\documentclass[leqno]{surv-l}

\usepackage{iftexml}
\usepackage{epigraph}

\usepackage{hyperref}
\usepackage{amssymb,graphicx}
\usepackage[mathscr]{euscript}
\usepackage{amsmidx,morefloats}

\theoremstyle{plain}
\newtheorem{theorem}{Theorem}
\newtheorem{exercise}{Exercise}[chapter]
\newtheorem{lemma}{Lemma}
\newtheorem*{theorem*}{Theorem}
\newtheorem*{lemma*}{Lemma}
\newtheorem*{corollary*}{Corollary}
\newtheorem{unsec}{H-B}
\theoremstyle{definition}
\newtheorem{example}[theorem]{Example}
\newtheorem*{definition*}{Definition}
\newtheorem{definition}{Definition}
\newtheorem*{claim*}{Claim}

\renewcommand{\thepart}{\Roman{part}}
\renewcommand{\thelemma}{\Alph{lemma}}

\makeindex{Subject}
\makeindex{Author}

\def\chaptername{Part}
\renewcommand{\thechapter}{\Roman{chapter}}
\setcounter{tocdepth}{0}

\makeatletter
\newcommand\addtocline[2]{\let\@secnumber\@empty\@tocwrite{#1}{#2}}
\ifTeXML\DeclareRobustCommand{\SkipTocEntry}[4]{}\fi%
\makeatother

\begin{document}

\refsectiontrue

\frontmatter

\tableofcontents

\mainmatter

%%%%%%intro

\chapter*{Introduction}

\section*{Background}

In this book we have attempted to collect most of the consequences of the axiom of choice that have been discovered during the past 90 years. Our goal is to present these statements and the relationships between them in an easily accessible format. Since 1904, when Zermelo introduced the axiom of choice, hundreds of its consequences have appeared in articles by almost as many authors. These articles have explored the complicated pattern of implications between the various consequences. With the introduction of independence proofs by Fraenkel, Mostowski, Specker and Cohen, the pattern of relationships became still more intricate. The results in this area are scattered over a wide spectrum of journals and books. As a result a particular theorem may be proved and reproved by several different authors. For example, Shannon\index{Author}{Shannon} \cite{bib:1988} asked whether or not the axiom of choice for countable families of countable sets implies the countable union theorem. In Howard\index{Author}{Howard} \cite{bib:1992} Shannon's question was answered in the negative. Shortly afterward Norbert Brunner, in his review of Howard \cite{bib:1992}, pointed out that the result had already been obtained by Felgner\index{Author}{Felgner} in \cite{bib:1974}. This book was begun with the primary purpose of preventing such occurrences. We hope it will also serve as a reference for non-set theorist who need to know the strength of some particular statement relative to the axiom of choice and its consequences.

In some ways our work is similar to that of Herman and Jean Rubin\index{Author}{H. Rubin/J. Rubin} (Rubin/Rubin \cite{bib:1985}) which contains equivalents of the axiom of choice. However, our task is complicated by the fact that the statements we consider are not, in general, equivalent. We are concerned with sentences in the language of set theory which can be proved using the axiom of choice but which are not theorems of set theory with that axiom omitted (with the exception of form 0 described below). We call these sentences \emph{forms of the axiom of choice}. The forms are numbered starting with form 0. Form 1 is the axiom of choice and form 0 is $0=0$, provable in set theory without any form of the axiom of choice. The remaining forms are numbered in the order we encountered them as we reviewed the literature.

We consider two versions of set theory. The first is Zermelo-Fraenkel set theory\index{Subject}{Zermelo-Fraenkel set theory} which we denote by $\mathrm{ZF}$. For our purposes any one of the axiomatizations of the theory $\mathrm{ZF}$ will be satisfactory. To be specific we shall use the one given in Jech\index{Author}{Jech} \cite{bib:1973b}, page 35. The second version of set theory is denoted by $\mathrm{ZF}^{0}$\index{Subject}{$\mathrm{ZF}^{0}$}. The theory $\mathrm{ZF}^{0}$ is the following modification of $\mathrm{ZF}$. A constant symbol $\emptyset$ is added to the language together with the axiom $(\forall x)(x\not\in\emptyset)$. In addition the axioms of extensionality (A1 in Jech \cite{bib:1973b}) and foundation (or regularity) (A8 in Jech \cite{bib:1973b}) are modified to allow the existence of \emph{atoms} (or \emph{urelements}). Each atom $a$ has the property that $a\neq\emptyset$ and $\neg(\exists x)(x\in a)$. Let set$(x)$ be the formula $x=\emptyset \vee(\exists y)(y\in x)$. Then the modified version of the extensionality axiom mentioned above is
\begin{equation*}
(\forall X)(\forall Y)[(\mathrm{set(X)}\wedge \mathrm{set}(Y))\rightarrow((\forall u)(u\in X\leftrightarrow u\in Y)\rightarrow X=Y)].
\end{equation*}
The modified version of the axiom of regularity is
\begin{equation*}
(\forall S)[(S\neq\emptyset \wedge \mathrm{set}(S) \wedge(\forall u\in S)(\mathrm{set}(u)))\rightarrow(\exists x\in S)(x\cap S=\emptyset)].
\end{equation*}

Corresponding to these two theories are two types of models: Cohen models\index{Subject}{Cohen models} of $\mathrm{ZF}$ which use \emph{forcing} to extend a model of $\mathrm{ZF}$ in which the axiom of choice holds to a model in which it fails and Fraenkel-Mostowski (FM) permutation models\index{Subject}{permutation models} of $\mathrm{ZF}^{0}$. Starting with a model of $\mathrm{ZF}^{0}$ in which the axiom of choice is true, and in which the class of atoms is non-empty, $\mathrm{FM}$ models use permutations of the atoms to obtain a submodel in which the axiom of choice is false. Both types of models are described in more detail in part~\ref{pt03:part03} of the book.

Our policy on notation and terminology is (for the most part) to use that of the original articles and books, or standard terminology. However, we are careful to distinguish between the terms ``denumerable'' and ``countable''. We use the term ``denumerable''\index{Subject}{denumerable set} to mean ``countably infinite''. A set $x$ is \emph{denumerable} or \emph{countably infinite} if $x\approx(\omega)$, and a set is \emph{countable}\index{Subject}{countable set} if it is either finite or denumerable.

\section*{Description of the parts of the book}

This book is divided into seven parts

Part~\ref{pt01:part01} is the numerical list of forms. For example, the listing for form 3 is

\noindent \textbf{FORM 3.} $2m=m$: For all infinite cardinals $m,\,2m=m$. Tarski\index{Author}{Tarski} \cite{bib:1954a}.

\textbf{[3 A]}\index{Author}{Sierpi\'{n}ski} For all\index{Subject}{countable axiom of choice} cardinals $m$ and $n,\,n\leq m$ and $m$ infinite $\rightarrow m+n=m$. Halpern/Howard\index{Author}{Halpern/Howard} \cite{bib:1970}.

\noindent \textbf{[3 B]} $\aleph_{0}\cdot m=m$: For all infinite cardinals $m,\,\aleph_{0}\cdot m=m$. Halpern/Howard \cite{bib:1970}.

\noindent \textbf{[3 C]} Definability of cardinal addition as the least upper bound: (For all cardinals $x,\,y$ and $z$) ($x+y=z$ iff ($x\leq z$ and $y\leq z$ and $(\forall u)(x, y\leq u\rightarrow z\leq u)$). H\"{a}ussler\index{Author}{H\"{a}ussler} \cite{bib:1983} and Tarski \cite{bib:1949a}.

\noindent \textbf{[3 D]} Every pair of cardinals has a least upper bound in the usual cardinal ordering and for all cardinals $\kappa$ and $\gamma$, if $n\kappa \leq\gamma$ for all $n\in\omega$, then $\aleph_{0}\kappa\leq\gamma$. Hickman\index{Author}{Hickman} \cite{bib:1979b}.

\noindent \textbf{[3 E]} $E(IV, V)$: For all infinite cardinals $m,\,m+1=m$ implies $2m=m$. Howard/Yorke\index{Author}{Howard/Yorke} \cite{bib:1989} and note 94.

\noindent \textbf{[3 F]} $E(III, V)$: For every set $X$, if $\mathcal{P}(X)$ is Dedekind infinite then $2|X|=|X|$. Howard/Yorke \cite{bib:1989} and note 94.

\noindent \textbf{[3 G]} $E(II, V)$: For every set $X$, if there is a non-empty family of subsets of $X$ which is linearly ordered by $\subseteq$ and has no maximal element then $2|X|=|X|$. Howard/Yorke \cite{bib:1989} and note 94.

\noindent \textbf{[3 H]} $E(Ia, V)$: For every set $X$, if $2|X|>|X|$ then $X$ is amorphous. Howard/Yorke \cite{bib:1989} and note 94.

In most cases a form's listing contains its number, a name for the form, a precise statement of the form and a reference to the bibliography. Following each form in part~\ref{pt01:part01} is a list of statements known to be equivalent to the form. We also refer to these statements as forms. In the example given above, forms [3 A] through [3 H] are all equivalent to form 3. The listing of each equivalent form ends with a reference to an article (or note) where the proof of equivalence can be found.

Part~\ref{pt02:part02} is the topical list of forms. This listing of the forms by subject matter is the part of the book which is most similar to H. Rubin/J. Rubin\index{Author}{H. Rubin/J. Rubin} \cite{bib:1985}. The listing of each form in part~\ref{pt02:part02} contains the same information as its listing in part~\ref{pt01:part01}. One difference between part~\ref{pt01:part01} and part~\ref{pt02:part02} is that in part~\ref{pt02:part02} a form and its equivalents need not occur together. For example, form 14 (the Boolean Prime Ideal Theorem) is in the section of algebraic forms while its equivalent form [14 P] (Alaoglu's Theorem) is in the section of forms from analysis.

Part~\ref{pt03:part03}, titled ``Models'', contains a brief description of each model of set theory. These descriptions include the forms that have been proved true in the model and the forms that have been proved false together with a reference to the proof. Each description ends with a list of all forms known to be true in the model and all forms known to be false. For example, in model $\mathcal{M}1$ form 14 (the Boolean Prime Ideal Theorem) was proved to be true in Halpern/Levy\index{Author}{Halpern/Levy} \cite{bib:1971}. These facts are included in the description of $\mathcal{M}1$. Since form 14 implies form 30 (the Ordering principle), form 30 is also true in the model and therefore form 30 occurs in the final list of forms which are true in the model. There are two sections in part~\ref{pt03:part03}. Section 1 contains Cohen models for $\mathrm{ZF}$. In most cases the models are denoted by $\mathcal{M}k$, where $k$ is a positive integer. Section 2 contains models of $\mathrm{ZF}^{0}$. Our notation for these models is $\mathcal{N}k$, where $k$ is a positive integer.

Part~\ref{pt04:part04}, titled ``Notes'', is a collection of notes which are referred to in other parts of the book. Whenever we consider notation or terminology to be nonstandard we add a note explaining it. Our policy on proofs is to omit them if they are readily available in the literature. However, if a proof is not, in our judgement, readily available, we have included it in a note. Other notes include a survey of the relationships between the (infinitely many) versions of the axiom of choice for sets of $n$ element sets, $n$ a positive integer (note 15), a brief exposition of the transfer theorems of Jech/Sochor and Pincus (note 103), and various properties of and relations between forms (notes 18 and 120).

Part~\ref{pt05:part05} is called ``References for Relationships Between Forms''. To explain part~\ref{pt05:part05} it is helpful to consider an example. A typical line from part~\ref{pt05:part05} is below.

43\quad 63\quad (3)\quad Felgner\index{Author}{Felgner} \cite{bib:1971a} (M2)

The first two numbers are form numbers. The third number which occurs in parentheses is a code number for the status of the implication ``form 43 implies form 63''. The number 3 is the code for ``not provable in Zermelo-Fraenkel set theory''. The other code numbers that we use are
\begin{enumerate}
\item[0] for ``The status of the implication is unknown.''
\item[1] for ``The implication is provable.''
\item[2] for ``The implication is provable and this fact follows from implications whose code numbers are 1.''
\item[3] for ``The implication is not provable in $\mathrm{ZF}$.''
\item[4] for ``The implication is not provable in $\mathrm{ZF}$ and this fact follows from implications whose code numbers are 1 or 3.''
\item[5] for ``The implication is not provable in $\mathrm{ZF}^{0}$.''
\item[6] for ``The implication is not provable in $\mathrm{ZF}^{0}$ and this fact follows from implications whose code numbers are odd.''
\item[7] for ``The status of the implication is unknown and we have found the question asked somewhere in the literature.''
\end{enumerate}

In part~\ref{pt05:part05} the implication codes\index{Subject}{implication codes} are odd numbers for the most part. Following the code number in a typical line of part~\ref{pt05:part05} is a reference to a bibliography item and possibly a note and, in the case of a non-implication, a model number. In the example, Felgner\index{Author}{Felgner} \cite{bib:1971a} refers to a bibliography item in part VI and M2 refers to the model $\mathcal{M}2$ from part~\ref{pt03:part03}.

Part~\ref{pt06:part06}, our bibliography, is a listing of all articles and books we have used arranged alphabetically by author.

Part VII consists of two $(n+1)$ by $(n+1)$ matrices where $n$ is the last form number. (Currently $n=383$.) We refer to these matrices as table 1\index{Subject}{table 1, description} and table 2\index{Subject}{table 2, description}.

Printed copies of table 1 and table 2 are not included. They may be
accessed using the software that is included. In addition the
program maketex.c in the software package will create \TeX\ files
containing copies of table 1 and table 2 which may then be printed.
Tables 1 and 2 are also available at either of the authors' web
sites: \url{http://www.math.purdue.edu/~jer/} or
\url{http://www.emunix.emich.edu/~phoward/}

The natural numbers $0$ through $n$ are the row and column headings in both table 1 and table 2 and are to be thought of as form numbers. Each entry in the body of table 1 is one of the code numbers $0$ through 7 used in part~\ref{pt05:part05}. The code number in row $n$ and column $k$ of table 1 describes the status of the implication ``form $n$ implies form $k$''. Odd entries (1, 3, 5 or 7) mean that the implication or non-implication (or open problem) has a reference listed in part~\ref{pt05:part05}. For example, the entry in row 43, column 39 of table 1 is 3. This means that form 43 does not imply form 39 in $\mathrm{ZF}$ and there is a reference in part~\ref{pt05:part05} to this result. The entry in table 1 at position 8, 40 (that is, row 8, column 40) is 4. This means that form 8 does not imply form 40 in $\mathrm{ZF}$ and this non-implication follows from implications and non-implications that have references in part~\ref{pt05:part05}.

When an even entry is encountered in table 1 it is possible, by studying table 1, to find implications and non-implications from which this even entry follows and which are referenced in part~\ref{pt05:part05}. For example, an examination of table 1 reveals that the entries at positions 43, 8 and 40, 39 are both 1 and that (as noted earlier) the entry at position 43, 39 is a 3. This means that form 43 implies form 8, form 40 implies form 39 and form 43 does not imply form 39 in $\mathrm{ZF}$. Further all these results are referenced in part~\ref{pt05:part05}. It follows from these three results that form 8 does not imply form 40 and hence we have ``reduced'' the non-implication ``form 8 does not imply form 40'' to results that have references in part~\ref{pt05:part05}.

The purpose of table 2 is to provide a direct path to this reduction. In table 2 the non-zero entries are form numbers. These non-zero entries occur only in positions where the corresponding entry in table 1 is 2, 4, or 6. If the table 1 entry at position $m,\,k$ is 2, then the entry in table 2 is a form number $j$ such that form $m$ implies form $j$ and form $j$ implies form $k$. Thus the implication ``form $m$ implies form $k$'' is ``reduced'' to the implications ``form $m$ implies form $j$ ''  and ``form $j$ implies form $k$''. If the table 1 entries at positions $m,\,j$ and $j,\,k$ are both $1$'s then the original implication (``form $m$ implies form $k$'') has been reduced to implications with references in part~\ref{pt05:part05}. If either of the table 1 entries at positions $m,\,j$ or $j,\,k$ is a 2 then the implication corresponding to the 2 can be reduced again using table 2. Similarly, if the entry at position $m,\,k$ of table 1 is a 4 or a 6, the table 2 entry at position $m,\, k$ will be a form number $j$ of a form such that one of two things happens: Either form $j$ implies form $m$ and form $j$ does not imply form $k$ or form $k$ implies form $j$ and form $m$ does not imply form $j$. These two possibilities are illustrated in the diagram below.
\begin{equation*}
\begin{array}{ll}
\mathrm{form}\ j\quad {\not\rightarrow}\quad \mathrm{form}\ k &\qquad\qquad\qquad\qquad\qquad\qquad \mathrm{form}\ k\\
\quad \downarrow & \qquad\qquad\qquad\qquad\qquad\qquad\quad\downarrow\\
\mathrm{form}\ m & \qquad \quad\qquad\mathrm{form}\ m\quad \not\rightarrow\quad\mathrm{form}\ j
\end{array}
\end{equation*}
\noindent In either case the non-implication ``form $m$ does not imply form $k$'' has been ``reduced''. Using this method with the 4 at position 8, 40 in table 1 we see that the table 2 entry at position 8, 40 is 231. Referring back to table 1, we see that the entry in position 8, 231 is a 3 and the entry at position 40, 231 is a 1. The non-implication ``form 8 does not imply form 40'' has been reduced to the two facts `` form 8 does not imply form 231'' and ``form 40 implies form 231'' both of which have references in part~\ref{pt05:part05}.

\section*{Software}

This section contains a description\index{Subject}{software description} of the software. The next section contains detailed instructions for its use. All programs and computer files mentioned in this section are on the diskette enclosed inside the back cover. Using the software that is included will eliminate the need for calculations using table 2. The enclosed diskette contains two data files: rfb1.tex and models.dat. The purpose of these files is described later in this section.

The software also includes several programs each of which is described below. There are two versions of each program: A source file with extension .c suitable for compiling on any machine with a c compiler and an executable version with extension .exe which can be run directly on any $\mathrm{PC}$ which runs MS DOS. Here are the descriptions of the c programs
\begin{enumerate}
\item[1.] mk\_book.c -- This program must be run before any of the others. (From the file rfb1.tex it creates two matrices of size 384 by 384 (which are versions of table 1 and table 2) and stores them in the two files book1 and book2. These files (book1 and book2) are used by the other c programs to extract information about implications among the forms. Each of book1 and book2 uses about .5 megabytes of memory.
\end{enumerate}

\textbf{AFTER RUNNING mk-book.c ANY OF THE FOLLOWING CAN BE USED:}
\begin{enumerate}
\item[2.] weakac.c -- This is the main program for obtaining information about the relationships between forms. When you run the program the data files book1 and book2 will be read into memory and then the main menu will appear on the screen. Then menu is show below.

\item[]MAIN MENU

\begin{enumerate}
\item[1.] see the meanings of the implication codes

type 1 $<$enter$>$
\item[2.] see what is known about the implication !`form n implies form k?` for any two form numbers n and k

type 2 $<$enter$>$
\item[3.] get information about a single form

type 3 $<$enter$>$
\item[4.] see table 1 restricted to any subset of the forms

type 4 $<$enter$>$
\item[5.] see a submatrix of table 1 consisting of consecutive rows and columns 

type 5 $<$enter$>$
\item[6.] see how an implication or non-implication follows from those with a reference in part~\ref{pt05:part05}

type 6 $<$enter$>$
\item[7.] see the models in which a form is true and those in which it is false

type 7 $<$enter$>$
\item[8.] exit the program

type 8 $<$enter$>$
\end{enumerate}
In order for option 7 to work the file models.dat must be present.

\item[3.] maketex.c -- A program which will create \TeX\ files for table 1 and table 2 from book1 and book2 if you want a hard copy of either of these.
\end{enumerate}

\section*{Using the Software}

\textbf{The Setup.} There are three steps in the setup procedure:
\begin{enumerate}
\item[1.] Copy the files from the floppy disk to your computer's memory.
\item[2.] Compile the c programs using a c language compiler. (This step is not necessary if you are working on a computer running MS DOS. We have included complied executable DOS files for each of the c programs.)
\item[3.] Run the program \textbf{mk{\_}book}. This program creates the data files used by the other programs.
\end{enumerate}

The program mk{\_}book takes between 5 and 10 minutes to run. It creates the data files book1 and book2 from the file rfb1.tex. The files copied from the floppy disk together with the data files book1 and book2 occupy about 2.3 megabytes of memory.

To carry out the setup procedure on a DOS machine, put the floppy disk in the disk drive. Assuming that this is drive a:, the following instructions given at the DOS prompt $(\mathrm{C}:\backslash >)$ will create a directory called weakac and copy the files from the disk into that directory.

\noindent md $\mathrm{c}:\backslash $weakac

\noindent copy $\mathrm{a}:^{\ast}\cdot^{\ast}\mathrm{c}:\backslash $weakac

\noindent To complete the setup on a DOS machine change to the weakac directory and run mk\_book with the following two commands.

\noindent cd c:$\backslash$ weakac

\noindent mk\_book

\noindent Answer y to the question ``Do you want to continue?''.

On a non-DOS machine a c language compiler is required. You must compile the c language programs mk\_book.c, weakac.c, and maketex.c. On a UNIX machine, for example, the command

\noindent cc -o mk\_book mk\_book.c

\noindent will compile the c program mk\_book.c and store the executable program in the file mk\_book. The other programs should be compiled similarly. To complete the setup type

\noindent mk\_book

\noindent at the UNIX prompt to run the program mk\_book.

\textbf{Running the Other Programs.} The program \textbf{weakac} is the program that is used for getting information about the relationships between the various consequences of $\mathrm{AC}$. It is described in the previous section. To run it type

\noindent weakac

\noindent at the DOS or UNIX prompt. Follow the instructions on the screen. In particular you are offered the option of writing the program's output in a file.

The program \textbf{maketex} will create \TeX\ files for either table 1 or table 2. At the DOS or UNIX prompt type

\noindent maketex

\noindent You will be prompted with the lines:

\noindent ``Enter the mane of the file to be converted to

\noindent TeX. This will usually be book1 or book2.''

To make a \TeX\ file for table 1 type book1. To make a \TeX\ file for table 2 type book2. You will then be prompted with the question:

\noindent ``To what file name should the TeX file be

\noindent written? Must be different from the input file.''

\noindent At this point you choose the name for the output file. It may be any valid file name other that book1 or book2. For example table1.tex or table2.tex are suitable. A few more lines appear on the screen and you are prompted with the question ``Today's date?'' Enter the date in any format. For example 8/10/97 or August 10, 1997 are suitable responses. The final question is

\noindent ``Will the output be table 1 or table 2?

\noindent respond 1 or 2.''

\noindent Make the appropriate response and the \TeX\ file for table 1 or table 2 will be created. When processed by a \TeX\ program and printed, each of these tables is 144 pages long.

\section*{Acknowledgements}

We have received many helpful suggestions and comments from colleagues. In particular we thank Andreas Blass, Norbert Brunner, Robert Cowen, Kyriakos Keremedis, Marianne Morillon, David Pincus, Herman and Arthur Rubin, Eric Schechter and John Truss. We would also like to thank the referee whose report we found extremely helpful. Our thanks also to Andrew Smith. His translations of several articles into English made life much easier for us. We are grateful to Leonard Lipshitz and the Mathematics Department at Purdue University for providing an opportunity for the authors to work side by side during the Fall 1994 semester. The publishers representatives, Sarah Donnelly and Sergei Gelfand, were very helpful. We should also like to thank Betty Gick, a secretary in the Purdue Mathematics Department, who helped us with \TeX. Finally we are grateful to our spouses Herman Rubin and Gale Howard for their patience and encouragement.

\section*{Final Note}

We have tried to give the original sources of the results. However, this is not always possible. When we give references like Jech\index{Author}{Jech} \cite{bib:1973b}, Mathias\index{Author}{Mathias} \cite{bib:1979}, G. Moore\index{Author}{G. Moore} \cite{bib:1982}, H. Rubin/J. Rubin\index{Author}{H. Rubin/J. Rubin} \cite{bib:1985}, etc., it means that information about the results can be found there. It does not necessarily mean that it is the original source. We have also made an effort to check the accuracy of our results, especially when the source is a preprint. As far as we can tell, our results are consistent. However, because of the amount of data that we have, it is virtually impossible to prevent some error from sneaking in. We just hope the errors are not serious and we encourage readers to contact us when they encounter errors or software problems. We would also appreciate hearing of new results in this area. We shall add all additions and corrections to our web pages.

\begin{table}[h]
{\begin{tabular*}{\textwidth}{@{\extracolsep{\fill}}ll@{}}
Paul Howard & Jean E. Rubin\\
Department of Mathematics & Department of Mathematics\\
Eastern Michigan University & Purdue University\\
Ypsilanti, MI 48197 & West Lafayette, IN 47907\\
email phoward@emunix.emich.edu & email jer@math.purdue.edu\\
\url{http://www.emunix.emich.edu/~phoward/} & \url{http://www.math.purdue.edu/~jer/}\\
\end{tabular*}}
{}
\end{table}

%%%%%%%part01

\part{Numerical List of Forms \label{pt01:part01}}

\textbf{FORM 0.} 0 = 0.

[\textbf{$\boldsymbol{0}$ A}] Cardinal successors 2: For every cardinal $m$ there is a cardinal $n$ such that $m<n$ and $(\neg\exists p)(m<p<n)$. Tarski\index{Author}{Tarski} \cite{bib:1954a} and Jech\index{Author}{Jech} \cite{bib:1966a}.

[\textbf{$\boldsymbol{0}$ B}] If $X$ is finite then $X^{\omega}$ is compact. Brunner\index{Author}{Brunner} \cite{bib:1984b}.

[\textbf{$\boldsymbol{0}$ C}] Closed subspaces of limited amorphous spaces are limited amorphous. (A space is \emph{limited amorphous} if each amorphous subset is compact.) Brunner \cite{bib:1984b}.

[\textbf{$\boldsymbol{0}$ D}] Products of limited amorphous spaces are limited amorphous. (A space is \emph{limited amorphous} if each amorphous subset is compact.) Brunner \cite{bib:1984b}

[\textbf{$\boldsymbol{0}$ E}] The box product of a well ordered family of limited amorphous spaces is limited amorphous. (A space is \emph{limited amorphous} if each amorphous subset is compact.) Brunner \cite{bib:1984b} and note 52.

[\textbf{$\boldsymbol{0}$ F}] Cancellation for surjective cardinal equivalence: For every $k\in \omega -\{0\},\,(\forall x\forall y)(kx=^{\ast} ky$ implies $x=^{\ast}y)$ (The expression $x=^{\ast}y$ means that there is a function $f$ from $x$ onto $y$ and a function $g$ from $y$ onto $x$.) Truss\index{Author}{Truss} \cite{bib:1984}.

[\textbf{$\boldsymbol{0}$ G}] Cancellation for injective cardinal inequality: For every $k\in\omega -\{0\},\,(\forall x\forall y)(kx\leq ky$ implies $x\leq y)$ ($\leq$ is the usual (injective) cardinal ordering.) Lindenbaum/Tarski\index{Author}{Lindenbaum/Tarski} \cite{bib:1926}.

[\textbf{$\boldsymbol{0}$ H}] F. Riesz Theorem:\index{Subject}{Riesz theorem} A Hilbert space is finite dimensional iff its closed unit ball is compact. Brunner \cite{bib:1983b}.

[\textbf{$\boldsymbol{0}$ I}] The following are equivalent in a Hilbert space:

\noindent(a) The closed unit ball is sequentially compact.

\noindent(b) Each unconditionally convergent series converges absolutely.

\noindent(c) Each ortho-normal system is Dedekind finite.

\noindent Brunner \cite{bib:1983b} and note 94.

[\textbf{$\boldsymbol{0}$ J}] Zermelo's Fixed Point Theorem:\index{Subject}{Zermelo's fixed point theorem} For every partially ordered set $(X, \leq)$, if every well ordered subset has a least upper bound then every $f:X\rightarrow X$ satisfying $\forall t\in X,\,t\leq f(t)$ has a fixed point. Manka\index{Author}{Manka} \cite{bib:1988a}, Zermelo\index{Author}{Zermelo} \cite{bib:1908b} and Abian\index{Author}{Abian} \cite{bib:1980}.

[\textbf{$\boldsymbol{0}$ K}] For all partially ordered sets $(X, \leq)$, if there is a $\displaystyle \sup$ function $\sigma$ on the well ordered subsets of $X$ then every $f:X\rightarrow X$ satisfying $\forall t\in X,\,t\leq f(t)$ has a fixed point. Manka \cite{bib:1988a} and note 38.

[\textbf{$\boldsymbol{0}$ L}] Caristi's Fixed Point Theorem:\index{Subject}{Caristi's fixed point theorem} If $(X, \rho)$ is a complete metric space and $\phi$: $X\rightarrow \mathbb{R}$ is bounded above and upper semi-continuous then in the Br{\o}ndsted ordering $(x\leq y$ iff $\rho(x, y)\leq\phi(y)-\phi(x))$ every $f:X\rightarrow X$ satisfying $\forall t\in X,\,t \leq f(t)$ has a fixed point. Caristi\index{Author}{Caristi} \cite{bib:1976}, Manka \cite{bib:1988a}, and note 38.

[\textbf{$\boldsymbol{0}$ M}] A $T_{1}$ topological space is Dedekind finite if and only if every subspace has a dense Dedekind finite subset. Brunner\index{Author}{Brunner} \cite{bib:1982d} and note 94.

[\textbf{$\boldsymbol{0}$ N}] If $X$ and $Y$ are topological spaces with dense Dedekind finite subsets so is $X\times Y$. Brunner \cite{bib:1982d} and note 94.

[\textbf{$\boldsymbol{0}$ O}] A Lindel\"{o}f space with a dense, Dedekind finite subset is compact. Brunner \cite{bib:1982d} and note 94.

[\textbf{$\boldsymbol{0}$ P}] A Lindel\"{o}f space with a Dedekind finite basis is compact. Brunner \cite{bib:1982d} and note 94.

[\textbf{$\boldsymbol{0}$ Q}] For all cardinal numbers $m$ and $n$ and Dedekind finite cardinals $p$, if $p+m=p+n$ then $m=n$. Tarski\index{Author}{Tarski} \cite{bib:1949b} and note 94

[\textbf{$\boldsymbol{0}$ R}] Every infinite Dedekind finite field is the direct limit of a strictly increasing $\omega$-sequence of finite subfields. Hickman\index{Author}{Hickman} \cite{bib:1982}, notes 45 and 94.

[\textbf{$\boldsymbol{0}$ S}] If $(P, \leq)$ is a quasi-order such that $P$ is non-empty, well orderable and every antichain is finite then there is a $P$-generic filter. Shannon\index{Author}{Shannon} \cite{bib:1990} and note 47.

[\textbf{$\boldsymbol{0}$ T}] If a set $S$ is hereditarily countable\index{Subject}{hereditarily countable set} (that is, every $y\in TC(S)$ is countable) then the rank of $S$ is $<\omega_{2}$. Jech\index{Author}{Jech} \cite{bib:1982}.

[\textbf{$\boldsymbol{0}$ U}] For any class of similar (universal) algebras $K,\,HSP\, K=Mod\,Eq\,K$. (The class of homomorphic images of subalgebras of products of non-empty families of algebras in $K$ is the same as the class of algebras which satisfy the set of equations holding in every algebra in $K$.) Andreka/Nemeti\index{Author}{Andreka/Nemeti} \cite{bib:1981} and note 50.

[\textbf{$\boldsymbol{0}$ V}] For every class $K$ of similar algebras $SP\,K\supseteq PS\,K$. (For any class $K$ of similar algebras, $S\,K$ is the class of all subalgebras of elements of $K$ (not closed under isomorphism) and $P\,K$ is the class of algebras isomorphic to direct products of algebras in $K$). Andreka/Nemeti \cite{bib:1980}.

[\textbf{$\boldsymbol{0}$ W}] For every class $K$ of similar algebras $SP^{r}K\supseteq P^{r}S\,K$. (For any class $K$ of similar algebras, $S\,K$ is the class of all subalgebras of elements of $K$ (not closed under isomorphism) and $P^{r}K$ is the class of all algebras isomorphic to reduced direct products of algebras in $K$.) Andreka/Nemeti \cite{bib:1980}.

[\textbf{$\boldsymbol{0}$ X}] $DC(\mathbf{\Sigma}_{2}^{1})$: If $P\subseteq {^{\omega}}\omega \times {^{\omega}}\omega$ and $P$ has domain ${^{\omega}}\omega,\, P\in\boldsymbol{\Sigma}_{2}^{1}$, then for any $y\in {^{\omega}}\omega$ there is a sequence $\langle x_{k}: k\in \omega\rangle$ of elements of ${^{\omega}}\omega$ with $x_{0}=y$ and $\langle x_{k},\,x_{k+1}\rangle\in P$ for all $k\in\omega$. Kanovei\index{Author}{Kanovei} \cite{bib:1979} and note 61.

[\textbf{$\boldsymbol{0}$ Y}] There is no free complete Boolean algebra over $D$ if $|D|\geq\aleph_{0}$. Stavi\index{Author}{Stavi} \cite{bib:1975}, Gaifman\index{Author}{Gaifman} \cite{bib:1970}, Hales\index{Author}{Hales} \cite{bib:1964}, and note 89.

[\textbf{$\boldsymbol{0}$ Z}] Suppose $X$ and $M$ are sets, $F$ is a function from $\mathcal{P}(X)-\{\emptyset\}$ into $M$ and $G$ is a function on $M$ such that $\forall z\in \mathcal{P}(X)-\{\emptyset\},\,G(F(z))\in \mathcal{P}(z)-\{\emptyset\}$. Then there are unique $W$ and $R$ such that $W\subseteq M,\,R\subseteq W\times W$ and $R$ is an irreflexive well ordering of $W$, the range of $G|W$ is a partition of $X$, and $\forall w\in W,\,w=F(X -(\bigcup\{G(\nu) :\nu\in W$ and $\nu\, R\,w\}))$. Kruse\index{Author}{Kruse} \cite{bib:1974}.

[\textbf{$\boldsymbol{0}$ AA}] $CT_{1}$: Restricted Compactness Theorem for Propositional Logic I: If $\Sigma$ is a set of formulas in a propositional language such that every finite subset of $\Sigma$ is satisfiable and if each propositional variable occurs in at most one formula in $\Sigma$, then $\Sigma$ is satisfiable. Kolany/Wojtylak\index{Author}{Kolany/Wojtylak} \cite{bib:1991}.

[\textbf{$\boldsymbol{0}$ AB}] If $X\subseteq \mathbb{R}$ is perfect, then $|X|=2^{\aleph_{0}}$. Truss\index{Author}{Truss} \cite{bib:1974b} and note 28.

[\textbf{$\boldsymbol{0}$ AC}] There are $2^{\aleph_{0}}$ perfect subsets of $\mathbb{R}$. Truss\index{Author}{Truss} \cite{bib:1974b} and note 28.

[\textbf{$\boldsymbol{0}$ AD}] For all ordinals $\beta$, well ordered cardinals $\gamma\geq 1$ and all $n\in\omega$, there is a well ordered cardinal $\kappa$ such that $\kappa \rightarrow(\beta)_{\gamma}^{n}$. Erd\"{o}s/Rado\index{Author}{Erd\"{o}s/Rado} \cite{bib:1952} and note 97.

[\textbf{$\boldsymbol{0}$ AE}] The number of well ordered subsets of a set: For all $x,\,|x|<|s(x)|$ where $s(x)=$\{ $y\subseteq x:y$ can be well ordered\}. Tarski\index{Author}{Tarski} \cite{bib:1939}.

[\textbf{$\boldsymbol{0}$ AF}] The number of well orderings of a set: For all $x,\,|x|<|w(x)|$ where $w(x)$ is the set of all well orderings of subsets of $x$. Russ \cite{bib:1973d}.

[\textbf{$\boldsymbol{0}$ AG}] For any $n\in\omega,\, n\neq 0,\, |s(x)|\not\leq|x^{n}|$ and $|x^{n}|<|w(x)|$ if $|x|$ exceeds some small finite value depending on $n$. $(s(x)$ is the set of well orderable subsets of $x$ and $w(x)$ is the set of well orderings of subsets of $x$.) Truss \cite{bib:1973d}

[\textbf{$\boldsymbol{0}$ AH}] No infinite well orderable set is amorphous.

[\textbf{$\boldsymbol{0}$ AI}] Turinici's Fixed Point Theorem\index{Subject}{Turinici's fixed point theorem} for Hausdorff Spaces: If $(X, \leq)$ is a directed, partially ordered set and $\tau$ is a Hausdorff topology on $X$ such that

\noindent(a) $(X, \leq)$ is upper semi-continuous with respect to $\tau$,

\noindent(b) Every well ordered subset of $(X, \leq)$ has a unique limit as a net, then every function $f:X\rightarrow X$ such that $\forall x\in X,\,x\leq f(x)$ has a fixed point. Manka\index{Author}{Manka} \cite{bib:1988b} and note 38.

[\textbf{$\boldsymbol{0}$ AJ}] $\aleph_{\alpha+1}$ is not an inaccessible cardinal for any ordinal $\alpha$. Drake\index{Author}{Drake} \cite{bib:1974} and note 20.

[\textbf{$\boldsymbol{0}$ AK}] Every separable metric space is second countable.

[\textbf{$\boldsymbol{0}$ AL}] For every set $A$, there is an $\aleph$ which cannot be covered by fewer than $|A|$ sets each of cardinality $<|A|$. Note 131.

[\textbf{$\boldsymbol{0}$ AM}] Assume $V$ is a separable vector space, $D$ is a countable dense subset of $V$, and $p:V\rightarrow \mathbb{R}$ satisfies $p(x+y)\leq p(x)+p(y)$, and $(\forall t\geq 0)(\forall x\in V)(p(tx) =tp(x))$. Also assume that $p$ satisfies the following condition:
\begin{equation*}\tag{$^{\ast}$}
(\forall\epsilon>0)(\forall x\in V)(\exists y\in D)\ \text{such that}\ p(x-y)<\epsilon\ \text{and}\ p(y-x)<\epsilon.
\end{equation*}

\noindent Moreover, if $f$ is a linear function from a subspace $S$ of $V$ into $\mathbb{R}$ which satisfies $(\forall x\in S)(f(x)\leq p(x))$, then $f$ can be extended to $f^{\ast}:V\rightarrow \mathbb{R}$ so that $f^{\ast}$ is linear and $(\forall x\in V)(f^{\ast}(x)\leq p(x))$. Note 143.

\noindent\textbf{FORM 1.} $C(\infty, \infty)$: The Axiom of Choice:\index{Subject}{axiom of choice} Every set of non-empty sets has a choice function.

[\textbf{1 A}] In every vector space, every generating set contains a basis. Halpern\index{Author}{Halpern} \cite{bib:1966}.

\noindent [\textbf{1 B}] Every commutative ring with unit has a maximal ideal. Hodges\index{Author}{Hodges} \cite{bib:1979}.

\noindent[\textbf{1C}] For all infinite cardinals $m,\,UT(m, m, m)$: If every member of an infinite set $A$ has power $|A|$ then $|\bigcup A|=|A|$. (Note that [1 C] implies $m^{2}=m$ for all infinite cardinals $m$.) Howard\index{Author}{Howard} \cite{bib:1994}.

\noindent[\textbf{1 D}] Cardinal successors 3: For every cardinal $m$ there is a cardinal $n$ such that $m<n$ and $(\forall p)(m<p\rightarrow n\leq p)$. Tarski \cite{bib:1954a} and Jech\index{Author}{Jech} \cite{bib:1966a}.

\noindent[\textbf{1 E}] $K(0)$: For every set $S$ there is an ordinal $\alpha$ and a one to one function $f:S\rightarrow\alpha$. (See form $81(n)$).

\noindent[\textbf{1 F}] For all ordinals $\alpha,\,DC(\aleph_{\alpha})$: Given a relation $R$ such that for every subset $Y$ of a set $X$ with $|Y|<\aleph_{\alpha}$, there is an $x\in X$ with $Y\ R\,x$, then there is a function $f:\aleph_{\alpha}\rightarrow X$ such that $\forall\beta<\aleph_{\alpha},\, \{f(\gamma):\gamma<\beta\}\,R\,f(\beta).)$ Levy\index{Author}{Levy} \cite{bib:1964} and note 133.

\noindent[\textbf{1 G}]\index{Author}{Shannon} Form 163 + form 8. Brunner\index{Author}{Brunner} \cite{bib:1982a} and note 94.

\noindent[\textbf{1 H}] Form 164 + form 67. Brunner \cite{bib:1982a} and note 94.

\noindent[\textbf{1 I}] In every metric space ($X,\,d$) there is an $\epsilon$-lattice for some $\epsilon>0$ which is not a strict upper bound of $d$. Brunner \cite{bib:1984e} and note 17.

\noindent[\textbf{1 J}] Form 49 + form 51. Morris\index{Author}{Morris} \cite{bib:1969} and note 121.

\noindent[\textbf{1 K}] Form 133 + form 64. Clear

\noindent[\textbf{1 L}] Hajnal's Free Set Principle:\index{Subject}{Hajnal's free set principle} For every set mapping $f:X\rightarrow[X]^{<\lambda}$ where $\lambda$ is a well orderable cardinal and $\lambda<|X|$ there is a free set of cardinality $|X|.\, ([X]^{<\lambda}$ is the set of subsets of $X$ of cardinality $<\lambda$. Brunner \cite{bib:1989} and note 22.

\noindent[\textbf{1 M}]\index{Subject}{ramification lemma} The ramification Lemma: If $S:X^{\leq\theta}\rightarrow \mathcal{P}(E)$ is a ramification system, then for each $g\in E$ there is an $f$ which is maximal (with respect to inclusion) in $\{h\in X^{\leq\theta}:g\in S(h)\}$. Brunner \cite{bib:1989} and note 22.

\noindent[\textbf{1 N}] For every set mapping $f:X\rightarrow[X]^{WO}$ (well ordered subsets of $X$) there is a co-well-orderable free subset $H$. ($H$ is \emph{co-well-orderable} if $X\,\backslash\, H$ is well orderable). Brunner \cite{bib:1989} and note 22.

\noindent[\textbf{1 O}] $A(D1, C)$: For every $T_{2}$ topological space $(X, T)$, if $X$ is covered by a well ordered family of closed, discrete sets then each family $U\subseteq T$ of pairwise disjoint sets is well ordered. Brunner \cite{bib:1983d} and note 26.

\noindent[\textbf{1 P}] $A(P)$: (Depends on $P$ for $P\in\{W1, W2, B1, L1, L2, S, C, H1, A1\})$ Every $T_{2}$ topological space has property $P$. Brunner \cite{bib:1983d} and note 26.

\noindent[\textbf{1 Q}] For any partial order $(X, \leq)$ there is a sup function on the bounded (above), well ordered subsets of $X$. Manka\index{Author}{Manka} \cite{bib:1988a} and note 38.

\noindent[\textbf{1 R}] For any partial order $(X, \leq)$ there is a sup function on the bounded subsets of $X$. Manka \cite{bib:1988a} and note 38.

\noindent[\textbf{1 S}] For any partial order $(X, \leq)$ if every well ordered subset of $X$ is bounded then there is a sup function on the well ordered subsets of $X$. Manka \cite{bib:1988a} and note 38.

\noindent[\textbf{1 T}] For every directed partial order $(X,\, \leq)$ there is a sup function on the bounded, well ordered subsets of $X$. Manka \cite{bib:1988a} and note 38.

\noindent[\textbf{1 U}] For every directed partial order, if every well ordered subset of $X$ is bounded above then there is a sup function on the well ordered subsets of $X$. Manka \cite{bib:1988a} and note 38.

\noindent[\textbf{1 V}] For every partially ordered set $(X, \leq)$, if every well ordered subset is bounded above then every function $f:X\rightarrow X$ satisfying $\forall t\in X,\,t\leq f(t)$ has a fixed point. Manka \cite{bib:1988a} and Abian\index{Author}{Abian} \cite{bib:1980}.

\noindent[\textbf{1 W}] Form 133 + form 88. Brunner \cite{bib:1982a}.

\noindent[\textbf{1 X}] Form 15 + form 91. Brunner \cite{bib:1982a}.

\noindent[\textbf{1 Y}] Refined Dual Cantor-Bernstein Theorem:\index{Subject}{Cantor-Bernstein theorem} If $f:X\rightarrow Y$ and $g:Y\rightarrow X$ are onto then there is a bijection $h:X\rightarrow Y$ such that $h\subseteq f \cup g^{-1}$. Banaschewski/Moore\index{Author}{Banaschewski/Moore} \cite{bib:1990}.

\noindent[\textbf{1 Z}] If $(P, \leq)$ is a quasi-order and $P\neq\emptyset$ and every antichain is finite, then there is a $P$ generic filter which is well orderable. Shannon \cite{bib:1990} and note 47.

\noindent[\textbf{1 AA}] Assume $f:S\rightarrow R$ where $f$ is surjective and $R\subseteq S$. Then every mapping $g$ of $S$ into $R$ has the form $g=f\circ h$ for some $h:S\rightarrow S$. Smith\index{Author}{Smith} \cite{bib:1982}.

\noindent[\textbf{1 AB}] Assume $f:S\rightarrow R$ where $f$ is surjective and $R\subseteq S$. Then the mappings of $S$ onto $R$ are precisely the mappings of the form $f\circ h$ where $h:S\rightarrow S$ and $(\forall y\in R)(\exists x\in S)(h(x)\in f^{-1}(\{y\}))$. Smith \cite{bib:1982}.

\noindent[\textbf{1 AC}] Assume $f:S\rightarrow R$ where $R\subseteq S$. Then the retractions of $S$ onto $R$ (that is, the mappings $g : S\rightarrow R$ which are the identity on $R$) are precisely the mappings $f\circ h$ where $h:S\rightarrow S$ and $(\forall y\in R)(h(y)\in f^{-1}(\{y\}))$. Smith \cite{bib:1982}.

\noindent[\textbf{1 AD}] For every class $K$ of similar algebras $SP^{g}K \supseteq P^{g}S\,K$. (For any class $K$ of similar algebras, $S\,K$ is the class of all subalgebras of elements of $K$ (not closed under isomorphism) and $P^{g}K$ is the class of direct products of elements of $K$ not closed under isomorphism). Andreka/Nemeti\index{Author}{Andreka/Nemeti} \cite{bib:1980}.

\noindent[\textbf{1 AE}] There is a set $S$ such that $SVC$ holds with $S$ and such that $\mathcal{P}(S)-\{\emptyset\}$ has a choice function. Blass\index{Author}{Blass} \cite{bib:1979} and note 59. Also see form 191.

\noindent[\textbf{1 AF}] Form 191 + form 8 + form 193. Blass \cite{bib:1979} and note 60.

\noindent[\textbf{1 AG}] Form 192 + form 191. Blass \cite{bib:1979} and note 60.

\noindent[\textbf{1 AH}] $E(V, VI)$: For all infinite cardinals $m,\,2m=m$ implies $m^{2}=m$. Howard/Yorke\index{Author}{Howard/Yorke} \cite{bib:1989} and note 94.

\noindent[\textbf{1 AI}] $E(VI, VII)$: For all infinite cardinals $m,\,m^{2}=m$ implies there is an ordinal $\alpha$ such that $m=\aleph_{\alpha}$. Howard/Yorke \cite{bib:1989} and note 94.

\noindent[\textbf{1 AJ}] Let $\mathcal{F}$ be a family of mappings of a poset $(X, \leq)$ into itself such that $\forall f\in \mathcal{F},\,\forall x\in X,\,f(x)\leq x$. If for some element $e$ of $X$ each chain in $X$ containing $e$ has a lower bound, then the family $\mathcal{F}$ has a common fixed point. Kasahara\index{Author}{Kasahara} \cite{bib:1976}.

\noindent[\textbf{1 AK}] There is an ordinal $\beta$ such that for all cardinals $x,\,\alpha(x+\aleph(x), x\cdot\aleph(x))< \beta$. Truss\index{Author}{Truss} \cite{bib:1973b} and note 98.

\noindent[\textbf{1 AL}] There is an ordinal $\beta$ such that for all cardinals $x,\,\alpha(x, x^{2})<\beta$. Truss \cite{bib:1973b} and note 98.

\noindent[\textbf{1 AM}] There is an ordinal $\beta$ such that for no cardinal $x$ is there a strictly increasing sequence of cardinals $\{y_{\gamma}:\gamma<\beta\}$ such that $y_{\gamma}^{2}=x$. Truss \cite{bib:1973b} and note 98.

\noindent[\textbf{1 AN}] Every lattice ($L, \cap,\,\cup, \subseteq\rangle$ of sets (that is, $L$ is a family of sets and $\cap,\,\cup$ and $\subseteq$ are the usual operations) with a greatest element has a maximal, proper ideal. Bell/Fremlin\index{Author}{Bell}\index{Author}{Bell/Fremlin} \cite{bib:1972}.

\noindent[\textbf{1 AO}] Any continuous surjection between compact spaces, where the range space is Hausdorff, has an irreducible restriction to a closed subset of its domain. Bell \cite{bib:1988} and note 114.

\noindent[\textbf{1 AP}] For any type $\tau$ of the form $\tau =(M)$, every homomorphism of an algebra of type $\tau$ is closed. Diener\index{Author}{Diener} \cite{bib:1989} and note 115.

\noindent[\textbf{1 AQ}] For any type $\tau$ of the form $\tau =(M)$, every homomorphism of an algebra \textbf{A} of type $\tau$ into an algebra \textbf{B} of type $\tau$ maps \textbf{A} onto a subalgebra of \textbf{B}. Diener \cite{bib:1989} and note 105.

\noindent[\textbf{1 AR}] Let $S$ be any family of disjoint non-empty sets and let $\tau =(\bigcup S)$, then any homomorphism from $\mathfrak{P}(\tau;\bigcup S)$ into $\mathfrak{P}(\tau;S)$ maps $\mathfrak{P}(\tau;\bigcup S)$ onto a subalgebra of $\mathfrak{P}(\tau;S)$. Diener\index{Author}{Diener} \cite{bib:1989} and note 115.

\noindent[\textbf{1 AS}] Turinici's Fixed Point Theorem:\index{Subject}{Turinici's fixed point theorem} If $(X, \leq)$ is a directed, partially ordered set and $\tau$ is a topology on $X$ such that

\noindent(a) $(X, \leq)$ is upper semicontinuous with respect to $\tau$,

\noindent(b) Every well ordered subset of $(X, \leq)$ is convergent as a net, then every function $f:X\rightarrow X$ such that $\forall x\in X,\,x\leq f(x)$ has a fixed point. Manka\index{Author}{Manka} \cite{bib:1988b} and note 38.

\noindent[\textbf{1 AT}] For all cardinal numbers $\kappa_{1}$ and $\kappa_{2}$, if a graph is both $\kappa_{1}$ and $\kappa_{2}$ colorable then it is also $\kappa$ colorable for some $\kappa$ such that $\kappa \leq^{\ast}\kappa_{1}$ and $\kappa \leq^{\ast}\kappa_{2}$. ($\leq^{\ast}$ is the surjective cardinal ordering). Galvin/Komj\'{a}th\index{Author}{Galvin/Komj\'{a}th} \cite{bib:1991} and note 118.

\noindent[\textbf{1 AU}] Every graph has an irreducible, good coloring. Galvin/Komj\'{a}th \cite{bib:1991} and note 118.

\noindent[\textbf{1 AV}] Every graph of the form $G(A, H(A))$ has an irreducible, good coloring. ($H$ is the Hartogs aleph function.) Galvin/Komj\'{a}th \cite{bib:1991} and note 118.

\noindent[\textbf{1 AW}] Every graph has a chromatic number. Galvin/Komj\'{a}th \cite{bib:1991} and note 118.

\noindent[\textbf{1 AX}] Birkhoff\index{Author}{Birkhoff} Representation Theorem:\index{Subject}{Birkhoff representation theorem} Every algebra is isomorphic to a subdirect product of subdirectly irreducible algebras. Birkhoff \cite{bib:1944}, Gr\"{a}tzer\index{Author}{Gr\"{a}tzer} \cite{bib:1979} and \cite{bib:1986}.

\noindent[\textbf{1 AY}] For each $\lambda\in\Lambda$ let $S_{\lambda}$ be a subset of some topological space $X_{\lambda}$. Then $\mathrm{cl}(\prod_{\lambda\in\Lambda}S_{\lambda})=\prod_{\lambda\in\Lambda}\mathrm{cl}\,(S_{\lambda})$, where the first ``cl'' is closure in the product topology and the second ``cl'' is closure in $X_{\lambda}$. Schechter\index{Author}{Schechter} \cite{bib:1992}.

\noindent[\textbf{1 AZ}] For each $\lambda\in\Lambda$, let $X_{\lambda}$ be a complete uniform space. Then, the product $\prod_{\lambda\in\Lambda}X_{\lambda}$ equipped with the product uniform structure, is also complete. (See, for example, Kelley\index{Author}{Kelley} \cite{bib:1955} pp 176ff for definitions.) Schechter \cite{bib:1992}.

\noindent[\textbf{1 BA}] (Depends on the prime \emph{p}) Every group has a maximal $p$-subgroup. Howard/Yorke\index{Author}{Howard/Yorke} \cite{bib:1987} and note 24.

\noindent[\textbf{1 BB}] Strong Nielsen-Schreier Theorem\index{Subject}{Nielsen-Schreier theorem} II: If $G$ is a group which is freely generated by $X$ and $U$ is a subgroup of $G$, then there is a subset $A$ of $G$ which freely generates $U$ and is level with respect to $X$. Howard\index{Author}{Howard} \cite{bib:1985} and note 129.

\noindent[\textbf{1 BC}] Assume $P$ is a partially ordered set in which every non-empty chain has an upper bound and $f:P\rightarrow P$ satisfies $\forall x\in f(P)\cup\{a\in P:a$ is an upper bound for some chain in $f(P)\},\, x\leq f(f(x))$. Then $f$ has a fixed apex $u$ (that is, there is some $v\in P$ such that $f(u)=v$ and $f(v)=u$.). Taskovi\'{c}\index{Author}{Taskovi\'{c}} \cite{bib:1988}, \cite{bib:1992a}.

\noindent[\textbf{1 BD}] Let $P$ be a partially ordered set in which every increasing sequence has an upper bound and $f:P\rightarrow P$ satisfies $x\leq f(f(x))$ for all $x\in f(P)\cup\{x\in P:x$ is an upper bound for some increasing sequence in $f(P)\}$. Then $f$ has a fixed apex $u$ (that is, $\exists v\in P$, such that $f(u)=v$ and $f(v)=u$). (For purposes of [1 BD], ``sequence'' means a function whose domain is an ordinal.) Taskovi\'{c} \cite{bib:1992b}.

\noindent[\textbf{1 BE}] Every unique factorization domain can be mapped homomorphically onto a non-zero subdirectly irreducible ring. Rav\index{Author}{Rav} \cite{bib:1988}, Gr\"{a}tzer\index{Author}{Gr\"{a}tzer} \cite{bib:1979} pp 122-124 and note 84.

\noindent[\textbf{1 BF}] $(\mathrm{D}T_{0})$: Every topological space $X$ has a $T_{0}$ subspace that is dense in $X$. $(Y$ is \emph{dense in} $X$ if there is no non-empty open subset $O\subseteq X$ such that $O\cap Y=\emptyset.)$ McCarten\index{Author}{McCarten} \cite{bib:1988} and note 106.

\noindent[\textbf{1 BG}] $(\mathrm{C}T_{0})$: Every topological space $X$ has a $T_{0}$ subspace that is codense in $X$. ($Y$ is \emph{codense in} $X$ if there is no non-empty closed subset $C\subseteq X$ such that $C\cap Y=\emptyset.)$ McCarten \cite{bib:1988} and note 106.

\noindent[\textbf{1 BH}] $(\mathrm{T}T_{0})$: Every topological space $X$ has a $T_{0}$ subspace that is thick in $X$. ($Y$ is \emph{thick in} $X$ if there is no non-empty open and closed (clopen) subset $H\subseteq X$ such that $H\cap Y=\emptyset$.) McCarten \cite{bib:1988} and note 106.

\noindent[\textbf{1 BI}] $(\mathrm{M}T_{0})$: Every topological space has a maximal $T_{0}$ subspace. Schnare\index{Author}{Schnare} \cite{bib:1968} and note 106.

\noindent[\textbf{1 BJ}] In every vector space over the two element field, every generating set contains a basis. Keremedis\index{Author}{Keremedis} \cite{bib:1996a}.

\noindent[\textbf{1 BK}] Every open cover $\mathcal{U}$ of a metric space ($X,d$) can be written as a well ordered union $\bigcup\{U_{\alpha}:\alpha\in\gamma\}$ where $\gamma$ is an ordinal and each $U_{\alpha}$ is pairwise disjoint. Howard/Keremedis/Rubin/Stanley\index{Author}{Howard/Keremedis/Rubin/Stanley} \cite{bib:1997}.

\noindent[\textbf{1 BL}] Every family of closed sets in a $T_{2}$ topological space includes a maximal subfamily with the finite intersection property. Keremedis \cite{bib:1996b} and H. Rubin/J. Rubin\index{Author}{H. Rubin/J. Rubin} \cite{bib:1985}, p 60,
$M14$.

\noindent[\textbf{1 BM}] For every family $\{(X_{i}, T_{i}):i\in k\}$ of $T_{1}$ topological spaces, the box topology contains an element $\prod_{i\in k}O_{i}$, such that $\emptyset \neq O_{i}$ and $O_{i}\neq X_{i}$ for all $i\in k$. Keremedis \cite{bib:1996b} and note 77.

\noindent[\textbf{1 BN}] For every family $\{(X_{i}, T_{i}):i\in k\}$ of disjoint $T_{1}$ topological spaces, the free union contains an open set $O$ such that $\emptyset \neq O\cap X_{i}$ and $O\cap X_{i}\neq X_{i}$ for all $i\in k$. Keremedis \cite{bib:1996b} and note 77.

\noindent[\textbf{1 BO}] For every $T_{1}$ topological space $(X, T)$ and every base (subbase) $B$ for $X$ there exists a $\subseteq$-maximal cellular family $Q\subseteq B$. Keremedis \cite{bib:1996b} and note 77.

\noindent[\textbf{1 BP}] For every dense in itself topological space $(X, T)$ and every base $B$ for $X$ there exists a nest $N$ included in $B$. Keremedis \cite{bib:1996b} and note 77.

\noindent[\textbf{1 BQ}] For every dense in itself topological space $(X, T)$ and every base $B$ for $X$ there exists a tower $R$ included in $B$. Keremedis \cite{bib:1996b} and note 77.

\noindent[\textbf{1 BR}] For every $T_{2}$ space $(X, T)$ and every subbase $\mathcal{B}$ for $X$, if $\mathcal{C}$ and $\mathcal{D}$ are cellular families of $X$ included in $\mathcal{B}$, then there is a cellular family $\mathcal{E} \subseteq \mathcal{B}$ such that $\mathcal{C}$ and $\mathcal{D}$ are equipollent to subsets of $\mathcal{E}$. Keremedis \cite{bib:1996b} and note 77.

\noindent[\textbf{1 BS}] For every Abelian group $(G, +)$, for every $H\subseteq G$, the infinite system of equations
\begin{equation*}
x_{i}+y_{i}=a_{i},\ i\in I,\  a_{i}\in G
\end{equation*}
has a solution in $H$ iff each of its equations has a solution in $H$. Keremedis \cite{bib:1996c}.

\noindent[\textbf{1 BT}] For every Boolean ring $(B, +, \cdot)$ and every $H\subseteq B$ which is closed under + there exists a $\subseteq$-maximal ideal $Q\subseteq B$ such that $H\cap Q=\{0\}$. Keremedis \cite{bib:1996c}.

\noindent[\textbf{1 BU}] In every vector space over $\mathbb{Q}$, every generating set includes a basis. Blass\index{Author}{Blass} \cite{bib:1984a} and note 137.

\noindent[\textbf{1 BV}] If $(X, T)$ is a $T_{4}$ topological space and $G=\{G_{i}:i\in K\},\,|G_{i}|\geq 2$, is a locally finite family of compact sets with non-empty pairwise disjoint interiors, then there is a continuous real valued function $f$ on $X$ which assumes at least two distinct values on each $G_{i}$. Keremedis\index{Author}{Keremedis} \cite{bib:1995} and note 43.

\noindent[\textbf{1 BW}] If $(X, T)$ is a $T_{4}$ topological space and $U=\{U_{i}:i\in K\},\,|U_{i}|\geq 2$, is a locally finite family of open non-empty pairwise disjoint sets, then there is a continuous real valued function $f$ on $X$ which assumes at least two distinct values on each $U_{i}$. Keremedis \cite{bib:1995} and note 43.

\noindent[\textbf{1 BX}] For every Abelian group $G$ and every subgroup $H$ of $G$ there exists a set of representatives for the quotient $G/H$. Keremedis \cite{bib:1996c}.

\noindent[\textbf{1 BY}] If $x$ is a set of sets of cardinality at least 2, then there exists a function $f$ such that for each $u\in x,\,f(u)$ is a finite, non-empty, proper subset of $u$. (Form [1 BY] implies [62 E] $(KW(\infty, <\aleph_{0}))$ and 67 (MC) and $62+67\rightarrow \mathrm{AC}$.) Blass \cite{bib:1979}.

\noindent[\textbf{1 BZ}] Vector Space Multiple Choice: For every family $V=\{V_{i}:i\in K\}$ of non-trivial vector spaces there is a family $F=\{F_{i}:i\in K\}$ such that for each $i\in K,\,F_{i}$ is a non-empty, finite, independent subset of $V_{i}$. Note 127.

\noindent[\textbf{1 CA}] Strong Nielsen-Schreier Theorem\index{Subject}{Nielsen-Schreier theorem} I: If $G$ is a group which is freely generated by $X$ and $U$ is a subgroup of $G$, then there is a subset $A$ of $G$ which freely generates $U$ and has the Nielsen property with respect to $X$. Howard\index{Author}{Howard} \cite{bib:1985} and note 129.

\noindent[\textbf{1 CB}] For every infinite cardinal $\kappa$, if $(X, T)$ is a $T_{2}$ space having a base $B$ such that no $Q\subseteq B,\,|Q|<\kappa$, is a tower, then for every family $\mathcal{D}=\{D_{i}:i\in\kappa\}$ of dense open sets of $X$ there is a filter $F\subseteq T$ such that for every $i\in \kappa$ there is a $b\in F$ with $b\subseteq D_{i}$. Keremedis \cite{bib:1996b} and note 77.

\noindent[\textbf{1 CC}] $PC(\aleph_{0}, \mathrm{odd}, \infty)\, +\,MC_{\omega^{+}}$. $PC(\aleph_{0}, \mathrm{odd}, \infty)$ is ``Every denumerable family of odd sized (finite) sets has an infinite subfamily with a choice function.'' and $MC_{\omega^{+}}$ is ``For every natural number $n\geq 1$ and for every set of non-empty sets, there is a function $f$ such that for each $u\in x,\,f(u)$ is a non-empty finite subset of $u$ such that $|f(u)|$ and $n$ are relatively prime.'' Note 133.

\noindent[\textbf{1 CD}] In a commutative ring, every proper ideal can be extended to a maximal ideal. Krull\index{Author}{Krull} \cite{bib:1929}. ([1 CD] implies [1 B]).

\noindent [\textbf{1 CE}] Every open cover $\mathcal{U}$ of a metric space ($X,d$) can be written as a well ordered union $\bigcup\{U_{\alpha}:\alpha\in\gamma\}$ where $\gamma$ is an ordinal and each $U_{\alpha}$ is discrete. Howard/Keremedis/Rubin/Stanley\index{Author}{Howard/Keremedis/Rubin/Stanley} \cite{bib:1997} and note 141.

\noindent[\textbf{1 CF}] If a topological space is $cwH$, then it is $cwH(B)$ for every basis $B$. Howard/Keremedis/Rubin/Rubin\index{Author}{Howard/Keremedis/Rubin/Rubin} \cite{bib:1997b} and note 141.

\noindent[\textbf{1 CG}] The disjoint union of $cwH$ spaces is $cwH(B)$ for every basis $B$. Howard/Keremedis/Rubin/Rubin \cite{bib:1997b} and note 141.

\noindent[\textbf{1 CH}] If a topological space is $cwN$, then it is $cwH(B)$ for any basis $B$. Howard/Keremedis/Rubin/Rubin \cite{bib:1997b} and note 141.

\noindent\textbf{FORM 2.} Existence of successor cardinals: For every cardinal $m$ there is a cardinal $n$ such that $m<n$ and $(\forall p<n)(p\leq m)$. Tarski\index{Author}{Tarski} \cite{bib:1954a}.

\noindent\textbf{FORM 3.} $2m=m$: For all infinite cardinals $m,\,2m=m$. Tarski \cite{bib:1954a}.

[\textbf{3 A}] For all cardinals $m$ and $n,\,n\leq m$ and $m$ infinite $\rightarrow m+n=m$. Halpern/Howard\index{Author}{Halpern/Howard} \cite{bib:1970}.

\noindent[\textbf{3 B}] $\aleph_{0}\cdot m=m$: For all infinite cardinals $m,\,\aleph_{0}\cdot m=m$. Halpern/Howard \cite{bib:1970}.

\noindent[\textbf{3 C}] Definability of cardinal addition as the least upper bound: (For all cardinals $x,\,y$ and $z$) $(x+y=z$ iff $(x\leq z$ and $y\leq z$ and $(\forall u)(x, y\leq u\rightarrow z\leq u))$. H\"{a}ussler\index{Author}{H\"{a}ussler} \cite{bib:1983} and Tarski \cite{bib:1949a}.

\noindent[\textbf{3 D}] Every pair of cardinals has a least upper bound in the usual cardinal ordering and for all cardinals $\kappa$ and $\gamma$, if $n\kappa \leq\gamma$ for all $n\in \omega$, then $\aleph_{0}\kappa\leq\gamma$. Hickman\index{Author}{Hickman} \cite{bib:1979b}.

\noindent[\textbf{3 E}] $E(IV, V)$: For all infinite cardinals $m,\,m+1=m$ implies $2m=m$. Howard/Yorke\index{Author}{Howard/Yorke} \cite{bib:1989} and note 94.

\noindent[\textbf{3 F}] $E(III, V)$: For every set $X$, if $\mathcal{P}(X)$ is Dedekind infinite then $2|X|=|X|$. Howard/Yorke \cite{bib:1989} and note 94.

\noindent[\textbf{3 G}] $E(II, V)$: For every set $X$, if there is a non-empty family of subsets of $X$ which is linearly ordered by $\subseteq$ and has no maximal element then $2|X|=|X|$. Howard/Yorke \cite{bib:1989} and note 94.

\noindent[\textbf{3 H}] $E(Ia, V)$: For every set $X$, if $2|X|>|X|$ then $X$ is amorphous. Howard/Yorke \cite{bib:1989}, notes 85(H) and 94.

\noindent\textbf{FORM 4.} Every infinite set is the union of some disjoint family of denumerable subsets. (Denumerable means $\cong\aleph_{0}$.) Whitehead\index{Author}{Whitehead} \cite{bib:1902}.

\noindent\textbf{FORM 5.} $C(\aleph_{0}, \aleph_{0}, \mathbb{R})$: Every denumerable set of non-empty denumerable subsets of $\mathbb{R}$ has a choice function. Sageev\index{Author}{Sageev} \cite{bib:1975}.

\noindent\textbf{FORM 6.} $UT(\aleph_{0}, \aleph_{0}, \aleph_{0}, \mathbb{R})$: The union of a denumerable family of denumerable subsets of $\mathbb{R}$ is denumerable. G. Moore\index{Author}{G. Moore} \cite{bib:1982} and note 3.

[\textbf{6 A}] Every uncountable subset of $\mathbb{R}$ contains a condensation point. G. Moore \cite{bib:1982} p 205 and 324, Sierpi\'{n}ski \cite{bib:1918}.

\noindent[\textbf{6 B}] Every uncountable subset of $\mathbb{R}$ has two condensation points. G. Moore \cite{bib:1982} p 205 and 324, Sierpi\'{n}ski \cite{bib:1918}.

\noindent\textbf{FORM 7.} There is no infinite decreasing sequence of cardinals. Schoenflies\index{Author}{Schoenflies} \cite{bib:1908}.

\noindent\textbf{FORM 8.} $C(\aleph_{0}, \infty)$: Every denumerable family of non-empty sets has a choice function. Sierpi\'{n}ski \cite{bib:1918}.

[\textbf{8 A}] $PC(\infty, \infty, \infty)$: Every infinite family of non-empty sets has an infinite subfamily with a choice function. Brunner\index{Author}{Brunner} \cite{bib:1983b}.

\noindent[\textbf{8 B}] $PC(\aleph_{0}, \infty, \infty)$: Every denumerable family of non-empty sets has an infinite subfamily with a choice function. Monro\index{Author}{Monro} \cite{bib:1975}.

\noindent[\textbf{8 C}] Every $\sigma$-compact space is Lindel\"{o}f. Brunner \cite{bib:1982b} and note 43.

\noindent[\textbf{8 D}] Every $\sigma$-compact, locally compact $T_{2}$ space is weakly Lindel\"{o}f. Brunner \cite{bib:1982b} and note 43.

\noindent[\textbf{8 E}] In a metric space every sequentially continuous real valued function is continuous. Brunner \cite{bib:1982c}.

\noindent[\textbf{8 F}] Martin's Axiom\index{Subject}{Martin's axiom} $(\aleph_{0})$: Whenever $(P\leq)$ is a non-empty, ccc quasi-order (ccc means every anti-chain is countable) and $\mathcal{D}$ is a family of $\leq\aleph_{0}$ dense subsets of $P$, then there is a $\mathcal{D}$ generic filter $G$ in $P$. Kunen\index{Author}{Kunen} \cite{bib:1980}, Shannon\index{Author}{Shannon} \cite{bib:1990} and note 47.

\noindent[\textbf{8 G}] Let $R$ be a relation and $R^{\ast}$, the transitive closure of $R$. $(x\,R^{\ast}\,y$ iff there exists a finite ascending $R$-chain connecting $x$ to $y$.) If every ascending $R$-chain is finite, then every ascending $R^{\ast}$-chain is finite. Diener\index{Author}{Diener} \cite{bib:1994}.

\noindent[\textbf{8 H}] A disjoint union of a countable set of compact $T_{2}$ spaces is Lindel\"{o}f. Keremedis\index{Author}{Keremedis} \cite{bib:1996d} and note 132.

\noindent[\textbf{8 I}] A disjoint union of a countable set of Lindel\"{o}f $T_{2}$ spaces is Lindel\"{o}f. Keremedis \cite{bib:1996d} and note 132.

\noindent[\textbf{8 J}] Hemicompact $T_{2}$ spaces are Lindel\"{o}f. (A \emph{hemicompact} space is a space $(X, T)$ that can be written as a disjoint union $\bigcup\{K_{n}:n\in\omega\}$ such that (i) each $K_{n}$ is a compact set in $X$; (ii) $K_{N}\subseteq K_{n+1}$ for all $n\in\omega$; and (iii) for every compact set $K$ in $X$, there is an $n\in\omega$ such that $K\subseteq K_{n}$.) Keremedis \cite{bib:1996d} and note 132.

\noindent[\textbf{8 K}] Hemilindel\"{o}f $T_{2}$ spaces are Lindel\"{o}f. (In [8 J], in the definition of hemicompact replace ``compact'' by ``Lindel\"{o}f'' to get the definition of hemilindel\"{o}f.) Keremedis \cite{bib:1996d} and note 132.

\noindent[\textbf{8 L}] Every second countable topological space is separable. Keremedis \cite{bib:1996d} and note 132.

\noindent[\textbf{8 M}] Form $131+$ form 32. Keremedis \cite{bib:1996d} and note 132.

\noindent[\textbf{8 N}] $DP$: If $(X, T)$ is a topological space having a countable $\pi$-base, then for every family $D=\{D_{i}:i\in\omega\}$ of dense open sets of $X$, there is a countable dense set $S\subseteq X$ such that for all $i\in\omega$ and for all but finitely many $s\in S,\,s\in D_{i}$. Keremedis \cite{bib:1996d} and note 132.

\noindent[\textbf{8 O}] $DUCCC$ + form 131. $DUCCC$ is ``If $X$ is a denumerable family of disjoint topological spaces each having the ccc property, then the disjoint union of the spaces in $X$ has the ccc property.'' $(Y$ has the $ccc$ \emph{property} if every collection of disjoint open sets is countable.) Keremedis \cite{bib:1996d} and note 132.

\noindent[\textbf{8 P}] \emph{DUSS} $+$ form 131. \emph{DUSS} is ``The disjoint union of a denumerable family of separable topological spaces is separable.'' Keremedis \cite{bib:1996d} and note 132.

\noindent[\textbf{8 Q}] $DUSC$ + form 131. $DUSC$ is ``The disjoint union of a denumerable family of second countable topological spaces is second countable.'' Keremedis \cite{bib:1996d} and note 132.

\noindent[\textbf{8 R}] \emph{SCDUL} + form 131. \emph{SCDUL} is ``The denumerable disjoint union of second countable spaces is Lindel\"{o}f.'' Keremedis \cite{bib:1996d} and note 132.

\noindent[\textbf{8 S}] $PC(WO, \infty, \infty)$: Every infinite well ordered set of non-empty sets has an infinite subset with a choice function. $([8\ \mathrm{A}]\rightarrow[8\ \mathrm{S}]\rightarrow[8\ \mathrm{B}])$

\noindent[\textbf{8 T}] $PUT(\aleph_{0}, \infty, \aleph_{0})$: The union of a denumerable set of pairwise disjoint non-empty sets has a denumerable subset. $([8\ \mathrm{A}]\rightarrow[8\ \mathrm{T}]\rightarrow[8\ \mathrm{B}])$

\noindent[\textbf{8 U}] $PUT(WO, \infty, WO)$: The union of an infinite well ordered set of pairwise disjoint non-empty sets has an infinite well ordered subset. $([8\ \mathrm{U}]\leftrightarrow[8\ \mathrm{T}])$

\noindent\textbf{FORM 9.} Finite $\Leftrightarrow$ Dedekind finite: $W_{\aleph_{0}}$ (Jech\index{Author}{Jech} \cite{bib:1973b}): $E(I, IV)$ (Howard/Yorke\index{Author}{Howard/Yorke} \cite{bib:1989}): Every Dedekind finite set is finite. Dedekind\index{Author}{Dedekind} \cite{bib:1888} and note 94.

[\textbf{9 A}] For all infinite $x$, if $x$ is Dedekind finite, then $\mathcal{P}(x)$ is is Dedekind finite). Notes 8 and 94.

\noindent[\textbf{9 B}] $E(II, IV)$ (Howard/Yorke\index{Author}{Howard/Yorke} \cite{bib:1989}): $(\forall x)(x$ is $T$-finite if and only if $x$ is Dedekind finite). Howard/Yorke \cite{bib:1989} and note 94.

\noindent[\textbf{9 C}] The image of a Dedekind finite set is Dedekind finite. Jech\index{Author}{Jech} \cite{bib:1973b} p 161 prob 11.2, notes 11 and 94.

\noindent[\textbf{9 D}] The union of a Dedekind finite family of finite sets is Dedekind finite. Jech \cite{bib:1973b} p 161 prob 11.4, notes 12 and 94.

\noindent[\textbf{9 E}] No Dedekind finite set can be mapped onto $\aleph_{0}$. (See form 160.) Monro\index{Author}{Monro} \cite{bib:1975} and note 94.

\noindent[\textbf{9 F}] There is no subset $A$ of the class of Dedekind finite cardinals such that $A$ under the usual cardinal ordering is order isomorphic to $\langle \mathbb{R},\,\leq\rangle$. Tarski\index{Author}{Tarski} \cite{bib:1965}, Jech \cite{bib:1973b} p 161 prob 6, and note 94.

\noindent[\textbf{9 G}] Every infinite set has a denumerable subset. Cantor\index{Author}{Cantor} \cite{bib:1895} and G. Moore\index{Author}{G. Moore} \cite{bib:1982} p 9.

\noindent[\textbf{9 H}] If $A$ is Dedekind finite and $B$ is Dedekind infinite, then $A\precsim B$. G. Moore \cite{bib:1982} p 28 and note 94.

\noindent[\textbf{9 I}] If $\emptyset \in K$ and for all $A\in K$ and for all $b\not\in A,\,A\, \cup\,\{b\}\in K$, then every Dedekind finite set is in $K$. G. Moore \cite{bib:1982} p 28 and note 94.

\noindent[\textbf{9 J}] The union of a Dedekind infinite family of non-empty sets is Dedekind infinite. G. Moore \cite{bib:1982} p 130 and note 94.

\noindent[\textbf{9 K}] If $A$ is uncountable and $B$ is countable, then $|A-B|=|A|$. G. Moore \cite{bib:1982} p 201.

\noindent[\textbf{9 L}] If $A$ is uncountable and $B$ is countable, then $|A\cup B|=|A|$. G. Moore \cite{bib:1982} p 202.

\noindent[\textbf{9 M}] If $\aleph_{0}<|A|$ and $\aleph_{0}=|B|$, then $\aleph_{0}<|A-B|$. G. Moore \cite{bib:1982} p 202.

\noindent\textbf{[9 N]} If $\aleph_{0}\leq^{\ast}m$, then $\aleph_{0}\leq m$. Lindenbaum/Tarski\index{Author}{Lindenbaum/Tarski} \cite{bib:1926}, G. Moore \cite{bib:1982} p 216, and note 69.

\noindent[\textbf{9 O}] $m\leq\aleph_{0}$ or $\aleph_{0}\leq m$. Lindenbaum/Tarski \cite{bib:1926} and G. Moore \cite{bib:1982} p 216.

\noindent[\textbf{9 P}] If $m+\aleph_{0}<n+\aleph_{0}$, then $m<n$. Lindenbaum/Tarski \cite{bib:1926} and G. Moore \cite{bib:1982} p 216.

\noindent[\textbf{9 Q}] If $m+\aleph_{0}=n+\aleph_{0}$, then $m=n$ or $m,\,n\leq\aleph_{0}$. Lindenbaum/Tarski \cite{bib:1926} and G. Moore \cite{bib:1982} p 216.

\noindent[\textbf{9 R}] $E(III, IV)$: Every Dedekind finite set has a Dedekind finite power set. Howard/Yorke \cite{bib:1989} and note 94.

\noindent[\textbf{9 S}] $E(Ia, IV)$: Every Dedekind finite set is amorphous. Howard/Yorke \cite{bib:1989}, notes 57 and 94.

\noindent[\textbf{9 T}] Restricted Choice: For every infinite set $X$ there is an infinite $Y\subseteq X$ such that the collection of non-empty subsets of $Y$ has a choice function. De la Cruz/Di Prisco\index{Author}{De la Cruz/Di Prisco} \cite{bib:1996}.

\noindent\textbf{FORM 10.} $C(\aleph_{0}, <\aleph_{0})$: Every denumerable family of non-empty finite sets has a choice function.

[\textbf{10 A}] $UT(\aleph_{0}, <\aleph_{0}, \aleph_{0})$: The union of denumerably many pairwise disjoint finite sets is denumerable. Shannon\index{Author}{Shannon} \cite{bib:1988} p 569.

\noindent[\textbf{10 B}] Tychonoff's\index{Subject}{Tychonoff's theorem} Compactness Theorem for Countable Families of Finite $T_{2}$ Spaces: The product of a countable family of finite $T_{2}$ spaces is compact. Krom\index{Author}{Krom} \cite{bib:1981}.

\noindent[\textbf{10 C}] Every set mapping $f:X\rightarrow[X]^{WO}$ on a Dedekind finite, infinite set $X$ has arbitrarily large, finite free subsets. Brunner\index{Author}{Brunner} \cite{bib:1989}, notes 22 and 94.

\noindent[\textbf{10 D}] In every Hilbert space with a Dedekind finite orthonormal basis, the closed unit ball is sequentially compact. Brunner \cite{bib:1983b}
and note 94.

\noindent[\textbf{10 E}] $PC(\aleph_{0}, <\aleph_{0}, \infty)$: Every denumerable family of finite sets has an infinite subfamily with a choice function. Brunner
\cite{bib:1983b}.

\noindent[\textbf{10 F}] K\"{o}nig's Lemma:\index{Subject}{K\"{o}nig's lemma} Every $\omega$ tree has an infinite chain. Pincus\index{Author}{Pincus} \cite{bib:1972c}, note 21, and note 35.

\noindent[\textbf{10 G}] For all metric spaces $X$ and $Y$ and continuous functions $f$ from $X$ onto $Y$ such that $f^{-1}(y)$ is finite for all $y\in Y$, if $X$ has a dense Dedekind finite subset, then so does $Y$. Brunner \cite{bib:1982d} and note 94.

\noindent[\textbf{10 H}] $\sigma$-compact, paracompact $T_{2}$ spaces are weakly Lindel\"{o}f\index{Author}{Lindel\"{o}f}. Brunner \cite{bib:1982b} and note 43.

\noindent[\textbf{10 I}] Assume $(P, \leq)$ is a partial order such that $P$ is a denumerable union of finite sets and contains a non-empty subset $Q$ relative to which depths are defined for all $x\in \mathrm{P}$. If all $Q$-strong antichains of $P$ are finite then for each denumerable family $\mathcal{D}$ of dense sets there is a $\mathcal{D}$ generic filter. Shannon \cite{bib:1990} and note 47.

\noindent[\textbf{10 J}] Let $S$ be a denumerable set of propositional sentences over an infinite set of propositional letters (in a propositional language where conjunction and disjunction are operators on unordered finite sets of sentences). Then $S$ has a model if and only if every finite subset of $S$ has a model. Lolli\index{Author}{Lolli} \cite{bib:1977}.

\noindent[\textbf{10 K}] Let $S$ be a consistent denumerable set of sentences in a first order language (where conjunction and disjunction are operators on unordered finite sets of sentences). Then $S$ has a model. Note 19.

\noindent[\textbf{10 L}] Let $S$ be a denumerable set of sentences in a first order language (where conjunction and disjunction are operators on unordered finite sets of sentences) such that every finite subset has a model. Then $S$ has a model. Note 19.

\noindent[\textbf{10 M}] $PUT(\aleph_{0}, <\aleph_{0}, \aleph_{0})$: The union of denumerably many pairwise disjoint finite sets has a denumerable subset.

\noindent[\textbf{10 N}] $PUT(WO, <\aleph_{0}, WO)$: The union of an infinite well ordered set of pairwise disjoint finite sets has an infinite well ordered subset.

\noindent[\textbf{10 O}] $PC(WO, <\aleph_{0}, \infty)$: Every infinite well ordered family of finite sets has an infinite subset with a choice function.

\noindent\textbf{FORM 11.} $\mathbb{R}$ is hereditarily Lindel\"{o}f. Lindel\"{o}f
\cite{bib:1905}.

\noindent\textbf{FORM 12.} $\mathbb{R}$ is Lindel\"{o}f. Young\index{Author}{Young} \cite{bib:1903}.

\noindent\textbf{FORM 13.} Every Dedekind finite subset of $\mathbb{R}$ is finite. G. Moore\index{Author}{G. Moore} \cite{bib:1982} and note 94.

[\textbf{13 A}] $\mathbb{R}$ has no dense Dedekind finite subset. Brunner \cite{bib:1982d} and note 94.

\noindent[\textbf{13 B}] There is a separable topological space $X$ and a set $A\subseteq \mathbb{R}$ such that $X^{A\times \omega}$ has no dense, Dedekind finite subset. Brunner\index{Author}{Brunner} \cite{bib:1982d} and note 94.

\noindent\textbf{FORM 14.} BPI: Every Boolean algebra has a prime ideal. M. Stone\index{Author}{M. Stone} \cite{bib:1936}.

[\textbf{14 A}] Ultrafilter Theorem: Every proper filter over a set $S$ (in $\mathcal{P}(S)$) can be extended to an ultrafilter. Jech\index{Author}{Jech} \cite{bib:1973b} p 17.

\noindent[\textbf{14 B}] Stone Representation Theorem:\index{Subject}{Stone representation theorem} Every Boolean algebra is isomorphic to a set algebra. Jech \cite{bib:1973b} p 27 prob 2.12 and M. Stone \cite{bib:1936}.

\noindent[\textbf{14 C}] Every commutative ring with unit has a prime ideal. Scott\index{Author}{Scott} \cite{bib:1954}.

\noindent[\textbf{14 D}] Every ring with unit has a (two sided) prime ideal. (An ideal $P$ in a ring $R$ is prime if $P\neq R$ and for any ideals $A$ and $B$ in $R$, if $AB\subseteq P$ then $A\subseteq P$ or $B\subseteq P.)$ Blass\index{Author}{Blass} \cite{bib:1986} and Banaschewski\index{Author}{Banaschewski} \cite{bib:1985}.

\noindent[\textbf{14 E}] Almost Maximal Ideal Theorem:\index{Subject}{almost maximal ideal theorem} Every bounded, non-trivial, distributive lattice has an almost maximal ideal. (Definitions: If $D$ is a bounded non-trivial distributive lattice, $I(D)$ is the lattice of ideals of $D$. An \emph{almost maximal ideal} in $D$ is an ideal $P$ of $D$ such that $s(P)=P$ where $s$ is the operation on $I(D)$ defined as follows: If $e$ is the unit of $I(D)$, (i) For $a\in I(D)$ we say $x\in I(D)$ is $a$-\emph{small} if $x\vee b=e$ implies $a\vee b=e$ for all $b\in I(D)$. (ii) For any $a\in I(D)$, $s(a)$ is the join of all a-small elements of $I(D))$. Banaschewski/Harting\index{Author}{Banaschewski/Harting} \cite{bib:1985} and Blass \cite{bib:1986}.

\noindent[\textbf{14 F}] Every non-trivial complete distributive lattice with compact unit has a prime ideal. Banaschewski \cite{bib:1985} and note 29.

\noindent[\textbf{14 G(n)}] $(n\in \omega, n\geq 3)\,P_{n}$: For every graph $G$ if every finite subgraph of $G$ is $n$-colorable then $G$ is $n$-colorable. L\"{a}uchli\index{Author}{L\"{a}uchli} \cite{bib:1971} and note 117.

\noindent[\textbf{14 H}] Compactness Theorem for First Order Logic: If $\Sigma$ is a set of formulas in a first order language such that every finite subset of $\Sigma$ has a model then $\Sigma$ has a model. Henkin\index{Author}{Henkin} \cite{bib:1954a} and Jech \cite{bib:1973b} p 17.

\noindent[\textbf{14 I}] Consistency Principle:\index{Subject}{consistency principle} For every binary mess $M$ on a set $S$ there is a function $f$ on $S$ which is consistent with $M$. Jech \cite{bib:1973b} p 17 and note 109.

\noindent[\textbf{14 J}] Tychonoff's\index{Subject}{Tychonoff's theorem} Compactness Theorem for Compact Hausdorff Spaces: The product of compact Hausdorff spaces is compact. (For many compactness theorems $T$, it is clear that [14 J] implies $T$ implies [14 K].) {\L}o\'{s}/Ryll-Nardzewski\index{Author}{{\L}o\'{s}/Ryll-Nardzewski}
\cite{bib:1954} and Rubin/Scott\index{Author}{Rubin/Scott} \cite{bib:1954}.

\noindent[\textbf{14 K}] Compactness of the Generalized Cantor Space: $\{0,1\}^{I}$ is compact for any I. (For many compactness theorems $T$, it is clear that [14 J] implies $T$ implies [14 K].) Jech \cite{bib:1973b} p 28, prob 17 and Mycielski\index{Author}{Mycielski} \cite{bib:1964}.

\noindent[\textbf{14 L}] Stone-\v{C}ech Compactification Theorem:\index{Subject}{Stone-\v{C}ech compactification theorem} For every topological space $X$ there is a pair $(Y, j)$ such that $Y$ is a compact topological space and $j$ is a continuous surjection from $X$ onto $Y$ such that for all topological spaces $Z$ such that $Z$ is compact and completely regular and for all continuous $f:X\rightarrow Z$, there is a unique continuous $h:Y\rightarrow Z$ such that $f=h\circ j$. Rubin/Scott \cite{bib:1954} and note 71.

\noindent[\textbf{14 M}] R. Cowen's Generalization of Konig's Lemma:\index{Subject}{K\"{o}nig's lemma} Let $T$ be a collection of locally finite trees such that for any finite set of levels of $T$, there is a consistent set of vertices piercing those levels. Then there is a consistent set of vertices piercing the entire set of levels of T. Cowen\index{Author}{Cowen} \cite{bib:1977b} and note 21.

\noindent[\textbf{14 N(n)}] $(n\in\omega,\,n\geq 2)$ [14 M] restricted to families of trees where no tree in the family has a vertex with more than $n$ successors. Cowen\index{Author}{Cowen} \cite{bib:1977b}, Howard\index{Author}{Howard} \cite{bib:1984a}, and note 21.

\noindent[\textbf{14 O}]\index{Subject}{Rado's selection lemma} Let $\{A(i):i\in I\}$\index{Subject}{FIP extension principle} be a collection of finite sets and $R$ a symmetric binary relation on $\bigcup_{i\in I}A(i)$. Suppose that for every finite $W\subseteq I$, there is an $R$-consistent choice function for $\{A(i):i\in W\}$. Then there is an $R$-consistent choice function for $\{A(i):i\in I\}$. {\L}o\'{s}/Ryll-Nardzewski\index{Author}{{\L}o\'{s}/Ryll-Nardzewski}
 \cite{bib:1951} and note 109.

\noindent[\textbf{14 P($n$)}] (For $n\geq 3)$ [14 O] restricted to $\{A(i):i\in I\}$ such that for all $i\in I,\,|A(i)|\leq n$. Cowen \cite{bib:1977b} and note 109. (For $n=2$ see form 141.)

\noindent[\textbf{14 Q}] Alaoglu's Theorem:\index{Subject}{Alaoglu's theorem} The unit ball is compact in the weak$^{\ast}$ topology in the dual of a Banach space. (See [52 E].) Johnstone\index{Author}{Johnstone} \cite{bib:1984} and note 23.

\noindent[\textbf{14 R}] Every coherent local lattice is spatial. Banaschewski\index{Author}{Banaschewski} \cite{bib:1981} and note 29.

\noindent[\textbf{14 S}] Every compact, regular local lattice is spatial. Banaschewski \cite{bib:1981} and note 29.

\noindent[\textbf{14 T}] Every regular commutative ring with unit has a maximal ideal. Banaschewski \cite{bib:1983} and note 29.

\noindent[\textbf{14 U}] Every closed proper ideal in a Gelfand algebra is contained in some maximal ideal. Banaschewski \cite{bib:1983} and note 29.

\noindent[\textbf{14 V}] FIP Extension Principle: Let $\mathcal{F}\subseteq 2^{X}$ have the finite intersection property, then $\mathcal{F}\subseteq U$ for some ultrafilter $U$ on $X$. Schechter\index{Author}{Schechter} \cite{bib:1996a}.

\noindent[\textbf{14 W}] Universal Subnet Theorem: Any net has a universal subnet. Schechter \cite{bib:1996a} and note 31.

\noindent[\textbf{14 X}] Cowen-Engler Lemma:\index{Subject}{Cowen-Engler lemma} Let $\Lambda$ and $X$ be sets. Let $\Phi$ be a family of functions from subsets of $\Lambda$ into X. Assume
\begin{enumerate}
\item[(a)] $\Phi(\lambda)=\{f(\lambda):f\in\Phi$ and $\lambda\in$ dom $f\}$ is finite for each $\lambda\in\Lambda$,
\item[(b)] each finite $S\subseteq\Lambda$ is the domain of at least one $f\in\Phi$, and
\item[(c)] $\Phi$ has finite character, (that is, a function $f$ from a subset of $\Lambda$ into $X$ is in $\Phi$ if and only if every finite subset of $f$ is in $\Phi$).
\end{enumerate}

\noindent Then at least one $f\in\Phi$ has domain $\Lambda$. Schechter \cite{bib:1996a}. (See [14 I].)

\noindent[\textbf{14 Y}] Form 99 + Form 62. Note 33.

\noindent[\textbf{14 Z}] Tychonoff's\index{Subject}{Tychonoff's theorem} Compactness Theorem for Families of Finite Spaces: If $\{X_{i}:i\in I\}$ is a family of finite topological spaces, then $\prod_{i\in I}X_{i}$ is compact. (Cleary [14 Z] implies [14 K] and is implied by [14 J].) Jech\index{Author}{Jech} \cite{bib:1973b} p 28 prob 17, A.A. Lewis\index{Author}{Lewis} \cite{bib:1983}, and Wolk\index{Author}{Wolk} \cite{bib:1967}.

\noindent[\textbf{14 AA}] Alexander's Subbase Lemma:\index{Subject}{Alexander's subbase lemma} A topological space is compact if and only if there is a subbase $S$ for the topology such that every cover by elements of $S$ has a finite subcover. Kelley\index{Author}{Kelley} \cite{bib:1955} p 143, A.A. Lewis \cite{bib:1983}, Wolk \cite{bib:1967}, and note 34.

\noindent[\textbf{14 AB}] Tychonoff's Compactness Theorem for Sober Spaces: A product of compact, sober topological spaces is compact. Blass\index{Author}{Blass} \cite{bib:1986}, Johnstone \cite{bib:1984} and note 37.

\noindent[\textbf{14 AC}] (Depends on $k, n\in \omega$ with $k\geq 2,\,n\geq 2$ and $n+k\geq 5$.) $P(k, n)$: If $X$ is a set and $P$ is a property of subsets of $X$ of $n$ character (that is, $\forall y\subseteq X(P(y)$ iff $\forall z\subseteq y(|z|\leq n\rightarrow P(z))))$, then if every finite subset of $X$ can be partitioned into $k$ or fewer $P$-sets (that is, sets $z$ such that $P(z))$, then $X$ can be partitioned into $k$ or fewer $P$-sets. Cowen\index{Author}{Cowen} \cite{bib:1982}.

\noindent[\textbf{14 AD}]\index{Subject}{Rado's selection lemma} (Depends on $k\in\omega, k\geq 2)\ P(k)$: If $X$ is a set and $P$ is a property of subsets of $X$ of finite character (that is, $\forall y\subseteq X(P(y)$ iff $\forall z\subseteq y(z$ finite $\rightarrow P(z))))$, then if every finite subset of $X$ can be partitioned into $k$ or fewer $P$-sets (that is, sets $z$ such that $P(z)$), then $X$ can be partitioned into $k$ or fewer $P$-sets. Cowen \cite{bib:1982}.

\noindent[\textbf{14 AE}] For any distributive lattice $D$ with zero and unit, any ideal $J\subseteq D$ and $a\not\in J$, there is a prime filter $F$ such that $a\in F$ and $F\cap J=\emptyset$. Banaschewski\index{Author}{Banaschewski} \cite{bib:1981}.

\noindent[\textbf{14 AF}] Disjoint Transversal axiom\index{Subject}{linking axiom} (2,3): If $X$ is a set, $\mathcal{U}$ is a subset of $\{y\subseteq X:|y|\leq 2\}$ and $\mathcal{V}$ is a subset of $\{y\subseteq X:|y|\leq 3\}$ and if for each two finite subsets $\mathcal{U}_{0}\subseteq \mathcal{U}$ and $\mathcal{V}_{0}\subseteq \mathcal{V}$ there are transversals $Y_{1}$ and $Y_{2}$ for $\mathcal{U}_{0}$ and $\mathcal{V}_{0}$ respectively such that $Y_{1}\cap Y_{2}=\emptyset$, then there are disjoint transversals $X_{1}$ and $X_{2}$ for $\mathcal{U}$ and $\mathcal{V}$ respectively. $(Y\subseteq X$ is a transversal for $\mathcal{U}$ if for all $y\in \mathcal{U},\, y\cap Y\neq(\emptyset.)$ Schrijver\index{Author}{Schrijver} \cite{bib:1978}.

\noindent[\textbf{14 AG}] Topological Rado Lemma: Let $\{T_{i}:i\in I\}$ be a family of compact Hausdorff spaces and suppose for each finite $F\subseteq I$, there is a set $\Phi(F)$ of partial choice functions on $I$ whose domains include $F$ and such that (a) $F_{1}\subseteq F_{2}$ implies $\Phi(F_{2})\subseteq\Phi(F_{1})$ and (b) the set $\{\phi|F:\phi\in\Phi(F)\}$ is closed in $\prod\{T_{i}:i\in F\}$ in the product topology. Then there is a choice function $f$ on $I$ satisfying for any finite $F\subseteq I,\,\exists\phi\in\Phi(F)$ such that $f$ and $\phi$ coincide on $F$. Rav\index{Author}{Rav} \cite{bib:1977}.

\noindent[\textbf{14 AH}] Principle of Consistent Choices: Let $\{T_{i}:i\in I\}$ be a family of compact Hausdorf spaces and suppose $R$ is a symmetric binary relation defined for all pairs of elements belonging to distinct $T_{i}$ and such that for all $i\neq j$ in $I$, the set $E(i, j)=\{(x, y):x\ R\ y$ and $x\in T_{i}$ and $y\in T_{j}\}$ is closed in $T_{i}\times T_{j}$. Furthermore suppose that for every finite subset $F\subseteq I$, with at least two elements, $\exists\phi\in\prod\{T_{i}:i\in F\}$ such that $\phi(i)\ R\ \phi(j)$ for all $i\neq j$ in $F$. Then $\exists f\in\prod\{T_{i}:i\in I\}$ such that $f(i)\ R\ f(j)$ for all $i\neq j$ in $I$. Rav \cite{bib:1977}.

\noindent[\textbf{14 AI}] Engelers Lemma:\index{Subject}{Engelers lemma} Let $\mathcal{E}$ be a family of functions whose domains are subsets of a set $S$ with range $\subseteq\{0,1\}$ and such that
\begin{enumerate}
\item[(a)] For every finite $F\subseteq S,\,\exists\phi\in \mathcal{E}$ with domain F.
\item[(b)] The restriction of any function in $\mathcal{E}$ to any subset of its domain is in $\mathcal{E}$.
\item[(c)] If $\phi$ is any function with domain $\phi\subseteq S$ such that the restriction of $\phi$ to any finite subset of its domain is in $\mathcal{E}$ then $\phi\in \mathcal{E}$.
\end{enumerate}
Then $\exists f\in \mathcal{E}$ whose domain is $S$. Rav \cite{bib:1977}.

\noindent[\textbf{14 AJ}] Robinson's Valuation Lemma:\index{Subject}{Robinson's valuation lemma} Let $\{\phi_{x}:x\in X\}$ be a family of partial valuations on a set $S$. $(\phi$ is a partial valuation on $S$ if the domain of $\phi \subseteq S$ and the range of $\phi \subseteq\{0,1\}.)$ Suppose $\mathfrak{B}$ is a filter base on $X$ such that for every finite $F\subseteq S$ and for every $B\in \mathfrak{B}$ there is $x\in B$ such that $F\subseteq$ dom $\phi_{x}$. Then there is a total valuation $f$ on $S$ (that is, dom $f=S)$ such that for all finite $F\subseteq S$ and all $B\in \mathfrak{B},\,\exists x\in B$ such that $F\subseteq$ dom $\phi_{x}$ and $f|F=\phi_{x}|F$. $(\mathfrak{B}$ is a filter base means $\forall B_{1},\,B_{2}\in \mathfrak{B},\,\exists B_{3}\in \mathfrak{B}$ such that $B_{3}\subseteq B_{1}\cap B_{2}$.) Rav \cite{bib:1977} and Cowen \cite{bib:1973}. Compare with [14
I].

\noindent[\textbf{14 AK}] Let $X$ be a set and $\mathcal{D}=\{Y_{i}:i\in I\}$ a family of finite pairwise disjoint subsets of $\mathcal{P}(X)$. Suppose for every finite subset $\{a, b, \ldots, k\}$ of $I,\,\exists s\in I$ such
that $Y_{s}\subseteq Y_{a}\wedge Y_{b}\wedge\ldots\wedge Y_{k}$ (If $Y$ and $Z$ are finite pairwise disjoint sets then $Y\wedge Z= \{A \cap B:A\in Y$ and $B\in Z$ and $A \cap B\neq\emptyset\}.)$ Then there is a choice function $\mu$ whose domain includes $\mathcal{D}$ such that the family $\{\mu(Y_{i}):i\in I\}$ (where $\mu(Y_{i})\in Y_{i})$ of subsets of $X$ is closed under finite intersections. Rav\index{Author}{Rav} \cite{bib:1977}.

\noindent[\textbf{14 AL}]\index{Subject}{Rado's selection lemma} The radical of a proper ideal $A$ in a commutative ring $R$ with identity is the intersection of all prime ideals in $R$ containing \emph{A}. (The \emph{radical} of $A= \{x\in R$: for some $m\in \omega,\, x^{m}\in A\}.)$ Rav \cite{bib:1977} and note 80.

\noindent[\textbf{14 AM}] Let $A$ be a subring of a commutative ring $R$ and $p$ a prime ideal in $A$ such that $p=Rp\cap A$. Then there is a prime ideal $J$ in $R$ such that $p=J\cap A$. Rav \cite{bib:1977} and note 80.

\noindent[\textbf{14 AN}] In every commutative ring with identity, every proper ideal is included in some prime ideal. Rav \cite{bib:1977}.

\noindent[\textbf{14 AO}] Suppose $R$ is a commutative ring, $A$ is a proper ideal in $R$, and $S$ is a multiplicative semigroup in $R$ not meeting $A$. Then there is a prime ideal $p$ in $R$ such that $A\subseteq p$ and $p\cap S=\emptyset$. Rav \cite{bib:1977} and note 80.

\noindent[\textbf{14 AP}] Suppose to each finite subset $F$ of a set I there corresponds a set $\Phi(F)$ of functions whose domains are subsets of $I$ including $F$ and such that (a) $F_{1}\subseteq F_{2}$ implies $\Phi(F_{2})\subseteq\Phi(F_{2})$ and (b) $\forall i\in I,\,\big\{\phi(i):\phi\in\bigcup\{\Phi(F):F$ finite and $F\subseteq I\}\big\}$ is finite. Then there is a function $f$ with domain $I$ such that for all finite $F\subseteq I,\,\exists\phi\in\Phi(F)$ such that $\phi$ and $f$ coincide on F. Rav \cite{bib:1977}, Cowen\index{Author}{Cowen} \cite{bib:1973} and note 80.

\noindent[\textbf{14 AQ}] Strong Rado Lemma 1: Let $T$ be a family of finite sets and let $\mathcal{F}$ be a collection of finite subsets of $T$ satisfying

\noindent(1) $(\forall K\in T)(\exists S\in \mathcal{F})(K\in S)$

\noindent(2) $(\forall S_{1}, S_{2}\in \mathcal{F})(\exists S_{3}\in \mathcal{F})(S_{1}\cup
S_{2}\subseteq S_{3})$

\noindent Assume there is a function $\phi$ with domain $\mathcal{F}$ such that for each $S\in \mathcal{F},\,\phi(S)= \phi_{S}$ is a function with domain $S$ and satisfies $(\forall K\in S)(\phi_{S}(K)\in K)$. Then there is a function $f$ with domain $T$ and with the following property: For all $S\in \mathcal{F}$ there is an $S^{\prime}\in \mathcal{F}$ such that $S\subseteq S^{\prime}$ and $f$ and $\phi_{S^{\prime}}$ agree on $S$. Howard\index{Author}{Howard} \cite{bib:1993}.

\noindent[\textbf{14 AR}] Patching Lemma: Let $L$ be a local system of $S$ (that is, $L\subseteq \mathcal{P}(S)$ and for each finite $E\subseteq S$, there exists $H\in L$ such that $E\subseteq H$). Let $F$ be a set and $n$ a positive integer. Suppose that for each $H\in L$, there is a function $f_{H}:H^{n}\rightarrow F$ and $\{f_{H}(x):H\in L\}$ is finite for each $x\in S^{n}$. Then there is a function $f:S^{n}\rightarrow F$ such that for any finite subset $K\subseteq S^{n}$, there is an $H\in L$ with $K\subseteq H^{n}$ and $f$ and $f_{H}$ agree on K. Hickin/Plotkin\index{Author}{Hickin/Plotkin} \cite{bib:1976} and note 83.

\noindent[\textbf{14 AS}] Let $\mathfrak{A}$ and $\mathfrak{B}$ be (universal) algebras of type $r$ where $\mathfrak{A}$ and $\mathfrak{B}$ have underlying sets $A$ and $B$ respectively and $B$ is finite. Suppose for every finite $C\subseteq A$ there is a partial $r$ homomorphism from $\mathfrak{A}$ into $\mathfrak{B}$ with domain $C$. Then there is an $r$ homomorphism from $\mathfrak{A}$ into $\mathfrak{B}$ . van Benthem\index{Author}{Benthem} \cite{bib:1975}.

\noindent[\textbf{14 AT}] Assume $D$ and $E$ are sets and for all $d\in D,\,d$ is finite. Suppose that for every finite $F\subseteq D$, there is an $S$ such that $(\forall f\in F)(f\cap S\in E)$. Then there is an $S$ such that $(\forall d\in D)(d\cap S\in E)$. van Benthem \cite{bib:1975}.

\noindent[\textbf{14 AU}] If $Z$ is a real vector lattice with order unit 1 then the set of all positive linear functionals $u$ on $Z$ with $u(1)=1$ is a convex set that has an extreme point. See H. Rubin/J. Rubin\index{Author}{H. Rubin/J. Rubin} \cite{bib:1985}, pages 177 and 182 for definitions. Also see Gardiner\index{Author}{Gardiner} \cite{bib:1974}.

\noindent[\textbf{14 AV}] Every Boolean algebra has a 2-valued measure. Pincus\index{Author}{Pincus} \cite{bib:1972c} and note 147.

\noindent[\textbf{14 AW}] The Compactness Theorem for Propositional Logic: If $\Sigma$ is a set of formulas in a propositional language such that every finite subset of $\Sigma$ is satisfiable, then $\Sigma$ is satisfiable. Henkin\index{Author}{Henkin} \cite{bib:1954b}.

\noindent[\textbf{14 AX}] Structural Completeness Theorem for Propositional Logic: Assume $At$ is a set of propositional variables and $S$ is the set of propositional formulas generated by \emph{At}. Assume $X\subseteq S$ and $a\in S$ satisfy the following: For any function $e:At\rightarrow S$ if $Cn(\{b[e]:b\in X\})\subseteq Cn(\emptyset)$ then $a[e]\in Cn(\emptyset)$. Then $a\in Cn(X)$. $(Cn(Y)$ is the deductive closure of $Y\ \cup$ all substitutions of axioms for the classical predicate calculus and for $a\in S,\,a[e]$ is $a$ with each propositional variable $p$ replaced by $e(p).)$ Pogorzelski/Prucnal\index{Author}{Pogorzelski/Prucnal} \cite{bib:1974}.

\noindent[\textbf{14 AY}] Let $F$ be a class of partial functions on $I$ of finite character such that (a) $\{f(\nu):f \in F\}$ is finite for each $\nu\in I$ and (b) for each finite $W \subseteq I$, there exists $f\in F$ with dom$(f)=W$. Then $F$ contains a function with domain $I$. ($F$ is of finite character means $f\in F$ iff for all finite $w\subseteq I, f|w\in F.)$ Cowen\index{Author}{Cowen} \cite{bib:1973}.

\noindent[\textbf{14 AZ}] Let $\Phi=\{\phi_{t}\}_{t\in T}$ be a set of partial valuations on a set $S$ (that is, dom$(\phi_{t})\subseteq S$ and range$(\phi_{t})\subseteq\{0,1\})$ such that for all finite $U\subseteq S$ there is a $t \in T$ for which $U\subseteq$ dom$(\phi_{t})$. Then there is a valuation $\Psi$ with dom$(\Psi)=S$ and such that for every finite $U\subseteq S$, there is a $t\in T$ with $U\subseteq$ dom$(\phi_{t})$ and $\Psi|U=\phi_{t}|U$. Cowen \cite{bib:1973}.

\noindent[\textbf{14 BA}] Monteiro's\index{Author}{Monteiro} Theorem\index{Subject}{Monteiro's theorem} for Complete Atomic Boolean Algebras: Let $K$ be a complete atomic Boolean algebra, $A$ a subalgebra of the Boolean algebra $B$ and $f$ a homomorphism from $A$ to $K$. Let $d$ be a semimorphism from $B$ into $K$ (that is, $d:B\rightarrow K,\, d(0)=0,\,d(1)=1$ and for all $x,\,y\in B,\, d(x\vee y)=d(x)\vee d(y)\ )$ with $f(x)\leq d(x)$ for all $x\in A$. Then there is a homomorphism $h$ from $B$ to $K$ such that $h(x)\leq d(x)$ for all $x\in B$ and $h/A=f$. Monteiro \cite{bib:1965} and Bacsich\index{Author}{Bacsich} \cite{bib:1972}.

\noindent[\textbf{14 BB}] Fra\"{\i}ss\'{e}'s Theorem:\index{Subject}{Fra\"{i} ss\'{e}'s theorem} Let $\mathcal{E}$ be inductively ordered (that is, for any $E,\,E^{\prime}\in \mathcal{E},\,\exists E^{\prime\prime}\in \mathcal{E}$ such that $E\cup E^{\prime}\subseteq E^{\prime\prime}$.) For any $E\in \mathcal{E}$, assume $\mathcal{R}_{E}$ is a finite set of $m$-ary relations on $E$. If for any $E,\,E^{\prime}\in \mathcal{E}$ with $E\subseteq E^{\prime}$ we have $\forall R\in \mathcal{R}_{E^{\prime}},\,R/E\in \mathcal{R}_{E}$ then there is an $m$-ary relation $S$ on $\bigcup \mathcal{E}$ such that $S|E\in \mathcal{R}_{E}$ for any $E\in \mathcal{E}$. B\'{e}n\'{e}jam\index{Author}{B\'{e}n\'{e}jam} \cite{bib:1970}. See [14 Y].

\noindent[\textbf{14 BC}] If $m_{0}$ is an additive, real valued measure on a subalgebra $\mathcal{B}_{0}$ of a Boolean algebra $\mathcal{B}$ then there exists a measure $m$ on $\mathcal{B}$ such that $m=m_{0}$ on $\mathcal{B}_{0}$ and the range of $m$ is contained in the closure of the range of $m_{0}$. Tarski\index{Author}{Tarski} \cite{bib:1930} and note 147.

\noindent[\textbf{14 BD}] Let $H$ be a hypergraph. If every finite subhypergraph has an independent vertex cover then $H$ has an independent vertex cover. Cowen \cite{bib:1990} and note 117.

\noindent[\textbf{14 BE}] Let $H$ be a hypergraph. If every finite subhypergraph is 2-colorable then $H$ is 2-colorable. Cowen \cite{bib:1990} and note 117.

\noindent[\textbf{14 BF}] Let $H$ be a hypergraph whose edges contain at most three vertices. If every finite subhypergraph of $H$ has an independent vertex cover then $H$ has an independent vertex cover. Cowen \cite{bib:1990} and note 117.

\noindent[\textbf{14 BG}] Let $H$ be a hypergraph whose edges contain at most three vertices. If every finite subhypergraph of $H$ is 2-colorable then $H$ is 2-colorable. Cowen\index{Author}{Cowen} \cite{bib:1990} and note 117.

\noindent[\textbf{14 BH}] The sum of finite frames is a spatial frame. Paseka\index{Author}{Paseka} \cite{bib:1989} and note 29.

\noindent[\textbf{14 BI}] Every compact, conjunctive frame is spatial. Paseka \cite{bib:1989} and note 29.

\noindent[\textbf{14 BJ}] The sum of compact, regular frames is spatial. Paseka \cite{bib:1989} and note 29.

\noindent[\textbf{14 BK}] The sum of 4-element Boolean algebras is spatial. Paseka \cite{bib:1989} and note 29.

\noindent[\textbf{14 BL}] Every compact, 0 dimensional frame is spatial. Paseka \cite{bib:1989} and note 29.

\noindent[\textbf{14 BM}] A system of equations over a finite ring $R$ has a solution in $R$ if and only if every finite sub-system has a solution in $R$. Abian\index{Author}{Abian} \cite{bib:1972} and note 30.

\noindent[\textbf{14 BN}] A system of equations over a finite field $F$ has a solution in $F$ if and only if every finite sub-system has a solution in $F$. Cowen/Emerson\index{Author}{Cowen/Emerson} \cite{bib:1995} and note 30.

\noindent[\textbf{14 BO}] A system of equations over a finite structure $(A, +, \cdot)$ has a solution in $A$, if and only if every finite sub-system has a solution in \emph{A}. note 30.

\noindent[\textbf{14 BP}] Every non-empty bounded distributive lattice has a prime ideal. Morillon\index{Author}{Morillon} \cite{bib:1988}.

\noindent[\textbf{14 BQ}] Every proper ideal in a bounded distributive lattice is the intersection of the prime ideals containing it. Morillon \cite{bib:1988} and note 79. (Compare with [14 AE].)

\noindent[\textbf{14 BR}] Every $\vee$-compact non trivial lattice has a principal prime ideal. Morillon \cite{bib:1988}, \cite{bib:1985} and note 71.

\noindent[\textbf{14 BS}] Every bounded distributive lattice has a prime ideal choice function. Morillon \cite{bib:1988}, \cite{bib:1985} and note 71.

\noindent[\textbf{14 BT}] If $(A_{i})_{i\in I}$ is a family of bounded, distributive, $\vee$-compact lattices and if $(J_{i})_{i\in I}$ is a family such that for all $i\in I,\,J_{i}$ is a proper ideal in $A_{i}$ then there is a family $(P_{i})_{i\in I}$ such that for all $i\in I,\,P_{i}$ is a prime principal ideal in $A_{i}$ containing $J_{i}$. Morillon \cite{bib:1985}, \cite{bib:1988} and note 71.

\noindent[\textbf{14 BU}] If $(A_{i})_{i\in I}$ is a family of distributive non trivial bounded lattices, then there is a family $(f_{i})_{i\in I}$ such that each $f_{i}$ is a prime ideal choice function on $A_{i}$. Morillon \cite{bib:1985}, \cite{bib:1988} and note 71.

\noindent[\textbf{14 BV}] For every set $E$ there is an effectively $E$-adequate ultrafilter. Morillon \cite{bib:1988}, \cite{bib:1985} and note 71.

\noindent[\textbf{14 BW}] For every set $E$ there is an ultrafilter choice function on $E$. Morillon \cite{bib:1985}, \cite{bib:1988} and note 71.

\noindent[\textbf{14 BX}] If $(E_{i})_{i\in I}$ is a family of sets and $(F_{i})_{i\in I}$ is a family such that for each $i,\,F_{i}$ is a filter in $\mathcal{P}(E_{i})$, then there is a family $(U_{i})_{i\in I}$ such that for each $i\in I,\,U_{i}$ is an ultrafilter in $\mathcal{P}(E_{i})$ containing $F_{i}$. Morillon \cite{bib:1985}, \cite{bib:1988} and note 71.

\noindent[\textbf{14 BY}] For every topological space $X$, every proper ideal in the normed $\mathbb{R}$-algebra of continuous bounded functions from $X$ into $\mathbb{R}$ is contained in a maximal ideal. Morillon\index{Author}{Morillon} \cite{bib:1988}.

\noindent[\textbf{14 BZ}] Every proper ideal in a non-trivial normed $\mathbb{R}$-algebra or $\mathbb{C}$-algebra with identity is contained in a maximal ideal. Morillon \cite{bib:1988}.

\noindent[\textbf{14 CA}] If $(A_{i})_{i\in I}$ is a family of non trivial commutative normed algebras with identity then there is a family $(J_{i})_{i\in I}$ such that for each $i\in I,\,J_{i}$ is a maximal ideal in $A_{i}$. Morillon \cite{bib:1988}.

\noindent[\textbf{14 CB}] If $X$ is a topological space and $I$ is a proper ideal in the algebra of all bounded continuous functions from $X$ to $\mathbb{R}$ such that every function in $I$ vanishes at at least one point of $X$ then $I$ is contained in a maximal ideal. Morillon \cite{bib:1988}.

\noindent[\textbf{14 CC}] A projective limit of compact topological spaces is compact. Morillon \cite{bib:1988} and note 71.

\noindent[\textbf{14 CD}] A projective limit of a directed projective system of compact topological spaces is compact. Morillon \cite{bib:1988} and note 71.

\noindent[\textbf{14 CE}] A projective limit of a directed projective system of compact nonempty topological spaces is non-empty. Morillon \cite{bib:1988} and note 71.

\noindent[\textbf{14 CF}] A projective limit of a directed projective system of non-empty sets is non-empty. Morillon \cite{bib:1988} and note 71.

\noindent[\textbf{14 CG}] A projective limit of a directed projective system of finite sets is compact. Morillon \cite{bib:1988} and note 71.

\noindent[\textbf{14 CH}] For every bounded, distributive lattice $(A, \wedge, \vee)$ in which $0$ is $\wedge$-compact, for all $a\in A,\,a\neq 0$ and for all sequences $(X_{n})_{n\in\omega}$ of subsets of $A$, there is a prime filter in $A$ which contains $a$ and respects each $X_{n}$. Morillon \cite{bib:1988} and note 71.

\noindent[\textbf{14 CI}] 3-Sat: Restricted Compactness Theorem for Propositional Logic II: If $\Sigma$ is a set of formulas in a propositional language such that every finite subset of $\Sigma$ is satisfiable and if every formula in $\Sigma$ is a disjunction of at most three literals, then $\Sigma$ is satisfiable. (A \emph{literal} is a propositional variable or its negation.) Cowen\index{Author}{Cowen} \cite{bib:1990}.

\noindent[\textbf{14 CJ}] The Completeness Theorem for Propositional Logic: Every consistent set of formulas is satisfiable. Henkin\index{Author}{Henkin} \cite{bib:1954a} and \cite{bib:1954b}.

\noindent[\textbf{14 CK}] The Completeness Theorem for First Order Logic: Every consistent set of formulas is satisfiable. Henkin \cite{bib:1954a} and \cite{bib:1954b}.

\noindent[\textbf{14 CL}] Gelfand Extreme Point Theorem:\index{Subject}{Gelfand extreme point theorem} If $A$ is a non-trivial Gelfand algebra then the closed unit ball in the dual of $A$ has an extreme point $e$ with $e\neq 0$. Morillon \cite{bib:1986}, notes 23 and 29.

\noindent\textbf{FORM 15.} $KW(\infty, \infty)$ (KW),\index{Subject}{Banach-Tarski paradox} The Kinna-Wagner Selection Principle: For every set $M$ there is a function $f$ such that for all $A\in M$, if $|A|>1$ then $\emptyset \neq f(A)\subsetneq A$. (See form $81(n)$). Kinna/Wagner\index{Author}{Kinna/Wagner} \cite{bib:1955}.

[\textbf{15 A}] $K(1)$: $(\forall M)(\exists\alpha$ an ordinal) $(|M|=|A|$ for some $A\subseteq \mathcal{P}(\alpha))$. (See form $81(n).)$ Jech \cite{bib:1973b}
p 53 prob 4.12 and Kinna/Wagner \cite{bib:1955}.

\noindent[\textbf{15 B}] SDO: For any family $X$ of infinite sets, there is a function $f$ with domain $X$ such that $\forall y\in X,\,f(y)$ is a dense linear ordering of $y$. Pincus\index{Author}{Pincus} \cite{bib:1997}.

\noindent[\textbf{15 C}] SUO: For any family $X$ of infinite sets, there is a function $f$ with domain $X$ such that $\forall y\in X,\,f(y)$ is an unbounded linear ordering of $y$. (An ordering is unbounded if there is neither a least element nor a greatest element.) Pincus\index{Author}{Pincus} \cite{bib:1997}.

\noindent[\textbf{15 D}]\index{Author}{Sierpi\'{n}ski} SDUO: For any family $X$ of infinite sets, there is a function $f$ with domain $X$ such that $\forall y\in X,\,f(y)$ is a dense unbounded linear ordering of $y$. Pincus \cite{bib:1997}.

\noindent\textbf{FORM 16.} $C(\aleph_{0}, \leq 2^{\aleph_{0}})$: Every denumerable collection of non-empty sets each with power $\leq 2^{\aleph_{0}}$ has a choice function. Sierpi\'{n}ski \cite{bib:1918}.

\noindent\textbf{FORM 17.} Ramsey's Theorem\index{Subject}{Ramsey's theorem} I: If $A$ is an infinite set and the family of all 2 element subsets of $A$ is partitioned into 2 sets $X$ and $Y$, then there is an infinite subset $B\subseteq A$ such that all 2 element subsets of $B$ belong to $X$ or all 2 element subsets of $B$ belong to $Y$. (Also, see form 325.) Jech\index{Author}{Jech} \cite{bib:1973b} p 164 prob 11.20 and Ramsey\index{Author}{Ramsey} \cite{bib:1929}.

\noindent\textbf{FORM 18.} $PUT(\aleph_{0},2, \aleph_{0})$: The union of a denumerable family of pairwise disjoint pairs has a denumerable subset. Brunner\index{Author}{Brunner} \cite{bib:1984f}.

[\textbf{18 A}] $PC(\aleph_{0},2, \infty)$: Every denumerable set of pairs has an infinite subset with a choice function.

\noindent[\textbf{18 B}] The union of a denumerable family of pairs has a denumerable subset. Note 135.

\noindent[\textbf{18 C}] $PC(WO, 2, \infty)$: Every infinite well ordered set of pairs has an infinite subset with a choice function.

\noindent[\textbf{18 D}] $PUT(WO, 2, WO)$: The union of an infinite well ordered family of pairwise disjoint pairs has an infinite well ordered subset. ([18D $\leftrightarrow$ 18)

\noindent\textbf{FORM 19.} A real function is analytically representable if and only if it is in Baire's classification.\index{Subject}{Baire's classification} G. Moore\index{Author}{G. Moore} \cite{bib:1982} (2.3.1).

\noindent\textbf{FORM 20.} If $\{A_{x}:x\in S\}$ and $\{B_{x}:x\in S\}$ are families of pairwise disjoint sets and $|A_{x}|=|B_{x}|$ for all $x\in S$, then $I\bigcup_{x\in S}A_{x}|=|\bigcup_{x\in S}B_{x}|$. G. Moore \cite{bib:1982} (1.4.12 and 1.7.8).

\noindent\textbf{FORM 21.} If $S$ is well ordered, $\{A_{x}:x\in S\}$ and $\{B_{x}:x\in S\}$ are families of pairwise disjoint sets, and $|A_{x}|=|B_{x}|$ for all $x\in S$, then $|\bigcup_{x\in S}A_{x}|=|\bigcup_{x\in S}B_{x}|$.
G. Moore \cite{bib:1982}.

\noindent\textbf{FORM 22.} $UT(2^{\aleph_{0}},2^{\aleph_{0}},2^{\aleph_{0}})$: If every member of an infinite set of cardinality $2^{\aleph_{0}}$ has power $2^{\aleph_{0}}$, then the union has power $2^{\aleph_{0}}$. Borel\index{Author}{Borel} \cite{bib:1898} and G. Moore \cite{bib:1982}.

\noindent\textbf{FORM 23.} $(\forall\alpha)(UT(\aleph_{\alpha}, \aleph_{\alpha}, \aleph_{\alpha}))$: For every ordinal $\alpha$, if $A$ and every member of $A$ has cardinality $\aleph_{\alpha}$, then $|\bigcup A|=\aleph_{\alpha}$. G. Moore \cite{bib:1982}.

\noindent\textbf{FORM 24.} $C(\aleph_{0},2^{(2^{\aleph_{0}})})$: Every denumerable collection of non-empty sets each with power $2^{(2^{\aleph_{0}})}$ has a choice function. G. Moore \cite{bib:1982}.

\noindent\textbf{FORM 25.} $\aleph_{\beta+1}$ is regular for all ordinals $\beta$. Feferman/Levy\index{Author}{Feferman/Levy} \cite{bib:1963} and Gitik\index{Author}{Gitik} \cite{bib:1980}.

\noindent\textbf{FORM 26.} $UT(\aleph_{0},2^{\aleph_{0}},2^{\aleph_{0}})$: The union of denumerably many sets each of power $2^{\aleph_{0}}$ has power $2^{\aleph_{0}}$. Borel \cite{bib:1898}.

\noindent\textbf{FORM 27.} $(\forall\alpha)(UT(\aleph_{0}, \aleph_{\alpha}, \aleph_{\alpha}))$: The union of denumerably many sets each of power $\aleph_{\alpha}$ has power $\aleph_{\alpha}$. G. Moore\index{Author}{G. Moore} \cite{bib:1982} p 36.

\noindent\textbf{FORM 28.}\index{Author}{Sierpi\'{n}ski} $C(\aleph_{0},2^{\aleph_{0}})$: Every denumerable collection of non-empty sets each with power $2^{\aleph_{0}}$ has a choice function. Sierpi\'{n}ski \cite{bib:1918}.

\noindent\textbf{FORM 29.} If $|S|=\aleph_{0}$ and $\{A_{x}:x\in S\}$ and $\{B_{x}:x\in S\}$ are families of pairwise disjoint sets and $|A_{x}|=|B_{x}|$ for all $x\in S$, then $|\bigcup_{x\in S}A_{x}|=|\bigcup_{x\in S}B_{x}|$. G. Moore \cite{bib:1982} p 324.

\noindent\textbf{FORM 30.} Ordering Principle:\index{Subject}{ordering principle} Every set can be linearly ordered. Schr\"{o}der\index{Author}{Schr\"{o}der} \cite{bib:1988}.

[\textbf{30 A}] If $R$ is a partial ordering in which the set of predecessors of each element is linearly ordered, then there is a linear ordering extending $R$. Felgner\index{Author}{Felgner} \cite{bib:1969}.

\noindent[\textbf{30 B}] ADO: Every infinite set differs finitely from a set with a dense linear ordering: for all infinite $x$, there is a $y\subseteq x$ such that $y$ has a dense linear ordering and $x-y$ is finite. Pincus\index{Author}{Pincus} \cite{bib:1997}.

\noindent[\textbf{30 C}] SADO: For every set $F$ there is a function $f$ with domain $F$ such that for all $x\in F,\, f(x)\subseteq x,\,f(x)$ has a dense linear ordering, and $x-f(x)$ is finite. Pincus \cite{bib:1997}.

\noindent[\textbf{30 D}] UO: There exists an unbounded linear ordering on every infinite set. (An ordering is unbounded if there is neither a least element nor a greatest element.) Pincus \cite{bib:1997}.

\noindent[\textbf{30 E}] ADUO: For all infinite $x$, there exists $y\subseteq x$ such that $y$ has a dense unbounded linear ordering and $x-y$ is finite. Pincus \cite{bib:1997}.

\noindent[\textbf{30 F}] SADUO: For every set $F$ there is a function $f$ with domain $F$ such that for all $x\in F,\, f(x)\subseteq x,\,f(x)$ has a dense unbounded linear ordering, and $x-f(x)$ is finite. Pincus \cite{bib:1997}.

\noindent[\textbf{30 G}] Weak Hausdorff Lemma:\index{Subject}{Hausdorff lemma} For every set $X$, every $\subseteq$-chain in $\mathcal{P}(X)$ is contained in a maximal $\subseteq$-chain. Note 130.

\noindent \textbf{FORM 31.} $UT(\aleph_{0}, \aleph_{0}, \aleph_{0})$: The countable union theorem: The union of a denumerable set of denumerable sets is denumerable. Cantor\index{Author}{Cantor} \cite{bib:1878}.

\noindent\textbf{FORM 32.} $C(\aleph_{0}, \leq\aleph_{0})$: Every denumerable set of non-empty countable sets has a choice function. Sierpi\'{n}ski \cite{bib:1918}.

[\textbf{32 A}] $C(\aleph_{0}, \aleph_{0})$. Every denumerable set of denumerable sets has a choice function.

\noindent[\textbf{32 B}] $PC(\aleph_{0}, \aleph_{0}, \infty)$: Every denumerable set of denumerable sets has an infinite subset with a choice function.

\noindent[\textbf{32 C}] $PC(WO, \aleph_{0}, \infty)$: Every infinite well ordered set of denumerable sets has an infinite subset with a choice function.

\noindent\textbf{\textbf{FORM 33($n$).}} If $n\in \omega-\{0,1\},\,C(LO, n)$: Every linearly ordered set of $n$ element sets has a choice function. Truss\index{Author}{Truss} \cite{bib:1973a}.

\noindent\textbf{FORM 34.} $\aleph_{1}$ is regular. Russell\index{Author}{Russell} \cite{bib:1906} and note 3.

[\textbf{34 A}] $\exists\alpha< \omega_{2}$ such that there are no hereditarily countable sets of rank $\alpha$. Jech\index{Author}{Jech} \cite{bib:1982} and note 48.

\noindent[\textbf{34 B}] $\omega_{1}$ with the order topology is not paracompact. Good/Tree\index{Author}{Good/Tree} \cite{bib:1995} and note 107.

\noindent[\textbf{34 C}]\index{Author}{Sierpi\'{n}ski} $\omega_{1}$ with the order topology has the property that every infinite subset has a limit point in $\omega_{1}$. Good/Tree \cite{bib:1995} and note 107.

\noindent[\textbf{34 D}] $\omega_{1}$ with the order topology does not contain a countable discrete family of open subsets. Good/Tree \cite{bib:1995} and note 107.

\noindent[\textbf{34 E}] $\omega_{1}$ with the order topology is not Lindel\"{o}f. Good/Tree \cite{bib:1995} and note 107.

\noindent\textbf{FORM 35.} The union of countably many meager subsets of $\mathbb{R}$ is meager. (Meager sets are the same as sets of the first category.) Feferman/Levy\index{Author}{Feferman/Levy}\index{Author}{Levy} \cite{bib:1963}, Jech\index{Author}{Jech} \cite{bib:1973b} p 7 prob 1.7, and note 28.

\noindent\textbf{FORM 36.} If $A\subseteq \mathbb{R}^{n}$ and $A \bigcap B$ is countable for every bounded $B$ then $A$ is countable. G. Moore\index{Author}{G. Moore} \cite{bib:1982} p 36.

\noindent\textbf{FORM 37.} Lebesgue measure is countably additive. Sierpi\'{n}ski \cite{bib:1916}.

\noindent\textbf{FORM 38.} $\mathbb{R}$ is not the union of a countable family of countable sets. Feferman/Levy \cite{bib:1963} and note 3.

[\textbf{38 A}] $UT(\aleph_{0}, \aleph_{0}, \neq 2^{\aleph_{0}})$: A denumerable union of denumerable sets does not have cardinality $2^{\aleph_{0}}$.

\noindent\textbf{FORM 39.} $C(\aleph_{1}, \infty)$: Every set $A$ of non-empty sets such that $|A|=\aleph_{1}$ has a choice function. Sierpi\'{n}ski \cite{bib:1918}, G. Moore \cite{bib:1982} p. 202.

\noindent\textbf{FORM 40.} $C(WO, \infty)$: Every well orderable set of non-empty sets has a choice function. G. Moore \cite{bib:1982} p 325.

[\textbf{40 A}] $A(L2, L1)$: For every $T_{2}$ topological space $(X, T)$, if each open cover of $X$ has a well ordered subcover, then each open cover has a well ordered refinement. Brunner\index{Author}{Brunner} \cite{bib:1983d} and note 26.

\noindent[\textbf{40 B}] For all ordinals $\alpha$, the $\aleph_{\alpha}$ partition principle holds: For every ordinal $\alpha$ and every cardinal $\kappa$, if $\aleph_{\alpha}\leq^{\ast}\kappa$, then $\aleph_{\alpha}\leq \kappa$. Pelc\index{Author}{Pelc} \cite{bib:1978}, Banaschewski/Moore\index{Author}{Banaschewski/Moore} \cite{bib:1990}, and note 69.

\noindent[\textbf{40 C}] For each idemmultiple set $x\ (2\times x\approx x)$, if there exists a surjection $f$ mapping $x$ onto an ordinal $\lambda$ such that for each $\psi<\lambda,\,f^{-1}[\{\psi\}]$ is Dedekind infinite, then $\lambda\precsim x$. Higasikawa\index{Author}{Higasikawa} \cite{bib:1995} and note 94.

\noindent\textbf{FORM 41}. $W_{\aleph_{1}}$: For every cardinal $m,\,m\leq\aleph_{1}$ or $\aleph_{1}\leq m$. Jech \cite{bib:1973b}.

\noindent\textbf{FORM 42.} L\"{o}wenheim-Skolem Theorem:\index{Subject}{L\"{o}wenheim-Skolem theorem} If a countable family of first order sentences is satisfiable in a set $M$ then it is satisfiable in a countable subset of $M$. G. Moore \cite{bib:1982} p 251.

\noindent\textbf{FORM 43.} $DC(\omega)$ (DC), Principle of Dependent Choices: If $S$ is a relation on a non-empty set $A$ and $(\forall x\in A)(\exists y\in A)(xSy)$ then there is a sequence $a(0),a(1),a(2),\ldots$ of elements of $A$ such that $(\forall n\in \omega) (a(n)\,S\,a(n+1))$. Tarski\index{Author}{Tarski} \cite{bib:1948} p 96 , Levy \cite{bib:1964} p 136 and note 54.

[\textbf{43 A}] Tarski's Lemma: Let $\{x(n) : n\in \omega\}$ be a denumerable collection of subsets of a Boolean algebra $B$ each of which has a meet in B. Then for any $a\neq 0$ in $B$ there is a proper filter $F$ of $B$ that decides $\bigwedge x(n)$ for all $n$ and has $a\in F$. ($F$ \emph{decides} $\bigwedge x(n)$ if either $\bigwedge x(n)\in F$ or $t\not\in F$ for some $t\in x(n).)$ Goldblatt\index{Author}{Goldblatt} \cite{bib:1985}.

\noindent[\textbf{43 B}]\index{Subject}{Tarski's lemma} Countable Henkin Principle:\index{Subject}{Henkin principle} Let $\Phi$ be the class of formulas in a language that includes the language of the sentential calculus. Let $I$ be a countable set of inferences, and $\Gamma$ a consistent subset of $\Phi$ such that $\Gamma \cup\{\psi\}$ respects $I$ for each $\psi\in\Phi$. Then for any $\phi \in\Phi$, if $\phi$ is consistent with $\Gamma$ then $\Gamma$ has a finitely consistent extension that contains $\phi$ and decides $I$. Goldblatt \cite{bib:1985} and note.

\noindent[\textbf{43 C}] Existence of a generic set: If $\mathcal{D}$ is a denumerable collection of dense sets in a partial order $(P, \leq)$ then for all $p\in P$ there is a $\mathcal{D}$ generic $G$ with $p\in G$. Goldblatt \cite{bib:1985} and note 47.

\noindent[\textbf{43 D}] If $\mathcal{D}$ is a denumerable collection of dense subsets of a partial order $(P, \leq)$ then for all $p\in P$ there is a sequence $f:\omega \rightarrow P$ with $f(0)=p$ and ($\forall n\in\omega)(f(n)\geq f(n+1))$ and range $(f) \cap D\neq\emptyset$ for all $D\in \mathcal{D}$. Goldblatt \cite{bib:1985}.

\noindent[\textbf{43 E}] The Baire Category Theorem for \v{C}ech Complete-I Spaces: Every \v{C}ech complete space is Baire. Goldblatt \cite{bib:1985} and note 28.

\noindent[\textbf{43 F}] Baire Category Theorem for Sequentially Complete Metric Spaces: Every sequentially complete metric space is Baire. Blair\index{Author}{Blair} \cite{bib:1977}, Brunner\index{Author}{Brunner} \cite{bib:1983c}, Goldblatt \cite{bib:1985}, and note 28.

\noindent[\textbf{43 G}] Form 106 + form 154. Brunner \cite{bib:1983c} \cite{bib:1984b}.

\noindent[\textbf{43 H}] Every tree in which every element has at least one successor has a branch. Blass\index{Author}{Blass} \cite{bib:1979} and note note 21.

\noindent[\textbf{43 I}] Form [106 A] + form 10. Blass \cite{bib:1979}.

\noindent[\textbf{43 J}] Baire Category Theorem for Cantor Complete Metric Spaces: Every Cantor complete metric space is Baire. Brunner \cite{bib:1983c} and note 28.

\noindent[\textbf{43 K}] Let $M$ be a non-empty set partially ordered by $<$ without $<$-minimal elements. Then there is a subset $C$ of $M$ which is linearly ordered by $<$ such that $C$ has no $<$-minimal element. Note 146.

\noindent[\textbf{43 L}] $Z_{\omega}$: If every chain in a partially ordered set $P$ is finite, then $P$ contains a maximal element. Tarski\index{Author}{Tarski} \cite{bib:1948}, Wolk\index{Author}{Wolk} \cite{bib:1983}, and note 54.

\noindent[\textbf{43 M}] If $(X, \rho)$ is a complete metric space and $\phi:X\rightarrow \mathbb{R}$ is bounded above and upper semi-continuous then in the Br{\o}ndsted ordering, $(x\leq y$ iff $\rho(x, y)\leq \phi(y)-\phi(x))$, there is a maximal element. Brunner \cite{bib:1987a} and note 38.

\noindent[\textbf{43 N}] $RM1,\,\aleph_{0}$: The representation theorem for multi-algebras with $\aleph_{0}$ unary operations: Assume $(A, F)$ is a multi-algebra with $\aleph_{0}$ unary operations (and no other operations), then there is an algebra $(B, G)$ with $\aleph_{0}$ unary operations and an equivalence relation $E$ on $B$ such that $(B/E, G/E)$ and $(A, F)$ are isomorphic multi-algebras. Note 50 and H\"{o}ft/Howard\index{Author}{H\"{o}ft/Howard} \cite{bib:1981}.

\noindent[\textbf{43 O}] (Depends on $n\in\omega,\ n\geq 1)\ RM1,\,n$: The representation theorem for multi-algebras with $n$ unary operations: Assume $(A, F)$ is a multi-algebra with $n$ unary operations (and no other operations), then there is an algebra $(B, G)$ with $n$ unary operations and an equivalence relation $E$ on $B$ such that $(B/E, G/E)$ and $(A, F)$ are isomorphic multi-algebras. H\"{o}ft/Howard \cite{bib:1981} and note 50.

\noindent[\textbf{43 P}] The ultraproduct of a well founded structure by a countably complete ultrafilter is well founded. Spector\index{Author}{Spector} \cite{bib:1980} and note 54.

\noindent[\textbf{43 Q}] Matrix$(\omega)$: For every atomless poset $(P, \leq)$ and every family $D=\{D_{i}:i\in \omega\}$ of dense open subsets of $P$, there exists a matrix $C=\{C_{i}:i\in \omega\},\, |\bigcup C|=\omega$ of $P$ for $D$. Keremedis\index{Author}{Keremedis} \cite{bib:1996d} and note 132.

\noindent[\textbf{43 R}] Any partial order $(A, \leq)$ with no infinite descending chains is well founded, (that is, for all $X$ such that $\emptyset \neq X\subseteq A, X$ has a minimal element.) Spector\index{Author}{Spector} \cite{bib:1980} and note 54.

\noindent[\textbf{43 S}] If $S$ is a relation on $A$ such that $(\forall x\in A)(\exists y\in A)(x\,S\, y)$ and $x_{0}$ is any element of A, then there is a sequence $a(1),a(2),\ldots$ of elements of $A$ such that $a(n)\,S\,a(n+1)$ for all $n\in \omega$ and such that $a(0)=x_{0}$. Notes 101 and 54.

\noindent[\textbf{43 T}] A sequentially complete metric space is not meager. Note 28.

\noindent[\textbf{43 U}] Baire Category Theorem for Filter Complete Metric Spaces: Every filter complete metric space is Baire, Morillon\index{Author}{Morillon} \cite{bib:1988} and note 28.

\noindent[\textbf{43 V}] Every convex-compact set in a Hausdorff topological vector space with a basis of convex open sets is Baire. Fossy/Morillon\index{Author}{Fossy/Morillon} \cite{bib:1995} and note 23.

\noindent\textbf{FORM 44.} $DC(\aleph_{1})$: Given a relation $R$ such that for every subset $Y$ of a set $X$ with $|Y|<\aleph_{1}$ there is an $x\in X$ with $Y\ R\,x$, then there is a function $f:\aleph_{1}\rightarrow X$ such that $(\forall\beta<\aleph_{1})(\{f(\gamma):\gamma<b\}\,R\,f(\beta))$. Levy\index{Author}{Levy} \cite{bib:1964}.

\noindent \textbf{FORM 45($n$).} If $n\in \omega-\{0,1\},\,C(\infty, n)$: Every set of $n$-element sets has a choice function. Mostowski\index{Author}{Mostowski} \cite{bib:1945}, Truss\index{Author}{Truss} \cite{bib:1973a}.

\noindent\textbf{FORM 46($K$).} If $K$ is a finite subset of $\omega-\{0,1\},\,C(\infty, K)$: For every $n\in K$, every set of $n$-element sets has a choice function. Mostowski \cite{bib:1945}, Truss \cite{bib:1973a}.

\noindent\textbf{FORM 47($n$).} If $n\in \omega-\{0, 1\},\,C(WO, n)$: Every well ordered collection of $n$-element sets has a choice function. Mostowski \cite{bib:1945}, Truss \cite{bib:1973a}.

\noindent\textbf{FORM 48($K$).} If $K$ is a finite subset of $\omega-\{0, 1\},\,C(WO, K)$: For every $n\in K,\,C(WO, n)$. Mostowski \cite{bib:1945}, Truss \cite{bib:1973a}.

\noindent\textbf{FORM 49.} Order Extension Principle:\index{Subject}{order extension principle} Every partial ordering can be extended to a linear ordering. Tarski\index{Author}{Tarski} \cite{bib:1924} p 78.

\noindent\textbf{FORM 50.} Sikorski's\index{Author}{Sikorski} Extension Theorem:\index{Subject}{Sikorski's extension theorem} Every homomorphism of a subalgebra $B$ of a Boolean algebra $A$ into a complete Boolean algebra $B^{\prime}$ can be extended to a homomorphism of $A$ into $B^{\prime}$. Sikorski \cite{bib:1964}, p 141.

[\textbf{50 A}] Every complete Boolean algebra is injective. Banaschewski\index{Author}{Banaschewski} \cite{bib:1988}.

\noindent[\textbf{50 B}] Monteiro's\index{Author}{Monteiro} Theorem:\index{Subject}{Monteiro's theorem} Let $K$ be a complete Boolean algebra, $A$ a subalgebra of the Boolean algebra $B$ and $f$ a homomorphism from $A$ to $K$. Let $d$ be a semimorphism from $B$ into $K$ (that is, $d:B\rightarrow K, d(0)=0,\,d(1)=1$ and for all $x$ and $y$ in $B,\,d(x\vee y)=d(x)\vee d(y))$ with $f(x)\leq d(x)$ for all $x\in A$. Then there is a homomorphism $h$ from $B$ to $K$ such that $h(x)\leq d(x)$ for all $x\in B$ and $h/A=f$. Monteiro \cite{bib:1965} and Bacsich\index{Author}{Bacsich} \cite{bib:1972}.

\noindent[\textbf{50 C}] Every DeMorgan compact regular frame is injective in the category of all such frames. Banaschewski \cite{bib:1988}, notes 29 and 60.

\noindent[\textbf{50 D}] For any extension $M\supseteq L$ of a compact regular frame $L$, there is an $s\in M$ maximal with respect to the property that $x\vee s=1\rightarrow x=1$ for all $x\in L$. Banaschewski \cite{bib:1988}, notes 29 and 60.

\noindent[\textbf{50 E}] Form 14 + form 297. Banaschewski \cite{bib:1988} and Bell\index{Author}{Bell} \cite{bib:1988}.

\noindent[\textbf{50 F}] For any complete Boolean algebra $B$ and any subalgebra $A$ of $B$, there is a $U\in V^{(B)}$ such that $V^{(B)}\models ``U$ is an ultrafilter on $A$ and $U^{B}\subseteq U\mbox{''}$. $(V^{(B)}$ is the Boolean valued universe constructed from $B,\,U^{B}=\{(\hat{x}, x):x\in B\}$ is the canonical ultrafilter on $B$ and $x\mapsto\hat{x}$ is the canonical embedding of $V$ into $V^{(B)}.)$ Bell\index{Author}{Bell} \cite{bib:1977} and \cite{bib:1983}.

\noindent[\textbf{50 G}]\index{Author}{Sierpi\'{n}ski} For any complete Boolean algebra $B$ and any $C\in V^{(B)}$ such that $V^{(B)}\models ``C$ is a Boolean algebra.'' there is a $U\in V^{(B)}$ such that $V^{(B)}\models ``U$ is an ultrafilter in \emph{C}.'' See [50 F] for definitions. Bell \cite{bib:1983}.

\noindent[\textbf{50 H}] If $m:A\rightarrow B$ is a monomorphism between Boolean algebras, then there is an $m$-minimal epimorphism $p:B\rightarrow C$ to some Boolean algebra \emph{C}. $(p$ is $m$-minimal if (1) $p\circ m$ is monic and (2) if $q:C\rightarrow D$ is any epimorphism to a Boolean algebra $D$ such that $q\circ p\circ m$ is monic, then $q$ is an isomorphism.) Bell \cite{bib:1988}.

\noindent[\textbf{50 I}] If $A$ is a subalgebra of a Boolean algebra $B$, then there is a $\subseteq$-maximal proper filter $F$ in $B$ such that $F\cap A=\{1\}$. Bell \cite{bib:1988}

\noindent[\textbf{50 J}] Form 299 + form 14. Bell \cite{bib:1988}.

\noindent[\textbf{50 K}] Form 300 + form 14. Bell \cite{bib:1988}.

\noindent[\textbf{50 L}] Form 301 + form 14. Bell \cite{bib:1988}.

\noindent[\textbf{50 M}] Form 302 + form 14. Bell \cite{bib:1988}.

\noindent\textbf{FORM 51.} Cofinality Principle:\index{Subject}{cofinality principle} Every linear ordering has a cofinal sub well ordering. Sierpi\'{n}ski \cite{bib:1918} p 117.

[\textbf{51 A}] For all linear orders $(X, \leq)$, there is a $\displaystyle \sup$ function on the bounded, well ordered subsets of $X$. Manka\index{Author}{Manka} \cite{bib:1988a} and note 38.

\noindent[\textbf{51 B}] $Z(L, W)$: Every non-empty linearly ordered set in which every well ordered subset has an upper bound has a maximal element. Harper/Rubin \cite{bib:1976}, Manka \cite{bib:1988a}, H. Rubin/J. Rubin\index{Author}{Harper/Rubin}\index{Author}{H. Rubin/J. Rubin} \cite{bib:1985} p 37 $M1(L, W)$, and note 39.

\noindent[\textbf{51 C}] For every linearly ordered set $(X, \leq)$, if every well ordered subset of $X$ is bounded above then there is a sup function on the well ordered subsets of \emph{X}. Manka \cite{bib:1988a} and note 38.

\noindent\textbf{FORM 52.} Hahn-Banach Theorem:\index{Subject}{Hahn-Banach theorem} If $V$ is a real vector space and $p:V\rightarrow \mathbb{R}$ satisfies $p(x+y)\leq p(x)+p(y)$ and $(\forall t>0)(p(tx)=tp(x))$ and $S$ is a subspace of $V$ and $f:S\rightarrow \mathbb{R}$ is linear and satisfies $(\forall x\in S)(f(x)\leq p(x))$ then $f$ can be extended to $f^{\ast}:V\rightarrow \mathbb{R}$ such that $f^{\ast}$ is linear and $(\forall x\in V)(f^{\ast}(x)\leq p(x))$.
L\'{o}s/Ryll-Nardzewski\index{Author}{{\L}o\'{s}/Ryll-Nardzewski}
 \cite{bib:1951}.

[\textbf{52 A}] On every Boolean algebra there is an additive real valued measure. Luxemburg\index{Author}{Luxemburg} \cite{bib:1969} and note 147.

\noindent[\textbf{52 B}] For every proper filter $\mathcal{F}$ of elements of a Boolean algebra $\mathcal{B}$, there exists an additive, real valued measure $m$ on $\mathcal{B}$ such that $m(x)=1$ for all $x\in \mathcal{F}$. Luxemburg \cite{bib:1969} and note 147.

\noindent[\textbf{52 C}] For every non-empty set X and every proper ideal $\mathcal{I}$ of in $\mathcal{P}(X)$ there is a measure $m$ on $\mathcal{P}(X)$ such that $m(x)=0$ for all $x\in \mathcal{I}$. Luxemburg \cite{bib:1969} and note 147.

\noindent[\textbf{52 D}] Let $\mathcal{B}_{0}$ be a subalgebra of a Boolean algebra $\mathcal{B}$ and $m_{0}$ an additive, real valued measure on $\mathcal{B}_{0}$. Then there is a real valued, additive measure $m$ on $\mathcal{B}$ such that $m=m_{0}$ on $\mathcal{B}_{0}$ and the range of $m$ is contained in the closed convex hull of the range of $m_{0}$. Luxemburg\index{Author}{Luxemburg} \cite{bib:1969} and note 147.

\noindent[\textbf{52 E}] The unit ball is convex-compact in the weak$^\ast$ topology in the dual of a Banach space. (See [14 Q].) Luxemburg \cite{bib:1969}, H. Rubin/J. Rubin \cite{bib:1985} p 178, and note 23.

\noindent[\textbf{52 F}] If $X$ is a topological vector space then for every continuous sublinear functional $p$ on $E$ there is a linear functional $f$ on $E$ such that $f\leq p$. Fossy/Morillon\index{Author}{Fossy/Morillon} \cite{bib:1995}.

\noindent[\textbf{52 G}] If $X$ is a normed vector space then for every continuous sublinear functional $p$ on $E$ there is a linear functional $f$ on $E$ such that $f\leq p$. Fossy/Morillon \cite{bib:1995}.

\noindent[\textbf{52 H}] Existence of Real Valued Banach Limits: For any directed set $(\Delta, \sqsubseteq)$ there exists a (real valued) Banach limit. Schechter\index{Author}{Schechter} \cite{bib:1996b} and note 31.

\noindent[\textbf{52 I}] Real Valued Convex Domination Theorem: For every real vector space $X$, every subspace $X_{0}\subseteq X$, every linear $\lambda_{0}:X_{0}\rightarrow \mathbb{R}$ and every convex $p:X\rightarrow \mathbb{R}$, if $\lambda_{0}\leq p$ on $X_{0}$ then $\lambda_{0}$ can be extended to a linear map $\lambda:X\rightarrow \mathbb{R}$ such that $\lambda\leq p$ on $X$. Schechter \cite{bib:1996b} and note 31.

\noindent[\textbf{52 J}] Support Theorem for Real Valued Convex Functions: For every real vector space $X$, every convex function $p:X\rightarrow \mathbb{R}$ is a pointwise maximum of affine functions. Schechter \cite{bib:1996b} and note 31.

\noindent[\textbf{52 K}] Sandwich Theorem for Real Valued Functions: For every real vector space $X$, if $e:X\rightarrow \mathbb{R}$ is a concave function, $g:X\rightarrow \mathbb{R}$ is a convex function, and $e\leq g$ everywhere on $X$, then there exists an affine function $f: X\rightarrow \mathbb{R}$ satisfying $e\leq f\leq g$ everywhere on $X$. Schechter \cite{bib:1996b} and note 31.

\noindent\textbf{FORM 53.} For all infinite cardinals $m,\,m^{2}\leq 2^{m}$. Mathias\index{Author}{Mathias} \cite{bib:1979} prob 1336.

\noindent\textbf{FORM 54.} For all infinite cardinals $m,\,m\ \mathrm{adj}\ 2^{m}$ implies $m$ is an aleph. ($m\ \mathit{adj}\ n$ iff $m<n$ and $(\neg\exists p)(m<p<n).)$ Mathias \cite{bib:1979} prob 1337, thm 1338 and prob 1339.

\noindent\textbf{FORM 55.} For all infinite cardinals $m$ and $n$, if $m\cdot n=m+n$ then, $m\leq n$ or $n\leq m$. Mathias \cite{bib:1979} p 125.

\noindent\textbf{FORM 56.} $\aleph(2^{\aleph_{0}})\neq\aleph_{\omega}$. ($\aleph(2^{\aleph_{0}})$ is Hartogs' aleph, the least $\aleph$ not $\leq 2^{\aleph_{0}}$.) Mathias \cite{bib:1979} p 125.

\noindent\textbf{FORM 57.} If $x$ and $y$ are Dedekind finite sets then either $|x|\leq|y|$ or $|y|\leq|x|$. Mathias \cite{bib:1979} p 125 and note 94.

\noindent\textbf{FORM 58.} There is an ordinal $\alpha$ such that $\aleph(2^{\aleph_{\alpha}})\neq\aleph_{\alpha+1}.\
 (\aleph(2^{\aleph_{\alpha}})$ is Hartogs' aleph, the least $\aleph$ not $\leq 2^{\aleph_{\alpha}}$.) Mathias \cite{bib:1979} p 126.

\noindent\textbf{FORM 59($\leq$).} If $(A, \leq)$ is a partial ordering that is not a well ordering, then there is no set $B$ such that $(B, \leq)$ (the usual injective cardinal ordering on $B$) is isomorphic to $(A, \leq)$. Mathias \cite{bib:1979} p 120.

\noindent\textbf{FORM 60.} $C(\infty, WO)$: Every set of non-empty, well orderable sets has a choice function. Jech\index{Author}{Jech} \cite{bib:1973b}, G. Moore\index{Author}{G. Moore} \cite{bib:1982} p 125.

[\textbf{60 A}] $(\forall\alpha)C(\infty, \aleph_{\alpha})$: For every ordinal $\alpha$, every family of sets each with power $\aleph_{\alpha}$ has a choice function.

\noindent\textbf{FORM 61.} $(\forall n\in\omega, n\geq 2)(C(\infty, n))$: For each $n\in\omega, n\geq 2$, every set of $n$ element sets has a choice function. Levy\index{Author}{Levy} \cite{bib:1962}.

\noindent\textbf{FORM 62.}\index{Author}{Royden} $C(\infty, <\aleph_{0})$: Every set of non-empty finite sets has a choice function.

[\textbf{62 A}] (Depends on the two primes $p_{1}\neq p_{2}$) Form $308(p_{1})\
+$ form $308(p_{2})$. Howard/Yorke\index{Author}{Howard/Yorke} \cite{bib:1987}.

\noindent[\textbf{62 B}] (Depends on the prime $p_{1}$) Form $45(p_{1})\
+$ form $308(p_{1})$. Howard/Yorke \cite{bib:1987}.

\noindent[\textbf{62 C}] For every prime $p$, if $\{G_{y}:y\in Y\}$ is a set of finite groups, then the weak direct product $\prod_{y\in Y}G_{y}$ has a maximal $p$-subgroup. Howard/Yorke \cite{bib:1987} and note 24.

\noindent[\textbf{62 D}] For every prime $p$, if $Y$ is a set of non-empty, finite sets, then the weak direct product $\prod_{y\in Y}S_{y}$ has a maximal $p$-subgroup. $(S_{y}$ is the symmetric group on $y.)$ Howard/Yorke \cite{bib:1987} and note 24.

\noindent [\textbf{62 E}] $KW(\infty, <\aleph_{0})$, The Kinna-Wagner Selection Principle for families of finite sets: For every set $M$ of finite sets there is a function $f$ such that for all $A\in M$, if $|A|>1$ then $\emptyset \neq f(A)\subsetneq A$. Note 70.

\noindent\textbf{FORM 63.} $SPI$: Weak ultrafilter principle: Every infinite set has a non-trivial ultrafilter. Jech\index{Author}{Jech} \cite{bib:1973b} p 172 prob 8.5.

\noindent\textbf{FORM 64.} $E(I, Ia)$ (Howard/Yorke \cite{bib:1989}): There are no amorphous sets. Note 94 and Levy \cite{bib:1958}, notes 57 and 94.

[\textbf{64 A}] An amorphous power of a compact $T_{2}$ space is compact. Brunner\index{Author}{Brunner} \cite{bib:1984b}.

\noindent[\textbf{64 B}] Finite products of metric $C$ spaces are $C$. (A space is $C$ if each open covering has an amorphous refinement.) Brunner \cite{bib:1984b}.

\noindent[\textbf{64 C}] $C$ spaces are compact. (A space is $a\ C$ \emph{space}
if each open covering has an amorphous refinement.) Brunner \cite{bib:1984b}.

\noindent[\textbf{64 D}] Metric $C$ spaces are limited amorphous. (A space is $C$ if every open covering has an amorphous refinement. A space is \emph{limited amorphous} if each amorphous subset is relatively compact.) Brunner \cite{bib:1984b}.

\noindent\textbf{FORM 65.} The Krein-Milman Theorem:\index{Subject}{Krein-Milman theorem} Let $K$ be a compact convex set in a locally convex topological vector space $X$. Then $K$ has an extreme point. (An extreme point is a point which is not an interior point of any line segment which lies in $K$.) Royden \cite{bib:1963} and H. Rubin/J. Rubin\index{Author}{H. Rubin/J. Rubin} \cite{bib:1985}, p. 177.

\noindent\textbf{FORM 66.} Every vector space over a field has a basis. Hamel\index{Author}{Hamel} \cite{bib:1905}.

\noindent\textbf{FORM 67.} $MC(\infty, \infty)$ (MC), The Axiom of Multiple Choice: For every set $M$ of non-empty sets there is a function $f$ such that $(\forall x\in M)(\emptyset\neq f(x)\subseteq x$ and $f(x)$ is finite).

[\textbf{67 A}] $A(A2, H2)$: For every $T_{2}$ topological space $(X, T)$ if $(X, T)$ is A2, then $(X, T)$ is hereditarily $A2$. $((X, T)$ is A2 means if $U\subseteq T$ covers $X$ then $\exists f:X\rightarrow T$ such that for all $x\in X,\,x\in f(x)$ and $f$'' $X$ refines $U$.) Brunner \cite{bib:1983d} and note 26.

\noindent[\textbf{67 B}] $MW$: Each set can be covered by a well ordered family of finite sets. Levy \cite{bib:1962}.

\noindent[\textbf{67 C}] $A(P)$: (Depends on $P\in\{D1, B2, B3, H2, A2\}.)$ Every $T_{2}$ topological space has property $P$. Brunner\index{Author}{Brunner} \cite{bib:1983d} and note 26.

\noindent[\textbf{67 D}] For every Abelian group $A$, every torsion free divisible subgroup of $A$ is a direct summand of $A$. Armbrust\index{Author}{Armbrust} \cite{bib:1972} and note 24.

\noindent[\textbf{67 E}] For any algebra \textbf{A} with finitary operations (that is, \textbf{A} is of type $\tau$ where $\tau =(K_{i})_{i\in I}$ with $K_{i}$ finite for each $i\in I$) and for any generating subset $Z$ of \textbf{A}, there is a function $h:A\rightarrow \mathcal{P}_{\omega}(Z)$ ($A$ is the domain of \textbf{A} and $\mathcal{P}_{\omega}(Z)$ is the set of all finite subsets of $Z$.) such that $\forall a\in A,\ a$ is in the subalgebra of \textbf{A} generated by $h(a)$. Diener\index{Author}{Diener} \cite{bib:1989} and note 115.

\noindent[\textbf{67 F}] For any algebra \textbf{A} with one unary operation $f$ and any generating set $Z$ of \textbf{A}, there is a function $h:A\rightarrow \mathcal{P}_{\omega}(Z)$ ($A$ is the domain of \textbf{A} and $\mathcal{P}_{\omega}(Z)$ is the set of all finite subsets of $Z$.) such that $\forall a\in A,\,a$ is in the subalgebra of \textbf{A} generated by $h(a)$. Diener \cite{bib:1989} and note 115.

\noindent[\textbf{67 G}] For every $T_{1}$ topological space $(X, T)$ and every cellular family $\mathcal{O}= \{O_{i}:i\in k\}\subseteq T-\{\emptyset\}$, there is a closed discrete set $\mathcal{I}$ such that for every $i\in k,\, \mathcal{I}\cap O_{i}\neq\emptyset$. Keremedis\index{Author}{Keremedis} \cite{bib:1996b} and note 77.

\noindent[\textbf{67 H}] Let $(X, T)$ be a $T_{4}$ topological space and let $\{G_{i}:i\in K\}$ be a locally finite family of pairwise disjoint, non-empty, open subsets of $X$. Then there exists a continuous real valued function $f:X\rightarrow \mathbb{R}$ such that for each $i\in K,\,f$ vanishes on the boundary of $G_{i}$ and $f(x_{i})=1$ for some $x_{i}\in G_{i}$. Keremedis
\cite{bib:1995} and note 43.

\noindent[\textbf{67 I}] If $(X, T)$ is a $T_{4}$ topological space and $U=\{U_{i}:i\in K\}$ is a family of non-empty open sets, then there is a family $F=\{f_{i}:i\in K\}$ of non-negative, non-constant, continuous real valued functions on $X$ such that each $f_{i}$ vanishes on $X-U_{i}$. Keremedis \cite{bib:1995}.

\noindent[\textbf{67 J}] If $(X, T)$ is a $T_{4}$ topological space and $U=\{U_{i}:i\in K\}$ is a family of non-empty open sets, then there is a family $V=\{V_{i}:i\in K\}$ of non-empty closed (compact) sets such that $U_{i}\supseteq V_{i}$ for each $i\in K$. Keremedis \cite{bib:1995}.

\noindent[\textbf{67 K}] DUN: The disjoint union of normal topological spaces is normal. Howard/Keremedis/Rubin/Rubin\index{Author}{Howard/Keremedis/Rubin/Rubin} \cite{bib:1997a} and note 141.

\noindent[\textbf{67 L}] DUU: If $C$ is a collection of pairwise disjoint topological spaces each of which satisfies Urysohn's lemma, then the disjoint union of the spaces in $C$ satisfies Urysohn's lemma. Howard/Keremedis/Rubin/Rubin \cite{bib:1997a} and note 141.

\noindent[\textbf{67 M}] DUT: If $C$ is a collection of pairwise disjoint topological spaces each of which satisfies the Tietze-Urysohn extension theorem, then the disjoint union of the spaces in $C$ satisfies the Tietze-Urysohn extension theorem.
Howard/Keremedis/Rubin/Rubin \cite{bib:1997a} and note 141.

\noindent[\textbf{67 N}] DUP: The disjoint union of paracompact spaces is paracompact. Howard/ Keremedis/Rubin/Rubin \cite{bib:1997b} and note 141.

\noindent[\textbf{67 O}] DUMET: The disjoint union of metacompact spaces is metacompact. Howard/Keremedis/Rubin/Rubin \cite{bib:1997b} and note 141.

\noindent[\textbf{67 P}] DUPFCS: The disjoint union of PFCS spaces is PFCS. ($X$ is a \emph{PFCS space} iff every point finite open cover of $X$ is shrinkable.) Howard/Keremedis/ Rubin/Rubin \cite{bib:1997b} and note 141.

\noindent[\textbf{67 Q}] DUPN: The disjoint union of paracompact spaces is normal. Howard/ Keremedis/Rubin/Rubin\index{Author}{Howard/Keremedis/Rubin/Rubin} \cite{bib:1997b} and note 141.

\noindent[\textbf{67 R}] DUcw$H$: The disjoint union of cwH spaces is cwH. Howard/Keremedis/ Rubin/Rubin \cite{bib:1997b} and note 141.

\noindent[\textbf{67 S}] DUcwN: The disjoint union of cwN spaces is cwN. Howard/Keremedis/ Rubin/Rubin \cite{bib:1997b} and note 141.

\noindent[\textbf{67 T}] DUcwNN: The disjoint union of cwN spaces is normal. Howard/Keremedis/Rubin/Rubin \cite{bib:1997b} and note 141.

\noindent[\textbf{67 U}] Every open cover $\mathcal{U}$ of a metric space ($X,\,d$) can be written as a well ordered union $\bigcup\{U_{\alpha}:\alpha\in\gamma\}$ where $\gamma$ is an ordinal and each $U_{\alpha}$ is locally finite. Howard/Keremedis/Rubin/Stanley\index{Author}{Howard/Keremedis/Rubin/Stanley} \cite{bib:1997} and note 141.

\noindent[\textbf{67 V}] Every open cover $\mathcal{U}$ of a metric space ($X,\,d$) can be written as a well ordered union $\bigcup\{U_{\alpha}:\alpha\in\gamma\}$ where $\gamma$ is an ordinal and each $U_{\alpha}$ is point finite. Howard/Keremedis/Rubin/Stanley \cite{bib:1997} and note 141.

\noindent\textbf{FORM 68.} Nielsen-Schreier Theorem:\index{Subject}{Nielsen-Schreier theorem} Every subgroup of a free group is free. Jech\index{Author}{Jech} \cite{bib:1973b} p 12.

\noindent\textbf{FORM 69.} Every field has an algebraic closure. Jech \cite{bib:1973b}
p 13.

\noindent\textbf{FORM 70.} There is a non-trivial ultrafilter on $\omega$. Jech \cite{bib:1973b} prob 5.24.

[\textbf{70 A}] Every Dedekind infinite set has a non-principal ultrafilter.

\noindent[\textbf{70 B}] Weak Universal Net Theorem: There is a universal net in $\omega$ which is not eventually constant. Schechter\index{Author}{Schechter} \cite{bib:1996a} and note 31.

\noindent[\textbf{70 C}] There is a 2-valued measure on $\mathcal{P}(\omega)$. Pincus\index{Author}{Pincus} \cite{bib:1972c} and note 147.

\noindent\textbf{FORM 71($\alpha$).} $W_{\aleph_{\alpha}}:(\forall x)(|x|\leq\aleph_{\alpha}$ or $|x|\geq\aleph_{\alpha})$. Pelc\index{Author}{Pelc} \cite{bib:1978}, Jech \cite{bib:1973b}, page 119 and Levy\index{Author}{Levy} \cite{bib:1964}.

\noindent\textbf{FORM 72.} Artin-Schreier Theorem:\index{Subject}{Artin-Schreier theorem} Every field in which $-1$ is not the sum of squares can be ordered. (The ordering, $\leq$, must satisfy (a) $a\leq b\rightarrow a+c\leq b+c$ for all $c$ and (b) $c\geq 0$ and $a\leq b\rightarrow a\cdot c\leq b\cdot c.\rangle$ Tarski\index{Author}{Tarski} \cite{bib:1954b}.

\noindent\textbf{FORM 73.} For every $A\subseteq \mathbb{R}$ and $x\in \mathbb{R}$ the following definitions are equivalent:

(1) $x$ is in the closure of $A$ iff every neighborhood of $x$ intersects $A$.

\noindent(2) $x$ is in the closure of $A$ iff there is a sequence $\{x_{n}\}\subseteq A$ such that $\displaystyle \lim x_{n}=x$. Jech \cite{bib:1973b} p 21.

\noindent[\textbf{73 A}] For all $A\subseteq \mathbb{R},\,x\in A$ and $f: A\rightarrow \mathbb{R}$ the following are equivalent:

\noindent(1) $(\forall\epsilon>0)(\exists\delta>0)(\forall y\in A)(|y-x|<\delta$ implies $|f(y)-f(x)|<\epsilon)$

\noindent(2) Whenever $\{x_{n}\}\subseteq A$ and $\lim x_{n}=x$ then $\lim f(x_{n})=f(x)$. Notes 4, 5 and 6.

\noindent[\textbf{73 B}] $(\forall f:\mathbb{R} \rightarrow \mathbb{R})(\forall x\in \mathbb{R})$ the following are equivalent:

\noindent(1) $(\forall\epsilon>0)(\exists\delta>0)((\forall y\in \mathbb{R})(|y-x|<\delta\rightarrow|f(y)-f(x)|<\epsilon)$

\noindent(2) Whenever $\displaystyle \lim x_{n}=x$, then $\lim f(x_{n})=f(x)$.
Note 10.

\noindent\textbf{FORM 74.} For every $A\subseteq \mathbb{R}$ the following are equivalent:

(1) A is closed and bounded.

\noindent(2) Every sequence $\{x_{n}\}\subseteq A$ has a convergent subsequence with limit in A. Jech \cite{bib:1973b} p 21.

\noindent\textbf{FORM 75.} Every subspace of a separable metric space is separable. Jech\index{Author}{Jech} \cite{bib:1973b} p 21.

\noindent\textbf{FORM 76.}\index{Author}{Sierpi\'{n}ski} $MC_{\omega}(\infty, \infty) (\omega\text{-}\mathrm{MC}$): For every family $X$ of pairwise disjoint nonempty sets, there is a function $f$ such that for each $x\in X,\, \mathrm{f(x)}$ is a non-empty countable subset of $x$. Howard/Keremedis/Rubin/Rubin\index{Author}{Howard/Keremedis/Rubin/Rubin} \cite{bib:1997b}.

[\textbf{76 A}] Every open cover $\mathcal{U}$ of a metric space can be written as a well ordered union $\bigcup\{U_{\alpha}:\alpha\in\gamma\}$ where $\gamma$ is an ordinal and each $U_{\alpha}$ is locally countable. Howard/Keremedis/Rubin/Stanley\index{Author}{Howard/Keremedis/Rubin/Stanley} \cite{bib:1997} and note 141.

\noindent[\textbf{76 B}] Every open cover $\mathcal{U}$ of a metric space can be written as a well ordered union $\bigcup\{U_{\alpha}:\alpha\in\gamma\}$ where $\gamma$ is an ordinal and each $U_{\alpha}$ is point countable. Howard/Keremedis/Rubin/Stanley \cite{bib:1997} and note 141.

\noindent\textbf{FORM 77.} A linear ordering of a set $P$ is a well ordering if and only if $P$ has no infinite descending sequences. Jech \cite{bib:1973b} p 23.

[\textbf{77 A}] If a linearly ordered set $(X, \leq)$ has no greatest element then it has a subset of order type $\omega$. Sierpi\'{n}ski \cite{bib:1918}.

\noindent\textbf{FORM 78.} Urysohn's Lemma:\index{Subject}{Urysohn's lemma} If $A$ and $B$ are disjoint closed sets in a normal space $S$, then there is a continuous $f:S\rightarrow[0, 1]$ which is 1 everywhere in $A$ and $0$ everywhere in $B$. Urysohn\index{Author}{Urysohn} \cite{bib:1925} pp 290-292.

\noindent\textbf{FORM 79.} $\mathbb{R}$ can be well ordered. Hilbert\index{Author}{Hilbert} \cite{bib:1900} p 263.

[\textbf{79 A}] $C(\infty, \infty, \mathbb{R})$: Every set whose elements are non-empty subsets of $\mathbb{R}$ has a choice function. Schechter\index{Author}{Schechter} \cite{bib:1996a}.

\noindent\textbf{FORM 80.} $C(\aleph_{0},2)$: Every denumerable set of pairs has a choice function. Mostowski\index{Author}{Mostowski} \cite{bib:1948}.

[\textbf{80 A}] $UT(\aleph_{0},2, \aleph_{0})$: The union of denumerably many pairs is denumerable.

\noindent\textbf{FORM 81($n$).} (For $n\in \omega\ K(n)$: For every set $S$ there is an ordinal $\alpha$ and a one to one function $f:S\rightarrow \mathcal{P}^{n}(\alpha)$. $(\mathcal{P}^{0}(X)=X$ and $\mathcal{P}^{n+1}(X)=\mathcal{P}(\mathcal{P}^{n}(X))$. $(K(0)$ is equivalent to form 1 (AC) and $K(1)$ is equivalent to the selection principle (form 15)). Monro\index{Author}{Monro} \cite{bib:1972}, \cite{bib:1973b} and note 90.

\noindent\textbf{FORM 82.} $E(I, III)$ (Howard/Yorke\index{Author}{Howard/Yorke} \cite{bib:1989}): If $X$ is infinite then $\mathcal{P}(X)$ is Dedekind infinite. ($X$ is finite $\Leftrightarrow \mathcal{P}(X)$ is Dedekind finite.) Levy\index{Author}{Levy} \cite{bib:1958} and note 94.

[\textbf{82 A}]. Every infinite set can be mapped onto $\omega$. Howard/Yorke \cite{bib:1989}, Monro \cite{bib:1972}.

\noindent[\textbf{82 B}] $E(Ia, III)$: For every set $X$, if $\mathcal{P}(X)$ is Dedekind finite then $X$ is amorphous. Howard/Yorke \cite{bib:1989}, notes 57 and 94.

\noindent\textbf{FORM 83.} $E(I, II)$ (Howard/Yorke \cite{bib:1989}): $T$-finite is equivalent to finite. Tarski\index{Author}{Tarski} \cite{bib:1924} and note 94.

[\textbf{83 A}] $E(Ia, II)$: Every $T$-finite set is amorphous. Howard/Yorke \cite{bib:1989}, notes 57 and 94.

\noindent\textbf{FORM 84.} $E(II, III)$ (Howard/Yorke \cite{bib:1989}): $(\forall x)(x$ is $T$-finite if and only if $\mathcal{P}(x)$ is Dedekind finite). Howard/Yorke \cite{bib:1989} and note 94.

\noindent\textbf{FORM 85.} $C(\infty, \aleph_{0})$: Every family of denumerable sets has a choice function. Jech \cite{bib:1973b} p 115 prob 7.13.

[\textbf{85 A}] $C(\infty, \leq\aleph_{0})$: Every family of non-empty countable sets has a choice function.

\noindent\textbf{FORM 86($\alpha$).} $C(\aleph_{\alpha}, \infty)$: If $X$ is a set of non-empty sets such that $|X|=\aleph_{\alpha}$, then $X$ has a choice function. Levy\index{Author}{Levy} \cite{bib:1964}.

[\textbf{86 A($\alpha$)}] If $(P, \leq)$ is a tree of height $\lambda\leq\aleph_{\alpha}$ which has fewer than $\lambda$ branches of height $<\lambda$ then one of the following holds for the upside down tree $(P, \geq)$:
\begin{enumerate}
\item[(a)] There is a cardinal $\mu<\lambda$ such that $(P, \geq)$ does not contain a strong antichain of cardinality $\mu$.
\item[(b)] $(P, \geq)$ contains a strong antichain of cardinality $\lambda$.
\item[(c)] $(P, \geq)$ contains a chain of cardinality $\lambda$.
\end{enumerate}
Shannon\index{Author}{Shannon} \cite{bib:1992} and note 119.

[\textbf{86 B($\alpha$)}] The frame envelope of any $\aleph_{\alpha}$ frame is $\aleph_{\alpha}$-Lindel\"{o}f. Banaschewski\index{Author}{Banaschewski} \cite{bib:1997} and note 29.

\noindent\textbf{FORM 87($\alpha$).} $DC(\aleph_{\alpha})$: Given a relation $R$ such that for every subset $Y$ of a set $X$ with $|Y|<\aleph_{\alpha}$, there is an $x\in X$ with $Y\,R\,x$ then there is a function $f:\aleph_{\alpha}\rightarrow X$ such that $(\forall\beta<\aleph_{\alpha})\ \{f(\gamma):\gamma<\beta\}\,R\,f(\beta)$. Levy \cite{bib:1964}.

[\textbf{87 A($\alpha$)}] $Z_{\aleph_{\alpha}}$: Let $P$ be a partially ordered set in which every well ordered chain has type $<\aleph_{\alpha}$. If every well ordered chain in $P$ has an upper bound in $P$, then $P$ contains a maximal element. Wolk\index{Author}{Wolk} \cite{bib:1983}.

\noindent[\textbf{87 B($\alpha$)}] $Z_{\aleph_{\alpha}}^{\ast}$: Let $P$ be a well-founded partially ordered set (that is, every chain in $P$ is well ordered.) in which every chain has type $<\aleph_{\alpha}$. If every chain in $P$ has an upper bound in $P$, then $P$ contains a maximal element. Wolk \cite{bib:1983}.

\noindent\textbf{FORM 88.} $C(\infty, 2)$: Every family of pairs has a choice function.

[\textbf{88 A}] $P(2,2)$: If $X$ is a set and $P$ is a property of subsets of $X$ of 2 character (that is, $\forall y\subseteq X(P(y)$ iff $\forall z\subseteq y(|z|\leq 2\rightarrow P(z))$)) then if every finite subset of $X$ can be partitioned into two or fewer $P$-sets (that is, sets $z$ such that $P(z)$) then $X$ can be partitioned into two or fewer $P$-sets. Cowen\index{Author}{Cowen} \cite{bib:1982}.

\noindent[\textbf{88 B}] $G_{2}$: A graph is 2-colorable if every finite subgraph is 2-colorable. Mycielski\index{Author}{Mycielski} \cite{bib:1961} and note 117.

\noindent[\textbf{88 C}] If $m$ is a cardinal number and $\mathcal{L}$ is a lattice isomorphic to the lattice of subalgebras of some unary universal algebra (a unary universal algebra is one with only unary or nullary operations) and $\alpha$ is an automorphism of $\mathcal{L}$ of order 2 (that is, $\alpha^{2}$ is the identity) then there is a unary algebra $\mathfrak{A}$ with $|$ domain $\mathfrak{A}|\geq m$ and an isomorphism $\rho$ from $\mathcal{L}$ onto the lattice of subalgebras of $\mathfrak{A}^{2}$ with $\rho(\alpha(x))=(\rho(x))^{-1}(=\{(s, t):(t, s)\in\rho(x)\})$ for all $x\in \mathcal{L}$. Lampe \cite{bib:1974}.

\noindent[\textbf{88 D}] For any cardinal $m$, there is a unary algebra $\mathfrak{A}=\langle A,\, F\rangle$ (only unary and nullary operations) and a proper subalgebra $\mathfrak{D}$ of $\mathfrak{A}^{2}$ such that $\mathfrak{A}^{2}$ has exactly four subalgebras, $\mathfrak{A},\,\emptyset,\,\mathfrak{D}$ and $\mathfrak{D}^{\ast}=\{(s, t):(t, s)\in \mathfrak{D}\}$. Lampe\index{Author}{Lampe} \cite{bib:1974}.

\noindent\textbf{FORM 89.} Antichain Principle:\index{Subject}{antichain principle} Every partially ordered set has a maximal antichain. Jech\index{Author}{Jech} \cite{bib:1973b} p 133 and Kurepa\index{Author}{Kurepa} \cite{bib:1953}.

[\textbf{89 A}] If $\forall y\in X,\,y=\langle u,\, R\rangle$ where $R$ is a partial order on $u$ then there is a function $f$ on $X$ such that $\forall\langle u,\,R\rangle\in X,\,f(\langle u, R\rangle)$ is a maximal antichain in $u$. Felgner\index{Author}{Felgner} \cite{bib:1969} and Harper/Rubin\index{Author}{Harper/Rubin} \cite{bib:1976}.

\noindent\textbf{FORM 90.} $LW$: Every linearly ordered set can be well ordered. Jech\index{Author}{Jech} \cite{bib:1973b} p 133 and Felgner/Jech\index{Author}{Felgner/Jech} \cite{bib:1973}.

\noindent\textbf{FORM 91.}\index{Author}{Sierpi\'{n}ski} $PW$:\index{Subject}{antichain principle} The power set of a well ordered set can be well ordered. Jech \cite{bib:1973b} p 133 and H. Rubin/J. Rubin\index{Author}{H. Rubin/J. Rubin} \cite{bib:1985}.

[\textbf{91 A}] $A(S, W1)$: For every $T_{2}$ topological space ($X,\, T)$ if $X$ has a dense well ordered subset, then $T$ is well ordered. (Clear) Brunner\index{Author}{Brunner} \cite{bib:1983d} and note 26.

\noindent\textbf{FORM 92.} Every subset of $\mathbb{R}$ is separable. Jech \cite{bib:1973b} p 142.

\noindent\textbf{FORM 93.} There is a non-measurable subset of $\mathbb{R}$. Sierpi\'{n}ski \cite{bib:1918}.

\noindent\textbf{FORM 94.} $C(\aleph_{0}, \infty, \mathbb{R})$: Every denumerable family of non-empty sets of reals has a choice function. Jech \cite{bib:1973b} p 148 prob 10.1.

[\textbf{94 A}] The discrete topology on $\omega$ is Lindel\"{o}f. Brunner \cite{bib:1982d} and note 40.

\noindent[\textbf{94B}] There is a Lindel\"{o}f space with a closed, discrete infinite subset. Brunner \cite{bib:1982d} and note 40.

\noindent\textbf{FORM 95 ($F$).} Existence of Complementary Subspaces over a Field $F$: If $F$ is a field, then every vector space $V$ over $F$ has the property that if $S\subseteq V$ is a subspace of $V$, then there is a subspace $S^{\prime}\subseteq V$ such that $S\cap S^{\prime}=\{0\}$ and $S\cup S^{\prime}$ generates $V$. Bleicher\index{Author}{Bleicher} \cite{bib:1964}, \cite{bib:1965}, H. Rubin/J. Rubin \cite{bib:1985}, pp 119ff, and Jech \cite{bib:1973b}, p 148 prob 10.4.

\noindent\textbf{FORM 96.} L\"{o}wig's\index{Author}{L\"{o}wig} Theorem.\index{Subject}{L\"{o}wig's theorem} If $B_{1}$ and $B_{2}$ are both bases for the vector space $V$ then $|B_{1}|=|B_{2}|$. L\"{o}wig \cite{bib:1934}.

\noindent\textbf{FORM 97.} Cardinal Representatives:\index{Subject}{cardinal representatives} For every set $A$ there is a function $c$ with domain $\mathcal{P}(A)$ such that for all $x,\,y\in \mathcal{P}(A)$, (i) $c(x)=c(y\rangle\leftrightarrow x\approx y$ and (ii) $c(x)\approx x$. Jech \cite{bib:1973b} p 154.

[\textbf{97 A}] For every set $\{X_{i}:i\in I\}$ of local cardinal numbers, there is a set $\{x_{i}:i\in I\}$ of sets such that for each $i\in I$ there exists a $y_{i}\in X_{i}$ such that $x_{i}\approx y_{i}$. ($X$ is a \emph{local cardinal number} if for all $x, y\in X, x\approx y$.) Higasikawa\index{Author}{Higasikawa} \cite{bib:1995}.

\noindent\textbf{FORM 98.} The set of all finite subsets of a Dedekind finite set is Dedekind finite. Jech \cite{bib:1973b} p 161 prob 11.5 and note 94.

\noindent\textbf{FORM 99.} Rado's\index{Author}{Rado} Selection Lemma:\index{Subject}{Rado's selection lemma} Let $\{K(\lambda):\lambda\in\Lambda\}$ be a family of finite subsets (of $X$) and suppose for each finite $S\subseteq\Lambda$ there is a function $\gamma(S):S\rightarrow X$ such that $(\forall\lambda\in S)(\gamma(S)(\lambda)\in K(\lambda))$. Then there is an $f:\Lambda\rightarrow X$ such that for every finite $S\subseteq\Lambda$ there is a finite $T$ such that $S\subseteq T\subseteq\Lambda$ and such that $f$ and $\gamma(T)$ agree on S. Rado \cite{bib:1949}.

[\textbf{99 A}] Weak Rado Lemma: Let $\Lambda$ be a set. Suppose that $\gamma$ is a function with domain all finite subsets of $\Lambda$ such that for each finite $S\subseteq\Lambda,\,\gamma(S):S\rightarrow\{0,1\}$. Then there is a $\phi$ defined on $\Lambda$ with the property that for every finite $S\subseteq\Lambda$ there is a finite $T\subseteq\Lambda$ such that $S\subseteq T$ and $\gamma(T)$ and $\phi$ agree on $S$. Schechter\index{Author}{Schechter}
\cite{bib:1996a} and Howard\index{Author}{Howard} \cite{bib:1993}.

\noindent[\textbf{99 B}] Rado's Lemma restricted to families of two element sets: Let $\{k(\lambda):\lambda\in \Lambda\}$ be a family of 2 element subsets of a set $X$ and suppose that for each finite $S\subseteq\Lambda$, there is a function $\gamma(S)$ from $S$ into $X$ such that $(\forall\lambda\in S)(\gamma(S)(\lambda)\in k(\lambda))$. Then there is a function from $\Lambda$ into $X$ such that for every finite $S\subseteq\Lambda$, there is a finite $T$ with $S\subseteq T\subseteq\Lambda$ such that $f$ and $\gamma(T)$ agree on $S$. Rav\index{Author}{Rav} \cite{bib:1977} and Howard\index{Author}{Howard} \cite{bib:1993}.

\noindent[\textbf{99 C}]\index{Subject}{Rado's selection lemma}\index{Author}{Sierpi\'{n}ski} Strong Rado Lemma 2: Let $T$ be a family of finite sets and let $\mathcal{F}$ be a collection of finite subsets of $T$ satisfying

\noindent(1) $(\forall K\in T)(\exists S\in \mathcal{F})(K\in S)$

\noindent(2) $(\forall S_{1}, S_{2}\in \mathcal{F})(S_{1}\cup S_{2}\in \mathcal{F})$

\noindent Assume there is a function $\phi$ with domain $\mathcal{F}$ such that for each $S\in \mathcal{F},\,\phi(S)= \phi_{S}$ is a function with domain $S$ and satisfies $(\forall K\in S)(\phi_{S}(K)\in K)$. Then there is a function $f$ with domain $T$ and with the following property: For all $S\in \mathcal{F}$ there is an $S^{\prime}\in \mathcal{F}$ such that $S\subseteq S^{\prime}$ and $f$ and $\phi_{S^{\prime}}$ agree on $S$. Howard \cite{bib:1993}.

\noindent\textbf{FORM 100.} Weak Partition Principle:\index{Subject}{partition principle} For all sets $x$ and $y$, if $x\precsim^{\ast}y$, then it is not the case that $y\prec x$. Lindenbaum/Tarski\index{Author}{Lindenbaum/Tarski} \cite{bib:1926} and note 69.

[\textbf{100 A}] For every set $A$ and function $f$ it is false that $A\prec f^{\prime\prime}A$. G. Moore\index{Author}{G. Moore} \cite{bib:1982}, p. 327.

\noindent[\textbf{100 B}] For all sets $x$ and $y$, if $x\precsim^{\ast}y$ and $y\precsim x$, then $x\approx y$. Higasikawa\index{Author}{Higasikawa} \cite{bib:1995} and note 69.

\noindent\textbf{FORM 101.} Partition Principle: If $S$ is a partition of $M$, then $S\precsim M$. Sierpi\'{n}ski \cite{bib:1947}.

[\textbf{101 A}] For all sets $x$ and $y,\,x\precsim^{\ast}y$ implies $x\precsim y$. Lindenbaum/Tarski \cite{bib:1926} and note 69.

\noindent\textbf{FORM 102.} For all Dedekind finite cardinals $p$ and $q$, if $p^{2}=q^{2}$ then $p=q$. Jech\index{Author}{Jech} \cite{bib:1973b} p 162 prob 11.12 and note 94.

\noindent\textbf{FORM 103.} If $(P, <)$ is a linear ordering and $|P|>\aleph_{1}$ then some initial segment of $P$ is uncountable. Jech \cite{bib:1973b} p 164 prob 11.21.

\noindent\textbf{FORM 104.} There is a regular uncountable aleph. Jech \cite{bib:1966b} p 165 prob 11.26.

\noindent\textbf{FORM 105.} There is a partially ordered set $(A, \leq)$ such that for no set $B$ is $(B, \leq)$ (the ordering on $B$ is the usual injective cardinal ordering) isomorphic to $(A, \leq)$. Forti/Honsell\index{Author}{Forti/Honsell} \cite{bib:1985} and Jech \cite{bib:1966c}.

\noindent\textbf{FORM 106.} Baire Category Theorem for Compact Hausdorff Spaces: Every compact Hausdorff space is Baire. Goldblatt\index{Author}{Goldblatt} \cite{bib:1985}, G. Moore \cite{bib:1982}, and note 28.

[\textbf{106 A}] $DMC$: Dependent Multiple Choice: Every tree in which every element has at least one immediate successor has a subtree whose levels are finite and in which every element has at least one immediate successor. Blass\index{Author}{Blass} \cite{bib:1979}, Brunner\index{Author}{Brunner} \cite{bib:1983c}, Fossy/Morillon\index{Author}{Fossy/Morillon} \cite{bib:1995} and note 21.

\noindent[\textbf{106 B}] Weak DC: If $R$ is a binary relation on a non-empty set $E$ such that $(\forall x\in E)(\exists y\in E)(x\,R\,y)$, then there is a sequence $\langle F_{n}\rangle_{n\in\omega}$ of non-empty finite subsets of $E$ such that $(\forall n\in \omega)(\forall x\in F_{n})(\exists y\in F_{n+1})(x\,R\,y)$. (This statement remains equivalent to form 106 if we replace ``binary relation'' by ``transitive binary relation''.) Fossy/Morillon \cite{bib:1995}.

\noindent[\textbf{106 C}] Every scattered compact space is a Baire space. Fossy/Morillon \cite{bib:1995} and note 28.

\noindent[\textbf{106 D}] For every compact $\mathrm{T}_{2}$ space $(X, T)$ and every family $D=\{D_{i}:i\in \omega\}$ of dense open sets of $X$ there is a regular filter base $\mathcal{F}\subseteq T$ (that is, if $F,\,G\in \mathcal{F}$, then there exists $Q\in \mathcal{F}$ such that $\overline{Q}\subseteq F\cap G$) such that for all $i\in\omega$, for all but finitely many $F\in \mathcal{F},\,\overline{F}\subset D_{i}$. Fossy/Morillon\index{Author}{Fossy/Morillon} \cite{bib:1995}

\noindent\textbf{FORM 107.} M. Hall's\index{Author}{Hall} Theorem:\index{Subject}{Hall's theorem} Let $\{S(\alpha):\alpha\in A\}$ be a collection of finite subsets (of a set $X$) then if
\begin{equation*}
\tag{$\ast$} \text{for each finite}\ F\subseteq A\ \text{there is an injective choice function on}\ F
\end{equation*}

\noindent then there is an injective choice function on $A$. (That is, a 1-1 function $f$ such that $(\forall\alpha\in A)(f(\alpha)\in S(\alpha)).)$ (According to a theorem of P. Hall $(\ast)$ is equivalent to $|\bigcup_{\alpha\in F}S(\alpha)|\geq|F|$. P. Hall's theorem does not require the axiom of choice.) M. Hall \cite{bib:1948}.

\noindent\textbf{FORM 108.} There is an ordinal $\alpha$ such that $2^{\aleph_{\alpha}}$ is not the union of a denumerable set of denumerable sets. Gitik\index{Author}{Gitik} \cite{bib:1985}.

\noindent\textbf{FORM 109.} Every field $F$ and every vector space $V$ over $F$ has the property that each linearly independent set $A\subseteq V$ can be extended to a basis. Bleicher\index{Author}{Bleicher} \cite{bib:1964}, \cite{bib:1965}, Blass\index{Author}{Blass} \cite{bib:1984a}, H. Rubin/J. Rubin\index{Author}{H. Rubin/J. Rubin} \cite{bib:1985}, pp 119 ff.

\noindent\textbf{FORM 110.} Every vector space over $\mathbb{Q}$ has a basis. Blass \cite{bib:1984a}.

\noindent\textbf{FORM 111.} $UT(WO, 2, WO)$: The union of an infinite well ordered set of 2-element sets is an infinite well ordered set.

[\textbf{111 A}] $C(WO, 2)$: Every well ordered family of 2-element sets has a choice function.

\noindent\textbf{FORM 112.} $MC(\infty, LO)$: For every family $X$ of non-empty sets each of which can be linearly ordered there is a function $f$ such that for all $y\in X,\,f(y)$ is a non-empty finite subset of $y$. Howard/Rubin\index{Author}{Howard/Rubin} \cite{bib:1977}.

\noindent\textbf{FORM 113.} Tychonoff's\index{Subject}{Tychonoff's theorem} Compactness Theorem for Countably Many Spaces: The product of a countable set of compact spaces is compact. Brunner\index{Author}{Brunner} \cite{bib:1985a}.

\noindent\textbf{FORM 114.} Every A-bounded $T_{2}$ topological space is weakly Loeb. ($A$-\emph{bounded} means amorphous subsets are relatively compact. \emph{Weakly Loeb} means the set of non-empty closed subsets has a multiple choice function.) Brunner \cite{bib:1985a}.

\noindent\textbf{FORM 115.} The product of weakly Loeb $T_{2}$ spaces is weakly Loeb. (See form 114 for the definition of weakly Loeb.) Brunner \cite{bib:1985a}.

\noindent\textbf{FORM 116.} Every compact $T_{2}$ space is weakly Loeb. (See form 114 for the definition of weakly Loeb.) Brunner \cite{bib:1985a}.

\noindent\textbf{FORM 117.} If $\kappa$ is a measurable cardinal, then $\kappa$ is the $\kappa$th inaccessible cardinal. Keisler/Tarski\index{Author}{Keisler/Tarski} \cite{bib:1964} and note 20.

\noindent \textbf{FORM 118.} Every linearly orderable topological space is normal. Birkhoff\index{Author}{Birkhoff} \cite{bib:1967} p 241.

[\textbf{118 A}] For each complete linear order, the set of non-empty convex subsets has a choice function. van Douwen\index{Author}{Douwen} \cite{bib:1985}.

\noindent[\textbf{118 B}] For each complete linear order, each pairwise disjoint collection of non-empty, convex, open subsets has a choice function. van Douwen \cite{bib:1985}.

\noindent[\textbf{118 C}] Every linearly orderable space is collectionwise normal. van Douwen\index{Author}{Douwen} \cite{bib:1985} and note 71.

\noindent[\textbf{118 D}] Every linearly orderable space is collectionwise Hausdorff. van Douwen \cite{bib:1985} and note 71.

\noindent[\textbf{118 E}] Every linearly orderable space satisfies Urysohn's lemma, van Douwen \cite{bib:1985} and note 141.

\noindent[\textbf{118 F}] Every linearly orderable space satisfies the Tietze-Urysohn extension theorem, van Douwen \cite{bib:1985} and note 141.

\noindent[\textbf{118 G}] Each complete linear order admits an extreme choice, that is, if $(X, \leq)$ is a linear order then there is a choice function $f$ on the set of non-empty open interval of $X$ with the property that for all $(a, b, c)\in X^{3}$ with $a<b<c,\,f((a, c))\in\{f((a, b)), b, f((b, c))\}$. Morillon\index{Author}{Morillon} \cite{bib:1991a}.

\noindent[\textbf{118 H}] For each linear order $(X, \leq)$, there is an ordinal $\alpha$ and a strictly increasing injective map $j$ from $(X, \leq)$ to $\{0,1\}^{\alpha}$ with the lexicographic order. Morillon \cite{bib:1987}.

\noindent[\textbf{118 I}] Every linearly ordered set of linear orders, each of which is non-empty and conditionally complete, has a choice function. Morillon \cite{bib:1988}, van Douwen \cite{bib:1985}, and note 71.

\noindent[\textbf{118 J}] If $(X, \leq)$ is a linear, complete order, then the order topology on $X$ is hereditarily normal. Morillon \cite{bib:1988}, van Douwen \cite{bib:1985}, and note 71.

\noindent[\textbf{118 K}] If $(X, \leq)$ is a linear order, then the order topology on $X$ is highly separated. Morillon \cite{bib:1988} and note 71.

\noindent[\textbf{118 L}] If $X$ is a linear, complete order, then the order topology on $X$ is highly separated. Morillon \cite{bib:1988} and note 71.

\noindent[\textbf{118 M}] If $(X, \leq)$ is a linear order, then the family $\{A(a, b):a, b\in X\wedge a<b\}$ has a choice function. $(A(a, b)$ denotes the set of continuous, order preserving functions from $X$ into $[0,1]$ such that $f(a)=0$ and $f(b)=1.)$ Morillon \cite{bib:1988} and note 71.

\noindent[\textbf{118 N}] If $(X, \leq)$ is a linear order then the family $\{A^{\prime}(a, b):a, b\in X\wedge a<b\}$ has a choice function. $(A^{\prime}(a, b)$ denotes the set of continuous functions from $X$ into $[0,1]$ such that $f(a)=0$ and $f(b)=1.)$ Morillon \cite{bib:1988} and note 71.

\noindent[\textbf{118 O}] If $(X, \leq)$ is a linear order, then the order topology on $X$ is effectively normal. Morillon \cite{bib:1988} and note 71.

\noindent[\textbf{118 P}] If $(X, \leq)$ is a complete linear order, then the order topology on $X$ is effectively normal. Morillon \cite{bib:1988} and note 71.

\noindent[\textbf{118 Q}] If $(X, \leq)$ is a linear order, then the order topology on $X$ is hereditarily normal. Morillon \cite{bib:1988} and note 71.

\noindent[\textbf{118 R}] If $(X, \leq)$ is a linear order, then the order topology on $X$ is monotonely normal. Morillon \cite{bib:1987} and note 71.

\noindent[\textbf{118 S}] If $(X, \leq)$ is a complete linear order then there is a choice function $\ast$ for the family of non-empty open intervals of $X$ such that if $a,\,b,\,c\in X$ and either $(a, b)\neq\emptyset$ or $(b, c)\neq\emptyset$ then $\ast((a, c))\in\{\ast((a, b)), \ast((b, c)), b, b^{+}\}$ (where $b^{+}$ is the successor of $b$ if it exists and $b$ otherwise). Morillon \cite{bib:1987} and note 71.

\noindent[\textbf{118 T}] Every complete linear order has an increasing choice function. Morillon \cite{bib:1987} and note 71.

\noindent[\textbf{118 U}] Every linearly ordered space is monotonically normal. (A space $(X, <)$ is \emph{monotonically normal} if there is an operator $V$ which assigns to each $x\in X$ and basic open neighborhood $U$ of $x$ a basic open neighborhood $V(x, U)$ of $x$ such that $V(x, U)\cap V(x^{\prime}, U^{\prime})\neq\emptyset$ implies $x\in U^{\prime}$ or $x^{\prime}\in U$.) Good/Tree\index{Author}{Good/Tree} \cite{bib:1995}.

\noindent [\textbf{118 V}]\index{Subject}{van Douwen's choice principle} If $(L, \leq)$ is a linear order and $G=\{[l_{i}, r_{i}]:i\in K\}$ is a locally finite family of closed intervals of $L$ whose interiors are non-empty and pairwise disjoint, then there exists a continuous real valued function $f:L\rightarrow \mathbb{R}$ such that $|f^{-1}(0) \cap [l_{i}, r_{i}]|\geq 2$ and $|f^{-1}(1) \cap [l_{i}, r_{i}]|\geq 1$ for all $i\in K$. Keremedis\index{Author}{Keremedis} \cite{bib:1995} and note 43.

\noindent[\textbf{118 W}] If $(L, \leq)$ is a linear order and $G=\{[l_{i}, r_{i}]:i\in K\}$ is a locally finite family of closed intervals of $L$ whose interiors are non-empty and pairwise disjoint, then there exists a continuous real valued function $f:L\rightarrow \mathbb{R}$ such that for each $i\in K,\,f(l_{i})=f(r_{i})=0$ and $f(x)>0$ for all $x\in(l_{i}, r_{i})$. Keremedis
\cite{bib:1995} and note 43.

\noindent[\textbf{118 X}] If $(L, \leq)$ is a complete linear order and $G=\{(l_{i}, r_{i}):i\in K\}$ is a family of open, non-empty, pairwise disjoint intervals of $L$, then there exists a set $\{f_{i}:i\in K\}$ of non constant continuous real valued functions on $L$ such that for each $i\in K,\,f_{i}$ vanishes outside of $(l_{i}, r_{i})$. Keremedis \cite{bib:1995}.

\noindent\textbf{FORM 119.} $C(\aleph_{0}$, uniformly orderable with order type of the integers): Suppose $\{A_{i}:i\in\omega\}$ is a set and there is a function $f$ such that for each $i\in \omega,\,f(i)$ is an ordering of $A_{i}$ of type $\omega^{\ast}+\omega$ (the usual ordering of the integers), then $\{A_{i}:i\in \omega\}$ has a choice function. van Douwen\index{Author}{Douwen} \cite{bib:1985}.

\noindent\textbf{FORM 120($K$).} If $K\subseteq\omega -\{0,1\},\, C(LO, K)$: Every linearly ordered set of non-empty sets each of whose cardinality is in $K$ has a choice function. Truss \cite{bib:1973a}.

\noindent\textbf{FORM 121.} $C(LO, <\aleph_{0})$: Every linearly ordered set of non-empty finite sets has a choice function. Truss\index{Author}{Truss} \cite{bib:1973a}.

\noindent\textbf{FORM 122.} $C(WO, <\aleph_{0})$:  Every well ordered set of non-empty finite sets has a choice function. Truss \cite{bib:1973a}.

[\textbf{122 A}] $UT(WO, <\aleph_{0}, WO)\ (U_{\aleph_{0}})$: The union of a well ordered set of finite sets is well orderable. Brunner/Howard\index{Author}{Brunner}\index{Author}{Brunner/Howard} \cite{bib:1992}. (See note 27 for more information on $UT(WO, \kappa, WO)$, $\kappa$ a well ordered cardinal. The proof that 122 implies [122 A] is similar to the proof given in Shannon\index{Author}{Shannon} \cite{bib:1988} p 569.)

\noindent\textbf{FORM 123.} $SPI^{\ast}$: Uniform weak ultrafilter principle: For each family $F$ of infinite sets $\exists f$ such that $\forall x\in F,\,f(x)$ is a non-principal ultrafilter on $x$. Halpern\index{Author}{Halpern} \cite{bib:1972}

[\textbf{123 A}] Form 70 + form 344. Morillon\index{Author}{Morillon} \cite{bib:1988}.

\noindent\textbf{FORM 124.} Every operator on a Hilbert space with an amorphous base is the direct sum of a finite matrix and a scalar operator. (A set is amorphous if it is not the union of two disjoint infinite sets.) Brunner \cite{bib:1984a}.

\noindent\textbf{FORM 125.} There does not exist an infinite, compact connected $p$ space. (A $p$ \emph{space} is a $T_{2}$ space in which the intersection of any well orderable family of open sets is open.) Brunner \cite{bib:1984c}.

\noindent\textbf{FORM 126.} $MC(\aleph_{0}, \infty)$, Countable axiom of multiple choice: For every denumerable set $X$ of non-empty sets there is a function $f$ such that for all $y\in X,\,f(y)$ is a non-empty finite subset of $y$. Brunner\index{Author}{Brunner} \cite{bib:1984d} and Blass\index{Author}{Blass} \cite{bib:1979}.

[\textbf{126 A}] Every positive functional on a $B$-lattice is continuous. Brunner \cite{bib:1984d} and note 16.

\noindent[\textbf{126 B}] Form 131 + form 350. Keremedis\index{Author}{Keremedis} \cite{bib:1996d} and note 132.

\noindent[\textbf{126 C}] $FSCB$ + form 131. FSCB is ``Every separable first countable space $(X, T)$ has a denumerable $\pi$-base.'' Keremedis \cite{bib:1996d} and note 132.

\noindent[\textbf{126 D}] For every denumerable family $A$ of disjoint non-empty sets there is an infinite set $C\subset\bigcup A$ such that for every $a\in A,\, 0\leq|C\cap a|<\omega$. Keremedis \cite{bib:1996d} and note 132.

\noindent[\textbf{126 E}] If $(X, T)$ is a topological space, then for every family $D=\{D_{i}:i\in \omega\}$ of dense open sets of $X$, there is a countable dense set $S\subseteq X$ such that for all $i\in\omega$ and for all but finitely many $s\in S,\,s\in D_{i}$. Keremedis \cite{bib:1996d} and note 132.

\noindent[\textbf{126 F}] Every denumerable family of sets having the sfip also has a pseudointersection. Keremedis \cite{bib:1996d} and note 132.

\noindent[\textbf{126 G}] DUN$(w)$: The countable disjoint union of normal topological spaces is normal. Howard/Keremedis/Rubin/Rubin\index{Author}{Howard/Keremedis/Rubin/Rubin} \cite{bib:1997a} and note 141.

\noindent[\textbf{126 H}] DUU$(w)$: If $C$ is a countable collection of pairwise disjoint topological spaces each of which satisfies Urysohn's lemma, then the disjoint union of the spaces in $C$ satisfies Urysohn's lemma. Howard/Keremedis/Rubin/Rubin \cite{bib:1997a} and note 141.

\noindent[\textbf{126 I}] DUT$(\omega)$: If $C$ is a countable collection of pairwise disjoint topological spaces each of which satisfies the Tietze-Urysohn extension theorem, then the disjoint union of the spaces in $C$ satisfies the Tietze-Urysohn extension theorem. Howard/Keremedis/Rubin/Rubin \cite{bib:1997a} and note 141.

\noindent[\textbf{126 J}] $F_{\sigma}$ subsets of paracompact spaces are paracompact. Howard/Keremedis/Rubin/Rubin \cite{bib:1997b} and note 141.

\noindent[\textbf{126 K}] $F_{\sigma}$ subsets of metacompact spaces are metacompact. Howard/Keremedis/Rubin/Rubin \cite{bib:1997b} and note 141.

\noindent\textbf{FORM 127.} An amorphous power of a compact $T_{2}$ space, which as a set is well orderable, is well orderable. Brunner \cite{bib:1984b}.

\noindent\textbf{FORM 128.} Aczel's Realization Principle:\index{Subject}{Aczel's Realization Principle} On every infinite set there is a Hausdorff topology with an infinite set of non-isolated points. Brunner \cite{bib:1984f}.

\noindent\textbf{FORM 129.} For every infinite set $A,\,A$ admits a partition into sets of order type $\omega^{\ast}+\omega$. (For every infinite $A$, there is a set $\{\langle C_{j}, <j\rangle:j\in J\}$ such that $\{C_{j}:j\in J\}$ is a partition of $A$ and for each $j\in J, <_{j}$ is an ordering of $C_{j}$ of type $\omega^{\ast}+\omega.)$ Brunner \cite{bib:1984f}.

\noindent\textbf{FORM 130.} $\mathcal{P}(\mathbb{R})$ is well orderable. Brunner \cite{bib:1984f}.

\noindent\textbf{FORM 131.} $MC_{\omega}(\aleph_{0}, \infty)$: For every denumerable family $X$ of pairwise disjoint non-empty sets, there is a function $f$ such that for each $x\in X,\,\mathrm{f(x)}$ is a non-empty countable subset of $x$. Howard/Keremedis/Rubin/Rubin \cite{bib:1997b}.

\noindent\textbf{FORM 132.} $PC(\infty<\aleph_{0}, \infty)$: Every infinite family of finite sets has an infinite subfamily with a choice function. Blass \cite{bib:1977a} and Kleinberg\index{Author}{Kleinberg} \cite{bib:1969}.

[\textbf{132 A}] Every set mapping $f:X\rightarrow[X]^{WO}$ (well orderable subsets of $X$) on a Dedekind finite, infinite set $X$ has an infinite free subset. Brunner\index{Author}{Brunner} \cite{bib:1989}. notes 22 and 94.

\noindent\textbf{FORM 133.} Every set is either well orderable or has an infinite amorphous subset. Brunner \cite{bib:1982a}.

\noindent\textbf{FORM 134.} If $X$ is an infinite $T_{1}$ space and $X^{Y}$ is $T_{5}$, then $Y$ is countable. $(T_{5}$ is ``hereditarily $T_{4}$''.) A. H. Stone\index{Author}{A. H. Stone} \cite{bib:1948}.

\noindent\textbf{FORM 135.} If $X$ is a $T_{2}$ space with at least two points and $X^{Y}$ is hereditarily metacompact then $Y$ is countable. (A space is \emph{metacompact} if every open cover has an open point finite refinement. If $B$ and $B^{\prime}$ are covers of a space $X$, then $B^{\prime}$ is a \emph{refinement} of $B$ if $(\forall x\in B^{\prime})(\exists y\in B)(x\subseteq y)$. $B$ is \emph{point finite} if $(\forall t\in X)$ there are only finitely many $x\in B$ such that $t\in x$.) van Douwen\index{Author}{Douwen} \cite{bib:1980}.

\textbf{FORM 136($k$).} Surjective Cardinal Cancellation (depends on $k\in \omega-\{0\})$: For all cardinals $x$ and $y,\,kx\leq^{\ast}ky$ implies $x\leq^{\ast}y$. Truss\index{Author}{Truss} \cite{bib:1984} and note 69.

\noindent\textbf{FORM 137($k$).} Suppose $k\in \omega-\{0\}$. If $f$ is a 1-1 map from $k\times X$ into $k\times Y$ then there are partitions $X=\bigcup_{i<k}X_{i}$ and $Y=\bigcup_{i<k}Y_{i}$ of $X$ and $Y$ such that $f$ maps $\bigcup_{i<k}(\{i\}\times X_{i})$ onto $\bigcup_{i<k}(\{i\}\times Y_{i})$. Truss \cite{bib:1984}.

[\textbf{137 A($k$)}] For all $X$ and $Y\subseteq \mathbb{R}$ form 137($k$) holds. Truss \cite{bib:1984}

\noindent \textbf{FORM 138($k$).} Suppose $k\in \omega$. If $f$ is a partial map from $k\times Y$ onto $k\times X$ (that is, the domain is a subset of $k\times Y$), then there are partitions $X=\bigcup_{i<k}X_{i}$ and $Y=\bigcup_{i<k}Y_{i}$ of $X$ and $Y$ such that $f$ maps $\bigcup_{i<k}(\{i\}\times Y_{i})$ onto $\bigcup_{i<k}(\{i\}\times X_{i})$. Truss \cite{bib:1984}.

[\textbf{138 A($k$)}] For all standard $X$ and $Y$, form $138(k)$ holds. Truss \cite{bib:1984}. This is equivalent to $138(k)$ in $\mathrm{ZF}^{0}$.

\noindent\textbf{FORM 139.} Using the discrete topology on 2, $2^{\mathcal{P}(\omega)}$ is compact. Truss \cite{bib:1984}.

\noindent\textbf{FORM 140.} Let $\Omega$ be the set of all (undirected) infinite cycles of reals (Graphs whose vertices are real numbers, connected, no loops and each vertex adjacent to exactly two others). Then there is a function $f$ on $\Omega$ such that for all $s\in\Omega,\,f(s)$ is a direction along $s$. Truss \cite{bib:1984}.

[\textbf{140 A}] (Form 137 with $k=2$): If $f$ is a 1-1 map from $2\times X$ into $2\times Y$ then there are partitions $X=X_{0}\cup X_{1}$ and $Y=Y_{0}\cup Y_{1}$ of $X$ and $Y$ such that $f$ maps $(\{0\}\times X_{0})\cup(\{1\}\times X_{1})$ onto $(\{0\}\times Y_{0})\cup(\{1\}\times Y_{1})$. Truss \cite{bib:1984}.

\noindent\textbf{FORM 141.} $[14\ \mathrm{P}(n)]$ with $n=2$: Let $\{A(i):i\in I\}$ be a collection of sets such that $\forall i\in I,\ |A(i)|\leq 2$ and suppose $R$ is a symmetric binary relation on $\bigcup_{i\in I}A(i)$ such that for all finite $W\subseteq I$ there is an $R$ consistent choice function for $\{A(i):i\in W\}$. Then there is an $R$ consistent choice function for $\{A(i):i\in I\}$. Cowen\index{Author}{Cowen} \cite{bib:1977b} and note 109.

\noindent\textbf{FORM 142.} $\neg PB$: There is a set of reals without the property of Baire. Jech\index{Author}{Jech} \cite{bib:1973b} p 7 and note 28.

\noindent\textbf{FORM 143.} $H(C, TR)$: If $(X, R)$ is a connected relation $(u\neq v\rightarrow u\ R\ v$ or $v\ R\ u)$ then $X$ contains a $\subseteq$-maximal transitive subset. Harper/Rubin\index{Author}{Harper/Rubin} \cite{bib:1976} and note 39.

\noindent\textbf{FORM 144.} Every set is almost well orderable. Keisler\index{Author}{Keisler} \cite{bib:1970} and note 25.

[\textbf{144 A}] Every $p$ space is discrete. (A $p$ \emph{space} is a $T_{2}$ space in which the intersection of any well orderable family of open sets is open). Brunner\index{Author}{Brunner} \cite{bib:1984c}.

\noindent\textbf{FORM 145.} Compact $P_{0}$-spaces are Dedekind finite. (A $P_{0}$-\emph{space} is a topological space in which the intersection of a countable collection of open sets is open.) Brunner \cite{bib:1984c} and note 94.

\noindent\textbf{FORM 146.} $A(F, A1)$: For every $T_{2}$ topological space $(X, T)$, if $X$ is a continuous finite to one image of an A1 space then $(X, T)$ is an A1 space. $((X, T)$ is A1 means if $U\subseteq T$ covers $X$ then $\exists f:X\rightarrow U$ such that $(\forall x\in X)(x\in f(x)).)$ Brunner \cite{bib:1983d} and note 26.

\noindent\textbf{FORM 147.} $A(D2)$: Every $T_{2}$ topological space $(X, T)$ can be covered by a well ordered family of discrete sets. Brunner \cite{bib:1983d} and note 26.

\noindent\textbf{FORM 148.} $A(W2, B1)$: For every $T_{2}$ topological space $(X, T)$, if $X$ is well ordered, then $X$ has a well ordered base. Brunner \cite{bib:1983d} and note 26.

\noindent\textbf{FORM 149.} $A(F)$: Every $T_{2}$ topological space is a continuous, finite to one image of an A1 space. Brunner \cite{bib:1983d} and note 26.

\noindent\textbf{FORM 150.} $PC(\infty, \aleph_{0}, \infty)$: Every infinite set of denumerable sets has an infinite subset with a choice function.

\noindent\textbf{FORM 151.} $UT(WO, \aleph_{0}, WO)\ (U_{\aleph_{1}})$: The union of a well ordered set of denumerable sets is well orderable. Brunner/Howard\index{Author}{Brunner/Howard} \cite{bib:1992}. (If $\kappa$ is a well ordered cardinal, see note 27 for $UT(WO, \kappa, WO).)$

\noindent\textbf{FORM 152.} $D_{\aleph_{0}}$: Every non-well-orderable set is the union of a pairwise disjoint, well orderable family of denumerable sets. Brunner/Howard \cite{bib:1992}. (See note 27 for $D_{\kappa},\,\kappa$ a well ordered cardinal.)

\noindent\textbf{FORM 153.} The closed unit ball of a Hilbert space is compact in the weak topology. Fossy/Morillon\index{Author}{Fossy/Morillon} \cite{bib:1995} and Rudin\index{Author}{Rudin} \cite{bib:1991}.

\noindent\textbf{FORM 154.} Tychonoff's\index{Subject}{Tychonoff's theorem} Compactness Theorem for Countably Many $T_{2}$ Spaces: The product of countably many $T_{2}$ compact spaces is compact. Brunner \cite{bib:1983c}.

\noindent\textbf{FORM 155.} $LC$: There are no non-trivial L\"{a}uchli continua. (A \emph{L\"{a}uchli continuum} is a strongly connected continuum. \emph{Continuum} $\equiv$ compact, connected, Hausdorff space; and \emph{strongly connected} $\equiv$ every continuous real valued function is constant.) Brunner \cite{bib:1983a} and L\"{a}uchli\index{Author}{L\"{a}uchli} \cite{bib:1962}.

\noindent\textbf{FORM 156.} Theorem of Gelfand and Kolmogoroff: Two compact $T_{2}$ spaces are homeomorphic if their rings of real valued continuous functions are isomorphic. Brunner \cite{bib:1983a} and Gelfand/Kolmogoroff\index{Author}{Gelfand/Kolmogoroff} \cite{bib:1939}.

\noindent\textbf{FORM 157.} Theorem of Goodner:\index{Subject}{Goodner's theorem} A compact $T_{2}$ space is extremally disconnected (the closure of every open set is open) if and only if each non-empty subset of $C(X)$ (set of continuous real valued functions on $X$) which is pointwise bounded has a supremum. Goodner\index{Author}{Goodner} \cite{bib:1950}.

\noindent\textbf{FORM 158.} In every Hilbert space $H$, if the closed unit ball is sequentially compact, then $H$ has an orthonormal basis. Brunner \cite{bib:1983b}.

\noindent\textbf{FORM 159.}\index{Author}{Sierpi\'{n}ski} The regular cardinals are cofinal in the class of ordinals. Blass\index{Author}{Blass} \cite{bib:1983b}.

[\textbf{159 A}]\index{Subject}{Gelfand-Kolmogoroff theorem} Every variety has an initial algebra. Blass \cite{bib:1983b} and note 32.

\noindent[\textbf{159 B}] For every variety $V$ and set $X,\,V$ has a free algebra on $X$. Blass \cite{bib:1983b} and note 32.

\noindent\textbf{FORM 160.} No Dedekind finite set can be mapped onto an aleph. (See [9 E].) Monro\index{Author}{Monro} \cite{bib:1975} and note 94.

\noindent\textbf{FORM 161.} Definability of cardinal addition in terms of $\leq$: There is a first order formula whose only non-logical symbol is $\leq$ (for cardinals) that defines cardinal addition. H\"{a}ussler\index{Author}{H\"{a}ussler} \cite{bib:1983}.

\noindent\textbf{FORM 162.} Non-existence of infinite units: There is no infinite cardinal number $A$ such that $A+A>A$ and for all cardinals $x$ and $y,\,x+y=A\rightarrow x=A$ or $y=A$. H\"{a}ussler \cite{bib:1983}.

\noindent\textbf{FORM 163.} Every non-well-orderable set has an infinite, Dedekind finite subset. Brunner\index{Author}{Brunner} \cite{bib:1982a} and note 94.

\noindent\textbf{FORM 164.} Every non-well-orderable set has an infinite subset with a Dedekind finite power set. Brunner \cite{bib:1982a}, notes 42 and 94.

\noindent\textbf{FORM 165.} $C(WO, WO)$: Every well ordered family of non-empty, well orderable sets has a choice function.

\noindent\textbf{FORM 166.} $PC(\infty, 2, \infty)$: Every infinite family of pairs has an infinite subfamily with a choice function. Brunner \cite{bib:1982a}.

\noindent\textbf{FORM 167.} $PKW(\aleph_{0}, \geq 2, \infty)$, Partial Kinna-Wagner Principle: For every denumerable family $F$ such that for all $x\in F,\,|x|\geq 2$, there is an infinite subset $H\subseteq F$ and a function $f$ such that for all $x\in H,\,\emptyset \neq f(x)\subsetneq x$. Brunner \cite{bib:1982a}.

\noindent\textbf{FORM 168.} Dual Cantor-Bernstein\index{Author}{Bernstein} Theorem:\index{Subject}{Cantor-Bernstein theorem} $(\forall x)(\forall y)(|x|\leq^{\ast}|y|$ and $|y|\leq^{\ast}|x|$ implies $|x|=|y|)$. Banaschewski/Moore\index{Author}{Banaschewski/Moore} \cite{bib:1990} and note 69.

\noindent\textbf{FORM 169.} There is an uncountable subset of $\mathbb{R}$ without a perfect subset. Banaschewski/Moore \cite{bib:1990}, Bernstein \cite{bib:1908}, and note 28.

\noindent\textbf{FORM 170.} $\aleph_{1}\leq 2^{\aleph_{0}}$. Bernstein \cite{bib:1908} and Banaschewski/Moore \cite{bib:1990}.

[\textbf{170 A}] There is a function that assigns to each denumerable set $D$ of real numbers a real number not in $D$. Sierpi\'{n}ski \cite{bib:1954}.

\noindent[\textbf{170 B}] There is an uncountable, well orderable set of reals. Raisonnier\index{Author}{Raisonnier} \cite{bib:1982} and note 3.

\noindent\textbf{FORM 171.} If $(P, \leq)$ is a partial order such that $P$ is the denumerable union of finite sets and all antichains in $P$ are finite then for each denumerable family $\mathcal{D}$ of dense sets there is a $\mathcal{D}$ generic filter. Shannon\index{Author}{Shannon} \cite{bib:1990} and note 47.

\noindent\textbf{FORM 172.} For every infinite set $S$, if $S$ is hereditarily countable (that is, every $y\in TC(S)$ is countable) then $|TC(S)|=\aleph_{0}$. Jech\index{Author}{Jech} \cite{bib:1982}.

\noindent\textbf{FORM 173.} MPL: Metric spaces are para-Lindel\"{o}f. Howard/Keremedis/Rubin/ Stanley\index{Author}{Howard/Keremedis/Rubin/Stanley} \cite{bib:1997} and note 141.

[\textbf{173 A}] MML: Metric spaces are meta-Lindel\"{o}f. Howard/Keremedis/Rubin/ Stanley \cite{bib:1997} and note 141.

\noindent\textbf{FORM 174($\alpha$).} $RM1,\,\aleph_{\alpha}$: The representation theorem for multi-algebras with $\aleph_{\alpha}$ unary operations: Assume $(A, F)$ is a multi-algebra with $\aleph_{\alpha}$ unary operations (and no other operations). Then there is an algebra $(B, G)$ with $\aleph_{\alpha}$ unary operations and an equivalence relation $E$ on $B$ such that $(B/E, G/E)$ and $(A, F)$ are isomorphic multi-algebras. H\"{o}ft/Howard\index{Author}{H\"{o}ft/Howard} \cite{bib:1981} and note 50.

\noindent\textbf{FORM 175.} Transitivity Condition: For all sets $x$, there is a set $u$ abd a function $f$ such that $u$ is transitive and $f$ is a one to one function from $x$ onto $u$. von Rimscha\index{Author}{Rimscha} \cite{bib:1982}.

\noindent\textbf{FORM 176.} Every infinite, locally finite group has an infinite Abelian subgroup. (\emph{Locally finite} means every finite subset generates a finite subgroup.) Plotkin\index{Author}{Plotkin} \cite{bib:1981}.

\noindent\textbf{FORM 177.} An infinite box product of regular $T_{1}$ spaces, each of cardinality greater than 1, is neither first countable nor connected. Brunner\index{Author}{Brunner} \cite{bib:1981a}, Knight\index{Author}{Knight} \cite{bib:1964}, and note 52.

\noindent\textbf{FORM 178($n, N$).} If $n\in\omega$, $n\geq 2$ and $N\subseteq\{1,2, \ldots, n-1\},\, N\neq\emptyset,\, MC(\infty, n, N)$: If $X$ is any set of $n$-element sets then there is a function $f$ with domain $X$ such that for all $A\in X,\,f(A)\subseteq A$ and $|f(A)|\in N$. Zuckerman\index{Author}{Zuckerman} \cite{bib:1969a}, \cite{bib:1981}, and note 53.

\noindent\textbf{FORM 179($\epsilon$).} Suppose $\epsilon>0$ is an ordinal. $\forall x, x\in W(\epsilon)$. David\index{Author}{David} \cite{bib:1980} and note 25.

\noindent \textbf{FORM 180.} Every Abelian group has a divisible hull. (If $A$ and $B$ are groups, $B$ \emph{is a divisible hull of} $A$ means $B$ is a divisible group, $A$ is a subgroup of $B$ and for every non-zero $b\in B,\, \exists n\in\omega$ such that $0\neq nb\in A$.) Fuchs \cite{bib:1970}, Theorem~24.4 p 107, Hodges\index{Author}{Hodges} \cite{bib:1980}, and note 24.

\noindent\textbf{FORM 181.} $C(2^{\aleph_{0}}, \infty)$: Every set $X$ of non-empty sets such that $|X|=2^{\aleph_{0}}$ has a choice function. Freyd\index{Author}{Freyd} \cite{bib:1980}.

\noindent\textbf{FORM 182.} There is an aleph whose cofinality is greater than $\aleph_{0}$. Gitik\index{Author}{Gitik} \cite{bib:1980}.

\noindent\textbf{FORM 183($\alpha$).} There are no $\aleph_{\alpha}$ minimal sets. That is, there are no sets $X$ such that

(1) $|X|$ is incomparable with $\aleph_{\alpha}$

\noindent(2) $\aleph_{\beta}<|X|$ for every $\beta<\alpha$ and

\noindent(3) $\forall Y\subseteq X,\,|Y|<\aleph_{\alpha}$ or $|X-Y|<\aleph_{\alpha}$.

\noindent Hickman\index{Author}{Hickman} \cite{bib:1980b} and note 56.

\noindent\textbf{FORM 184.} Existence of a double uniformization: For all $X$ and $Y$, for all $E\subseteq X\times Y$, if there is an infinite cardinal $\kappa$ satisfying:

\noindent(1) $\forall y_{0}\in Y,\,|\{(x, y_{0}):x\in X\}\cap E|=\kappa$ and

\noindent(2) $\forall x_{0}\in X,\,|\{(x_{0}, y):y\in Y\}\cap E|=\kappa$,

\noindent then $\exists U\subseteq E$ such that for all $(x, y)\in E,\,\exists!z\in X$ such that $(z, y)\in U$ and $\exists!z\in Y$ such that $(x, z)\in U$. ($U$ is called a \emph{double uniformization} of $E$.) Kaniewski/Rogers\index{Author}{Kaniewski/Rogers} \cite{bib:1980}.

\noindent\textbf{FORM 185.} Every linearly ordered Dedekind finite set is finite. Hickman \cite{bib:1979a} and note 94.

[\textbf{185 A}] If $X$ is a linearly ordered set and there is neither an increasing nor a decreasing injection $f:\omega\rightarrow X$, then $X$ is finite. Hickman\index{Author}{Hickman} \cite{bib:1979a}.

\noindent\textbf{FORM 186.} Every pair\index{Subject}{double uniformization} of cardinal numbers has a least upper bound (in the usual cardinal ordering.) Hickman \cite{bib:1979b}.

\noindent\textbf{FORM 187.} Every pair of cardinal numbers has a greatest lower bound (in the usual cardinal ordering.) Hickman \cite{bib:1979b}.

\noindent\textbf{FORM 188.} $EP$ \emph{Ab}: For every Abelian group $A$ there is a projective Abelian group $G$ and a homomorphism from $G$ onto $A$. Blass\index{Author}{Blass} \cite{bib:1979} and note 60.

\noindent\textbf{FORM 189.} $EI$ \emph{Ab}: For every Abelian group $A$ there is an injective Abelian group $G$ and a one to one homomorphism from $A$ into $G$. Blass \cite{bib:1979} and note 60.

\noindent\textbf{FORM 190.} There is a non-trivial injective Abelian group. Blass \cite{bib:1979} and note 60.

\noindent\textbf{FORM 191.} $SVC$: There is a set $S$ such that for every set $a$, there is an ordinal $\alpha$ and a function from $S\times\alpha$ onto $a$. Blass \cite{bib:1979}.

\noindent\textbf{FORM 192.} $EP$ sets: For every set $A$ there is a projective set $X$ and a function from $X$ onto $A$. Blass \cite{bib:1979} and note 60.

[\textbf{192 A}] $(\forall x)(\exists y)$ ($x\leq^{\ast}y$ and $y$ indexed choice holds). Blass in Mathias\index{Author}{Mathias} \cite{bib:1977b}, notes 60 and 69.

\noindent \textbf{FORM 193.} $EFP$ \emph{Ab}: Every Abelian group is a homomorphic image of a free projective Abelian group. Blass \cite{bib:1979} and note 60.

\noindent\textbf{FORM 194.} $C(\varPi_{2}^{1})$ or $AC(\varPi_{2}^{1})$: If $P\in \omega \times^{\omega} \omega,\,P$ has domain $\omega$, and $P$ is in $\varPi_{2}^{1}$, then there is a sequence of elements $\langle x_{k}: k\in \omega\rangle$ of $^{\omega}\omega$ with $\langle k,\,x_{k}\rangle\in P$ for all $k\in\omega$. Kanovei\index{Author}{Kanovei} \cite{bib:1979} and note 61.

\noindent\textbf{FORM 195.} Every general linear system has a linear global reaction. Note 62 and Harazisville (see Kharazishvili\index{Author}{Kharazishvili} \cite{bib:1979}).

\noindent\textbf{FORM 196($\alpha$).} $\aleph_{\alpha}$ and $\aleph_{\alpha+1}$ are not both measurable. Bull\index{Author}{Bull} \cite{bib:1978} and note 20.

\noindent\textbf{FORM 197.} $\mathbb{R}^{2}$ is the union of three sets $C$ with the property that for all $x\in C$ there is a straight line $L$ such that $L\cap C=\{x\}$. Davies \cite{bib:1978}.

\noindent\textbf{FORM 198.} For every set $S$, if the only linearly orderable subsets of $S$ are the finite subsets of $S$, then either $S$ is finite or $S$ has an amorphous subset. Hickman \cite{bib:1978b} and note 57.

\noindent\textbf{FORM 199($n$).} (For $n\in \omega-\{0, 1\})$ If all $\varSigma_{n}^{1}$, Dedekind finite subsets of $^{\omega}\omega$ are finite, then all $\varPi_{n}^{1}$ Dedekind finite subsets of $^{\omega}\omega$ are finite. Kanovei \cite{bib:1978}, notes 57 and 61.

\noindent\textbf{FORM 200.} For all infinite $x,\,|2^{x}|=|x!|$. Dawson/Howard\index{Author}{Dawson/Howard} \cite{bib:1976} and note 64.

\noindent\textbf{FORM 201.} Linking Axiom for Boolean Algebras: Every Boolean algebra has a maximal linked system. ($L\subseteq B$ is \emph{linked} if $a\wedge b\neq 0$ for all $a$ and $b\in L$.) Schrijver\index{Author}{Schrijver} \cite{bib:1978}.

[\textbf{201 A}] Strong Linking Axiom for Boolean Algebras: Each linked system in a Boolean algebra is contained in some maximal linked system. (If $B$ is a Boolean algebra, $L\subseteq B$ is \emph{linked} if $a\wedge b\neq 0$ for all $a$ and $b\in L$.) Schrijver \cite{bib:1978}.

\noindent[\textbf{201 B}] Disjoint Transversal axiom (2,2): If $X$ is a set and $\mathcal{U}$ and $\mathcal{V}$ are subsets of $\{y\subseteq X:|y|\leq 2\}$ and if for each two finite subsets $\mathcal{U}_{0}\subseteq \mathcal{U}$ and $\mathcal{V}_{0}\subseteq \mathcal{V}$ there are transversals $Y_{1}$ and $Y_{2}$ for $\mathcal{U}_{0}$ and $\mathcal{V}_{0}$ respectively such that $Y_{1}\cap Y_{2}=\emptyset$, then there are disjoint transversals $X_{1}$ and $X_{2}$ for $\mathcal{U}$ and $\mathcal{V}$ respectively. $(Y\subseteq X$ is a \emph{transversal} for $\mathcal{U}$ if for all $y\in CalU,\, y\cap Y\neq\emptyset$.) Schrijver\index{Author}{Schrijver} \cite{bib:1978}.

\noindent\textbf{FORM 202.} $C(LO, \infty)$: Every linearly ordered family of non-empty sets has a choice function. Truss\index{Author}{Truss} \cite{bib:1978}.

[\textbf{202 A}] $H(TR\& C, P)$: If $(X, R)$ is a transitive, connected relation. Then $X$ has a $\subseteq$-maximal partially ordered subset. Harper/Rubin\index{Author}{Harper/Rubin} \cite{bib:1976} and note 39.

\noindent\textbf{FORM 203.} $C$(disjoint,$\subseteq \mathbb{R})$: Every partition of $\mathcal{P}(\omega)$ into non-empty subsets has a choice function. Truss \cite{bib:1978}.

\noindent\textbf{FORM 204.} For every infinite $X$, there is a function from $X$ onto $2X$. Truss \cite{bib:1978}.

\noindent\textbf{FORM 205.} For all cardinals $m$ and $n$, if $m\leq^{\ast}n$ and $\neg(n\leq^{\ast}m)$ then there is a cardinal $k<n$ such that $m\leq^{\ast}k$. Pelc\index{Author}{Pelc} \cite{bib:1978} and note 69.

\noindent\textbf{FORM 206.} The existence of a non-principal ultrafilter: There exists an infinite set $X$ and a non-principal ultrafilter on $X$. Blass\index{Author}{Blass} \cite{bib:1977b}.

\noindent\textbf{FORM 207($\alpha$).} $UT(\aleph_{\alpha}, \aleph_{\alpha}, <2^{\aleph_{\alpha}})$: The union of $\aleph_{\alpha}$ sets each of cardinality $\aleph_{\alpha}$ has cardinality less than $2^{\aleph_{\alpha}}$. Figura\index{Author}{Figura} \cite{bib:1977}.

\noindent\textbf{FORM 208.} For all ordinals $\alpha,\,\aleph_{\alpha+1}\leq 2^{\aleph_{\alpha}}$. Figura \cite{bib:1977}.

\noindent\textbf{FORM 209.} There is an ordinal $\alpha$ such that for all $X$, if $X$ is a denumerable union of denumerable sets then $\mathcal{P}(X)$ cannot be partitioned into $\aleph_{\alpha}$ non-empty sets. Morris\index{Author}{Morris} \cite{bib:1970}.

\noindent\textbf{FORM 210.} The commutator subgroup of a free group is free. L\"{a}uchli\index{Author}{L\"{a}uchli} \cite{bib:1962}.

\noindent\textbf{FORM 211.} $DCR$: Dependent choice for relations on $\mathbb{R}$: If $R\subseteq \mathbb{R} \times \mathbb{R}$ satisfies $(\forall x\in \mathbb{R})(\exists y\in \mathbb{R})(x\,R\,y)$ then there is a sequence $\langle x(n): n\in\omega\rangle$ of real numbers such that $(\forall n\in \omega)(x(n)\, R\, x(n+1))$. Mathias\index{Author}{Mathias}
\cite{bib:1977a}.

\noindent\textbf{FORM 212.} $C(2^{\aleph_{0}}, \subseteq \mathbb{R})$: If $R$ is a relation on $\mathbb{R}$ such that for all $x\in \mathbb{R}$, there is a $y\in \mathbb{R}$ such that $x\,R\,y$, then there is a function $f:\mathbb{R} \rightarrow \mathbb{R}$ such that for all $x\in \mathbb{R},\,x\,R\,f(x)$. Mathias \cite{bib:1977a}.

\noindent\textbf{FORM 213.} $C(\infty, \aleph_{1})$: If $(\forall y\in X)(|y|=\aleph_{1})$ then $X$ has a choice function. Pincus\index{Author}{Pincus} \cite{bib:1977a}.

\noindent\textbf{FORM 214.} $Z(\omega)$: For every family $A$ of infinite sets, there is a function $f$ such that for all $y\in A,\,f(y)$ is a non-empty subset of $y$ and $|f(y)|=\aleph_{0}$. Pincus \cite{bib:1977a}.

\noindent\textbf{FORM 215.} If $(\forall y\subseteq X)(y$ can be linearly ordered implies $y$ is finite), then $X$ is finite. Pincus \cite{bib:1972a}.

\noindent\textbf{FORM 216.} Every infinite tree has either an infinite chain or an infinite antichain. Pincus \cite{bib:1972a} and note 21.

\noindent\textbf{FORM 217.} Every infinite partially ordered set has either an infinite chain or an infinite antichain. Pincus \cite{bib:1972a}.

\noindent\textbf{FORM 218.} $(\forall n\in\omega -\{0\})MC(\infty, \infty$, relatively prime to $n$): $\forall n\in \omega-\{0\}$, if $X$ is a set of non-empty sets, then there is a function $f$ such that for all $x\in X,\,f(x)$ is a non-empty, finite subset of $x$ and $|f(x)|$ is relatively prime to $n$. Bleicher\index{Author}{Bleicher} \cite{bib:1965}.

[\textbf{218 A}]\index{Author}{Pineus/Solovay} Existence of Complementary Subspaces: For every field $F$, every vector space $V$ over $F$ has the property that if $S\subseteq V$ is a subspace of $V$, then there is a subspace $S^{\prime}\subseteq V$ such that $S\cap S^{\prime}=\{0\}$ and $S\cup S^{\prime}$ generates $V$. Bleicher
\cite{bib:1964}, \cite{bib:1965}, H. Rubin/J. Rubin\index{Author}{H. Rubin/J. Rubin} \cite{bib:1985}, pp 123ff, Theorems 6.36, 6.37, and 6.38.

\noindent\textbf{FORM 219.} $(\forall n\in \omega -\{0\})MC(\infty, WO$, relatively prime to $n)$: For all non-zero $n\in \omega$, if $X$ is a set of non-empty well orderable sets, then there is a function $f$ such that for all $x\in X,\, f(x)$ is a non-empty, finite subset of $x$, and $|f(x)|$ is relatively prime to $n$. Pincus\index{Author}{Pincus} \cite{bib:1972a}.

\noindent\textbf{FORM 220($p$).} Suppose $p\in\omega$ and $p$ is a prime. Any two elementary Abelian $p$-groups (all non-trivial elements have order $p$) of the same cardinality are isomorphic. Hickman\index{Author}{Hickman} \cite{bib:1977b}.

\noindent\textbf{FORM 221.} For all infinite $X$, there is a non-principal measure on $\mathcal{P}(X)$. Pincus/Solovay \cite{bib:1977} and note 147.

\noindent\textbf{FORM 222.} There is a non-principal measure on $\mathcal{P}(\omega)$. Pincus/Solovay \cite{bib:1977} and note 147.

\noindent\textbf{FORM 223.} There is an infinite set $X$ and a non-principal measure on $\mathcal{P}(X)$. Pincus/Solovay \cite{bib:1977} and note 147.

\noindent\textbf{FORM 224.} There is a partition of the real line into $\aleph_{1}$ Borel sets $\{B_{\alpha}:\alpha<\aleph_{1}\}$ such that for some $\beta<\aleph_{1},\,\forall\alpha<\aleph_{1},\, B_{\alpha}\in G_{\beta}$. $(G_{\beta}$ for $\beta<\aleph_{1}$ is defined by induction, $G_{0}=$\{$A: A$ is an open subset of $\mathbb{R}$\} and for $\beta>0$,

$G_{\beta}=\{\bigcap_{i=0}^{\infty}A_{i}:(\forall i\in w)(\exists\xi<\beta)(A_{i}\in G_{\xi})\}\ \mathrm{if}\ \beta\ \mathrm{is\ even\ and}$

\noindent $G_{\beta}=\{\bigcup_{i=0}^{\infty}A_{i}:(\forall i \in \omega)(\exists\xi<\beta)(A_{i}\in G_{\xi})\}\ \mathrm{if}\ \beta\ \mathrm{is\ odd}.)$

\noindent Stern\index{Author}{Stern} \cite{bib:1979}.

\noindent\textbf{FORM 225.} Every proper filter on $\omega$ can be extended to an ultrafilter. Rav\index{Author}{Rav} \cite{bib:1977}.

[\textbf{225 A}] There exists a compact, extremally disconnected topological space having a Dedekind infinite set of clopen subsets. Morillon\index{Author}{Morillon} \cite{bib:1993}, notes 94 and 114.

\noindent[\textbf{225 B}] There exists a compact, $\sigma$-extremally disconnected topological space having a Dedekind infinite set of clopen subsets. Morillon \cite{bib:1993}, notes 94 and 114.

\noindent\textbf{FORM 226.} Let $R$ be a commutative ring with identity, $B$ a proper subring containing 1 and $q$ a prime ideal in $B$. Then there is a subring $A$ of $R$ and a prime ideal $p$ in $A$ such that (a) $B\subseteq A$ (b) $q=B\cap p$ (c) $R-p$ is multiplicatively closed and (d) if $A\neq R$, then $R-A$ is multiplicatively closed. Rav \cite{bib:1977}.

\noindent\textbf{FORM 227.} For all groups $G$, if every finite subgroup of $G$ can be fully ordered then $G$ can be fully ordered. Rav \cite{bib:1977} and note 79.

\noindent\textbf{FORM 228.} Every torsion free Abelian group can be fully ordered. Rav \cite{bib:1977} and note 79.

\noindent\textbf{FORM 229.} If $(G, \circ, \leq)$ is a partially ordered group, then $\leq$ can be extended to a linear order on $G$ if and only if for every finite set $\{a_{1}, \ldots, a_{n}\}\subseteq G$, with $a_{i}\neq$ the identity for $i=1$ to $n$, the signs $\epsilon_{1},\,\ldots,\,\epsilon_{n}\
(\epsilon_{i}=\pm 1)$ can be chosen so that $P\cap S(a_{1}^{\epsilon_{1}}, \ldots, a_{n}^{\epsilon_{n}})=\emptyset$ (where $S(b_{1}, \ldots, b_{n})$ is the normal sub-semi-group of $G$ generated by $b_{1},\,\ldots,\,b_{n}$ and $P=\{g\in G:e\leq g\}$ where $e$ is the identity of $G$.) Rav\index{Author}{Rav} \cite{bib:1977}.

\noindent\textbf{FORM 230.} $L^{1}=HOD$.\index{Subject}{Artin-Schreier theorem} Szczepaniak\index{Author}{Szczepaniak} \cite{bib:1977} and note 82.

\noindent\textbf{FORM 231.} $UT(WO, WO, WO)$: The union of a well ordered collection of well orderable sets is well orderable. Tsukada\index{Author}{Tsukada} \cite{bib:1977}.

[\textbf{231 A}] Every almost well orderable set ($W$-set) is well orderable. Note 25.

\noindent\textbf{FORM 232.} Every metric space ($X, d$) has a $\sigma$-point finite base. Howard/ Keremedis/Rubin/Stanley\index{Author}{Howard/Keremedis/Rubin/Stanley} \cite{bib:1997} and note 141.

[\textbf{232 A}] Every metric space has a $\gamma$-point finite base, for some ordinal $\gamma$. Howard/Keremedis/Rubin/Stanley \cite{bib:1997} and note 141.

\noindent[\textbf{232 B}] Every metric space has a $\sigma$-locally finite base. Howard/Keremedis/ Rubin/Stanley \cite{bib:1997} and note 141.

\noindent[\textbf{232 C}] Every metric space has a $\gamma$-locally finite base, for some ordinal $\gamma$. Howard/Keremedis/Rubin/Stanley \cite{bib:1997} and note 141.

\noindent[\textbf{232 D}] Every metric space ($X, d$) has a $\sigma$-discrete base. Howard/Keremedis/Rubin/Stanley \cite{bib:1997} and note 141.

\noindent[\textbf{232 E}] Every metric space ($X, d$) has a $\gamma$-discrete base, for some ordinal $\gamma$. Howard/Keremedis/Rubin/Stanley \cite{bib:1997} and note 141.

\noindent[\textbf{232 F}] Every metric space ($X, d$) has a $\sigma$-disjoint base. Howard/Keremedis/ Rubin/Stanley \cite{bib:1997} and note 141.

\noindent [\textbf{232 G}] Every metric space ($X, d$) has a $\gamma$-disjoint base for some ordinal $\gamma$. Howard/Keremedis/Rubin/Stanley \cite{bib:1997} and note 141.

\noindent\textbf{FORM 233.} If a field has an algebraic closure it is unique up to isomorphism. Mathias\index{Author}{Mathias} \cite{bib:1977b}.

\noindent\textbf{FORM 234.} There is a non-Ramsey set: There is a set $A$ of infinite subsets of $\omega$ such that for every infinite subset $N$ of $\omega,\,N$ has a subset which is in $A$ and a subset which is not in $A$. Prikry\index{Author}{Prikry} \cite{bib:1976} and Mathias \cite{bib:1968}.

\noindent\textbf{FORM 235.} If $V$ is a vector space and $B_{1}$ and $B_{2}$ are bases for $V$ then $|B_{1}|$ and $|B_{2}|$ are comparable. Gross\index{Author}{Gross} \cite{bib:1976}.

\noindent\textbf{FORM 236.} If $V$ is a vector space with a basis and $S$ is a linearly independent subset of $V$ such that no proper extension of $S$ is a basis for $V$, then $S$ is a basis for $V$. Gross \cite{bib:1975}, \cite{bib:1976}.

[\textbf{236 A}] If $V$ is a vector space with a basis, then every linearly independent subset is contained in a basis.

\noindent\textbf{FORM 237.} The order of any group is divisible by the order of any of its subgroups, (i.e., if $H$ is a subgroup of $G$ then there is a set $A$ such that $|H\times A|=|G|.)$ Hickman\index{Author}{Hickman} \cite{bib:1976}.

\noindent\textbf{FORM 238.} Every elementary Abelian group (that is, for some prime $p$ every non identity element has order \emph{p}) is the direct sum of cyclic subgroups. Hickman \cite{bib:1976}.

\noindent\textbf{FORM 239.} For all $X$, if $X\neq\emptyset$, then there is a binary operation on $X$ that makes $X$ a group. Hickman\index{Author}{Hickman} \cite{bib:1976}.

\noindent\textbf{FORM 240.} If a group $G$ satisfies ``every ascending chain of subgroups is finite,'' then every subgroup of $G$ is finitely generated. Hickman \cite{bib:1976}.

\noindent\textbf{FORM 241.} Every algebraic closure of $\mathbb{Q}$ has a real closed subfield. Hodges\index{Author}{Hodges} \cite{bib:1976a}, L\"{a}uchli\index{Author}{L\"{a}uchli} \cite{bib:1962} and note 84.

\noindent\textbf{FORM 242.} There is, up to an isomorphism, at most one algebraic closure of $\mathbb{Q}$. Hodges \cite{bib:1976a}, L\"{a}uchli \cite{bib:1962} and note 84.

\noindent\textbf{FORM 243.} Every principal ideal domain is a unique factorization domain. Hodges \cite{bib:1976a} and note 84.

\noindent\textbf{FORM 244.} Every principal ideal domain has a maximal ideal. Hodges \cite{bib:1976a} and note 84.

\noindent\textbf{FORM 245.} There is a function $f:\omega_{1}\rightarrow\omega_{1}^{\omega}$ such that for every $\alpha,\, 0<\alpha<\omega_{1},\,f(\alpha)$ is a function from $\omega$ onto $\alpha$. Litman\index{Author}{Litman} \cite{bib:1976}.

[\textbf{245 A}] There is a function $g:\omega_{1}\rightarrow \omega_{1}^{\omega}$ such that for all limit ordinals $\alpha<\omega_{1},\,g(\alpha)$ is an increasing function from $\omega$ into $\alpha$ with the range of $g$ cofinal in $\alpha$. Litman \cite{bib:1976}

\noindent[\textbf{245 B}] Every closed set A in $\omega_{1}$ which contains only limit ordinals is a derivative (that is, $\exists B\subseteq\omega_{1}$ such that $A=\{x:x$ is a limit and $\{t:t\in B\}$ is cofinal in $x\}.)$ Litman \cite{bib:1976}.

\noindent\textbf{FORM 246.} The monadic theory theory $MT(\omega_{1}, <)$ of $\omega_{1}$ is recursive. Litman \cite{bib:1976} and note 85.

\noindent\textbf{FORM 247.} Every atomless Boolean algebra is Dedekind infinite. Plotkin\index{Author}{Plotkin} \cite{bib:1976}, notes 86 and 94.

\noindent\textbf{FORM 248.} For any $\kappa,\,\kappa$ is the cardinal number of an infinite complete Boolean algebra if and only if $\kappa^{\aleph_{0}}=\kappa$. Plotkin \cite{bib:1976} and note 86.

\noindent\textbf{FORM 249.} If $T$ is an infinite tree in which every element has exactly 2 immediate successors then $T$ has an infinite branch. Truss\index{Author}{Truss} \cite{bib:1975}, note 87 and note 21.

\noindent\textbf{FORM 250.} $(\forall n\in \omega-\{0,1\})(C(WO, n))$: For every natural number $n\geq 2$, every well ordered family of $n$ element sets has a choice function. Truss \cite{bib:1975}.

\noindent\textbf{FORM 251.} The additive groups $(\mathbb{R}, +)$ and $(\mathbb{C}, +)$ are isomorphic. Ash\index{Author}{Ash} \cite{bib:1975}.

\noindent\textbf{FORM 252.} The additive groups of $\mathbb{Q} \oplus \mathbb{R}$ and $\mathbb{R}$ are isomorphic. Ash \cite{bib:1975}.

\noindent\textbf{FORM 253.} {\L}o\'{s}' Theorem: If $M=\langle A,\,R_{j}\rangle_{j\in J}$ is a relational system, $X$ any set and $\mathcal{F}$ an ultrafilter in $\mathcal{P}(X)$, then $M$ and $M^{X}/\mathcal{F}$ are elementarily equivalent. Howard\index{Author}{Howard} \cite{bib:1975}.

\noindent\textbf{FORM 254.} $Z(D, R, l)$: Every directed relation $(P, R)$ in which ramified subsets have least upper bounds, has a maximal element. Harper/Rubin\index{Author}{Harper/Rubin} \cite{bib:1976} and note 39.

\noindent\textbf{FORM 255.} $Z(D, R)$: Every directed relation $(P, R)$ in which every ramified subset $A$ has an upper bound, has a maximal element. Harper/Rubin \cite{bib:1976} and note 39.

\noindent\textbf{FORM 256.} $Z(P, F)$: Every partially ordered set $(X, R)$ in which every forest $A$ has an upper bound, has a maximal element. Harper/Rubin\index{Author}{Harper/Rubin} \cite{bib:1976} and note 39.

\noindent\textbf{FORM 257.} $Z(TR, P)$:\index{Subject}{{\L}os' theorem}
 Every transitive relation $(X, R)$ in which every partially ordered subset has an upper bound, has a maximal element. Harper/Rubin \cite{bib:1976} and note 39.

\noindent\textbf{FORM 258.} $Z(D, L)$: Every directed relation $(X, R)$ in which linearly ordered subsets have upper bounds, has a maximal element. Harper/Rubin \cite{bib:1976} and note 39.

\noindent\textbf{FORM 259.} $Z(TR\& C, W)$: If ($X,\, R)$ is a transitive and connected relation in which every well ordered subset has an upper bound, then $(X, R)$ has a maximal element. Harper/Rubin \cite{bib:1976} and note 39.

\noindent\textbf{FORM 260.} $Z(TR\& C, P)$: If $(X, R)$ is a transitive and connected relation in which every partially ordered subset has an upper bound, then $(X, R)$ has a maximal element. Harper/Rubin \cite{bib:1976} and note 39.

\noindent\textbf{FORM 261.} $Z(TR, T)$: Every transitive relation $(X, R)$ in which every subset which is a tree has an upper bound, has a maximal element. Harper/Rubin \cite{bib:1976} and note 39.

\noindent\textbf{FORM 262.} $Z(TR, R)$: Every transitive relation $(X, R)$ in which every ramified subset $A$ has an upper bound, has a maximal element. Harper/Rubin \cite{bib:1976} and note 39.

\noindent\textbf{FORM 263.} $H(AS\& C, P)$: Every every relation $(X, R)$ which is antisymmetric and connected contains a $\subseteq$-maximal partially ordered subset. Harper/Rubin \cite{bib:1976} and note 39.

\noindent\textbf{FORM 264.} $H(C, P)$: Every connected relation $(X, R)$ contains a $\subseteq$-maximal partially ordered set. Harper/Rubin \cite{bib:1976} and note 39.

\noindent\textbf{FORM 265.} $H(A, TR)$: Every relation $(X, R)$ contains a $\subseteq$-maximal transitive subset. Harper/Rubin \cite{bib:1976} and note 39.

\noindent\textbf{FORM 266.} $H(AS, P)$: Every antisymmetric relation contains $\subseteq$-maximal partially ordered subset. Harper/Rubin \cite{bib:1976} and note 39.

\noindent\textbf{FORM 267.} There is no infinite, free complete Boolean algebra. Stavi\index{Author}{Stavi} \cite{bib:1975} and note 89.

\noindent\textbf{FORM 268}. If $\mathcal{L}$ is a lattice isomorphic to the lattice of subalgebras of some unary universal algebra (a unary universal algebra is one with only unary or nullary operations) and $\alpha$ is an automorphism of $\mathcal{L}$ of order 2 (that is, $\alpha^{2}$ is the identity) then there is a unary algebra $\mathfrak{A}$ and an isomorphism $\rho$ from $\mathcal{L}$ onto the lattice of subalgebras of $\mathfrak{A}^{2}$ with $\rho(\alpha(x))=(\rho(x))^{-1}(=\{(s, t):(t, s)\in\rho(x)\}\rangle$ for all $x\in \mathcal{L}$. Lampe\index{Author}{Lampe} \cite{bib:1974}.

\noindent\textbf{FORM 269}. For every cardinal $m$, there is a set $A$ such that $2^{|A|^{2}}\geq m$ and there is a choice function on the collection of 2-element subsets of $A$. Lampe \cite{bib:1974}.

\noindent\textbf{FORM 270}. $CT_{\mathrm{fin}}$: The compactness theorem for propositional logic restricted to sets of formulas in which each variable occurs only in a finite number of formulas. Kolany/Wojtylak\index{Author}{Kolany/Wojtylak} \cite{bib:1991}.

[\textbf{270 A}] $C(\infty, <\aleph_{0})\wedge CT_{2}\ (CT_{2}$ is the compactness theorem for propositional logic restricted to sets of formulas in which each variable occurs in at most two formulas. See form $271(n).)$ Kolany/Wojtylak\index{Author}{Kolany}\index{Author}{Kolany/Wojtylak} \cite{bib:1991}.

\noindent[\textbf{270 B}] $F_{\mathrm{fin}}$: For any family $\mathcal{F}$ of pairwise disjoint, non-empty, finite sets if $R$ is a symmetric binary relation on $\bigcup \mathcal{F}$ which satisfies (i) the set $\{b\in\bigcup \mathcal{F}:\neg a\,R\,b\}$ is finite for every $a\in\bigcup \mathcal{F}$ and (ii) every finite subfamily of $\mathcal{F}$ has an $R$-consistent choice function, then $\mathcal{F}$ has an $R$-consistent choice function. Kolany/Wojtylak \cite{bib:1991}.

\noindent [\textbf{270 C(n)}] $F_{n}^{\mathrm{fin}}$ ($n\geq 3$ a natural number): For any family $\mathcal{F}$ of pairwise disjoint, $n$ element sets if $R$ is a symmetric binary relation on $\bigcup \mathcal{F}$ which satisfies (i) the set $\{b\in\bigcup \mathcal{F}:\neg a\ R\ b\}$ is finite for every $a\in\bigcup \mathcal{F}$ and (ii) every finite subfamily of $\mathcal{F}$ has an $R$-consistent choice function, then $\mathcal{F}$ has an $R$-consistent choice function. Kolany \cite{bib:1992}.

\noindent\textbf{FORM 271 ($n$)}. If $n\in \omega-\{0,1\},\,CT_{n}$: The compactness theorem for propositional logic restricted to sets of formulas in which each variable occurs in at most $n$ formulas. Kolany/Wojtylak \cite{bib:1991}.

\noindent\textbf{FORM 272}. There is an $X\subseteq \mathbb{R}$ such that neither $X$ nor $\mathbb{R} -X$ has a perfect subset. Truss\index{Author}{Truss} \cite{bib:1974b} and note 28.

\noindent\textbf{FORM 273}. There is a subset of $\mathbb{R}$ which is not Borel. Truss \cite{bib:1974b}.

\noindent\textbf{FORM 274}. There is a cardinal number $x$ and an $n\in \omega$ such that $\neg (x\ \mathrm{adj}^{n}\,x^{2}$). (The expression ``$x\,\mathrm{adj}^{n}\, ya^{\prime\prime}$ means there are cardinals $z_{0},\,\ldots,\,z_{n}$ such that $z_{0}=x$ and $z_{n}=y$ and for all $i,\,0\leq i<n,\,z_{i}<z_{i+1}$ and if $z_{i}<z\leq z_{i+1}$, then $z=z_{i+1}$.) (Compare with [$0$ A]). Truss \cite{bib:1974b}.

\noindent\textbf{FORM 275}. The sequence of cardinals $\langle\aleph_{n}: n\in \omega\rangle$ has a unique minimal upper bound. Truss \cite{bib:1974b}.

\noindent\textbf{FORM 276}. $E(V^{\prime\prime}, III)$: For every set $A,\,\mathcal{P}(A)$ is Dedekind finite if and only if $A=\emptyset$ or $2|\mathcal{P}(A)|>|\mathcal{P}(A)|$. Howard/Spi\u{s}iak\index{Author}{Howard/Spi\u{s}iak} \cite{bib:1994} and note 94.

\noindent\textbf{FORM 277.} $E(D, VII)$: Every non-well-orderable cardinal is decomposable. Monro\index{Author}{Monro} \cite{bib:1973a} and note 94.

\noindent\textbf{FORM 278}. In an integral domain $R$, if every ideal is finitely generated then $R$ has a maximal proper ideal. Hodges\index{Author}{Hodges} \cite{bib:1974} and note 45 E.

\noindent\textbf{FORM 279}. The Closed Graph Theorem for operations between Fr\'{e}chet Spaces: Suppose $X$ and $Y$ are Fr\'{e}chet spaces, $T:X\rightarrow Y$ is linear and $G=\{(x, Tx): x\in X\}$ is closed in $X\times Y$. Then $T$ is continuous. Rudin\index{Author}{Rudin} \cite{bib:1991} p. 51, Wright\index{Author}{Wright} \cite{bib:1973}, note 96.

\noindent\textbf{FORM 280.} There is a complete separable metric space with a subset which does not have the Baire property. Wright \cite{bib:1973} and note 28.

\noindent\textbf{FORM 281.} There is a Hilbert space $H$ and an unbounded linear operator on $H$. Wright \cite{bib:1973}, note 96.

\noindent\textbf{FORM 282.} $\omega\not\rightarrow (\omega)^{\omega}$. Erd\"{o}s/Rado\index{Author}{Erd\"{o}s/Rado} \cite{bib:1952}, Kleinberg/Seiferas\index{Author}{Kleinberg/Seiferas} \cite{bib:1973} and note 97.

\noindent\textbf{FORM 283.} Cardinality of well ordered subsets: For all $n\in \omega$ and for all infinite $x,\,|x^{n}|<|s(x)|$ where $s(x)$ is the set of all well orderable subsets of $x$. Truss\index{Author}{Truss} \cite{bib:1973d}.

\noindent\textbf{FORM 284.} A system\index{Subject}{closed graph theorem} of linear equations over a field $F$ has a solution in $F$ if and only if every finite sub-system has a solution in $F$. Cowen/Emerson\index{Author}{Cowen/Emerson} \cite{bib:1995}.

\noindent\textbf{FORM 285.} Let $E$ be a set and $f:E\rightarrow E$, then $f$ has a fixed point if and only if $E$ is not the union of three mutually disjoint sets $E_{1},\,E_{2}$ and $E_{3}$ such that $E_{i}\cap f(E_{i})=\emptyset$ for $i=1,2,3$. Makowski/Wi\'{s}niewski\index{Author}{Makowski/Wi\'{s}niewski}\index{Author}{Wi\'{s}niewski} \cite{bib:1969}, Abian\index{Author}{Abian} \cite{bib:1968}, Baker\index{Author}{Baker} \cite{bib:1964}, and Kenyon\index{Author}{Kenyon} \cite{bib:1963}.

\noindent\textbf{FORM 286.} Extended Krein-Milman Theorem:\index{Subject}{Krein-Milman theorem} Let K be a quasicompact (sometimes called convex-compact), convex subset of a locally convex topological vector space, then K has an extreme point. H. Rubin/J. Rubin\index{Author}{H. Rubin/J. Rubin} \cite{bib:1985}, p. 177--178 and note 23.

\noindent\textbf{FORM 287.} The Hahn-Banach Theorem\index{Subject}{Hahn-Banach theorem} for Separable Normed Linear Spaces: Assume $V$ is a separable normed linear space and $p:V\rightarrow \mathbb{R}$ satisfies $p(x+y)\leq p(x)+p(y)$ and $(\forall t\geq 0)(\forall x\in V)(p(tx)=tp(x))$ and assume $f$ is a linear function from a subspace $S$ of $V$ into $\mathbb{R}$ which satisfies $(\forall x\in S)(f(x)\leq p(x))$, then $f$ can be extended to $f^{\ast}:V\rightarrow \mathbb{R}$ so that $f^{\ast}$ is linear and $(\forall x\in V)(f^{\ast}(x)\leq p(x))$. Note 145.

\noindent\textbf{FORM 288($n$).} If $n\in\omega -\{0,1\}, C(\aleph_{0}, n)$: Every denumerable set of $n$-element sets has a choice function. Wi\'{s}niewski \cite{bib:1972}.

\noindent\textbf{FORM 289.} If $S$ is a set of subsets of a countable set and $S$ is closed under chain unions, then $S$ has a $\subseteq$-maximal element. Smullyan\index{Author}{Smullyan} \cite{bib:1994} and note 102.

\noindent\textbf{FORM 290.} For all infinite $x,\,|2^{x}|=|x^{x}|$. Pincus\index{Author}{Pincus} \cite{bib:1978} and note 64.

\noindent\textbf{FORM 291.} For all infinite $x,\,|x!|=|x^{x}|$. Pincus \cite{bib:1978} and note 64.

\noindent\textbf{FORM 292.} $MC(LO, \infty)$: For each linearly ordered family of non-empty sets $X$, there is a function $f$ such that for all $x\in X\
f(x)$ is non-empty, finite subset of $x$). Howard/Rubin\index{Author}{Howard/Rubin} \cite{bib:1977}.

\noindent\textbf{FORM 293.} For all sets $x$ and $y$, if $x$ can be linearly ordered and there is a mapping of $x$ onto $y$, then $y$ can be linearly ordered.

\noindent\textbf{FORM 294.} Every linearly ordered $W$-set is well orderable. Morris\index{Author}{Morris} \cite{bib:1969} and note 25.

\noindent\textbf{FORM 295.} DO: Every infinite set has a dense linear ordering. Gonzalez\index{Author}{Gonzalez} \cite{bib:1995a}.

\textbf{[295 A]} DUO: Every infinite set has a dense unbounded linear ordering. (An ordering is \emph{unbounded} if there is neither a least element nor a greatest element.) Pincus \cite{bib:1997}.

\noindent\textbf{FORM 296.} Part-$\infty$: Every infinite set is the disjoint union of infinitely many infinite sets. Gonzalez \cite{bib:1995b}, Pincus \cite{bib:1997}.

\noindent\textbf{FORM 297.} Extremally disconnected compact Hausdorff spaces are projective in the category of all compact Hausdorff spaces. Banaschewski\index{Author}{Banaschewski} \cite{bib:1988}, Bell\index{Author}{Bell} \cite{bib:1988}, notes 60 and 114.

\noindent\textbf{FORM 298.} Every compact Hausdorff space has a Gleason cover. Bell\index{Author}{Bell} \cite{bib:1988}, Gleason\index{Author}{Gleason} \cite{bib:1958}, and note 114.

\noindent\textbf{FORM 299.} Any extremally disconnected compact Hausdorff space is projective in the category of Boolean topological spaces. Bell \cite{bib:1988}, notes 60 and 114.

\noindent\textbf{FORM 300.} Any continuous surjection between extremally disconnected compact Hausdorff spaces has an irreducible restriction to a closed subset of its domain. Bell \cite{bib:1988} and note 114.

\noindent\textbf{FORM 301.} Any continuous surjection between Boolean spaces has an irreducible restriction to a closed subset of its domain. Bell \cite{bib:1988} and note 114.

\noindent\textbf{FORM 302.} Any continuous surjection between compact Hausdorff spaces has an irreducible restriction to a closed subset of its domain. Bell \cite{bib:1988} and note 114.

\noindent\textbf{FORM 303.} If $B$ is a Boolean algebra, $S\subseteq B$ and $S$ is closed under $\wedge$, then there is a $\subseteq$-maximal proper ideal $I$ of $B$ such that $I \cap S=\{0\}$. H. Rubin/J. Rubin\index{Author}{H. Rubin/J. Rubin} \cite{bib:1985}.

\noindent\textbf{FORM 304.} There does not exist a $T_{2}$ topological space $X$ such that every infinite subset of $X$ contains an infinite compact subset. Brunner\index{Author}{Brunner} \cite{bib:1985c}.

\noindent\textbf{FORM 305.} There are $2^{\aleph_{0}}$ Vitali equivalence classes.\index{Subject}{Vitali equivalence classes} (\emph{Vitali equivalence classes} are equivalence classes of the real numbers under the relation $x\equiv y\leftrightarrow(\exists q\in \mathbb{Q})(x-y=q).)$. Kanovei\index{Author}{Kanovei} \cite{bib:1991}.

\noindent\textbf{FORM 306.} The set of Vitali equivalence classes is linearly orderable. (\emph{Vitali equivalence classes} are equivalence classes of the real numbers under the relation $x\equiv y\leftrightarrow(\exists q\in \mathbb{Q})(x-y=q).)$. Kanovei \cite{bib:1991}.

\noindent\textbf{FORM 307.} If $m$ is the cardinality of the set of Vitali equivalence classes, then $H(m)=H(2^{\aleph_{0}})$, where $H$ is Hartogs aleph function and the \emph{Vitali equivalence classes} are equivalence classes of the real numbers under the relation $x\equiv y\leftrightarrow(\exists q\in \mathbb{Q})(x-y=q)$. Kanovei \cite{bib:1991}.

\noindent \textbf{FORM 308($p$).} If $p$ is a prime and if $\{G_{y}:y\in Y\}$ is a set of finite groups, then the weak direct product $\prod_{y\in Y}G_{y}$ has a maximal $p$-subgroup. Howard/Yorke\index{Author}{Howard/Yorke} \cite{bib:1987} and note 24.

\textbf{[308 A($p$)]} If $p$ is a prime and if $Y$ is a set of non-empty, finite sets, then the weak direct product $\prod_{y\in Y}S_{y}$ has a maximal $p$-subgroup. $(S_{y}$ is the symmetric group on $y$.) Howard/Yorke \cite{bib:1987} and note 24.

\noindent \textbf{[308 B($p$)]} If $p$ is a prime and if $Y$ is a set of non-empty finite sets such that $(\forall y\in Y)(\gcd(|y|,p)=1)$, then $Y$ has a choice function. (This is similar to form $46(K)$, but in $46(K)$ the set $K$ is required to be finite.) Howard/Yorke \cite{bib:1987}.

\noindent \textbf{[308 C($p$)]} If $p$ is a prime and if $Y$ is a set of non-empty finite sets such that $(\forall y\in Y)(|y|\equiv 1\mod p)$, then $Y$ has a choice function. (This is similar to form $46(K)$, but in $46(K)$ the set $K$ is required to be finite.) Howard/Yorke \cite{bib:1987}.

\noindent\textbf{FORM 309.} The Banach-Tarski Paradox: There are three finite partitions $\{P_{1},\,\ldots,\,P_{n}\}$, $\{Q_{1},\ldots, Q_{r}\}$ and $\{S_{1}, \ldots, S_{n}, T_{1}, \ldots, T_{r}\}$ of $B^{3}=\{x\in \mathbb{R}^{3}:|x|\leq 1\}$ such that $P_{i}$ is congruent to $S_{i}$ for $1\leq i\leq n$ and $Q_{i}$ is congruent to $T_{i}$ for $1\leq i\leq r$. Banach/Tarski\index{Author}{Banach/Tarski} \cite{bib:1924}.

\noindent\textbf{FORM 310.} The Measure Extension Theorem: Suppose that $\mathcal{A}_{0}$ is a subring (that is, $a, b\in \mathcal{A}_{0}\rightarrow a\vee b\in \mathcal{A}_{0}$ and $a-b\in \mathcal{A}_{0})$ of a Boolean algebra $\mathcal{A}$ and $\mu$ is a measure on $\mathcal{A}_{0}$ (that is, $\mu : \mathcal{A}\rightarrow[0, \infty],\, \mu(a\vee b)=\mu(a)+\mu(b)$ for $a\wedge b=0$, and $\mu(0)=0.)$ then there is a measure on $\mathcal{A}$ that extends $\mu$. Wagon\index{Author}{Wagon} \cite{bib:1985}.

\noindent\textbf{FORM 311.} Abelian groups are amenable. $(G$ is \emph{amenable} if there is a finitely additive measure $\mu$ on $\mathcal{P}(G)$ such that $\mu(G)=1$ and $\forall A\subseteq G,\,\forall g\in G,\,\mu(gA)=\mu(A).)$
Wagon \cite{bib:1985}.

\noindent\textbf{FORM 312.} A subgroup of an amenable group is amenable. $(G$ is \emph{amenable} if there is a finitely additive measure $\mu$ on $\mathcal{P}(G)$ such that $\mu(G)=1$ and $\forall A\subseteq G,\,\forall g\in G,\,\mu(gA)=\mu(A).)$ Wagon \cite{bib:1985}.

\noindent\textbf{FORM 313.} $\mathbb{Z}$ (the set of integers under addition) is amenable. $(G$ is \emph{amenable} if there is a finitely additive measure $\mu$ on $\mathcal{P}(G)$ such that $\mu(G)=1$ and $\forall A\subseteq G,\,\forall g\in G,\,\mu(gA)=\mu(A).)$ Wagon \cite{bib:1985}.

\noindent\textbf{FORM 314.} Every permutation on a non-empty set can be expressed as a product of two reflections. (A reflection is a permutation $\phi$ such that $\phi^{2}$ is the identity.) Degen\index{Author}{Degen} \cite{bib:1988}.

\noindent\textbf{FORM 315.} $\Omega=\omega_{1}$, where $\Omega=\{\alpha\in$ On: $(\forall\beta\leq\alpha)(\beta=0\vee(\exists\gamma)(\beta=\gamma+1)\vee$
there is a sequence $\langle\gamma_{n}: n\in \omega\rangle$ such that for each $n\ \gamma_{n}<\beta$ and $\beta=\bigcup_{n<\omega}\gamma_{n}$\} Gitik\index{Author}{Gitik} \cite{bib:1980}.

\noindent\textbf{FORM 316.} If a linearly ordered set $(A, \leq)$ has the fixed point property then $(A, \leq)$ is complete. $((A, \leq)$ has the \emph{fixed point property} if every function $f:A\rightarrow A$ satisfying $(x\leq y\Rightarrow f(x)\leq f(y))$ has a fixed point, and $((A, \leq)$ is \emph{complete} if every subset of $A$ has a least upper bound.) H\"{o}ft/Howard\index{Author}{H\"{o}ft/Howard} \cite{bib:1994}.

\noindent\textbf{FORM 317.} Weak Sikorski Theorem: If $B$ is a complete, well orderable Boolean algebra and $f$ is a homomorphism of the Boolean algebra $A^{\prime}$ into $B$ where $A^{\prime}$ is a subalgebra of the Boolean algebra $A$, then $f$ can be extended to a homomorphism of $A$ into $B$. Howard\index{Author}{Howard} \cite{bib:1973}.

[\textbf{317 A}] If $B$ is a complete, well orderable Boolean algebra and $A$ is a Boolean algebra of which $B$ is a subalgebra, then there is an ideal $I$ of $A$ satisfying $I\cap B =\{0\}$ and $I$ is maximal among those ideals $J$ of $A$ satisfying $J\cap B =\{0\}$. Howard \cite{bib:1973}.

\noindent\textbf{FORM 318.} $\aleph_{1}$ is not measurable. Jech\index{Author}{Jech} \cite{bib:1968a} and note 20.

\noindent\textbf{FORM 319.} Measurable cardinals are inaccessible. Drake\index{Author}{Drake} \cite{bib:1974} and note 20.

\noindent\textbf{FORM 320.} No successor cardinal, $\aleph_{\alpha+1}$, is measurable. Drake \cite{bib:1974} and note 20.

\noindent\textbf{FORM 321.} There does not exist an ordinal $\alpha$ such that $\aleph_{\alpha}$ is weakly compact and $\aleph_{\alpha+1}$ is measurable. Bull\index{Author}{Bull} \cite{bib:1978} and note 20.

\noindent\textbf{FORM 322.}\index{Author}{Brunner} $KW(WO, \infty)$, The Kinna-Wagner Selection Principle for a well ordered family of sets: For every well ordered set $X$ there is a function $f$ such that for all $A\in M$, if $|A|>1$ then $\emptyset \neq f(A)\subsetneq A$. (See form 15). Kinna/Wagner \cite{bib:1955}, Brunner \cite{bib:1982a}.

\noindent\textbf{FORM 323.} $KW(\infty, WO)$, The Kinna-Wagner\index{Author}{Kinna/Wagner} Selection Principle for a family of well orderable sets: For every set $X$ of well orderable sets there is a function $f$ such that for all $A\in X$, if $|A|>1$ then $\emptyset \neq f(A)\subsetneq A$. (See form 15.) Howard/A. Rubin/J. Rubin\index{Author}{Howard/A. Rubin/J. Rubin} \cite{bib:1997}.

\noindent\textbf{FORM 324.}\index{Subject}{Sikorski's extension theorem} $KW(WO, WO)$, The Kinna-Wagner Selection Principle for a well ordered family of well orderable sets: For every well ordered set $X$ of well orderable sets, there is a function $f$ such that for all $A\in X$, if $|A|>1$ then $\emptyset \neq f(A)\subsetneq A$. (See form 15.) Howard/A. Rubin/J. Rubin \cite{bib:1997}.

\noindent\textbf{FORM 325.} Ramsey's\index{Author}{Ramsey} Theorem II: $\forall n,\,m\in \omega$, if A is an infinite set and the family of all $m$ element subsets of $A$ is partitioned into $n$ sets $S_{j},\,1\leq j\leq n$, then there is an infinite subset $B\subseteq A$ such that all $m$ element subsets of $B$ belong to the same $S_{j}$. (Also, see form 17.) Ramsey \cite{bib:1929}.

\noindent\textbf{FORM 326.} 2-SAT: Restricted Compactness Theorem for Propositional Logic III: If $\Sigma$ is a set of formulas in a propositional language such that every finite subset of $\Sigma$ is satisfiable and if every formula in $\Sigma$ is a disjunction of at most two literals, then $\Sigma$ is satisfiable. (A \emph{literal} is a propositional variable or its negation.) Wojtylak\index{Author}{Wojtylak}
\cite{bib:1995}.

\noindent\textbf{FORM 327.} $KW(WO, <\aleph_{0})$, The Kinna-Wagner Selection Principle for a well ordered family of finite sets: For every well ordered set $X$ of finite sets there is a function $f$ such that for all $A\in X$, if $|A|>1$ then $\emptyset \neq f(A)\subsetneq A$. (See form 15.) Howard/A. Rubin/J. Rubin \cite{bib:1997}.

\noindent\textbf{FORM 328.} $MC(WO, \infty)$: For every well ordered set $X$ such that for all $x\in X,\,|x|\geq 1$, there is a function $f$ such that and for every $x\in X,\,f(x)$ is a finite, non-empty subset of $x$. (See form 67.) Howard/A. Rubin/J. Rubin \cite{bib:1997}.

\noindent\textbf{FORM 329.} $MC(\infty, WO)$: For every set $M$ of well orderable sets such that for all $x\in X,\,|x|\geq 1$, there is a function $f$ such that for every $x\in X,\,f(x)$ is a finite, non-empty subset of $x$. (See form 67.) Howard/A. Rubin/J. Rubin \cite{bib:1997}.

\noindent\textbf{FORM 330.} $MC(WO, WO)$: For every well ordered set $X$ of well orderable sets such that for all $x\in X,\, |x|\geq 1$, there is a function $f$ such that for every $x\in X,\,f(x)$ is a finite, non-empty subset of $x$. (See form 67.) Howard/A. Rubin/J. Rubin \cite{bib:1997}.

\noindent\textbf{FORM 331.} If $(X_{i})_{i\in I}$ is a family of compact non-empty topological spaces then there is a family $(F_{i})_{i\in I}$ such that $\forall i\in I,\,F_{i}$ is an irreducible closed subset of $X_{i}$. Morillon\index{Author}{Morillon} \cite{bib:1988} and note 71.

\noindent\textbf{FORM 332.} A product of non-empty compact sober topological spaces is nonempty. Morillon \cite{bib:1988} and note 71.

\noindent\textbf{FORM 333.} $MC(\infty, \infty, \mathrm{odd})$: For every set $X$ of sets such that for all $x\in X,\,|x|\geq 1$, there is a function $f$ such that for every $x\in X,\,f(x)$ is a finite, non-empty subset of $x$ and $|f(x)|$ is odd. Keremedis\index{Author}{Keremedis} \cite{bib:1996a}.

[\textbf{333 A}] In every vector space $B$ over the two element field, every subspace of $B$ has a complementary subspace. Keremedis \cite{bib:1996a}.

\noindent[\textbf{333 B}]\index{Subject}{Ramsey's theorem} Every quotient group of an Abelian group each of whose non-unit elements has order 2 has a set of representatives. Keremedis \cite{bib:1996a}.

\noindent\textbf{FORM 334.} $MC(\infty, \infty, \mathrm{even})$: For every set $X$ of sets such that for all $x\in X,\,|x|\geq 2$, there is a function $f$ such that for every $x\in X,\,f(x)$ is a finite, non-empty subset of $x$ and $|f(x)|$ is even. Keremedis\index{Author}{Keremedis} \cite{bib:1996a}.

\noindent\textbf{FORM 335($n$).} Every quotient group of an Abelian group each of whose elements has order $\leq n$ has a set of representatives. Keremedis \cite{bib:1996c}.

\noindent\textbf{FORM 336($n$).} (For $n\in \omega,\,n\geq 2.)$ For every infinite set $X$, there is an infinite $Y\subseteq X$ such that the set of all $n$-element subsets of $Y$ has a choice function. Note 111.

\noindent\textbf{FORM 337.} Every separable metric space is Lindel\"{o}f. Good/Tree\index{Author}{Good/Tree} \cite{bib:1995}.

\noindent\textbf{FORM 338.} Every second countable metric space is Lindel\"{o}f. Good/Tree \cite{bib:1995}.

\noindent\textbf{FORM 339.} Every second countable metric space is separable. Good/Tree \cite{bib:1995}.

\noindent\textbf{FORM 340.} Every Lindel\"{o}f metric space is separable. Good/Tree \cite{bib:1995}.

\noindent\textbf{FORM 341.} Every Lindel\"{o}f metric space is second countable. Good/Tree \cite{bib:1995}.

\noindent\textbf{FORM 342($n$).} (For $n\in \omega, n\geq 2.)\ PC(\infty, n, \infty)$: Every infinite family of $n$-element sets has an infinite subfamily with a choice function. (See form 166.) Note 111.

\noindent\textbf{FORM 343.} A product of non-empty, compact $T_{2}$ topological spaces is non-empty. Morillon\index{Author}{Morillon} \cite{bib:1988}.

\noindent\textbf{FORM 344.} If $(E_{i})_{i\in I}$ is a family of non-empty sets, then there is a family $(U_{i})_{i\in I}$ such that $\forall i\in I,\,U_{i}$ is an ultrafilter on $E_{i}$. Morillon \cite{bib:1988}.

\noindent\textbf{FORM 345.} Rasiowa-Sikorski Axiom: If $(B, \wedge, \vee)$ is a Boolean algebra, $a$ is a non-zero element of $B$, and $\{X_{n}:n\in \omega\}$ is a denumerable set of subsets of $B$ then there is a maximal prime filter $F$ of $B$ such that $a\in F$ and for each $n\in \omega$, if $X_{n}\subseteq F$ and $\bigwedge X_{n}$ exists then $\bigwedge X_{n}\in F$. Morillon \cite{bib:1988}.

[\textbf{345 A}] Form 14 + form 43. Morillon \cite{bib:1988}.

\noindent[\textbf{345 B}] Form 14 + form 106. Morillon \cite{bib:1988}.

\noindent[\textbf{345 C}] Every bounded, distributive lattice has the strong Rasiowa-Sikorski property. Morillon \cite{bib:1988} and note 71.

\noindent[\textbf{345 D}] Every bounded, distributive lattice has the Rasiowa-Sikorski property. Morillon \cite{bib:1988} and note 71.

\noindent[\textbf{345 E}] Every Boolean algebra has the Rasiowa-Sikorski property. Morillon \cite{bib:1988} and note 71.

\noindent\textbf{FORM 346.} Vector Space Kinna-Wagner Principle: For every family $V=\{V_{i}:i\in K\}$ of non-trivial vector spaces there is a family $F=\{F_{i}:i\in K\}$ such that for each $i\in K,\,F_{i}$ is a non-empty, independent subset of $V_{i}$. Note 127.

\noindent\textbf{FORM 347.} Idemmultiple Partition Principle: If $y$ is idemmultiple $(2\times y\approx y)$ and $x\precsim^{\ast} y$, then $x\precsim y$. Higasikawa\index{Author}{Higasikawa} \cite{bib:1995} and note 69.

\noindent\textbf{FORM 348.} If $G$ is a group and $X$ and $Y$ both freely generate $G$ then $|X|=|Y|$. Federer/Jonsson \cite{bib:1950} and Howard \cite{bib:1985}.

\noindent\textbf{FORM 349.}\index{Author}{Sierpi\'{n}ski} $MC(\infty, \aleph_{0})$: For every set $X$ of non-empty denumerable sets there is a function $f$ such that for all $x\in X,\,f(x)$ is a finite, non-empty subset of $x$.

\noindent\textbf{FORM 350.}\index{Subject}{Rasiowa-Sikorski property}\index{Author}{Federer/Jonsson}\index{Author}{Howard} $MC(\aleph_{0}, \aleph_{0})$: For every denumerable set $X$ of non-empty denumerable sets there is a function $f$ such that for all $x\in X,\,f(x)$ is a finite, non-empty subset of $x$. Keremedis\index{Author}{Keremedis} \cite{bib:1996d} and note 132.

\noindent\textbf{FORM 351.} A countable product of metrizable spaces is metrizable. Keremedis \cite{bib:1996d} and note 132.

\noindent\textbf{FORM 352.} A countable product of second countable spaces is second countable. Keremedis \cite{bib:1996d} and note 132.

\noindent\textbf{FORM 353.} A countable product of first countable spaces is first countable. Keremedis \cite{bib:1996d} and note 132.

\noindent\textbf{FORM 354.} A countable product of separable $T_{2}$ spaces is separable. Keremedis \cite{bib:1996d} and note 132.

\noindent\textbf{FORM 355.} $KW(\aleph_{0}, \infty)$, The Kinna-Wagner Selection Principle for a denumerable family of sets: For every denumerable set $X$ there is a function $f$ such that for all $A\in X$, if $|A|>1$ then $\emptyset \neq f(A)\subsetneq A$.

\noindent\textbf{FORM 356.} $KW(\infty, \aleph_{0})$, The Kinna-Wagner Selection Principle for a family of denumerable sets: For every set $X$ of denumerable sets there is a function $f$ such that for all $A\in X$, if $|A|>1$ then $\emptyset \neq f(A)\subsetneq A$.

\noindent\textbf{FORM 357.} $KW(\aleph_{0}, \aleph_{0})$, The Kinna-Wagner Selection Principle for a denumerable family of denumerable sets: For every denumerable set $X$ of denumerable sets there is a function $f$ such that for all $A\in X$, if $|A|>1$ then $\emptyset \neq f(A)\subsetneq A$.

\noindent\textbf{FORM 358.} $KW(\aleph_{0}, <\aleph_{0})$, The Kinna-Wagner Selection Principle for a denumerable family of finite sets: For every denumerable set $X$ of finite sets there is a function $f$ such that for all $A\in X$, if $|A|>1$ then $\emptyset \neq f(A)\subsetneq A$.

\noindent\textbf{FORM 359.} If $\{A_{x}:x\in S\}$ and $\{B_{x}:x\in S\}$ are families of pairwise disjoint sets and $|A_{x}|\leq|B_{x}|$ for all $x\in S$, then $|\bigcup_{x\in S}A_{x}|\leq |\bigcup_{x\in S}B_{x}|$. Higasikawa\index{Author}{Higasikawa} \cite{bib:1995}.

\noindent\textbf{FORM 360.} Every second countable topological space is Lindel\"{o}f. G. Moore\index{Author}{G. Moore} \cite{bib:1982} p 236.

\noindent\textbf{FORM 361.} In $\mathbb{R}$, the union of a denumerable number of analytic sets is analytic. G. Moore \cite{bib:1982}, pp 181 and 325.

\noindent\textbf{FORM 362.} In $\mathbb{R}$, every Borel set is analytic. G. Moore \cite{bib:1982}, pp 181 and 325.

\noindent\textbf{FORM 363.} There are exactly $2^{\aleph_{0}}$ Borel sets in $\mathbb{R}$. G. Moore \cite{bib:1982}, p 325.

\noindent\textbf{FORM 364.} In $\mathbb{R}$, there is a measurable set that is not Borel. G. Moore \cite{bib:1982}, p 325.

\noindent\textbf{FORM 365.} For every uncountable set $A$, if $A$ has the same cardinality as each of its uncountable subsets then $|A|=\aleph_{1}$. Sierpi\'{n}ski \cite{bib:1918}.

\noindent\textbf{FORM 366.} There is a discontinuous function $f:\mathbb{R} \rightarrow \mathbb{R}$ such that for all real $x$ and $y,\,f(x+y)=f(x)+f(y)$. Hamel\index{Author}{Hamel} \cite{bib:1905}.

\noindent\textbf{FORM 367.} There is a Hamel basis for $\mathbb{R}$ as a vector space over $\mathbb{Q}$. Hamel \cite{bib:1905}.

\noindent\textbf{FORM 368.} The set of all denumerable subsets of $\mathbb{R}$ has power $2^{\aleph_{0}}$. Sierpi\'{n}ski \cite{bib:1917}.

\noindent\textbf{FORM 369.}\index{Author}{Sierpi\'{n}ski} If\index{Author}{Hamel} $\mathbb{R}$ is partitioned into two sets, at least one of them has cardinality $2^{\aleph_{0}}$. Luzin/Sierpi\'{n}ski\index{Author}{Luzin/Sierpi\'{n}ski} \cite{bib:1917}.

\noindent\textbf{FORM 370.} Weak Gelfand Extreme Point Theorem: If $A$ is a non-trivial Gelfand algebra then the closed unit ball in the dual of $A$ has an extreme point $e$. Morillon\index{Author}{Morillon} \cite{bib:1986}, notes 23 and 29.

[\textbf{370 A}] Measured Boolean Prime Ideal Theorem: If $\mathcal{B}$ is a Boolean algebra and there is a measure on $\mathcal{B}$ then there is an ultrafilter on $\mathcal{B}$. Morillon \cite{bib:1989}, \cite{bib:1990} and note 147.

\noindent\textbf{FORM 371.} There is an infinite, compact, Hausdorff, extremally disconnected topological space. Morillon \cite{bib:1993} and note 114.

\noindent\textbf{FORM 372.} Generalized Hahn-Banach Theorem: Assume that $X$ is a real vector space, $(Z,\preccurlyeq)$ is a Dedekind complete ordered vector space and $X_{0}$ is a subspace of $X$. If $\lambda_{0}:X_{0}\rightarrow Z$ is linear and $p:X\rightarrow Z$ is sublinear and if $\lambda_{0}\preccurlyeq p$ on $X_{0}$ then $\lambda_{0}$ can be extended to a linear map $\lambda:X\rightarrow Z$ such that $\lambda\preccurlyeq p$ on $X$. Schechter \cite{bib:1996b} and note 31.

[\textbf{372 A}] Existence of Banach Limits: For any directed set $(\Delta, \sqsubseteq)$ and any Dedekind complete ordered vector space $(Z,\preccurlyeq)$ there exists a $Z$-valued Banach limit. Schechter\index{Author}{Schechter} \cite{bib:1996b} and note 31.

\noindent[\textbf{372 B}] Generalized Convex Domination Theorem: For every real vector space $X$, every Dedekind complete ordered vector space $(Z, \preccurlyeq)$, every subspace $X_{0}\subseteq X$, every linear $\lambda_{0}:X_{0}\rightarrow Z$ and every convex $p:X\rightarrow Z$, if $\lambda_{0}\preccurlyeq p$ on $X_{0}$ then $\lambda_{0}$ can be extended to a linear map $\lambda:X\rightarrow Z$ such that $\lambda\preccurlyeq p$ on $X$. Schechter \cite{bib:1996b} and note 31.

\noindent[\textbf{372 C}] Generalized Support Theorem: For every real vector space $X$ and every Dedekind complete ordered vector space $(Z, \preccurlyeq)$, every convex function $p:X\rightarrow Z$ is a pointwise maximum of affine functions. Schechter \cite{bib:1996b} and note 31.

\noindent[\textbf{372 D}] Generalized Sandwich Theorem: For every real vector space $X$ and every Dedekind complete ordered vector space $(Z,\preccurlyeq)$, if $e:X\rightarrow Z$ is a concave function, $g:X\rightarrow Z$ is a convex function, and $e\leq g$ everywhere on $X$, then there exists an affine function $f:X\rightarrow Z$ satisfying $e\leq f\leq g$ everywhere on $X$. Schechter \cite{bib:1996b} and note 31.

\noindent\textbf{FORM 373($n$).} (For $n\in w, n\geq 2$.) $PC(\aleph_{0}, n, \infty)$: Every denumerable set of $n$-element sets has an infinite subset with a choice function.

\textbf{[373 A($n$)]} (For $n\in \omega, n\geq 2.)\ PC(WO, n, \infty)$: Every infinite well ordered set of $n$-element sets has an infinite subset with a choice function.

\noindent\textbf{[373 B($n$)]} $PUT(\aleph_{0}, n, \aleph_{0})$ for $n\in\omega -\{0,1\}$: The union of a denumerable set of pairwise disjoint $n$-element sets has a denumerable subset.

\noindent\textbf{[373 C($n$)]} $PUT(WO, n, WO)$ for $n\in w-\{0, 1\}$: The union of an infinite well ordered set of pairwise disjoint $n$-element sets has an infinite well ordered subset.

\noindent\textbf{FORM 374($n$).} $UT(\aleph_{0}, n, \aleph_{0})$ for $n\in \omega-\{0,1\}$: The union of a denumerable set of $n$-element sets is denumerable.

\textbf{[374 A($n$)]} For each $i,\,2\leq i\leq n\in\omega -\{0,1\}, C(\aleph_{0}, i)$: Every denumerable set of $i$-element sets has a choice function.

\noindent\textbf{FORM 375.}\index{Subject}{convex domination theorem}\index{Subject}{Gelfand extreme point theorem}\index{Subject}{Hahn-Banach theorem}\index{Subject}{measure extension theorem}\index{Subject}{sandwich theorem}\index{Subject}{support theorem} Tietze-Urysohn\index{Author}{Urysohn} Extension Theorem: If $(X, T)$ is a normal topological space, $A$ is closed in $X$, and $f:A\rightarrow[0,1]$ is continuous, then there exists a continuous function $g:X\rightarrow[0, 1]$ which extends $f$. Urysohn \cite{bib:1924}.

\noindent\textbf{FORM 376.} Restricted Kinna Wagner Principle: For every infinite set $X$ there is an infinite subset $Y$ of $X$ and a function $f$ such that for every $z\subseteq Y$, if $|z|\geq 2$ then $f(z)$ is a non-empty proper subset of $z$. De la Cruz/Di Prisco\index{Author}{De la Cruz/Di Prisco} \cite{bib:1996}.

\noindent\textbf{FORM 377.} Restricted Ordering Principle: For every infinite set $X$ there is an infinite subset $Y$ of $X$ such that $Y$ can be linearly ordered. De la Cruz/Di Prisco \cite{bib:1996}.

\noindent\textbf{FORM 378.} Restricted Choice for Families of Well Ordered Sets: For every infinite set $X$ there is an infinite subset $Y$ of $X$ such that the family of non-empty well orderable subsets of $Y$ has a choice function. De la Cruz/Di Prisco \cite{bib:1996}.

[\textbf{378A}] Restricted Choice for Families of Finite Sets: For every infinite set $X$ there is an infinite subset $Y$ of $X$ such that the family of non-empty Finite subsets of $Y$ has a choice function. De la Cruz/Di Prisco \cite{bib:1996}.

\noindent\textbf{FORM 379.} $PKW(\infty, \infty, \infty)$: For every infinite family $X$ of sets each of which has at least two elements, there is an infinite subfamily $Y$ of $X$ and a function $f$ such that for all $y\in Y,\, f(y)$ is a non-empty proper subset of $y$. De la Cruz/Di Prisco \cite{bib:1996}.

\noindent\textbf{FORM 380.} $PC(\infty, WO, \infty)$: For every infinite family of non-empty well orderable sets, there is an infinite subfamily $Y$ of $X$ which has a choice function. De la Cruz/Di Prisco \cite{bib:1996}.

\noindent\textbf{FORM 381.} DUM: The disjoint union of metrizable spaces is metrizable. Howard/Keremedis/Rubin/Rubin\index{Author}{Howard/Keremedis/Rubin/Rubin} \cite{bib:1997b}.

\noindent\textbf{FORM 382.} DUMN: The disjoint union of metrizable spaces is normal. Howard/Keremedis/Rubin/Rubin \cite{bib:1997b}.

\noindent\textbf{FORM 383.} MP: Every metric space is paracompact. Howard/Keremedis/Rubin/Stanley\index{Author}{Howard/Keremedis/Rubin/Stanley} \cite{bib:1997} and note 141.

[\textbf{383 A}] Every open cover $\mathcal{U}$ of a metric space ($X, d$) has a refinement $\mathcal{V}$ which covers $X$ and which can be written as a well ordered union $\bigcup\{V_{\alpha}:\alpha\in\gamma\}$ where each $V_{\alpha}$ is locally finite. Howard/Keremedis/Rubin/Stanley \cite{bib:1997} and note 141.

\noindent[\textbf{383 B}] MM: Every metric space is metacompact. Howard/Keremedis/Rubin/Stanley \cite{bib:1997} and note 141.

\noindent[\textbf{383 C}] Every open cover $\mathcal{U}$ of a metric space ($X, d$) has a refinement $\mathcal{V}$ which covers $X$ and which can be written as a well ordered union $\bigcup\{V_{\alpha}:\alpha\in\gamma\}$ where each $V_{\alpha}$ is point finite. Howard/Keremedis/Rubin/Stanley \cite{bib:1997} and note 141.

%%%%%%%%%%%part02

\part{Topical List of Forms \label{pt02:part02}}

\section*{Table of Contents}

\begin{enumerate}
\item[] ALGEBRAIC FORMS 

\begin{enumerate}
\item[I.] Groups 

\item[II.] Lattices 

\item[III.] Rings 

\begin{enumerate}
\item[A.] Boolean Algebras\index{Subject}{Tietze-Urysohn theorem} 

\item[B.] Fields (Including Algebraic Properties of $\mathbb{R}$ and $\mathbb{Q}$) 

\item[C.] General Rings 
\end{enumerate}

\item[IV.] Universal Algebra 

\item[V.] Vector Spaces 
\end{enumerate}

\item[] ANALYSIS (FORMS FROM ANALYSIS) 

\item[] CARDINAL NUMBER FORMS 

\begin{enumerate}
\item[I.] Alephs and their Properties 

\item[II.] Cardinal Arithmetic 

\item[III.] Cardinality of Unions 

\item[IV.] Definitions of Finite Including Forms Involving Dedekind Finite 

\item[V.] Ordering Cardinals 

\item[VI.] Inaccessible Cardinals 

\item[VII.] Miscellaneous Cardinal Number Forms 
\end{enumerate}

\item[] CHOICE FORMS 

\begin{enumerate}
\item[I.] Choosing Single Elements 

\item[II.] Dependent Choice 

\item[III.] Multiple Choice and Kinna-Wagner Principles 

\item[IV.] Partial Choice $PC$ 

\item[V.] Conditional Choice Including Variations of Rado's Lemma 

\item[VI.] Miscellaneous Choice 
\end{enumerate}

\item[] FIXED POINT THEOREMS 

\item[] GRAPH THEORETIC FORMS 

\item[] LOGIC (FORMS FROM LOGIC) 

\item[] MAXIMAL PRINCIPLES

\item[] MEASURES (FORMS INVOLVING MEASURES ON SETS) 

\item[] MISCELLANEOUS FORMS 

\item[] ORDERING PRINCIPLES INCLUDING PROPERTIES OF PARTIAL ORDERS 

\begin{enumerate}
\item[I.] Linear And Well Orderings 

\item[II.] Versions Of Martin's Axiom 

\item[III.] Trees

\item[IV.] Other Partial Ordering Properties
\end{enumerate}

\item[] TOPOLOGICAL FORMS (INCLUDING PROPERTIES OF $\mathbb{R}$)

\begin{enumerate}
\item[I.] Baire Category Type Theorems

\item[II.] Product Theorems

\item[III.] General Topology

\item[IV.] Linearly Orderable Spaces

\item[V.] Compact Spaces

\item[VI.] Properties Of $\mathbb{R}$

\item[VII.] Metric Spaces
\end{enumerate}
\end{enumerate}
\addtocontents{toc}{\SkipTocEntry}
\setcounter{chapter}{1}
\chapter{Topical List of Forms}\label{part02}

\section*{ALGEBRAIC FORMS}

\textbf{I. Groups}

[\textbf{1 BA}] (Depends on the prime $p$) Every group has a maximal $p$-subgroup. Howard/Yorke\index{Author}{Howard/Yorke} \cite{bib:1987d} and note 24.

\noindent [\textbf{1 BB}] Strong Nielsen-Schreier Theorem\index{Subject}{Nielsen-Schreier theorem} II: If $G$ is a group which is freely generated by $X$ and $U$ is a subgroup of $G$, then there is a subset $A$ of $G$ which freely generates $U$ and is level with respect to $X$. Howard\index{Author}{Howard} \cite{bib:1985n} and note 129.

\noindent [\textbf{1 BS}] For every Abelian group $(G,\,+)$, for every $H\subseteq G$, the infinite system of equations
\begin{equation*}
x_{i}+y_{i}=a_{i},\,i\in I,\,a_{i}\in G
\end{equation*}
has a solution in $H$ iff each of its equations has a solution in $H$. Keremedis\index{Author}{Keremedis} \cite{bib:1996c1}.

\noindent [\textbf{1 BX}] For every Abelian group $G$ and every subgroup $H$ of $G$ there exists a set of representatives for the quotient $G/H$. Keremedis \cite{bib:1996c1}.

\noindent [\textbf{1 CA}] Strong Nielsen-Schreier Theorem I: If $G$ is a group which is freely generated by $X$ and $U$ is a subgroup of $G$, then there is a subset $A$ of $G$ which freely generates $U$ and has the Nielsen property with respect to $X$. Howard \cite{bib:1985n} and note 129.

\noindent [\textbf{62 A}] (Depends on the two primes $p_{1}\neq p_{2}$) Form $308(p_{1})+$ form $308(p_{2})$. Howard/Yorke \cite{bib:1987d}.

\noindent [\textbf{62 B}] (Depends on the prime $p_{1}$) Form $45(p_{1})+$ form $308(p_{1})$. Howard/Yorke \cite{bib:1987d}.

\noindent [\textbf{62 C}] For every prime $p$, if $\{G_{y}:y\in Y\}$ is a set of finite groups, then the weak direct product $\prod_{y\in Y}G_{y}$ has a maximal $p$-subgroup. Howard/Yorke \cite{bib:1987d} and note 24.

\noindent [\textbf{62 D}] For every prime $p$, if $Y$ is a set of non-empty, finite sets, then the weak direct product $\prod_{y\in Y}S_{y}$ has a maximal $p$-subgroup. $(S_{y}$ is the symmetric group on $y$.) Howard/Yorke \cite{bib:1987d} and note 24.

\noindent [\textbf{67 D}] For every Abelian group $A$, every torsion free divisible subgroup of $A$ is a direct summand of $A$. Armbrust\index{Author}{Armbrust} \cite{bib:1972a} and note 24.

\noindent \textbf{68.} Nielsen-Schreier Theorem: Every subgroup of a free group is free. Jech\index{Author}{Jech} \cite{bib:1973b} p12.

\noindent \textbf{176.} Every infinite, locally finite group has an infinite Abelian subgroup. (\emph{Locally finite} means every finite subset generates a finite subgroup.) Plotkin\index{Author}{Plotkin} \cite{bib:1981i}.

\noindent \textbf{180.} Every Abelian group has a divisible hull. (If $A$ and $B$ are groups, $B$ \emph{is a divisible hull of} $A$ means $B$ is a divisible group, $A$ is a subgroup of $B$ and for every non-zero $b\in B,\,\exists n\in w$ such that $0\neq nb\in A$.) Fuchs\index{Author}{Fuchs} \cite{bib:1970e}, Theorem 24.4 p 107, Hodges\index{Author}{Hodges} \cite{bib:1980j}, and note 24.

\noindent \textbf{188.} $EP$ \emph{Ab}: For every Abelian group $A$ there is a projective Abelian group $G$ and a homomorphism from $G$ onto $A$. Blass \cite{bib:1979} and note 60.

\noindent \textbf{189.} $EI$ \emph{Ab}: For every Abelian group $A$ there is an injective Abelian group $G$ and a one to one homomorphism from $A$ into $G$. Blass\index{Author}{Blass} \cite{bib:1979} and note 60.

\noindent \textbf{190.} There is a non-trivial injective Abelian group. Blass \cite{bib:1979} and note 60.

\noindent \textbf{193.} $EFP$ \emph{Ab}: Every Abelian group is a homomorphic image of a free projective Abelian group. Blass \cite{bib:1979} and note 60.

\noindent \textbf{210.} The commutator subgroup of a free group is free. L\"{a}uchli\index{Author}{Lauchli@L\"{a}uchli} \cite{bib:1962d}.

\noindent \textbf{220($p$).} Suppose $p\in\omega$ and $p$ is a prime. Any two elementary Abelian $p$-groups (all non-trivial elements have order \emph{p}) of the same cardinality are isomorphic. Hickman\index{Author}{Hickman} \cite{bib:1977b2}.

\noindent \textbf{227.} For all groups $G$, if every finite subgroup of $G$ can be fully ordered then $G$ can be fully ordered. Rav \cite{bib:1977k} and note 79.

\noindent \textbf{228.} Every torsion free Abelian group can be fully ordered. Rav\index{Author}{Rav} \cite{bib:1977k} and note 79.

\noindent \textbf{229.} If $(G,\,\circ,\,\leq)$ is a partially ordered group, then $\leq$ can be extended to a linear order on $G$ if and only if for every finite set $\{a_{1},\,\ldots,\,a_{n}\}\subseteq G$, with $a_{i}\neq$ the identity for $i=1$ to $n$, the signs $\epsilon_{1},\,\ldots,\,\epsilon_{n}\,(\epsilon_{i}=\pm 1)$ can be chosen so that $P\cap S(a_{1}^{\epsilon_{1}},\,\ldots,\,a_{n}^{\epsilon_{n}})=\emptyset$ (where $S(b_{1},\, \ldots,\,b_{n})$ is the normal sub-semi-group of $G$ generated by $b_{1},\,\ldots,\,b_{n}$ and $P=\{g\in G:e\leq g\}$ where $e$ is the identity of $G$.) Rav \cite{bib:1977k}.

\noindent \textbf{237.} The order of any group is divisible by the order of any of its subgroups, (i.e., if $H$ is a subgroup of $G$ then there is a set $A$ such that $|H\times A|=|G|$.) Hickman \cite{bib:1976j}.

\noindent \textbf{238.} Every elementary Abelian group (that is, for some prime $p$ every non identity element has order \emph{p}) is the direct sum of cyclic subgroups. Hickman \cite{bib:1976j}.

\noindent \textbf{239.} For all $X$, if $X\neq\emptyset$, then there is a binary operation on $X$ that makes $X$ a group. Hickman \cite{bib:1976j}.

\noindent \textbf{240.} If a group $G$ satisfies ``every ascending chain of subgroups is finite,'' then every subgroup of $G$ is finitely generated. Hickman \cite{bib:1976j}.

\noindent \textbf{251.} The additive groups $(\mathbb{R},\,+)$ and $(\mathbb{C},\,+)$ are isomorphic. Ash\index{Author}{Ash} \cite{bib:1975}.

\noindent \textbf{252.} The additive groups of $\mathbb{Q} \oplus \mathbb{R}$ and $\mathbb{R}$ are isomorphic. Ash \cite{bib:1975}.

\noindent \textbf{308($p$).} If $p$ is a prime and if $\{G_{y}:y\in Y\}$ is a set of finite groups, then the weak direct product $\prod_{y\in Y}G_{y}$ has a maximal $p$-subgroup. Howard/Yorke \cite{bib:1987} and note 24.

\noindent \textbf{[308A($p$)]} If $p$ is a prime and if $Y$ is a set of non-empty, finite sets, then the weak direct product $\prod_{y\in Y}S_{y}$ has a maximal $p$-subgroup ($S_{y}$ is the symmetric group on $y$.) Howard/Yorke\index{Author}{Howard/Yorke} \cite{bib:1987} and note 24.

\noindent \textbf{311.} Abelian groups are amenable. ($G$ is \emph{amenable} if there is a finitely additive measure $\mu$ on $\mathcal{P}(G)$ such that $\mu(G)=1$ and $\forall A\subseteq G,\,\forall g\in G,\,\mu(gA)=\mu(A)$.) Wagon \cite{bib:1985u}.

\noindent \textbf{312.} A subgroup of an amenable group is amenable. $(G$ is \emph{amenable} if there is a finitely additive measure $\mu$ on $\mathcal{P}(G)$ such that $\mu(G)=1$ and $\forall A\subseteq G,\,\forall g\in G,\,\mu(gA)=\mu(A).)$ Wagon \cite{bib:1985u}.

\noindent \textbf{313.} $\mathbb{Z}$ (the set of integers under addition) is amenable. ($G$ is \emph{amenable} if there is a finitely additive measure $\mu$ on $\mathcal{P}(G)$ such that $\mu(G)=1$ and $\forall A\subseteq G,\,\forall g\in G,\,\mu(gA)=\mu(A)$.) Wagon \cite{bib:1985u}.

\noindent \textbf{314.} Every permutation on a non-empty set can be expressed as a product of two reflections. (A reflection is a permutation $\phi$ such that $\phi^{2}$ is the identity.) Degen\index{Author}{Degen} \cite{bib:1988c}.

\noindent [\textbf{333 B}] Every quotient group of an Abelian group each of whose non-unit elements has order 2 has a set of representatives. Keremedis\index{Author}{Keremedis} \cite{bib:1996a}.

\noindent \textbf{335 ($n$).} Every quotient group of an Abelian group each of whose elements has order $\leq n$ has a set of representatives. Keremedis \cite{bib:1996c1}.

\noindent \textbf{348.} If $G$ is a group and $X$ and $Y$ both freely generate $G$ then $|X|=|Y|$. Federer/Jonsson\index{Author}{Federer/Jonsson} \cite{bib:1950} and Howard\index{Author}{Howard} \cite{bib:1985n}.

\textbf{II. Lattices}

\noindent [\textbf{1 AN}] Every lattice $\langle L, \cap,\,\cup, \subseteq\rangle$ of sets (that is, $L$ is a family of sets and $\cap,\,\cup$ and $\subseteq$ are the usual operations) with a greatest element has a maximal, proper ideal. Bell/Fremlin\index{Author}{Bell/Fremlin} \cite{bib:1972c}.

\noindent [\textbf{14 E}] Almost Maximal Ideal Theorem: Every bounded, non-trivial, distributive lattice has an almost maximal ideal. (Definitions: If $D$ is a bounded non-trivial distributive lattice, $I(D)$ is the lattice of ideals of $D$. An \emph{almost maximal ideal} in $D$ is an ideal $P$ of $D$ such that $s(P)=P$ where $s$ is the operation on $I(D)$ defined as follows: If $e$ is the unit of $I(D)$, (i) For $a\in I(D)$ we say $x\in I(D)$ is $a$-\emph{small} if $x\vee b=e$ implies $a\vee b=e$ for all $b\in I(D)$. (ii) For any $a\in I(D)$, $s(a)$ is the join of all a-small elements of $I(D)$). Banaschewski/Harting\index{Author}{Banaschewski/Harting} \cite{bib:1985d} and Blass\index{Author}{Blass} \cite{bib:1986}.

\noindent [\textbf{14 F}] Every non-trivial complete distributive lattice with compact unit has a prime ideal. Banaschewski \cite{bib:1985c} and note 29.

\noindent [\textbf{14 R}] Every coherent local lattice is spatial. Banaschewski\index{Author}{Banaschewski} \cite{bib:1981b} and note 29.

\noindent [\textbf{14 S}] Every compact, regular local lattice is spatial. Banaschewski \cite{bib:1981b} and note 29.

\noindent [\textbf{14 AE}] For any distributive lattice $D$ with zero and unit, any ideal $J\subseteq D$ and $a\not\in J$, there is a prime filter $F$ such that $a\in F$ and $F\cap J=\emptyset$. Banaschewski\index{Author}{Banaschewski} \cite{bib:1981b}.

\noindent [\textbf{14 AU}] If $Z$ is a real vector lattice with order unit 1 then the set of all positive linear functionals $u$ on $Z$ with $u(1)=1$ is a convex set that has an extreme point. See H. Rubin/J. Rubin\index{Author}{H. Rubin/J. Rubin} \cite{bib:1985q}, pages 177 and 182 for definitions. Also see Gardiner\index{Author}{Gardiner} \cite{bib:1974i}.

\noindent [\textbf{14 BH}] The sum of finite frames is a spatial frame. Paseka\index{Author}{Paseka} \cite{bib:1989s} and note 29.

\noindent [\textbf{14 BI}] Every compact, conjunctive frame is spatial. Paseka \cite{bib:1989s} and note 29.

\noindent [\textbf{14 BJ}] The sum of compact, regular frames is spatial. Paseka \cite{bib:1989s} and note 29.

\noindent [\textbf{14 BL}] Every compact, 0 dimensional frame is spatial. Paseka \cite{bib:1989s} and note 29.

\noindent [\textbf{14 BP}] Every non-empty bounded distributive lattice has a prime ideal. Morillon\index{Author}{Morillon} \cite{bib:1988h}.

\noindent [\textbf{14 BQ}] Every proper ideal in a bounded distributive lattice is the intersection of the prime ideals containing it. Morillon \cite{bib:1988h} and note 79. (Compare with [14 AE]).

\noindent [\textbf{14 BR}] Every $\vee$-compact non trivial lattice has a principal prime ideal. Morillon \cite{bib:1988h}, \cite{bib:1985p} and note 71.

\noindent [\textbf{14 BS}] Every bounded distributive lattice has a prime ideal choice function. Morillon \cite{bib:1988h}, \cite{bib:1985p} and note~71.

\noindent [\textbf{14 BT}] If $(A_{i})_{i\in I}$ is a family of bounded, distributive, $\vee$-compact lattices and if $(J_{i})_{i\in I}$ is a family such that for all $i\in I,\,J_{i}$ is a proper ideal in $A_{i}$ then there is a family $(P_{i})_{i\in I}$ such that for all $i\in I,\,P_{i}$ is a prime principal ideal in $A_{i}$ containing $J_{i}$. Morillon \cite{bib:1985p}, \cite{bib:1988h} and note 71.

\noindent [\textbf{14 BU}] If $(A_{i})_{i\in I}$ is a family of distributive non trivial bounded lattices, then there is a family $(f_{i})_{i\in I}$ such that each $f_{i}$ is a prime ideal choice function on $A_{i}$. Morillon \cite{bib:1985p}, \cite{bib:1988h} and note 71.

\noindent [\textbf{14 CH}] For every bounded, distributive lattice $(A, \wedge, \vee)$ in which $0$ is $\wedge$-compact, for all $a\in A,\,a\neq 0$ and for all sequences $(X_{n})_{n\in\omega}$ of subsets of $A$, there is a prime filter in $A$ which contains $a$ and respects each $X_{n}$. Morillon \cite{bib:1988h} and note 71.

\noindent [\textbf{50 C}] Every DeMorgan compact regular frame is injective in the category of all such frames. Banaschewski \cite{bib:1988a}, notes 29 and 60.

\noindent [\textbf{50 D}] For any extension $M\supseteq L$ of a compact regular frame $L$, there is an $s\in M$ maximal with respect to the property that $x\vee s=1\rightarrow x=1$ for all $x\in L$. Banaschewski \cite{bib:1988a}, notes 29 and 60.

\noindent [\textbf{86 B($\alpha$)}] The frame envelope of any $\aleph_{\alpha}$ frame is $\aleph_{\alpha}$-Lindel\"{o}f. Banaschewski \cite{bib:1997} and note 29.

\noindent [\textbf{88 C}] If $m$ is a cardinal number and $\mathcal{L}$ is a lattice isomorphic to the lattice of subalgebras of some unary universal algebra (a unary universal algebra is one with only unary or nullary operations) and $\alpha$ is an automorphism of $\mathcal{L}$ of order 2 (that is, $\alpha^{2}$ is the identity) then there is a unary algebra $\mathfrak{A}$ with $|\ \mathrm{domain}\ \mathfrak{A}|\geq m$ and an isomorphism $\rho$ from $\mathcal{L}$ onto the lattice of subalgebras of $\mathfrak{A}^{2}$ with $\rho(\alpha(x))=(\rho(x))^{-1}(=\{(s,\,t):(t,\,s)\in \rho(x)\})$ for all $x\in \mathcal{L}$. Lampe\index{Author}{Lampe} \cite{bib:1974m}.

\noindent \textbf{268.} If $\mathcal{L}$ is a lattice isomorphic to the lattice of subalgebras of some unary universal algebra (a unary universal algebra is one with only unary or nullary operations) and $\alpha$ is an automorphism of $\mathcal{L}$ of order 2 (that is, $\alpha^{2}$ is the identity) then there is a unary algebra $\mathfrak{A}$ and an isomorphism $\rho$ from $\mathcal{L}$ onto the lattice of subalgebras of $\mathfrak{A}^{2}$ with $\rho(\alpha(x))=(\rho(x))^{-1}(=\{(s,\,t):(t,\,s)\in\rho(x)\})$ for all $x\in \mathcal{L}$. Lampe \cite{bib:1974m}.

\noindent [\textbf{345 C}] Every bounded, distributive lattice has the strong Rasiowa-Sikorski property\index{Subject}{Rasiowa-Sikorski property}. Morillon\index{Author}{Morillon} \cite{bib:1988h} and note 71.

\noindent [\textbf{345 D}] Every bounded, distributive lattice has the Rasiowa-Sikorski property. Morillon \cite{bib:1988h} and note 71.

\textbf{III. Rings}

\textbf{A. Boolean Algebras}

\noindent [\textbf{$\boldsymbol{0}$ Y}] There is no free complete Boolean algebra over $D$ if $|D|\geq\aleph_{0}$. Stavi\index{Author}{Stavi} \cite{bib:1975k}, Gaifman\index{Author}{Gaifman} \cite{bib:1970f}, Hales\index{Author}{Hales} \cite{bib:1964c}, and note 89.

\noindent [\textbf{1 BT}] For every Boolean ring $(B,\,+,\,\cdot)$ and every $H\subseteq B$ which is closed under $+$ there exists a $\subseteq$-maximal ideal $Q\subseteq B$ such that $H\cap Q=\{0\}$. Keremedis\index{Author}{Keremedis} \cite{bib:1996c1}.

\noindent \textbf{14.} BPI: Every Boolean algebra has a prime ideal. M. Stone \cite{bib:1936}.

\noindent [\textbf{14 B}] Stone Representation Theorem:\index{Subject}{Stone representation theorem} Every Boolean algebra is isomorphic to a set algebra. Jech\index{Author}{Jech} \cite{bib:1973b} p 27 prob 2.12 and M. Stone \cite{bib:1936}.

\noindent [\textbf{14 AV}] Every Boolean algebra has a 2-valued measure. Pincus\index{Author}{Pincus} \cite{bib:1972c1} and note 147.

\noindent [\textbf{14 BA}] Monteiro's Theorem\index{Subject}{Monteiro's theorem} for Complete Atomic Boolean Algebras: Let $K$ be a complete atomic Boolean algebra, $A$ a subalgebra of the Boolean algebra $B$ and $f$ a homomorphism from $A$ to $K$. Let $d$ be a semimorphism from $B$ into $K$ (that is, $d\,:\,B\rightarrow K,\,d(0)=0,\,d(1)=1$ and for all $x,\,y\in B,\, d(x\vee y)=d(x)\vee d(y)$ ) with $f(x)\leq d(x)$ for all $x\in A$. Then there is a homomorphism $h$ from $B$ to $K$ such that $h(x)\leq d(x)$ for all $x\in B$ and $h/A=f$. Monteiro\index{Author}{Monteiro} \cite{bib:1965g} and Bacsich\index{Author}{Bacsich} \cite{bib:1972b}.

\noindent [\textbf{14 BC}] If $m_{0}$ is an additive, real valued measure on a subalgebra $\mathcal{B}_{0}$ of a Boolean algebra $\mathcal{B}$ then there exists a measure $m$ on $\mathcal{B}$ such that $m=m_{0}$ on $\mathcal{B}_{0}$ and the range of $m$ is contained in the closure of the range of $m_{0}$. Tarski\index{Author}{Tarski} \cite{bib:1930} and note 147.

\noindent [\textbf{14 BK}] The sum of 4-element Boolean algebras is spatial. Paseka\index{Author}{Paseka} \cite{bib:1989s} and note 29.

\noindent [\textbf{43 A}] Tarski's Lemma:\index{Subject}{Tarski's lemma} Let $\{x(n):n\in \omega\}$ be a denumerable collection of subsets of a Boolean algebra $B$ each of which has a meet in B. Then for any $a\neq 0$ in $B$ there is a proper filter $F$ of $B$ that decides $\bigwedge x(n)$ for all $n$ and has $a\in F$. ($F$ \emph{decides} $\bigwedge x(n)$ if either $\bigwedge x(n)\in F$ or $t\not\in F$ for some $t\in x(n)$.) Goldblatt\index{Author}{Goldblatt} \cite{bib:1985m}.

\noindent \textbf{50.} Sikorski's Extension Theorem:\index{Subject}{Sikorski's extension theorem} Every homomorphism of a subalgebra $B$ of a Boolean algebra $A$ into a complete Boolean algebra $B^{\prime}$ can be extended to a homomorphism of $A$ into $B^{\prime}$. Sikorski\index{Author}{Sikorski} \cite{bib:1964p}, p 141.

\noindent [\textbf{50 A}] Every complete Boolean algebra is injective. Banaschewski\index{Author}{Banaschewski} \cite{bib:1988a}.

\noindent [\textbf{50 B}] Monteiro's Theorem:\index{Subject}{Monteiro's theorem} Let $K$ be a complete Boolean algebra, $A$ a subalgebra of the Boolean algebra $B$ and $f$ a homomorphism from $A$ to $K$. Let $d$ be a semimorphism from $B$ into $K$ (that is, $d:B\rightarrow K,\,d(0)=0,\,d(1)=1$ and for all $x$ and $y$ in $B,\,d(x\vee y)=d(x)\vee d(y))$ with $f(x)\leq d(x)$ for all $x\in A$. Then there is a homomorphism $h$ from $B$ to $K$ such that $h(x)\leq d(x)$ for all $x\in B$ and $h/A=f$. Monteiro\index{Author}{Monteiro} \cite{bib:1965g} and Bacsich\index{Author}{Bacsich} \cite{bib:1972b}.

\noindent [\textbf{50 E}] Form 14 + form 297. Banaschewski \cite{bib:1988a} and Bell\index{Author}{Bell} \cite{bib:1988b}.

\noindent [\textbf{50 F}] For any complete Boolean algebra $B$ and any subalgebra $A$ of $B$, there is a $U\in V^{(B)}$ such that $V^{(B)}\models\ \text{``} U$ is an ultrafilter on $A$ and $U^{B}\subseteq U$''. ($V^{(B)}$ is the Boolean valued universe constructed from $B,\,U^{B}=\{(\hat{x},\,x):x\in B\}$ is the canonical ultrafilter on $B$ and $x\mapsto\hat{x}$ is the canonical embedding of $V$ into $V^{(B)}$.) Bell \cite{bib:1977}  and \cite{bib:1983a}.

\noindent [\textbf{50 G}] For any complete Boolean algebra $B$ and any $C\in V^{(B)}$ such that $V^{(B)}\models \text{``}C$ is a Boolean algebra.'' there is a $U\in V^{(B)}$ such that $V^{(B)}\models \text{``}U$ is an ultrafilter in $C$.'' See [50 F] for definitions. Bell \cite{bib:1983a}.

\noindent [\textbf{50 H}] If $m:A\rightarrow B$ is a monomorphism between Boolean algebras, then there is an $m$-minimal epimorphism $p:B\rightarrow C$ to some Boolean algebra $C$. ($p$ is $m$-minimal if (1) $p\circ m$ is monic and (2) if $q:C\rightarrow D$ is any epimorphism to a Boolean algebra $D$ such that $q\circ p\circ m$ is monic, then $q$ is an isomorphism.) Bell \cite{bib:1988b}.

\noindent [\textbf{50 I}] If $A$ is a subalgebra of a Boolean algebra $B$, then there is a $\subseteq$-maximal proper filter $F$ in $B$ such that $F\cap A=\{1\}$. Bell \cite{bib:1988b}

\noindent [\textbf{50 J}] Form 299$+$ form 14. Bell \cite{bib:1988b}.

\noindent [\textbf{50 K}] Form 300$+$ form 14. Bell \cite{bib:1988b}.

\noindent [\textbf{50 L}] Form $301+$ form 14. Bell \cite{bib:1988b}.

\noindent [\textbf{50 M}] Form $302+$ form 14. Bell \cite{bib:1988b}.

\noindent [\textbf{52 A}] On every Boolean algebra there is an additive real valued measure. Luxemburg\index{Author}{Luxemburg} \cite{bib:1969d} and note 147.

\noindent [\textbf{52 B}] For every proper filter $\mathcal{F}$ of elements of a Boolean algebra $\mathcal{B}$, there exists an additive, real valued measure $m$ on $\mathcal{B}$ such that $m(x)=1$ for all $x\in \mathcal{F}$. Luxemburg \cite{bib:1969d} and note 147.

\noindent [\textbf{52 C}] For every non-empty set X and every proper ideal $\mathcal{I}$ of in $\mathcal{P}(X)$ there is a measure $m$ on $\mathcal{P}(X)$ such that $m(x)=0$ for all $x\in \mathcal{I}$. Luxemburg \cite{bib:1969d} and note 147.

\noindent [\textbf{52 D}] Let $\mathcal{B}_{0}$ be a subalgebra of a Boolean algebra $\mathcal{B}$ and $m_{0}$ an additive, real valued measure on $\mathcal{B}_{0}$. Then there is a real valued, additive measure $m$ on $\mathcal{B}$ such that $m=m_{0}$ on $\mathcal{B}_{0}$ and the range of $m$ is contained in the closed convex hull of the range of $m_{0}$. Luxemburg\index{Author}{Luxemburg} \cite{bib:1969d} and note 147.

\noindent [\textbf{123 A}] Form 70$+$ form 344. Morillon\index{Author}{Morillon} \cite{bib:1988h}.

\noindent \textbf{201.} Linking Axiom for Boolean Algebras:\index{Subject}{linking axiom} Every Boolean algebra has a maximal linked system. ($L\subseteq B$ is \emph{linked} if $a\wedge b\neq 0$ for all $a$ and $b\in L$.) Schrijver\index{Author}{Schrijver} \cite{bib:1978j}.

\noindent [\textbf{201 A}] Strong Linking Axiom for Boolean Algebras: Each linked system in a Boolean algebra is contained in some maximal linked system. (If $B$ is a Boolean algebra, $L\subseteq B$ is \emph{linked} if $a\wedge b\neq 0$ for all $a$ and $b\in L$.) Schrijver \cite{bib:1978j}.

\noindent \textbf{247.} Every atomless Boolean algebra is Dedekind infinite. Plotkin\index{Author}{Plotkin} \cite{bib:1976p}, notes 86 and 94.

\noindent \textbf{248.} For any $\kappa,\,\kappa$ is the cardinal number of an infinite complete Boolean algebra if and only if $\kappa^{\aleph_{0}}=\kappa$. Plotkin \cite{bib:1976p} and note 86.

\noindent \textbf{267.} There is no infinite, free complete Boolean algebra. Stavi\index{Author}{Stavi} \cite{bib:1975k} and note 89.

\noindent \textbf{303.} If $B$ is a Boolean algebra, $S\subseteq B$ and $S$ is closed under $\wedge$, then there is a $\subseteq$-maximal proper ideal $I$ of $B$ such that $I \cap S=\{0\}$. H. Rubin/J. Rubin\index{Author}{H. Rubin/J. Rubin} \cite{bib:1985q}.

\noindent \textbf{310.} The Measure Extension Theorem:\index{Subject}{measure extension theorem} Suppose that $\mathcal{A}_{0}$ is a subring (that is, $a,\,b\in \mathcal{A}_{0}\rightarrow a\vee b\in \mathcal{A}_{0}$ and $a-b\in \mathcal{A}_{0}$) of a Boolean algebra $\mathcal{A}$ and $\mu$ is a measure on $\mathcal{A}_{0}$ (that is, $\mu\,:\,\mathcal{A}\rightarrow[0,\,\infty],\,\mu(a\vee b)=\mu(a)+\mu(b)$ for $a\wedge b=0$, and $\mu(0)=0$.) then there is a measure on $\mathcal{A}$ that extends $\mu$. Wagon\index{Author}{Wagon} \cite{bib:1985u}.

\noindent \textbf{317.} Weak Sikorski Theorem:\index{Subject}{Sikorski's extension theorem} If $B$ is a complete, well orderable Boolean algebra and $f$ is a homomorphism of the Boolean algebra $A^{\prime}$ into $B$ where $A^{\prime}$ is a subalgebra of the Boolean algebra $A$, then $f$ can be extended to a homomorphism of $A$ into $B$. Howard \cite{bib:1973h}.

\noindent [\textbf{317 A}] If $B$ is a complete, well orderable Boolean algebra and $A$ is a Boolean algebra of which $B$ is a subalgebra, then there is an ideal $I$ of $A$ satisfying $I\cap B =\{0\}$ and $I$ is maximal among those ideals $J$ of $A$ satisfying $J\cap B=\{0\}$. Howard\index{Author}{Howard} \cite{bib:1973h}.

\noindent \textbf{344.} If $(E_{i})_{i\in I}$ is a family of non-empty sets, then there is a family $(U_{i})_{i\in I}$ such that $\forall i\in I,\,U_{i}$ is an ultrafilter on $E_{i}$. Morillon \cite{bib:1988h}.

\noindent \textbf{345.} Rasiowa-Sikorski Axiom:\index{Subject}{Rasiowa-Sikorski axiom} If $(B,\,\wedge,\,\vee)$ is a Boolean algebra, $a$ is a non-zero element of $B$, and $\{X_{n}:n\in w\}$ is a denumerable set of subsets of $B$ then there is a maximal prime filter $F$ of $B$ such that $a\in F$ and for each $n\in \omega$, if $X_{n}\subseteq F$ and $\bigwedge X_{n}$ exists then $\bigwedge X_{n}\in F$. Morillon~\cite{bib:1988h}.

\noindent [\textbf{345 A}] Form $14+$ form 43. Morillon \cite{bib:1988h}.

\noindent [\textbf{345 B}] Form $14+$ form 106. Morillon \cite{bib:1988h}.

\noindent [\textbf{345 E}] Every Boolean algebra has the Rasiowa-Sikorski property. Morillon \cite{bib:1988h} and note 71.

\noindent [\textbf{370 A}] Measured Boolean Prime Ideal Theorem: If $\mathcal{B}$ is a Boolean algebra and there is a measure on $\mathcal{B}$ then there is an ultrafilter on $\mathcal{B}$. Morillon\index{Author}{Morillon} \cite{bib:1989r}, \cite{bib:1990e} and note 147.

\textbf{B. Fields (Including Algebraic Properties of} $\mathbb{R}$ \textbf{and} $\mathbb{Q}$) (See note 45.)

\noindent [\textbf{$\boldsymbol{0}$ R}] Every infinite Dedekind finite field is the direct limit of a strictly increasing $\omega$-sequence of finite subfields. Hickman\index{Author}{Hickman} \cite{bib:1982f}, notes 45 and 94.

\noindent [\textbf{14 BN}] A system of equations over a finite field $F$ has a solution in $F$ if and only if every finite sub-system has a solution in $F$. Cowen/Emerson\index{Author}{Cowen/Emerson} \cite{bib:1995} and note 30.

\noindent \textbf{69.} Every field has an algebraic closure. Jech\index{Author}{Jech} \cite{bib:1973b} p~13.

\noindent \textbf{72.} Artin-Schreier Theorem: Every field in which $-1$ is not the sum of squares can be ordered. (The ordering, $\leq$, must satisfy (a) $a\leq b\rightarrow a+c\leq b+c$ for all $c$ and (b) $c\geq 0$ and $a\leq b\rightarrow a\cdot c\leq b\cdot\mathrm{c}$.) Tarski\index{Author}{Tarski} [\textbf{1954b}].

\noindent \textbf{233.} If a field has an algebraic closure it is unique up to isomorphism. Mathias\index{Author}{Mathias} [\textbf{1977b}].

\noindent \textbf{241.} Every algebraic closure of $\mathbb{Q}$ has a real closed subfield. Hodges [\textbf{1976a}], L\"{a}uchli\index{Author}{Lauchli@L\"{a}uchli} [\textbf{1962}] and note 84.

\noindent \textbf{242.} There is, up to an isomorphism, at most one algebraic closure of $\mathbb{Q}$. Hodges\index{Author}{Hodges} [\textbf{1976a}], L\"{a}uchli [\textbf{1962}] and note 84.

\noindent \textbf{284.} A system of linear equations over a field $F$ has a solution in $F$ if and only if every finite sub-system has a solution in $F$. Cowen/Emerson \cite{bib:1995}.

\noindent \textbf{367.} There is a Hamel basis for $\mathbb{R}$ as a vector space over $\mathbb{Q}$. Hamel\index{Author}{Hamel} [\textbf{1905}].

\textbf{C. General Rings}

\noindent [\textbf{1 B}] Every commutative ring with unit has a maximal ideal. Hodges [\textbf{1979}].

\noindent [\textbf{1 CD}] In a commutative ring, every proper ideal can be extended to a maximal ideal. Krull\index{Author}{Krull} [\textbf{1929}]. ([1 CD] implies [1 B]).

\noindent [\textbf{14 C}] Every commutative ring with unit has a prime ideal. Scott\index{Author}{Scott} [\textbf{1954}].

\noindent [\textbf{14 D}] Every ring with unit has a (two sided) prime ideal. (An ideal $P$ in a ring $R$ is prime if $P\neq R$ and for any ideals $A$ and $B$ in $R$, if $AB\subseteq P$ then $A\subseteq P$ or $B\subseteq P$.) Blass\index{Author}{Blass} [\textbf{1986}] and Banaschewski\index{Author}{Banaschewski} [\textbf{1985}].

\noindent [\textbf{14 T}] Every regular commutative ring with unit has a maximal ideal. Banaschewski [\textbf{1983}] and note 29.

\noindent [\textbf{14 U}] Every closed proper ideal in a Gelfand algebra is contained in some maximal ideal. Banaschewski\index{Author}{Banaschewski} [\textbf{1983}] and note 29.

\noindent [\textbf{14 AL}] The radical of a proper ideal $A$ in a commutative ring $R$ with identity is the intersection of all prime ideals in $R$ containing $A$. (The \emph{radical} of $A= \{x\in R:\ \text{for some}\ m\in\omega,\, x^{m}\in A\}.)$ Rav\index{Author}{Rav} [\textbf{1977}] and note 80.

\noindent [\textbf{14 AM}] Let $A$ be a subring of a commutative ring $R$ and $p$ a prime ideal in $A$ such that $p=Rp\cap A$. Then there is a prime ideal $J$ in $R$ such that $p=J\cap A$. Rav [\textbf{1977}] and note 80.

\noindent [\textbf{14 AN}] In every commutative ring with identity, every proper ideal is included in some prime ideal. Rav\index{Author}{Rav} [\textbf{1977}].

\noindent [\textbf{14 AO}] Suppose $R$ is a commutative ring, $A$ is a proper ideal in $R$, and $S$ is a multiplicative semigroup in $R$ not meeting $A$. Then there is a prime ideal $p$ in $R$ such that $A\subseteq p$ and $p\cap S=\emptyset$. Rav [\textbf{1977}] and note 80.

\noindent [\textbf{14 BM}] A system of equations over a finite ring $R$ has a solution in $R$ if and only if every finite sub-system has a solution in $R$. Abian\index{Author}{Abian} [\textbf{1972}] and note 30.

\noindent [\textbf{14 BO}] A system of equations over a finite structure $(A,\,+,\,\cdot)$ has a solution in $A$, if and only if every finite sub-system has a solution in $A$. note 30.

\noindent \textbf{226.} Let $R$ be a commutative ring with identity, $B$ a proper subring containing 1 and $q$ a prime ideal in $B$. Then there is a subring $A$ of $R$ and a prime ideal $p$ in $A$ such that (a) $B\subseteq A$ (b) $q=B\cap p$ (c) $R-p$ is multiplicatively closed and (d) if $A\neq R$, then $R-A$ is multiplicatively closed. Rav [\textbf{1977}].

\noindent \textbf{243.} Every principal ideal domain is a unique factorization domain. Hodges\index{Author}{Hodges} [\textbf{1976a}] and note 84.

\noindent \textbf{244.} Every principal ideal domain has a maximal ideal. Hodges [\textbf{1976a}] and note 84.

\noindent \textbf{278.} In an integral domain $R$, if every ideal is finitely generated then $R$ has a maximal proper ideal. Hodges [\textbf{1974}] and note 45 E.

\textbf{IV. Universal Algebra}

\noindent [\textbf{$\boldsymbol{0}$ U}] For any class of similar (universal) algebras $K,\,HSP\,K=Mod\,Eq\,K$. (The class of homomorphic images of subalgebras of products of non-empty families of algebras in $K$ is the same as the class of algebras which satisfy the set of equations holding in every algebra in $K$.) Andreka/Nemeti\index{Author}{Andreka/Nemeti} \cite{bib:1981b} and note 50.

\noindent [\textbf{$\boldsymbol{0}$ V}] For every class $K$ of similar algebras $SP\,K\supseteq PS\,K$. (For any class $K$ of similar algebras, $S\,K$ is the class of all subalgebras of elements of $K$ (not closed under isomorphism) and $P\,K$ is the class of algebras isomorphic to direct products of algebras in $K$). Andreka/Nemeti~[\textbf{1980}].

\noindent [\textbf{$\boldsymbol{0}$ W}] For every class $K$ of similar algebras $SP^{r}K\supseteq P^{r}S\,K$. (For any class $K$ of similar algebras, $S\,K$ is the class of all subalgebras of elements of $K$ (not closed under isomorphism) and $P^{r}K$ is the class of all algebras isomorphic to reduced direct products of algebras in $K$.) Andreka/Nemeti [\textbf{1980}].

\noindent [\textbf{1 AD}] For every class $K$ of similar algebras $SP^{g}K\supseteq P^{g}S\,K$. (For any class $K$ of similar algebras, $S\,K$ is the class of all subalgebras of elements of $K$ (not closed under isomorphism) and $P^{g}K$ is the class of direct products of elements of $K$ not closed under isomorphism). Andreka/Nemeti [\textbf{1980}].

\noindent [\textbf{1 AP}] For any type $\tau$ of the form $\tau =(M)$, every homomorphism of an algebra of type $\tau$ is closed. Diener\index{Author}{Diener} [\textbf{1989}] and note 115.

\noindent [\textbf{1 AQ}] For any type $\tau$ of the form $\tau =(M)$, every homomorphism of an algebra $\mathbf{A}$ of type $\tau$ into an algebra $\mathbf{B}$ of type $\tau$ maps $\mathbf{A}$ onto a subalgebra of $\mathbf{B}$. Diener [\textbf{1989}] and note 105.

\noindent [\textbf{1 AR}] Let $S$ be any family of disjoint non-empty sets and let $\tau =(\bigcup S)$, then any homomorphism from $\mathfrak{P}(\tau;\bigcup\,S)$ into $\mathfrak{P}(\tau;S)$ maps $\mathfrak{P}(\tau;\bigcup\,S)$ onto a subalgebra of $:\mathfrak{P}(\tau;S)$. Diener\index{Author}{Diener} [\textbf{1989}] and note 115.

\noindent [\textbf{1 AX}] Birkhoff Representation Theorem:\index{Subject}{Birkhoff representation theorem} Every algebra is isomorphic to a subdirect product of subdirectly irreducible algebras. Birkhoff\index{Author}{Birkhoff} [\textbf{1944}], Gr\"{a}tzer [\textbf{1979}] and [\textbf{1986}].

\noindent [\textbf{1 BE}] Every unique factorization domain can be mapped homomorphically onto a non-zero subdirectly irreducible ring. Rav\index{Author}{Rav} [\textbf{1988}], Gr\"{a}tzer\index{Author}{Gratzer@Gr\"{a}tzer} [\textbf{1979}] pp 122-124 and note 84.

\noindent [\textbf{14 AS}] Let $\mathfrak{A}$ and $\mathfrak{B}$ be (universal) algebras of type $r$ where $\mathfrak{A}$ and $\mathfrak{B}$ have underlying sets $A$ and $B$ respectively and $B$ is finite. Suppose for every finite $C\subseteq A$ there is a partial $r$ homomorphism from $\mathfrak{A}$ into $\mathfrak{B}$ with domain $C$. Then there is an $r$ homomorphism from $\mathfrak{A}$ into $\mathfrak{B}$ . van Benthem\index{Author}{Benthem} [\textbf{1975}].

\noindent [\textbf{43 N}] $RM1,\,\aleph_{0}$: The representation theorem for multi-algebras with $\aleph_{0}$ unary operations: Assume $(A,\,F)$ is a multi-algebra with $\aleph_{0}$ unary operations (and no other operations), then there is an algebra $(B,\,G)$ with $\aleph_{0}$ unary operations and an equivalence relation $E$ on $B$ such that $(B/E,\,G/E)$ and $(A,\,F)$ are isomorphic multi-algebras. Note 50 and H\"{o}ft/Howard\index{Author}{Hoft/Howard@H\"{o}ft/Howard} \cite{bib:1981b}.

\noindent [\textbf{43 O}] (Depends on $n\in\omega,\,n\geq 1)\,RM1,\,n$: The representation theorem for multi-algebras with $n$ unary operations: Assume $(A,\,F)$ is a multi-algebra with $n$ unary operations (and no other operations), then there is an algebra $(B,\,G)$ with $n$ unary operations and an equivalence relation $E$ on $B$ such that $(B/E,\,G/E)$ and $(A,\,F)$ are isomorphic multi-algebras. H\"{o}ft/Howard\index{Author}{Hoft/Howard@H\"{o}ft/Howard} \cite{bib:1981b} and note 50.

\noindent [\textbf{67 E}] For any algebra $\mathbf{A}$ with finitary operations (that is, $\mathbf{A}$ is of type $\tau$ where $\tau=(K_{i})_{i\in I}$ with $K_{i}$ finite for each $i\in I$) and for any generating subset $Z$ of $\mathbf{A}$, there is a function $h:A\rightarrow \mathcal{P}_{\omega}(Z)$ ($A$ is the domain of $\mathbf{A}$ and $\mathcal{P}_{\omega}(Z)$ is the set of all finite subsets of $Z$.) such that $\forall a\in A,\,a$ is in the subalgebra of $\mathbf{A}$ generated by $h(a)$. Diener [\textbf{1989}] and note 115.

\noindent [\textbf{67 F}] For any algebra $\mathbf{A}$ with one unary operation $f$ and any generating set $Z$ of $\mathbf{A}$, there is a function $h:A\rightarrow \mathcal{P}_{\omega}(Z)$ ($A$ is the domain of $\mathbf{A}$ and $\mathcal{P}_{\omega}(Z)$ is the set of all finite subsets of $Z$.) such that $\forall a\in A,\,a$ is in the subalgebra of $\mathbf{A}$ generated by $h(a)$. Diener [\textbf{1989}] and note 115.

\noindent [\textbf{88 C}] If $m$ is a cardinal number and $\mathcal{L}$ is a lattice isomorphic to the lattice of subalgebras of some unary universal algebra (a unary universal algebra is one with only unary or nullary operations) and $\alpha$ is an automorphism of $\mathcal{L}$ of order 2 (that is, $\alpha^{2}$ is the identity) then there is a unary algebra $\mathfrak{A}$ with $|\,\mathrm{domain}\,(\mathfrak{A})|\geq m$ and an isomorphism $\rho$ from $\mathcal{L}$ onto the lattice of subalgebras of $\mathfrak{A}^{2}$ with $\rho(\alpha(x))=(\rho(x))^{-1}(=\{(s,\,t):(t,\, s)\in\rho(x)\})$ for all $x\in \mathcal{L}$. Lampe\index{Author}{Lampe} \cite{bib:1974m}.

\noindent [\textbf{88 D}] For any cardinal $m$, there is a unary algebra $\mathfrak{A} =\langle A,\, F\rangle$ (only unary and nullary operations) and a proper subalgebra $\mathfrak{D}$ of $\mathfrak{A}^{2}$ such that $\mathfrak{A}^{2}$ has exactly four subalgebras, $\mathfrak{A},\,\emptyset,\,\mathfrak{D}$ and $\mathfrak{D}^{\ast}=\{(s,\,t):(t,\,s)\in \mathfrak{D}\}$. Lampe \cite{bib:1974m}.

\noindent [\textbf{159 A}] Every variety has an initial algebra. Blass\index{Author}{Blass} [\textbf{1983b}] and note 32.

\noindent [\textbf{159 B}] For every variety $V$ and set $X,\,V$ has a free algebra on $X$. Blass [\textbf{1983b}] and note 32.

\noindent \textbf{174($\alpha$).} $RM1,\,\aleph_{\alpha}$: The representation theorem for multi-algebras with $\aleph_{\alpha}$ unary operations: Assume $(A,\,F)$ is a multi-algebra with $\aleph_{\alpha}$ unary operations (and no other operations). Then there is an algebra $(B,\,G)$ with $\aleph_{\alpha}$ unary operations and an equivalence relation $E$ on $B$ such that $(B/E,\,G/E)$ and $(A,\, F)$ are isomorphic multi-algebras. H\"{o}ft/Howard \cite{bib:1981b} and note 50.

\noindent \textbf{268.} If $\mathcal{L}$ is a lattice isomorphic to the lattice of subalgebras of some unary universal algebra (a \emph{unary universal algebra} is one with only unary or nullary operations) and $\alpha$ is an automorphism of $\mathcal{L}$ of order 2 (that is, $\alpha^{2}$ is the identity), then there is a unary algebra $\mathfrak{A}$ and an isomorphism $\rho$ from $\mathcal{L}$ onto the lattice of subalgebras of $\mathfrak{A}^{2}$ with $\rho(\alpha(x))=(\rho(x))^{-1}(=\{(s,\,t):(t,\,s)\in\rho(x)\})$ for all $x\in \mathcal{L}$. Lampe\index{Author}{Lampe} \cite{bib:1974m}.

\textbf{V. Vector Spaces}

\noindent [\textbf{$\boldsymbol{0}$ AM}] Assume $V$ is a separable vector space, $D$ is a countable dense subset of $V$, and $p:V\rightarrow \mathbb{R}$ satisfies $p(x+y)\leq p(x)+p(y)$, and $(\forall t\geq 0)(\forall x\in V)(p(tx)=tp(x))$. Also assume that $p$ satisfies the following condition:
\begin{equation*}
\tag{$^{\ast}$} (\forall\epsilon>0)(\forall x\in V)(\exists y\in D)\ \text{such that}\ p(x-y)<\epsilon\ \mathrm{and}\ p(y-x)<\epsilon.
\end{equation*}
Moreover, if $f$ is a linear function from a subspace $S$ of $V$ into $\mathbb{R}$ which satisfies $(\forall x\in S)(f(x)\leq p(x))$, then $f$ can be extended to $f^{\ast}:V\rightarrow \mathbb{R}$ so that $f^{\ast}$ is linear and $(\forall x\in V)(f^{\ast}(x)\leq p(x))$. Note 143.

\noindent [\textbf{1 A}] In every vector space, every generating set contains a basis. Halpern\index{Author}{Halpern} [\textbf{1966}].

\noindent [\textbf{1 BJ}] In every vector space over the two element field, every generating set contains a basis. Keremedis\index{Author}{Keremedis} [\textbf{1996a}].

\noindent [\textbf{1 BU}] In every vector space over $\mathbb{Q}$, every generating set includes a basis. Blass\index{Author}{Blass} [\textbf{1984a}] and note 137.

\noindent [\textbf{1 BZ}] Vector Space Multiple Choice: For every family $V=\{V_{i}:i\in K\}$ of non-trivial vector spaces there is a family $F=\{F_{i}:i\in K\}$ such that for each $i\in K,\,F_{i}$ is a non-empty, finite, independent subset of $V_{i}$. Note 127.

\noindent [\textbf{14 AU}] If $Z$ is a real vector lattice with order unit 1 then the set of all positive linear functionals $u$ on $Z$ with $u(1)=1$ is a convex set that has an extreme point. See H. Rubin/J. Rubin\index{Author}{H. Rubin/J. Rubin} \cite{bib:1985q}, pages 177 and 182 for definitions. Also see Gardiner\index{Author}{Gardiner} \cite{bib:1974i}.

\noindent [\textbf{43 V}] Every convex-compact set in a Hausdorff topological vector space with a basis of convex open sets is Baire. Fossy/Morillon\index{Author}{Fossy/Morillon} [\textbf{1995}] and note 23.

\noindent \textbf{52.} Hahn-Banach Theorem:\index{Subject}{Hahn-Banach theorem} If $V$ is a real vector space and $p:V\rightarrow \mathbb{R}$ satisfies $p(x+y)\leq p(x)+p(y)$ and $(\forall t>0)(p(tx)=tp(x))$ and $S$ is a subspace of $V$ and $f:S\rightarrow \mathbb{R}$ is linear and satisfies $(\forall x\in S)(f(x)\leq p(x))$ then $f$ can be extended to $f^{\ast}:V\rightarrow \mathbb{R}$ such that $f^{\ast}$ is linear and $(\forall x\in V)(f^{\ast}(x)\leq p(x))$. L\'{o}s/Ryll-Nardzewski\index{Author}{Los/Ryll-Nardzewski@{\L}o\'{s}/Ryll-Nardzewski}
 [\textbf{1951}].

\noindent [\textbf{52 F}] If $X$ is a topological vector space then for every continuous sublinear functional $p$ on $E$ there is a linear functional $f$ on $E$ such that $f\leq p$. Fossy/Morillon\index{Author}{Fossy/Morillon} [\textbf{1995}].

\noindent [\textbf{52 G}] If $X$ is a normed vector space then for every continuous sublinear functional $p$ on $E$ there is alinear functional $f$ on $E$ such that $f\leq p$. Fossy/Morillon [\textbf{1995}].

\noindent [\textbf{52 H}] Existence of Real Valued Banach Limits: For any directed set $(\Delta,\,\sqsubseteq)$ there exists a (real valued) Banach limit. Schechter\index{Author}{Schechter} [\textbf{1996b}] and note 31.

\noindent [\textbf{52 I}] Real Valued Convex Domination Theorem: For every real vector space $X$, every subspace $X_{0}\subseteq X$, every linear $\lambda_{0}:X_{0}\rightarrow \mathbb{R}$ and every convex $p:X\rightarrow \mathbb{R}$, if $\lambda_{0}\leq p$ on $X_{0}$ then $\lambda_{0}$ can be extended to a linear map $\lambda:X\rightarrow \mathbb{R}$ such that $\lambda\leq p$ on $X$. Schechter [\textbf{1996b}] and note 31.

\noindent [\textbf{52 J}] Support Theorem for Real Valued Convex Functions: For every real vector space $X$, every convex function $p:X\rightarrow \mathbb{R}$ is a pointwise maximum of affine functions. Schechter [\textbf{1996b}] and note 31.

\noindent [\textbf{52 K}] Sandwich Theorem for Real Valued Functions: For every real vector space $X$, if $e:X\rightarrow \mathbb{R}$ is a concave function, $g:X\rightarrow \mathbb{R}$ is a convex function, and $e\leq g$ everywhere on $X$, then there exists an affine function $f:X\rightarrow \mathbb{R}$ satisfying $e\leq f\leq g$ everywhere on $X$. Schechter [\textbf{1996b}] and note 31.

\noindent \textbf{65.} The Krein-Milman Theorem:\index{Subject}{Krein-Milman theorem} Let $K$ be a compact convex set in a locally convex topological vector space $X$. Then $K$ has an extreme point. (An extreme point is a point which is not an interior point of any line segment which lies in $K$.) Royden\index{Author}{Royden} [\textbf{1963}] and H. Rubin/J. Rubin\index{Author}{H. Rubin/J. Rubin} \cite{bib:1985q}, p. 177.

\noindent \textbf{66.} Every vector space over a field has a basis. Hamel\index{Author}{Hamel} [\textbf{1905}].

\noindent \textbf{95 ($F$).} Existence of Complementary Subspaces over a Field $F$: If $F$ is a field, then every vector space $V$ over $F$ has the property that if $S\subseteq V$ is a subspace of $V$, then there is a subspace $S^{\prime}\subseteq V$ such that $S\cap S^{\prime}=\{0\}$ and $SUS^{\prime}$ generates $V$. Bleicher\index{Author}{Bleicher} [\textbf{1964}], [\textbf{1965}], H. Rubin/J. Rubin \cite{bib:1985q}, pp 119ff, and Jech\index{Author}{Jech} \cite{bib:1973b}, p 148 prob 10.4.

\noindent \textbf{96.} L\"{o}wig's Theorem.\index{Subject}{L\"{o}wig's theorem} If $B_{1}$ and $B_{2}$ are both bases for the vector space $V$ then $|B_{1}|=|B_{2}|$. L\"{o}wig\index{Author}{Lowig@L\"{o}wig} [\textbf{1934}].

\noindent \textbf{109.} Every field $F$ and every vector space $V$ over $F$ has the property that each linearly independent set $A\subseteq V$ can be extended to a basis. Bleicher [\textbf{1964}], [\textbf{1965}], Blass\index{Author}{Blass} [\textbf{1984a}], H. Rubin/J. Rubin \cite{bib:1985q}, pp 119ff.

\noindent \textbf{110.} Every vector space over $\mathbb{Q}$ has a basis. Blass [\textbf{1984a}].

\noindent \textbf{195.} Every general linear system has a linear global reaction. Note 62 and Harazisville (see Kharazishvili\index{Author}{Kharazishvili} [\textbf{1979}]).

\noindent [\textbf{218 A}] Existence of Complementary Subspaces: For every field $F$, every vector space $V$ over $F$ has the property that if $S\subseteq V$ is a subspace of $V$, then there is a subspace $S^{\prime}\subseteq V$ such that $S\cap S^{\prime}=\{0\}$ and $SUS^{\prime}$ generates $V$. Bleicher [\textbf{1964}], [\textbf{1965}], H. Rubin/J. Rubin \cite{bib:1985q}, pp 123ff, Theorems 6.36, 6.37, and 6.38.

\noindent \textbf{235.} If $V$ is a vector space and $B_{1}$ and $B_{2}$ are bases for $V$ then $|B_{1}|$ and $|B_{2}|$ are comparable. Gross\index{Author}{Gross} [\textbf{1976}].

\noindent \textbf{236.} If $V$ is a vector space with a basis and $S$ is a linearly independent subset of $V$ such that no proper extension of $S$ is a basis for $V$, then $S$ is a basis for $V$. Gross [\textbf{1975}], [\textbf{1976}].

\noindent [\textbf{236 A}] If $V$ is a vector space with a basis, then every linearly independent subset is contained in a basis.

\noindent \textbf{286.} Extended Krein-Milman Theorem:\index{Subject}{Krein-Milman theorem} Let $\mathrm{K}$ be a quasicompact (sometimes called convex-compact), convex subset of a locally convex topological vector space, then $\mathrm{K}$ has an extreme point. H. Rubin/J. Rubin\index{Author}{H. Rubin/J. Rubin} \cite{bib:1985q}, p. 177--178 and note 23.

\noindent \textbf{287.} The Hahn-Banach Theorem\index{Subject}{Hahn-Banach theorem} for Separable Normed Linear Spaces: Assume $V$ is a separable normed linear space and $p:V\rightarrow \mathbb{R}$ satisfies $p(x+y)\leq p(x)+ p(y)$ and $(\forall t\geq 0)(\forall x\in V)(p(tx)=tp(x))$ and assume $f$ is a linear function from a subspace $S$ of $V$ into $\mathbb{R}$ which satisfies $(\forall x\in S)(f(x)\leq p(x))$, then $f$ can be extended to $f^{\ast}:V\rightarrow \mathbb{R}$ so that $f^{\ast}$ is linear and $(\forall x\in V)(f^{\ast}(x)\leq p(x))$. Note 145.

\noindent [\textbf{333 A}] In every vector space $B$ over the two element field, every subspace of $B$ has a complementary subspace. Keremedis\index{Author}{Keremedis} [\textbf{1996a}].

\noindent \textbf{346.} Vector Space Kinna-Wagner Principle: For every family $V=\{V_{i}:i\in K\}$ of non-trivial vector spaces there is a family $F=\{F_{i}:i\in K\}$ such that for each $i\in K,\,F_{i}$ is a non-empty, independent subset of $V_{i}$. Note 127.

\noindent \textbf{367.} There is a Hamel\index{Author}{Hamel} basis for $\mathbb{R}$ as a vector space over $\mathbb{Q}$. Hamel [\textbf{1905}].

\noindent \textbf{372.} Generalized Hahn-Banach Theorem: Assume that $X$ is a real vector space, $(Z,\preccurlyeq)$ is a Dedekind complete ordered vector space and $X_{0}$ is a subspace of $X$. If $\lambda_{0}:X_{0}\rightarrow Z$ is linear and $p:X\rightarrow Z$ is sublinear and if $\lambda_{0}\preccurlyeq p$ on $X_{0}$ then $\lambda_{0}$ can be extended to a linear map $\lambda:X\rightarrow Z$ such that $\lambda\preccurlyeq p$ on $X$. Schechter\index{Author}{Schechter} [\textbf{1996b}] and note 31.

\noindent [\textbf{372 A}] Existence of Banach Limits: For any directed set $(\triangle,\,\sqsubseteq)$ and any Dedekind complete ordered vector space $(Z,\preccurlyeq)$ there exists a $Z$-valued Banach limit. Schechter [\textbf{1996b}] and note 31.

\noindent [\textbf{372 B}] Generalized Convex Domination Theorem:\index{Subject}{convex domination theorem} For every real vector space $X$, every Dedekind complete ordered vector space $(Z,\preccurlyeq)$, every subspace $X_{0}\subseteq X$, every linear $\lambda_{0}:X_{0}\rightarrow Z$ and every convex $p:X\rightarrow Z$, if $\lambda_{0}\preccurlyeq p$ on $X_{0}$ then $\lambda_{0}$ can be extended to a linear map $\lambda:X\rightarrow Z$ such that $\lambda\preccurlyeq p$ on $X$. Schechter [\textbf{1996b}] and note 31.

\noindent [\textbf{372 C}] Generalized Support Theorem:\index{Subject}{support theorem} For every real vector space $X$ and every Dedekind complete ordered vector space $(Z,\preccurlyeq)$, every convex function $p:X\rightarrow Z$ is a pointwise maximum of affine functions. Schechter [\textbf{1996b}] and note 31.

\noindent [\textbf{372 D}] Generalized Sandwich Theorem: For every real vector space $X$ and every Dedekind complete ordered vector space $(Z,\,\preccurlyeq)$, if $e:X\rightarrow Z$ is a concave function, $g:X\rightarrow Z$ is a convex function, and $e\leq g$ everywhere on $X$, then there exists an affine function $f:X\rightarrow Z$ satisfying $e\leq f\leq g$ everywhere on $X$. Schechter\index{Author}{Schechter} [\textbf{1996b}] and note 31.

\section*{ANALYSIS (FORMS FROM ANALYSIS)}

\noindent [\textbf{$\boldsymbol{0}$ H}] F. Riesz Theorem:\index{Subject}{Riesz theorem} A Hilbert space is finite dimensional iff its closed unit ball is compact. Brunner\index{Author}{Brunner} [\textbf{1983b}].

\noindent [\textbf{$\boldsymbol{0}$ I}] The following are equivalent in a Hilbert space:
\begin{enumerate}
\item[(a)] The closed unit ball is sequentially compact.
\item[(b)] Each unconditionally convergent series converges absolutely.
\item[(c)] Each ortho-normal system is Dedekind finite.
\end{enumerate}
Brunner [\textbf{1983b}] and note 94.

\noindent [\textbf{$\boldsymbol{0}$ AM}] Assume $V$ is a separable vector space, $D$ is a countable dense subset of $V$, and $p:V\rightarrow \mathbb{R}$ satisfies $p(x+y)\leq p(x)+p(y)$, and $(\forall t\geq 0)(\forall x\in V)(p(tx)=tp(x))$. Also assume that $p$ satisfies the following condition:
\begin{equation*}
\tag{$^{\ast}$} (\forall\epsilon>0)(\forall x\in V)(\exists y\in D)\ \text{such that}\ p(x-y)<\epsilon\ \mathrm{and}\ p(y-x)<\epsilon.
\end{equation*}
Moreover, if $f$ is a linear function from a subspace $S$ of $V$ into $\mathbb{R}$ which satisfies $(\forall x\in S)(f(x)\leq p(x))$, then $f$ can be extended to $f^{\ast}:V\rightarrow \mathbb{R}$ so that $f^{\ast}$ is linear and $(\forall x\in V)(f^{\ast}(x)\leq p(x))$. Note 143.

\noindent [\textbf{8 E}] In a metric space every sequentially continuous real valued function is continuous. Brunner [\textbf{1982c}].

\noindent [\textbf{10 D}] In every Hilbert space with a Dedekind finite orthonormal basis, the closed unit ball is sequentially compact. Brunner [\textbf{1983b}] and note 94.

\noindent [\textbf{14 Q}] Alaoglu's Theorem:\index{Subject}{Alaoglu's theorem} The unit ball is compact in the weak$^{\ast}$ topology in the dual of a Banach space. (See [52 E].) Johnstone\index{Author}{Johnstone} [\textbf{1984}] and note 23.

\noindent [\textbf{14 U}] Every closed proper ideal in a Gelfand algebra is contained in some maximal ideal. Banaschewski [\textbf{1983}] and note 29.

\noindent [\textbf{14 BY}] For every topological space $X$, every proper ideal in the normed $\mathbb{R}$-algebra of continuous bounded functions from $X$ into $\mathbb{R}$ is contained in a maximal ideal. Morillon\index{Author}{Morillon} \cite{bib:1988h}.

\noindent [\textbf{14 BZ}] Every proper ideal in a non-trivial normed $\mathbb{R}$-algebra or $\mathbb{C}$-algebra with identity is contained in a maximal ideal. Morillon \cite{bib:1988h}.

\noindent [\textbf{14 CA}] If $(A_{i})_{i\in I}$ is a family of non trivial commutative normed algebras with identity then there is a family $(J_{i})_{i\in I}$ such that for each $i\in I,\,J_{i}$ is a maximal ideal in $A_{i}$. Morillon \cite{bib:1988h}.

\noindent [\textbf{14 CB}] If $X$ is a topological space and $I$ is a proper ideal in the algebra of all bounded continuous functions from $X$ to $\mathbb{R}$ such that every function in $I$ vanishes at at least one point of $X$ then $I$ is contained in a maximal ideal. Morillon \cite{bib:1988h}.

\noindent [\textbf{14 CL}] Gelfand Extreme Point Theorem:\index{Subject}{Gelfand extreme point theorem} If $A$ is a non-trivial Gelfand algebra then the closed unit ball in the dual of $A$ has an extreme point $e$ with $e\neq 0$. Morillon [\textbf{1986}], notes 23 and 29.

\noindent [\textbf{52 E}] The unit ball is convex-compact in the weak$^{\ast}$ topology in the dual of a Banach space. (See [14 Q].) Luxemburg\index{Author}{Luxemburg} \cite{bib:1969d}, H. Rubin /J. Rubin \cite{bib:1985q} p 178, and note 23.

\noindent \textbf{65.} The Krein-Milman Theorem:\index{Subject}{Krein-Milman theorem} Let $K$ be a compact convex set in a locally convex topological vector space $X$. Then $K$ has an extreme point. (An extreme point is a point which is not an interior point of any line segment which lies in $K$.) Royden\index{Author}{Royden} [\textbf{1963}] and H. Rubin/J. Rubin\index{Author}{H. Rubin/J. Rubin} \cite{bib:1985q}, p. 177.

\noindent [\textbf{118 V}] If $(L,\,\leq)$ is a linear order and $G=\{[l_{i},\,r_{i}]:i\in K\}$ is a locally finite family of closed intervals of $L$ whose interiors are non-empty and pairwise disjoint, then there exists a continuous real valued function $f:L\rightarrow \mathbb{R}$ such that $|f^{-1}(0)\cap[l_{i},\,r_{i}]|\geq 2$ and $|f^{-1}(1) \cap [l_{i},\,r_{i}]|\geq 1$ for all $i\in K$. Keremedis\index{Author}{Keremedis} [\textbf{1995}] and note 43.

\noindent [\textbf{118 W}] If $(L,\,\leq)$ is a linear order and $G=\{[l_{i},\,r_{i}]:i\in K\}$ is a locally finite family of closed intervals of $L$ whose interiors are non-empty and pairwise disjoint, then there exists a continuous real valued function $f:L\rightarrow \mathbb{R}$ such that for each $i\in K,\,f(l_{i})=f(r_{i})=0$ and $f(x)>0$ for all $x\in(l_{i},\,r_{i})$. Keremedis [\textbf{1995}] and note 43.

\noindent [\textbf{118 X}] If $(L,\,\leq)$ is a complete linear order and $G=\{(l_{i},\,r_{i}):i\in K\}$ is a family of open, non-empty, pairwise disjoint intervals of $L$, then there exists a set $\{f_{i}:i\in K\}$ of non constant continuous real valued functions on $L$ such that for each $i\in K,\,f_{i}$ vanishes outside of $(l_{i},\, r_{i})$. Keremedis [\textbf{1995}].

\noindent \textbf{124.} Every operator on a Hilbert space with an amorphous base is the direct sum of a finite matrix and a scalar operator. (A set is amorphous if it is not the union of two disjoint infinite sets.) Brunner\index{Author}{Brunner} [\textbf{1984a}].

\noindent [\textbf{126 A}] Every positive functional on a $B$-lattice is continuous. Brunner\index{Author}{Brunner} [\textbf{1984d}] and note 16.

\noindent \textbf{153.} The closed unit ball of a Hilbert space is compact in the weak topology. Fossy/Morillon\index{Author}{Fossy/Morillon} [\textbf{1995}] and Rudin\index{Author}{Rudin} [\textbf{1991}].

\noindent \textbf{156.} Theorem of Gelfand and Kolmogoroff:\index{Subject}{Gelfand-Kolmogoroff theorem} Two compact $T_{2}$ spaces are homeomorphic if their rings of real valued continuous functions are isomorphic. Brunner [\textbf{1983a}] and Gelfand/Kolmogoroff\index{Author}{Gelfand/Kolmogoroff} [\textbf{1939}].

\noindent \textbf{157.} Theorem of Goodner:\index{Subject}{Goodner's theorem} A compact $T_{2}$ space is extremally disconnected (the closure of every open set is open) if and only if each non-empty subset of $C(X)$ (set of continuous real valued functions on $X$) which is pointwise bounded has a supremum. Goodner\index{Author}{Goodner} [\textbf{1950}].

\noindent \textbf{158.} In every Hilbert space $H$, if the closed unit ball is sequentially compact, then $H$ has an orthonormal basis. Brunner [\textbf{1983b}].

\noindent \textbf{279.} The Closed Graph Theorem\index{Subject}{closed graph theorem} for operations between Fr\'{e}chet Spaces: Suppose $X$ and $Y$ are Fr\'{e}chet spaces, $T:X\rightarrow Y$ is linear and $G=\{(x,\,Tx): x\in X\}$ is closed in $X\times Y$. Then $T$ is continuous. Rudin [\textbf{1991}] p. 51, Wright \cite{bib:1973h}, note 96.

\noindent \textbf{281.} There is a Hilbert space $H$ and an unbounded linear operator on $H$. Wright \cite{bib:1973h}, note 96.

\noindent \textbf{286.} Extended Krein-Milman Theorem:\index{Subject}{Krein-Milman theorem} Let $\mathrm{K}$ be a quasicompact (sometimes called convex-compact), convex subset of a locally convex topological vector space, then $\mathrm{K}$ has an extreme point. H. Rubin/J. Rubin\index{Author}{H. Rubin/J. Rubin} \cite{bib:1985q}, p. 177--178 and note 23.

\noindent \textbf{287.} The Hahn-Banach Theorem\index{Subject}{Hahn-Banach theorem} for Separable Normed Linear Spaces: Assume $V$ is a separable normed linear space and $p:V\rightarrow \mathbb{R}$ satisfies $p(x+y)\leq p(x)+ p(y)$ and $(\forall t\geq 0)(\forall x\in V)(p(tx)=tp(x))$ and assume $f$ is a linear function from a subspace $S$ of $V$ into $\mathbb{R}$ which satisfies $(\forall x\in S)(f(x)\leq p(x))$, then $f$ can be extended to $f^{\ast}:V\rightarrow \mathbb{R}$ so that $f^{\ast}$ is linear and $(\forall x\in V)(f^{\ast}(x)\leq p(x))$. Note 145.

\noindent \textbf{370.} Weak Gelfand Extreme Point Theorem:\index{Subject}{Gelfand extreme point theorem} If $A$ is a non-trivial Gelfand algebra then the closed unit ball in the dual of $A$ has an extreme point $e$. Morillon\index{Author}{Morillon} [\textbf{1986}], notes 23 and 29.

\section*{CARDINAL NUMBER FORMS}

\textbf{I. Alephs and their Properties}

\noindent [\textbf{$\boldsymbol{0}$ AL}] For every set $A$, there is an $\aleph$ which cannot be covered by fewer than $|A|$ sets each of cardinality $<|A|$. Note 131.

\noindent \textbf{25.} $\aleph_{\beta+1}$ is regular for all ordinals $\beta$. Feferman/Levy\index{Author}{Feferman/Levy}\index{Author}{Levy} [\textbf{1963}] and Gitik\index{Author}{Gitik} [\textbf{1980}].

\noindent \textbf{34.} $\aleph_{1}$ is regular. Russell\index{Author}{Russell} [\textbf{1906}] and note 3.

\noindent [\textbf{40 B}] For all ordinals $\alpha$, the $\aleph_{\alpha}$ partition principle holds: For every ordinal $\alpha$ and every cardinal $\kappa$, if $\aleph_{\alpha}\leq^{\ast}\kappa$, then $\aleph_{\alpha}\leq\kappa$. Pelc\index{Author}{Pelc} [\textbf{1978}], Banaschewski/Moore\index{Author}{Banaschewski/Moore} [\textbf{1990}], and note 69.

\noindent \textbf{41.} $W_{\aleph_{1}}$: For every cardinal $m,\,m\leq\aleph_{1}$ or $\aleph_{1}\leq m$. Jech \cite{bib:1973b}.

\noindent \textbf{54.} For all infinite cardinals $m,\,m$ adj $2^{m}$ implies $m$ is an aleph. ($m\ adj\ n$ iff $m<n$ and $(\neg\exists p)(m<p<n)$.) Mathias\index{Author}{Mathias} [\textbf{1979}] prob 1337, thm 1338 and prob 1339.

\noindent \textbf{56.} $\aleph(2^{\aleph_{0}})\neq\aleph_{\omega}$. ($\aleph(2^{\aleph_{0}})$ is Hartogs' aleph, the least $\aleph$ not $\leq 2^{\aleph_{0}}$.) Mathias [\textbf{1979}] p 125.

\noindent \textbf{58.} There is an ordinal $\alpha$ such that $\aleph(2^{\aleph_{\alpha}})\neq\aleph_{\alpha+1}$. $(\aleph(2^{\aleph_{\alpha}})$ is Hartogs' aleph, the least $\aleph$ not $\leq 2^{\aleph_{\alpha}}$.) Mathias [\textbf{1979}] p 126.

\noindent \textbf{71($\alpha$).} $W_{\aleph_{\alpha}}:(\forall x)(|x|\leq\aleph_{\alpha}$ or $|x|\geq\aleph_{\alpha}$). Pelc [\textbf{1978}], Jech\index{Author}{Jech} \cite{bib:1973b}, page 119 and Levy [\textbf{1964}].

\noindent \textbf{104.} There is a regular uncountable aleph. Jech [\textbf{1966b}] p 165 prob 11.26.

\noindent \textbf{108.} There is an ordinal $\alpha$ such that $2^{\aleph_{\alpha}}$ is not the union of a denumerable set of denumerable sets. Gitik [\textbf{1985}].

\noindent \textbf{159.} The regular cardinals are cofinal in the class of ordinals. Blass\index{Author}{Blass} [\textbf{1983b}].

\noindent \textbf{170.} $\aleph_{1}\leq 2^{\aleph_{0}}$. Bernstein\index{Author}{Bernstein} [\textbf{1908}] and Banaschewski/Moore [\textbf{1990}].

\noindent \textbf{182.} There is an aleph whose cofinality is greater than $\aleph_{0}$. Gitik [\textbf{1980}].

\noindent \textbf{183($\alpha$).} There are no $\aleph_{\alpha}$ minimal sets. That is, there are no sets $X$ such that
\begin{enumerate}
\item[(1)] $|X|$ is incomparable with $\aleph_{\alpha}$
\item[(2)] $\aleph_{\beta}<|X|$ for every $\beta<\alpha$ and
\item[(3)] $\forall Y\subseteq X,\,|Y|<\aleph_{\alpha}$ or $|X-Y|<\aleph_{\alpha}$.
\end{enumerate}
Hickman\index{Author}{Hickman} [\textbf{1980b}] and note 56.

\noindent \textbf{208.} For all ordinals $\alpha,\,\aleph_{\alpha+1}\leq 2^{\aleph_{\alpha}}$. Figura\index{Author}{Figura} [\textbf{1977}].

\noindent \textbf{209.} There is an ordinal $\alpha$ such that for all $X$, if $X$ is a denumerable union of denumerable sets then $\mathcal{P}(X)$ cannot be partitioned into $\aleph_{\alpha}$ non-empty sets. Morris\index{Author}{Morris} [\textbf{1970}].

\noindent \textbf{245.} There is a function $f:\omega_{1}\rightarrow\omega_{1}^{\omega}$ such that for every $\alpha,\, 0<\alpha<\omega_{1},\,f(\alpha)$ is a function from $\omega$ onto $\alpha$. Litman\index{Author}{Litman} [\textbf{1976}].

\noindent [\textbf{245 A}] There is a function $g:\omega_{1}\rightarrow\omega_{1}^{\omega}$ such that for all limit ordinals $\alpha<\omega_{1},\,g(\alpha)$ is an increasing function from $\omega$ into $\alpha$ with the range of $g$ cofinal in $\alpha$. Litman [\textbf{1976}]

\noindent [\textbf{245 B}] Every closed set A in $\omega_{1}$ which contains only limit ordinals is a derivative (that is, $\exists B\subseteq\omega_{1}$ such that $A=\{x:x\ \text{is a limit and}\ \{t:t\in B\}\ \text{is cofinal in}\ x\}$.) Litman [\textbf{1976}].

\noindent \textbf{246.} The monadic theory theory $MT(\omega_{1},\,<)$ of $\omega_{1}$ is recursive. Litman [\textbf{1976}] and note 85.

\noindent \textbf{275.} The sequence of cardinals $\langle\aleph_{n}: n\in\omega\rangle$ has a unique minimal upper bound. Truss\index{Author}{Truss} [\textbf{1974b}].

\noindent \textbf{315.} $\Omega=\omega_{1}$, where $\Omega=\{\alpha\in\ \mathrm{On}: (\forall\beta\leq\alpha)(\beta=0\vee(\exists\gamma)(\beta=\gamma+ 1) \vee$ there is a sequence $\langle\gamma_{n}: n\in\omega\rangle$ such that for each $n\ \gamma_{n}<\beta$ and $\beta=\bigcup\nolimits_{n<\omega}\gamma_{n}\}$ Gitik\index{Author}{Gitik} [\textbf{1980}].

\noindent \textbf{365.} For every uncountable set $A$, if $A$ has the same cardinality as each of its uncountable subsets then $|A|=\aleph_{1}$. Sierpi\'{n}ski\index{Author}{Sierpinski@Sierpi\'{n}ski} [\textbf{1918}].

\textbf{II. Cardinal Arithmetic}

\noindent [\textbf{$\boldsymbol{0}$ A}] Cardinal successors 2: For every cardinal $m$ there is a cardinal $n$ such that $m<n$ and $(\neg\exists p)(m<p<n)$. Tarski\index{Author}{Tarski} [\textbf{1954a}] and Jech\index{Author}{Jech} [\textbf{1966a}].

\noindent [\textbf{$\boldsymbol{0}$ F}] Cancellation for surjective cardinal equivalence: For every $k\in w-\{0\},\,(\forall x\forall y)(kx=^{\ast} ky\ \mathrm{implies}\ x=^{\ast} y)$ (The expression $x=^{\ast}y$ means that there is a function $f$ from $x$ onto $y$ and a function $g$ from $y$ onto $x$.) Truss [\textbf{1984}].

\noindent [\textbf{$\boldsymbol{0}$ G}] Cancellation for injective cardinal inequality: For every $k\in \omega-\{0\},\,(\forall x\forall y)(kx\leq ky\ \mathrm{implies}\ x\leq y)$ ($\leq$ is the usual (injective) cardinal ordering.) Lindenbaum/Tarski\index{Author}{Lindenbaum/Tarski} [\textbf{1926}].

\noindent [\textbf{$\boldsymbol{0}$ Q}] For all cardinal numbers $m$ and $n$ and Dedekind finite cardinals $p$, if $p+m=p+n$ then $m=n$. Tarski [\textbf{1949b}] and note 94

\noindent [\textbf{1 D}] Cardinal successors 3: For every cardinal $m$ there is a cardinal $n$ such that $m<n$ and $(\forall p)(m<p\rightarrow n\leq p)$. Tarski [\textbf{1954a}] and Jech [\textbf{1966a}].

\noindent \textbf{2.} Existence of successor cardinals: For every cardinal $m$ there is a cardinal $n$ such that $m<n$ and $(\forall p<n)(p\leq m)$. Tarski [\textbf{1954a}].

\noindent \textbf{3.} $2m=m$: For all infinite cardinals $m,\,2m=m$. Tarski [\textbf{1954a}].

\noindent [\textbf{3 A}] For all cardinals $m$ and $n,\,n\leq m$ and $m$ infinite $\rightarrow m+n=m$. Halpern/Howard [\textbf{1970}].

\noindent [\textbf{3 B}] $\aleph_{0}\cdot m=m$: For all infinite cardinals $m,\,\aleph_{0}\cdot m=m$. Halpern/Howard\index{Author}{Halpern/Howard} [\textbf{1970}].

\noindent [\textbf{3 C}] Definability of cardinal addition as the least upper bound: (For all cardinals $x,\,y$ and $z$) ($x+y=z$ iff ($x\leq z$ and $y\leq z$ and $(\forall u)(x,\,y\leq u\rightarrow z\leq u)$). H\"{a}ussler\index{Author}{Haussler@H\"{a}ussler} [\textbf{1983}] and Tarski\index{Author}{Tarski} [\textbf{1949a}].

\noindent [\textbf{9 F}] There is no subset $A$ of the class of Dedekind finite cardinals such that $A$ under the usual cardinal ordering is order isomorphic to $\langle \mathbb{R}, \leq\rangle$. Tarski [\textbf{1965}], Jech\index{Author}{Jech} \cite{bib:1973b} p 161 prob 6, and note 94.

\noindent [\textbf{9 K}] If $A$ is uncountable and $B$ is countable, then $|A-B|=|A|$. G. Moore\index{Author}{G. Moore} [\textbf{1982}] p 201.

\noindent [\textbf{9 L}] If $A$ is uncountable and $B$ is countable, then $|A\cup B|=|A|$. G. Moore [\textbf{1982}] p 202.

\noindent [\textbf{9 M}] If $\aleph_{0}<|A|$ and $\aleph_{0}=|B|$, then $\aleph_{0}<|A-B|$. G. Moore [\textbf{1982}] p 202.

\noindent \textbf{[9 N]} If $\aleph_{0}\leq^{\ast}m$, then $\aleph_{0}\leq m$. Lindenbaum/Tarski\index{Author}{Lindenbaum/Tarski} [\textbf{1926}], G. Moore [\textbf{1982}] p 216, and note 69.

\noindent [\textbf{9 O}] $m\leq\aleph_{0}$ or $\aleph_{0}\leq m$. Lindenbaum/Tarski [\textbf{1926}] and G. Moore [\textbf{1982}] p 216,

\noindent [\textbf{9 P}] If $m+\aleph_{0}<n +\aleph_{0}$, then $m<n$. Lindenbaum/Tarski [\textbf{1926}] and G. Moore [\textbf{1982}] p 216.

\noindent [\textbf{9 Q}] If $m+\aleph_{0}=n+\aleph_{0}$, then $m=n$ or $m,\,n\leq\aleph_{0}$. Lindenbaum/Tarski [\textbf{1926}] and G. Moore [\textbf{1982}] p 216.

\noindent \textbf{53.} For all infinite cardinals $m,\,m^{2}\leq 2^{m}$. Mathias [\textbf{1979}] prob 1336.

\noindent \textbf{55.} For all infinite cardinals $m$ and $n$, if $m\cdot n=m+n$ then, $m\leq n$ or $n\leq m$. Mathias\index{Author}{Mathias} [\textbf{1979}] p 125.

\noindent \textbf{136($k$).} Surjective Cardinal Cancellation (depends on $k\in\omega-\{0\}$): For all cardinals $x$ and $y,\,kx\leq^{\ast}ky$ implies $x\leq^{\ast}y$. Truss\index{Author}{Truss} [\textbf{1984}] and note 69.

\noindent \textbf{161.} Definability of cardinal addition in terms of $\leq$: There is a first order formula whose only non-logical symbol is $\leq$ (for cardinals) that defines cardinal addition. H\"{a}ussler~[\textbf{1983}].

\noindent \textbf{162.} Non-existence of infinite units: There is no infinite cardinal number $A$ such that $A+A>A$ and for all cardinals $x$ and $y,\,x+y=A\rightarrow x=A$ or $y=A$. H\"{a}ussler~[\textbf{1983}].

\noindent \textbf{200.} For all infinite $x,\,|2^{x}|=|x!|$. Dawson/Howard\index{Author}{Dawson/Howard} [\textbf{1976}] and note 64.

\noindent \textbf{204.} For every infinite $X$, there is a function from $X$ onto $2X$. Truss [\textbf{1978}].

\noindent \textbf{277.} $E(D,\,VII)$: Every non-well-orderable cardinal is decomposable. Monro\index{Author}{Monro} [\textbf{1973a}] and note 94.

\noindent \textbf{283.} Cardinality of well ordered subsets: For all $n\in\omega$ and for all infinite $x,\,|x^{n}|<|s(x)|$ where $s(x)$ is the set of all well orderable subsets of $x$. Truss [\textbf{1973d}].

\noindent \textbf{290.} For all infinite $x,\,|2^{x}|=|x^{x}|$. Pincus [\textbf{1978}] and note 64.

\noindent \textbf{291.} For all infinite $x,\, |x!|=|x^{x}|$. Pincus\index{Author}{Pincus} [\textbf{1978}] and note 64.

\textbf{III. Cardinality of Unions}

\noindent [\textbf{1 C}] For all infinite cardinals $m,\,UT(m,\,m,\,m)$: If every member of an infinite set $A$ has power $|A|$ then $|\bigcup A|=|A|$. (Note that [1 C] implies $m^{2}=m$ for all infinite cardinals m.) Howard\index{Author}{Howard} [\textbf{1994}].

\noindent \textbf{6.} $UT(\aleph_{0},\,\aleph_{0},\,\aleph_{0},\,\mathbb{R})$: The union of a denumerable family of denumerable subsets of $\mathbb{R}$ is denumerable. G. Moore\index{Author}{G. Moore} [\textbf{1982}] and note 3.

\noindent [\textbf{8 T}] $PUT(\aleph_{0},\,\infty,\,\aleph_{0})$: The union of a denumerable set of pairwise disjoint non-empty sets has a denumerable subset. $([8\ \mathrm{A}]\rightarrow[8\ \mathrm{T}]\rightarrow[8\ \mathrm{B}])$

\noindent [\textbf{8 U}] $PUT(WO,\,\infty,\,WO)$: The union of an infinite well ordered set of pairwise disjoint non-empty sets has an infinite well ordered subset. $([8\ \mathrm{U}]\leftrightarrow[8\ \mathrm{T}])$

\noindent [\textbf{10 A}] $UT(\aleph_{0},\,<\aleph_{0},\,\aleph_{0})$: The union of denumerably many pairwise disjoint finite sets is denumerable. Shannon\index{Author}{Shannon} [\textbf{1988}] p 569.

\noindent [\textbf{10 M}] $PUT(\aleph_{0},\,<\aleph_{0},\,\aleph_{0})$: The union of denumerably many pairwise disjoint finite sets has a denumerable subset.

\noindent [\textbf{10 N}] $PUT(WO,\,<\aleph_{0},\,WO)$: The union of an infinite well ordered set of pairwise disjoint finite sets has an infinite well ordered subset.

\noindent \textbf{18.} \emph{PUT} $(\aleph_{0},2,\,\aleph_{0})$: The union of a denumerable family of pairwise disjoint pairs has a denumerable subset. Brunner\index{Author}{Brunner} [\textbf{1984f}].

\noindent [\textbf{18 B}] The union of a denumerable family of pairs has a denumerable subset. Note 135.

\noindent [\textbf{18 D}] $PUT(WO,\,2,\,WO)$: The union of an infinite well ordered family of pairwise disjoint pairs has an infinite well ordered subset. $([18\ \mathrm{D}\leftrightarrow 18)$

\noindent \textbf{20.} If $\{A_{x}:x\in S\}$ and $\{B_{x}:x\in S\}$ are families of pairwise disjoint sets and $|A_{x}|=|B_{x}|$ for all $x\in S$, then $|\bigcup_{x\in S}A_{x}|= |\bigcup_{x\in S}B_{x}|$. G. Moore [\textbf{1982}] (1.4.12 and 1.7.8).

\noindent \textbf{21.} If $S$ is well ordered, $\{A_{x}:x\in S\}$ and $\{B_{x}:x\in S\}$ are families of pairwise disjoint sets, and $|A_{x}|=|B_{x}|$ for all $x\in S$, then $|\bigcup_{x\in S}A_{x}|=|\bigcup_{x\in S}B_{x}|$. G. Moore [\textbf{1982}].

\noindent \textbf{22.} $UT(2^{\aleph_{0}},2^{\aleph_{0}},2^{\aleph_{0}})$: If every member of an infinite set of cardinality $2^{\aleph_{0}}$ has power $2^{\aleph_{0}}$, then the union has power $2^{\aleph_{0}}$. Borel\index{Author}{Borel} [\textbf{1898}] and G. Moore [\textbf{1982}].

\noindent \textbf{23.} $(\forall\alpha)(UT(\aleph_{\alpha},\,\aleph_{\alpha},\,\aleph_{\alpha}))$: For every ordinal $\alpha$, if $A$ and every member of $A$ has cardinality $\aleph_{\alpha}$, then $|\bigcup A|=\aleph_{\alpha}$. G. Moore [\textbf{1982}].

\noindent \textbf{26.} $UT(\aleph_{0},2^{\aleph_{0}},2^{\aleph_{0}})$: The union of denumerably many sets each of power $2^{\aleph_{0}}$ has power $2^{\aleph_{0}}$. Borel [\textbf{1898}].

\noindent \textbf{27.} $(\forall\alpha)(UT(\aleph_{0},\,\aleph_{\alpha},\,\aleph_{\alpha}))$: The union of denumerably many sets each of power $\aleph_{\alpha}$ has power $\aleph_{\alpha}$. G. Moore [\textbf{1982}] p 36.

\noindent \textbf{29.} If $|S|=\aleph_{0}$ and $\{A_{x}:x\in S\}$ and $\{B_{x}:x\in S\}$ are families of pairwise disjoint sets and $|A_{x}|=|B_{x}|$ for all $x\in S$, then $|\bigcup_{x\in S}A_{x}|=|\bigcup_{x\in S}B_{x}|$. G. Moore [\textbf{1982}] p 324.

\noindent \textbf{31.} $UT(\aleph_{0},\,\aleph_{0},\,\aleph_{0})$: The countable union theorem: The union of a denumerable set of denumerable sets is denumerable. Cantor\index{Author}{Cantor} [\textbf{1878}].

\noindent [\textbf{38 A}] $UT(\aleph_{0},\,\aleph_{0},\,\neq 2^{\aleph_{0}})$: A denumerable union of denumerable sets does not have cardinality $2^{\aleph_{0}}$.

\noindent [\textbf{80 A}] $UT(\aleph_{0},2,\,\aleph_{0})$: The union of denumerably many pairs is denumerable.

\noindent \textbf{111.} $UT(WO,\,2,\,WO)$: The union of an infinite well ordered set of 2-element sets is an infinite well ordered set.

\noindent [\textbf{122 A}] $UT(WO,\,<\aleph_{0},\,WO)(U_{\aleph_{0}})$: The union of a well ordered set of finite sets is well orderable. Brunner/Howard\index{Author}{Brunner/Howard} [\textbf{1992}]. (See note 27 for more information on $UT(WO,\,\kappa,\,WO),\,\kappa$ a well ordered cardinal. The proof that 122 implies [122 A] is similar to the proof given in Shannon\index{Author}{Shannon} [\textbf{1988}] p 569.)

\noindent \textbf{151.} $UT(WO,\,\aleph_{0},\,WO)(U_{\aleph_{1}})$: The union of a well ordered set of denumerable sets is well orderable. Brunner/Howard [\textbf{1992}]. (If $\kappa$ is a well ordered cardinal, see note 27 for $UT(WO,\,\kappa,\, WO)$.)

\noindent \textbf{207($\alpha$).} $UT(\aleph_{\alpha},\,\aleph_{\alpha},\,<2^{\aleph_{\alpha}})$: The union of $\aleph_{\alpha}$ sets each of cardinality $\aleph_{\alpha}$ has cardinality less than $2^{\aleph_{\alpha}}$. Figura\index{Author}{Figura}~[\textbf{1977}].

\noindent \textbf{231.} $UT(WO,\,WO,\,WO)$: The union of a well ordered collection of well orderable sets is well orderable. Tsukada~[\textbf{1977}].

\noindent \textbf{296.} Part-$\infty$: Every infinite set is the disjoint union of infinitely many infinite sets. Gonzalez\index{Author}{Gonzalez} [\textbf{1995b}], Pincus\index{Author}{Pincus} [\textbf{1997}].

\noindent \textbf{359.} If $\{A_{x}:x\in S\}$ and $\{B_{x}:x\in S\}$ are families of pairwise disjoint sets and $|A_{x}|\leq|B_{x}|$ for all $x\in S$, then $|\bigcup_{x\in S}A_{x}|\leq |\bigcup_{x\in S}B_{x}|$. Higasikawa\index{Author}{Higasikawa} [\textbf{1995}].

\noindent \textbf{[373 $\mathbf{B}(n)$]} $PUT(\aleph_{0},\,n,\,\aleph_{0})$ for $n\in \omega-\{0,1\}$: The union of a denumerable set of pairwise disjoint $n$-element sets has a denumerable subset.

\noindent \textbf{[373 $\mathbf{C}(n)$]} $PUT(WO,\,n,\,WO)$ for $n\in\omega -\{0,1\}$: The union of an infinite well ordered set of pairwise disjoint $n$-element sets has an infinite well ordered subset.

\noindent \textbf{374($n$).} $UT(\aleph_{0},\,n,\,\aleph_{0})$ for $n\in\omega -\{0,1\}$: The union of a denumerable set of $n$-element sets is denumerable.

\textbf{IV. Definitions of Finite Including Forms Involving Dedekind Finite}

\noindent [\textbf{1 AH}] $E(V,\,VI)$: For all infinite cardinals $m,\,2m=m$ implies $m^{2}=m$. Howard/Yorke\index{Author}{Howard/Yorke} [\textbf{1989}] and note 94.

\noindent [\textbf{1 AI}] $E(VI,\,VII)$: For all infinite cardinals $m,\,m^{2}=m$ implies there is an ordinal $\alpha$ such that $m=\aleph_{\alpha}$. Howard/Yorke [\textbf{1989}] and note 94.

\noindent [\textbf{3 E}] $E(IV,\,V)$: For all infinite cardinals $m,\,m+1=m$ implies $2m=m$. Howard/Yorke [\textbf{1989}] and note 94.

\noindent [\textbf{3 F}] $E(III,\,V)$: For every set $X$, if $\mathcal{P}(X)$ is Dedekind infinite then $2|X|=|X|$. Howard/Yorke [\textbf{1989}] and note~94.

\noindent [\textbf{3 G}] $E(II,\,V)$: For every set $X$, if there is a non-empty family of subsets of $X$ which is linearly ordered by $\subseteq$ and has no maximal element then $2|X|=|X|$. Howard/Yorke [\textbf{1989}] and note 94.

\noindent [\textbf{3 H}] $E(Ia,\,V)$: For every set $X$, if $2|X|>|X|$ then $X$ is amorphous. Howard Yorke [\textbf{1989}], notes 85(H) and 94.

\noindent \textbf{9.} Finite $\Leftrightarrow$ Dedekind finite: $W_{\aleph_{0}}$ (Jech \cite{bib:1973b}): $E(I,\,IV)$ (Howard/Yorke\index{Author}{Howard/Yorke} [\textbf{1989}]): Every Dedekind finite set is finite. Dedekind\index{Author}{Dedekind} [\textbf{1888}] and note 94.

\noindent [\textbf{9 A}] For all infinite $x$, if $x$ is Dedekind finite, then $\mathcal{P}(x)$ is is Dedekind finite). Notes 8 and 94.

\noindent [\textbf{9 B}] $E(II,\,IV)$ (Howard/Yorke [\textbf{1989}]): $(\forall x)$($x$ is $T$-finite if and only if $x$ is Dedekind finite). Howard/Yorke [\textbf{1989}] and note 94.

\noindent [\textbf{9 C}] The image of a Dedekind finite set is Dedekind finite. Jech\index{Author}{Jech} \cite{bib:1973b} p 161 prob 11.2, notes 11 and 94.

\noindent [\textbf{9 D}] The union of a Dedekind finite family of finite sets is Dedekind finite. Jech \cite{bib:1973b} p 161 prob 11.4, notes 12 and 94.

\noindent [\textbf{9 G}] Every infinite set has a denumerable subset. Cantor\index{Author}{Cantor} [\textbf{1895}] and G. Moore\index{Author}{G. Moore} [\textbf{1982}] p 9.

\noindent [\textbf{9 H}] If $A$ is Dedekind finite and $B$ is Dedekind infinite, then $A\precsim B$. G. Moore [\textbf{1982}] p 28 and note 94.

\noindent [\textbf{9 I}] If $\emptyset \in K$ and for all $A\in K$ and for all $b\not\in A,\,A\cup\{b\}\in K$, then every Dedekind finite set is in $K$. G. Moore [\textbf{1982}] p 28 and note 94.

\noindent [\textbf{9 J}] The union of a Dedekind infinite family of non-empty sets is Dedekind infinite. G. Moore [\textbf{1982}] p 130 and note 94.

\noindent [\textbf{9 R}] $E(III,\,IV)$: Every Dedekind finite set has a Dedekind finite power set. Howard/Yorke [\textbf{1989}] and note 94.

\noindent [\textbf{9 S}] $E(Ia,\, IV)$: Every Dedekind finite set is amorphous. Howard/Yorke [\textbf{1989}], notes 57 and 94.

\noindent \textbf{13.} Every Dedekind finite subset of $\mathbb{R}$ is finite. G. Moore [\textbf{1982}] and note 94.

\noindent [\textbf{13 A}] $\mathbb{R}$ has no dense Dedekind finite subset. Brunner\index{Author}{Brunner} [\textbf{1982d}] and note 94.

\noindent \textbf{57.} If $x$ and $y$ are Dedekind finite sets then either $|x|\leq|y|$ or $|y|\leq|x|$. Mathias\index{Author}{Mathias} [\textbf{1979}] p 125 and note 94.

\noindent \textbf{64.} $E(I,\,Ia)$ (Howard/Yorke [\textbf{1989}]): There are no amorphous sets. Note 94 and Levy\index{Author}{Levy} [\textbf{1958}], notes 57 and 94.

\noindent \textbf{82.} $E(I,\,III)$ (Howard/Yorke [\textbf{1989}]): If $X$ is infinite then $\mathcal{P}(X)$ is Dedekind infinite, ($X$ is finite $\Leftrightarrow \mathcal{P}(X)$ is Dedekind finite.) Levy [\textbf{1958}] and note 94.

\noindent [\textbf{82 A}]. Every infinite set can be mapped onto $\omega$. Howard/Yorke [\textbf{1989}], Monro\index{Author}{Monro} [\textbf{1972}].

\noindent [\textbf{82 B}] $E(Ia,\,III)$: For every set $X$, if $\mathcal{P}(X)$ is Dedekind finite then $X$ is amorphous. Howard/Yorke [\textbf{1989}], notes 57 and 94.

\noindent \textbf{83.} $E(I,\,II)$ (Howard/Yorke [\textbf{1989}]): $T$-finite is equivalent to finite. Tarski\index{Author}{Tarski} [\textbf{1924}] and note 94.

\noindent [\textbf{83 A}] $E(Ia,\,II)$: Every $T$-finite set is amorphous. Howard/Yorke [\textbf{1989}], notes 57 and 94.

\noindent \textbf{84.} $E(II,\,III)$ (Howard/Yorke [\textbf{1989}]): $(\forall x)$ ($x$ is $T$-finite if and only if $\mathcal{P}(x)$ is Dedekind finite). Howard/Yorke\index{Author}{Howard/Yorke} [\textbf{1989}] and note 94.

\noindent \textbf{98.} The set of all finite subsets of a Dedekind finite set is Dedekind finite. Jech\index{Author}{Jech} \cite{bib:1973b} p 161 prob 11.5 and note~94.

\noindent \textbf{102.} For all Dedekind finite cardinals $p$ and $q$, if $p^{2}=q^{2}$ then $p=q$. Jech \cite{bib:1973b} p 162 prob 11.12 and note 94.

\noindent \textbf{163.} Every non-well-orderable set has an infinite, Dedekind finite subset. Brunner\index{Author}{Brunner} [\textbf{1982a}] and note 94.

\noindent \textbf{164.} Every non-well-orderable set has an infinite subset with a Dedekind finite power set. Brunner [\textbf{1982a}], notes 42 and 94.

\noindent \textbf{185.} Every linearly ordered Dedekind finite set is finite. Hickman\index{Author}{Hickman} [\textbf{1979a}] and note 94.

\noindent [\textbf{185 A}] If $X$ is a linearly ordered set and there is neither an increasing nor a decreasing injection $f:\omega \rightarrow X$, then $X$ is finite. Hickman [\textbf{1979a}].

\noindent \textbf{198.} For every set $S$, if the only linearly orderable subsets of $S$ are the finite subsets of $S$, then either $S$ is finite or $S$ has an amorphous subset. Hickman\index{Author}{Hickman} [\textbf{1978b}] and note 57.

\noindent \textbf{199($n$).} (For $n\in\omega-\{0,1\})$ If all $\varSigma_{n}^{1}$, Dedekind finite subsets of $^{\omega}\omega$ are finite, then all $\varPi_{n}^{1}$ Dedekind finite subsets of $^{\omega}\omega$ are finite. Kanovei\index{Author}{Kanovei} [\textbf{1978}], notes 57 and 61.

\noindent \textbf{215.} If $(\forall y\subseteq X)$ ($y$ can be linearly ordered implies $y$ is finite), then $X$ is finite. Pincus\index{Author}{Pincus} [\textbf{1972a}].

\noindent \textbf{276.} $E(V^{\prime\prime},\,III)$: For every set $A,\,\mathcal{P}(A)$ is Dedekind finite if and only if $A=\emptyset$ or $2|\mathcal{P}(A)|>|\mathcal{P}(A)|$. Howard/Spi\u{s}iak\index{Author}{Howard/Spisiak@Howard/Spi\u{s}iak} [\textbf{1994}] and note 94.

\noindent \textbf{277.} $E(D,\,VII)$: Every non-well-orderable cardinal is decomposable. Monro\index{Author}{Monro} [\textbf{1973a}] and note 94.

\textbf{V. Ordering Cardinals}

\noindent [\textbf{$\boldsymbol{0}$ AE}] The number of well ordered subsets of a set: For all $x,\,|x|<|s(x)|$ where $s(x)=\{y\subseteq x:y\ \text{can be well ordered}\ \}$. Tarski\index{Author}{Tarski} [\textbf{1939}].

\noindent [\textbf{$\boldsymbol{0}$ AF}] The number of well orderings of a set: For all $x,\,|x|<|w(x)|$ where $w(x)$ is the set of all well orderings of subsets of $x$. Truss\index{Author}{Truss} [\textbf{1973d}].

\noindent [\textbf{$\boldsymbol{0}$ AG}] For any $n\in \omega,\,n\neq 0,\,|s(x)|\not\leq|x^{n}|$ and $|x^{n}|<|w(x)|$ if $|x|$ exceeds some small finite value depending on $n$. ($s(x)$ is the set of well orderable subsets of $x$ and $w(x)$ is the set of well orderings of subsets of $x$.) Truss [\textbf{1973d}]

\noindent [\textbf{1 Y}] Refined Dual Cantor-Bernstein Theorem:\index{Subject}{Cantor-Bernstein theorem} If $f:X\rightarrow Y$ and $g:Y\rightarrow X$ are onto then there is a bijection $h:X\rightarrow Y$ such that $h\subseteq f\cup g^{-1}$. Banaschewski/Moore\index{Author}{Banaschewski/Moore} [\textbf{1990}].

\noindent [\textbf{1 AK}] There is an ordinal $\beta$ such that for all cardinals $x,\,\alpha(x+\aleph(x),\,x\cdot\aleph(x))< \beta$. Truss [\textbf{1973b}] and note 98.

\noindent [\textbf{1 AL}] There is an ordinal $\beta$ such that for all cardinals $x,\,\alpha(x,\,x^{2})<\beta$. Truss [\textbf{1973b}] and note 98.

\noindent [\textbf{1 AM}] There is an ordinal $\beta$ such that for no cardinal $x$ is there a strictly increasing sequence of cardinals $\{y_{\gamma}:\gamma<\beta\}$ such that $y_{\gamma}^{2}=x$. Truss\index{Author}{Truss} [\textbf{1973b}] and note 98.

\noindent [\textbf{3 C}] Definability of cardinal addition as the least upper bound: ($\forall$ cardinals $x,\,y$ and $z$) ($x+y=z$ iff ($x\leq z$ and $y\leq z$ and $(\forall u)(x,\,y\leq u\rightarrow z\leq u)$)). H\"{a}ussler\index{Author}{Haussler@H\"{a}ussler} [\textbf{1983}] and Tarski\index{Author}{Tarski} [\textbf{1949a}].

\noindent [\textbf{3 D}] Every pair of cardinals has a least upper bound in the usual cardinal ordering and for all cardinals $\kappa$ and $\gamma$, if $n\kappa \leq\gamma$ for all $n\in w$, then $\aleph_{0}\kappa\leq\gamma$. Hickman\index{Author}{Hickman} [\textbf{1979b}].

\noindent \textbf{7.} There is no infinite decreasing sequence of cardinals. Schoenflies\index{Author}{Schoenflies} [\textbf{1908}].

\noindent [\textbf{9 E}] No Dedekind finite set can be mapped onto $\aleph_{0}$. (See form 160.) Monro\index{Author}{Monro} [\textbf{1975}] and note 94.

\noindent [\textbf{40 B}] For all ordinals $\alpha$, the $\aleph_{\alpha}$ partition principle holds: For every ordinal $\alpha$ and every cardinal $\kappa$, if $\aleph_{\alpha}\leq^{\ast}\kappa$, then $\aleph_{\alpha}\leq\kappa$. Pelc\index{Author}{Pelc} [\textbf{1978}], Banaschewski/Moore\index{Author}{Banaschewski/Moore} [\textbf{1990}], and note 69.

\noindent \textbf{59($\leq$).} If $(A,\,\leq)$ is a partial ordering that is not a well ordering, then there is no set $B$ such that $(B,\,\leq)$ (the usual injective cardinal ordering on $B$) is isomorphic to $(A,\,\leq)$. Mathias\index{Author}{Mathias} [\textbf{1979}] p 120.

\noindent \textbf{100.} Weak Partition Principle:\index{Subject}{partition principle} For all sets $x$ and $y$, if $x\precsim^{\ast} y$, then it is not the case that $y\prec x$. Lindenbaum/Tarski\index{Author}{Lindenbaum/Tarski} [\textbf{1926}] and note 69.

\noindent [\textbf{100 A}] For every set $A$ and function $f$ it is false that $A\prec f^{\prime\prime}A$. G. Moore\index{Author}{G. Moore} [\textbf{1982}], p. 327.

\noindent [\textbf{100 B}] For all sets $x$ and $y$, if $x\precsim^{\ast}y$ and $y\precsim x$, then $x\approx y$. Higasikawa\index{Author}{Higasikawa} [\textbf{1995}] and note 69.

\noindent \textbf{101.} Partition Principle: If $S$ is a partition of $M$, then $s\precsim\ M$. Sierpi\'{n}ski\index{Author}{Sierpinski@Sierpi\'{n}ski} [\textbf{1947}].

\noindent [\textbf{101 A}] For all sets $x$ and $y,\,x\precsim^{\ast}y$ implies $x\precsim y$. Lindenbaum/Tarski [\textbf{1926}] and note 69.

\noindent \textbf{103.} If $(P,\,<)$ is a linear ordering and $|P|>\aleph_{1}$ then some initial segment of $P$ is uncountable. Jech \cite{bib:1973b} p 164 prob 11.21.

\noindent \textbf{105.} There is a partially ordered set $(A,\,\leq)$ such that for no set $B$ is $(B,\,\leq)$ (the ordering on $B$ is the usual injective cardinal ordering) isomorphic to $(A,\,\leq)$. Forti/Honsell\index{Author}{Forti/Honsell} [\textbf{1985}] and Jech\index{Author}{Jech} [\textbf{1966c}].

\noindent \textbf{159.} The regular cardinals are cofinal in the class of ordinals. Blass\index{Author}{Blass} [\textbf{1983b}].

\noindent \textbf{160.} No Dedekind finite set can be mapped onto an aleph. (See [9 E].) Monro [\textbf{1975}] and note 94.

\noindent \textbf{161.} Definability of cardinal addition in terms of $\leq$: There is a first order formula whose only non-logical symbol is $\leq$ (for cardinals) that defines cardinal addition. H\"{a}ussler~[\textbf{1983}].

\noindent \textbf{162.} Non-existence of infinite units: There is no infinite cardinal number $A$ such that $A+A>A$ and for all cardinals $x$ and $y,\,x+y=A\rightarrow x=A$ or $y=A$. H\"{a}ussler~[\textbf{1983}].

\noindent \textbf{168.} Dual Cantor-Bernstein Theorem:\index{Subject}{Cantor-Bernstein theorem} $(\forall x)(\forall y)(|x|\leq^{\ast}|y|\ \mathrm{and}\ |y|\leq^{\ast}|x|\ \mathrm{implies}\ |x|=|y|)$. Banaschewski/Moore\index{Author}{Banaschewski/Moore} [\textbf{1990}] and note 69.

\noindent \textbf{186.} Every pair of cardinal numbers has a least upper bound (in the usual cardinal ordering.) Hickman\index{Author}{Hickman} [\textbf{1979b}].

\noindent \textbf{187.} Every pair of cardinal numbers has a greatest lower bound (in the usual cardinal ordering.) Hickman [\textbf{1979b}].

\noindent [\textbf{192 A}] $(\forall x)(\exists y)$($x\leq^{\ast}y$ and $y$ indexed choice holds). Blass in Mathias\index{Author}{Mathias} [\textbf{1977b}], notes 60 and 69.

\noindent \textbf{204.} For every infinite $X$, there is a function from $X$ onto $2X$. Truss\index{Author}{Truss} [\textbf{1978}].

\noindent \textbf{205.} For all cardinals $m$ and $n$, if $m\leq^{\ast}n$ and $\neg(n\leq^{\ast}m)$ then there is a cardinal $k<n$ such that $m\leq^{\ast}k$. Pelc\index{Author}{Pelc} [\textbf{1978}] and note 69.

\noindent \textbf{208.} For all ordinals $\alpha,\,\aleph_{\alpha+1}\leq 2^{\aleph_{\alpha}}$. Figura\index{Author}{Figura} [\textbf{1977}].

\noindent \textbf{274.} There is a cardinal number $x$ and an $n\in\omega$ such that $\neg(x\ \mathrm{adj}^{n}\,x^{2})$. (The expression ``$x\ \mathrm{adj}^{n}\ ya^{\prime\prime}$ means there are cardinals $z_{0},\,\ldots,\,z_{n}$ such that $z_{0}=x$ and $z_{n}=y$ and for all $i,\,0\leq i<n,\,z_{i}<z_{i+1}$ and if $z_{i}<z\leq z_{i+1}$, then $z=z_{i+1}$.) (Compare with [0 A]). Truss [\textbf{1974b}].

\noindent \textbf{283.} Cardinality of well ordered subsets: For all $n\in \omega$ and for all infinite $x,\,|x^{n}|<|s(x)|$ where $s(x)$ is the set of all well orderable subsets of $x$. Truss [\textbf{1973d}].

\noindent \textbf{347.} Idemmultiple Partition Principle: If $y$ is idemmultiple $(2\times y\approx y)$ and $x\precsim^{\ast} y$, then $x\precsim y$. Higasikawa\index{Author}{Higasikawa} [\textbf{1995}] and note 69.

\textbf{VI. Inaccessible Cardinals}

\noindent [\textbf{$\boldsymbol{0}$ AJ}] $\aleph_{\alpha+1}$ is not an inaccessible cardinal for any ordinal $\alpha$. Drake\index{Author}{Drake} [\textbf{1974}] and note 20.

\noindent \textbf{117.} If $\kappa$ is a measurable cardinal, then $\kappa$ is the $\kappa$th inaccessible cardinal. Keisler/Tarski\index{Author}{Keisler/Tarski} [\textbf{1964}] and note 20.

\noindent \textbf{196($\alpha$).} $\aleph_{\alpha}$ and $\aleph_{\alpha+1}$ are not both measurable. Bull\index{Author}{Bull} [\textbf{1978}] and note 20.

\noindent \textbf{318.} $\aleph_{1}$ is not measurable. Jech\index{Author}{Jech} [\textbf{1968a}] and note 20.

\noindent \textbf{319.} Measurable cardinals are inaccessible. Drake [\textbf{1974}] and note 20.

\noindent \textbf{320.} No successor cardinal, $\aleph_{\alpha+1}$, is measurable. Drake [\textbf{1974}] and note 20.

\noindent \textbf{321.} There does not exist an ordinal $\alpha$ such that $\aleph_{\alpha}$ is weakly compact and $\aleph_{\alpha+1}$ is measurable. Bull [\textbf{1978}] and note 20.

\textbf{VII. Miscellaneous Cardinal Number Forms}

\noindent [\textbf{15 A}] $K(1):(\forall M)$($\exists\alpha$ an ordinal) ($|M|=|A|$ for some $A\subseteq \mathcal{P}(\alpha)$). (See form 81($n$).) Jech \cite{bib:1973b} p 53 prob 4.12 and Kinna/Wagner\index{Author}{Kinna/Wagner} [\textbf{1955}].

\noindent [\textbf{40 C}] For each idemmultiple set $x\,(2\times x\approx x)$, if there exists a surjection $f$ mapping $x$ onto an ordinal $\lambda$ such that for each $\psi<\lambda,\,f^{-1}[\{\psi\}]$ is Dedekind infinite, then $\lambda\precsim\ x$. Higasikawa [\textbf{1995}] and note 94.

\noindent \textbf{81($n$).} (For $n\in \omega)\ K(n)$: For every set $S$ there is an ordinal $\alpha$ and a one to one function $f:S\rightarrow \mathcal{P}^{n}(\alpha)$. ($\mathcal{P}^{0}(X)=X$ and $\mathcal{P}^{n+1}(X)=\mathcal{P}(\mathcal{P}^{n}(X)$). ($K(0)$ is equivalent to form 1 (AC) and $K(1)$ is equivalent to the selection\index{Subject}{Rado's selection lemma} principle (form 15)). Monro\index{Author}{Monro} [\textbf{1972}], [\textbf{1973b}] and note 90.

\noindent \textbf{97.} Cardinal Representatives:\index{Subject}{cardinal representatives}\index{Subject}{countable axiom of choice} For every set $A$ there is a function $c$ with domain $\mathcal{P}(A)$ such that for all $x,\,y\in \mathcal{P}(A)$, (i) $c(x)=c(y)\leftrightarrow x\approx y$ and (ii) $c(x)\approx x$. Jech\index{Author}{Jech} \cite{bib:1973b} p~154.

\noindent [\textbf{97 A}] For every set $\{X_{i}:i\in I\}$ of local cardinal numbers, there is a set $\{x_{i}:i\in I\}$ of sets such that for each $i\in I$ there exists a $y_{i}\in X_{i}$ such that $x_{i}\approx y_{i}$. ($X$ is a \emph{local cardinal number} if for all $x,\,y\in X,\,x\approx y.)$ Higasikawa\index{Author}{Higasikawa} [\textbf{1995}].

\noindent \textbf{175.} Transitivity Condition: For all sets $x$, there is a set $u$ abd a function $f$ such that $u$ is transitive and $f$ is a one to one function from $x$ onto $u$. von Rimscha\index{Author}{Rimscha} [\textbf{1982}].

\noindent \textbf{269.} For every cardinal $m$, there is a set $A$ such that $2^{|A|^{2}}\geq m$ and there is a choice function on the collection of 2-element subsets of $A$. Lampe\index{Author}{Lampe} \cite{bib:1974m}.

\noindent \textbf{296.} Part-$\infty$: Every infinite set is the disjoint union of infinitely many infinite sets. Gonzalez\index{Author}{Gonzalez} [\textbf{1995b}], Pincus\index{Author}{Pincus} [\textbf{1997}].

\section*{CHOICE FORMS}

\textbf{I. Choosing Single Elements}

\noindent \textbf{1.} $C(\infty,\,\infty)$: The Axiom of Choice:\index{Subject}{axiom of choice} Every set of non-empty sets has a choice function.

\noindent \textbf{5.} $C(\aleph_{0},\,\aleph_{0},\,\mathbb{R})$: Every denumerable set of non-empty denumerable subsets of $\mathbb{R}$ has a choice function. Sageev\index{Author}{Sageev} [\textbf{1975}].

\noindent \textbf{8.} $C(\aleph_{0},\,\infty)$: Every denumerable family of non-empty sets has a choice function. Sierpi\'{n}ski\index{Author}{Sierpinski@Sierpi\'{n}ski} [\textbf{1918}].

\noindent [\textbf{8 M}] Form 131+ form 32. Keremedis\index{Author}{Keremedis} [\textbf{1996d}] and note~132.

\noindent \textbf{10.} $C(\aleph_{0},\,<\aleph_{0})$: Every denumerable family of non-empty finite sets has a choice function.

\noindent [\textbf{14 Y}] Form 99+ Form 62. Note 33.

\noindent [\textbf{14 BS}] Every bounded distributive lattice has a prime ideal choice function. Morillon\index{Author}{Morillon} \cite{bib:1988h}, \cite{bib:1985p} and note~71.

\noindent [\textbf{14 BT}] If $(A_{i})_{i\in I}$ is a family of bounded, distributive, $\vee$-compact lattices and if $(J_{i})_{i\in I}$ is a family such that for all $i\in I,\,J_{i}$ is a proper ideal in $A_{i}$ then there is a family $(P_{i})_{i\in I}$ such that for all $i\in I,\,P_{i}$ is a prime principal ideal in $A_{i}$ containing $J_{i}$. Morillon \cite{bib:1985p}, \cite{bib:1988h} and note 71.

\noindent [\textbf{14 BU}] If $(A_{i})_{i\in I}$ is a family of distributive non trivial bounded lattices then there is a family $(f_{i})_{i\in I}$ such that each $f_{i}$ is a prime ideal choice function on $A_{i}$. Morillon \cite{bib:1985p}, \cite{bib:1988h} and note 71.

\noindent [\textbf{14 BW}] For every set $E$ there is an ultrafilter choice function on $E$. Morillon \cite{bib:1985p}, \cite{bib:1988h} and note 71.

\noindent [\textbf{14 BX}] If $(E_{i})_{i\in I}$ is a family of sets and $(F_{i})_{i\in I}$ is a family such that for each $i,\,F_{i}$ is a filter in $\mathcal{P}(E_{i})$, then there is a family $(U_{i})_{i\in I}$ such that for each $i\in I,\,U_{i}$ is an ultrafilter in $\mathcal{P}(E_{i})$ containing $F_{i}$. Morillon\index{Author}{Morillon} \cite{bib:1985p}, \cite{bib:1988h} and note~71.

\noindent \textbf{16.} $C(\aleph_{0},\,\leq 2^{\aleph_{0}})$: Every denumerable collection of non-empty sets each with power $\leq 2^{\aleph_{0}}$ has a choice function. Sierpi\'{n}ski\index{Author}{Sierpinski@Sierpi\'{n}ski} [\textbf{1918}].

\noindent \textbf{24.} $C(\aleph_{0},2^{(2^{\aleph_{0}})})$: Every denumerable collection of non-empty sets each with power $2^{(2^{\aleph_{0}})}$ has a choice function. G. Moore\index{Author}{G. Moore} [\textbf{1982}].

\noindent \textbf{28.} $C(\aleph_{0},2^{\aleph_{0}})$: Every denumerable collection of non-empty sets each with power $2^{\aleph_{0}}$ has a choice function. Sierpi\'{n}ski [\textbf{1918}].

\noindent \textbf{32.} $C(\aleph_{0},\,\leq\aleph_{0})$: Every denumerable set of non-empty countable sets has a choice function. Sierpi\'{n}ski [\textbf{1918}].

\noindent [\textbf{32 A}] $C(\aleph_{0},\,\aleph_{0})$. Every denumerable set of denumerable sets has a choice function.

\noindent \textbf{33($n$).} If $n\in \omega-\{0,1\},\,C(LO,\,n)$: Every linearly ordered set of $n$ element sets has a choice function. Truss\index{Author}{Truss} [\textbf{1973a}].

\noindent \textbf{39.} $C(\aleph_{1},\,\infty)$: Every set $A$ of non-empty sets such that $|A|=\aleph_{1}$ has a choice function. Sierpi\'{n}ski [\textbf{1918}], G. Moore [\textbf{1982}] p. 202.

\noindent \textbf{40.} $C(WO,\,\infty)$: Every well orderable set of non-empty sets has a choice function. G. Moore [\textbf{1982}] p 325.

\noindent \textbf{45($n$).} If $n\in\omega-\{0,1\},\,C(\infty,\,n)$: Every set of $n$-element sets has a choice function. Mostowski\index{Author}{Mostowski} [\textbf{1945}], Truss [\textbf{1973a}].

\noindent \textbf{46 ($K$)}. If $K$ is a finite subset of $\omega-\{0,\,1\},\,C(\infty,\,K)$: For every $n\in K$, every set of $n$-element sets has a choice function. Mostowski [\textbf{1945}], Truss [\textbf{1973a}].

\noindent \textbf{47($n$).} If $n\in \omega-\{0,1\},\,C(WO,\,n)$: Every well ordered collection of $n$-element sets has a choice function. Mostowski [\textbf{1945}], Truss [\textbf{1973a}].

\noindent \textbf{48($K$).} If $K$ is a finite subset of $\omega-\{0,\,1\},\,C(WO,\,K)$: For every $n\in K,\,C(WO,\,n)$. Mostowski [\textbf{1945}], Truss [\textbf{1973a}].

\noindent \textbf{60.} $C(\infty,\,WO)$: Every set of non-empty, well orderable sets has a choice function. Jech\index{Author}{Jech} \cite{bib:1973b}, G. Moore [\textbf{1982}] p 125.

\noindent [\textbf{60 A}] $(\forall\alpha)C(\infty,\,\aleph_{\alpha})$: For every ordinal $\alpha$, every family of sets each with power $\aleph_{\alpha}$ has a choice function.

\noindent \textbf{61.} $(\forall n\in\omega,\,n\geq 2)(C(\infty,\,n))$: For each $n\in \omega,\,n\geq 2$, every set of $n$ element sets has a choice function. Levy\index{Author}{Levy} [\textbf{1962}].

\noindent \textbf{62.} $C(\infty,\,<\aleph_{0})$: Every set of non-empty finite sets has a choice function.

\noindent [\textbf{79 A}] $C(\infty,\,\infty,\,\mathbb{R})$: Every set whose elements are non-empty subsets of $\mathbb{R}$ has a choice function. Schechter\index{Author}{Schechter} [\textbf{1996a}].

\noindent \textbf{80.} $C(\aleph_{0},2)$: Every denumerable set of pairs has a choice function. Mostowski [\textbf{1948}].

\noindent \textbf{85.} $C(\infty,\,\aleph_{0})$: Every family of denumerable sets has a choice function. Jech\index{Author}{Jech} \cite{bib:1973b} p 115 prob 7.13.

\noindent [\textbf{85 A}] $C(\infty,\,\leq\aleph_{0})$: Every family of non-empty\index{Author}{Freyd} countable sets has a choice function.

\noindent \textbf{86($\alpha$).} $C(\aleph_{\alpha},\,\infty)$: If $X$ is a set of non-empty sets such that $|X|=\aleph_{\alpha}$, then $X$ has a choice function. Levy\index{Author}{Levy} [\textbf{1964}].

\noindent \textbf{88.} $C(\infty,\,2)$: Every family of pairs has a choice function.

\noindent \textbf{94.} $C(\aleph_{0},\,\infty,\,\mathbb{R})$: Every denumerable family of non-empty sets of reals has a choice function. Jech \cite{bib:1973b} p 148 prob 10.1.

\noindent [\textbf{111 A}] $C(WO,\,2)$: Every well ordered family of 2-element sets has a choice function.

\noindent [\textbf{118 I}] Every linearly ordered set of linear orders, each of which is non-empty and conditionally complete, has a choice function. Morillon\index{Author}{Morillon} \cite{bib:1988h}, van Douwen\index{Author}{Douwen} [\textbf{1985}], and note 71.

\noindent [\textbf{118 M}] If $(X,\,\leq)$ is a linear order, then the family $\{A(a,\,b):a,\,b\in X\wedge a<b\}$ has a choice function. ($A(a,\,b)$ denotes the set of continuous, order preserving functions from $X$ into $[0,1]$ such that $f(a)=0$ and $f(b)=1$.) Morillon \cite{bib:1988h} and note 71.

\noindent [\textbf{118 N}] If $(X,\,\leq)$ is a linear order then the family $\{A^{\prime}(a,\,b):a,\,b\in X\wedge a<b\}$ has a choice function. ($A^{\prime}(a, b)$ denotes the set of continuous functions from $X$ into $[0,1]$ such that $f(a)=0$ and $f(b)=1.)$ Morillon \cite{bib:1988h} and note 71.

\noindent [\textbf{118 S}] If $(X,\,\leq)$ is a complete linear order then there is a choice function $\ast$ for the family of non-empty open intervals of $X$ such that if $a,\,b,\,c\in X$ and either $(a,\,b)\neq\emptyset$ or $(b,\, c)\neq\emptyset$ then $\ast((a,\,c))\in\{\ast((a,\,b)), \ast((b,\,c)), b,\,b^{+}\}$ (where $b^{+}$ is the successor of $b$ if it exists and $b$ otherwise). Morillon [\textbf{1987}] and note 71.

\noindent [\textbf{118 T}] Every complete linear order has an increasing choice function. Morillon [\textbf{1987}] and note 71.

\noindent \textbf{119.} $C(\aleph_{0},\text{uniformly orderable with order type of the integers})$: Suppose $\{A_{i}: i\in\omega\}$ is a set and there is a function $f$ such that for each $i\in\omega,\,f(i)$ is an ordering of $A_{i}$ of type $\omega^{\ast}+\omega$ (the usual ordering of the integers), then $\{A_{i}:i\in\omega\}$ has a choice function. van Douwen\index{Subject}{van Douwen's choice principle} [\textbf{1985}].

\noindent \textbf{120($K$).} If $K\subseteq\omega-\{0,1\},\,C(LO,\,K)$: Every linearly ordered set of non-empty sets each of whose cardinality is in $K$ has a choice function. Truss\index{Author}{Truss} [\textbf{1973a}].

\noindent \textbf{121.} $C(LO,\,<\aleph_{0})$: Every linearly ordered set of non-empty finite sets has a choice function. Truss [\textbf{1973a}].

\noindent \textbf{122.} $C(WO,\,<\aleph_{0})$: Every well ordered set of non-empty finite sets has a choice function. Truss [\textbf{1973a}].

\noindent \textbf{165.} $C(WO,\,WO)$: Every well ordered family of non-empty, well orderable sets has a choice function.

\noindent \textbf{181.} $C(2^{\aleph_{0}},\,\infty)$: Every set $X$ of non-empty sets such that $|X|=2^{\aleph_{0}}$ has a choice function. Freyd [\textbf{1980}].

\noindent \textbf{194.} $C(\varPi_{2}^{1})$ or $AC(\varPi_{2}^{1})$: If $P\in\omega \times ^{\omega}\omega,\,P$ has domain $\omega$, and $P$ is in $\varPi_{2}^{1}$, then there is a sequence of elements $\langle x_{k}: k\in \omega\rangle$ of $^{\omega}\omega$ with $\langle k,\,x_{k}\rangle\in P$ for all $k\in\omega$. Kanovei\index{Author}{Kanovei} [\textbf{1979}] and note 61.

\noindent \textbf{202.} $C(LO,\,\infty)$: Every linearly ordered family of non-empty sets has a choice function. Truss\index{Author}{Truss} [\textbf{1978}].

\noindent \textbf{203.} $C(\mathrm{disjoint},\subseteq \mathbb{R})$: Every partition of $\mathcal{P}(\omega)$ into non-empty subsets has a choice function. Truss [\textbf{1978}].

\noindent \textbf{212.} $C(2^{\aleph_{0}},\,\subseteq \mathbb{R})$: If $R$ is a relation on $\mathbb{R}$ such that for all $x\in \mathbb{R}$, there is a $y\in \mathbb{R}$ such that $x\,R\,y$, then there is a function $f:\mathbb{R} \rightarrow \mathbb{R}$ such that for all $x\in \mathbb{R},\,x\,R\,f(x)$. Mathias\index{Author}{Mathias} [\textbf{1977a}].

\noindent \textbf{213.} $C(\infty,\,\aleph_{1})$: If $(\forall y\in X)(|y|=\aleph_{1})$ then $X$ has a choice function. Pincus\index{Author}{Pincus} [\textbf{1977a}].

\noindent \textbf{250.} $(\forall n\in\omega\,-\{0,1\})(C(WO,\,n))$: For every natural number $n\geq 2$, every well ordered family of $n$ element sets has a choice function. Truss [\textbf{1975}].

\noindent \textbf{288($n$).} If $n\in\omega\,-\{0,1\},\,C(\aleph_{0},\,n)$: Every denumerable set of $n$-element sets has a choice function. Wi\'{s}niewski~[\textbf{1972}].

\noindent [\textbf{308} $\mathbf{B}(p)$] If $p$ is a prime and if $Y$ is a set of non-empty finite sets such that $(\forall y\in Y)(\mathrm{gcd}(|y|,\,p)=1)$, then $Y$ has a choice function. (This is similar to form 46$(K)$, but in 46$(K)$ the set $K$ is required to be finite.) Howard/Yorke\index{Author}{Howard/Yorke} [\textbf{1987}].

\noindent [\textbf{308} $\mathbf{C}(p)$] If $p$ is a prime and if $Y$ is a set of non-empty finite sets such that $(\forall y\in Y)(|y|\equiv 1\mod\,p)$, then $Y$ has a choice function. (This is similar to form 46$(K)$, but in 46$(K)$ the set $K$ is required to be finite.) Howard/Yorke~[\textbf{1987}].

\noindent [\textbf{374} $\mathbf{A}(n)$] For each $i,\,2\leq i\leq n\in \omega-\{0,1\},\,C(\aleph_{0},\,i)$: Every denumerable set of $i$-element sets has a choice function.

\textbf{II. Dependent Choice}

\noindent [\textbf{$\boldsymbol{0}$ X}] $DC(\boldsymbol{\Sigma}_{2}^{1})$: If $P\subseteq {^{\omega}\omega} \times{^{\omega}\omega}$ and $P$ has domain ${^{\omega}\omega},\,P\in\boldsymbol{\Sigma}_{2}^{1}$, then for any $y\in{^{\omega}\omega}$ there is a sequence $\langle x_{k}: k\in \omega\rangle$ of elements of $^{\omega}\omega$ with $x_{0}=y$ and $\langle x_{k},\,x_{k+1}\rangle\in P$ for all $k\in\omega$. Kanovei [\textbf{1979}] and note 61.

\noindent [\textbf{1 F}] For all ordinals $\alpha,\,DC(\aleph_{\alpha})$: Given a relation $R$ such that for every subset $Y$ of a set $X$ with $|Y|<\aleph_{\alpha}$, there is an $x\in X$ with $YR\,x$, then there is a function $f:\aleph_{\alpha}\rightarrow X$ such that $\forall\beta<\aleph_{\alpha},\, \{f(\gamma):\gamma<\beta\}\,R\,f(\beta)$.) Levy\index{Author}{Levy} [\textbf{1964}] and note 133.

\noindent \textbf{43.} $DC(\omega)$ ($DC$), Principle of Dependent Choices: If $S$ is a relation on a non-empty set $A$ and $(\forall x\in A)(\exists y\in A)(xSy)$ then there is a sequence $a(0), a(1), a(2),\ldots$ of elements of $A$ such that $(\forall n\in\omega) (a(n)\,S\,a(n+1))$. Tarski\index{Author}{Tarski} [\textbf{1948}] p 96, Levy [\textbf{1964}] p 136 and note 54.

\noindent [\textbf{43 S}] If $S$ is a relation on $A$ such that $(\forall x\in A)(\exists y\in A)(x\,S\,y)$ and $x_{0}$ is any element of $\mathrm{A}$, then there is a sequence $a(1),a(2),\ldots$ of elements of $A$ such that $a(n)\,S\,a(n+1)$ for all $n\in \omega$ and such that $a(0)=x_{0}$. Notes 101 and~54.

\noindent \textbf{44.} $DC(\aleph_{1})$: Given a relation $R$ such that for every subset $Y$ of a set $X$ with $|Y|<\aleph_{1}$ there is an $x\in X$ with $Y\,R\,x$, then there is a function $f:\aleph_{1}\rightarrow X$ such that $(\forall\beta<\aleph_{1})(\{f(\gamma):\gamma<b\}\,R\,f(\beta))$. Levy [\textbf{1964}].

\noindent \textbf{87($\alpha$).} $DC(\aleph_{\alpha})$: Given a relation $R$ such that for every subset $Y$ of a set $X$ with $|Y|<\aleph_{\alpha}$, there is an $x\in X$ with $Y\,R\,x$ then there is a function $f:\aleph_{\alpha}\rightarrow X$ such that $(\forall\beta<\aleph_{\alpha})\{f(\gamma):\gamma<\beta\}\,R\,f(\beta)$. Levy\index{Author}{Levy} [\textbf{1964}].

\noindent [\textbf{106 A}] $DMC$: Dependent Multiple Choice: Every tree in which every element has at least one immediate successor has a subtree whose levels are finite and in which every element has at least one immediate successor. Blass\index{Author}{Blass} [\textbf{1979}], Brunner\index{Author}{Brunner} [\textbf{1983c}], Fossy/Morillon\index{Author}{Fossy/Morillon} [\textbf{1995}] and note~21.

\noindent [\textbf{106 B}] Weak $\mathrm{DC}$: If $R$ is a binary relation on a non-empty set $E$ such that $(\forall x\in E)(\exists y\in E)(x\,R\,y)$, then there is a sequence $\langle F_{n}\rangle_{n\in\omega}$, of non-empty finite subsets of $E$ such that ($\forall n\in\omega) (\forall x\in F_{n})(\exists y\in F_{n+1})(x\,R\,y)$. (This statement remains equivalent to form 106 if we replace ``binary relation'' by ``transitive binary relation''.) Fossy/Morillon~[\textbf{1995}].

\noindent \textbf{211.} $DCR$: Dependent choice for relations on $\mathbb{R}$: If $R\subseteq \mathbb{R}\ \times \mathbb{R}$ satisfies $(\forall x\in \mathbb{R})(\exists y\in \mathbb{R})(x\,R\,y)$ then there is a sequence $\langle x(n): n\in \omega\rangle$ of real numbers such that $(\forall n\in\omega)(x(n)\,R\,x(n+1))$. Mathias\index{Author}{Mathias} [\textbf{1977a}].

\textbf{III. Multiple Choice and Kinna-Wagner Principles}

\noindent [\textbf{1 E}] $K(0)$: For every set $S$ there is an ordinal $\alpha$ and a one to one function $f:S\rightarrow\alpha$. (See form 81$(n)$).

\noindent [\textbf{1 BY}] If $x$ is a set of sets of cardinality at least 2, then there exists a function $f$ such that for each $u\in x,\,f(u)$ is a finite, non-empty, proper subset of $u$. (Form [1 BY] implies [62 E] $(KW(\infty,\,<\aleph_{0})$) and 67 (MC) and $62+67\rightarrow \mathrm{AC}$.) Blass [\textbf{1979}].

\noindent [\textbf{1 CC}] $PC(\aleph_{0},\,\mathrm{odd},\,\infty) +MC_{\omega^{+}}$. $PC(\aleph_{0},\,\mathrm{odd},\,\infty)$ is ``Every denumerable family of odd sized (finite) sets has an infinite subfamily with a choice function.'' and $MC_{\omega^{+}}$ is ``For every natural number $n\geq 1$ and for every set of non-empty sets, there is a function $f$ such that for each $u\in x,\,f(u)$ is a non-empty finite subset of $u$ such that $|f(u)|$ and $n$ are relatively prime.'' Note 133.

\noindent [\textbf{8 M}] Form 131+ form 32. Keremedis\index{Author}{Keremedis} [\textbf{1996d}] and note 132.

\noindent [\textbf{8 O}] $DUCCC\ +$ form 131. $DUCCC$ is ``If $X$ is a denumerable family of disjoint topological spaces each having the ccc property, then the disjoint union of the spaces in $X$ has the ccc property.'' ($Y$ has the \emph{ccc property} if every collection of disjoint open sets is countable.) Keremedis [\textbf{1996d}] and note 132.

\noindent [\textbf{8 P}] $DUSS\ +$ form 131. $DUSS$ is ``The disjoint union of a denumerable family of separable topological spaces is separable.'' Keremedis [\textbf{1996d}] and note 132.

\noindent [\textbf{8 Q}] $DUSC\ +$ form 131. $DUSC$ is ``The disjoint union of a denumerable family of second countable topological spaces is second countable.'' Keremedis [\textbf{1996d}] and note~132.

\noindent [\textbf{8 R}] $SCDUL\ +$ form 131. $SCDUL$ is ``The denumerable disjoint union of second countable spaces is Lindel\"{o}f.'' Keremedis [\textbf{1996d}] and note 132.

\noindent [\textbf{14 AT}] Assume $D$ and $E$ are sets and for all $d\in D,\,d$ is finite. Suppose that for every finite $F\subseteq D$, there is an $S$ such that $(\forall f\in F)(f\cap\,S\in E)$. Then there is an $S$ such that $(\forall d\in D)(d\cap\,S\in E)$. van Benthem\index{Author}{Benthem} [\textbf{1975}].

\noindent \textbf{15.} $KW(\infty,\,\infty)$ (KW), The Kinna-Wagner Selection Principle: For every set $M$ there is a function $f$ such that for all $A\in M$, if $|A|>1$ then $\emptyset \neq f(A)\subsetneq A$. (See form $81(n)$). Kinna/Wagner\index{Author}{Kinna/Wagner} [\textbf{1955}].

\noindent [\textbf{62 E}] $KW(\infty,\,<\aleph_{0})$, The Kinna-Wagner Selection Principle for families of finite sets: For every set $M$ of finite sets there is a function $f$ such that for all $A\in M$, if $|A|>1$ then $\emptyset \neq f(A)\subsetneq A$. Note 70.

\noindent \textbf{67.} $MC(\infty,\,\infty)$ (MC), The Axiom of Multiple Choice: For every set $M$ of non-empty sets there is a function $f$ such that ($\forall x\in M)(\emptyset\neq f(x)\subseteq x$ and $f(x)$ is finite).

\noindent \textbf{76.} $MC_{\omega}(\infty,\,\infty)\ (\omega\text{-}\mathrm{MC})$: For every family $X$ of pairwise disjoint non-empty sets, there is a function $f$ such that for each $x\in X,\,\mathrm{f}(\mathrm{x})$ is a non-empty countable subset of $x$. Howard/Keremedis/Rubin/Rubin\index{Author}{Howard/Keremedis/Rubin/Rubin} [\textbf{1997b}].

\noindent [\textbf{106 A}] $DMC$: Dependent Multiple Choice: Every tree in which every element has at least one immediate successor has a subtree whose levels are finite and in which every element has at least one immediate successor. Blass\index{Author}{Blass} [\textbf{1979}], Brunner\index{Author}{Brunner} [\textbf{1983c}], Fossy/Morillon\index{Author}{Fossy/Morillon} [\textbf{1995}] and note 21.

\noindent [\textbf{106 D}] For every compact $\mathrm{T}_{2}$ space $(X,\,T)$ and every family $D=\{D_{i}:i\in\omega\}$ of dense open sets of $X$ there is a regular filter base $\mathcal{F}\subseteq T$ (that is, if $F,\,G\in \mathcal{F}$, then there exists $Q\in \mathcal{F}$ such that $\overline{Q}\subseteq F\cap G$) such that for all $i\in\omega$, for all but finitely many $F\in \mathcal{F},\,\overline{F}\subset D_{i}$. Fossy/Morillon~[\textbf{1995}]

\noindent \textbf{112.} $MC(\infty,\,LO)$: For every family $X$ of non-empty sets each of which can be linearly ordered there is a function $f$ such that for all $y\in X,\,f(y)$ is a non-empty finite subset of $y$. Howard/Rubin\index{Author}{Howard/Rubin} [\textbf{1977}].

\noindent \textbf{126.} $MC(\aleph_{0},\,\infty)$, Countable axiom of multiple choice: For every denumerable set $X$ of non-empty sets there is a function $f$ such that for all $y\in X,\,f(y)$ is a non-empty finite subset of $y$. Brunner [\textbf{1984d}] and Blass [\textbf{1979}].

\noindent [\textbf{126 B}] Form 131+ form 350. Keremedis\index{Author}{Keremedis} [\textbf{1996d}] and note 132.

\noindent [\textbf{126 C}] $FSCB\ +$ form 131. FSCB is ``Every separable first countable space $(X,\,T)$ has a denumerable $\pi$-base.'' Keremedis [\textbf{1996d}] and note 132.

\noindent \textbf{131.} $MC_{\omega}(\aleph_{0},\,\infty)$: For every denumerable family $X$ of pairwise disjoint nonempty sets, there is a function $f$ such that for each $x\in X,\,\mathrm{f}(\mathrm{x})$ is a non-empty countable subset of $x$. Howard/Keremedis/Rubin/Rubin [\textbf{1997b}].

\noindent \textbf{167.} $PKW(\aleph_{0},\,\geq 2,\,\infty)$, Partial Kinna-Wagner Principle: For every denumerable family $F$ such that for all $x\in F,\,|x|\geq 2$, there is an infinite subset $H\subseteq F$ and a function $f$ such that for all $x\in H,\,\emptyset \neq f(x)\subsetneq x$. Brunner [\textbf{1982a}].

\noindent \textbf{178($n,\,N$).} If $n\in \omega,\,n\geq 2$ and $N\subseteq\{1,2,\,\ldots,\,n-1\},\, N\neq\emptyset,\,MC(\infty,\, n,\,N)$: If $X$ is any set of $n$-element sets then there is a function $f$ with domain $X$ such that for all $A\in X,\,f(A)\subseteq A$ and $|f(A)|\in N$. Zuckerman\index{Author}{Zuckerman} [\textbf{1969a}], \cite{bib:1981b}, and note 53.

\noindent \textbf{214.} $Z(\omega)$: For every family $A$ of infinite sets, there is a function $f$ such that for all $y\in A,\, f(y)$ is a non-empty subset of $y$ and $|f(y)|=\aleph_{0}$. Pincus [\textbf{1977a}].

\noindent \textbf{218.} $(\forall n\in \omega-\{0\})MC$($\infty,\,\infty$, relatively prime to $n$): $\forall n\in \omega-\{0\}$, if $X$ is a set of non-empty sets, then there is a function $f$ such that for all $x\in X,\,f(x)$ is a non-empty, finite subset of $x$ and $|f(x)|$ is relatively prime to $n$. Bleicher\index{Author}{Bleicher} [\textbf{1965}].

\noindent \textbf{219.} $(\forall n\in\omega\,-\{0\})MC$($\infty,\,WO$, relatively prime to $n$): For all non-zero $n\in \omega$, if $X$ is a set of non-empty well orderable sets, then there is a function $f$ such that for all $x\in X,\, f(x)$ is a non-empty, finite subset of $x$, and $|f(x)|$ is relatively prime to $n$. Pincus\index{Author}{Pincus} [\textbf{1972a}].

\noindent \textbf{292.} $MC(LO,\,\infty)$: For each linearly ordered family of non-empty sets $X$, there is a function $f$ such that for all $x\in Xf(x)$ is non-empty, finite subset of $x$). Howard/Rubin\index{Author}{Howard/Rubin} $[\textbf{1977}]$.

\noindent \textbf{322.} $KW(WO,\,\infty)$, The Kinna-Wagner Selection Principle for a well ordered family of sets: For every well ordered set $X$ there is a function $f$ such that for all $A\in M$, if $|A|>1$ then $\emptyset \neq f(A)\subsetneq A$. (See form 15). Kinna/Wagner\index{Author}{Kinna/Wagner} [\textbf{1955}], Brunner\index{Author}{Brunner} [\textbf{1982a}].

\noindent \textbf{323.} $KW(\infty,\,WO)$, The Kinna-Wagner Selection Principle for a family of well orderable sets: For every set $X$ of well orderable sets there is a function $f$ such that for all $A\in X$, if $|A|>1$ then $\emptyset \neq f(A)\subsetneq A$. (See form 15.) Howard/A. Rubin/J. Rubin\index{Author}{Howard/A. Rubin/J. Rubin} [\textbf{1997}].

\noindent \textbf{324.} $KW(WO,\,WO)$, The Kinna-Wagner Selection Principle for a well ordered family of well orderable sets: For every well ordered set $X$ of well orderable sets, there is a function $f$ such that for all $A\in X$, if $|A|>1$ then $\emptyset \neq f(A)\subsetneq A$. (See form 15.) Howard/A. Rubin/J. Rubin [\textbf{1997}].

\noindent \textbf{327.} $KW(WO,\,<\aleph_{0})$, The Kinna-Wagner Selection Principle for a well ordered family of finite sets: For every well ordered set $X$ of finite sets there is a function $f$ such that for all $A\in X$, if $|A|>1$ then $\emptyset \neq f(A)\subsetneq A$. (See form 15.) Howard/A. Rubin/J. Rubin [\textbf{1997}].

\noindent \textbf{328.} $MC(WO, \infty)$; For every well ordered set $X$ such that for all $x\in X,\,|x|\geq 1$, there is a function $f$ such that and for every $x\in X,\,f(x)$ is a finite, non-empty subset of $x$. (See form 67.) Howard/A. Rubin/J. Rubin [\textbf{1997}].

\noindent \textbf{329.} $MC(\infty, WO)$: For every set $M$ of well orderable sets such that for all $x\in X,\,|x|\geq 1$, there is a function $f$ such that for every $x\in X,\,f(x)$ is a finite, non-empty subset of $x$. (See form 67.) Howard/A. Rubin/J. Rubin~[\textbf{1997}].

\noindent \textbf{330.} $MC(WO,\,WO)$: For every well ordered set $X$ of well orderable sets such that for all $x\in X,\, |x|\geq 1$, there is a function $f$ such that for every $x\in X,\,f(x)$ is a finite, non-empty subset of $x$. (See form 67.) Howard/A. Rubin/J. Rubin [\textbf{1997}].

\noindent \textbf{333.} $MC(\infty,\,\infty,\,\mathrm{odd})$: For every set $X$ of sets such that for all $x\in X,\,|x|\geq 1$, there is a function $f$ such that for every $x\in X,\,f(x)$ is a finite, non-empty subset of $x$ and $|f(x)|$ is odd. Keremedis\index{Author}{Keremedis} [\textbf{1996a}].

\noindent \textbf{334.} $MC$($\infty,\,\infty$, even): For every set $X$ of sets such that for all $x\in X,\,|x|\geq 2$, there is a function $f$ such that for every $x\in X,\,f(x)$ is a finite, non-empty subset of $x$ and $|f(x)|$ is even. Keremedis\index{Author}{Keremedis} [\textbf{1996a}].

\noindent \textbf{349.} $MC(\infty,\,\aleph_{0})$: For every set $X$ of non-empty denumerable sets there is a function $f$ such that for all $x\in X,\,f(x)$ is a finite, non-empty subset of $x$.

\noindent \textbf{350.} $MC(\aleph_{0},\,\aleph_{0})$: For every denumerable set $X$ of non-empty denumerable sets there is a function $f$ such that for all $x\in X,\,f(x)$ is a finite, non-empty subset of $x$. Keremedis [\textbf{1996d}] and note 132.

\noindent \textbf{355.} $KW(\aleph_{0},\,\infty)$, The Kinna-Wagner Selection Principle for a denumerable family of sets: For every denumerable set $X$ there is a function $f$ such that for all $A\in X$, if $|A|>1$ then $\emptyset \neq f(A)\subsetneq A$.

\noindent \textbf{356.} $KW(\infty,\,\aleph_{0})$, The Kinna-Wagner Selection Principle for a family of denumerable sets: For every set $X$ of denumerable sets there is a function $f$ such that for all $A\in X$, if $|A|>1$ then $\emptyset \neq f(A)\subsetneq A$.

\noindent \textbf{357.} $KW(\aleph_{0},\,\aleph_{0})$, The Kinna-Wagner Selection Principle for a denumerable family of denumerable sets: For every denumerable set $X$ of denumerable sets there is a function $f$ such that for all $A\in X$, if $|A|>1$ then $\emptyset \neq f(A)\subsetneq A$.

\noindent \textbf{358.} $KW(\aleph_{0},\,<\aleph_{0})$, The Kinna-Wagner Selection Principle for a denumerable family of finite sets: For every denumerable set $X$ of finite sets there is a function $f$ such that for all $A\in X$, if $|A|>1$ then $\emptyset \neq f(A)\subsetneq A$.

\noindent \textbf{376.} Restricted Kinna Wagner Principle: For every infinite set $X$ there is an infinite subset $Y$ of $X$ and a function $f$ such that for every $z\subseteq Y$, if $|z|\geq 2$ then $f(z)$ is a non-empty proper subset of $z$. De la Cruz/Di Prisco\index{Author}{De la Cruz/Di Prisco} [\textbf{1996}].

\textbf{IV. Partial Choice $PC$}

\noindent [\textbf{1 CC}] $PC(\aleph_{0},\,\mathrm{odd},\,\infty) +MC_{\omega^{+}}$. $PC(\aleph_{0},\,\mathrm{odd},\,\infty)$ is ``Every denumerable family of odd sized (finite) sets has an infinite subfamily with a choice function.'' and $MC_{\omega^{+}}$ is ``For every natural number $n\geq 1$ and for every set of non-empty sets, there is a function $f$ such that for each $u\in x,\,f(u)$ is a non-empty finite subset of $u$ such that $|f(u)|$ and $n$ are relatively prime.'' Note 133.

\noindent [\textbf{8 A}] $PC(\infty,\,\infty,\,\infty)$: Every infinite family of non-empty sets has an infinite subfamily with a choice function. Brunner\index{Author}{Brunner} [\textbf{1983b}].

\noindent [\textbf{8 B}] $PC(\aleph_{0},\,\infty,\,\infty)$: Every denumerable family of non-empty sets has an infinite subfamily with a choice function. Monro\index{Author}{Monro} [\textbf{1975}].

\noindent [\textbf{8 S}] $PC(WO,\,\infty,\,\infty)$: Every infinite well ordered set of non-empty sets has an infinite subset with a choice function. $([8\ \mathrm{A}]\rightarrow[8\ \mathrm{S}]\rightarrow[8\ \mathrm{B}])$

\noindent [\textbf{9 T}] Restricted Choice: For every infinite set $X$ there is an infinite $Y\subseteq X$ such that the collection of non-empty subsets of $Y$ has a choice function. De la Cruz/Di Prisco~[\textbf{1996}].

\noindent [\textbf{10 E}] $PC(\aleph_{0},\,<\aleph_{0},\,\infty)$: Every denumerable family of finite sets has an infinite subfamily with a choice function. Brunner [\textbf{1983b}].

\noindent [\textbf{10 O}] $PC(WO,\,<\aleph_{0},\,\infty)$: Every infinite well ordered family of finite sets has an infinite subset with a choice function.

\noindent [\textbf{18 A}] $PC(\aleph_{0},2,\,\infty)$: Every denumerable set of pairs has an infinite subset with a choice function.

\noindent [\textbf{18 C}] $PC(WO,\,2,\,\infty)$: Every infinite well ordered set of pairs has an infinite subset with a choice function.

\noindent [\textbf{32 B}] $PC(\aleph_{0},\,\aleph_{0},\,\infty)$: Every denumerable set of denumerable sets has an infinite subset with a choice function.

\noindent [\textbf{32 C}] $PC(WO,\,\aleph_{0},\,\infty)$: Every infinite well ordered set of denumerable sets has an infinite subset with a choice function.

\noindent [\textbf{126 D}] For every denumerable family $A$ of disjoint non-empty sets there is an infinite set $C\subset\bigcup A$ such that for every $a\in A,\, 0\leq|C\cap a|<\omega$. Keremedis\index{Author}{Keremedis} [\textbf{1996d}] and note 132.

\noindent \textbf{132.} $PC(\infty<\aleph_{0},\,\infty)$: Every infinite family of finite sets has an infinite subfamily with a choice function. Blass\index{Author}{Blass} [\textbf{1977a}] and Kleinberg\index{Author}{Kleinberg} [\textbf{1969}].

\noindent \textbf{150.} $PC(\infty,\,\aleph_{0},\,\infty)$: Every infinite set of denumerable sets has an infinite subset with a choice function.

\noindent \textbf{166.} $PC(\infty,\,2,\,\infty)$: Every infinite family of pairs has an infinite subfamily with a choice function. Brunner\index{Author}{Brunner} [\textbf{1982a}].

\noindent \textbf{167.} $PKW(\aleph_{0},\,\geq 2,\,\infty)$, Partial Kinna-Wagner Principle: For every denumerable family $F$ such that for all $x\in F,\,|x|\geq 2$, there is an infinite subset $H\subseteq F$ and a function $f$ such that for all $x\in H,\,\emptyset \neq f(x)\subsetneq x$. Brunner [\textbf{1982a}].

\noindent \textbf{336($n$).} (For $n\in\omega,\,n\geq 2$.) For every infinite set $X$, there is an infinite $Y\subseteq X$ such that the set of all $n$-element subsets of $Y$ has a choice function. Note 111.

\noindent \textbf{342($n$).} (For $n\in\omega,\,n\geq 2$.) $PC(\infty,\,n,\,\infty)$: Every infinite family of $n$-element sets has an infinite subfamily with a choice function. (See form 166.) Note 111.

\noindent \textbf{373($n$).} (For $n\in\omega,\,n\geq 2$.) $PC(\aleph_{0},\,n,\,\infty)$: Every denumerable set of $n$-element sets has an infinite subset with a choice function.

\noindent [\textbf{373 A($n$)}] (For $n\in \omega,\,n\geq 2$.) $PC(WO,\,n,\,\infty)$: Every infinite well ordered set of $n$-element sets has an infinite subset with a choice function.

\noindent \textbf{378.} Restricted Choice for Families of Well Ordered Sets: For every infinite set $X$ there is an infinite subset $Y$ of $X$ such that the family of non-empty well orderable subsets of $Y$ has a choice function. De la Cruz/Di Prisco [\textbf{1996}].

\noindent [\textbf{378 A}] Restricted Choice for Families of Finite Sets: For every infinite set $X$ there is an infinite subset $Y$ of $X$ such that the family of non-empty finite subsets of $Y$ has a choice function. De la Cruz/Di Prisco\index{Author}{De la Cruz/Di Prisco}  [\textbf{1996}].

\noindent \textbf{379.} $PKW(\infty,\,\infty,\,\infty)$: For every infinite family $X$ of sets each of which has at least two elements, there is an infinite subfamily $Y$ of $X$ and a function $f$ such that for all $y\in Y,\,f(y)$ is a non-empty proper subset of $y$. De la Cruz/Di Prisco [\textbf{1996}].

\noindent \textbf{380.} $PC(\infty,\,WO,\,\infty)$: For every infinite family of non-empty well orderable sets, there is an infinite subfamily $Y$ of $X$ which has a choice function. De la Cruz/Di Prisco~[\textbf{1996}].

\textbf{V. Conditional Choice Including Variations of Rado's Lemma}

\noindent [\textbf{14 I}] Consistency Principle:\index{Subject}{consistency principle} For every binary mess $M$ on a set $S$ there is a function $f$ on $S$ which is consistent with $M$. Jech\index{Author}{Jech} \cite{bib:1973b} p 17 and note 109.

\noindent [\textbf{14 O}] Let $\{A(i):i\in I\}$ be a collection of finite sets and $R$ a symmetric binary relation on $\bigcup_{i\in I}A(i)$. Suppose that for every finite $W\subseteq I$, there is an $R$-consistent choice function for $\{A(i):i\in W\}$. Then there is an $R$-consistent choice function for $\{A(i):i\in I\}$. {\L}o\'{s}/Ryll-Nardzewski\index{Author}{Los/Ryll-Nardzewski@{\L}o\'{s}/Ryll-Nardzewski}
 [\textbf{1951}] and note 109.

\noindent [\textbf{14 $\mathbf{P}(n)$}] (For $n\geq 3$) [14 O] restricted to $\{A(i):i\in I\}$ such that for all $i\in I,\,|A(i)|\leq n$. Cowen\index{Author}{Cowen} [\textbf{1977b}] and note 109. (For $n=2$ see form 141.)

\noindent [\textbf{14 Y}] Form $99$ + form 62. Note 33.

\noindent [\textbf{14 AG}] Topological Rado\index{Author}{Rado} Lemma:\index{Subject}{Rado's selection lemma} Let $\{T_{i}:i\in I\}$ be a family of compact Hausdorff spaces and suppose for each finite $F\subseteq I$, there is a set $\Phi(F)$ of partial choice functions on $I$ whose domains include $F$ and such that (a) $F_{1}\subseteq F_{2}$ implies $\Phi(F_{2})\subseteq\Phi(F_{1})$ and (b) the set $\{\phi|F:\phi\,\in\Phi(F)\}$ is closed in $\prod\{T_{i}:i\in F\}$ in the product topology. Then there is a choice function $f$ on $I$ satisfying for any finite $F\subseteq I,\,\exists\phi\in\Phi(F)$ such that $f$ and $\phi$ coincide on $F$. Rav\index{Author}{Rav} [\textbf{1977}].

\noindent [\textbf{14 AH}] Principle of Consistent Choices: Let $\{T_{i}:i\in I\}$ be a family of compact Hausdorf spaces and suppose $R$ is a symmetric binary relation defined for all pairs of elements belonging to distinct $T_{i}$ and such that for all $i\neq j$ in $I$, the set $E(i,j)=\{(x,\,y):x\,R\,y\ \mathrm{and}\ x\in T_{i}\ \mathrm{and}\ y\in T_{j}\}$ is closed in $T_{i}\times T_{j}$. Furthermore suppose that for every finite subset $F\subseteq I$, with at least two elements, $\exists\phi\in\prod\{T_{i}:i\in F\}$ such that $\phi(i)\,R\,\phi(j)$ for all $i\neq j$ in $F$. Then $\exists f\in\prod\{T_{i}:i\in I\}$ such that $f(i)\,R\,f(j)$ for all $i\neq j$ in $I$. Rav [\textbf{1977}].

\noindent [\textbf{14 AI}] Engelers Lemma:\index{Subject}{Engelers lemma} Let $\mathcal{E}$ be a family of functions whose domains are subsets of a set $S$ with range $\subseteq\{0,1\}$ and such that
\begin{enumerate}
\item[(a)] For every finite $F\subseteq S,\,\exists\phi\in \mathcal{E}$ with domain F.
\item[(b)] The restriction of any function in $\mathcal{E}$ to any subset of its domain is in $\mathcal{E}$.
\item[(c)] If $\phi$ is any function with domain $\phi \subseteq S$ such that the restriction of $\phi$ to any finite subset of its domain is in $\mathcal{E}$ then $\phi \in \mathcal{E}$.
\end{enumerate}
Then $\exists f\in \mathcal{E}$ whose domain is $S$. Rav [\textbf{1977}].

\noindent [\textbf{14 AJ}] Robinson's Valuation Lemma:\index{Subject}{Robinson's valuation lemma} Let $\{\phi_{x}:x\in X\}$ be a family of partial valuations on a set $S$. ($\phi$ is a partial valuation on $S$ if the domain of $\phi\subseteq S$ and the range of $\phi \subseteq\{0,1\}.)$ Suppose $\mathfrak{B}$ is a filter base on $X$ such that for every finite $F\subseteq S$ and for every $B\in \mathfrak{B}$ there is $x\in B$ such that $F\subseteq\ \mathrm{dom}\ \phi_{x}$. Then there is a total valuation $f$ on $S$ (that is, $\mathrm{dom}\,f=S$) such that for all finite $F\subseteq S$ and all $B\in \mathfrak{B},\,\exists x\in B$ such that $F\subseteq\ \mathrm{dom}\,\phi_{x}$ and $f|F=\phi_{x}|F$. ($\mathfrak{B}$ is a filter base means $\forall B_{1},\,B_{2}\in \mathfrak{B},\,\exists B_{3}\in \mathfrak{B}$ such that $B_{3}\subseteq B_{1}\cap B_{2}$.) Rav [\textbf{1977}] and Cowen [\textbf{1973}]. Compare with [14 I].

\noindent [\textbf{14 AK}] Let $X$ be a set and $\mathcal{D}=\{Y_{i}:i\in I\}$ a family of finite pairwise disjoint subsets of $\mathcal{P}(X)$. Suppose for every finite subset $\{a,\,b,\,\ldots,\,k\}$ of $I,\,\exists s\in I$ such that $Y_{s}\subseteq Y_{a}\wedge Y_{b}\wedge\ldots\wedge Y_{k}$ (If $Y$ and $Z$ are finite pairwise disjoint sets then $Y\wedge Z= \{A \cap B:A\in Y\ \mathrm{and}\ B\in Z\ \mathrm{and}\ A \cap B\neq\emptyset\}$.) Then there is a choice function $\mu$ whose domain includes $\mathcal{D}$ such that the family $\{\mu(Y_{i}):i\in I\}$ (where $\mu(Y_{i})\in Y_{i})$ of subsets of $X$ is closed under finite intersections. Rav\index{Author}{Rav} [\textbf{1977}].

\noindent [\textbf{14 AP}] Suppose to each finite subset $F$ of a set I there corresponds a set $\Phi(F)$ of functions whose domains are subsets of $I$ including $F$ and such that (a) $F_{1}\subseteq F_{2}$ implies $\Phi(F_{2})\subseteq\Phi(F_{2})$ and (b) $\forall i\in I,\,\{\phi(i):\phi\in\bigcup\{\Phi(F):F$ finite and $F\subseteq I\}\big\}$ is finite. Then there is a function $f$ with domain $I$ such that for all finite $F\subseteq I,\,\exists\phi\in\Phi(F)$ such that $\phi$ and $f$ coincide on F. Rav\index{Author}{Rav} [\textbf{1977}], Cowen\index{Author}{Cowen} [\textbf{1973}] and note 80.

\noindent [\textbf{14 AQ}] Strong Rado Lemma 1: Let $T$ be a family of finite sets and let $\mathcal{F}$ be a collection of finite subsets of $T$ satisfying
\begin{enumerate}
\item[(1)] $(\forall K\in T)(\exists S\in \mathcal{F})(K\in S)$
\item[(2)] $(\forall S_{1},\,S_{2}\in \mathcal{F})(\exists S_{3}\in \mathcal{F})(S_{1}\cup S_{2}\subseteq S_{3})$
\end{enumerate}
Assume there is a function $\phi$ with domain $\mathcal{F}$ such that for each $S\in \mathcal{F},\,\phi(S)= \phi_{S}$ is a function with domain $S$ and satisfies $(\forall K\in S)(\phi_{S}(K)\in K)$. Then there is a function $f$ with domain $T$ and with the following property: For all $S\in \mathcal{F}$ there is an $S^{\prime}\in \mathcal{F}$ such that $S\subseteq S^{\prime}$ and $f$ and $\phi_{S^{\prime}}$ agree on $S$. Howard\index{Author}{Howard} [\textbf{1993}].

\noindent [\textbf{14 AR}] Patching Lemma:\index{Subject}{patching lemma} Let $L$ be a local system of $S$ (that is, $L\subseteq \mathcal{P}(S)$ and for each finite $E\subseteq S$, there exists $H\in L$ such that $E\subseteq H$). Let $F$ be a set and $n$ a positive integer. Suppose that for each $H\in L$, there is a function $f_{H}:H^{n}\rightarrow F$ and $\{f_{H}(x):H\in L\}$ is finite for each $x\in S^{n}$. Then there is a function $f:S^{n}\rightarrow F$ such that for any finite subset $K\subseteq S^{n}$, there is an $H\in L$ with $K\subseteq H^{n}$ and $f$ and $f_{H}$ agree on K. Hickin/Plotkin\index{Author}{Hickin/Plotkin} [\textbf{1976}] and note 83.

\noindent [\textbf{14 AY}] Let $F$ be a class of partial functions on $I$ of finite character such that (a) $\{f(\nu):f \in F\}$ is finite for each $\nu\in I$ and (b) for each finite $W \subseteq I$, there exists $f\in F$ with $\mathrm{dom}(f)=W$. Then $F$ contains a function with domain $I$. ($F$ is of finite character means $f\in F$ iff for all finite $\omega\subseteq I,\,f|w\in F$.) Cowen [\textbf{1973}].

\noindent [\textbf{14 AZ}] Let $\Phi=\{\phi_{t}\}_{t\in T}$ be a set of partial valuations on a set $S$ (that is, $\mathrm{dom}(\phi_{t})\subseteq S$ and range$(\phi_{t})\subseteq\{0,1\}$) such that for all finite $U\subseteq S$ there is a $t\in T$ for which $U\subseteq\ \mathrm{dom}(\phi_{t})$. Then there is a valuation $\Psi$ with $\mathrm{dom}(\Psi)=S$ and such that for every finite $U\subseteq S$, there is a $t\in T$ with $U\subseteq\, \mathrm{dom}(\phi_{t})$ and $\Psi|U=\phi_{t}|U$. Cowen [\textbf{1973}].

\noindent \textbf{99.} Rado's Selection Lemma:\index{Subject}{Rado's selection lemma} Let $\{K(\lambda):\lambda\in\Lambda\}$ be a family of finite subsets (of $X$) and suppose for each finite $S\subseteq\Lambda$ there is a function $\gamma(S):S\rightarrow X$ such that $(\forall\lambda\in S)(\gamma(S)(\lambda)\in K(\lambda))$. Then there is an $f:\Lambda\rightarrow X$ such that for every finite $S\subseteq\Lambda$ there is a finite $T$ such that $S\subseteq T\subseteq\Lambda$ and such that $f$ and $\gamma(T)$ agree on S. Rado [\textbf{1949}].

\noindent [\textbf{99 A}] Weak Rado Lemma: Let $\Lambda$ be a set. Suppose that $\gamma$ is a function with domain all finite subsets of $\Lambda$ such that for each finite $S\subseteq\Lambda,\,\gamma(S):S\rightarrow\{0,1\}$. Then there is a $\phi$ defined on $\Lambda$ with the property that for every finite $S\subseteq\Lambda$ there is a finite $T\subseteq\Lambda$ such that $S\subseteq T$ and $\gamma(T)$ and $\phi$ agree on $S$. Schechter\index{Author}{Schechter} [\textbf{1996a}] and Howard [\textbf{1993}].

\noindent [\textbf{99 B}] Rado's Lemma restricted to families of two element sets: Let $\{k(\lambda):\lambda\in \Lambda\}$ be a family of 2 element subsets of a set $X$ and suppose that for each finite $S\subseteq\Lambda$, there is a function $\gamma(S)$ from $S$ into $X$ such that $(\forall\lambda\in S)(\gamma(S)(\lambda)\in k(\lambda))$. Then there is a function from $\Lambda$ into $X$ such that for every finite $S\subseteq\Lambda$, there is a finite $T$ with $S\subseteq T\subseteq\Lambda$ such that $f$ and $\gamma(T)$ agree on $S$. Rav [\textbf{1977}] and Howard [\textbf{1993}].

\noindent [\textbf{99 C}] Strong Rado Lemma 2: Let $T$ be a family of finite sets and let $\mathcal{F}$ be a collection of finite subsets of $T$ satisfying
\begin{enumerate}
\item[(1)] $(\forall K\in T)(\exists S\in \mathcal{F})(K\in S)$
\item[(2)] $(\forall S_{1},\,S_{2}\in \mathcal{F})(S_{1}\cup S_{2}\in \mathcal{F})$
\end{enumerate}
Assume there is a function $\phi$ with domain $\mathcal{F}$ such that for each $S\in \mathcal{F},\,\phi(S)= \phi_{S}$ is a function with domain $S$ and satisfies $(\forall K\in S)(\phi_{S}(K)\in K)$. Then there is a function $f$ with domain $T$ and with the following property: For all $S\in \mathcal{F}$ there is an $S^{\prime}\in \mathcal{F}$ such that $S\subseteq S^{\prime}$ and $f$ and $\phi_{S^{\prime}}$ agree on $S$. Howard\index{Author}{Howard} [\textbf{1993}].

\noindent \textbf{107.} M. Hall's Theorem:\index{Subject}{Hall's theorem} Let $\{S(\alpha):\alpha\in A\}$ be a collection of finite subsets (of a set $X$) then if
\begin{equation*}
\tag{$^{\ast}$} \text{for each finite}\ F\subseteq A\ \text{there is an injective choice function on}\ F
\end{equation*}
then there is an injective choice function on $A$. (That is, a 1-1 function $f$ such that $(\forall\alpha\in A)(f(\alpha)\in S(\alpha))$.) (According to a theorem of P. Hall $(\ast)$ is equivalent to $|\bigcup_{\alpha\in F}S(\alpha)|\geq|F|$. P. Hall's theorem does not require the axiom of choice.) M. Hall\index{Author}{Hall} [\textbf{1948}].

\noindent \textbf{141.} $[14\mathrm{P}(n)]$ with $n=2$: Let $\{A(i):i\in I\}$ be a collection of sets such that $\forall i\in I,\, |A(i)|\leq 2$ and suppose $R$ is a symmetric binary relation on $\bigcup_{i\in I}A(i)$ such that for all finite $W\subseteq I$ there is an $R$ consistent choice function for $\{A(i):i\in W\}$. Then there is an $R$ consistent choice function for $\{A(i):i\in I\}$. Cowen\index{Author}{Cowen} [\textbf{1977b}] and note 109.

\noindent \textbf{184.} Existence of a double uniformization: For all $X$ and $Y$, for all $E\subseteq X\times Y$, if there is an infinite cardinal $\kappa$ satisfying:
\begin{enumerate}
\item[(1)] $\forall y_{0}\in Y,\,|\{(x,\,y_{0}):x\in X\}\cap E|=\kappa$ and
\item[(2)] $\forall x_{0}\in X,\,|\{(x_{0},\,y):y\in Y\}\cap E|=\kappa$,
\end{enumerate}
then $\exists U\subseteq E$ such that for all $(x,\,y)\in E,\,\exists!z\in X$ such that $(z,\,y)\in U$ and $\exists!z\in Y$ such that $(x,\,z)\in U$. ($U$ is called a \emph{double uniformization}\index{Subject}{double uniformization} of $E$.) Kaniewski/Rogers\index{Author}{Kaniewski/Rogers} [\textbf{1980}].

\noindent [\textbf{270 B}] $F_{\mathrm{fin}}$: For any family $\mathcal{F}$ of pairwise disjoint, non-empty, finite sets if $R$ is a symmetric binary relation on $\bigcup \mathcal{F}$ which satisfies (i) the set $\{b\in\bigcup \mathcal{F}:\neg a\,R\,b\}$ is finite for every $a\in\bigcup \mathcal{F}$ and (ii) every finite subfamily of $\mathcal{F}$ has an $R$-consistent choice function, then $\mathcal{F}$ has an $R$-consistent choice function. Kolany/Wojtylak\index{Author}{Kolany/Wojtylak} [\textbf{1991}].

\noindent [\textbf{270 $\mathbf{C}(\mathbf{n})$}] $F_{n}^{\mathrm{fin}}$ ($n\geq 3$ a natural number): For any family $\mathcal{F}$ of pairwise disjoint, $n$ element sets if $R$ is a symmetric binary relation on $\bigcup \mathcal{F}$ which satisfies (i) the set $\{b\in\bigcup\,\mathcal{F}:\neg a\,R\,b\}$ is finite for every $a\in\bigcup \mathcal{F}$ and (ii) every finite subfamily of $\mathcal{F}$ has an $R$-consistent choice function, then $\mathcal{F}$ has an $R$-consistent choice function. Kolany\index{Author}{Kolany} [\textbf{1992}].

\textbf{VI. Miscellaneous Choice}

\noindent [\textbf{14 BW}] For every set $E$ there is an ultrafilter choice function on $E$. Morillon\index{Author}{Morillon} \cite{bib:1985p}, \cite{bib:1988h} and note 71.

\noindent \textbf{140.} Let $\Omega$ be the set of all (undirected) infinite cycles of reals (Graphs whose vertices are real numbers, connected, no loops and each vertex adjacent to exactly two others). Then there is a function $f$ on $\Omega$ such that for all $s\in\Omega,\,f(s)$ is a direction along $s$. Truss\index{Author}{Truss} [\textbf{1984}].

\noindent [\textbf{170 A}]\index{Author}{Taskovi\'{c}} There is a function that assigns to each denumerable set $D$ of real numbers a real number not in $D$. Sierpi\'{n}ski\index{Author}{Sierpinski@Sierpi\'{n}ski} [\textbf{1954}].

\noindent \textbf{269.} For every cardinal $m$, there is a set $A$ such that $2^{|A|^{2}}\geq m$ and there is a choice function on the collection of 2-element subsets of $A$. Lampe\index{Author}{Lampe} \cite{bib:1974m}.

\section*{FIXED POINT THEOREMS}

\noindent [\textbf{$\boldsymbol{0}$ J}] Zermelo's Fixed Point Theorem:\index{Subject}{Zermelo's fixed point theorem} For every partially ordered set $(X,\,\leq)$, if every well ordered subset has a least upper bound then every $f:X\rightarrow X$ satisfying $\forall t\in X,\,t\leq f(t)$ has a fixed point. Manka\index{Author}{Manka} [\textbf{1988a}], Zermelo\index{Author}{Zermelo} [\textbf{1908b}] and Abian\index{Author}{Abian} [\textbf{1980}].

\noindent [\textbf{$\boldsymbol{0}$ K}] For all partially ordered sets $(X,\,\leq)$, if there is a sup function $\sigma$ on the well ordered subsets of $X$ then every $f:X\rightarrow X$ satisfying $\forall t\in X,\,t\leq f(t)$ has a fixed point. Manka [\textbf{1988a}] and note 38.

\noindent [\textbf{$\boldsymbol{0}$ L}] Caristi's Fixed Point Theorem:\index{Subject}{Caristi's fixed point theorem} If $(X,\,\rho)$ is a complete metric space and $\phi:X\rightarrow \mathbb{R}$ is bounded above and upper semi-continuous then in the Br{\o}ndsted ordering $(x\leq y\,\mathrm{iff}\, p(x,\,y)\leq\phi(y)-\phi(x))$ every $f:X\rightarrow X$ satisfying $\forall t\in X,\,t\leq f(t)$ has a fixed point. Caristi\index{Author}{Caristi} [\textbf{1976}], Manka [\textbf{1988a}], and note 38.

\noindent [\textbf{$\boldsymbol{0}$ AI}] Turinici's Fixed Point Theorem\index{Subject}{Turinici's fixed point theorem} for Hausdorff Spaces: If $(X, \leq)$ is a directed, partially ordered set and $\tau$ is a Hausdorff topology on $X$ such that

\noindent (a) $(X,\,\leq)$ is upper semi-continuous with respect to $\tau$,

\noindent (b) Every well ordered subset of $(X,\,\leq)$ has a unique limit as a net, then every function $f:X\rightarrow X$ such that $\forall x\in X,\,x\leq f(x)$ has a fixed point. Manka [\textbf{1988b}] and note~38.

\noindent [\textbf{1 V}] For every partially ordered set $(X,\,\leq)$, if every well ordered subset is bounded above then every function $f:X\rightarrow X$ satisfying $\forall t\in X,\,t\leq f(t)$ has a fixed point. Manka [\textbf{1988a}] and Abian [\textbf{1980}].

\noindent [\textbf{1 AJ}] Let $\mathcal{F}$ be a family of mappings of a poset $(X,\,\leq)$ into itself such that $\forall f\in \mathcal{F},\,\forall x\in X,\,f(x)\leq x$. If for some element $e$ of $X$ each chain in $X$ containing $e$ has a lower bound, then the family $\mathcal{F}$ has a common fixed point. Kasahara\index{Author}{Kasahara} [\textbf{1976}].

\noindent [\textbf{1 AS}] Turinici's Fixed Point Theorem: If $(X,\,\leq)$ is a directed, partially ordered set and $\tau$ is a topology on $X$ such that

\noindent (a) $(X,\,\leq)$ is upper semicontinuous with respect to $\tau$,

\noindent (b) Every well ordered subset of $(X,\,\leq)$ is convergent as a net, then every function $f:X\rightarrow X$ such that $\forall x\in X,\,x\leq f(x)$ has a fixed point. Manka [\textbf{1988b}] and note~38.

\noindent [\textbf{1 BC}] Assume $P$ is a partially ordered set in which every non-empty chain has an upper bound and $f:P\rightarrow P$ satisfies $\forall x\in f(P)\cup\{a\in P: a\ \text{is an upper bound for some chain in}\ f(P)\},\, x\leq f(f(x))$. Then $f$ has a fixed apex $u$ (that is, there is some $v\in P$ such that $f(u)=v$ and $f(v)=u$.). Taskovi\'{c} [\textbf{1988}], [\textbf{1992a}].

\noindent [\textbf{1 BD}] Let $P$ be a partially ordered set in which every increasing sequence has an upper bound and $f:P\rightarrow P$ satisfies $x\leq f(f(x))$ for all $x\in f(P)\cup\{x\in P:x$ is an upper bound for some increasing sequence in $f(P)\}$. Then $f$ has a fixed apex $u$ (that is, $\exists v\in P$, such that $f(u)=v$ and $f(v) =u$). (For purposes of [1 BD], ``sequence'' means a function whose domain is an ordinal.) Taskovi\'{c}\index{Author}{Taskovi\'{c}} [\textbf{1992b}].

\noindent \textbf{285.} Let $E$ be a set and $f:E\rightarrow E$, then $f$ has a fixed point if and only if $E$ is not the union of three mutually disjoint sets $E_{1},\,E_{2}$ and $E_{3}$ such that $E_{i}\cap f(E_{i})=\emptyset$ for $i=1,2,3$. Makowski/Wi\'{s}niewski\index{Author}{Makowski/Wisniewski@Makowski/Wi\'{s}niewski} [\textbf{1969}], Abian\index{Author}{Abian} [\textbf{1968}], Baker\index{Author}{Baker} [\textbf{1964}], and Kenyon\index{Author}{Kenyon} [\textbf{1963}].

\noindent \textbf{316.} If a linearly ordered set $(A,\,\leq)$ has the fixed point property then $(A,\,\leq)$ is complete. ($(A,\,\leq)$ has the \emph{fixed point property} if every function $f:A\rightarrow A$ satisfying $(x\leq y\Rightarrow f(x)\leq f(y)$) has a fixed point, and ($(A,\,\leq)$ is \emph{complete} if every subset of $A$ has a least upper bound.) H\"{o}ft/Howard\index{Author}{Hoft/Howard@H\"{o}ft/Howard} [\textbf{1994}].

\section*{GRAPH THEORETIC FORMS}

\noindent [\textbf{1 AT}] For all cardinal numbers $\kappa_{1}$ and $\kappa_{2}$, if a graph is both $\kappa_{1}$ and $\kappa_{2}$ colorable then it is also $\kappa$ colorable for some $\kappa$ such that $\kappa \leq^{\ast}\kappa_{1}$ and $\kappa \leq^{\ast}n_{2}$. ($\leq^{\ast}$ is the surjective cardinal ordering). Galvin/Komj\'{a}th\index{Author}{Galvin/Komjath@Galvin/Komj\'{a}th} [\textbf{1991}] and note 118.

\noindent [\textbf{1 AU}] Every graph has an irreducible, good coloring. Galvin/Komj\'{a}th [\textbf{1991}] and note 118.

\noindent [\textbf{1 AV}] Every graph of the form $G(A,\,H(A))$ has an irreducible, good coloring. ($H$ is the Hartogs aleph function.) Galvin/Komj\'{a}th [\textbf{1991}] and note 118.

\noindent [\textbf{1 AW}] Every graph has a chromatic number. Galvin/Komj\'{a}th [\textbf{1991}] and note 118.

\noindent [\textbf{14} $\mathbf{G}(\mathbf{n})$] $(n\in \omega,\,n\geq 3)\,P_{n}$: For every graph $G$ if every finite subgraph of $G$ is $n$-colorable then $G$ is $n$-colorable. L\"{a}uchli\index{Author}{Lauchli@L\"{a}uchli} [\textbf{1971}] and note 117.

\noindent [\textbf{14 M}] R. Cowen's Generalization of Konig's Lemma:\index{Subject}{K\"{o}nig's lemma} Let $T$ be a collection of locally finite trees such that for any finite set of levels of $T$, there is a consistent set of vertices piercing those levels. Then there is a consistent set of vertices piercing the entire set of levels of T. Cowen\index{Author}{Cowen} [\textbf{1977b}] and note 21.

\noindent [\textbf{14} $\mathbf{N}(\mathbf{n})$] $(n\in\omega,\, n\geq 2)$ [14 M] restricted to families of trees where no tree in the family has a vertex with more than $n$ successors. Cowen [\textbf{1977b}], Howard\index{Author}{Howard} [\textbf{1984a}], and note 21.

\noindent [\textbf{14 BD}] Let $H$ be a hypergraph. If every finite subhypergraph has an independent vertex cover then $H$ has an independent vertex cover. Cowen [\textbf{1990}] and note 117.

\noindent [\textbf{14 BE}] Let $H$ be a hypergraph. If every finite subhypergraph is 2-colorable then $H$ is 2-colorable. Cowen [\textbf{1990}] and note 117.

\noindent [\textbf{14 BF}] Let $H$ be a hypergraph whose edges contain at most three vertices. If every finite subhypergraph of $H$ has an independent vertex cover then $H$ has an independent vertex cover. Cowen\index{Author}{Cowen} [\textbf{1990}] and note 117.

\noindent [\textbf{14 BG}] Let $H$ be a hypergraph whose edges contain at most three vertices.
If every finite subhypergraph of $H$ is 2-colorable then $H$ is 2-colorable. Cowen [\textbf{1990}] and note~117.

\noindent [\textbf{88 B}] $G_{2}$: A graph is 2-colorable if every finite subgraph is 2-colorable. Mycielski\index{Author}{Mycielski} [\textbf{1961}] and note 117.

\section*{LOGIC (FORMS FROM LOGIC)}

\noindent [\textbf{$\boldsymbol{0}$ AA}] $CT_{1}$: Restricted Compactness Theorem for Propositional Logic I: If $\Sigma$
is a set of formulas in a propositional language such that every finite subset of $\Sigma$ is satisfiable and if each propositional variable occurs in at most one formula in $\Sigma$, then $\Sigma$ is satisfiable. Kolany/Wojtylak\index{Author}{Kolany/Wojtylak} [\textbf{1991}].

\noindent [\textbf{10 J}] Let $S$ be a denumerable set of propositional sentences over an infinite set of propositional letters (in a propositional language where conjunction and disjunction are operators on unordered finite sets of sentences). Then $S$ has a model if and only if every finite subset of $S$ has a model. Lolli\index{Author}{Lolli} [\textbf{1977}].

\noindent [\textbf{10 K}] Let $S$ be a consistent denumerable set of sentences in a first order language (where conjunction and disjunction are operators on unordered finite sets of sentences). Then $S$ has a model. Note 19.

\noindent [\textbf{10 L}] Let $S$ be a denumerable set of sentences in a first order language (where
conjunction and disjunction are operators on unordered finite sets of sentences)
such that every finite subset has a model. Then $S$ has a model. Note 19.

\noindent [\textbf{14 H}] Compactness Theorem for First Order Logic: If $\Sigma$ is a set of formulas
in a first order language such that every finite subset of $\Sigma$ has a model then $\Sigma$ has a model. Henkin\index{Author}{Henkin} [\textbf{1954a}] and Jech\index{Author}{Jech} \cite{bib:1973b} p 17.

\noindent [\textbf{14 AW}] The Compactness Theorem for Propositional Logic: If $\Sigma$ is a set of formulas in a propositional language such that every finite subset of $\Sigma$ is satisfiable, then $\Sigma$ is satisfiable. Henkin [\textbf{1954b}].

\noindent [\textbf{14 AX}] Structural Completeness Theorem for Propositional Logic: Assume \emph{At} is a set of propositional variables and $S$ is the set of propositional formulas generated by \emph{At}. Assume $X \subseteq S$ and $a\in S$ satisfy the following: For any function $e : At \rightarrow S$ if $Cn(\{b[e] : b \in X\}) \subseteq\ Cn(\emptyset)$ then $a[e] \in Cn(\emptyset)$. Then $a\in Cn(X)$. ($Cn(Y)$ is the deductive closure of $Y\cup$ all substitutions of axioms for the classical predicate calculus and for $a\in S,\, a[e]$ is $a$ with each propositional variable $p$ replaced by $e(p)$.) Pogorzelski/Prucnal\index{Author}{Pogorzelski/Prucnal} [\textbf{1974}].

\noindent [\textbf{14 CI}] 3-Sat: Restricted Compactness Theorem for Propositional Logic II: If $\Sigma$ is a set of formulas in a propositional language such that every finite subset of $\Sigma$ is satisfiable and if every formula in $\Sigma$ is a disjunction of at most three literals, then $\Sigma$ is satisfiable. (A \emph{literal} is a propositional variable or its negation.) Cowen [\textbf{1990}].

\noindent [\textbf{14 CJ}] The Completeness Theorem for Propositional Logic: Every consistent set of formulas is satisfiable. Henkin [\textbf{1954a}] and [\textbf{1954b}].

\noindent [\textbf{14 CK}] The Completeness Theorem for First Order Logic: Every consistent set of formulas is satisfiable. Henkin [\textbf{1954a}] and [\textbf{1954b}].

\noindent \textbf{42.} L\"{o}wenheim-Skolem Theorem:\index{Subject}{L\"{o}wenheim-Skolem theorem} If a countable family of first order sentences
is satisfiable in a set $M$ then it is satisfiable in a countable subset of $M$. G. Moore\index{Author}{G. Moore} [\textbf{1982}] p~251.

\noindent [\textbf{43 B}] Countable Henkin Principle:\index{Subject}{Henkin principle} Let $\Phi$ be the class of formulas in a language that includes the language of the sentential calculus. Let $I$ be a countable set of inferences, and $\Gamma$ a consistent subset of $\Phi$ such that $\Gamma\cup\{\psi\}$ respects $I$ for each $\psi\in\Phi$. Then for any $\phi\in\Phi$, if $\phi$ is consistent with $\Gamma$ then $\Gamma$ has a finitely consistent extension that contains $\phi$ and decides $I$. Goldblatt\index{Author}{Goldblatt} \cite{bib:1985m} and note.

\noindent [\textbf{43 P}] The ultraproduct of a well founded structure by a countably complete ultrafilter is well founded. Spector\index{Author}{Spector} [\textbf{1980}] and note 54.

\noindent \textbf{253.} Lo\'{s}' Theorem:\index{Subject}{{\L}os' theorem}
 If $M = \langle A,R_{j}\rangle _{j\in J}$ is a relational system, $X$ any set and $\mathcal{F}$ an ultrafilter in $\mathcal{P}(X)$ then $M$ and $M^{X}/\mathcal{F}$ are elementarily equivalent.
Howard\index{Author}{Howard} [\textbf{1975}].

\noindent \textbf{270.} $CT_{\mathrm{fin}}$: The compactness theorem for propositional logic restricted to sets
of formulas in which each variable occurs only in a finite number of formulas.
Kolany/Wojtylak\index{Author}{Kolany/Wojtylak}\index{Author}{Wojtylak} [\textbf{1991}].

\noindent [\textbf{270 A}] $C(\infty, < \aleph_{0})\wedge CT_{2}$ ($CT_{2}$ is the compactness theorem for propositional logic restricted to sets of formulas in which each variable occurs in at most two formulas. See form 271($n$).) Kolany/Wojtylak [\textbf{1991}].

\noindent \textbf{271 ($n$).} If $n\in\omega -\{0,1\}, CT_{n}$: The compactness theorem for propositional logic restricted to sets of formulas in which each variable occurs in at most $n$ formulas. Kolany/Wojtylak [\textbf{1991}].

\noindent \textbf{326.} 2-SAT: Restricted Compactness Theorem for Propositional Logic III: If $\Sigma$ is a set of formulas in a propositional language such that every finite subset of $\Sigma$ is satisfiable and if every formula in $\Sigma$ is a disjunction of at most two literals, then $\Sigma$ is satisfiable. (A \emph{literal} is a propositional variable or its negation.) Wojtylak [\textbf{1995}].

\section*{MAXIMAL PRINCIPLES}

\noindent [\textbf{1 B}] Every commutative ring with unit has a maximal ideal. Hodges\index{Author}{Hodges} [\textbf{1979}].

\noindent [\textbf{1 M}] The ramification Lemma:\index{Subject}{ramification lemma} If $S: X^{\leq\theta}\rightarrow \mathcal{P}(E)$ is a ramification system, then for each $g\in E$ there is an $f$ which is maximal (with respect to inclusion) in $\{h\in X^{\leq\theta} : g\in S(h)\}$. Brunner\index{Author}{Brunner} [\textbf{1989}] and note 22.

\noindent [\textbf{1 AN}] Every lattice $\langle L,\cap,\cup,\subseteq\rangle$ of sets (that is, $L$ is a family of sets and $\cap,\,\cup$
and $\subseteq$ are the usual operations) with a greatest element has a maximal, proper ideal. Bell/Fremlin\index{Author}{Bell/Fremlin} [\textbf{1972}].

\noindent [\textbf{1 BA}] (Depends on the prime $p$) Every group has a maximal $p$-subgroup. Howard/Yorke\index{Author}{Howard/Yorke} [\textbf{1987}] and note 24.

\noindent [\textbf{1 BI}] $(\mathrm{M}T_{0})$: Every topological space has a maximal $T_{0}$ subspace. Schnare\index{Author}{Schnare} [\textbf{1968}] and note 106.

\noindent [\textbf{1 BO}] For every $T_{1}$ topological space $(X,\,T)$ and every base (subbase) $B$ for $X$ there exists a $\subseteq$-maximal cellular family $Q\subseteq B$. Keremedis\index{Author}{Keremedis} [\textbf{1996b}] and note 77.

\noindent [\textbf{1 BT}] For every Boolean ring $(B,+,\cdot)$ and every $H\subseteq B$ which is closed under + there exists a $\subseteq$-maximal ideal $Q\subseteq B$ such that $H\cap Q = \{0\}$. Keremedis\index{Author}{Keremedis} \cite{bib:1996c1}.

\noindent [\textbf{1 CD}] In a commutative ring, every proper ideal can be extended to a maximal ideal. Krull\index{Author}{Krull} [\textbf{1929}]. ([1 CD] implies [1 B]).

\noindent [\textbf{14 A}] Ultrafilter Theorem: Every proper filter over a set $S$ (in $\mathcal{P}(S)$) can be extended to an ultrafilter. Jech\index{Author}{Jech} \cite{bib:1973b} p 17.

\noindent [\textbf{14 E}] Almost Maximal Ideal Theorem: Every bounded, non-trivial, distributive lattice has an almost maximal ideal. (Definitions: If $D$ is a bounded non-trivial distributive lattice, $I(D)$ is the lattice of ideals of $D$. An \emph{almost maximal ideal} in $D$ is an ideal $P$ of $D$ such that $s(P) = P$ where s is the operation on $I(D)$ defined as follows: If $e$ is the unit of $I(D)$, (i) For $a \in I(D)$ we say $x \in I(D)$ is \emph{a-small} if $x \vee b = e$ implies $a \vee b = e$ for all $b \in I(D)$. (ii) For any $a\in I(D), s(a)$ is the join of all a-small elements of $I(D)$). Banaschewski/Harting\index{Author}{Banaschewski/Harting} [\textbf{1985}] and Blass\index{Author}{Blass} [\textbf{1986}].

\noindent [\textbf{14 T}] Every regular commutative ring with unit has a maximal ideal. Banaschewski\index{Author}{Banaschewski} [\textbf{1983}] and note 29.

\noindent [\textbf{14 U}] Every closed proper ideal in a Gelfand algebra is contained in some maximal ideal. Banaschewski [\textbf{1983}] and note 29.

\noindent [\textbf{14 V}] FIP Extension Principle:\index{Subject}{FIP extension principle} Let $\mathcal{F} \subseteq 2^{X}$ have the finite intersection property, then $\mathcal{F}\subseteq U$ for some ultrafilter $U$ on $X$. Schechter\index{Author}{Schechter} [\textbf{1996a}].

\noindent [\textbf{14 BV}] For every set $E$ there is an effectively $E$-adequate ultrafilter. Morillon \cite{bib:1988h}, \cite{bib:1985p} and note 71.

\noindent [\textbf{14 BW}] For every set $E$ there is an ultrafilter choice function on $E$. Morillon\index{Author}{Morillon} \cite{bib:1985p}, \cite{bib:1988h} and note 71.

\noindent [\textbf{14 BX}] If $(E_{i})_{i\in I}$ is a family of sets and $(F_{i})_{i\in I}$ is a family such that for each $i, F_{i}$ is a filter in $\mathcal{P}(E_{i})$ then there is a family $(U_{i})_{i\in I}$ such that for each $i\in I,\,U_{i}$ is an ultrafilter in $\mathcal{P}(E_{i})$ containing $F_{i}$. Morillon \cite{bib:1985p}, \cite{bib:1988h} and note 71.

\noindent [\textbf{14 BY}] For every topological space X, every proper ideal in the normed $\mathbb{R}$-algebra of continuous bounded functions from $X$ into $\mathbb{R}$ is contained in a maximal ideal. Morillon \cite{bib:1988h}.

\noindent [\textbf{14 BZ}] Every proper ideal in a non-trivial normed $\mathbb{R}$-algebra or $\mathbb{C}$-algebra with identity is contained in a maximal ideal. Morillon \cite{bib:1988h}.

\noindent [\textbf{14 CA}] If $(A_{i})_{i\in I}$ is a family of non trivial commutative normed algebras with identity then there is a family $(J_{i})_{i\in I}$ such that for each $i\in I, J_{i}$ is a maximal ideal in $A_{i}$. Morillon \cite{bib:1988h}.

\noindent [\textbf{14 CB}] If $X$ is a topological space and $I$ is a proper ideal in the algebra of all bounded continuous functions from $X$ to $\mathbb{R}$ such that every function in $I$ vanishes at at least one point of $X$ then $I$ is contained in a maximal ideal.
Morillon \cite{bib:1988h}.

\noindent [\textbf{30 G}] Weak Hausdorff Lemma:\index{Subject}{Hausdorff lemma} For every set $X$, every $\subseteq$-chain in $\mathcal{P}(X)$ is contained in a maximal $\subseteq$-chain. Note~130.

\noindent [\textbf{43 L}] $Z_{\omega}$: If every chain in a partially ordered set $P$ is finite, then $P$ contains a maximal element. Tarski\index{Author}{Tarski} [\textbf{1948}], Wolk\index{Author}{Wolk} [\textbf{1983}], and note 54.

\noindent [\textbf{43 M}] If $(X, \rho)$ is a complete metric space and $\phi : X\rightarrow \mathbb{R}$ is bounded above and upper semi-continuous then in the Br{\o}ndsted ordering, $(x < y\ \mathrm{iff}\ \rho(x, y)\leq
\phi(y)-\phi(x))$, there is a maximal element. Brunner\index{Author}{Brunner} [\textbf{1987a}] and note 38.

\noindent [\textbf{50 I}] If $A$ is a subalgebra of a Boolean algebra $B$, then there is a $\subseteq$-maximal proper filter $F$ in $B$ such that $F\cap A =\{1\}$. Bell\index{Author}{Bell} \cite{bib:1988b}

\noindent [\textbf{51 B}] $Z(L,\,W)$: Every non-empty linearly ordered set in which every well
ordered subset has an upper bound has a maximal element. Harper/Rubin\index{Author}{Harper/Rubin}
[\textbf{1976}], Manka\index{Author}{Manka} [\textbf{1988a}], H. Rubin/J. Rubin\index{Author}{H. Rubin/J. Rubin} \cite{bib:1985q} p 37 $M1(L,\,W)$, and note 39.

\noindent [\textbf{62 C}] For every prime $p$, if $\{G_{y}: y\in Y\}$ is a set of finite groups, then the weak direct product $\prod_{y \in Y} G_{y}$ has a maximal $p$-subgroup. Howard/Yorke\index{Author}{Howard/Yorke} [\textbf{1987}] and note 24.

\noindent [\textbf{62 D}] For every prime $p$, if $Y$ is a set of non-empty, finite sets, then the weak direct product $\prod _{y\in Y} S_{y}$ has a maximal $p$-subgroup. ($S_{y}$ is the symmetric group on $y$.) Howard/Yorke [\textbf{1987}] and note 24.

\noindent \textbf{63.} $SPI$: Weak ultrafilter principle: Every infinite set has a non-trivial ultrafilter. Jech\index{Author}{Jech} \cite{bib:1973b} p 172 prob 8.5.

\noindent \textbf{70.} There is a non-trivial ultrafilter on $\omega$. Jech \cite{bib:1973b} prob 5.24.

\noindent [\textbf{70 A}] Every Dedekind infinite set has a non-principal ultrafilter.

\noindent [\textbf{87 A($\alpha$)}] $Z_{\aleph_{\alpha}}$: Let $P$ be a partially ordered set in which every well ordered chain has type $<\aleph_{\alpha}$. If every well ordered chain in $P$ has an upper bound in $P$, then $P$ contains a maximal element. Wolk\index{Author}{Wolk} [\textbf{1983}].

\noindent [\textbf{87 B($\alpha$)}] $Z_{\aleph_{\alpha}}^{\ast}$: Let $P$ be a well-founded partially ordered set (that is, every
chain in $P$ is well ordered.) in which every chain has type $<\aleph_{\alpha}$. If every chain in $P$ has an upper bound in $P$, then $P$ contains a maximal element. Wolk [\textbf{1983}].

\noindent \textbf{89.} Antichain Principle: Every partially ordered set has a maximal antichain. Jech \cite{bib:1973b} p 133 and Kurepa\index{Author}{Kurepa} [\textbf{1953}].

\noindent [\textbf{89 A}] If $\forall y\in X, y = \langle u,R\rangle$ where $R$ is a partial order on $u$ then there is a function $f$ on $X$ such that $\forall\langle u, R\rangle \in X, f(\langle u,R\rangle)$ is a maximal antichain in $u$. Felgner\index{Author}{Felgner} [\textbf{1969}] and Harper/Rubin [\textbf{1976}].

\noindent \textbf{123.} $SPI^{\ast}$: Uniform weak ultrafilter principle: For each family $F$ of infinite sets $\exists f$ such that $\forall x \in F, f(x)$ is a non-principal ultrafilter on $x$. Halpern\index{Author}{Halpern} [\textbf{1972}]

\noindent \textbf{143.} $H(C,TR)$: If $(X,\,R)$ is a connected relation$ (u\neq v\rightarrow u\,R\,v\ \mathrm{or}\ v\,R\,u)$
then $X$ contains a $\subseteq$-maximal transitive subset. Harper/Rubin [\textbf{1976}] and note 39.

\noindent \textbf{201.} Linking Axiom for Boolean Algebras:\index{Subject}{linking axiom} Every Boolean algebra has a maximal linked system. ($L \subseteq B$ is \emph{linked} if $a \wedge b\neq 0$ for all $a$ and $b\in L$.) Schrijver\index{Author}{Schrijver} \cite{bib:1978j}.

\noindent [\textbf{201 A}] Strong Linking Axiom for Boolean Algebras: Each linked system in a
Boolean algebra is contained in some maximal linked system. (If $B$ is a Boolean algebra, $L\subseteq B$ is \emph{linked} if $a\wedge b\neq 0$ for all $a$ and $b\in L$.) Schrijver\index{Author}{Schrijver} \cite{bib:1978j}.

\noindent [\textbf{202 A}] $H(TR\&C,P)$: If $(X,\,R)$ is a transitive, connected relation. Then $X$ has a $\subseteq$-maximal partially ordered subset. Harper/Rubin\index{Author}{Harper/Rubin} [\textbf{1976}] and note 39.

\noindent \textbf{206.} The existence of a non-principal ultrafilter: There exists an infinite set $X$ and a non-principal ultrafilter on $X$. Blass\index{Author}{Blass} [\textbf{1977b}].

\noindent \textbf{225.} Every proper filter on $\omega$ can be extended to an ultrafilter. Rav\index{Author}{Rav} [\textbf{1977}].

\noindent \textbf{244.} Every principal ideal domain has a maximal ideal. Hodges\index{Author}{Hodges} [\textbf{1976a}] and
note 84.

\noindent \textbf{254.} $Z(D, P, l)$: Every directed relation $(P, R)$ in which ramified subsets have
least upper bounds, has a maximal element. Harper/Rubin [\textbf{1976}] and note
39.

\noindent \textbf{255.} $Z(D, R)$: Every directed relation $(P,\,R)$ in which every ramified subset $A$ has an upper bound, has a maximal element. Harper/Rubin [\textbf{1976}] and note 39.

\noindent \textbf{256.} $Z(P,F)$: Every partially ordered set $(X,\,R)$ in which every forest $A$ has an upper bound, has a maximal element. Harper/Rubin [\textbf{1976}] and note 39.

\noindent \textbf{257.} $Z(TR,\,P)$: Every transitive relation $(X,\,R)$ in which every partially ordered subset has an upper bound, has a maximal element. Harper/Rubin [\textbf{1976}] and note 39.

\noindent \textbf{258.} $Z(D,\,L)$: Every directed relation $(X,\,R)$ in which linearly ordered subsets have upper bounds, has a maximal element. Harper/Rubin [\textbf{1976}] and note 39.

\noindent \textbf{259.} $Z(TR\&C, W)$: If $(X,\,R)$ is a transitive and connected relation in which every well ordered subset has an upper bound, then $(X,\,R)$ has a maximal element. Harper/Rubin [\textbf{1976}] and note 39.

\noindent \textbf{260.} $Z(TR\&C,P)$: If $(X,\,R)$ is a transitive and connected relation in which every partially ordered subset has an upper bound, then $(X,\,R)$ has a maximal element. Harper/Rubin [\textbf{1976}] and note 39.

\noindent \textbf{261.} $Z(TR,\,T)$: Every transitive relation $(X,\,R)$ in which every subset which is a tree has an upper bound, has a maximal element. Harper/Rubin [\textbf{1976}] and note 39.

\noindent \textbf{262.} $Z(TR,\,R)$: Every transitive relation $(X, R)$ in which every ramified subset $A$ has an upper bound, has a maximal element. Harper/Rubin [\textbf{1976}] and note 39.

\noindent \textbf{263.} $H(AS\&C,P)$: Every every relation $(X,\,R)$ which is antisymmetric and connected contains a $\subseteq$-maximal partially ordered subset. Harper/Rubin [\textbf{1976}] and note 39.

\noindent \textbf{264.} $H(C, P)$: Every connected relation $(X,\,R)$ contains a $\subseteq$-maximal partially ordered set. Harper/Rubin [\textbf{1976}] and note 39.

\noindent \textbf{265.} $H(A,\,TR)$: Every relation $(X,\,R)$ contains a $\subseteq$-maximal transitive subset. Harper/Rubin [\textbf{1976}] and note~39.

\noindent \textbf{266.} $H(AS,\,P)$: Every antisymmetric relation contains $\subseteq$-maximal partially ordered subset. Harper/Rubin\index{Author}{Harper/Rubin} [\textbf{1976}] and note 39.

\noindent \textbf{278.} In an integral domain $R$, if every ideal is finitely generated then $R$ has a maximal proper ideal. Hodges\index{Author}{Hodges} [\textbf{1974}] and note 45 E.

\noindent \textbf{289.} If $S$ is a set of subsets of a countable set and $S$ is closed under chain unions, then $S$ has a $\subseteq$-maximal element. Smullyan\index{Author}{Smullyan} [\textbf{1994}] and note 102.

\noindent \textbf{303.} If $B$ is a Boolean algebra, $S\subseteq B$ and $S$ is closed under $\wedge$, then there is a $\subseteq$-maximal proper ideal $I$ of $B$ such that $I \cap S =\{0\}$. H. Rubin/J. Rubin\index{Author}{H. Rubin/J. Rubin} \cite{bib:1985q}.

\noindent \textbf{308($p$).} If $p$ is a prime and if $\{G_{y}:y\in Y\}$ is a set of finite groups, then the weak direct product $\prod_{y\in Y}G_{y}$ has a maximal $p$-subgroup. Howard/Yorke [\textbf{1987}] and note 24.

\noindent [\textbf{308 A($p$)}] If $p$ is a prime and if $Y$ is a set of non-empty, finite sets, then the weak direct product $\prod_{y\in Y}S_{y}$ has a maximal $p$-subgroup. $(S_{y}$ is the symmetric group on $y$.) Howard/Yorke\index{Author}{Howard/Yorke} [\textbf{1987}] and note 24.

\noindent [\textbf{317 A}] If $B$ is a complete, well orderable Boolean algebra and $A$ is a Boolean algebra of which $B$ is a subalgebra, then there is an ideal $I$ of $A$ satisfying $I\cap B =\{0\}$ and $I$ is maximal among those ideals $J$ of $A$ satisfying $J\cap B=\{0\}$. Howard\index{Author}{Howard} \cite{bib:1973h}.

\noindent \textbf{345.} Rasiowa-Sikorski Axiom:\index{Subject}{Rasiowa-Sikorski axiom} If $(B,\,\wedge,\,\vee)$ is a Boolean algebra, $a$ is a non-zero element of $B$, and $\{X_{n}:n\in \omega\}$ is a denumerable set of subsets of $B$ then there is a maximal prime filter $F$ of $B$ such that $a\in F$ and for each $n\in\omega$, if $X_{n}\subseteq F$ and $\bigwedge X_{n}$ exists then $\bigwedge X_{n}\in F$. Morillon\index{Author}{Morillon} \cite{bib:1988h}.

\section*{MEASURES (FORMS INVOLVING MEASURES ON SETS)}

\noindent [\textbf{14 AV}] Every Boolean algebra has a 2-valued measure. Pincus\index{Author}{Pincus} \cite{bib:1972c1} and note 147.

\noindent [\textbf{14 BC}] If $m_{0}$ is an additive, real valued measure on a subalgebra $\mathcal{B}_{0}$ of a Boolean algebra $\mathcal{B}$ then there exists a measure $m$ on $\mathcal{B}$ such that $m=m_{0}$ on $\mathcal{B}_{0}$ and the range of $m$ is contained in the closure of the range of $m_{0}$. Tarski\index{Author}{Tarski} \cite{bib:1930} and note 147.

\noindent \textbf{37.} Lebesgue measure is countably additive. Sierpi\'{n}ski\index{Author}{Sierpinski@Sierpi\'{n}ski} [\textbf{1916}].

\noindent [\textbf{52 A}] On every Boolean algebra there is an additive real valued measure. Luxemburg\index{Author}{Luxemburg} \cite{bib:1969d} and note 147.

\noindent [\textbf{52 B}] For every proper filter $\mathcal{F}$ of elements of a Boolean algebra $\mathcal{B}$, there exists an additive, real valued measure $m$ on $\mathcal{B}$ such that $m(x)=1$ for all $x\in \mathcal{F}$. Luxemburg \cite{bib:1969d} and note 147.

\noindent [\textbf{52 C}] For every non-empty set X and every proper ideal $\mathcal{I}$ of in $\mathcal{P}(X)$ there is a measure $m$ on $\mathcal{P}(X)$ such that $m(x)=0$ for all $x\in \mathcal{I}$. Luxemburg \cite{bib:1969d} and note 147.

\noindent [\textbf{52 D}] Let $\mathcal{B}_{0}$ be a subalgebra of a Boolean algebra $\mathcal{B}$ and $m_{0}$ an additive, real valued measure on $\mathcal{B}_{0}$. Then there is a real valued, additive measure $m$ on $\mathcal{B}$ such that $m=m_{0}$ on $\mathcal{B}_{0}$ and the range of $m$ is contained in the closed convex hull of the range of $m_{0}$. Luxemburg\index{Author}{Luxemburg} \cite{bib:1969d} and note 147.

\noindent [\textbf{70 C}] There is a 2-valued measure on $\mathcal{P}(\omega)$. Pincus\index{Author}{Pincus} \cite{bib:1972c1} and note 147.

\noindent \textbf{93.} There is a non-measurable subset of $\mathbb{R}$. Sierpi\'{n}ski\index{Author}{Sierpinski@Sierpi\'{n}ski} [\textbf{1918}].

\noindent \textbf{221.} For all infinite $X$, there is a non-principal measure on $\mathcal{P}(X)$. Pincus/Solovay [\textbf{1977}] and note 147.

\noindent \textbf{222.} There is a non-principal measure on $\mathcal{P}(\omega)$. Pincus/Solovay\index{Author}{Pincus/Solovay} [\textbf{1977}] and note 147.

\noindent \textbf{223.} There is an infinite set $X$ and a non-principal measure on $\mathcal{P}(X)$. Pincus/Solovay [\textbf{1977}] and note 147.

\noindent \textbf{310.} The Measure Extension Theorem: Suppose that $\mathcal{A}_{0}$ is a subring (that is, $a,\,b\in \mathcal{A}_{0}\rightarrow a\vee b\in \mathcal{A}_{0}$ and $a-b\in \mathcal{A}_{0}$) of a Boolean algebra $\mathcal{A}$ and $\mu$ is a measure on $\mathcal{A}_{0}$ (that is, $\mu\,:\,\mathcal{A}\rightarrow[0,\,\infty],\,\mu(a\vee b)=\mu(a)+\mu(b)$ for $a\wedge b=0$, and $\mu(0)=0$.) then there is a measure on $\mathcal{A}$ that extends $\mu$. Wagon \cite{bib:1985u}.

\noindent \textbf{311.} Abelian groups are amenable. ($G$ is \emph{amenable} if there is a finitely additive measure $\mu$ on $\mathcal{P}(G)$ such that $\mu(G)=1$ and $\forall A\subseteq G,\,\forall g\in G,\,\mu(gA)=\mu(A)$.) Wagon~{bib:1985u}.

\noindent \textbf{312.} A subgroup of an amenable group is amenable. ($G$ is \emph{amenable} if there is a finitely additive measure $\mu$ on $\mathcal{P}(G)$ such that $\mu(G)=1$ and $\forall A\subseteq G,\,\forall g\in G,\,\mu(gA)=\mu(A)$.) Wagon \cite{bib:1985u}.

\noindent \textbf{313.} $\mathbb{Z}$ (the set of integers under addition) is amenable. ($G$ is \emph{amenable} if there is a finitely additive measure $\mu$ on $\mathcal{P}(G)$ such that $\mu(G)=1$ and $\forall A\subseteq G,\,\forall g\in G,\,\mu(gA)=\mu(A)$.) Wagon \cite{bib:1985u}.

\noindent [\textbf{370 A}] Measured Boolean Prime Ideal Theorem: If $\mathcal{B}$ is a Boolean algebra and there is a measure on $\mathcal{B}$ then there is an ultrafilter on $\mathcal{B}$. Morillon\index{Author}{Morillon} \cite{bib:1989r}, \cite{bib:1990e} and note~147.

\section*{MISCELLANEOUS FORMS}

\textbf{0.} $0=0$.

\noindent [\textbf{$\boldsymbol{0}$ T}] If a set $S$ is hereditarily countable\index{Subject}{hereditarily countable set} (that is, every $y\in TC(S)$ is countable) then the rank of $S$ is $<\omega_{2}$. Jech\index{Author}{Jech} [\textbf{1982}].

\noindent [\textbf{$\boldsymbol{0}$ AD}] For all ordinals $\beta$, well ordered cardinals $\gamma\geq 1$ and all $n\in \omega$, there is a well ordered cardinal $\kappa$ such that $\kappa \rightarrow(\beta)_{\gamma}^{n}$. Erd\"{o}s/Rado\index{Author}{Erdos/Rado@Erd\"{o}s/Rado} [\textbf{1952}] and note 97.

\noindent [\textbf{1 G}] Form $163+$ form 8. Brunner\index{Author}{Brunner} [\textbf{1982a}] and note 94.

\noindent [\textbf{1 H}] Form $164+$ form 67. Brunner [\textbf{1982a}] and note 94.

\noindent [\textbf{1 L}] Hajrial's Free Set Principle:\index{Subject}{Hajnal's free set principle} For every set mapping $f:X\rightarrow[X]^{<\lambda}$ where $\lambda$ is a well orderable cardinal and $\lambda<|X|$ there is a free set of cardinality $|X|.\, ([X]^{<\lambda}$ is the set of subsets of $X$ of cardinality $<\lambda$. Brunner [\textbf{1989}] and note 22.

\noindent [\textbf{1 M}] The ramification Lemma:\index{Subject}{ramification lemma} If $S:X^{\leq\theta}\rightarrow \mathcal{P}(E)$ is a ramification system, then for each $g\in E$ there is an $f$ which is maximal (with respect to inclusion) in $\{h\in X^{\leq\theta}:g\in S(h)\}$. Brunner\index{Author}{Brunner} [\textbf{1989}] and note 22.

\noindent [\textbf{1 N}] For every set mapping $f:X\rightarrow[X]^{WO}$ (well ordered subsets of $X$) there is a co-well-orderable free subset $H$. ($H$ is \emph{co-well-orderable} if $X\backslash H$ is well orderable). Brunner [\textbf{1989}] and note 22.

\noindent [\textbf{1 W}] Form 133 + form 88. Brunner [\textbf{1982a}].

\noindent [\textbf{1 X}] Form 15 + form 91. Brunner [\textbf{1982a}].

\noindent [\textbf{1 Y}] Refined Dual Cantor-Bernstein Theorem:\index{Subject}{Cantor-Bernstein theorem} If $f:X\rightarrow Y$ and $g:Y\rightarrow X$ are onto then there is a bijection $h:X\rightarrow Y$ such that $h\subseteq f\cup g^{-1}$. Banaschewski/Moore\index{Author}{Banaschewski/Moore} [\textbf{1990}].

\noindent [\textbf{1 AA}] Assume $f:S\rightarrow R$ where $f$ is surjective and $R\subseteq S$. Then every mapping $g$ of $S$ into $R$ has the form $g=f\circ h$ for some $h:S\rightarrow S$. Smith [\textbf{1982}].

\noindent [\textbf{1 AB}] Assume $f:S\rightarrow R$ where $f$ is surjective and $R\subseteq S$. Then the mappings of $S$ onto $R$ are precisely the mappings of the form $f\circ h$ where $h:S\rightarrow S$ and $(\forall y\in R)(\exists x\in S)(h(x)\in f^{-1}(\{y\}))$. Smith\index{Author}{Smith} [\textbf{1982}].

\noindent [\textbf{1 AC}] Assume $f:S\rightarrow R$ where $R\subseteq S$. Then the retractions of $S$ onto $R$ (that is, the mappings $g\,:\,S\rightarrow R$ which are the identity on $R$) are precisely the mappings $f\circ h$ where $h:S\rightarrow S$ and $(\forall y\in R)(h(y)\in f^{-1}(\{y\}))$. Smith [\textbf{1982}].

\noindent [\textbf{1 AE}] There is a set $S$ such that $SVC$ holds with $S$ and such that $\mathcal{P}(S)-\{\emptyset\}$ has a choice function. Blass\index{Author}{Blass} [\textbf{1979}] and note 59. Also see form 191.

\noindent [\textbf{1 AF}] Form 191 + form 8 + form 193. Blass [\textbf{1979}] and note 60.

\noindent [\textbf{1 AG}] Form 192 + form 191. Blass [\textbf{1979}] and note 60.

\noindent \textbf{4.} Every infinite set is the union of some disjoint family of denumerable subsets. (Denumerable means $\cong\aleph_{0}$.) Whitehead\index{Author}{Whitehead} [\textbf{1902}].

\noindent [\textbf{10 C}] Every set mapping $f:X\rightarrow[X]^{WO}$ on a Dedekind finite, infinite set $X$ has arbitrarily large, finite free subsets. Brunner [\textbf{1989}], notes 22 and 94.

\noindent [\textbf{14 I}] Consistency Principle:\index{Subject}{consistency principle} For every binary mess $M$ on a set $S$ there is a function $f$ on $S$ which is consistent with $M$. Jech\index{Author}{Jech} \cite{bib:1973b} p 17 and note 109.

\noindent [\textbf{14 X}] Cowen-Engler Lemma:\index{Subject}{Cowen-Engler lemma} Let $\Lambda$ and $X$ be sets. Let $\Phi$ be a family of functions from subsets of $\Lambda$ into $X$. Assume
\begin{enumerate}
\item[(a)] $\Phi(\lambda)=\{f(\lambda): f\in\Phi\ \mathrm{and}\ \lambda\in \mathrm{dom}\ f\}$ is finite for each $\lambda\in\Lambda$,
\item[(b)] each finite $S\subseteq\Lambda$ is the domain of at least one $f\in\Phi$, and
\item[(c)] $\Phi$ has finite character, (that is, a function $f$ from a subset of $\Lambda$ into $X$ is in $\Phi$ if and only if every finite subset of $f$ is in $\Phi$).
\end{enumerate}
Then at least one $f\in\Phi$ has domain $\Lambda$. Schechter\index{Author}{Schechter} [\textbf{1996a}]. (See 14 I.)

\noindent [\textbf{14 AC}] (Depends on $k,\,n\in\omega$ with $k\geq 2,\,n\geq 2$ and $n+k\geq 5$.) $P(k,\,n)$: If $X$ is a set and $P$ is a property of subsets of $X$ of $n$ character (that is, $\forall y\subseteq X(P(y)$ iff $\forall z\subseteq y(|z|\leq n\rightarrow P(z)))$), then if every finite subset of $X$ can be partitioned into $k$ or fewer $P$-sets (that is, sets $z$ such that $P(z)$) then $X$ can be partitioned into $k$ or fewer $P$-sets. Cowen\index{Author}{Cowen} [\textbf{1982}].

\noindent [\textbf{14 AD}] (Depends on $k\in\omega,\,k\geq 2$) $P(k)$: If $X$ is a set and $P$ is a property of subsets of $X$ of finite character (that is, $\forall y\subseteq X(P(y)\ \mathrm{iff}\ \forall z\subseteq y(z\ \mathrm{finite}\ \rightarrow P(z)))$), then if every finite subset of $X$ can be partitioned into $k$ or fewer $P$-sets (that is, sets $z$ such that $P(z)$), then $X$ can be partitioned into $k$ or fewer $P$-sets. Cowen [\textbf{1982}].

\noindent [\textbf{14 AF}] Disjoint Transversal axiom (2,3): If $X$ is a set, $\mathcal{U}$ is a subset of $\{y\subseteq X:|y|\leq 2\}$ and $\mathcal{V}$ is a subset of $\{y\subseteq X:|y|\leq 3\}$ and if for each two finite subsets $\mathcal{U}_{0}\subseteq \mathcal{U}$ and $\mathcal{V}_{0}\subseteq \mathcal{V}$ there are transversals $Y_{1}$ and $Y_{2}$ for $\mathcal{U}_{0}$ and $\mathcal{V}_{0}$ respectively such that $Y_{1}\cap Y_{2}=\emptyset$, then there are disjoint transversals $X_{1}$ and $X_{2}$ for $\mathcal{U}$ and $\mathcal{V}$ respectively. ($Y\subseteq X$ is a transversal for $\mathcal{U}$ if for all $y\in \mathcal{U},\, y\cap Y\neq\emptyset$.) Schrijver\index{Author}{Schrijver} \cite{bib:1978j}.

\noindent [\textbf{14 AT}] Assume $D$ and $E$ are sets and for all $d\in D,\,d$ is finite. Suppose that for every finite $F\subseteq D$, there is an $S$ such that $(\forall f\in F)(f\cap\,S\in E)$. Then there is an $S$ such that $(\forall d\in D)(d\cap\,S\in E)$. van Benthem\index{Author}{Benthem} [\textbf{1975}].

\noindent [\textbf{14 BB}] Fra\"{\i}ss\'{e}'s Theorem:\index{Subject}{Fra\"{i} ss\'{e}'s theorem} Let $\mathcal{E}$ be inductively ordered (that is, for any $E,\,E^{\prime}\in \mathcal{E},\,\exists E^{\prime\prime}\in \mathcal{E}$ such that $E\cup E^{\prime}\subseteq E^{\prime\prime}$.) For any $E\in \mathcal{E}$, assume $\mathcal{R}_{E}$ is a finite set of $m$-ary relations on $E$. If for any $E,\,E^{\prime}\in \mathcal{E}$ with $E\subseteq E^{\prime}$ we have $\forall R\in \mathcal{R}_{E^{\prime}},\,R/E\in \mathcal{R}_{E}$ then there is an $m$-ary relation $S$ on $\bigcup \mathcal{E}$ such that $S|E\in \mathcal{R}_{E}$ for any $E\in \mathcal{E}$. B\'{e}n\'{e}jam\index{Author}{B\'{e}n\'{e}jam} [\textbf{1970}]. See [14 Y].

\noindent [\textbf{14 CF}] A projective limit of a directed projective system of non-empty sets is non-empty. Morillon\index{Author}{Morillon} \cite{bib:1988h} and note~71.

\noindent \textbf{17.} Ramsey's Theorem\index{Subject}{Ramsey's theorem} I: If $A$ is an infinite set and the family of all 2 element subsets of $A$ is partitioned into 2 sets $X$ and $Y$, then there is an infinite subset $B\subseteq A$ such that all 2 element subsets of $B$ belong to $X$ or all 2 element subsets of $B$ belong to $Y$. (Also, see form 325.) Jech\index{Author}{Jech} \cite{bib:1973b} p 164 prob 11.20 and Ramsey\index{Author}{Ramsey} [\textbf{1929}].

\noindent [\textbf{34 A}] $\exists\alpha<\omega_{2}$ such that there are no hereditarily countable sets of rank $\alpha$. Jech [\textbf{1982}] and note 48.

\noindent [\textbf{43 I}] Form [106 A] + form 10. Blass\index{Author}{Blass} [\textbf{1979}].

\noindent [\textbf{43 P}] The ultraproduct of a well founded structure by a countably complete ultrafilter is well founded. Spector\index{Author}{Spector} [\textbf{1980}] and note 54.

\noindent [\textbf{88 A}] $P(2,2)$: If $X$ is a set and $P$ is a property of subsets of $X$ of 2 character (that is, $\forall y\subseteq X(P(y)\ \mathrm{iff}\,\forall z\subseteq y(|z|\leq 2\rightarrow P(z)))$) then if every finite subset of $X$ can be partitioned into two or fewer $P$-sets (that is, sets $z$ such that $P(z)$) then $X$ can be partitioned into two or fewer $P$-sets. Cowen [\textbf{1982}].

\noindent [\textbf{126 F}] Every denumerable family of sets having the sfip also has a pseudo-intersection. Keremedis\index{Author}{Keremedis} [\textbf{1996d}] and note 132.

\noindent [\textbf{132 A}] Every set mapping $f:X\rightarrow[X]^{WO}$ (well orderable subsets of $X$) on a Dedekind finite, infinite set $X$ has an infinite free subset. Brunner\index{Author}{Brunner} [\textbf{1989}]. notes 22 and 94.

\noindent \textbf{137($k$).} Suppose $k\in\omega -\{0\}$. If $f$ is a 1-1 map from $k\times X$ into $k\times Y$ then there are partitions $X=\bigcup_{i<k}X_{i}$ and $Y=\bigcup_{i<k}Y_{i}$ of $X$ and $Y$ such that $f$ maps $\bigcup_{i<k}(\{i\}\times X_{i})\ \mathrm{onto}\ \bigcup_{i<k}(\{i\}\times Y_{i})$. Truss [\textbf{1984}].

\noindent [\textbf{137 A($k$)}] For all $X$ and $Y\subseteq \mathbb{R}$ form 137($k$) holds. Truss\index{Author}{Truss} [\textbf{1984}]

\noindent \textbf{138($k$).} Suppose $k\in\omega$. If $f$ is a partial map from $k\times Y$ onto $k\times X$ (that is, the domain is a subset of $k\times Y$), then there are partitions $X= \bigcup_{i<k}X_{i}$ and $Y=\bigcup_{i<k}Y_{i}$ of $X$ and $Y$ such that $f$ maps $\bigcup_{i<k}(\{i\}\times Y_{i})$ onto $\bigcup_{i<k}(\{i\}\times X_{i})$. Truss [\textbf{1984}].

\noindent [\textbf{138 A($k$)}] For all standard $X$ and $Y$, form 138$(k)$ holds. Truss [\textbf{1984}]. This is equivalent to 138$(k)$ in $\mathrm{ZF}^{0}$.

\noindent [\textbf{140 A}] (Form 137 with $k=2$): If $f$ is a 1-1 map from $2\times X$ into $2\times Y$ then there are partitions $X=X_{0}\cup X_{1}$ and $Y=Y_{0}\cup Y_{1}$ of $X$ and $Y$ such that $f$ maps $(\{0\}\times X_{0})\cup (\{1\}\times X_{1})$ onto $(\{0\}\times Y_{0})\cup (\{1\}\times Y_{1})$. Truss [\textbf{1984}].

\noindent \textbf{152.} $D_{\aleph_{0}}$: Every non-well-orderable set is the union of a pairwise disjoint, well orderable family of denumerable sets. Brunner/Howard\index{Author}{Brunner/Howard} [\textbf{1992}]. (See note 27 for $D_{\kappa},\,\kappa$ a well ordered cardinal.)

\noindent \textbf{172.} For every infinite set $S$, if $S$ is hereditarily countable (that is, every $y\in TC(S)$ is countable) then $|TC(S)|=\aleph_{0}$. Jech\index{Author}{Jech} [\textbf{1982}].

\noindent \textbf{179($\epsilon$).} Suppose $\epsilon>0$ is an ordinal. $\forall x,\,x\in W(\epsilon)$. David\index{Author}{David} [\textbf{1980}] and note 25.

\noindent \textbf{184.} Existence of a double uniformization: For all $X$ and $Y$, for all $E\subseteq X\times Y$, if there is an infinite cardinal $\kappa$ satisfying:
\begin{enumerate}
\item[(1)] $\forall y_{0}\in Y,\,|\{(x,\,y_{0}):x\in X\}\cap E|=\kappa$ and
\item[(2)] $\forall x_{0}\in X,\,|\{(x_{0},\,y):y\in Y\}\cap E|=\kappa$,
\end{enumerate}
then $\exists U\subseteq E$ such that for all $(x,\,y)\in E,\,\exists!z\in X$ such that $(z,\,y)\in U$ and $\exists!z\in Y$ such that $(x,\,z)\in U$. ($U$ is called a \emph{double uniformization} of $E$.) Kaniewski/Rogers\index{Author}{Kaniewski/Rogers} [\textbf{1980}].

\noindent \textbf{191.} $SVC$: There is a set $S$ such that for every set $a$, there is an ordinal $\alpha$ and a function from $S\times\alpha$ onto $a$. Blass [\textbf{1979}].

\noindent \textbf{192.} $EP$ sets: For every set $A$ there is a projective set $X$ and a function from $X$ onto $A$. Blass\index{Author}{Blass} [\textbf{1979}] and note 60.

\noindent [\textbf{201 B}] Disjoint Transversal axiom (2,2): If $X$ is a set and $\mathcal{U}$ and $\mathcal{V}$ are subsets of $\{y\subseteq X:|y|\leq 2\}$ and if for each two finite subsets $\mathcal{U}_{0}\subseteq \mathcal{U}$ and $\mathcal{V}_{0}\subseteq \mathcal{V}$ there are transversals $Y_{1}$ and $Y_{2}$ for $\mathcal{U}_{0}$ and $\mathcal{V}_{0}$ respectively such that $Y_{1}\cap Y_{2}=\emptyset$, then there are disjoint transversals $X_{1}$ and $X_{2}$ for $\mathcal{U}$ and $\mathcal{V}$ respectively ($Y\subseteq X$ is a \emph{transversal} for $\mathcal{U}$ if for all $y\in CalU,\, y\cap Y\neq\emptyset$.) Schrijver\index{Author}{Schrijver} \cite{bib:1978j}.

\noindent \textbf{230.} $L^{1}=HOD$. Szczepaniak\index{Author}{Szczepaniak} [\textbf{1977}] and note 82.

\noindent \textbf{234.} There is a non-Ramsey set: There is a set $A$ of infinite subsets of $\omega$ such that for every infinite subset $N$ of $\omega,\,N$ has a subset which is in $A$ and a subset which is not in $A$. Prikry\index{Author}{Prikry} [\textbf{1976}] and Mathias\index{Author}{Mathias} [\textbf{1968}].

\noindent \textbf{282.} $\omega{\not\rightarrow} (\omega)^{\omega}$. Erd\"{o}s/Rado\index{Author}{Erdos/Rado@Erd\"{o}s/Rado} [\textbf{1952}], Kleinberg/Seiferas\index{Author}{Kleinberg/Seiferas} [\textbf{1973}] and note 97.

\noindent \textbf{309.} The Banach-Tarski Paradox: There are three finite partitions $\{P_{1},\,\ldots,\,P_{n}\},\,\{Q_{1},\, \ldots,\,Q_{r}\}$ and $\{S_{1},\,\ldots,\,S_{n},\,T_{1},\,\ldots,\,T_{r}\}$ of $B^{3}=\{x\in \mathbb{R}^{3}:|x|\leq 1\}$ such that $P_{i}$ is congruent to $S_{i}$ for $1\leq i\leq n$ and $Q_{i}$ is congruent to $T_{i}$ for $1\leq i\leq r$. Banach/Tarski\index{Author}{Banach/Tarski} [\textbf{1924}].

\noindent \textbf{325.} Ramsey's Theorem\index{Subject}{Ramsey's theorem} II: $\forall n,\,m\in \omega$, if\index{Subject}{Banach-Tarski paradox} A is an infinite set and the family of all $m$ element subsets of $A$ is partitioned into $n$ sets $S_{j},\,1\leq j\leq n$, then there is an infinite subset $B\subseteq A$ such that all $m$ element subsets of $B$ belong to the same $S_{j}$. (Also, see form 17.) Ramsey\index{Author}{Ramsey} [\textbf{1929}].

\section*{ORDERING PRINCIPLES INCLUDING PROPERTIES OF PARTIAL ORDERS}

\textbf{I. Linear And Well Orderings}

\noindent [\textbf{$\boldsymbol{0}$ Z}] Suppose $X$ and $M$ are sets, $F$ is a function from $\mathcal{P}(X)-\{\emptyset\}$ into $M$ and $G$ is a function on $M$ such that $\forall z\in \mathcal{P}(X)-\{\emptyset\},\,G(F(z))\in \mathcal{P}(z)-\{\emptyset\}$. Then there are unique $W$ and $R$ such that $W\subseteq M,\,R\subseteq W\times W$ and $R$ is an irreflexive well ordering of $W$, the range of $G|W$ is a partition of $X$, and $\forall w\in W,\,w=F(X-\,(\bigcup\{G(\nu)\,:\nu\in W\,\mathrm{and}\,\nu\,R\,w\})$). Kruse\index{Author}{Kruse} [\textbf{1974}].

\noindent [\textbf{$\boldsymbol{0}$ AH}] No infinite well orderable set is amorphous.

\noindent [\textbf{1 J}] Form 49 + form 51. Morris\index{Author}{Morris} [\textbf{1969}] and note 121.

\noindent [\textbf{1 K}] Form 133 + form 64. Clear

\noindent [\textbf{8 G}] Let $R$ be a relation and $R^{\ast}$, the transitive closure of $R$. ($x\,R^{\ast}\,y$ iff there exists a finite ascending $R$-chain connecting $x$ to $y$.) If every ascending $R$-chain is finite, then every ascending $R^{\ast}$-chain is finite. Diener\index{Author}{Diener} [\textbf{1994}].

\noindent [\textbf{15 B}] SDO: For any family $X$ of infinite sets, there is a function $f$ with domain $X$ such that $\forall y\in X,\,f(y)$ is a dense linear ordering of $y$. Pincus\index{Author}{Pincus} [\textbf{1997}].

\noindent [\textbf{15 C}] SUO: For any family $X$ of infinite sets, there is a function $f$ with domain $X$ such that $\forall y\in X,\,f(y)$ is an unbounded linear ordering of $y$. (An ordering is unbounded if there is neither a least element nor a greatest element.) Pincus [\textbf{1997}].

\noindent [\textbf{15 D}] SDUO: For any family $X$ of infinite sets, there is a function $f$ with domain $X$ such that $\forall y\in X,\,f(y)$ is a dense unbounded linear ordering of $y$. Pincus [\textbf{1997}].

\noindent \textbf{30.} Ordering Principle:\index{Subject}{ordering principle} Every set can be linearly ordered. Schr\"{o}der\index{Author}{Schroder@Schr\"{o}der} [\textbf{1988}].

\noindent [\textbf{30 A}] If $R$ is a partial ordering in which the set of predecessors of each element is linearly ordered, then there is a linear ordering extending $R$. Felgner\index{Author}{Felgner} [\textbf{1969}].

\noindent [\textbf{30 B}] ADO: Every infinite set differs finitely from a set with a dense linear ordering: for all infinite $x$, there is a $y\subseteq x$ such that $y$ has a dense linear ordering and $x-y$ is finite. Pincus [\textbf{1997}].

\noindent [\textbf{30 C}] SADO: For every set $F$ there is a function $f$ with domain $F$ such that for all $x\in F,\, f(x)\subseteq x,\,f(x)$ has a dense linear ordering, and $x-f(x)$ is finite, Pincus [\textbf{1997}].

\noindent [\textbf{30 D}] UO: There exists an unbounded linear ordering on every infinite set. (An ordering is unbounded if there is neither a least element nor a greatest element.) Pincus~[\textbf{1997}].

\noindent [\textbf{30 E}] ADUO: For all infinite $x$, there exists $y\subseteq x$ such that $y$ has a dense unbounded linear ordering and $x-y$ is finite. Pincus [\textbf{1997}].

\noindent [\textbf{30 F}] SADUO: For every set $F$ there is a function $f$ with domain $F$ such that for all $x\in F,\, f(x)\subseteq x,\,f(x)$ has a dense unbounded linear ordering, and $x-f(x)$ is finite. Pincus\index{Author}{Pincus} [\textbf{1997}].

\noindent \textbf{49.} Order Extension Principle:\index{Subject}{order extension principle} Every partial ordering can be extended to a linear ordering. Tarski\index{Author}{Tarski} [\textbf{1924}] p 78.

\noindent \textbf{51.} Cofinality Principle:\index{Subject}{cofinality principle} Every linear ordering has a cofinal sub well ordering. Sierpi\'{n}ski\index{Author}{Sierpinski@Sierpi\'{n}ski} [\textbf{1918}] p 117.

\noindent [\textbf{51 A}] For all linear orders $(X,\,\leq)$, there is a sup function on the bounded, well ordered subsets of $X$. Manka\index{Author}{Manka} [\textbf{1988a}] and note 38.

\noindent [\textbf{51 C}] For every linearly ordered set $(X,\,\leq)$, if every well ordered subset of $X$ is bounded above then there is a sup function on the well ordered subsets of $X$. Manka [\textbf{1988a}] and note 38.

\noindent [\textbf{67 B}] $MW$: Each set can be covered by a well ordered family of finite sets. Levy\index{Author}{Levy} [\textbf{1962}].

\noindent \textbf{77.} A linear ordering of a set $P$ is a well ordering if and only if $P$ has no infinite descending sequences. Jech\index{Author}{Jech} \cite{bib:1973b} p 23.

\noindent [\textbf{77 A}] If a linearly ordered set $(X,\,\leq)$ has no greatest element then it has a subset of order type $\omega$. Sierpi\'{n}ski [\textbf{1918}].

\noindent \textbf{79.} $\mathbb{R}$ can be well ordered. Hilbert\index{Author}{Hilbert} [\textbf{1900}] p 263.

\noindent \textbf{90.} $LW$: Every linearly ordered set can be well ordered. Jech \cite{bib:1973b} p 133 and Felgner/Jech\index{Author}{Felgner/Jech} [\textbf{1973}].

\noindent \textbf{91.} $PW$: The power set of a well ordered set can be well ordered. Jech \cite{bib:1973b} p 133 and H. Rubin/J. Rubin\index{Author}{H. Rubin/J. Rubin} \cite{bib:1985q}.

\noindent \textbf{103.} If $(P,\,<)$ is a linear ordering and $|P|>\aleph_{1}$ then some initial segment of $P$ is uncountable. Jech \cite{bib:1973b} p 164 prob 11.21.

\noindent [\textbf{118 G}] Each complete linear order admits an extreme choice, that is, if $(X,\,\leq)$ is a linear order then there is a choice function $f$ on the set of non-empty open interval of $X$ with the property that for all $(a,\,b,\,c)\in X^{3}$ with $a<b<c,\,f((a,\,c))\in\{f((a,\,b)),\,b,\,f((b,\,c))\}$. Morillon~[\textbf{1991a}].

\noindent [\textbf{118 H}] For each linear order $(X,\,\leq)$, there is an ordinal $\alpha$ and a strictly increasing injective map $j$ from $(X,\,\leq)$ to $\{0,1\}^{\alpha}$ with the lexicographic order. Morillon [\textbf{1987}].

\noindent [\textbf{118 I}] Every linearly ordered set of linear orders each of which is non-empty and conditionally complete has a choice function. Morillon\index{Author}{Morillon} \cite{bib:1988h}, van Douwen\index{Author}{Douwen} [\textbf{1985}], and note 71.

\noindent [\textbf{118 J}] If $(X,\,\leq)$ is a linear, complete order, then the order topology on $X$ is hereditarily normal. Morillon \cite{bib:1988h}, van Douwen [\textbf{1985}], and note 71.

\noindent [\textbf{118 K}] If $(X,\,\leq)$ is a linear order, then the order topology on $X$ is highly separated. Morillon \cite{bib:1988h} and note 71.

\noindent [\textbf{118 L}] If $X$ is a linear, complete order, then the order topology on $X$ is highly separated. Morillon \cite{bib:1988h} and note 71.

\noindent [\textbf{118 M}] If $(X,\,\leq)$ is a linear order, then the family $\{A(a,\,b):a,\,b\in X\wedge a<b\}$ has a choice function. ($A(a,\,b)$ denotes the set of continuous, order preserving functions from $X$ into $[0,1]$ such that $f(a)=0$ and $f(b)=1$.) Morillon\index{Author}{Morillon} \cite{bib:1988h} and note 71.

\noindent [\textbf{118 N}] If $(X,\,\leq)$ is a linear order, then the family $\{A^{\prime}(a,\,b):a,\,b\in X\wedge a<b\}$ has a choice function. ($A^{\prime}(a,\,b)$ denotes the set of continuous functions from $X$ into $[0,1]$ such that $f(a)=0$ and $f(b)=1.)$ Morillon \cite{bib:1988h} and note 71.

\noindent [\textbf{118 O}] If $(X,\,\leq)$ is a linear order, then the order topology on $X$ is effectively normal. Morillon \cite{bib:1988h} and note 71.

\noindent [\textbf{118 P}] If $(X,\,\leq)$ is a complete linear order, then the order topology on $X$ is effectively normal. Morillon \cite{bib:1988h} and note 71.

\noindent [\textbf{118 Q}] If $(X,\,\leq)$ is a linear order, then the order topology on $X$ is hereditarily normal. Morillon \cite{bib:1988h} and note 71.

\noindent [\textbf{118 R}] If $(X,\,\leq)$ is a linear order, then the order topology on $X$ is monotonely normal. Morillon [\textbf{1987}] and note 71.

\noindent [\textbf{118 S}] If $(X,\,\leq)$ is a complete linear order, then there is a choice function $\ast$ for the family of non-empty open intervals of $X$ such that if $a,\,b,\,c\in X$ and either $(a,\,b)\neq\emptyset$ or $(b,\, c)\neq\emptyset$, then $\ast((a,\,c))\in\{\ast((a,\,b)),\,\ast((b,\,c)),\,b,\,b^{+}\}$ (where $b^{+}$ is the successor of $b$ if it exists and $b$ otherwise). Morillon [\textbf{1987}] and note 71.

\noindent [\textbf{118 T}] Every complete linear order has an increasing choice function. Morillon [\textbf{1987}] and note 71.

\noindent \textbf{129.} For every infinite set $A,\,A$ admits a partition into sets of order type $\omega^{\ast}+\omega$. (For every infinite $A$, there is a set $\{\langle C_{j},\,<j\rangle:j\in J\}$ such that $\{C_{j}:j\in J\}$ is a partition of $A$ and for each $j\in J, <j$ is an ordering of $C_{j}$ of type $\omega^{\ast}+\omega$.) Brunner\index{Author}{Brunner} [\textbf{1984f}].

\noindent \textbf{130.} $\mathcal{P}(\mathbb{R})$ is well orderable. Brunner [\textbf{1984f}].

\noindent \textbf{133.} Every set is either well orderable or has an infinite amorphous subset. Brunner [\textbf{1982a}].

\noindent \textbf{163.} Every non-well-orderable set has an infinite, Dedekind finite subset. Brunner [\textbf{1982a}] and note 94.

\noindent \textbf{164.} Every non-well-orderable set has an infinite subset with a Dedekind finite power set. Brunner [\textbf{1982a}], notes 42 and 94.

\noindent [\textbf{170 B}] There is an uncountable, well orderable set of reals. Raisonnier\index{Author}{Raisonnier} [\textbf{1982}] and note 3.

\noindent \textbf{185.} Every linearly ordered Dedekind finite set is finite. Hickman\index{Author}{Hickman} [\textbf{1979a}] and note 94.

\noindent [\textbf{185 A}] If $X$ is a linearly ordered set and there is neither an increasing nor a decreasing injection $f:\omega\rightarrow X$, then $X$ is finite. Hickman [\textbf{1979a}].

\noindent \textbf{198.} For every set $S$, if the only linearly orderable subsets of $S$ are the finite subsets of $S$, then either $S$ is finite or $S$ has an amorphous subset. Hickman [\textbf{1978b}] and note 57.

\noindent \textbf{227.} For all groups $G$, if every finite subgroup of $G$ can be fully ordered then $G$ can be fully ordered. Rav [\textbf{1977}] and note 79.

\noindent \textbf{228.} Every torsion free Abelian group can be fully ordered. Rav\index{Author}{Rav} [\textbf{1977}] and note 79.

\noindent \textbf{229.} If $(G,\,\circ,\,\leq)$ is a partially ordered group, then $\leq$ can be extended to a linear order on $G$ if and only if for every finite set $\{a_{1},\,\ldots,\,a_{n}\}\subseteq G$, with $a_{i}\neq$ the identity for $i=1$ to $n$, the signs $\epsilon_{1},\,\ldots,\,\epsilon_{n}\,(\epsilon_{i}=\pm 1)$ can be chosen so that $P\cap S(a_{1}^{\epsilon_{1}},\,\ldots,\,a_{n}^{\epsilon_{n}})=\emptyset$ (where $S(b_{1},\,\ldots,\, b_{n})$ is the normal sub-semi-group of $G$ generated by $b_{1},\,\ldots,\,b_{n}$ and $P=\{g\in G:e\leq g\}$ where $e$ is the identity of $G$.) Rav~[\textbf{1977}].

\noindent [\textbf{231 A}] Every almost well orderable set ($W$-set) is well orderable. Note 25.

\noindent \textbf{293.} For all sets $x$ and $y$, if $x$ can be linearly ordered and there is a mapping of $x$ onto $y$, then $y$ can be linearly ordered.

\noindent \textbf{294.} Every linearly ordered $W$-set is well orderable. Morris\index{Author}{Morris} [\textbf{1969}] and note 25.

\noindent \textbf{295.} DO: Every infinite set has a dense linear ordering. Gonzalez\index{Author}{Gonzalez} [\textbf{1995a}].

\noindent [\textbf{295 A}] DUO: Every infinite set has a dense unbounded linear ordering. (An ordering is \emph{unbounded} if there is neither a least element nor a greatest element.) Pincus\index{Author}{Pincus} [\textbf{1997}].

\noindent \textbf{316.} If a linearly ordered set $(A,\,\leq)$ has the fixed point property then $(A,\,\leq)$ is complete. ($(A,\,\leq)$ has the \emph{fixed point property} if every function $f:A\rightarrow A$ satisfying $(x\leq y\Rightarrow f(x)\leq f(y)$) has a fixed point, and ($(A,\,\leq)$ is \emph{complete} if every subset of $A$ has a least upper bound.) H\"{o}ft/Howard\index{Author}{Hoft/Howard@H\"{o}ft/Howard} [\textbf{1994}].

\noindent \textbf{377.} Restricted Ordering Principle: For every infinite set $X$ there is an infinite subset $Y$ of $X$ such that $Y$ can be linearly ordered. De la Cruz/Di Prisco\index{Author}{De la Cruz/Di Prisco}  [\textbf{1996}].

\section*{{II}. Versions Of Martin's Axiom}

\noindent [\textbf{$\boldsymbol{0}$ S}] If $(P,\,\leq)$ is a quasi-order such that $P$ is non-empty, well orderable and every antichain is finite then there is a $P$-generic filter . Shannon\index{Author}{Shannon} [\textbf{1990}] and note 47.

\noindent [\textbf{1 Z}] If $(P,\,\leq)$ is a quasi-order and $P\neq\emptyset$ and every antichain is finite, then there is a $P$ generic filter which is well orderable. Shannon [\textbf{1990}] and note~47.

\noindent [\textbf{8 F}] Martin's Axiom\index{Subject}{Martin's axiom} $(\aleph_{0})$: Whenever $(P\leq)$ is a non-empty, ccc quasi-order (ccc means every anti-chain is countable) and $\mathcal{D}$ is a family of $\leq\aleph_{0}$ dense subsets of $P$, then there is a $\mathcal{D}$ generic filter $G$ in $P$. Kunen\index{Author}{Kunen} [\textbf{1980}], Shannon [\textbf{1990}] and note 47.

\noindent [\textbf{10 I}] Assume $(P,\,\leq)$ is a partial order such that $P$ is a denumerable union of finite sets and contains a non-empty subset $Q$ relative to which depths are defined for all $x\in \mathrm{P}$. If all $Q$-strong antichains of $P$ are finite then for each denumerable family $\mathcal{D}$ of dense sets there is a $\mathcal{D}$ generic filter. Shannon\index{Author}{Shannon} \cite{bib:1990f} and note 47.

\noindent [\textbf{43 C}] Existence of a generic set: If $\mathcal{D}$ is a denumerable collection of dense sets in a partial order $(P,\,\leq)$ then for all $p\in P$ there is a $\mathcal{D}$ generic $G$ with $p\in G$. Goldblatt\index{Author}{Goldblatt} \cite{bib:1985m} and note 47.

\noindent [\textbf{43 D}] If $\mathcal{D}$ is a denumerable collection of dense subsets of a partial order $(P,\,\leq)$ then for all $p\in P$ there is a sequence $f:\omega \rightarrow P$ with $f(0)=p$ and $(\forall n\in \omega)(f(n)\geq f(n+1))$ and range $(f)\cap D\neq\emptyset$ for all $D\in \mathcal{D}$. Goldblatt~\cite{bib:1985m}.

\noindent \textbf{171.} If $(P,\,\leq)$ is a partial order such that $P$ is the denumerable union of finite sets and all antichains in $P$ are finite then for each denumerable family $\mathcal{D}$ of dense sets there is a $\mathcal{D}$ generic filter. Shannon \cite{bib:1990f} and note 47.

\textbf{III. Trees}

\noindent [\textbf{10 F}] K\"{o}nig's Lemma:\index{Subject}{K\"{o}nig's lemma} Every $\omega$ tree has an infinite chain. Pincus\index{Author}{Pincus} \cite{bib:1972c1}, note 21, and note 35.

\noindent [\textbf{43 H}] Every tree in which every element has at least one successor has a branch. Blass\index{Author}{Blass} [\textbf{1979}] and note note 21.

\noindent [\textbf{86 A($\alpha$)}] If $(P,\,\leq)$ is a tree of height $\lambda\leq\aleph_{\alpha}$ which has fewer than $\lambda$ branches of height $<\lambda$ then one of the following holds for the upside down tree $(P,\, \geq)$:
\begin{enumerate}
\item[(a)] There is a cardinal $\mu<\lambda$ such that $(P, \geq)$ does not contain a strong antichain of cardinality $\mu$.
\item[(b)] $(P,\,\geq)$ contains a strong antichain of cardinality $\lambda$.
\item[(c)] $(P,\,\geq)$ contains a chain of cardinality $\lambda$.
\end{enumerate}
Shannon [\textbf{1992}] and note 119.

\noindent [\textbf{106 A}] $DMC$: Dependent Multiple Choice: Every tree in which every element has at least one immediate successor has a subtree whose levels are finite and in which every element has at least one immediate successor. Blass [\textbf{1979}], Brunner\index{Author}{Brunner} [\textbf{1983c}], Fossy/Morillon\index{Author}{Fossy/Morillon} [\textbf{1995}] and note~21.

\noindent \textbf{216.} Every infinite tree has either an infinite chain or an infinite antichain. Pincus [\textbf{1972a}] and note 21.

\noindent \textbf{249.} If $T$ is an infinite tree in which every element has exactly 2 immediate successors then $T$ has an infinite branch. Truss\index{Author}{Truss} [\textbf{1975}], note 87 and note 21.

\textbf{IV. Other Partial Ordering Properties (See note 71.)}

\noindent [\textbf{$\boldsymbol{0}$ J}] Zermelo's Fixed Point Theorem:\index{Subject}{Zermelo's fixed point theorem} For every partially ordered set $(X,\,\leq)$, if every well ordered subset has a least upper bound then every $f:X\rightarrow X$ satisfying $\forall t\in X,\,t\leq f(t)$ has a fixed point. Manka\index{Author}{Manka} [\textbf{1988a}], Zermelo\index{Author}{Zermelo} [\textbf{1908b}] and Abian\index{Author}{Abian} [\textbf{1980}].

\noindent [\textbf{$\boldsymbol{0}$ K}] For all partially ordered sets $(X,\,\leq)$, if there is a sup function $\sigma$ on the well ordered subsets of $X$ then every $f:X\rightarrow X$ satisfying $\forall t\in X,\,t\leq f(t)$ has a fixed point. Manka [\textbf{1988a}] and note 38.

\noindent [\textbf{$\boldsymbol{0}$ L}] Caristi's Fixed Point Theorem: If $(X,\,\rho)$ is a complete metric space and $\phi$: $X\rightarrow \mathbb{R}$ is bounded above and upper semi-continuous then in the Br{\o}ndsted ordering $(x\leq y\,\mathrm{iff}\,\rho(x,\,y)\leq\phi(y)-\phi(x))$ every $f:X\rightarrow X$ satisfying $\forall t\in X,\,t \leq f(t)$ has a fixed point. Caristi\index{Author}{Caristi} [\textbf{1976}], Manka [\textbf{1988a}], and note 38.

\noindent [\textbf{1 Q}] For any\index{Subject}{Caristi's fixed point theorem} partial order $(X,\,\leq)$ there is a $\sup$ function on the bounded (above), well ordered subsets of $X$. Manka [\textbf{1988a}] and note 38.

\noindent [\textbf{1 R}] For any partial order $(X,\,\leq)$ there is a $\sup$ function on the bounded subsets of $X$. Manka [\textbf{1988a}] and note~38.

\noindent [\textbf{1 S}] For any partial order $(X,\,\leq)$ if every well ordered subset of $X$ is bounded then there is a sup function on the well ordered subsets of $X$. Manka [\textbf{1988a}] and note 38.

\noindent [\textbf{1 T}] For every directed partial order $(X,\,\leq)$ there is a sup function on the bounded, well ordered subsets of $X$. Manka\index{Author}{Manka} [\textbf{1988a}] and note 38.

\noindent [\textbf{1 U}] For every directed partial order, if every well ordered subset of $X$ is bounded above then there is a sup function on the well ordered subsets of $X$. Manka [\textbf{1988a}] and note~38.

\noindent [\textbf{1 V}] For every partially ordered set $(X,\,\leq)$, if every well ordered subset is bounded above then every function $f:X\rightarrow X$ satisfying $\forall t\in X,\,t\leq f(t)$ has a fixed point. Manka [\textbf{1988a}] and Abian\index{Author}{Abian} [\textbf{1980}].

\noindent [\textbf{1 AJ}] Let $\mathcal{F}$ be a family of mappings of a poset $(X,\,\leq)$ into itself such that $\forall f\in \mathcal{F},\,\forall x\in X,\,f(x)\leq x$. If for some element $e$ of $X$ each chain in $X$ containing $e$ has a lower bound, then the family $\mathcal{F}$ has a common fixed point. Kasahara\index{Author}{Kasahara} [\textbf{1976}].

\noindent [\textbf{14 W}] Universal Subnet Theorem: Any net has a universal subnet. Schechter\index{Author}{Schechter} [\textbf{1996a}] and note 31.

\noindent [\textbf{43 K}] Let $M$ be a non-empty set partially ordered by $<$ without $<$-minimal elements. Then there is a subset $C$ of $M$ which is linearly ordered by $<$ such that $C$ has no $<$-minimal element. Note 146.

\noindent [\textbf{43 L}] $Z_{\omega}$: If every chain in a partially ordered set $P$ is finite, then $P$ contains a maximal element. Tarski\index{Author}{Tarski} [\textbf{1948}], Wolk\index{Author}{Wolk} [\textbf{1983}], and note 54.

\noindent [\textbf{43 M}] If $(X,\,\rho)$ is a complete metric space and $\phi:X\rightarrow \mathbb{R}$ is bounded above and upper semi-continuous then in the Br{\o}ndsted ordering, $(x\leq y\,\mathrm{iff}\,\rho(x,\,y)\leq \phi(y)-\phi(x))$, there is a maximal element. Brunner\index{Author}{Brunner} [\textbf{1987a}] and note 38.

\noindent [\textbf{43 Q}] Matrix$(\omega)$: For every atomless poset $(P,\,\leq)$ and every family $D=\{D_{i}:i\in \omega\}$ of dense open subsets of $P$, there exists a matrix $C=\{C_{i}:i\in \omega\},\, |\bigcup C|=\omega$ of $P$ for $D$. Keremedis\index{Author}{Keremedis} [\textbf{1996d}] and note 132.

\noindent [\textbf{43 R}] Any partial order $(A,\,\leq)$ with no infinite descending chains is well founded, (that is, for all $X$ such that $\emptyset\,\neq X\subseteq A,\,X$ has a minimal element.) Spector\index{Author}{Spector} [\textbf{1980}] and note 54.

\noindent \textbf{59($\leq$).} If $(A,\,\leq)$ is a partial ordering that is not a well ordering, then there is no set $B$ such that $(B,\,\precsim)$ (the usual injective cardinal ordering on $B$) is isomorphic to $(A,\,\leq)$. Mathias\index{Author}{Mathias} [\textbf{1979}] p 120.

\noindent [\textbf{70 B}] Weak Universal Net Theorem: There is a universal net in $\omega$ which is not eventually constant. Schechter [\textbf{1996a}] and note 31.

\noindent [\textbf{87 A($\alpha$)}] $Z_{\aleph_{\alpha}}$: Let $P$ be a partially ordered set in which every well ordered chain has type $<\aleph_{\alpha}$. If every well ordered chain in $P$ has an upper bound in $P$, then $P$ contains a maximal element. Wolk [\textbf{1983}].

\noindent [\textbf{87 B($\alpha$)}] $Z_{\aleph_{\alpha}}^{\ast}$: Let $P$ be a well-founded partially ordered set, (that is, every chain in $P$ is well ordered), in which every chain has type $<\aleph_{\alpha}$. If every chain in $P$ has an upper bound in $P$, then $P$ contains a maximal element. Wolk [\textbf{1983}].

\noindent \textbf{89.} Antichain Principle:\index{Subject}{antichain principle} Every partially ordered set has a maximal antichain. Jech\index{Author}{Jech} \cite{bib:1973b} p 133 and Kurepa\index{Author}{Kurepa} [\textbf{1953}].

\noindent \textbf{105.} There is a partially ordered set $(A,\,\leq)$ such that there is no set $(B,\,\precsim)$ ($\precsim$, the usual injective cardinal ordering) isomorphic to $(A,\,\leq)$. Forti/Honsell\index{Author}{Forti/Honsell} [\textbf{1985}] and Jech [\textbf{1966c}].

\noindent \textbf{144.} Every set is almost well orderable. Keisler\index{Author}{Keisler} [\textbf{1970}] and note 25.

\noindent \textbf{217.} Every infinite partially ordered set has either an infinite chain or an infinite antichain. Pincus\index{Author}{Pincus} [\textbf{1972a}].

\section*{TOPOLOGICAL FORMS (INCLUDING PROPERTIES OF $\mathbb{R}$)}

\textbf{I. Baire Category Type Theorems (See note 28.)}

\noindent [\textbf{43 E}] The Baire Category Theorem for \v{C}ech Complete-I Spaces: Every \v{C}ech complete space is Baire. Goldblatt\index{Author}{Goldblatt} \cite{bib:1985m} and note 28.

\noindent [\textbf{43 F}] Baire Category Theorem for Sequentially Complete Metric Spaces: Every sequentially complete metric space is Baire. Blair\index{Author}{Blair} [\textbf{1977}], Brunner\index{Author}{Brunner} [\textbf{1983c}], Goldblatt \cite{bib:1985m}, and note 28.

\noindent [\textbf{43 J}] Baire Category Theorem for Cantor Complete Metric Spaces: Every Cantor complete metric space is Baire. Brunner [\textbf{1983c}] and note 28.

\noindent [\textbf{43 T}] A sequentially complete metric space is not meager. Note 28.

\noindent [\textbf{43 U}] Baire Category Theorem for Filter Complete Metric Spaces: Every filter complete metric space is Baire, Morillon\index{Author}{Morillon} \cite{bib:1988h} and note 28.

\noindent [\textbf{43 V}] Every convex-compact set in a Hausdorff topological vector space with a basis of convex open sets is Baire. Fossy/Morillon [\textbf{1995}] and note 23.

\noindent \textbf{106.} Baire Category Theorem for Compact Hausdorff Spaces: Every compact Hausdorff space is Baire. Goldblatt \cite{bib:1985m}, G. Moore\index{Author}{G. Moore} [\textbf{1982}], and note 28.

\noindent [\textbf{106 C}] Every scattered compact space is a Baire space. Fossy/Morillon\index{Author}{Fossy/Morillon} [\textbf{1995}] and note 28.

\textbf{II. Product Theorems}

\noindent [\textbf{$\boldsymbol{0}$ B}] If $X$ is finite then $X^{\omega}$ is compact. Brunner [\textbf{1984b}].

\noindent [\textbf{$\boldsymbol{0}$ D}] Products of limited amorphous spaces are limited amorphous. (A space is \emph{limited amorphous} if each amorphous subset is compact.) Brunner [\textbf{1984b}]

\noindent [\textbf{$\boldsymbol{0}$ E}] The box product of a well ordered family of limited amorphous spaces is limited amorphous. (A space is \emph{limited amorphous} if each amorphous subset is compact.) Brunner [\textbf{1984b}] and note 52.

\noindent [\textbf{1 AY}] For each $\lambda\in\Lambda$ let $S_{\lambda}$ be a subset of some topological space $X_{\lambda}$. Then $\mathrm{cl}( \prod_{\lambda\in\Lambda}S_{\lambda})=\prod_{\lambda\in\Lambda}\ \mathrm{cl}\ (S_{\lambda})$, where the first ``$\mathrm{cl}$'' is closure in the product topology and the second ``$\mathrm{cl}$'' is closure in $X_{\lambda}$. Schechter [\textbf{1992}].

\noindent [\textbf{1 AZ}] For each $\lambda\in\Lambda$, let $X_{\lambda}$ be a complete uniform space. Then, the product $\prod_{\lambda\in\Lambda}X_{\lambda}$ equipped with the product uniform structure, is also complete. (See, for example, Kelley\index{Author}{Kelley} [\textbf{1955}] pp 176ff for definitions.) Schechter\index{Author}{Schechter} [\textbf{1992}].

\noindent [\textbf{10 B}] Tychonoff's\index{Subject}{Tychonoff's theorem} Compactness Theorem for Countable Families of Finite $T_{2}$ Spaces: The product of a countable family of finite $T_{2}$ spaces is compact. Krom\index{Author}{Krom} \cite{bib:1981b}.

\noindent [\textbf{14 J}] Tychonoff's Compactness Theorem for Compact Hausdorff Spaces: The product of compact Hausdorff spaces is compact. (For many compactness theorems $T$, it is clear that [14 J] implies $T$ implies [14 K].) {\L}o\'{s}/Ryll-Nardzewski\index{Author}{Los/Ryll-Nardzewski@{\L}o\'{s}/Ryll-Nardzewski}
 [\textbf{1954}] and Rubin/Scott\index{Author}{Rubin/Scott} [\textbf{1954}].

\noindent [\textbf{14 K}] Compactness of the Generalized Cantor Space: $\{0,1\}^{I}$ is compact for any I. (For many compactness theorems $T$, it is clear that [14 J] implies $T$ implies [14 K].) Jech\index{Author}{Jech} \cite{bib:1973b} p 28, prob 17 and Mycielski\index{Author}{Mycielski} [\textbf{1964}].

\noindent [\textbf{14 Z}] Tychonoff's Compactness Theorem for Families of Finite Spaces: If $\{X_{i}:i\in I\}$ is a family of finite topological spaces, then $\prod_{i\in I}X_{i}$ is compact. (Cleary [14 Z] implies [14 K] and is implied by [14 J].) Jech \cite{bib:1973b} p 28 prob 17, A.A. Lewis\index{Author}{Lewis} [\textbf{1983}], and Wolk\index{Author}{Wolk} [\textbf{1967}].

\noindent [\textbf{14 AB}] Tychonoff's Compactness Theorem for Sober Spaces: A product of compact, sober topological spaces is compact. Blass\index{Author}{Blass} [\textbf{1986}], Johnstone\index{Author}{Johnstone} [\textbf{1984}] and note 37.

\noindent [\textbf{43 G}] Form 106 + form 154. Brunner\index{Author}{Brunner} [\textbf{1983c}] [\textbf{1984b}].

\noindent [\textbf{64 A}] An amorphous power of a compact $T_{2}$ space is compact. Brunner [\textbf{1984b}].

\noindent [\textbf{64 B}] Finite products of metric $C$ spaces are $C$. (A space is $C$ if each open covering has an amorphous refinement.) Brunner [\textbf{1984b}].

\noindent \textbf{113.} Tychonoff's Compactness Theorem for Countably Many Spaces: The product of a countable set of compact spaces is compact. Brunner [\textbf{1985a}].

\noindent \textbf{115.} The product of weakly Loeb $T_{2}$ spaces is weakly Loeb. (See form 114 for the definition of weakly Loeb.) Brunner [\textbf{1985a}].

\noindent \textbf{127.} An amorphous power of a compact $T_{2}$ space, which as a set is well orderable, is well orderable. Brunner [\textbf{1984b}].

\noindent \textbf{134.} If $X$ is an infinite $T_{1}$ space and $X^{Y}$ is $T_{5}$, then $Y$ is countable. ($T_{5}$ is ``hereditarily $T_{4}$''.) A. H. Stone\index{Author}{A. H. Stone} [\textbf{1948}].

\noindent \textbf{135.} If $X$ is a $T_{2}$ space with at least two points and $X^{Y}$ is hereditarily metacompact then $Y$ is countable. (A space is \emph{metacompact} if every open cover has an open point finite refinement. If $B$ and $B^{\prime}$ are covers of a space $X$, then $B^{\prime}$ is a \emph{refinement} of $B$ if $(\forall x\in B^{\prime})(\exists y\in B)(x\subseteq y)$. $B$ is \emph{point finite} if $(\forall t\in X)$ there are only finitely many $x\in B$ such that $t\in x$.) van Douwen\index{Author}{Douwen} [\textbf{1980}].

\noindent \textbf{139.} Using the discrete topology on $2,\, 2^{\mathcal{P}(\omega)}$ is compact. Truss\index{Author}{Truss} [\textbf{1984}].

\noindent \textbf{154.} Tychonoff's\index{Subject}{Tychonoff's theorem} Compactness Theorem for Countably Many $T_{2}$ Spaces: The product of countably many $T_{2}$ compact spaces is compact. Brunner [\textbf{1983c}].

\noindent \textbf{177.} An infinite box product of regular $T_{1}$ spaces, each of cardinality greater than 1, is neither first countable nor connected. Brunner\index{Author}{Brunner} [\textbf{1981a}], Knight\index{Author}{Knight} [\textbf{1964}], and note 52.

\noindent \textbf{332.} A product of non-empty compact sober topological spaces is non-empty. Morillon\index{Author}{Morillon} \cite{bib:1988h} and note 71.

\noindent \textbf{343.} A product of non-empty, compact $T_{2}$ topological spaces is non-empty. Morillon \cite{bib:1988h}.

\noindent \textbf{351.} A countable product of metrizable spaces is metrizable. Keremedis\index{Author}{Keremedis} [\textbf{1996d}] and note 132.

\noindent \textbf{352.} A countable product of second countable spaces is second countable. Keremedis [\textbf{1996d}] and note 132.

\noindent \textbf{353.} A countable product of first countable spaces is first countable. Keremedis [\textbf{1996d}] and note 132.

\noindent \textbf{354.} A countable product of separable $T_{2}$ spaces is separable. Keremedis [\textbf{1996d}] and note 132.

\textbf{III. General Topology}

\noindent [\textbf{$\boldsymbol{0}$ C}] Closed subspaces of limited amorphous spaces are limited amorphous. (A space is \emph{limited amorphous} if each amorphous subset is compact.) Brunner [\textbf{1984b}].

\noindent [\textbf{$\boldsymbol{0}$ M}] A $T_{1}$ topological space is Dedekind finite if and only if every subspace has a dense Dedekind finite subset. Brunner [\textbf{1982d}] and note 94.

\noindent [\textbf{$\boldsymbol{0}$ N}] If $X$ and $Y$ are topological spaces with dense Dedekind finite subsets so is $X\times Y$. Brunner [\textbf{1982d}] and note 94.

\noindent [\textbf{1 O}] $A(D1,\,C)$: For every $T_{2}$ topological space $(X,\,T)$, if $X$ is covered by a well ordered family of closed, discrete sets then each family $U\subseteq T$ of pairwise disjoint sets is well ordered. Brunner [\textbf{1983d}] and note 26.

\noindent [\textbf{1 P}] $A(P)$: (Depends on $P$ for $P\in\{W1,\,W2,\,B1,\,L1,\,L2$, $S,\,C,\,H1,\,A1\}$) Every $T_{2}$ topological space has property $P$. Brunner [\textbf{1983d}] and note 26.

\noindent [\textbf{1 BF}] $(\mathrm{D}T_{0})$: Every topological space $X$ has a $T_{0}$ subspace that is dense in $X$. ($Y$ is \emph{dense in} $X$ if there is no non-empty open subset $O\subseteq X$ such that $O\cap Y=\emptyset$.) McCarten\index{Author}{McCarten} [\textbf{1988}] and note 106.

\noindent [\textbf{1 BG}] $(\mathrm{C}T_{0})$: Every topological space $X$ has a $T_{0}$ subspace that is codense in $X$. ($Y$ is \emph{codense in} $X$ if there is no non-empty closed subset $C\subseteq X$ such that $C\cap Y=\emptyset$.) McCarten [\textbf{1988}] and note 106.

\noindent [\textbf{1 BH}] $(\mathrm{T}T_{0})$: Every topological space $X$ has a $T_{0}$ subspace that is thick in $X$. ($Y$ is \emph{thick in} $X$ if there is no non-empty open and closed (clopen) subset $H\subseteq X$ such that $H\cap Y=\emptyset$.) McCarten [\textbf{1988}] and note 106.

\noindent [\textbf{1 BI}] $(MT_{0})$: Every topological space has a maximal $T_{0}$ subspace. Schnare\index{Author}{Schnare} [\textbf{1968}] and note 106.

\noindent [\textbf{1 BL}] Every family of closed sets in a $T_{2}$ topological space includes a maximal subfamily with the finite intersection property. Keremedis [\textbf{1996b}] and H. Rubin/J. Rubin\index{Author}{H. Rubin/J. Rubin} \cite{bib:1985q}, p 60, $M14$.

\noindent [\textbf{1 BM}] For every family $\{(X_{i},\,T_{i}):i\in k\}$ of $T_{1}$ topological spaces, the box topology contains an element $\prod_{i\in k}O_{i}$, such that $\emptyset \neq O_{i}$ and $O_{i}\neq X_{i}$ for all $i\in k$. Keremedis [\textbf{1996b}] and note 77.

\noindent [\textbf{1 BN}] For every family $\{(X_{i},\,T_{i}):i\in k\}$ of disjoint $T_{1}$ topological spaces, the free union contains an open set $O$ such that $\emptyset \neq O\cap X_{i}$ and $O\cap X_{i}\neq X_{i}$ for all $i\in k$. Keremedis\index{Author}{Keremedis} [\textbf{1996b}] and note 77.

\noindent [\textbf{1 BO}] For every $T_{1}$ topological space $(X,\,T)$ and every base (subbase) $B$ for $X$ there exists a $\subseteq$-maximal cellular family $Q\subseteq B$. Keremedis [\textbf{1996b}] and note 77.

\noindent [\textbf{1 BP}] For every dense in itself topological space $(X,\,T)$ and every base $B$ for $X$ there exists a nest $N$ included in $B$. Keremedis [\textbf{1996b}] and note 77.

\noindent [\textbf{1 BQ}] For every dense in itself topological space $(X,\,T)$ and every base $B$ for $X$ there exists a tower $R$ included in $B$. Keremedis [\textbf{1996b}] and note 77.

\noindent [\textbf{1 BR}] For every $T_{2}$ space $(X,\,T)$ and every subbase $\mathcal{B}$ for $X$, if $\mathcal{C}$ and $\mathcal{D}$ are cellular families of $X$ included in $\mathcal{B}$, then there is a cellular family $\mathcal{E} \subseteq \mathcal{B}$ such that $\mathcal{C}$ and $\mathcal{D}$ are equipollent to subsets of $\mathcal{E}$. Keremedis [\textbf{1996b}] and note 77.

\noindent [\textbf{1 BW}] If $(X,\,T)$ is a $T_{4}$ topological space and $U=\{U_{i}:i\in K\},\,|U_{i}|\geq 2$, is a locally finite family of open non-empty pairwise disjoint sets, then there is a continuous real valued function $f$ on $X$ which assumes at least two distinct values on each $U_{i}$. Keremedis [\textbf{1995}] and note 43.

\noindent [\textbf{1 CB}] For every infinite cardinal $\kappa$, if $(X,\,T)$ is a $T_{2}$ space having a base $B$ such that no $Q\subseteq B,\,|Q|<\kappa$, is a tower, then for every family $\mathcal{D}=\{D_{i}:i\in\kappa\}$ of dense open sets of $X$ there is a filter $F\subseteq T$ such that for every $i\in\kappa$ there is a $b\in F$ with $b\subseteq D_{i}$. Keremedis [\textbf{1996b}] and note 77.

\noindent [\textbf{1 CF}] If a topological space is $cwH$, then it is $cwH(B)$ for every basis $B$. Howard/Keremedis/Rubin/Rubin [\textbf{1997b}] and note 141.

\noindent [\textbf{1 CG}] The disjoint union of $cwH$ spaces is $cwH(B)$ for every basis $B$. Howard/Keremedis/Rubin/Rubin\index{Author}{Howard/Keremedis/Rubin/Rubin} [\textbf{1997b}] and note 141.

\noindent [\textbf{1 CH}] If a topological space is $cwN$, then it is $cwH(B)$ for any basis $B$. Howard/Keremedis/Rubin/Rubin [\textbf{1997b}] and note 141.

\noindent [\textbf{8 C}] Every $\sigma$-compact space is Lindel\"{o}f. Brunner [\textbf{1982b}] and note 43.

\noindent [\textbf{8 D}] Every $\sigma$-compact, locally compact $T_{2}$ space is weakly Lindel\"{o}f. Brunner\index{Author}{Brunner} [\textbf{1982b}] and note 43.

\noindent [\textbf{8 I}] A disjoint union of a countable set of Lindel\"{o}f $T_{2}$ spaces is Lindel\"{o}f. Keremedis [\textbf{1996d}] and note 132.

\noindent [\textbf{8 J}] Hemicompact $T_{2}$ spaces are Lindel\"{o}f. (A \emph{hemicompact} space is a space $(X,\,T)$ that can be written as a disjoint union $\bigcup\{K_{n}:n\in\omega\}$ such that (i) each $K_{n}$ is a compact set in $X$; (ii) $K_{N}\subseteq K_{n+1}$ for all $n\in\omega$; and (iii) for every compact set $K$ in $X$, there is an $n\in\omega$ such that $K\subseteq K_{n}$.) Keremedis [\textbf{1996d}] and note 132.

\noindent [\textbf{8 K}] Hemilindel\"{o}f $T_{2}$ spaces are Lindel\"{o}f. (In [8 J], in the definition of hemicompact replace ``compact'' by ``Lindel\"{o}f'' to get the definition of hemilindel\"{o}f.) Keremedis [\textbf{1996d}] and note 132.

\noindent [\textbf{8 L}] Every second countable topological space is separable. Keremedis [\textbf{1996d}] and note 132.

\noindent [\textbf{8 N}] $DP$: If $(X,\,T)$ is a topological space having a countable $\pi$-base, then for every family $D=\{D_{i}:i\in\omega\}$ of dense open sets of $X$, there is a countable dense set $S\subseteq X$ such that for all $i\in \omega$ and for all but finitely many $s\in S,\,s\in D_{i}$. Keremedis\index{Author}{Keremedis} [\textbf{1996d}] and note 132.

\noindent [\textbf{8 O}] $DUCCC\ +$ form 131. $DUCCC$ is ``If $X$ is a denumerable family of disjoint topological spaces each having the ccc property, then the disjoint union of the spaces in $X$ has the ccc property.'' $(Y$ has the $ccc$ \emph{property} if every collection of disjoint open sets is countable.) Keremedis [\textbf{1996d}] and note 132.

\noindent [\textbf{8 P}] $DUSS\ +$ form 131. $DUSS$ is ``The disjoint union of a denumerable family of separable topological spaces is separable.'' Keremedis [\textbf{1996d}] and note 132.

\noindent [\textbf{8 Q}] $DUSC\ +$ form 131. $DUSC$ is ``The disjoint union of a denumerable family of second countable topological spaces is second countable.'' Keremedis [\textbf{1996d}] and note~132.

\noindent [\textbf{8 R}] $SCDUL\ +$ form 131. $SCDUL$ is ``The denumerable disjoint union of second countable spaces is Lindel\"{o}f.'' Keremedis [\textbf{1996d}] and note 132.

\noindent [\textbf{10 H}] $\sigma$-compact, paracompact $T_{2}$ spaces are weakly Lindel\"{o}f. Brunner\index{Author}{Brunner} [\textbf{1982b}] and note 43.

\noindent [\textbf{13 B}] There is a separable topological space $X$ and a set $A\subseteq \mathbb{R}$ such that $X^{A\times\omega}$ has no dense, Dedekind finite subset. Brunner [\textbf{1982d}] and note 94.

\noindent [\textbf{14 BY}] For every topological space $X$, every proper ideal in the normed $\mathbb{R}$-algebra of continuous bounded functions from $X$ into $\mathbb{R}$ is contained in a maximal ideal. Morillon\index{Author}{Morillon} \cite{bib:1988h}.

\noindent [\textbf{14 CB}] If $X$ is a topological space and $I$ is a proper ideal in the algebra of all bounded continuous functions from $X$ to $\mathbb{R}$ such that every function in $I$ vanishes at at least one point of $X$, then $I$ is contained in a maximal ideal. Morillon \cite{bib:1988h}.

\noindent [\textbf{40 A}] $A(L2,\,L1)$: For every $T_{2}$ topological space $(X,\, T)$, if each open cover of $X$ has a well ordered subcover, then each open cover has a well ordered refinement. Brunner [\textbf{1983d}] and note 26.

\noindent [\textbf{50 L}] Form 301 + form 14. Bell\index{Author}{Bell} \cite{bib:1988b}.

\noindent [\textbf{67 A}] $A(A2,\,H2)$: For every $T_{2}$ topological space $(X,\,T)$ if $(X,\,T)$ is A2, then $(X,\,T)$ is hereditarily $A2$. ($(X,\,T)$ is A2 means if $U\subseteq T$ covers $X$ then $\exists f:X\rightarrow T$ such that for all $x\in X,\,x\in f(x)$ and $f$''$ X$ refines $U$.) Brunner [\textbf{1983d}] and note 26.

\noindent [\textbf{67 C}] $A(P)$: (Depends on $P\in\{D1,\,B2,\,B3,\,H2,\,A2\}$.) Every $T_{2}$ topological space has property $P$. Brunner\index{Author}{Brunner} [\textbf{1983d}] and note 26.

\noindent [\textbf{67 G}] For every $T_{1}$ topological space $(X,\,T)$ and every cellular family $\mathcal{O}= \{O_{i}:i\in k\}\subseteq T-\{\emptyset\}$, there is a closed discrete set $\mathcal{I}$ such that for every $i\in k,\, \mathcal{I}\cap O_{i}\neq\emptyset$. Keremedis [\textbf{1996b}] and note 77.

\noindent [\textbf{67 I}] If $(X,\,T)$ is a $T_{4}$ topological space and $U=\{U_{i}:i\in K\}$ is a family of non-empty open sets, then there is a family $F=\{f_{i}:i\in K\}$ of non-negative, non-constant, continuous real valued functions on $X$ such that each $f_{i}$ vanishes on $X-U_{i}$. Keremedis\index{Author}{Keremedis} [\textbf{1995}].

\noindent [\textbf{67 J}] If $(X,\,T)$ is a $T_{4}$ topological space and $U=\{U_{i}:i\in K\}$ is a family of non-empty open sets, then there is a family $V=\{V_{i}:i\in K\}$ of non-empty closed (compact) sets such that $U_{i}\supseteq V_{i}$ for each $i\in K$. Keremedis [\textbf{1995}].

\noindent [\textbf{67 K}] DUN: The disjoint union of normal topological spaces is normal. Howard/Keremedis/Rubin/Rubin [\textbf{1997a}] and note 141.

\noindent [\textbf{67 L}] DUU: If $C$ is a collection of pairwise disjoint topological spaces each of which satisfies Urysohn's lemma\index{Subject}{Urysohn's lemma}, then the disjoint union of the spaces in $C$ satisfies Urysohn's\index{Author}{Urysohn} lemma. Howard/Keremedis/Rubin/Rubin\index{Author}{Howard/Keremedis/Rubin/Rubin} [\textbf{1997a}] and note 141.

\noindent [\textbf{67 M}] DUT: If $C$ is a collection of pairwise disjoint topological spaces each of which satisfies the Tietze-Urysohn extension theorem, then the disjoint union of the spaces in $C$ satisfies the Tietze-Urysohn extension theorem. Howard/ Keremedis/Rubin/Rubin [\textbf{1997a}] and note 141.

\noindent [\textbf{67 N}] DUP: The disjoint union of paracompact spaces is paracompact. Howard/Keremedis/Rubin/Rubin [\textbf{1997b}] and note 141.

\noindent [\textbf{67 O}] DUMET: The disjoint union of metacompact spaces is metacompact. Howard/Keremedis/Rubin/Rubin [\textbf{1997b}] and note 141.

\noindent [\textbf{67 P}] DUPFCS: The disjoint union of PFCS spaces is PFCS. ($X$ is a \emph{PFCS space} iff every point finite open cover of $X$ is shrinkable.) Howard/Keremedis/Rubin/Rubin [\textbf{1997b}] and note 141.

\noindent [\textbf{67 Q}] DUPN: The disjoint union of paracompact spaces is normal. Howard/Keremedis/Rubin/Rubin [\textbf{1997b}] and note 141.

\noindent [\textbf{67 R}] DUcwH: The disjoint union of cwH spaces is cwH. Howard/Keremedis/ Rubin/Rubin [\textbf{1997b}] and note 141.

\noindent [\textbf{67 S}] DUcwN: The disjoint union of cwN spaces is cwN. Howard/Keremedis/ Rubin/Rubin [\textbf{1997b}] and note 141.

\noindent [\textbf{67 T}] DUcwNN: The disjoint union of cwN spaces is normal. Howard/ Keremedis/Rubin/Rubin [\textbf{1997b}] and note~141.

\noindent \textbf{78.} Urysohn's Lemma: If $A$ and $B$ are disjoint closed sets in a normal space $S$, then there is a continuous $f:S\rightarrow[0,\,1]$ which is 1 everywhere in $A$ and $0$ everywhere in $B$. Urysohn [\textbf{1925}] pp 290-292.

\noindent [\textbf{91 A}] $A(S,\,W1)$: For every $T_{2}$ topological space $(X,\,T)$ if $X$ has a dense well ordered subset, then $T$ is well ordered. (Clear) Brunner [\textbf{1983d}] and note 26.

\noindent [\textbf{94 A}] The discrete topology on $\omega$ is Lindel\"{o}f. Brunner [\textbf{1982d}] and note 40.

\noindent [\textbf{94 B}] There is a Lindel\"{o}f space with a closed, discrete infinite subset. Brunner\index{Author}{Brunner} [\textbf{1982d}] and note 40.

\noindent \textbf{114.} Every A-bounded $T_{2}$ topological space is weakly Loeb. (A-\emph{bounded} means amorphous subsets are relatively compact. \emph{Weakly Loeb} means the set of nonempty closed subsets has a multiple choice function.) Brunner [\textbf{1985a}].

\noindent [\textbf{126 C}] $FSCB\ +$ form 131. \emph{FSCB} is ``Every separable first countable space $(X,\,T)$ has a denumerable $\pi$-base.'' Keremedis\index{Author}{Keremedis} [\textbf{1996d}] and note 132.

\noindent [\textbf{126 E}] If $(X,\,T)$ is a topological space, then for every family $D=\{D_{i}:i\in \omega\}$ of dense open sets of $X$, there is a countable dense set $S\subseteq X$ such that for all $i\in\omega$ and for all but finitely many $s\in S,\,s\in D_{i}$. Keremedis [\textbf{1996d}] and note 132.

\noindent [\textbf{126 G}] DUN$(\omega)$: The countable disjoint union of normal topological spaces is normal. Howard/Keremedis/Rubin/Rubin [\textbf{1997a}] and note 141.

\noindent [\textbf{126 H}] DUU$(\omega)$: If $C$ is a countable collection of pairwise disjoint topological spaces each of which satisfies Urysohn's lemma, then the disjoint union of the spaces in $C$ satisfies Urysohn's lemma. Howard/Keremedis/Rubin/Rubin\index{Author}{Howard/Keremedis/Rubin/Rubin} [\textbf{1997a}] and note 141.

\noindent [\textbf{126 I}] DUT$(\omega)$: If $C$ is a countable collection of pairwise disjoint topological spaces each of which satisfies the Tietze-Urysohn extension theorem, then the disjoint union of the spaces in $C$ satisfies the Tietze-Urysohn extension theorem. Howard/Keremedis/Rubin/Rubin [\textbf{1997a}] and note 141.

\noindent [\textbf{126 J}] $F_{\sigma}$ subsets of paracompact spaces are paracompact. Howard/Keremedis/Rubin/Rubin [\textbf{1997b}] and note~141.

\noindent [\textbf{126 K}] $F_{\sigma}$ subsets of metacompact spaces are metacompact. Howard/Keremedis/Rubin/Rubin [\textbf{1997b}] and note~141.

\noindent \textbf{128.} Aczel's Realization Principle:\index{Subject}{Aczel's Realization Principle} On every infinite set there is a Hausdorff topology with an infinite set of non-isolated points. Brunner [\textbf{1984f}].

\noindent [\textbf{144 A}] Every $p$ space is discrete. (A $p$ \emph{space} is a $T_{2}$ space in which the intersection of any well orderable family of open sets is open). Brunner [\textbf{1984c}].

\noindent \textbf{146.} $A(F,\,A1)$: For every $T_{2}$ topological space $(X,\,T)$, if $X$ is a continuous finite to one image of an A1 space then $(X,\,T)$ is an A1 space. ($(X,\,T)$ is A1 means if $U\subseteq T$ covers $X$ then $\exists f:X\rightarrow U$ such that $(\forall x\in X)(x\in f(x))$.) Brunner [\textbf{1983d}] and note 26.

\noindent \textbf{147.} $A(D2)$: Every $T_{2}$ topological space $(X,\,T)$ can be covered by a well ordered family of discrete sets. Brunner [\textbf{1983d}] and note 26.

\noindent \textbf{148.} $A(W2,\,B1)$: For every $T_{2}$ topological space $(X,\,T)$, if $X$ is well ordered, then $X$ has a well ordered base. Brunner [\textbf{1983d}] and note 26.

\noindent \textbf{149.} $A(F)$: Every $T_{2}$ topological space is a continuous, finite to one image of an A1 space. Brunner [\textbf{1983d}] and note 26.

\noindent \textbf{301.} Any continuous surjection between Boolean spaces has an irreducible restriction to a closed subset of its domain. Bell\index{Author}{Bell} \cite{bib:1988b} and note 114.

\noindent \textbf{360.} Every second countable topological space is Lindel\"{o}f. G. Moore\index{Author}{G. Moore} [\textbf{1982}] p 236.

\noindent \textbf{371.} There is an infinite, compact, Hausdorff, extremally disconnected topological space. Morillon\index{Author}{Morillon} [\textbf{1993}] and note~114.

\noindent \textbf{375.} Tietze-Urysohn Extension Theorem: If $(X,\,T)$ is a normal topological space, $A$ is closed in $X$, and $f:A\rightarrow[0,\,1]$ is continuous, then there exists a continuous function $g:X\rightarrow[0,\, 1]$ which extends $f$. Urysohn\index{Author}{Urysohn} [\textbf{1924}].

\textbf{IV. Linearly Orderable Spaces}

\noindent [\textbf{34 B}] $\omega_{1}$ with the order topology is not paracompact. Good/Tree\index{Author}{Good/Tree} [\textbf{1995}] and note 107.

\noindent [\textbf{34 C}] $\omega_{1}$ with the order topology has the property that every infinite subset has a limit point in $\omega_{1}$. Good/Tree [\textbf{1995}] and note 107.

\noindent [\textbf{34 D}] $\omega_{1}$ with the order topology does not contain a countable discrete family of open subsets. Good/Tree [\textbf{1995}] and note 107.

\noindent [\textbf{34 E}] $\omega_{1}$ with the order topology is not Lindel\"{o}f. Good/Tree [\textbf{1995}] and note 107.

\noindent \textbf{118.} Every linearly orderable topological space is normal. Birkhoff\index{Author}{Birkhoff} [\textbf{1967}] p 241.

\noindent [\textbf{118 A}] For each complete linear order, the set of non-empty convex subsets has a choice function, van Douwen\index{Author}{Douwen} [\textbf{1985}].

\noindent [\textbf{118 B}] For each complete linear order, each pairwise disjoint collection of non-empty, convex, open subsets has a choice function, van Douwen [\textbf{1985}].

\noindent [\textbf{118 C}] Every linearly orderable space is collectionwise normal. van Douwen [\textbf{1985}] and note 71.

\noindent [\textbf{118 D}] Every linearly orderable space is collectionwise Hausdorff. van Douwen [\textbf{1985}] and note 71.

\noindent [\textbf{118 E}] Every linearly orderable space satisfies Urysohn's lemma, van Douwen [\textbf{1985}] and note 141.

\noindent [\textbf{118 F}] Every linearly orderable space satisfies the Tietze-Urysohn extension theorem, van Douwen [\textbf{1985}] and note~141.

\noindent [\textbf{118 J}] If $(X,\,\leq)$ is a complete linear order, then the order topology on $X$ is hereditarily normal. Morillon \cite{bib:1988h}, van Douwen [\textbf{1985}], and note 71.

\noindent [\textbf{118 K}] If $(X,\,\leq)$ is a linear order, then the order topology on $X$ is highly separated. Morillon \cite{bib:1988h} and note 71.

\noindent [\textbf{118 L}] If $(X,\,\leq)$ is a complete linear order, then the order topology on $X$ is highly separated. Morillon \cite{bib:1988h} and note 71.

\noindent [\textbf{118 O}] If $(X,\,\leq)$ is a linear order, then the order topology on $X$ is effectively normal. Morillon \cite{bib:1988h} and note 71.

\noindent [\textbf{118 R}] If $(X,\,\leq)$ is a linear order, then the order topology on $X$ is monotonely normal. Morillon [\textbf{1987}] and note 71.

\noindent [\textbf{118 U}] Every linearly ordered space is monotonically normal. (A space, $(X,\,<)$ is \emph{monotonically normal} if there is an operator $V$ which assigns to each $x\in X$ and basic open neighborhood $U$ of $x$ a basic open neighborhood $V(x,\,U)$ of $x$ such that $V(x,\,U) \cap V(x^{\prime},\,U^{\prime})\neq\emptyset$ implies $x\in U^{\prime}$ or $x^{\prime}\in U$.) Good/Tree [\textbf{1995}].

\noindent \textbf{336.} $\omega_{1}$ with the order topology is not paracompact. Good/Tree\index{Author}{Good/Tree} [\textbf{1995}] and note 43.

\textbf{V. Compact Spaces}

\noindent [\textbf{$\boldsymbol{0}$ O}] A Lindel\"{o}f space with a dense, Dedekind finite subset is compact. Brunner\index{Author}{Brunner} [\textbf{1982d}] and note 94.

\noindent [\textbf{$\boldsymbol{0}$ P}] A Lindel\"{o}f space with a Dedekind finite basis is compact. Brunner [\textbf{1982d}] and note 94.

\noindent [\textbf{1 AO}] Any continuous surjection between compact spaces, where the range space is Hausdorff, has an irreducible restriction to a closed subset of its domain. Bell\index{Author}{Bell} \cite{bib:1988b} and note 114.

\noindent [\textbf{1 BV}] If $(X,\,T)$ is a $T_{4}$ topological space and $G=\{G_{i}:i\in K\},\,|G_{i}|\geq 2$, is a locally finite family of compact sets with non-empty pairwise disjoint interiors, then there is a continuous real valued function $f$ on $X$ which assumes at least two distinct values on each $G_{i}$. Keremedis [\textbf{1995}] and note 43.

\noindent [\textbf{8 H}] A disjoint union of a countable set of compact $T_{2}$ spaces is Lindel\"{o}f. Keremedis\index{Author}{Keremedis} [\textbf{1996d}] and note 132.

\noindent [\textbf{14 L}] Stone-\v{C}ech Compactification Theorem: For every topological space $X$ there is a pair $(Y,\,j)$ such that $Y$ is a compact topological space and $j$ is a continuous surjection from $X$ onto $Y$ such that for all topological spaces $Z$ such that $Z$ is compact and completely regular and for all continuous $f:X\rightarrow Z$, there is a unique continuous $h:Y\rightarrow Z$ such that $f=h\circ j$. Rubin/Scott\index{Author}{Rubin/Scott} [\textbf{1954}] and note 71.

\noindent [\textbf{14 AA}] Alexander's Subbase Lemma:\index{Subject}{Alexander's subbase lemma} A topological space is compact if and only if there is a subbase $S$ for the topology such that every cover by elements of $S$ has a finite subcover. Kelley\index{Author}{Kelley} [\textbf{1955}] p 143, A.A. Lewis\index{Author}{Lewis} [\textbf{1983}], Wolk\index{Author}{Wolk} [\textbf{1967}], and note 34.

\noindent [\textbf{14 AG}] Topological Rado Lemma: Let $\{T_{i}:i\in I\}$ be a family of compact Hausdorff spaces and suppose for each finite $F\subseteq I$, there is a set $\Phi(F)$ of partial choice functions on $I$ whose domains include $F$ and such that (a) $F_{1}\subseteq F_{2}$ implies $\Phi(F_{2})\subseteq\Phi(F_{1})$ and (b) the set $\{\phi|F:\phi\in\Phi(F)\}$ is closed in $\prod\{T_{i}:i\in F\}$ in the product topology. Then there is a choice function $f$ on $I$ satisfying for any finite $F\subseteq I,\,\exists\phi\in\Phi(F)$ such that $f$ and $\phi$ coincide on $F$. Rav\index{Author}{Rav} [\textbf{1977}].

\noindent [\textbf{14 AH}] Principle of Consistent Choices: Let $\{T_{i}:i\in I\}$ be a family of compact Hausdorf spaces and suppose $R$ is a symmetric binary relation defined for all pairs of elements belonging to distinct $T_{i}$ and such that for all $i\neq j$ in $I$, the set $E(i,\,j)=\{(x,\,y):x\,R\,y\ \mathrm{and}\ x\in T_{i}\ \mathrm{and}\ y\in T_{j}\}$ is closed in $T_{i}\times T_{j}$. Furthermore suppose that for every finite subset $F\subseteq I$, with at least
two elements, $\exists\phi\in\prod\{T_{i}:i\in F\}$ such that $\phi(i)\,R\,\phi(j)$ for all $i\neq j$ in $F$. Then $\exists f\in\prod\{T_{i}:i\in I\}$ such that $f(i)\,R\,f(j)$ for all $i\neq j$ in $I$. Rav [\textbf{1977}].

\noindent [\textbf{14 CC}] A projective limit of compact topological spaces is compact. Morillon\index{Author}{Morillon} \cite{bib:1988h} and note 71.

\noindent [\textbf{14 CD}] A projective\index{Subject}{Goodner's theorem} limit of a directed projective system of compact topological spaces is compact. Morillon \cite{bib:1988h} and note 71.

\noindent [\textbf{14 CE}] A projective limit of a directed projective system of compact nonempty topological spaces is non-empty. Morillon \cite{bib:1988h} and note 71.

\noindent [\textbf{14 CG}] A projective limit of a directed projective system of finite sets is compact. Morillon \cite{bib:1988h} and note 71.

\noindent [\textbf{50 E}] Form 14 + form 297. Banaschewski\index{Author}{Banaschewski} \cite{bib:1988a} and Bell \cite{bib:1988b}.

\noindent [\textbf{50 J}] Form 299 + form 14. Bell\index{Author}{Bell} \cite{bib:1988b}.

\noindent [\textbf{50 K}] Form 300 + form 14. Bell \cite{bib:1988b}.

\noindent [\textbf{50 M}] Form 302 + form 14. Bell \cite{bib:1988b}.

\noindent [\textbf{64 C}] $C$ spaces are compact. (A space is \emph{a} $C$ \emph{space} if each open covering has an amorphous refinement.) Brunner\index{Author}{Brunner} [\textbf{1984b}].

\noindent [\textbf{67 H}] Let $(X,\,T)$ be a $T_{4}$ topological space and let $\{G_{i}:i\in K\}$ be a locally finite family of pairwise disjoint, non-empty, open subsets of $X$. Then there exists a continuous real valued function $f:X\rightarrow \mathbb{R}$ such that for each $i\in K,\,f$ vanishes on the boundary of $G_{i}$ and $f(x_{i})=1$ for some $x_{i}\in G_{i}$. Keremedis\index{Author}{Keremedis} [\textbf{1995}] and note 43.

\noindent [\textbf{106 D}] For every compact $T_{2}$ space $(X,\,T)$ and every family $D=\{D_{i}:i\in \omega\}$ of dense open sets of $X$ there is a regular filter base $\mathcal{F}\subseteq T$ (that is, if $F,\,G\in \mathcal{F}$, then there exists $Q\in \mathcal{F}$ such that $\overline{Q}\subseteq F\cap G$) such that $\forall i\in \omega$, for all but finitely many $F\in \mathcal{F},\,\overline{F}\subset D_{i}$. Fossy/Morillon\index{Author}{Fossy/Morillon} [\textbf{1995}].

\noindent \textbf{116.} Every compact $T_{2}$ space is weakly Loeb. (See form 114 for the definition of weakly Loeb.) Brunner [\textbf{1985a}].

\noindent \textbf{125.} There does not exist an infinite, compact connected $p$ space. (A $p$ \emph{space} is a $T_{2}$ space in which the intersection of any well orderable family of open sets is open.) Brunner [\textbf{1984c}].

\noindent \textbf{145.} Compact $P_{0}$-spaces are Dedekind finite. (A $P_{0}$-\emph{space} is a topological space in which the intersection of a countable collection of open sets is open.) Brunner [\textbf{1984c}] and note 94.

\noindent \textbf{155.} $LC$: There are no non-trivial L\"{a}uchli continua. (A \emph{L\"{a}uchli continuum} is a strongly connected continuum. \emph{Continuum} $\equiv$ compact, connected, Hausdorff space; and \emph{strongly connected} $\equiv$ every continuous real valued function is constant.) Brunner [\textbf{1983a}] and L\"{a}uchli\index{Author}{Lauchli@L\"{a}uchli} [\textbf{1962}].

\noindent \textbf{156.} Theorem of Gelfand and Kolmogoroff:\index{Subject}{Gelfand-Kolmogoroff theorem} Two compact $T_{2}$ spaces are homeomorphic if their rings of real valued continuous functions are isomorphic. Brunner [\textbf{1983a}] and Gelfand/Kolmogoroff\index{Author}{Gelfand/Kolmogoroff} [\textbf{1939}].

\noindent \textbf{157.} Theorem of Goodner\index{Author}{Goodner}: A compact $T_{2}$ space is extremally disconnected (the closure of every open set is open) if and only if each non-empty subset of $C(X)$ (set of continuous real valued functions on $X$) which is pointwise bounded, has a supremum. Goodner [\textbf{1950}].

\noindent [\textbf{225 A}] There exists\index{Subject}{Baire's classification} a compact, extremally disconnected topological space having a Dedekind infinite set of clopen subsets. Morillon [\textbf{1993}], notes 94 and 114.

\noindent [\textbf{225 B}] There exists\index{Author}{Feferman/Levy} a compact, $\sigma$-extremally disconnected topological space having a Dedekind infinite set of clopen subsets. Morillon [\textbf{1993}], notes 94 and 114.

\noindent \textbf{297.} Extremally\index{Author}{Jech} disconnected compact Hausdorff spaces are projective in the category of all compact Hausdorff spaces. Banaschewski\index{Author}{Banaschewski} \cite{bib:1988a}, Bell\index{Author}{Bell} \cite{bib:1988b}, notes 60 and~114.

\noindent \textbf{298.} Every compact Hausdorff space has a Gleason\index{Author}{Gleason} cover. Bell \cite{bib:1988b}, Gleason [\textbf{1958}], and note 114.

\noindent \textbf{299.} Any extremally disconnected compact Hausdorff space is projective in the category of Boolean topological spaces. Bell \cite{bib:1988b}, notes 60 and 114.

\noindent \textbf{300.} Any continuous surjection between extremally disconnected compact Hausdorff spaces has an irreducible restriction to a closed subset of its domain. Bell \cite{bib:1988b} and note 114.

\noindent \textbf{302.} Any continuous surjection between compact Hausdorff spaces has an irreducible restriction to a closed subset of its domain. Bell \cite{bib:1988b} and note 114.

\noindent \textbf{304.} There does not exist a $T_{2}$ topological space $X$ such that every infinite subset of $X$ contains an infinite compact subset. Brunner\index{Author}{Brunner} [\textbf{1985c}].

\noindent \textbf{331.} If $(X_{i})_{i\in I}$ is a family of compact non-empty topological spaces then there is a family $(F_{i})_{i\in I}$ such that $\forall i\in I,\,F_{i}$ is an irreducible closed subset of $X_{i}$. Morillon \cite{bib:1988h} and note 71.

\textbf{VI. Properties of} $\mathbb{R}$

\noindent [\textbf{$\boldsymbol{0}$ AB}] If $X\subseteq \mathbb{R}$ is perfect, then $|X|=2^{\aleph_{0}}$. Truss\index{Author}{Truss} [\textbf{1974b}] and note 28.

\noindent [\textbf{$\boldsymbol{0}$ AC}] There are $2^{\aleph_{0}}$ perfect subsets of $\mathbb{R}$. Truss [\textbf{1974b}] and note 28.

\noindent [\textbf{6 A}] Every uncountable subset of $\mathbb{R}$ contains a condensation point. G. Moore\index{Author}{G. Moore} [\textbf{1982}] p 205 and 324, Sierpi\'{n}ski\index{Author}{Sierpinski@Sierpi\'{n}ski} [\textbf{1918}].

\noindent [\textbf{6 B}] Every uncountable subset of $\mathbb{R}$ has two condensation points. G. Moore [\textbf{1982}] p 205 and 324, Sierpi\'{n}ski [\textbf{1918}].

\noindent \textbf{11.} $\mathbb{R}$ is hereditarily Lindel\"{o}f\index{Author}{Lindelof@Lindel\"{o}f}. Lindel\"{o}f [\textbf{1905}].

\noindent \textbf{12.} $\mathbb{R}$ is Lindel\"{o}f. Young\index{Author}{Young} [\textbf{1903}].

\noindent \textbf{13.} Every Dedekind finite subset of $\mathbb{R}$ is finite. G. Moore [\textbf{1982}] and note 94.

\noindent [\textbf{13 A}] $\mathbb{R}$ has no dense Dedekind finite subset. Brunner [\textbf{1982d}] and note 94.

\noindent \textbf{19.} A real function is analytically representable if and only if it is in Baire's classification. G. Moore\index{Author}{G. Moore} [\textbf{1982}]~(2.3.1).

\noindent \textbf{35.} The union of countably many meager subsets\index{Author}{Davies} of $\mathbb{R}$ is meager. (Meager sets are the same as sets of the first category.) Feferman/Levy\index{Author}{Feferman/Levy} [\textbf{1963}], Jech\index{Author}{Jech} \cite{bib:1973b} p 7 prob 1.7, and note 28.

\noindent \textbf{36.} If $A\subseteq \mathbb{R}^{n}$ and $A \bigcap B$ is countable for every bounded $B$ then $A$ is countable. G. Moore [\textbf{1982}] p 36.

\noindent \textbf{37.} Lebesgue measure is countably additive. Sierpi\'{n}ski~[\textbf{1916}].

\noindent \textbf{38.} $\mathbb{R}$ is not the union of a countable family of countable sets. Feferman/Levy [\textbf{1963}] and note 3.

\noindent \textbf{73.} For every $A\subseteq \mathbb{R}$ and $x\in \mathbb{R}$ the following definitions are equivalent:

(1) $x$ is in the closure of $A$ iff every neighborhood of $x$ intersects $A$.

(2) $x$ is in the closure of $A$ iff there is a sequence $\{x_{n}\}\subseteq A$ such that $\lim x_{n}=x$. Jech \cite{bib:1973b} p 21.

\noindent [\textbf{73 A}] For all $A\subseteq \mathbb{R},\,x\in A$ and $f:A\rightarrow \mathbb{R}$ the following are equivalent:

(1) $(\forall\epsilon>0)(\exists\delta>0)(\forall y\in A)(|y-x|<\delta$ implies $|f(y)-f(x)|<\epsilon)$

(2) Whenever $\{x_{n}\}\subseteq A$ and $\lim x_{n}=x$ then $\lim f(x_{n})=f(x)$.
Notes 4, 5 and 6.

\noindent [\textbf{73 B}] $(\forall f:\mathbb{R} \rightarrow \mathbb{R})(\forall x\in \mathbb{R})$ the following are equivalent:

(1) $(\forall\epsilon>0)(\exists\delta>0)((\forall y\in \mathbb{R}\rangle(|y-x|<\delta\rightarrow|f(y)-f(x)|<\epsilon)$

(2) Whenever $\lim x_{n}=x$, then $\lim f(x_{n})=f(x)$.
Note~10.

\noindent \textbf{74.} For every $A\subseteq \mathbb{R}$ the following are equivalent:

(1) $A$ is closed and bounded.

(2) Every sequence $\{x_{n}\}\subseteq A$ has a convergent subsequence with limit in A. Jech \cite{bib:1973b} p 21.

\noindent \textbf{79.} $\mathbb{R}$ can be well ordered. Hilbert\index{Author}{Hilbert} [\textbf{1900}] p 263.

\noindent [\textbf{79 A}] $C(\infty,\,\infty,\,\mathbb{R})$: Every set whose elements are non-empty subsets of $\mathbb{R}$ has a choice function. Schechter\index{Author}{Schechter} [\textbf{1996a}].

\noindent \textbf{92.} Every subset of $\mathbb{R}$ is separable. Jech \cite{bib:1973b} p 142.

\noindent \textbf{93.} There is a non-measurable subset of $\mathbb{R}$. Sierpi\'{n}ski\index{Author}{Sierpinski@Sierpi\'{n}ski} [\textbf{1918}].

\noindent \textbf{130.} $\mathcal{P}(\mathbb{R})$ is well orderable. Brunner\index{Author}{Brunner} [\textbf{1984f}].

\noindent [\textbf{137 A($k$)}] For all $X$ and $Y\subseteq \mathbb{R}$ form 137$(k)$ holds. Truss\index{Author}{Truss} [\textbf{1984}].

\noindent \textbf{142.} $\neg PB$: There is a set of reals without the property of Baire. Jech \cite{bib:1973b} p 7 and note 28.

\noindent \textbf{169.} There is an uncountable subset of $\mathbb{R}$ without a perfect subset. Banaschewski/Moore\index{Author}{Banaschewski/Moore} [\textbf{1990}], Bernstein\index{Author}{Bernstein} [\textbf{1908}], and note 28.

\noindent [\textbf{170 A}] There is a function that assigns to each denumerable set $D$ of reals a real number not in $D$. Sierpi\'{n}ski~[\textbf{1954}].

\noindent [\textbf{170 B}] There is an uncountable, well orderable set of reals. Raisonnier\index{Author}{Raisonnier} [\textbf{1982}] and note 3.

\noindent \textbf{197.} $\mathbb{R}^{2}$ is the union of three sets $C$ with the property that for all $x\in C$ there is a straight line $L$ such that $L\cap C=\{x\}$. Davies [\textbf{1978}].

\noindent \textbf{203.} $C(\mathrm{disjoint},\subseteq \mathbb{R})$: Every partition of $\mathcal{P}(w)$ into non-empty subsets has a choice function. Truss\index{Author}{Truss} [\textbf{1978}].

\noindent \textbf{211.} $DCR$: Dependent\index{Author}{Sierpinski@Sierpi\'{n}ski} Choice\index{Author}{Hamel} for Relations on $\mathbb{R}$: If $R\subseteq \mathbb{R} \times \mathbb{R}$ satisfies $(\forall x\in \mathbb{R})(\exists y\in \mathbb{R})(x\ R\ y)$ then there is a sequence $\langle x(n): n\in \omega\rangle$ of real numbers such that $(\forall n\in \omega)(x(n)\,R\,x(n+1))$. Mathias\index{Author}{Mathias} [\textbf{1977a}].

\noindent \textbf{212.} $C(2^{\aleph_{0}},\,\subseteq \mathbb{R})$: If $R$ is a relation on $\mathbb{R}$ such that for all $x\in \mathbb{R}$, there is a $y\in \mathbb{R}$ such that $x\,R\,y$, then there is a function $f:\mathbb{R} \rightarrow \mathbb{R}$ such that for all $x\in \mathbb{R},\,x\,R\,f(x)$. Mathias [\textbf{1977a}].

\noindent \textbf{224.} There is a partition of the real line into $\aleph_{1}$ Borel sets $\{B_{\alpha}:\alpha<\aleph_{1}\}$ such that for some $\beta<\aleph_{1},\,\forall\alpha<\aleph_{1},\,B_{\alpha}\in G_{\beta}$. ($G_{\beta}$ for $\beta<\aleph_{1}$ is defined by induction, $G_{0}=\{A: A\ \text{is an open subset of}\ \mathbb{R}\}$ and for $\beta>0,\,G_{\beta}=\{\bigcup_{i=0}^{\infty}A_{i}:(\forall i\in\omega)(\exists\xi<\beta)(A_{i}\in G_{\xi})\}$ if $\beta$ is even and $G_{\beta}=\{\bigcap_{i=0}^{\infty}A_{i}:(\forall i\in \omega)(\exists\xi<\beta)(A_{i}\in G_{\xi})\}$ if $\beta$ is odd.) Stern\index{Author}{Stern} [\textbf{1979}].

\noindent \textbf{251.} The additive groups $(\mathbb{R},\,+)$ and $(\mathbb{C},\,+)$ are isomorphic. Ash [\textbf{1975}].

\noindent \textbf{252.} The additive groups of $\mathbb{Q} \oplus \mathbb{R}$ and $\mathbb{R}$ are isomorphic. Ash\index{Author}{Ash} [\textbf{1975}].

\noindent \textbf{272.} There is an $X\subseteq \mathbb{R}$ such that neither $X$ nor $\mathbb{R} -X$ has a perfect subset. Truss [\textbf{1974b}] and note 28.

\noindent \textbf{273.} There is a subset of $\mathbb{R}$ which is not Borel. Truss [\textbf{1974b}].

\noindent \textbf{305.} There are $2^{\aleph_{0}}$ Vitali equivalence classes\index{Subject}{Vitali equivalence classes}. (\emph{Vitali equivalence classes} are equivalence classes of the real numbers under the relation $x\equiv y\leftrightarrow(\exists q\in\mathbb{Q})(x-y=q)$.). Kanovei\index{Author}{Kanovei} [\textbf{1991}].

\noindent \textbf{306.} The set of Vitali equivalence classes is linearly orderable. (\emph{Vitali equivalence classes} are equivalence classes of the real numbers under the relation $x\equiv y\leftrightarrow(\exists q\in \mathbb{Q})(x-y=q)$.). Kanovei [\textbf{1991}].

\noindent \textbf{307.} If $m$ is the cardinality of the set of Vitali equivalence classes, then $H(m)= H(2^{\aleph_{0}})$, where $H$ is Hartogs aleph function and the \emph{Vitali equivalence classes} are equivalence classes of the real numbers under the relation $x\equiv y\leftrightarrow(\exists q\in\mathbb{Q})(x-y=q)$. Kanovei [\textbf{1991}].

\noindent \textbf{309.} The Banach-Tarski Paradox:\index{Subject}{Banach-Tarski paradox} There are three finite partitions $\{P_{1},\,\ldots,\,P_{n}\},\,\{Q_{1},\, \ldots,\,Q_{r}\}$ and $\{S_{1},\,\ldots,\,S_{n},\,T_{1},\,\ldots,\,T_{r}\}$ of $B^{3}=\{x\in \mathbb{R}^{3}:|x|\leq 1\}$ such that $P_{i}$ is congruent to $S_{i}$ for $1\leq i\leq n$ and $Q_{i}$ is congruent to $T_{i}$ for $1\leq i\leq r$. Banach/Tarski\index{Author}{Banach/Tarski} [\textbf{1924}].

\noindent \textbf{361.} In $\mathbb{R}$, the union of a denumerable number of analytic sets is analytic. G. Moore\index{Author}{G. Moore} [\textbf{1982}], pp 181 and 325.

\noindent \textbf{362.} In $\mathbb{R}$, every Borel set is analytic. G. Moore [\textbf{1982}], pp 181 and 325.

\noindent \textbf{363.} There are exactly $2^{\aleph_{0}}$ Borel sets in $\mathbb{R}$. G. Moore [\textbf{1982}], p 325.

\noindent \textbf{364.} In $\mathbb{R}$, there is a measurable set that is not Borel. G. Moore [\textbf{1982}], p 325.

\noindent \textbf{366.} There is a discontinuous function $f:\mathbb{R} \rightarrow \mathbb{R}$ such that for all real $x$ and $y,\,f(x+y)=f(x)+f(y)$. Hamel [\textbf{1905}].

\noindent \textbf{367.} There is a Hamel\index{Author}{Jech} basis for $\mathbb{R}$ as a vector space over $\mathbb{Q}$. Hamel [\textbf{1905}].

\noindent \textbf{368.} The set of all denumerable subsets of $\mathbb{R}$ has power $2^{\aleph_{0}}$. Sierpi\'{n}ski [\textbf{1917}].

\noindent \textbf{369.} If $\mathbb{R}$ is partitioned into two sets, at least one of them has cardinality $2^{\aleph_{0}}$. Luzin/Sierpi\'{n}ski\index{Author}{Luzin/Sierpinski@Luzin/Sierpi\'{n}ski} [\textbf{1917}].

\textbf{VII. Metric Spaces}

\noindent [\textbf{$\boldsymbol{0}$ L}] Caristi's Fixed Point Theorem: If $(X,\,\rho)$ is a complete metric space and $\phi:X\rightarrow \mathbb{R}$ is bounded above and upper semi-continuous then in the Br{\o}ndsted ordering $(x\leq y\,\mathrm{iff}\, \rho(x,\,y)\leq\phi(y)-\phi(x))$ every $f:X\rightarrow X$ satisfying $\forall t\in X,\,t\leq f(t)$ has a fixed point. Caristi\index{Author}{Caristi} [\textbf{1976}], Manka\index{Author}{Manka} [\textbf{1988a}], and note 38.

\noindent [\textbf{$\boldsymbol{0}$ AK}] Every separable metric space is second countable.

\noindent [\textbf{1 I}] In every metric space $(X,\,d)$ there is an $\epsilon$-lattice for some $\epsilon>0$ which is not a strict upper bound of $d$. Brunner\index{Author}{Brunner} [\textbf{1984e}] and note 17.

\noindent [\textbf{1 BK}] Every open cover $\mathcal{U}$ of a metric space $(X,\,d)$ can be written as a well ordered union $\bigcup\{U_{\alpha}:\alpha\in\gamma\}$ where $\gamma$ is an ordinal and each $U_{\alpha}$ is pairwise disjoint. Howard/Keremedis/Rubin/Stanley\index{Author}{Howard/Keremedis/Rubin/Stanley} [\textbf{1997}].

\noindent [\textbf{1 CE}] Every open cover $\mathcal{U}$ of a metric space $(X,\,d)$ can be written as a well ordered union $\bigcup\{U_{\alpha}:\alpha\in\gamma\}$ where $\gamma$ is an ordinal and each $U_{\alpha}$ is discrete. Howard/Keremedis/Rubin/Stanley [\textbf{1997}] and note 141.

\noindent [\textbf{10 G}] For all metric spaces $X$ and $Y$ and continuous functions $f$ from $X$ onto $Y$ such that $f^{-1}(y)$ is finite for all $y\in Y$, if $X$ has a dense Dedekind finite subset, then so does $Y$. Brunner [\textbf{1982d}] and note 94.

\noindent [\textbf{43 F}] Baire Category Theorem for Sequentially Complete Metric Spaces: Every sequentially complete metric space is Baire. Blair\index{Author}{Blair} [\textbf{1977}], Brunner [\textbf{1983c}], Goldblatt\index{Author}{Goldblatt} \cite{bib:1985m} and note 28.

\noindent [\textbf{43 J}] Baire Category Theorem for Cantor Complete Metric Spaces: Every Cantor complete metric space is Baire. Brunner [\textbf{1983c}] and note 28.

\noindent [\textbf{43 M}] If $(X,\,\rho)$ is a complete metric space and $\phi:X\rightarrow \mathbb{R}$ is bounded above and upper semi-continuous then in the Br{\o}ndsted ordering, $(x\leq y\,\mathrm{iff}\,\rho(x,\,y)\leq \phi(y)-\phi(x)$), there is a maximal element. Brunner [\textbf{1987a}] and note 38.

\noindent [\textbf{64 D}] Metric $C$ spaces are limited amorphous. (A space is $C$ if every open covering has an amorphous refinement. A space is \emph{limited amorphous} if each amorphous subset is relatively compact.) Brunner [\textbf{1984b}].

\noindent [\textbf{67 U}] Every open cover $\mathcal{U}$ of a metric space $(X,\,d)$ can be written as a well ordered union $\bigcup\{U_{\alpha}:\alpha\in\gamma\}$ where $\gamma$ is an ordinal and each $U_{\alpha}$ is locally finite. Howard/Keremedis/Rubin/Stanley [\textbf{1997}] and note 141.

\noindent [\textbf{67 V}] Every open cover $\mathcal{U}$ of a metric space ($X$, \emph{d}) can be written as a well ordered union $\bigcup\{U_{\alpha}:\alpha\in\gamma\}$ where $\gamma$ is an ordinal and each $U_{\alpha}$ is point finite. Howard/Keremedis/Rubin/Stanley [\textbf{1997}] and note 141.

\noindent \textbf{75.} Every subspace of a separable metric space is separable. Jech \cite{bib:1973b} p 21.

\noindent [\textbf{76 A}] Every open cover $\mathcal{U}$ of a metric space can be written as a well ordered union $\bigcup\{U_{\alpha}:\alpha\in\gamma\}$ where $\gamma$ is an ordinal and each $U_{\alpha}$ is locally countable. Howard/Keremedis/Rubin/Stanley [\textbf{1997}] and note 141.

\noindent [\textbf{76 B}] Every open cover $\mathcal{U}$ of a metric space can be written as a well ordered union $\bigcup\{U_{\alpha}:\alpha\in\gamma\}$ where $\gamma$ is an ordinal and each $U_{\alpha}$ is point countable. Howard/Keremedis/Rubin/Stanley\index{Author}{Howard/Keremedis/Rubin/Stanley} [\textbf{1997}] and note 141.

\noindent \textbf{173.} MPL: Metric spaces are para-Lindel\"{o}f. Howard/Keremedis/Rubin/Stanley [\textbf{1997}] and note 141.

\noindent [\textbf{173 A}] MML: Metric spaces are meta-Lindel\"{o}f. Howard/Keremedis/Rubin/ Stanley [\textbf{1997}] and note 141.

\noindent \textbf{232.} Every metric space $(X,\,d)$ has a $\sigma$-point finite base. Howard/Keremedis/ Rubin/Stanley [\textbf{1997}] and note 141.

\noindent [\textbf{232 A}] Every metric space has a $\gamma$-point finite base, for some ordinal $\gamma$. Howard/Keremedis/Rubin/Stanley [\textbf{1997}] and note 141.

\noindent [\textbf{232 B}] Every metric space has a $\sigma$-locally finite base. Howard/Keremedis/ Rubin/Stanley [\textbf{1997}] and note 141.

\noindent [\textbf{232 C}] Every metric space has a $\gamma$-locally finite base, for some ordinal $\gamma$. Howard/Keremedis/Rubin/Stanley [\textbf{1997}] and note 141.

\noindent [\textbf{232 D}] Every metric space $(X,\,d)$ has a $\sigma$-discrete base. Howard/Keremedis/ Rubin/Stanley [\textbf{1997}] and note 141.

\noindent [\textbf{232 E}] Every metric space $(X,\,d)$ has a $\gamma$-discrete base, for some ordinal $\gamma$. Howard/Keremedis/Rubin/Stanley [\textbf{1997}] and note 141.

\noindent [\textbf{232 F}] Every metric space $(X,\,d)$ has a $\sigma$-disjoint base. Howard/Keremedis/ Rubin/Stanley [\textbf{1997}] and note 141.

\noindent [\textbf{232 G}] Every metric space $(X,\,d)$ has a $\gamma$-disjoint base for some ordinal $\gamma$. Howard/Keremedis/Rubin/Stanley [\textbf{1997}] and note 141.

\noindent \textbf{280.} There is a complete separable metric space with a subset which does not have the Baire property. Wright\index{Author}{Wright} [\textbf{1973}] and note 28.

\noindent \textbf{337.} Every separable metric space is Lindel\"{o}f. Good/Tree [\textbf{1995}].

\noindent \textbf{338.} Every second countable metric space is Lindel\"{o}f. Good/Tree\index{Author}{Good/Tree} [\textbf{1995}].

\noindent \textbf{339.} Every second countable metric space is separable. Good/Tree [\textbf{1995}].

\noindent \textbf{340.} Every Lindel\"{o}f metric space is separable. Good/Tree [\textbf{1995}].

\noindent \textbf{341.} Every Lindel\"{o}f metric space is second countable. Good/Tree [\textbf{1995}].

\noindent \textbf{381.} DUM: The disjoint union of metrizable spaces is metrizable. Howard/Keremedis/Rubin/Rubin\index{Author}{Howard/Keremedis/Rubin/Rubin} [\textbf{1997b}].

\noindent \textbf{382.} DUMN: The disjoint union of metrizable spaces is normal. Howard/Keremedis/Rubin/Rubin [\textbf{1997b}].

\noindent \textbf{383.} MP: Every metric space is paracompact. Howard/Keremedis/Rubin/Stanley [\textbf{1997}] and note 141.

\noindent [\textbf{383 A}] Every open cover $\mathcal{U}$ of a metric space $(X,\,d)$ has a refinement $\mathcal{V}$ which covers $X$ and which can be written as a well ordered union $\bigcup\{V_{\alpha}:\alpha\in\gamma\}$ where each $V_{\alpha}$ is locally finite. Howard/Keremedis/Rubin/Stanley [\textbf{1997}] and note 141.

\noindent [\textbf{383 B}] MM: Every metric space is metacompact. Howard/Keremedis/Rubin/Stanley\index{Author}{Howard/Keremedis/Rubin/Stanley} [\textbf{1997}] and note 141.

\noindent [\textbf{383 C}] Every open cover $\mathcal{U}$ of a metric space ($X$, \emph{d}) has a refinement $\mathcal{V}$ which covers $X$ and which can be written as a well ordered union $\bigcup\{V_{\alpha}:\alpha\in\gamma\}$ where each $V_{\alpha}$ is point finite. Howard/Keremedis/Rubin/Stanley [\textbf{1997}] and note 141.

%%%%%%%%part03

\part{Models\label{pt03:part03}}

\section*{Table of Contents}

\begin{enumerate}
\item[] Introduction

\item[\S 1:] Cohen Models, Models of ZF

\begin{enumerate}
\item[la.] $\mathcal{M}1$: Cohen's original model.

(14, 15, 31, 60, 118, 128, 163, 165, 191, and 277 are true, but 13,
17, 50, 65, 106, 131, 144, 253, 289, 299, 300, 337, 338, and 339 are
false.)

\item[lb.] $\mathcal{M}1 (\langle(\omega_{1}\rangle)$: Cohen/Pincus Model.

(43, 97, 214, and 295 are true, but 20, 39, 85, 91, 151, 163, 328,
and 329 are false.)

\item[2a.] $\mathcal{M}2$: Feferman's model.

(40, 142, 165, and 204 are true, but 88, 152, 163, 203, 222, and 274
are false.)

\item[2b.] $\mathcal{M}2(\kappa)$: Feferman/Pincus Model.

(222 is true, but 52 and 91 are false.)

\item[2c.] $\mathcal{M}2(\langle(\omega_{2}\rangle)$: Feferman/Truss Model.

(181 is true, but 163 and 203 are false.)

\item[3.] $\mathcal{M}3$: Mathias' model.

(15, 60, and 165 are true, but 49, 91, 99, 144, and 152 are false.)

\item[4.] $\mathcal{M}4$: Pincus' Model I.

(52 is true, but 65, 88, 91, and 358 are false.)

\item[5.] $\mathcal{M}5(\aleph)$: Solovay's Model.

(43 and 307 are true, but 93, 163, 169, 224, 234, 277, 280, 281, 312,
and 318 are false.)

\item[6.] $\mathcal{M}6$: Sageev's Model I.

(3 and 295 are true, but 5, 126, 163, and 330 are false.)

\item[7a.] $\mathcal{M}7$: Cohen's Second Model.

(191 is true, but 18, 91, 340, 341, and 358 are false.)

\item[7b.] $\mathcal{M}7(n)$: Generalization of $\mathcal{M}7$.

($47(n)$ and 91 are false.)

\item[8.] $\mathcal{M} 8$: Apter's Model.

(8, 25, and 91 are false.)

\item[9.] $\mathcal{M} 9$: Feferman/Levy Model.

(38 is false.)

\item[10.] $\mathcal{M}10$: Derrick/Drake Model.

(43, 56, 91, and 275 are false.)

\item[11.] $\mathcal{M}11$: Forti/Honsell Model.

(15, 60, and 165 are true, but 91, 105, 144, 152, and 163 are false.)

\item[12.] $\mathcal{M}12(\aleph)$: Truss' Model I.

(6 and 231 are true, but 30, 34, 106, 152, 169, 273, and 328 are
false.)

\item[13.] $\mathcal{M}13$: Feferman/Solovay Model.

(40 and 165 are true, but 44, 91, 152, and 163 are false.)

\item[14.] $\mathcal{M}14$: Morris' Model.

(30 and 51 are true, but 49, 91, 99, and 144 are false.)

\item[15.] $\mathcal{M}15$: Feferman/Blass Model.

(144 and 253 are true, but 206 is false.)

\item[16.] $\mathcal{M}16(n)$: Monro's Model I.

($15,\,81(n)$, and 91 are false.)

\item[17.] $\mathcal{M}17$: Gitik's Model.

(91, 182, and 315 are false.)

\item[18.] $\mathcal{M}18$: Shelah's Model I.

(43 and 93 are true, but 142 and 163 are false.)

\item[19.] $\mathcal{M}19(\aleph)$: Monro's Model II.

(91 and 160 are false.)

\item[20.] $\mathcal{M}20$: Felgner's Model I.

(30, 60, and 165 are true, but 31, 91, 131, 144, 152, and 328 are false.)

\item[21.] $\mathcal{M}21$: Felgner's Model II.

(8 is true, but 91, 163, and 231 are false.)

\item[22.] $\mathcal{M}22$: Plotkin's Model I.

(9, 91, and 176 are false.)

\item[23.] $\mathcal{M}23$: Hodges' Model.

(91 and 180 are false.)

\item[24.] $\mathcal{M} 24$: Blass' Model.

(91 and 190 are false.)

\item[25.] $\mathcal{M}25$: Freyd's Model.

(40 and 165 are true, but 91, 152, 163, and 181 are false.)

\item[26.] $\mathcal{M}26$: Kanovei's Model I.

($199(n)$ is false.)

\item[27.] $\mathcal{M}27$: Pincus/Solovay Model I.

(43, 52, and 253 are true, but 163, 206, 286, and 371 are false.)

\item[28.] $\mathcal{M}28$: Morris' Model.

(91 and 209 are false.)

\item[29.] $\mathcal{M} 29$: Pincus' Model II.

(200, 214, 290, 291, and 295 are true, but 3, 91, 118, 151, 163, 350, and 357 are false.)

\item[30.] $\mathcal{M}30$: Pincus/Solovay Model II.

(43, 253, and 312 are true, but 163 and 223 are false.)

\item[31.] $\mathcal{M}31$: Szczepaniak's Model.

(91 and 230 are false.)

\item[32.] $\mathcal{M}32$: Sageev's Model II.

(30 and 57 are true, but 9, 91, 126, and 328 are false.)

\item[33.] $\mathcal{M}33$: Plotkin's Model II.

(9, 88, 91, 247, and 248 are false.)

\item[34.] $\mathcal{M} 34(\aleph_{1})$: Pincus' Model III.

(15 and 85 are true, but 91, 163, 213, and 330 are false.)

\item[35.] $\mathcal{M}35(\epsilon)$: David's Model.

($179(\epsilon)$ is true, but 91 and $179(\beta)$ are false.)

\item[36.] $\mathcal{M}36$: Figura's Model.

(170 is true, but 34 and 91 are false.)

\item[37.] $\mathcal{M} 37$: Monro's Model III.

(371 is true, but 64 and 91 are false.)

\item[38.] $\mathcal{M}38$: Shelah's Model II.

(43, 142, and 169 are true, but 93 and 163 are false.)

\item[39.] $\mathcal{M}39(\kappa,\,\lambda)$: Kanovei's model II.

(43 is true, but 91, 163, and 307 are false.)

\item[40.] $\mathcal{M}40(\kappa)$: Pincus' Model IV.

(60, $87(\alpha)$ for $\aleph_{\alpha}<\kappa$, 165, 295, and 345 are true, but 51, 91, 144, 152, 163, 253, and 286 are false.)

\item[41.] $\mathcal{M} 41$ Kanovei's Model III.

(194 is false.)

\item[42.] $\mathcal{M} 42$ Bull's Model.

(39 and 86 are true, but 91, 163, and 321 are false.)

\item[43.] $\mathcal{M}43$: Pincus' Model V.

(43, 62, and 214 are true, but 30, 91, 99, 151, 163, 328, and 329 are false.)

\item[44.] $\mathcal{M}44$: Pincus' Model VI.

(43, 214, and 295 are true, but 15, 91, 118, 151, 163, 328, and 329 are false.)

\item[45.] $\mathcal{M}45$: Pincus' Model VII.

(43, 214, and 295 are true, but 49, 91, 99, 151, 163, 328, and 329

are false.)

\item[46.] $\mathcal{M}46(m,\,M)$: Pincus' Model VIII.

$(43,\,46(M)$, and 214 are true, but $45(n)$, 91, 151, 163, and 329 are false.)

\item[47.] $\mathcal{M}47(n,\,M)$: Pincus' Model IX.

$(43,\,46(\neg M)$, 63, and 214 are true, but $46(n\in K)$, $47(n)$, 91, 163, 329, and 324 are false.)
\end{enumerate}

\item[\S 2:] Fraenkel-Mostowski Models, Models of $\mathrm{ZF}^{0}$

\begin{enumerate}
\item[1.] $\mathcal{N}1$: The Basic Fraenkel Model.

(6, 11, 16, 23, 24, 37, 73, 89, 92, 112, 114, 115, 116, 127, 130, 133,
134, 135, 147, 191, 217, 232, 233, 263, 304, 305, 309, 313, 322, 325,
361, 363, and 380 are true, but 15, 53, 64, 69, 76, 126, 128, 131,
146, 177, 200, 239, 267, 278, 292, and 344 are false.)

\item[2a.] $\mathcal{N}2$: The Second Fraenkel Model.

(6, 11, 37, 65, 73, 92, 99, 130, 163, 191, 305, 313, 334, 361, and 363
are true, but 18, $45(n), 98,128,154,163,164,198,344$, and 358
are false.)

\item[2b.] $\mathcal{N}2(n)$: A generalization of $\mathcal{N}2$.

(6, 11, 37, 67, 73, 92, 130, 191, 305, 313, 361, and 363 are true,
but $46(n\in K)$, 154, 164, 344, and $373(n)$ are false.)

\item[2c.] $\mathcal{N}2^{\ast}(3)$: Howard's variation of $\mathcal{N}2(3)$.

(6, 11, 37, 73, 92, 130, 141, 191, 305, 313, 333, 361, and 363 are
true, but 154, 164, 288, 334, 344, and 358 are false.)

\item[2d.] $\mathcal{N}2(n,\,M)$: Mostowski's variation of $\mathcal{N}2(n)$.

$(6,\,11,\,37,\,46(\neg M)$, 73, 91, 92, 130, 191, 273. 305, 309, 313, 361,
363, 368, and 369 are true, but 15, $46(n\in K)$, $47(n)$, 192, and 344
are false.)

\item[2e.] $\mathcal{N}2(\aleph_{\alpha})$: Jech's Model.

$(6,\,11,\,37,\,44\,(\mathrm{if}\ \aleph_{\alpha}\ \mathrm{is\ regular\ and}\ \alpha>1), 63(\mathrm{if}\,\alpha>0)$, 73, $87(\beta)$
$(\mathrm{if}\ \alpha>\beta)$, 91, 92, 130, 135 (if $\aleph_{\alpha}$ is regular and $\alpha>1$), 191, 273,
305, 309, 313, 361, 363, 368, and 369 are true, but 3,15, $45(n)$, 88,
163, 192, 193, 327, and 344 are false.)

\item[2f.] $\mathcal{N}2$(LO): van Douwen's Model.

(6, 11, 37, 73, 91, 92, 130, 191, 305, 309, 313, 361, and 363 are
true, but 62, 119, and 344 are false.)

\item[3.] $\mathcal{N}3$: Mostowski's Linearly Ordered Model.

(6, 11, 16, 23, 24, 37, 60, 73, 91, 92, 128, 130, 164, 165, 191, 305,
309, 313, 317, 361, 363, 368, 369, and 381 are true, but 15, 76, 84,
97, 118, 125, 126, 131, 147, 155, 156, 157, 200, 253, 290, 295, 296,
304, and 355 are false.)

\item[4.] $\mathcal{N}4$: The Mathias/Pincus Model I.

(6, 11, 37, 64, 73, 90, 92, 130, 191, 305, 309, 313, 361, and 363 are
true, but 83, 89, 114, and 133 are false.)

\item[5.] $\mathcal{N}5$: The Mathias/Pincus Model II (an extension of $\mathcal{N}4$).

(11, 16, 23, 24, 30, 37, 60, 73, 91, 92, 130, 165, 191, 305, 309, 313,
361, and 363 are true, but 49, 84, 99, 118, 131, and 144 are false.)

\item[6.] $\mathcal{N}6$: Levy's Model I.

(6, 11, 37, 73, 92, 130, 191, 218, 305, 313, 361, and 363 are true,
but 154, 164, 171, $308(p)$, 314, 334, 344, and 358 are false.)

\item[7.] $\mathcal{N}7$: L\"{a}uchli's Model I.

(11, 16, 23, 24, 37, 60, 73, 91, 92, 123, 130, 165, 191, 305, 309, 313,
361, and 363 are true, but 30, 99, 133, 152, 192, and 328 are false.)

\item[8.] $\mathcal{N}8$: L\"{a}uchli's Model II.

(6, 11, 37, 73, 91, 92, 130, 191, 305, 309, 313, 361, and 363 are
true, but 15 and 78 (Urysohn's Lemma) are false.)

\item[9.] $\mathcal{N}9$: Halpern/Howard Model.

(3, 6, 11, 37, 63, 73, 91, 92, 130, 191, 305, 309, 313, 361, and 363
are true, but 88, 133, 192, and 344 are false.)

\item[10.] $\mathcal{N}10$: H\"{o}ft/Howard/Mostowski Model.

(6, 9, 11, 37, 63, 73, 91, 92, 130, 191, 305, 309, 313, 361, and 363
are true, but 77, 88, 106, and 344 are false.)

\item[11.] $\mathcal{N}11$: Jech's Model.

(6, 11, 37, 73, 91, 92, 130, 191, 305, 309, 313, 361, and 363 are
true, but 15, $59(\precsim)$, and 192 are false.)

\item[12a.] $\mathcal{N}12(\aleph_{1})$: A variation of Fraenkel's model, $\mathcal{N}1$.

(11, 23, 63, 130, 165, 187, 191, 202, 305, 309, and 313 are true,
but 3, 41, $45(n)$, 67, 88, 100, 133, 152, 163, 192, 193, and 344 are
false.)

\item[12b.] $\mathcal{N}12(\aleph_{2})$: Another variation of $\mathcal{N}1$.

(11, 23, 40, 44, 63, 91, 130, 135, 165, 191, 305, 309, and 313 are
true, but 62, 67, 100, 133, 152, 161, 163, 192, 193, and 344 are
false.)

\item[12c.] $\mathcal{N}12(\aleph_{\alpha})$:A generalization of $\mathcal{N}12(\aleph_{1})$.

$(11,\,63,\,87(\beta)$ if $\aleph_{\beta}<\mathrm{cf}(\aleph_{\alpha})$, 91, 130, 191, 305, 309, and 313 are
true, but 3, 15, 133, 163, 192, and 193 are false.)

\item[13.]  $\mathcal{N}13$: L\"{a}uchli/Jech Model.

(6, 11, 37, 73, 91, 92, 130, 191, 305, 309, 313, 361, and 363 are
true, but 9, 15, and 235 are false.)

\item[14.] $\mathcal{N}14$: Morris/Jech Model.

(6, 11, 37, 43, 63, 73, 91, 92, 130, 191, 305, 309, 313, 361, and 363
are true, but 15, 103, 133, 163, and 192 are false.)

\item[15.] $\mathcal{N}15$: Brunner/Howard Model I.

(8, 11, 43, 63, 91, 122, 130, 152, 191, 305, 309, 313, and 363 are
true, but 3, 62, 151, 163, 192, 193, 324, 330, and 344 are false.)

\item[16a.] $\mathcal{N}16$: Jech/Levy/Pincus Model.

(6, 11, 37, 41, 63, 73, 91, 92, 130, 191, 305, 309, 313, 361, and 363
are true, but 8, 15, 106, and 133 are false.)

\item[16b.] $\mathcal{N}16(\aleph_{\gamma})$: Levy's Model II.

$(6,\,11,\,37,\,63,\,71(\alpha)$ for $\alpha<\gamma$, 73, 91, 92, 130, 191, 305, 309, 313,
361, and 363 are true, but 8, 15, $71(\gamma), 106$, and 133 are false.)

\item[17.] $\mathcal{N}17$: Brunner/Howard Model II.

(11, 31, 37, 73, 91, 92, 130, 191, 305, 309, 361, and 363 are true,
but 27, 28, 62, 106, 131, 163, 165, and 344 are false.)

\item[18.] $\mathcal{N}18$: Howard's Model I.

(6, 11, 32, 37, 73, 91, 92, 130, 191, 305, 309, 313, 361, and 363 are
true, but 31, 62, 131, 163, 165, and 344 are false.)

\item[19.] $\mathcal{N}19(_{\sim}^{\prec})$: Tsukada's Model.

(11, 16, 23, 24, 37, 60, 73, 91, 92, 130, 165, 191, 305, 317, 361, and
363 are true, but 15, 51, $59(_{\sim}^{\prec})$, 118, 144, 152, 192, 253, and 328
are false.)

\item[20.] $\mathcal{N}20$: Truss' Model II.

(6, 11, 37, 73, 91, 92, 130, 191, 305, 309, 313, 361, and 363 are
true, but 15, $136(k)$, and 192 are false.)

\item[21.] $\mathcal{N}21(\aleph_{\alpha+1})$: Jensen's Model.

$(11,\,39\,(\mathrm{for}\,\alpha\geq 1)$, 63, $71(\beta)\,(\mathrm{for}\,\beta<\alpha+1),\,86(\beta)$ (for $\beta<\alpha+1$),
91, 130, 191, 305, 309, 313, and 363 are true, but 15, 43, 106, 133,
163, and 193 are false.)

\item[22.] $\mathcal{N}22(K,p)$: Makowski/Wi\'{s}niewski/Mostowski Model.

$(6,\,11,\,37,\,45(n) (for\,K=\{n\}),\,46(K)$, 73, 88 (for $K=\{2\})$, 91,
92, 130, 191, 305, $308(p), 309,313,361$, and 363 are true, but 285,
$288(p)$, and 344 are false.)

\item[23.] $\mathcal{N}$23: Howard/Mostowski Linearly Ordered Model.

(8, 11, 63, 91, 130, 191, 305, 309, 313, and 363 are true, but 88,
133, 163, 192, 193, and 344 are false.)

\item[24a.] $\mathcal{N}24$: Hickman's Model I.

(11, 16, 23, 24, 37, 73, 92, 121, 130, 133, 191, 305, 309, 313, 361,
and 363 are true, but 124, 127, 131, 166, and 344 are false.)

\item[24b.] $\mathcal{N}24(n)$: An extension of $\mathcal{N}24$ to $n$-element sets, $n>1$.

(11, 16, 23, 24, 37, 73, 92, 121, 130, 133, 191, 305, 309, 313, 361,
and 363 are true, but 64, 131, and 344 are false.)

\item[24c.] $\mathcal{N}24(n,\,LO)$: Truss' Model III.

(6, 11, 37, 73, 91, 92, 122, 130, 191, 305, 309, 313, 361, and 363
are true, but $33(n)$, 67, 192, and 344 are false.)

\item[25.] $\mathcal{N}25$: Brunner's Model I.

(6, 11, 37, 73, 91, 92, 130, 191, 305, 309, 313, 361, and 363 are
true, but 15, 158, and 192 are false.)

\item[26.] $\mathcal{N}26$: Brunner/Pincus Model, a variation of $\mathcal{N}2$.

(11, 16, 23, 24, 37, 73, 92, 130, 133, 191, 305, 309, 313, 361, 363,
and 380 are true, but 64, 131, 167, and 344 are false.)

\item[27.] $\mathcal{N}27$: Hickman's Model II.

(6, 11, 37, 43, 63, 73, 91, 92, 130, 191, 305, 309, 313, 361, and 363
are true, but 15, 133, 163, $183(\alpha)$, and 192 are false.)

\item[28.] $\mathcal{N}28$: Blass' Permutation Model.

(6, 11, 37, 73, 91, 92, 130, 305, 309, 313, 361, and 363 are true,
but 129, 190, and 293 are false.)

\item[29.] $\mathcal{N}29$: Dawson/Howard Model.

(11, 16, 23, 24, 37, 60, 73, 91, 92, 130, 165, 191, 305, 317, 361, and
363 are true, but 15, 51, 118, 126, 131, 144, 152, 167, 200, 253, and
291 are false.)

\item[30.] $\mathcal{N}$30: L\"{a}uchli's Model III.

(6, 11, 37, 73, 91, 92, 130, 191, 305, 309, 313, 361, and 363 are
true, but 15, 192, and 210 are false.)

\item[31.] $\mathcal{N}31$: L\"{a}uchli's Model IV.

(6, 11, 37, 73, 91, 92, 130, 191, 305, 309, 313, 361, and 363 are
true, but 15, 241, 243, and 244 are false.)

\item[32.] $\mathcal{N}32$: Hickman's Model III.

(6, 11, 37, 73, 91, 92, 130, 191, 305, 309, 313, 361, and 363 are
true, but 15, 192, 237, 238, and 240 are false.)

\item[33.] $\mathcal{N}33$: Howard/H.Rubin/J.Rubin Model.

(11, 23, 40, 63, 91, 130, 165, 191, 305, 309, and 313 are true, but
3, 62, 133, 152, 163, 192, 193, 260, 344, 349, and 356 are false.)

\item[34.] $\mathcal{N}34$: Pincus' Model X.

(6, 11, 37, 73, 91, 92, 130, 191, 216, 305, 309, 313, 361, and 363
are true, but 15 and 217 are false.)

\item[35.] $\mathcal{N}$35: Truss' Model IV.

(6, 11, 37, 73, 91, 92, 130, 191, 250, 305, 309, 313, 361, and 363
are true, but 249 and 344 are false.)

\item[36.] $\mathcal{N}36(\beta)$: Brunner/Howard Model III.

$(11,\,39,\,43,\,63,\,86(\beta)$, 91, 122, 130, 191, 305, 309, 313, and 363 are
true, but 15, 67, 133, 152, 163, 192, and 193 are false.)

\item[37.] $\mathcal{N}$37: A variation of Blass' model, $\mathcal{N}28$.

(6, 11, 37, 73, 91, 92, 130, 191, 305, 309, 313, 316, 361, and 363
are true, but 51, 129, 192, and 293 are false.)

\item[38.] $\mathcal{N}38$: Howard/Rubin Model I.

(11, 23, 60, 91, 113, 130, 165, 191, 295, 305, 317, and 363 are true,
but 3, 15, 51, 106, 118, 144, 152, 163, 181, 253, and 328 are false.)

\item[39.] $\mathcal{N}39$: Howard's Model II.

(11, 16, 23, 24, 37, 73, 91, 92, 130, 191, 231, 305, 309, 313, 361,
and 363 are true, but 15, 29, 131, 152, and 328 are false.)

\item[40.] $\mathcal{N}$40: Howard/Rubin Model II.

(11, 23, 60, 91, 130, 165, 191, 295, 305, 317, and 345 are true, but
3, 15, 51, 118, 144, 152, 163, 181, 192, 193, 253, and 328 are false.)

\item[41.] $\mathcal{N}41$: Another variation of $\mathcal{N}3$.

(6, 9, 11, 37, 63, 73, 91, 92, 122, 130, 191, 305, 309, 313, 361, and
363 are true, but 88, 163, 294, 344, 350, and 357 are false.)

\item[42.] $\mathcal{N}42(p)$: Hickman's Model IV.

(6, 11, 37, 73, 91, 92, 130, 191, 305, 309, 313, 361, and 363 are
true, but 15, 192, and $220(p)$ are false.)

\item[43.] $\mathcal{N}$43: Brunner's Model II.

(6, 11, 37, 73, 91, 92, 130, 191, 305, 309, 313, 361, and 363 are
true, but 64, 80, and 344 are false.)

\item[44.] $\mathcal{N}44$: Gross' model.

(6, 11, 37, 73, 91, 92, 130, 191, 305, 309, 313, 361, and 363 are
true, but 15, 192, and 236 are false.)

\item[45.] $\mathcal{N}45(p)$: Howard/Rubin Model III.

(6, 11, 37, 73, 91, 92, 130, 191, 305, 309, 313, 361, and 363 are
true, but 15, 192, and $220(p)$ are false.)

\item[46.] $\mathcal{N}46$: H\"{o}ft/Howard Model I.

(6, 11, 37, 73, 84, 91, 92, 130, 191, 305, 309, 313, 361, and 363 are
true, but 15 and 185 are false.)

\item[47.] $\mathcal{N}47$: H\"{o}ft/Howard Model II.

(6, 11, 37, 73, 77, 91, 92, 130, 191, 305, 309, 313, 361, and 363 are
true, but 15, 192, and 316 are false.)

\item[48.] $\mathcal{N}48$: Pincus' Model XI.

(11, 16, 23, 24, 37, 60, 73, 91, 92, 130, 165, 191, 295, 305, 317, 361,
and 363 are true, but 51, 106, 118, 126, 131, 144, 253, 296, and
328 are false.)

\item[49.] $\mathcal{N}49$: De la Cruz/Di Prisco Model.

(6, 9, 11, 37, 63, 73, 91, 92, 130, 191, 305, 309, 313, 361, and 363
are true, but $47(n)$, 106, 163, 167, 344, and 380 are false.)

\item[50.] $\mathcal{N}50(E)$: Brunner's Model III.

(6, 11, 37, 67, 73,92, 130, 191, 305, 313, 322, 361, and 363 are true,
but 10, 154, 164, $342(n)$ (for $n$ as described above) and 344 are
false.)

\item[51.] $\mathcal{N}51$: W\c{e}glorz/Brunner Model.

(6, 11, 37, 64, 73, 90, 92, 130, 191, 305, 309, 313, 361, and 363 are
true, but 9, 30, 221, and 344 are false.)

\item[52.] $\mathcal{N}52$: Felgner/Truss Model.

(6, 11, 37, 49, 73, 91, 92, 130, 191, 305, 309, 313, 361, and 363 are
true, but 14, 15, 51, 99, 118, 192, and 253 are false.)

\item[53.] $\mathcal{N}53$:Good/Tree/Watson Model I.

(6, 9, 11, 37, 63, 73, 91, 92, 130, 273, 305, 309, 313, 361, 363, 368,
369 is true, but 88, 133, 192, 344, 382, and 383 are false.)

\item[54.] $\mathcal{N}54$: Good/Tree/Watson Model II.

(6, 11, 37, 73, 91, 92, 130, 273, 305, 309, 313, 361, 363, 368, 369 is
true, but 15, 173, and 192, are false.)
\end{enumerate}
\end{enumerate}
\addtocontents{toc}{\SkipTocEntry}
\setcounter{chapter}{2}
\chapter{Models}\label{part03}

\section*{Introduction}

We consider two types of models: Models of Zermelo-Fraenkel set theory, ZF, in section I and models of $\mathrm{ZF}^{0}$ in section II. These two versions of set theory are described in the introduction to the book. (In what follows, AC is the Axiom of Choice and GCH is the Generalized Continuum Hypothesis.) Our models of ZF are Cohen models\index{Subject}{Cohen models} (or some modern variation) which are obtained from models of ZFC = ZF +AC by using the method of \emph{forcing} to construct an extension which is a model of ZF without AC. The models of $\mathrm{ZF}^{0}$ that we use are the permutation models of Fraenkel\index{Author}{Fraenkel} and Mostowski (FM models). We start with a model of $\mathrm{ZF^{0}+ AC}$ which we refer to as the ground model and use a group of permutations of the atoms to construct a submodel in which AC is false. Details on the construction of models for ZF by the method of forcing can be found in Cohen\index{Author}{Cohen} \cite{bib:1966}, Felgner \cite{bib:1971a}, Jech\index{Author}{Jech} \cite{bib:1973b}, Kunen\index{Author}{Kunen} \cite{bib:1980}, and Solovay\index{Author}{Solovay} \cite{bib:1970}. For more details on permutation models\index{Subject}{permutation models} and the Fraenkel-Mostowski\index{Subject}{Fraenkel-Mostowski (FM) models} method we refer the reader to Felgner \cite{bib:1971a}, Fraenkel\index{Author}{Felgner} \cite{bib:1922}, Jech \cite{bib:1973b}, and Levy\index{Author}{Levy} \cite{bib:1963a}.

We refer to the statements in the numerical list of forms both by their name (e.g. Boolean Prime Ideal Theorem or BPI) and by their number (e.g. BPI is form 14). Sometimes we refer to a form by a symbol. For example, $C(\infty,\,\infty)$ is AC in the form: ``Every set of non-empty sets has a choice function.'' (form 1); $C(\aleph_{0},\,<\aleph_{0})$ is AC for a denumerable set of finite sets (10); $C(WO,\,\infty)$ is AC for a well ordered set of non-empty sets (40); $UT(\aleph_{0},\,\aleph_{0},\,\aleph_{0})$ is ``The union of a denumerable number of denumerable sets is denumerable.'' (The countable union theorem, 31); etc. These and other abbreviations are defined in both the numerical list of forms and the topical list of forms.

\section{Cohen Models, Models of ZF. ($\mathcal{M}$)}\label{pt03:sec1}

There are some weak forms of AC which are true in all Cohen models and some which are false in all Cohen models. These are described briefly below. More information can be found in the notes to which we refer.

Several of our weak forms of AC are known to imply AC in the theory ZF (but not in the theory $\mathrm{ZF}^{0}$). More detailed information can be found in notes 18 and 75. Since AC is false in all of our Cohen models of ZF all forms which imply AC in ZF must be false. In the theory $\mathrm{ZF}^{0}$, all of these forms are known to imply form 91. Therefore, in the description of each Cohen model, we will say ``91 is false'' and refer to note 75.

Secondly, there are several implications involving more than two forms that are known to hold in $\mathrm{ZF}^{0}.$ These are listed with references in note 120. Most are of the form A+B implies C, where A,\,B, and C are form numbers. When it follows from these implications that a certain form is true or false in a model we will, in the description of the model, assert that the form is true (or false) and give a reference to the part of note 120 from which the truth or falsity follows.

\noindent $\mathcal{M}1$: Cohen's original model. Add a denumerable number of generic reals (subsets of $\omega$), $a_{1},\,a_{2},\, \cdots$, along with the set $b$ containing them. It follows from the work of Halpern and Levy that the Boolean Prime Ideal Theorem (BPI,14) and the Kinna-Wagner Selection Principle (15) are true in this model. Cohen\index{Author}{Cohen} proved that in this model there is an infinite subset of the reals that has no denumerable subset (13 is false). Also, the Krein-Milman Theorem (KM, 65) is false because BPI + KM $\rightarrow$ AC and Haddad and Morillon proved that in $\mathcal{M}$1 every orderable topological space is normal (118 is true). Jech\index{Author}{Jech} has shown that the Axiom of Choice for a set of well orderable sets $(C(\infty,\,WO)$, 60) is true and the countable union theorem $(UT(\aleph_{0},\,\aleph_{0},\,\aleph_{0})$, 31) is true. Blass\index{Author}{Blass} proved that Ramsey's Theorem (17) is false. Brunner \cite{bib:1982a} has shown that every set that cannot be well ordered has an infinite Dedekind finite subset (163 is true), but the interval [0, 1] of real numbers is not well orderable and has no infinite subset with a Dedekind finite power set (164 is false). It is shown in note 102 that there is a set $S$ consisting of finite unions of sets in $b$, the set of generic reals, that is closed with respect to chain unions, but has no $\subseteq$-maximal element (289 is false). Brunner\index{Author}{Brunner} \cite{bib:1984f} shows that in this model on every infinite set there is a Hausdorff topology with an infinite set of non-isolated points (128 is true). Blass \cite{bib:1979} has shown that in every symmetric model there is a set $X$ such that for each set $a$, there is an ordinal $\alpha$ and a function $f$ mapping $X\times\alpha$ onto $a$ (191 is true). (See note 59.) Since BPI is true, it follows from Howard\index{Author}{Howard} \cite{bib:1975} that {\L}os' Theorem (253) is false. Monro\index{Author}{Monro} proves that every cardinal that cannot be well ordered is decomposable (i.e. is the sum of two smaller positive cardinals), so 277 is true. In Jech \cite{bib:1973b} (prob 5.20), it is shown that the power set of an infinite set is Dedekind infinite (82 is true). David\index{Author}{David} shows that no infinite subset of $b$ is almost well orderable so 144 is false. See note 25 for definitions. By the results of Bell\index{Author}{Bell} \cite{bib:1983}, \cite{bib:1988} and Monro \cite{bib:1972} the following three forms are false in $\mathcal{M}$1: Form 50 (Sikorski's extension theorem on Boolean homomorphisms), form 299 (Any extremally disconnected compact Hausdorff space is projective in the category of Boolean topological spaces.) and form 300 (Any continuous surjection between extremally disconnected Hausdorff spaces has an irreducible restriction to a closed subset of its domain.) See notes 60 and 114 for definitions. Since 154 is true (14 implies 154) and 43 is false (43 implies 13), form 106 (the Baire Category Theorem for compact $T_{2}$ spaces) is false because 106 + 154 implies 43. (See Brunner \cite{bib:1984b} and note 95.) In this model, $\mathbb{R}$ with the order topology is a separable metric space that is not Lindel\"{o}f (337 is false). Also, there is a dense set $K\subseteq \mathbb{R}$ which has no denumerable subset such that $K$ is a second countable metric space that is neither Lindel\"{o}f nor separable (338 and 339 are false). (See Good/Tree \cite{bib:1995}.)

\noindent $\mathcal{M}1 \models$ 14, 15, 31, 60, 82, 118, 128, 163, 165, 191, and 277, but 13, 17, 50, 65, 106, 131, 144, 164, 253, 289, 299, 300, 337, 338, and 339 are false. References include Bell \cite{bib:1983}, \cite{bib:1988}, Blass \cite{bib:1979}, \cite{bib:1977a}, Brunner \cite{bib:1982a}, \cite{bib:1984b}, \cite{bib:1984f}, Cohen \cite{bib:1963}, \cite{bib:1966}, David \cite{bib:1980}, Good/Tree\index{Author}{Good/Tree} \cite{bib:1995}, Halpern/Levy\index{Author}{Halpern/Levy} \cite{bib:1971}, Haddad/Morillon\index{Author}{Haddad/Morillon} \cite{bib:1990}, Howard \cite{bib:1975}, Monro \cite{bib:1972}, \cite{bib:1973a}, Sageev\index{Author}{Sageev} \cite{bib:1981}, Shoenfield\index{Author}{Shoenfield} \cite{bib:1967}, Smullyan\index{Author}{Smullyan} \cite{bib:1994} and note 120(30 and 49). See, also, Felgner\index{Author}{Felgner} \cite{bib:1971a} (pp 100--103, 131--145), and Jech \cite{bib:1973b} (Chapt. 5, probs 20, 22).

(A minimal list of forms implying each form known to be true in the model: 14, 15, 31, 60, 118, 128, 163, 165, 191, 277.

A minimal list of forms at least one of which is implied by each form known to be false in the model: 13, 17, 50, 65, 131, 144, 253, 289, 299, 300, 337, 338, 339.

A list of all forms known to be true in the model: 0, 5, 6, 10, 14, 15, 18, 19, 30, 31, 32, 33, 34, 35, 36, 38, 45, 46, 47, 48, 49, 52, 60, 61, 62, 63, 64, 69, 70, 72, 80, 82, 83, 84, 85, 88, 93, 96, 99, 102, 104, 107, 108, 111, 118, 119, 120, 121, 122, 123, 127, 128, 132, 137, 139, 140, 141, 142, 146, 150, 151, 153, 154, 163, 165, 178, 182, 189, 190, 191, 198, 201, 206, 209, 213, 216, 221, 222, 223, 225, 226, 227, 228, 229, 231, 233, 235, 241, 242, 249, 250, 268, 269, 270, 271, 276, 277, 280, 283, 285, 287, 288, 293, 294, 295, 296, 298, 308, 309, 311, 313, 314, 322, 323, 324, 326, 327, 329, 330, 331, 332, 336, 342, 343, 344, 349, 350, 355, 356, 357, 358, 370, 371, 373, 374, 376, 377, 378, 379, 380, (383).

A\index{Author}{Tsukada} list of all forms known to be false in the model: 1, 2, 3, 4, 7, 8, 9, 11, 12, 13, 16, 17, 39, 40, 41, 43, 44, 50, 51, 65, 66, 67, 71, 73, 75, 76, 77, 79, 86, 87, 89, 90, 91, 92, 94, 95, 100, 101, 106, 109, 112, 113, 114, 126, 129, 130, 131, 133, 144, 147, 149, 152, 161, 164, 168, 174, 175, 179, 181, 185, 188, 192, 193, 202, 203, 211, 212, 214, 218, 253, 255, 256, 257, 258, 259, 260, 261, 262, 264, 286, 289, 292, 297, 299, 300, 301, 302, 303, 316, 325, 328, 333, 334, 335, 337, 338, 339, 345, 347, 359, 360, (383).)

$\mathcal{M}1(\langle(\omega_{1}\rangle)$: Cohen/Pincus Model. Pincus\index{Author}{Pincus} extends the methods of Cohen and adds a generic $\omega_{1}$-sequence, $\langle I_{\alpha}:\alpha\in\omega_{1}\rangle$, of denumerable sets, where $I_{0}$ is a denumerable set of generic reals, each $I_{\alpha+1}$ is a generic set of enumerations of $I_{\alpha}$, and for a limit ordinal $\lambda,\,I_{\lambda}$ is a generic set of choice functions for $\{I_{\alpha}:\alpha<\lambda\}$. Pincus proves that the Principle of Dependent Choices (43) is true, for each family of infinite sets $x$, there is a function $f$ such that for each $y\in x,\,f(y)\subseteq y$ and $|f(y)|=\aleph_{0}(214)$, is true, every set can be linearly ordered (30) is true, there is a class of cardinal representatives, hence form 97 is true, but $C(\aleph_{1},\,\aleph_{0})$ is false because $<I_{\alpha}:\alpha\in \omega_{1}>$ has no choice function. Therefore, forms 39 $(C(\aleph_{1},\,\infty))$ and 85 $(C(\infty,\,\aleph_{0}))$ are false. Since form 64 (There is no amorphous set.) is true (30 implies 64), form 133 (Every set is either well orderable or has an infinite amorphous subset.) is false. Form 163 (Every set is either well orderable or has an infinite Dedekind finite subset.) is also false because 43 implies 8 and 8 + 163 implies AC (Brunner\index{Author}{Brunner} \cite{bib:1982a}). Since form 43 implies form 296 (Every infinite set is the disjoint union of infinitely many infinite sets.) it follows from Pincus \cite{bib:1997} that form 295 (Every infinite set has a dense linear ordering.) is also true. (30 + 296 implies 295.) Form 62 $(C(\infty,\, <\aleph_{0}))$ is true because 295 implies 62 and form 60 $(C(\infty,\,WO))$ is false because 60 implies 85. Therefore, form 329 $(MC(\infty,\,WO))$ is false because 62 + 329 $\rightarrow$ 60.

$\mathcal{M}1(\langle\omega_{1}\rangle) \models$ 30, 43, 97, 214, and 295, but 20, 39, 85, 91, 133, 151, 163, 328, and 329 are false. References Pincus \cite{bib:1977a}, \cite{bib:1997}, Brunner \cite{bib:1982a}, note 18, and note 120(32, 43, and 47).

(A minimal list of forms implying each form known to be true in the model: 43, 97, 214, 295.

A minimal list of forms at least one of which is implied by each form known to be false in the model: 39, 85, 91, 151, 163, 328, 329.

A list of all forms known to be true in the model: 0, 4, 5, 6, 8, 9, 10, 13, 16, 17, 18, 19, 24, 26, 27, 28, 29, 30, 31, 32, 33, 34, 35, 36, 37, 38, 42, 43, 45, 46, 47, 48, 57, 61, 62, 64, 73, 74, 75, 76, 77, 78, 80, 82, 83, 84, 88, 92, 93, 94, 97, 98, 102, 104, 106, 108, 111, 113, 119, 120, 121, 122, 124, 126, 127, 128, 131, 132, 140, 142, 146, 150, 152, 154, 155, 173, 178, 182, 185, 194, 198, 199, 209, 211, 214, 216, 217, 243, 249, 250, 268, 269, 273, 276, 279, 280, 282, 283, 285, 288, 293, 295, 296, 304, 308, 325, 327, 336, 337, 338, 339, 340, 341, 342, 350, 351, 352, 353, 354, 355, 357, 358, 360, 361, 362, 363, 364, 373, 374, 375, 376, 377, 378, 380, (383).

A list of all forms known to be false in the model: 1, 20, 21, 23, 39, 40, 44, 60, 66, 67, 85, 89, 90, 91, 95, 101, 109, 112, 114, 133, 147, 149, 151, 163, 164, 202, 213, 218, 231, 255, 256, 257, 258, 259, 260, 261, 262, 264, 286, 292, 328, 329, 333, 334, 335, 347, 359, (383).)

$\mathcal{M}2$: Feferman's model. Add a denumerable number of generic reals to the base model, but do not collect them. Feferman\index{Author}{Feferman} has shown that in this model there is an infinite set of pairs with no choice function (88 is false), there is a partition of $\mathcal{P}(\omega)$ with no choice function (203 is false) and the Boolean Prime Ideal Theorem (14) is false. Truss \cite{bib:1978}\index{Author}{Truss} has shown that in this model $C(WO,\, \infty)$ (40) is true. (He gives credit to Solovay for the proof.) In Truss \cite{bib:1978} it is also shown that $\mathcal{P}(\omega)$ does not have a non-trivial ultrafilter (70 is false), and for every infinite set $x$ there is a mapping of $x$ onto $2\times x$ (204) is true. (See note 65.) Pincus\index{Author}{Pincus} shows that there is a set of reals that does not have the property of Baire (142 is true), but $\mathcal{P}(\omega)$ has no nontrivial real valued measure which is 0 on finite sets (222 is false). (The proof of form 142 requires that the Continuum Hypothesis holds in the outer model.) Since 40 implies 8 $(C(\aleph_{0},\,\infty))$, it follows from Brunner \cite{bib:1982a} that in this model there is a set that cannot be well ordered and does not have an infinite Dedekind finite subset, (163 is false). (Form 8 plus form 163 iff AC.) Truss \cite{bib:1974b} proves that there is a sequence $x_{1},\,x_{2},\,\cdots$ of cardinals such that for each $n,\, x_{n}=\aleph_{n+1}+2^{\aleph_{0}}$ adj $\aleph_{n+1}+\aleph_{2}\cdot2^{\aleph_{0}}$ adj $\aleph_{n+1}+\aleph_{3}\cdot2^{\aleph_{0}}$ adj $\cdots$ adj $\aleph_{n+1}+\aleph_{n+1}\cdot 2^{\aleph_{0}}=x_{n}^{2}$. (The symbol $x$ adj $y$ means $x<y$ and there is nothing between $x$ and $y$.) It follows that 274 is false.

$\mathcal{M}2 \models$ 40, 142, 165, and 204, but 70, 88, 152, 163, 203, 222, and 274 are false. References Brunner \cite{bib:1982a}, Feferman \cite{bib:1965}, Pincus \cite{bib:1972c}, Truss \cite{bib:1978}, \cite{bib:1974b}, note 65, and note 120(34 and 46). (See also Felgner\index{Author}{Felgner} \cite{bib:1971a} pp 160-166 where credit for some of the proofs is given to D. Scott and R. Solovay.)

(A minimal list of forms implying each form known to be true in the model: 40, 142, 165, 204.

A minimal list of forms at least one of which is implied by each form known to be false in the model: 88, 152, 163, 203, 222, 274.

A list of all forms known to be true in the model: 0, 5, 6, 8, 9, 10, 13, 16, 17, 18, 19, 24, 26, 27, 28, 29, 31, 32, 34, 35, 36, 37, 38, 39, 40, 42, 43, 47, 48, 57, 58, 64, 73, 74, 75, 77, 78, 80, 82, 83, 84, 86, 92, 94, 98, 104, 106, 108, 111, 113, 119, 122, 124, 126, 127, 128, 131, 132, 142, 150, 151, 154, 155, 165, 182, 185, 194, 196, 198, 199, 204, 208, 209, 211, 216, 217, 231, 243, 249, 250, 273, 279, 280, 282, 288, 294, 296, 304, 322, 324, 325, 327, 328, 330, 336, 337, 338, 339, 340, 341, 342, 350, 351, 352, 353, 354, 355, 357, 358, 360, 361, 362, 363, 364, 373, 374, 375, 376, 377, 378, 380, (383).

A list of all forms known to be false in the model: 1, 14, 15, 30, 49, 50, 52, 60, 61, 62, 63, 66, 67, 68, 70, 79, 85, 88, 89, 90, 91, 95, 99, 107, 109, 112, 114, 123, 130, 133, 147, 149, 152, 163, 164, 201, 202, 203, 213, 214, 218, 221, 222, 225, 264, 270, 274, 284, 287, 292, 295, 303, 311, 313, 317, 323, 326, 331, 332, 333, 334, 335, 343, 345, 346, 372, (383).)

$\mathcal{M}2(\kappa)$: Feferman/Pincus Model. This is an extension of $\mathcal{M}2$ in which there are $\kappa$ generic sets, where $\kappa$ is a regular cardinal. ($\mathcal{M}2$ is $\mathcal{M}2(\aleph_{0}).)$ Pincus proves that the Hahn-Banach Theorem (52) is false in this model. If $\kappa$ is sufficiently large, then there is a non-trivial measure on $\mathcal{P}(\omega)$ which is 0 on finite sets (222 is true). $\kappa$ has to be large enough so that form 222 is boundable in sets of rank $\mu$, where $\mu$ is less than $\kappa.$

$\mathcal{M}2(\kappa) \models$ 222, but 52, and 91 are false. References Pincus\index{Author}{Pincus} \cite{bib:1972c} and note 18.

(A minimal list of forms implying each form known to be true in the model: 222.

A minimal list of forms at least one of which is implied by each form known to be false in the model: 52, 91.

A list of all forms known to be true in the model: 0, 142, 222, 223, 280, (383).

A list of all forms known to be false in the model: 1, 14, 50, 52, 66, 67, 89, 90, 91, 95, 109, 112, 114, 133, 147, 149, 164, 202, 218, 264, 292, 303, 317, 333, 334, 335, 345, 372, (383).)

$\mathcal{M}2(\langle\omega_{2}\rangle)$: Feferman/Truss\index{Author}{Truss} Model. This is another extension of $\mathcal{M}2$. The set of all sequences of generic reals of length less than $\omega_{2}$ are added to a countable transitive model of ZFC + V = L. Truss proves that $C(|\mathbb{R}|,\,\infty)$ (181) is true, but there is a partition of $\mathcal{P}(\omega)$ that has no choice function (203 is false). Since the Axiom of choice is false, it follows that all forms that imply AC in ZF are also false. In particular, $C(LO,\,\infty)$ (202), is false. (See the first paragraph of section I.) Form 163 (every set is either well orderable or has an infinite Dedekind finite subset) is also false because 181 implies 8 and 8 + 163 implies AC (Brunner\index{Author}{Brunner} \cite{bib:1982a}).

$\mathcal{M}2(\langle \omega_{2}\rangle) \models$ 181, but 163, 202, and 203 are false. References Truss \cite{bib:1978} and Brunner \cite{bib:1982a}.

(A minimal list of forms implying each form known to be true in the model: 181.

A minimal list of forms at least one of which is implied by each form known to be false in the model: 163, 203.

A list of all forms known to be true in the model: 0, 5, 6, 8, 9, 10, 13, 16, 17, 18, 19, 24, 26, 27, 28, 29, 31, 32, 34, 35, 36, 37, 38, 57, 64, 73, 74, 75, 80, 82, 83, 84, 92, 94, 98, 104, 108, 119, 124, 126, 127, 128, 131, 132, 150, 181, 182, 185, 194, 198, 199, 209, 216, 217, 249, 282, 288, 296, 304, 325, 336, 337, 338, 339, 340, 341, 342, 350, 351, 352, 353, 354, 355, 357, 358, 360, 361, 362, 373, 374, 376, 377, 378, 380, (383).

A list of all forms known to be false in the model: 1, 66, 67, 79, 89, 90, 91, 95, 109, 112, 114, 130, 133, 147, 149, 163, 164, 202, 203, 218, 264, 292, 333, 334, 335, (383).)

$\mathcal{M}3$: Mathias' model. Mathias proves that the FM model $\mathcal{N}4$ can be transformed into a model of ZF, $\mathcal{M}3$. $\mathcal{M}3$ is an extension of $\mathcal{M}1$ in which a universal homogeneous partial ordering is defined on $b$, the set of generic reals. (See Jech\index{Author}{Jech} \cite{bib:1973b} p 101 for definitions.) Mathias then proves that the partial ordering on $b$ cannot be extended to a linear ordering (49 is false), but every set can be linearly ordered (30 is true). Also, in this model, every set of well orderable sets has a choice function (60) and the Kinna-Wagner Selection Principle (15) are true. Form 60 implies 294 (Every linearly ordered $W$-set is well orderable.), and 30 + 294 + 144 (Every set is a $W$-set.) implies AC. It follows that 144 is false. (See note 25 for definitions.) Since form 64 (There is no amorphous set.) is true (15 implies 64), form 133 (Every set is either well orderable or has an infinite amorphous subset.) is false because 64 + 133 implies AC. The Boolean Prime Ideal Theorem (14) is false because 14 implies 49, and $C(\infty,\,<\aleph_{0})$ (62) is true because 15 implies 62. It was shown by Blass (see note 33) that 62 + 99 (Rado's Selection Lemma) implies 14, so form 99 is false.

$\mathcal{M}3 \models$ 15, 60, and 165, but 49, 91, 99, 144, 152, and 133 are false. References Mathias\index{Author}{Mathias} \cite{bib:1967}, notes 25, 33 and 120(6, 34 and 43). Also, see Felgner\index{Author}{Felgner} \cite{bib:1971a} (pp 114--128) and Jech\index{Author}{Jech} \cite{bib:1973b} (p 113, prob 8).

(A minimal list of forms implying each form known to be true in the model: 15, 60, 165.

A minimal list of forms at least one of which is implied by each form known to be false in the model: 49, 91, 99, 144, 152.

A list of all forms known to be true in the model: 0, 5, 10, 15, 18, 30, 32, 33, 45, 46, 47, 48, 60, 61, 62, 64, 80, 82, 83, 84, 85, 88, 93, 102, 111, 119, 120, 121, 122, 127, 132, 140, 142, 146, 150, 151, 165, 178, 198, 213, 216, 231, 249, 250, 268, 269, 270, 271, 276, 280, 283, 285, 288, 293, 294, 295, 296, 308, 314, 322, 323, 324, 327, 329, 330, 336, 342, 349, 350, 355, 356, 357, 358, 373, 374, 376, 377, 378, 379, 380, (383).

A list of all forms known to be false in the model: 1, 14, 49, 50, 66, 67, 89, 90, 91, 95, 99, 109, 112, 114, 133, 144, 147, 149, 152, 164, 179, 202, 214, 218, 264, 292, 303, 317, 333, 334, 335, 345, (383).)

$\mathcal{M}4$: Pincus' Model I. This model has many of the properties of the model given in Cohen\index{Author}{Cohen} \cite{bib:1966} (p 143). The ground model, $\mathcal{M}$, is a model of ZF + V = L. Forcing conditions are functions from countable subsets of $\omega_{1}$ to $\{0, 1\}$. $\mathcal{M}4$ is a symmetric submodel of the corresponding generic extension of $\mathcal{M}$. Pincus\index{Author}{Pincus} proves that the Hahn-Banach Theorem (52) is true, but there is an infinite set with no denumerable subset (9 is false), and there is a denumerable family of pairs that has no choice function. It follows that $KW(\aleph_{0},\,<\aleph_{0})$ (358) and $C(\infty,\,2)$ (88) are false. Pincus also proves that the Krein-Milman Theorem (65) is false. (Since BPI (form 14) implies 358, BPI is also false.)

$\mathcal{M}4 \models$ 52, but 65, 88, 91, and 358 are false. References Pincus \cite{bib:1972c} and note 18.

(A minimal list of forms implying each form known to be true in the model: 52.

A minimal list of forms at least one of which is implied by each form known to be false in the model: 65, 88, 91, 358.

A list of all forms known to be true in the model: 0, 52, 93, 142, 221, 222, 223, 280, 287, 309, (383).

A list of all forms known to be false in the model: 1, 2, 3, 4, 7, 8, 9, 10, 14, 15, 16, 17, 20, 21, 23, 27, 28, 29, 30, 31, 32, 39, 40, 41, 43, 44, 49, 50, 60, 61, 62, 65, 66, 67, 68, 71, 85, 86, 87, 88, 89, 90, 91, 95, 100, 101, 107, 109, 112, 113, 114, 121, 122, 123, 129, 132, 133, 147, 149, 150, 151, 152, 153, 161, 164, 165, 168, 174, 181, 192, 201, 202, 213, 214, 218, 231, 255, 256, 257, 258, 259, 260, 261, 262, 264, 270, 284, 286, 292, 295, 303, 317, 322, 323, 324, 325, 326, 327, 331, 332, 333, 334, 335, 343, 345, 346, 347, 352, 354, 355, 358, 359, 376, 377, 378, 380, (383).)

$\mathcal{M}5(\aleph)$: Solovay's Model. An inaccessible cardinal $\aleph$ is collapsed to $\aleph_{1}$ in the outer model and then $\mathcal{M}5(\aleph)$ is the smallest model containing the ordinals and $\mathbb{R}$. (The method of collapsing cardinals that is used in the model and some of the proofs are due to A. Levy. An inaccessible cardinal is used to construct the model. See, for example, Drake\index{Author}{Drake} \cite{bib:1974}, Kanamori/Magador\index{Author}{Kanamori/Magador} \cite{bib:1978} or Solovay/Reinhardt/Kanamori\index{Author}{Solovay/Reinhardt/Kanamori} \cite{bib:1978} for information about large cardinals.) In this model, The Principle of Dependent Choices (43) is true and every set of reals is Lebesgue measurable (93 is false). Moreover, if $\aleph$ is a measurable cardinal, then so is $\aleph_{1}$ (318 is false). (Also, see Jech\index{Author}{Jech} \cite{bib:1968a}.) (It has been shown by Shelah\index{Author}{Shelah} \cite{bib:1984} and Raisonnier \cite{bib:1984} that if ZF + $C(\aleph_{0},\,\infty)(8)+\neg\ 93$ is consistent (43 implies 8), then so is ZFC + ``there exists an inaccessible cardinal''. Thus, the existence of an inaccessible cardinal is necessary for this result.) Moreover, in $\mathcal{M}5(\aleph)$, every set of reals has the Baire property (142 is false) and the cardinal $2^{\aleph_{0}}$ is not well orderable and is not decomposable (the sum of two smaller positive cardinals) so 277 is false. Also, the Hahn-Banach Theorem (52) and the Axiom of Choice for pairs (88) are false. Every uncountable set of reals contains a perfect subset (169 is false). The following two statements are also false in this model: There is a complete separable metric space with a non meager subset (280) and there is a Hilbert space $H$ and an unbounded linear operator on $H$ (281). (See also note 96.) Stern\index{Author}{Stern} proves that in this model, every infinite sequence of Borel sets of the same rank is countable (224 is false). (Truss\index{Author}{Truss} has shown that if the same construction is used with $\aleph$ a limit cardinal then 169 is still false, but 43 is also false. See $\mathcal{M}12(\aleph).)$ Mathias\index{Author}{Mathias} has shown that there is no non-Ramsey set (234 is false). According to Kanovei\index{Author}{Kanovei} \cite{bib:1991}, if $m$ is the cardinality of the set of Vitali equivalence classes, then, in $\mathcal{M}5(\aleph)$, $H(m)=H(2^{\aleph_{0}})$, where $H$ is Hartogs' function (307 is true). $(H(x)$ is the smallest ordinal that is not similar to a subset of $x$ and the Vitali equivalence classes are equivalence classes of the real numbers under the relation $x\equiv y\leftrightarrow(\exists q\in \mathbb{Q})(x-y=q).)$ Since 43 implies 8 $(C(\aleph_{0},\,\infty))$, it follows from Brunner\index{Author}{Brunner} \cite{bib:1982a} that in this model there is a set that cannot be well ordered and does not have an infinite Dedekind finite subset, (163 is false). (Form 8 plus form 163 iff AC.) In Wagon\index{Author}{Wagon} \cite{bib:1985}, it is shown that $\mathbb{Z}$ is not amenable in $\mathcal{M}5(\aleph)$ (313 is false) and that there is an amenable group with a non-amenable subgroup in $\mathcal{M}5(\aleph)$ (312 is false). ($G$ is amenable if there is a finitely additive measure $\mu$ on $\mathcal{P}(G)$ such that $\mu(G)=1$ and $\forall A\subseteq G,\,\forall g\in G,\,\mu(gA)=\mu(A).)$

$\mathcal{M}5(\aleph) \models$ 43 and 307, but 88, 93, 142, 163, 169, 224, 234, 277, 280, 281, 312, 313, and 318 are false. References Solovay\index{Author}{Solovay} \cite{bib:1965}, \cite{bib:1970}, Brunner \cite{bib:1982a}, Jech \cite{bib:1968a}, Kanovei \cite{bib:1991}, Monro\index{Author}{Monro} \cite{bib:1973a}, Mathias \cite{bib:1968}, Prikry\index{Author}{Prikry} \cite{bib:1976}, Raisonnier \cite{bib:1984}, Shelah \cite{bib:1984}, Stern \cite{bib:1979}, Truss \cite{bib:1974b}, Wright\index{Author}{Wright} \cite{bib:1973}, Wagon \cite{bib:1985}, and note 20.

(A minimal list of forms implying each form known to be true in the model: 43, 307.

A minimal list of forms at least one of which is implied by each form known to be false in the model: 93, 163, 169, 224, 234, 277, 280, 281, 312, 318.

A list of all forms known to be true in the model: 0, 5, 6, 8, 9, 10, 13, 16, 17, 18, 19, 24, 26, 27, 28, 29, 31, 32, 34, 35, 36, 37, 38, 42, 43, 57, 64, 73, 74, 75, 77, 78, 80, 82, 83, 84, 92, 94, 98, 104, 106, 108, 113, 119, 124, 126, 127, 128, 131, 132, 150, 154, 155, 182, 185, 194, 198, 199, 209, 211, 216, 217, 243, 249, 273, 279, 282, 288, 296, 304, 307, 325, 336, 337, 338, 339, 340, 341, 342, 350, 351, 352, 353, 354, 355, 357, 358, 360, 361, 362, 363, 364, 373, 374, 375, 376, 377, 378, 380, (383).

A list of all forms known to be false in the model: 1, 14, 15, 30, 49, 50, 52, 60, 61, 62, 63, 66, 67, 68, 70, 79, 85, 88, 89, 90, 91, 93, 95, 99, 100, 101, 107, 109, 112, 114, 117, 123, 130, 133, 142, 147, 149, 163, 164, 168, 169, 170, 184, 195, 201, 202, 203, 213, 218, 221, 222, 224, 225, 234, 251, 252, 264, 270, 277, 280, 281, 284, 287, 292, 295, 303, 305, 306, 309, 310, 311, 312, 313, 317, 318, 319, 320, 323, 326, 331, 332, 333, 334, 335, 343, 345, 346, 359, 366, 367, 368, 372, (383).)

$\mathcal{M}6$: Sageev's Model I. Using iterated forcing, Sageev\index{Author}{Sageev} constructs $\mathcal{M}6$ by adding a denumerable number of generic tree-like structures to the ground model, a model of ZF + V = L. He proves that in $\mathcal{M}6$, every set can be linearly ordered (30 is true) and the $2m=m$ principle (3) holds, but the Axiom of Choice for a denumerable set of denumerable subsets of the reals (5) is false. Gonzalez has shown that every infinite set has a dense linear ordering in $\mathcal{M}6$ (295 is true). Form 3 implies form 10 $(C(\aleph_{0},\,<\aleph_{0}))$ and form 231 $(UT(WO,\,WO,\,WO))$ implies form 5. Thus, form 10 is true and form 231 is false. It is shown in note 123 that 10 + 163 (Every non well orderable set has an infinite Dedekind finite subset.) implies 231. Consequently, form 163 is false. (Halpern/Howard\index{Author}{Halpern/Howard} \cite{bib:1970} have shown using a permutation model that the $2m=m$ principle does not imply the Axiom of Choice in $\mathrm{ZF}^{0}$ See $\mathcal{N}9$.) Form 165 $(C(WO,\,WO))$ is false because 165 implies 5 and form 122 $(C(WO,\,<\aleph_{0}))$ is true because 295 implies 122. Therefore, form 330 $(MC(WO,\,WO))$ is false because 122 + 330 $\rightarrow$ 165. In addition, 126 $(MC(\aleph_{0},\,\infty))$ is false because 30 is true and 5 is false. $(8\,(C(\aleph_{0},\, \infty))$ implies 5 and 30 + 126 implies 8.)

$\mathcal{M}6 \models$ 3, 30, and 295, but 5, 126, 163, and 330 are false. References Sageev \cite{bib:1975}, Gonzalez\index{Author}{Gonzalez} \cite{bib:1995a}, notes 2(8), 120(9), and 123.

(A minimal list of forms implying each form known to be true in the model: 3, 295.

A minimal list of forms at least one of which is implied by each form known to be false in the model: 5, 126, 163, 330.

A list of all forms known to be true in the model: 0, 3, 4, 9, 10, 13, 17, 18, 30, 33, 45, 46, 47, 48, 53, 57, 61, 62, 64, 80, 82, 83, 84, 88, 93, 98, 102, 111, 120, 121, 122, 124, 127, 128, 132, 140, 142, 146, 178, 185, 186, 198, 199, 204, 216, 217, 249, 250, 268, 269, 276, 280, 283, 285, 288, 293, 295, 296, 304, 308, 325, 327, 336, 342, 358, 373, 374, 376, 377, 378, (383).

A list of all forms known to be false in the model: 1, 5, 6, 8, 16, 20, 21, 23, 27, 28, 29, 31, 32, 39, 40, 43, 44, 60, 66, 67, 79, 85, 86, 87, 89, 90, 91, 94, 95, 101, 106, 109, 112, 113, 114, 126, 130, 133, 147, 149, 150, 163, 164, 165, 174, 181, 188, 192, 193, 202, 203, 212, 213, 218, 255, 256, 257, 258, 259, 260, 261, 262, 264, 286, 292, 328, 329, 330, 333, 334, 335, 345, 347, 352, 354, 359, (383).)

$\mathcal{M}7$: Cohen's Second Model. In this model there are two denumerable sets $U = \{U_{i}:i\in \omega\}$ and $V=\{V_{i}:i\in \omega\}$ such that for each $i\in \omega,\,U_{i}$ and $V_{i}$ cannot be distinguished. Therefore, in this model there is a denumerable set of pairs whose union has no denumerable subset, so form 18 is false, and $KW(\aleph_{0},\,<\aleph_{0})$ (358) is also false. Blass\index{Author}{Blass} \cite{bib:1979} has shown that in every symmetric model there is a set $X$ such that for each set $a$, there is an ordinal $\alpha$ and a function $f$ mapping $X\times\alpha$ onto $a$ (191 is true). (See note 59.) It is shown in Good/Tree\index{Author}{Good/Tree} \cite{bib:1995} that in this model there is a Lindel\"{o}f metric space that is neither second countable nor separable (340 and 341 are false).

$\mathcal{M}7 \models$ 191, but 18, 91, 340, 341, and 358 are false. References Blass \cite{bib:1979}, Cohen\index{Author}{Cohen} \cite{bib:1965}, \cite{bib:1963}, \cite{bib:1966}, Good/Tree \cite{bib:1995}, Jech\index{Author}{Jech} \cite{bib:1973b} (pp 68ff), and note 18.

(A minimal list of forms implying each form known to be true in the model: 191.

A minimal list of forms at least one of which is implied by each form known to be false in the model: 18, 91, 340, 341, 358.

A list of all forms known to be true in the model: 0, 182, 189, 190, 191, (383).

A list of all forms known to be false in the model: 1, 2, 3, 4, 7, 8, 9, 10, 14, 15, 16, 17, 18, 20, 21, 23, 27, 28, 29, 30, 31, 32, 39, 40, 41, 43, 44, 49, 50, 60, 61, 62, 66, 67, 68, 71, 80, 85, 86, 87, 88, 89, 90, 91, 95, 100, 101, 107, 109, 111, 112, 113, 114, 121, 122, 123, 129, 132, 133, 147, 149, 150, 151, 152, 153, 161, 164, 165, 166, 167, 168, 174, 181, 192, 201, 202, 213, 214, 218, 231, 250, 255, 256, 257, 258, 259, 260, 261, 262, 264, 270, 284, 286, 292, 295, 303, 317, 322, 323, 324, 325, 326, 327, 331, 332, 333, 334, 335, 340, 341, 343, 345, 346, 347, 352, 354, 355, 358, 359, 376, 377, 378, 380, (383).)

$\mathcal{M}7(n)$: Generalization of $\mathcal{M}7$. Model $\mathcal{M}7$ can be generalized to $n$ denumerable sets for $1<n\in \omega$, then the Axiom of Choice for a denumerable number of $n$ element sets, $C(\aleph_{0},\,n)$, is false for $1<n<\omega.$ Therefore, $C(WO,\,n)\ (47(n))$ is also false for $1<n\in \omega$.

In $\mathcal{M}7(n)$, forms $47(n)$ and 91 are false. Reference note 18.

(A minimal list of forms implying each form known to be true in the model: 0.

A minimal list of forms at least one of which is implied by each form known to be false in the model: 47, 91.

A list of all forms known to be true in the model: 0, (383).

A list of all forms known to be false in the model: 1, 14, 15, 20, 21, 23, 30, 33, 40, 45, 47, 49, 50, 60, 61, 62, 66, 67, 68, 85, 89, 90, 91, 95, 101, 107, 109, 112, 114, 121, 122, 123, 133, 147, 149, 151, 164, 165, 202, 213, 218, 231, 250, 255, 256, 257, 258, 259, 260, 261, 262, 264, 270, 284, 286, 292, 295, 303, 317, 322, 323, 324, 327, 331, 332, 333, 334, 335, 343, 345, 346, 347, 359, (383).)

$\mathcal{M}8$: Apter's Model. Suppose $\mathcal{M} \models$ ZFC $+$ ``There are cardinals $\kappa <\delta< \lambda$ such that $\kappa$ is a supercompact limit of supercompact cardinals, $\lambda$ is a measurable cardinal, and $\delta$ is $\lambda$ supercompact.'' (See, for example, Drake\index{Author}{Drake} \cite{bib:1974}, Kanamori/Magador\index{Author}{Kanamori/Magador} \cite{bib:1978}, or Solovay/Reinhardt/Kanamori\index{Author}{Solovay/Reinhardt/Kanamori} \cite{bib:1978} for information about large cardinals.) $\mathcal{M}8$ is constructed by first forcing over the ground model $\mathcal{M}$, constructing an inner model $\mathcal{M}^{\prime}$, doing an additional forcing argument over $\mathcal{M}^{\prime}$, and then constructing the final inner model $\mathcal{M}8$. In this model $\aleph_{n}$, for $3\leq n\leq\omega$ are singular cardinals so 25 ($\aleph_{\alpha+1}$ is regular for all ordinals $\alpha$.) is false. Also, the Axiom of Choice for a denumerable number of sets, 8, is false. (Supercompact and measurable cardinals are used to construct the model.)

In $\mathcal{M}8$, forms 8, 25, and 91 are false. References Apter\index{Author}{Apter} \cite{bib:1985a}, notes 18 and 20.

(A minimal list of forms implying each form known to be true in the model: 0.

A minimal list of forms at least one of which is implied by each form known to be false in the model: 8, 25, 91.

A list of all forms known to be true in the model: 0, (383).

A list of all forms known to be false in the model: 1, 8, 20, 21, 23, 25, 39, 40, 43, 44, 51, 66, 67, 86, 87, 89, 90, 91, 95, 101, 109, 112, 113, 114, 133, 147, 149, 164, 174, 181, 192, 202, 218, 255, 256, 257, 258, 259, 260, 261, 262, 264, 286, 292, 333, 334, 335, 345, 347, 359, (383).)

$\mathcal{M}9$: Feferman/Levy Model. Assume the ground model, $\mathcal{M}$, satisfies ZF + GCH (the Generalized Continuum Hypothesis). Let $P$, the set of forcing conditions, be the set of all finite sets $p$ of triples $(n,\,i,\, \alpha)$ that satisfy: (i) $n\in \omega,\,i\in\omega, \alpha\in\omega_{n}$; and (ii) $(n,\,i,\,\alpha)$, $(n,\,i,\,\beta)\in p$ implies $\alpha=\beta$. Let $G$ be a generic subset of $P$ and let $\mathcal{M}[G]$ be the corresponding generic extension of $\mathcal{M}$. $\mathcal{M}9$ is a symmetric submodel of $\mathcal{M}[G]$. In this model, $\mathbb{R}\subseteq\bigcup\{R_{n}:n\in \omega\}$, where $|R_{n}|=\aleph_{n+1}$ for $n\in \omega$, and $|\aleph_{n}|=|\aleph_{0}|$ for each $n\in \omega$. Therefore, the set of real numbers is a countable union of countable sets (38 is false).

In $\mathcal{M}$9, form 38 is false. References Feferman/Levy\index{Author}{Feferman/Levy} \cite{bib:1963}, Cohen\index{Author}{Cohen} \cite{bib:1966} (p 143), Jech\index{Author}{Jech} \cite{bib:1973b} (p 142).

(A minimal list of forms implying each form known to be true in the model: 0.

A minimal list of forms at least one of which is implied by each form known to be false in the model: 38.

A list of all forms known to be true in the model: 0, (383).

A list of all forms known to be false in the model: 1, 6, 8, 16, 20, 21, 23, 25, 27, 28, 29, 31, 34, 36, 37, 38, 39, 40, 43, 44, 51, 66, 67, 79, 86, 87, 89, 90, 91, 95, 101, 109, 112, 113, 114, 130, 133, 147, 149, 164, 170, 172, 174, 181, 192, 202, 218, 245, 255, 256, 257, 258, 259, 260, 261, 262, 264, 286, 292, 333, 334, 335, 345, 347, 359, (383).)

$\mathcal{M}10$: Derrick/Drake\index{Author}{Derrick/Drake} Model. Let $\mathcal{M}$ be a model of ZF
+ GCH. Add to $\mathcal{M}$ generic functions $f_{n}$ for each $n\in \omega$, where $f_{n}:\omega_{n}\rightarrow \mathcal{P}(w)$, but do not add $\{f_{n}:n\in \omega\}$. It is shown that in this model $\aleph(2^{\aleph_{0}})=\aleph_{\omega}$ (56 is false). $(\aleph(x)$ is the smallest ordinal that is not similar to a subset of $x$, Hartogs' function.) (In fact, for every $\alpha>0$, it is consistent with ZF that $\aleph(2^{\aleph_{0}})=\aleph_{\alpha}.)$ It is shown by Truss\index{Author}{Truss} that in $\mathcal{M}10$ the sequence $\{\aleph_{n}:n\in \omega\}$ has two distinct minimal upper bounds $2^{\aleph_{0}}$ and $\aleph_{\omega}$. (275 is false). Truss also proves that there is a well ordered set of sets that does not have a choice function (40 is false) and that the Principle of Dependent Choices (43) is false.

In $\mathcal{M}10$, forms 40, 43, 56, 91 and 275 are false. References Derrick/Drake \cite{bib:1967b}, Mathias\index{Author}{Mathias} \cite{bib:1979} p 125, Truss \cite{bib:1974b}, note 18.

(A minimal list of forms implying each form known to be true in the model: 0.

A minimal list of forms at least one of which is implied by each form known to be false in the model: 43, 56, 91, 275.

A list of all forms known to be true in the model: 0, (383).

A list of all forms known to be false in the model: 1, 40, 43, 44, 56, 66, 67, 87, 89, 90, 91, 95, 101, 109, 112, 114, 133, 147, 149, 164, 174, 192, 202, 218, 255, 256, 257, 258, 259, 260, 261, 262, 264, 275, 286, 292, 333, 334, 335, 345, 347, 359, (383.))

$\mathcal{M}11$: Forti/Honsell Model. Using a model of ZF + V = L for the ground model, the authors construct a generic extension, $\mathcal{M}$, using Easton\index{Author}{Easton} forcing which adds $\kappa$ generic subsets to each regular cardinal $\kappa$. (See Easton \cite{bib:1970}.) Then, $\mathcal{M}11$ is a symmetric submodel of $\mathcal{M}$. In this model, the Kinna-Wagner Principle (15) and the Axiom of Choice for a family of well orderable sets (60) are true, but for any partially ordered set $(X,\,\precsim)$ there is a bijection $f:X\rightarrow Y$ such that for all $x,\, y\in X,\,x\precsim\ y\leftrightarrow|f(x)|\leq|f(y)|$ (105 is false). Since form 64 (There is no amorphous set.) is true (30 implies 64), form 133 (Every set is either well orderable or has an infinite amorphous subset.) is false. Form 163 (Every set is either well orderable or has an infinite Dedekind finite subset.) is also false because 43 implies 8 and 8 + 163 implies AC (Brunner\index{Author}{Brunner} \cite{bib:1982a}). Form 15 implies the Ordering Principle (30) and form 60 implies that every linearly ordered set that is almost well orderable is well orderable (294). (See note 25 for definitions.) However, 30 + 144 (Every set is almost well orderable.) + 231 iff AC. Thus, form 144 is false.

$\mathcal{M}11 \models$ 15, 60, and 165, but 91, 105, 133, 144, 152, and 163 are false. References Forti/Honsell\index{Author}{Forti/Honsell} \cite{bib:1985}, Brunner\index{Author}{Brunner} \cite{bib:1982a}, Easton\index{Author}{Easton} \cite{bib:1970}, notes 25 and 120(6, 34, and 43).

(A minimal list of forms implying each form known to be true in the model: 15, 60, 165.

A minimal list of forms at least one of which is implied by each form known to be false in the model: 91, 105, 144, 152, 163.

A list of all forms known to be true in the model: 0, 5, 10, 15, 18, 30,
32, 33, 45, 46, 47, 48, 60, 61, 62, 64, 80, 82, 83, 84, 85, 88, 93, 102, 111, 119, 120, 121, 122, 127, 132, 140, 142, 146, 150, 151, 165, 178, 198, 213, 216, 231, 249, 250, 268, 269, 270, 271, 276, 280, 283, 285, 288, 293, 294, 295, 296, 308, 314, 322, 323, 324, 327, 329, 330, 336, 342, 349, 350, 355, 356, 357, 358, 373, 374, 376, 377, 378, 379, 380, (383).

A list of all forms known to be false in the model: 1, 59, 66, 67, 89, 90, 91, 95, 105, 109, 112, 114, 133, 144, 147, 149, 152, 163, 164, 179, 202, 214, 218, 264, 292, 333, 334, 335, (383).)

$\mathcal{M}12(\aleph)$: Truss'\index{Author}{Truss} Model I. This is a variation of Solovay's model, $\mathcal{M}5(\aleph)$ in which $\aleph$ is singular. (Truss uses the method of collapsing cardinal which is due to A. Levy, but, unlike $\mathcal{M}5$, $\aleph$ is not an inaccessible cardinal.) Let $\mathcal{M}$ be a countable model of ZF + V = L and let $\aleph$ be a limit cardinal in $\mathcal{M}$. For each ordinal $\alpha$, the set of conditions $Q^{\alpha}$ is the set of finite sets $p$ of triples, $(\beta,\,n,\,\gamma)$ where $\gamma<\beta<\alpha$ and $n\in \omega$ such that if $(\beta,\,n,\,\gamma_{1})$, $(\beta,\,n,\,\gamma_{2})\in p$, then $\gamma_{1}=\gamma_{2}$. $Q_{\alpha}$ is the set of finite sets $p$ of pairs of the form $(n,\,\gamma)$ where $\gamma<\alpha$ and $n\in \omega$ such that if $(n,\,\gamma_{1})$, $(n,\, \gamma_{2})\in p$ then $\gamma_{1}=\gamma_{2}$.
A function $f: \omega\rightarrow\alpha$ is called an $\mathcal{M}$-\emph{generic collapsing map} if the set of finite subsets of $\{(n,\,\gamma):f(n)=\gamma\}$ is an $\mathcal{M}$-generic subset of $Q_{\alpha}$. Let $G$ be an $\mathcal{M}$-generic subset of $Q^{\aleph}$ then, $G_{\alpha}$, the projection of $G$ onto $Q_{\alpha}$ is an $\mathcal{M}$-generic subset of $Q_{\alpha}$, for each $\alpha< \aleph$. Let $f_{\alpha}=\bigcup G_{\alpha}$. Then $\mathcal{M}12(\aleph)$ is the smallest model of ZF containing the same ordinals as $\mathcal{M}$ and each $\langle f_{\beta}:\beta<\alpha\rangle$ for $\alpha<\aleph$. In $\mathcal{M}12(\aleph)$, $\aleph=\omega_{1}$. If $\aleph$ is singular, $\omega_{1}$ is singular so 34 is false. On the other hand, the union of a denumerable set of denumerable sets of reals is denumerable (6 is true) and the union of a well ordered set of well orderable sets can be well ordered (231 is true). $\mathcal{P}(\mathcal{P}(\omega))$ cannot be linearly ordered, so 30 is false. Truss also proves that the Principle of Dependent Choices (43) is false, every set of reals is Borel (273 is false), and every uncountable set of reals has a perfect subset (169 is false). Form 40 $(C(WO,\,\infty))$ is false because 40 implies 31 and form 122 $(C(WO,\, <\aleph_{0}))$ is true because 231 implies 122. Therefore, form 328 $(MC(WO,\,\infty))$ is false because 122 + 328 $\rightarrow$ 40

$\mathcal{M}12(\aleph)\models$ 6 and 231, but 30, 31, 34, 43, 106, 152, 169, 273, and 328 are false. References Truss \cite{bib:1974b} and note 43 and 120(43 and 55).

(A minimal list of forms implying each form known to be true in the model: 6, 231.

A minimal list of forms at least one of which is implied by each form known to be false in the model: 30, 34, 106, 152, 169, 273, 328.

A list of all forms known to be true in the model: 0, 5, 6, 10, 18, 35, 36,
38, 47, 48, 80, 108, 111, 122, 151, 216, 231, 249, 250, 288, 294, 327, 358, 373, 374, (383).

A list of all forms known to be false in the model: 1, 8, 14, 15, 16, 20, 21, 23, 25, 27, 28, 29, 30, 31, 34, 39, 40, 43, 44, 49, 50, 51, 66, 67, 86, 87, 89, 90, 95, 101, 106, 109, 112, 113, 114, 133, 149, 152, 169, 172, 174, 181, 188, 192, 193, 202, 214, 218, 245, 255, 256, 257, 258, 259, 260, 261, 262, 264, 273, 286, 292, 295, 303, 317, 328, 333, 334, 335, 345, 347, 359, 363, 364, (383).)

$\mathcal{M}13$: Feferman/Solovay Model. This model is an extension of $\mathcal{M}2$ in which there are $\omega_{1}$ generic real numbers, but no set to collect them. Solovay has shown that in this model, the Axiom of Choice for a well ordered family of sets (40) and the Principle of Dependent Choices (DC, 43) are true, while the extension of DC to $\omega_{1}$ (44) is false. Since 40 implies 8 $(C(\aleph_{0},\,\infty))$, it follows from Brunner \cite{bib:1982a} that in this model there is a set that cannot be well ordered and does not have an infinite Dedekind finite subset, (163 is false). (Form 8 plus form 163 iff AC.)

$\mathcal{M}13 \models$ 40, 43, and 165, but 44, 67, 91, 152 and 163 are false. References Felgner\index{Author}{Felgner} \cite{bib:1971a} p 166, Brunner\index{Author}{Brunner} \cite{bib:1982a}, notes 18 and 120(41, 34 and 46).

(A minimal list of forms implying each form known to be true in the model: 40, 165.

A minimal list of forms at least one of which is implied by each form known to be false in the model: 44, 91, 152, 163.

A list of all forms known to be true in the model: 0, 5, 6, 8, 9, 10, 13,
16, 17, 18, 19, 24, 26, 27, 28, 29, 31, 32, 34, 35, 36, 37, 38, 39, 40, 42, 43, 47, 48, 57, 58, 64, 73, 74, 75, 77, 78, 80, 82, 83, 84, 86, 92, 94, 98, 104, 106, 108, 111, 113, 119, 122, 124, 126, 127, 128, 131, 132, 150, 151, 154, 155, 165, 182, 185, 194, 196, 198, 199, 208, 209, 211, 216, 217, 231, 243, 249, 250, 273, 279, 282, 288, 294, 296, 304, 322, 324, 325, 327, 328, 330, 336, 337, 338, 339, 340, 341, 342, 350, 351, 352, 353, 354, 355, 357, 358, 360, 361, 362, 363, 364, 373, 374, 375, 376, 377, 378, 380, (383).

A list of all forms known to be false in the model: 1, 44, 66, 67, 89, 90, 91, 95, 109, 112, 114, 133, 147, 149, 152, 163, 164, 202, 214, 218, 264, 292, 333, 334, 335, (383).)

$\mathcal{M}14$: Morris'\index{Author}{Morris} Model. This is an extension of Mathias' model, $\mathcal{M}3$. Morris shows that the cofinal subsets of the partially ordered set $b$, the set of generic reals, can be systematically well ordered. He then proves that in this model, the Cofinality Principle (51) and the Ordering Principle (30) are true, but the Order Extension Principle (49) is false. Forms 30 and 294 (Every linearly ordered $W$-set is well orderable.) are both true in this model (51 implies 294). Since 30 + 294 + 144 (Every set is a $W$-set.) implies AC, it follows that 144 is false. (See note 25 for definitions.) Moreover, form 64 (There is no amorphous set.) is true (30 implies 64), so form 133 (Every set is either well orderable or has an infinite amorphous subset.) is false because 64 + 133 implies AC. The Boolean Prime Ideal Theorem (14) is false because 14 implies 49, and $C(\infty,\,<\aleph_{0})(62)$ is true because 30 implies 62. It was shown by Blass (see note 33) that 62 +
99 (Rado's Selection Lemma) implies 14, so form 99 is false.

$\mathcal{M}14 \models$ 30 and 51, but 49, 91, 99, 133, and 144 are false. References Morris \cite{bib:1969}, notes 18, 25 and 33.

(A minimal list of forms implying each form known to be true in the model: 30, 51.

A minimal list of forms at least one of which is implied by each form known to be false in the model: 49, 91, 99, 144.

A list of all forms known to be true in the model: 0, 10, 13, 18, 19, 25,
30, 33, 34, 38, 45, 46, 47, 48, 51, 58, 61, 62, 64, 77, 80, 83, 84, 88, 93, 102, 104, 108, 111, 120, 121, 122, 127, 132, 140, 142, 146, 178, 182, 185, 198, 199, 208, 216, 249, 250, 268, 269, 276, 280, 283, 285, 288, 293, 294, 308, 315, 316, 327, 336, 342, 358, 373, 374, 377, 378, (383).

A list of all forms known to be false in the model: 1, 14, 49, 50, 66, 67, 89, 90, 91, 95, 99, 109, 112, 114, 133, 144, 147, 149, 164, 179, 202, 218, 264, 292, 303, 317, 333, 334, 335, 345, (383).)

$\mathcal{M}15$: Feferman/Blass Model. Blass\index{Author}{Blass} constructs a model similar to Feferman's model, $\mathcal{M}2$. Suppose $\mathcal{M}$ = ZF + V = L. By forcing with finite partial functions from $\omega \times \omega$ into 2, a generic extension $\mathcal{M}[G]$ of $\mathcal{M}$ is obtained that contains an $\omega$-sequence of generic reals $a_{n}\subseteq \omega$. For each $x\subseteq\omega$,
let $\delta(x)$ be the set of reals whose symmetric difference with $x$ is finite. Let $f$ be the function in $\mathcal{M}[G]$ such that for all $n\in\omega,\,f(n)=\{\delta(a_{n}),\, \delta(\omega-a_{n})\}$. Let $S=\bigcup_{n\in\omega}(\delta(a_{n})\cup\delta(\omega-a_{n}))\cup\{f\}$. $\mathcal{M}15$ is the submodel of $\mathcal{M}[G]$ that is hereditarily ordinal definable from $S$ and $\bigcup\ S$. In this model, every set is almost well orderable (144 is true) and all ultrafilters are principal (206 is false). Howard\index{Author}{Howard} has shown that if form 206 is false, then {\L}os' Theorem (253) is true.

$\mathcal{M}15 \models$ 144 and 253, but 206 is false. References Blass \cite{bib:1977b}, Feferman\index{Author}{Feferman} \cite{bib:1965}, Howard \cite{bib:1975}.

(A minimal list of forms implying each form known to be true in the model: 144, 253.

A minimal list of forms at least one of which is implied by each form known to be false in the model: 206.

A list of all forms known to be true in the model: 0, 125, 144, 253, (383).

A list of all forms known to be false in the model: 1, 14, 50, 63, 66, 67, 70, 79, 89, 90, 91, 95, 99, 109, 112, 114, 123, 130, 133, 147, 149, 164, 202, 206, 218, 225, 264, 292, 303, 317, 333, 334, 335, 345, (383).)

$\mathcal{M}16(n)$: Monro's Model I. Monro has constructed an $\omega$-sequence of models of ZF such that,

$(\ast)$ For each ordinal $\alpha,\,\mathcal{P}^{n}(\alpha)$ is the same in the $n$th model and the $n + 1$st model.

\noindent $(\mathcal{P}^{0}(\alpha)=\alpha$, and $\mathcal{P}^{n+1}(\alpha)=\mathcal{P}(\mathcal{P}^{n}(\alpha)).)$ He starts with a countable transitive model $\mathcal{M} \models$ ZF + V
= L. Let $H_{\kappa}(A,\,B)$ be the set of all functions $f\subseteq A\times B$ such that $|\mathrm{dom}(f)|<\kappa$. Let $\mathcal{M}^{0}=\mathcal{M}$ and $J^{0}=\omega$. Then $\mathcal{M}^{k+1}$ is obtained from $\mathcal{M}^{k}$ as follows. Let $G^{k+1}$ be $H_{\omega}(J^{k}\times J^{k},\,2)$-generic over $\mathcal{M}_{k}$. For $r\in J^{k}$, let $G_{r}^{k+1}= \{(s,\,a):(r,\,s,\,a)\in\bigcup(G^{k+1}\}$. $(G_{r}^{k+1}:J^{k}\rightarrow 2.)$ Then $J^{k+1}=\{G_{r}^{k+1}:r\in J^{k}\}$ and $\mathcal{M}^{k+1}=\mathcal{M}^{k}[G^{k+1}]$. Let $K(n)$ be the statement: For every set $x$ there is an ordinal $\alpha$ and a 1-1 function $f:x\longrightarrow \mathcal{P}^{n}(\alpha)$. $(K(n)$ is form $81(n)$; $K(0)$ is equivalent to AC; and $K(1)$ is equivalent to the Kinna-Wagner Selection Principle, form 15.) Monro proves that if $\mathcal{M}_{1}$ and $\mathcal{M}_{2}$ are two models such that for each $\alpha,\,\mathcal{P}^{n}(\alpha)$ is the same in each model and, in addition, $\mathcal{M}_{1}\neq \mathcal{M}_{2}$, then $K(n-1)$ is false in each model. Since, each successive pair of models in the sequence constructed has the property $(\ast)$, it follows that for each $n\in \omega$, there is a model of ZF in which 81$(n)$ is false. Monro also proves that for each $n\in \omega$, $\mathcal{M}^{n}\models K(n+1)$. Consequently, it follows that for all $n>0,\,K(n+1) \not\rightarrow\ K(n-1)$.

In $\mathcal{M}16(n)$, forms 15, 81$(n)$, and 91 are false. References Monro \cite{bib:1972}, \cite{bib:1973b}, note 18.

(A minimal list of forms implying each form known to be true in the model: 0.

A minimal list of forms at least one of which is implied by each form known to be false in the model: 15, 81, 91.

A list of all forms known to be true in the model: 0, (383).

A list of all forms known to be false in the model: 1, 15, 66, 67, 81, 89, 90, 91, 95, 109, 112, 114, 133, 147, 149, 164, 202, 218, 264, 292, 333, 334, 335, (383).)

$\mathcal{M}17$: Gitik's Model. Using the assumption that for every ordinal $\alpha$ there is a strongly compact cardinal $\kappa$ such that $\kappa >\alpha$, Gitik\index{Author}{Gitik} extends the universe $V$ by a filter $G$ generic over a proper class of forcing conditions. (Inaccessible cardinals are used in the construction.) In $V[G]$, he constructs a symmetric submodel $\mathcal{M}17$ in which every infinite set is a countable union of sets of smaller cardinality. Thus, there is no regular uncountable cardinal (104 is false), there is no aleph whose cofinality is greater than $\aleph_{0}$ (182 is false), and $\Omega=$ On (315 is false). See note 55. In $\mathcal{M}17$, forms 91, 104, 182, and 315 are false. References Gitik \cite{bib:1980}, \cite{bib:1985}, notes 18 and 55.

(A minimal list of forms implying each form known to be true in the model: 0.

A minimal list of forms at least one of which is implied by each form known to be false in the model: 91, 182, 315.

A list of all forms known to be true in the model: 0, (383).

A list of all forms known to be false in the model: 1, 8, 16, 20, 21, 23, 25, 27, 28, 29, 31, 34, 39, 40, 43, 44, 51, 66, 67, 86, 87, 89, 90, 91, 95, 101, 104, 109, 112, 113, 114, 133, 147, 149, 159, 164, 172, 174, 181, 182, 191, 192, 202, 218, 245, 255, 256, 257, 258, 259, 260, 261, 262, 264, 286, 292, 315, 333, 334, 335, 345, 347, 359, (383).)

$\mathcal{M}18$: Shelah's Model I. Shelah modified Solovay's model, $\mathcal{M}5$, and constructed a model without using an inaccessible cardinal in which the Principle of Dependent Choices (43) is true and every set of reals has the property of Baire (142 is false). Stern\index{Author}{Stern} has shown that in this model there is a non-measurable set of real numbers (93 is true). Since 43 implies 8 $(C(\aleph_{0},\,\infty))$, it follows from Brunner\index{Author}{Brunner} \cite{bib:1982a} that in this model there is a set that cannot be well ordered and does not have an infinite Dedekind finite subset, (163 is false). (Form 8 plus form 163 iff AC.)

$\mathcal{M}18 \models$ 43 and 93, but 142 and 163 are false. References Shelah \cite{bib:1984}, Solovay\index{Author}{Solovay} \cite{bib:1970}, Stern \cite{bib:1985}, and Brunner \cite{bib:1982a}.

(A minimal list of forms implying each form known to be true in the model: 43, 93.

A minimal list of forms at least one of which is implied by each form known to be false in the model: 142, 163.

A list of all forms known to be true in the model: 0, 5, 6, 8, 9, 10, 13,
16, 17, 18, 19, 24, 26, 27, 28, 29, 31, 32, 34, 35, 36, 37, 38, 42, 43, 57, 64, 73, 74, 75, 77, 78, 80, 82, 83, 84, 92, 93, 94, 98, 104, 106, 108, 113, 119, 124, 126, 127, 128, 131, 132, 150, 154, 155, 182, 185, 194, 198, 199, 209, 211, 216, 217, 243, 249, 273, 279, 282, 288, 296, 304, 325, 336, 337, 338, 339, 340, 341, 342, 350, 351, 352, 353, 354, 355, 357, 358, 360, 361, 362, 363, 364, 373, 374, 375, 376, 377, 378, 380, (383).

A list of all forms known to be false in the model: 1, 14, 15, 30, 49, 50, 52, 60, 61, 62, 63, 66, 67, 68, 70, 79, 85, 88, 89, 90, 91, 95, 99, 107, 109, 112, 114, 123, 130, 133, 142, 147, 149, 163, 164, 201, 202, 213, 218, 221, 222, 225, 264, 270, 284, 287, 292, 295, 303, 310, 311, 313, 317, 323, 326, 331, 332, 333, 334, 335, 343, 345, 346, 372, (383).)

$\mathcal{M}19(\aleph)$: Monro's Model II. Let $\mathcal{M}$ be a countable transitive model of ZF + V = L and let $\aleph$ be a regular cardinal in $\mathcal{M}$. A forcing condition is a partial function, $p$, of cardinality less than $\aleph$ such that $p:\aleph\times\aleph\rightarrow 2$. Let $G$ be a generic subset of the set of forcing conditions and let $\mathcal{M}[G]$ be the corresponding generic extension.

Let $p(\mu)=\{(\nu,\,i):(\mu,\,\nu,\,i)\in p\}$, for each $\mu<\aleph$, let $G_{\mu}=\{p(\mu):p\in G\}$, and $G^{\ast}=\{G_{\mu}:\mu<\aleph\}$. Then $\mathcal{M}19(\aleph)$ is the set of all sets in $\mathcal{M}[G]$ that are constructible from $G^{\ast}.$ Monro\index{Author}{Monro} proves that $\mathcal{M}19(\aleph)$ has the same ordinals and alephs as $\mathcal{M}$ and that $G^{\ast}$ is an infinite Dedekind finite set in $\mathcal{M}19(\aleph)$. He also proves that if $x$ is any set that cannot be well ordered, then x can be mapped onto $\aleph$. Thus, $G^{\ast}$ can be mapped onto $\aleph$ so form 160 is false.

In $\mathcal{M}19(\aleph)$, forms 91 and 160 are false. References Monro \cite{bib:1975} and note 18.

(A minimal list of forms implying each form known to be true in the model: 0.

A minimal list of forms at least one of which is implied by each form known to be false in the model: 91, 160.

A list of all forms known to be true in the model: 0, (383).

A list of all forms known to be false in the model: 1, 66, 67, 89, 90, 91, 95, 109, 112, 114, 133, 147, 149, 160, 164, 202, 218, 264, 292, 333, 334, 335, (383).)

$\mathcal{M}20$: Felgner's Model I. Let $\mathcal{M}$ be a model of ZF +
V = L. Felgner\index{Author}{Felgner} defines forcing conditions that force $\aleph_{\omega}$ in $\mathcal{M}$ to be $\aleph_{1}$. He then proves that in the resulting model, $\mathcal{M}20$, every set can be linearly ordered (30 is true) and the union of a well ordered set of well orderable sets can be well ordered (231 is true), but there is a denumerable set of denumerable sets which does not have a denumerable union (31 is false). Also, in this model every set of well orderable sets has a choice function (60 is true). Form 60 implies 294 (Every linearly ordered $W$-set is well orderable.), and 30 + 294 + 144 (Every set is a $W$-set.) implies AC. It follows that 144 is false. (See note 25 for definitions.) Since form 64 (There is no amorphous set.) is true (30 implies 64), form 133 (Every set is either well orderable or has an infinite amorphous subset.) is false. Form 40 $(C(WO,\,\infty))$ is false because 40 implies 31 and form 122 $(C(WO,\,<\aleph_{0}))$ is true because 30 implies 122. Therefore, form 328 $(MC(WO,\,\infty))$ is false because 122 + 328 $\rightarrow$ 40.

$\mathcal{M}20 \models$ 30, 60, 165, and 231, but 31, 91, 131, 133, 144, 152, and 328 are false. References Felgner \cite{bib:1974}, notes 18, 25, and 120(34, 43, and 49).

(A minimal list of forms implying each form known to be true in the model: 30, 60, 165.

A minimal list of forms at least one of which is implied by each form known to be false in the model: 31, 91, 131, 144, 152.

A list of all forms known to be true in the model: 0, 5, 10, 18, 30, 32,
33, 45, 46, 47, 48, 60, 61, 62, 64, 80, 83, 85, 88, 93, 102, 111, 119, 120, 121, 122, 127, 132, 140, 142, 146, 150, 151, 165, 178, 198, 213, 216, 231, 249, 250, 268, 269, 270, 271, 276, 280, 283, 285, 288, 293, 294, 308, 314, 323, 324, 327, 329, 330, 336, 342, 349, 350, 356, 357, 358, 373, 374, 377, 378, 380, (383).

A list of all forms known to be false in the model: 1, 8, 16, 20, 21, 23, 27, 28, 29, 31, 39, 40, 43, 44, 66, 67, 76, 86, 87, 89, 90, 91, 95, 101, 106, 109, 112, 113, 114, 126, 131, 133, 144, 147, 149, 152, 164, 174, 179, 181, 188, 192, 193, 202, 214, 218, 255, 256, 257, 258, 259, 260, 261, 262, 264, 286, 292, 328, 333, 334, 335, 345, 347, 359, (383).)

$\mathcal{M}21$: Felgner's Model II. Suppose $\mathcal{M}\models$ ZF + V
= L. Define $B=\{f:(\exists\alpha< \omega_{1})f:\alpha\rightarrow\omega\}$. $\mathcal{M}21$ is constructed using forcing conditions which are finite partial functions from $\omega\times B$ to $\{0, 1\}$. In this model, there is a well ordered set of well orderable sets whose union cannot be well ordered (231 is false), but the Axiom of Choice for a denumerable set (8) is true. It follows from Brunner\index{Author}{Brunner} \cite{bib:1982a} that in this model there is a set that cannot be well ordered and does not have an infinite Dedekind finite subset, (163 is false). (Form 8 plus form 163 iff AC.)

$\mathcal{M}21\models$ 8, but 91, 163, and 231 are false. References Felgner\index{Author}{Felgner} \cite{bib:1974}, Brunner\index{Author}{Brunner} \cite{bib:1982a}, and note 18.

(A minimal list of forms implying each form known to be true in the model: 8.

A minimal list of forms at least one of which is implied by each form known to be false in the model: 91, 163, 231.

A list of all forms known to be true in the model: 0, 5, 6, 8, 9, 10, 13, 16, 17, 18, 19, 24, 26, 27, 28, 29, 31, 32, 34, 35, 36, 37, 38, 57, 64, 73, 74, 75, 80, 82, 83, 84, 92, 94, 98, 104, 108, 119, 124, 126, 127, 128, 131, 132, 150, 182, 185, 194, 198, 199, 209, 216, 217, 249, 282, 288, 296, 304, 325, 336, 337, 338, 339, 340, 341, 342, 350, 351, 352, 353, 354, 355, 357, 358, 360, 361, 362, 373, 374, 376, 377, 378, 380, (383).

A list of all forms known to be false in the model: 1, 40, 60, 66, 67, 89, 90, 91, 95, 101, 109, 112, 114, 133, 147, 149, 163, 164, 202, 218, 231, 255, 256, 257, 258, 259, 260, 261, 262, 264, 286, 292, 333, 334, 335, 347, 359, (383).)

$\mathcal{M}22$: Plotkin's Model I. Let $T$ be a complete first order theory with equality which has infinite models and is $\aleph_{0}$-categorical. Let $\mathcal{M}$ be a countable standard transitive model of ZFC which contains a countable model $\mathfrak{A}$ of $T$. It is shown in Plotkin\index{Author}{Plotkin} \cite{bib:1969} that there is a generic extension of $\mathcal{M}$ that contains an isomorphic copy, $\mathfrak{A}^{\prime}$, of $\mathfrak{A}$, such that in the symmetric submodel constructed, $\mathfrak{A}^{\prime}$ is infinite, Dedekind finite, and the relations of $T$ are definable. Let $p$ be an odd prime and let $T$ be the first order theory of an infinite extra-special $p$-group of exponent $p$. (See note 24 for definitions.) It is shown in Felgner \cite{bib:1975} that this theory is complete and $\aleph_{0}$-categorical. $\mathcal{M}22$ is the resulting model. Plotkin \cite{bib:1981} proves the countable extra-special $p$-groups of exponent $p$ in $\mathcal{M}22$ are all locally finite and have no infinite Abelian subgroups. (Locally finite means that every finite subset generates a finite subgroup.) Thus, form 176 is false. (Since the groups constructed are infinite and Dedekind finite, form 9 is also false.)

In $\mathcal{M}22$, forms 9, 91 and 176 are false. References Plotkin \cite{bib:1969}, \cite{bib:1981}, Felgner \cite{bib:1975}, notes 18 and 24.

(A minimal list of forms implying each form known to be true in the model: 0.

A minimal list of forms at least one of which is implied by each form known to be false in the model: 9, 91, 176.

A list of all forms known to be true in the model: 0, (383).

A list of all forms known to be false in the model: 1, 2, 3, 4, 7, 8, 9, 39, 40, 41, 43, 44, 66, 67, 71, 86, 87, 89, 90, 91, 95, 100, 101, 109, 112, 113, 114, 129, 133, 147, 149, 152, 161, 164, 168, 174, 176, 181, 192, 202, 214, 218, 255, 256, 257, 258, 259, 260, 261, 262, 264, 286, 292, 333, 334, 335, 345, 347, 359, (383).)

$\mathcal{M}23$: Hodges' Model. Let $\mathcal{M}$ be a countable transitive model of ZFC + V = L and let $\kappa$ be a regular cardinal in $\mathcal{M}$. Let $A=\sum\nolimits_{\alpha<\kappa}A_{\alpha}$, where each $A_{\alpha}$ is a copy of the cyclic group $Z(4)$. Using forcing conditions of the form $p:A\times\kappa \times\kappa \rightarrow 2$, where $p$ is a partial function with domain of cardinality less than $\kappa$, it is shown in Hodges\index{Author}{Hodges} \cite{bib:1976b} that there is a generic extension $\mathcal{M}[G]$ in which there is a generic copy of $A$. It is also shown in Hodges \cite{bib:1976b} that $\mathcal{M}[G]$ can be cut down to a transitive model of ZF which contains $A$ and all elements are definable from elements of $\mathcal{M}$ and sequences of generic sets of length less than $\kappa$. This is the model $\mathcal{M}23$, It is shown in Hodges \cite{bib:1980} that in $\mathcal{M}23,\,A$ is an Abelian group with no divisible hull (180 is false). (If $A$ and $B$ are groups, ``$B$ is a divisible hull of $A$'' means $B$ is a divisible group, $A$ is a subgroup of $B$ and for every non-zero $b\in B,\,\exists n\in w$ such that $0\neq nb\in A$. An Abelian group $B$ is divisible if $n$ divides $a$ for all $a\in B$ and positive $n\in\omega.)$

In $\mathcal{M}23$, forms 91 and 180 are false. References Hodges\index{Author}{Hodges} \cite{bib:1976b}, \cite{bib:1980}, and note 18.

(A minimal list of forms implying each form known to be true in the model: 0.

A minimal list of forms at least one of which is implied by each form known to be false in the model: 91, 180.

A list of all forms known to be true in the model: 0, (383).

A list of all forms known to be false in the model: 1, 66, 67, 89, 90, 91, 95, 109, 112, 114, 133, 147, 149, 164, 180, 202, 218, 264, 292, 333, 334, 335, (383).)

\noindent $\mathcal{M}24$: Blass' Model. Let $\mathcal{M}$ be a countable transitive model of ZFC + V = L. For each regular cardinal $\xi$ adjoin generic sets $b_{\xi}^{i}\subseteq\xi$ for each $i\in \mathbb{Z}$ using Easton forcing. (See Easton\index{Author}{Easton} \cite{bib:1970}. There is a proper class of generic sets.) Call the resulting forcing extension $\mathcal{M}^{\prime}$. In $\mathcal{M}^{\prime}$, let $a_{\xi}^{i}$ be the set of all subsets of $\xi$ whose symmetric difference with $b_{\xi}^{i}$ is a constructible set of cardinality less than $\xi$. For each regular $\xi$ and each $i\in \mathbb{Z}$, define $S(a_{\xi}^{i})=a_{\xi}^{i+1}.$ Let $\mathcal{M}24$ be the submodel of $\mathcal{M}^{\prime}$ consisting of all sets that are hereditarily ordinal definable from $S$ and the $b_{\xi}^{i}$'s. Blass\index{Author}{Blass} \cite{bib:1979} has shown that in this model there is no non-trivial injective Abelian group. (See note 60 for definitions.)

In $\mathcal{M}24$, forms 91 and 190 are false. References Blass \cite{bib:1979}, Easton \cite{bib:1970}, notes 18 and 60.

(A minimal list of forms implying each form known to be true in the model: 0.

A minimal list of forms at least one of which is implied by each form known to be false in the model: 91, 190.

A list of all forms known to be true in the model: 0, (383).

A list of all forms known to be false in the model: 1, 66, 67, 89, 90, 91, 95, 109, 112, 114, 133, 147, 149, 164, 189, 190, 191, 202, 218, 264, 292, 333, 334, 335, (383).)

$\mathcal{M}25:$ Freyd's\index{Author}{Freyd} Model. Using topos-theoretic methods due to Fourman\index{Author}{Fourman}, Freyd constructs a Boolean-valued model of ZF in which every well ordered family of sets has a choice function (40 is true), but $C(|\mathbb{R}|,\,\infty)(181)$ is false. Freyd describes an algebraic theory with two constants, three unary operations, and two binary operations. (See Freyd \cite{bib:1980}, p 112.) Let $F$ be the free algebra generated by the two constants and let $\mathcal{B}(F)$ be the corresponding category of Boolean sheaves. The Fourman interpretation arising from $\mathcal{B}(F)$ is the desired model. (See, for example, MacLane\index{Author}{MacLane} \cite{bib:1971} and Johnstone\index{Author}{Johnstone} \cite{bib:1977} for background information on category theory.)

Since 40 implies 8 $(C(\aleph_{0},\,\infty))$, it follows from Brunner\index{Author}{Brunner} \cite{bib:1982a} that in this model there is a set that cannot be well ordered and does not have an infinite Dedekind finite subset, (163 is false). (Form 8 plus form 163 iff AC.)

$\mathcal{M}25 \models$ 40 and 165, but 67, 91, 152, 163, and 181 are false. References Freyd \cite{bib:1980}, Fourman \cite{bib:1980}, MacLane \cite{bib:1971}, Johnston\index{Author}{Johnston} \cite{bib:1984}, Brunner \cite{bib:1982a}, notes 18 and 120(24, 41, and 46).

(A minimal list of forms implying each form known to be true in the model: 40, 165.

A minimal list of forms at least one of which is implied by each form known to be false in the model: 91, 152, 163, 181.

A list of all forms known to be true in the model: 0, 5, 6, 8, 9, 10, 13, 16, 17, 18, 19, 24, 26, 27, 28, 29, 31, 32, 34, 35, 36, 37, 38, 39, 40, 42, 43, 47, 48, 57, 58, 64, 73, 74, 75, 77, 78, 80, 82, 83, 84, 86, 92, 94, 98, 104, 106, 108, 111, 113, 119, 122, 124, 126, 127, 128, 131, 132, 150, 151, 154, 155, 165, 182, 185, 194, 196, 198, 199, 208, 209, 211, 216, 217, 231, 243, 249, 250, 273, 279, 282, 288, 294, 296, 304, 322, 324, 325, 327, 328, 330, 336, 337, 338, 339, 340, 341, 342, 350, 351, 352, 353, 354, 355, 357, 358, 360, 361, 362, 363, 364, 373, 374, 375, 376, 377, 378, 380, (383).

A list of all forms known to be false in the model: 1, 66, 67, 89, 90, 91, 95, 109, 112, 114, 133, 147, 149, 152, 163, 164, 181, 202, 214, 218, 264, 292, 333, 334, 335, (383).)

$\mathcal{M}26$: Kanovei's Model I. Starting with a model of ZF + V = L and using forcing techniques due to Jensen\index{Author}{Jensen} \cite{bib:1968}, Kanovei\index{Author}{Kanovei} constructs a model of ZF in which there is an infinite Dedekind finite set $A$ of generic reals that is in the class $\Pi_{n}^{1}$, but there are no infinite Dedekind finite subsets of $\mathbb{R}$ in the class $\varSigma_{n}^{1}$, where $n\in \omega,\,n\geq 2$. Thus, form $199(n)$ is false.

In $\mathcal{M}26$, form $199(n)$ is false. References Kanovei \cite{bib:1978}, Jensen \cite{bib:1968}, and note 61.

(A minimal list of forms implying each form known to be true in the model: 0.

A minimal list of forms at least one of which is implied by each form known to be false in the model: 199.

A list of all forms known to be true in the model: 0, (383).

A list of all forms known to be false in the model: 1, 2, 3, 4, 7, 8, 9, 11, 12, 13, 16, 39, 40, 41, 43, 44, 51, 66, 67, 71, 73, 75, 77, 79, 86, 87, 89, 90, 91, 92, 94, 95, 100, 101, 106, 109, 112, 113, 114, 126, 129, 130, 133, 147, 149, 152, 161, 164, 168, 174, 175, 181, 185, 188, 192, 193, 199, 202, 203, 211, 212, 214, 218, 255, 256, 257, 258, 259, 260, 261, 262, 264, 286, 292, 316, 328, 333, 334, 335, 345, 347, 359, (383).)

$\mathcal{M}27$: Pincus/Solovay Model I. Let $\mathcal{M}_{1}$ be a model of ZFC + V = L. $R\in 2^{x}$ is said to be \emph{random} if $R\in b$ for every constructible Borel set $b$ which has Lebesgue measure 1. Let $\mathcal{M}_{2}=L(R)$, where $R$ is a random element of $2^{\omega^{L}_{1}}.$ (Models like this were also discussed in Sacks\index{Author}{Sacks} \cite{bib:1969} and Solovay\index{Author}{Solovay} \cite{bib:1971}.) In $\mathcal{M}_{2}$, for each $\alpha\in \omega_{1}$, let $r_{\alpha}=R|\alpha\cdot \omega$. Let $G_{\alpha}^{\ast}$ be the set of all subsets of $\alpha\cdot \omega$ which are finite unions of half open intervals of the form $[\gamma\cdot\omega,\,(\gamma+1)\cdot \omega)$, and $G_{\alpha}$ is the set of subsets of $\alpha\cdot \omega$ which differ finitely from a member of $G_{\alpha}^{\ast}$ . $G_{\alpha}$ is an Abelian group under the operation of symmetric difference. Let $F$ be the function on $\omega_{1}$ such that for each $\alpha\in \omega_{1},\,F(\alpha)$ is the orbit of $r_{\alpha}$ under the action of the group $G_{\alpha}$. Then, $\mathcal{M}27$ is the smallest model of ZF containing $F$. Solovay proves that the Principle of Dependent Choices (43) and the Hahn-Banach Theorem (52) are true in $\mathcal{M}27$, but all ultrafilters are principal in this model so 206 is false. Since 43 implies 8 $(C(\aleph_{0},\,\infty))$, it follows from Brunner\index{Author}{Brunner} \cite{bib:1982a} that in this model there is a set that cannot be well ordered and does not have an infinite Dedekind finite subset, (163 is false). (Form 8 plus form 163 iff AC.) Howard\index{Author}{Howard} \cite{bib:1975} has shown that if form 206 is false, then {\L}o\'{s}' Theorem (253) is true. Morillon\index{Author}{Morillon} \cite{bib:1993} has shown that if form 9 is true and form 206 is false then there does not exist an infinite, compact, Hausdorff, extremally disconnected topological space so form 371 is false.

$\mathcal{M}27 \models$ 43, 52, and 253, but 163, 206, 286, and 371 are false. References Pincus/Solovay\index{Author}{Pincus/Solovay} \cite{bib:1977}, Brunner \cite{bib:1982a}, Howard\index{Author}{Howard} \cite{bib:1975}, Morillon\index{Author}{Morillon} \cite{bib:1993}, Sacks\index{Author}{Sacks} \cite{bib:1969}, Solovay\index{Author}{Solovay} \cite{bib:1971}, and note 120(26).

(A minimal list of forms implying each form known to be true in the model: 43, 52, 253.

A minimal list of forms at least one of which is implied by each form known to be false in the model: 163, 206, 286, 371.

A list of all forms known to be true in the model: 0, 5, 6, 8, 9, 10, 13, 16, 17, 18, 19, 24, 26, 27, 28, 29, 31, 32, 34, 35, 36, 37, 38, 42, 43, 52, 57, 64, 73, 74, 75, 77, 78, 80, 82, 83, 84, 92, 93, 94, 98, 104, 106, 108, 113, 119, 124, 126, 127, 128, 131, 132, 142, 150, 154, 155, 182, 185, 194, 198, 199, 209, 211, 216, 217, 221, 222, 223, 243, 249, 253, 273, 279, 280, 282, 287, 288, 296, 304, 309, 325, 336, 337, 338, 339, 340, 341, 342, 350, 351, 352, 353, 354, 355, 357, 358, 360, 361, 362, 363, 364, 373, 374, 375, 376, 377, 378, 380, (383).

A list of all forms known to be false in the model: 1, 14, 50, 63, 66, 67, 70, 79, 89, 90, 91, 95, 99, 109, 112, 114, 123, 130, 133, 147, 149, 163, 164, 202, 206, 218, 225, 264, 286, 292, 303, 317, 333, 334, 335, 345, 371, (383).)

$\mathcal{M}28$: Morris'\index{Author}{Morris} Model. Morris constructs a generic extension of a countable standard model of ZFC in which there is a proper class of generic sets. (The generic sets are added at each level $R_{\alpha}=\{x$: rank of $x$ is $<\alpha\}$.) In this model, for each ordinal $\alpha$, there is a set $X$ such that $X$ is a denumerable union of denumerable sets and $\mathcal{P}(X)$ can be partioned into $\aleph_{\alpha}$ nonempty sets.

In $\mathcal{M}28$, forms 91 and 209 are false. References Morris \cite{bib:1970}, and note 18.

(A minimal list of forms implying each form known to be true in the model: 0.

A minimal list of forms at least one of which is implied by each form known to be false in the model: 91, 209.

A list of all forms known to be true in the model: 0, (383).

A list of all forms known to be false in the model: 1, 8, 16, 20, 21, 22, 23, 24, 26, 27, 28, 29, 31, 39, 40, 43, 44, 66, 67, 86, 87, 89, 90, 91, 95, 101, 109, 112, 113, 114, 133, 147, 149, 164, 174, 181, 192, 202, 207, 209, 218, 255, 256, 257, 258, 259, 260, 261, 262, 264, 286, 292, 333, 334, 335, 345, 347, 359, (383).)

$\mathcal{M}29$: Pincus' Model II. Pincus\index{Author}{Pincus} constructs a generic extension $M[I]$ of a model $M$ of ZF + class choice + GCH in which $I=\bigcup_{n\in\omega}I_{n},\,I_{-1}=2$ and $I_{n+1}$ is a denumerable set of independent functions from $\omega$ onto $I_{n}$. In this model $\{I_{n}:n\in\omega\}$ is a denumerable set of denumerable sets which has neither a choice function (32 is false), nor a Kinna-Wagner selection function $(KW(\aleph_{0},\,\aleph_{0})$, 357 is false), nor a multiple choice function $(MC(\aleph_{0},\, \aleph_{0})$, 350 is false); $|I|<|2\times I|$ (3 is false); every set can be linearly ordered (30 is true); every infinite set has a denumerable subset (9 is true); and for all infinite $x,\, |x!|=|2^{x}|=|x^{x}|$ (200, 290, and 291 are true). Pincus also proves that if $x$ is a set of infinite sets, then there is a function $f$ with domain $x$ such that for each $y\in x,\,f(y)\subseteq y$, and $|f(y)|=\aleph_{0}$ (214 is true). In this model form 10 $(C(\aleph_{0},\,<\aleph_{0}))$ is true (30 implies 10) and form 165 $(C(WO,\,WO))$ is false (165 implies 32). Thus, form 163 (Every non well orderable set has an infinite Dedekind finite subset.) is false because Brunner\index{Author}{Brunner} \cite{bib:1982a} shows that 10 + 163 implies 165. Since form 30 is true and form 214 implies form 296 (Every infinite set is the disjoint union of infinitely many infinite sets.), it follows from Pincus \cite{bib:1997} that form 295 (Every infinite set has a dense linear ordering.) is true.

$\mathcal{M}29 \models$ 9, 30, 200, 214, 290, 291, and 295, but 3, 32, 91, 118, 151, 163, 350, and 357 are false. References Pincus \cite{bib:1974a}, \cite{bib:1978}, \cite{bib:1997}, Brunner \cite{bib:1982a}, notes 18, 64, and 120(16 and 43).

(A minimal list of forms implying each form known to be true in the model: 200, 214, 290, 291, 295.

A minimal list of forms at least one of which is implied by each form known to be false in the model: 3, 91, 118, 151, 163, 350, 357.

A list of all forms known to be true in the model: 0, 4, 9, 10, 13, 17, 18, 30, 33, 45, 46, 47, 48, 57, 61, 62, 64, 76, 80, 82, 83, 84, 88, 93, 98, 102, 111, 120, 121, 122, 124, 127, 128, 131, 132, 140, 142, 146, 152, 173, 178, 185, 198, 199, 200, 214, 216, 217, 249, 250, 268, 269, 276, 280, 283, 285, 288, 290, 291, 293, 295, 296, 304, 308, 325, 327, 336, 342, 358, 373, 374, 376, 377, 378, (383).

A list of all forms known to be false in the model: 1, 2, 3, 8, 15, 16, 20, 21, 23, 27, 28, 29, 31, 32, 39, 40, 43, 44, 60, 66, 67, 85, 86, 87, 89, 90, 91, 95, 101, 106, 109, 112, 113, 114, 118, 126, 133, 147, 149, 150, 151, 163, 164, 165, 174, 181, 188, 192, 193, 202, 213, 218, 231, 255, 256, 257, 258, 259, 260, 261, 262, 264, 286, 292, 322, 323, 324, 328, 329, 330, 333, 334, 335, 345, 347, 349, 350, 351, 352, 353, 354, 355, 356, 357, 359, (383).)

$\mathcal{M}30$: Pincus\index{Author}{Pincus}/Solovay Model II. In this construction, an $\omega_{1}$ sequence of generic reals is added to a model of ZFC in such a way that the Principle of Dependent Choices (43) is true, but no nonprincipal measure exists (223 is false). (The construction is similar to that given for $\mathcal{M}27$.) Since 43 implies 8 $(C(\aleph_{0},\,\infty))$, it follows from Brunner\index{Author}{Brunner} \cite{bib:1982a} that in this model there is a set that cannot be well ordered and does not have an infinite Dedekind finite subset (163 is false). (Form 8 plus form 163 iff AC.) Howard\index{Author}{Howard} \cite{bib:1975} has shown that if form 206 (There exists a nonprincipal ultrafilter.) is false (206 implies 223), then {\L}os' Theorem (253) is true. By results of Wagon\index{Author}{Wagon} \cite{bib:1985}, the amenable groups in $\mathcal{M}30$ are the finite groups and therefore 312 (A subgroup of an amenable group is amenable.) is true and 313 ($\mathbb{Z}$ (the set of integers under addition) is amenable.) is false. $(G$ is amenable if there is a finitely additive measure $\mu$ on $\mathcal{P}(G)$ such that $\mu(G)=1$ and $\forall A\subseteq G,\, \forall g\in G,\,\mu(gA)=\mu(A).)$

$\mathcal{M}30 \models$ 43, 253, and 312, but 163, 223, and 313 are false. References Pincus/Solovay\index{Author}{Pincus/Solovay} \cite{bib:1977}, Brunner \cite{bib:1982a}, Howard \cite{bib:1975}, and Wagon \cite{bib:1985}.

(A minimal list of forms implying each form known to be true in the model: 43, 253, 312.

A minimal list of forms at least one of which is implied by each form known to be false in the model: 163, 223.

A list of all forms known to be true in the model: 0, 5, 6, 8, 9, 10, 13, 16, 17, 18, 19, 24, 26, 27, 28, 29, 31, 32, 34, 35, 36, 37, 38, 42, 43, 57, 64, 73, 74, 75, 77, 78, 80, 82, 83, 84, 92, 94, 98, 104, 106, 108, 113, 119, 124, 126, 127, 128, 131, 132, 150, 154, 155, 182, 185, 194, 198, 199, 209, 211, 216, 217, 243, 249, 253, 273, 279, 282, 288, 296, 304, 312, 325, 336, 337, 338, 339, 340, 341, 342, 350, 351, 352, 353, 354, 355, 357, 358, 360, 361, 362, 363, 364, 373, 374, 375, 376, 377, 378, 380, (383).

A list of all forms known to be false in the model: 1, 14, 50, 52, 63, 66, 67, 70, 79, 89, 90, 91, 95, 99, 109, 112, 114, 123, 130, 133, 147, 149, 163, 164, 202, 206, 218, 221, 222, 223, 225, 264, 287, 292, 303, 311, 313, 317, 333, 334, 335, 345, 372, (383).)

$\mathcal{M}31$: Szczepaniak's\index{Author}{Szczepaniak} Model. Two models of ZF, $M_{1}$ and $M_{2},\,M_{1}\subseteq M_{2}$, are constructed with the same ordinals (so $L^{1}$ in $M_{1}$ is the same as $L^{1}$ in $M_{2}$), and a generic real $a\in M_{1}$ so that $a\not\in HOD$ in $M_{1}$, but $a\in HOD$ in $M_{2}$. Thus, in at least one of these models $L^{1}\neq HOD$ and that is the model we call $\mathcal{M}31$. (See note 82 for definitions.)

In $\mathcal{M}31$, forms 91 and 230 are false. References Szczepaniak \cite{bib:1977}, notes 18 and 82.

(A minimal list of forms implying each form known to be true in the model: 0.

A minimal list of forms at least one of which is implied by each form known to be false in the model: 91, 230.

A list of all forms known to be true in the model: 0, (383).

A list of all forms known to be false in the model: 1, 66, 67, 89, 90, 91, 95, 109, 112, 114, 133, 147, 149, 164, 202, 218, 230, 264, 292, 333, 334, 335, (383).)

$\mathcal{M}32$: Sageev's Model II. Starting with a model $\mathcal{M}$ of ZF + V = L, Sageev constructs a sequence of models $\mathcal{M}\ \subseteq N_{0}\subseteq N_{1}\subseteq\cdots\subseteq N_{\kappa}$ where $\kappa$ is an inaccessible cardinal, $N_{0}$ is Cohen's model $\mathcal{M}1$, and $N_{\kappa}$ is $\mathcal{M}32$. In this model, every set can be linearly ordered (30 is true); if $\Delta$ is the set of Dedekind finite cardinals, then $2^{\aleph_{0}}<|\Delta|\leq 2^{2^{\aleph_{0}}}$ (9 is false) and each pair of cardinals in $\Delta$ is comparable (57 is true). Sageev\index{Author}{Sageev} also shows that $C(\infty,\,<\aleph_{0})$ (62) is true. Since form 64 (There is no amorphous set.) is true (30 implies 64), form 133 (Every set is either well orderable or has an infinite amorphous subset.) is false. Form 40 $(C(WO,\,\infty))$ is false because 40 implies 9 and form 122 $(C(WO,\,<\aleph_{0}))$ is true because 30 implies 122. Therefore, form 328 $(MC(WO,\,\infty))$ is false because 122 + 328 $\rightarrow$ 40. In addition, 126 $(MC(\aleph_{0},\,\infty))$ is false because 30 is true and 8 is false. $(8(C(\aleph_{0},\,\infty))$ implies 9 and 30 + 126 implies 8).

$\mathcal{M}32 \models$ 30, 57, and 62, but 9, 91, 126, 133, and 328 are false. References Sageev \cite{bib:1981}, notes 18 and 120(9).

(A minimal list of forms implying each form known to be true in the model: 30, 57.

A minimal list of forms at least one of which is implied by each form known to be false in the model: 9, 91, 126.

A list of all forms known to be true in the model: 0, 10, 18, 30, 33, 45, 46, 47, 48, 57, 61, 62, 64, 80, 83, 88, 93, 102, 111, 120, 121, 122, 127, 132, 140, 142, 146, 178, 198, 216, 249, 250, 268, 269, 276, 280, 283, 285, 288, 293, 308, 327, 336, 342, 358, 373, 374, 377, 378, (383).

A list of all forms known to be false in the model: 1, 2, 3, 4, 7, 8, 9, 39, 40, 41, 43, 44, 66, 67, 71, 86, 87, 89, 90, 91, 95, 100, 101, 106, 109, 112, 113, 114, 126, 129, 133, 147, 149, 152, 161, 164, 168, 174, 181, 188, 192, 193, 202, 214, 218, 255, 256, 257, 258, 259, 260, 261, 262, 264, 286, 292, 328, 333, 334, 335, 345, 347, 359, (383).)

$\mathcal{M}33$: Plotkin's Model II. The construction is similar to the construction of $\mathcal{M}22.$ An atomless Boolean algebra in a countable standard modell $\mathcal{M}$ of ZF + V = L is embedded in a generic extension, $\mathcal{M}33$, of $\mathcal{M}$ in such a way that the isomorphic copy, $B$, of the Boolean algebra has the property that it is infinite, atomless, and Dedekind finite (9 and 247 are false) and $|B|^{\aleph_{0}}>|B|$ (248 is false). Plotkin also proves that $C(\infty,\,2)$ (88) is false in this model. Additional properties of $B$ include the following: Every filter in $B$ is principal; $B$ has no infinite free subalgebra; $B$ is not the direct union of complete weakly homogeneous Boolean algebras; $B$ is not inversely reducible. See Plotkin \cite{bib:1976} for definitions.

In $\mathcal{M}33$, forms 9, 88, 91, 247 and 248 are false. References Plotkin\index{Author}{Plotkin} \cite{bib:1969}, \cite{bib:1976}, and note 18.

(A minimal list of forms implying each form known to be true in the model: 0.

A minimal list of forms at least one of which is implied by each form known to be false in the model: 9, 88, 91, 247, 248.

A list of all forms known to be true in the model: 0, (383).

A list of all forms known to be false in the model: 1, 2, 3, 4, 7, 8, 9, 14, 15, 30, 39, 40, 41, 43, 44, 49, 50, 60, 61, 62, 66, 67, 68, 71, 85, 86, 87, 88, 89, 90, 91, 95, 100, 101, 107, 109, 112, 113, 114, 123, 129, 133, 147, 149, 152, 161, 164, 168, 174, 181, 192, 201, 202, 213, 214, 218, 247, 248, 255, 256, 257, 258, 259, 260, 261, 262, 264, 270, 284, 286, 292, 295, 303, 317, 323, 326, 331, 332, 333, 334, 335, 343, 345, 346, 347, 359, (383).)

$\mathcal{M}34(\aleph_{1})$: Pincus' Model III. Pincus\index{Author}{Pincus} proves that Cohen's model $\mathcal{M}1$ can be extended by adding $\aleph_{1}$ generic sets along with the set $b$ containing them and well orderings of all countable subsets of $b$. The model is constructed in such a way that every collection of denumerable sets has a choice function (85 is true), but there is a denumerable collection of sets of cardinality $\aleph_{1}$ which does not have a choice function. Therefore, both $C(WO,\,WO)$ (165) and $C(\infty,\, \aleph_{1})$ (213) are false. Pincus also shows that the Kinna-Wagner Principle (15) and the Ordering Principle (30) are true. Since form 64 (There is no amorphous set.) is true (30 implies 64), form 133 (Every set is either well orderable or has an infinite amorphous subset.) is false. Form 165 is false and form 122 $(C(WO,\,<\aleph_{0}))$ is true because 15 implies 122. Therefore, form 330 $(MC(WO,\,WO))$ is false because 122 + 330 $\rightarrow$ 165.

$\mathcal{M}34 \models$ 15, 30 and 85, but 91, 133, 163, 165, 213, and 330 are false. References Pincus \cite{bib:1977a}, notes 74 and 120(6 and 30).

(A minimal list of forms implying each form known to be true in the model: 15, 85.

A minimal list of forms at least one of which is implied by each form known to be false in the model: 91, 163, 213, 330.

A list of all forms known to be true in the model: 0, 5, 10, 15, 18, 30, 32, 33, 45, 46, 47, 48, 61, 62, 64, 80, 82, 83, 84, 85, 88, 93, 102, 111, 119, 120, 121, 122, 127, 132, 140, 142, 146, 150, 178, 198, 216, 249, 250, 268, 269, 270, 271, 276, 280, 283, 285, 288, 293, 295, 296, 308, 314, 322, 323, 324, 327, 336, 342, 349, 350, 355, 356, 357, 358, 373, 374, 376, 377, 378, 379, (383).

A list of all forms known to be false in the model: 1, 40, 60, 66, 67, 89, 90, 91, 95, 101, 109, 112, 114, 133, 147, 149, 163, 164, 165, 202, 213, 218, 255, 256, 257, 258, 259, 260, 261, 262, 264, 286, 292, 328, 329, 330, 333, 334, 335, 347, 359, (383).)

$\mathcal{M}35(\epsilon)$: David's Model. In Cohen's model $\mathcal{M}1$, define sets $B_{n}=\{x\subset\omega:|x\,\Delta\,a_{n}|<\omega \vee|x\,\Delta\, \omega -a_{n}|<\omega\}$ (where $\Delta$ is the symmetric difference). Let $C$ be the function on $\omega$ such that for each $n\in\omega$, $C(n)=B_{n}$. The symmetric submodel, $N$, of $\mathcal{M}1$ which is the set of all sets that are hereditarily ordinal definable from $\{C\}\cup\{a_{n}:n\in\omega\}$ has the property that $(\forall x\in N)(x\in W(1))$, but it is not the case that all sets are in $W(0)$, the set of all well orderable sets. (See note 25.) David proves that this procedure can be generalized so that for any ordinal $\epsilon$ there is a model, $\mathcal{M}35(\epsilon)$, such that for each set $x$ in the model, $x\in W(\epsilon)$, and $\epsilon$ is the smallest ordinal with this property. Thus, form $179(\epsilon)$ is true, so every set is almost well orderable (144) is true, but form $179(\beta)$, $\beta<\epsilon$, is false.

$\mathcal{M}35(\epsilon)\models$ 144 and $179(\epsilon)$, but 91 and $179(\beta)$, $\beta<\epsilon$, are false. References David\index{Author}{David} \cite{bib:1980}, notes 18 and 25.

(A minimal list of forms implying each form known to be true in the model: 179.

A minimal list of forms at least one of which is implied by each form known to be false in the model: 91, 179.

A list of all forms known to be true in the model: 0, 125, 144, 179, (383).

A list of all forms known to be false in the model: 1, 66, 67, 89, 90, 91, 95, 109, 112, 114, 133, 147, 149, 164, 179, 202, 218, 264, 292, 333, 334, 335, (383).)

$\mathcal{M}36$: Figura's Model. Starting with a countable, standard model, $\mathcal{M}$, of ZFC + 2$^{\aleph_{0}}=\aleph_{\omega+1}$, Figura\index{Author}{Figura} uses forcing conditions that are functions from a subset of $\omega\times\omega$ to $\omega_{\omega}$ to construct a symmetric extension of $\mathcal{M}$ in which there is an uncountable well ordered subset of the reals (170 is true), but $\aleph_{1}=\aleph_{\omega}$ so $\aleph_{1}$ is singular (34 is false). (In fact, Figura proves this result for any two cardinals $\kappa$ (replacing $\aleph_{0})$ and $\mu$ (replacing $\aleph_{\omega)}$) such that cf$(\kappa)=\kappa \leq \mathrm{cf}(\mu)<\mu$ and $2^{\kappa}=\mu^{+}$ in $\mathcal{M}$. See note 3.)

$\mathcal{M}36 \models$ 170, but 34 and 91 are false. References Figura \cite{bib:1981}, notes 3 and 18.

(A minimal list of forms implying each form known to be true in the model: 170.

A minimal list of forms at least one of which is implied by each form known to be false in the model: 34, 91.

A list of all forms known to be true in the model: 0, 38, 93, 108, 170, (383).

A list of all forms known to be false in the model: 1, 8, 16, 20, 21, 23, 25, 27, 28, 29, 31, 34, 39, 40, 43, 44, 51, 66, 67, 86, 87, 89, 90, 91, 95, 101, 109, 112, 113, 114, 133, 147, 149, 164, 172, 174, 181, 192, 202, 218, 245, 255, 256, 257, 258, 259, 260, 261, 262, 264, 286, 292, 333, 334, 335, 345, 347, 359, (383).)

$\mathcal{M}37$: Monro's Model III. This is a generic extension of $\mathcal{M}1$ in which there is an amorphous set (64 is false) and $C(\infty,\,2)$ (88) is false. The model is constructed as follows: If $b$ is the countable set of generic reals from which $\mathcal{M}1$ is constructed, define a partial ordering $(Q,\,\leq)$ in $\mathcal{M}1$ where $Q$ is the set of ordered pair $(A,\,p)$ such that $A$ is a finite subset of $b$ and $\rho$: $A^{2}\rightarrow\{0, 1\}$ satisfies $\{(a,\,b):\rho(a,\,b)=1\}$ is an equivalence relation on $A$. The ordering $\leq$ is defined by $(B,\,\sigma)\leq(A,\,\rho)$ if and only if $A\subseteq B$ and $\sigma$ extends $\rho$. Let $F$ be $Q$ generic over $\mathcal{M}1$. Then Monro's model, $\mathcal{M}37$, is $\mathcal{M}1[F]$. In Morillon\index{Author}{Morillon} \cite{bib:1993} it is shown that if form 64 is false then there is an infinite, compact, Hausdorff, extremally disconnected topological space and therefore 371 is true.

$\mathcal{M}37 \models$ 371 but 64, 88, and 91 are false. References Monro\index{Author}{Monro} \cite{bib:1972}, Morillon \cite{bib:1993}, and note 18.

(A minimal list of forms implying each form known to be true in the model: 371.

A minimal list of forms at least one of which is implied by each form known to be false in the model: 64, 91.

A list of all forms known to be true in the model: 0, 371, (383).

A list of all forms known to be false in the model: 1, 2, 3, 4, 7, 8, 9, 14, 15, 30, 39, 40, 41, 43, 44, 45, 49, 50, 57, 60, 61, 62, 64, 66, 67, 68, 71, 82, 83, 85, 86, 87, 88, 89, 90, 91, 95, 100, 101, 106, 107, 109, 112, 113, 114, 123, 126, 129, 133, 147, 149, 152, 161, 164, 168, 174, 181, 188, 192, 193, 201, 202, 213, 214, 215, 218, 255, 256, 257, 258, 259, 260, 261, 262, 264, 270, 284, 286, 292, 295, 296, 303, 317, 323, 326, 328, 331, 332, 333, 334, 335, 336, 343, 345, 346, 347, 359, 376, 377, 378, (383).)

$\mathcal{M}38$: Shelah's Model II. In a model of ZFC + $``\kappa$ is a strongly inaccessible cardinal'', Shelah\index{Author}{Shelah} uses Levy's method of collapsing cardinals to collapse $\kappa$ to $\aleph_{1}$ similarly to Solovay\index{Author}{Solovay} \cite{bib:1970}. (See $\mathcal{M}5$ and $\mathcal{M}12(\aleph)$.) He then builds a special forcing notion and constructs a model in which the Principle of Dependent Choices (43) is true and every set of reals is Lebesgue measurable (93 is false) just as in $\mathcal{M}5$. However, unlike $\mathcal{M}5$, in Shelah's model there is a set of reals without the Baire property (142 is true) and there is an uncountable set of reals with no perfect set (169 is true). Since 43 implies 8 $(C(\aleph_{0},\,\infty))$, it follows from Brunner\index{Author}{Brunner} \cite{bib:1982a} that in this model there is a set that cannot be well ordered and does not have an infinite Dedekind finite subset, (163 is false). (Form 8 plus form 163 iff AC.)

$\mathcal{M}38 \models$ 43, 142, and 169, but 93 and 163 are false. References Shelah \cite{bib:1985}, Solovay \cite{bib:1970}, and Brunner \cite{bib:1982a}.

(A minimal list of forms implying each form known to be true in the model: 43, 142, 169.

A minimal list of forms at least one of which is implied by each form known to be false in the model: 93, 163.

A list of all forms known to be true in the model: 0, 5, 6, 8, 9, 10, 13, 16, 17, 18, 19, 24, 26, 27, 28, 29, 31, 32, 34, 35, 36, 37, 38, 42, 43, 57, 64, 73, 74, 75, 77, 78, 80, 82, 83, 84, 92, 94, 98, 104, 106, 108, 113, 119, 124, 126, 127, 128, 131, 132, 142, 150, 154, 155, 169, 182, 185, 194, 198, 199, 209, 211, 216, 217, 243, 249, 273, 279, 280, 282, 288, 296, 304, 325, 336, 337, 338, 339, 340, 341, 342, 350, 351, 352, 353, 354, 355, 357, 358, 360, 361, 362, 363, 364, 373, 374, 375, 376, 377, 378, 380, (383).

A list of all forms known to be false in the model: 1, 14, 15, 30, 49, 50, 52, 60, 61, 62, 63, 66, 67, 68, 70, 79, 85, 88, 89, 90, 91, 93, 95, 99, 100, 101, 107, 109, 112, 114, 123, 130, 133, 147, 149, 163, 164, 168, 170, 184, 195, 201, 202, 203, 213, 218, 225, 251, 252, 264, 270, 284, 292, 295, 303, 305, 306, 309, 317, 323, 326, 331, 332, 333, 334, 335, 343, 345, 346, 359, 366, 367, 368, 372, (383).)

$\mathcal{M}39(\kappa,\,\lambda)$: Kanovei's model II. This model depends on the two cardinals $\kappa <\lambda$ such that both $\kappa$ and $\lambda$ have cofinality $>\omega$ and neither $\kappa$ nor $\lambda$ can be written as $\theta^{+}$ where $\theta$ is a cardinal of countable cofinality and such that $\aleph_{2}\leq \kappa$.

Suppose $\kappa$ and $\lambda$ are cardinals in a model $\mathcal{M}$ of ZFC + V = L. Assume that $\kappa$ and $\lambda$ satisfy the properties of the preceding paragraph in $\mathcal{M}$. Let $P=\{p:p$ is a function such that $\mathrm{dom}(p)\subseteq\lambda\times \omega$ is finite and $\mathrm{ran}(p)\subseteq\{0, 1\}\}$ and let $G$ be $P$ generic over $M$. For each $\zeta <\lambda$, let $a_{\zeta}(k)=i$ if and only if $(\exists p\in G)(p(\zeta,\,k)=i)$. $a_{\zeta}$ is a function from $\omega$ to $\{0, 1\}$ which we identify with a real number between 0 and 1 whose binary expansion has digits $a_{\zeta}(0), a_{\zeta}(1), a_{\zeta}(2),\ldots$ and use $[a_{\zeta}]$ to denote the Vitali equivalence class of this real number. (Vitali equivalence classes are equivalence classes or real numbers under the relation $x\equiv y\leftrightarrow(\exists q\in \mathbb{Q})(x-y=q).)$ Now define the following
\begin{enumerate}
\item[1.] $B=\{b\in\mathcal{M}:b$ is a bijection from $\lambda$ onto $\lambda$\}.
\item[2.] $Z=\{z\subseteq\lambda\times \omega:z$ is finite\}
\item[3.] For $b\in B,\,z\in Z$ and $\zeta <\lambda\ za_{\zeta}(k)=\left\{\begin{array}{ll}
    a_{\zeta}(k) & \mathrm{for}\,(\zeta,\,k)\not\in z\\
    1-a_{\zeta}(k) & \mathrm{for}\,(\zeta,\,k)\in z \end{array}\right.$
\item[4.] $zb\,|\,u=\langle za_{b(\zeta)}:\zeta \in u\rangle$
\item[5.] $|zb\,|\,u|=\langle|za_{b(\zeta)}|:\zeta \in u\rangle$
\item[6.] $C_{\theta}=\{u\subseteq\lambda:u\in \mathcal{M}$ has cardinality $<\theta$ in $\mathcal{M}$\}
\item[7.] $W_{\kappa}=\{zb\,|\,u:z\in Z\wedge b\in B\wedge u\in C_{\kappa}\}$
\item[8.] $[W_{\lambda}]=\{|zb\,|\,u|:z\in Z \wedge b\in B\wedge u\in C_{\lambda}\}$
\end{enumerate}

The model $\mathcal{M}39(\kappa,\,\lambda)$ is HOD$(W_{\kappa}\cup[W_{\lambda}]\cup\{W_{\kappa},\,[W_{\lambda}]\})$. In Kanovei \cite{bib:1991} it is shown that the Principle of Dependent Choices (43) is true and that $H(2^{\aleph_{0}})< H(m)$ where $m$ is the cardinality of the set of Vitali equivalence classes and $H$ is Hartogs' function (307 is false). $(H(x)$ is the smallest ordinal that is not similar to a subset of $x$.) Since 43 implies 8 $(C(\aleph_{0},\,\infty))$, it follows from Brunner\index{Author}{Brunner} \cite{bib:1982a} that in this model there is a set that cannot be well ordered and does not have an infinite Dedekind finite subset, (163 is false). (Form 8 plus form 163 iff AC.)

$\mathcal{M}39(\kappa,\,\lambda)\models$ 43 but 91, 163, and 307 are false. References Kanovei \cite{bib:1991}, Brunner \cite{bib:1982a}, and note 18.

(A minimal list of forms implying each form known to be true in the model: 43.

A minimal list of forms at least one of which is implied by each form known to be false in the model: 91, 163, 307.

A list of all forms known to be true in the model: 0, 5, 6, 8, 9, 10, 13, 16, 17, 18, 19, 24, 26, 27, 28, 29, 31, 32, 34, 35, 36, 37, 38, 42, 43, 57, 64, 73, 74, 75, 77, 78, 80, 82, 83, 84, 92, 94, 98, 104, 106, 108, 113, 119, 124, 126, 127, 128, 131, 132, 150, 154, 155, 182, 185, 194, 198, 199, 209, 211, 216, 217, 243, 249, 273, 279, 282, 288, 296, 304, 325, 336, 337, 338, 339, 340, 341, 342, 350, 351, 352, 353, 354, 355, 357, 358, 360, 361, 362, 363, 364, 373, 374, 375, 376, 377, 378, 380, (383).

A list of all forms known to be false in the model: 1, 66, 67, 89, 90, 91, 95, 109, 112, 114, 133, 147, 149, 163, 164, 202, 218, 264, 292, 305, 307, 333, 334, 335, (383).)

$\mathcal{M}40(\kappa)$: Pincus' Model IV. The ground model $\mathcal{M}$, is a model of ZF + the class form of AC. In $\mathcal{M},\,\kappa$ is a regular cardinal larger than $\aleph_{0}$ such that $2^{\kappa}=\kappa^{+}$. Pincus\index{Author}{Pincus} constructs a generic extension of $\mathcal{M}$ using $\kappa^{+}$ generic iterations. To assure that all the sets added in this manner have canonical supports, he enumerates the sets at each level using hereditarily almost disjoint sets. (See Pincus \cite{bib:1977a} p 125 for definitions.) The construction is made is such a way that at each level DC (43) holds, but AC does not. In Pincus \cite{bib:1977a} it is shown that in this model, every set can be linearly ordered (30 is true), $C(\infty,\,WO)$ (60) is true; and for all $\aleph_{\alpha}<\kappa,\,\mathrm{DC}(\aleph_{\alpha})\ (87(\alpha))$ is true. Pincus \cite{bib:1977b} proves that the Boolean Prime Ideal Theorem (14) is also true. Since 87 implies 8 $(C(\aleph_{0},\,\infty))$, it follows from Brunner \cite{bib:1982a} that in this model there is a set that cannot be well ordered and does not have an infinite Dedekind finite subset, (163 is false). (Form 8 plus form 163 iff AC.) Form 14 implies 49 (Every partial ordering can be extended to a linear ordering.) and 49 + 51 (Every linear ordering has a cofinal sub well ordering.) iff AC. (See Morris\index{Author}{Morris} \cite{bib:1969} and note 121.) Thus, form 51 is false. Since form 30 is true and form $87(\alpha)$ implies form 296 (Every infinite set is the disjoint union of infinitely many infinite sets.), it follows from Pincus \cite{bib:1997} that form 295 (Every infinite set has a dense linear ordering.) is true. Since 43 + 14 is equivalent to 345 ([345 A]), form 345 is true.

$\mathcal{M}40(\kappa)\models$ 14, 30, 43, 60, 87$(\alpha)$ for $\aleph_{\alpha}<\kappa$, 165, 295, and 345, but 51, 91, 144, 152, 163, 253, and 286 are false. References Pincus \cite{bib:1977a}, \cite{bib:1977b}, \cite{bib:1997}, Brunner \cite{bib:1982a}, Morris \cite{bib:1969}, notes 18, 120(26, 28, 34, 43, and 45) and 121.

(A minimal list of forms implying each form known to be true in the model: 60, 87, 165, 295, 345.

A minimal list of forms at least one of which is implied by each form known to be false in the model: 51, 91, 144, 152, 163, 253, 286.

A list of all forms known to be true in the model: 0, 5, 6, 8, 9, 10, 13, 14, 16, 17, 18, 19, 24, 26, 27, 28, 29, 30, 31, 32, 33, 34, 35, 36, 37, 38, 42, 43, 45, 46, 47, 48, 49, 52, 57, 60, 61, 62, 63, 64, 69, 70, 71, 72, 73, 74, 75, 77, 78, 80, 82, 83, 84, 85, 86, 87, 88, 92, 93, 94, 96, 98, 99, 102, 104, 106, 107, 108, 111, 113, 119, 120, 121, 122, 123, 124, 126, 127, 128, 131, 132, 137, 139, 140, 141, 142, 146, 150, 151, 153, 154, 155, 165, 174, 178, 182, 183, 185, 194, 196, 198, 199, 201, 206, 209, 211, 213, 216, 217, 221, 222, 223, 225, 226, 227, 228, 229, 231, 233, 235, 241, 242, 243, 249, 250, 268, 269, 270, 271, 273, 276, 279, 280, 282, 283, 285, 287, 288, 293, 294, 295, 296, 298, 304, 308, 309, 311, 313, 314, 323, 324, 325, 326, 327, 329, 330, 331, 332, 336, 337, 338, 339, 340, 341, 342, 343, 344, 345, 349, 350, 351, 352, 353, 354, 355, 356, 357, 358, 360, 361, 362, 363, 364, 370, 371, 373, 374, 375, 376, 377, 378, 380, (383).

A list of all forms known to be false in the model: 1, 51, 66, 67, 89, 90, 91, 95, 109, 112, 114, 133, 144, 147, 149, 152, 163, 164, 179, 202, 214, 218, 253, 256, 259, 261, 264, 286, 292, 333, 334, 335, (383).)

$\mathcal{M}41$ Kanovei's Model III. Let $\mathbb{P}$ be the set of conditions from the model in Jensen\index{Author}{Jensen} \cite{bib:1968}. $\mathbb{P}$ is a subset of the set of all constructible perfect trees. $\mathcal{M}41$ is the symmetric part of $L[\boldsymbol{\alpha}]$ where $\boldsymbol{\alpha}=\langle\alpha_{n}: n\in\omega\rangle$ is $\mathbb{P}^{w}$ generic. In Kanovei\index{Author}{Kanovei} \cite{bib:1976} it is shown that there is a $\varPi_{2}^{1}$ subset $P$ of $\omega \times{^{\omega}\omega}$ with dom $P=\omega$ and for which no sequence $\langle x_{n}:n\in\omega\rangle$ with $\langle k,\,x_{k}\rangle\in P$ for all $k\in \omega$ exists (form 194 is false).

In $\mathcal{M}41$, form 194 is false. Reference Kanovei \cite{bib:1976}.

(A minimal list of forms implying each form known to be true in the model: 0.

A minimal list of forms at least one of which is implied by each form known to be false in the model: 194.

A list of all forms known to be true in the model: 0, (383).

A list of all forms known to be false in the model: 1, 8, 16, 39, 40, 43, 44, 66, 67, 79, 86, 87, 89, 90, 91, 94, 95, 101, 109, 112, 113, 114, 130, 133, 147, 149, 164, 174, 181, 192, 194, 202, 203, 212, 218, 255, 256, 257, 258, 259, 260, 261, 262, 264, 286, 292, 333, 334, 335, 345, 347, 359, (383).)

$\mathcal{M}42$ Bull's Model. Let $\mathcal{M}$ be a countable transitive model of ZFC + ``There are uncountable regular cardinals $\aleph_{\alpha}<\aleph_{\beta}<\aleph_{\gamma}$ such that $\aleph_{\alpha}$ is $\aleph_{\gamma}$-supercompact; $\aleph_{\beta}$ is the first measurable cardinal greater than $\aleph_{\alpha}$; and $\aleph_{\gamma}=|2^{\aleph_{\beta}}|.$'' Using backward Easton forcing (which is due to Silver), Bull constructs a generic extension of $\mathcal{M}$. Then using collapsing functions, he constructs an inner model, $\mathcal{M}42$, in which $\aleph_{\beta}$ is collapsed to $\aleph_{\alpha+1}$. He proves that $\aleph_{\alpha}$ is an uncountable, weakly compact cardinal and $\aleph_{\alpha+1}$ is measurable (321 is false). Bull\index{Author}{Bull} also proves that $C(\aleph_{\alpha},\,\infty)$ (86) is true in this model. Since $\alpha\geq 1$, form 39 $(C(\aleph_{1},\,\infty))$ is also true. Form 163 (Every set is either well orderable or has an infinite Dedekind finite subset.) is false because 39 implies 8 and 8 + 163 implies AC (Brunner\index{Author}{Brunner} \cite{bib:1982a}).

$\mathcal{M}42 \models$ 39 and 86, but 91, 163, and 321 are false. References Bull \cite{bib:1978}, Brunner \cite{bib:1982a}, notes 18 and 20.

(A minimal list of forms implying each form known to be true in the model: 39, 86.

A minimal list of forms at least one of which is implied by each form known to be false in the model: 91, 163, 321.

A list of all forms known to be true in the model: 0, 5, 6, 8, 9, 10, 13, 16, 17, 18, 19, 24, 26, 27, 28, 29, 31, 32, 34, 35, 36, 37, 38, 39, 57, 64, 73, 74, 75, 80, 82, 83, 84, 86, 92, 94, 98, 104, 108, 119, 124, 126, 127, 128, 131, 132, 150, 182, 185, 194, 196, 198, 199, 209, 216, 217, 249, 282, 288, 296, 304, 325, 336, 337, 338, 339, 340, 341, 342, 350, 351, 352, 353, 354, 355, 357, 358, 360, 361, 362, 373, 374, 376, 377, 378, 380, (383).

A list of all forms known to be false in the model: 1, 66, 67, 89, 90, 91, 95, 109, 112, 114, 117, 133, 147, 149, 163, 164, 202, 218, 264, 292, 319, 320, 321, 333, 334, 335, (383).)

$\mathcal{M}43$: Pincus' Model V. This is the model of Pincus\index{Author}{Pincus} \cite{bib:1977a}, Theorem 2.1. (A). It uses Pincus' general method for adding dependent choice (43). Beginning with a model $\mathcal{M}$ of ZF + ``The class form of choice'', Pincus adds three generic $\omega_{1}$ sequences: $\langle I_{\alpha}:\alpha< \omega_{1}\rangle,\,\langle A_{\alpha}\,:\,\alpha<\omega_{1}\rangle$ and $\langle \mathfrak{A}_{\alpha}:\alpha<\omega_{1}\rangle$. The sets $I_{\alpha},\,\alpha<\omega_{1}$ are obtained as in $\mathcal{M}1(\langle \omega_{1}\rangle)$, $A_{\alpha}$ is a generic set of disjoint subsets of $I_{\alpha}$ and $\mathfrak{A}_{\alpha}$ is a generically added universal homogeneous structure for the theory of finite choice operators with domain $A_{\alpha}$. (See Pincus \cite{bib:1977a} for definitions.) In Pincus \cite{bib:1977a}, it is shown that in $\mathcal{M}43$ Dependent Choice (43) is true, $C(\infty,\, <\aleph_{0})$ (62) is true, ``For every collection $A$ of sets each of which has cardinality at least $\aleph_{0}$, there is a function $f$ with domain $A$ such that $\forall x\in A,\,f(x)\subseteq x$ and $|f(x)|=\aleph_{0}$.'' (214) is true but the Ordering Principle (30) is false.

$\mathcal{M}43 \models$ 43, 62 and 214, but 30, 67, 91, 99, 133, 151, 163, 328, and 329 are false. References Pincus \cite{bib:1977a}, notes 18 and 120(1, 4, 29, 31, 32, 43, and 48).

(A minimal list of forms implying each form known to be true in the model: 43, 62, 214.

A minimal list of forms at least one of which is implied by each form known to be false in the model: 30, 91, 99, 151, 163, 328, 329.

A list of all forms known to be true in the model: 0, 4, 5, 6, 8, 9, 10, 13, 16, 17, 18, 19, 24, 26, 27, 28, 29, 31, 32, 33, 34, 35, 36, 37, 38, 42, 43, 45, 46, 47, 48, 57, 61, 62, 64, 73, 74, 75, 76, 77, 78, 80, 82, 83, 84, 88, 92, 93, 94, 98, 102, 104, 106, 108, 111, 113, 119, 120, 121, 122, 124, 126, 127, 128, 131, 132, 140, 142, 146, 150, 152, 154, 155, 173, 178, 182, 185, 194, 198, 199, 209, 211, 214, 216, 217, 243, 249, 250, 268, 269, 273, 276, 279, 280, 282, 283, 285, 288, 296, 304, 308, 325, 327, 336, 337, 338, 339, 340, 341, 342, 350, 351, 352, 353, 354, 355, 357, 358, 360, 361, 362, 363, 364, 373, 374, 375, 376, 377, 378, 380, (383).

A list of all forms known to be false in the model: 1, 14, 15, 20, 21, 23, 30, 40, 49, 50, 60, 66, 67, 89, 90, 91, 95, 99, 101, 109, 112, 114, 133, 147, 149, 151, 163, 164, 202, 218, 231, 255, 256, 257, 258, 259, 260, 261, 262, 264, 286, 292, 295, 303, 317, 328, 329, 333, 334, 335, 345, 347, 359, (383).)

$\mathcal{M}44$: Pincus' Model VI. This is the model of Pincus \cite{bib:1977a}, Theorem 2.1. (B). It uses Pincus' general method for adding dependent choice (43). Beginning with a model $\mathcal{M}$ of ZF + ``The class form of choice'', Pincus adds three generic $\omega_{1}$ sequences: $\langle I_{\alpha}:\alpha<\omega_{1}\rangle,\, \langle A_{\alpha}:\alpha<\omega_{1}\rangle$ and $\langle \mathfrak{A}_{\alpha}:\alpha<\omega_{1}\rangle$. The sets $I_{\alpha},\,\alpha<\omega_{1}$ are obtained as in $\mathcal{M}1(\langle \omega_{1}\rangle)$, $A_{\alpha}$ is a generic set of disjoint subsets of $I_{\alpha}$ and $\mathfrak{A}_{\alpha}$ is a generically added universal homogeneous structure for the theory dense linear orderings without first or last element (that is, $\mathfrak{A}_{\alpha}$ is a copy of the rationals with the usual ordering) with domain $A_{\alpha}$. (See Pincus \cite{bib:1977a} for definitions.)

In Pincus \cite{bib:1977a}, it is shown that in $\mathcal{M}44$ the Ordering Principle (30) is true, $f$ with domain $A$ such that $\forall x\in A,\,f(x)\subseteq x$ and $|f(x)|=\aleph_{0}$.'' (214) is true but the Kinna-Wagner Principle (15) is false. Form 43 implies form 296 (Every infinite set is the disjoint union of infinitely many infinite sets.) and it is shown in Pincus \cite{bib:1997} that 30 + 296 implies 295 (Every infinite set has a dense linear ordering.) Thus, form 295 is true.

$\mathcal{M}44 \models$ 30, 43, 214, and 295, but 15, 67, 91, 118, 133, 151, 163, 328, and 329 are false. References Pincus\index{Author}{Pincus} \cite{bib:1977a}, \cite{bib:1997}, notes 18 and 120(1, 4, 16, 31, 32, 43, and 48).

(A minimal list of forms implying each form known to be true in the model: 43, 214, 295.

A minimal list of forms at least one of which is implied by each form known to be false in the model: 15, 91, 118, 151, 163, 328, 329.

A list of all forms known to be true in the model: 0, 4, 5, 6, 8, 9, 10, 13, 16, 17, 18, 19, 24, 26, 27, 28, 29, 30, 31, 32, 33, 34, 35, 36, 37, 38, 42, 43, 45, 46, 47, 48, 57, 61, 62, 64, 73, 74, 75, 76, 77, 78, 80, 82, 83, 84, 88, 92, 93, 94, 98, 102, 104, 106, 108, 111, 113, 119, 120, 121, 122, 124, 126, 127, 128, 131, 132, 140, 142, 146, 150, 152, 154, 155, 173, 178, 182, 185, 194, 198, 199, 209, 211, 214, 216, 217, 243, 249, 250, 268, 269, 273, 276, 279, 280, 282, 283, 285, 288, 293, 295, 296, 304, 308, 325, 327, 336, 337, 338, 339, 340, 341, 342, 350, 351, 352, 353, 354, 355, 357, 358, 360, 361, 362, 363, 364, 373, 374, 375, 376, 377, 378, 380, (383).

A list of all forms known to be false in the model: 1, 15, 20, 21, 23, 40, 60, 66, 67, 89, 90, 91, 95, 101, 109, 112, 114, 118, 133, 147, 149, 151, 163, 164, 202, 218, 231, 255, 256, 257, 258, 259, 260, 261, 262, 264, 286, 292, 328, 329, 333, 334, 335, 347, 359, (383).)

$\mathcal{M}45$: Pincus' Model VII. This is the model of Pincus \cite{bib:1977a}, Theorem 2.1. (C). It uses Pincus' general method for adding dependent choice (43). Beginning with a model $\mathcal{M}$ of ZF + ``The class form of choice'', Pincus adds three generic $\omega_{1}$ sequences: $\langle I_{\alpha}:\alpha<\omega_{1}\rangle,\, \langle A_{\alpha}:\alpha <\omega_{1}\rangle$ and $\langle \mathfrak{A}_{\alpha}:\alpha<\omega_{1}\rangle$. The sets $I_{\alpha},\,\alpha<\omega_{1}$ are obtained as in $\mathcal{M}1(\langle\omega_{1}\rangle), A_{\alpha}$ is a generic set of disjoint subsets of $I_{\alpha}$ and $\mathfrak{A}_{\alpha}$ is a generically added universal homogeneous structure for the theory of partial orders together with a dense linear order both with domain $A_{\alpha}$. (See Pincus \cite{bib:1977a} for definitions.) In Pincus \cite{bib:1977a}, it is shown that in $\mathcal{M}45$ the Ordering Principle (30) is true, Dependent Choice (43) is true, ``For every collection $A$ of sets each of which has cardinality at least $\aleph_{0}$, there is a function $f$ with domain $A$ such that $\forall x\in A,\,f(x)\subseteq x$ and $|f(x)|=\aleph_{0}$.'' (214) is true but the Order Extension Principle (49) is false. Form 43 implies form 296 (Every infinite set is the disjoint union of infinitely many infinite sets.) and it is shown in Pincus \cite{bib:1997} that 30 + 296 implies 295 (Every infinite set has a dense linear ordering.) Thus, form 295 is true.

$\mathcal{M}45 \models$ 30, 43, 214, and 295, but 49, 67, 91, 99, 133, 151, 163, 328, and 329 are false. References Pincus \cite{bib:1977a}, \cite{bib:1997}, notes 18 and 120(1, 4, 29, 31, 32, 43, and 48).

(A minimal list of forms implying each form known to be true in the model: 43, 214, 295.

A minimal list of forms at least one of which is implied by each form known to be false in the model: 49, 91, 99, 151, 163, 328, 329.

A list of all forms known to be true in the model: 0, 4, 5, 6, 8, 9, 10, 13, 16, 17, 18, 19, 24, 26, 27, 28, 29, 30, 31, 32, 33, 34, 35, 36, 37, 38, 42, 43, 45, 46, 47, 48, 57, 61, 62, 64, 73, 74, 75, 76, 77, 78, 80, 82, 83, 84, 88, 92, 93, 94, 98, 102, 104, 106, 108, 111, 113, 119, 120, 121, 122, 124, 126, 127, 128, 131, 132, 140, 142, 146, 150, 152, 154, 155, 173, 178, 182, 185, 194, 198, 199, 209, 211, 214, 216, 217, 243, 249, 250, 268, 269, 273, 276, 279, 280, 282, 283, 285, 288, 293, 295, 296, 304, 308, 325, 327, 336, 337, 338, 339, 340, 341, 342, 350, 351, 352, 353, 354, 355, 357, 358, 360, 361, 362, 363, 364, 373, 374, 375, 376, 377, 378, 380, (383).

A list of all forms known to be false in the model: 1, 14, 20, 21, 23, 40, 49, 50, 60, 66, 67, 89, 90, 91, 95, 99, 101, 109, 112, 114, 133, 147, 149, 151, 163, 164, 202, 218, 231, 255, 256, 257, 258, 259, 260, 261, 262, 264, 286, 292, 303, 317, 328, 329, 333, 334, 335, 345, 347, 359, (383).)

$\mathcal{M}46(m,\,M)$: Pincus'\index{Author}{Pincus} Model VIII. This model depends on the natural number $m$ and the set of natural numbers $M$ which must satisfy Mostowski's condition:
\begin{enumerate}
\item[$S(M,\,m)$:] For every decomposition $m=p_{1}+\ldots+p_{s}$ of $m$ into a sum of primes at least one $p_{i}$ divides an element of $M$. (See Mostowski\index{Author}{Mostowski} \cite{bib:1945} and note 15.)
\end{enumerate}

It is the model of Pincus \cite{bib:1977a}, Theorem 2.1. (D) and uses Pincus' general method for adding dependent choice (43). Beginning with a model $\mathcal{M}$ of ZF + ``The class form of choice''. Pincus adds three generic $\omega_{1}$ sequences: $\langle I_{\alpha}:\alpha<\omega_{1}\rangle,\,\langle A_{\alpha}:\alpha<\omega_{1}\rangle$ and $\langle \mathfrak{A}_{\alpha}:\alpha<\omega_{1}\rangle$. The sets $I_{\alpha},\,\alpha<\omega_{1}$ are obtained as in $\mathcal{M}1(\langle \omega_{1}\rangle)$, $A_{\alpha}$ is a generic set of disjoint subsets of $I_{\alpha}$ and $\mathfrak{A}_{\alpha}$ is a generically added copy of the structure introduced in Pincus \cite{bib:1974b} with domain $A_{\alpha}$. (See Pincus \cite{bib:1977a} for definitions.) In Pincus \cite{bib:1977a}, it is shown that in $\mathcal{M}46(m,\,M)$ Dependent Choice (43) is true, ``For every collection $A$ of sets each of which has cardinality at least $\aleph_{0}$, there is a function $f$ with domain $A$ such that $\forall x\in A,\,f(x)\subseteq x$ and $|f(x)|=\aleph_{0}.$'' (214) is true, ``Every set of sets, all of whose cardinalities are in $M$, has a choice function.'' $(46(M))$ is true but $C(\infty,\,n)\ (45(n))$ is false.

$\mathcal{M}46(n,\,M)\models$ 43, 46$(M)$, and 214, but $45(n)$, 91, 133, 151, 163, and 329 are false. References Pincus \cite{bib:1974b}, \cite{bib:1977a}, notes 18 and 120(1, 4, 43, and 48).

(A minimal list of forms implying each form known to be true in the model: 43, 46, 214.

A minimal list of forms at least one of which is implied by each form known to be false in the model: 45, 91, 151, 163, 329.

A list of all forms known to be true in the model: 0, 4, 5, 6, 8, 9, 10, 13, 16, 17, 18, 19, 24, 26, 27, 28, 29, 31, 32, 34, 35, 36, 37, 38, 42, 43, 46, 48, 57, 64, 73, 74, 75, 76, 77, 78, 80, 82, 83, 84, 92, 94, 98, 104, 106, 108, 113, 119, 120, 124, 126, 127, 128, 131, 132, 150, 152, 154, 155, 173, 182, 185, 194, 198, 199, 209, 211, 214, 216, 217, 243, 249, 273, 279, 282, 288, 296, 304, 325, 336, 337, 338, 339, 340, 341, 342, 350, 351, 352, 353, 354, 355, 357, 358, 360, 361, 362, 363, 364, 373, 374, 375, 376, 377, 378, 380, (383).

A list of all forms known to be false in the model: 1, 14, 15, 20, 21, 23, 30, 40, 45, 49, 50, 60, 61, 62, 66, 67, 68, 85, 89, 90, 91, 95, 101, 107, 109, 112, 114, 123, 133, 147, 149, 151, 163, 164, 202, 213, 218, 231, 255, 256, 257, 258, 259, 260, 261, 262, 264, 270, 284, 286, 292, 295, 303, 317, 323, 329, 331, 332, 333, 334, 335, 343, 345, 346, 347, 359, (383).)

$\mathcal{M}47(n,\,M)$: Pincus' Model IX. This is the model of Pincus \cite{bib:1977a}, Theorem 2.1. (E). In Pincus \cite{bib:1977a} it is shown that $\mathcal{N}2(n,\,M)$ may be transferred to a model of ZF, $\mathcal{N}2(n,\,M)T$, in which $46(\neg M)$ is true, but $46(n\in K)$ and $47(n)$ are false. Then $\mathcal{N}2(n,\,M)T$ can be modified so that, in addition, the Principle of Dependent Choices (43) is true and for each set $x$ of infinite sets, there is a function $f$ such that for each $y\in x,\,f(y)\subseteq y$ and $|f(y)|=\aleph_{0}$ (214) is true. We shall call the modified model $\mathcal{M}47(n,\,M)$. Since 43 implies 9 (Every Dedekind finite set is finite.) and 70 (There is a non-trivial ultra filter on $\omega$.) is true in every FM model, it follows that 63 (Every infinite set has a non-trivial ultrafilter.) is also true. Since form 64 (There is no amorphous set.) is true (30 implies 64), form 133 (Every set is either well orderable or has an infinite amorphous subset.) is false. In this model form 10 $(C(\aleph_{0},\,<\aleph_{0}))$ is true (30 implies 10) and form 165 $(C(WO,\,WO))$ is false (165 implies 32). Thus, form 163 (Every non well orderable set has an infinite Dedekind finite subset.) is false because Brunner\index{Author}{Brunner} \cite{bib:1982a} shows that 10 + 163 implies 165.

$\mathcal{M}47(n,\,M)\models$ 43, 46$(\neg M)$, 63, and 214, but $46(n\in K)$, $47(n)$, 91, 163, 329, and 344 are false. References Pincus\index{Author}{Pincus} \cite{bib:1977a}, Brunner \cite{bib:1982a}, and notes 18 and 120(48 and 56).

(A minimal list of forms implying each form known to be true in the model: 43, 46, 63, 214.

A minimal list of forms at least one of which is implied by each form known to be false in the model: 46, 47, 91, 163, 329, 344.

A list of all forms known to be true in the model: 0, 4, 5, 6, 8, 9, 10, 13, 16, 17, 18, 19, 24, 26, 27, 28, 29, 31, 32, 34, 35, 36, 37, 38, 42, 43, 46, 48, 57, 63, 64, 70, 73, 74, 75, 76, 77, 78, 80, 82, 83, 84, 92, 93, 94, 98, 104, 106, 108, 113, 119, 120, 124, 126, 127, 128, 131, 132, 142, 150, 152, 154, 155, 173, 182, 185, 194, 198, 199, 206, 209, 211, 214, 216, 217, 221, 222, 223, 243, 249, 273, 279, 280, 282, 288, 296, 304, 325, 336, 337, 338, 339, 340, 341, 342, 350, 351, 352, 353, 354, 355, 357, 358, 360, 361, 362, 363, 364, 373, 374, 375, 376, 377, 378, 380, (383).

A list of all forms known to be false in the model: 1, 14, 15, 20, 21, 23, 30, 33, 40, 45, 46, 47, 49, 50, 60, 61, 62, 66, 67, 68, 85, 89, 90, 91, 95, 101, 107, 109, 112, 114, 121, 122, 123, 133, 147, 149, 151, 163, 164, 165, 202, 213, 218, 231, 250, 255, 256, 257, 258, 259, 260, 261, 262, 264, 270, 284, 286, 292, 295, 303, 317, 322, 323, 324, 327, 329, 331, 332, 333, 334, 335, 343, 344, 345, 346, 347, 359, (383).)

\section{Fraenkel-Mostowski Models, Models of $\mathrm{ZF}^{0}\ (\mathcal{N})$}\label{pt03:sec02}\index{Subject}{Fraenkel-Mostowski (FM) models}

The models we give here are all permutation models (FM models) of $\mathrm{ZF}^{0}.$ We shall describe the models by specifying the set $A$ of urelements\index{Subject}{urelements} (or atoms)\index{Subject}{atoms} in the ground model; the group of permutations, $\mathcal{G}$ of $A$; and the set of supports\index{Subject}{supports}, $S$, an ideal of subsets of $A$. (On occasion, we shall give a filter of subgroups of $\mathcal{G}$ rather than $S$. See, for example, Jech\index{Author}{Jech} \cite{bib:1973b}, Chapter 4, for details.)

In all the models of $\mathrm{ZF}^{0}$ given below all the standard properties of the reals are true, PW (91) (The power set of a well ordered set can be well ordered.) is true, but the Axiom of Choice is false and the Kinna-Wagner Principle (15) is false. Blass\index{Author}{Blass} \cite{bib:1979} has shown that in every permutation model with a set of atoms there is a set $X$ such that for each set $a$ there is an ordinal $\alpha$ and a function $f$ mapping $X\times\alpha$ onto $a$ (191 is true). Blass also shows that form 192 (For every set $A$ there is a projective set $X$ and a function $f$ mapping $A$ onto $X$.) + 191 implies AC, so 192 is false in every permutation model with a set of atoms. (See notes 59 and 60.) Thus, each of the following is true in each model: 5, 6, 11, 12, 13, 35, 36, 37, 38, 70, 73, 74, 79, 91, 92, 93, 94, 108, 130, 142, 145, 169, 170, 194, 197, 199, 203, 206, 211, 212, 222, 223, 251, 252, 272, 273, 280, 289, 305, 306, 307, 309, 313, 361, 362, 363, 364, 366, 367, 368, and 369. (Each one of these 53 forms is implied by one of the following 16 forms: 6, 11, 37, 73, 91, 92, 130, 273, 305, 309, 313, 361, 363, 368, 369.) However, 15 is false in each permutation model. In addition, 182, 189, 190, and 191 are true in every permutation model in which the class of atoms is a set and 192 is false in all such models.

Many of the results obtained for $\mathrm{ZF}^{0}$ models can be transferred to ZF models because of the work of Jech/Sochor\index{Author}{Jech/Sochor} \cite{bib:1966a}, \cite{bib:1966b}, and Pincus\index{Author}{Pincus} \cite{bib:1971}, \cite{bib:1972a}, \cite{bib:1974c}, \cite{bib:1977a}. (See note 103 for a summary of the transfer results.) If a result is transferable\index{Subject}{transferable}, one of the numerals 3 or 4 will occur in the appropriate position in Table 1. In the corresponding place in part V, REFERENCES FOR RELATIONSHIPS BETWEEN FORMS, we will give the $\mathrm{ZF}^{0}$ model with a ``T'' following its number. (The following subset of the above list of forms that are true in every permutation model are boundable statements so they are transferable: 5, 6, 11, 12, 13, 35, 36, 37, 38, 70, 73, 74, 79, 92, 93, 94, 130, 142, 169, 170, 194, 197, 199, 203, 206, 211, 212, 222, 223, 251, 252, 272, 273, 280, 289, 305, 306, 307, 309, and 313. (Each one of these 42 forms is implied by one of the following 12 forms: 6, 11, 37, 73, 92, 130, 170, 273, 305, 309, 313.) In addition, the following forms are transferable: 8, 9, 10, 14, 18, 17, 30, 31, 32, 39, 40, 43, 44, 45, 46, 47, 48, 52, 60, 61, 64, 71, 80, 82, 83, 85, 86, 87, 88, 98, 119, 120, 121, 122, 123, 128, 132, 163, 165, 166, 178, 185, 213, 216, 217, 250, 295, 336, 342, 350, 357, 358, 373, 374, 376, 377, and 378. See note 103.) Since the lists of forms given above occur in note 18, we shall give a reference to that note when, in the description of a model, we assert that one of the forms listed above is true (or false) in the model.

Also, there are several implications involving more than two forms that are known to hold in $\mathrm{ZF}^{0}.$ These are listed with references in note 120. Most are of the form A + B implies C, where A, B, and C are form numbers. When it follows from these implications that a certain form is true or false in a model we will, in the description of the model, assert that the form is true (or false) and give a reference to the appropriate part of note 120.

Finally, in note 2 we give a list of implications which hold in every FM model and involve more than two forms. When the truth of falsity of a form in an FM model follows from one of these, we shall refer to the relevant part of note 2.

\noindent $\mathcal{N}1$: The Basic Fraenkel Model. The set of atoms, $A$ is denumerable; $\mathcal{G}$ is the group of all permutations on $A$; and $S$ is the set of all finite subsets of $A$. In $\mathcal{N}1,\,A$ is infinite, but neither $A$ nor $\mathcal{P}(A)$ have a denumerable subset (9 and 82 are false); $A$ is amorphous (64 is false); $|A|^{2}\not\leq 2^{|A|}$ (53 is false); and $MC_{\omega}(\aleph_{0},\,\infty)$ (131) is false because if $F_{n}$ is the set of all $n$-element subsets of $A$, then $\{F_{n}:n\in\omega\}$ has no countable multiple choice function. However, $A$ (Antichain Condition 89) is true; the Weak Ultra Filter Principle (63) is true, but the Boolean Prime Ideal Theorem (14) is false; Ramsey's Theorem (17) is true; every set that cannot be well ordered has an infinite subset in one-to-one correspondence with a subset of $A$ (133 is true) (Blass\index{Author}{Blass} \cite{bib:1977a}); there is a free group that has a subgroup that is not free (68 is false); and there is a field that does not have an algebraic closure (69 is false). (See L\"{a}uchli\index{Author}{Lauchli@L\"{a}uchli} \cite{bib:1962} and Jech\index{Author}{Jech} \cite{bib:1973b} p145ff for these last two results.) Also, AC for a set of pairs (88) is false.

In Brunner\index{Author}{Brunner} \cite{bib:1984b} it is shown that 17 + 133 implies 134, so 134 is true. It is shown in note 123 that 133 implies 231 $(UT(WO,\,WO,\,WO))$. In every FM model 231 implies 23 $((\forall\alpha)UT(\aleph_{\alpha},\,\aleph_{\alpha},\,\aleph_{\alpha}))$. It follows that 23 is also true. It is shown in note 88 that the Countable Union Theorem (31) is true. Brunner \cite{bib:1984b} has shown that an amorphous power of a compact $T_{2}$ space is compact (127 is true) and that if $X$ is a $T_{2}$ space with at least two points and $X^{Y}$ is hereditarily metacompact, then $Y$ is countable (135 is true). Also, in this paper, Brunner shows that the Baire Category Theorem for compact Hausdorff spaces (106) is false (127 + 106 implies 64). In Howard/Rubin\index{Author}{Howard/Rubin} \cite{bib:1977} it is shown that $MC(\infty,\,LO)$ (112) is true, but $MC(LO,\,\infty)$ (292) is false. Brunner \cite{bib:1985a} shows that a product of weakly Loeb $T_{2}$-spaces is weakly Loeb (115 is true) and that every compact $T_{2}$-space is weakly Loeb (116 is true). In Brunner \cite{bib:1983d} it is shown that there is a $T_{2}$-space $X$ which is $F$ (a continuous finite-to-one image of an $A_{1}$-space), but not $A_{1}$ (for every open covering $O$ of $X$ there is a mapping $f:X\rightarrow O$ such that for each $x\in X,\,x\in f(x))$ (146 is false) and also that every $T_{2}$ topological space can be covered by a well ordered family of discrete sets (147 is true). Brunner \cite{bib:1984f} shows that in this model there is an infinite $T_{2}$-space which does not have an infinite set of non-isolated points (128 is false). (Note that the statement that $\mathcal{P}(\mathbb{R})$ can be well ordered (130) is true in every model of $\mathrm{ZF}^{0}$.)

Brunner \cite{bib:1981a} shows that the box product $\mathbb{R}^{A}$ is metrizable and connected (177 is false). (See note 52.) Dawson/Howard show that $|A!|$ and $|2^{A}|$ are incomparable (200 is false). (See note 64.) Hickman proves that in this model it is not possible to construct a binary operation on $A$ which makes $A$ a group (239 is false). Stavi proves that there exist an infinite, free, complete Boolean algebra on $A$(267 is false). (See note 89 for definitions.) Hodges proves that if $R$ is the polynomial ring over $A$ with coefficients in $\mathbb{Q}$ then $R$ is an integral domain such that every ideal in $R$ is finitely generated, but $R$ has no maximal proper ideal (278 is false). It is shown in Harper/Rubin that in this model every set on which there is a relation $R$ which is anti-symmetric and connected has a maximal subset on which $R$ is transitive (263 is true). In Brunner \cite{bib:1985a} it is shown that in this model for $T_{2}$ spaces, $A$-bounded is equivalent to weakly Loeb (114 is true). In this model, every infinite partially ordered set has either an infinite chain or an infinite antichain (217 is true). (See note 105.)

In note 41 we show that in $\mathcal{N}1$ algebraic closures of fields are unique when they exist (233 is true). In note 116 we show that in $\mathcal{N}1$ every $\mathrm{T}_{2}$ topological space contains an infinite subset which has no infinite compact subsets (304) is true. In note 66 we show that $KW(WO,\,\infty)$ (322) is true. (This result also appears in Brunner\index{Author}{Brunner} \cite{bib:1982a}, Proposition 4.5.) The set of all finite subsets of $A$ that have at least two elements is a set of well orderable sets for which there is no Kinna-Wagner selection function. Therefore, $KW(\infty,\,WO)$ (323) is false. Ramsey's Theorem II (325) is true (see note 46). Form 40 $(C(WO,\,\infty))$ is false because 40 implies 64 and form 122 $(C(WO,\,<\aleph_{0}))$ is true because 133 implies 122. Therefore, form 328 $(MC(WO,\,\infty))$ is false because 122 + 328 $\rightarrow$ 40. De la Cruz and Di Prisco have shown that every infinite collection of non-empty well orderable sets has an infinite subfamily with a choice function (380 is true). It is shown in Howard/Keremedis/Rubin/Stanley\index{Author}{Howard/Keremedis/Rubin/Stanley} \cite{bib:1997} that form 232 (Every metric space $(X, d)$ has a $\sigma$-discrete basis.) is true. Since form 165 $(C(WO,\,WO))$ is true, (133 implies 165) it follows from note 2(8 and 9) that 16 and 24 are also true.

\noindent $\mathcal{N}1 \models$ 6, 11, 16, 17, 23, 24, 31, 37, 63, 73, 89, 91, 92, 112, 114, 115, 116, 127, 130, 133, 134, 135, 147, 191, 217, 232, 233, 263, 273, 304, 305, 309, 313, 322, 325, 361, 363, 368, 369 and 380, but 15, 53, 64, 68, 69, 76, 106, 126, 128, 131, 146, 177, 200, 239, 267, 278, 292, 323, 328, and 344 are false. References include Blass\index{Author}{Blass} \cite{bib:1977a}, Brunner \cite{bib:1981a}, \cite{bib:1982a}, \cite{bib:1983d}, \cite{bib:1984b}, \cite{bib:1984f}, \cite{bib:1985a}, Dawson/Howard\index{Author}{Dawson/Howard} \cite{bib:1976}, De la Cruz/Di Prisco\index{Author}{De la Cruz/Di Prisco}  \cite{bib:1996}, Fraenkel\index{Author}{Fraenkel} \cite{bib:1922}, Felgner/Jech\index{Author}{Felgner/Jech} \cite{bib:1973}, Felgner\index{Author}{Felgner} \cite{bib:1971a}, Jech\index{Author}{Jech} \cite{bib:1973b}, Jech/Sochor\index{Author}{Jech/Sochor} \cite{bib:1966a}, Halpern\index{Author}{Halpern} \cite{bib:1964}, Harper/Rubin\index{Author}{Harper/Rubin} \cite{bib:1976}, Hickman\index{Author}{Hickman} \cite{bib:1976}, Hodges\index{Author}{Hodges} \cite{bib:1974}, Howard/Keremedis/Rubin/Stanley \cite{bib:1997}, Howard/Rubin\index{Author}{Howard/Rubin} \cite{bib:1977}, L\"{a}uchli\index{Author}{Lauchli@L\"{a}uchli} \cite{bib:1962}, Levy\index{Author}{Levy} \cite{bib:1958}, Pincus\index{Author}{Pincus} \cite{bib:1969}, Specker\index{Author}{Specker} \cite{bib:1957}, Stavi\index{Author}{Stavi} \cite{bib:1975}, notes 2(8,9), 18, 41, 46, 52, 64, 66, 88, 89, 105, 109, 116, 120(56) and 128.

(A minimal list of forms implying each form known to be true in the model: 11, 16, 23, 24, 73, 89, 92, 112, 114, 115, 116, 127, 130, 133, 134, 135, 147, 191, 217, 232, 233, 263, 304, 305, 309, 313, 322, 325, 361, 363, 380.

A minimal list of forms at least one of which is implied by each form known to be false in the model: 53, 64, 69, 128, 131, 146, 177, 200, 239, 267, 278, 292, 344.

A list of all forms known to be true in the model: 0, 5, 6, 10, 11, 12, 13, 16, 17, 18, 19, 23, 24, 25, 26, 27, 28, 31, 32, 34, 35, 36, 37, 38, 47, 48, 51, 58, 63, 70, 73, 74, 77, 79, 80, 84, 89, 90, 91, 92, 93, 94, 104, 108, 111, 112, 114, 115, 116, 118, 119, 122, 124, 127, 130, 132, 133, 134, 135, 142, 145, 147, 151, 173, 182, 185, 189, 190, 191, 194, 197, 199, 203, 206, 207, 208, 209, 211, 212, 216, 217, 221, 222, 223, 231, 232, 233, 241, 242, 249, 250, 251, 252, 254, 263, 272, 273, 280, 288, 289, 294, 304, 305, 306, 307, 309, 313, 315, 316, 322, 324, 325, 327, 342, 350, 355, 357, 358, 361, 362, 363, 364, 366, 367, 368, 369, 371, 373, 374, 380, 383, (383).

A list of all forms known to be false in the model: 1, 2, 3, 4, 7, 8, 9, 14, 15, 30, 39, 40, 41, 43, 44, 45, 49, 50, 53, 57, 60, 61, 62, 64, 66, 67, 68, 69, 71, 76, 82, 83, 85, 86, 87, 88, 95, 100, 101, 106, 107, 109, 113, 123, 126, 128, 129, 131, 146, 149, 152, 161, 168, 174, 177, 181, 188, 192, 193, 200, 201, 202, 213, 214, 215, 218, 239, 255, 256, 257, 258, 259, 260, 261, 262, 264, 267, 270, 278, 284, 286, 292, 295, 296, 303, 317, 323, 326, 328, 331, 332, 333, 334, 335, 336, 343, 344, 345, 346, 347, 359, 376, 377, 378, (383).)

\noindent $\mathcal{N}2$: The Second Fraenkel Model. $A=\{a_{i}:i\in\omega\}$. Let $B=\{\{a_{2i},\,a_{2i+1}\}:i\in \omega\}$. $\mathcal{G}$ is the group of permutations on $A$ that leaves $B$ point-wise fixed. $S$ is the set of all finite subsets of $A$. In $\mathcal{N}2$, MC (67) is true. (In fact, $MC(\infty,\,\infty,$ even) (334) is true. See note 93.) However, the set $B$ is a denumerable set of pairs that has no choice function $(C(\aleph_{0},2)$, form 80 is false) and has no Kinna-Wagner selection function $(KW(\aleph_{0},\, <\aleph_{0})$, form 358, is false). $A$ is infinite, but has no countably infinite subset (9 is false); $B$ does have a countably infinite subset (98 is false), but $\bigcup B$ does not (18 is false); the Hahn-Banach Theorem (52) and the Krein-Milman Theorem (65) are true, but the Boolean Prime Ideal Theorem (14) is false.

It also follows from results of Brunner\index{Author}{Brunner} \cite{bib:1984b} that in $\mathcal{N}2$, there is a compact $T_{2}$ topological space $X$ such that $X^{\omega}$ is not compact, thus 154 is false. (See note 95.) Brunner \cite{bib:1982a} shows that in this model every set that cannot be well ordered has an infinite Dedekind finite subset (163 is true). Howard\index{Author}{Howard} proves that Rado's Selection Lemma (99) is true. Hickman\index{Author}{Hickman} proves that the only subsets of $A$ that can be linearly ordered are finite, but $A$ has no amorphous subset (198 is false). However, no infinite set in this model is $T$ finite (83 is true). (See note 105.) The set of all $n$-element subsets of $A$ has no choice function for each $n\geq 2\ (45(n)$ is false).

Since 106 (the Baire Category Theorem for compact Hausdorff spaces) is true (67 implies 106) and 43 (Dependent Choices) is false (43 implies 18), it follows that the Tychonoff Theorem for countably many $T_{2}$ spaces (154) is false because it is shown in Brunner \cite{bib:1983c} that 106 + 154 implies 43. Form 130 $(\mathcal{P}(\mathbb{R})$ is well orderable.) is true in every FM model, 18 is false in $\mathcal{N}2$, and Brunner \cite{bib:1984f} has shown that 130 + 128 (On every infinite set there is a $T_{2}$ topology with an infinite set of non-isolated points.) implies 18. Thus, 128 must be false in $\mathcal{N}2$.

\noindent $\mathcal{N}2 \models$ 6, 11, 37, 65, 67, 73, 83, 92, 99, 130, 163, 191, 273, 305, 313, 334, 361, and 363, but 18, 45$(n)$, 80, 98, 128, 154, 163, 164, 198, 344, and 358 are false. References Brunner \cite{bib:1982a}, \cite{bib:1983c}, \cite{bib:1984b}, \cite{bib:1984f}, Fraenkel\index{Author}{Fraenkel} \cite{bib:1922}, Hickman \cite{bib:1978b}, Howard \cite{bib:1984b}, Jech\index{Author}{Jech} \cite{bib:1973b}, Pincus\index{Author}{Pincus} \cite{bib:1972a}, \cite{bib:1972b}, notes 18, 93, 95, 105, and 120(2 and 56).

(A minimal list of forms implying each form known to be true in the model: 6, 11, 37, 65, 73, 92, 99, 130, 163, 191, 305, 313, 334, 361, 363.

A minimal list of forms at least one of which is implied by each form known to be false in the model: 18, 45, 98, 128, 154, 163, 164, 198, 344, 358.

A list of all forms known to be true in the model: 0, 5, 6, 11, 12, 13, 19, 25, 34, 35, 36, 37, 38, 51, 52, 58, 64, 65, 67, 70, 73, 74, 76, 77, 78, 79, 82, 83, 84, 89, 90, 91, 92, 93, 94, 99, 104, 106, 108, 112, 114, 118, 119, 125, 126, 127, 130, 131, 142, 144, 145, 147, 155, 163, 173, 182, 185, 189, 190, 191, 194, 197, 199, 203, 206, 208, 211, 212, 221, 222, 223, 225, 232, 251, 252, 254, 272, 273, 280, 287, 289, 292, 294, 305, 306, 307, 309, 313, 315, 316, 328, 329, 330, 334, 349, 350, 361, 362, 363, 364, 366, 367, 368, 369, 371, 375, 381, 382, 383, (383).

A list of all forms known to be false in the model: 1, 2, 3, 4, 7, 8, 9, 10, 14, 15, 16, 17, 18, 20, 21, 23, 27, 28, 29, 30, 31, 32, 39, 40, 41, 43, 44, 45, 49, 50, 60, 61, 62, 68, 71, 80, 85, 86, 87, 88, 95, 98, 100, 101, 107, 109, 111, 113, 121, 122, 123, 128, 129, 132, 133, 150, 151, 152, 153, 154, 161, 163, 164, 165, 166, 167, 168, 174, 181, 192, 198, 201, 202, 213, 214, 215, 218, 231, 250, 255, 256, 257, 258, 259, 260, 261, 262, 264, 270, 284, 286, 295, 303, 317, 322, 323, 324, 325, 326, 327, 331, 332, 333, 335, 343, 344, 345, 346, 347, 352, 354, 355, 358, 359, 376, 377, 378, 380, (383).)

\noindent $\mathcal{N}2(n)$: A generalization of $\mathcal{N}2$. This is a generalization of $\mathcal{N}2$ in which there is a denumerable set of $n$ element sets for $n\in \omega-\{0, 1\}$. $A=\bigcup B$, where $B$ is a set of pairwise disjoint $n$ element sets $b_{i}$ with $i\in\omega$; $\mathcal{G}$ is the group of all permutations that leave $B$ point-wise fixed $(\phi(b_{i})=b_{i}$ for each $i\in\omega)$; and $S$ is the set of all finite subsets of $A$. In this model, the Axiom of Multiple Choice (67) is true, but $C(WO,\,n)\ (47(n))$, $C(\aleph_{0},\,n)\ (288(n))$ and $PC(\aleph_{0},\,n,\,\infty)\ (373(n))$ are false. Let $K$ be a non-empty set of positive integers, then for $n\in K,\,C(\infty,\,K)$ is false. For $n\in K$, we shall denote $C(\infty,\,K)$ by $46(n\in K)$.

\noindent $\mathcal{N}2(n) \models$ 6, 11, 37, 67, 73, 92, 130, 191, 273, 305, 313, 361, and 363, but $46(n\in K)$, $47(n)$, 154, 164, 288$(n)$, 344, and 373$(n)$ are false. References Fraenkel\index{Author}{Fraenkel} \cite{bib:1922}, Mostowski\index{Author}{Mostowski} \cite{bib:1945}, Zuckerman\index{Author}{Zuckerman} \cite{bib:1971}, notes 15, 18, and 120(2, 10, and 56).

(A minimal list of forms implying each form known to be true in the model: 6, 11, 37, 67, 73, 92, 130, 191, 305, 313, 361, 363.

A minimal list of forms at least one of which is implied by each form known to be false in the model: 46, 154, 164, 344, 373.

A list of all forms known to be true in the model: 0, 5, 6, 11, 12, 13, 19, 25, 34, 35, 36, 37, 38, 51, 52, 58, 64, 67, 70, 73, 74, 76, 77, 78, 79, 82, 83, 84, 89, 90, 91, 92, 93, 94, 104, 106, 108, 112, 114, 118, 119, 125, 126, 127, 130, 131, 142, 144, 145, 147, 155, 173, 182, 185, 189, 190, 191, 194, 197, 199, 203, 206, 208, 211, 212, 221, 222, 223, 232, 251, 252, 254, 272, 273, 280, 287, 289, 292, 294, 305, 306, 307, 309, 313, 315, 316, 328, 329, 330, 349, 350, 361, 362, 363, 364, 366, 367, 368, 369, 371, 375, 381, 382, 383, (383).

A list of all forms known to be false in the model: 1, 2, 3, 4, 7, 8, 9, 10, 14, 15, 16, 17, 20, 21, 23, 27, 28, 29, 30, 31, 32, 33, 39, 40, 41, 43, 44, 45, 46, 47, 49, 50, 60, 61, 62, 68, 71, 85, 86, 87, 100, 101, 107, 109, 113, 121, 122, 123, 129, 132, 133, 150, 151, 152, 153, 154, 161, 164, 165, 168, 174, 181, 192, 202, 213, 214, 218, 231, 250, 255, 256, 257, 258, 259, 260, 261, 262, 264, 270, 284, 286, 288, 295, 303, 317, 322, 323, 324, 325, 327, 331, 332, 335, 343, 344, 345, 346, 347, 352, 354, 359, 373, 374, 376, 377, 378, 380, (383).)

\noindent $\mathcal{N}2^{\ast}(3)$: Howard's\index{Author}{Howard} variation\index{Author}{Fraenkel} of $\mathcal{N}2(3)$. $A=\bigcup B$, where $B$ is a set of pairwise disjoint 3 element sets, $T_{i}=\{a_{i},\,b_{i},\,c_{i}\}$. For each $i\in\omega$ define a function $\eta_{i}:T_{i}\rightarrow T_{i}$ such that $\eta_{i}:a_{i}\mapsto b_{i}\mapsto c_{i}$. $\mathcal{G}$ is the group of all permutations $\phi$ of $A$ such that for each $i\in \omega,\,\phi|T_{i}$ is either the identity, $\eta_{i}$, or $\eta_{i}^{2}$. $S$ is the set of all finite subsets of $A$. In this model, $B$ is a denumerable set of triples that has no choice function and has no Kinna-Wagner selection function so $(\forall n\in \omega-\{0, 1\})C(\aleph_{0},\,n)$ (288) is false and $KW(\aleph_{0},\,<\aleph_{0})$ (358) is false. (Since the Boolean Prime Ideal Theorem (14) implies 288, form 14 is also false.) Howard shows that if $P$ is any set of pairs in the model and $R$ is a symmetric relation on $\bigcup P$ such that every finite subset of $P$ has an $R$-consistent choice function, (A choice function $f$ on $X$ is $R$-consistent if for all $u,\,v\in X,\,u\neq v$ implies $f(u)Rf(v)$.) then $P$ has an $R$-consistent choice function (141 is true). It is shown in note 93 that $MC(\infty,\,\infty,\mathrm{odd})$ (333) is true. Keremedis\index{Author}{Keremedis} \cite{bib:1996a} proves that $333 + 334(MC(\infty,\,\infty,\mathrm{even})$ iff AC, so 334 is false.

\noindent $\mathcal{N}2^{\ast}(3)\models$ 6, 11, 37, 73, 92, 130, 141, 191, 273, 305, 313, 333, 361, and 363, but 154, 164, 288, 334, 344, and 358 are false. References Howard \cite{bib:1984a}, Keremedis \cite{bib:1996a}, and notes 18, 93, and 120(2, 10, and 56).

(A minimal list of forms implying each form known to be true in the model: 6, 11, 37, 73, 92, 130, 141, 191, 305, 313, 333, 361, 363.

A minimal list of forms at least one of which is implied by each form known to be false in the model: 154, 164, 288, 334, 344, 358.

A list of all forms known to be true in the model: 0, 5, 6, 11, 12, 13, 18, 19, 25, 34, 35, 36, 37, 38, 51, 52, 58, 64, 67, 70, 73, 74, 76, 77, 78, 79, 80, 82, 83, 84, 88, 89, 90, 91, 92, 93, 94, 104, 106, 108, 112, 114, 118, 119, 125, 126, 127, 130, 131, 140, 141, 142, 144, 145, 147, 155, 173, 182, 185, 189, 190, 191, 194, 197, 199, 203, 206, 208, 211, 212, 221, 222, 223, 232, 251, 252, 254, 268, 269, 272, 273, 276, 280, 287, 289, 292, 294, 305, 306, 307, 309, 313, 315, 316, 328, 329, 330, 333, 349, 350, 361, 362, 363, 364, 366, 367, 368, 369, 371, 375, 381, 382, 383, (383).

A list of all forms known to be false in the model: 1, 2, 3, 4, 7, 8, 9, 10, 14, 15, 16, 17, 20, 21, 23, 27, 28, 29, 30, 31, 32, 33, 39, 40, 41, 43, 44, 45, 47, 49, 50, 60, 61, 62, 68, 71, 85, 86, 87, 100, 101, 107, 109, 113, 121, 122, 123, 129, 132, 133, 150, 151, 152, 153, 154, 161, 164, 165, 168, 174, 181, 192, 202, 213, 214, 218, 231, 250, 255, 256, 257, 258, 259, 260, 261, 262, 264, 270, 284, 286, 288, 295, 303, 317, 322, 323, 324, 325, 327, 331, 332, 334, 335, 343, 344, 345, 346, 347, 352, 354, 355, 358, 359, 374, 376, 377, 378, 380, (383).)

\noindent $\mathcal{N}2(n,\,M)$: Mostowski's variation of $\mathcal{N}2(n)$. $A,\,B$, and $S$ are the same as in $\mathcal{N}2(n)$. In this model, $\mathcal{G}$ is the group of all permutations of $A$ such that for each $\phi \in \mathcal{G},\, \phi(b_{i})=b_{i}$ (as in $\mathcal{N}2(n)$), but in addition $\phi$ is different from the identity on $b_{i}$ for at most a finite number of $i\in\omega$. Let $K$ be a non-empty set of positive integers. The condition $M(K,\, n)$ holds if for every decomposition of $n$ into a sum of primes, $n=p_{1}+p_{2}+\cdots+p_{m}$, there is a sequence of $m$ non-negative integers, $q_{1},\,q_{2},\,\cdots,\,q_{m}$, such that $p_{1}q_{1}+p_{2}q_{2}+\cdots+p_{m}q_{m}\in K$. Mostowski proves that $C(\infty,\,K)$ (46) is true in this model if $M(K,\,n)$ is false. We shall denote $C(\infty,\,K)$ by $46(\neg M)$ when condition $M(K,\,n)$ is false. Also, it is clear that $C(WO,\,n)\ (47(n))$ and $C(\infty,\,K)\ (46(n\in K))$ are false in this model.

\noindent $\mathcal{N}2(n,\,M) \models$ 6, 11, 37, 46$(\neg M)$, 73, 91, 92, 130, 191, 273, 305, 309, 313, 361, 363, 368, and 369, but 15, 46$(n\in K)$, 47$(n)$, 192, and 344 are false. References Fraenkel \cite{bib:1922}, Mostowski \cite{bib:1945}, Pincus\index{Author}{Pincus} \cite{bib:1977a}, notes 73 and 120(56).

(A minimal list of forms implying each form known to be true in the model: 6, 11, 37, 46, 73, 91, 92, 130, 191, 305, 309, 313, 361, 363.

A minimal list of forms at least one of which is implied by each form known to be false in the model: 46, 47, 192, 344.

A list of all forms known to be true in the model: 0, 5, 6, 11, 12, 13, 35, 36, 37, 38, 46, 48, 70, 73, 74, 79, 91, 92, 93, 94, 108, 120, 130, 142, 145, 182, 189, 190, 191, 194, 197, 199, 203, 206, 211, 212, 222, 223, 251, 252, 272, 273, 280, 289, 305, 306, 307, 309, 313, 361, 362, 363, 364, 366, 367, 368, 369, 371, (383).

A list of all forms known to be false in the model: 1, 14, 15, 20, 21, 23, 30, 33, 40, 45, 46, 47, 49, 50, 60, 61, 62, 68, 85, 101, 107, 109, 121, 122, 123, 133, 151, 165, 192, 202, 213, 218, 231, 250, 255, 256, 257, 258, 259, 260, 261, 262, 264, 270, 284, 286, 295, 303, 317, 322, 323, 324, 327, 331, 332, 343, 344, 345, 346, 347, 359, (383).)

\noindent $\mathcal{N}2(\aleph_{\alpha})$: Jech's Model. This is an extension of $\mathcal{N}2$ in which $A=\{a_{\gamma}:\gamma\in \omega_{\alpha}\}$; $B$ is the corresponding set of $\aleph_{\alpha}$ pairs of elements of $A$; $\mathcal{G}$ is the group of all permutations on $A$ that leave $B$ point-wise fixed; and $S$ is the set of all subsets of $A$ of cardinality less than $\aleph_{\alpha}$. $(\mathcal{N}2(\aleph_{0})$ is $\mathcal{N}2$.) If $\aleph_{\alpha}$ is regular and $\alpha>1$, then $DC(\omega_{1})$ (44) is true in this model; and if $\alpha>\beta$, then $DC(\omega_{\beta})\ (87(\beta))$ is true. (Therefore, the Principle of Dependent Choices (DC, 43) is true if $\alpha>0$.) On the other hand, $C(WO,\,<\aleph_{0})$ (122) is false for all $\alpha\geq 0,\,C(\infty,\,n)$ (45$(n)$ and 88) is false for each $n\geq 2$ and $\alpha\geq 0$, and $KW(WO,\,<\aleph_{0})$ (327) is false for all $\alpha\geq 0$.

Since BPI (14) implies 88, 14 is also false. Brunner\index{Author}{Brunner} \cite{bib:1982a} has shown that for all $\alpha>0$ there is a set that cannot be well ordered and does not have an infinite Dedekind finite subset (163 is false) because 44 implies 8 $(C(\aleph_{0},\,\infty))$ and 8 plus 163 iff AC. Also, Jech\index{Author}{Jech} (\cite{bib:1973b}, p 132, prob 5) notes that if $\alpha\geq 1$, since there is a non-trivial ultrafilter on $\omega$ (70 is true) and every infinite set has a countably infinite subset (9 is true), the weak ultrafilter principal (63) is true. Form 44 implies forms 39 $(C(\aleph_{1},\,\infty))$ and 41 (Every cardinal is comparable with $\aleph_{1}$.). van Douwen\index{Author}{Douwen} \cite{bib:1980} has shown that 39 + 41 implies 135 (If $X$ is a $T_{2}$ space with at least two points and $X^{Y}$ is hereditarily metacompact, then $Y$ is countable.). Therefore, form 135 is also true. Blass\index{Author}{Blass} \cite{bib:1979} has shown that 191 (There is a set $X$ such that for each set $a$ there is an ordinal $\alpha$ and a function $f$ mapping $X\times\alpha$ onto $a$.) is true in every permutation model with a set of atoms. He has also shown that 8 $(C(\aleph_{0},\,\infty)) +191+193$ (Every Abelian group is a homomorphic image of a free projective Abelian group.) iff AC. Thus, form 193 must be false if $\aleph_{\alpha}$ is regular and $\alpha>1$. (See note 60 for definitions.) It follows from note 104 that the $2m=m$ principle (3) is false in this model.

\noindent $\mathcal{N}2(\aleph_{\alpha})\models$ 6, 11, 37, 43 (for $\alpha>0$), 44 (if $\aleph_{\alpha}$ is regular and $\alpha>1$), 63 (if $\alpha>0$), 73, 87$(\beta)$ (if $\alpha>\beta)$, 91, 130, 135 (if $\aleph_{\alpha}$ is regular and $\alpha>1$), 191, 273, 305, 309, 313, 361, 363, 368, and 369, but 3, 15, 40, $45(n)$, 88, 122, 163, 192, 193, 327, and 344 are false. References Blass \cite{bib:1979}, Brunner \cite{bib:1982a}, Jech \cite{bib:1966b}, \cite{bib:1973b}, Mostowski\index{Author}{Mostowski} \cite{bib:1948}, van Douwen \cite{bib:1980}, notes 18, 60, 104, and 120(56).

(A minimal list of forms implying each form known to be true in the model: 11, 44, 63, 87, 91, 130, 135, 191, 305, 309, 313.

A minimal list of forms at least one of which is implied by each form known to be false in the model: 3, 45, 88, 163, 192, 193, 327, 344.

A list of all forms known to be true in the model: 0, 5, 6, 8, 9, 10, 11, 12, 13, 16, 17, 18, 19, 24, 26, 27, 28, 29, 31, 32, 34, 35, 36, 37, 38, 39, 41, 42, 43, 44, 57, 63, 64, 70, 71, 73, 74, 75, 77, 78, 79, 80, 82, 83, 84, 86, 87, 91, 92, 93, 94, 98, 104, 106, 108, 113, 119, 124, 126, 127, 128, 130, 131, 132, 135, 142, 145, 150, 154, 155, 174, 182, 183, 185, 189, 190, 191, 194, 196, 197, 198, 199, 203, 206, 209, 211, 212, 216, 217, 221, 222, 223, 243, 249, 251, 252, 272, 273, 279, 280, 282, 288, 289, 296, 304, 305, 306, 307, 309, 313, 325, 336, 337, 338, 339, 340, 341, 342, 350, 351, 352, 353, 354, 355, 357, 358, 360, 361, 362, 363, 364, 365, 366, 367, 368, 369, 371, 373, 374, 375, 376, 377, 378, 380, (383).

A list of all forms known to be false in the model: 1, 2, 3, 14, 15, 20, 21, 23, 30, 40, 45, 49, 50, 60, 61, 62, 68, 85, 88, 95, 101, 107, 109, 121, 122, 123, 133, 151, 163, 165, 192, 193, 201, 202, 213, 218, 231, 255, 256, 257, 258, 259, 260, 261, 262, 264, 270, 284, 286, 295, 303, 317, 322, 323, 324, 326, 327, 331, 332, 333, 335, 343, 344, 345, 346, 347, 359, (383).)

\noindent $\mathcal{N}2(\mathrm{LO}):$ van Douwen's Model. This model is another variation of $\mathcal{N}2$. The set of atoms, $A=\bigcup_{n\in \omega}B_{n}$ where the $B_{n}$'s are pairwise disjoint sets each of which is ordered like the integers, $\mathbb{Z}$. $\mathcal{G}$ is the group of all permutations which fix each $B_{n}$ and also preserve the ordering on each $B_{n},\,n\in\omega$. $S$ is the set of all finite subsets of $A$. In $\mathcal{N}2(\mathrm{LO})$, $\{B_{n}:n\in\omega\}$ is a denumerable set of sets each of which is ordered like $\mathbb{Z}$ which has no choice function, so form 119 is false.

\noindent $\mathcal{N}2(\mathrm{LO}) \models$ 6, 11, 37, 73, 91, 92, 130, 191, 273, 305, 309, 313, 361, 363, 368, and 369, but 15, 62, 119, 323, and 344 are false. References van Douwen\index{Author}{Douwen} \cite{bib:1985}, notes 2(3, 6, and 56) and 18.

(A minimal list of forms implying each form known to be true in the model: 6, 11, 37, 73, 91, 92, 130, 191, 305, 309, 313, 361, 363.

A minimal list of forms at least one of which is implied by each form known to be false in the model: 62, 119, 344.

A list of all forms known to be true in the model: 0, 5, 6, 11, 12, 13, 35, 36, 37, 38, 70, 73, 74, 79, 91, 92, 93, 94, 108, 130, 142, 145, 182, 189, 190, 191, 194, 197, 199, 203, 206, 211, 212, 222, 223, 251, 252, 272, 273, 280, 289, 305, 306, 307, 309, 313, 361, 362, 363, 364, 366, 367, 368, 369, 371, (383).

A list of all forms known to be false in the model: 1, 8, 14, 15, 16, 20, 21, 23, 27, 28, 29, 30, 31, 32, 39, 40, 43, 44, 49, 50, 60, 62, 66, 67, 68, 85, 86, 87, 89, 90, 95, 101, 107, 109, 112, 113, 114, 118, 119, 123, 133, 149, 150, 165, 174, 181, 192, 202, 213, 218, 255, 256, 257, 258, 259, 260, 261, 262, 264, 270, 286, 292, 295, 303, 317, 323, 331, 332, 333, 334, 335, 343, 344, 345, 346, 347, 352, 354, 359, 375, (383).)

\noindent $\mathcal{N}3$: Mostowski's Linearly Ordered Model. $A$ is countably infinite; $\precsim$ is a dense linear ordering on $A$ without first or last elements $((A,\,\precsim)\cong(\mathbb{Q},\,\leq));\mathcal{G}$ is the group of all order automorphisms on $(A,\,\precsim)$; and $S$ is the set of all finite subsets of $A$. In $\mathcal{N}3,\, PW$ (91) is true, but LW (90) is false. The set $A$ is infinite, but neither $A$ nor $\mathcal{P}(A)$ have a countably infinite subset (9 and 82 are false); every $T$ finite set is finite (83 is true); $A$ is not $T$ finite, but $\mathcal{P}(A)$ is Dedekind finite (84 is false); every set can be linearly ordered (30 is true); the Boolean Prime Ideal Theorem (14) is true; AC for a family of well orderable sets (60) is true; the Cardinal Representative Theorem (97) is false; and the Baire category theorem for compact Hausdorff spaces (106) is false. Also, $MC_{\omega}(\aleph_{0},\,infty)$ (131) is false because if $F_{n}$ is the set of all $n$-element subsets of $A$, then $\{F_{n}:n\in\omega\}$ has no countable multiple choice function.

It is shown in note 88 that the Countable Union Theorem (31) is true. Form 23 is true because 60 implies $UT(WO,\,WO,\,WO)$ (231) and, in every FM model, 231 implies $(\forall\alpha)UT(\aleph_{\alpha},\,\aleph_{\alpha},\,\aleph_{\alpha})$ (23). Brunner \cite{bib:1984c} has shown that there does exist an infinite connected compact P-space (125 is false) and in Brunner\index{Author}{Brunner} \cite{bib:1982a} it is shown that every infinite set that cannot be well ordered has an infinite subset with a Dedekind finite power set (164 is true). In Brunner \cite{bib:1983a} it is shown that in $\mathcal{N}3$ there is a non-trivial L\"{a}uchli continuum (155 is false), there are two compact $T_{2}$ spaces which have isomorphic rings of real valued continuous functions, but are not homeomorphic (156 is false), and there is a compact $T_{2}$, space, $X$, with the property that every non-empty subset of $C(X)$ (the set of continuous real valued functions on $X$) which is point-wise bounded has a supremum, but $X$ is not extremely disconnected (the closure of every open set is open) (157 is false). In Brunner \cite{bib:1983d} it is shown that if $a,\,b\in A,\,a\prec b$, then the closed interval $[a,\, b]$ is a compact $T_{2}$ space that cannot be covered by a well ordered set of discrete sets (147 is false). Dawson/Howard show that $|A!|>|2^{A}|$ (200 is false), and since it is clear that $|A!|\leq|A^{A}|$, 290 is false. (See note 64).

Gonzalez has shown that there is an infinite set in $\mathcal{N}3$ which does not have a dense linear ordering in $\mathcal{N}3$ (295 is false). Since form 295 is false and 30 is true it follows from Pincus\index{Author}{Pincus} \cite{bib:1997} that form 296 (Every infinite set is the disjoint union of infinitely many infinite sets.) is false. (30 + 296 implies 295.) Using the fact that the Kinna-Wagner Selection Principle (15) is false and that the Ordering Principle (30) is true, Krom proves that there is an ordered topological space that is not normal (118 is false). Brunner \cite{bib:1985c} comments that the order topology on $A$ is $T_{2}$ and each infinite subset of $A$ contains a closed interval which is compact (304 is false). In Brunner \cite{bib:1982a}, it is shown that in this model, every set is either well orderable or contains a copy of sub-intervals of $A$. It is noted in Brunner\index{Author}{Brunner} \cite{bib:1984f} that an interval with the order-topology is a $T$-2 space with an infinite number of non-isolated points so it follows that form 128 is true. Howard\index{Author}{Howard} \cite{bib:1973} has shown that in every FM model 14 implies the Weak Sikorski Theorem (317), so 317 is true,

For each $i\in\omega-\{0\}$, let $c_{i}$ be the set of all $i$-element subsets of $A$. Then the set $\{c_{i}:i\in \omega-\{0\}\}$ is a denumerable set that has no multiple choice function $(MC(\aleph_{0},\, \infty)$ (126) is false) and no Kinna-Wagner Selection function $(KW(\aleph_{0},\,\infty)$ (355) is false). It is shown in Howard/Keremedis/Rubin/Rubin \cite{bib:1997b} that form 381 (DUM) is true. Since form 165 $(C(WO,\,WO))$ is true, (60 implies 165) it follows from note 2(8 and 9) that 16 and 24 are also true.

\noindent $\mathcal{N}3\models$ 6, 11, 14, 16, 23, 24, 31, 37, 60, 73, 83, 91, 92, 128, 130, 164, 165, 191, 273, 305, 309, 313, 317, 361, 363, 368, 369, and 381, but 15, 84, 90, 97, 106, 118, 125, 126, 131, 147, 155, 156, 157, 200, 253, 290, 295, 296, 304, and 355 are false. References Mostowski\index{Author}{Mostowski} \cite{bib:1939}, Brunner \cite{bib:1982a}, \cite{bib:1983a}, \cite{bib:1983c}, \cite{bib:1983d}, \cite{bib:1984c}, \cite{bib:1984f}, \cite{bib:1985c}, Dawson/Howard\index{Author}{Dawson/Howard} \cite{bib:1976}, Gonzalez\index{Author}{Gonzalez} \cite{bib:1995a}, Halpern\index{Author}{Halpern} \cite{bib:1964}, Howard \cite{bib:1973}, Howard/Keremedis/Rubin/Rubin\index{Author}{Howard/Keremedis/Rubin/Rubin} \cite{bib:1997b}, Howard/Yorke\index{Author}{Howard/Yorke} \cite{bib:1989}, Jech\index{Author}{Jech} \cite{bib:1973b}, Krom\index{Author}{Krom} \cite{bib:1986}, L\"{a}uchli\index{Author}{Lauchli@L\"{a}uchli} \cite{bib:1964}, Levy\index{Author}{Levy} \cite{bib:1958}, Pincus\index{Author}{Pincus} \cite{bib:1969}, \cite{bib:1997}, Sageev\index{Author}{Sageev} \cite{bib:1981}, notes 2(8,9), 18 and 120(34 and 45).

(A minimal list of forms implying each form known to be true in the model: 11, 16, 23, 24, 60, 73, 92, 128, 130, 164, 165, 191, 305, 317, 361, 363, 381.

A minimal list of forms at least one of which is implied by each form known to be false in the model: 84, 97, 118, 125, 131, 147, 155, 156, 157, 200, 253, 290, 295, 296, 304, 355.

A list of all forms known to be true in the model: 0, 5, 6, 10, 11, 12, 13, 14, 16, 18, 19, 23, 24, 25, 26, 27, 28, 30, 31, 32, 33, 34, 35, 36, 37, 38, 45, 46, 47, 48, 49, 52, 60, 61, 62, 63, 64, 69, 70, 72, 73, 74, 79, 80, 83, 85, 88, 91, 92, 93, 94, 96, 99, 102, 104, 107, 108, 111, 119, 120, 121, 122, 123, 127, 128, 130, 132, 137, 139, 140, 141, 142, 145, 146, 150, 151, 153, 154, 164, 165, 178, 182, 189, 190, 191, 194, 197, 198, 199, 201, 203, 206, 207, 209, 211, 212, 213, 216, 221, 222, 223, 225, 226, 227, 228, 229, 231, 233, 235, 241, 242, 249, 250, 251, 252, 268, 269, 270, 271, 272, 273, 276, 280, 283, 285, 287, 288, 289, 293, 294, 298, 305, 306, 307, 308, 309, 311, 313, 314, 315, 317, 323, 324, 326, 327, 329, 330, 331, 332, 336, 342, 343, 344, 349, 350, 356, 357, 358, 361, 362, 363, 364, 366, 367, 368, 369, 370, 371, 373, 374, 377, 378, 380, 381, 382, (383).

A list of all forms known to be false in the model: 1, 2, 3, 4, 7, 8, 9, 15, 39, 40, 41, 43, 44, 51, 66, 67, 71, 76, 77, 78, 82, 84, 86, 87, 89, 90, 95, 97, 100, 101, 106, 109, 112, 113, 114, 118, 125, 126, 129, 131, 133, 144, 147, 149, 152, 155, 156, 157, 161, 168, 174, 179, 181, 185, 188, 192, 193, 200, 202, 214, 218, 253, 255, 256, 257, 258, 259, 260, 261, 262, 264, 286, 290, 292, 295, 296, 304, 316, 322, 328, 333, 334, 335, 345, 347, 355, 359, 375, (383).)

\noindent $\mathcal{N}4$: The Mathias/Pincus Model I. $A$ is countably infinite; $\precsim$ is a universal homogeneous partial ordering on $A$ (See Jech \cite{bib:1973b} p 101 for definitions.); $\mathcal{G}$ is the group of all order automorphisms on $(A,\,\precsim)$; and $S$ is the set of all finite subsets of $A$. In $\mathcal{N}4$, every linearly ordered set can be well ordered (90) is true but the Antichain Principle (89) is false. There are no amorphous sets (64 is true) and the set of atoms is an infinite $T$ finite set (83 is false). (See note 105.) Since form 64 is true and AC is false, form 133 (Every set is either well orderable or has an infinite amorphous subset.) is false. The set of atoms with the discrete topology is a $T_{2}$-space which is not weakly Loeb (114 is false). Since form 118 (Every orderable topological space is normal.) is true (90 implies 118), it follows from Krom \cite{bib:1986} that not every set can be linearly ordered (30 is false). (The Kinna-Wagner Principle (15) is false in every FM model and Krom proves that $\neg 15+30\rightarrow\neg 118$.)

\noindent $\mathcal{N}4 \models$ 6, 11, 37, 64, 73, 90, 92, 130, 191, 273, 305, 309, 313, 361, and 363, but 30, 83, 89, 114, and 133 are false. References Mathias\index{Author}{Mathias} \cite{bib:1967}, Pincus\index{Author}{Pincus} \cite{bib:1969}, Brunner\index{Author}{Brunner} \cite{bib:1985a}, Felgner/Jech\index{Author}{Felgner/Jech} \cite{bib:1973}, Jech\index{Author}{Jech} \cite{bib:1973b}, Krom\index{Author}{Krom} \cite{bib:1986}, notes 18 and 105.

(A minimal list of forms implying each form known to be true in the model: 6, 11, 37, 64, 73, 90, 92, 130, 191, 305, 309, 313, 361, 363.

A minimal list of forms at least one of which is implied by each form known to be false in the model: 83, 89, 114, 133.

A list of all forms known to be true in the model: 0, 5, 6, 11, 12, 13, 19, 25, 34, 35, 36, 37, 38, 51, 58, 64, 70, 73, 74, 77, 79, 84, 90, 91, 92, 93, 94, 104, 108, 118, 119, 127, 130, 142, 145, 182, 185, 189, 190, 191, 194, 197, 199, 203, 206, 208, 211, 212, 222, 223, 251, 252, 272, 273, 280, 289, 294, 305, 306, 307, 309, 313, 315, 316, 361, 362, 363, 364, 366, 367, 368, 369, 371, (383).

A list of all forms known to be false in the model: 1, 2, 3, 4, 7, 8, 9, 14, 15, 30, 39, 40, 41, 43, 44, 49, 50, 66, 67, 71, 82, 83, 86, 87, 89, 95, 100, 101, 106, 109, 113, 114, 126, 129, 133, 149, 152, 161, 168, 174, 181, 188, 192, 193, 202, 214, 215, 218, 255, 256, 257, 258, 259, 260, 261, 262, 264, 286, 295, 303, 317, 328, 333, 334, 335, 345, 347, 359, (383).)

\noindent $\mathcal{N}5$: The Mathias/Pincus Model II (an extension of $\mathcal{N}4$). $A$ is countably infinite, $\precsim$ and $\leq$ are universal homogeneous partial and linear orderings, respectively, on $A$, (See Jech \cite{bib:1973b} p 101 for definitions.); $\mathcal{G}$ is the group of all order automorphisms on $(A,\precsim,\,\leq)$; and $S$ is the set of all finite subsets of $A$. In $\mathcal{N}5,\precsim$ cannot be extended to a linear ordering on $A$ (49 is false), but every set can be linearly ordered (30 is true). Also, the Boolean Prime Ideal Theorem (14) is false. Since 30 is true, a set is $T$ finite iff it is finite (Levy\index{Author}{Levy} \cite{bib:1958}), so $A$ is not $T$ finite, but $\mathcal{P}(A)$ is Dedekind finite (84 is false). ($T$ finite is equivalent to II-finite from note 94.) Using the fact that the Kinna-Wagner Selection Principle (15) is false and that the Ordering Principle (30) is true, Krom\index{Author}{Krom} proves that there is an ordered topological space that is not normal (118 is false). Since form 14 is false and form 30 implies form 62 $(C(\infty,\,<\aleph_{0}))$, Blass has shown that Rado's Selection Principle (99) is false. (See note 33 for the proof that 62 + 99 implies 14.) Howard\index{Author}{Howard} \cite{bib:1973} has shown that in every FM model, 62 implies 60 $(C(\infty,\,WO))$, so 60 is true. Moreover, in every FM model, 60 + 40 $(C(WO,\,\infty))$ implies AC, so 40 is false. (See Howard/Rubin\index{Author}{Howard/Rubin} \cite{bib:1995b}.) Since form 165 $(C(WO,\,WO))$ is true, (60 implies 165) it follows from note 2(8 and 9) that 16 and 24 are also true.

\noindent $\mathcal{N}5\models$ 6, 11, 16, 23, 24, 30, 37, 60, 73, 91, 92, 130, 165, 191, 273, 305, 309, 313, 361, and 363, but 14, 40, 49, 84, 99, 118, 131, and 144 are false. References Krom \cite{bib:1986}, Mathias \cite{bib:1967}, Pincus \cite{bib:1969}, Felgner/Jech \cite{bib:1973}, Jech \cite{bib:1973b}, Howard \cite{bib:1973}, Howard/Rubin \cite{bib:1995b}, Levy \cite{bib:1958}, notes 2(2,8,9), 18, 33, 94, and 120(28, 34 and 49).

(A minimal list of forms implying each form known to be true in the model: 11, 16, 23, 24, 30, 60, 73, 91, 92, 130, 165, 191, 305, 309, 313, 361, 363.

A minimal list of forms at least one of which is implied by each form known to be false in the model: 49, 84, 99, 118, 131, 144.

A list of all forms known to be true in the model: 0, 5, 6, 10, 11, 12, 13, 16, 18, 19, 23, 24, 25, 26, 27, 28, 30, 31, 32, 33, 34, 35, 36, 37, 38, 45, 46, 47, 48, 60, 61, 62, 64, 70, 73, 74, 79, 80, 83, 85, 88, 91, 92, 93, 94, 102, 104, 108, 111, 119, 120, 121, 122, 127, 130, 132, 140, 142, 145, 146, 150, 151, 165, 178, 182, 189, 190, 191, 194, 197, 198, 199, 203, 206, 207, 209, 211, 212, 213, 216, 222, 223, 231, 249, 250, 251, 252, 268, 269, 270, 271, 272, 273, 276, 280, 283, 285, 288, 289, 293, 294, 305, 306, 307, 308, 309, 313, 314, 315, 323, 324, 327, 329, 330, 336, 342, 349, 350, 356, 357, 358, 361, 362, 363, 364, 366, 367, 368, 369, 371, 373, 374, 377, 378, 380, (383).

A list of all forms known to be false in the model: 1, 2, 3, 4, 7, 8, 9, 14, 15, 39, 40, 41, 43, 44, 49, 50, 51, 66, 67, 71, 76, 77, 82, 84, 86, 87, 89, 90, 95, 99, 100, 101, 106, 109, 112, 113, 114, 118, 126, 129, 131, 133, 144, 149, 152, 161, 168, 174, 179, 181, 185, 188, 192, 193, 202, 214, 218, 255, 256, 257, 258, 259, 260, 261, 262, 264, 286, 292, 303, 316, 317, 328, 333, 334, 335, 345, 347, 359, (383).)

\noindent $\mathcal{N}6$: Levy's Model I. $A=\{a_{n}:n\in \omega\}$ and $A=\bigcup\{P_{n}:n\in\omega\}$, where $P_{0}=\{a_{0}\},\, P_{1}=\{a_{1},\,a_{2}\},\,P_{2}=\{a_{3},\,a_{4},\,a_{5}\},\,P_{3}=\{a_{6},\,a_{7},\,a_{8},\,a_{9},\, a_{10}\},\,\cdots$; in general for $n>0,\,|P_{n}|=p_{n}$, where $p_{n}$ is the $n$th prime. $\mathcal{G}$ is the group generated by $\{\pi_{n}:n\in\omega\}$, where, if $P_{n}=\{a_{m+1},\,a_{m+2},\,\cdots,\, a_{m+p_{n}}\}$, then

\begin{equation*}
\pi_{n}:a_{m+1}\mapsto a_{m+2}\mapsto\cdots a_{m+p_{n}}\mapsto a_{m+1}\ \mathrm{and}\ \pi_{n}(x)=x,\ \mathrm{for\ all}\,x\not\in P_{n}.
\end{equation*}
$S$ is the set of all finite subsets of $A$. Levi proves in $\mathcal{N}6$ that there is no choice function on $\{P_{n}:n\in\omega\}$ so $C(\aleph_{0},\,<\aleph_{0})$ (10) is false, but for each $n\in\omega, n>0,\, C(\infty,\,n)$ (61) is true. (It is also clear that there is no Kinna/Wagner selection function on $\{P_{n}:n\in\omega\}$ so $KW(\aleph_{0},\,<\aleph_{0})$ (358) is also false.) Levy\index{Author}{Levy}, also proves the axiom of Multiple Choice (67) is true. However, Bleicher has shown that $(\forall n\in \omega)MC(\infty,\, \infty$, relatively prime to $n$) (218) is equivalent to 61 + 67, so 218 is also true. Since 218 implies 333 $(MC(\infty,\,\infty,\mathrm{odd}))$ and Keremedis has shown that 333 + 334$(MC(\infty,\,\infty,\mathrm{even}))\leftrightarrow$ AC, it follows that 334 is false. Since 218 is true, and therefore, [218 A] (Existence of Complementary Subspaces) is also true, it follows that 95$(F)$ (Existence of Complementary Subspaces over a Field $F$) is true.

Shannon proves if $T=\{f:\exists nf$ is a choice function on $\{P_{0},\,P_{1},\,\cdots,\,P_{n}\}$\} with the partial order $f\leq g$ iff $g\subseteq f$, then $T$ is a denumerable union of finite sets, all antichains are finite, and there is a denumerable family of dense sets for which there in no generic filter (171 is false). See note 47 for definitions. In Howard/Yorke\index{Author}{Howard/Yorke} \cite{bib:1987} it is shown that for any prime $p$, there is a set of finite groups $\{G_{y}:y\in Y\}$ such that the weak direct product has no maximal $p$-subgroup $(308(p)$ is false in $\mathcal{N}6$ for any prime $p$). Since form 106 (Baire category Theorem for compact Hausdorff spaces) is true (218 implies 106) and 43 (Principle of Dependent Choices) is false (43 implies 171), form 154 (the Tychonoff theorem for countably many $T_{2}$ spaces.) must be false. (106 + 154 implies 43. See Brunner\index{Author}{Brunner} \cite{bib:1983c}) Degen has shown that if $\rho_{i}=\pi_{i}/P_{i},\,\rho_{i}$ is not defined on $A-P_{i}$, and $\phi =\bigcup_{i\in\omega}\rho_{i}$, then $\phi$ cannot be expressed as the product of two reflections. ($\phi$ is a reflection if $\phi^{2}=\phi$.) Consequently, form 314 (Every permutation on a non-empty set can be expressed as a product of two reflections.) is false.

\noindent $\mathcal{N}6\models$ 6, 11, 37, 61, 67, 73, 92, 95, 130, 191, 218, 273, 305, 313, 361, and 363, but 10, 154, 164, 171, $308(p)$, 314, 334, 344, and 358, are false. References Bleicher \cite{bib:1965}, Brunner \cite{bib:1984b}, Degen\index{Author}{Degen} \cite{bib:1988}, Howard/Yorke\index{Author}{Howard/Yorke} \cite{bib:1987}, Jech\index{Author}{Jech} \cite{bib:1973b} (Theorem~7.11 and prob 7.15), Keremedis\index{Author}{Keremedis} \cite{bib:1996a}, Levy\index{Author}{Levy} \cite{bib:1962}, H. Rubin/J. Rubin\index{Author}{H. Rubin/J. Rubin} \cite{bib:1985}, Shannon \cite{bib:1990}, notes 18, 35, 95, and 120(2 and 56).

(A minimal list of forms implying each form known to be true in the model: 6, 11, 37, 73, 92, 130, 191, 218, 305, 313, 361, 363.

A minimal list of forms at least\index{Author}{Bleicher} one of which is implied by each form known to be false in the model: 154, 164, 171, 308, 314, 334, 344, 358.

A list of all forms known to be true in the model: 0, 5, 6, 11, 12, 13, 18, 19, 25, 33, 34, 35, 36, 37, 38, 45, 46, 47, 48, 51, 52, 58, 61, 64, 67, 70, 73, 74, 76, 77, 78, 79, 80, 82, 83, 84, 88, 89, 90, 91, 92, 93, 94, 95, 104, 106, 108, 111, 112, 114, 118, 119, 120, 125, 126, 127, 130, 131, 140, 142, 144, 145, 147, 155, 173, 182, 185, 189, 190, 191, 194, 197, 199, 203, 206, 208, 211, 212, 218, 221, 222, 223, 232, 250, 251, 252, 254, 268, 269, 272, 273, 276, 280, 287, 288, 289, 292, 294, 305, 306, 307, 309, 313, 315, 316, 328, 329, 330, 333, 336, 342, 349, 350, 361, 362, 363, 364, 366, 367, 368, 369, 371, 373, 375, 381, 382, 383, (383).

A list of all forms known to be false in the model: 1, 2, 3, 4, 7, 8, 9, 10, 14, 15, 16, 17, 20, 21, 23, 27, 28, 29, 30, 31, 32, 39, 40, 41, 43, 44, 49, 50, 60, 62, 68, 71, 85, 86, 87, 100, 101, 107, 113, 121, 122, 123, 129, 132, 133, 150, 151, 152, 153, 154, 161, 164, 165, 168, 171, 174, 181, 192, 202, 213, 214, 231, 255, 256, 257, 258, 259, 260, 261, 262, 264, 270, 286, 295, 303, 308, 314, 317, 322, 323, 324, 325, 327, 331, 332, 334, 343, 344, 345, 346, 347, 352, 354, 355, 358, 359, 376, 377, 378, 380, (383).)

\noindent $\mathcal{N}7$: L\"{a}uchli's Model I. $A$ is countably infinite. The group $\mathcal{G}$ is a permutation group on $A$ with the following properties:
\begin{enumerate}
\item[1.] If $E_{1}$ and $E_{2}$ are finite subsets of $A$ and if $\phi\in \mathcal{G}$ is the identity on $E_{1}\cap E_{2}$, then there exist $\pi_{i}$ and $\rho_{i}$ in $\mathcal{G},\,i=1, 2,\cdots n$, such that $\phi=\pi_{1}\cdot \rho_{1}\cdot\pi_{2}\cdot \rho_{2}\cdots\pi_{n}\cdot \rho_{n}$ where $\pi_{i}$ is the identity on $E_{1}$ and $\rho_{i}$ is the identity on $E_{2}$.
\item[2.] If $E$ is a finite subset of $A$ and $\phi\in \mathcal{G}$ is such that $\phi(E)=E$, then $\phi$ is the identity on $E$.
\item[3.] For each finite $E\subset A$, there are distinct elements $a,\,b,\,c\in A$ and $\phi,\,\rho \in \mathcal{G}$ which are the identity on $E$ such that $\phi:a\mapsto b\mapsto c$ and $\rho:c\mapsto a\mapsto b$.
\end{enumerate}
$S$ is the set of all finite subsets of $A$. L\"{a}uchli proves that in this model, AC for a set of finite sets (62) and AC for a set of well orderable sets (60) are true, but the Ordering Principle (30) is false. Since 60 implies $UT(WO,\,WO,\,WO)$ (231) and in every FM model 231 implies $(\forall\alpha)UT(\aleph_{\alpha},\,\aleph_{\alpha},\,\aleph_{\alpha})$, form 23 is true. Also, 60 implies that there is no amorphous set (64 is true), and 64 implies 133 (Every set is either well orderable or has an infinite amorphous subset.) is false. Pincus proves that there is a function which associates each infinite set with a nonprincipal ideal in its power set (123 is true). Moreover, Rado's Lemma (99) is false in this model because BPI (14) is false (14 implies 30) and 99 + 62 implies 14 (note 33). Since 60 is true, AC for a well ordered set of sets (40) is false. (See Howard/Rubin\index{Author}{Howard/Rubin} \cite{bib:1995b}.) Form 40 is false and form 122 $(C(WO,\,<\infty))$ is true because 60 implies 122. Therefore, form 328 $(MC(WO,\,\infty))$ is false because 122 + 328 $\rightarrow$ 40. Since form 165 $(C(WO,\,WO))$ is true, (60 implies 165) it follows from note 2(8 and 9) that 16 and 24 are also true.

\noindent $\mathcal{N}7\models$ 11, 16, 23, 24, 37, 60, 73, 91, 92, 123, 130, 165, 191, 273, 305, 309, 313, 361, 363, 368, and 369, but 30, 40, 99, 133, 152, 192, and 328 are false. References L\"{a}uchli \cite{bib:1964}, Jech\index{Author}{Jech} \cite{bib:1973b} (p 104, Thm 7.9), Pincus\index{Author}{Pincus} \cite{bib:1976}, Howard/Rubin\index{Author}{Howard/Rubin} \cite{bib:1995b}, notes 2(8,9), 18, 33, and 120(34 and 47).

(A minimal list of forms implying each form known to be true in the model: 11, 16, 23, 24, 60, 73, 91, 92, 123, 130, 165, 191, 305, 309, 313, 361, 363.

A minimal list of forms at least one of which is implied by each form known to be false in the model: 30, 99, 133, 152, 192, 328.

A list of all forms known to be true in the model: 0, 5, 6, 10, 11, 12, 13, 16, 18, 19, 23, 24, 25, 26, 27, 28, 31, 32, 33, 34, 35, 36, 37, 38, 45, 46, 47, 48, 60, 61, 62, 63, 64, 70, 73, 74, 79, 80, 85, 88, 91, 92, 93, 94, 102, 104, 108, 111, 119, 120, 121, 122, 123, 127, 130, 132, 140, 142, 145, 146, 150, 151, 165, 178, 182, 189, 190, 191, 194, 197, 199, 203, 206, 207, 209, 211, 212, 213, 216, 221, 222, 223, 231, 249, 250, 251, 252, 268, 269, 270, 271, 272, 273, 276, 280, 283, 285, 288, 289, 294, 305, 306, 307, 308, 309, 313, 314, 315, 323, 324, 327, 329, 330, 331, 332, 336, 342, 343, 344, 349, 350, 356, 357, 358, 361, 362, 363, 364, 366, 367, 368, 369, 371, 373, 374, 378, 380, (383).

A list of all forms known to be false in the model: 1, 14, 15, 30, 40, 49, 50, 66, 67, 95, 99, 101, 109, 133, 149, 152, 192, 202, 214, 218, 255, 256, 257, 258, 259, 260, 261, 262, 264, 286, 295, 303, 317, 328, 333, 334, 335, 345, 347, 359, (383).)

\noindent $\mathcal{N}8$: L\"{a}uchli's Model II. $A$ is countably infinite, ordered like the rational numbers; $\mathcal{G}$ is the group of all order automorphisms of $A$; and $S$ is the set of all subsets $E$ of $A$ such that $E$ has at most a finite number of accumulation points and every infinite subset of $E$ has an accumulation point. L\"{a}uchli\index{Author}{Lauchli@L\"{a}uchli} shows that Urysohn's Lemma (78) is false in this model.

\noindent $\mathcal{N}8\models$ 6, 11, 37, 73, 91, 92, 130, 191, 273, 305, 309, 313, 361, 363, 368, and 369, but 15 and 78 are false. References L\"{a}uchli\index{Author}{Lauchli@L\"{a}uchli} \cite{bib:1962}, Jech/Sochor\index{Author}{Jech/Sochor} \cite{bib:1966a}, and note 18.

(A minimal list of forms implying each form known to be true in the model: 6, 11, 37, 73, 91, 92, 130, 191, 305, 309, 313, 361, 363.

A minimal list of forms at least one of which is implied by each form known to be false in the model: 15, 78.

A list of all forms known to be true in the model: 0, 5, 6, 11, 12, 13, 35, 36, 37, 38, 70, 73, 74, 79, 91, 92, 93, 94, 108, 130, 142, 145, 182, 189, 190, 191, 194, 197, 199, 203, 206, 211, 212, 222, 223, 251, 252, 272, 273, 280, 289, 305, 306, 307, 309, 313, 361, 362, 363, 364, 366, 367, 368, 369, 371, (383).

A list of all forms known to be false in the model: 1, 15, 40, 43, 44, 66, 67, 78, 87, 95, 101, 106, 109, 149, 174, 188, 192, 193, 202, 218, 255, 256, 257, 258, 259, 260, 261, 262, 264, 286, 333, 334, 335, 345, 347, 359, 375, (383).)

\noindent $\mathcal{N}9$: Halpern/Howard Model. $A$ is a set of atoms with the structure of the set $\{s:s:\omega\rightarrow \omega\wedge(\exists n)(\forall j>n)(s_{j}=0)\}$. We identify $A$ with this set to simplify the description of $\mathcal{G}$. $\mathcal{G}$ is the group of all permutations $\phi$ on $A$ such that $\{a\in A$: $\phi(a)\neq a\}$ is bounded; and $S=\{A_{0}^{n}:n\in \omega\}$, where $A_{0}^{n}$ is the set of all sequences $s$ such that $s_{j}=0$ for all $j\geq n$. (The set $x$ has support $A_{0}^{n}$ if all permutations in $\mathcal{G}$ that leave $A_{0}^{n}$ pointwise fixed and leave the first $n$ coordinates of each $s\in A$ pointwise fixed, also leave $x$ fixed.) In this model, the axiom of choice for 2-element sets (88) is false, but the $2m=m$ principle (3) is true. Form 3 implies that there are no amorphous sets (64) is true and 64 implies that 133 (Every set is either well orderable or has an infinite amorphous subset.) is false. Since 3 implies 9 (Every Dedekind finite set is finite.) and 70 (There is a non-trivial ultrafilter on $\omega$.) is true in every FM model, it follows that 63 (Every infinite set has a non-trivial ultrafilter.) is also true.

\noindent $\mathcal{N}9\models$ 3, 6, 11, 37, 63, 73, 91, 92, 130, 191, 273, 305, 309, 313, 361, 363, 368, and 369, but 88, 133, 192, and 344 are false. References Halpern/Howard\index{Author}{Halpern/Howard} \cite{bib:1974}, notes 18 and 120(56).

(A minimal list of forms implying each form known to be true in the model: 3, 6, 11, 37, 63, 73, 91, 92, 130, 191, 305, 309, 313, 361, 363.

A minimal list of forms at least one of which is implied by each form known to be false in the model: 88, 133, 192, 344.

A list of all forms known to be true in the model: 0, 3, 4, 5, 6, 9, 10, 11, 12, 13, 17, 18, 35, 36, 37, 38, 53, 57, 63, 64, 70, 73, 74, 79, 80, 82, 83, 84, 91, 92, 93, 94, 98, 108, 124, 127, 128, 130, 132, 142, 145, 182, 185, 186, 189, 190, 191, 194, 197, 198, 199, 203, 204, 206, 211, 212, 216, 217, 221, 222, 223, 249, 251, 252, 272, 273, 280, 288, 289, 296, 304, 305, 306, 307, 309, 313, 325, 336, 342, 358, 361, 362, 363, 364, 366, 367, 368, 369, 371, 373, 374, 376, 377, 378, (383).

A list of all forms known to be false in the model: 1, 14, 15, 30, 49, 50, 60, 61, 62, 68, 85, 88, 95, 107, 109, 123, 133, 192, 201, 213, 218, 270, 284, 295, 303, 317, 323, 326, 331, 332, 333, 335, 343, 344, 345, 346, (383).)

\noindent $\mathcal{N}10$: H\"{o}ft/Howard\index{Author}{Hoft/Howard@H\"{o}ft/Howard}/Mostowski Model. (The model is a variation of $\mathcal{N}3$.) $A$ as ordered by $\precsim$ has the same order type as the rationals; $\mathcal{G}$ is the group of all order automorphisms of $A$; $S$ is the set of all subsets $E$ of $A$ that satisfy the following three conditions:
\begin{enumerate}
\item[1.] $E$ is well ordered by $\precsim$.
\item[2.] $E$ is bounded in $A$.
\item[3.] If $b:\alpha\rightarrow E$ is an order preserving bijection from $\alpha$ onto $E$ and if $\lambda<\alpha$ is a limit ordinal then $\{b(\beta):\beta<\gamma\}$ has no least upper bound in $(A,\,\precsim)$.
\end{enumerate}
In H\"{o}ft/Howard \cite{bib:1994} it is shown that, in $\mathcal{N}10$, every Dedekind finite set is finite (9 is true), but $(A,\,\precsim)$ is a linearly ordered set with no infinite descending sequences that cannot be well ordered (77 is false). In addition (see note 113) there is an infinite set of pairs with no choice function (88 is false). Since 70 (There is a nontrivial ultrafilter on $\omega$.) is true in every FM model and 9 + 70 implies 63 (Every infinite set has a non-trivial ultrafilter.), 63 is also true.

\noindent $\mathcal{N}10 \models$ 6, 9, 11, 37, 63, 73, 91, 92, 130, 191, 273, 305, 309, 313, 361, 363, 368, and 369, but 77, 88, 106 and 344 are false. References H\"{o}ft/Howard \cite{bib:1994}, notes 18, 113, and 120(55 and 56).

(A minimal list of forms implying each form known to be true in the model: 6, 9, 11, 37, 63, 73, 91, 92, 130, 191, 305, 309, 313, 361, 363.

A minimal list of forms at least one of which is implied by each form known to be false in the model: 77, 88, 106, 344.

A list of all forms known to be true in the model: 0, 5, 6, 9, 10, 11, 12, 13, 17, 18, 35, 36, 37, 38, 57, 63, 64, 70, 73, 74, 79, 80, 82, 83, 84, 91, 92, 93, 94, 98, 108, 124, 127, 128, 130, 132, 142, 145, 182, 185, 189, 190, 191, 194, 197, 198, 199, 203, 206, 211, 212, 216, 217, 221, 222, 223, 249, 251, 252, 272, 273, 280, 288, 289, 296, 304, 305, 306, 307, 309, 313, 325, 336, 342, 358, 361, 362, 363, 364, 366, 367, 368, 369, 371, 373, 374, 376, 377, 378, (383).

A list of all forms known to be false in the model: 1, 14, 15, 30, 40, 43, 44, 49, 50, 51, 60, 61, 62, 66, 67, 68, 77, 85, 87, 88, 89, 90, 95, 101, 106, 107, 109, 112, 114, 123, 133, 149, 174, 188, 192, 193, 201, 202, 213, 218, 255, 256, 257, 258, 259, 260, 261, 262, 264, 270, 284, 286, 292, 295, 303, 316, 317, 323, 326, 331, 332, 333, 334, 335, 343, 344, 345, 346, 347, 359, (383).)

\noindent $\mathcal{N}11$: Jech's Model. Let $(I,\,\precsim)$ be a partially ordered set in the kernel (in the base model without atoms). Let $A=\{a_{i,n}:i\in I,\,n\in \omega\}$. For each $p\subseteq I$, let $S_{p}=\{a_{i,n}:i\in p,\, n\in \omega\}$. $\mathcal{G}$ is the group of all permutations of $A$ that leave $S_{\{i\}}$ fixed, for each $i\in I$. $S$ is the set of all finite subsets of $A$. In this model, there is a set of cardinals that is isomorphic to $(I,\,\precsim)$. Thus, form $59(\precsim)$ is false.

\noindent $\mathcal{N}11\models$ 6, 11, 37, 73, 91, 92, 130, 191, 273, 305, 309, 313, 361, 363, 368, and 369, but 15, 59$(_{\sim}^{\prec})$, and 192 are false. References Jech\index{Author}{Jech} \cite{bib:1966c} and note 18.

(A minimal list of forms implying each form known to be true in the model: 6, 11, 37, 73, 91, 92, 130, 191, 305, 309, 313, 361, 363.

A minimal list of forms at least one of which is implied by each form known to be false in the model: 15, 59, 192.

A list of all forms known to be true in the model: 0, 5, 6, 11, 12, 13, 35, 36, 37, 38, 70, 73, 74, 79, 91, 92, 93, 94, 108, 130, 142, 145, 182, 189, 190, 191, 194, 197, 199, 203, 206, 211, 212, 222, 223, 251, 252, 272, 273, 280, 289, 305, 306, 307, 309, 313, 361, 362, 363, 364, 366, 367, 368, 369, 371, (383).

A list of all forms known to be false in the model: 1, 15, 59, 192, (383).)

\noindent $\mathcal{N}12(\aleph_{1})$: A variation of Fraenkel's model, $\mathcal{N}1$. The cardinality of $A$ is $\aleph_{1}$, $\mathcal{G}$ is the group of all permutations on $A$, and $S$ is the set of all countable subsets of $A$. In $\mathcal{N}12(\aleph_{1})$, every Dedekind finite set is finite (9 is true), but the $2m=m$ principle (3) is false. Also, AC for a well ordered set of sets (40) is true, but not all cardinals are comparable with $\aleph_{1}$ (41 is false) and for each $n\geq 2$ the set of $n$-element subsets of $A$ has no choice function (45$(n)$ is false). A. Rubin\index{Author}{A. Rubin} shows that each pair of cardinals has a greatest lower bound (187 is true). (See note 108.) Truss\index{Author}{Truss} proves that every linearly ordered set has a choice function (202 is true). Since 202 implies $UT(WO,\,WO,\,WO)$ (231) and in every FM model 231 implies $(\forall\alpha)UT(\aleph_{\alpha},\,\aleph_{\alpha},\,\aleph_{\alpha})$, form 23 is true. Since 202 implies 8 $(C(\aleph_{0},\,\infty))$, it follows from Brunner\index{Author}{Brunner} \cite{bib:1982a} that in this model there is a set that cannot be well ordered and does not have an infinite Dedekind finite subset, (163 is false). (Form 8 plus form 163 iff AC.) Since 70 (There is a non-trivial ultrafilter on $\omega$.) is true in every FM model and 9 + 70 implies 63 (Every infinite set has a non-trivial ultra filter.), 63 is also true. Blass\index{Author}{Blass} \cite{bib:1979} has shown that 191 (There is a set $X$ such that for each set $a$ there is an ordinal $\alpha$ and a function $f$ mapping $X\times\alpha$ onto $a$.) is true in every permutation model with a set of atoms. He has also shown that 8 $(C(\aleph_{0},\,\infty))+191+193$ (Every Abelian group is a homomorphic image of a free projective Abelian group.) iff AC. Since 202 implies 8, form 193 must be false. (See note 60 for definitions.) The set of all pairs of atoms is a set of well orderable sets that has no choice function. Thus, $C(\infty,\,2)$ (88) and $KW(\infty,\,WO)$ (323) are false. Since $C(WO,\,<\aleph_{0})$ (122) is true (202 implies 122) it follows that MC (67) is false because 67 + 122 implies AC. (See [67 B].) It is shown in Higasikawa\index{Author}{Higasikawa} \cite{bib:1995} that the Idemmultiple Partition Principle, form 347, is false.

\noindent $\mathcal{N}12(\aleph_{1})\models$ 9, 11, 23, 40, 63, 130, 165, 187, 191, 202, 305, 309, and 313, but 3, 41, 45$(n)$, 67, 88, 133, 152, 163, 192, 193, 323, 344, and 347 are false. See references for $\mathcal{N}1$, A. Rubin \cite{bib:1978}, Higasikawa \cite{bib:1995}, Jech \cite{bib:1973b} (pp 127 and 131), Truss \cite{bib:1978}, Brunner \cite{bib:1982a}, Blass \cite{bib:1979}, notes 18, 60, 108, and 120(4, 34, 46, and 56).

(A minimal list of forms implying each form known to be true in the model: 11, 23, 63, 130, 165, 187, 191, 202, 305, 309, 313.

A minimal list of forms at least one of which is implied by each form known to be false in the model: 3, 41, 45, 67, 88, 133, 152, 163, 192, 193, 344, 347.

A list of all forms known to be true in the model: 0, 5, 6, 8, 9, 10, 11, 12, 13, 16, 17, 18, 19, 23, 24, 25, 26, 27, 28, 29, 31, 32, 33, 34, 35, 36, 37, 38, 39, 40, 42, 43, 47, 48, 57, 58, 63, 64, 70, 73, 74, 75, 77, 78, 79, 80, 82, 83, 84, 86, 91, 92, 93, 94, 98, 104, 106, 108, 111, 113, 119, 120, 121, 122, 124, 126, 127, 128, 130, 131, 132, 142, 145, 150, 151, 154, 155, 165, 181, 182, 185, 186, 187, 189, 190, 191, 194, 196, 197, 198, 199, 202, 203, 206, 207, 208, 209, 211, 212, 216, 217, 221, 222, 223, 231, 243, 249, 250, 251, 252, 260, 272, 273, 279, 280, 282, 288, 289, 294, 296, 304, 305, 306, 307, 309, 313, 315, 322, 324, 325, 327, 328, 330, 336, 337, 338, 339, 340, 341, 342, 350, 351, 352, 353, 354, 355, 357, 358, 360, 361, 362, 363, 364, 366, 367, 368, 369, 371, 373, 374, 375, 376, 377, 378, 380, (383).

A list of all forms known to be false in the model: 1, 2, 3, 14, 15, 30, 41, 44, 45, 49, 50, 60, 61, 62, 66, 67, 68, 85, 88, 95, 101, 107, 109, 123, 133, 149, 152, 163, 192, 193, 201, 213, 214, 218, 270, 284, 295, 303, 317, 323, 326, 331, 332, 333, 334, 335, 343, 344, 345, 346, 347, 359, (383).)

\noindent $\mathcal{N}12(\aleph_{2})$: Another variation of $\mathcal{N}1$. Change ``$\aleph_{1}$'' to $``\aleph_{2}$'' in $\mathcal{N}12(\aleph_{1})$ above. Thus, $A$ has cardinality $\aleph_{2}$, $\mathcal{G}$ is the group of all permutations of $A$, and $S$ is the set of all subsets of $A$ of cardinality at most $\aleph_{1}$. It follows from Jech\index{Author}{Jech} \cite{bib:1973b}(p 127, Theorem 8.9), that the Axiom of Choice for a well ordered family of sets (40) is true, the Principle of Dependent Choices for $\omega_{1}$ sets (44) is true, but the $2m=m$ principle (3) is false. Since 40 implies $UT(WO,\,WO,\,WO)$ (231) and in every FM model 231 implies $(\forall\alpha)UT(\aleph_{\alpha},\, \aleph_{\alpha},\,\aleph_{\alpha})$, form 23 is true. H\"{a}ussler\index{Author}{Haussler@H\"{a}ussler} \cite{bib:1983} has shown that cardinal addition is not definable in terms of $\leq$ in this model. Thus, 161 is false. Since 40 implies 8 $(C(\aleph_{0},\,\infty))$, it follows from Brunner\index{Author}{Brunner} \cite{bib:1982a} that in this model there is a set that cannot be well ordered and does not have an infinite Dedekind finite subset, (163 is false). (Form 8 plus form 163 iff AC.) Since 40 implies 9 (Every Dedekind finite set is finite.) and 70 (There is a non-trivial ultrafilter on $\omega$.) is true in every FM model, it follows that 63 (Every infinite set has a non-trivial ultrafilter.) is also true. Form 44 implies forms 39 $(C(\aleph_{1},\,\infty))$ and 41 (For each cardinal $m,\,m\leq\aleph_{1}$ or $\aleph_{1}\leq m$.). van Douwen\index{Author}{Douwen} \cite{bib:1980} has shown that 39 + 41 implies 135 (If $X$ is a $T_{2}$ space with at least two points and $X^{Y}$ is hereditarily metacompact, then $Y$ is countable.), so form 135 is also true. Blass\index{Author}{Blass} \cite{bib:1979} has shown that 191 (There is a set $X$ such that for each set $a$ there is an ordinal $\alpha$ and a function $f$ mapping $X\times\alpha$ onto $a$.) is true in every permutation model with a set of atoms. He has also shown that 8 $(C(\aleph_{0},\,\infty))+191+193$ (Every Abelian group is a homomorphic image of a free projective Abelian group.) iff AC. Since 40 implies 8, form 193 must be false. (See note 60 for definitions.) It is shown in Howard/Rubin\index{Author}{Howard/Rubin} \cite{bib:1995b} that in every FM model $40 + 60\ (C(\infty,\,WO))$ implies AC, so 60 is false. Howard\index{Author}{Howard} \cite{bib:1973} proves that in every FM model 62 $(C(\infty,\,<\aleph_{0}))$ implies 60, so 62 is also false. Since $C(WO,\,<\aleph_{0})$ (122) is true (40 implies 122) it follows that MC (67) is false because $67 + 122$ implies AC. (See [67 B].) It is shown in Higasikawa\index{Author}{Higasikawa} \cite{bib:1995} that the Idemmultiple Partition Principle, form 347, is false.

\noindent $\mathcal{N}12(\aleph_{2})\models$ 11, 23, 40, 44, 63, 91, 130, 135, 165, 191, 305, 309, 313, 368, and 369, but 60, 62, 67, 133, 152, 161, 163, 192, 193, 344, and 347 are false. References Jech \cite{bib:1973b} (p 127), H\"{a}ussler \cite{bib:1983} (p 197), Blass \cite{bib:1979}, Brunner \cite{bib:1982a}, van Douwen\index{Author}{Douwen} \cite{bib:1980}, Higasikawa\index{Author}{Higasikawa} \cite{bib:1995}, Howard\index{Author}{Howard} \cite{bib:1973}, Howard/Rubin\index{Author}{Howard/Rubin} \cite{bib:1995b}, notes 18, 60, and 120(4, 46, and 56).

(A minimal list of forms implying each form known to be true in the model: 11, 23, 40, 44, 63, 91, 130, 135, 165, 191, 305, 309, 313.

A minimal list of forms at least one of which is implied by each form known to be false in the model: 62, 67, 133, 152, 161, 163, 192, 193, 344, 347.

A list of all forms known to be true\index{Author}{Haussler@H\"{a}ussler} in the model: 0, 5, 6, 8, 9, 10, 11, 12, 13, 16, 17, 18, 19, 23, 24, 25, 26, 27, 28, 29, 31, 32, 34, 35, 36, 37, 38, 39, 40, 41, 42, 43, 44, 47, 48, 57, 58, 63, 64, 70, 73, 74, 75, 77, 78, 79, 80, 82, 83, 84, 86, 91, 92, 93, 94, 98, 104, 106, 108, 111, 113, 119, 122, 124, 126, 127, 128, 130, 131, 132, 135, 142, 145, 150, 151, 154, 155, 165, 182, 185, 189, 190, 191, 194, 196, 197, 198, 199, 203, 206, 207, 208, 209, 211, 212, 216, 217, 221, 222, 223, 231, 243, 249, 250, 251, 252, 272, 273, 279, 280, 282, 288, 289, 294, 296, 304, 305, 306, 307, 309, 313, 315, 322, 324, 325, 327, 328, 330, 336, 337, 338, 339, 340, 341, 342, 350, 351, 352, 353, 354, 355, 357, 358, 360, 361, 362, 363, 364, 365, 366, 367, 368, 369, 371, 373, 374, 375, 376, 377, 378, 380, (383).

A list of all forms known to be false in the model: 1, 14, 15, 30, 49, 50, 60, 62, 66, 67, 68, 85, 95, 101, 107, 109, 123, 133, 149, 152, 161, 163, 192, 193, 213, 214, 218, 270, 295, 303, 317, 323, 331, 332, 333, 334, 335, 343, 344, 345, 346, 347, 359, (383).)

\noindent $\mathcal{N}12(\aleph_{\alpha})$: A generalization of $\mathcal{N}12(\aleph_{1})$. Replace ``$\aleph_{1}$'' by ``$\aleph_{\alpha}$'' where $\aleph_{\alpha}$ is a singular cardinal. Thus, $|A|=\aleph_{\alpha}$; $\mathcal{G}$ is the group of all permutations on $A$; and $S$ is the set of all subsets of $A$ with cardinality less than $\aleph_{\alpha}$. Jech\index{Author}{Jech} shows that $DC(\aleph_{\beta})$ holds for all $\aleph_{\beta}<\mathrm{cf}(\aleph_{\alpha})$ ($87(\beta)$ for $\aleph_{\beta}<\mathrm{cf}(\aleph_{\alpha})$ is true); $W_{\aleph_{\beta}}$ (For all cardinals $m,\, m\leq\aleph_{\beta}$ or $\aleph_{\beta}\leq m.)$ holds for all $\beta<\alpha$ ($71(\beta)$ is true for $\beta<\alpha$); but $71(\alpha)$ and $C(\infty,\mathrm{cf}\aleph_{\alpha})$ are false. It follows from note 104 that the $2m=m$ principle (3) is false. Since 87 implies 8 $(C(\aleph_{0},\,\infty))$, it follows from Brunner\index{Author}{Brunner} \cite{bib:1982a} that in this model there is a set that cannot be well ordered and does not have an infinite Dedekind finite subset, which implies that 163 is false. (Form 8 plus form 163 iff AC.) Since 87$(\beta)$ implies 9 (Every Dedekind finite set is finite.) and 70 (There is a non-trivial ultrafilter on $\omega$.) is true in every FM model, it follows that 63 (Every infinite set has a non-trivial ultrafilter.) is also true. Blass\index{Author}{Blass} \cite{bib:1979} has shown that 191 (There is a set $X$ such that for each set $a$ there is an ordinal $\alpha$ and a function $f$ mapping $X\times\alpha$ onto $a$.) is true in every permutation model with a set of atoms. He has also shown that 8 + 191 + 193 (Every Abelian group is a homomorphic image of a free projective Abelian group.) iff AC. Thus, form 193 must be false if $\aleph_{\beta}<\mathrm{cf}(\aleph\alpha)$. (See note 60 for definitions.)

\noindent $\mathcal{N}12(\aleph_{\alpha})\models$ 11, 63, 71$(\alpha)$, $87(\beta)$ if $\aleph_{\beta}<\mathrm{cf}(\aleph_{\alpha})$, 91, 130, 191, 305, 309, 313, 368, and 369, but 3, 15, 133, 163, 192, and 193 are false. References Jech \cite{bib:1973b} (p 124, Thm 8.6), Monro\index{Author}{Monro} \cite{bib:1974}, Blass \cite{bib:1979}, Brunner \cite{bib:1982a}, notes 18, 60, 104, and 120(4).

(A minimal list of forms implying each form known to be true in the model: 11, 63, 87, 91, 130, 191, 305, 309, 313.

A minimal list of forms at least one of which is implied by each form known to be false in the model: 3, 15, 133, 163, 192, 193.

A list of all forms known to be true in the model: 0, 5, 6, 8, 9, 10, 11, 12, 13, 16, 17, 18, 19, 24, 26, 27, 28, 29, 31, 32, 34, 35, 36, 37, 38, 42, 43, 57, 63, 64, 70, 71, 73, 74, 75, 77, 78, 79, 80, 82, 83, 84, 86, 87, 91, 92, 93, 94, 98, 104, 106, 108, 113, 119, 124, 126, 127, 128, 130, 131, 132, 142, 145, 150, 154, 155, 174, 182, 183, 185, 189, 190, 191, 194, 196, 197, 198, 199, 203, 206, 209, 211, 212, 216, 217, 221, 222, 223, 243, 249, 251, 252, 272, 273, 279, 280, 282, 288, 289, 296, 304, 305, 306, 307, 309, 313, 325, 336, 337, 338, 339, 340, 341, 342, 350, 351, 352, 353, 354, 355, 357, 358, 360, 361, 362, 363, 364, 366, 367, 368, 369, 371, 373, 374, 375, 376, 377, 378, 380, (383).

A list of all forms known to be false in the model: 1, 2, 3, 15, 133, 163, 192, 193, (383).)

\noindent $\mathcal{N}13$: L\"{a}uchli/Jech Model. $A=B_{1}\cup B_{2}$, where $B_{1}=\bigcup\{A_{j1}:j\in \omega\}$, and $B_{2}=\bigcup\{A_{j2}:j\in \omega\}$, and each $A_{ji}$ is a 6-element set. Consider the following permutations on 6 elements: $\alpha_{1}=(12)(34)(5)(6)$, $\beta_{1}=(13)(24)(5)(6)$, $\gamma_{1}=(14)(23)(5)(6)$, $\alpha_{2}=\alpha_{1}$, $\beta_{2}=(12)(3)(4)(56)$, and $\gamma_{2}=(1)(2)(34)(56)$. For each $j,\,\alpha$ is the permutation on $A_{j1}\cup A_{j2}$ that acts like $\alpha_{1}$ on $A_{j1}$ and $\alpha_{2}$ on $A_{j2}$; similarly for $\beta$ and $\gamma$. $\mathcal{G}$ is the group of all permutations on $A$ that preserves $A_{j1}$ and $A_{j2}$ and acts on $A_{j1}\cup A_{j2}$ like the identity, $\alpha,\,\beta$, or $\gamma$. $S$ is the set of all finite subsets of $A$. In $\mathcal{N}13$, there are two isomorphic vector spaces, one with $B_{1}$ as a basis and the other with $B_{2}$ as a basis, but $|B_{1}|\neq|B_{2}|$ because $B_{1}$ contains a countably infinite subset, but $B_{2}$ does not. And in L\"{a}uchli \cite{bib:1962} it is shown that $|B_{1}|$ and $|B_{2}|$ are incomparable.

\noindent $\mathcal{N}13\models$ 6, 11, 37, 73, 91, 92, 130, 191, 273, 305, 309, 313, 361, 363, 368, and 369, but 9, 15, and 235 are false. References L\"{a}uchli \cite{bib:1962}, Jech/Sochor\index{Author}{Jech/Sochor} \cite{bib:1966a}, Jech\index{Author}{Jech} \cite{bib:1973b} (p 149 prob 5), and note 18.

(A minimal list of forms implying each form known to be true in the model: 6, 11, 37, 73, 91, 92, 130, 191, 305, 309, 313, 361, 363.

A minimal list of forms at least one of which is implied by each form known to be false in the model: 9, 15, 235.

A list of all forms known to be true in the model: 0, 5, 6, 11, 12, 13, 35, 36, 37, 38, 70, 73, 74, 79, 91, 92, 93, 94, 108, 130, 142, 145, 182, 189, 190, 191, 194, 197, 199, 203, 206, 211, 212, 222, 223, 251, 252, 272, 273, 280, 289, 305, 306, 307, 309, 313, 361, 362, 363, 364, 366, 367, 368, 369, 371, (383).

A list of all forms known to be false in the model: 1, 2, 3, 4, 7, 8, 9, 14, 15, 39, 40, 41, 43, 44, 50, 71, 86, 87, 96, 100, 101, 107, 113, 129, 152, 161, 168, 174, 181, 192, 202, 214, 235, 255, 256, 257, 258, 259, 260, 261, 262, 264, 286, 303, 317, 345, 347, 359, (383).)

\noindent $\mathcal{N}14$: Morris/Jech Model. $A=\bigcup\{A_{\alpha}:\alpha<\omega_{1}\}$, where the $A_{\alpha}$'s are pairwise disjoint, each is countably infinite, and each is ordered like the rationals; $\mathcal{G}$ is the group of all permutations on $A$ that leave each $A_{\alpha}$ fixed and preserve the ordering on each $A_{\alpha}$; and $S=\{B_{\gamma}:\gamma<\omega_{1}\}$, where $B_{\gamma}=\bigcup\{A_{\alpha}:\alpha<\gamma\}$. Let $P=A\cup\omega_{1}$ and order $P$ so that the elements in $A_{\alpha}$ precede those in $A_{\beta}$ if $\alpha<\beta$, the elements in each $A_{\alpha}$ are ordered like the rationals as above, and each $\alpha\in \omega_{1}$ precedes $A_{\alpha}$ and follows all $A_{\beta}$ for $\beta<\alpha$. Then $P$ is a linearly ordered set such that $|P|>\aleph_{1}$, but each initial segment of $P$ is countable. Thus, form 103 is false. In note 144 we show dependent choice (form 43) is true.

\noindent $\mathcal{N}14\models$ 6, 11, 37, 43, 63, 73, 91, 92, 130, 191, 273, 305, 309, 313, 361, 363, 368, and 369, but 15, 103, 133, 163, and 192 are false. References Jech \cite{bib:1973b} (p 164 prob 21), notes 18, 144, and 120(1, 4, and 20).

(A minimal list of forms implying each form known to be true in the model: 11, 43, 63, 91, 130, 191, 305, 309, 313.

A minimal list of forms at least one of which is implied by each form known to be false in the model: 15, 103, 133, 163, 192.

A list of all forms known to be true in the model: 0, 5, 6, 8, 9, 10, 11, 12, 13, 16, 17, 18, 19, 24, 26, 27, 28, 29, 31, 32, 34, 35, 36, 37, 38, 42, 43, 57, 63, 64, 70, 73, 74, 75, 77, 78, 79, 80, 82, 83, 84, 91, 92, 93, 94, 98, 104, 106, 108, 113, 119, 124, 126, 127, 128, 130, 131, 132, 142, 145, 150, 154, 155, 182, 185, 189, 190, 191, 194, 197, 198, 199, 203, 206, 209, 211, 212, 216, 217, 221, 222, 223, 243, 249, 251, 252, 272, 273, 279, 280, 282, 288, 289, 296, 304, 305, 306, 307, 309, 313, 325, 336, 337, 338, 339, 340, 341, 342, 350, 351, 352, 353, 354, 355, 357, 358, 360, 361, 362, 363, 364, 366, 367, 368, 369, 371, 373, 374, 375, 376, 377, 378, 380, (383).

A list of all forms known to be false in the model: 1, 15, 103, 133, 163, 192, (383).)

\noindent $\mathcal{N}15$: Brunner/Howard\index{Author}{Howard} Model I. $A=\{a_{i,\alpha}:i\in \omega\wedge\alpha\in \omega_{1}\}$. Let $\Gamma$ be the group of all even permutations on $\omega$. $\mathcal{G}=\{g:(\forall\alpha\in\omega_{1})(\exists\gamma\in\Gamma)(\forall i\in \omega)g(a_{i,\alpha})=a_{\gamma(i),\alpha}\}$. $S$ is the set of all countable subsets of $A$. In this model, $|A|<|A|+|A|$ (3 is false). The axiom of choice for a denumerable set of sets (8) is true and it is also true that every infinite set is the union of a disjoint well orderable family of denumerable sets (152 is true). If $B_{\alpha}=\{a_{i,\alpha}:i\in \omega\}$, then $\bigcup\{B_{\alpha}:\alpha\in \omega_{1}\}$ cannot be well ordered (151 is false), but the union of a well ordered family of finite sets can be well ordered ([122 A] is true). $C(\infty,\,WO)$ (60) implies 151 and in every FM model $C(\infty,\,<\aleph_{0})$ (62) implies 60 (Howard \cite{bib:1973}). Thus, 62 is false. It follows from Brunner\index{Author}{Brunner} \cite{bib:1982a} that in this model there is a set that cannot be well ordered and does not have an infinite Dedekind finite subset, (163 is false). (Form 8 plus form 163 iff AC.) Since 8 implies 9 (Every Dedekind finite set is finite.) and 70 (There is a non-trivial ultrafilter on $\omega$.) is true in every FM model, it follows that 63 (every infinite set has a non-trivial ultrafilter) is also true. Blass\index{Author}{Blass} \cite{bib:1979} has shown that 191 (There is a set $X$ such that for each set $a$ there is an ordinal $\alpha$ and a function $f$ mapping $X\times\alpha$ onto $a$.) is true in every permutation model with a set of atoms. He has also shown that 8 + 191 + 193 (Every Abelian group is a homomorphic image of a free projective Abelian group.) iff AC. Thus, 193 is false. (See note 60 for definitions.) $KW(WO,\,WO)$ (324) and $MC(WO,\, WO)$ (330) are false because $\{B_{\alpha}:\alpha\in \omega_{1}\}$ is a well ordered set of well orderable sets that has no Kinna-Wagner selection function and no multiple choice function. In note 144 we show dependent choice (form 43) is true.

\noindent $\mathcal{N}15\models$ 8, 11, 43, 63, 91, 122, 130, 152, 191, 273, 305, 309, 313, 363, 368, and 369, but 3, 62, 67, 151, 163, 192, 193, 324, 330, and 344 are false. References Brunner/Howard\index{Author}{Brunner/Howard} \cite{bib:1992}, Brunner \cite{bib:1982a}, Blass \cite{bib:1979}, Howard \cite{bib:1973}, notes 18, 60, 120(32, 34, 41, and 56), and 144.

(A minimal list of forms implying each form known to be true in the model: 11, 43, 63, 91, 122, 130, 152, 191, 305, 309, 313.

A minimal list of forms at least one of which is implied by each form known to be false in the model: 3, 62, 151, 163, 192, 193, 324, 330, 344.

A list of all forms known to be true in the model: 0, 4, 5, 6, 8, 9, 10, 11, 12, 13, 16, 17, 18, 19, 24, 26, 27, 28, 29, 31, 32, 34, 35, 36, 37, 38, 42, 43, 47, 48, 57, 63, 64, 70, 73, 74, 75, 77, 78, 79, 80, 82, 83, 84, 91, 92, 93, 94, 98, 104, 106, 108, 111, 113, 119, 122, 124, 126, 127, 128, 130, 131, 132, 142, 145, 150, 152, 154, 155, 182, 185, 189, 190, 191, 194, 197, 198, 199, 203, 206, 209, 211, 212, 216, 217, 221, 222, 223, 243, 249, 250, 251, 252, 272, 273, 279, 280, 282, 288, 289, 296, 304, 305, 306, 307, 309, 313, 325, 327, 336, 337, 338, 339, 340, 341, 342, 350, 351, 352, 353, 354, 355, 357, 358, 360, 361, 362, 363, 364, 366, 367, 368, 369, 371, 373, 374, 375, 376, 377, 378, 380, (383).

A list of all forms known to be false in the model: 1, 2, 3, 14, 15, 20, 21, 23, 30, 40, 49, 50, 60, 62, 66, 67, 68, 85, 95, 101, 107, 109, 123, 133, 149, 151, 163, 165, 192, 193, 202, 213, 218, 231, 255, 256, 257, 258, 259, 260, 261, 262, 264, 270, 286, 295, 303, 317, 322, 323, 324, 328, 329, 330, 331, 332, 333, 334, 335, 343, 344, 345, 346, 347, 359, (383).)

\noindent $\mathcal{N}16$: Jech/Levy/Pincus Model. $A$ has cardinality $\aleph_{\omega}$; $\mathcal{G}$ is the group of all permutations on $A$; and $S$ is the set of all subsets of $A$ of cardinality less that $\aleph_{\omega}$. In this model, every infinite set has a countably infinite subset (9 is true), but not every set of denumerable sets has a choice function (8 is false). Also, in $\mathcal{N}16$, for every cardinal $\kappa <\aleph_{\omega}$ every cardinal is comparable to $\kappa$ (41 is true), but there are cardinals that are not comparable to $\aleph_{\omega}$. Form 41 implies that there are no amorphous sets (64 is true), but 64 implies that 133 (Every set is either well orderable or has an infinite amorphous subset.) is false. Since 70 (There is a non-trivial ultrafilter on $\omega.)$ is true in every FM model and 9 + 70 implies 63 (Every infinite set has a non-trivial ultrafilter.), 63 is also true.

\noindent $\mathcal{N}16\models$ 6, 9, 11, 37, 41, 63, 73, 91, 92, 130, 191, 273, 305, 309, 313, 361, 363, 368, and 369, but 8, 15, 106, and 133 are false. References Levy\index{Author}{Levy} \cite{bib:1964}, Pincus\index{Author}{Pincus} \cite{bib:1969}, Jech\index{Author}{Jech} \cite{bib:1973b} (p 124), notes 18 and 120(55).

(A minimal list of forms implying each form known to be true in the model: 6, 11, 37, 41, 63, 73, 91, 92, 130, 191, 305, 309, 313, 361, 363.

A minimal list of forms at least one of which is implied by each form known to be false in the model: 8, 15, 106, 133.

A list of all forms known to be true in the model: 0, 5, 6, 9, 10, 11, 12, 13, 17, 18, 35, 36, 37, 38, 41, 57, 63, 64, 70, 73, 74, 79, 80, 82, 83, 84, 91, 92, 93, 94, 98, 108, 124, 127, 128, 130, 132, 142, 145, 182, 185, 189, 190, 191, 194, 197, 198, 199, 203, 206, 211, 212, 216, 217, 221, 222, 223, 249, 251, 252, 272, 273, 280, 288, 289, 296, 304, 305, 306, 307, 309, 313, 325, 336, 342, 358, 361, 362, 363, 364, 365, 366, 367, 368, 369, 371, 373, 374, 376, 377, 378, (383).

A list of all forms known to be false in the model: 1, 8, 15, 39, 40, 43, 44, 66, 67, 86, 87, 95, 101, 106, 109, 113, 133, 149, 174, 181, 188, 192, 193, 202, 218, 255, 256, 257, 258, 259, 260, 261, 262, 264, 286, 333, 334, 335, 345, 347, 359, (383).)

\noindent $\mathcal{N}16(\aleph_{\gamma})$: Levy's Model II. This is an extension of $\mathcal{N}16$ in which $A$ has cardinality $\aleph_{\gamma}$ where $\mathrm{cf}(\aleph_{\gamma})=\aleph_{0}$; $\mathcal{G}$ is the group of all permutations on $A$; and $S$ is the set of all subsets of $A$ of cardinality less that $\aleph_{\gamma}$. Just as in $\mathcal{N}16$ form 9 (Every Dedekind finite set is finite.) is true, but 8 $(C(\aleph_{0},\,\infty))$ is false. Also, for every cardinal $\aleph_{\alpha},\,\alpha<\gamma$, every cardinal is comparable to $\aleph_{\alpha}$ (71($\alpha$) is true), but there are cardinals that are not comparable to $\aleph_{\gamma}$ ($71(\gamma)$ is false). Form $71(\alpha)$ implies that there are no amorphous sets (64 is true) and 64 implies that 133 (Every set is either well orderable or has an infinite amorphous subset.) is false. Since 70 (There is a non-trivial ultrafilter on $\omega$.) is true in every FM model, and 9 + 70 implies 63 (Every infinite set has a non-trivial ultra filter.), it follows that 63 is also true.

\noindent $\mathcal{N}16(\aleph_{\gamma}\models$ 6, 9, 11, 37, 63, 71$(\alpha)$ for $\alpha<\gamma$, 73, 91, 92, 130, 191, 273, 305, 309, 313, 361, 363, 368, and 369, but 8, 15, $71(\gamma)$, 106, and 133 are false. References Levy \cite{bib:1964}, Pincus \cite{bib:1969}, Jech \cite{bib:1973b} (p 124), notes 18 and 120(55).

(A minimal list of forms implying each form known to be true in the model: 6, 11, 37, 63, 71, 73, 91, 92, 130, 191, 305, 309, 313, 361, 363.

A minimal list of forms at least one of which is implied by each form known to be false in the model: 8, 15, 71, 106, 133.

A list of all forms known to be true\index{Author}{Jech} in the model: 0, 5, 6, 9, 10, 11, 12, 13, 17, 18, 35, 36, 37, 38, 57, 63, 64, 70, 71, 73, 74, 79, 80, 82, 83, 84, 91, 92, 93, 94, 98, 108, 124, 127, 128, 130, 132, 142, 145, 182, 183, 185, 189, 190, 191, 194, 197, 198, 199, 203, 206, 211, 212, 216, 217, 221, 222, 223, 249, 251, 252, 272, 273, 280, 288, 289, 296, 304, 305, 306, 307, 309, 313, 325, 336, 342, 358, 361, 362, 363, 364, 366, 367, 368, 369, 371, 373, 374, 376, 377, 378, (383).

A list of all forms known to be false in the model: 1, 8, 15, 39, 40, 43, 44, 66, 67, 71, 86, 87, 95, 101, 106, 109, 113, 133, 149, 174, 181, 188, 192, 193, 202, 218, 255, 256, 257, 258, 259, 260, 261, 262, 264, 286, 333, 334, 335, 345, 347, 359, (383).)

\noindent $\mathcal{N}17$: Brunner/Howard\index{Author}{Brunner/Howard} Model II. $A=\{a_{\alpha,i}:\alpha\in \omega_{1}\wedge i\in\omega\}$. Let $\Gamma$ be the group of all permutations on $\omega_{1}$. $\mathcal{G}\ =\{g:(\exists\gamma\in\Gamma)(\forall\alpha\in \omega_{1})(\forall i\in \omega)g(a_{\alpha,i})=a_{\gamma(\alpha),i}\}$. $S$ is the set of all finite subsets of $A$. In this model, the union of a denumerable number of denumerable sets is denumerable (31 is true), but $C(\aleph_{0},2^{\aleph_{0}})$ (28) is false. If we assume that $2^{\aleph_{0}}=\aleph_{1}$ in the ground model, then $C(\aleph_{0},\,\aleph_{1})$ is also false. Therefore, $UT(\aleph_{0},\,\aleph_{1},\, \aleph_{1})$ (27) (which implies $C(\aleph_{0},\,\aleph_{1})$) is also false. Form 23 $((\forall\alpha)UT(\aleph_{\alpha},\,\aleph_{\alpha},\,\aleph_{\alpha}))$ is false because 23 implies 27. Since 231 $(UT(WO,\,WO,\,WO)$ implies 23 in every FM model,\index{Author}{Pincus} it follows that 231 is false.

\noindent $\mathcal{N}17\models$ 11, 31, 37, 73, 91, 92, 130, 191, 273, 305, 309, 13, 361, 363, 368, and 369, but 15, 27, 28, 62, 106, 131, 163, 165, 231, 323, and 344 are false. References Brunner/Howard \cite{bib:1992}, notes 2(3, 5, and 6), 18 and 120(30 49, 55, and 56).

(A minimal list of forms implying\index{Author}{Levy} each form known to be true in the model: 11, 31, 37, 73, 91, 92, 130, 191, 305, 309, 361, 363.

A minimal list of forms at least one of which is implied by each form known to be false in the model: 27, 28, 62, 131, 163, 165, 231, 344.

A list of all forms known to be true in the model: 0, 5, 6, 10, 11, 12, 13, 18, 19, 31, 32, 34, 35, 36, 37, 38, 70, 73, 74, 79, 80, 91, 92, 93, 94, 104, 108, 119, 130, 142, 145, 182, 189, 190, 191, 194, 197, 199, 203, 206, 209, 211, 212, 216, 222, 223, 249, 251, 252, 272, 273, 280, 288, 289, 305, 306, 307, 309, 350, 357, 358, 361, 362, 363, 364, 366, 367, 368, 369, 371, 373, 374, (383).

A list of all forms known to be false in the model: 1, 8, 14, 15, 16, 20, 21, 23, 27, 28, 29, 30, 39, 40, 43, 44, 49, 50, 60, 62, 66, 67, 68, 76, 85, 86, 87, 95, 101, 106, 107, 109, 113, 123, 126, 131, 133, 149, 163, 165, 174, 181, 188, 192, 193, 202, 213, 214, 218, 231, 255, 256, 257, 258, 259, 260, 261, 262, 264, 270, 286, 295, 303, 317, 323, 328, 331, 332, 333, 334, 335, 343, 344, 345, 346, 347, 359, (383).)

\noindent $\mathcal{N}18$: Howard's Model I. Let $B=B_{n}:n\in\omega$ where the $B_{n}$'s are pairwise disjoint and each is countably infinite and let $A=\bigcup B$. For each $n\in \omega$, let $F_{n}$ be a non-principal ultrafilter in $\mathcal{P}(B_{n})$ and let $I_{n}$ be the corresponding non-principal maximal ideal. $\mathcal{G}$ is the set of all permutations $\sigma$ of $A$ such that
\begin{enumerate}
\item[1.] $\sigma(B_{n})=B_{n}$, for each $n\in \omega$.
\item[2.] For each $n\in \omega$ there is a $u\in F_{n}$ such that $\sigma$ fixes $u$ pointwise.
\end{enumerate}
$S$ is the set of all sets of the form $C\cup\bigcup_{n\in\omega}g(n)$ where $C$ is a finite union of the $B_{n}$'s and $g:\omega \rightarrow \mathcal{P}(A)$ satisfies $g(n)\in I_{n}$ for each $n\in \omega$. In $\mathcal{N}18,\,\{B_{n}:n\in \omega\}$ is a denumerable collection of denumerable sets whose union is not denumerable (31 is false), but AC for a denumerable number of denumerable sets ([32 A]) is true. Form 23 $((\forall\alpha)UT(\aleph_{\alpha},\,\aleph_{\alpha},\,\aleph_{\alpha}))$ is false because 23 implies 31. Since 231 $(UT(WO,\,WO,\,WO)$ implies 23 in every FM model, it follows that 231 is false.

\noindent $\mathcal{N}18\models$ 6, 11, 32, 37, 73, 91, 92, 130, 191, 273, 305, 309, 313, 361, 363, 368, and 369, but 15, 31, 62, 131, 163, 165, 231, 323, and 344 are false. References Howard\index{Author}{Howard} \cite{bib:1992}, notes 2(3, 5, and 6), 18 and 120 (30, 49, and 56).

(A minimal list of forms implying each form known to be true in the model: 6, 11, 32, 37, 73, 91, 92, 130, 191, 305, 309, 313, 361, 363.

A minimal list of forms at least one of which is implied by each form known to be false in the model: 31, 62, 131, 163, 165, 231, 344.

A list of all forms known to be true in the model: 0, 5, 6, 10, 11, 12, 13, 18, 32, 35, 36, 37, 38, 70, 73, 74, 79, 80, 91, 92, 93, 94, 108, 119, 130, 142, 145, 182, 189, 190, 191, 194, 197, 199, 203, 206, 211, 212, 216, 222, 223, 249, 251, 252, 272, 273, 280, 288, 289, 305, 306, 307, 309, 313, 350, 357, 358, 361, 362, 363, 364, 366, 367, 368, 369, 371, 373, 374, (383).

A list of all forms known to be false in the model: 1, 8, 14, 15, 16, 20, 21, 23, 27, 28, 29, 30, 31, 39, 40, 43, 44, 49, 50, 60, 62, 66, 67, 68, 76, 85, 86, 87, 95, 101, 106, 107, 109, 113, 123, 126, 131, 133, 149, 163, 165, 174, 181, 188, 192, 193, 202, 213, 214, 218, 231, 255, 256, 257, 258, 259, 260, 261, 262, 264, 270, 286, 295, 303, 317, 323, 328, 331, 332, 333, 334, 335, 343, 344, 345, 346, 347, 359, (383).)

\noindent $\mathcal{N}19(\precsim)$: Tsukada's Model. Let $(P,\,\precsim)$ be a partially ordered set that is not well ordered; Let $Q$ be a countably infinite set, disjoint from $P$; and let $I=P\cup Q$. Let $A=$\{ $a_{iq}:i\in I,\,q$ a rational number\}; let $\mathcal{G}$ be the set of all permutations $\pi$ of $A$ such that $\pi(a_{iq})=a_{i\pi_{i}(q)}$ where for each $i\in I,\,\pi_{i}$ is an order automorphism of the rationals. $S$ is the set of all finite subsets of $A$. In this model, the Boolean Prime Ideal Theorem (14) is true and every set of well orderable sets has a choice function (60) is true. (Since 60 implies $UT(WO,\,WO,\, WO)$ (231) and in every FM model 231 implies $(\forall\alpha)UT(\aleph_{\alpha},\,\aleph_{\alpha},\,\aleph_{\alpha})$ (23), form 23 is true.) However, the set $(P,\,\precsim)$ can be embedded into the natural ordering of the cardinals so 59$(\precsim)$ is false. AC for a well ordered family of sets (40) is false. (In every FM model, 40 + 60 $\rightarrow$ AC, Howard/Rubin\index{Author}{Howard/Rubin} \cite{bib:1995b}.) Using the fact that the Kinna-Wagner Selection Principle (15) is false and that the Ordering Principle (30) is true, Krom proves that there is an ordered topological space that is not normal (118 is false). Form 14 implies 49 (Every partial ordering can be extended to a linear ordering.) and 49 + 51 (Every linear ordering has a cofinal sub-well ordering) iff AC. (See Morris \cite{bib:1969} and note 121.) Thus, form 51 is false. Howard \cite{bib:1973} has shown that in every FM model 14 implies the Weak Sikorski Theorem (317), so 317 is true. Form 40 is false and form 122 $(C(WO,\,<\infty))$ is true because 60 implies 122. Therefore, form 328 $(MC(WO,\,\infty))$ is false because 122 + 328 $\rightarrow$ 40. Since form 165 $(C(WO,\,WO))$ is true, (60 implies 165) it follows from note 2(8 and 9) that 16 and 24 are also true.

\noindent $\mathcal{N}19(\precsim)\models$ 11, 14, 16, 23, 24, 37, 60, 73, 91, 92, 130, 165, 191, 273, 305, 317, 361, 363, 368, and 369, but 15, 40, 51, 59$(\precsim)$, 118, 144, 153, 192, 253, and 328 are false. References Tsukada \cite{bib:1977}, Howard \cite{bib:1973}, Howard/Rubin \cite{bib:1995b}, Krom \cite{bib:1986}, Morris\index{Author}{Morris} \cite{bib:1969}, notes 2(8,9), 18, 107, 120(28, 34, 45, and 47), and 121.

(A minimal list of forms implying each form known to be true in the model: 11, 16, 23, 24, 60, 73, 91, 92, 130, 165, 191, 305, 317, 361, 363.

A minimal list of forms at least one of which is implied by each form known to be false in the model: 15, 51, 59, 118, 144, 153, 192, 253, 328.

A list of all forms known to be true in the model: 0, 5, 6, 10, 11, 12, 13, 14, 16, 18, 19, 23, 24, 25, 26, 27, 28, 30, 31, 32, 33, 34, 35, 36, 37, 38, 45, 46, 47, 48, 49, 52, 60, 61, 62, 63, 64, 69, 70, 72, 73, 74, 79, 80, 83, 85, 88, 91, 92, 93, 94, 96, 99, 102, 104, 107, 108, 111, 119, 120, 121, 122, 123, 127, 130, 132, 137, 139, 140, 141, 142, 145, 146, 150, 151, 153, 154, 165, 178, 182, 189, 190, 191, 194, 197, 198, 199, 201, 203, 206, 207, 209, 211, 212, 213, 216, 221, 222, 223, 225, 226, 227, 228, 229, 231, 233, 235, 241, 242, 249, 250, 251, 252, 268, 269, 270, 271, 272, 273, 276, 280, 283, 285, 287, 288, 289, 293, 294, 298, 305, 306, 307, 308, 309, 311, 313, 314, 315, 317, 323, 324, 326, 327, 329, 330, 331, 332, 336, 342, 343, 344, 349, 350, 356, 357, 358, 361, 362, 363, 364, 366, 367, 368, 369, 370, 371, 373, 374, 377, 378, 380, (383).

A list of all forms\index{Author}{Krom} known to be false\index{Author}{Howard/Rubin} in the model: 1, 14, 15, 40, 50, 51, 59, 66, 67, 89, 90, 95, 101, 109, 112, 114, 118, 133, 144, 149, 153, 179, 192, 202, 218, 253, 255, 256, 257, 258, 259, 260, 261, 262, 264, 286, 292, 303, 317, 328, 333, 334, 335, 345, 347, 359, (383).)

\noindent $\mathcal{N}20$: Truss'\index{Author}{Truss} Model\index{Author}{Morris} II. Let $X=\{a(i,\,k,\,l):i\in 2,\,k\in \mathbb{Z},\,l\in \omega\},\,Y=\{a(i,j,\,k,\,l):i,\,j\in 2,\,k\in \mathbb{Z},\,i\in \omega\}$ and $A$ is the disjoint union of $X$ and $Y$. The set of atoms is countable with order type $(2\times \mathbb{Z}\times \omega)\cup(2^{2}\times \mathbb{Z}\times \omega)$. Let $V(k,\,l)=\{a(i,\,k,\,l):i\in 2\},\,W(k,\, l)=\{a(i,j,\,k,\,l):i,\,j\in 2\},\,S=\{V(k,\,l):k\in \mathbb{Z},\,l\in \omega\},\, T=\{W(k,\,l):k\in \mathbb{Z},\,l\in \omega\}$, and $U=\{\{V(k,\,l),\,W(k,\,l)\}:k\in \mathbb{Z},\,l\in \omega\}\cup\{\{V(k+1,\,l),\,W(k,\,l)\}$: $k\in \mathbb{Z},\,l\in \omega\}$. Using forcing, construct a generic choice function $F$ on the set of all pairs of elements in $S\cup T$. Using $F$, it is possible to construct a function $f$ that maps $2\times Y$ onto $2\times X$. $\mathcal{G}$ is the group of all permutations of $A$ that leaves $f,\,X,\,Y,\,S,\,T$, and $U$ fixed. The set of supports is the finite subsets of \emph{A}. (This is a Fraenkel-Mostowski-Cohen model.) In this model, there is a mapping of $2\times Y$ onto $2\times X$, but no mapping of $Y$ onto $X$.

\noindent $\mathcal{N}20\models$ 6, 11, 37, 73, 91, 92, 130, 191, 273, 305, 309, 313, 361, 363, 368, and 369, but 15, 136$(k)$, for all $k\in \omega-\{0, 1\}$, and 192 are false. References Truss \cite{bib:1984} and note 18.

(A minimal list of forms implying each form known to be true in the model: 6, 11, 37, 73, 91, 92, 130, 191, 305, 309, 313, 361, 363.

A minimal list of forms at least one of which is implied by each form known to be false in the model: 15, 136, 192.

A list of all forms known to be true in the model: 0, 5, 6, 11, 12, 13, 35, 36, 37, 38, 70, 73, 74, 79, 91, 92, 93, 94, 108, 130, 142, 145, 182, 189, 190, 191, 194, 197, 199, 203, 206, 211, 212, 222, 223, 251, 252, 272, 273, 280, 289, 305, 306, 307, 309, 313, 361, 362, 363, 364, 366, 367, 368, 369, 371, (383).

A list of all forms known to be false in the model: 1, 15, 136, 138, 192, (383).)

\noindent $\mathcal{N}21(\aleph_{\alpha+1})$: Jensen's Model. We assume $\aleph_{\alpha+1}$ is a regular cardinal. Let $A= \{(a_{i_{1}},\,a_{i_{2}},\, \ldots,\,a_{i_{n}}):n\in \omega\wedge(\forall j\leq n)(i_{j}<\aleph_{\alpha+1})\}$. $A$ is partially ordered by extension. Let $\mathcal{G}$ be the group of all permutations of $A$ that preserves the partial ordering, and let $S$ be the set of all subsets $E$ of $A$ that satisfy the following properties:
\begin{enumerate}
\item[(1)] The cardinality of $E$ is less than $\aleph_{\alpha+1}$;
\item[(2)] If $s\subseteq t$ and $t\in E$, then $s\in E$; and
\item[(3)] There is no function $f:\omega\rightarrow\aleph_{\alpha+1}\,f|n\in E$ for all $n\in\omega$.
\end{enumerate}
In this model, the Principle of Dependent Choices (43) is false. The Axiom of choice for a family of $\aleph_{\beta}$ sets $(86(\beta))$ is true for all $\beta<\alpha+1$ and, therefore, the Axiom of Choice for a family of $\aleph_{1}$ sets (39) is true for all $\alpha\geq 1$. For all cardinals $\kappa,\,\kappa$ is comparable with $\aleph_{\beta}(71(\beta))$ is true for each $\beta<\alpha+1$. Since 39 implies 9 (Every Dedekind finite set is finite.) and 70 (There is a non-trivial ultrafilter on $\omega$.) is true in every FM model, it follows that 63 (Every infinite set has a non-trivial ultrafilter.) is also true. Since 39 implies 8 $(C(\aleph_{0},\,\infty))$, it follows from Brunner\index{Author}{Brunner} \cite{bib:1982a} that in this model there is a set that cannot be well ordered and does not have an infinite Dedekind finite subset, (163 is false). (Form 8 plus form 163 iff AC.) Blass\index{Author}{Blass} \cite{bib:1979} has shown that 191 (There is a set $X$ such that for each set $a$ there is an ordinal $\alpha$ and a function $f$ mapping $X\times\alpha$ onto $a$.) is true in every permutation model with a set of atoms. He has also shown that 8 + 191 + 193 (Every Abelian group is a homomorphic image of a free projective Abelian group.) iff AC. Since 8 and 191 are true, 193 must be false. (See note 60 for definitions.)

\noindent $\mathcal{N}21(\aleph_{\alpha+1})\models$ 11, 39 (for $\alpha\geq 1)$, 63, 71$(\beta)$ (for $\beta<\alpha+1)$, 86$(\beta)$ (for $\beta<\alpha+1)$, 91, 130, 191, 273, 305, 309, 313, 363, 368, and 369, but 15, 43, 106, 133, 163, and 193 are false. References Jensen\index{Author}{Jensen} \cite{bib:1966}, \cite{bib:1967}, Blass \cite{bib:1979}, Brunner \cite{bib:1982a}, and note 60. See, also, Jech\index{Author}{Jech} \cite{bib:1973b} (p 130) and Felgner\index{Author}{Felgner} \cite{bib:1971a} (pp 151-159), notes 18 and 120(4 and 55).

(A minimal list of forms implying each form known to be true in the model: 11, 39, 63, 71, 86, 91, 130, 191, 305, 309, 313, 363.

A minimal list of forms at least one of which is implied by each form known to be false in the model: 15, 106, 133, 163.

A list of all forms known to be true in the model: 0, 5, 6, 8, 9, 10, 11, 12, 13, 16, 17, 18, 19, 24, 26, 27, 28, 29, 31, 32, 34, 35, 36, 37, 38, 39, 57, 63, 64, 70, 71, 73, 74, 75, 79, 80, 82, 83, 84, 86, 91, 92, 93, 94, 98, 104, 108, 119, 124, 126, 127, 128, 130, 131, 132, 142, 145, 150, 182, 183, 185, 189, 190, 191, 194, 196, 197, 198, 199, 203, 206, 209, 211, 212, 216, 217, 221, 222, 223, 249, 251, 252, 272, 273, 280, 282, 288, 289, 296, 304, 305, 306, 307, 309, 313, 325, 336, 337, 338, 339, 340, 341, 342, 350, 351, 352, 353, 354, 355, 357, 358, 360, 361, 362, 363, 364, 366, 367, 368, 369, 371, 373, 374, 376, 377, 378, 380, (383).

A list of all forms known to be false in the model: 1, 15, 40, 43, 44, 66, 67, 87, 95, 101, 106, 109, 133, 149, 163, 174, 188, 192, 193, 202, 218, 255, 256, 257, 258, 259, 260, 261, 262, 264, 286, 333, 334, 335, 345, 347, 359, (383).)

\noindent $\mathcal{N}22(K,p)$: Makowski/Wi\'{s}niewski/Mostowski Model. $K$ is a finite subset of $\omega$. Let $p$ be a prime that is greater than every element of $K$. Let $A=\bigcup\{A_{i}:i\in \omega\}$ where The $A_{i}$'s are pairwise disjoint and each has cardinality $p$. For each $i\in\omega$, let $g_{i}$ be a permutation of $A_{i}$ which is a $p$-cycle. Let $\mathcal{G}$ be the group of all permutations of $A$ which are the identity on all but a finite number of the $A_{i}$'s and are equal to some power of $g_{i}$ on the others. Let $S$ be the set of all finite subsets of $A$. In this model, $C(\infty,\,K)$ (46$(K)$ is true. Thus, $C(\infty,\, n)(45(n))$ is true if $K=\{n\}$, and $C(\infty,\,2)$ (88) is true if $K=\{2\}$. However, there is a set $E$ and a function $f:E\rightarrow E$ such that $f$ has no fixed point and $E$ is not the union of three mutually disjoint sets, $E_{1},\,E_{2}$, and $E_{3}$, such that $E_{i}\cap f(E_{i})=\emptyset$, for $i=1,2,3$ (285 is false). In Howard/Yorke \cite{bib:1987}, it is shown that every group has a maximal $p$-subgroup in $\mathcal{N}22(K,p)$ (form $308(p)$ is true) and that $C(\aleph_{0},p)$ is false (form $288(p)$ is false).

\noindent $\mathcal{N}22(K,p)\models$ 6, 11, 37, 45$(n)$ (for $K=\{n\})$, $46(K)$, 73, 88 (for $K=\{2\})$, 91, 92, 130, 191, 273, 305, 308$(p)$, 309, 313, 361, 363, 368, and 369, but 285, 288$(p)$, and 344 are false. References Howard\index{Author}{Howard}/Yorke \cite{bib:1987}, Makowski/Wi\'{s}niewski \cite{bib:1969}, Mostowski \cite{bib:1945}, and note 120(56).

(A minimal list of forms implying each form known to be true in the model: 6, 11, 37, 45, 46, 73, 88, 91, 92, 130, 191, 305, 308, 309, 313, 361, 363.

A minimal list of forms at least one of which is implied by each form known to be false in the model: 285, 288, 344.

A list of all forms known to be true in the model: 0, 5, 6, 11, 12, 13, 18, 33, 35, 36, 37, 38, 45, 46, 47, 48, 64, 70, 73, 74, 79, 80, 88, 91, 92, 93, 94, 108, 120, 127, 130, 140, 142, 145, 182, 189, 190, 191, 194, 197, 199, 203, 206, 211, 212, 222, 223, 251, 252, 268, 269, 272, 273, 276, 280, 288, 289, 305, 306, 307, 308, 309, 313, 361, 362, 363, 364, 366, 367, 368, 369, 371, 373, (383).

A list of all forms known to be false in the model: 1, 2, 3, 4, 7, 8, 9, 10, 14, 15, 16, 17, 20, 21, 23, 27, 28, 29, 30, 31, 32, 33, 39, 40, 41, 43, 44, 45, 47, 49, 50, 60, 61, 62, 68, 71, 85, 86, 87, 100, 101, 107, 109, 113, 121, 122, 123, 129, 132, 133, 150, 151, 152, 153, 161, 165, 168, 174, 181, 192, 202, 213, 214, 218, 231, 250, 255, 256, 257, 258, 259, 260, 261, 262, 264, 270, 284, 285, 286, 288, 295, 303, 317, 322, 323, 324, 325, 327, 331, 332, 335, 343, 344, 345, 346, 347, 352, 354, 359, 374, 376, 377, 378, 380, (383).)

\noindent $\mathcal{N}23$: Howard/Mostowski\index{Author}{Mostowski}\index{Author}{Makowski/Wisniewski@Makowski/Wi\'{s}niewski} Linearly Ordered Model. This model is a modification of $\mathcal{N}3$. $A$ is countably infinite; $<$ is a dense linear ordering on $A$ without first or last elements $((A,\,<)\cong(\mathbb{Q},\, <))$; $\mathcal{G}$ is the group of all order automorphisms on $(A,\,<)$; and $S$ is the set of all subsets of $A$ that are well ordered by $<$. P. Howard\index{Author}{Howard/Yorke} has shown that AC for pairs (88) is false in this model and A. Rubin has shown that AC for a denumerable number of sets (8) is true. See note 100. It follows from Brunner\index{Author}{Brunner} \cite{bib:1982a} that in this model there is a set that cannot be well ordered and does not have an infinite Dedekind finite subset, (163 is false). (Form 8 plus form 163 iff AC.) Since 8 implies 9 (Every Dedekind finite set is finite.) and 70 (There is a non-trivial ultrafilter on $\omega$.) is true in every FM model, it follows that 63 (Every infinite set has a non-trivial ultrafilter.) is also true. Blass\index{Author}{Blass} \cite{bib:1979} has shown that 191 (There is a set $X$ such that for each set $a$ there is an ordinal $\alpha$ and a function $f$ mapping $X\times\alpha$ onto $a$.) is true in every permutation model with a set of atoms. He has also shown that 8 + 191 + 193 (Every Abelian group is a homomorphic image of a free projective Abelian group.) iff AC. Thus, 193 is false. (See note 60 for definitions.)

\noindent $\mathcal{N}23 \models$ 8, 11, 63, 91, 130, 191, 273, 305, 309, 313, 363, 368, and 369, but 88, 133, 163, 192, 193, and 344 are false. References Brunner \cite{bib:1982a}, Blass \cite{bib:1979}, notes 18, 60, 100, and 120(4 and 56).

(A minimal list of forms implying each form known to be true in the model: 8, 11, 63, 91, 130, 191, 305, 309, 313, 363.

A minimal list of forms at least one of which is implied by each form known to be false in the model: 88, 133, 163, 192, 193, 344.

A list of all forms known to be true in the model: 0, 5, 6, 8, 9, 10, 11, 12, 13, 16, 17, 18, 19, 24, 26, 27, 28, 29, 31, 32, 34, 35, 36, 37, 38, 57, 63, 64, 70, 73, 74, 75, 79, 80, 82, 83, 84, 91, 92, 93, 94, 98, 104, 108, 119, 124, 126, 127, 128, 130, 131, 132, 142, 145, 150, 182, 185, 189, 190, 191, 194, 197, 198, 199, 203, 206, 209, 211, 212, 216, 217, 221, 222, 223, 249, 251, 252, 272, 273, 280, 282, 288, 289, 296, 304, 305, 306, 307, 309, 313, 325, 336, 337, 338, 339, 340, 341, 342, 350, 351, 352, 353, 354, 355, 357, 358, 360, 361, 362, 363, 364, 366, 367, 368, 369, 371, 373, 374, 376, 377, 378, 380, (383).

A list of all forms known to be false in the model: 1, 14, 15, 30, 49, 50, 60, 61, 62, 68, 85, 88, 95, 107, 109, 123, 133, 163, 192, 193, 201, 213, 218, 270, 284, 295, 303, 317, 323, 326, 331, 332, 333, 335, 343, 344, 345, 346, (383).)

\noindent $\mathcal{N}24$: Hickman's Model I. This model is a variation of $\mathcal{N}2$. Let $A=\{a_{i}: i\in \omega\}\cup\{b_{i}:i\in \omega\}$, a union of pairwise disjoint denumerable sets; $\mathcal{G}$ is the group of all permutations on $A$ such that for each $\sigma\in \mathcal{G}$, and for each $i\in \omega,\,\sigma(\{a_{i},\,b_{i}\})= \{a_{j},\, b_{j}\}$, for some $j\in\omega$; and $S$ is the set of all finite subsets of $A$. Then both $A$ and $P=\{\{a_{i},\,b_{i}\}:i\in \omega\}$ are amorphous sets in $\mathcal{N}24$ (64 is false), so it follows from Brunner\index{Author}{Brunner} \cite{bib:1984a}, that in this model there is a Hilbert Space $H$ and an operator on $H$ which is not the direct sum of a finite matrix and a scalar operator (124 is false). Also, every denumerable family of finite sets has a choice function (10 is true) and every linearly ordered family of finite sets has a choice function (121 is true), but it is not the case that every infinite family of pairs has an infinite subfamily with a choice function (166 is false). There is an amorphous power of a compact $T_{2}$ space which is well orderable, but not compact (127 is false). In Brunner \cite{bib:1982a}, it is shown that if a set cannot be well ordered, then it has an amorphous subset (133 is true). It follows from Brunner \cite{bib:1982a} that MC (67) is false because 133 implies 164 (Every non well orderable set has an infinite subset with a Dedekind finite power set.) and 67 + 164 iff AC. It is shown in note 123 that 133 implies 231 $(UT(WO,\,WO,\,WO))$. In every FM model 231 implies 23 $((\forall\alpha)UT(\aleph_{\alpha},\,\aleph_{\alpha},\,\aleph_{\alpha}))$. It follows that 23 is also true. Since form 165 $(C(WO,\,WO))$ is true, (133 implies 165) it follows from note 2(8 and 9) that 16 and 24 are also true.

\noindent $\mathcal{N}24\models$ 10, 11, 16, 23, 24, 37, 73, 92, 121, 130, 133, 191, 273, 305, 309, 313, 361, and 363, but 64, 67, 124, 127, 131, 132, 166, and 344 are false. References Hickman\index{Author}{Hickman} \cite{bib:1976}, Brunner\index{Author}{Brunner} \cite{bib:1982a}, \cite{bib:1983b}, \cite{bib:1984a}, \cite{bib:1984b}, notes 2(8,9), 18, and 120(49 and 56).

(A minimal list of forms implying each form known to be true in the model: 11, 16, 23, 24, 73, 92, 121, 130, 133, 191, 305, 309, 313, 361, 363.

A minimal list of forms at least one of which is implied by each form known to be false in the model: 124, 127, 131, 132, 166, 344.

A list of all forms known to be true in the model: 0, 5, 6, 10, 11, 12, 13, 16, 18, 19, 23, 24, 25, 26, 27, 28, 31, 32, 33, 34, 35, 36, 37, 38, 47, 48, 51, 58, 63, 70, 73, 74, 77, 79, 80, 84, 90, 91, 92, 93, 94, 104, 108, 111, 118, 119, 120, 121, 122, 130, 133, 142, 145, 151, 182, 185, 189, 190, 191, 194, 197, 199, 203, 206, 207, 208, 209, 211, 212, 216, 221, 222, 223, 231, 249, 250, 251, 252, 272, 273, 280, 288, 289, 294, 305, 306, 307, 309, 313, 315, 316, 327, 350, 357, 358, 361, 362, 363, 364, 366, 367, 368, 369, 371, 373, 374, (383).

A list of all forms known to be false in the model: 1, 2, 3, 4, 7, 8, 9, 14, 15, 17, 30, 39, 40, 41, 43, 44, 45, 49, 50, 57, 60, 61, 62, 64, 66, 67, 68, 71, 76, 82, 83, 85, 86, 87, 88, 95, 100, 101, 106, 107, 109, 113, 123, 124, 126, 127, 129, 131, 132, 149, 150, 152, 161, 166, 168, 174, 181, 188, 192, 193, 201, 202, 213, 214, 215, 218, 255, 256, 257, 258, 259, 260, 261, 262, 264, 270, 284, 286, 295, 296, 303, 317, 323, 325, 326, 328, 331, 332, 333, 334, 335, 336, 343, 344, 345, 346, 347, 359, 376, 377, 378, 380, (383).)

\noindent $\mathcal{N}24(n)$: An extension of $\mathcal{N}24$ to $n$-element sets, $n>1.\,A=\bigcup B$, where $B=\{b_{i}:i\in \omega\}$ is a pairwise disjoint set of $n$-element sets, $\mathcal{G}$ is the group of all permutations of $A$ which are permutations of $B$; and $S$ is the set of all finite subsets of $A$. The set $B$ is a set of $n$-element sets which does not have a choice function (45($n$) is false).

Using similar arguments as in $\mathcal{N}24$, it follows that forms 10, 121, and 133 are true, but 64 is false. It follows from Brunner \cite{bib:1982a} that MC (67) is false because 133 implies 164 and 67 + 164 iff AC. It is shown in note 123 that 133 implies 231 $(UT(WO,\,WO,\,WO))$. In every FM model 231 implies 23 $((\forall\alpha)UT(\aleph_{\alpha},\,\aleph_{\alpha},\,\aleph_{\alpha}))$. It follows that 23 is also true. Since form 165 $(C(WO,\,WO))$ is true, (133 implies 165) it follows from note 2(8 and 9) that 16 and 24 are also true.

\noindent $\mathcal{N}24(n)\models$ 10, 11, 16, 23, 24, 37, 73, 92, 121, 130, 133, 191, 273, 305, 309, 313, 361, and 363, but 45$(n)$, 64, 67, 131, and 344 are false. References Brunner \cite{bib:1982a}, Hickman\index{Author}{Hickman} \cite{bib:1976}, Truss\index{Author}{Truss} \cite{bib:1973a}, notes 2(8,9), 18, and 120(49 and 56).

(A minimal list of forms implying each form known to be true in the model: 11, 16, 23, 24, 73, 92, 121, 130, 133, 191, 305, 309, 313, 361, 363.

A minimal list of forms at least one of which is implied by each form known to be false in the model: 64, 131, 344.

A list of all forms known to be true in the model: 0, 5, 6, 10, 11, 12, 13, 16, 18, 19, 23, 24, 25, 26, 27, 28, 31, 32, 33, 34, 35, 36, 37, 38, 47, 48, 51, 58, 63, 70, 73, 74, 77, 79, 80, 84, 90, 91, 92, 93, 94, 104, 108, 111, 118, 119, 120, 121, 122, 130, 133, 142, 145, 151, 182, 185, 189, 190, 191, 194, 197, 199, 203, 206, 207, 208, 209, 211, 212, 216, 221, 222, 223, 231, 249, 250, 251, 252, 272, 273, 280, 288, 289, 294, 305, 306, 307, 309, 313, 315, 316, 327, 350, 357, 358, 361, 362, 363, 364, 366, 367, 368, 369, 371, 373, 374, (383).

A list of all forms known to be false in the model: 1, 2, 3, 4, 7, 8, 9, 14, 15, 30, 39, 40, 41, 43, 44, 45, 49, 50, 57, 60, 61, 62, 64, 66, 67, 68, 71, 76, 82, 83, 85, 86, 87, 88, 95, 100, 101, 106, 107, 109, 113, 123, 126, 129, 131, 149, 152, 161, 168, 174, 181, 188, 192, 193, 201, 202, 213, 214, 215, 218, 255, 256, 257, 258, 259, 260, 261, 262, 264, 270, 284, 286, 295, 296, 303, 317, 323, 326, 328, 331, 332, 333, 334, 335, 336, 343, 344, 345, 346, 347, 359, 376, 377, 378, (383).)

\noindent $\mathcal{N}24(n,\,LO)$: Truss' Model III. This is a variation of $\mathcal{N}24(n)$ in which the set $B$ is linearly ordered. Suppose that $n>1$ and let $A=\{a_{qi}:q\in \mathbb{Q},\,i\in n\}$. For each $q\in \mathbb{Q}$, let $b_{q}=\{a_{qi}:i\in n\}$ and let $B=\{b_{q}:q\in \mathbb{Q}\}$. We define a linear ordering, $\precsim$, on $B$ so that for all $p,\,q\in \mathbb{Q},\, b_{p}\precsim b_{q}$ iff $p\leq q$. $\mathcal{G}$ is the group of all permutations on $A$ which are order automorphisms on $(B,\precsim)$ and $S$ is the set of all finite subsets of $A$. In this model $B$ is a linearly ordered set of $n$-element sets that does not have a choice function (33$(n)$ is false). However, every well ordered set of finite sets has a choice function (122 is true). Since 122 + 67 (MC) implies AC (see [67 B]), it follows that MC is false.

\noindent $\mathcal{N}24(n,\,LO)\models$ 6, 11, 37, 73, 91, 92, 122, 130, 191, 273, 305, 309, 313, 361, 363, 368, and 369, but 33$(n)$, 67, 192, and 344 are false. References Truss \cite{bib:1973a}, notes 18 and 120(56).

(A minimal list of forms implying each form known to be true in the model: 6, 11, 37, 73, 91, 92, 122, 130, 191, 305, 309, 313, 361, 363.

A minimal list of forms at least one of which is implied by each form known to be false in the model: 33, 67, 192, 344.

A list of all forms known to be true in the model: 0, 5, 6, 10, 11, 12, 13, 18, 35, 36, 37, 38, 47, 48, 70, 73, 74, 79, 80, 91, 92, 93, 94, 108, 111, 122, 130, 142, 145, 182, 189, 190, 191, 194, 197, 199, 203, 206, 211, 212, 216, 222, 223, 249, 250, 251, 252, 272, 273, 280, 288, 289, 305, 306, 307, 309, 313, 327, 358, 361, 362, 363, 364, 366, 367, 368, 369, 371, 373, 374, (383).

A list of all forms known to be false in the model: 1, 14, 15, 30, 33, 45, 49, 50, 60, 61, 62, 66, 67, 68, 85, 95, 107, 109, 121, 123, 149, 192, 202, 213, 218, 264, 270, 284, 295, 303, 317, 323, 331, 332, 333, 334, 335, 343, 344, 345, 346, (383).)

\noindent $\mathcal{N}25$: Brunner's Model I. The set of atoms, $A$, is equipped with the structure of the Hilbert space $\ell_{2},\, \mathcal{G}$ is the group of all permutations on $A$ that preserve the norm (unitary operators), and $S$ is the set of all finite subsets of $A$. Brunner\index{Author}{Brunner} has shown that in this model $A$ is a Hilbert space in which the unit ball is sequentially compact, but $A$ has no orthonormal basis (158 is false).

\noindent $\mathcal{N}25\models$ 6, 11, 37, 73, 91, 92, 130, 191, 273, 305, 309, 313, 361, 363, 368, and 369, but 15, 158, and 192 are false. References Brunner \cite{bib:1983b}, and note 18.

(A minimal list of forms implying each form known to be true in the model: 6, 11, 37, 73, 91, 92, 130, 191, 305, 309, 313, 361, 363.

A minimal list of forms at least one of which is implied by each form known to be false in the model: 15, 158, 192.

A list of all forms known to be true in the model: 0, 5, 6, 11, 12, 13, 35, 36, 37, 38, 70, 73, 74, 79, 91, 92, 93, 94, 108, 130, 142, 145, 182, 189, 190, 191, 194, 197, 199, 203, 206, 211, 212, 222, 223, 251, 252, 272, 273, 280, 289, 305, 306, 307, 309, 313, 361, 362, 363, 364, 366, 367, 368, 369, 371, (383).

A list of all forms known to be false in the model: 1, 15, 158, 192, (383).)

\noindent $\mathcal{N}26$: Brunner/Pincus\index{Author}{Pincus} Model, a variation of $\mathcal{N}2$. The set of atoms $A=\bigcup_{n\in\omega}P_{n}$, where the $P_{n}$'s are pairwise disjoint denumerable sets; $\mathcal{G}$ is the set of all permutations $\sigma$ on $A$ such that $\sigma(P_{n})=P_{n}$, for all $n\in \omega$; and $S$ is the set of all finite subsets of $A$. Each $P_{n}$ for $n\in \omega$ is amorphous (64 is false) and each set that cannot be well ordered has an infinite subset in 1-1 correspondence with $A$. Thus, Brunner proves that every set is either well orderable or has an infinite amorphous subset (133 is true). It is clear that the partial Kinna-Wagner Principle, $PKW(\aleph_{0},\,\infty,\,\infty)$ (167), the Kinna-Wagner Principle, $KW(\aleph_{0},\,\infty)$ (355), and the multiple-choice axiom, $MC_{\omega}(\aleph_{0},\,\infty)$ (131), are all false for the set $\{P_{n}:n\in \omega\}$. It is shown in note 123 that 133 implies 231 $(UT(WO,\,WO,\,WO))$. In every FM model 231 implies 23 $((\forall\alpha)UT(\aleph_{\alpha},\,\aleph_{\alpha},\,\aleph_{\alpha}))$. It follows that 23 is also true. De la Cruz and Di Prisco have shown that every infinite collection of non-empty well orderable sets has an infinite subfamily with a choice function (380 is true). Since form 165 $(C(WO,\,WO))$ is true, (133 implies 165) it follows from note 2(8 and 9) that 16 and 24 are also true.

\noindent \noindent $\mathcal{N}26\models$ 11, 16, 23, 24, 37, 73, 92, 130, 133, 191, 273, 305, 309, 313, 361, 363, and 380 but 64, 131, 167, 344, and 355, are false. References Brunner \cite{bib:1982a}, De la Cruz/Di Prisco\index{Author}{De la Cruz/Di Prisco}  \cite{bib:1996}, Pincus \cite{bib:1976}, notes 2(8,9), 18, and 120(56).

(A minimal list of forms implying each form known to be true in the model: 11, 16, 23, 24, 73, 92, 130, 133, 191, 305, 309, 313, 361, 363, 380.

A minimal list of forms at least one of which is implied by each form known to be false in the model: 64, 131, 167, 344, 355.

A list of all forms known to be true in the model: 0, 5, 6, 10, 11, 12, 13, 16, 18, 19, 23, 24, 25, 26, 27, 28, 31, 32, 34, 35, 36, 37, 38, 47, 48, 51, 58, 63, 70, 73, 74, 77, 79, 80, 84, 90, 91, 92, 93, 94, 104, 108, 111, 118, 119, 122, 130, 132, 133, 142, 145, 151, 182, 185, 189, 190, 191, 194, 197, 199, 203, 206, 207, 208, 209, 211, 212, 216, 221, 222, 223, 231, 249, 250, 251, 252, 272, 273, 280, 288, 289, 294, 305, 306, 307, 309, 313, 315, 316, 327, 342, 350, 357, 358, 361, 362, 363, 364, 366, 367, 368, 369, 371, 373, 374, 380, (383).

A list of all forms known to be false in the model: 1, 2, 3, 4, 7, 8, 9, 14, 15, 30, 39, 40, 41, 43, 44, 45, 49, 50, 57, 60, 61, 62, 64, 66, 67, 68, 71, 76, 82, 83, 85, 86, 87, 88, 95, 100, 101, 106, 107, 109, 113, 123, 126, 129, 131, 149, 152, 161, 167, 168, 174, 181, 188, 192, 193, 201, 202, 213, 214, 215, 218, 255, 256, 257, 258, 259, 260, 261, 262, 264, 270, 284, 286, 295, 296, 303, 317, 322, 323, 326, 328, 331, 332, 333, 334, 335, 336, 343, 344, 345, 346, 347, 355, 359, 376, 377, 378, (383).)

\noindent $\mathcal{N}27$: Hickman's Model II. Let $A$ be a set with cardinality $\aleph_{1}$ such that $A= \{(a_{\alpha},\, b_{\beta}):\alpha<\omega,\,\beta<\omega_{1}\}$. Let $<$ be the reverse lexicographic order on $A$, i.e., $(a_{\alpha},\,b_{\beta})<(a_{\gamma},\,b_{\delta})$ iff $\beta<\delta$, or $\beta=\delta$ and $\alpha<\gamma$. $\mathcal{G}$ is the group of all automorphisms on $A$ that preserve $<$, and $S$ is the set of all countable subsets of $A$. Hickman\index{Author}{Hickman} \cite{bib:1980b} has shown that if $\aleph_{\alpha}$ is a successor cardinal, then there exists an $\aleph_{\alpha}$ minimal set ($183(\alpha)$ is false for successor cardinals). In note 144 we show dependent choice (form 43) is true.

\noindent $\mathcal{N}27\models$ 6, 11, 37, 43, 63, 73, 91, 92, 130, 191, 273, 305, 309, 313, 361, 363, 368, and 369, but 15, 133, 163, 183$(\alpha)$, for $\aleph_{\alpha}$ a successor cardinal, and 192 are false. References Hickman \cite{bib:1980b}, notes 18, 56, 120(1, 4, and 20) and 144.

(A minimal list of forms implying each form known to be true in the model: 11, 43, 63, 91, 130, 191, 305, 309, 313.

A minimal list of forms at least one of which is implied by each form known to be false in the model: 15, 133, 163, 183, 192.

A list of all forms known to be true in the model: 0, 5, 6, 8, 9, 10, 11, 12, 13, 16, 17, 18, 19, 24, 26, 27, 28, 29, 31, 32, 34, 35, 36, 37, 38, 42, 43, 57, 63, 64, 70, 73, 74, 75, 77, 78, 79, 80, 82, 83, 84, 91, 92, 93, 94, 98, 104, 106, 108, 113, 119, 124, 126, 127, 128, 130, 131, 132, 142, 145, 150, 154, 155, 182, 185, 189, 190, 191, 194, 197, 198, 199, 203, 206, 209, 211, 212, 216, 217, 221, 222, 223, 243, 249, 251, 252, 272, 273, 279, 280, 282, 288, 289, 296, 304, 305, 306, 307, 309, 313, 325, 336, 337, 338, 339, 340, 341, 342, 350, 351, 352, 353, 354, 355, 357, 358, 360, 361, 362, 363, 364, 366, 367, 368, 369, 371, 373, 374, 375, 376, 377, 378, 380, (383).

A list of all forms known to be false in the model: 1, 15, 71, 87, 133, 163, 183, 192, (383).)

\noindent $\mathcal{N}28$: Blass' Permutation Model. The set $A=\{a_{\xi}^{i}:i\in \mathbb{Z},\,\xi\,\in\aleph_{1}\}$. For each $\xi \in\aleph_{1}$, let $g_{\xi}$ be the permutation of $A$ such that
\begin{equation*}
g_{\xi}(a_{\eta}^{i})=\left\{\begin{matrix}
a_{\eta}^{i+1} & \mathrm{if}\,\eta=\xi\\ a_{\eta}^{i} & \mathrm{if}\,\eta\,\neq\xi
\end{matrix}\right.
\end{equation*}
Let $\mathcal{G}$ be the group of permutations generated by all the $g_{\xi}$'s for $\xi \in\aleph_{1}$. Let
\begin{equation*}
S=\{\{a_{\xi}^{i}:i\in \mathbb{Z},\,\xi \in E\}:E\ \mathrm{a\,finite\,subset\,of}\,\aleph_{1}\}.
\end{equation*}
Let $\mathcal{N}$ be the permutation model generated by $A,\,\mathcal{G}$, and $S$. For any set $x,\,x$ is called \emph{small} if $A \cap TC(x)$ is countable. ($TC(x)$ is the transitive closure of $x$.) Let the universe of $\mathcal{N}28$ be the class of all small sets in $\mathcal{N}$. Blass shows that $\mathcal{N}28$ is a model of $\mathrm{ZF}^{0}$ in which $A$ is a proper class and he also shows that there is no non-trivial injective Abelian group in this model (190 is false). Blass also proves that form 190 is true in every permutation model in which the class of atoms is a set. Let $A^{\prime}=\Big\{a_{\xi}^{i}: i\in \mathbb{Z}\wedge\xi \in\aleph_{0}\Big\}$. $A^{\prime}$ is a denumerable subset of the atoms (in the ground model) and is therefore a set in the model. Further, $A^{\prime}$ is linearly orderable in $\mathcal{N}28$. For each $\xi \in\aleph_{0}$, let $O_{\xi}=\{a_{\xi}^{2i}:i\in \mathbb{Z}\}$ and $E_{\xi}=\{a_{\xi}^{2i+1}:i\in \mathbb{Z}\}$. The set $B=\{O_{\xi}:\xi\,\in\aleph_{0}\}\cup\{E_{\xi}:\xi\, \in\aleph_{0}\}$ cannot be linearly ordered in $\mathcal{N}28$ so form 293 is false. Form 129 is false because the set $B$ cannot be partitioned into sets of order type $^{\ast} \omega+\omega$. (See model $\mathcal{N}37$.)

\noindent $\mathcal{N}28\models$ 6, 11, 37, 73, 91, 92, 130, 273, 305, 309, 313, 361, 363, 368, and 369, but 129, 190, and 293 are false. References Blass\index{Author}{Blass} \cite{bib:1979}, and note 18.

(A minimal list of forms implying each form known to be true in the model: 6, 11, 37, 73, 91, 92, 130, 305, 309, 313, 361, 363.

A minimal list of forms at least one of which is implied by each form known to be false in the model: 129, 190, 293.

A list of all forms known to be true in the model: 0, 5, 6, 11, 12, 13, 35, 36, 37, 38, 70, 73, 74, 79, 91, 92, 93, 94, 108, 130, 142, 145, 194, 197, 199, 203, 206, 211, 212, 222, 223, 251, 252, 272, 273, 280, 289, 305, 306, 307, 309, 313, 361, 362, 363, 364, 366, 367, 368, 369, 371, (383).

A list of all forms known to be false in the model: 1, 14, 15, 30, 49, 50, 129, 189, 190, 191, 293, 295, 303, 317, 345, (383).)

\noindent $\mathcal{N}29$: Dawson/Howard\index{Author}{Dawson/Howard} Model. Let $A=\bigcup\{B_{n};n\in \omega\}$ is a disjoint union, where each $B_{n}$ is denumerable and ordered like the rationals by $\leq_{n}$. Thus, for each $n\in \omega,\,(B_{n},\, \leq_{n})\cong(\mathbb{Q},\,\leq)$. $\mathcal{G}$ is the group of all permutations on $A$ such that for all $n\in \omega$, and all $\phi\in \mathcal{G},\,\phi$ is an order automorphism of $(B_{n},\,\leq_{n})$ and $S$ is the set of all finite subsets of $A$. In Dawson/Howard it is shown that $|A!|<|2^{A}|$, thus form 200 $((\forall x)(|2^{x}|=|x!|))$ is false. Also, since $|2^{A}|$ is clearly $\leq|A^{A}|$, form 291 $((\forall x)(|x!|=|x^{x}|))$ is also false. In note 136 we prove that BPI (14) is true. Thus, $C(\infty,\, <\aleph_{0})$ (62) is true. Howard\index{Author}{Howard} \cite{bib:1973} proves that in every FM model 62 implies 60 $(C(\infty,\,WO))$. However, the set $\{B_{n}:n\in \omega\}$ is a denumerable set that has no multiple choice function $(MC(\aleph_{0},\,\infty)$ (126) is false) nor does it have an infinite subset with a Kinna-Wagner selection function $(PKW(\aleph_{0},\,\geq 2,\,\infty)$ (167) is false). Since form 165 $(C(WO,\,WO))$ is true, (60 implies 165) it follows from note 2(8 and 9) that 16 and 24 are also true.

\noindent $\mathcal{N}29\models$ 6, 11, 14, 16, 23, 24, 37, 60, 73, 91, 92, 130, 165, 191, 273, 305, 317, 361, 363, 368, and 369, but 15, 51, 118, 126, 131, 133, 144, 152, 167, 200, 253, and 291 are false. References Dawson/Howard \cite{bib:1976}, Howard \cite{bib:1973}, notes 2(3, 8, 9,), 18, 64, 120(3, 4, 6, 16, 28, 34, 45, 47, and 49), and 136.

(A minimal list of forms implying each form known to be true in the model: 11, 16, 23, 24, 60, 73, 91, 92, 130, 165, 191, 305, 317, 361, 363.

A minimal list of forms at least one of which is implied by each form known to be false in the model: 15, 51, 118, 131, 144, 152, 167, 200, 253, 291.

A list of all forms known to be true in the model: 0, 5, 6, 10, 11, 12, 13, 14, 16, 18, 19, 23, 24, 25, 26, 27, 28, 30, 31, 32, 33, 34, 35, 36, 37, 38, 45, 46, 47, 48, 49, 52, 60, 61, 62, 63, 64, 69, 70, 72, 73, 74, 79, 80, 83, 85, 88, 91, 92, 93, 94, 96, 99, 102, 104, 107, 108, 111, 119, 120, 121, 122, 123, 127, 130, 132, 137, 139, 140, 141, 142, 145, 146, 150, 151, 153, 154, 165, 178, 182, 189, 190, 191, 194, 197, 198, 199, 201, 203, 206, 207, 209, 211, 212, 213, 216, 221, 222, 223, 225, 226, 227, 228, 229, 231, 233, 235, 241, 242, 249, 250, 251, 252, 268, 269, 270, 271, 272, 273, 276, 280, 283, 285, 287, 288, 289, 293, 294, 298, 305, 306, 307, 308, 309, 311, 313, 314, 315, 317, 323, 324, 326, 327, 329, 330, 331, 332, 336, 342, 343, 344, 349, 350, 356, 357, 358, 361, 362, 363, 364, 366, 367, 368, 369, 370, 371, 373, 374, 377, 378, 380, (383).

A list of all forms known to be false in the model: 1, 8, 15, 39, 40, 43, 44, 51, 66, 67, 76, 86, 87, 89, 90, 95, 101, 106, 109, 112, 113, 114, 118, 126, 131, 133, 144, 149, 152, 167, 174, 179, 181, 188, 192, 193, 200, 202, 214, 218, 253, 255, 256, 257, 258, 259, 260, 261, 262, 264, 286, 291, 292, 328, 333, 334, 335, 345, 347, 359, (383).)

\noindent $\mathcal{N}30$: L\"{a}uchli's Model III. The set $A$ is denumerable; $\mathcal{G}$ is the group generated by the set of transpositions on $A$; and $S$ is the set of all finite subsets of $A$. The free group generated by $A$ does not have a free commutator subgroup (210 is false).

\noindent $\mathcal{N}30\models$ 6, 11, 37, 73, 91, 92, 130, 191, 273, 305, 309, 313, 361, 363, 368, and 369, but 15, 192, and 210 are false. References L\"{a}uchli\index{Author}{Lauchli@L\"{a}uchli} \cite{bib:1962}, and note 18.

(A minimal list of forms implying each form known to be true in the model: 6, 11, 37, 73, 91, 92, 130, 191, 305, 309, 313, 361, 363.

A minimal list of forms at least one of which is implied by each form known to be false in the model: 15, 192, 210.

A list of all forms known to be true in the model: 0, 5, 6, 11, 12, 13, 35, 36, 37, 38, 70, 73, 74, 79, 91, 92, 93, 94, 108, 130, 142, 145, 182, 189, 190, 191, 194, 197, 199, 203, 206, 211, 212, 222, 223, 251, 252, 272, 273, 280, 289, 305, 306, 307, 309, 313, 361, 362, 363, 364, 366, 367, 368, 369, 371, (383).

A list of all forms known to be false in the model: 1, 15, 68, 192, 210, (383).)

\noindent $\mathcal{N}31$: L\"{a}uchli's Model IV. The set $A$ is denumerable. A field, $F$ is constructed with base set $A$ which is isomorphic to the algebraic closure of $\mathbb{Q}$. $\mathcal{G}$ is the set of all permutations on $A$ that are automorphisms on $F$. $S$ is the set of all finite subsets of $A$. L\"{a}uchli proves that in this model, $F$ has no real closed subfield so 241 is false. Also, Hodges\index{Author}{Hodges} proves that if $R$ is the ring of integers in $F$, then $R$ is a principal ideal domain that is not a unique factorization domain and has no maximal ideal (243 and 244 are false).

\noindent $\mathcal{N}31\models$ 6, 11, 37, 73, 91, 92, 130, 191, 273, 305, 309, 313, 361, 363, 368, and 369, but 15, 241, 243, and 244 are false. References Hodges \cite{bib:1976a}, L\"{a}uchli \cite{bib:1962}, notes 18 and 84.

(A minimal list of forms implying each form known to be true in the model: 6, 11, 37, 73, 91, 92, 130, 191, 305, 309, 313, 361, 363.

A minimal list of forms at least one of which is implied by each form known to be false in the model: 15, 241, 243, 244.

A list of all forms known to be true in the model: 0, 5, 6, 11, 12, 13, 35, 36, 37, 38, 70, 73, 74, 79, 91, 92, 93, 94, 108, 130, 142, 145, 182, 189, 190, 191, 194, 197, 199, 203, 206, 211, 212, 222, 223, 251, 252, 272, 273, 280, 289, 305, 306, 307, 309, 313, 361, 362, 363, 364, 366, 367, 368, 369, 371, (383).

A list of all forms known to be false in the model: 1, 14, 15, 40, 43, 44, 50, 87, 101, 174, 192, 202, 233, 241, 242, 243, 244, 255, 256, 257, 258, 259, 260, 261, 262, 264, 278, 286, 303, 317, 345, 347, 359, (383).)

\noindent $\mathcal{N}32$: Hickman's Model III. This is a variation of $\mathcal{N}1$. Let $A$ be denumerable. Construct an elementary Abelean $p$-group $G$ on $A$ where $p$ is a prime. (Every non-identity element has order $p$.) Let $\mathcal{G}$ be the group of automorphisms of $G$, and let $S$ be the set of finite subsets of $A$. Hickman proves that if $H$ is any nontrivial subgroup of $G$, then the order of $H$ does not divide the order of $G$ (237 is false). Thus, it follows that $G$ is not the direct sum of cyclic groups (238 is false). Hickman also shows that in $G$ every ascending chain of subgroups is finite, but $G$ is not finitely generated (240 is false).

\noindent $\mathcal{N}32\models$ 6, 11, 37, 73, 91, 92, 130, 191, 273, 305, 309, 313, 361, 363, 368, and 369, but 15, 192, 237, 238, and 240 are false. References Hickman\index{Author}{Hickman} \cite{bib:1976}, and note 18.

(A minimal list of forms implying each form known to be true in the model: 6, 11, 37, 73, 91, 92, 130, 191, 305, 309, 313, 361, 363.

A minimal list of forms at least one of which is implied by each form known to be false in the model: 15, 192, 237, 238, 240.

A list of all forms known to be true in the model: 0, 5, 6, 11, 12, 13, 35, 36, 37, 38, 70, 73, 74, 79, 91, 92, 93, 94, 108, 130, 142, 145, 182, 189, 190, 191, 194, 197, 199, 203, 206, 211, 212, 222, 223, 251, 252, 272, 273, 280, 289, 305, 306, 307, 309, 313, 361, 362, 363, 364, 366, 367, 368, 369, 371, (383).

A list of all forms known to be false in the model: 1, 15, 192, 237, 238, 240, (383).)

\noindent $\mathcal{N}33$: Howard\index{Author}{Howard}/H.Rubin/J.Rubin\index{Author}{Howard/H.Rubin/J.Rubin} Model. $A$ is countably infinite; $\precsim$ is a dense linear ordering on $A$ without first or last elements $((A,\,\precsim)\cong(\mathbb{Q},\,\leq))$; $\mathcal{G}$ is the group of all order automorphisms on $(A,\,\precsim)$; and $S$ is the set of all bounded subsets of $A$. It is shown that the axiom of choice for a well ordered set of non-empty sets (40) is true. (Since 40 implies $UT(WO,\,WO,\,WO)$ (231) and in every FM model 231 implies $(\forall\alpha)UT(\aleph_{\alpha},\, \aleph_{\alpha},\,\aleph_{\alpha})$ (23), form 23 is true.) However, a set $x\subseteq A\times A$ and a transitive relation $R$ on $x$ are constructed so that every subset of $x$ on which $R$ is antisymmetric has an upper bound, but $x$ has no maximal element (260 is false). AC for a set of well orderable sets (60) is false. (In every FM model, 40 + 60 $\rightarrow$ AC, Howard/Rubin \cite{bib:1995b}.) In every FM model $C(\infty,\,<\aleph_{0})$ (62) implies 60 (Howard \cite{bib:1973}). Thus, 62 is false. Since 40 implies 8 $(C(\aleph_{0},\,\infty))$, it follows from Brunner\index{Author}{Brunner} \cite{bib:1982a} that in this model there is a set that cannot be well ordered and does not have an infinite Dedekind finite subset, (163 is false). (Form 8 plus form 163 iff AC.) Since 8 implies 9 (Every Dedekind finite set is finite.) and 70 (There is a non-trivial ultrafilter on $\omega$.) is true in every FM model, it follows that 63 (Every infinite set has a non-trivial ultrafilter.) is also true. Blass\index{Author}{Blass} \cite{bib:1979} has shown that 191 (There is a set $X$ such that for each set $a$ there is an ordinal $\alpha$ and a function $f$ mapping $X\times\alpha$ onto $a$.) is true in every permutation model with a set of atoms. He has also shown that 8 + 191 + 193 (Every Abelian group is a homomorphic image of a free projective Abelian group.) iff AC. Since 40 implies 8, form 193 must be false. (See note 60 for definitions.) Form 3 $(2m=m)$ is false because $2|A|\neq|A|$. Forms 349 $(MC(\infty,\,\aleph_{0}))$ and 356 $(KW(\infty,\,\aleph_{0}))$ are also false because the set of all infinite bounded intervals of $A$ is a set of denumerable sets which has no multiple choice function and no Kinna-Wagner selection function.

\noindent $\mathcal{N}33\models$ 11, 23, 40, 63, 91, 130, 165, 191, 305, 309, 313, 368, and 369, but 3, 60, 62, 133, 152, 163, 192, 193, 260, 344, 349 and 356 are false. References Howard/H.Rubin/J.Rubin \cite{bib:1973}, Blass \cite{bib:1979}, Brunner \cite{bib:1982a}, Howard \cite{bib:1973}, Howard/Rubin\index{Author}{Howard/Rubin} \cite{bib:1995b}, notes 18, 60, 107, and 120(4, 34, 46, 56).

(A minimal list of forms implying each form known to be true in the model: 11, 23, 40, 63, 91, 130, 165, 191, 305, 309, 313.

A minimal list of forms at least one of which is implied by each form known to be false in the model: 3, 62, 133, 152, 163, 192, 193, 260, 344, 349, 356.

A list of all forms known to be true in the model: 0, 5, 6, 8, 9, 10, 11, 12, 13, 16, 17, 18, 19, 23, 24, 25, 26, 27, 28, 29, 31, 32, 34, 35, 36, 37, 38, 39, 40, 42, 43, 47, 48, 57, 58, 63, 64, 70, 73, 74, 75, 77, 78, 79, 80, 82, 83, 84, 86, 91, 92, 93, 94, 98, 104, 106, 108, 111, 113, 119, 122, 124, 126, 127, 128, 130, 131, 132, 142, 145, 150, 151, 154, 155, 165, 182, 185, 189, 190, 191, 194, 196, 197, 198, 199, 203, 206, 207, 208, 209, 211, 212, 216, 217, 221, 222, 223, 231, 243, 249, 250, 251, 252, 272, 273, 279, 280, 282, 288, 289, 294, 296, 304, 305, 306, 307, 309, 313, 315, 322, 324, 325, 327, 328, 330, 336, 337, 338, 339, 340, 341, 342, 350, 351, 352, 353, 354, 355, 357, 358, 360, 361, 362, 363, 364, 366, 367, 368, 369, 371, 373, 374, 375, 376, 377, 378, 380, (383).

A list of all forms known to be false in the model: 1, 2, 3, 14, 15, 30, 49, 50, 60, 62, 66, 67, 68, 85, 95, 107, 109, 123, 133, 149, 152, 163, 192, 193, 202, 213, 214, 218, 255, 256, 257, 258, 259, 260, 261, 262, 264, 270, 295, 303, 317, 323, 329, 331, 332, 333, 334, 335, 343, 344, 345, 346, 349, 356, (383).)

\noindent $\mathcal{N}34$: Pincus'\index{Author}{Pincus} Model X. For each $q\in \mathbb{Q}$, let $C_{q}=\{a_{q},\,b_{q}\}$, a pair of atoms and let $A=\bigcup_{q\in \mathbb{Q}}C_{q}$. Let $\mathscr{C} =\{C_{q}:q\in \mathbb{Q}\,\}$ and let $\preceq$ be the ordering $\mathscr{C}$ inherits from $\mathbb{Q}$. (I. e., $C_{q}\preceq C_{r}\Leftrightarrow q\leq r$ where $\leq$ is the usual ordering on $\mathbb{Q}$.) Let $\mathcal{G}$ be the group of permutations of $A$ generated by the the two groups
\begin{align*}
G_{1}=\{\phi:&(\exists f)[f\,\mathrm{is\,an\,order\,automorphism\,of}\,\mathbb{Q}\\
&\wedge(\forall q\in \mathbb{Q})(\phi(a_{q})=a_{f(q)}\wedge\phi(b_{q})=b_{f(q)})]\}
\end{align*}
and
\begin{align*}
G_{2}=\{\psi:&\psi\ \mathrm{is\ a\ finite\ product\ of\ transposition},\\
&\psi=(a_{q_{1}},\,b_{q_{1}})\circ\cdots \circ(a_{q_{n}},\,b_{a_{n}})\ \mathrm{for\,some}\,q_{1},\ldots,\,q_{n}\in \mathbb{Q}\}.
\end{align*}
$G_{1}$ could also be described as $\{\phi:\phi:A\rightarrow A$ is one to one and onto and (the extension of) $\phi$ is an order isomorphism of $(\mathscr{C},\,\preceq)\}$. $S$ is the set of finite subsets of $A$. In this model, each infinite tree has either an infinite chain or an infinite antichain (216 is true), but there is an infinite partially ordered set that has neither an infinite chain nor an infinite antichain (217) is false. (See note 105.)

\noindent $\mathcal{N}34 \models$ 6, 11, 37, 73, 91, 92, 130, 191, 216, 273, 305, 309, 313, 361, 363, 368, and 369, but 15 and 217 are false. References Pincus \cite{bib:1972a}, notes 18 and 105.

(A minimal list of forms implying each form known to be true in the model: 6, 11, 37, 73, 91, 92, 130, 191, 216, 305, 309, 313, 361, 363.

A minimal list of forms at least one of which is implied by each form known to be false in the model: 15, 217.

A list of all forms known to be true in the model: 0, 5, 6, 11, 12, 13, 35, 36, 37, 38, 70, 73, 74, 79, 91, 92, 93, 94, 108, 130, 142, 145, 182, 189, 190, 191, 194, 197, 199, 203, 206, 211, 212, 216, 222, 223, 251, 252, 272, 273, 280, 289, 305, 306, 307, 309, 313, 361, 362, 363, 364, 366, 367, 368, 369, 371, (383).

A list of all forms known to be false in the model: 1, 2, 3, 4, 7, 8, 9, 15, 39, 40, 41, 43, 44, 71, 86, 87, 100, 101, 113, 129, 152, 161, 168, 174, 181, 192, 202, 214, 217, 255, 256, 257, 258, 259, 260, 261, 262, 264, 286, 345, 347, 359, (383).)

\noindent $\mathcal{N}35$: Truss' Model IV. The set of atoms, $A$, is denumerable and each element of $A$ is associated with a finite sequence of zeros and ones. We define a partial ordering $\precsim$ on $A$ such that if $a_{i},\, a_{j}\in A,\,a_{i}\precsim a_{j}$ if $a_{i}$ is an initial segment of $a_{j}$. Then, $(A,\,\precsim)$ is a tree. $\mathcal{G}$ is the group of permutation of $A$ which is generated by the set $\{\sigma_{a}:a\in A\}$ where if $\tau$ is the transposition (0, 1) and $a\precsim b$ then if $b=\langle a,\,i,\, b^{\prime}\rangle$, then $\sigma_{a}(b)=\langle a,\,\tau(i), b^{\prime}\rangle$, otherwise, $\sigma_{a}(b)=b$. For this model, it is easier to describe the filter $F$ of subgroups of $G$ rather than the ideal of supports $S$. $F$ is generated by $\{\mathcal{G}_{m}:m\in \omega\}$ where $\mathcal{G}_{0}=\mathcal{G}$, and $\mathcal{G}_{m+1}=\bigcap\{K:K$ is a subgroup of $\mathcal{G}_{m}\wedge|\mathcal{G}_{m}$: $K|=2\}$. In this model $(A,\,\precsim)$ is a binary tree with no infinite branch (249) is false, but for each $n\in \omega$ the Axiom of Choice for a well ordered family of $n$-element sets (250) is true.

\noindent $\mathcal{N}35\models$ 6, 11, 37, 73, 91, 92, 130, 191, 250, 273, 305, 309, 313, 361, 363, 368, and 369, but 249 and 344 are false. References Truss\index{Author}{Truss} \cite{bib:1975}, notes 18, 87, and 120(56).

(A minimal list of forms implying each form known to be true in the model: 6, 11, 37, 73, 91, 92, 130, 191, 250, 305, 309, 313, 361, 363.

A minimal list of forms at least one of which is implied by each form known to be false in the model: 249, 344.

A list of all forms known to be true in the model: 0, 5, 6, 11, 12, 13, 18, 35, 36, 37, 38, 47, 48, 70, 73, 74, 79, 80, 91, 92, 93, 94, 108, 111, 130, 142, 145, 182, 189, 190, 191, 194, 197, 199, 203, 206, 211, 212, 222, 223, 250, 251, 252, 272, 273, 280, 288, 289, 305, 306, 307, 309, 313, 361, 362, 363, 364, 366, 367, 368, 369, 371, 373, (383).

A list of all forms known to be false in the model: 1, 2, 3, 4, 7, 8, 9, 10, 14, 15, 16, 17, 20, 21, 23, 27, 28, 29, 30, 31, 32, 39, 40, 41, 43, 44, 49, 50, 60, 62, 68, 71, 85, 86, 87, 100, 101, 107, 113, 121, 122, 123, 129, 132, 133, 150, 151, 152, 153, 161, 165, 168, 174, 181, 192, 202, 213, 214, 231, 249, 255, 256, 257, 258, 259, 260, 261, 262, 264, 270, 286, 295, 303, 317, 323, 325, 331, 332, 343, 344, 345, 346, 347, 352, 354, 359, 376, 377, 378, 380, (383).)

\noindent $\mathcal{N}36(\beta)$: Brunner/Howard\index{Author}{Brunner/Howard} Model III. This model is a modification of $\mathcal{N}15$. $A= \{a_{i,\alpha}:i\in \omega\wedge\alpha\in \omega_{\beta+1}\}$. Let $\Gamma$ be the group of all permutations on $\omega$. Then $\mathcal{G}=\{g:(\exists\gamma\in\Gamma)(\forall i\in\omega))(\forall\alpha\in \omega_{\beta+1})g(a_{i,\alpha})=a_{\gamma(i),\alpha}\}$. $S$ is the set of all subsets of $A$ of cardinality at most $\aleph_{\beta}$. In this model $C(\aleph_{\beta},\,\infty)(86(\beta)$ is true and, therefore, $C(\aleph_{1},\,\infty)$ (39) is true if $\beta\geq 1$. The union of a denumerable set of finite sets is well orderable ([122 A] is true), but there is an infinite set that is not the union of a disjoint well orderable family of denumerable sets (152 is false). Since 39 implies 8 $(C(\aleph_{0},\, \infty))$, it follows from Brunner\index{Author}{Brunner} \cite{bib:1982a} that in this model there is a set that cannot be well ordered and does not have an infinite Dedekind finite subset, (163 is false). (Form 8 plus form 163 iff AC.) Since 8 implies 9 (Every Dedekind finite set is finite.) and 70 (There is a non-trivial ultrafilter on $\omega$.) is true in every FM model, it follows that 63 (Every infinite set has a non-trivial ultrafilter.) is also true. Blass\index{Author}{Blass} \cite{bib:1979} has shown that 191 (There is a set $X$ such that for each set $a$ there is an ordinal $\alpha$ and a function $f$ mapping $X\times\alpha$ onto $a$.) is true in every permutation model with a set of atoms. He has also shown that 8 + 191 + 193 (Every Abelian group is a homomorphic image of a free projective Abelian group.) iff AC. Thus, form 193 must be false. (See note 60 for definitions.) In note 144 we show dependent choice (form 43) if true.

\noindent $\mathcal{N}36(\beta)\models$ 11, 39, (if $\beta\geq 1)$, 43, 63, 86$(\beta)$, 91, 122, 130, 191, 273, 305, 309, 313, 363, 368, and 369, but 15, 67, 133, 152, 163, 192, and 193 are false. References Brunner/Howard \cite{bib:1992}, Blass \cite{bib:1979}, Brunner \cite{bib:1982a}, notes 18, 60, 120(4 and 41), and 144.

(A minimal list of forms implying each form known to be true in the model: 11, 39, 43, 63, 86, 91, 122, 130, 191, 305, 309, 313.

A minimal list of forms at least one of which is implied by each form known to be false in the model: 15, 67, 133, 152, 163, 192, 193.

A list of all forms known to be true in the model: 0, 5, 6, 8, 9, 10, 11, 12, 13, 16, 17, 18, 19, 24, 26, 27, 28, 29, 31, 32, 34, 35, 36, 37, 38, 39, 42, 43, 47, 48, 57, 63, 64, 70, 73, 74, 75, 77, 78, 79, 80, 82, 83, 84, 86, 91, 92, 93, 94, 98, 104, 106, 108, 111, 113, 119, 122, 124, 126, 127, 128, 130, 131, 132, 142, 145, 150, 154, 155, 182, 185, 189, 190, 191, 194, 196, 197, 198, 199, 203, 206, 209, 211, 212, 216, 217, 221, 222, 223, 243, 249, 250, 251, 252, 272, 273, 279, 280, 282, 288, 289, 296, 304, 305, 306, 307, 309, 313, 325, 327, 336, 337, 338, 339, 340, 341, 342, 350, 351, 352, 353, 354, 355, 357, 358, 360, 361, 362, 363, 364, 366, 367, 368, 369, 371, 373, 374, 375, 376, 377, 378, 380, (383).

A list of all forms known to be false in the model: 1, 15, 66, 67, 95, 109, 133, 149, 152, 163, 192, 193, 214, 218, 333, 334, 335, (383).)

\noindent $\mathcal{N}37$: A variation of Blass' model, $\mathcal{N}28$. Let $A=\{a_{i,j}:i\in \omega,\,j\in \mathbb{Z}\}$. $\mathcal{G}$ is the group of permutations of $A$ generated by $\{\phi_{i}:i\in\omega\}$, where $\phi_{i}(a_{i,j})=a_{i,j+1}$ and if $k\neq i,\,\phi_{i}(a_{k,j})=a_{k,j}$. $S$ is the set of finite subsets of $A$. For each $i\in\omega$, let $E_{i}=\{a_{i,2j}:j\in \mathbb{Z}\}$ and $O_{i}=\{a_{i,2j+1}:j\in \mathbb{Z}\}$. Then $B=\{E_{i}:i\in\omega\}\cup\{O_{i}:i\in \omega\}$ is a partition of $A$. We define a linear ordering $\precsim$ on $A$ such that $a_{i,j}\precsim a_{k,l}$ if $i<k$, or $i=k$ and $j\leq l$, the lexicographical ordering. Note that $\{a_{i,0}\}$ is a support for $E_{i}$ and $O_{i}$, while the empty set is a support for $(A,\,\precsim)$ and $B$. Thus, $A$ can be linearly ordered, but $B$ cannot be linearly ordered (293 is false). $B$ is also an infinite set which does not admit a partition into sets of order type $^{\ast} \omega+\omega$ so form 129 is false. In H\"{o}ft/Howard\index{Author}{Hoft/Howard@H\"{o}ft/Howard} \cite{bib:1981} it is shown that $(A,\precsim)$ has no cofinal well ordered subset (form 51 is false) and that every linear ordering with the fixed point property is complete (316 is true). $((A,\, \precsim)$ has the \emph{fixed point} property if every function $f:A\rightarrow A$ satisfying $(x\precsim y\Rightarrow f(x)\precsim f(y))$ has a fixed point, and $((A,\precsim)$ is \emph{complete} if every non-empty subset of $A$ has a least upper bound.)

\noindent $\mathcal{N}37 \models$ 6, 11, 37, 73, 91, 92, 130, 191, 273, 305, 309, 313, 316, 361, 363, 368, and 369, but 51, 129, 192, and 293 are false. References H\"{o}ft/Howard \cite{bib:1994}, and note 18.

(A minimal list of forms implying each form known to be true in the model: 6, 11, 37, 73, 91, 92, 130, 191, 305, 309, 313, 316, 361, 363.

A minimal list of forms at least one of which is implied by each form known to be false in the model: 51, 129, 192, 293.

A list of all forms known to be true in the model: 0, 5, 6, 11, 12, 13, 35, 36, 37, 38, 70, 73, 74, 77, 79, 84, 91, 92, 93, 94, 108, 130, 142, 145, 182, 185, 189, 190, 191, 194, 197, 199, 203, 206, 211, 212, 222, 223, 251, 252, 272, 273, 280, 289, 305, 306, 307, 309, 313, 316, 361, 362, 363, 364, 366, 367, 368, 369, 371, (383).

A list of all forms known to be false in the model: 1, 14, 15, 30, 49, 50, 51, 66, 67, 89, 90, 95, 109, 112, 114, 129, 133, 149, 192, 218, 256, 259, 261, 292, 293, 295, 303, 317, 333, 334, 335, 345, (383).)

\noindent $\mathcal{N}38$: Howard/Rubin Model I. Let $(A,\,\leq)$ be an ordered set of atoms which is order isomorphic to $\mathbb{Q}^{\omega}$, the set of all functions from $\omega$ into $\mathbb{Q}$ ordered by the lexicographic ordering. That is, if $a,\,b\in \mathbb{Q}^{\omega}$ then $a<b$ if and only if there is some $n\in \omega$ such that $(\forall j<n)(a_{j}=b_{j})$ and $a_{n}<b_{n}$.

For each $b\in A$ and $n\in \omega$,
\begin{enumerate}
\item[1.] $A_{b}^{n}=\{a\in A|a_{i}=b_{i}$ for $0\leq i\leq n\}$ is the $n-$\emph{level block containing} $b$.
($A_{b}^{n}$ will not be in the model we are defining.)
\item[2.] The sequence $\langle b_{n+1},\,b_{n+2},\ldots\rangle$ is the \emph{position of} $b$ \emph{in its} $n-$\emph{level block}.
\item[3.] $\mathscr{B}^{n}=\{A_{a}^{n}|a\in A\}$ is the set of $n-$level blocks.
\item[4.] $\leq_{n}$ is the relation on $\mathscr{B}^{n}$ define by $A_{a}^{n}\leq_{n}A_{b}^{n}$ if and only if $a\leq b$.
\item[5.] Let $f$ be an order automorphism of $(\mathscr{B}^{n},\,\leq_{n})$. (See lemmas A and B in Howard/Rubin \cite{bib:1996}.) We define $\phi_{f}$ to be the unique order automorphism of $(A,\,\leq)$ which satisfies $\phi_{f}^{\prime\prime}A_{a}^{n}=f(A_{a}^{n})$ for all $a\in A$ and such that for all $a\in A,\,a$ and $\phi_{f}(a)$ have the same position in their $n$-level blocks. (By 2 above, this means that $(\forall a\in A)(\forall i>n)(a_{i}=\phi_{f}(a)_{i})$.)
\item[6.] For $n\in \omega,\,G_{n}$ is the group $\{\phi_{f}:f$ is an order automorphism of $(\mathscr{B}^{n},\,\leq_{n})\}$.
\end{enumerate}

$\mathcal{G}$ is the group $\bigcup_{n\in\omega}G_{n}$. (Note that for $n<m,\,G_{n}\subseteq G_{m}$.) $S$ is the set of all $E\subseteq A$ which satisfy the following conditions:
\begin{enumerate}
\item[(a)] $E$ is well ordered by the ordering $\leq$ on $A$.
\item[(b)] $(\forall n\in \omega)$ ($ E\cap A_{a}^{n}\neq\emptyset$ for only finitely many $a\in A$)
\item[(c)] $E$ is countable.
\end{enumerate}

It is shown in Howard/Rubin\index{Author}{Howard/Rubin} \cite{bib:1996} that the Boolean Prime Ideal Theorem (14) and the Axiom of Choice for a denumerable family of sets (8) are true, but both the Principle of Dependent Choices (43) and the Axiom of Choice for families of cardinality $2^{\aleph_{0}}$ (181) are false. By using a slight modification of the proof of the Tychonoff Compactness Theorem (see, for example, Kelley\index{Author}{Kelley} \cite{bib:1955}, p 143), it can be shown that 14 + 8 implies that the product of a countable number of compact spaces is compact (113). It follows from note 104, that the $2m=m$ principle (3) is false. Using the fact that the Kinna-Wagner Selection Principle (15) is false and that the Ordering Principle (30) is true (14 implies 30), Krom proves that there is an ordered topological space that is not normal (118 is false). Since 8 is true, it follows from Brunner\index{Author}{Brunner} \cite{bib:1982a} that in this model there is a set that cannot be well ordered and does not have an infinite Dedekind finite subset, (163 is false). (Form 8 plus form 163 iff AC.) Since form [154 A] (If $X$ is a compact $T_{2}$ topological space, then $X^{\omega}$ is compact.) is true (14 implies 154) and 43 is false, 106 (the Baire Category Theorem for compact $T_{2}$ spaces) is false because 106 + 154 implies 43. (See Brunner \cite{bib:1984b} and note 95.) Form 14 implies 49 (Every partial ordering can be extended to a linear ordering.) and 49 + 51 (Every linear ordering has a cofinal sub well ordering.) iff AC. (See Morris\index{Author}{Morris} \cite{bib:1969} and note 121.) Thus, form 51 is false. Since form 14 implies form 30 (Ordering Principle) and form 113 implies form 296 (Every infinite set is the disjoint union of infinitely many infinite sets.), it follows from Pincus\index{Author}{Pincus} \cite{bib:1997} that form 295 (Every infinite set has a dense linear ordering.) is true. Form 14 implies 62 $(C(\infty,\,<\aleph_{0}))$ and Howard\index{Author}{Howard} \cite{bib:1973} has shown that in every FM model, 62 implies 60 $(C(\infty,\,WO))$. Thus, 60 is true. It is also proved in Howard \cite{bib:1973} that in every FM model 14 implies the Weak Sikorski Theorem (317), so 317 is also true. Form 40 $(C(WO,\,\infty))$ is false because 40 implies 106 and form 122 $(C(WO,\,<\aleph_{0}))$ is true because 60 implies 122. Therefore, form 328 $(MC(WO,\,\infty))$ is false because 122 + 328 $\rightarrow$ 40.

\noindent $\mathcal{N}38\models$ 8, 11, 14, 23, 60, 91, 113, 130, 165, 191, 273, 295, 305, 317, 363, 368, and 369, but 3, 15, 43, 51, 106, 118, 144, 152, 163, 181, 253, and 328 are false. References Howard/Rubin \cite{bib:1996}, Kelley \cite{bib:1955}, Krom \cite{bib:1986}, Blass\index{Author}{Blass} \cite{bib:1979}, Brunner \cite{bib:1982a}, \cite{bib:1984b}, Howard \cite{bib:1973}, Morris \cite{bib:1969}, Pincus \cite{bib:1997}, notes 2(3), 18, 60, 95, 104, 120(28, 34, 45, and 47) and 121.

(A minimal list of forms implying each form known to be true in the model: 11, 23, 60, 91, 113, 130, 165, 191, 295, 305, 317, 363.

A minimal list of forms at least one of which is implied by each form known to be false in the model: 3, 15, 51, 106, 118, 144, 152, 163, 181, 253, 328.

A list of all forms known to be true in the model: 0, 5, 6, 8, 9, 10, 11, 12, 13, 14, 16, 17, 18, 19, 23, 24, 25, 26, 27, 28, 29, 30, 31, 32, 33, 34, 35, 36, 37, 38, 45, 46, 47, 48, 49, 52, 57, 60, 61, 62, 63, 64, 69, 70, 72, 73, 74, 75, 79, 80, 82, 83, 84, 85, 88, 91, 92, 93, 94, 96, 98, 99, 102, 104, 107, 108, 111, 113, 119, 120, 121, 122, 123, 124, 126, 127, 128, 130, 131, 132, 137, 139, 140, 141, 142, 145, 146, 150, 151, 153, 154, 165, 178, 182, 185, 189, 190, 191, 194, 197, 198, 199, 201, 203, 206, 207, 209, 211, 212, 213, 216, 217, 221, 222, 223, 225, 226, 227, 228, 229, 231, 233, 235, 241, 242, 249, 250, 251, 252, 268, 269, 270, 271, 272, 273, 276, 280, 282, 283, 285, 287, 288, 289, 293, 294, 295, 296, 298, 304, 305, 306, 307, 308, 309, 311, 313, 314, 315, 317, 323, 324, 325, 326, 327, 329, 330, 331, 332, 336, 337, 338, 339, 340, 341, 342, 343, 344, 349, 350, 351, 352, 353, 354, 355, 356, 357, 358, 360, 361, 362, 363, 364, 366, 367, 368, 369, 370, 371, 373, 374, 376, 377, 378, 380, (383).

A list of all forms known to be false in the model: 1, 2, 3, 15, 40, 43, 44, 51, 66, 67, 87, 89, 90, 95, 101, 106, 109, 112, 114, 118, 133, 144, 149, 152, 163, 174, 179, 181, 188, 192, 193, 202, 214, 218, 253, 255, 256, 257, 258, 259, 260, 261, 262, 264, 286, 292, 328, 333, 334, 335, 345, 347, 359, (383).)

$\mathcal{N}39$: Howard's Model II. $A$ is denumerable and is a disjoint union $\bigcup_{i\in\omega}B_{i}\cup\bigcup_{i\in\omega}C_{i}$, where for all $i\in\omega$, $|B_{i}|=|C_{i}|=\aleph_{0}$. Let $f_{i}$ be a fixed bijection from $B_{i}$ onto $C_{i}$ for each $i\in \omega$. $\mathcal{G}$ is the group of all permutations $\phi$ of $A$ such that:
\begin{enumerate}
\item[1.] $\phi(B_{i})=B_{i}$ and $\phi(C_{i})=C_{i})$, $\forall i\in \omega$.
\item[2.] $\{a\in A:\phi(a)\neq a\}$ is finite.
\item[3.] $\phi|B_{i}$ and $\phi|C_{i}$ are even permutations, $\forall i\in \omega$.
\end{enumerate}
For each finite $x\subseteq \omega$ and each finite $E\subseteq A$, define
\begin{equation*}
G(x,\,E)=\{\phi\in \mathrm{fix}(E):(\forall i\in x)(\phi(f_{i})=f_{i})\}
\end{equation*}
Let $\Gamma$ be the filter of subgroups of $\mathcal{G}$ generated by the $G(x,\,E)$, where $x$ is a finite subset of $\omega$ and $E$ is a finite subset of $A$. (The filter is easier to describe for this model than the supports.) $\mathcal{N}39$ is the model determined by $A,\,\mathcal{G}$, and the filter $\Gamma$. It is shown in Howard \cite{bib:1994} that $|B_{i}|=|C_{i}|$ for all $i\in \omega$, but $|\bigcup_{i\in \omega}B_{i}|\neq |\bigcup_{i\in\omega}C_{i}|$, so 29 is false. Howard\index{Author}{Howard} also proves that the union of a well ordered collection of well orderable sets is well orderable $(UT(WO,\,WO,\,WO)$, form 231) is true. However, in every FM model, 231 implies $(\forall\alpha)UT(\aleph_{\alpha},\,\aleph_{\alpha},\,\aleph_{\alpha})$ so form 23 is also true. $C(WO,\,<\aleph_{0})$ (122) is true (231 implies 122) and 122 + 67 (MC) implies AC (see [67 B]), so it follows that MC is false. Since form 165 $(C(WO,\,WO))$ is true, (231 implies 165) it follows from note 2(8 and 9) that 16 and 24 are also true.

\noindent \noindent $\mathcal{N}39\models$ 11, 16, 23, 24, 37, 73, 91, 92, 130, 191, 231, 273, 305, 309, 313, 361, 363, 368, and 369, but 15, 29, 67, 131, 152, and 328 are false. References Howard \cite{bib:1994}, notes 2(8, 9), 18 and 120(33, 47, and 49).

(A minimal list of forms implying each form known to be true in the model: 11, 16, 23, 24, 73, 91, 92, 130, 191, 231, 305, 309, 313, 361, 363.

A minimal list of forms at least one of which is implied by each form known to be false in the model: 15, 29, 131, 152.

A list of all forms known to be true in the model: 0, 5, 6, 10, 11, 12, 13, 16, 18, 19, 23, 24, 25, 26, 27, 28, 31, 32, 34, 35, 36, 37, 38, 47, 48, 70, 73, 74, 79, 80, 91, 92, 93, 94, 104, 108, 111, 119, 122, 130, 142, 145, 151, 182, 189, 190, 191, 194, 197, 199, 203, 206, 207, 209, 211, 212, 216, 222, 223, 231, 249, 250, 251, 252, 272, 273, 280, 288, 289, 294, 305, 306, 307, 309, 313, 315, 327, 350, 357, 358, 361, 362, 363, 364, 366, 367, 368, 369, 371, 373, 374, (383).

A list of all forms known to be false in the model: 1, 8, 15, 20, 21, 29, 39, 40, 43, 44, 66, 67, 76, 86, 87, 95, 101, 106, 109, 113, 126, 131, 149, 152, 174, 181, 188, 192, 193, 202, 214, 218, 255, 256, 257, 258, 259, 260, 261, 262, 264, 286, 328, 333, 334, 335, 345, 347, 359, (383).)

\noindent $\mathcal{N}40$: Howard/Rubin\index{Author}{Howard/Rubin} Model II. A variation of $\mathcal{N}38$. $A$ and $\mathcal{G}$ are the same as in $\mathcal{N}38$, but $S$ is the set of all subsets of $A$ that satisfy conditions (b) and (c) in the definition of $S$ in $\mathcal{N}38$. Thus, we no longer require supports to be well ordered by $\leq$. $\mathcal{N}40$ has many of the same properties as $\mathcal{N}38$, BPI (14) and $C(\aleph_{0},\,\infty)$ (8) are true, but $C(2^{\aleph_{0}},\,\infty)$ (181) and the $2m=m$ principle (3) are false. (See note 104.) Form 113 (Tychonoff's Compactness Theorem for countably many spaces) is true here also, because 14 + 8 implies 113. However, in this model, the Principle of Dependent Choices (43) is true. See Howard/Rubin \cite{bib:1996} for proofs. Using the fact that the Kinna-Wagner Selection Principle (15) is false and the Ordering Principle (30) is true (14 implies 30), Krom\index{Author}{Krom} proves that there is an ordered topological space that is not normal (118 is false). Since 8 is true, it follows from Brunner\index{Author}{Brunner} \cite{bib:1982a} that in this model there is a set that cannot be well ordered and does not have an infinite Dedekind finite subset, (163 is false). (Form 8 plus form 163 iff AC.) Form 14 implies 49 (Every partial ordering can be extended to a linear ordering.) and 49 + 51 (Every linear ordering has a cofinal sub well ordering.) iff AC. (See Morris\index{Author}{Morris} \cite{bib:1969} and note 121.) Thus, form 51 is false. Blass\index{Author}{Blass} \cite{bib:1979} has shown that 191 (There is a set $X$ such that for each set $a$ there is an ordinal $\alpha$ and a function $f$ mapping $X\times\alpha$ onto $a$.) is true in every permutation model with a set of atoms. He has also shown that 8 + 191 + 193 (Every Abelian group is a homomorphic image of a free projective Abelian group.) iff AC. Thus, 193 is false. (See note 60 for definitions.) Since form 14 implies form 30 (Ordering Principle) and form 43 implies form 296 (Every infinite set is the disjoint union of infinitely many infinite sets.), it follows from Pincus\index{Author}{Pincus} \cite{bib:1997} that form 295 (Every infinite set has a dense linear ordering.) is true. Form 14 implies 62 $(C(\infty,\,<\aleph_{0}))$ and Howard\index{Author}{Howard} \cite{bib:1973} has shown that in every FM model, 62 implies 60 $(C(\infty,\,WO))$. Thus, 60 is true. Moreover, in every FM model, 60 + 40 $(C(WO,\,\infty))$ implies AC, so 40 is false. (See Howard/Rubin \cite{bib:1995b}.) It is also shown in Howard \cite{bib:1973} that in every FM model 14 implies the Weak Sikorski Theorem (317), so 317 is true. Form 40 is false and form 122 $(C(WO,\,<\infty))$ is true because 60 implies 122. Therefore, form 328 $(MC(WO,\,\infty))$ is false because 122 + 328 $\rightarrow$ 40. Since 43 + 14 is equivalent to 345 ([345 A]), form 345 is true.

\noindent $\mathcal{N}40\models$ 8, 11, 14, 23, 43, 60, 91, 113, 130, 165, 191, 295, 305, 317, 345, 368, and 369, but 3, 15, 40, 51, 118, 144, 152, 163, 181, 192, 193, 253, and 328 are false. References Blass \cite{bib:1979}, Brunner \cite{bib:1982a}, Howard \cite{bib:1973}, Howard/Rubin \cite{bib:1995b},\cite{bib:1996}, Krom \cite{bib:1986}, Pincus \cite{bib:1997}, notes 2(3), 18, 60, and 120(28, 34, 45, and 47).

(A minimal list of forms implying each form known to be true in the model: 11, 23, 60, 91, 130, 165, 191, 295, 305, 317, 345.

A minimal list of forms at least one of which is implied by each form known to be false in the model: 3, 15, 51, 118, 144, 152, 163, 181, 192, 193, 253, 328.

A list of all forms known to be true in the model: 0, 5, 6, 8, 9, 10, 11, 12, 13, 14, 16, 17, 18, 19, 23, 24, 25, 26, 27, 28, 29, 30, 31, 32, 33, 34, 35, 36, 37, 38, 42, 43, 45, 46, 47, 48, 49, 52, 57, 60, 61, 62, 63, 64, 69, 70, 72, 73, 74, 75, 77, 78, 79, 80, 82, 83, 84, 85, 88, 91, 92, 93, 94, 96, 98, 99, 102, 104, 106, 107, 108, 111, 113, 119, 120, 121, 122, 123, 124, 126, 127, 128, 130, 131, 132, 137, 139, 140, 141, 142, 145, 146, 150, 151, 153, 154, 155, 165, 178, 182, 185, 189, 190, 191, 194, 197, 198, 199, 201, 203, 206, 207, 209, 211, 212, 213, 216, 217, 221, 222, 223, 225, 226, 227, 228, 229, 231, 233, 235, 241, 242, 243, 249, 250, 251, 252, 268, 269, 270, 271, 272, 273, 276, 279, 280, 282, 283, 285, 287, 288, 289, 293, 294, 295, 296, 298, 304, 305, 306, 307, 308, 309, 311, 313, 314, 315, 317, 323, 324, 325, 326, 327, 329, 330, 331, 332, 336, 337, 338, 339, 340, 341, 342, 343, 344, 345, 349, 350, 351, 352, 353, 354, 355, 356, 357, 358, 360, 361, 362, 363, 364, 366, 367, 368, 369, 370, 371, 373, 374, 375, 376, 377, 378, 380, (383).

A list of all forms known to be false in the model: 1, 2, 3, 15, 40, 51, 66, 67, 89, 90, 95, 101, 109, 112, 114, 118, 133, 144, 149, 152, 163, 179, 181, 192, 193, 202, 214, 218, 253, 255, 256, 257, 258, 259, 260, 261, 262, 264, 286, 292, 328, 333, 334, 335, 347, 359, (383).)

\noindent $\mathcal{N}41$: Another variation of $\mathcal{N}3$. $A=\bigcup\{B_{n};n\in\omega\}$ is a disjoint union, where each $B_{n}$ is denumerable and ordered like the rationals by $\leq_{n}$. Thus, for each $n \in \omega,\,(B_{n},\, \leq_{n})\cong(\mathbb{Q},\,\leq)$. $\mathcal{G}$ is the group of all permutations on $A$ such for all $n\in\omega$, and all $\phi\in \mathcal{G},\,\phi$ is an order automorphism of $(B_{n},\,\leq_{n})$. $S$ is the set of all finite unions of the $B_{n}$'s. It follows that $A$ can be linearly ordered in this model, for if $a\in B_{m}$ and $b\in B_{n},\,a\leq b$ if $m<n$ or $m=n$ and $a\leq_{n}b$. For each $n \in \omega,\,B_{n}$ can be well ordered so $A$ is the union of a well ordered set of well orderable sets. Since $A$ can be linearly ordered and not well ordered, it follows that 294 is false. (See note 25 for definitions.) Let $C_{j}$ be the set of all pairs of countable subsets of $B_{j}$, and let $D=\bigcup_{j\in\omega}C_{j}$. Then $D$ has no choice function. (Let $\{E_{j},\,O_{j}\}\in C_{j}$ such that $E_{j}=\{a_{2i}^{j}:i\in \mathbb{Z}\}$ and $O_{j}=\{a_{2i+1}^{j}:i\in \mathbb{Z}\}$, where $\mathbb{Z}$ is the set of integers. (If $m,\,n\in \mathbb{Q},\,a_{m}^{j}\leq a_{n}^{j}$ iff $m\leq n$.) If $f$ were a choice function on $D$ with support $F$, then there is a $\phi\in$ fix $(F)$ and a $j\in \omega$ such that $\phi$ interchanges $E_{j}$ and $O_{j}$.) Thus, $C(\infty,\,2)$ (88) is false. Also, in this model every infinite set has a countably infinite subset (9) is true and the axiom of choice for well orderable families of finite sets (122) is true. See note 112. Since 70 (There is a non-trivial ultrafilter on $\omega$) is true in every FM model and 9 + 70 implies 63 (Every infinite set has a non-trivial ultrafilter.) it follows that 63 is true. Form 10 $(C(\aleph_{0},\,<\aleph_{0}))$ is true (9 implies 10) and form 231 $(UT(WO,\,WO,\,WO))$ is false (231 implies 294). It follows that form 163 (Every non well orderable set has an infinite Dedekind finite subset) is false. (10 + 163 implies 231. See note 123.) The set $\{B_{n}:n\in\omega\}$ is a denumerable set of denumerable sets that has no Kinna-Wagner selection function and no multiple choice function so $KW(\aleph_{0},\,\aleph_{0})$ (357) and $MC(\aleph_{0},\,\aleph_{0})$ (350) are both false.

\noindent $\mathcal{N}41\models$ 6, 9, 11, 37, 63, 73, 91, 92, 122, 130, 191, 273, 305, 309, 313, 361, 363, 368, and 369, but 88, 163, 294, 344, 350, and 357 are false. References notes 18, 112, 120(56) and 123.

(A minimal list of forms implying each form known to be true in the model: 6, 9, 11, 37, 63, 73, 91, 92, 122, 130, 191, 305, 309, 313, 361, 363.

A minimal list of forms at least one of which is implied by each form known to be false in the model: 88, 163, 294, 344, 350, 357.

A list of all forms known to be true in the model: 0, 5, 6, 9, 10, 11, 12, 13, 17, 18, 35, 36, 37, 38, 47, 48, 57, 63, 64, 70, 73, 74, 79, 80, 82, 83, 84, 91, 92, 93, 94, 98, 108, 111, 122, 124, 127, 128, 130, 132, 142, 145, 182, 185, 189, 190, 191, 194, 197, 198, 199, 203, 206, 211, 212, 216, 217, 221, 222, 223, 249, 250, 251, 252, 272, 273, 280, 288, 289, 296, 304, 305, 306, 307, 309, 313, 325, 327, 336, 342, 358, 361, 362, 363, 364, 366, 367, 368, 369, 371, 373, 374, 376, 377, 378, (383).

A list of all forms known to be false in the model: 1, 8, 14, 15, 16, 20, 21, 23, 27, 28, 29, 30, 31, 32, 39, 40, 43, 44, 49, 50, 51, 60, 61, 62, 66, 67, 68, 85, 86, 87, 88, 89, 90, 95, 101, 106, 107, 109, 112, 113, 114, 123, 126, 133, 149, 150, 163, 165, 174, 181, 188, 192, 193, 201, 202, 213, 218, 231, 255, 256, 257, 258, 259, 260, 261, 262, 264, 270, 284, 286, 292, 294, 295, 303, 317, 322, 323, 324, 326, 328, 329, 330, 331, 332, 333, 334, 335, 343, 344, 345, 346, 347, 349, 350, 351, 352, 353, 354, 355, 356, 357, 359, (383).)

\noindent $\mathcal{N}42(p)$: Hickman's Model IV. This model is an extension of $\mathcal{N}32$. Let $A$ be the disjoint union of two denumerable sets $B_{0}$ and $B_{1}$, where $B_{0}$ and $B_{1}$ are both Abelian $p$-groups (each non-identity element has order \emph{p}) for some prime $p\geq 5$. Let $\mathcal{G}$ be the group of all permutations $\phi$ of $A$ such that for $i=0, 1,\,\phi(B_{i})=B_{i}$ and $\phi$ is an automorphism on $B_{i}$. Moreover, for each $\phi\in \mathcal{G}$ and each $b\in B_{0},\,(\ast) \phi(f(b))=f(\phi(b))$, where $f$ is a bijection from $B_{0}$ onto $B_{1}$ in the outer model which is proved to exist if $p\geq 5$. $S$ is the set of all finite subsets of $A$. Hickman\index{Author}{Hickman} proves using condition $(\ast)$ that $f$ is in $\mathcal{N}42(p)$ so that $|B_{0}|=|B_{1}|$, but $B_{0}$ and $B_{1}$ are not isomorphic ($220(p)$ is false).

\noindent $\mathcal{N}42(p)\models$ 6, 11, 37, 73, 91, 92, 130, 191, 273, 305, 309, 313, 361, 363, 368, and 369, but 15, 192, and 220$(p)$, for $p\geq 5$, are false. References Hickman \cite{bib:1977b}, note 18.

(A minimal list of forms implying each form known to be true in the model: 6, 11, 37, 73, 91, 92, 130, 191, 305, 309, 313, 361, 363.

A minimal list of forms at least one of which is implied by each form known to be false in the model: 15, 192, 220.

A list of all forms known to be true in the model: 0, 5, 6, 11, 12, 13, 35, 36, 37, 38, 70, 73, 74, 79, 91, 92, 93, 94, 108, 130, 142, 145, 182, 189, 190, 191, 194, 197, 199, 203, 206, 211, 212, 222, 223, 251, 252, 272, 273, 280, 289, 305, 306, 307, 309, 313, 361, 362, 363, 364, 366, 367, 368, 369, 371, (383).

A list of all forms known to be false in the model: 1, 15, 192, 220, (383).)

\noindent $\mathcal{N}43$: Brunner's Model II. The set of atoms $A=\bigcup\{P_{n}:n\in\omega\}$, where $|P_{n}|=n+1$ for each $n\in \omega$ and the $P_{n}$'s are pairwise disjoint. $\mathcal{G}$ is the group of all permutations $\phi$, on $A$ that $\phi(P_{n})=P_{n}$, for each $n\in \omega$. $S$ is the set of all subsets, $E$, of $A$ such that $\sup\{|E\cap\,P_{n}|:n\in \omega\}$ is finite. Brunner\index{Author}{Brunner} proves that if $X=\prod\{P_{n}:n\in \omega\}$ and $F$ is an ultrafilter on $\omega$, then the ultraproduct $X/F$ is amorphous (form 64 is false) even though every Dedekind finite subset of $X$ is finite. It is shown in note 122 that both $C(\aleph_{0},2)$ (80) and $MC(\infty,\, \infty)$ (67) are false.

\noindent $\mathcal{N}43\models$ 6, 11, 37, 73, 91, 92, 130, 191, 273, 305, 309, 313, 361, 363, 368, and 369, but 64, 67, 80, and 344 are false. References Brunner \cite{bib:1986}, notes 18, 120(56), and 122.

(A minimal list of forms implying each form known to be true in the model: 6, 11, 37, 73, 91, 92, 130, 191, 305, 309, 313, 361, 363.

A minimal list of forms at least one of which is implied by each form known to be false in the model: 64, 80, 344.

A list of all forms known to be true in the model: 0, 5, 6, 11, 12, 13, 35, 36, 37, 38, 70, 73, 74, 79, 91, 92, 93, 94, 108, 130, 142, 145, 182, 189, 190, 191, 194, 197, 199, 203, 206, 211, 212, 222, 223, 251, 252, 272, 273, 280, 289, 305, 306, 307, 309, 313, 361, 362, 363, 364, 366, 367, 368, 369, 371, (383).

A list of all forms known to be false in the model: 1, 2, 3, 4, 7, 8, 9, 10, 14, 15, 16, 17, 20, 21, 23, 27, 28, 29, 30, 31, 32, 39, 40, 41, 43, 44, 45, 49, 50, 57, 60, 61, 62, 64, 66, 67, 68, 71, 80, 82, 83, 85, 86, 87, 88, 95, 100, 101, 106, 107, 109, 111, 113, 121, 122, 123, 126, 129, 132, 133, 149, 150, 151, 152, 153, 161, 165, 168, 174, 181, 188, 192, 193, 201, 202, 213, 214, 215, 218, 231, 250, 255, 256, 257, 258, 259, 260, 261, 262, 264, 270, 284, 286, 295, 296, 303, 317, 322, 323, 324, 325, 326, 327, 328, 331, 332, 333, 334, 335, 336, 343, 344, 345, 346, 347, 352, 354, 359, 376, 377, 378, 380, (383).)

\noindent $\mathcal{N}44$: Gross' model. $A$ is a vector space over a finite field with basis $B=\bigcup_{i\in \omega}B_{i}$ where the $B_{i}$ are pairwise disjoint and $|B_{i}|=4$ for each $i\in \omega$. Let $B_{i}= \{b_{i1},\, \ldots,\,b_{i4}\}.\,G$ is the group of vector space automorphisms $\phi$ of $A$ such that there is a permutation $\pi$ of $\omega$ and for each $i\in \omega$ there is a permutation $\alpha_{i}$ of \{1, 2, 3, 4\} in the group generated by the permutations (1, 2), (3, 4), and (1, 3)(2, 4) such that for all $i\in \omega$ and $j\in\{1,2,3,4\},\,\phi(b_{ij})=b_{\pi(i)\alpha_{i}(j)}$. $S$ is the set of finite subsets of $A$. In Gross\index{Author}{Gross} \cite{bib:1975} it is shown that, in $\mathcal{N}44$, there is a vector space $V$ with a basis and a subset $T\subseteq V$ such that no finite extension of $T$ is a basis and no infinite extension of $T$ is independent. (Hence, form 236 is false).

\noindent $\mathcal{N}44\models$ 6, 11, 37, 73, 91, 92, 130, 191, 273, 305, 309, 313, 361, 363, 368, and 369, but 15, 192, and 236 are false. References Gross \cite{bib:1975}, \cite{bib:1976}, and note 18.

(A minimal list of forms implying each form known to be true in the model: 6, 11, 37, 73, 91, 92, 130, 191, 305, 309, 313, 361, 363.

A minimal list of forms at least one of which is implied by each form known to be false in the model: 15, 192, 236.

A list of all forms known to be true in the model: 0, 5, 6, 11, 12, 13, 35, 36, 37, 38, 70, 73, 74, 79, 91, 92, 93, 94, 108, 130, 142, 145, 182, 189, 190, 191, 194, 197, 199, 203, 206, 211, 212, 222, 223, 251, 252, 272, 273, 280, 289, 305, 306, 307, 309, 313, 361, 362, 363, 364, 366, 367, 368, 369, 371, (383).

A list of all forms known to be false in the model: 1, 15, 109, 192, 236, (383).)

\noindent $\mathcal{N}45(p)$: Howard/Rubin Model III. Let $p$ be a prime. The set of atoms $A$ is a disjoint union,
\begin{equation*}
A=\bigcup_{i=1}^{p}B_{i}\cup\bigcup_{i=1}^{p}C_{i}\cup D
\end{equation*}
where $B_{i}=\{b_{n,\alpha}^{i}:n\,\in\omega\,\wedge\alpha<\aleph_{1}\},\, C_{i}=\{c_{n,\alpha}^{i}:n\in\omega\wedge\alpha<\aleph_{1}\}$, and $D=\{d_{n,\alpha}:n\in \omega\wedge\alpha<\aleph_{1}\}$. For each $\lambda<\aleph_{1}$, let
\begin{equation*}
A_{\lambda}=\bigcup_{i=1}^{p}B_{i,\lambda}\cup\bigcup_{i=1}^{p}C_{i,\lambda}\cup D_{\lambda}
\end{equation*}
where $B_{i,\lambda}=\{b_{n,\lambda}^{i}:n\in \omega\},\,C_{i,\lambda}=\{c_{n,\lambda}^{i}:n\in \omega\}$ and $D_{\lambda}=\{d_{n,\lambda}:n\in \omega\}$. Let $\phi_{\lambda}$ be the permutation of $A$ which is the product of the cycles,
\begin{equation*}
\phi_{\lambda}=\prod_{n\in \omega}(b_{n,\lambda}^{0},\,b_{n,\lambda}^{1},\,\cdots,\, b_{n,\lambda}^{p-1})(c_{n,\lambda}^{0},\,c_{n,\lambda}^{1},\,\cdots,\,c_{n,\lambda}^{p-1}).
\end{equation*}
(Note that $\phi_{\lambda}(d_{n,\lambda})=d_{n,\lambda}$ for each $n\in\omega$ and $\phi_{\lambda}(a)=a,$ for each $a\in A-A_{\lambda}.)$ The group $\mathcal{G}$ is the group generated by the $\phi_{\lambda}$ for $\lambda<\aleph_{1}$ and $S$ is the set of all finite subsets of $A$.

An elementary Abelian $p$-group is an Abelian group such that every nonidentity element has order $p$. (Note that an elementary Abelian $p$-group is just a vector space over the $p$-element field $\{0, 1,\, \cdots,p-1\}.)$ Let $G_{1}$ be the vector space over the $p$-element field $\{0, 1,\,\cdots,p-1\}$ with basis $\bigcup_{i=1}^{p}B_{i,\lambda}$ and $G_{2}$ be the vector space over the same $p$-element field with basis $\bigcup_{i=1}^{p}C_{i}\cup D$. It is shown that $|G_{1}|=|G_{2}|$, but $G_{1}$ and $G_{2}$ are not isomorphic. Thus, form 220$(p)$ is false in this model. (This extends the result of Hickman\index{Author}{Hickman} \cite{bib:1977b} to $p=2$ and 3. See $\mathcal{N}42(p)$.)

\noindent $\mathcal{N}45(p)\models$ 6, 11, 37, 73, 91, 92, 130, 191, 273, 305, 309, 313, 361, 363, 368, and 369, but 15, 192, and 220$(p)$ are false. References Howard/Rubin\index{Author}{Howard/Rubin} \cite{bib:1995a}, and note 18.

(A minimal list of forms implying each form known to be true in the model: 6, 11, 37, 73, 91, 92, 130, 191, 305, 309, 313, 361, 363.

A minimal list of forms at least one of which is implied by each form known to be false in the model: 15, 192, 220.

A list of all forms known to be true in the model: 0, 5, 6, 11, 12, 13, 35, 36, 37, 38, 70, 73, 74, 79, 91, 92, 93, 94, 108, 130, 142, 145, 182, 189, 190, 191, 194, 197, 199, 203, 206, 211, 212, 222, 223, 251, 252, 272, 273, 280, 289, 305, 306, 307, 309, 313, 361, 362, 363, 364, 366, 367, 368, 369, 371, (383).

A list of all forms known to be false in the model: 1, 15, 192, 220, (383).)

\noindent $\mathcal{N}46$: H\"{o}ft/Howard\index{Author}{Hoft/Howard@H\"{o}ft/Howard} Model I. The set $A$ of atoms is denumerable and ordered by $\leq$ so that $(A,\, \leq)$ is order isomorphic to the rationals. $A$ is written as the union $A=D_{1}\cup D_{2}\cup D_{3}$ of three dense disjoint subsets. $G$ is the group of all order automorphisms $\phi$ of $A$ such that $\phi(D_{i})=D_{i}$ for $i=1,2,3$. $S$ is the set of all subsets $E$ of $A$ which satisfy the following three conditions: (a) $E\cap D_{1}$ is finite, (b) $E\cap D_{2}$ is well ordered by $\leq$ and (c) If $b:\alpha\rightarrow E\cap D_{2}$ is an order preserving bijection from an ordinal $\alpha$ onto $E\cap D_{2}$ and $\lambda\leq\alpha$ is a limit ordinal, then the least upper bound of $\{b(\gamma):\gamma<\lambda\}$ in $(A,\,\leq)$ exists and is in $D_{3}$. In H\"{o}ft/Howard \cite{bib:1994}, it is shown that in $\mathcal{N}46,\,D_{1}$ is an infinite linearly orderable set which is Dedekind finite (form 185 is false) and that $\forall X$, if $\mathcal{P}(X)$ is Dedekind finite then every subset of $\mathcal{P}(X)$ which is linearly ordered by $\subseteq$ has a maximum element (form 84 is true).

\noindent $\mathcal{N}46 \models$ 6, 11, 37, 73, 84, 91, 92, 130, 191, 273, 305, 309, 313, 361, 363, 368, and 369, but 15 and 185 are false. References H\"{o}ft/Howard \cite{bib:1994}, note 18.

(A minimal list of forms implying each form known to be true in the model: 6, 11, 37, 73, 84, 91, 92, 130, 191, 305, 309, 313, 361, 363.

A minimal list of forms at least one of which is implied by each form known to be false in the model: 15, 185.

A list of all forms known to be true in the model: 0, 5, 6, 11, 12, 13, 35, 36, 37, 38, 70, 73, 74, 79, 84, 91, 92, 93, 94, 108, 130, 142, 145, 182, 189, 190, 191, 194, 197, 199, 203, 206, 211, 212, 222, 223, 251, 252, 272, 273, 280, 289, 305, 306, 307, 309, 313, 361, 362, 363, 364, 366, 367, 368, 369, 371, (383).

A list of all forms known to be false in the model: 1, 2, 3, 4, 7, 8, 9, 15, 39, 40, 41, 43, 44, 51, 66, 67, 71, 77, 86, 87, 89, 90, 95, 100, 101, 106, 109, 112, 113, 114, 126, 129, 133, 149, 152, 161, 168, 174, 181, 185, 188, 192, 193, 202, 214, 218, 255, 256, 257, 258, 259, 260, 261, 262, 264, 286, 292, 316, 328, 333, 334, 335, 345, 347, 359, (383).)

\noindent $\mathcal{N}47$: H\"{o}ft/Howard Model II. This model is similar to $\mathcal{N}33$. The atoms $A$ are ordered by $\leq$ so that they have order type that of the real numbers $\mathbb{R} (|A|=2^{\aleph_{0}})$. $\mathcal{G}$ is the group of all order automorphisms of $(A,\,\leq)$ and $S$ is the set of bounded subsets of $A$. It is shown in H\"{o}ft/Howard \cite{bib:1994} that in $\mathcal{N}47$ a linear ordering with no infinite descending sequences is a well ordering (form 77 is true) and that $(A,\,\leq)$ has the fixed point property but is not complete (since $A$ has no least upper bound). It follows that form 316 is false.

\noindent $\mathcal{N}47\models$ 6, 11, 37, 73, 77, 91, 92, 130, 191, 273, 305, 309, 313, 361, 363, 368, and 369, but 15, 192, and 316 are false. References H\"{o}ft/Howard\index{Author}{Hoft/Howard@H\"{o}ft/Howard} \cite{bib:1994} and note 18.

(A minimal list of forms implying\index{Author}{Brunner} each form known to be true in the model: 6, 11, 37, 73, 77, 91, 92, 130, 191, 305, 309, 313, 361, 363.

A minimal list of forms at least one of which is implied by each form known to be false in the model: 15, 192, 316.

A list of all forms known to be true in the model: 0, 5, 6, 11, 12, 13, 35, 36, 37, 38, 70, 73, 74, 77, 79, 84, 91, 92, 93, 94, 108, 130, 142, 145, 182, 185, 189, 190, 191, 194, 197, 199, 203, 206, 211, 212, 222, 223, 251, 252, 272, 273, 280, 289, 305, 306, 307, 309, 313, 361, 362, 363, 364, 366, 367, 368, 369, 371, (383).

A list of all forms known to be false in the model: 1, 15, 51, 66, 67, 89, 90, 95, 109, 112, 114, 133, 149, 192, 218, 256, 259, 261, 292, 316, 333, 334, 335, (383).)

\noindent $\mathcal{N}48$: Pincus' Model XI. $\mathcal{A}= (A,\,<,\,C_{0},\,C_{1},\,\ldots)$ is called an \emph{ordered colored set} (OC set) if $<$ is a linear ordering on $A$ and the $C_{i}$, for $i\in \omega$ are subsets of $A$ such that for each $a\in A$ there is exactly one $n\in \omega$ such that $a\in C_{n}$. If $a\in C_{n},\,C_{n}$ is called the \emph{color} of $a$. In this model we take the set of atoms to be a denumerable, universal, homogeneous OC set. (See Jech\index{Author}{Jech} \cite{bib:1973b} p101 for definitions.) It follows that $<$ is a dense unbounded linear ordering such that between any two distinct elements of $A$ there are elements of every color. $\mathcal{G}$ is the group of automorphisms of $\mathcal{A}$ and $S$ is the set of all finite subsets of $A$. Pincus proves that each $C_{i}$, where $i\in \omega$, is an infinite set which is not the union of an infinite number of pairwise disjoint infinite sets (296 is false). Pincus also proves that there is a dense linear ordering on every set (295 is true) and the Boolean Prime Ideal Theorem (14) is true. Form 14 implies 49 (Every partial ordering can be extended to a linear ordering.) and 49 + 51 (Every linear ordering has a cofinal sub well ordering.) iff AC. (See Morris\index{Author}{Morris} \cite{bib:1969} and note 121.) Thus, form 51 is false. Form 154 (Tychonoff's compactness theorem for countably many $T_{2}$ spaces) is true (14 implies 154) and 43 is false (43 implies 296). Thus, it follows from Brunner\index{Author}{Brunner} \cite{bib:1983c} and note 95 that 106 (the Baire category theorem for compact Hausdorff spaces) is false (106 + 154 iff 43). Form 14 implies the Ordering Principle (30) and the Kinna-Wagner Principle (15) is false in every FM model. It is shown in Krom\index{Author}{Krom} \cite{bib:1986} that there is an orderable topological space that is not normal (118 is false). Form 14 implies 62 $(C(\infty,\,<\aleph_{0}))$ and Howard\index{Author}{Howard} \cite{bib:1973} has shown that in every FM model 62 implies 60 $(C(\infty,\,WO))$. Thus, 60 is true. It is also shown in Howard \cite{bib:1973} that in every FM model 14 implies the Weak Sikorski Theorem (317), so 317 is true. Form 30 (Ordering Principle) $+$ form 126 $(MC(\aleph_{0},\,\infty))$ implies form 8 $(C(\aleph_{0},\, \infty))$. Since 30 is true and 8 is false, it follows that 126 is false. Since form 165 $(C(WO,\,WO))$ is true, (60 implies 165) it follows from note 2(8 and 9) that 16 and 24 are also true.

\noindent $\mathcal{N}48\models$ 6, 11, 14, 16, 23, 24, 37, 60, 73, 91, 92, 130, 165, 191, 273, 295, 305, 317, 361, 363, 368, and 369, but 51, 106, 118, 126, 131, 144, 253, 296, and 328 are false. References Pincus \cite{bib:1997}, Brunner \cite{bib:1983c}, Howard \cite{bib:1973}, Krom \cite{bib:1986}, notes 2(3, 8, 9), 18, 95, and 120(28, 32, 34, 45, and 49).

(A minimal list of forms implying each form known to be true in the model: 11, 16, 23, 24, 60, 73, 91, 92, 130, 165, 191, 295, 305, 317, 361, 363.

A minimal list of forms at least one of which is implied by each form known to be false in the model: 51, 118, 131, 144, 253, 296.

A list of all forms\index{Author}{Krom} known to be true in the model: 0, 5, 6, 10, 11, 12, 13, 14, 16, 18, 19, 23, 24, 25, 26, 27, 28, 30, 31, 32, 33, 34, 35, 36, 37, 38, 45, 46, 47, 48, 49, 52, 60, 61, 62, 63, 64, 69, 70, 72, 73, 74, 79, 80, 83, 85, 88, 91, 92, 93, 94, 96, 99, 102, 104, 107, 108, 111, 119, 120, 121, 122, 123, 127, 130, 132, 137, 139, 140, 141, 142, 145, 146, 150, 151, 153, 154, 165, 178, 182, 189, 190, 191, 194, 197, 198, 199, 201, 203, 206, 207, 209, 211, 212, 213, 216, 221, 222, 223, 225, 226, 227, 228, 229, 231, 233, 235, 241, 242, 249, 250, 251, 252, 268, 269, 270, 271, 272, 273, 276, 280, 283, 285, 287, 288, 289, 293, 294, 295, 298, 305, 306, 307, 308, 309, 311, 313, 314, 315, 317, 323, 324, 326, 327, 329, 330, 331, 332, 336, 342, 343, 344, 349, 350, 356, 357, 358, 361, 362, 363, 364, 366, 367, 368, 369, 370, 371, 373, 374, 377, 378, 380, (383).

A list of all forms known to be false\index{Author}{Pincus} in the model: 1, 2, 3, 4, 7, 8, 9, 15, 39, 40, 41, 43, 44, 51, 66, 67, 71, 76, 86, 87, 89, 90, 95, 100, 101, 106, 109, 112, 113, 114, 118, 126, 129, 131, 133, 144, 149, 152, 161, 168, 174, 179, 181, 188, 192, 193, 202, 214, 218, 253, 255, 256, 257, 258, 259, 260, 261, 262, 264, 286, 292, 296, 328, 333, 334, 335, 345, 347, 359, (383).)

\noindent $\mathcal{N}49$: De la Cruz/Di Prisco Model\index{Author}{De la Cruz/Di Prisco}. Let $A=\{a(i,p):i\in \omega\wedge p\in \mathbb{Q}/\mathbb{Z}\}$. $(\mathbb{Q}/\mathbb{Z}$ denotes the additive group.) For each $i\in \omega$, let $A_{i}=\{a(i,p):p\in \mathbb{Q}/\mathbb{Z}\}$ and for each $i\in \omega$ and $p\in \mathbb{Q}/\mathbb{Z}$, let $\pi_{i}^{p}: A\rightarrow A$ be defined by $\pi_{i}^{p}(a(j,\,q))= \left\{\begin{array}{ll}
a(j,p+q) & \mathrm{if}\,j=i\\ a(j,\,q) & \mathrm{otherwise}
\end{array}\right.$. $G$ is the group generated by the permutations $\pi_{i}^{p}$ for $i\in \omega$ and $p\in \mathbb{Q}/\mathbb{Z}$. $S$ is the set of finite supports. De la Cruz and Di Prisco have shown that every infinite set in this model is Dedekind infinite (9 is true) and that $\{A_{i}:i\in \omega\}$ is a countable family of well orderable sets such that no infinite subfamily has Kinna-Wagner selection function. It follows that 167 and 380 are false. In addition De la Cruz and Di Prisco have shown that for every $n\in\omega, n\geq 2$, there is a set of $n$ elements sets in the model with no choice function (47(n) is false).

\noindent $\mathcal{N}49\models$ 6, 9, 11, 37, 63, 73, 91, 92, 130, 191, 305, 309,
313, 361, and 363, but 47(n), 106, 163, 167, 344 and 380 are false. References De la Cruz/Di Prisco \cite{bib:1996}, notes 18 and 120(20, 23, 55 and 56).

(A minimal list of forms implying each form known to be true in the model: 6, 9, 11, 37, 63, 73, 91, 92, 130, 191, 305, 309, 313, 361, 363.

A minimal list of forms at least one of which is implied by each form known to be false in the model: 47, 106, 163, 167, 344, 380.

A list of all forms known to be true in the model: 0, 5, 6, 9, 10, 11, 12,
13, 17, 18, 35, 36, 37, 38, 57, 63, 64, 70, 73, 74, 79, 80, 82, 83, 84, 91, 92, 93, 94, 98, 108, 124, 127, 128, 130, 132, 142, 145, 182, 185, 189, 190, 191, 194, 197, 198, 199, 203, 206, 211, 212, 216, 217, 221, 222, 223, 249, 251, 252, 272, 273, 280, 288, 289, 296, 304, 305, 306, 307, 309, 313, 325, 336, 342, 358, 361, 362, 363, 364, 366, 367, 368, 369, 371, 373, 374, 376, 377, 378, (383).

A list of all forms known to be false in the model: 1, 8, 14, 15, 20, 21, 23, 30, 33, 39, 40, 43, 44, 45, 47, 49, 50, 60, 61, 62, 66, 67, 68, 85, 86, 87, 95, 101, 106, 107, 109, 113, 121, 122, 123, 133, 149, 151, 163, 165, 167, 174, 181, 188, 192, 193, 202, 213, 218, 231, 250, 255, 256, 257, 258, 259, 260, 261, 262, 264, 270, 284, 286, 295, 303, 317, 322, 323, 324, 327, 331, 332, 333, 334, 335, 343, 344, 345, 346, 347, 359, 380, (383).)

\noindent $\mathcal{N}50(E)$: Brunner's Model III. $E$ is a finite set of prime numbers. For each $p\in E$ and $n\in \omega$, let $A_{p,n}$ be a set of atoms of cardinality $p^{n}$. $A$ is the disjoint union $A=\bigcup_{p\in E,n\in\omega}A_{p,n}$. For each $p\in E$ and $n\in \omega$, let $\eta_{p,n}$ be a permutation of $A_{p,n}$ consisting of a single cycle of length $p^{n}$. For each $p\in E$ let $\eta_{p}$ be the permutation of $A$ defined by $\eta_{p}(a)=\eta_{p,n}(a)$ if $a\in A_{p,n}$ and $\eta_{p}(a)=a$ for $a\not\in\bigcup_{n\in\omega}A_{p,n}$.
$\mathcal{G}$ is the group of permutations of $A$ generated by $\{\eta_{p}:p\in E\}$. $S$ is the ideal of finite subsets of $A$. This is the model of the lemma on page 72 of Brunner\index{Author}{Brunner} \cite{bib:1990}. In Brunner \cite{bib:1990} it is shown that $MC(\infty,\,\infty)$ (form 67) and $KW(WO, \infty)$ (form 322) are true and that $PC(\aleph_{0},\,<\aleph_{0},\,\infty)$ (form [10
E]) and $PC(\infty, n,\,\infty)$ (form $45(n))$ (for $n=\sum\nolimits_{p\in E}n_{p}\cdot p$ where $n_{p}\in \omega$ for $p\in E$) are false.

\noindent $\mathcal{N}50(E)\models$ 6, 11, 37, 67, 73, 92, 130, 191, 273, 305, 313,
322, 361, and 363 but 10, 154, 164, 342(n) (for $n$ as described above) and 344 are false. References Brunner \cite{bib:1990}, notes 18 and 120(2, 10, and 56).

(A minimal list of forms implying each form known to be true in the model: 6, 11, 37, 67, 73, 92, 130, 191, 305, 313, 322, 361, 363.

A minimal list of forms at least one of which is implied by each form known to be false in the model: 10, 154, 164, 342, 344.

A list of all forms known to be true in the model: 0, 5, 6, 11, 12, 13, 18,
19, 25, 34, 35, 36, 37, 38, 47, 48, 51, 52, 58, 64, 67, 70, 73, 74, 76, 77, 78, 79, 80, 82, 83, 84, 89, 90, 91, 92, 93, 94, 104, 106, 108, 111, 112, 114, 118, 119, 125, 126, 127, 130, 131, 142, 144, 145, 147, 155, 173, 182, 185, 189, 190, 191, 194, 197, 199, 203, 206, 208, 211, 212, 221, 222, 223, 232, 250, 251, 252, 254, 272, 273, 280, 287, 288, 289, 292, 294, 305, 306, 307, 309, 313, 315, 316, 322, 324, 327, 328, 329, 330, 349, 350, 355, 357, 358, 361, 362, 363, 364, 366, 367, 368, 369, 371, 373, 375, 381, 382, 383, (383).

A list of all forms known to be false in the model: 1, 2, 3, 4, 7, 8, 9, 10, 14, 15, 16, 17, 20, 21, 23, 27, 28, 29, 30, 31, 32, 39, 40, 41, 43, 44, 49, 50, 60, 61, 62, 68, 71, 85, 86, 87, 100, 101, 107, 109, 113, 121, 122, 123, 129, 132, 133, 150, 151, 152, 153, 154, 161, 164, 165, 168, 174, 181, 192, 202, 213, 214, 218, 231, 255, 256, 257, 258, 259, 260, 261, 262, 264, 270, 284, 286, 295, 303, 317, 323, 325, 331, 332, 342, 343, 344, 345, 346, 347, 352, 354, 359, 376, 377, 378, 380, (383).)

\noindent $\mathcal{N}51$: W\c{e}glorz\index{Author}{W\c{e}glorz}/Brunner Model. Let $A$ be denumerable and $\mathcal{G}$ be the group of all permutations of $A$. Let $\mathcal{F}$ be the normal filter generated by the groups $\mathcal{G}_{B}$, where $B\subseteq A$ and $\mathcal{G}_{B}=\{\sigma\in \mathcal{G}:\sigma(B)=B\}$. In this model there is no non-principal measure on $A$, so form 221 is false, but forms 64 (There are no amorphous sets.) and 90 (Every linearly ordered set can be well ordered.) are true. (See note 128.)

\noindent $\mathcal{N}51\models$ 6, 11, 37, 64, 73, 90, 92, 130, 191, 273, 305, 309,
313, 361, and363, but 9, 30, 49, 192, 221, and 344 are false. References W\c{e}glorz \cite{bib:1969}, notes 18, 120(3, 6, 16, 20, and 56) and 128.

(A minimal list of forms implying each form known to be true in the model: 6, 11, 37, 64, 73, 90, 92, 130, 191, 305, 309, 313, 361, 363.

A minimal list of forms at least one of which is implied by each form known to be false in the model: 9, 30, 221, 344.

A list of all forms known to be true in the model: 0, 5, 6, 11, 12, 13, 19,
25, 34, 35, 36, 37, 38, 51, 58, 64, 70, 73, 74, 77, 79, 84, 90, 91, 92, 93, 94, 104, 108, 118, 119, 127, 130, 142, 145, 182, 185, 189, 190, 191, 194, 197, 199, 203, 206, 208, 211, 212, 222, 223, 251, 252, 272, 273, 280, 289, 294, 305, 306, 307, 309, 313, 315, 316, 361, 362, 363, 364, 366, 367, 368, 369, 371, (383).

A list of all forms known to be false in the model: 1, 2, 3, 4, 7, 8, 9, 14, 15, 30, 39, 40, 41, 43, 44, 49, 50, 52, 63, 66, 67, 71, 86, 87, 95, 100, 101, 109, 113, 123, 129, 133, 149, 152, 161, 168, 174, 181, 192, 202, 214, 218, 221, 255, 256, 257, 258, 259, 260, 261, 262, 264, 286, 295, 303, 317, 333, 334, 335, 344, 345, 347, 359, 372, (383).)

\noindent $\mathcal{N}52$: Felgner/Truss Model. Let $(\mathcal{B},\,\prec)$ be a countable universal homogeneous linearly ordered Boolean algebra, (i.e., $<$ is a linear ordering extending the Boolean partial ordering on $B$). Let $A=\{a_{x}:x\in \mathcal{B}\}$; let $\mathcal{G}$ be the group of all permutations on $A$ which preserve the structure (a linearly ordered Boolean algebra); and let $S$ be the set of finite subsets of $A$. It is shown in Felgner/Truss \cite{bib:1996} that $A$ has no prime ideal so form 14 is false. It is also shown that every partial ordering can be extended to a linear ordering so form 49 is true.

\noindent $\mathcal{N}52\models$ 6, 11, 37, 49, 73, 91, 92, 130, 191, 305, 309, 313,
361, and 363, but 14, 15, 51, 99, 118, 192, and 253 are false. References Felgner/Truss\cite{bib:1996}, notes 18, 120(3, 16, 29, and 45).

(A minimal list of forms implying each form known to be true in the model: 6, 11, 37, 49, 73, 91, 92, 130, 191, 305, 309, 313, 361, 363.

A minimal list of forms at least one of which is implied by each form known to be false in the model: 15, 51, 99, 118, 192, 253.

A list of all forms known to be true in the model: 0, 5, 6, 10, 11, 12, 13,
18, 30, 33, 35, 36, 37, 38, 45, 46, 47, 48, 49, 61, 62, 64, 70, 73, 74, 79, 80, 83, 88, 91, 92, 93, 94, 102, 108, 111, 120, 121, 122, 127, 130, 132, 140, 142, 145, 146, 178, 182, 189, 190, 191, 194, 197, 198, 199, 201, 203, 206, 211, 212, 216, 222, 223, 249, 250, 251, 252, 268, 269, 272, 273, 276, 280, 283, 285, 288, 289, 293, 305, 306, 307, 308, 309, 313, 326, 327, 336, 342, 358, 361, 362, 363, 364, 366, 367, 368, 369, 371, 373, 374, 377, 378, (383).

A list of all forms known to be false in the model: 1, 14, 15, 50, 51, 66, 67, 89, 90, 95, 99, 109, 112, 114, 118, 133, 149, 192, 218, 253, 256, 259, 261, 292, 303, 317, 333, 334, 335, 345, (383).)

$\mathcal{N}53$: Good/Tree/Watson\index{Author}{Good/Tree/Watson} Model I. Let $A=\bigcup\{Q_{n}:n\in \omega\}$, where $Q_{n}=\{a_{n,q}$: $q\in \mathbb{Q}\}$. Suppose $<$ is the lexicographic ordering on $A$, i.e.
\begin{equation*}
a_{n,q}<a_{m,p}\,\mathrm{if}\ n<m,\ \mathrm{or}\ n=m\ \mathrm{and}\ q<p.
\end{equation*}
The group of permutations $\mathcal{G}$, is the group of all permutations on $A$ which are a rational translation on $Q_{n}$, i.e. if $\phi\in G$, then $\phi|Q_{n}(a_{n,q})=a_{n,q+r_{n}}$ for some $r_{n}\in \mathbb{Q}$, and supports are finite. Let $d$ be the metric on $A$ given by:
\begin{equation*}
d(a_{n,q},\,a_{m,p})=1\ \mathrm{if}\ n\neq m
\end{equation*}
and
\begin{equation*}
d(a_{n,q},\,a_{m,p})=|q-p|/(1+|q-p|)\ \mathrm{if}\ n=m.
\end{equation*}
It is shown in Good/Tree/Watson \cite{bib:1996} that $(A,\,d)$ is a linearly ordered, zero dimensional, metric space, but is not paracompact (MP, form 383 is false). Each $Q_{n}$, being a second countable metric space, is paracompact, and the disjoint union topology on $A$ coincides With the metric topology which is normal. Hence, $A$ can be considered as a normal non paracompact disjoint union of paracompact spaces. Thus, DUP ([67 N]) is false. It is shown in Howard\index{Author}{Howard}/Keremedis/Rubin/Rubin \cite{bib:1997b} that MM (383), DUMN (382), and $C(\infty,\,2)$ (88) are each false, but every infinite set does have a countably infinite subset so form 9 is true.

\noindent $\mathcal{N}53\models$ 6, 9, 11, 37, 63, 73, 91, 92, 130, 273, 305, 309,
313, 361, and 363, but 88, 133, 192, 344, 382, and 383 are false. References Good/Tree/Watson\index{Author}{Good/Tree/Watson} \cite{bib:1996}, Howard/Keremedis/Rubin/Rubin\index{Author}{Howard/Keremedis/Rubin/Rubin} \cite{bib:1997b}, notes 18 and 120(4, 20, and 56).

(A minimal list of forms implying each form known to be true in the model: 6, 9, 11, 37, 63, 73, 91, 92, 130, 305, 309, 313, 361, 363.

A minimal list of forms at least one of which is implied by each form known to be false in the model: 88, 133, 192, 344, 382, 383.

A list of all forms known to be true in the model: 0, 5, 6, 9, 10, 11, 12,
13, 17, 18, 35, 36, 37, 38, 57, 63, 64, 70, 73, 74, 79, 80, 82, 83, 84, 91, 92, 93, 94, 98, 108, 124, 127, 128, 130, 132, 142, 145, 185, 194, 197, 198, 199, 203, 206, 211, 212, 216, 217, 221, 222, 223, 249, 251, 252, 272, 273, 280, 288, 289, 296, 304, 305, 306, 307, 309, 313, 325, 336, 342, 358, 361, 362, 363, 364, 366, 367, 368, 369, 371, 373, 374, 376, 377, 378, (383).

A list of all forms known to be false in the model: 1, 14, 15, 30, 49, 50, 60, 61, 62, 66, 67, 68, 85, 88, 95, 107, 109, 123, 133, 149, 192, 201, 213, 218, 232, 270, 284, 295, 303, 317, 323, 326, 331, 332, 333, 334, 335, 343, 344, 345, 346, 381, 382, 383, (383).)

\noindent $\mathcal{N}54$: Good/Tree/Watson Model II. This model is a variation of $\mathcal{N}53$. The set of atoms $A=\bigcup\{Q_{n}:\,n\in \omega\}$, where $Q_{n}=\{a_{n,q}:q\in \mathbb{R}\}$. We let $<$ be the lexicographic ordering on $A$. The group of permutations $\mathcal{G}$, is the group of all permutations on $A$ which are a translation on $Q_{n}$, that is, if $\phi\in \mathcal{G}$, then $\phi|Q_{n}(a_{n,q})=a_{n,q+r_{n}}$ for some $r_{n}\in \mathbb{R}$. Supports are finite. Let $d$ be the metric on $A$ given by:
\begin{equation*}
d(a_{n,q},\,a_{m,p})=1\,\mathrm{if}\ n\neq m
\end{equation*}
and
\begin{equation*}
d(a_{n,q},\,a_{m,p})=|q-p|/(1+|q-p|)\,\mathrm{if}\,n=m.
\end{equation*}
It is shown in Howard/Keremedis/Rubin/Stanley \cite{bib:1997} that $(A,\,d)$ is a metric space that is not meta-Lindel\"{o}f (173 is false).

\noindent $\mathcal{N}54\models$ 6, 11, 37, 73, 91, 92, 130, 273, 305, 309, 313, 361, and 363, but 15, 173, and 192 are false. References Good/Tree/Watson \cite{bib:1996}, Howard/Keremedis/Rubin/Stanley\index{Author}{Howard/Keremedis/Rubin/Stanley} \cite{bib:1997}, and note 18.

(A minimal list of forms implying each form known to be true in the model: 6, 11, 37, 73, 91, 92, 130, 305, 309, 313, 361, 363.

A minimal list of forms at least one of which is implied by each form known to be false in the model: 15, 173, 192.

A list of all forms known to be true in the model: 0, 5, 6, 11, 12, 13, 35,
36, 37, 38, 70, 73, 74, 79, 91, 92, 93, 94, 108, 130, 142, 145, 194, 197, 199, 203, 206, 211, 212, 222, 223, 251, 252, 272, 273, 280, 289, 305, 306, 307, 309, 313, 361, 362, 363, 364, 366, 367, 368, 369, 371, (383).

A list of all forms known to be false in the model: 1, 15, 66, 67, 76, 95, 109, 149, 173, 192, 214, 218, 232, 333, 334, 335, 383, (383).)

%%%%%%%%part04

\part{Notes \label{pt04:part04}}

{Table of Contents}

\begin{enumerate}
\item[0.] Open problems that don't fit into the table format

\item[1.] Independence results involving more that 2 forms

\item[2.] Implications that hold in every Fraenkel-Mostowski model

\item[3.] Results from Mathias \cite{bib:1979} and Figura \cite{bib:1981}

\item[4.] A proof that 73 implies [73A]

\item[5.] A proof that [73 A] implies 74

\item[6.] A proof that [73 A] implies 73

\item[7.] A proof that 43 implies 112

\item[8.] A proof that 9 and [9 A] are equivalent

\item[9.] References for forms 71, 86 and 87

\item[10.] A proof that [73 A] is equivalent to [73 B]

\item[11.] A proof that 9 is equivalent to [9 C]

\item[12.] A proof that 9 is equivalent to [9 D]

\item[13.] A note on the result $39\, {\not\rightarrow}\, 43$

\item[14.] A note on the result $40\, {\not\rightarrow}\,44$

\item[15.] Relations between the finite axioms of choice: $C(\infty,\,n)$, $C(LO,\,n)$,
$C(WO,\,n)$, $C(\infty,\,\in Z)$, $C(LO,\,\in Z)$, $C(WO,\,\in Z)$ and $C(\aleph_{0},\,n)$
(Forms 45, 33, 47, 46, 120, 48, and 288 respectively) and generalizations

\item[16.] Definitions from Day \cite{bib:1973} for form 126: $B$ lattice, wedge in a
linear space, ordered linear space, cone in a linear space, vector
lattice

\item[17.] Definition of $\epsilon$ lattice from Brunner \cite{bib:1984e}

\item[18.] A list of special forms including the following: 1. A list of forms
that are transferable; 2. A list of forms whose negation may not be
transferable; 3. A list of forms that are true in every permutation
model; 4. A list of forms that are false in every permutation model;
5. A list of forms that are false in every Cohen model; and 6. A
list of forms that depend on parameters

\item[19.] The equivalence of [10 F], [10 K], [10 L] and 10

\item[20.] Definitions and consistency results for large cardinal axioms

\item[21.] Definitions for [10 F], [14 M], [14 N], [43 H], [106 A], 216, and 249:
Trees, properties of trees, and families of locally finite trees

\item[22.] Definitions for [1 L], [1 M], [1 M], [10 C], and [132 A]: set mapping,
free subset (relative to a set mapping), and ramification system

\item[23.] Definitions for [14 Q] and [52 E]: Weak$^\ast$ topology on the dual of a
Banach space, convex-compact subset of a Banach space

\item[24.] Definitions from group theory used in $\mathcal{M}22$ and forms [62 C], [62
$D],\,[67 D]$, 180, and $[308 \mathrm{A}(p)]$: Frattini subgroup, $p$-group, group
of exponent $p$, elementary Abelian group, special $p$-group, extra-special $p$-group, divisible and torsion free groups

\item[25.] Definitions for forms 149 and 179: $W(\alpha)$ for $\alpha$ an ordinal, the class
of almost well ordered sets

\item[26.] A summary of topological definitions and results from Brunner
\cite{bib:1983d}

\item[27.] Notes on $UT(WO,\,\kappa,\,WO)$, $C(WO,\,\kappa)$ and $D_{\kappa}$ from Brunner/
Howard \cite{bib:1992}

\item[28.] Definitions for Baire category type theorems and a summary of
results from Brunner \cite{bib:1983c}

\item[29.] Definitions for forms [14 R]~-- [14 U], [14 BH]~-- [14 BL], and [86 B]: Local lattice, frame, complete Heyting algebra, spatial lattice,
possible properties of local lattices (compact, coherent, regular, 0
dimensional, conjunctive), irreducible, closed set in a topological
space, generic point for a closed set, sober topological space, sum
of two frames, regular ring, Gelfand algebra, frame envelope, $\aleph_{\alpha^{-}}$
frame, $\aleph_{\alpha}$-Lindel\"{o}f frame

\item[30.] A proof that forms 14, [14 BM], [14 BN], and [14 BO] are equivalent

\item[31.] Definitions for forms [14 W] and [70 A] from Schechter \cite{bib:1996a}:
Net, subnet, universal net on a set

\item[32.] Definitions for forms [159 A] and [159 B] from Blass \cite{bib:1983b}: similarity type, variety, free algebra in a variety, initial algebra in a
variety

\item[33.] A proof that $99+62$ implies 14

\item[34.] A proof that [14 AA] implies [14Z]

\item[35.] A proof that form 10 is equivalent to form [10 F]

\item[36.] A proof that 284 implies 61

\item[37.] Definitions for [14 AB]: Sober topological space

\item[38.] Definitions from Manka \cite{bib:1988a} and \cite{bib:1988b}: Supremal subset of
a poset, sup function, upper semi-continuous function

\item[39.] Summary of definitions and results from Harper/Rubin \cite{bib:1976} having to do with possible properties of relations

\item[40.] A proof that 94, [94 A], and [94 B] are equivalent

\item[41.] A proof that 233 is true in $\mathcal{N}1$

\item[42.] Conditions under which a topological space is well orderable from
Brunner \cite{bib:1982a}

\item[43.] Topological definitions and results from Brunner \cite{bib:1982b}

\item[44.] A note on $57+62{\not\rightarrow} 9$

\item[45.] References to consult for set theory without the axiom of choice

\item[46.] Ramsey's Theorem, form 325

\item[47.] Definitions from Shannon \cite{bib:1990}: Generic filter in a partial ordering

\item[48.] The equivalence of 34 and [34 A]

\item[49.] Transitivity conditions from von Rimscha \cite{bib:1982}

\item[50.] Definitions from universal algebra (including multi-algebras) from
H\"{o}ft/Howard \cite{bib:1981} and Andreka/Nemeti \cite{bib:1981}

\item[51.] An incorrect proof that [10 F] implies 17

\item[52.] A result from Brunner \cite{bib:1981a} on the box product of topological
spaces

\item[53.] Relationships between versions of 178 from Zuckerman \cite{bib:1969a}
and \cite{bib:1981}

\item[54.] The equivalence of 43, [43 R], and [43 S]

\item[55.] Results from Gitik \cite{bib:1980} and \cite{bib:1985} on cofinalities of ordinals
without the axiom of choice

\item[56.] Generalizations of amorphous sets from Hickman \cite{bib:1980b}

\item[57.] The structure of amorphous sets. (See Truss \cite{bib:1995}.)

\item[58.] Results on Dedekind finite cardinals from Hickman \cite{bib:1979b}

\item[59.] Notes on form 191 from Blass \cite{bib:1979}

\item[60.] Definitions from category theory

\item[61.] Notes from Kanovei \cite{bib:1979} on $AC(K)$ and $DC(K)$ where $K$ is
either in the analytic or projective hierarchy

\item[62.] Definitions from Harazisville (see Kharazishvili \cite{bib:1979}): General
linear system and linear global reaction

\item[63.] References for class forms of $AC$

\item[64.] Notes from Dawson/Howard \cite{bib:1976} and Pincus \cite{bib:1978} on form 200

\item[65.] Notes from Truss \cite{bib:1978} on $\mathcal{M}2$

\item[66.] A proof that 322 is true in $\mathcal{N}1$

\item[67.] A proof that 203 implies 94

\item[68.] A proof that 3 implies 53

\item[69.] A proof that [40 B] implies 208

\item[70.] A proof that [62 E] implies 62

\item[71.] Definitions from Morillon \cite{bib:1988}

\item[72.] A proof that $207(\beta)$ implies 209

\item[73.] Results from Pincus \cite{bib:1977a} on adding dependent choice (43) to
models of ZF

\item[74.] Results from Pincus \cite{bib:1977a} on $C(\infty,\,<\aleph_{\alpha})$ and $C(\infty,\,\aleph_{\alpha})$

\item[75.] Forms equivalent to AC in ZF

\item[76.] Proofs that 126 implies 82 and 185

\item[77.] Definitions for [1 BL]~-- [1 BR], [1 CB], and [67 G]

\item[78.] Extending filters to ultrafilters and 99

\item[79.] Definitions for 227 and 228: Positive cone in a partially ordered
group, fully ordered group

\item[80.] Proofs that [14 AL], [14 AM], [14 AO], and [14 AP] imply 14

\item[81.] A proof that 30 implies 83

\item[82.] Definitions for 230 from Myhill/Scott \cite{bib:1971}: $L$ and $HOD$

\item[83.] A proof that [14 AR] implies 14

\item[84.] Standard definitions from abstract algebra

\item[85.] Definition of the monadic theory $MT(\omega_{1},\,<)$ from 246

\item[86.] Properties of Boolean algebras not provable in ZF from Plotkin
\cite{bib:1976}

\item[87.] Weakenings of K\"{o}nig's lemma from Truss \cite{bib:1975}

\item[88.] A proof that 31 does not imply 9

\item[89.] Definition of a free complete Boolean algebra for 267 and [0Y]

\item[90.] Results on $81(n)$ from Monro \cite{bib:1972}

\item[91.] A proof that 269 is false in $\mathcal{N}l$

\item[92.] Relationships between $271(m)$ and $45(n)$ from Kolany/Wojtylak
\cite{bib:1991}

\item[93.] A proof that 334 is true in $\mathcal{N}2$ and that 333 is true in $\mathcal{N}2^{\ast}(3)$

\item[94.] Relationships between definitions of finite

\item[95.] A proof that 154 is not a theorem of $\mathrm{ZF}^{0}$

\item[96.] Some topological consequences of 43 and $\neg 280$ and definitions for 279 and 281 from Wright \cite{bib:1973}: Locally convex topological vector space, invariant metric on a vector space, Fr\'{e}chet space

\item[97.] Definitions for forms 282 and [0 AD] and related results from Kleinberg/Seiferas \cite{bib:1973}: $\kappa \rightarrow(\beta)_{\gamma}^{\alpha}$

\item[98.] Definitions for forms [1 AK], [1 AL] and [1 AM] from Truss \cite{bib:1973a}: $\aleph(x)$, and $\aleph(x,\,y)$

\item[99.] A proof that 65 does not imply 14

\item[100.] A proof that 88 is false and 8 is true in $\mathcal{N}23$

\item[101.] A proof that 43 implies [43 S]

\item[102.] A proof that 289 is false in $\mathcal{M}1$

\item[103.] A summary of the transfer results from Jech/Sochor \cite{bib:1966a} and
\cite{bib:1966b} and Pincus \cite{bib:1972a}, \cite{bib:1977a} and \cite{bib:1971}

\item[104.] A condition insuring the failure of 3 in a permutation model

\item[105.] Proofs that 64 is true and 83 false in $\mathcal{N}4,\ 83$ is true and 215 is false
in $\mathcal{N}2,\ 216$ is true and 217 false in $\mathcal{N}34,\ 217$ is true and 9 is false
in $\mathcal{N}1$

\item[106.] Results from McCarten \cite{bib:1988} and Schnare \cite{bib:1968} relating to
forms [1 BE], [1 BF], [1 BG], and [1 BH]

\item[107.] Proofs that each of [34 B], [34 C] [34 D], and [34 E] is equivalent to 34

\item[108.] A proof that 187 is true in $\mathcal{N}12(\aleph_{1})$

\item[109.] Definitions for the conditional choice forms

\item[110.] A proof that 36 implies 38

\item[111.] Relationships between the forms $336(n)$ and $342(n)$

\item[112.] A proof that 9 and 122 are true in $\mathcal{N}41$

\item[113.] A proof that 88 is false in $\mathcal{N}10$

\item[114.] Definitions from Bell \cite{bib:1988}: Irreducible function on topological
spaces, Boolean topological space, extremally disconnected space, Gleason cover

\item[115.] Definitions from Diener \cite{bib:1989}: Algebra of type $\tau$, closed homomorphism between algebras, Peano algebra

\item[116.] A proof that 304 is true in $\mathcal{N}1$

\item[117.] Definitions from Cowen \cite{bib:1990} for forms [14 BD]~-- [14 BG]: Hypergraph, subhypergraph, independent vertex set, vertex cover,
$n$-coloring

\item[118.] Definitions from Galvin/Komj\'{a}th \cite{bib:1991} for [1 AT]~-- [1 AW]: Chromatic number and irreducible coloring of a graph

\item[119.] Definitions related to partial orders and quasi orders from Shannon \cite{bib:1992} for [86$(\alpha)$ A]

\item[120.] Some relationships between forms which are provable in $\mathrm{ZF}^{0}$ and are not readily available from table 1

\item[121.] A proof that $49+51$ implies AC

\item[122.] A proof that 67 and 80 are false in $\mathcal{N}43$

\item[123.] Proofs that $10+163$ implies 231 and that 133 implies 231

\item[124.] A proof that 314 implies 119

\item[125.] Forms of the Hahn-Banach theorem from Buskes \cite{bib:1993}

\item[126.] Definitions from logic for [43 B] from Goldblatt \cite{bib:1985}

\item[127.] Proofs that 346 implies [62 E] and [1 BZ] implies 1

\item[128.] A proof that forms 64 and 90 are true and form 221 is false in $\mathcal{N}51$

\item[129.] Definitions from free groups for forms [1 BB], [1 CA], 68, and 348

\item[130.] A proof that 30 and [30 G] are equivalent

\item[131.] A proof of [0 AL] in $\mathrm{ZF}^{0}$

\item[132.] Definitions and results from Keremedis \cite{bib:1996d}

\item[133.] Proofs that [1 F] and [1 CC] are equivalent to 1

\item[134.] A proof that 100 implies 369

\item[135.] A proof that 18 is equivalent to [18 B]

\item[136.] A proof that 14 is true in $\mathcal{N}29$

\item[137.] A proof that [1 BU] implies AC

\item[138.] A proof that 375 implies 119

\item[139.] A proof that 67 implies 375

\item[140.] A proof that 214 implies 152

\item[141.] Definitions for forms [1 CE]~-- [1 CH], [67 K]~-- [67 T], [126 G]~-- [126 K], 173, [173 A], 232, [232 A]~-- [232 G], 383, and [383 A]~-- [383 C]

\item[142.] A proof that 234 implies 282

\item[143.] A proof that [0 AM] is provable in $\mathcal{ZF}^{0}$

\item[144.] A proof that 43 holds in FM models with countable supports

\item[145.] A proof that 287 implies 222

\item[146.] A proof that [43 W] is equivalent to 43

\item[147.] Definitions for forms involving measures on Boolean algebras

\item[148.] A proof that 154 and [154 A] are equivalent
\end{enumerate}

\addtocontents{toc}{\SkipTocEntry}
\chapter{Notes}\label{part04}

\section*{NOTE 0}
\begin{enumerate}
\item[ ] Open problems\index{Subject}{open problems} that don't conveniently fit into the table format.
\item[1.] Does ZF\index{Subject}{implications true in every FM model} $\vdash$ (form 43 (DC) + form 85 $(C(\infty,\,\aleph_{0}))\rightarrow$ form 213 $(C(\infty,\, \aleph_{1}))$? Pincus\index{Author}{Pincus}\index{Author}{Pineus/Solovay} \cite{bib:1977a}.

\item[2.] Does ZF $\vdash$ (form 40 $(C(WO,\,\infty)$) + form 30 (OP) ) $\rightarrow$ AC? Pincus \cite{bib:1977a}.
\item[3.] Does ZF $\vdash$ form 91 (PW) + form 132 $PC(\infty,\,<\aleph_{0},\,\infty)\rightarrow$ AC? Brunner\index{Author}{Brunner} \cite{bib:1984b}.
\item[4.] Does ZF $\vdash$ form 336(5) (Every infinite set has an infinite subset $Y$ such that $\{t\in Y:|t|=5\}$ has a choice function) $\rightarrow$ form 342(5) $(PC(\infty,\,5,\,\infty))$? Note 111.
\end{enumerate}

\section*{NOTE 1}

In this note we list independence results\index{Subject}{independence results} which are not readily available from table 1. In each case we give the form numbers, the relationship, the reference and the model or models.
\begin{enumerate}
\item[1.] $14+43\not\rightarrow 1$, (ZF) Mathias\index{Author}{Mathias} \cite{bib:1979}, p 160 (due to Pincus), $\mathcal{M}40(\kappa),\, \mathcal{N}40$.
\item[2.] $52 + 65 \not\rightarrow 1\ (\mathrm{ZF}^{0})$, Pincus \cite{bib:1972b},
$\mathcal{N}2$.
\item[3.] $14+295 \not\rightarrow 296$ (ZF), Pincus \cite{bib:1997}, $\mathcal{N}48\mathrm{T}$.
\item[4.] $14+295\not\rightarrow 15$ (ZF), Pincus \cite{bib:1997}, $\mathcal{N}48\mathrm{T}$.
\item[5.] $14+15\not\rightarrow 1$ (ZF), Pincus \cite{bib:1997}, $\mathcal{M}1$.
\item[6.] $130+18 \not\rightarrow 128$ (ZF), Brunner \cite{bib:1984f}, $\mathcal{N}2\mathrm{T}$.
\item[7.] $130+18+128 \not\rightarrow 9$ (ZF), Brunner \cite{bib:1984f},
$\mathcal{N}3\mathrm{T}$.
\item[8.] $18 + 128 \not\rightarrow 13$ (ZF), Brunner \cite{bib:1984f}, $\mathcal{M}1$.
\item[9.] $163+67 \not\rightarrow 10\ (\mathrm{ZF}^{0})$, Brunner \cite{bib:1982a},
$\mathcal{N}2$.
\item[10.] $163+ 67 \not\rightarrow 164\ (\mathrm{ZF}^{0})$, Brunner \cite{bib:1982a},
$\mathcal{N}2$.
\item[11.] $164+60+91 \not\rightarrow 133\ (\mathrm{ZF}^{0})$, Brunner \cite{bib:1982a},
$\mathcal{N}3$.
\item[12.] $133+17+89\not\rightarrow 62\ (\mathrm{ZF}^{0})$, Brunner \cite{bib:1982a},
$\mathcal{N}1$.
\item[13.] $62+57 \not\rightarrow 9$ (ZF), Sageev\index{Author}{Sageev} \cite{bib:1981}, $\mathcal{M} 32$. See note 44 for a related result.
\item[14.] $52+43 \not\rightarrow 206$ (ZF), Pincus/Solovay \cite{bib:1977}, $\mathcal{M}27$.
\end{enumerate}

\section*{NOTE 2}

In this note we list implications that hold in every Fraenkel-Mostowski model. None of the implications listed are known to be theorems of ZF. In each case we list the implication and a reference if one is available.
\begin{enumerate}
\item[1.] $231\rightarrow 23$, this is reasonably clear.
\item[2.] $40+60\rightarrow 1$, Howard/Rubin \cite{bib:1995b}.
\item[3.] $62\rightarrow 60$, Howard\index{Author}{Howard} \cite{bib:1973}.
\item[4.] $14\rightarrow 317$, Howard \cite{bib:1973}.
\item[5.] $165\rightarrow 231$.
\end{enumerate}

\begin{proof}
If $X$\index{Author}{Brunner/Rubin} is a set of well orderable sets in a Fraenkel-Mostowski model $\mathcal{N}$, then for each $y\in X$, the set $y^{\prime}$ of well orderings of $y$ is well orderable in $\mathcal{N}$ since any permutation fixing $y$ pointwise also fixes $y^{\prime}$ pointwise. Therefore, by form 165 $(C(WO,\,WO))$, the set $X^{\prime}=\{y^{\prime}:y\in X\}$ has a choice function. Using this choice function and a well ordering of $X$ a well ordering of $\bigcup X$ can be constructed.
\end{proof}

\begin{enumerate}
\item[6.] $323\rightarrow 60$, because $323\rightarrow 62$ is a theorem of $\mathrm{ZF}^{0}$, and, by 3 above, $62\rightarrow 60$ in every Fraenkel-Mostowski model. Howard\index{Author}{Howard} \cite{bib:1973}.
\item[7.] $8+67\leftrightarrow 1$, Brunner/Rubin \cite{bib:1986}.
\item[8.] $165\rightarrow 24$, K. Keremedis. (In every FM model $\mathcal{P}(\mathbb{R})$ can be well ordered, form 130 is true, and $|\mathcal{P}(\mathbb{R})|=2^{(2^{\aleph_{0}})}.)$
\item[9.] $165\rightarrow 16$. (In every FM model, $\mathbb{R}$ can be well ordered, form 79 is true, and $|\mathbb{R}|=2^{\aleph_{0}}.)$
\item[10.] $165\rightarrow 212$. (See 9.)
\end{enumerate}

\section*{NOTE 3}

From Mathias\index{Author}{Mathias} \cite{bib:1979} p 124:

\begin{theorem*}
(Solovay) $Con\text{(}ZF\text{)}\rightarrow Con\text{(}ZF+\text{43}\ \text{(}DC(\omega)\text{)}\
\text{+}$ Every endomorphism of the additive group of the real line is continuous)
\end{theorem*}

\begin{theorem*}
(Levy) $Con\text{(}ZF\text{)}\rightarrow Con\text{(}ZF \text{+}\ \text{cf}\text{(}\omega_{1})=\omega +\mathbb{R}$ is the denumerable union of denumerable sets $\text{+}\ \neg\Sigma_{3}^{1}-AC$).
\end{theorem*}

In Figura\index{Author}{Figura} \cite{bib:1981} a summary of these results is given. In addition Figura generalizes form 34 ($\aleph_{1}$ is regular), form 38 ($\mathbb{R}$ is not a countable union of countable sets) and form 170 ($\mathbb{R}$ has an uncountable well orderable subset.) These generalizations for regular cardinals $\kappa$, are respectively, (in the notation of Figura \cite{bib:1981}):

$\neg CP(\kappa):\kappa^{+}$ is regular.

$\neg WOP(\kappa):\mathcal{P}(\kappa)$ has a well ordered subset of power not $\leq \kappa$.

$\neg DP(\kappa):\mathcal{P}(\kappa)$ is not the union of $\kappa$ sets of power $\leq \kappa$.

\noindent In summary, the results are
\begin{enumerate}
\item[1.] $DP(\kappa)$ implies $WOP(\kappa)$ and $DP(\kappa)$ implies $CP(\kappa)$.
\item[2.] $\mathrm{Con}(\mathrm{ZF}+CP(\kappa)\,+WOP(\kappa)+DP(\kappa))$
\item[3.] $\mathrm{Con(ZF}+CP(\kappa)+\neg WOP(\kappa)\,+\neg DP(\kappa))$
\item[4.] $\mathrm{Con(ZF}+CP(\kappa)\,+WOP(\kappa)\,+\neg DP(\kappa))$
\item[5.] $\mathrm{Con(ZF}+WOP(\kappa)\,+\neg CP(\kappa)+\neg DP(\kappa))$
\item[6.] $\mathrm{Con(ZF}+\{\kappa:DP(\kappa)\}$ is a proper class)
\end{enumerate}

\section*{NOTE 4}

Form 73 $\rightarrow$ [73 A]. (Form 73 is: For every $A\subseteq \mathbb{R}$ and $x\in \mathbb{R}$ the following definitions are equivalent:
\begin{enumerate}
\item[(1)] $x$ is in the closure of $A$ iff every neighborhood of $x$ intersects $A$.
\item[(2)] $x$ is in the closure of $A$ if there is a sequence $\{x_{n}\}\subseteq A$ such that $\lim x_{n}=x$.
\end{enumerate}
and form [73 A] is: For all $A\subseteq \mathbb{R},\,x\in A$ and $f: A\rightarrow \mathbb{R}$ the following are equivalent:
\begin{enumerate}
\item[(1)] $(\forall\epsilon)(\exists\delta)(\forall y\in A)(|y-x|<\delta$ implies $|f(y)-f(x)|<\epsilon)$
\item[(2)] Whenever $\{x_{n}\}\subseteq A$ and $\lim x_{n}=x$ then $\lim f(x_{n})=f(x).)$
\end{enumerate}

\begin{proof}
Assume Form 73. Fix $A\subseteq \mathbb{R}$ and $x\in A$. The fact that part (1) of [73 A] implies part (2) of [73 A] is clear. Assume (2) of [73 A] and $\neg(1)$ of [73 A]. Then $\exists\epsilon\,\forall\delta\,\exists y\in A$ such that $|y-x|<\delta$ and $|f(x)-f(y)|\geq\epsilon$. Let $B=\{y\in A:|f(x)-f(y)|\geq\epsilon\}$. Every neighborhood of $x$ intersects $B$. By form 73, $x= \lim x_{n}$ for some sequence $\{x_{n}\}\subseteq B(\subseteq A)$. But $\lim f(x_{n})\neq f(x)$ contradicting part (2) of [73 A].
\end{proof}

\section*{NOTE 5}

[73 A] $\rightarrow$ form 74. (See note 4 for a statement of [73 A]. Form 74 is: For every $A\subseteq \mathbb{R}$ the following are equivalent:
\begin{enumerate}
\item[(1)] $A$ is closed and bounded.
\item[(2)] Every sequence $\{x_{n}\}\subseteq A$ has a convergent subsequence with limit in A.)
\end{enumerate}

\begin{proof}
Assume [73 A]. The usual proof that part (1) of form 74 implies part (2) doesn't require any choice. So assume $A\subseteq \mathbb{R}$ and $A$ satisfies (2) of 74 but not (1) of 74. Assume first that $A$ is bounded. Then $A$ is not closed so there is an $x_{0}\not\in A$ such that every neighborhood of $x_{0}$ intersects $A$. Define $f: A\cup \{x_{0}\}\rightarrow \mathbb{R}$ by $f(x)=1$ if $x\in A$ and $f(x)=0$ if $x=x_{0}$. Then (1) of [73 A] fails (with $\epsilon=1)$ so there is a sequence $\{x_{n}\}\subseteq A \cup\{x_{0}\}$ such that $\lim f(x_{n})\neq f(x_{0})$. This means we can assume $\{x_{n}\}\subseteq A$ (replacing $\{x_{n}\}$ by a subsequence if necessary). By (2) of form 74, $\lim x_{n}\in A$ that is $x_{0}\in A$, a contradiction. Now assume that $A$ is unbounded. Let $f$ be a monotone, increasing function from the interval $(-1,1)$ one to one, onto $\mathbb{R}$ which satisfies $(\forall a,\,b\in(-1,1))(|a-b|\leq|f(a)-f(b)|)$. (For example, $f(x)=\frac{2x}{1-x^{2}}$ would work.) Then $f^{-1}(A)$ is bounded and satisfies (2) of 74 because if $\{x_{n}\}$ is a sequence in $f^{-1}(A)$, then $\{f(x_{n})\}$ is a sequence in $A$. Since $A$ satisfies (2) of 74, $\{f(x_{n})\}$ has a convergent subsequence with limit in $A$. Using the property which we required $f$ to have we infer that $f^{-1}$ applied to the convergent subsequence mentioned above gives a convergent subsequence of $\{x_{n}\}$ with limit in $f^{-1}(A)$. Since we have proved (2) of 74 implies (1) of 74 for bounded sets we conclude that $f^{-1}(A)$ is closed. Hence there are numbers $a$ and $b$ such that $-1<a<b<1$ and $(\forall x\in f^{-1}(A))(a<x<b)$. Since $f$ is monotone increasing $(\forall x\in A)(f(a)<x<f(b))$. Hence $A$ is bounded. 
\end{proof}

\section*{NOTE 6}

[73 A] $\rightarrow$ form 73. (See note 4 for a statements of forms 73 and [73 A].)

\begin{proof}
Assume [73 A] and assume $A\subseteq \mathbb{R}$ and $x_{0}\in \mathbb{R}$ satisfy part (1) of 73 but not (2) of 73 (This means $x_{0}\not\in A$). Then
\begin{equation*}
\text{For every sequence}\ \{x_{n}\}\ \mathrm{from}\ A,\ \lim x_{n}\neq x_{0} \qquad\qquad\qquad (\ast)
\end{equation*}
 Define $f:A \cup\{x_{0}\}\rightarrow \mathbb{R}$ by $f(x)=1$ if $x\in A$ and $f(x)=0$ if $x=x_{0}$. By $(\ast)$, if $\{x_{n}\}$ converges to $x_{0}$ and $\{x_{n}\}\subseteq A \cup\{x_{0}\}$ then $\{x_{n}\}$ is eventually constant and equal to $x_{0}$. Therefore, $\lim f(x_{n})=f(x_{0})$. Therefore, (2) of [73 A] holds (with $A$ of [73 A] equal to $A \cup\{x_{0}\}$ of the present proof). Hence (1) of [73 A] holds, which means $(\exists\delta>0)(\forall x\in A\cup\{x_{0}\})(|x-x_{0}|<\delta\rightarrow|f(x)-f(x_{0})|<1)$. But by (1) of 73, the set $B=\{x:|x-x_{0}|<\delta\}\cap A\neq\emptyset$. Choose $y\in B$, then $y\neq x_{0},\, y\in A$ and $|1-0|=|f(y)-f(x)|<1$, a contradiction. 
\end{proof}

\section*{NOTE 7}

In this note we prove that form 43 (Dependent Choice) implies form 113 (A countable product of compact topological spaces is compact.) Our proof is based on the proof of theorem 6.1 in Good/Tree\index{Author}{Good/Tree} \cite{bib:1995}. Let $\{X_{\alpha}:\alpha<\omega\}$ be a countable set of compact topological spaces. Let $X=\prod_{\alpha\in\omega}X_{\alpha}$ and for $\beta<\omega$ define $p_{\beta}:X\rightarrow\prod_{\alpha\leq\beta}X_{\alpha}$ by $p_{\beta} (x_{0},\,x_{1},\,\ldots) =(x_{0},\,x_{1},\, \ldots,\,x_{\beta})$. To show that $X$ is compact let $\mathcal{C}$ be a collection of closed sets in $X$ with the finite intersection property. We will use Dependent Choice to choose a sequence $(x_{0},\,x_{1},\,\ldots)$ with the property that for all $\alpha<\omega,\,(x_{0},\,x_{1},\,\ldots,\, x_{\alpha})\in\bigcap\{\overline{p_{\alpha}(C)}:C\in \mathcal{C}\}$. Then we will prove, as in Good/Tree \cite{bib:1995}, that
$(x_{1},\,x_{2},\,\ldots) \in\bigcap \mathcal{C}$.

\begin{definition*}
If $x_{0}\in X_{0}$ and $\alpha<\kappa$ we say $x_{0}$ \emph{can be extended to level} $\alpha$ if there is a sequence $(x_{1},\,\ldots,\,x_{\alpha})$ such that $(x_{0},\,x_{1},\,\ldots,\, x_{\alpha})\in\bigcap\{\overline{p_{\alpha}(C)}:C\in \mathcal{C}\}$. Similarly, if $\alpha<\beta<\kappa$ and $(x_{0},\,x_{1},\, \ldots,\,x_{\alpha})\in\bigcap\{\overline{p_{\alpha}(C)}:C\in \mathcal{C}\}$ then we say $(x_{0},\,x_{1},\,\ldots,\,x_{\alpha})$ can be extended to level $\beta$ if there is a sequence $(y_{\alpha+1},\,\ldots,\,y_{\beta})$ such that $(x_{0},\,x_{1},\,\ldots,\,x_{\alpha},\,y_{\alpha+1},\,\ldots,\, y_{\beta})\in\cap\{\overline{p_{\beta}(C)}:C\in \mathcal{C}\}$.
\end{definition*}

\begin{definition*}
If $\beta\leq\gamma<\kappa$ define
\begin{equation*}
p_{\beta,\gamma}:\prod_{\alpha\leq\gamma}X_{\alpha}\rightarrow\prod_{\alpha\leq\beta}X_{\alpha}
\end{equation*}
by
\begin{equation*}
p_{\beta,\gamma}\,(x_{0},\,x_{1},\,\ldots,\,x_{\beta},\,\ldots,\,x_{\gamma})=(x_{0},\,x_{1},\,\ldots,\, x_{\beta}).
\end{equation*}

For every $\alpha<\kappa$, let $X_{0,\alpha}=\{x\in X_{0}:x$ can be extended to level $\alpha$\}. $X_{0,\alpha}$ is closed in $X_{0}$ because $X_{0,\alpha}=p_{0,\alpha} (\bigcap\{\overline{p_{\alpha}(C)}:C\in \mathcal{C}\})$ and
\begin{enumerate}
\item[A.] $\prod_{0\leq\beta\leq\alpha}X_{\beta}=X_{0}\times\prod_{1\leq\beta\leq\alpha}X_{\beta}$ is compact.
\item[B.] $\bigcap\{\overline{p_{\alpha}(C)}:C\in \mathcal{C}\}$ is closed in $\prod_{0\leq\beta\leq\alpha}X_{\beta}$.
\item[C.] If $Y$ and $Z$ are compact topological spaces and $D$ is closed in $Y\times Z$ then the projection of $D$ onto $Y$ is closed.
\end{enumerate}

We also note that $X_{0,\alpha}\neq\emptyset$ since $\bigcap\{\overline{p_{\alpha}(C)}:C\in \mathcal{C}\}$ is non-empty (by the compactness of $\prod_{0\leq\beta\leq\alpha}X_{\beta}$). We can conclude that the family $\{X_{0,\alpha}:\alpha<\kappa\}$ has the finite intersection property since $X_{0,\alpha_{1}} \cap\ldots\cap X_{0,\alpha_{k}}=X_{0,\alpha}$ where $\alpha=\max\{\alpha_{1},\,\ldots,\,\alpha_{k}\}$. Since $X_{0}$ is compact $\bigcap_{\alpha<\kappa}X_{0,\alpha}\neq\emptyset$. Choose $x_{0} \in\bigcap_{\alpha<\kappa}X_{0,\alpha}$.

Assume that $x_{0},\,\ldots,\,x_{\beta}$ have been chosen so that $(x_{0},\,\ldots,\,x_{\beta})$ is extendible to level $\alpha$ for every $\alpha$ such that $\beta<\alpha<\kappa$. For every such $\alpha$, let $X_{\beta+1,\alpha}=\{x\in X_{\beta+1}:(x_{0},\, \ldots,\,x_{\beta},\,x)$ is extendible to level $\alpha$\}.
\end{definition*}

\begin{claim*}
$X_{\beta+1,\alpha}$ is closed.
\end{claim*}

\begin{proof}
First note that
\begin{equation*}
X_{\beta+1,\alpha}=\{x\in X_{\beta+1}:(x_{0},\,\ldots,\,x_{\beta},\,x)\in p_{\beta+1,\alpha}(\bigcap\{\overline{p_{\alpha}(C)}:C\in \mathcal{C}\}).
\end{equation*}
$\bigcap\{\overline{p_{\alpha}(C)}:C\in \mathcal{C}\}$ is closed in $\prod_{\gamma\leq\alpha}X_{\gamma}=\Big(\prod_{\gamma\leq\beta+1}X_{\gamma}\Big)\times\Big(\coprod_{\beta+1<\gamma\leq\alpha}X_{\gamma}\Big)$. Therefore, by C. above, $p_{\beta+1,\alpha} (\bigcap\{\overline{p_{\alpha}(C)}:C\in \mathcal{C}\})$ is closed. Assume $y\in X_{\beta+1}$ is in the closure of $X_{\beta+1,\alpha}$. We will complete the proof of the claim by showing that $y\in X_{\beta+1,\alpha}$. Let $N$ be a neighborhood of $(x_{0},\,\ldots,\,x_{\beta},\,y)$ in $\prod_{\gamma\leq\beta+1}X_{\gamma}$. Then the projection of $N$ onto $X_{\beta+1}$ is a neighborhood of $y$. Therefore there is an $x\in X_{\beta+1,\alpha}$ such that $x$ is in this projection and therefore, the sequence $(x_{0},\,\ldots,\, x_{\beta},\,x)$ is in $N$. Also $(x_{0},\,\ldots,\,x_{\beta},\,x)$ is in $p_{\beta+1,\alpha}(\bigcap\{\overline{p_{\alpha}(C)}:C\in \mathcal{C}\})$. This shows that $(x_{0},\,\ldots,\, x_{\beta},\,y)$ is in the closure of $p_{\beta+1,\alpha}(\bigcap\{\overline{p_{\alpha}(C)}:C\in \mathcal{C}\})$ and hence $(x_{0},\, \ldots,\,x_{\beta},\,y)$ is in $p_{\beta+1,\alpha}(\bigcap\{\overline{p_{\alpha}(C)}:C\in \mathcal{C}\})$. $y$ is therefore in $X_{\beta+1,\alpha}$.
\end{proof}

As above, the family $\{X_{\beta+1,\alpha}:\beta<\alpha<\kappa\}$ has the finite intersection property. Since $X_{\beta+1}$ is compact, $\bigcap_{\beta<\alpha<\kappa}X_{\beta+1,\alpha}\neq\emptyset$. Choose $x_{\beta+1}\in\bigcap_{\beta<\alpha<\kappa}X_{\beta+1,\alpha}$.

The argument that $(x_{0},\,x_{1},\,\ldots)$ is in $\bigcap \mathcal{C}$ is by contradiction. Suppose the assertion is false. Then there is an open set $U=\prod_{\alpha<\omega}U_{\alpha}$ and $C\in \mathcal{C}$ such that $(x_{0},\,x_{1},\, \ldots) \in U$ and $U\cap C=\emptyset$. Choose $\gamma\in \omega$ so that $U_{\alpha}=X_{\alpha}$ whenever $\alpha>\gamma$, and let $V=\prod_{\alpha\leq\gamma}U_{\alpha}$. $V$ is open in $\prod_{\alpha\leq\gamma}X_{\alpha}$. Further $(x_{0},\,x_{1},\,\ldots,\,x_{\gamma})\in V$ hence $V\cap p_{\gamma}(C)\neq\emptyset$.
This means that there is some sequence $(y_{0},\,y_{1},\,\ldots) \in C$ with $(y_{0},\,y_{1},\,\ldots,\,y_{\gamma})\in V$. Which (by definition of $V$) implies that $(y_{0},\, y_{1},\,\ldots) \in U$. This contradicts $U\cap C=\emptyset$.

\section*{NOTE 8}

[9 A] (The statement
\begin{equation*}
(\forall x)((x\ \text{infinite and}\ x\ \text{Dedekind finite) implies}\ \mathcal{P}(x)\ \text{is Dedekind finite)}\qquad (\ast))
\end{equation*}
is equivalent to form 9 (Dedekind finite $=$ finite).

\begin{proof}
Assume $(\ast)$ and suppose that $x$ is a set which is infinite and not Dedekind infinite. Then by $(\ast)$, $Y=\mathcal{P}(x)$ is infinite and Dedekind finite. Applying $(\ast)$ again, $z=\mathcal{P}(\mathcal{P}(x))$ is infinite and Dedekind finite. But $\{A_{n}:n\in \omega\}$ is a denumerable subset of $z$ where $A_{n}=\{t\subseteq x:|t|=n\}$, a contradiction. Conversely assuming form 9, the hypotheses of $(\ast)$ are never true.
\end{proof}

\section*{NOTE 9}

For the relationships between form 71 $(W_{\aleph_{\alpha}}$, Every set is
comparable with $\aleph_{\alpha}.)$, form 86 $(C(\aleph_{\alpha},\,\infty))$ and form 87 $(DC(\aleph_{\alpha}))$ see Jech \cite{bib:1973b}, chapter 8 and the papers he refers to (Levy\index{Author}{Levy} \cite{bib:1964}, Mostowski\index{Author}{Mostowski} \cite{bib:1958}, Jech\index{Author}{Jech} \cite{bib:1966b}, Pincus\index{Author}{Pincus} \cite{bib:1969} and Jensen\index{Author}{Jensen} \cite{bib:1967}). Also see Pincus \cite{bib:1972a}, pp 731ff.

\section*{NOTE 10}

[73 A] is equivalent to [73 B]. (See note 4 for a statement of [73 A]. Form [73 B] is: $(\forall f:\mathbb{R}\,\rightarrow \mathbb{R})(\forall x\in \mathbb{R})$ the following are equivalent:
\begin{enumerate}
\item[(1)] $(\forall\epsilon>0)(\exists\delta>0)((\forall y)(|y-x|<\delta\rightarrow|f(y)-f(x)|<\epsilon)$
\item[(2)] Whenever $\lim x_{n}=x$, then $\lim f(x_{n})=f(x).)$
\end{enumerate}

\begin{proof}
Clearly [73 A] implies [73 B]. Assume [73 B] and assume that [73 A] fails. Then $\exists A\subseteq \mathbb{R}\, \exists x_{0}\in A$ and $\exists f:A\rightarrow \mathbb{R}$ such that (2) of [73 A] holds but (1) of [73 A] fails. Define $g:\mathbb{R} \rightarrow \mathbb{R}$ by $g(x)=f(x)$ if $x\in A$ and $g(x)=f(x_{0})$ if $x\not\in A$. Then (2) of [73 B] holds for $g$ so (1) of [73 B] holds for $g$. This implies that (1) of [73 A] holds for $f$.
\end{proof}

\section*{NOTE 11}

Form 9 (Dedekind finite $=$ finite) is equivalent to [9 C] (The image of a Dedekind finite set is Dedekind finite.)

\begin{proof}
Clearly form 9 implies [9 C]. By note 8 (above) to prove [9 C] implies 9, it suffices to prove [9 C] implies [9 A]. Assume [9 C] and that $x$ is an infinite, Dedekind finite set and that $\mathcal{P}(x)$ is Dedekind infinite. Then there is a one to one function $f:\omega \rightarrow \mathcal{P}(x)$.
\end{proof}

\begin{lemma*}
We may assume that $(\forall i,\,j\in \omega)(i\neq j$ implies $f(i) \cap f(j)=\emptyset)$.
\end{lemma*}

\begin{proof}
If $g:\omega\rightarrow \mathcal{P}(x)$ is one to one, then we define $f:\omega\rightarrow \mathcal{P}(x)$ by induction so that for each $i\in\omega,\,f(i)\cap f(j)=\emptyset$ for all $j<i$, and such that the set algebra generated by $\{g(k)\cap(x-\bigcup_{j\leq i}f(j)):k>j\}$ is infinite, $f(0)$ is defined as follows: Either the set algebra generated by $\{g(k)\,\cap\,g(0):k>0\}$ is infinite or the set algebra generated by $\{g(k)\,\cap\,(x-g(0)):k>0\}$ is infinite. If the latter is true then $f(0)=g(0)$. If the latter is false then $f(0)=x-g(0)$. Suppose that $f(i)$ has been defined for every $i<n$ satisfying the required conditions. Then for some $k$ (and we take the least such) $g(k) \cap (x-\bigcup_{j<n}f(j))\neq\emptyset$ (call this set $A$) and $(x-\bigcup_{j<n}f(j))-g(k)\neq\emptyset$ (call this set $B$). Since $A\cup B=(x-\bigcup_{j<n}f(j))$, one of the two set algebras, either that generated by $\{g(k)\,\cap\,A:k>n\}$ or that generated by $\{g(k)\,\cap\, B:k>n\}$, is infinite. If the latter is the case, let $f(k)=A$. If the latter fails, let $f(k)=B$. This completes the proof of the lemma.
\end{proof}

Now define $h:X\rightarrow \mathcal{P}(x)$ by
\begin{equation*}
h(t)=\left\{\begin{array}{ll}
f(i), & \mathrm{if}\ t\in f(i)\\
\emptyset\ & \mathrm{otherwise}.
\end{array}\right.
\end{equation*}
By [9 C] the image of $X$ under $h$ is Dedekind finite, a contradiction.

\section*{NOTE 12}

Form 9 (Dedekind finite $=$ finite) and [9 D] (The union of a Dedekind finite family of finite sets is Dedekind finite) are equivalent.

\begin{proof}
It is clear that form 9 implies [9 D]. Since 9 and [9 A] are equivalent (see note 8.) it suffices to show that [9 D] implies [9 A]. Assume [9 D] and that $x$ is a set which is Dedekind finite such that $\mathcal{P}(x)$ is Dedekind infinite. Then there is a one to one function $f:\omega\rightarrow \mathcal{P}(x)$ and, as in note 11, we may assume that $(\forall i,\,j\in \omega)(i\neq j$ implies $f(i) \cap f(j)=\emptyset)$. It is also no loss of generality to assume $(\forall i\in \omega)(f(i)\neq\emptyset)$. Let
\begin{equation*}
W=\{\{i,\,t\}:i\in\omega\wedge t\in f(i)\}.
\end{equation*}
Then $W$ is Dedekind finite since $x$ is. But $\omega\subseteq\bigcup W$, a contradiction. $\square$(See Jech\index{Author}{Jech} \cite{bib:1973b} p 161 prob 11.4.)
\end{proof}

\section*{NOTE 13}

The fact that form 39 does not imply form 43 follows from the theorem of Jensen mentioned in Mathias\index{Author}{Mathias} \cite{bib:1979} p. 123 and proved in Felgner\index{Author}{Felgner} \cite{bib:1971a}:

\begin{theorem*}
If $\alpha$\index{Subject}{Gauntt's condition $L(n,\, Z)$} is an infinite cardinal in a countable standard model $M$ of $ZF+V=L$ then $M$ can be extended to an model $N$ of $ZF$ with the same ordinals and (well ordered) cardinals as $M$ in which $C(\alpha,\,\infty)$ (form 39) is true and $DC(\omega)$ (form 43) is false.
\end{theorem*}

\section*{NOTE 14}

Form 40 does not imply form 44 by the following theorem which is proved in Jech\index{Author}{Jech} \cite{bib:1966b}:

\begin{theorem*}
Let $M$ be a countable standard model for $ZF+AC$ and $\alpha$ a regular, infinite cardinal in $M$. Then there is a model $N$ of $ZF$ containing $M$ with the same ordinals as $M$ such that in $N,\,\neg DC(\alpha)$ (see form 87), $\neg C(\alpha,\,2)$ and there is a cardinal incomparable with $\alpha$. But for every $\beta<\alpha,\,DC(\beta)$ and $C(\beta,\,\infty)$ hold and $\beta$ is comparable with every cardinal.
\end{theorem*}

\section*{NOTE 15}

Finite axioms of choice:\index{Subject}{finite axioms of choice} $C(\infty,\,n),C(LO,\,n),C(WO,\,n),C(\infty,\,\in Z),C(LO, \in Z),C(WO,\,\in Z)$ and $C(\aleph_{0},\,n)$, (forms 45, 33, 47, 46, 120, 48, and 288 respectively) and generalizations have been studied extensively in Blass\index{Author}{Blass} \cite{bib:1984b}, Conway\index{Author}{Conway} \cite{bib:1971}, Gauntt\index{Author}{Gauntt} \cite{bib:1970} and \cite{bib:1968}, Mostowski\index{Author}{Mostowski} \cite{bib:1958}, Truss\index{Author}{Truss} \cite{bib:1973a}, Wi\'{s}niewski\index{Author}{Wi\'{s}niewski} \cite{bib:1972} and Zuckerman\index{Author}{Zuckerman} \cite{bib:1971}. Here is a summary of the results. (See Truss \cite{bib:1973a} for more details.)

Until further notice $n\in \omega$ and $Z$ is a finite subset of $\omega$. The first two of the following conditions on the pair $(n,\,Z)$ were introduced by Mostowski\index{Subject}{Mostowski's condition $D(n,\,Z)$} in \cite{bib:1958} and the third by Gauntt in \cite{bib:1970}.

$D(n,\,Z)$. For any subgroup $G$ of $S_{n}$ (the symmetric group on n) without fixed points, there is a subgroup $H$ of $G$ and proper subgroups $K_{1},\,\ldots,\,K_{r}$of $H$ such that $\Sigma|H:K_{i}|\in Z$.

$K(n,\,Z)$.\index{Subject}{Mostowski's condition $K(n,\,Z)$} For any subgroup $G$ of $S_{n}$ without fixed points, there is a subgroup $H$ of $G^{\omega}$ (the restricted direct product of countably many copies of $G$) and proper subgroups $K_{1},\,\ldots,\, K_{r}$ of $H$ such that $\Sigma|H:K_{i}|\in Z$

$L(n,\,Z)$. For any subgroup $G$ of $S_{n}$ without fixed points, there are proper subgroups $K_{1},\, \ldots,\,K_{r}$ of $G$ such that $\Sigma|G:K_{i}|\in Z$.

Blass \cite{bib:1984b} has noticed that $L(n,\,Z)$ is can be formulated as follows:

Every group that can act without fixed points on an $n$ element set can also act without fixed points on a $k$ element set for at least one $k\in Z$.

\noindent The equivalence of the two conditions is based on the following:

\begin{lemma*}
Suppose that $H$ is a group and $A$ is a set of cardinality $k$ then $H$ acts without fixed points on $A\Leftrightarrow$ there are proper subgroups $J_{1},\,\ldots,\, J_{r}$ of $H$ (not necessarily distinct) such that $\Sigma|H:J_{i}|=k$.
\end{lemma*}

\begin{proof}
(Outline) $(\Rightarrow)$ If $H$ acts without fixed points on $A$, let $A_{1},\,\ldots,\,A_{r}$ be the orbits of elements of A under the action. Then $|A_{i}|\geq 2$ and $\Sigma|A_{i}|=k$. If, for each $i,\, 1\leq i\leq r$ we choose $a_{i}\in A_{i}$ and let $J_{i}$ be the subgroup of $H$ consisting of those elements that fix $a_{i}$ then $J_{i}$ is proper and $|H:J_{i}|=|A_{i}|$.

\noindent$(\Leftarrow)$ Let $A$ be the disjoint union of the left cosets of the $J_{i},\,1\leq i\leq r$. Then $H$ acts on $A$ without fixed points.
\end{proof}

\begin{theorem*}
$D(n,\,Z)\Leftrightarrow ZF\vdash C(\infty,\,Z)\rightarrow C(\infty,\,n)$. $(\Leftarrow$ Gauntt\index{Author}{Gauntt} \emph{\cite{bib:1970}}, $\Rightarrow$ Mostowski\index{Author}{Mostowski} \emph{\cite{bib:1958}})
\end{theorem*}

\begin{theorem*}
$L(n,\,Z)\Leftrightarrow ZF\vdash C(WO,\,Z)\rightarrow C(WO,\,n)$. (Gauntt \emph{\cite{bib:1968}})
\end{theorem*}

\begin{theorem*}
$K(n,\,Z)\Leftrightarrow ZF\vdash C(\infty,\,Z)\rightarrow C(WO,\,n)$. ($\Rightarrow$ Truss\index{Author}{Truss} \emph{\cite{bib:1973a}}, $\Leftarrow$ Mostowski \emph{\cite{bib:1958}})
\end{theorem*}

\begin{theorem*}
$L(n,\,Z)\Leftrightarrow ZF\vdash C(LO,\,Z)\rightarrow C(LO,\,n)$. (Truss \emph{\cite{bib:1973a}})
\end{theorem*}

\begin{theorem*}
$ZF\vdash C(LO,\,Z)\rightarrow C(WO,\,n)\Leftrightarrow ZF\vdash C(WO,\,Z)\rightarrow C(WO,\,n)$. (Truss \emph{\cite{bib:1973a}})
\end{theorem*}

\begin{theorem*}
$ZF\vdash C(\infty,\,Z)\rightarrow C(LO,\,n)\Leftrightarrow ZF\vdash C(\infty,\,Z)\rightarrow C(\infty,\,n)$. (Truss \emph{\cite{bib:1973a}})
\end{theorem*}

\noindent Blass\index{Author}{Blass} \cite{bib:1984b} has introduced the following forms each depending on the set I:

$C(I,\,n)$. Every $I$ indexed family of $n$ element sets has a choice function. (Depends on $n\in \omega$.)

$C(I,\,\in Z)$. $(\forall n\in Z)(C(I,\,n))$. (Depends on finite $Z\subseteq \omega$.) and has shown

\begin{theorem*}
$L(n,\,Z)\Leftrightarrow ZF\vdash(\forall I)(C(I,\,\in Z)\rightarrow C(I,\,n))$.
\end{theorem*}

We now turn to some generalizations of the finite axioms of choice due to Truss \cite{bib:1973a}. Some definitions are necessary. For any set $X$ let $e(X)$ be the set of nonempty finite subsets of $X$ and define $e(0,\, X)=X$ and $e(n+1,\,X)=e(e(n,\,X))$. (We assume that $e(k,\,X) \cap e(j,\,X)=\emptyset$ for $j\neq k$. If this is not the case we replace $e(n,\,X)$ by $\{e(n,\,X)\}\times e(n,\,X)$ for each $n\in \omega$.) Let $V(m,\, X)$ be the set of equivalence classes on $e(m,\,X)$ under the equivalence relation $aRb$ iff $\exists$ a permutation $\phi$ of $X$ such that (the extension of) $\phi(a)=\mathrm{b}$. If $|X|=n\in\omega$ then there is a natural 1-1 correspondence between $V(m,\,n)$ and $V(m,\,X).(\nu\rightarrow f(\nu)$ where $f$ is (the extension of) any one to one correspondence between $X$ and $n$.) For $\nu\in V(m,\,n)$ we denote the corresponding equivalence class in $V(m,\,X)$ by $\nu(X)$. Let $V=\bigcup\{V(m,\,n):m,\,n\in \omega,\, n\neq 0\}$. For any $Z\subseteq V$, Truss defines the following choice axioms: ($Z$ no longer denotes a finite subset of $\omega$.)

$C(\infty,\,Z)$. For any set $Y,\,\{\nu(X):\nu\in Z$ and $X\in Y$ and $\nu\in V(m,\,|X|)$ for some $m\in \omega$\} has a choice function.

$C(LO,\,Z)$. For any linearly ordered set $Y$, same as above.

$C(WO,\,Z)$. For any well ordered set $Y$, same as above.

If $\nu\in V$, then we let $C(\infty,\,\nu)$, $C(LO,\,\nu)$ and $C(WO,\,\nu)$ denote respectively $C(\infty,\,\{\nu\})$, $C(LO,\,\{\nu\})$ and $C(WO,\,\{\nu\})$.

\noindent For example, if $Z=\{\nu\}$ where $\nu$ is the set of all pairs of elements of $n$ then $C(\infty,\, \{\nu\})\equiv C(\infty,\,\nu)$ is ``For every set $Y$ of $n$ elements sets there is a function $f$ such that $\forall X\in Y,\,f(X)$ is a two element subset of X.''

Necessary and sufficient conditions are given for
\begin{enumerate}
\item[(A)] $\mathrm{ZF}\vdash C(\infty,\,Z_{1})\wedge C(\mathrm{LO},\,Z_{2})\wedge C(WO,\,Z_{3})\rightarrow C(\infty,\,\nu)$
\item[(B)] $\mathrm{ZF}\vdash C(\infty,\,Z_{1})\wedge C(LO,\,Z_{2})\wedge C(WO,\,Z_{3})\rightarrow C(LO,\,\nu)$ and
\item[(C)] $\mathrm{ZF}\vdash C(\infty,\,Z_{1})\wedge C(LO,\,Z_{2})\wedge C(WO,\,Z_{3})\rightarrow C(WO,\,\nu)$.
\end{enumerate}

They are

$A(Z,\nu)$. If $\nu\in V(m,\,n)$ and $G$ is a group of permutations of $n$ which moves every member of $\nu$ then there is an $X\in e(k,\,n)$ for some $k$ such that there is and a $w\in Z$ such that $w\in V(n^{\prime},\,|X|)$ for some $n^{\prime}$ and for each $\zeta \in w(X)$, $H(X) \cap G$ is not a subset of $H(\zeta)$. $(\mathrm{If}\,\zeta\,\in e(n,\,X)$ then $H(\zeta)=\{\sigma\in S(X):\sigma(\zeta)=\zeta\}$.)

\begin{theorem*}
$A(Z_{1},\,\nu)$ is equivalent to ($A$).
\end{theorem*}

$B(Z_{1},\,Z_{2},\,\nu)$. If $G$ is a group of permutations of $n$ (where
$\nu\in V(j,\,n)$ for some $j$) which moves every member of $\nu$, there is an $X$ such that either
\begin{enumerate}
\item[] $X\in e(k,\,n)$ for some $k$, and there is some $w\in Z_{1}$ a such that $w\in V(k,\,|X|)$ for some $k$ and for each $\zeta \in w(X)$, $H(X) \cap G$ is not a subset of $H(\zeta)$ or
\item[] $X\in e(k,\,n\cdot m)$ for some $k$ and $m\in \omega$ and there is a $w\in Z_{2}$ such that $w\in V(j,\,|X|)$ for some $j\in \omega$ and for each $\zeta\in w(X),\, G^{m}$ is not a subset of $H(\zeta)$.
\end{enumerate}

\begin{theorem*}
$B(Z_{1},\,Z_{2},\,\nu)$ is equivalent to ($B$).
\end{theorem*}

$C(Z_{1},\,Z_{2},\,\nu)$. If $G$ is a group of permutations of $n$ (where
$\nu\in V(j,\,n)$ for some $j\in \omega)$ which moves every member of $\nu$, then there is an $X\in e(j,\,n\cdot m)$ for some $j,\,m\in\omega$ such that either
\begin{enumerate}
\item[] for some $w\in Z_{1},\,w\in V(j,\,|X|)$ for some $j\in \omega$ and for each $\zeta \in w(X)$, $G^{m}\cap H(X)$ is not a subset of $H(\zeta)$
or
\item[] for some $w\in Z_{2},\,w\in V(j,\,|X|)$ for some $j\in \omega$ and $H(X)=G^{m}$ and for each $\zeta \in w(X)$, $G^{m}$ is not a subset of $H(\zeta)$.
\end{enumerate}

\begin{theorem*}
$C(Z_{1},\,Z_{2}\cup Z_{3},\,\nu)$ is equivalent to ($C$).
\end{theorem*}

\noindent See note 73 for adding dependent choice to the independence results above.

The finite axioms of choice have been generalized in Conway\index{Author}{Conway} \cite{bib:1971} as follows:

\begin{definition*}
Let $k$ be a natural number.
\begin{enumerate}
\item[1.] An ordered $k$-tuple $(A,\,B,\,C,\,\ldots)$ of disjoint finite sets of cardinalities $a,\,b,\,\ldots$ respectively is called an $(a,\,b,\,c,\,\ldots)$ set.
\item[2.] $[a,\,b,\,c,\,\ldots]$ is the axiom: For every collection of $(a,\,b,\,c,\,\ldots)$ sets there is a function $f$ defined on the collection such that $f(A,\,B,\,C,\,\ldots) \in A\cup B \cup C \cup\cdots$
\item[3.] $\langle a,\, b,\,c,\,\ldots\rangle$ is the additive semi-group generated by $a,\,b,\,c,\,\ldots$ (I. e., the closure of $\{a,\,b,\,c,\,\ldots\}$ under addition.)
\item[4.] If $S=\langle a,\,b,\,c,\,\ldots\rangle$ then $[S]$ denotes the axiom $[a,\,b,\,c,\,\ldots]$. (by result I below, if $\langle a,\,b,\,c,\,\ldots\rangle$ and $\langle\alpha,\,\beta,\,\gamma\ldots\rangle$ are equal then $[a,\,b,\,c,\,\ldots]$ and $[\alpha,\,\beta,\,\gamma,\,\ldots]$ are equivalent.)

\item[5.] If $G$ is any finite group, $\langle G\rangle\subseteq \mathbb{N}$ is the semi-group generated by the indices in $G$ of all proper subgroups of $G$.
\end{enumerate}
\end{definition*}

\begin{theorem*}
(Some parts due to Mostowski)
\begin{enumerate}
\item[\emph{I.}] If $\langle\alpha,\,\beta,\,\gamma,\,\ldots\rangle\subseteq\langle a,\,b,\,c,\,\ldots\rangle$ then $[\alpha,\, \beta,\,\gamma,\,\ldots]\Rightarrow[a,\,b,\,c,\,\ldots]$.
\item[\emph{II.}] Let $S_{1},\,\ldots,\,S_{N},\,S$ be semi-groups ($\subseteq \mathbb{N}$), then for the implication
\begin{equation*}\tag{$^{\ast}$}
[S_{1}]\wedge[S_{2}]\wedge\cdots\wedge[S_{N}]\Rightarrow[S]
\end{equation*}
to hold it suffices that every group $G$ with $S\subseteq\langle G\rangle$ has some subgroup $H$ for which $S_{i}\subseteq\langle H\rangle$ for some $i,\,1\leq i\leq N$.
\item[\emph{III.}] (Mostowski) For ($\ast$) to hold it is necessary that every Abelian group $G$ with $S\subseteq\langle G\rangle$ have some subgroup $H$ with $S_{i}\subseteq\langle H\rangle$ for some $i,\,1\leq i\leq N$.
\item[IV.] The condition of II is necessary and sufficient for effective implication, that is, for the existence of a formula $F(f_{1},\,\ldots,\,f_{N})=f$ which defines a choice function $f$ for the class of all $S$ sets whenever we substitute such choice functions for the classes of all $S_{1}$ sets, $S_{2}$ sets, $\ldots$, $S_{N}$ sets for $f_{1},\,f_{2},\ldots,f_{N}$.
\end{enumerate}
\end{theorem*}

The following are proved in Wi\'{s}niewski\index{Author}{Wi\'{s}niewski} \cite{bib:1972}:

\begin{theorem*}
Let $P$ be the additive semi-group of positive integers generated by a set $Z$ of prime numbers ($P$ consists of sums of the form $p_{1}+p_{2}+\cdots+p_{k}$ where $p_{i},\, i=1,\,\ldots,\,k$ are in $Z$.) Then there is a model $\mathcal{N}$ of $ZF^{0}$ in which $C(\infty,\,n)$ is true for all $n\not\in P$ and $C(\aleph_{0},\,n)$ is false for all $n\in P$.
\end{theorem*}

\begin{theorem*}
Let $G$ be a finite group, $Q(G)$ the additive semi-group of positive integers generated by the indices of proper subgroups of $G$ in $G$ and $R(G)$ the additive semi-group of positive integers generated by the orders of non-identity elements of G. Then there is a model $\mathcal{N}$ of $ZF^{0}$ in which $C(LO,\,n)$ is true for all $n\not\in Q(G)$, $C(\infty,\, n)$ is false for all $n\not\in R(G)$ and $C(\aleph_{0},\,n)$ is false for all $n\in Q(G)$.
\end{theorem*}

\begin{theorem*}
 Let $p$ and $q$ be primes such that $p_{1}=(p^{q}-1)/(p-1)$ is prime. Then there is a model $\mathcal{N}$ of $ZF^{0}$ in which
\begin{enumerate}
\item[\emph{1.}] $C(LO,\,n)$ is true for $n$ of the form $kp_{1}+rp^{q}$
\item[\emph{2.}] $C(\infty,\,n)$ is false for $n$ of the form $kp+rq$
\item[\emph{3.}] $C(\aleph_{0},\,n)$ is false for $n$ of the form $kp_{1}+rp^{q}$.
\end{enumerate}
\end{theorem*}

\begin{theorem*}
If for every prime $p$ there are infinitely many primes of the form $(p^{q}-1)/(p-1)$ then there is a model $\mathcal{N}$ of $ZF^{0}$ where $C(\aleph_{0},\,n)$ is true for every $n$ and $C(\infty,\,n)$ is false for every $n>1$.
\end{theorem*}

\textsc{Note.} The Mersenne Prime conjecture\index{Subject}{Mersenne prime conjecture} (there are infinitely many primes of the form $2^{n}-1$) is a special case of the hypotheses of this theorem.

Along the same lines Zuckerman\index{Author}{Zuckerman} \cite{bib:1971} proves

\begin{theorem*}
Let $Z$ be any set of primes and $P$ the additive semigroup of positive integers generated by Z. Then there is a model $\mathcal{N}$ of $ZF^{0}$ in which form 67 (the axiom of multiple choice, $MC(\infty,\,\infty)$) is true and for every $n\in \omega,\,n>1,\,C(\infty,\,n)$ is true if and only if $n\in P$.
\end{theorem*}

\section*{NOTE 16}

This note contains the definitions necessary to understand the forms in Brunner\index{Author}{Brunner} \cite{bib:1984d}, that is, form 126.

\begin{definition*}\
\begin{enumerate}
\item[1.] A $B$-\emph{lattice}\index{Subject}{$B$-lattice} is a vector lattice\index{Subject}{vector lattice} (see below for the definition of vector lattice) together with a Frechet-complete (Cauchy sequences are convergent) Riez-norm $\Vert\cdot\Vert$ (i.e., (i) $-y\leq x\leq y\rightarrow\Vert x\Vert\leq\Vert y\Vert$ (ii) if $\Vert x\Vert\leq 1$ there is a $y\geq 0$ such that $-y\leq x\leq y$ and $\Vert y\Vert\leq 1.)$
\item[2.] A \emph{positive linear functional} on $X$ is a linear mapping $f: X\rightarrow \mathbb{R}$ such that $f(x)\geq 0$ when $x\geq 0$.
\item[3.] A \emph{wedge}\index{Subject}{wedge}\index{Subject}{wedge in a linear space} $W$ in a real linear space $L$ is a convex set such that $t\subseteq W$ for every $t\geq 0\ (t\in \mathbb{R})$. A \emph{cone}\index{Subject}{cone in a linear space} is a wedge which contains no line through 0.
\item[4.] An \emph{ordered linear space}\index{Subject}{ordered linear space} (OSL) is a linear space $L$ in which some wedge $W$, called the positive wedge has been, specified. The partial ordering $\leq$ on $L$ is defined by $y\leq x$ iff $x-y\in W$. (Note that this differs from the definition of \emph{ordered vector space} in note 31.)

\item[ ] (Facts: (a) $\leq$ is reflexive \& transitive (b) translation and multiplication by positive numbers preserves order, multiplication by negative numbers reverses order, (c) ($x\geq y$ and $y\geq x$ implies $y=x$) if and only if the wedge is a cone. (d) $x\geq 0$ iff $x\in W$. If a relation satisfies (a) and (b) then (d) defines a wedge.)
\item[5.] Let $V$ be a partially ordered linear space with a positive wedge (defined). $V$ is a \emph{vector lattice} if the positive wedge in $V$ is a cone and if each set $\{x,\,x^{\prime}\}$ of two elements has a least upper bound $x\vee x^{\prime}$. These definitions are taken primarily from Day\index{Author}{Day} \cite{bib:1973}.
\end{enumerate}
\end{definition*}

\section*{NOTE 17}

\begin{definition*}
An $\epsilon$ lattice\index{Subject}{$\epsilon$ lattice} in a metric space ($X$,dist) is a maximal subset $L$ of $X$ which satisfies $(\forall x,\,y\in L)(\mathrm{dist}(x,\,y)\geq\epsilon$ for $x\neq y)$. From Brunner\index{Author}{Brunner} \cite{bib:1984e}. See [1 I].
\end{definition*}

\section*{NOTE 18}

This note contains a list of special forms including the following: 1. A list of forms that are transferable;\index{Subject}{transferable forms} 2. A list of forms whose negation may not be transferable; 3. A list of forms that are true in every permutation model; 4. A list of forms that are false in every permutation model; 5. A list of forms that are false in every Cohen model; and 6. A list of forms that depend on parameters.

\noindent 1. LIST OF FORMS THAT ARE TRANSFERABLE.

5, 6, 8, 9, 10, 11, 12, 13, 14, 17, 18, 30, 31, 32, 35, 36, 37, 38, 39, 40,
43, 44, 45, 46, 47, 48, 52, 60, 61, 64, 70, 71, 73, 74, 79, 80, 82, 83, 85, 86, 87, 88, 92, 93, 94, 98, 119, 120, 121, 122, 123, 128, 130, 132, 142, 163, 165, 166, 169, 170, 178, 185, 194, 197, 199, 203, 206, 211, 212, 213, 216, 217, 222, 223, 250, 251, 252, 272, 273, 280, 289, 295, 305, 306, 307, 309, 313, 325, 336, 342, 350, 357, 358, 361, 362, 363, 364, 366, 367, 368, 369, 373, 374, 376, 377, 378.

\noindent 2. LIST OF FORMS WHOSE NEGATION MAY NOT BE TRANSFERABLE.

69, 97, 104, 105, 108, 175, 182, 188, 189, 190, 191, 192, 193, 206, 209, 223, 267, 269, 270, 274, 280, 281,

\noindent 3. LIST OF FORMS THAT ARE TRUE IN EVERY PERMUTATION MODEL.

5, 6, 11, 12, 13, 35, 36, 37, 38, 70, 73, 74, 79, 91, 92, 93, 94, 108, 130, 142, 145, 169, 170, 194, 197, 199, 203, 206, 211, 212, 222, 223, 251, 252, 272, 273, 280, 289, 305, 306, 307, 309, 313, 361, 362, 363, 364, 366, 367, 368, 369. (Each one of these 53 forms is implied by one of the following 16 forms: 6, 11, 37, 73, 91, 92, 130, 273, 305, 309, 313, 361, 363, 368, 369.) (All but 91 are transferable. ) In addition, 182, 189, 190, and 191 are true in every permutation model in which the class of atoms is a set. (Each one of these four forms is implied by 191.)

\noindent4. LIST OF\index{Subject}{forms false in every FM model}\index{Subject}{forms true in every FM model} FORMS THAT ARE FALSE IN EVERY PERMUTATION MODEL.

15, 192 (if the collection of atoms is a set).

\noindent5. LIST OF FORMS THAT ARE FALSE IN EVERY COHEN MODEL.\index{Subject}{forms false in every Cohen model} (Each of these forms implies AC in ZF.)

66, 67, 89, 90, 91, 95, 109, 112, 114, 133, 147, 149, 164, 202, 218, 264, 292, 333, 334, 335 (Each one of these 20 forms implies form 91 and each is implied by one of the following 9 forms: 66, 95, 109, 133, 149, 218, 264, 334, 335.)

\noindent6. LIST OF FORMS THAT DEPEND ON PARAMETERS.

33, 45, 46, 47, 48, 59, 71, 81, 86, 87, 95, 120, 136, 137, 138, 174, 178, 179, 183, 196, 199, 207, 220, 271, 288, 308, 335, 336, 342, 373, 374.

\section*{NOTE 19}

It is clear that K\"{o}nig's lemma\index{Subject}{K\"{o}nig's lemma}, [10 F] implies [10 K] and using the usual proof that the Completeness Theorem implies the Compactness Theorem, it follows that [10 K] implies [10 L]. To show that [10 L] implies 10, $C(\aleph_{0},\,<\aleph_{0})$, let $X=\{x_{i}: i\in \omega\}$ be a denumerable set of finite sets, where each $x_{i}=\{x_{i,0},\,x_{i,1},\,\cdots x_{i,n_{i}}\}$. Let $\mathcal{L}$ be a first order language with a single function symbol $f$ and constant symbols for each $x_{i}\in X$ and each element $x_{i,j}$ of $\bigcup X$. Let $S$ be the following set of sentence in $\mathcal{L}$: For each $i\in\omega$, the conjunction of

$\wedge (\{x_{i,j}\neq x_{i,k}:j,\,k\leq n_{i}$ and $j\neq k\})$ and

$\vee\,(\{f(x_{i})=x_{i,0},\,\ldots,\,f(x_{i})=x_{i,n_{i}}\})$.

Each finite subset of $S$ has a model, so [10 L] implies that the denumerable set of sentences has a model. In that model, $f$ is a choice function on $X$.

\section*{NOTE 20}

We give some definitions and properties of inaccessible cardinals\index{Subject}{inaccessible cardinal}. (Proofs of the results given below can be found in Drake\index{Author}{Drake} \cite{bib:1974}.)

\begin{definition*}
A cardinal $\kappa$ is \emph{(strongly) inaccessible} if $\kappa$ is regular and for all $\mu<\kappa,\, 2^{\mu}<\kappa$. (An inaccessible cardinal is also called a \emph{strong limit cardinal}.)
\end{definition*}

\begin{lemma*}
No successor cardinal, $\aleph_{\alpha+1}$, is inaccessible. (This is form [0 AJ].)
\end{lemma*}

\begin{definition*}
$\mu$ is called a 2-\emph{valued}, $\kappa$-\emph{additive (complete) measure on} $x$ if $\mu$ is a function from $x$ to $\{0,1\}$ such that
\begin{enumerate}
\item[(a)] $\mu(x)=1$;
\item[(b)] If $y=\bigcup_{\alpha<\kappa}x_{\alpha}\subseteq x$ is a disjoint union, then $\mu(y)=\sum\nolimits_{\alpha\in \kappa}\mu(x_{\alpha})$.
\end{enumerate}
(The measure is called \emph{non-trivial} if the measure of a one-element set is 0.)
\end{definition*}

\begin{definition*}
A cardinal $\kappa$ is said to be \emph{measurable} if $\kappa >\aleph_{0}$ and $\kappa$ has a 2-valued $\kappa$-additive non-trivial measure.
\end{definition*}

\begin{theorem*}
($ZFC$) Every measurable cardinal\index{Subject}{measurable cardinal} is inaccessible. (Form 319.)
\end{theorem*}

\begin{definition*}
A filter is $\kappa$-\emph{additive (complete)} if the intersection of $<\kappa$ members of the filter is in the filter.
\end{definition*}

\begin{theorem*}
A cardinal $\kappa$ is measurable iff there is a non-principal ultrafilter on $\kappa$ which is $\kappa$-additive.
\end{theorem*}

\begin{definition*}
A measure $\mu$ on $\kappa$ is \emph{normal} if $\mu$ is a non-trivial $\kappa$-additive measure on $\kappa$ such that for any function $f:\kappa \rightarrow\kappa$ that is regressive (i.e., such that $f(\alpha)<\alpha$ for all $\alpha<\kappa,\,\alpha\neq 0$) there is a $\beta<\kappa$ such that $\mu(\{\alpha<\kappa:f(\alpha)=\beta\})=1$ (i.e., every regressive function on $\kappa$ is constant on a set of measure
1).
\end{definition*}

\begin{definition*}
If $X=\langle X_{\alpha}:\alpha<\kappa\rangle$ is $\kappa$-sequence of subsets of $\kappa$, the \emph{diagonal intersection}\index{Subject}{diagonal intersection}, $D(X)=\bigcap_{\alpha<\kappa}(X_{\alpha}\cup
\alpha+1)=\{\beta<\kappa:\forall\alpha<\beta(\beta\in X_{\alpha}\}$.
\end{definition*}

\begin{lemma*}
$\mu$ is normal iff the diagonal intersection of any $\kappa$-sequence of sets of measure 1 is also of measure 1.
\end{lemma*}

\begin{theorem*}
(ZFC) If $\kappa$ is measurable then $\kappa$ has a normal measure.
\end{theorem*}

\begin{definition*}
If $\kappa$ and $\lambda$ are cardinals, $\mathcal{L}_{\kappa,\lambda}$ is the language obtained by adding infinite conjunctions and disjunctions of cardinality less than $\kappa$ and infinite quantifications of blocks of variables of cardinal less than $\lambda$. (The ordinary first-order language is equivalent to $\mathcal{L}_{\omega,\omega}.$)
\end{definition*}

\begin{definition*}
\begin{enumerate}
\item[(a)] A cardinal $\kappa >\aleph_{0}$ is \emph{weakly compact}\index{Subject}{weakly compact cardinal} if the language $\mathcal{L}_{\kappa,\kappa}$ is $(\kappa,\,\kappa)$-compact, i.e., if $\Gamma$ is a set of sentences of $\mathcal{L}_{\kappa,\kappa}$ with cardinality $\leq\kappa$ and every subset of $\Gamma$ of cardinality $<\kappa$ has a model, then $\Gamma$ has a model.
\item[(b)] A cardinal $\kappa >\aleph_{0}$ is \emph{strongly compact} if $\mathcal{L}_{\kappa,\kappa}$ is compact, i.e., if $\Gamma$ is any set of formulas of $\mathcal{L}_{\kappa,\kappa}$ such that every subset of $\Gamma$ of cardinality $<\kappa$ has a model, then $\Gamma$ has a model.
\end{enumerate}
\end{definition*}

\begin{definition*}
\begin{enumerate}
\item[(a)] For $n\in \omega,\,[X]^{n}$ is the set of all $n$-element subsets of $X$, and $[X]^{<\omega}$, is the set of all finite subsets of $X$.
\item[(b)] The partition relation\index{Subject}{partition relation} $\kappa \rightarrow(\kappa)^{n}$ holds iff for every function $f:[\kappa]^{n}\rightarrow 2$ there is an unbounded subset $X$ of $\kappa$ such that $f/[X]^{2}$ is constant. (The set $X$ is called \emph{homogeneous for} $f$.)
\item[(c)] The partition relation $\kappa \rightarrow(\kappa)^{<\omega}$ holds iff for every function $f:[\kappa]^{<\omega}\rightarrow 2$ there is an unbounded subset $X$ of $\kappa$ such that $f$ is constant on $[X]^{n}$ for every $n\in \omega$.
\end{enumerate}
\end{definition*}

\begin{theorem*}
(ZFC) If $\kappa$ is inaccessible, then $\kappa$ is weakly compact iff $\kappa \rightarrow(\kappa)^{2}$.
\end{theorem*}

\begin{lemma*}
(ZFC) Every measurable cardinal is weakly compact.
\end{lemma*}

\noindent(Bull\index{Author}{Bull} \cite{bib:1978} defines a measurable cardinal as a cardinal that has a normal measure. With this definition, ZF $\vdash$ Every measurable cardinal is weakly compact.)

\begin{theorem*}
(ZFC) If $\kappa$ is a regular infinite cardinal, $\kappa$ is strongly compact iff for any $\kappa$-complete filter $J$ on a set $X,\,J$ can be extended to a $\kappa$-complete ultrafilter on $X$.
\end{theorem*}

\begin{lemma*}
(ZFC) Every strongly compact cardinal\index{Subject}{strongly compact cardinal} is measurable.
\end{lemma*}

\begin{definition*}
\begin{enumerate}
\item[(a)] $\mathcal{P}_{\kappa}(\lambda)$ is the set of all subsets of $\lambda$ of cardinality $<\kappa$.
\item[(b)] Let $X=\mathcal{P}_{\kappa}(\lambda)$,
\begin{enumerate}
\item[(i)] For each $\xi <\lambda$, let $C_{\xi}=\{x\in X:\xi\,\in x\}$. An ultrafilter extending the filter $\{C_{\xi}:\xi\,<\lambda\}$ is called a \emph{regular} ultrafilter\index{Subject}{regular ultrafilter}. (The measure corresponding to the regular ultrafilter is called a \emph{regular} measure\index{Subject}{regular measure}.)
\item[(ii)] A measure $\mu$ on $X$ is called \emph{normal} if it is $\kappa$-additive, regular and such that if $f(x)\in x$ for each $x\in X$, then for some $\xi <\lambda,\,\mu(\{x\in X:f(x)=\xi\})=1$.
\end{enumerate}
\end{enumerate}
\end{definition*}

\begin{definition*}
\begin{enumerate}
\item[(a)] If $\kappa$ and $\lambda$ are cardinals such that $\aleph_{0}<\kappa \leq\lambda,\,\kappa$ is $\lambda$-\emph{supercompact} if there is a normal measure\index{Subject}{normal measure} on $\mathcal{P}_{\kappa}(\lambda)$.
\item[(b)] If $\kappa >\aleph_{0},\,\kappa$ is \emph{supercompact} if $\kappa$ is $\lambda$-supercompact for every $\lambda\geq\kappa$.
\end{enumerate}
\end{definition*}

\begin{lemma*}
(ZFC) Every supercompact cardinal\index{Subject}{supercompact cardinal} is strongly compact.
\end{lemma*}

\begin{lemma*}
Form 319 implies form 320. (Use [0 AJ].)
\end{lemma*}

\begin{definition*}
$\kappa$ is a \emph{Ramsey cardinal}\index{Subject}{Ramsey cardinal} iff $\kappa \rightarrow(\kappa)^{<\omega}$.
\end{definition*}

\begin{lemma*}
Every measurable cardinal is a Ramsey cardinal.
\end{lemma*}

\begin{definition*} A cardinal $\kappa$ is called a \emph{Rowbottom cardinal}\index{Subject}{Rowbottom cardinal} iff, whenever $\mathfrak{A} = \langle A,\,U,\,\cdots\rangle$ is a structure of countable length, where $|A|=\kappa,\,U\subset A$, and $|U|<\kappa$, then
$\mathfrak{A}$ has an elementary substructure $\mathfrak{B} =\langle B,\,U^{\prime},\cdots\rangle$ with $|\mathfrak{B}|=\kappa$ and $|U^{\prime}|\leq\aleph_{0}$.
\end{definition*}

\begin{lemma*}
Every Ramsey cardinal is a Rowbottom cardinal.
\end{lemma*}

\begin{definition*} A cardinal $\kappa$ is said to be \emph{almost huge} iff there is an $n\in \omega$ and a cardinal $\lambda>\kappa$, a $\kappa$-additive normal ultrafilter $U$ on $2^{\lambda}$, and a sequence $\kappa =\lambda_{0}<\lambda_{1}<\cdots<\lambda_{n}=\lambda$ such that for each $i<n,\,\{x\subseteq\lambda:x\cap\, \lambda_{i+1}=\lambda_{i}\}\in U$.
\end{definition*}

\begin{lemma*}
Every almost huge cardinal\index{Subject}{almost huge cardinal} is supercompact.
\end{lemma*}

In a series of papers (Apter\index{Author}{Apter} \cite{bib:1985b}, \cite{bib:1988}, \cite{bib:1990}, \cite{bib:1992} and Apter/Henle\index{Author}{Apter/Henle} \cite{bib:1991}) various consistency results about large cardinals are obtained. We give some examples below.
\begin{enumerate}
\item[1.] Con(ZFC + There exists cardinals $\kappa <\lambda$ such $\kappa$ is $\lambda$-supercompact and $\lambda$ is measurable) $\rightarrow$ Con(ZF + $\mathrm{cof}(\kappa)=\omega+\kappa$ is a Rowbottom cardinal $+\ \kappa^{+}$ is a measurable cardinal).
\item[2.] Con(ZFC + There exists an almost huge cardinal) $\rightarrow$ Con(ZF
+ $\aleph_{1}$ is measurable + All uncountable cardinals are singular).
\item[3.] Con(ZFC + There exists an almost huge cardinal) $\rightarrow$ Con(ZF
+ $\neg C(\aleph_{0},\,\infty)\ +$ An uncountable cardinal is measurable iff it is the successor of a singular cardinal).
\item[4.] Con(ZFC + There exists an almost huge cardinal) $\rightarrow$ Con(ZF $+\neg C(\aleph_{0},\,\infty)+ \forall\alpha\,\aleph_{\alpha}$ or $\aleph_{\alpha+1}$ is measurable).
\item[] In Bull\index{Author}{Bull} \cite{bib:1978}, besides form 321, there are the following consistency results. Let $S$ be the following statement:
\end{enumerate}

$S$. There are cardinals $\aleph_{\alpha}<\aleph_{\beta}<\aleph_{\gamma}$ such that $\aleph_{\alpha}$ is $\aleph_{\gamma}$-supercompact; $\aleph_{\beta}$ is the first measurable cardinal greater than $\aleph_{\alpha}$; and $\aleph_{\gamma}=|2^{\aleph_{\beta}}|$.
\begin{enumerate}
\item[1.] Con(ZFC + $S$) $\rightarrow$ Con(ZF + form 43 + $\aleph_{1}$ is weakly compact + $\aleph_{2}$ is measurable).
\end{enumerate}

$S^{\prime}$.: There\index{Subject}{ramification system} are cardinals $\aleph_{\alpha}<\aleph_{\beta}$ such that $\aleph_{\beta}$ is the first weakly compact cardinal greater than $\aleph_{\alpha}$ and $\aleph_{\alpha}$ is $\aleph_{\beta}$-supercompact.
\begin{enumerate}
\item[2.] Con(ZFC + $S^{\prime}$) $\rightarrow$ Con(ZF + $C(\aleph_{\alpha},\,\infty)+\aleph_{\alpha}$ and $\aleph_{\alpha+1}$ are both weakly compact. (In fact, 2 may be extended to a finite sequence of weakly compact cardinals.)
\end{enumerate}

\section*{NOTE 21}

Definitions for forms [10 F], [14 M], [14 N], [43 H], [106 A], 216, and 249.

\begin{definition*}
\begin{enumerate}
\item[1.] A \emph{tree}\index{Subject}{tree} is a connected undirected graph without circuits one of whose vertices is designated as the origin.
\item[] Note that a tree may also be defined as a partially ordered set $(P,\,\leq)$ with a least element and with the property that for any element $x\in P$ the set of predecessors of $x$ is a finite set which is linearly ordered by $\leq$. Every tree in this sense is also a tree in the sense of definition 1 (under the obvious translation) and vice versa. Note also that the definition of \emph{tree} for form [86 A($\alpha$)] in note 119 differs from the definition given above.

\item[2.] The number of vertices on the unique path connecting a vertex $v$ with the origin is the \emph{level}\index{Subject}{level in a tree} of $v,\,l(v)$. (Thus the set of vertices of a tree can be decomposed into an at most denumerable set of levels.)
\item[3.] A vertex $v^{\prime}$ is a \emph{successor} of a vertex $v$ if $v$ and $v^{\prime}$ are connected by an edge and $l(v^{\prime})=l(v)+1$.

\item[4.] A tree is \emph{finite} if its set of vertices is finite and \emph{locally finite} if each vertex has only finitely many successors.
\item[5.] A \emph{branch} in a tree\index{Subject}{branch in a tree} is a maximal path beginning at the origin. If $v$ and $v^{\prime}$ are on the same branch, then $v^{\prime}$ \emph{dominates} $v$ if $l(v^{\prime})\geq l(v)$.
\item[] Konig's lemma states that any infinite locally finite tree\index{Subject}{locally finite tree} has an infinite branch.
\item[6.] An $\omega$-\emph{tree}\index{Subject}{$\omega$-tree} is a locally finite tree with at least one vertex of level $n$ for each $n\in \omega$.
\item[7.] A \emph{subtree}\index{Subject}{subtree} of a tree $(P,\,\leq)$ is a subset of $P$ which is a tree under the induced ordering and which is closed under predecessor.
\item[8.] An \emph{antichain} in a tree\index{Subject}{antichain in a tree} $(P,\,\leq)$ is a subset of $P$ whose elements are pairwise incomparable. (Note that this differs from the definition of \emph{antichain} used in form [86 A($\alpha$)] and given in note 119.)
\item[9.] Let $T$ be a collection of locally finite trees (not necessarily pairwise disjoint).
\item[] By a \emph{vertex} or a \emph{level} of $T$, we mean a vertex or a level of some tree in $T$.
\item[10.] If $v$ and $v^{\prime}$ are vertices of $T$, then $v^{\prime}$ \emph{dominates} $v$ in $T$ if $v^{\prime}$ dominates $v$ in some tree in T.
\item[11.] Let $S$ be a set of vertices of $T$. $S$ \emph{pierces} a level $l$ of $T$ if $|S\cap l|=1$.
\item[12.] $S$ is \emph{consistent} if for every $v,\,v^{\prime}$ in $S$ there is a $v^{\prime\prime}$ which dominates them both in T.
\end{enumerate}
\end{definition*}

\section*{NOTE 22}

Definitions for forms [1 L], [1 M], [1 N], [10 C] and [132 A]. A \emph{set mapping}\index{Subject}{set mapping} $f$ is a function $f:X\rightarrow \mathcal{P}(X)$ such that for all $x\in X,\,x\not\in f(x)$. A \emph{free subset}\index{Subject}{free subset} of $X$ (relative to $f$) is a subset $H$ of $X$ such that for all $x,\,y\in H,\,x\not\in f(y)$. A \emph{ramification system} is a triple $\langle X,\,E,\,S\rangle$ where $X$ and
$E$ are sets and $S\,:\,X^{\leq\theta}\rightarrow \mathcal{P}(E)$ (where $\theta$ is an ordinal and $X^{\leq\theta}$ denotes all functions $t:\alpha\rightarrow X$ where $\alpha\leq\theta$) satisfies (i) $S(\emptyset)=E$ (ii) If $f,\, g\in\mathrm{dom}(S)$ and $f\subseteq g$, then $S(g)\subseteq S(f)$ and (iii) If $f\in\mathrm{dom}(S)$ and $\eta=\mathrm{dom}(f)$ is a limit ordinal then $S(f)=\bigcap_{\alpha<\eta}S(f/\alpha)$.

\section*{NOTE 23}

Definitions for forms [14 Q] and [52 E]: If $E$ is a Banach space\index{Subject}{Banach space} (complete, normed, linear space) then $E^{\ast}$ (the \emph{dual} of $E$) is the space of all bounded linear functionals on $E$. (That is, all $f:E\rightarrow\,\mathbb{R}$
such that $(\exists M)(\forall x\in E)(|f(x)|\leq\Vert x\Vert M)$. $E^{\ast}$ is a Banach space if we define $\Vert f\Vert=\sup_{x\neq 0}\Big(\frac{|f(x)|}{\Vert x\Vert}\Big)$. Each $x\in E$ can be thought of as a linear functional on $E^{\ast}$ if we define $x(f)=f(x)$ for all $f$ in $E^{\ast}$. The \emph{weak}$^{\ast}$ \emph{topology}\index{Subject}{weak$^{\ast}$ topology} on $E^{\ast}$ is the weakest topology that makes all these linear functionals continuous. Sets of the form $w_{x,\epsilon}=\{f\in E^{\ast}:|f(x)|<\epsilon\}$ form a basis for this topology. If $X\subseteq E,\,X$ is \emph{convex-compact}\index{Subject}{convex-compact subset of a topological
vector space} if whenever $F_{i}$, for $i\in I$, are closed convex subsets of $X$ and $\{X\,\cap\,F_{i}:i\in I\}$ has the finite intersection property, then $\bigcap_{i\in I}(X\cap\,F_{i})\neq\emptyset$. Note that the definition of \emph{convex-compact} may be given in any topological vector space.

\section*{NOTE 24}

This note contains some definitions from group theory that are used in $\mathcal{M}22$, and forms [62 C], [62 D], [67 D], 180, and 308.

\begin{definition*}
Let $G$ be a group and let $p$ be a prime.
\begin{enumerate}
\item[(a)] The \emph{Frattini subgroup of}\index{Subject}{Frattini subgroup} $G,\,\Phi(G)$, is the intersection of all maximal subgroups of $G$.
\item[(b)] Let $Z(G)$ be the center of $G$ and $G^{\prime}$, the commutator subgroup of $G$.
\item[(c)] $G$ is said to be a $p$-\emph{group}\index{Subject}{$p$-group} if for all $x\in G$ there is an $n\in \omega$ such that $x^{p^{n}}=1$. $G$ is said to have \emph{exponent} $p$ if for every $x\in G,\,x^{p}=1$. An Abelian group is said to be \emph{elementary} if there is a prime $p$ such that every non-identity element has order $p$.
$\qquad G$ is said to be \emph{divisible}\index{Subject}{divisible group} if for each $a\in G$ and for each integer $m\neq 0$ the equation $x^{m}=a$ has a solution in $G$ and $G$ is called $n$-\emph{divisible}\index{Subject}{$n$-divisible group} if the equation $x^{m}=a$ has a solution for all $m<n$.
\item[(d)] A factor $H/K$ of a series of $G$ is said to be a \emph{central factor}\index{Subject}{central factor of a group} of $G$ if $K$ is normal in $G$ and $H/K$ is a subgroup of $Z(G/K)$.
\item[(e)] $G$ is said to be \emph{nilpotent}\index{Subject}{nilpotent group} if it has a series all of whose factors are central factors of $G$.
\item[(f)] $Z_{n}(G)$ is defined as follows: $Z_{0}(G)=1$ and for each integer $n>0$,
\begin{equation*}
Z_{n}(G)/Z_{n-1}(G)=Z(G/Z_{n-1}(G)).
\end{equation*}
The \emph{class} of $G$ is the least integer $n$ such that $Z_{n}(G)=G$.
\item[(g)] $G$ is a \emph{special} $p$-\emph{group} if $G$ is a $p$-group, $G^{\prime}=Z(H)=\Phi(H)$, and $G$ is either elementary Abelian\index{Subject}{elementary Abelian group} or nilpotent of class 2.
\item[(h)] $G$ is an \emph{extra-special} $p$-\emph{group}\index{Subject}{special $p$-group} if $G$ is special and $G^{\prime}$ is cyclic of order $p$.
\item[(i)] A group is \emph{divisible} if each element is divisible by each positive integer.
\item[(j)] A group is \emph{torsion free}\index{Subject}{torsion free group} if the only element of finite order is the identity element.
\end{enumerate}
\end{definition*}

\section*{NOTE 25}

For forms 144 and 179: Let $W(0)$ be the class of well orderable sets and let $W(\alpha+1)=\{\bigcup Q:Q\subseteq W(\alpha)$ and $Q\in W(0)\}$ and $W(\lambda)=\bigcup\{W(\beta):\beta<\lambda\}$ for limit ordinals $\lambda$. Then $W$ (the class of almost well orderable sets\index{Subject}{almost well orderable sets}) is defined to be $\bigcup\{W(\alpha):\alpha\in\, On\}$. See Keisler\index{Author}{Keisler} \cite{bib:1970}. Blass\index{Author}{Blass} \cite{bib:1977b} has shown that If $P\subseteq W$ and $P\in W$, then $\bigcup P\in W$. Regarding the form(s) 179 $(\forall x,\,x\in W(\alpha))$, R. David\index{Author}{David} \cite{bib:1980} shows that for ordinals $\alpha$ and $\beta$ with $0<\alpha<\beta$, the statement $\forall x,\,x\in W(\beta)$ does not imply $\forall x,\,x\in W(\alpha)$. (Clearly,
$\forall x,\,x\in W(\alpha)$ implies $\forall x,\,x\in W(\beta).)$ It is clear that 30 + 294 (every linearly ordered $W$-set\index{Subject}{$W$-set} can be well ordered) implies that every $W$-set can be well ordered. Thus, 30 + 294 + 144 (every set is a $W$-set) implies AC.

\section*{NOTE 26}

A summary of the definitions and results from Brunner\index{Author}{Brunner} \cite{bib:1983d}. Brunner considers 16 properties a topological space $(X,\,\underline{X})$ may have:
\begin{enumerate}
\item[W1] : $\underline{X}$ is w.o. (well ordered);
\item[W2] : $X$ is w.o.;
\item[D1] : $X$ is covered by a w.o. family of closed, discrete sets.
\item[D2] : $X$ is covered by a w.o. family of discrete sets;
\item[B1] : $X$ has a w.o. base
\item[B2] : $\exists f:X\times W\rightarrow\underline{X}$ such that $W$ is w.o. and $f^{\prime}\{x\}\times W)$ is a neighborhood base at $x$ for each $x\in X$.
\item[B3] : $\exists f:X\times W\rightarrow\underline{X}$ such that $W$ is w.o. and $\forall x\in X,\,\{x\}=\bigcap f^{\prime}\{x\}\times W)$
\item[L1] : Each open cover of $X$ has a w.o. subcover;
\item[L2] : Each open cover of $X$ has a w.o. refinement;
\item[S] : $X$ has a dense, w.o. subset;
\item[C] : Each family $O\subseteq\underline{X}$ of pairwise disjoint sets is w.o.
\item[A1] : If $O\subseteq\underline{X}$ covers $X,\,\exists f:X\rightarrow O$ such that $\forall x\in X,\, x\in f(x)$;
\item[A2] : If $O\subseteq\underline{X}$ covers $X,\,\exists f:X\rightarrow\underline{X}$ s.t. $\forall x\in X,\,x\in f(x)$ and $f^{\prime}X$ refines $O$;
\item[H1] : $X$ is hereditarily $A1$;
\item[H2] : $X$ is hereditarily $A2$;
\item[F] : $X$ is a continuous finite to one image of an $A1$ space.
\end{enumerate}

\begin{definition*}
If $P$ is one of the above then $A(P)$ is the statement: ``Every topological space $(X,\, \underline{X})$ is P.'' $A(P,\,Q)$ is the statement: ``If a topological space satisfies $P$ then it satisfies $Q$.''
\end{definition*}

All of the $A(P)$'s occur in the list of forms, but not all the $A(P,\,Q)$'s occur.

The table below, taken from Brunner \cite{bib:1983d}, summarizes the results. In the table, $0$ means that $A(P,\, Q)$ ($P$ the row heading and $Q$ the column heading) is provable in $\mathrm{ZF}^{0}$. $N$ means $A(P,\,Q)$ is not provable in $\mathrm{ZF}^{0}$. $AC$ (form 1), $MC$ (form 67) and $PW$ (form 91) mean respectively that $A(P,\,Q)$ is equivalent to $AC,\,MC$ or $PW$ in $\mathrm{ZF}^{0}$. And $IMC$ or $IPW$ mean that $A(P,\,Q)$ implies $MC$ or $PW$ respectively in $\mathrm{ZF}^{0}$.

\begin{table}[h]
\begin{tabular}{l|llllllllllllllll}
 & W1 & W2 & D1 & D2 & B1 & B2 & B3 & L1 & L2 & S & C & H1 & H2 & A1 & F & A2 \\ \hline
W1 & 0 & 0 & 0 & 0 & 0 & 0 & 0 & 0 & 0 & 0 & 0 & 0 & 0 & 0 & 0 & 0\\
W2 & PW & 0 & 0 & 0 & ? & ? & 0 & N & ? & 0 & 0 & N & ? & N & N & ? \\
D1 & AC & AC & 0 & 0 & AC & N & 0 & AC & AC & AC & AC & AC & N & AC & MC & N \\
D2 & AC & AC & IMC & 0 & AC & MC & MC & AC & AC & AC & AC & AC & MC & AC & MC & MC\\
B1 & PW & IPW & IPW & IPW & 0 & 0 & 0 & N & 0 & N & 0 & N & 0 &  N & N & 0 \\
B2 & AC & AC & IPW & IPW & AC & 0 & 0 & AC & AC & AC & AC & AC & 0 & AC & MC & 0 \\ B3 & AC & AC & IPW & IPW & AC & N & 0 & AC & AC & AC & AC & AC & N & AC & MC & N \\ L1 & AC & AC & MC & N & AC & MC & MC & 0 & 0 & AC & AC & AC & MC & 0 & 0 & 0 \\
L2 & AC &  AC &  MC &  IPW &  AC &  MC &  MC &  N & 0 & AC & AC & AC & MC & N & N & 0 \\
S & PW & IPW & IPW & IPW & N & N & ? & N & N  & 0 & 0 & N & N & N & N & N\\ C & IPW & PW & PW & PW & N & N & N & N & N & N & 0 & N & N & N & N & N \\ H1 & N & N & N & ? & N & N & N & N & N & N & N & 0 & 0 & 0 & 0 & 0 \\
H2 & AC & AC & AC & IPW & AC & MC & MC & AC & AC & AC & AC & AC & 0 & AC & MC & 0\\
A1 & AC & AC & MC & N & AC & MC & MC & N & N & AC & AC & AC & MC & 0 & 0 & 0 \\
F & AC & AC & MC & N & AC & MC & MC & N & N & AC & AC & AC & MC & N & 0 & 0 \\
A2 & AC & AC &  MC & IPW & AC & MC & MC & AC & AC & AC & AC & AC & MC & AC & MC & 0 \\
\end{tabular}
\end{table}

\begin{theorem*}
\begin{enumerate}
\item[\emph{(i)}] $ZF\vdash A(P)\leftrightarrow AC$ for all $P$ in the list.
\item[\emph{(ii)}] $ZF^{0}\vdash A(P)\leftrightarrow AC$ for $P\in\{W1,\,W2,\,B1,\,L1,\,L2,\,S,\,C,\,H1,\,A1\}$.
\item[\emph{(iii)}] $ZF^{0}\vdash A(P)\leftrightarrow MC$ (form 67) for $P\in\{D1,\,B2,\,B3,\,H2,\, A2\}$.
\item[\emph{(iv)}] $ZF^{0}\vdash AC\rightarrow A(F)\rightarrow MC$ (form 67) $\rightarrow A(D2)\rightarrow PW$ (form 91).
\end{enumerate}
\end{theorem*}

\section*{NOTE 27}

Notes on $UT(WO,\,\kappa,\,WO),C(WO,\,\kappa)$ and $D_{\kappa}$ from Brunner/Howard\index{Author}{Brunner}\index{Author}{Brunner/Howard} \cite{bib:1992}.

\begin{definition*}
$D_{\kappa}$ or $D(\kappa)$ is the sentence ``Every non-well orderable set is the union of a disjoint, well orderable family of $\kappa$ element sets.
\end{definition*}

(Note that $UT(WO,\,\kappa,\,WO)$, $C(WO,\,\kappa)$, and $D(\kappa)$ all depend on the well ordered, infinite cardinal $\kappa$ and that $UT(WO,\,<\aleph_{0},\,WO)$, $UT(WO,\,<\aleph_{1},\,WO)$, and $D(\aleph_{0})$ are forms [122 A], 151 and 152 respectively.)

\begin{theorem*}
($ZF^{0}$) $UT(WO,\,\kappa,\,WO)$ implies $C(WO,\,\kappa)$ but does not imply $C(\aleph_{0},\, \kappa^{+})$ nor does it imply $D(\kappa^{+})$.
\end{theorem*}

\begin{theorem*}
($ZF^{0}$) $D(\kappa^{+})$ does not imply $D(\kappa)$.
\end{theorem*}

\section*{NOTE 28}

In this note we give definitions for the various versions of the Baire category theorem and summarize the results from Brunner \cite{bib:1983c}.

\begin{definition*}
Let $(X,\,T)$ be a topological space.
\begin{enumerate}
\item[1.] A set $Y\subseteq X$ is \emph{nowhere dense} if the closure of $Y$ has empty interior.
\item[2.] A set $Y\subseteq X$ is \emph{meager}\index{Subject}{meager set} or \emph{of the first category}\index{Subject}{first category set} if $Y$ is a countable union of nowhere dense sets\index{Subject}{nowhere dense set}.
\item[3.] A set $Y\subseteq X$ is \emph{perfect}\index{Subject}{perfect set} if $Y$ is closed, non-empty and has no isolated points.
\item[4.] A set $Y\subseteq X$ has the \emph{Baire property}\index{Subject}{Baire property} if $(Y\backslash U)\cup(U\backslash Y)$ is meager for some open set $U$.
\item[5.] $(X,\,T)$ is \emph{Baire} or is a \emph{Baire space}\index{Subject}{Baire space} if the intersection of each countable sequence of dense, open sets in $X$ is dense in $X$.
\item[6.] $(X,\,T)$ is \emph{sequentially complete} if every sequence has a convergent subsequence.
\item[7.] $(X,\,T)$ is \emph{\v{Cech complete}}-$I$ if it satisfies
\begin{enumerate}
\item[(i)] $(X,\,T)$ is regular (every neighborhood of a point $x\in S$ contains a closed neighborhood of $x$) and
\item[(ii)] there is a denumerable collection $\{\mathcal{C}_{n}:n\in \omega\}$ of open covers of $X$ such that for any collection $F$ of closed subsets of $X$ with the finite intersection property, if for each $n,\,F$ contains a subset of diameter less than $\mathcal{C}_{n}$ (that is, $(\exists u\in F)(\exists c\in \mathcal{C}_{n})(u\subseteq c))$, then $F$ has a non-empty intersection.
\end{enumerate}
(This is the definition of \v{C}ech complete used in form [43 E].)
\item[8.] $(X,\,T)$ is \emph{\v{C}ech complete-II} if ($X,\, T)$ is homeomorphic to a $G_{\delta}$ set in a compact, $T_{2}$ space. ($ G_{\delta}\equiv$ countable intersection of open sets.)
\item[] (This is the definition of \v{C}ech complete used in $K9$ below. $K9$ is equivalent to form 106.)
\item[9.] ($X,\, T)$ is \emph{co-compact}\index{Subject}{co-compact space} if $\exists$ a family $F$ of closed sets such that
\begin{enumerate}
\item[(i)] If $G\subseteq F$ and $G$ has the finite intersection property, then $\bigcap G\neq\emptyset$ and
\item[(ii)] If $x\in \mathcal{O}$, with $\mathcal{O}$ open, then $\exists A\in F$ with $x\in A^{\circ}$ (the interior of $A$) and $A\subseteq$ closure $(\mathcal{O}).)$
\end{enumerate}
\item[10.] $(X,\,T)$ is \emph{scattered}\index{Subject}{scattered space} or \emph{Boolean}\index{Subject}{Boolean space} if it is Hausdorff and has a basis of clopen sets.
\end{enumerate}
\end{definition*}

\begin{definition*}
Let ($X, d$) be a metric space.
\begin{enumerate}
\item[1.] ($X, d$) is \emph{complete} or \emph{Frechet complete}\index{Subject}{Frechet complete space} if Cauchy sequences converge.
\item[2.] ($X, d$) is \emph{Cantor complete} if $\bigcap_{n\in\omega}A(n)\neq\emptyset$ for every sequence $(A(n))_{n\in\omega}$ of non-empty closed sets satisfying
\begin{enumerate}
\item[(i)] $(\forall n\in \omega)(A(n+1)\subseteq A(n))$ and
\item[(ii)] $\lim_{n\rightarrow\infty}$ diameter$(A(n))=0$.
\end{enumerate}
\item[3.] ($X, d$) is \emph{Kuratowski complete} if $\bigcap_{n\in\omega}A(n)\neq\emptyset$ for each nested sequence $(A(n))_{n\in\omega}$ of non-empty closed sets for which $a(A(n))\rightarrow 0$ where for $A\subseteq X,\, a(A)=\inf\{\max\{\mathrm{diameter}(B):B\in Y\}:Y$ is a finite covering of $A$\}.
\item[4.] ($X, d$) is \emph{filter complete}\index{Subject}{filter complete space} if every Cauchy filter converges.
\begin{enumerate}
\item[$K1$] : Every Frechet Complete metric space is Baire. (This is what is usually known as the Baire Category Theorem and is [43 F] in the list of forms.)
\item[$K2$] : Every Cantor complete metric space is Baire. ([43 J])
\item[$K3$] : Every metric Kuratowski complete space\index{Subject}{Kuratowski complete space} is Baire.
\item[$K4$] : Every sequentially compact metric space is Baire.
\item[$K5$] : Every countably compact (Every countable open cover has a finite subcover.) metric space is Baire.
\item[$K6$] : Banach spaces (complete, normed, linear spaces) are Baire.
\item[$K7$] : Countably compact, regular spaces are Baire.
\item[$K8$] : Compact $T_{2}$-spaces are Baire. (This is form 106)
\item[$K9$] : \v{C}ech complete\index{Subject}{\v{C}ech complete space}-II spaces are Baire.
\item[$K10$] : Co-compact, $T_{2}$ spaces are Baire.
\item[K11] : Separable, Frechet complete, metric spaces are Baire.
\item[$K12$] : Separable, countably compact, regular spaces are Baire.
\end{enumerate}
\end{enumerate}

The following results are proved in $\mathrm{ZF}^{0}:K1$ and $K2$ are equivalent to $DC$. $DC$ (form 43) implies $K9$ and $K10$. $K1$ implies $K6$. $DMC$ (Dependent Multiple choice, Form [106 A]) implies $K3$ implies $K5$. $K4$ for Cantor complete spaces\index{Subject}{Cantor complete space} implies $DF=F$ (form 9). $K11$ is provable in $\mathrm{ZF}^{0}$. $DMC$ (form [106 A]) implies $K7$ implies $K8$. $DMC$ implies $K9$ implies $K8$. $DMC$ implies $K10$ implies $K8$. $PW$ (form 91) implies $K12$.

The following independences are shown to hold in $\mathrm{ZF}^{0}:MC\not\rightarrow K4,\, MC\not\rightarrow K6$, BPI (Form 14) $\not\rightarrow K8$ (Form 106), $MC_{\mathrm{LO}}$ (form 112) $\not\rightarrow K7,\,C(\aleph_{0},\,\infty) \not\rightarrow K8$ and none of the following are provable in ZF: $Ki,\,i=1,2,4,6,7,8,9,10$.

The following questions are asked:

Are the following provable in $\mathrm{ZF}^{0}$?: $K3,\,K5,\,K12$.

Are $Ki,\,7\leq i\leq 10$, equivalent to $DMC$ (form [106 A]).

In Fossy/Morillon\index{Author}{Fossy/Morillon} \cite{bib:1995} it is shown that $K8$ implies $DMC$ from which it follows that $Ki,\,7\leq i\leq 10$ are all equivalent to $DMC$.
\end{definition*}

\section*{NOTE 29}

Definitions for forms [14 R] through [14 U], [14 BH] through [14 BL], and [86 B]. See Banaschewski\index{Author}{Banaschewski} \cite{bib:1981}, \cite{bib:1983}, \cite{bib:1997}, Paseka\index{Author}{Paseka} \cite{bib:1989}, and Vickers\index{Author}{Vickers} \cite{bib:1989}.

\begin{definition*}
\begin{enumerate}
\item[1.] A \emph{local lattice}\index{Subject}{local lattice} (or \emph{frame}\index{Subject}{frame} or \emph{complete Heyting algebra}\index{Subject}{Heyting algebra}) is a complete lattice $L$ satisfying the distribution law: $x\bigvee_{i\in I}x_{i}=\bigvee_{i\in I}(x\wedge x_{i})$, for all $x,\,x_{i}\in L$ and $\{x_{i}:i\in I\}\subseteq L$. (Example: $\mathcal{D}X$ - the open sets of a topological
space)
\item[2.] A lattice $L$ is \emph{spatial}\index{Subject}{spatial lattice} if for some topological space $X,\,L$ is isomorphic to $\mathcal{D}X$. (Theorem. Every complete Boolean algebra is a local lattice. A complete Boolean algebra is a spatial iff it is atomic.)
\item[3.] $c\in L$ ($L$ a local lattice) is \emph{compact}\index{Subject}{compact local lattice} if and only if $c\leq \bigvee_{i\in I}x_{i}$ implies $c\leq x_{i_{1}}\vee\ldots\vee x_{i_{n}}$, for some indices $i_{1},\,\ldots,\,i_{n}$.
\item[4.] $L$ is \emph{compact} iff its top element $e$ is compact.
\item[5.] $L$ is \emph{coherent}\index{Subject}{coherent local lattice} iff (i) Each element of $L$ is the join of compact elements and (ii) the meet of any finite set, including the empty set, of compact elements is compact.
\item[6.] $L$ is \emph{regular}\index{Subject}{regular local lattice} iff for each $x\in L,\,x$ is the join of all $y\ll x$, where $y\ll x$ means $(\exists z\in L)(y\wedge z=0$ and $x\vee z=e)$.
\item[7.] $L$ is 0 \emph{dimensional} iff for each $x\in L,\,x=\bigvee\{y\in L:y^{\ast}\vee y=e$ and $y\leq x\}$. $(y^{\ast}$ is the pseudo-complement
of $y,\,y^{\ast}=\bigvee\{c\in L\,:\,c\wedge a=0\}.)$
\item[8.] $L$ is \emph{conjunctive}\index{Subject}{conjunctive local lattice} iff $\forall a,\,b\in L$, such that $a<b,\,\exists c\in L$ such that $a\vee c<e$ and $b\vee c=e$.
\item[9.] A closed set $F$ of a topological space $X$ is \emph{irreducible}\index{Subject}{irreducible set in a topological space} iff whenever $G_{i},\,1\leq i\leq n$, are closed sets with $F\subseteq \bigcup_{i=1}^{n}G_{i}$, then $F\subseteq G_{i}$ for some $i$.
\item[10.] If $F$ is a closed set in a topological space $X$, then $x\in X$ is a \emph{generic point for} $F$ iff $F$ is the closure of $\{x\}$.
\item[11.] A topological space is \emph{sober}\index{Subject}{sober topological space} iff every irreducible closed set has a unique generic point\index{Subject}{generic point}.
\item[12.] The \emph{sum of two frames} $L_{1}$ and $L_{2}$ is the lattice whose set is the Cartesian product of the two frames with operations done componentwise.
\item[13.] A ring $R$ is \emph{regular} iff $(\forall a)(\exists b)(aba=a)$. (It is known that every prime ideal in a commutative regular ring\index{Subject}{regular ring} is maximal.)
\item[14.] A \emph{Gelfand algebra}\index{Subject}{Gelfand algebra} is a commutative Banach algebra $A$ with unit 1 over the real numbers $\mathbb{R}$ such that $\Vert a^{2}\Vert=\Vert a\Vert^{2}$ and $1+a^{2}$ is invertible for all $a\in A$.
\item[15.] A $\aleph_{\alpha}$-\emph{frame} is a partially ordered set $A$ with zero (0) and unit (e) in which every pair of elements has a meet and every subset of cardinality $\leq\aleph_{\alpha}$ has a join and which satisfies the distributive law: $a\wedge \bigvee S=\bigvee\{a\wedge t:t\in S\}$ for any $a\in A$ and subset $S$ of $A$ of cardinality $\leq\aleph_{\alpha}$.
\item[16.] A \emph{cover} of a lattice $L$ is a subset $S$ of $L$ such that $\bigvee S=e$ where $e$ is the unit in the lattice. A \emph{subcover} of $S$ is a subset of $S$ which is still a cover.
\item[17.] A lattice $L$ is $\aleph_{\alpha}$-\emph{Lindel\"{o}f} if every cover of $L$ has a subcover of cardinality $\leq\aleph_{\alpha}$.
\item[18.] For any $\aleph_{\alpha}$-frame $A$ and ideal $J$ of $A$, let $k_{0}(A)=\{\bigvee S:S\subseteq A\wedge|S|\leq\aleph_{\alpha}\}$. Then \{$J:J$ is an ideal in $A$ and $k_{0}(J)=J$\} ordered by $\subseteq$ is a frame and is called \emph{the frame envelope of}\index{Subject}{frame envelope} $A$.
\end{enumerate}
\end{definition*}

\section*{NOTE 30}

This note contains the proof that forms 14 (The Boolean Prime Ideal Theorem), [14 BM], [14 BN], and [14 BO] are equivalent. ([14 BM]: A system of equations over a finite ring $R$ has a solution in $R$ if and only if every finite sub-system has a solution in $R$. [14 BN]: A system of equations over a finite field $F$ has a solution in $F$ if and only if every finite sub-system has a solution in $F$. [14 BO]: A system of equations over a finite structure $(A,\,+,\,\cdot)$ has a solution in $A$, if and only if every finite sub-system has a solution in $A$.) First, it is clear that [14 BO] implies [14 BM] implies [14 BN]. We shall show that [14 J] implies [14 BO] and that [14 BN] implies 14.

\begin{theorem}\label{part04:thm1}
[14 J] implies [14 BO].
\end{theorem}

\begin{proof}
Let $A$ be a finite set with two operations + and $\cdot$ defined on it. Suppose $\Sigma$ is a set of equations over $A$ such that every finite subset of $\Sigma$ has a solution. Let $V$ be the set of variables which occur in $\Sigma$. Since $A$ is finite the discrete topology on $A$ is compact. By [14 J], the product $P=\prod_{i\in V} A$ is compact. For each finite subset $\sigma$ of $\Sigma$, let $x_{\sigma}=\{f\in P:f$ is a solution for $\sigma$\}. Let $B=\{x_{\sigma}:\sigma$ is a finite subset of $\Sigma$\}. The set $B$ is a set of closed subsets of $P$ with the finite intersection property. Since $P$ is compact, the intersection of $B$ is non-empty. Any element of the intersection is a solution of the system $\Sigma$.
\end{proof}

\begin{theorem}\label{part04:thm2}
[14 BN] implies 14.
\end{theorem}

\begin{proof}
Let $B$ be a Boolean algebra and for each $b\in B$, let $x_{b}$ be a variable. For $b\in B$, let $b^{\prime}$ be the complement of $b$. For each $b,\,c\in B$, consider the following set of equations $\Sigma$ over the two element field $\{0,1\}$:
\begin{equation*}
x_{b}+x_{b}^{\prime}=1,\qquad\quad x_{b\cap c}=x_{b}\cdot x_{c}.
\end{equation*}
It is easy to see that each finite subset of $\Sigma$ has a solution. [14 BN] implies that $\Sigma$ has a solution. Therefore, the set $\{b\in B:x_{b}=1\}$ is an ultrafilter.
\end{proof}

\section*{NOTE 31}

Definitions for forms [14 W], [70 A], [52 H] through [52 L], 372 and [372
A] through [372 D]. These are modifications of definitions from Schechter\index{Author}{Schechter} \cite{bib:1996a} and \cite{bib:1996b}.
\begin{enumerate}
\item[1.] A \emph{pre-order}\index{Subject}{pre-order} is a reflexive, transitive relation.
\item[2.] A \emph{directed set}\index{Subject}{directed set} $(\mathbb{D},\,\leq)$ is a non-empty, preordered set in which each finite subset has an upper bound.
\item[3.] A \emph{net}\index{Subject}{net} in a set $X$ is any function from a directed set into $X$. $(x:\mathbb{D}\rightarrow X$, usually denoted $(x(\delta):\delta\in \mathbb{D})$ or $(x_{\delta}))$.
\item[4.] If $x:\mathbb{D}\rightarrow X$ is a net and $S$ is any set we say $x$ is \emph{eventually in} $S$ if $(\exists\delta_{o}\in \mathbb{D})(\forall\delta\in \mathbb{D})(\delta\geq\delta_{o}\rightarrow x(\delta)\in S)$ and we say $x$ is \emph{frequently in} $S$ if $(\forall\delta\in \mathbb{D})(\exists\delta^{\prime}\in \mathbb{D})(\delta^{\prime}\geq\delta$ and $x(\delta^{\prime})\in S)$
\item[5.] $x$ is \emph{frequently constant} if for some $c,\,x$ is frequently in $\{c\}$.
\item[6.] Assume $(x(\alpha):\alpha\in \mathbb{A})$ and $(y(\beta):\beta\in \mathbb{B})$ are two nets in $X$, then $(y(\beta))$ is a \emph{subnet}\index{Subject}{subnet} of $(x(\alpha))$ (in the sense of Aarnes and Andenees - (This is not the usual definition of subnet.)) if $(\forall S\subseteq X)$(If ($x(\alpha))$ is eventually in $S$, then $(y(\beta))$ is eventually in $S$). (Which is equivalent to $(\forall S\subseteq X)($ If $(y(\beta))$ is frequently in $S$ then $(x(\alpha))$ is frequently in $S$). )
\item[7.] A \emph{universal net}\index{Subject}{universal net} (or \emph{ultranet})\index{Subject}{ultranet} on a set $X$ is a net $(x(\delta):\delta\in \mathbb{D})$ with the property that $(\forall S\subseteq X)$ ($x$ is eventually in $S$ or $x$ is eventually in $X\backslash S$).
\item[8.] A partially ordered set is \emph{Dedekind complete}\index{Subject}{Dedekind complete partially ordered set} if each subset that has an upper bound has a least upper bound. (Equivalently, if each subset that has a lower bound has a greatest lower bound.)
\item[9.] An \emph{ordered vector space}\index{Subject}{ordered vector space} is a real vector space $Z$ equipped with a partial ordering $\preccurlyeq$ which also satisfies:
\begin{enumerate}
\item[ ] $x\preccurlyeq y$ implies $x+z\preccurlyeq y+z$
\item[ ] $x\succcurlyeq 0$ and $r\geq 0$ imply $rx\succcurlyeq 0$.
\end{enumerate}

(Note: This differs from the definition of \emph{ordered linear space} in note 16.)

\noindent In definitions 10 through 13, $X$ is a vector space, $(Z,\preccurlyeq)$ is an ordered vector space and $p:X\rightarrow Z$.
\item[10.] $p$ is \emph{sublinear}\index{Subject}{sublinear function} if $p(x+y)\preccurlyeq p(x)+p(y)$ and for all $c\in[0,\,+\infty)$, $p(cx)=cp(x)$.
\item[11.] $p$ is \emph{convex}\index{Subject}{convex function} if $p(cx+(1-c)y)\preccurlyeq cp(x)+(1-c)p(y)$ for all $c\in[0,1]$.
\item[12.] $p$ is \emph{concave}\index{Subject}{concave function} if $p(cx+(1-c)y)\succcurlyeq cp(x)+(1-c)p(y)$ for all $c\in[0,1]$.
\item[13.] $p$ is \emph{affine}\index{Subject}{affine function} if $p(cx+(1-c)y)=cp(x)+(1-c)p(y)$ for all $c\in[0,1]$.
\item[14.] Let $(\Delta,\,\sqsubseteq)$ be a directed set and $(Z,\preccurlyeq)$ be a Dedekind complete ordered vector space. Let $B(\Delta,\,Z)$ be the set of functions $f$ from $\Delta$ to $Z$ such that the range of $f$ has both an upper bound and a lower bound. For $f\in B(\Delta,\,Z)$ let
\begin{enumerate}
\item[]$\liminf(f)=\sup_{\alpha\in\Delta}\inf_{\alpha\sqsubseteq\beta}f(\beta)$ and
\item[] $\limsup(f)=\inf_{\alpha\in\Delta}\sup_{\alpha\sqsubseteq\beta}f(\beta)$
\end{enumerate}
(Since $Z$ is Dedekind complete both $\liminf(f)$ and $\limsup(f)$ exist in $Z$ and $\lim\inf(f)\preccurlyeq \limsup(f).)$ We define a $Z$-\emph{valued Banach limit on}\index{Subject}{$Z$-valued Banach limit on} $\Delta$ to be a linear map LIM: $B(\Delta,\,Z)\rightarrow Z$ which satisfies
\begin{equation*}
(\forall f\in B(\Delta,\,Z)) (\liminf(f)\preccurlyeq \mathrm{LIM}\ (f)\preccurlyeq \limsup(f)
\end{equation*}
\end{enumerate}

\section*{NOTE 32}

Notes for forms [159 A] and [159 B] from Blass\index{Author}{Blass} \cite{bib:1983b}.

A \emph{similarity type}\index{Subject}{similarity type} $\tau$ is a set $J$ of operation symbols and for each $j\in J$, an arity $I_{j}$ ($I_{j}$ indexes the argument places of $j$). An \emph{algebra of type} $\tau$\index{Subject}{algebra of type $\tau$} is a set $A$ and for each $j\in J$, an operation $j_{A}$ from $A^{I_{j}}$ into $A$. The following is a theorem of ZF. (See Blass \cite{bib:1983b}.)

\begin{theorem*}
For any similarity type $\tau$ and set $V$ there is an algebra $F(V)$ of type $\tau$ (called the $\tau$ free algebra generated by $V$) with the property that $V\subseteq F(V)$ and for any $A$ of type $\tau$ and any $\eta:V\rightarrow A$ there is a unique homomorphism $\alpha:F(V)\rightarrow A$ with $\alpha(v)=v$ for all $v\in V$.
\end{theorem*}

\begin{definition*}
A $\tau$ \emph{identity} is a triple $(V,\,w,\,w^{\prime})$ where $V$ is a set (of variables) and $w$ and $w^{\prime}$ are words in the free $\tau$ algebra $F(V)$.

Given a similarity type $\tau$ and a set of identities $E$, the \emph{variety}\index{Subject}{variety} determined by $\tau$ and $E$ is the class of all algebras of type $\tau$ where all the identities in $E$ are true. $((V,\,w,\,w^{\prime})$ is true in $A$ if for every function $\eta:V\rightarrow A,\,\alpha(w)=\alpha(w^{\prime})$ where $\alpha$ is the unique homomorphism extending $\eta$.)

If $V$ is a variety and $X$ is a set then a \emph{free algebra}\index{Subject}{free algebra in a variety} in the variety on $X$ is an algebra $A$ and a function $\eta:X\rightarrow A$ such that for any $A^{\prime}$ in the variety, and $\eta^{\prime}:X\rightarrow A^{\prime}$, there is a unique homomorphism $\alpha:A\rightarrow A^{\prime}$ with $\eta^{\prime}=\alpha\eta$.

If $X=\emptyset$ then $A$ is called an \emph{initial algebra} in the variety.
\end{definition*}

\section*{NOTE 33}

The following proof that form 99 (Rado's selection lemma) and form 62 $(C(\infty, < \aleph_{0}))$ imply the Boolean prime ideal theorem (form 14) is due to A. Blass.

Assume Rado's lemma and $C(\infty,\,<\aleph_{0})$ and let $B$ be a Boolean algebra. Let $X=\{\{x,\overline{x}\}:x\in B\}$ where $\overline{x}$ denotes the complement of $x$ in B. For each finite subset $A$ of $X$, let $f_{A}$ be a choice function on $A$ chosen so that $\{f_{A}(Y):Y\in A\}$ can be extended to a proper ideal in the subalgebra of $B$ generated by $\bigcup A$. (Choosing $f_{A}$ for each finite $A\subseteq X$ uses $C(\infty,\,<\aleph_{0}).)$ By Rado's lemma there is a choice function $f$ for $X$ such that for each finite $A\subseteq X$ there is a finite $B\subseteq X$ such that $A\subseteq B$ and $f_{B}$ agrees with $f$ on $A$. Let $I=\{f(Y):Y\in X\}$. It is reasonably easy to show that $I$ is a maximal proper ideal in $B$.

\section*{NOTE 34}

A proof that Alexander's subbase lemma\index{Subject}{Alexander's subbase lemma} [14 AA] (A topological space is compact if and only if there is a subbase $S$ for the topology such that every cover by elements of $S$ has a finite subcover.) implies [14 Z] (Tychonoff's Compactness Theorem for families of finite spaces.)

Let $\{X_{i}:i\in I\}$ be a family of non-empty finite sets and let $X=\prod_{i\in I}X_{i}$. If $X$ is empty we are done otherwise let $g\in X$ and for each $j\in I$ and $a\in X_{j}$ let $B_{j,a}=\{f\in X:f(j)\in X_{j}\backslash \{a\}\}$.
$\mathcal{B} =\{B_{a,j}:j\in I$ and $a\in X_{j}\}$ is a subbase for the topology on X. We show every cover of $X$ by elements of $\mathcal{B}$ has a finite subcover. Let $\mathcal{C}$ be such a cover. If for some $j\in I,\,B_{a,j}$ and $B_{a^{\prime},j}$ are in $\mathcal{C}$ with $a\neq a^{\prime}$, then these two sets are the required finite subcover. If not then we can define an element of $X$ that is not covered by $\mathcal{C}$ as follows: $f(j)=\left\{\begin{array}{ll}
a, & \mathrm{if}\ B_{j,a}\in C\\ g(j), & \mathrm{otherwise}
\end{array}\right. . $

\section*{NOTE 35}

Form 10 $(C(\aleph_{0},\,<\aleph_{0},\,\aleph_{0}))$ is equivalent to [10
F] (Every $\omega$-tree has an infinite chain.) The proof that 10 implies [10 F] uses the fact that form 10 is equivalent to $UT(\aleph_{0},\,<\aleph_{0},\,\aleph_{0})$ which implies an $\omega$-tree is denumerable. For the argument that [10
F] implies form 10, let $A$ be a denumerable set of finite sets and $T$ be the set of all choice functions on initial\index{Subject}{initial set in a variety} segments of $A$. Define $<$ on $T$ by $f<g$ iff $f$ is a proper subset of $g$. $T$ is an $\omega$ tree and an infinite chain in $T$ induces a choice function on $A$. This proof is due to Arthur and Jean Rubin. A proof of the equivalence of form 10 and [10
F] also occurs in Pelc\index{Author}{Pelc} \cite{bib:1978}.

\section*{NOTE 36}

In this note we prove that form 284 (Every finitely solvable system of linear equations over a field has a solution in the field.) implies $\forall n\in \omega,\,n\geq 2,\,C(\infty,\,n)$ (form 61). The proof is by mathematical induction on $n$, starting with the case $n=1$ which is clear. Assume that $C(\infty,\,k)$ is true for every $k<n$ and let $X$ be a collection of pairwise disjoint $n$ element sets. By the induction hypothesis there is a choice function $h$ for the set $Y=\{z:(\exists x\in X)(z\subseteq x)\wedge 1\leq|z|<n\}$. Let $p$ be the smallest prime that divides $n$, say $n=r\cdot p$. We define a system of linear equations over the field $\mathbb{Z}_{p}$ (the integers $\mod\,p$) as follows: We have one variable $t_{a}$ for each $a\in\bigcup X$ and for each $x=\{a_{1},\,\ldots,\,a_{n}\}\in X$ we include the equation $t_{a_{1}}+\cdots+t_{a_{n}}=1$. This system of equations is clearly finitely satisfiable hence by form 284 it has a solution $S:\{t_{a}:a\in\bigcup\,X\}\rightarrow \mathbb{Z}_{p}$. We now define a choice function $f$ on $X$. Let $x=\{a_{1},\,\ldots,\,a_{n}\}$ be an element of $X$. The solution $S$ determines a partition of $x$ into the sets $z_{m}=\{a\in x:S(t_{a})=m\},\,0\leq m\leq p-1$, some of which may be empty. However, there is no $m,\,0\leq m\leq p-1$, for which $x=z_{m}$ because the existence of such an $m$ would mean that $S(t_{a_{1}})=\cdots=S(t_{a_{n}})=m$. Hence $S(t_{a_{1}})+\cdots+S(t_{a_{n}})=m\cdot n=m\cdot(r\cdot p)=0$, which is a contradiction since $S$ is a solution to the equation $t_{a_{1}}+\cdots+t_{a_{n}}=1$.

Let $m_{0}$ be the least element of $\mathbb{Z}_{p}$ for which $z_{m_{0}}\neq\emptyset$. Since $z_{m_{0}}\neq x,\, |z_{m_{0}}|<n$ and we may therefore define $f(x)=h(z_{m_{0}})$.

\section*{NOTE 37}

Definitions for form [14 AB] (Tychonoff theorem for sober spaces).

\begin{definition*}
A topological space is \emph{sober}\index{Subject}{sober topological space} if every closed set $C\neq\emptyset$ that is not the union of two proper closed (not necessarily disjoint) subsets (such a $C$ is called \emph{irreducible}) is the closure of a unique singleton.
\end{definition*}

In Hausdorff spaces, irreducible closed sets are singletons and therefore

\begin{theorem*}
Hausdorff implies sober implies $T_{0}$.
\end{theorem*}

See Blass\index{Author}{Blass} \cite{bib:1986} and Johnstone\index{Author}{Johnstone} \cite{bib:1984}.

\section*{NOTE 38}

Definitions from Manka\index{Author}{Manka} \cite{bib:1988a} and \cite{bib:1988b}.

\begin{definition*}
Suppose that $(X,\,\leq)$ is a partial ordering
\begin{enumerate}
\item[1.] A set $\mathcal{A}\subseteq \mathcal{P}(X)$ is \emph{supremal in}\index{Subject}{supremal set} $(X,\,\leq)$ if
\begin{enumerate}
\item[(a)] $\forall S\in \mathcal{A}$ and for every upper bound $y$ of $S,\,S\cup\{y\}\in \mathcal{A}$ and
\item[(b)] $\forall S\in \mathcal{A}$ every initial segment of $S$ is in $\mathcal{L}$.
\end{enumerate}
\item[2.] If $\mathcal{A}$ is supremal in $(X,\,\leq)$ and $\sigma:\mathcal{A}\rightarrow X$ then $\sigma$ is a \emph{sup function}\index{Subject}{sup function} iff
\begin{enumerate}
\item[(a)] $\forall S\in \mathcal{A},\,\sigma(S)$ is an upper bound for $S$ and
\item[(b)] $\forall S\in \mathcal{A},\,\sigma(S)$ is the greatest element of $S$ if there is one.
\end{enumerate}
\item[3.] $(X,\,\leq)$ is \emph{directed} if each two element subset of $X$ has an upper bound.
\end{enumerate}
\end{definition*}

\begin{definition*}
If ($X, \rho$) is a metric space and $\phi:X\rightarrow \mathbb{R}$, then $\phi$ is \emph{upper semi-continuous}\index{Subject}{upper semi-continuous} iff for every convergent sequence $\{x_{n}\}$ in $(X,\,\rho)$ and for all $r\in \mathbb{R}$, the inequality $\phi(x_{n})\geq r$ for all $n\in\omega$ implies $\phi(\lim x_{n})\geq r$.
\end{definition*}

\begin{definition*}
If $(X,\,\leq)$ is a directed, partially ordered set and $\tau$ is a topology on $X$, then $(X,\,\leq)$ is \emph{upper-semi continuous with respect to} $\tau$ if for every $x_{0}$ in $X,\,\{y\in X:x_{0}\leq y\}$ is closed.
\end{definition*}

\section*{NOTE 39}

In this note the results of Harper/Rubin\index{Author}{Harper/Rubin} \cite{bib:1976} are summarized.

\begin{definition*}
Suppose that $R$ is a relation on $X$.
\begin{enumerate}
\item[1.] $R$ is \emph{connected} on $X$ if $(\forall u,\,v\in X)(u\neq v\rightarrow uRv$ or $vRu)$.
\item[2.] An element $u\in X$ is $R$-\emph{maximal} in $X$ if $(\forall v\in X)(uRv\rightarrow vRu)$. $R$-\emph{minimal} is defined similarly.
\item[3.] If $Y\subseteq X,\,u\in X$ is an $(R)$ \emph{upper bound} for $Y$ if $(\forall v\in Y)(v\neq u\rightarrow vRu)$. \emph{Lower bound} is defined similarly.
\item[4.] $X$ is \emph{directed (upwards)} by $R$ iff $R$ partially orders $X$ and every finite subset of $X$ has an upper bound.
\item[5.] If $y\in X$, the \emph{initial segment generated by} $y$ is $y^{\ast}=\{u\in X:uRy\}$.
\item[6.] $X$ is \emph{ramified}\index{Subject}{ramified set} by $R$ if $R$ partially orders $X$ and for every $y\in X,\,y^{\ast}$ is linearly ordered by $R$.
\item[7.] $(X,\,R)$ is a \emph{forest}\index{Subject}{forest} if $R$ is a partial ordering and $y^{\ast}$ is well ordered by $R$ for every $y\in X$.
\item[8.] $(X,\,R)$ is a \emph{tree} if $(X,\,R)$ is a forest and is directed downward by $R$.
\item[9.] $Q\subseteq X$ is \emph{quasi-cofinal}\index{Subject}{quasi-cofinal} in $X$ if $Q$ has no strict upper bound in $S$. (``Strict'' $=$ ``not in $Q$''.)
\item[10.] $Q\subseteq X$ is \emph{cofinal} in $X$ if $Q$ is linearly ordered by $R$ and $(\forall y\in X)(\exists z\in Q)(y\,R\,z)$.
\item[11.] $A\subseteq X$ is an \emph{antichain}\index{Subject}{antichain} if $\forall x,\,y\in A$, neither $xRy$ nor $yRx$ holds.
\item[] In Harper/Rubin \cite{bib:1976}, the following weak forms of $AC$ are studied:
\item[]$Z(Q,\,U)$: Every non-empty $Q$-ordered set in which every $U$-ordered subset has an upper bound, has a maximal element.
\item[]$C(Q,\,U)$: Every $Q$-ordered set contains a quasi-cofinal $U$-ordered subset.
\item[]$H(Q,\,U)$: Every $Q$-ordered set contains a $\subseteq$-maximal, $U$-ordered subset.
\item[]$H^{\prime}(Q,\,U)$: Every $Q$-ordered set has the property that each $U$-ordered subset can be extended to a $\subseteq$-maximal $U$-ordered subset.
\end{enumerate}
where $Q$ and $U$ range over the following properties: $A$: arbitrary (binary) relation, $TR$: transitive, $AS$: antisymmetric, $C$: connected, $P$: partially ordered, $W$: well ordered, $L$: linearly ordered, $D$: directed, $R$: ramified, $F$: forest, $T$: tree\index{Subject}{tree}, $AS$ \& $C$ and, $TR$ \& $C$.

In particular:

\noindent Form 51 is equivalent to $Z(L,\,U)$ where $U=W,\,F$ or $T$ and to $Z(R,\,F)$ and to $Z(R,\,T)$.

\noindent Form 143 is equivalent to $H(C,\,TR)$ and $H(C,\,TR$\&$C$),

\noindent Form 202 ([202 A]) is equivalent to $H(TR\& C,\,U)$ where $U=AS,\,AS$\&$C,\,P,\,L,\,D$, or $R$.

\noindent Form 255 is equivalent to $Z(P,\,R)$

\noindent Form 256 is equivalent to $Z(Q,\,U)$ where $Q\in\{P,\,D\}$ and $U\in\{F,\,T\}$.

\noindent Form 257 is equivalent to $Z(TR,\,AS)$ and $Z(TR,\,P)$.

\noindent Form 258 is equivalent to $Z(D,\,U)$ where $U\in\{C,\,AS\& C,\,TR\& C,\,L\}$.

\noindent Form 259 is equivalent to $Z(TR\& C,\,U)$ where $U\in\{W,\,F,\,T\}$.

\noindent Form 260 is equivalent to $Z(TR\& C,\,U)$ where $U\in\{AS,\,AS\& C,\,P,\,L,\,D,\,R\}$.

\noindent Form 261 is equivalent to $Z(TR,\,F)$ and $Z(TR,\,T)$.

\noindent Form 263 is equivalent to $H(AS\& C,\,U)$ where $U\in\{TR,\,TR\& C,\,P,\,L,\,D,\,R\}$.

\noindent Form 264 is equivalent to $H(C,\,U)$ where $U\in\{P,\,L,\,D,\,R\}$.

\noindent Form 266 is equivalent to $H(AS,\,TR)$ and $H(AS,\,P)$ and $\{AS,\,AS\&
C,\,P,\,L,\,D,\,R\}$.

\noindent Form 261 is equivalent to $Z(TR,\,F)$ and $Z(TR,\,T)$.

\noindent Form 263 is equivalent to $H(AS\& C,\,U)$ where $U\in\{TR,\,TR\& C,\,P,\,L,\,D,\,R\}$.

\noindent Form 264 is equivalent to $H(C,\,U)$ where $U\in\{P,\,L,\,D,\,R\}$.

\noindent Form 202 ([202 A]) is equivalent to $H(TR\& C,U)$ where $U=AS,\,AS\&
C,\,P,\,L,\,D$, or $R$.

\noindent Form 266 is equivalent to $H(AS,\,TR)$ and $H(AS,\,P)$.
\end{definition*}

\section*{NOTE 40}

Notes on [94 A] and [94 B]. Brunner\index{Author}{Brunner} \cite{bib:1982d} proves that [94 B] implies form 13 (lemma 6, p. 164) by showing that the negation of form 13 implies the negation of [94 A] which implies the negation of [94 B]. Since [94 A] clearly implies [94 B], it only remains to prove that [94 A] implies 94 (to complete the argument for the equivalence of 94, [94 A] and [94 B]). (Form
94: $C(\aleph_{0},\,\infty,\, \mathbb{R})$; [94 A]: The discrete topology on $\omega$ is Lindel\"{o}f; and [94 B]: There is a Lindel\"{o}f space with a closed, discrete infinite subset.)

Assume [94 A] and let $\{A_{i}:i\in \omega\}$ be a denumerable family where for all $i\in \omega,\, A_{i}\neq\emptyset$ and $A_{i}\subseteq \mathcal{P}(\omega)$. It suffices to find a choice function for the family $\{A_{i}:i\in \omega\}$ (identifying reals with sets of natural numbers).

Now using the family $\{A_{i}:i\in \omega\}$ we define a second family of sets of sets of natural numbers $\{B_{i}:i\in \omega\}$ so that
\begin{equation*}\tag{$\ast$}
\forall i\in \omega,\, 2^{i}\in\bigcup B_{i}\ \text{and for}\ j\in \omega,\,j\neq i\ \text{implies}\ 2^{i}\not\in\bigcup B_{j}.
\end{equation*}
More specifically, for each $i\in\omega$, let $m(i)$ be the least element of $\bigcup A_{i}$. Define $f_{i}:\{k\in \omega:k\geq m(i)\}\rightarrow \omega$ by $f_{i}(k)=[2(k-m(i))+1]2^{i}$ and let $B_{i}=\{f_{i}^{\prime\prime}x:x\in A_{i}\}$. It is clear that $\{B_{i}:i\in \omega\}$ satisfies $(\ast)$ and that it suffices to find a choice function for $\{B_{i}:i\in \omega\}$. Further $\bigcup_{i\in\omega}B_{i}\cup\big\{\{k\in \omega:k\not\in\bigcup_{i\in\omega}(\bigcup\,B_{i})\}\big\}$ is an open cover for $\omega$ (using the discrete topology). Since the topology is Lindel\"{o}f there is a countable subcover $\{x_{j}:j\in \omega\}$. Defining, for each $i\in \omega,\,F(B_{i})=x_{j}$ where $j$ is the least element of $\omega$ such that $2^{i}\in x_{j}$ provides a choice function for $\{B_{i}:i\in \omega\}$.

\section*{NOTE 41}

The following proof that form 233 is true in $\mathcal{N}1$ is due to David Pincus. Form 233 is ``Algebraic closures of fields, if they exist, are unique.''

\begin{lemma*}
If $K$ is an algebraically closed field, if $\pi$ is a non-trivial automorphism of $K$ satisfying $\pi^{2}=1_{K}$ (the identity on $K$), and if $i\in K$ is a square root of $-1$, then $\pi(i)=-i\neq i$.
\end{lemma*}

(H. L\"{a}uchli has given a proof of this lemma with the additional hypothesis that the characteristic of $K$ is not equal to 2. See Jech\index{Author}{Jech} \cite{bib:1973b}, p. 147.)

\begin{proof}
Let $L=\{x\in K:\pi(x)=x\}$. The order of the group of automorphisms generated by $\pi$ is 2. Therefore by Artin's theorem (theorem 1.8 on page 264 of Lang \cite{bib:1993}) the index of $L$ in $K$ is 2. By corollary 9.3, page 299 of Lang \cite{bib:1993} (If $L$ is a field and $K$ is the algebraic closure of $L$ and is of finite degree over $L$, then $K=L(i)$ where $i^{2}=-1$ and L has characteristic $0$.) $K=L(i)$, where $i^{2}=-1$ and $L$ has characteristic $0$, and hence $K$ has characteristic $0$.

Choose $x\in K$\index{Author}{Lang} such that $\pi(x)\neq x$. Let $y=\pi(x)-x$ so that $y\neq 0$ and $\pi(y)=-y$. Let $z$ be a square root of $y$. Then $\pi(z)$ is a square root of $\pi(y)=-y$ hence $\pi(z)=iz$ for some square root $i$ of --1. Since $y\neq 0,\,z\neq 0$. We can therefore compute
\begin{equation*}
\pi(i)=\pi\Big(\frac{\pi(z)}{z}\Big)=\frac{z}{\pi(z)}=\frac{1}{i}=-i\neq i.
\end{equation*}
The last inequality is a consequence of the fact that the characteristic of $K$ is not 2. Since $\pi$ has order 2, and $\pi(-i)=i$. This completes the proof of the lemma.
\end{proof}

Now let $(K,\,+,\,\cdot,\,0,1)$ be a field in $\mathcal{N}1$ with finite support $E\subseteq A$ where $A$ is the set of atoms. Assume that $K$ is algebraically closed. We will show that every element of $K$ has support $E$. This implies that all algebraically closed fields are well orderable in $\mathcal{N}1$ and therefore the standard proof of the uniqueness of algebraic closures (using AC) is valid in $\mathcal{N}1$.

The proof is by contradiction. Assume that $x\in K$ does not have support $E$. We will arrive at a contradiction in two steps. The first is to show that for some permutation $\pi$ of $A$ which fixes $E$ pointwise, $\pi(x)\neq x$ and $\pi^{2}$ is the identity. The permutation $\pi$ induces an automorphism of $(K,\,+,\cdot,\,0,1)$ and we can therefore apply the lemma to conclude that $\pi(i)=-i\neq i$ for some square root $i$ of -1 in $K$. The second step is to show that for every permutation $\pi$ of $A$ that fixes $E$ pointwise, $\pi(i)=i$ for every square root $i$ of -1 in $K$.

For step one, suppose that $x$ has support $E\cup F$ where $E\cap F=\emptyset$. Let $F^{\prime}$ be a subset of $A$ for which $|F^{\prime}|=|F|$ and $F^{\prime}\cap(E\cup
F)=\emptyset$. Let $\pi$ be a permutation of $A$ which interchanges the elements of $F$ and $F^{\prime}$, that is, $\pi$ is a product of cycles of length 2, $\pi=(a_{1},\,b_{1})(a_{2},\, b_{2})\cdots(a_{n},\,b_{n})$ where $F=\{a_{1},\,\ldots,\,a_{n}\}$ and $F^{\prime}=\{b_{1},\,\ldots,\,b_{n}\}$. Clearly $\pi^{2}$ is the identity and $\pi$ fixes $E$ pointwise. The proof that $\pi(x)\neq x$ is a standard permutation model argument and depends on the fact that $x$ does not have support $E$.

For step two, we note that if $i$ is a square root of --1 in K, then $i$ is a solution to the equation $x^{2}+1=0$ all of whose coefficients are fixed by any permutation that fixes $E$ pointwise. Therefore if $\eta$ is a permutation of $A$ that fixes $E$ pointwise, $\eta(i)$ is also a solution to $x^{2}+1=0$. Suppose that $i$ does not have support $E$. Let $E\cup F$ be a support of $i$ where $E\cap F=\emptyset$. By a standard permutation model argument, we can show that if $\pi$ and $\pi^{\prime}$ are any two permutations of $A$ that fix $E$ pointwise and $\pi(F) \cap \pi^{\prime}(F)=\emptyset$ then $\pi(x)\neq\pi^{\prime}(x)$. It follows that the equation $x^{2}+1=0$ has infinitely many solutions in $K$ which is a contradiction and the proof is complete.

\section*{NOTE 42}

In Brunner\index{Author}{Brunner} \cite{bib:1982a} is a study, in $\mathrm{ZF}^{0}$ without $AC$, of conditions under which a topological space is well orderable. These conditions are summarized below.

\begin{theorem*}
(2.4 from Brunner \emph{\cite{bib:1982a}}) Assume $X$ is a $T_{2}$ topological space
\begin{enumerate}
\item[\emph{(1)}] Assuming form 163, $X$ is well orderable if (a) $X$ is hereditarily separable or (b) $X$ is hereditarily Lindel\"{o}f and linearly orderable.
\item[(2)] Assuming form 164, $X$ is well orderable if (a) $X$ is hereditarily Lindel\"{o}f or (b) $X$ is first countable\index{Subject}{first countable} and Lindel\"{o}f or (c) $X$ is first countable and hereditarily Souslin\index{Subject}{Souslin space} or (d) $X$ is second countable\index{Subject}{second countable} or (e) $X$ is separable.

\item[(3)] Assuming form 133, $X$ is well orderable if (a) $X$ is hereditarily Souslin or (b) $X$ is linearly orderable. (Note: $X$ is linearly orderable means that the topology on $X$ is the order topology of some linear order, $X$ is well orderable means that the set $X$ is well orderable but the topology need not be the order topology of the well ordering.)
\end{enumerate}
\end{theorem*}

\begin{definition*}
\begin{enumerate}
\item[(1)] $X$ is \emph{first countable} if every point has a countable neighborhood base.
\item[(2)] $X$ is \emph{second countable} is $X$ has a countable basis.
\item[(3)] $X$ is \emph{Souslin} if every family of disjoint open sets is countable.
\item[(4)] $X$ is \emph{metacompact}\index{Subject}{metacompact} if every open cover has a point finite refinement. (See form 135 for a definition of point finite.)
\end{enumerate}
\end{definition*}

\begin{corollary*}
(2.5 in Brunner\index{Author}{Brunner} \emph{\cite{bib:1982a})} Assuming form 164, the following are equivalent for a metric space $X$: (1) $X$ is separable, (2) $X$ is second countable, (3) $X$ is Souslin and (4) $X$ is Lindel\"{o}f.
\end{corollary*}

\begin{theorem*}
(2.6) Under the assumption of form 133 and form 17 (Ramsey' s Theorem) a regular, Lindel\"{o}f, Souslin and hereditarily metacompact space is well orderable.
\end{theorem*}

\begin{theorem*}
(2.7) Assuming form 133, a $T_{2}$ topological space $X$ is well orderable if $X^{2}$ is hereditarily normal and $X$ is Lindel\"{o}f.
\end{theorem*}

\begin{theorem*}
(2.8) Assuming form 164, a $T_{2}$ space is well orderable if $X^{2}$ is hereditarily locally Lindel\"{o}f and $X$ is hereditarily Souslin.
\end{theorem*}

Finally some counter examples: In the model $\mathcal{M}1$ (The basic Cohen model) $X=[0,1]$ (the closed interval in $\mathbb{R}$) with the order topology is compact, metric, separable, second countable, hereditarily Souslin and linearly ordered, but $X$ is not well orderable.

In the model $\mathcal{N}3$ (The ordered Mostowski model of $\mathrm{ZF}^{0}$): Let $X=[a,\,b]$ be an interval in the set of atoms with the order topology. $X$ is compact, hereditarily normal (that is $T_{5}$) and $T_{2}$. Further, each finite power of $X$ is compact, $T_{4}$ and hereditarily Souslin, but $X$ is not well orderable.

In the model $\mathcal{N}1$ (The basic Fraenkel model with a denumerable set of atoms, using the group of all permutations of the atoms and finite supports): Let $X= [A]^{<\omega}$ (all finite subsets of the set of atoms) with a basis for the topology all set of the form $\{x:E\subseteq x$ and $x\cap F=\emptyset\}$ where $E$ and $F$ range over the finite subsets of U. The space $X$ is Souslin, perfectly regular and not well orderable.

\section*{NOTE 43}

These are definitions from Brunner \cite{bib:1982b} and results similar to the equivalence of [8 C] and [8 D] to form 8 and [10 H] to form 10. We also include some results from Brunner \cite{bib:1987b}.

\begin{definition*}\
\begin{enumerate}
\item[1.] A topological space is $\sigma$-\emph{compact}\index{Subject}{$\sigma$-compact} if it can be covered by a countable set of compact sets.
\item[2.] A space is \emph{weakly Lindel\"{o}f}\index{Subject}{weakly Lindel\"{o}f} if every open cover has a countable refinement. If $C$ is an open cover of a space $X$ a \emph{refinement} of $C$ is an open cover $C^{\prime}$ of $X$ such that for all $y\in C^{\prime}$, there is a $z$ in $C$ such that $y\subseteq z$.
\item[3.] A space $X$ is \emph{paracompact}\index{Subject}{paracompact} if $X$ is Hausdorff and every open covering of $X$ has an open, locally finite refinement. (An open covering $C$ of $X$ is \emph{locally finite} if for all $x\in X$, there is an open set containing $x$ that intersects only finitely many elements of $C$. In general, a family of subsets $S$ of $X$ is called locally finite if each point in $X$ has a neighborhood meeting a finite number of sets in $S.)$
\item[4.] An \emph{ultra metric group}\index{Subject}{ultra metric group} is a topological group with a metrizable topology using a metric $d$ that satisfies for all $x,\,y$, and $z$ in the group,
\begin{equation*}
d(x,\,y)\leq\max(d(x,\,z),\,d(y,\,z)).
\end{equation*}
The results from Brunner\index{Author}{Brunner} \cite{bib:1982b} are:
\item[a.] It is not possible to prove, without $AC$, that every weakly Lindel\"{o}f space is Lindel\"{o}f since in the basic Cohen model $\mathcal{M}1,\,\mathbb{R}$ is weakly Lindel\"{o}f but not Lindel\"{o}f.
\item[b.] The assertion that every $\sigma$-compact, locally compact space with countable neighborhood bases is Lindel\"{o}f implies form 32 $(C(\aleph_{0},\,\aleph_{0}))$.
\item[c.] In the basic Cohen model, $\mathcal{M}1,\,C(\aleph_{0},\,\aleph_{0})$ holds and there is a $\sigma$-compact locally compact space with countable neighborhood bases which is not Lindel\"{o}f.
\item[d.] $C(\aleph_{0},\,<\aleph_{0})$ is implied by ``Every $\sigma$-compact, ultra metric group without isolated points is Lindel\"{o}f.''
\item[e.] The assertion ``Every $\sigma$-compact, regular space without isolated points is Lindel\"{o}f'' implies form 9 (Dedekind finite = Finite).

\noindent The results from Brunner \cite{bib:1987b} are:
\item[a.] In $\mathcal{N}3,\,L$ is a paracompact and weakly Lindel\"{o}f topological space which is not Lindel\"{o}f. (In $\mathcal{N}3,\,L=(A\times \mathbb{Z})\cup \mathbb{Z}$, where the ordering on $L$ is defined as follows: for $a,\, b\in A$ and $m,\,n\in \mathbb{Z},\,(a,\,n)<(b,\,n)$ if $a<b$ and $(a,\,n)<m<(b,\,m)$ if $n<m$. The topology on $L$ is the order topology.)
\item[b.] In $\mathcal{M}1$, every $T_{1}$ Lindel\"{o}f space is compact.
\item[c.] In $\mathcal{M}17$, every countably compact space is compact.
\end{enumerate}
\end{definition*}

\section*{NOTE 44}

Sageev\index{Author}{Sageev} \cite{bib:1981} shows that the implication form 62 $(C(\infty,\,<\aleph_{0}))\
+$ form 57 (Any two Dedekind cardinals are comparable.) implies form 9 (Dedekind finite $=$ finite) is not provable in ZF. Ellentuck\index{Author}{Ellentuck} \cite{bib:1974} shows that $\mathrm{ZF}\vdash(62+57\rightarrow\langle \omega, \cdot, +\rangle$ is an elementary submodel of $\langle\Delta,\,\cdot, +\rangle)$ where $\Delta$ is the set of Dedekind finite cardinals $(\omega\subseteq\Delta)$ and $\cdot$ and $+$ are the usual cardinal operations.

\section*{NOTE 45}

For research on set theory without $AC$ the following references may be consulted:

A. Results on Dedekind finite, infinite fields in set theory without $AC$ see Hickman\index{Author}{Hickman} \cite{bib:1978a}, \cite{bib:1980a}, and \cite{bib:1982}.

B. The Dedekind cardinals without $AC$ see Ellentuck \cite{bib:1970}, \cite{bib:1973}, \cite{bib:1976} and Monro\index{Author}{Monro} \cite{bib:1975}.

C. Order types of linearly ordered Dedekind finite sets without $AC$ see Hickman\index{Author}{Hickman} \cite{bib:1979a}.

D. Ordered sets satisfying ``there is no infinite descending sequence'', see Hickman \cite{bib:1975}.

E. Possible properties of rings see Hodges\index{Author}{Hodges} \cite{bib:1974} where the negation of form 278 is shown to be consistent with ZF.

F. Estimates of the sizes of $\aleph(x)$ and $\aleph^{\ast}(x)$ in set theory without AC. $\aleph(x)$ is the smallest ordinal $\alpha$ such that $\neg(|\alpha|\leq|x|)$ and $\aleph^{\ast}(x)$ is the smallest ordinal $\alpha$ such that $\neg(|\alpha|\leq\ast|x|)$ where $\leq\ast$ is the surjective cardinal ordering. See Kruse\index{Author}{Kruse} \cite{bib:1974}.

G. Possible cardinality of the set of Dedekind finite cardinals in a Fraenkel-Mostowski model see A. Rubin/J. Rubin\index{Author}{A. Rubin/J. Rubin} \cite{bib:1974}.

H. Amorphous sets $X$ for which $\mathcal{P}\mathcal{P}\mathcal{P}(X)$ is large see Monro\index{Author}{Monro} \cite{bib:1973c}. ($X$ is amorphous if $X$ is infinite and $\forall Y\subseteq X,\,Y$ is finite or $X-Y$ is finite.)

I. Topology without the axiom of choice see Chandler\index{Author}{Chandler} \cite{bib:1972}.

J. Possible order types of convex sets of cardinal numbers see Truss\index{Author}{Truss} \cite{bib:1973e}.

K. Large cardinals in set theory without the axiom of choice see Spector\index{Author}{Spector} \cite{bib:1991} and Apter\index{Author}{Apter} \cite{bib:1992}.

L. Logic without $AC$ see Malitz\index{Author}{Malitz} \cite{bib:1992}.

M. Posets and lattices without $AC$ see Hickman \cite{bib:1977a}.

\section*{NOTE 46}

With regard to form 325 (Ramsey's Theorem\index{Subject}{Ramsey's theorem} II), all the arguments given in Blass\index{Author}{Blass} \cite{bib:1977a} can be modified slightly to show that all the properties of form 17 (Ramsey's Theorem I) proved by Blass also hold of 325. In particular
\begin{enumerate}
\item[1.] Form 9 (Dedekind finite $=$ finite) implies form 325.
\item[2.] Form 325 is true in $\mathcal{N}1$.
\item[3.] Form 325 is injectively boundable.
\end{enumerate}

\section*{NOTE 47}

Definitions from Shannon\index{Author}{Shannon} \cite{bib:1990}.

\begin{definition*}
Assume that $(P,\,\leq)$ is a quasi-order\index{Subject}{quasi-order} (reflexive and transitive).
\begin{enumerate}
\item[1.] $x$ and $y$ in $P$ are \emph{compatible}\index{Subject}{compatible} if $\exists z\in P$ such that $z\leq x$ and $z\leq \mathrm{y}$.
\item[2.] \emph{Incompatible} means not compatible.
\item[3.] $I\subseteq P$ is an \emph{antichain}\index{Subject}{antichain} if each pair of elements of $I$ are incompatible.
\item[4.] $D\subseteq P$ is \emph{dense}\index{Subject}{dense subset of a quasi-order} if $\forall x\in P,\,\exists y\in D$ such that $y\leq x$.
\item[5.] $F\subseteq P$ is a \emph{filter}\index{Subject}{filter in a quasi-order} if
\begin{enumerate}
\item[a)] $\forall x,\,y\in F\,\exists z\in F$ such that $z\leq x$ and $z\leq y$ and
\item[b)] $\forall x\in F\, \forall y\in P,\,x\leq y$ implies $y\in F$.
\end{enumerate}
\item[6.] If $\mathcal{D}\subseteq \mathcal{P}(P)$, a filter $F$ is $\mathcal{D}$ \emph{generic}\index{Subject}{generic filter} if $F\cap D\neq\emptyset$ for all $D$ in $\mathcal{D}$. $F$ is $P$ \emph{generic} if $F\cap D\neq\emptyset$ for all dense subsets $D$ of $P$.
\item[7.] A chain $C$ of $(P,\,\leq)$ is a \emph{covering chain}\index{Subject}{covering chain} iff for all $x\in C$, if $x$ is not maximal in $C$, then $\exists y\in C$ such that $y$ covers $x$ in $P$ (i.e., $y>x$ and there is no $z\in P$ such that $y>z>x$).
\item[8.] If $Q\subseteq P$ and $x\in P$, the \emph{depth} of $x$ (relative to $Q$) is the order type of the smallest well ordered (under the quasi-order) covering chain $C$ which begins at $x$ and ends at an element of $Q$. (The depth of $x$ is only defined if such a chain exists.)
\item[9.] Two elements $x$ and $y$ of $P$ are $Q$ \emph{strongly incompatible}\index{Subject}{strongly incompatible} if $x$ and $y$ are incompatible and the depth of $x$ is different from the depth of $y$.
\item[10.] $I\subseteq P$ is a $Q$ \emph{strong antichain}\index{Subject}{strong antichain} if any two elements are $Q$ strongly incompatible.
\end{enumerate}

Shannon\index{Author}{Shannon} \cite{bib:1990} proves that the following are equivalent for each ordinal $\alpha$:

(i) $C(\aleph_{\alpha},\, <\aleph_{0})$

(ii) If $(P,\,\leq)$ is a quasi-order such that $P\neq\emptyset$ and $P$ is a union of $\aleph_{\alpha}$ finite sets and every antichain is finite, then there is a well orderable, $P$ generic filter.

(iii) For some ordinal $\beta\geq\alpha$, for every quasi-order $(P,\,\leq)$ such that $P\neq\emptyset$ and such that $P$ is the union of $\aleph_{\alpha}$ finite sets and such that every antichain is finite then for every family $\mathcal{D}$ of $\aleph_{\beta}$ dense sets there is a well orderable $\mathcal{D}$ generic filter.
\end{definition*}

\section*{NOTE 48}

Jech\index{Author}{Jech} \cite{bib:1982} shows that [34 A] implies form 34. The proof that 34 implies [34 A] uses the fact that if $\aleph_{1}$ is regular then no set of rank $\omega_{1}$ can be countable.

\section*{NOTE 49}

In von Rimscha\index{Author}{Rimscha} \cite{bib:1982} three transitivity conditions are considered:

$Tr.\ \forall x\exists u\exists f$ such that $u$ is transitive and $f$ is a function from $f$ one to one and onto $u$. This is form 175.

$Tr^{\prime}.\ \forall x\exists u\exists f$ such that $u$ is transitive and $u\subseteq TC(x)$ and $f$ is a one to one function from $x$ onto $u$.

$Tr^{\prime\prime}.\ \forall x\exists u\exists f$ such that $u$ is transitive and $f$ is a one to one function from $x$ onto $u$ and $\forall s\in x,\,f(s)\in TC(s)$.

It is shown that in ZF, AC implies $Tr^{\prime\prime}$ and that $Tr^{\prime\prime}$ implies $Tr^{\prime}$ implies $Tr$. All proofs can be done in $\mathrm{ZF}^{0}$ with minor modifications. It is also shown that $\mathrm{ZF} \vdash Tr\rightarrow$ form 9 (Dedekind finite $=$ finite). This could be reformulated as $\mathrm{ZF}^{0}\vdash Tr\rightarrow$ all pure Dedekind finite sets are finite. We can therefore conclude that $Tr$ (form 175) implies Dedekind finite sets of reals are finite (form 13).

It is however possible to prove in $\mathrm{ZF}^{0}$ that $Tr^{\prime\prime}$ (form 174) implies Dedekind finite $=$ finite (form 9). We proceed as follows: Assume $Tr^{\prime\prime}$. We first show that every infinite subset of the atoms $A$ is Dedekind infinite. Assume $B\subseteq A$ and that $B$ is infinite. We will use $Tr^{\prime\prime}$ to get a choice function for $\mathcal{P}(B)$ from which it follows that $B$ has a denumerable subset. Let $\kappa$ be the least well ordered cardinal which is not $\leq|\mathcal{P}(B)|$. For each $a\in B$ and ordinal $\alpha$, define $a_{\alpha}$ by induction: $a_{0}=a,\,a_{\lambda}=\bigcup_{\alpha<\lambda}a_{\alpha}$ for limit ordinals $\lambda$ and $a_{\alpha+1}=a_{\alpha}\cup\{a_{\alpha}\}$.
$(a_{\alpha}$ is the ordinal $\alpha$ with every occurrence of $\emptyset$ replaced by $a$.) Then for each $\alpha,\, |a_{\alpha}|=|\alpha|$. Let $B^{\prime}=\{a_{\kappa}:a\in B\}$ and let $x=\mathcal{P}(B^{\prime})$. By $Tr^{\prime\prime}$ there is a transitive set $u$ and a function $f:x\rightarrow u$ which is one to one and onto and such that for each $s\in x,\,f(s)\in TC(s)$. Suppose that $C\subseteq B$ and let $s=\{a_{\kappa}:a\in C\}\in x$. Then
\begin{equation*}
TC(s)=\{s\}\cup s\cup \bigg(\bigcup_{a\in C}a_{\kappa}\bigg)
\end{equation*}
since $|TC(s)|\geq\kappa,\,s\not\in u$. $(|u|=|\mathcal{P}(B)|$ which is not $\geq\kappa.)$. So $f(s)\in\bigcup_{a\in C}a_{\kappa}$. Since $a\neq b$ implies $a_{\kappa}\cap b_{\kappa}=\emptyset,\, f(s)$ determines a unique element of $C$.

Now let $x$ be any infinite set. By $Tr^{\prime\prime}$ there is a transitive $u$ and a one to one function $f$ from $x$ onto $u$. If $u$ contains only finitely many atoms, then the von Rimscha\index{Author}{Rimscha} proof can be modified to show that $x$ has a denumerable subset. If $u$ contains an infinite set of atoms $B$, then by the first argument $B$ and therefore $x$ contains a denumerable subset.

\textsc{Problem}
(from von Rimscha \cite{bib:1982}): Are any of the implications $Tr\rightarrow Tr^{\prime}\rightarrow Tr^{\prime\prime}\rightarrow AC$ provable in ZF?

\section*{NOTE 50}

Definitions from Andreka/Nemeti\index{Author}{Andreka/Nemeti} \cite{bib:1981} and H\"{o}ft/Howard\index{Author}{H\"{o}ft/Howard} \cite{bib:1981}.

An \emph{algebra} is a pair $(A,\,F)$ where $A$ is a set and $F$ is a function from an index set $I$ such that for each $i\in I$, there is a $k_{i}\in \omega$ such that $F(i)$ is a function from $A^{k_{i}}$ into $A$. The function $\Delta$ with domain $I$ defined by $\Delta(i)=k_{i}$ is called the \emph{type of the algebra}\index{Subject}{type of an algebra} $(A,\,F)$. Two algebras of the same type are \emph{similar}.

\begin{definition*}
Let $K$ be any class of similar algebras
\begin{enumerate}
\item[1.] $S(K)=$ the class of subalgebras of algebras in $K$
\item[2.] $H(K)=$ homomorphic images of algebras in $K$
\item[3.] $P(K)=$ direct products of non-empty families of algebras in $K$
\item[4.] $HSP\ K=H(S(P(K)))$.
\item[5.] $Mod\,Eq\,K=$ the class of all algebras which are models of the set of equations which hold in every member of $K$.
\end{enumerate}

The fact that [0 U] is equivalent to form $0$ answers problem 31 in Gr\"{a}tzer\index{Author}{Gr\"{a}tzer}
\cite{bib:1979}.

A \emph{multi-algebra}\index{Subject}{multi-algebra} is a pair $(A,\,F)$ where $A$ is a set and $F$ is a function from an index set $I$ such that for each $i\in I$, there is a $k_{i}\in \omega$ such that $F(i)$ is a function from $A^{k_{i}}$ into $\mathcal{P}(A)-\{\emptyset\}$. We will denote $F(i)$ by $f_{i}$. The function $\Delta$ with domain $I$ defined by $\Delta(i)=k_{i}$ is called the \emph{type of the algebra}. Note that this doesn't require that constants, i. e., 0-ary operations, have to be singleton sets.

If $(A,\,F)$ and $(B,\,G)$ are algebras or multi-algebras of the same type then a function $\phi:A\rightarrow B$ is a \emph{homomorphism} if for every $i\in I$ and $\mathbf{a}= (a_{1},\,\ldots a_{k_{i}})\in A^{k_{i}},\,\phi(f_{i}(a_{1},\,\ldots a_{k_{i}}))=g_{i}(\phi(a_{1}),\,\ldots\phi(a_{k_{i}}))$. The function $\phi$ is an \emph{isomorphism} between multi-algebras if $\phi$ and $\phi^{-1}$ are bijective homomorphisms.
\end{definition*}

\begin{definition*}
Suppose that $(A,\,F)$ is a multi-algebra of type $\Delta$ and $(B,\,G)$ is an algebra of type $\Delta$. Let $E$ be an equivalence relation on $B$ and let $C/E$ be the set of equivalence classes, $(C/E,\,G/E)$ is the multi-algebra of type $\Delta$ defined by
\begin{equation*}
(G/E)(i)([a_{1}],\ldots[a_{k_{i}}])=\{[b]:b=g_{i}(c_{1},\,\ldots\,c_{k_{i}})\wedge[c_{j}]=[a_{j}]\ \mathrm{for}\ j=1,\,\ldots,\,k_{i}\}.
\end{equation*}
Here we have used $[a]$ to denote the $E$ equivalence class of $a\in C$.
\end{definition*}

\section*{NOTE 51}

In Drake\index{Author}{Drake} \cite{bib:1974}, p.203 there is a proof of Ramsey's theorem (form 17) which uses K\"{o}nig's lemma (form [10 F]). However this proof also uses the well ordering theorem and hence does not show that [10 F] implies form 17. (Which is in fact false. See Blass\index{Author}{Blass} \cite{bib:1977a}.) Therefore the proof in Lolli\index{Author}{Lolli} \cite{bib:1977} that Ramsey's theorem implies form 10 is invalid because it uses ``[10 F] implies 17''

\section*{NOTE 52}

In Brunner\index{Author}{Brunner} \cite{bib:1981a} it is shown that there is a model of ZF in which $\mathbb{R}^{A}$ is metrizable and connected, where A is a suitably chosen set and the box topology is used on the product. (If for each $i\in I,\,X_{i}$ is a topological space, then the \emph{box product}\index{Subject}{box product} is the topological space whose underlying set is the Cartesian product $\prod_{i\in I}X_{i}$ and whose topology has as a basis, sets of the form $\prod_{i\in I}U_{i}$, where $U_{i}$ is open in $X_{i}$.)

\section*{NOTE 53}

In Zuckerman\index{Author}{Zuckerman} \cite{bib:1969a} and \cite{bib:1981} the relationship between versions of form $178(n,\,N)$ are considered. (Form $178(n,\,N)$: If $n\in \omega,\,n\geq 2$ and $N\subseteq\{1,2,\,\ldots,\,n-1\},\, N\neq\emptyset,\, MC(\infty,\,n,\,N)$: If $X$ is any set of $n$-element sets then there is a function $f$ with domain $X$ such that $\forall A\in X,\,f(A)\subseteq A$ and $|f(A)|\in N.)$ A sample of the results follows. Using the notation of Zuckerman \cite{bib:1981}, $S(n,\,\ell)$ is $MC(\infty,\,n,\,\{\ell\})$ and $T(n,\,\ell)$ is $MC(\infty,\, n,\,\{1,2,\,\ldots,\,\ell\})$.

\begin{theorem*}
If $m\leq n$ and $n-\ell<m$ then $S(n,\,\ell)\rightarrow T(m,\,\ell)$.
\end{theorem*}

\begin{theorem*}
If $l<n$ and $k\geq 0$, then $S(n,\,\ell)\wedge S(kn+\ell,\,1)$ imply $S(n,\,1)$.
\end{theorem*}

\begin{theorem*}
Assume $K\geq 1$

a) $\ell<n,\,T(kn,\,\ell)\rightarrow T(n,\,\ell)$.

b) for $\ell$ not of the form $jn$ for any $j\leq k,\,S(kn,\,\ell)\rightarrow T(n,\,l)$.
\end{theorem*}

\begin{theorem*}
If $l<n$ and $k\geq 0$, then $S(n,\,\ell)\wedge S(kn+\ell,\,1)$ imply $S(n,\,1)$.
\end{theorem*}

\begin{theorem*}
Assume $K\geq 1$

a) $\ell<n,\,T(kn,\,\ell)\rightarrow T(n,\ell)$.

b) for $\ell$ not of the form $jn$ for any $j\leq k,\,S(kn,\,\ell)\rightarrow T(n,\, l)$.
\end{theorem*}

\section*{NOTE 54}

Spector\index{Author}{Spector} \cite{bib:1980} proves that form 43 (dependent choice) implies [43
P] which implies [43 R] and [43 R] clearly implies [43 L]. Diener\index{Author}{Diener} \cite{bib:1994} shows that if forms 43 and [43 S] are altered by restricting them to relations which are partial orderings, the resulting forms remain equivalent to form 43. In addition, if [43 R] is altered by removing the restriction to partial orderings the resulting form remains equivalent to form 43.

\section*{NOTE 55}

The following results are from Gitik\index{Author}{Gitik} \cite{bib:1980}. Let $\Omega=\{\alpha\in$ On: $(\forall\beta\leq\alpha)(\beta= 0\vee(\exists\gamma)(\beta=\gamma+1)\vee(\forall n\in \omega)(\exists\gamma_{n}\leq\beta)(\beta=\bigcup_{n\in\omega}\gamma_{n})\}$. Gitik shows, by making small changes in $\mathcal{M}17$, that each of the following are consistent:
\begin{enumerate}
\item[1.] ZF + $\Omega=$ On.
\item[2.] ZF + $\Omega=\omega_{\Omega}$.
\item[3.] ZF + $\Omega=\omega_{\alpha+1}$, for any $\alpha$.
\end{enumerate}
In fact, item 1 holds in $\mathcal{M}17$ itself. It is clear that ZFC $\vdash\Omega=\omega_{1}$ (form 315).

In Gitik \cite{bib:1985}, the following theorem is given: Suppose $M$ is a countable model of ZFC and $\kappa$ is an almost huge cardinal in $M$. (See note 20 for the definition of various types of inaccessible cardinals.) If $S$ is any subset of $\kappa$ consisting of nonlimit cardinals, then there is a model $\mathcal{M}(S)$ of ZF such that the regular cardinals are exactly $\{\aleph_{\alpha}:\alpha\in S\cup\{0\}\}$. In the case that $S=\emptyset$, then all uncountable $\aleph_{\alpha}$'s are singular. Note, however, that in $\mathcal{M}17$ all uncountable $\aleph_{\alpha}$'s are singular. Further, the construction of $\mathcal{M}17$ requires only the existence of a strongly compact cardinal which is weaker than the existence of an almost huge cardinal.

\section*{NOTE 56}

Hickman\index{Author}{Hickman} \cite{bib:1980b}\index{Subject}{amorphous set} generalizes the notion of amorphous as follows:

\begin{definition*}\
\begin{enumerate}
\item[1.] $X$ is an $\aleph_{\alpha}$ \emph{set} if $|X|$ is incomparable with $\aleph_{\alpha}$ and for all $\beta<\alpha,\,\aleph_{\beta}<|X|$.
\item[2.] $X$ is $\aleph_{\alpha}$ \emph{minimal} if $X$ is an $\aleph_{\alpha}$ set and for all $Y\subseteq X$, either $|Y|<\aleph_{\alpha}$ or $|X-Y|<\aleph_{\alpha}$, (abbreviated $X$ is $\alpha$-m).
\item[3.] $X$ is \emph{weakly} $\aleph_{\alpha}$ \emph{minimal}\index{Subject}{weakly $\aleph_{\alpha}$, minimal} if $X$ is an $\aleph_{\alpha}$ set and for all $Y\subseteq X$ at most one of $Y$ and $X-Y$ is an $\aleph_{\alpha}$ set, ($X$ is $\alpha$-wm).
\item[4.] $X$ is \emph{pseudo} $\aleph_{\alpha}$ \emph{minimal}\index{Subject}{pseudo $\aleph_{\alpha}$ minimal} if $X$ is an $\aleph_{\alpha}$ set and for all $Y\subseteq X$, exactly one of $Y$ and $Y-X$ is an $\aleph_{\alpha}$ set, ($X$ is $\alpha$-pm).
\end{enumerate}
\end{definition*}

The following are proved:

\setcounter{theorem}{0}
\begin{theorem}\label{part04:thm1a}
$X$ is $0$-m iff $X$ is $0$-wm iff $X$ is 0-pm.
\end{theorem}

\begin{theorem}\label{part04:thm2a}
(a) For all $\alpha,\,X$ is $(\alpha+1)$-pm implies $X$ is $(\alpha+1)$-wm and (b) $X$ is $(\alpha+1)$-wm iff $X$ is $(\alpha+1)$-m.
\end{theorem}

\begin{theorem}\label{part04:thm3}
If $\alpha$ is a limit ordinal $(\alpha\neq 0)$, then

\noindent (a) $X$ is $\alpha$-m implies $X$ is $\alpha$-wm.

\noindent (b) $X$ is $\alpha$-wm iff $X$ is $\alpha$-pm.
\end{theorem}

\begin{theorem}\label{part04:thm4}
For all ordinals $\alpha,\,Con(ZF+\exists X(X$ is $\alpha$-m$))$.
\end{theorem}

\begin{theorem}\label{part04:thm5}
(a) For all ordinals $\alpha$, the following is not theorem of $ZF$:
\begin{equation*}
\forall X(X\,is\,(\alpha+1)\text{-m}\rightarrow X\ is\ (\alpha+1)\text{-mp})
\end{equation*}
(b) For all limit ordinals $\alpha$, it is not a theorem of $ZF$ that
\begin{equation*}
\forall X(X\ is\ \alpha\text{-mp}\rightarrow X\ is\ \alpha\text{-m}).
\end{equation*}
\end{theorem}

\begin{theorem}\label{part04:thm6}
If $\aleph_{\alpha}$ is regular, then Con(ZF + there is a linearly ordered $(\alpha+1)$-pm set).
\end{theorem}

\begin{theorem}\label{part04:thm7}
For any $\alpha$, Con(ZF + there is a linearly ordered $(\alpha+1)$-m set).
\end{theorem}

\begin{theorem}\label{part04:thm8}
If $\alpha$ is a limit ordinal there are no linearly ordered $\alpha$-wm sets.
\end{theorem}

\textsc{Open Problem.}
The consistency of the existence of linearly ordered $(\alpha+1)$-pm set where $\aleph_{\alpha}$ is not regular.

\section*{NOTE 57}
Truss\index{Author}{Truss} \cite{bib:1995} studies the various structures an amorphous set\index{Subject}{amorphous set} can carry. (A set is \emph{amorphous} if it is infinite but is not the disjoint union of two infinite sets.) If $\pi$ is a partition of an amorphous set $U$ into infinitely many parts, then each part must be finite. (Otherwise,
if $X$ were an infinite part, then $U=X\cup (U-X)$, contradicting the fact that $U$ is amorphous.) Define a map $f$ from $U$ into $\omega$ such that for each $x\in U,\,f(x)=n$ where $x\in\xi \in\pi$ and $|\xi|=n$. Since $U$ is amorphous, the range of $f$ must be finite. Since $U=\bigcup\{f^{-1}(n): n\in$ range of $f$\}, there is a unique $n$ such that $f^{-1}(n)$ is infinite. Let $n(\pi)$ be this unique $n$, then $\{\xi\in\pi:|\xi|=n(\pi)\}$ is infinite.

Let $\Pi(U)$ be the set of finitary partitions of $U$ and let $n(U)= \sup\{n(\pi):\pi\in\Pi(U)\}$. If $n(U)$ is finite, $U$ is called \emph{bounded}, and if $n(U)$ infinite, $U$ is called \emph{unbounded}.

Two interesting results that Truss obtains are:

\begin{theorem*}
The following are consistent with $ZF$:
\begin{enumerate}
\item[1.] There is an amorphous set and all amorphous sets are bounded.
\item[2.] There is an amorphous set and all amorphous sets are unbounded.
\end{enumerate}
\end{theorem*}

See Truss \cite{bib:1995} for additional results.

\section*{NOTE 58}

J. L. Hickman\index{Author}{Hickman} \cite{bib:1979b} shows that if $\Gamma$ and $\Delta$ are Dedekind finite cardinals then the following are equivalent: (1) lub$\{\Gamma,\,\Delta\}$ exists, (2) $\Gamma$ and $\Delta$ are comparable and (3) gib $\{\Gamma,\,\Delta\}$ exists.

\section*{NOTE 59}

Form 191 is the statement: There is a set $S$ such that for every set $a$, there is an ordinal $\alpha$ and a function from $S\times\alpha$ onto $a$ and form 182 is: There is an aleph whose cofinality is greater than $\aleph_{0}$. In Blass\index{Author}{Blass} \cite{bib:1979} it is proved, assuming form 191, that $\forall A$ there is a limit ordinal $\alpha$ such that $A$ cannot be mapped cofinally onto $\alpha$. If $A$ cannot be mapped cofinally onto $\alpha$, then $A$ cannot be mapped cofinally onto $\aleph_{\alpha}$. Therefore form 191 implies form 182.

Blass also shows that if 191 holds with a set $S$ then form 14 (the Boolean Prime Ideal Theorem) holds if and only if there is an ultrafilter $U$ on $S^{<\omega}$ that is regular in the sense that $\{p\in S^{<\omega}: s\in$ range $p\}$ is in $U$ for every $s\in S$.

With regard to $(191,n)$, Blass shows that 191 holds in most known models of $\mathrm{ZF}^{0}$ or ZF. In particular,
\begin{enumerate}
\item[1.] Every permutation model whose class of atoms forms a set satisfies 191.
\item[2.] Every symmetric model satisfies 191. (A symmetric model is a model constructed from a ground model $M$ of ZFC by taking a complete Boolean algebra $B$ in $M$, an $M$ generic filter $G$ in $B$, a group $\mathcal{G}$ of automorphisms of $B$ and a normal filter of subgroups of $\mathcal{G}$. The model consists of members of $M[G]$ that have hereditarily symmetric names.)
\item[3.] If $V$ is any model of ZFC and $A\in V$, then the submodel $HOD(A)$ of sets hereditarily ordinal definable over A satisfies 191.
\end{enumerate}

\section*{NOTE 60}

Definitions from category theory, in particular for forms 188, 189, 190, 192 and 193 from Blass\index{Author}{Blass} \cite{bib:1979}.

\begin{definition*}
Assume $\mathbf{C}$ is a category
\begin{enumerate}
\item[1.] An object $p$ in $\mathbf{C}$ is \emph{projective}\index{Subject}{projective object in a category} (in $\mathbf{C}$) if for all objects $a$ and $b$ in $\mathbf{C}$, if $f\in \mathrm{hom}(a,\, b)$ is an epimorphism and $g\in \mathrm{hom}(p,\,b)$, then there is an $h\in \mathrm{hom}(p,\,a)$ such that $f\circ h=g$.
\item[2.] $p$ is \emph{injective}\index{Subject}{injective object in a category} if for all objects $a$ and $b$ in $\mathbf{C}$, if $f\in \mathrm{hom}(b,\,a)$ is a monomorphism and $g\in \mathrm{hom}(b,p)$ then there is an $h\in \mathrm{hom}(a,p)$ such that $h\circ f=g$.
\item[3.] An Abelian group $G$ is \emph{projective}\index{Subject}{projective Abelian group} if for all Abelian groups $A$ and $B$ if $f:A\rightarrow B$ is a homomorphism from $A$ onto $B$ and $g:G\rightarrow B$ is a homomorphism then there is a homomorphism $h:G\rightarrow A$ such that $f\circ h=g$.
\item[4.] An Abelian group $G$ is \emph{injective}\index{Subject}{injective Abelian group} if for all Abelian groups $A$ and $B$, if $f:B\rightarrow A$ is a one to one homomorphism and $g:B\rightarrow G$ is a homomorphism then there is a homomorphism $h:A\rightarrow G$ such that $h\circ f=g$.
\item[5.] A set $X$ is \emph{projective}\index{Subject}{projective set} if for all sets $A$ and $B$, for all $f$ from $A$ onto $B$ and for all $g$ from $X$ into $B$ there is a function $h$ from $X$ into $A$ such that $f\circ h=g$, (i.e., For any $X$-indexed family $\{y_{t}:t\in X\}$ of non-empty sets there is a choice function.)
\end{enumerate}
\end{definition*}

\section*{NOTE 61}

Kanovei\index{Author}{Kanovei} \cite{bib:1979} studies the relationships between $AC(K)$ and $DC(K)$ where $K$ is either in the analytic or projective hierarchy\index{Subject}{projective hierarchy}. We give definitions and summarize some of the results from Kanovei \cite{bib:1979}. $AC(K)$ is: For any $P\subseteq \omega\times {^{\omega}}\omega$ with domain $\omega$, if $P\in K$, then there is a sequence $\langle x_{k}: k\in \omega\rangle$ of elements of ${^{\omega}}\omega$ with $\langle k,\,x_{k}\rangle\in P$ for all $k\in\omega$. And

\noindent $DC(K)$ is: For every $P\subseteq {^{\omega}}\omega\times {^{\omega}}\omega$ with domain $^{\omega}\omega$, if $P\in K$ then for any $y\in{^{\omega}}\omega$ there is a sequence $\langle x_{k}: k\in\omega\rangle$ with $x_{0}=y$ and for all $k\in\omega, \langle x_{k},\,x_{k+1}\rangle\in P$. The sets in the analytic hierarchy\index{Subject}{analytic hierarchy} $(\varSigma_{n}^{1},\, \varPi_{n}^{1}$ and
$\varDelta_{n}^{1}$ for $n\in \omega)$ and the sets in the projective hierarchy $(\boldsymbol{\Sigma}_{n}^{1},\, \boldsymbol{\Pi}_{n}^{1}$ and $\boldsymbol{\Delta}_{n}^{1}$
for $n\in\omega)$ are defined as follows: For purposes of these definitions $^{k,\ell}\omega$ is ${^{k}}\omega\times{^{\ell}}(^{\omega}\omega)$ and a relation $R\subseteq {^{k,\ell}}\omega$ is of rank $(k,\,\ell)$.
\begin{enumerate}
\item[1.] $R$ of rank $(k,\,\ell)$ is \emph{arithmetic} if for all $\mathbf{m}\in {^{k}} \omega$ and all $\boldsymbol{\alpha}\in{^{\ell}}(^{\omega}\omega)$, $R(\mathbf{m},\,\boldsymbol{\alpha})$ iff $Q_{1}x_{1}\ldots Q_{r}x_{r}P(\mathbf{x},\,\mathbf{m},\,\boldsymbol{\alpha},)$ where $P\subseteq^{r+k,\ell}\omega$ is recursive, and each $Q_{i}$ is either
$\exists$ or $\forall$.
\item[2.] $R$ of rank $(k,\,\ell)$ is \emph{analytic} if $R(\mathbf{m},\,\boldsymbol{\alpha})$ iff $Q_{1}\beta_{1}\ldots Q_{r}\beta_{r}P(\mathbf{m},\,\boldsymbol{\alpha},\,\beta_{1},\,\ldots,\,\beta_{r})$ where $P$ is of rank $(k,\,\ell+r)$ and is arithmetic.
\item[3.] (The analytic hierarchy)
\begin{enumerate}
\item[(i)] $\varSigma_{0}^{1}=\varPi_{0}^{1}=$ the arithmetic relations.
\item[(ii)] $\varSigma_{r+1}^{1}=\{\exists\beta P(\mathbf{m},\,\beta,\,\boldsymbol{\alpha}):P\in\varSigma_{r}^{1}\wedge P$ has rank $(k,\, \ell+1)\}$
\item[(ii)] $\varPi_{n+1}^{1}=\{\forall\beta P(\mathbf{m},\,\beta,\,\boldsymbol{\alpha}):P\in\varSigma_{r}^{1}\wedge P$ has rank $(k,\, \ell+1)\}$
\item[(iv)] $\varDelta_{r}^{1}=\varSigma_{r}^{1}\cap \varPi_{r}^{1}$.
\end{enumerate}
\item[4.] $R$ of rank $(k,\,\ell)$ is \emph{recursive} in $\boldsymbol{\beta}$ iff the characteristic function $K_{R}(\mathbf{m},\, \boldsymbol{\alpha})$ of $R$ is $G(\mathbf{m},\,\boldsymbol{\alpha},\,\boldsymbol{\beta})$ for some partial recursive $G$.
\item[5.] Assume $\boldsymbol{\beta}$ has rank $(0,\,\ell)$ then $\varSigma_{r}^{1}[\beta]$ and $\varPi_{r}^{1}[\boldsymbol{\beta}]$ are defined by $R\in\varSigma_{r}^{1}[\boldsymbol{\beta}]$ iff $R(\mathbf{m},\,\boldsymbol{\alpha})\leftrightarrow S(\mathbf{m},\,\boldsymbol{\alpha},\,\boldsymbol{\beta})$ for some $S\in\varSigma_{r}^{1}$. $\varPi_{r}^{1}[\boldsymbol{\beta}]$ is defined similarly.
\item[6.] (The projective hierarchy)
\begin{enumerate}
\item[(i)] $\boldsymbol{\Sigma}_{r}^{1}=\bigcup\{\varSigma_{r}^{1}[\beta]:b\in
{^{\omega}}\omega\}$
\item[(ii)] $\boldsymbol{\Pi}_{r}^{1}=\bigcup\{\varPi_{r}^{1}[\beta]:b\in {^{\omega}}\omega\}$
\item[(iii)] $\boldsymbol{\Delta}_{r}^{1}=\boldsymbol{\Sigma}_{r}^{1}\cap\boldsymbol{\Pi} 1_{r}$.
\end{enumerate}
\end{enumerate}
In Kanovei\index{Author}{Kanovei} \cite{bib:1979}, it is shown for any $n\geq 2$ that, $AC(\varPi_{n+1}^{1})$ does not follow, in ZF from $DC(\boldsymbol{\Pi}_{n}^{1})$, that $DC(\boldsymbol{\Pi}_{n}^{1})$ implies $AC(\varPi_{n}^{1})$ and that $AC(\varPi_{n}^{1})$ does not imply $DC(\boldsymbol{\Pi}_{n}^{1})$. Kanovei \cite{bib:1976} has further results on $AC$ and the projective hierarchy. See forms [0 X], 194 and 199.

\section*{NOTE 62}

The following definitions for form 195 are from Harazisville (see Kharazishvili\index{Author}{Kharazishvili}
\cite{bib:1979}.)

\begin{definition*}
\begin{enumerate}
\item[1.] If $X$ and $Y$ are vector spaces over a field $K$ then a \emph{general linear system}\index{Subject}{general linear system} is a relation $S\subseteq X\times Y$ which has domain $X$ and is a subspace of $X\times Y$.
\item[2.] If $S$ is as in 1. then a \emph{linear global reaction of}\index{Subject}{linear global retraction} $S$ is a function $\Phi:C\times X\rightarrow Y$ such that
\begin{enumerate}
\item[(i)] $C$ is a vector space over $K$
\item[(ii)] $\Phi$ is a global reaction of $S$ (i.e., $(x,\,y)\in S$ iff $x\in X$ and $(\exists c)(c\in C$ and $y=\Phi(c,\,x)).)$
\item[(iii)] There are linear functions $\Phi_{1}:C\rightarrow Y$ and $\Phi_{2}:X\rightarrow Y$ such that $\Phi(c,\,x)=\Phi_{1}(c)+\Phi_{2}(x)$.
\end{enumerate}
\end{enumerate}
\end{definition*}

\section*{NOTE 63}

Since most class forms of AC\index{Subject}{class forms of AC} are not consequences of form 1 (AC) we are not including class forms of AC in this project. A summary of what is known about class forms of AC can be found in H. Rubin/J. Rubin\index{Author}{H. Rubin/ J. Rubin} \cite{bib:1985}, p. 185. Also see Felgner/Flannagan\index{Author}{Felgner/Flannagan}\index{Author}{Flannagan} \cite{bib:1978}, Felgner\index{Author}{Felgner} \cite{bib:1971c} and \cite{bib:1976}, Flannagan \cite{bib:1974}, Fleischer\index{Author}{Fleischer} \cite{bib:1973}, Harper/Rubin\index{Author}{Harper/Rubin} \cite{bib:1976}. part II and Howard/A. Rubin/J. Rubin\index{Author}{Howard/A. Rubin/J. Rubin} \cite{bib:1978}.

\section*{NOTE 64}

Form 200 is studied in Dawson/Howard\index{Author}{Dawson/Howard} \cite{bib:1976} where it is shown that
\begin{enumerate}
\item[1.] Each of
\begin{enumerate}
\item[(a)] $\exists$ infinite $x$ such that $2^{x}<x!$
\item[(b)] $\exists$ infinite $x$ such that $2^{x}>x!$ and
\item[(c)] $\exists$ infinite $x$ such that $2^{x}$ and $x!$ are incomparable
\end{enumerate}
are consistent with ZF. ($x!$ is the cardinality of the set of all one to one functions from $x$ onto $x$.).
\item[2.] for all $x,\,2x=x$ implies $2^{x}\leq x!$.
\end{enumerate}

In Pincus\index{Author}{Pincus} \cite{bib:1978} a model of ZF is constructed in which AC (form 1) is false, $\forall$ infinite $x,\, 2x=x$ (form 3) is false and $\forall$ infinite $x,\,2^{x}=x!$ is true. Also true in this mod
el are $\forall$ infinite $x,\,2^{x}=|x^{x}|$ and $\forall$ infinite $x,\,2^{x}=|x^{x}-x!|$.

\section*{NOTE 65}

Truss\index{Author}{Truss} \cite{bib:1978} shows that in Feferman's\index{Author}{Feferman} model, $\mathcal{M}2$ (Feferman \cite{bib:1965}), the following hold:
\begin{enumerate}
\item[1.] $\forall$ ordinals $\alpha,\,2^{\aleph_{\alpha}}=\aleph_{\alpha+1} \cdot 2^{\aleph_{0}}$
\item[2.] For all $X$, if $[X]^{2}$ has a choice function, then $\exists\alpha$ such that $|X|\leq 2^{\aleph_{\alpha}}$.
\item[3.] Form 40 $(C(WO,\,\infty))$.
\item[4.] $\mathcal{P}(\omega)$ has no non-principal prime ideal, (the negation of form 70).
\item[5.] $\neg C(\infty,\,2)$, (the negation of form 88).
\item[6.] There is a partition of $\mathcal{P}(\omega)$ into non-empty subsets without a choice function, (the negation of form 203).
\item[7.] Form 204, For every infinite $X$ there is a function from $X$ onto $2\times X$.
\end{enumerate}

\section*{NOTE 66}

In this note we prove that form 322 $(KW(WO,\,\infty))$ is true in $\mathcal{N}1$. (Brunner\index{Author}{Brunner} \cite{bib:1982a}, (Proposition~ 4.5) also has a proof of this result.)

Assume that $X$ is a well ordered set in $\mathcal{N}1$ and that $\forall y\in X,\,|y|\geq 2$. Let $E$ be a (finite) support of a well ordering of $X$ and let fix$(E)$ be the group of permutations of $A$ (the set of atoms) which fix $E$ pointwise. Then for all $\phi\in$ fix$(E)$, $\phi$ fixes $X$ pointwise. Choose an atom $a\in A-E$. We will show that for every $y\in X,\,E\cup\{a\}$ is a support of a non-empty proper subset $y^{\prime}$ of $y$. This will suffice since it follows that the function $y\mapsto y^{\prime}$ is in $\mathcal{N}1$.

Choose an element $y\in X$. If every element of $y$ has support $E$ then any nonempty proper subset of $y$ has support $E$ and therefore has support $E\cup\{a\}$. So we may assume that $y$ has some element $t$ such that $E$ is not a support of $t$. Choose a finite non-empty subset $F$ of $A$ such that $E\cap F=\emptyset$ and $E\cup F$ is a support of $t$ and such that $F$ is minimal in the sense that for no proper subset $F^{\prime}$ of $F$ is $E\cup F^{\prime}$ a support of $t$. We may assume without loss of generality that $a\in F$. (If $a\not\in F$ then we replace $t$ by $\psi(t)$ and $F$ by $\psi(F)$ where $\psi$ is a permutation of $A$ chosen so that $\psi\in$ fix$(E)$ and $a\in\psi(F).)$ Let $y'=\{\phi(t):\phi \in$ fix$(E\cup\{a\})\}$. Clearly $y^{\prime}$ has support $E\cup\{a\},\,y^{\prime}\subseteq y$ and $y^{\prime}\neq\emptyset$ (since $t\in y^{\prime}$). It only remains to show that $y^{\prime}\neq y$.

\begin{lemma*}
If $\beta$ and $\gamma$ are in fix$(E)$ and $\beta(F)\neq\gamma(F)$ then $\beta(t)\neq\gamma(t)$.
\end{lemma*}

\begin{proof}
Assume $b\in\beta(F)-\gamma(F)$. Since $F$ is minimal, there is no proper subset $B$ of $\beta(F)$ for which $E\cup B$ is a support of $\beta(t)$. Hence there is a $\psi\in$ fix$((E\cup \beta(F))-\{b\})$ such that $\psi(\beta(t))\neq\beta(t)$. We can find such a $\psi$ for which $\psi(b)\not\in\gamma(F)$. (If $\psi(b)\in\gamma(F)$ then we choose $d\in A-(\gamma(F)\cup\beta(F)\cup E)$. Since $\psi(\beta)\not\in\beta(F)$ the cycle $\rho=(d,\,\psi(b))$ of length two fixes $E\cup\beta(F)$ pointwise so that $\rho(\beta(t))=\beta(t)$. If we let $\psi^{\prime}=\rho \circ\psi$ then $\psi^{\prime}$ fixes $(E\cup \beta(F))-\{b\}$ pointwise. Further, since $\rho(\beta(t))=\beta(t)$ and $\psi(\beta(t))\neq\beta(t)$, we conclude that $\psi^{\prime}(\beta(t))=\rho(\psi(\beta(t)))\neq\beta(t)$. We can therefore replace $\psi$ by $\psi^{\prime}$.)

Let $\sigma$ be the cycle $(b,\,\psi(b))$. Since $\sigma$ agrees with $\psi$ on a support of $\beta(t)$, $\sigma(\beta(t))=\psi(\beta(t))\neq\beta(t)$. Since $\sigma$ fixes $\gamma(F)$ pointwise, $\sigma(\gamma(t))=\gamma(t)$. Therefore $\beta(t)\neq\gamma(t)$. This completes the proof of the lemma.
\end{proof}

Choose an $\eta\in$ fix$(E)$ so that $\eta(F) \cap F=\emptyset$. Then for every $\phi\in$ fix$(E\cup\{a\})$, $\eta(F)\neq\phi(F)$. It follows from the lemma that for all such $\phi,\,\eta(t)\neq\phi(t)$. Therefore $\eta(t)\in(y-y^{\prime})$.

\section*{NOTE 67}

\begin{theorem*} Form 203 (Every partition of $\mathcal{P}(\omega)$ into non-empty subsets has a choice function) implies form 94 ($C(\aleph_{0},\,\infty,\,\mathbb{R})$).
\end{theorem*}

\begin{proof}
Form 203 implies that every partition of $\mathcal{P}(\omega)\times \omega$ into non-empty subsets has a choice function. Assume $\{A_{i}:i\in \omega\}$ is a denumerable set of sets of reals. (We identify $\mathbb{R}$ with $\mathcal{P}(\omega)$.) For each $i\in \omega$, let $B_{i}=A_{i}\times\{i\}$. Then the family $\{B_{i}:i\in \omega\}$ is a disjoint family from $\mathcal{P}(\omega)\times \omega$ and therefore, by 203, has a choice function. This induces a choice function on $\{A_{i}:i\in \omega\}$.
\end{proof}

\section*{NOTE 68}

\begin{theorem*}
form 3 ($2m =m$ for infinite cardinals $m$) implies form 53 ($m^{2}\leq 2^{m}$ for infinite cardinals $m$).
\end{theorem*}

\begin{proof}
Assume that for all infinite cardinals $m,\,2m=m$ and that $X$ is an infinite set. Then $f(x,\, y)=\{(0,\,x),\,(1,\,y)\}$ is a 1-1 function from $X\times X$ into $\mathcal{P}(2\times X)$.
\end{proof}

\section*{NOTE 69}

\begin{definition*}
$|x|\leq^{\ast}|y|$\index{Subject}{$\leq^{\ast}$} (or $x\precsim^{\ast} y)$\index{Subject}{$\precsim^{\ast}$} if $y$ can be mapped onto $x$.
\end{definition*}

\begin{theorem*}
Form [40 B] implies form 208.
\end{theorem*}

\begin{proof}

[40 B] is ``For all ordinals $\alpha$ and all sets $x$ if $\aleph_{\alpha}\leq^{\ast}|x|$, then $\aleph_{\alpha}\leq|x|$.'' and 208 is ``$2^{\aleph_{\alpha}}\geq\aleph_{\alpha+1}$'' (Form 208 depends on the ordinal $\alpha$).) Assume [40 B]. As is shown in Kunen\index{Author}{Kunen} \cite{bib:1980}, chapter 1, $|\aleph_{\alpha}\times\aleph_{\alpha}|=\aleph_{\alpha}$ and therefore, $|\mathcal{P}(\aleph_{\alpha}\times\aleph_{\alpha})|= |\mathcal{P}(\aleph_{\alpha})|$. The function $f$ from $\mathcal{P}(\aleph_{\alpha}\times\aleph_{\alpha})$ defined by $f(R)=$ the order type of $R$, if $R$ is a well ordering and $f(R)=0$ otherwise is onto $\aleph_{\alpha+1}$. It follows that $\aleph_{\alpha+1}\leq^{\ast}|\mathcal{P}(\aleph_{\alpha})|$ and therefore by [40 B], $\aleph_{\alpha+1}\leq|\mathcal{P}(\aleph_{\alpha})|$.
\end{proof}

\section*{NOTE 70}

We give a proof that form [62 E] $(KW(\infty,\,<\aleph_{0}))$ implies form 62 $(C(\infty,\,<\aleph_{0}))$. Let $x$ be a set of finite sets and let $y=\{u:(\exists v\in x)u\subseteq v\}$. By [62 E], there is a Kinna-Wagner selection function, $f$, on $y$. We define a choice function $g$ on $x$ as follows: Let $u\in x$ be such that $|u|=n$. There exists an $m\in \omega,\,m<n$, such that $f^{m}(u)$ is a singleton. Define $g(u)$ to be the single element of $f^{m}(u)$. (Form 323, $KW(WO,\,\infty)$, clearly implies [62 E].)

\section*{NOTE 71}

In this note we give definitions from Morillon\index{Author}{Morillon} \cite{bib:1988} for forms [14 L], [14 BP] through [14 BX], [14 CC] through [14 CH], [118 I] through [118 T], 331, 332, 343, 344, and [345 C] through [345 E].

Let $X$ be a topological space, let $E=\{(A,\,B):A$ and $B$ are closed in $X$ and $A \cap B=\emptyset\}$ and let $F=\{(U,\,V):U$ and $V$ are open in $X$ and $U\cap V=\emptyset$\}.
\begin{enumerate}
\item[1.] $f:E\rightarrow F$ is a \emph{normality operator}\index{Subject}{normality operator} on $X$ if $\forall(A,\,B)\in E,\,f(A,\, B)=(U,\,V)$ where $A\subseteq U$ and $B\subseteq V$.
\item[2.] $X$ is \emph{effectively normal}\index{Subject}{effectively normal} if $X$ is $T_{2}$ and there is a normality operator on $X$.
\item[3.] $X$ is \emph{hereditarily normal}\index{Subject}{hereditarily normal} (or \emph{completely normal})\index{Subject}{completely normal} if all subspaces of $X$ are normal.
\item[4.] $X$ is \emph{hereditarily effectively normal} if all subspaces of $X$ are effectively normal.
\item[5.] A set $I\subseteq \mathcal{P}(X)$ is \emph{discrete} if $\forall x\in X,\,\exists$ a neighborhood $N$ of $x$ such that $N\cap A\neq\emptyset$ for at most one element $A\in I$.
\item[6.] $X$ is\index{Subject}{hereditarily effectively normal} \emph{collectionwise normal}\index{Subject}{collectionwise normal} if $X$ is $T_{2}$ and for every discrete $I\subseteq\{U:U$ is closed in $S\}$ there is a set $J$ of pairwise disjoint open subsets of $X$ such that for every $F\in I,\,\exists U\in J$ such that $F\subseteq U$.
\item[7.] Let $D=$\{ $I:I$ is a discrete collection of closed subsets of $X$\} and $O=\{J:J$ is a collection of pairwise disjoint open subsets of $X$\}. Then $f:D\rightarrow O$ is a \emph{collectionwise normality operator}\index{Subject}{collectionwise normality operator} on $X$ if $\forall I\in D$, if $f(I)=J$ then $\forall F\in I,\,\exists U\in J$ such that $F\subseteq U$. $N_{c}(X)$ denotes the set of collectionwise normality operators on $X$.
\item[8.] $X$ is \emph{effectively collectionwise normal}\index{Subject}{effectively collectionwise normal} if X is $T_{2}$ and $N_{c}(X)\neq\emptyset$.
\item[9.] $X$ is \emph{highly separated}\index{Subject}{highly separated} if $X$ is $T_{2}$ and $\{N_{c}(F):F\subseteq X\}$ has a choice function.
\item[10.] $X$ is \emph{functionally separated}\index{Subject}{functionally separated} if for all pairs $(A,\,B)$ of closed disjoint subsets of $X$ there is a function $f:X\rightarrow[0,\,1]$ such that $A\subseteq f^{-1}(0)$ and $B\subseteq f^{-1}(1)$.
\item[11.] $X$ is \emph{Tietze-Urysohn}\index{Subject}{Tietze-Urysohn space} if for every closed $F\subseteq X$ and every continuous function $f:F\rightarrow[0,\,1]$ there is a continuous function $g:X\rightarrow[0,\,1]$ which extends $f$.
\item[12.] Let $C=$\{ $g:g:X\rightarrow[0, 1]$ and $g$ is continuous\}. A function $f:E\rightarrow C$ is a \emph{Urysohn operator}\index{Subject}{Urysohn operator} on $X$ if $\forall(A,\,B)\in E$, if $U=f(A,\,B)$ then $A\subseteq U^{-1}(0)$ and $B\subseteq U^{-1}(1)$. $U(X)$ denotes the set of Urysohn operators on $X$.
\item[13.] Let $B=\{(F,\,U):F$ is a closed subset of $X$ and $U:F\rightarrow[0, 1]$ is continuous\}. A function $f:B\rightarrow C$ is a \emph{Tietze-Urysohn operator}\index{Subject}{Tietze-Urysohn operator} if $\forall(F,\,U)\in B$, if $f(F,\, U)=V$ then $V$ extends $U$.
\item[14.] $X$ is \emph{effectively functionally separated}\index{Subject}{effectively functionally separated} if $U(X)\neq\emptyset$.
$X$ is \emph{effectively Tietze-Urysohn}\index{Subject}{effectively Tietze-Urysohn} if T-U$(X$) $\neq\emptyset$. (Where T-U($X$) is the set of all Tietze-Urysohn operators on $X$.)
\item[15.] $X$ is \emph{highly functionally separated}\index{Subject}{highly functionally separated} if $\{U(A):A\subseteq X\}$ has a choice function.
\item[16.] $X$ is \emph{highly Tietze-Urysohn}\index{Subject}{highly Tietze-Urysohn} if \{T-U$(A):A\subseteq X\}$ has a choice function.
\item[17.] $X$ is \emph{collectionwise separated}\index{Subject}{collectionwise separated} (or \emph{collectionwise Hausdorff})\index{Subject}{collectionwise Hausdorff} if it is $T_{2}$ and for every discrete set $I$ of singletons from $X$ there is a set $J$ of open subsets of $X$ such that for every $F\in I$, there is a unique $O\in J$ containing (the singleton) $F$.

\noindent Assume that $(X,\,\leq)$ is a linear order.
\item[18.] $A\subseteq X$ is \emph{convex} if $\forall(x,\,y)\in A^{2},\,\forall u\in X$, if $x<u<y$ then $u\in A$.
\item[19.] $(X, \leq)$ is \emph{conditionally complete} if each non-empty subset with an upper bound has a least upper bound and each non-empty subset with a lower bound has a greatest lower bound.
\item[20.] $(X,\,\leq)$ is \emph{complete} if $X=\emptyset$ or every subset has a least upper bound and a greatest lower bound.
\item[21.] Let $X$ be a topological space, let $E$ be defined as above and let $\Omega(X)$ be the set of open subsets of $X$. A function $f:E\rightarrow\Omega(X)$ is a \emph{monotone normality operator}\index{Subject}{monotone normality operator} on $X$ if for every $(F,\,G)\in E$ the following conditions hold
\begin{enumerate}
\item[(i)] $F\subseteq f(F,\,G)\subseteq \overline{f(F,\,G)}\subseteq X-G$.
\item[(ii)] If $(F^{\prime},\,G^{\prime})\in E,\,F\subseteq F^{\prime}$ and $G^{\prime}\subseteq G$ then $f(F,\,G)\subseteq f(F^{\prime},\,G^{\prime})$.
\end{enumerate}
\item[22.] A topological space $X$ is \emph{monotonely normal} if it is $T_{1}$ and there is a monotone normality operator on $X$.
\item[23.] Let $(X,\,\leq)$ be a complete linear order and let $E$ be the set of non-empty open intervals of $X$. An \emph{increasing choice function} on $(X,\,\leq)$ is a function $\ast:E\rightarrow X$ such that
\begin{enumerate}
\item[(i)] $\forall I\in E,\,\ast(I)\in I$.
\item[(ii)] If $(a,\,b)\in E$ and $(b,\,c)\in E$ then $\ast((a,\,b))\leq\ast((a,\,c))\leq\ast((b,\,c))$.
\end{enumerate}
\item[24.] A topological\index{Subject}{increasing choice function} space $X$ is \emph{completely regular}\index{Subject}{completely regular} if for all open $U\subseteq X$ and all $p\in U$, there is a continuous $f:X\rightarrow[0,\,1]$ such that $f(p)=0$ and $\forall x\in X-U,\,f(x)=1$.
\item[25.] An ordered set $(T,\,\leq)$ is a \emph{lattice} if any two elements have a greatest lower bound and a least upper bound. A lattice is \emph{bounded} is it has a smallest and largest element.
\item[26.] An ideal $I$ in a lattice $(T,\,\leq)$ is \emph{prime} if $\forall x,\,y\in T$, if $x\wedge y\in I$ then either $x\in I$ or $y\in I$.
\item[27.] If $T$ is a bounded distributive lattice a function $f$ is a \emph{prime ideal choice function} on $T$ if for each proper ideal $I$ of $T,\,f(I)$ is a prime ideal containing $I$.
\item[28.] Assume that $T$ is a bounded distributive lattice in which every subset has a least upper bound then $T$ is $\vee$-\emph{compact}\index{Subject}{$\vee$-compact} if for all $F\subseteq T$ such that $\bigvee F=1$, there is a finite $F^{\prime}\subseteq F$ such that $\bigvee F^{\prime}=1.$ (If $(\Omega,\,\subseteq)$ is the lattice of open sets of a topological space $X$ then $X$ is compact if and only if $(\Omega,\,\subseteq)$ is $\vee$-compact.)
\item[29.] Let $E$ be a set. An \emph{ultra filter choice function}\index{Subject}{ultrafilter choice function} $f$ on $E$ is a function such that for all proper filters $I$ in $\mathcal{P}(E)$, $f(I)$ is an ultra filter containing $I$.
\item[30.] Let $E$ be a set and $U$ an ultra filter on a set $I$. $U$ is $E$-\emph{adequate} if for every filter $F$ in $\mathcal{P}(E)$ there is a function $g:I\rightarrow E$ such that $g(U)$ contains $F$.
\item[31.] A topological space is \emph{irreducible} if it is not empty and is not the union of two proper closed subsets.
\item[32.] A topological space $X$ is \emph{sober}\index{Subject}{sober topological space} if every closed irreducible subset of $X$ is the closure of a unique singleton.
\item[33.] A \emph{Banach algebra}\index{Subject}{Banach algebra} is a normed algebra in which every Cauchy filter converges.
\item[34.] Let $(I,\,\leq)$ be a preorder (reflexive and transitive), let $(E_{\alpha})_{\alpha\in I}$ be a family of sets and for each $(\alpha,\,\beta)\in I\times I$ let $f_{\alpha,\beta}$ be a function from $E_{\beta}$ to $E_{\alpha}$. Assume that the $f_{\alpha,\beta}$ satisfy the following conditions:
\begin{enumerate}
\item[] $(LP_{I})\ \alpha\leq\beta\leq\gamma\Rightarrow f_{\alpha,\gamma}=f_{\alpha,\beta}\circ f_{\beta,\gamma}$.
\item[] $(LP_{II})$ For all $\alpha\in I,\,f_{\alpha,\alpha}$ is the identity on $E_{\alpha}$.
\end{enumerate}
Then $((E_{\alpha})_{\alpha\in I},\,(f_{\alpha,\beta})_{(\alpha,\beta)\in I\times I})$ is called a \emph{projective system relative to the preorder}\index{Subject}{projective system} $(I,\,\leq)$. Let $E=\prod_{\alpha\in I}E_{\alpha}$ and for each pair $(\alpha,\,\beta)\in I\times I$ such that $\alpha\leq\beta$, let $E_{\alpha,\beta}=\{x\in E:f_{\alpha,\beta}(x_{\beta})=x_{\alpha}\}$. Then the set $Y=\bigcap_{\alpha<\beta}E_{\alpha,\beta}$ is called the \emph{projective limit}\index{Subject}{projective limit} of the projective system $((E_{\alpha})_{\alpha\in I},\,(f_{\alpha,\beta})_{(\alpha,\beta)\in I\times I})$. $Y$ is denoted $\lim_{\leftarrow}(E_{\alpha},\,f_{\alpha,\beta})$. In case the preorder $(I,\,\leq)$ is directed to the right (that is, $\forall(\alpha,\,\beta)\in I^{2},\,\exists\gamma\in I$ such that $\alpha\leq\gamma$ and $\beta\leq\gamma$.) we say that the projective system $((E_{\alpha})_{\alpha\in I},\, (f_{\alpha,\beta})_{(\alpha,\beta)\in I\times I})$ is directed.
\item[35.] If $((X_{\alpha})_{\alpha\in I},\,(f_{\alpha,\beta})_{(\alpha,\beta)\in I\times I})$ is a projective system relative to the preorder $(I, \leq)$ and if each $X_{\alpha}$ has a topology $C_{\alpha}$ such that $f_{\alpha,\beta}$ is continuous relative to the topologies $C_{\alpha}$ and $C_{\beta}$ then the system $((X_{\alpha})_{\alpha\in I},\,(f_{\alpha,\beta})_{(\alpha,\beta)\in I\times I})$ is called $a$ \emph{projective system of topological spaces}. In this case the projective limit of the system with the topology induced by the product topology on $\prod_{\alpha\in I}X_{\alpha}$ is called the \emph{the projective limit of the projective system of topological spaces}.
\item[36.] If $(A,\,\wedge,\,\vee)$ is a bounded, distributive lattice, $X$ is a subset of $A$ and $F$ is a filter in $A$, then $F$ \emph{respects} $X$ if $X\subseteq F$ and $\bigwedge X$ exists implies $\bigwedge X\in F$.
\item[37.] If $(A,\,\wedge,\,\vee)$ is a bounded, distributive lattice, $X$ is a subset of $A$ and $F$ is a filter in $A$, then $F$ \emph{strongly respects} $X$ if $X\subseteq F$ implies $\exists y\in F$ such that $\forall x\in X,\,y\leq x$.
\item[38.] If $X$ is a subset of a bounded, distributive lattice $(A,\,\wedge,\,\vee)$ then $X$ \emph{is of Boolean type} if for all $b\in A,\,[(\forall x\in X,\,b\vee x=1)$ and $(\forall y\leq X,\,y\leq b)]$ implies $b=1$. $(y\leq X$ means $(\forall x\in X)(y\leq x).)$
\item[39.] If $(A,\,\wedge,\,\vee)$ is a bounded, distributive lattice we say that \emph{A has the Rasiowa-Sikorski property}\index{Subject}{Rasiowa-Sikorski property} (respectively \emph{the strong Rasiowa-Sikorski property}) if for all $a,\,b\in A$ such that $a\not\leq b$ and all sequences $(X_{n})_{n\in\omega}$ of subsets of $A$ of Boolean type there is a prime filter containing $a$, not containing $b$ and respecting (respectively strongly respecting) each $X_{n}$.
\item[40.] If $(A,\,\wedge,\,\vee)$ is a bounded distributive lattice then an element $a\in A$ is $\vee$-\emph{compact} if for every subset $I\subseteq A,\,\bigvee I=a$ implies $(\exists J\subseteq I)$ ($J$ finite and $\bigvee J=a$).
\end{enumerate}

\section*{NOTE 72}

We prove that form $207(\beta)$ implies 209 for any ordinal $\beta$. (form $207(\beta)$ is $UT(\aleph_{\beta},\,\aleph_{\beta},\,<2^{\aleph_{\beta}})$ and 209 is ($\exists$ an ordinal $\alpha$) $(\forall X)$ (If $X$ is a denumerable union of denumerable sets then $\mathcal{P}(X)$ cannot be partitioned into $\aleph_{\alpha}$ non-empty sets.)) Assume form $207(\beta)$ is true. Then for any set $X$, if $X$ is a denumerable union of denumerable sets, then $|X|<2^{\aleph_{\beta}}$. Choose an ordinal $\alpha$ so that $\neg(\aleph_{\alpha}\leq |\mathcal{P}(\mathcal{P}(2^{\aleph_{\beta}}))|)$. For any $X$ which is the denumerable union of denumerable sets we have $|\mathcal{P}(\mathcal{P}(X))|\leq|\mathcal{P}(\mathcal{P}(2^{\aleph_{\beta}}))|$.
A partition of $\mathcal{P}(X)$ into $\aleph_{\alpha}$ non-empty sets would give $\aleph_{\alpha}\leq|\mathcal{P}(\mathcal{P}(X))|$.

Similarly, form 26 implies form 209 and form 31 implies form 209.

\section*{NOTE 73}

In Pincus\index{Author}{Pincus} \cite{bib:1977a} a general method for adding Dependent Choice (43) to models of ZF is described and the following models of ZF are constructed with the properties indicated:

\textsc{Model A} ($\mathcal{M}43$). Form 62 $(C(\infty,\,<\aleph_{0}))$ and form 43 (DC) are true and form 30 (OP) is false.

\textsc{Model B} ($\mathcal{M}44$). Forms 30 and 43 are true and form 15 (Kinna-Wagner Principle) is false.

\textsc{Model C} $(\mathcal{M}45)$. Forms 30 and 43 are true and form 49 (the order extension principle) is false.

\textsc{Model D} $(\mathcal{M}46(n,\,M))$. For each finite set $Z$ of positive natural numbers and natural number $n>1$, for which the following is false
\begin{enumerate}
\item[]$S(Z,\,n)$: For every decomposition $n=p_{1}+\ldots+p_{s}$ of $n$ into a sum of primes at least one $p_{i}$ divides and element of $Z$, (See Mostowski\index{Author}{Mostowski} \cite{bib:1945}. Also see note 15. Condition $S$ implies condition $D$ of note 15.)
\end{enumerate}

there is a model in which form 46 $(C(\infty,\,Z))$ and form 43 are true and form 45 $(C(\infty,\,n))$ is false.

\textsc{Model E} $(\mathcal{M}47(n,\,M))$. For every finite set $Z$ of positive integers and integer $n>1$, for which the following condition is false
\begin{enumerate}
\item[] $M(Z,\,n)$: For any decomposition $n=p_{1}+\ldots+p_{s}$ of $n$ into a sum of primes there are non-negative $q_{1},\,\ldots$, $q_{s}$ such that $q_{1}p_{1}+\ldots+q_{s}p_{s}\in Z$.
\end{enumerate}

There is a model in which form 46 $(C(\infty,\,Z))$ and form 43 are true and form 47 $(C(WO,\,n))$ is false. (See Mostowski \cite{bib:1945}. Also, condition $K$ of note 15 implies condition $M$.)

\section*{NOTE 74}

In Pincus\index{Author}{Pincus} \cite{bib:1977a} it is shown that $C(\infty,\,<\aleph_{\alpha})\rightarrow C(\infty,\,\aleph_{\alpha})$ is not a theorem of ZF for regular $\aleph_{\alpha}$. For singular $\aleph_{\alpha}$ the question is open.

\section*{NOTE 75}

Several weak\index{Author}{Felgner/Jech}\index{Author}{Mostowski} forms of AC are known to be equivalent to AC in ZF\index{Subject}{forms equivalent to AC in ZF}. These include:
\begin{enumerate}
\item[1.] MC: The Axiom of Multiple Choice, form 67;
\item[2.] LW: Every linearly ordered set can be well ordered, form 90;
\item[3.] PW: The power set of a well ordered set can be well ordered, form 91;
\item[4.] VB: Every vector space has a basis, form 66;
\item[5.] A: The Antichain Principle, form 89; and
\item[6.] $C(LO,\,\infty)$: Every linearly ordered family of non-empty sets has a choice function, form 202.
\end{enumerate}
Consequently everything that implies one of these six forms is equivalent to AC in ZF. That is, 95, 109, 112, 114, 133, 147, 149, 164, 218, 264, 292, 333, 334, and $335(n)$ for $n\geq 2$, also imply AC in ZF. In fact in $\mathrm{ZF}^{0}$ the following two chains of implications are provable: AC $\rightarrow$ VB $\rightarrow$ MC $\rightarrow$ A $\rightarrow$ LW $\rightarrow$ PW, and AC $\rightarrow C(LO,\,\infty)\rightarrow$ LW. With regard to the reverse implications, VB $\rightarrow$ AC, MC $\rightarrow$ VB and LW $\rightarrow C(LO,\,\infty)$ are open questions. The other reverse implications are known to be unprovable in $\mathrm{ZF}^{0}$. (See Jech\index{Author}{Jech} \cite{bib:1973b}, page 134, Felgner/Jech
\cite{bib:1973}, Blass\index{Author}{Blass} \cite{bib:1984a}, H. Rubin\index{Author}{H. Rubin} \cite{bib:1960}, H. Rubin/J. Rubin\index{Author}{H. Rubin/ J. Rubin} \cite{bib:1985}, Part I, 5, Brunner\index{Author}{Brunner} \cite{bib:1983d},\cite{bib:1984b},\cite{bib:1985a}, Howard/Rubin\index{Author}{Howard/Rubin} \cite{bib:1977}, Truss\index{Author}{Truss} \cite{bib:1978}, Keremedis\index{Author}{Keremedis} \cite{bib:1996a},\cite{bib:1996c}, and note 18.) It follows from the two chains of implications that each of the twenty forms mentioned above implies PW, form 91, in $\mathrm{ZF}^{0}$.

In Pincus \cite{bib:1972a}, page 740, it is shown that if $\neg\Phi$ is boundable (see Pincus \cite{bib:1972a}, p. 722 for the definition of boundable) and $\Phi\rightarrow AC$ is a theorem of ZF, then $\Phi\rightarrow$ form 91 (PW) is a theorem of $\mathrm{ZF}^{0}$.

\section*{NOTE 76}

In this note we prove that form 126 $(MC(\aleph_{0},\,\infty))$ implies forms 82 (If a set is infinite then its power set is Dedekind infinite.) and 185 (Every Dedekind finite, linearly ordered set is finite.)

\begin{lemma*}
$MC(\aleph_{0},\,\infty)$ implies that for every infinite set $X$ there is a sequence $(Y_{n})_{n\in\omega}$ of finite subsets of $X$ such that $\forall n\in \omega,\,Y_{n}\subsetneq Y_{n+1}$.
\end{lemma*}

\begin{proof}
Let $X$ be an infinite set and for each $n\in\omega$, let $X_{n}=\{Y\subseteq X:|Y|=n\}$. By $MC(\aleph_{0},\,\infty)$ there is a function $f$ such that $\forall n\in\omega,\,f(X_{n})$ is a finite non-empty subset of $X_{n}$. We define the sequence $(Y_{n})_{n\in\omega}$ by induction as follows: $Y_{0}=\emptyset$. Assume that $Y_{n}$ is defined and let $k$ be the least natural number such that $|Y_{n}|<k$. We let $Y_{n+1}=(\bigcup f(X_{k}))\cup Y_{n}$. Since $\mathrm{f}(X_{k})$ is a set of $k$ element sets $|Y_{n}|<k\leq |\bigcup f(X_{k})|\leq|Y_{n+1}|$.
\end{proof}

It is clear from the lemma that $MC(\aleph_{0},\,\infty)$ implies that the power set of an infinite set is Dedekind infinite. To show that $MC(\aleph_{0},\,\infty)$ implies form 185, let $X$ be a Dedekind finite set which is linearly ordered by $\leq$. The assumption that $X$ is infinite leads to a contradiction as follows: By the lemma we get a sequence $(Y_{n})_{n\in\omega}$ of finite subsets of $X$ such that $Y_{n}\subsetneq Y_{n+1}$ for all $n\in \omega$. If we define $t_{n}$ to be the $\leq$-least element of $Y_{n+1}-Y_{n}$ then $(t_{n})_{n\in \omega}$ is an infinite sequence of distinct elements of $X$ and therefore $X$ is Dedekind infinite.

\section*{NOTE 77}

In this note we include definitions from Keremedis\index{Author}{Keremedis} \cite{bib:1996b} for forms [1 BL] through [1 BR], [1CB], and [67 G].

In what follows $(X,\,T)$ is a topological space and $P\subset \mathcal{P}(X)$.

\begin{definition*}
A non-empty set $\mathcal{F}\subseteq P-\{\emptyset\}$ is called a $P$-\emph{filter} iff
\begin{enumerate}
\item[(1)] $F_{1},\,F_{2}\in \mathcal{F}$ implies $F_{1}\cap F_{2}\in \mathcal{F}$.
\item[(2)] $F\in \mathcal{F},\,F\subset F^{\prime}$ and $F^{\prime}\in P$ implies $F^{\prime}\in \mathcal{F}$.
\end{enumerate}
A $T$-filter is called an \emph{open filter}. $\mathcal{F}$ is called a \emph{filter base} if it just satisfies (1) above.
\end{definition*}

\begin{definition*}
A non-empty set $\mathcal{I}\subseteq P-\{X\}$ is called a $P$-\emph{ideal} iff
\begin{enumerate}
\item[(1)] $I_{1},\,I_{2}\in \mathcal{I}$ implies $I_{1}\cup I_{2}\in \mathcal{I}$.
\item[(2)] $I\in \mathcal{I},\,I^{\prime}\subset I$ and $I^{\prime}\in P$ implies $I^{\prime}\in \mathcal{I}$.
\end{enumerate}
If $F$ is the set of closed subsets of $X$, then a $F$-ideal is called a \emph{closed ideal}. $\mathcal{I}$ is called \emph{complete} if $\bigcup C\in \mathcal{I}$ for every $C\subseteq \mathcal{I}$.
\end{definition*}

\begin{definition*}
A disjoint set $C\subseteq T-\{\emptyset\}$ is called a \emph{cellular family
of} $X$.
\end{definition*}

\begin{definition*}
A \emph{nest} of $(X,\,T)$ is a maximal linearly ordered subset of the partially ordered set $(T-\{\emptyset\},\,\subseteq)$.
\end{definition*}

\begin{definition*}
A \emph{tower} $\mathcal{T}$ of $(X,\,T)$ is a linearly ordered subset of the partially ordered set $(T-\{\emptyset\},\,\subseteq)$ such that either the interior of $\bigcap \mathcal{T}$ is empty or the interior of $\bigcap \mathcal{T}$ does not properly contain any non-empty open set.
\end{definition*}

\begin{definition*}
If $\{(X_{i},\,T_{i}):i\in k\}$ is a family of topological spaces, then the \emph{box topology} on $\prod_{i\in k}X_{i}$ is the topology having a base of all sets of the form $\prod_{i\in k}O_{i}$ where $O_{i}\in T_{i}$ for all $i\in k$.
\end{definition*}

\begin{definition*}
If $\{(X_{i},\,T_{i}):i\in k\}$ is a family of disjoint topological spaces, then the \emph{free union} is the set $X=\bigcup\{X_{i}:i\in k\}$ together with the topology $T$ containing all sets $O\subseteq X$ such that $O\cap X_{i}\in T_{i}$ for all $i\in k$.
\end{definition*}

\section*{NOTE 78}

In Rav\index{Author}{Rav} \cite{bib:1977}, it is shown that under the assumption of $RL$ (form 99) and certain conditions on a set $X$ and filter $\mathcal{F}$ in $\mathcal{P}(X)$, $\mathcal{F}$ can be extended to an ultrafilter. Either of the following conditions suffices:
\begin{enumerate}
\item[1.] $\{A:A\subseteq \mathcal{P}(X)\wedge A$ finite $\wedge\forall x,\,y\in A(x\neq y\rightarrow x\cap y =\emptyset)\}$ has a choice function.
\item[2.] $\mathcal{F}$ has a choice function.
\end{enumerate}

\section*{NOTE 79}

Definitions for forms 227 and 228.

\begin{definition*}\
\begin{enumerate}
\item[1.] A \emph{partially ordered group} is a triple $(G,\,\circ,\,\leq)$ such that $(G,\,\circ)$ is a group and $(G,\,\leq)$ is a partially ordered set satisfying for all $x,\,y,\,a$ and $b$ in $G$ if $x\leq y$, then $axb\leq ayb$.
\item[2.] $P=\{a\in G:a\geq e\}$ is the \emph{positive cone}\index{Subject}{positive cone} (where $e$ is the identity in $G$).
\item[3.] If $P\cup P^{-1}=G$\index{Subject}{partially ordered group} then $G$ is \emph{fully ordered}\index{Subject}{fully ordered group}. (This is equivalent to saying that $\leq$ is a linear ordering.)
\item[4.] A group is \emph{torsion free}\index{Subject}{torsion free group} if the only element of finite order is the identity element.
\end{enumerate}
\end{definition*}

\section*{NOTE 80}

In Rav\index{Author}{Rav} \cite{bib:1977} forms [14 AL], [14 AM], [14 AO] and [14 AP] are shown to follow from form 14 (BPI). The arguments that they imply form 14 are fairly straightforward. ([14 AL]: The radical of a proper ideal $A$ in a commutative ring $R$ with identity (The radical of $A=\{x\in R$: for some $m\in \omega,\,x^{m}\in A\}.)$ is the intersection of all prime ideals in $R$ containing A; [14 AM]: Let $A$ be a subring of a commutative ring $R$ and $p$ a prime ideal in $A$ such that $p=Rp\cap A$. Then there is a prime ideal $J$ in $R$ such that $p=J\cap A;$ [14 AO]: Suppose $R$ is a commutative ring, $A$ is a proper ideal in $R$, and $S$ is a multiplicative semigroup in $R$ not meeting $A$. Then there is a prime ideal $p$ in $R$ such that $A\subseteq p$ and $p\cap S=\emptyset$; and [14 AP]: Suppose to each finite subset $F$ of a set I there corresponds a set $\Phi(F)$ of functions whose domains are subsets of $I$ including $F$ and such that (a) $F_{1}\subseteq F_{2}$ implies $\Phi(F_{2})\subseteq\Phi(F_{2})$ and (b) $\forall i\in I,\,\{\phi(i):\phi \in\bigcup$\{ $\Phi(F):F$ finite and $F\subseteq I\}\}$ is finite. Then there is a function $f$ with domain $I$ such that for all finite $F\subseteq I,\,\exists\phi\in\Phi(F)$ such that $\phi$ and $f$ coincide on F.)

\begin{theorem*}\quad
[14 AL] implies BPI.
\end{theorem*}

\begin{proof}
Suppose $B$ is a Boolean algebra with more than two elements. Choose $b\in B$ with $b\neq 0$ and $b\neq 1$. Let $A$ be the ideal generated by $b$ and apply [14 AL] to get a prime ideal in $B$.
\end{proof}

\begin{theorem*}\quad

[14 AM] implies BPI.
\end{theorem*}

\begin{proof}
Let $B$ be a Boolean algebra and let $r$ be a proper ideal in B. Let $A=\{x:x\in r$ or $\overline{x}\in r\}$. $A,\,B$ and $r$ satisfy the hypotheses of [14 AM]. The conclusion of [14 AM] gives a prime ideal $J$ in $B$ containing $r$.
\end{proof}

\begin{theorem*}\quad

[14 AO] implies BPI.
\end{theorem*}

\begin{proof}
If $B$ is a Boolean algebra and $A$ is a proper ideal in $B$ let $S=\{1\}\subseteq B$. $S$ is a multiplicative sub-semi-group of $B$ and therefore the hypotheses of [14 AO] are satisfied. The conclusion of [14 AO] gives a prime ideal in $B$.
\end{proof}

\begin{theorem*}\quad

[14 AP] implies BPI.
\end{theorem*}

\begin{proof}
Since it is shown in Cowen \cite{bib:1973} (see Rav \cite{bib:1977}, Corollary~ 2.3) that [14 AP] implies Rado's lemma (form 99) and it is known that Rado's lemma together with $C(\infty,\,<\aleph_{0})$ (form 62) imply BPI (see form [14
Y]), it suffices to prove that [14 AP] implies $C(\infty,\,<\aleph_{0})$. Let $\{x_{i}:i\in I\}$ be a family of pairwise disjoint finite sets. For each finite subset $F$ of $I$, let
\begin{equation*}
\Phi(F)=\{\phi: \mathrm{dom}\ \phi\ \text{is finite}\ \wedge F\subseteq \text{dom}\ \phi \subseteq I\wedge(\forall i\in \text{dom}\ \phi)(\phi(i)\in x_{i})\}
\end{equation*}
$\Phi$ and $I$ satisfy the hypotheses of [14 AP] and therefore by [14 AP] there is a function $f$ with domain $I$ such that $\forall$ finite $F\subseteq I,\,\exists\phi$ in $\Phi(F)$ such that $\phi$ and $f$ agree on $F$. Therefore, for all $i\in I,\,f(i)\in x_{i}$ and we have a choice function for $\{x_{i}:i\in I\}$.
\end{proof}

\section*{NOTE 81}

We argue that form 30 (the ordering principle) implies form 83 (If every subset of $\mathcal{P}(X)$ which is linearly ordered by $\subseteq$ has a maximum element then $X$ is finite.)

Assume $X$\index{Author}{Cowen}\index{Author}{Rav} is infinite and let $\leq$ be a linear ordering of $X$. There must be either a non-empty subset $Y$ of $X$ with no least element (under $\leq$) or a non-empty subset with no greatest element. (If every non-empty subset of $X$ has a least element then let $x_{1}$ be the least element of $X,\,x_{2}$ be the least element of $X-\{x_{1}\}$, etc. This produces a subset $\{x_{1},\,x_{2},\,\ldots\}$ with no greatest element.) If $Y$ is an infinite subset of $X$ with no greatest element, we get a linearly ordered subset $A$ of $\mathcal{P}(X)$ with no maximum by letting $A=\{\{x:x\in Y\wedge x\leq y\}:y\in Y\}$. We can proceed similarly if $Y$ is an infinite subset of $X$ with no smallest element.

\section*{NOTE 82}

We give the definitions for form 230 $(L^{1}=HOD)$\index{Subject}{$L^{1}$} from Myhill/Scott\index{Author}{Myhill/Scott} \cite{bib:1971}.

\begin{definition*}\
\begin{enumerate}
\item[1.] $DF^{0}(A)$ is the set of all subsets of $A$ definable by first order formulas in $\langle A, \in_{A}\rangle$ using parameters from $A$, i.e., $DF^{0}(A)$ is all sets of the form $\{x\in A:\Phi^{(A)} (x,\,a_{0},\,\ldots a_{n-1})\}$ where $\Phi$ is a formula of the type described above, $a_{0},\,\ldots,\, a_{n-1}$ are in $A$ and $\Phi^{(A)}$ means all variables range over $A$.
\item[2.] G\"{o}del's constructible sets are defined by $L_{\alpha}^{0}=\bigcup_{\beta<\alpha}DF^{0}(L_{\beta}^{0})$ and $L^{0}= \bigcup_{\alpha\in \mathrm{On}}L_{\alpha}^{0}$.
\item[3.] $DF^{1}(A)$ consists of all sets of the form $\{x\in A:\Phi(x,\,a_{0},\,\ldots,\,a_{n-1})\}$ where $a_{0},\,\ldots,\,a_{n-1}$ are in $A$ and the quantifiers are restricted to $A$ or $\mathcal{P}(A)$.
\item[4.] $L_{\alpha}^{1}=\bigcup_{\beta<\alpha}DF^{1}(L_{\beta}^{1})$ and $L^{1}=\bigcup_{\alpha\in \mathrm{On}}L_{\alpha}^{1}$.
\item[5.] $HOD$\index{Subject}{$HOD$} is the class of hereditarily ordinal definable sets.
\end{enumerate}
\end{definition*}

\section*{NOTE 83}

In Hickin/Plotkin\index{Author}{Hickin/Plotkin} \cite{bib:1976} it is shown that form 14 (BPI) implies [14 AR]. Since [14 AQ] follows from [14 AR] we get the equivalence of [14 AR] and 14.

\section*{NOTE 84}

This note contains (standard) algebraic definitions for forms 241 through 244.

\begin{definition*}\
\begin{enumerate}
\item[(1)] If $F$ is a subfield of $E,\,\alpha\in E$ is \emph{algebraic over} $F$ if $\exists a_{0},\,a_{1},\, \ldots,\,a_{n}\in F$ not all zero such that $a_{0}+a_{1}\alpha\ldots+a_{n}\alpha^{n}=0$.
\item[(2)] $E$ is \emph{algebraic over} $F$ if every element of $E$ is algebraic over $F$.
\item[(3)] A field $L$ is \emph{algebraically closed} if every polynomial of degree $\geq 1$ with coefficients in $L$ has a root in $L$.
\item[(4)] If $K$ is a field, an \emph{algebraic closure} of $K$ is an algebraically closed algebraic extension of $K$.
\item[(5)] A field $K$ is \emph{real} if --1 is not the sum of squares in $K$.
\item[(6)] A field is \emph{real closed} if it is real and any real, algebraic extension of $K$ equals $K$.
\item[(7)] A \emph{principal ideal domain} is an integral domain (commutative ring with identity and no zero divisors) in which every ideal is principal.
\item[(8)] A \emph{unique factorization domain}\index{Subject}{unique factorization domain} is an integral domain in which
\begin{enumerate}
\item[(i)] every non-zero,\index{Subject}{principal ideal domain}\index{Subject}{real closed field} non-unit element (unit $\equiv$ invertible) can be written $a= c_{1}c_{2}\ldots c_{n}$ with $c_{1},\,c_{2},\,\ldots,\,c_{n}$ irreducible. (That is, each $c_{i}$ is non-zero and non-unit and if $c_{i}=ab$ then either $a$ or $b$ is a unit.)
\item[(ii)] If\index{Subject}{algebraic closure}\index{Subject}{algebraic over a field} $a=c_{1}c_{2}\ldots c_{n}=d_{1}d_{2}\ldots d_{m}$ where $c_{i}$ and $d_{i}$ are irreducible then $n=m$ and there is a permutation $\sigma$ of $\{$1, 2, $\ldots,\,n\}$ such that $c_{i}=r_{i}d_{\sigma(i)}$ where $r_{i}$ is a unit, $i=1,\,\ldots,\,n$.
\end{enumerate}
\end{enumerate}
\end{definition*}

\section*{NOTE 85}

The monadic theory $MT(\omega_{1},\,<)$ mentioned in form 246 is defined as follows: The language has a countable set $x,\,y,\,\ldots$ of variables ranging over $\omega_{1}$, a countable set of variables $X,\,Y,\, \ldots$ ranging over subsets of $\omega_{1}$, a relation $<$ on $\omega_{1}$ and a relation $\in$ between individual variables and set variables. The formulas are built up with the usual logical connectives and quantifiers $\forall$ and $\exists$ on both types of variables. $MT(\omega_{1},\,<)$ is the set of all sentences true in $(\omega_{1},\,<)$

\section*{NOTE 86}

In Plotkin\index{Author}{Plotkin} \cite{bib:1976}, it is shown that several properties of Boolean algebras provable in ZFC are not provable in ZF alone. Forms 247 and 248 are examples.

\section*{NOTE 87}

In Truss\index{Author}{Truss} \cite{bib:1975} the following weakenings of K\"{o}nig's lemma are considered: For any subset $Z\subseteq\omega,\,KL(Z)$ is the formula ``If $T$ is an infinite tree such that for every element $a$ of $T$, the number of immediate successors of $a$ is in $Z$, then $T$ has an infinite branch.'' Form 249 is $KL(\{2\})$. Here is a summary of the results:
\begin{enumerate}
\item[1.] Form 250 $((\forall n\in \omega)(C(WO,\,n))\rightarrow KL(\{k\})$ is not a theorem of $\mathrm{ZF}^{0}$ for any $k>1$.
\item[2.] Form 80 $(C(\aleph_{0},2))\rightarrow KL(\{2\})$ is not a theorem of $\mathrm{ZF}^{0}$.
\item[3.] The following conditions are equivalent:
\begin{enumerate}
\item[(i)] $\mathrm{ZF}^{0}\vdash KL(Z)\rightarrow KL(\{n\})$
\item[(ii)] $\mathrm{ZF}^{0}\vdash KL(Z)\rightarrow C(\aleph_{0},\,n)$
\item[(iii)] $\mathrm{ZF}^{0}\vdash(\forall m\in Z)(KL(\{m\})))\rightarrow C(\aleph_{0},\,n)$
\item[(iv)] Gauntt's condition $L(n,\,Z)$, see note 15.
\end{enumerate}
\end{enumerate}

\section*{NOTE 88}

Either $\mathcal{N}1$ or $\mathcal{N}3$ may be used to prove that form 31 $(UT(\aleph^{0},\,\aleph^{0},\,\aleph^{0}))$ does not imply form 9 (Dedekind finite $=$ finite). Form 9 is known to be false in these models and a proof that form 31 holds is given below. The result is transferable to ZF by Pincus \cite{bib:1972a}.

Let $\mathcal{N}$ be either $\mathcal{N}1$ or $\mathcal{N}3$ with set of atoms $A$ and $\mathcal{G}$ the group of permutations used in defining the model. Assume that in $\mathcal{N},\,B$ is a denumerable set of denumerable sets. Let $e\subseteq A$ be a finite support for every $b\in B$. Such an $e$ exists since $B$ is well orderable in $\mathcal{N}$. Assume $b\in B$ and $x\in b$. The proof will be completed by showing that $e$ is a support of $x$. We argue by contradiction. Assume that $e$ is not a support of $x$. We will show that under this assumption $b$ has a subset of the same cardinality as an infinite subset of $A$ which contradicts the countability of $b$.

Let $e^{\prime}\supseteq e$ be a support of $x$. Then there is a $t\in e^{\prime}\backslash e$ and a $\phi\in \mathcal{G}$ which fixes $e^{\prime}\backslash \{t\}$ pointwise and such that $\phi(x)\neq x$. (It follows that $\phi(t)\neq t$.) It can be verified that
\begin{equation*}
\{\langle\psi(t),\psi(x)\rangle:\psi\in \mathcal{G}\ \text{and}\ \psi\ \text{fixes}\ e^{\prime}\backslash \{t\}\ \text{pointwise}\}
\end{equation*}
is a\index{Author}{Pincus} one to one function from an infinite subset of $A$ into $b$ which is in $\mathcal{N}$.

\section*{NOTE 89}

Definitions for form 267 and [0 Y].

\begin{definition*}
A \emph{free complete Boolean algebra over a set}\index{Subject}{free complete Boolean algebra} $D$ is a complete Boolean algebra $B$ together with a mapping $f:D\rightarrow B$ such that range$(f)$ generates $B$ as a complete Boolean algebra and for every complete Boolean algebra $C$ and function $g:D\rightarrow C$, there is a complete homomorphism $h:B\rightarrow C$ such that $g=h\circ f$
\end{definition*}

\section*{NOTE 90}

Monro\index{Author}{Monro} \cite{bib:1972} shows that for all $n\in \omega,\,\neg(\mathrm{ZF}\vdash K(n+1)\rightarrow K(n))$. $K(n)$ is form $81(n)$: For every set $S$ there is an ordinal $\alpha$ and a one to one function $f:S\rightarrow \mathcal{P}^{n}(\alpha)$. $(\mathcal{P}^{0}(X)=X$ and $\mathcal{P}^{n+1}(X)=\mathcal{P}(\mathcal{P}^{n}(X))$. $K(0)$ is form [1 E] and $K(1)$ is form [15 A]. Monro \cite{bib:1972} also says that in an unpublished model due to H. Stewart\index{Author}{Stewart} $K(2)$ holds, form 85 $(C(\infty,\, \aleph_{0}))$ fails and form 64 (there is no amorphous set) fails. ($A$ is amorphous if $A$ is infinite and is not the union of two non-empty, disjoint infinite sets.)

\section*{NOTE 91}

In this note we argue that form 269 is false in the first Fraenkel model $\mathcal{N}1$ of ZF$^{0}$. Form 269 is: ``For every cardinal $m$, there is a set $A$ such that $2^{|A|^{2}}\geq m$ and there is a choice function on the collection of 2 element subsets of $A$.'' It suffices to show that for any set $A$ in $\mathcal{N}1$ if the set $[A]^{2}$ of two element subset of $A$ has a choice function, then $A$ is well orderable in $\mathcal{N}1$. Assume that $A$ is such a set and assume that $E$ is a (finite) support of $A$ and a choice function $f$ on $[A]^{2}$. Assume $x\in A$. We show by contradiction that $E$ must be a support of $x$. Suppose not and let $F$ be a support of $x$ containing $E$. We claim there is a permutation $\psi$ in $\mathrm{fix}_{G}(E)$ and an element $y\in A$ such that $x\neq y,\,\psi(x)=y$ and $\psi(y)=x$. This would contradict our choice of $E$ as a support for a choice function on $[A]^{2}$ since $\psi$ fixes $\{x,\,y\}$ but moves both of its elements. To get such a $\psi$, let $\phi\in \mathrm{fix}_{G}(E)$ satisfy $\phi(x)\neq x$ and $\phi$ moves only finitely many atoms. Assume $F-E=\{t_{1},\,\ldots,\,t_{n}\}$ and choose a set of atoms $\{s_{1},\,\ldots,\, s_{n}\}$ disjoint from $F\cup\{t\in U:\phi(t)\neq t\}$. Let $\psi$ be the product of transpositions $\prod_{i=1}^{n}(t_{i},\,s_{i})$. The element $\psi(x)$ is different from $x$ since $\phi(x)\neq x$ and $\phi(\psi(x))=\psi(x)$. It is also clear that if $y=\psi(x)$, then $\psi(y)=x$.

\section*{NOTE 92}

In Kolany/Wojtylak \cite{bib:1991} some relationships between forms 271$(m)$ and 45$(n)$ are given. $(271(n)$ is $CT_{n}$: The for the propositional calculus restricted to sets of formulas where each variable occurs in at most $n$ formulas, and $45(n)$ is $C(\infty,\,n).)$ For example, $CT_{m}$ implies $C(\infty,\,n)$ if $m\geq n!+n(n-1)$. Also considered are forms $CT_{m}^{\ast}\
(CT_{m}$ restricted to formulas which are disjunctions of formulas of the form $p$ or $(\neg p)$ where $p$ is a propositional variable). For example, it is shown that $CT_{3}^{\ast}$ implies $C(\infty,\,2)$.

\section*{NOTE 93}

We give\index{Subject}{compactness theorem}\index{Author}{Kolany/Wojtylak} proofs due to A. Rubin\index{Author}{A. Rubin} that $MC(\infty,\,\infty$, even) (334) is true in $\mathcal{N}2$ and $MC(\infty,\, \infty$, odd) (333) is true in $\mathcal{N}2^{\ast}(3)$.

\setcounter{theorem}{0}
\begin{theorem}\label{part04:thm1b}
$MC(\infty,\,\infty,\,even)$ is true in $\mathcal{N}2$.
\end{theorem}

\begin{proof}
Let $M$ be a non-empty set in $\mathcal{N}2$, each of whose elements has at least two elements. The group $\mathcal{G}$ in $\mathcal{N}2$ is Abelian and each of its elements has order two. Let $\mathcal{F}$ be the filter determined by $S$, the set of finite subsets of the set of atoms, $A$. $(\mathcal{F}$ is generated by the set of all subgroups $H$ of $\mathcal{G}$ such that there is an $E\in S$ such that $H$ fixes $E$ pointwise.) The filter $\mathcal{F}$ is in the model because each element in $\mathcal{G}$ has order two so $\mathcal{F}$ has null support. In fact,
\begin{equation}
\label{part04:eqn1} (\forall H\in \mathcal{F})(\forall\sigma\in H)\sigma(H)=H.
\end{equation}
Define a linear ordering $\leq$ on $\mathcal{F}$ which extends $\supseteq$, that is, if $H_{1},\,H_{2}\in \mathcal{F}$ and $H_{1}\supseteq H_{2}$, then $H_{1}\leq H_{2}$. The relation $\leq$ is in the model because of (\ref{part04:eqn1}). (In fact, $\leq$ is a well ordering, but we do not need this.) Let $W$ be the set of all sets in the model that have null support. Every permutation in $\mathcal{G}$ leaves $W$ pointwise fixed so $W$ can be well ordered by a relation $R$.

Suppose $X\in M$. We shall give an algorithm to construct a subset of $X$ whose cardinality is a power of two. Let $Y=\{\mathcal{G}^{\prime\prime}\{x\}:x\in X\}$. Each element of $Y$ has null support so $Y$ can be well ordered by $R$. Moreover, each element of $Y$ is finite. (In fact, the cardinality of each element of $Y$ is a power of two because $|\mathcal{G}^{\prime\prime}\{x\}|=|\mathcal{G}/\mathrm{sym}_{\mathcal{G}}(x)|$, where $\mathrm{sym}_{\mathcal{G}}(x)=\{\sigma\in \mathcal{G}:\sigma(x)=x\}.)$ Let $y$ be the $R$-first element of $Y$ and let $x=y\cap X$. Let $\mathrm{fix}_{\mathcal{G}}(x)$ be the group of all $\sigma \in \mathcal{G}$ that fixes $x$ pointwise and let
\begin{equation*}
J=\{H\in \mathcal{F}:\mathrm{fix}_{\mathcal{G}} (x) \subseteq H\ \&\ (\exists z\in x)H^{\prime\prime}\{z\}\subseteq x\}
\end{equation*}
Since $\mathrm{fix}_{\mathcal{G}}(x)\in J,\, J\neq\emptyset$. Moreover, $J$ is finite and is linearly ordered by $\leq$. Let $H_{0}$ be the $\leq$-first element of $J$. We claim that if $z,\,w\in x$ and, $H_{0}^{\prime\prime}\{z\}$ and $H_{0}^{\prime\prime}\{w\}$ are both subsets of $x$, then $H^{\prime\prime}_{0}\{z\}=H^{\prime\prime}_{0}\{w\}$. The reason for this is as follows. By the definition of $y(x\subseteq y)$, there is a $\sigma \in \mathcal{G}$ such that $\sigma(z)=w$. Let $K=H_{0}\cup
H_{0}\sigma$. Then $K^{\prime\prime}\{z\}=H_{0}^{\prime\prime}\{z\}\cup H^{\prime\prime}_{0}\{w\}\subseteq x$. We have $\mathrm{fix}_{\mathcal{G}} (x)\subseteq H_{0}\subseteq K$ so $K\in J$, and $K\leq H_{0}$. However, $H_{0}$ is the $\leq$-first element of $J$, so $K=H_{0}$. Consequently, $H^{\prime\prime}_{0}\{z\}=H_{0}^{\prime\prime}\{w\}$.

Then, for each $X\in M$, we choose the unique $H_{0}^{\prime\prime}\{z\}$ which is a subset of $x\subseteq X$. $H_{0}^{\prime\prime}\{z\}$ has the same cardinality as the factor group, $H_{0}/\mathrm{sym}_{H_{0}}(z)$, which is a power of two.

Let $Z=\{u\subseteq X:X\in M\}$ and let $f$ be a multiple choice function on $Z$ such that for each $u\in Z,\,|f(u)|$ is a power of 2. We define a multiple choice function $g$ on $M$ satisfying $MC(\infty,\, \infty$, even) as follows. Suppose $X\in M$. If $|f(X)|>1$, then $g(X)=f(X)$. If $|f(X)|=1$, consider $u=f(X-\{f(X)\})$. If $|u|=1$, then define $g(X)=\{f(X),\,u\}$; and if $|u|>1$, define $g(X)=u$.
\end{proof}

\begin{theorem}\label{part04:thm2b}
$MC(\infty,\,\infty,\,odd)$ is true in $\mathcal{N}2^{\ast}(3)$.
\end{theorem}

\begin{proof}
In this model the group $\mathcal{G}$ is Abelian and each element has order three. Let $M$ be a set of non-empty sets. We shall show that there is a multiple choice function $f$ on $M$ such that for each $X\in M,\,|f(X)|$ is a power of three. The proof is the same as the proof of $MC(\infty,\,\infty$,
even) in Theorem~\ref{part04:thm1b} above up to and including the definition of $x=y\cap X$. Then we proceed as follows. For each $u\in x$, let $Q(u)=\{\sigma\in \mathcal{G}:\sigma(u)\in x\}$. Each $Q(u)$ has null support and is therefore in $W$. Let $Q_{0}$ be the $R$-first element of $\{Q(u):u\in x\}$ and let $z=\{u\in x:Q(u)=Q_{0}\}$. We claim that the cardinality of $z$ is a power of three. Fix $u\in z$ and let $K=\{\sigma\in \mathcal{G}:\sigma(u)\in z\}.\,K$ is a subgroup of $\mathcal{G}$ because $K=\{\sigma\in \mathcal{G}:(\sigma^{-1})^{\prime\prime}Q_{0}=Q_{0}\}$. Therefore, $|z|=|\mathcal{G}/$fix$_\mathcal{G}(u)|$ which is a power of three.
\end{proof}

\section*{NOTE 94}

In this note we summarize what is known about the relationships between possible definitions of ``finite''\index{Subject}{definitions of finite}.

Several authors have studied this topic they include Diel\index{Author}{Diel} \cite{bib:1974}, Howard/Yorke\index{Author}{Howard/Yorke} \cite{bib:1989}, Levy\index{Author}{Levy} \cite{bib:1958}, Spi\u{s}iak\index{Author}{Spi\u{s}iak} \cite{bib:1993}, Spi\u{s}iak/Vojt\'{a}\u{s}\index{Author}{Spi\u{s}iak/Vojt\'{a}\u{s}} \cite{bib:1988}, Tarski\index{Author}{Tarski} \cite{bib:1938b} and \cite{bib:1954a}, Truss\index{Author}{Truss} \cite{bib:1974a} and Hickman\index{Author}{Hickman} \cite{bib:1972}. The first eight definitions and their names are due to Tarski \cite{bib:1954a}

\begin{definition*}. A set $A$ is said to be
\begin{enumerate}
\item[] \emph{decomposable} if there are two pairwise disjoint non-empty sets $x$ and $y$ such that $x\prec A,\, y\prec A$, and $A=x\cup
y$. (Similarly, a cardinal number $m$ is said to be \emph{decomposable}, if there exist two non-zero cardinal number $p$ and $q$ such that $p<m,\,q<m$, and $m=p+q$.)

\item[]\emph{Dedekind finite}\index{Subject}{Dedekind finite} if it is not equivalent to a proper subset of itself. Otherwise, $A$ is \emph{Dedekind infinite}.

\item[] $T$-\emph{finite}\index{Subject}{$T$-finite} if every non-empty, monotone $X\subseteq \mathcal{P}(A)$ has a $\subseteq$-maximal element.
\item[] $I$-\emph{finite} if every non-void family of subsets of $A$ has an $\subseteq$-maximal element. (This is (equivalent to) the usual definition of finite.)
\item[] $Ia$-\emph{finite} if it is not the union of two disjoint sets neither of which is I-finite.
\item[] \emph{II-finite} if every non-empty $\subseteq$-monotone family of subsets of $A$ has a $\subseteq$-maximal element.
\item[] \emph{III-finite} if the power set of $A$ is Dedekind finite.
\item[] \emph{IV-finite} if $A$ is Dedekind finite.
\item[] \emph{V-finite} if $|A|=0$ or 2 $\cdot|A|>|A|$.

\item[] $VI$-\emph{finite} if $|A|=0$ or $|A|=1$ or $|A|^{2}>|A|$.
\item[] \emph{VII-finite} if $A$ is I-finite or not well orderable.
\end{enumerate}

In Tarski \cite{bib:1954a} and Levy \cite{bib:1958} it is shown that if a set is finite according to any definition in the above list then it is finite according to any following definition and none of the reverse implications hold.

In Tarski \cite{bib:1938b} the following definitions are considered:
\end{definition*}

\begin{definition*}
For each $n\in \omega$, A set $A$ is $T(n)$-finite if for every $S\subseteq \mathcal{P}(A)$, if every collection of $n$ non-empty, pairwise disjoint subsets of $A$ contains an element of $S$, then there are $2n+1$ elements $x(i)$, $1\leq i\leq 2n+1$ of $S$ such that $A\subseteq\bigcup\{x(i):1 \leq i\leq 2n+1\}$.
\end{definition*}

In Howard/Yorke \cite{bib:1989} the following definition is introduced:

\begin{definition*}
A set $A$ is $D$ finite if $|A|=0$ or $|A|=1$ or $A=C\cup D$ with $|C|<|A|$ and $|D|<|A|$. (Note that $A$ is $D$ finite if and only if $|A|=0$ or $|A|=1$ or $A$ is decomposable.
\end{definition*}

It is\index{Author}{Levy}\index{Author}{Tarski} shown there that IV $\rightarrow D$ and $D\rightarrow$ VII. Further, $V \not\rightarrow D$ and $D\not\rightarrow$ VI.

In Truss\index{Author}{Truss} \cite{bib:1974a} the following classes of Dedekind finite cardinals are considered:

\begin{definition*}\
\begin{enumerate}
\item[]$\Delta_{1}=\{x:x=y+z\rightarrow y$ or $z$ is finite\}. (We will say $A$ is $\Delta_{1}$ finite if $|A|\in\Delta_{1}$.
\item[] $\Delta_{1}$-finite is equivalent to Ia-finite.)
\item[] $\Delta_{2}=\{|A|$: any linearly ordered partition of $A$ is finite\}. $(\Delta_{2}$-finite is equivalent to II-finite.)
\item[] $\Delta_{3}=\{|A|$: any linearly ordered subset of $A$ is finite\}.
\item[] $\Delta_{4}=\{x:\neg(\aleph_{0}\leq^{\ast}x)\}$. ($\Delta_{4}$-finite is equivalent to III-finite.)
\item[] $\Delta_{5}=\{x:\neg(x+1\leq^{\ast}x)\}$.
\end{enumerate}

It is shown that I-finite $\rightarrow\Delta_{1}\rightarrow\Delta_{2}\rightarrow\Delta_{4}\rightarrow\Delta_{5}\rightarrow$ IV-finite and $\Delta_{2}\rightarrow\Delta_{3}\rightarrow$ IV-finite and that none of the implications are reversible. It is also shown that various combinations of equality and inequality between these classes are consistent. For example, it is shown that
\begin{equation*}
\text{ZF}+\omega=\Delta_{1}=\Delta_{2}=\Delta_{4}\neq\Delta_{5}\neq\ \text{IV-finite}\ +\Delta_{2}=\Delta_{3}\neq\ \text{IV-finite}
\end{equation*}
is consistent.

In Diel\index{Author}{Diel} \cite{bib:1974} two definitions are finite are considered: ``Almost finite''\index{Subject}{almost finite}  which is equivalent to III-finite and ``strongly Dedekind finite''\index{Subject}{strongly Dedekind finite} which is equivalent to $\Delta_{5}$-finite.

Finally in Howard/Spi\u{s}iak\index{Author}{Howard/Spi\u{s}iak}\index{Author}{Spi\u{s}iak} \cite{bib:1994}, Spi\u{s}iak \cite{bib:1993} and Spi\u{s}iak/Vojt\'{a}\u{s}\index{Author}{Spi\u{s}iak/Vojt\'{a}\u{s}} \cite{bib:1988} the following forms are considered:
\end{definition*}

\begin{definition*}
Let $F$ be one of I, Ia, II, III, IV, V, VI or VII. Then a set $A$ is $F^{\prime\prime}$-finite if $\mathcal{P}(A)$ is $F$-finite.
\end{definition*}

It is shown that each of $\mathrm{I}^{\prime\prime}, \mathrm{Ia}^{\prime\prime}, \mathrm{II}^{\prime\prime}$ and $\mathrm{III}^{\prime\prime}$ are equivalent to I and that the following implications hold: III $\rightarrow \mathrm{V}^{\prime\prime}\rightarrow \mathrm{IV},\mathrm{V}^{\prime\prime}\rightarrow \mathrm{VI}^{\prime\prime} \rightarrow \mathrm{V}$, and $\mathrm{VII}\rightarrow \mathrm{VI}^{\prime\prime}$. Further it is consistent (with $\mathrm{ZF}$) that there exists a set with any combination of the properties $\mathrm{III},\mathrm{V}^{\prime\prime},\mathrm{VI}^{\prime\prime}, \mathrm{IV}$ and $\mathrm{V}$ not excluded by the known implications described so far in this note. Finally, it is shown that form 88 $(C(\infty,\,2))$ implies that $\mathrm{V}^{\prime\prime}$ and III are equivalent.

In addition, in Howard/Yorke\index{Author}{Howard/Yorke} \cite{bib:1989}, if $J$ and $K$ are two possible definitions of finite then the assertion that $J$ and $K$ are equivalent is denoted by $E(J,\,K)$ and the strength of these weak forms of $AC$ is investigated. Several of these appear in the list of forms: $E(I,\,IV)$ is form 9, $E(I,\,Ia)$ is form 64, $E(I,\,III)$ is form 82, $E(I,\,II)$ is form 83, $E(II,\,III)$ is form 84, $E(II,\,IV)$ is form [9 B], $E(IV,\,V)$ is form [3 E], $E(V,\,VI)$ is form [1 AH], $E(VI,\,VII)$ is form [1 AI], $E(V^{\prime\prime},\,III)$ is form 276 and $E(D,\,VII)$ is form 277.

We note that $E(I,\,D)$ which is equivalent to ``\emph{every infinite set is decomposable}'' is equivalent to the axiom of choice since it is form CN 4F in H. Rubin/J. Rubin \cite{bib:1985}, p. 139.

The known relationships between the various definitions of finite are summarized in the following diagram:

\begin{figure}[h]
\includegraphics{part04-vend-scan-01.eps}
\end{figure}

\noindent(In addition\index{Author}{H. Rubin/ J. Rubin} $\mathrm{VI}^{\prime\prime}\not\rightarrow \mathrm{IV}, \mathrm{IV} \not\rightarrow \mathrm{VI}^{\prime\prime},\, \mathrm{V}\not\rightarrow D,\, D\not\rightarrow \mathrm{VI}$, and none of the arrows are reversible.)

This leaves\index{Author}{Howard/Yorke} several problems open. For example, what is the relationship between $\Delta_{3}$ and III or between $\Delta_{5}$ and $\mathrm{V}^{\prime\prime}$?

(In addition $\mathrm{VI}^{\prime\prime}\not\rightarrow \mathrm{IV}, \mathrm{IV} \not\rightarrow \mathrm{VI}^{\prime\prime},\,\mathrm{V} \not\rightarrow D,\, D\not\rightarrow \mathrm{VI}$, and none of the arrows are reversible.)

This leaves several problems open. For example, what is the relationship between $\Delta_{3}$ and III or between $\Delta_{5}$ and $\mathrm{V}^{\prime\prime}$?

\section*{NOTE 95}

Form [154 A] appears in Brunner\index{Author}{Brunner} \cite{bib:1984b} as $T_{\omega}$ and there Brunner refers to Brunner \cite{bib:1983c} for the result that form [154
A] is not provable in $\mathrm{ZF}^{0}$. In Brunner \cite{bib:1983c} (p 437 Satz 2.2 (5)), the result $(tych_{2}(\omega)$ + form 106) $\Leftrightarrow DC$ appears, that is,
\begin{equation*}\tag{$\ast$}
(\mathrm{form}\ 154 + \mathrm{form}\ 106)\Leftrightarrow \mathrm{form}\ 43.
\end{equation*}
Also in Brunner \cite{bib:1983c} (p 437 Satz 2.2 (1)) is the result
\begin{equation*}\tag{$\ast\ast$}
\mathrm{form}\ [106 A]\Rightarrow \mathrm{form}\ 106
\end{equation*}
In the $\mathrm{ZF}^{0}$ model $\mathcal{N}2$, form [106 A] is true and form 43 is false. Thus, combining $(\ast)$ and $(\ast\ast)$, we get the result that form 154 is not a theorem of $\mathrm{ZF}^{0}$.

\section*{NOTE 96}

In Wright \cite{bib:1973} it is shown that form $0\not\Rightarrow$ form 281 by showing that every linear operator on a Hilbert space is bounded in Solovay's model $\mathcal{M}5$. In this note we give the definitions for forms 279 and 281 and related results from Wright \cite{bib:1973}.

\begin{definition*}
Suppose $X$ is a topological vector space
\begin{enumerate}
\item[1.] $X$ is\index{Author}{Brunner} \emph{locally convex}\index{Subject}{locally convex} if there is a neighborhood base at $\overset{\rightarrow}{0}$ consisting of convex sets.
\item[2.] A metric $d$ on $X$ is \emph{invariant} if for all $x,\,y$ and $z$ in $X,\,d(x+z,\,y+z)=d(x,\,y)$.
\item[3.] $X$ is a \emph{Fr\'{e}chet space}\index{Subject}{Fr\'{e}chet space} if $X$ is locally convex and the topology on $X$ is induced by a complete, invariant metric\index{Subject}{invariant metric}.
\end{enumerate}

In the model $\mathcal{M}5$ (Solovay\index{Author}{Solovay}\cite{bib:1970}) form 43 (DC) and the negation of form 280 hold (the negation of 280 is ``Every subset of a complete, separable metric space has the Baire property.''). Wright\index{Author}{Wright} \cite{bib:1973} shows that the following are consistent by showing that they are theorems of $\mathrm{ZF}+DC+\neg \mathrm{form}\,280$:
\begin{enumerate}
\item[A.] If $X$ and $Y$ are separable, metrizable topological groups and $X$ is complete and $H:X\rightarrow Y$ is a group homomorphism, then $H$ is continuous.
\item[B.] If $X$ is a Fr\'{e}chet space and $Y$ is a separable, metrizable topological vector space and $T:X\rightarrow Y$ is linear, then $T$ is continuous.
\item[C.] Every linear functional on a Fr\'{e}chet space is continuous.
\item[D.] If $X$ and $Y$ are Fr\'{e}chet spaces, $T:X\rightarrow Y$ is linear and there are enough linear functionals on $Y$ to separate the points of $Y$ then $T$ is continuous.
\item[E.] If $H$ is a Hilbert space and $T:H\rightarrow H$ is linear, then $T$ is bounded.
\end{enumerate}
\end{definition*}

\section*{NOTE 97}

Definitions for forms 282 and [0 AD] and related results from Kleinberg/Seiferas\index{Author}{Kleinberg/Seiferas}
\cite{bib:1973}.

\begin{definition*}
Assume $\kappa \geq\omega$ and $\gamma\geq 1$ are well ordered cardinals and $\alpha$ and $\beta$ are ordinals. Then $\kappa \rightarrow(\beta)_{\gamma}^{\alpha}$\index{Subject}{$\kappa \rightarrow(\beta)_{\gamma}^{\alpha}$} denotes the partition relation ``For every function $F$ mapping the set of $\alpha$-sequences from $\kappa$ into $\gamma$ there exists a subset $X$ of $\kappa$ of order type $\beta$ such that the range of $F$ on the set of $\alpha$ sequences from $X$ has cardinality 1.'' If $\gamma=2$ we write $\kappa \rightarrow(\beta)^{\alpha}$. (An $\alpha$-sequence from $Y\subseteq On$ is an order preserving map from $\alpha$ to $Y$.) $\kappa\not\rightarrow (\beta)_{\gamma}^{\alpha}$ denotes the negation of $\kappa \rightarrow(\beta)_{\gamma}^{\alpha}$.
\end{definition*}

Form [0 AD] is: For any finite $n$, for all ordinals $\beta$ and well ordered cardinals $\gamma$, there is a well ordered cardinal $\kappa$ such that $\kappa \rightarrow(\beta)_{\gamma}^{n}$ and form 282 is $\omega\not\rightarrow(\omega)^{\omega}$. The axiom of choice implies for any $\alpha\geq \omega$, given any well ordered cardinal $\gamma>1$ and ordinal $\beta$ there is no $\kappa \geq \omega$ for which $\kappa \rightarrow(\beta)_{\gamma}^{\alpha}$. Also proved in Kleinberg/Seiferas \cite{bib:1973} are the following:

\begin{theorem*}
Assume $\gamma>1$ is a well ordered cardinals and $\alpha$ and $\beta$ are ordinals with $\alpha\leq\beta$ then
\begin{enumerate}
\item[\emph{1.}] $C(WO,\,\infty)$ (which is form 40) implies $\neg\exists\kappa(\kappa\rightarrow(\beta)_{\gamma}^{\alpha})$ provided
\begin{equation*}\tag{$\ast$}
\alpha\geq \omega+\omega\ \text{(ordinal addition) or}\ \beta\geq\lfloor\alpha\rfloor+\omega+\lceil\alpha\rceil\ \mathrm{or}\ \gamma\geq \omega
\end{equation*}
\item[\emph{2.}] If ($\ast$) fails then $\exists \kappa(\kappa\rightarrow(\beta)_{\gamma}^{\alpha})\Leftrightarrow \omega\rightarrow(\omega)^{\omega}$.
\item[\emph{3.}] Assume $C(WO,\,\infty)$ and ``$\mathbb{R}$ can be well ordered'' then for all $\alpha,\,\beta$ and $\gamma,\,\neg\exists \kappa(\kappa\rightarrow(\beta)_{\gamma}^{\alpha})$
\item[\emph{4.}] $C(\aleph_{\alpha},\,\infty)$ (which is form 86) implies $\aleph_{\alpha}\not\rightarrow(\omega+\omega)^{\omega+\omega}$.
\item[\emph{5.}] $C(\aleph_{\alpha},\,\infty)\Rightarrow\aleph_{\alpha}\not\rightarrow(\omega+\omega)^{\omega}$.
\item[\emph{6.}] $C(\aleph_{\alpha},\,\infty)\Rightarrow\aleph_{\alpha}\not\rightarrow (\omega)_{\omega}^{\omega}$. (So form 8 ($C(\aleph_{0},\,\infty)$) implies form 282.
\end{enumerate}
\end{theorem*}

\section*{NOTE 98}

Definitions for\index{Author}{Kleinberg/Seiferas} forms [1 AK], [1 AL] and [1 AM] from Truss\index{Author}{Truss} \cite{bib:1973a}.

\begin{definition*}
Assume $x$ and $y$ are cardinal numbers
\begin{enumerate}
\item[1.] $\aleph(x)$ is the least well ordered cardinal $\aleph$ such that $\aleph\not\leq x$.
\item[2.] $\alpha(x,\,y)$ is the least ordinal $\alpha$ such that whenever $\{z_{\gamma}:\gamma<\beta\}$ is a sequence of cardinals satisfying $x<z_{\gamma}<y$ for all $\gamma<\beta$ and $\gamma<\delta\rightarrow z_{\gamma}<z_{\delta}$ then $\beta<\alpha$.
\end{enumerate}
\end{definition*}

\section*{NOTE 99}

Form 65 (Krein-Milman Theorem)\index{Subject}{Krein-Milman theorem} does not imply form 14 (BPI) in $\mathrm{ZF}^{0}$. The argument is as follows: By results of Bell/Jellett\index{Author}{Bell/Jellett} \cite{bib:1971}, form 14 + form 65 implies form 1 (AC). By Pincus \cite{bib:1972b} and \cite{bib:1972c}, 65 is true and 1 is false in $\mathcal{N}2$ hence form 14 is false in $\mathcal{N}2$. Similarly, form 52 does not imply form 286 in ZF since form 52 + form 286 implies form 1 (Bell/Jellett \cite{bib:1971}) and by Pincus\index{Author}{Pincus} \cite{bib:1972b}, form 52 does not imply form 14 in ZF.

\section*{NOTE 100}

In this note we show that in $\mathcal{N}23\,C(\infty,\,2)$ (form 88) is false and $C(\aleph_{0},\,\infty)$ (form 126) is true. The proof that $C(\aleph_{0},\,\infty)$ is true is due to A. Rubin.

First we shall show that $C(\infty,\,2)$ is false. Let $B$ be the set of subsets $x$ of $A$ such that the order type of $x$ with respect to $<$ is $\omega$. Let $D$ be the set of equivalence classes of elements of $B$ under the relation $xRy$ iff the symmetric difference, $(x-y) \cup(y-x)$ is finite. (Denote the equivalence class of $x$ by $[x]$.) Let $P$ be the set of pairs of elements of $D$. $P$ has no choice function in the model because if $f$ is such a choice function and $E\subseteq A$ is a support of $f$ it follows that for all order preserving permutations $\phi$ of $A$ such that $\phi$ fixes $E$ pointwise and for all $[x],\, [y]$ in $D$, if $\phi$ fixes $\{[x],\,[y]\}$ then $\phi$ fixes both $[x]$ and $[y]$. We can construct a $\phi,\,x$ and $y$ for which this is false as follows: Since $E$ is well ordered there are elements $a,\,b\in A$ such that the interval $(a,\,b)$ contains no points of $E$. Let $\phi$ be an order preserving permutation of $A$ that fixes $E$ pointwise, fixes $a$ and $b$ but moves at least one element $t$ of $(a,\,b)$. We may assume (by replacing $\phi$ with $\phi^{-1}$ if necessary) that $t<\phi(t)$. Let $x=\{\phi^{2n}(t):n\in \omega\}$ and $y=\{\phi^{2n+1}(t):n\in\omega\}$. Then $\phi([x])=[y]$ and $\phi([y])=[x] ($and $[x]\neq[y])$.

Next, we shall show that $C(\aleph_{0},\,\infty)$ is true in $\mathcal{N}23$. Let $X=\{Y_{i}:i\in\omega\}$ be a denumerable set of non-empty sets in the model. Let $B$ be a common support for $X$ and the $Y_{i}$'s, $i\in \omega$. Let the open intervals of $A-B$ be denoted by $I_{\beta}$, for $\beta<\alpha$. For each $i\in\omega$, choose (in the outer model) $z_{i}\in Y_{i}$ and let $B_{i}$ be a support of $z_{i}$. Let $J_{\beta i}=I_{\beta}\cap B_{i}$. (Note: Whenever we choose a sequence of elements in $A$, we choose them in increasing order. Also, $\beta$ will range over $\alpha$, and $i$ and $j$ will range over $\omega$.)

We then define $K_{\beta ji}$ as follows: For each $\beta$ and $i$, choose $E_{\beta i}=\{e_{\beta ji}:j\in \omega\}$ cofinal in $I_{\beta}$ such that each $e_{\beta oi}<J_{\beta i}$ and $E_{\beta i}\cap \mathrm{cl}(B_{i})=\emptyset$. Then define $K_{\beta ji}= B_{i}\cap (e_{\beta ji},\,e_{\beta j+1i})$. (If $J_{\beta i}$ is not cofinal in $I_{\beta}$, choose $E_{\beta i}$ such that $e_{\beta oi}<J_{\beta i}<e_{\beta 1i}$ so there is only one non-empty $K_{\beta ji}$.)

For each $\beta$, choose $\{c_{\beta i}:i\in \omega\}$ cofinal in $I_{\beta}$ and choose $\{d_{\beta ji}:i\in \omega\}$ in the interval $(c_{\beta j},\,c_{\beta j+1})$. For each $\beta,\,j$, and $i$, let $f_{\beta ji}$ be a function mapping $K_{\beta ji}$ into the interval $(d_{\beta ji},\,d_{\beta ji+1})$ and let $L_{\beta ji}$ be the range of $f_{\beta ji}$. Next, let $f_{i}$ be an extension of $\bigcup_{\beta,j}f_{\beta ji}$ to a complete isomorphism of $A$ fixing $B$. We claim that $F=\{<Y_{i},\,f_{i}(z_{i})>:i\in \omega\}$ is a choice function for $X$ that is in the model. Since $f_{i}$ preserves $B$, we must have that $f_{i}(z_{i})\in f(Y_{i})=Y_{i}$. A support of $F$ is $\bigcup_{\beta,j,i}L_{\beta ji}$ which is a well ordered union of the well ordered sets $L_{\beta ji}$ and, therefore, can be well ordered.

\section*{NOTE 101}

The argument that form 43 (DC) implies [43 S] (If $S$ is a relation on $A$ such that $(\forall x\in A)(\exists y\in A)(x\,S\,y)$ and $x_{0}$ is any element of A, then there is a sequence $a(1),a(2),\ldots$ of elements of $A$ such that $a(n)\,S\,a(n+1)$ for all $n\in \omega$ and such that $a(0)=x_{0}.)$ is as follows: Let $S$ be a relation on a set $A$ satisfying the hypotheses of [43 S] and assume that $x_{0}\in A$. Let $B$ be the set of finite sequences $y=\langle y(i)\rangle_{i=0}^{n}$ of distinct elements of $A$ such that $y(0)=x_{0}$ and $y(i)\,S\,y(i+1)$ for $i=0$ to $n-1$. Define $R$ on $B$ by $y\,R\,z\Leftrightarrow y$ is a proper initial segment of $z$. Then the relation $R$ on $B$ satisfies the hypotheses of form 43. The conclusion of form 43 gives a sequence $y_{0}\,R\,y_{1}\,R\,y_{2}\,\ldots$ which gives an infinite sequence of elements of $A$ satisfying the conclusion of [43 S].

\section*{NOTE 102}

Form 289 (If $S$ is a set of subsets of a countable set and $S$ is closed under chain unions, then $S$ has a $\subseteq$-maximal element.) is false in $\mathcal{M}1$. Let $b$ be the generic set of generic reals in $\mathcal{M}1$, let $E(b)$ be the set of all finite subsets of $b$, and let $S=\{\cup x:x\in E(b)\}$. Since $b$ is generic, for any finite subset $\{a_{0},\,a_{1},\,\ldots,\,a_{n}\}$ of $b,\,a_{0}\not\subseteq \cup\{a_{1},\,\ldots,\,a_{n}\}$ It follows that if $x,\,y\in E(b)$, then $\cup x\subseteq \cup y$ iff $x\subseteq y$. Hence, any chain in $S$ corresponds to a chain in $E(b)$ and a $\subseteq$-maximal element of $S$ corresponds to a $\subseteq$-maximal element of $E(b)$. We claim that there is no infinite chain in $E(b)$. For suppose $x_{0}\subset x_{1}\subset x_{2}\cdots$ is an infinite chain in $E(b)$. Let $a_{n}$ be the smallest element in $x_{n+1}-x_{n}$ for $n\in \omega$. Then $\{a_{n}:n\in \omega\}$ is a denumerable subset of $b$, which is impossible because $b$ is Dedekind finite. Thus, all chains in $E(b)$ are finite so $E(b)$ is closed under chain unions. Consequently, $S$ is closed under chain unions. Since $b$ is infinite, $E(b)$ has no $\subseteq$-maximal element and, therefore, $S$ has no $\subseteq$-maximal element.

\section*{NOTE 103}

The results of Jech/Sochor \cite{bib:1966a} and \cite{bib:1966b} and Pincus \cite{bib:1971}, \cite{bib:1972a}, and \cite{bib:1977a} show that many independence results in $\mathrm{ZF}^{0}$ will transfer automatically to ZF. These transfer theorems are given in this note. We use the notation of Pincus.

\begin{definition*}
Assume $x$ is a set, $A$ is a class and $\Phi(x_{1},\,\ldots,\,x_{n})$ is a
formula in the language of $\mathrm{ZF}$ with variables among $x_{1},\,\ldots,\,x_{n}$.
\begin{enumerate}
\item[(1)] $|x|$ is the least ordinal $\alpha$ such that there is a one to one function from $x$ onto $\alpha$. (Note that this definition of $|x|$ is for this note only and differs from the definition of $|x|$ used elsewhere in this document. Also note that $|x|$ is defined only for well ordered sets $x$.)
\item[(2)] $|x|^{-}=\sup\{|y|$ : there is a surjection from $x$ onto $y$\}.
\item[(3)] $|x|_{-}=\sup\{|y|$ : there is an injection from $y$ to $x$\}.
\item[(4)] $R_{\beta}(x)$ is defined for ordinals $\beta$ by induction: $R_{0}(x)=x,\, R_{\alpha+1}(x)=R_{\alpha}(x)\cup \mathcal{P}(R_{\alpha}(x))$ and for limit ordinals $\lambda,\, R_{\lambda}(x)=\bigcup_{\gamma<\lambda}R_{\gamma}(x)$.
\item[(5)] $\Phi^{A}(x_{1},\,\ldots,\,x_{n})$ is $\Phi$ with quantifiers restricted to $A$. Similarly, if $\sigma$ is a term in $\mathrm{ZF}$ then $\sigma^{A}$ is the term defined by the same formula that defines $\sigma$ but with its quantifiers restricted to $A$.
\item[(6)] Assume that $V$ is a model of $\mathrm{ZF}$ and $V_{0}\subseteq V$ is a substructure of $V$ which is a model of $\mathrm{ZF}^{0}$ such that every element of a $V_{0}$ \emph{set} is an element of $V_{0}$. (But elements of $V_{0}$ atoms are not in $V_{0}$.) Assume also that $V$ and $V_{0}$ have the same ordinals, the same cofinality function, the same surjective cardinal function $|x|^{-}$ (on $V_{0}$ sets), the same injective cardinal function $|x|_{-}$ (on $V_{0}$ sets) and the same cardinality function $|x|$ (on $V_{0}$ sets.) Then the formula $\Phi(x_{1},\,\ldots,\,x_{n})$ is \emph{absolute} if for all $x_{1},\,\ldots,\,x_{n}$ in $V_{0},\, \Phi^{V_{0}} (x_{1},\,\ldots,\,x_{n})\leftrightarrow\Phi(x_{1},\,\ldots,\,x_{n})$. (This is sometimes called absolute for $V_{0},\,V$.)
\item[(7)] A term $\sigma(x_{1},\,\ldots,\,x_{n})$ is \emph{ordinal valued} if $\mathrm{ZF}^{0}\vdash\sigma(x_{1},\, \ldots,\,x_{n})\in On$.
\item[(8)] A formula $\Phi(x_{1},\,\ldots,\,x_{n})$ is \emph{boundable} if for some absolute ordinal valued term $r,\,\mathrm{ZF}^{0}\vdash\Phi(x_{1},\,\ldots,\,x_{n})\leftrightarrow\Phi^{R_{r}(x_{1}\cup\cdots\cup
x_{n})}(x_{1},\,\ldots,\, x_{n})$. Similarly, the term $\sigma(x_{1},\,\ldots,\,x_{n})$ is \emph{boundable} if
\begin{equation*}
\mathrm{ZF}^{0}\vdash\sigma(x_{1},\,\ldots,\,x_{n})=\sigma^{R_{r}(x_{1}\cup\cdots \cup x_{n})}(x_{1},\,\ldots,\,x_{n}).
\end{equation*}
\item[(9)] A sentence is \emph{boundable} if it is the existential closure of a boundable formula\index{Subject}{boundable formula}.
\end{enumerate}

The theorem of Jech/Sochor\index{Author}{Jech/Sochor} \cite{bib:1966b} is
\end{definition*}

\begin{theorem*}
A boundable sentence is transferable\index{Subject}{transfer theorems}.
\end{theorem*}

Many applications of the theorem use the following corollary:

\begin{corollary*}
If $\Phi$ is a sentence such that there is an ordinal $\alpha$ such that $ZF^{0}\vdash \Phi\leftrightarrow\Phi^{R_{\alpha}(A)}$ where $A$ is the set of atoms, then $\Phi$ is transferable. (Note that $\alpha$ must be a constant absolute ordinal valued term.)
\end{corollary*}

For example, form 5, $C(\aleph_{0},\,\aleph_{0},\,\mathbb{R})$ is the statement ``If $D$ is a denumerable collection of denumerable subsets of $\mathbb{R}$ then $\exists f: D\rightarrow\bigcup D$ such that $\forall d\in D,\,f(d)\in D$. This is equivalent to $(C(\aleph_{0},\, \aleph_{0)}\mathbb{R}))^{R_{\omega_{2}}(A)}$. Similarly, the following forms are boundable and therefore transferable: 6, 11, 12, 13, 35, 36, 37, 38, 70, 73, 74, 79, 92, 93, 94, 130, 142, 169, 170, 194, 197, 199, 203, 206, 211, 212, 222, 223, 251, 252, 272, 273, 280, 289, 305, 306, 307, 309, 313, 361, 362, 363, 364, 366, 367, 368, and 369. In addition all of these forms are true in all Fraenkel-Mostowski models since they make assertions about the standard part of the model only.

Pincus\index{Author}{Pincus} has strengthened the theorem of Jech/Sochor by the following:

\begin{definition*}
A formula $\Phi(y_{1},\,\ldots,\,y_{k})$ is \emph{surjectively boundable} if it is a con junction of formulas of the form
\begin{equation*}
\begin{split}
\Phi_{i}=\forall x_{1},\,\ldots,\, x_{n}[(|\bigcup_{j=1}^{n}x_{j}|^{-}\leq\sigma_{i}(y_{1},\,\ldots,\,y_{k})\wedge&
\bigcup_{j=1}^{n}x_{j}\cap TCL( \bigcup_{m=1}^{k}y_{m})=\emptyset)\\
&\rightarrow\Psi(x_{1},\,\ldots,\,x_{n},\,y_{1},\,\ldots,\,y_{k})],
\end{split}\tag{$\ast$}
\end{equation*}
where $\sigma(y_{1},\,\ldots,\,y_{k})$ and $\Psi(x_{1},\,\ldots,\,x_{n},\,y_{1},\,\ldots,\,y_{k})$ are boundable. A sentence is \emph{surjectively boundable}\index{Subject}{surjectively boundable} if it is the existential closure of a surjectively boundable formula.

A formula is \emph{injectively boundable}\index{Subject}{injectively boundable} if it is a conjunction of formulas of the form $(\ast)$ but with $\bigcup_{j=1}^{n} x_{j}|^{-}$ replaced by $\bigcup_{j=1}^{n}x_{j}|_{-}$ and a sentence is \emph{injectively boundable} if it is the existential closure of an injectively boundable formula.
\end{definition*}

\begin{theorem*}
An injectively boundable statement is transferable.
\end{theorem*}

Note that boundable statements are surjectively boundable and surjectively boundable statements are injectively boundable.

In many cases the following corollary suffices:

\begin{corollary*}
A formula $\Phi$ is transferable if $\Phi$ is a conjunction of sentences of the form
\begin{equation*}
\forall x_{1},\,\ldots,\,x_{n}(|\bigcup_{j=1}^{n} x_{j}|_{-}\leq\sigma_{i}\rightarrow\Psi_{i}(x_{1},\,\ldots,\, x_{n}))
\end{equation*}
where $\sigma_{i}$ is a boundable constant term and $\Psi$ is boundable.
\end{corollary*}

For example, form 9, Dedekind finite $\cong$ finite, is injectively boundable since it is equivalent to $\forall x$ ($|x|_{-}\leq \omega\rightarrow x$ is finite). Similarly, 18 and 128 are injectively boundable. Form 166 $(PC(\infty,\,2,\,\infty))$ is injectively boundable since it is equivalent to $(\forall x)(|x|_{-}\leq\omega\rightarrow$ every infinite family of pairs whose union is $x$ has an infinite subfamily with a choice function). Similarly, 132, 336, 342, 376, 377, and 378 are injectively boundable. Other injectively boundable forms are 10, 31, 32, 80, 119, 350, 357, 358, 373, and 374 (all using Howard/Solski\index{Author}{Howard/Solski} \cite{bib:1993}, lemma 3.5), 17 (Blass\index{Author}{Blass} \cite{bib:1977a}), 64 (Pincus\index{Author}{Pincus} \cite{bib:1972a}, 2B3), 71, 82, 83 (Pincus \cite{bib:1972a}, 2B4), 98, 163, 216, (Pincus \cite{bib:1972a},
2B5), 217 (Pincus \cite{bib:1972a} 2B6), and 325 (note 46).

In Pincus \cite{bib:1972a} the following transfer theorem also appears.

\begin{definition*}
A formula $\Psi(x,\,w_{1},\,\ldots,\,w_{n})$ is an \emph{onto property} if it is provable that
\begin{enumerate}
\item[(a)] If there is a function $f$ from $x$ onto $y$ and $\Psi(x, w_{1},\,\ldots,\, w_{n})$ then $\Psi(y, w_{1},\,\ldots w_{n})$.
\item[(b)] $x\in V_{0}\rightarrow(\Psi(x,\,w_{1},\,\ldots,\,w_{n})\leftrightarrow\Psi^{V_{0}}(x,\,w_{1},\,\ldots,\, w_{n}))$.
\item[] A term $\sigma(x,\,w_{1},\,\ldots w_{n})$ is \emph{invariant} if
\item[(a)] $\sigma(x,\,w_{1},\,\ldots,\,w_{n})=\{z\in R_{\alpha}(x):\Phi^{R_{\alpha}(x)}(z,\,x,\,w_{1},\,\ldots,\, w_{n})\}$ for some $\Phi$.
\item[(b)] $\sigma(x,\,w_{1},\,\ldots,\,w_{n})\neq\emptyset$ if $x$ is infinite
\item[(c)] There is a definable operator $\hat{\ }$ such that if $f:y\rightarrow x$ is one to one then $\hat{f}:\sigma(y,\,w_{1},\,\ldots,\,w_{n})\rightarrow\sigma(x,\,w_{1},\,\ldots,\,w_{n})$ is one to one.
\end{enumerate}

A formula $Z(\overset{\rightarrow}{\beta})$ with ordinal parameters $\overset{\rightarrow}{\beta}$ is a \emph{term-choice formula} if it is equivalent to a formula in the form
\begin{align*}
(\forall y)&[(\Psi(y,\,\overset{\rightarrow}{\beta})\wedge(\forall x\in y)[x\text{is well orderable}])\\
&\rightarrow(\{\sigma(x,\,\overset{\rightarrow}{\beta}):x\in y\wedge\sigma(x,\,\overset{\rightarrow}{\beta})\neq\emptyset\}\ \text{has a choice function})]
\end{align*}
where $\Psi(y,\,\overset{\rightarrow}{\beta})$ is onto and $\sigma(y,\,\overset{\rightarrow}{\beta})$ is invariant.

A sentence is a \emph{term-choice sentence} if it is equivalent to a statement in the form $\forall\overset{\rightarrow}{\beta}[\Omega(\overset{\rightarrow}{\beta})\wedge Z(\overset{\rightarrow}{\beta})]$ where $Z(\overset{\rightarrow}{\beta})$ is term choice and $\Omega(\overset{\rightarrow}{\beta})$ is absolute.
\end{definition*}

\begin{theorem*}
Class 2 statements are transferable where a class 2 statement is a finite conjunction of

(a) surjectively\index{Subject}{onto property} boundable statements,

(b) ``every set can be linearly ordered'',

(c) term-choice\index{Subject}{term-choice sentence} statements.
\end{theorem*}

The following forms are class 2 statements: forms 30, 45, 46, 47, 48, 60 and 165 (all from Pincus\index{Author}{Pincus} \cite{bib:1972a}, pp 736--737), 61, 85, 88, 122, 178, 185, 213, 250, and 295.

Another transfer theorem from Pincus \cite{bib:1972a}:

\begin{theorem*}
If $\Phi$ is a class 2 statement which is true in the Mostowski linear order model $(\mathcal{N}3)$ then $\Phi$ is $ZF$ consistent with the Boolean prime ideal theorem (form 14).
\end{theorem*}

Finally a transfer theorem from Pincus \cite{bib:1977a}.

\begin{theorem*}
If $\Phi$ is a conjunction of any of the following seven kinds of statements and $\Phi$ has a Fraenkel-Mostowski model, then $\Phi$ has a $ZF$ model:
\begin{enumerate}
\item[\emph{1.}] Injectively boundable statements.
\item[\emph{2.}] For any ordinal number $\alpha,\,C(<\aleph_{\alpha},\,\infty)$. (This includes, for any ordinal $\alpha,\,C(\aleph_{\alpha},\,\infty)$ which is equivalent to $C(<\aleph_{\alpha+1},\, \infty)$.)
\item[\emph{3.}] For every ordinal $\eta,\, DC^{<\eta}$. See the note below for definitions.
\item[\emph{4.}] The ordering principle, $OP$ (form 30)
\item[\emph{5.}] Term choice statements
\item[\emph{6.}] The Boolean prime ideal theorem (form 14).
\item[\emph{7.}] The Hahn-Banach theorem (form 52).
\end{enumerate}
\end{theorem*}

\textsc{Note.} $DC^{<\eta}$ is the statement $\forall\lambda<\eta,\,DC^{\lambda}$ and $DC^{\lambda}$ is ``If for every $\alpha<\lambda$ every $R$-admissible $\alpha$ sequence extends to an $R$-admissible $\alpha+1$ sequence then there is an $R$-admissible $\lambda$ sequence.'' (An $\alpha$ sequence is a function with domain $\alpha$ and an $\alpha$ sequence $\sigma$ is $R$-admissible if for all $\beta<\alpha,\,\sigma|\beta\, R\,\sigma(\beta)$.

Note also that $DC^{\aleph_{\alpha}}$ is equivalent to form 87 $(DC(\aleph_{\alpha}))$.

\begin{proof}
We show that $DC(\aleph_{\alpha})$ implies $DC^{\aleph_{\alpha}}$. The proof of the other impli cation is easier. Assume $R$ is a binary relation such that for every $\beta<\aleph_{\alpha}$ every $R$-admissible $\beta$ sequence of elements of $X$ extends to an $R$-admissible $\beta+1$ sequence of elements of $X$. Define the relation $R^{\prime}$ by $\sigma\, R^{\prime}\,t$ iff $(\sigma$ is an $R$-admissible $\beta$ sequence for some $\beta<\aleph_{\alpha}$ and the $\beta+1$ sequence $\sigma^{\prime}$ which extends $\sigma$ for which $\sigma^{\prime}(\beta)=t$ is $R$-admissible) or ($\sigma$ is not $R$-admissible). $R^{\prime}$ satisfies the hypotheses of $DC(\aleph_{\alpha})$ and therefore there is an $R^{\prime}$-admissible $\aleph_{\alpha}$ sequence $f$. $f$ is $R$-admissible for if not let $\beta$ be the least ordinal $<\aleph_{\alpha}$ for which $\neg(f|\beta\, R\,f(\beta))$. Since $f|\beta\, R^{\prime}\, f(\beta)$ it must be the case that $f|\beta$ is not $R$-admissible. Hence for some $\gamma<\beta,\,\neg((f|\beta)|\gamma\, R\,f(\gamma))$. Since $(f|\beta)|\gamma=f|\gamma$, this contradicts the definition of $\beta$.
\end{proof}

Forms that can be transferred using this theorem and which have not been mentioned in this note include 8, 14, 39, 40, 43, 44, 52, 86, and 87,

In Pincus \cite{bib:1974c}, Postscript 3, it is shown that 123 is transferable and according to Pincus \cite{bib:1976}, 120 and 121 are transferable.

\section*{NOTE 104}

If $N=N(A,\,\mathcal{G},\,S)$ is a permutation model, then the $2m=m$ principle (form 3) is false in $N$ if the following holds in $N$:
\begin{equation*}\tag{$\ast$}
\begin{split}
(\forall E\in S)&[A-E\neq\emptyset\\
&\wedge(\forall a\in A-E)(\forall b\in A-\{a\})(\exists\sigma\in \mathcal{G})(\forall c\in E)(\sigma(c)=c\\
&\wedge(\sigma(a)\neq a\wedge\sigma(b)=b)\\
&\vee(\sigma(a)=b\wedge\sigma(b)=a\wedge(\forall c\in A-\{a,\,b\})\sigma(c)=c))].
\end{split}
\end{equation*}

\begin{proof}
Suppose there is a $1-1$ function $f$ in $N$ mapping $2\times A$ into $A$. Let $E$ be a support of $f$ and let $a\in A-E$. Then, since $f$ is $1-1$, $f(i,\,a)\neq a$ for $i=0$ or $i=1$. Suppose that $f(i,\, a)=b\neq a$. It follows from $(\ast)$ that there is a $\sigma\in \mathcal{G}$ that leaves $E$ pointwise fixed (so $\sigma(f)=f)$ and $\sigma(a)=c\neq a$. If $\sigma(b)=b$, then $f(i,\,a)=b$ and $f(i,\,c)=b$, contradicting the fact that $f$ is $1- 1$. If $\sigma(a)=b$ and $\sigma(b)=a$, then $f(i,\,a)=b$ and $f(i,\, b)=a$. If $j\in\{0,1\}$ and $j\neq i$, then $f(j,\,a)=c$, where $c\neq a$ and $c\neq b$. Consequently, $f(j,\,b)=c$ again contradicting the fact that $f$ is $1-1$.
\end{proof}

Examples of models that satisfy $(\ast)$ are $\mathcal{N}1,\,\mathcal{N}2,\,\mathcal{N}2(\aleph_{\alpha})$, and $\mathcal{N}12(\aleph_{\alpha})$.

\section*{NOTE 105}

Proofs for the models given in 2B3--2B6 of Pincus\index{Author}{Pincus} \cite{bib:1972a}.

\setcounter{theorem}{0}
\begin{theorem}\label{part04:thm1c}
$\mathcal{N}4\models$ 64, but 83 is false (2B3)
\end{theorem}

(Form 64 is ``There are no amorphous sets. (Every infinite set is the union of two disjoint infinite sets.)'' and form 83 is ``$T$-finite is equivalent to finite.'' (A set $S$ is $T$-finite if every non-empty, monotone $X\subseteq \mathcal{P}(S)$ has a $\subseteq$-maximal element.).)

\begin{proof}
First we shall show that in $\mathcal{N}4$ there is no infinite $\subseteq$-chain in $\mathcal{P}(A)$. Suppose $C=\{S_{i}:i\in I\}$ is such a chain with support E. Then for all $\phi\in \mathrm{fix}(E)(=\{\psi\in \mathcal{G}:(\forall a\in A)\psi(a)=a\})$, for all $i\in I,\,\phi(S_{i})\in C$, and for all $i,\,j\in I,\,S_{i}\subseteq S_{j}$ iff $\phi(S_{i})\subseteq\phi(S_{j})$. There must be an $S_{j}\in C$ and a $\phi\in$ fix$(E)$ such that $\phi(S_{j})\neq S_{j}$. Otherwise, $C$ can be well ordered, which contradicts the fact that in $\mathcal{N}4,\, \mathcal{P}(A)$ is Dedekind finite. Suppose $\phi(S_{j})=S_{k}$. We may assume without loss of generality that $S_{j}\subset S_{k}$. It follows that $S_{k}$ is Dedekind infinite, contradicting the fact that $A$ is Dedekind finite.

Next, we shall show that every infinite set is the union of two disjoint infinite sets (64 is true). Let $x$ be an infinite set with support $E$. Suppose there is a $t\in x$ and a $\phi\in$ fix$(E)$ such that $\phi(t)\neq t$. (If no such $t$ exists, $x$ can be well ordered and is therefore the union of two disjoint infinite sets.) Let $F$ be a support of $t$, then $\phi(F)$ is a support of $\phi(t)$ and $F-E\neq\emptyset$. Let $a\in F-E$. Since $\precsim$ is a homogeneous universal partial ordering there are infinitely many elements $b\in A-(E\cup F)$ such that $a\prec b$ and infinitely many $c\in A-(E\cup
F)$ such that $c\prec a$. Let $B=\{b\in A-(E\cup F):a\prec b\}$ and $C=\{c\in A-(E\cup F):c\prec a\}$. Since $a\in F-E$, for each $b\in B$ and $c\in C$ there are $\phi,\,\psi\in \mathrm{fix}(E)$ such that $\phi(a)=b$ and $\psi(a)=c$. It follows that both of the following subsets of $x$ are infinite: $\{\phi(t):\phi\in$ fix$(E)\wedge(\exists b\in B)(\phi(a)=b)\}$ and $\{\psi(t):\psi\in$ fix$(E)\wedge(\exists c\in C)(\psi(a)=c)\}$. Thus, $x$ can be partitioned into two infinite sets.
\end{proof}

\begin{theorem}\label{part04:thm2c}
$\mathcal{N}2\models$ 83, but 215 is false (2B4).
\end{theorem}

(Form 83 is given in theorem 1 and form 215 is ``If $(\forall y\subseteq X)(y$ can be linearly ordered implies $y$ is finite), then $X$ is finite.'')

\begin{proof}
First we note that the set of atoms, $A$, has the property that no infinite $x\subseteq A$ can be linearly ordered so 215 is false. We shall show that for all infinite $x$ in $\mathcal{N}2,\,\mathcal{P}(x)$ has an infinite $\subseteq$-chain. Suppose $x$ is an infinite set with support $E$. Suppose $t\in x$ with support $F$, then $y=\{\phi(F):\phi\in$ fix$(E)\}$ is finite. Consequently, $z_{t}=\{\phi(t):\phi\in$ fix$(E)\}$ is finite. Hence, since $x$ is infinite, $Z=\{z_{t}:t\in x\}$ is infinite and it can be well ordered because each element has support $E$. Let $u_{0},\,u_{1},\,u_{2},\,\cdots$ be an infinite sequence of elements of $Z$. We can construct an infinite $\subseteq$-chain $\{U_{n}:n\in \omega\}$ in $\mathcal{P}(x)$ as follows: $U_{0}=u_{0}$ and $U_{n+1}=\bigcup_{i=0}^{n}u_{i}$.
\end{proof}

\begin{theorem}\label{part04:thm3a}
$\mathcal{N}34\models$ 216, but 217 is false (2B5).
\end{theorem}

(Form 216 is ``Every infinite tree has either an infinite chain or an infinite antichain.'' and form 217 is ``Every infinite partially ordered set has either an infinite chain or an infinite antichain.''.)

\begin{proof}
First we note that the relation $R$ on the atoms defined by $c\, R\, d$ if and only if $(\exists q,\,r\in \mathbb{Q})(q<r\wedge c\in C_{q}\wedge d\in C_{r})\vee c=d$ is a partial ordering in the model with no infinite chain and no infinite antichain in the model (217 is false). We shall prove that Every infinite tree $(X,\,\precsim)$ in the model has either an infinite chain (in the model) or an infinite antichain (in the model).

If $H$ is a subgroup of $\mathcal{G}$ and $E\subset A$, we will use the notation $\mathrm{fix}_{H}(E)$ to denote the subgroup $\{\phi\in H:(\forall x\in E)(\phi(x)=x)\}$ of $H$. $H$ will usually be $\mathcal{G}$ or $G_{1}$. One useful fact is that if $\phi\in G_{1}$ and $\psi\in G_{2}$ then there is a $\psi^{\prime}\in G_{2}$ such that $\phi\circ\psi=\psi^{\prime}\circ\phi$.

We will call a finite subset $D$ of $\mathcal{C}$ a $\ast$-support of $x$ if $\bigcup D$ is a support of $x$. The following lemma about dense linear orderings and their automorphisms will be useful.
\end{proof}

\begin{lemma}\label{part04:lemA}
Assume $a,\,b\in \mathbb{Q},\,a<b$ and that $E$ and $F$ are finite subsets of the interval $(a,\,b)$. Also assume that $\phi$ is an order automorphism of $(a,\,b)$ which fixes $E$ pointwise. Then there is a finite set $\{\phi_{1},\,\ldots,\,\phi_{n}\}$ of automorphisms of $(a,\,b)$ each of which fixes $E$ pointwise and such that
\begin{enumerate}
\item[\emph{(1)}] For $i=1,\,\ldots,\,n,\,\phi_{i}$ moves at most one element of $F$ and
\item[\emph{(2)}] $(\forall z\in F)(\phi_{1}\circ\cdots\circ\phi_{n}(z)=\phi(z))$.
\end{enumerate}
\end{lemma}

(Although we won't prove this lemma, the first step would be the case where $E=\emptyset,\,F=\{y,\,z\}$ with $y<z$ and $z<\phi(y)$. In this case we would choose $\phi_{2}$ so that $\phi_{2}(z)=z$ and $y<\phi_{2}(y)<z$. Then $\phi_{1}$ could be chosen so that $\phi_{1}(y)=y,\, \phi_{1}\circ\phi_{2}(y)=\phi(y)$ and $\phi_{1}(z)=\phi(z).)$

Assume $(X,\,\precsim)$ is a tree in $\mathcal{N}34$ with $\ast$-support $E$. If every element of $X$ has $\ast$-support $E$ then $X$ is well ordered in the model and therefore applying the axiom of choice in the ground model will give either an infinite chain or infinite antichain in the model. We therefore may assume that $(\exists t\in X)(\exists\eta\in \mathrm{fix}_{\mathcal{G}}(\bigcup E))(\eta(t)\neq t)$. Assume that $t$ has $\ast$-support $E\cup F$ where $E\cap F=\emptyset$. By a previous remark $\eta$ can be written $\psi \circ \phi$ where $\phi\in G_{1}$ and $\psi\in G_{2}$ and by Lemma~\ref{part04:lemA} there are permutations $\phi_{1},\,\ldots,\,\phi_{n}$ in $\mathrm{fix}_{G_{1}}(E)$ such that each $\phi_{i}$ moves at most one element of $F$ and $\forall z\in F,\, \phi_{1}\circ\cdots \circ\phi_{n}(z)=\phi(z)$.

\begin{claim*}
There is a $\gamma$ in $\mathrm{fix}_{G_{1}}(E)$ that moves at most one element of $F$ such that $\gamma(t)\neq t$.
\end{claim*}

\begin{proof}If $\phi(t)\neq t$ then for at least one $i,\,1\leq i\leq n,\,\phi_{i}(t)\neq t$ and the proof is finished by letting $\gamma=\phi_{i}$. Otherwise $\phi(t)=t$ and therefore $\psi(t)=\psi\circ\phi(t)=\eta(t)\neq t$. The permutation $\psi$ is a product of transpositions one of which, say $\tau_{q}=(a_{q},\,b_{q})$ moves $t$. Further $C_{q}=\{a_{q},\,b_{q}\}\in F$ otherwise $\tau_{q}$ fixes $\bigcup F$ pointwise and hence $\tau_{q}(t)=t$. Choose $\gamma\in \mathrm{fix}_{G_{1}}(\bigcup[(E\cup
F)-\{C_{q}\}])$ so that $\gamma(C_{q})\neq C_{q}$. Then $\gamma(t)$ has $\ast$-support $[(E\cup F)-\{C_{q}\}]\cup \{\gamma(C_{q})\}$ hence $\tau_{q}(\gamma(t))=\gamma(t)$ since $\tau_{q}$ fixes a support of $\gamma(t)$ pointwise. But $\tau_{q}(t)\neq t$ therefore $\gamma(t)\neq t$. This completes the proof of the claim.
\end{proof}

Using the notation of the claim, let $D=(E\cup F)-\{C_{q}\}$. And assume, as in the claim, that $\gamma\in G_{1}$ fixes $D$ pointwise and that $\gamma(t)\neq t$. We will show that the set $W=\{\beta(t):\beta\in \mathrm{fix}_{G}(\bigcup
D)\}$ (which is clearly in the model) is an infinite antichain in $(X,\,\precsim)$.

The first step is to argue by contradiction that $\{t,\,\gamma(t)\}$ is not a (two element) chain in $(X,\,\precsim)$. Assume $\{t,\,\gamma(t)\}$ is a chain and assume that $\gamma(t)\prec t$. (The argument is similar is $t\prec\gamma(t).)$ We also assume that $\gamma(C_{q})\prec C_{q}$ in the ordered set $(\mathcal{C},\, \preceq)$ (Again, an argument similar to the one we will give here will work if $C_{q}\prec\gamma(C_{q}).)$ The set $S=$\{ $\phi(C_{q}):\phi\in G_{1}$ and $\phi$ fixes $D$ pointwise\} is an interval in $(\mathcal{C},\, \preceq)$. We will argue that, under the above assumptions, that $T=\{\phi(t):\phi\in G_{1}$ and $\phi$ fixes $D$ pointwise\} under the ordering $\precsim$ is order isomorphic to $S$. The isomorphism $f$ is defined by $f(\phi(t))=\phi(C_{q})$. (The details of the proof that $f$ is an isomorphism will be left to the reader.) This is a contradiction since a tree cannot have dense linearly ordered subset.

Therefore $\{t,\,\gamma(t)\}$ is an antichain. There is only one element of $\mathrm{fix}_{G_{2}}(\bigcup D)$ that could move $t$, namely the transposition $\tau_{q}=(a_{q},\,b_{q})$. If $\tau_{q}(t)\neq t$ then $\{t,\, \tau_{q}(t)\}$ is also an antichain since $\tau_{q}^{2}$ is the identity permutation. In addition $\{\tau_{q}(t),\,\gamma(t)\}= \tau(\{t,\,\gamma(t)\}$ is an antichain because $\tau$ fixes $\precsim$. Similarly $\{t,\,\gamma(\tau_{q}(t))\}$ is an antichain.

Now assume $\beta_{1}(t)$ and $\beta_{2}(t)$ are in $W$. We first argue that if $\beta_{1}(t)\neq\beta_{2}(t)$ then $\{\beta_{1}(t),\,\beta_{2}(t)\}$ is an antichain. Since $\beta_{1}\in G,\,\beta_{1}=\phi_{1}\circ\psi_{1}$ where $\phi_{1}\in \mathrm{fix}_{G_{1}}(\bigcup D)$ and $\psi_{1}\in \mathrm{fix}_{G_{2}}(\bigcup D)$. Since the only transposition in $\mathrm{fix}_{G_{2}}$ that can possibly move $t$ is the transposition $\tau_{q}=(a_{q},\,b_{q})$ we can assume that $\psi_{1}$ is either $\tau_{q}$ or the identity. Therefore $\beta_{1}(t)=\phi_{1}(t)\circ\psi_{1}(t)$ where $\phi_{1}\in G_{1}$ and fixes $\bigcup D$ pointwise and $\psi_{1}$ is either $\tau_{q}$ or the identity. Similarly $\beta_{2}(t)=\phi_{2}(t)\circ\psi_{2}(t)$ where $\phi_{2}\in G_{1}$ and fixes $\bigcup D$ pointwise and $\psi_{2}$ is either $\tau_{q}$ or the identity. The proof now proceeds by cases.

\textsc{Case 1.} $\beta_{1}(t)=\phi_{1}(t)$ and $\beta_{2}(t)=\phi_{2}(t)$.

Since $\beta_{1}(t)\neq\beta_{2}(t)$ we have $\phi_{1}(C_{q})\neq\phi_{2}(C_{q})$. Therefore there is an $\alpha\in G_{1}$ which fixes $D$ pointwise and for which $\alpha(\{C_{q},\, \gamma(C_{q})\})=\{\phi_{1}(C_{q}),\,\phi_{2}(C_{q})\}$. Hence $\alpha(\{t,\,\gamma(t)\})=\{\phi_{1}(t),\, \phi_{2}(t)\}$ and therefore, since $\alpha$ fixes $\precsim,\,\{\beta_{1}(t),\,\beta_{2}(t)\}$ is an antichain.

\textsc{Case 2.} $\beta_{1}(t)=\phi_{1}\circ\tau_{q}(t)$ and $\beta_{2}(t)=\phi_{2}(t)$ where $\tau_{q}(t)\neq t$.

Assume that $\phi_{1}(C_{q})\prec\phi_{2}(C_{q})$ and that $C_{q}\prec\gamma(C_{q})$. (The arguments in the other possible cases are similar.) Then there is an $\alpha\in G_{1}$ which fixes $D$ pointwise and for which $\alpha(C_{q})=\phi_{1}(C_{q})$ and $\alpha(\gamma(C_{q}))=\phi_{2}(C_{q})$. Hence $\alpha(\tau_{q}(t))= \phi_{1}(\tau_{q}(t)$ and $\alpha(\gamma(t))=\phi_{2}(t)$. Therefore $\alpha(\{\tau_{q}(t), \gamma(t)\})=\{\beta_{1}(t),\,\beta_{2}(t)\}$. which must therefore be an antichain.

The other cases are similar. Since every two element subset of $W$ is an antichain, $W$ is an antichain.

It remains to show that $W$ is infinite. This can be done by showing that for all $\phi_{1},\,\phi_{2}\in \mathrm{fix}_{G_{1}}(\bigcup D)$, $\phi_{1}(t)=\phi_{2}(t)\leftrightarrow\phi_{1}(C_{q})=\phi_{2}(C_{q})$

\begin{theorem}\label{part04:thm4a}
$\mathcal{N}1\models$ 217, but 9 is false (2B6).
\end{theorem}

(Form 217 is given in theorem 3 and form 9 is ``Every Dedekind finite set is finite.''.)

\begin{proof}
The set of atoms is an infinite Dedekind finite set so 9 is false. We shall show that every infinite partially ordered set (poset) in $\mathcal{N}1$ either has an infinite chain or an infinite antichain. Let $(X,\, \leq)$ be a poset with support $E$. If every element of $X$ has support $E$, then $X$ can be well ordered and we are done. Assume there is a $t\in X$ and a $\psi\in \mathrm{fix}(E)$ such that $\psi(t)\neq t$. We may assume that $\psi$ moves only finitely many atoms. Let $F=E\cup K$ be a support of $t$, where $E\cap K=\emptyset$.
\end{proof}

We first argue that there is a $k\in K$ and $\phi\in \mathcal{G}$ which fixes $F-\{k\}$ pointwise and $\phi(t)\neq t$. Choose $\eta \in \mathrm{fix}(E)$ such that $\eta(K)$ is disjoint from $K\cup \{a\in A:\psi(a)\neq a\}$. We have that $\eta(t)\neq t$ because $\psi(t)\neq t$ but $\psi(\eta(t))=t$. Assume $K=\{k_{1},\,k_{2},\,\cdots,\,k_{n}\}$ and for $i=1,2,\,\cdots,\,n$, let $f_{i}$ be the permutation of $A$ that interchanges $k_{i}$ and $\eta(k_{i})$. Then $\prod_{i=1}^{n}f_{i}$ agrees with $\eta$ on a support of $t$, so $\gamma(t)= \eta(t)\neq t$. Therefore, for at least one of the $f_{i}$'s, $f_{i}(t)\neq t$.

Now, for $\gamma,\,\sigma\in \mathrm{fix}(F-\{k\})$, $\gamma(t)=\sigma(t)$ iff $\gamma(k)=\sigma(k)$. It follows that $y=\{\gamma(t):\gamma\in \mathrm{fix}(F-\{k\})\}$ is an infinite subset of $x$ that is in $\mathcal{N}1$. Moreover, $y$ is an antichain, for suppose $\gamma(t)<\sigma(t)$, where $\gamma(t),\sigma(t)\in y$. Let $\rho \in \mathrm{fix}(F-\{k\})$ be a permutation that interchanges $\gamma(t)$ and $\sigma(t)$. Then $\sigma(t)<\gamma(t)$ which is a contradiction.

\section*{NOTE 106}

This note contains results from McCarten\index{Author}{McCarten} \cite{bib:1988} and Schnare\index{Author}{Schnare} \cite{bib:1968} relating to forms [1 BF] (D$T_{0})$, [1 BG] (C$T_{0})$, [1 BH] (T$T_{0})$, and [1 BI] (M$T_{0})$.
\begin{enumerate}
\item[][1 BF] (D$T_{0}$): Every topological space $X$ has a $T_{0}$ subspace that is dense in $X$. ($Y$ is \emph{dense in} $X$ if there is no non-empty open subset $O\subseteq X$ such that $O\cap Y=\emptyset.)$
\item[] [1 BG] (C$T_{0}$): Every topological space $X$ has a $T_{0}$ subspace that is codense in $X$. ($Y$ is \emph{codense in} $X$ if there is no non-empty closed subset $C\subseteq X$ such that $C\cap Y=\emptyset.)$
\item[] [1 BH] (T$T_{0}$): Every topological space $X$ has a $T_{0}$ subspace that is thick in $X$. $(Y$ is \emph{thick in} $X$ if there is no non-empty open and closed (clopen) subset $H\subseteq X$ such that $H\cap Y=\emptyset$.)
\item[][1 BI] (M$T_{0}$): Every topological space has a maximal $T_{0}$ subspace.
\item[1.] Forms $\mathrm{D}T_{0}$ and $\mathrm{C}T_{0}$ are false if ``$T_{0}$'' is replaced by ``$T_{1}$''. (A counter example is the topological space $(\mathbb{R},\,\tau)$, where $\tau =\{G\subseteq \mathbb{R}:G=(a,\,\infty)$, $a\in \mathbb{R} \vee G= \emptyset \vee G=\mathbb{R}\}.)$ McCarten \cite{bib:1988}.
\item[2.] Form $\mathrm{T}T_{0}$ remains equivalent to $\mathrm{AC}$ if $T_{0}$ is replaced by $T_{n}$, for any $n$ such that $T_{n}$ implies $T_{0}$. McCarten \cite{bib:1988}.
\item[3.] $\mathrm{M}T_{0}$ is equivalent to $\mathrm{M}T_{1}$. Schnare \cite{bib:1968}.
\item[4.] Schnare\index{Author}{Schnare} \cite{bib:1968} shows that the statement $\mathrm{MP}$ is false if $\mathrm{P}=T_{2},\,T_{3},\,T_{4}$, regular, completely regular, normal, paracompact, metrizable, linearly disconnected or discrete. (It follows from the lemma that if P is a property that is possessed by every finite discrete space, but not by a countably infinite minimal $T_{1}$ space, then MP is false.)
\end{enumerate}

\section*{NOTE 107}

Good/Tree\index{Author}{Good/Tree}\index{Author}{McCarten} \cite{bib:1995} proves that each of [34 B] through [34 E] implies 34. In this note we supply the details of the converses.

In Kelley\index{Author}{Kelley} \cite{bib:1955}, pages 173--2 there is an outline of a proof that form [34 B] ($\omega_{1}$ with the order topology is not paracompact.) follows from $\mathrm{AC}$. We note that the proof requires only form 34 ($\omega_{1}$ is regular.) and therefore 34 implies [34 B].

Let $\mathcal{C}$ be the open cover of $\omega_{1}$ from Kelley's hint. That is, $\mathcal{C}=\{\{\beta:\beta<\alpha\}:\alpha<\omega_{1}\}$. We show that $\mathcal{C}$ has no open, locally finite refinement by contradiction. Assume that $\mathcal{A}$ is an open, locally finite refinement of $\mathcal{C}$. We define an increasing function $f:\omega\rightarrow \omega_{1}$ by induction: $f(0)=0$. Assume that $f(n)$ is defined. For each $\beta\leq f(n)$, the set $D_{\beta}=\{x\in \mathcal{A}:\beta\in x\}$ is finite since $\mathcal{A}$ is locally finite. It's also true that for each $x\in D_{\beta},\, \sup\{\gamma:\gamma\in x\}$ is less that $\omega_{1}$ since $x$ is a subset of one of the elements of $\mathcal{C}$. Let $g(\beta)$ be the maximum of $\{\sup(x):x\in D_{\beta}\}$.
$g(\beta)$ is also less than $\omega_{1}$. Now apply form 34 to the set $W=\{g(\beta):\beta\leq f(n)\}$. This is a countable set of ordinals each $<\omega_{1}$ and therefore $\sup W<\omega_{1}$. Let $f(n+1)$ be the least ordinal greater than $\sup W$. The important property of the function $f$ is that
\begin{equation*}\tag{$\ast$}
(\forall n\in \omega)(\forall j<n)(\forall x\in \mathcal{A})(f(j)\in x\rightarrow f(n)\not\in x)
\end{equation*}
The sequence $\langle f(n)\rangle_{n\in \omega}$ is strictly increasing so its $\sup$ (call the $\sup\, \beta$) is a limit ordinal. Further, by form 34, $\beta<\omega_{1}$. Since $\mathcal{A}$ is locally finite, there is an (open) neighborhood $y$ of $\beta$ such that $\{x\in \mathcal{A}:x\cap\,y\neq\emptyset\}$ is finite. Since $y$ is open and $\beta$ is a limit ordinal, there is an ordinal $\gamma<\beta$ such that the interval $(\gamma,\,\beta+1)\subseteq y$. For some natural number $k$, this interval must contain $f(n)$ for all $n\geq k$. For each $n\geq k$ we consider the sets $D_{f(n)}=\{x\in \mathcal{A}:f(n)\in x\}$. These sets are pairwise disjoint (by $(\ast)$), each is non-empty and for each element $x\in D_{f(n)},\, x\cap (\gamma,\, \beta+1)\neq\emptyset$. Therefore $\{x\in \mathcal{C}:x\cap\,y\neq\emptyset\}$ is infinite, a contradiction.

To show that 34 implies [34 C] $(\omega_{1}$ with the order topology has the property that every infinite subset has a limit point in $\omega_{1}$.), let $X$ be an infinite subset of $\omega_{1}$. Since $X$ is well ordered and infinite, it must have an initial segment of order type $\omega$. Form 34 implies that each such initial segment must be bounded, so the least upper bound of this set is a limit point of $X$ in $\omega_{1}$.

If $\omega_{1}$ is regular and contains a denumerable discrete family of open sets, then the set of smallest elements in each set would be an infinite set with no limit point, contradicting [34 C]. Thus, 34 implies [34 D] $(\omega_{1}$ with the order topology does not contain a countable discrete family of open subsets.).

The uncountable open cover of $\omega,\,\mathcal{C}=\{\{\beta:\beta<\alpha\}:\alpha<\omega_{1} \}$, described above does not contain a countable subcover. If it did we could construct a countable cofinal sequence in $\omega_{1}$, contradicting 34. Therefore, 34 implies [34 E] $(\omega_{1}$ with the order topology is not Lindel\"{o}f.).

\section*{NOTE 108}

We give the proof that each pair of cardinal numbers has a greatest lower bound (glb) (form 187) is true in $\mathcal{N}12(\aleph_{1})$. (The proof is due to A. Rubin.)\index{Author}{A. Rubin}

\begin{theorem*}
$\mathcal{N}12(\aleph_{1})\models$ every pair of cardinals has a greatest lower bound (187).
\end{theorem*}

\begin{proof}
Let $A^{[\omega]}$ be the set of all 1-1 functions from $\omega$ to $A$, the set of atoms. Let $Q_{0}$ be the set of all partitions of $A^{[\omega]}$ that have the empty set as a support. (Thus, for all $q\in Q_{0}$ if $u,\,v\in q$, then $|u|=|v|.)$ Since $Q_{0}$ can be well ordered, we can construct a subset $Q\subseteq Q_{0}$ such that for each $q\in Q_{0}$ there is exactly one $q^{\prime}\in Q$ such that $|q'|=|q|$. Let $D:Q_{0}\longrightarrow Q$ be such that for each $q\in Q_{0},\,D(q)$ is the unique member $q^{\prime}\in Q$ such that $|q'|=|q|$. Then, define $J$ on $Q_{0}$ such that for each $q\in Q_{0},\, J(q):q\approx D(q)$. Both $D$ and $J$ are in the model. (Clearly $D$ has empty support. The set $F=$\{ $\{f:f$ is a 1--1 function from $q$ onto $D(q)\}:q\in Q_{0}$\} is a well ordered family of sets and, therefore, has a choice function in $\mathcal{N}12(\aleph_{1})$. The function $J$ is a choice function on $F$.)

Suppose $X$ is a set with support $E$. Let $B$ and $C$ be pairwise disjoint denumerable sets of atoms that are disjoint from $E$. Let $H= \mathrm{fix}(E\cup
B)$ and let $T_{X}=\{\{\sigma(x):\sigma\in H\}:x\in X\}$. $T_{X}$ can be well ordered because $H$ leaves $T_{X}$ pointwise fixed. Since the axiom of choice for a well ordered family of sets is true in this model, we can construct a choice function $ch$ for $T_{X}$ with support $E\cup B\cup C$. Define functions $g$ and $h$ such that $g:\omega\approx C$ and $h:A\approx A-E\cup B$. Let $t\in T_{X}$ and $s\in A^{[\omega]}$. Let $\sigma_{s}$ be a permutation in $H$ such that for all $c\in C,\, \sigma_{s}(c)=h\circ s
\circ g^{-1}(c)$. For each $x\in t\in T$, define $j_{t}^{\prime}(x)=\{s\in A^{[\omega]}:x=\sigma_{s}(ch(t))\}$, then $\{j_{t}^{\prime}(x):x\in t\in T\}=q^{\prime}\in Q_{0}$ and $j_{t}^{\prime}:t\approx q^{\prime}$. Define the function $j_{t}$ so that $j_{t}=J(q^{\prime})\circ j_{t}^{\prime}$, then $j_{t}(t)=D(q^{\prime})\in Q$. Suppose $F: T\approx\alpha$, where $\alpha$ is an ordinal. For each $x\in X$, where $x\in t\in T$, define $f(x)=\langle F(t)$, $j_{t}(x)\rangle$, then $f:X  \approx\bigcup_{q\in Q}A_{q}\times q$, where $A_{q}$ is a set of ordinal numbers. Suppose for each $q\in Q,\,|A_{q}|=\alpha_{q}$, then $X \approx\bigcup_{q\in Q}\alpha_{q}\times q$. This construction is independent of the choice of $E,\,B$ and $C$. Now that we have this representation for each set in the model, we are ready to prove that form 187 is true.
\end{proof}

\begin{lemma*}
Suppose, using the construction above, $X \approx\bigcup_{q\in Q}\alpha_{q}\times q$ and $Y\approx  \bigcup_{q\in Q}\beta_{q}\times q$. Then, $X\precsim Y$ iff for all $q,\,\alpha_{q}\leq\beta_{q}$.
\end{lemma*}

\begin{proof}
Suppose $f:X\precsim Y$. Let $E$ be a support for $X$ and choose $B\subseteq A$ disjoint from $E$ such that $E\cup B$ is a support for $X,\,Y$, and $f$. It follows that for each $t_{X}\in T_{X},\, f(t_{X})=t_{Y}\in T_{Y}$ and $f:t_{X}\approx t_{Y}$. Thus, for each $q\in Q,\,f$ induces a 1-1 map from $\alpha_{q}$ into $\beta_{q}$. The converse is clear.
\end{proof}

Finally, if $X \approx\bigcup_{q\in Q}\alpha_{q}\times q$ and $Y \approx\bigcup_{q\in Q}\beta_{q}\times q$ are the representations of $X$ and $Y$, we claim that $Z= \bigcup_{q\in Q}\min(\alpha_{q},\,\beta_{q})\,\times\, q$ has the property that $|Z|= \mathrm{glb}(|X|,\,|Y|)$. Clearly, $|Z|\leq|X|$ and $|Z|\leq|Y|$. Suppose, $|Z_{0}|\leq|X|$ and $|Z_{0}|\leq|Y|$. Also, suppose $Z_{0} \approx\bigcup_{q\in Q}\gamma_{q}\times q$ is the representation of $Z_{0}$. By the lemma, $\gamma_{q}\leq\alpha_{q}$ and $\gamma_{q}\leq\beta_{q}$ for all $q$. Thus, $\gamma_{q}\leq\min(\alpha_{q},\,\beta_{q})$ for all $q$ so it follows that $|Z_{0}|\leq|Z|$.

\section*{NOTE 109}

Below are some definitions for forms involving conditional choice and variations of Rado's lemma.

\begin{definition*}\
\begin{enumerate}
\item[1.] $M$ is a \emph{binary mess on a set}\index{Subject}{binary mess} $S$ if $M$ is a set of functions from finite subsets of $S$ to $\{0,1\}$ such that for all finite $P\subseteq S$, there is a $t\in M$ with domain $P$ and for all $t\in M$, and for all finite $P\subseteq S,\,t|P\in M$.
\item[2.] If $M$ is a binary mess on $S$, a function $f:S\rightarrow\{0,\,1\}$ is \emph{consistent with} $M$ if $f|P\in M$ for all finite $P\subseteq S$.
\item[3.] If $f$ is a function and $R$ a symmetric binary relation, $f$ is $R$-\emph{consistent} if for all $x$ and $y$ in the domain of $f,\,f(x)Rf(y)$.
\end{enumerate}
\end{definition*}

\section*{NOTE 110}

In this note we prove the result from G. Moore\index{Author}{G. Moore} \cite{bib:1982} that form 36 (If $A\subseteq \mathbb{R}^{n}$ and $A \bigcap B$ is countable for every bounded $B$ then $A$ is countable.) implies form 38 ($\mathbb{R}$ is not the union of a countable family of countable sets.)

We prove the contrapositive. Assume that form 38 is false and therefore that $\mathbb{R} = \bigcup_{n\in\omega}A_{n}$ where each $A_{n}$ is countable. Let $f$ be any bijection between $\mathbb{R}$ and the interval $(0,1)$. For each $n\in \omega$ let $A_{n}^{\prime}=f^{\prime\prime}A_{n}+n$ (where for any set $X$ of reals and any real $t,\,X+t=\{x+t:x\in X\})$. Then
\begin{enumerate}
\item[(a)] $A_{n}^{\prime}\subseteq(n,\,n+1)$
\item[(b)] $|A_{n}^{\prime}|=|A_{n}|$ and therefore each $A_{n}$ is countable.
\item[(c)] There is an injection from $\mathbb{R}= \bigcup_{n\in \omega}A_{n}$ to $\bigcup_{n\in\omega}A_{n}^{\prime}(x\mapsto f(x)+n$ where $n$ is the least natural number such that $x\in A_{n}$).
\end{enumerate}

Let $A= \bigcup_{n\in\omega}A_{n}^{\prime}$. It follows from (a) that every bounded subset of $A$ is contained in some finite union $\bigcup_{n<N}A_{n}^{\prime}$. Hence, by (b), every bounded subset of $A$ is countable. By (c), $A$ is uncountable hence form 36 is false.

\section*{NOTE 111}

The relationships between the forms $336(n)$ and $342(n)$ were discussed by C. Montenegro in a talk at the Latin American Logic Symposium in Bogota, 1995. $(336(n)$ is: If $X$ is infinite then there is an infinite $Y\subseteq X$ such that $\{t\subseteq Y:|t|=n\}$ has a choice function and $342(n)$ is $PC(\infty,\,n,\, \infty).)$ He noted that for $n=2$ or 3, $336(n)$ clearly implies $342(n)$. He proved that 336(4) implies 342(4). The question of whether or not 336(5) implies 342(5) is open.

\section*{NOTE 112}

In this note we prove that in $\mathcal{N}41$, form 9 (Dedekind finite and finite are equivalent) is true and form 122 $(C(WO,\,<\aleph_{0}))$ is true.

\begin{lemma*}
Assume that $\psi$ is an order automorphism of $(\mathbb{Q},\,\leq)$ with the property that $\forall a\in \mathbb{Q},\, \lim_{n\rightarrow\infty}\psi^{n}(a)=\infty$.
 (That is, $(\forall a\in \mathbb{Q})(\forall b\in \mathbb{Q})(\exists n\in \omega)(\psi^{n}(a)\geq b)$.) Then there is an order automorphism $\eta$ of $(\mathbb{Q},\,\leq)$ such that for all $a\in \mathbb{Q},\,\eta\psi\eta^{-1}(a)=a+1$.
\end{lemma*}

\begin{proof}
It follows from the hypotheses that $\psi^{n}(0)<\psi^{n+1}(0)$ for each $n\in\omega$ and that $\mathbb{Q}$ is the disjoint union $\mathbb{Q} = \bigcup_{n\in\mathbb{Z}}[\psi^{n}(0),\, \psi^{n+1}(0))$. Let $\sigma$ be an order isomorphism from the interval $[0,\,\psi(0))$ onto the interval $[0,1)$. We define the function $\eta$ on $\mathbb{Q}$ as follows: Assume $a\in[\psi^{n}(0),\,\psi^{n+1}(0))$, then $\eta(a)=\sigma(\psi^{-n}(a))+ n$.

A straightforward argument shows that $\eta\ |\ [\psi^{n}(0),\,\psi^{n+1}(0))$ is an order automorphism from $[\psi^{n}(0),\,\psi^{n+1}(0))$ onto $[n,\,n+1)$ for all $n\in \mathbb{Z}$ and that for $b\in[n,\,n+1)$, $\eta^{-1}(b)=\psi^{n}\sigma^{-1}(b-n)$. It follows from the first assertion that $\eta$ is an order automorphism of $\mathbb{Q}$ and from the second we can calculate for $b\in[n,\,n+1)$:
\begin{align*}
\eta\psi\eta^{-1}(b)&=\eta\psi(\psi^{n}\sigma^{-1}(b-n))=\eta(\psi^{n+1}\sigma^{-1}(b-n))\\
&=\sigma(\psi^{-(n+1)}(\psi^{n+1}\sigma^{-1}(b-n)))+(n+1)\\
&\qquad(\mathrm{since}\ \psi^{n+1}(\sigma^{-1}(b-n))\in[\psi^{n+1}(0), \psi^{n+2}(0)))\\
&=b+1
\end{align*}
This completes the proof of the lemma.
\end{proof}

\begin{corollary*}
Any order automorphism $\psi$ of $(\mathbb{Q},\,\leq)$ satisfying
\begin{equation*}
(\forall a\in \mathbb{Q})(\forall b\in \mathbb{Q})(\exists n\in \omega)(b\leq\psi^{n}(a))
\end{equation*}
has an $n^{th}$ root for every $n\in \omega-\{0\}$. That is, there is an order automorphism $\psi^{\prime}$ of $(\mathbb{Q},\,\leq)$ such that $(\psi^{\prime})^{n}=\psi$.
\end{corollary*}

\begin{proof}
With $\eta$ as in the lemma, $\psi=\eta^{-1}\pi\eta$ where $\pi(a)=a+1$. Assume $n\in \omega-\{0\}$ and let $\pi_{n}(a)=a+\frac{1}{n}$. If we let $\psi^{\prime}=\eta^{-1}\pi_{n}\eta$, then $(\psi^{\prime})^{n}=\psi$.
\end{proof}

\begin{corollary*}
Any order automorphism $\psi$ of $(\mathbb{Q},\,\leq)$ which satisfies $(\forall a\in \mathbb{Q})(\forall b\in \mathbb{Q})(\exists n\in \omega)(\psi^{n}(a)\leq b)$ has an $n^{th}$ root for all $n\in \omega-\{0\}$.
\end{corollary*}

\begin{corollary*}
Every order automorphism of $(\mathbb{Q},\,\leq)$ has an $n^{th}$ root for all $n\in \omega-\{0\}$.
\end{corollary*}

\begin{proof}
Let $\phi$ be an order automorphism of $(\mathbb{Q},\,\leq)$ and for each $a\in \mathbb{Q}$ let $[a]=$\{ $b\in \mathbb{Q}:\exists n\in \mathbb{Z}$ such that $b$ is between $a$ and $\phi^{n}(a)$ \}. If $a$ is a fixed point of $\phi$ then $[a]=\{a\}$ otherwise $[a]$ is an open interval in the sense that if $c<d<e$ and $c$ and $e$ are in $[a]$, then $d\in[a]$ and $[a]$ has no first or last element. Further the set $\{[a]:a\in \mathbb{Q}\}$ is a partition of $\mathbb{Q}$. If $a$ is not a fixed point of $\phi$, then it is also the case that $\phi|[a]$ is an order automorphism of $[a]$ and either $(\forall b\in[a])(\forall c\in [a])(\exists n\in \omega)(c\leq\phi^{n}(b))$ (if $\phi(a)>a)$ or $(\forall b\in[a])(\forall c\in[a])(\exists n\in\omega)(\phi^{n}(b)\leq c)$ (if $\phi(a)<a)$. Therefore, by the previous two corollaries $\phi\ |\
[a]$ has an $n^{th}$ root which we denote by $(\phi|[a])^{\frac{1}{n}}$. Then $\psi^{\prime}$ defined by $\psi^{\prime}(a)=(\phi\ |\ [a])^{\frac{1}{n}}(a)$ is an $n^{th}$ root of $\psi$.
\end{proof}

It will follow from the lemma below that form 9 $(DF=F)$ and form 122 $(C(WO,\,<\aleph_{0}))$ are true in $\mathcal{N}41$.

\begin{lemma*}
If $x\in \mathcal{N}41$ has support $E=B_{0}\cup\cdots \cup B_{k}$ and there is some $\phi\in G$ which fixes $E$ pointwise but does not fix $x$ pointwise, then $x$ has a well orderable, infinite subset.
\end{lemma*}

\begin{proof}
Assume the hypotheses. Then there is some $y\in x$ and $\phi\in G$ such that $\phi$ fixes $E$ pointwise and $\phi(y)\neq y$. If we assume that $y$ has support $E^{\prime}= (B_{0}\cup\cdots B_{k})\cup(B_{k+1}\cup\cdots \cup B_{r})$ then we may also assume without loss of generality that for some $i$ with $k+1\leq i\leq r,\,\phi$ fixes $A-B_{i}$ pointwise.

We note that if $z$ is any element of $\mathcal{N}41$ of the form $z=\psi(y)$ where $\psi\in G$ and $\psi$ fixes $A-B_{i}$ pointwise, then $z$ has support $E^{\prime}$ (because $\psi(E^{\prime})=E^{\prime}$) and $z\in X$ (since $\psi$ fixes $E$ pointwise). We intend to show that there are infinitely many such elements $z$. Assume not. Let $W$ be the set of all such $z$ and assume $|W|=q\in \omega$. By the last corollary, the permutation $\phi$ has an $n^{th}$ root for every positive integer $n$ and we choose an $n$ which is divisible by all of the integers $1,2,\ldots ,q+1$. Let $\eta \in G$ satisfy $\eta^{n}=\phi$ (and we may assume that $\eta$ fixes $A-B_{i}$ pointwise). Clearly, for any $j,\,1\leq j\leq q+1,\,\eta^{j}(y)\neq y$ (otherwise $\eta^{n}(y)=\phi(y)=y$). It follows that for $1\leq j<m\leq q+1,\,\eta^{j}(y)\neq\eta^{m}(y)$ because equality would imply that $y=\eta^{m-j}(y)$. Therefore, $|W|\geq q+1$ a contradiction and $W$ is infinite. Since each $z\in W$ has support $E^{\prime},\,W$ can be well ordered. This complete the proof of the lemma.
\end{proof}

To show that $C(WO,\,<\aleph_{0})$ is true in $\mathcal{N}41$ let $Y$ be a well ordered set of finite sets in $\mathcal{N}41$. Suppose that $E$ is a support of a well ordering of $Y$. Then for each $x$ in $Y,\,E$ is a support of $x$. By the lemma, since $x$ is finite, $E$ is a support of every element of $x$. It follows that if $f$ is any choice function for $Y$ (in the ground model) then $E$ is a support of $f$.

\section*{NOTE 113}

\begin{theorem*}
$C(\infty,\,2)$ (form 88) is false in $\mathcal{N}10$.
\end{theorem*}

\begin{proof}
Let $C$ be the set of all order preserving functions $f:\omega\rightarrow A$ such that rng$(f)$ has an upper bound in $A$ but no least upper bound in $A$. This means that rng$(f)$ is a support and therefore $f$ is in $\mathcal{N}10$.
$C$ is in $\mathcal{N}10$ since it has empty support. Now define an equivalence relation $\sim$ on $C$ by $f\sim g\leftrightarrow(\exists n\in \omega)(f(k)= g(k+n)$ for all $k\in\omega$ or $g(k)=f(k+n)$ for all $k\in \omega$. Let $[f]_{\sim}$ denote the $\sim$ equivalence class of $f$ and let $X$ be the set of all pairs of elements of $\{[f]_{\sim}:f\in C\}$. The set $X$ has no choice function in $\mathcal{N}10$ for suppose that $E$ is a support of such a choice function. Choose a sequence $\langle a_{n}\rangle_{n\in \omega}$ which has order type $\omega$ in the ordering on $A$ and which lies in one of the intervals determined by $E$ in the ordering on $A$. There is a permutation $\psi\in \mathcal{G}$ which fixes $E$ pointwise and which satisfies $(\forall n\in\omega)(\psi(a_{n})=a_{n+1})$. Therefore, if we define $f:\omega\rightarrow A$ by $f(n)=a_{2n}$ and $g:\omega\rightarrow A$ by $g(n)=a_{2n+1}$, then $\psi([f]_{\sim})=[g]_{\sim}$ and $\psi([g]_{\sim})=[f]_{\sim}$. Therefore, no choice function with support $E$ for a set containing $\{[f]_{\sim},\,[g]_{\sim}\}$ exists.
\end{proof}

\section*{NOTE 114}

Definitions from Bell \cite{bib:1988}.

\begin{definition*}
Assume that $X$ and $Y$ are topological spaces.
\begin{enumerate}
\item[1.] If $f$ is a continuous surjection from $X$ onto $Y,\,f$ is \emph{irreducible} if the image of a proper closed subset of $X$ is a proper subset of $Y$.
\item[2.] $X$ is \emph{0-dimensional} if $X$ has a base of clopen sets.
\item[3.] $X$ is \emph{Boolean} if $X$ is compact, Hausdorff and $0$-dimensional.
\item[4.] $X$ is \emph{extremally disconnected} if the closure of any open set is open.
\item[5.] $X$ is $\sigma$-\emph{extremally disconnected}\index{Subject}{extremally disconnected}\index{Subject}{$\sigma$-extremally disconnected} if $X$ is $0$-dimensional and the algebra of clopen subsets of $X$ is $\sigma$-complete (that is, countably complete). (From Morillon\index{Author}{Morillon} \cite{bib:1993})
\item[6.] If $X$\index{Subject}{Boolean topological space}\index{Author}{Bell} is compact Hausdorff, a \emph{Gleason cover}\index{Subject}{Gleason cover} of $X$ is a pair $(E,\,\pi)$ where $E$ is an extremally disconnected compact Hausdorff space and $\pi$ is an irreducible\index{Subject}{irreducible function} continuous surjection $\pi:E\rightarrow X$.
\end{enumerate}
\end{definition*}

\section*{NOTE 115}

Definitions from Diener\index{Author}{Diener} \cite{bib:1989}.

\begin{definition*}\
\begin{enumerate}
\item[1.] A \emph{type} $\tau$ is a function $\tau =(K_{i})_{i\in I}$ where for each $i\in I,\,K_{i}$ is a set (not necessarily finite). If $I=\{i\}$ and $K_{i}=M$ we write $\tau =(M)$.
\item[2.] An \emph{algebra of type} $\tau$\index{Subject}{algebra of type $\tau$} is an ordered pair $\mathbf{A}=(A,\,(f_{i})_{i\in I})$ such that $(\forall i\in I)(f_{i}:A^{K_{i}}\rightarrow A)$.
\item[3.] A \emph{homomorphism} from an algebra \textbf{A} of type $\tau$ to an algebra \textbf{B} of type $\tau$ is closed if the image of every subalgebra of \textbf{A} is a subalgebra of \textbf{B}.
\item[4.] An algebra is a \emph{Peano algebra}\index{Subject}{Peano algebra} if
\begin{enumerate}
\item[a.] The ranges of the operations are pairwise disjoint.
\item[b.] All of the operations are injective.
\item[c.] The elements which are not values of any operation form a generating set.
\end{enumerate}
\item[5.] $\mathfrak{P}(\tau;Z)$ is the Peano algebra of type $\tau$ with generating set $Z$.
\end{enumerate}
\end{definition*}

In Diener \cite{bib:1989} it is shown that for any set $M$, if $\tau =(M)$, then the axiom of choice for families of cardinality $|M|$ (that is, $C(|M|,\,\infty)$) is equivalent to the the statement ``Every homomorphism of $\tau$ algebras is closed.'' See also forms [1 AP], [1 AQ], [1 AR], [67 E] and [67 F].

\section*{NOTE 116}

A proof that form 304 (There does not exist a $T_{2}$ topological space $X$ such that every infinite subset of $X$ contains an infinite compact subset.) is true in $\mathcal{N}1$.

\noindent Let $X$ be an infinite $\mathrm{T}_{2}$ topological space in $\mathcal{N}1$. We need to show that there is an infinite $Y\subseteq X\ (Y\in \mathcal{N}1)$ such that $Y$ has no infinite compact subsets (in $\mathcal{N}1)$. If $X$ is well orderable then it is clear that such a $Y$ exists. Assume that $X$ is not well orderable in $\mathcal{N}1$.

Choose an $x\in X$, a support $E$, an element $a\in A$, and a permutation $\phi\in G$ which fixes $E$ pointwise so that
\begin{enumerate}
\item[(a)] $E$ is a support of $X$ and its topology.
\item[(b)] $E\cup \{a\}$ is a support of $x$ and
\item[(c)] $\phi(x)\neq x$. (And therefore by (b) $\phi(a)\neq a$.)
\end{enumerate}

\begin{lemma*} For all $\psi,\,\eta \in G$ such that $\psi$ and $\eta$ fix $E$ pointwise, $\psi(x)=\eta(x)$ if and only if $\psi(a)=\eta(a)$.
\end{lemma*}

\begin{proof}
If $\psi(a)=\eta(a)$, then $\psi$ and $\eta$ agree on a support of $x$ and therefore $\psi(x)=\eta(x)$. Conversely, if $\psi(a)\neq\eta(a)$, then the permutation $\beta$, defined to be the product of the two transpositions $\beta=(a,\,\psi(a))(\phi(a),\,\eta(a))\in G$, fixes $E$ pointwise, agrees with $\psi$ on a support of $x$, and agrees with $\eta\phi^{-1}$ on a support of $\phi(x)$. It follows that $\beta(x)=\psi(x)$ and that $\beta(\phi(x))=\eta\phi^{-1}(\phi(x))=\eta(x)$. Since $x\neq\phi(x)$ we conclude that $\beta(x)\neq\beta(\phi(x))$, that is $\psi(x)\neq\eta(x)$.
\end{proof}

By the lemma, the set of pairs $F=\{(\psi(x),\,\psi(a)):\psi\in G$ and $\psi$ fixes $E$ pointwise\} is a one to one function (in the model $\mathcal{N}1$) from $Y=_{\mathrm{def}}\{\psi(x):\psi\in G$ and $\psi$ fixes $E$ pointwise\} onto $A-E$.

Since the topology on $X$ is $\mathrm{T}_{2}$ we can choose open sets $C$ and $D$ so that $x\in C,\,\phi(x)\in D$ and $C\cap D=\emptyset$. Since $Y$ can be put in a one to one correspondence with a subset of the atoms in the model, every subset of $Y$ in the model must be finite or cofinite. This means that at least one of $Y\cap C$ or $Y\cap D$ is finite. We may assume that $Y\cap C$ is finite by replacing $C$ by $C\cap\phi^{-1}(D)$ if necessary. Under this assumption we can conclude that $\mathcal{C}=\{\psi(C):\psi\in G$ and $\psi$ fixes $E$ pointwise\} is an open cover for $Y$. Further each element of $\mathcal{C}$ is finite. It follows that for any infinite subset $Z$ of $Y,\,\mathcal{C}$ is an open cover for $Z$ without a finite subcover.

\section*{NOTE 117}

Definitions from Cowen\index{Author}{Cowen} \cite{bib:1990} for forms [14 BD] through [14 BG]. A \emph{hypergraph}\index{Subject}{hypergraph} is a pair $H=\langle V,\, E\rangle$ where $V$ is a set (of vertices) and $E$ is a set of non-empty finite subsets of $V$. $E$ is called the set of edges. $H$ is finite if $V$ is finite. $K=\langle W,\, F\rangle$ is a \emph{subhypergraph}\index{Subject}{subhypergraph} of $H=\langle V,\, E\rangle$ if $W\subseteq V$ and $F\subseteq E$. $W\subseteq V$ is \emph{independent} in $H=\langle V,\, E\rangle$ if no two elements of $W$ belong to the same edge of $H$.
 $W\subseteq V$ is a \emph{vertex cover}\index{Subject}{vertex cover} of $H=\langle V,\, E\rangle$ if each element of $\mathrm{E}$ contains at least one element of $W$. An $n$-\emph{coloring} of $H=\langle V,\, E\rangle$ is a function $f:V\rightarrow\{0,\,1,\,\ldots,\,n-1\}$ such that $|f^{\prime\prime}e|>1$ for all $e\in E$ with $|e|>1$.

\section*{NOTE 118}

Definitions from Galvin/Komj\'{a}th\index{Author}{Galvin/Komj\'{a}th} \cite{bib:1991} for forms [1 AT] through [1 AW].

\begin{definition*}
Assume that $G=(V,\,E)$ is a graph.
\begin{enumerate}
\item[1.] $S\subseteq V$ is \emph{dependent} if $\{x,\,y\}\in S$ for some $\{x,\,y\}\in E$. Otherwise $S$ is \emph{independent}.
\item[2.] $f:V\rightarrow C$ is a \emph{good coloring} if $f^{-1}(c)$ is independent for each $c\in C$.
\item[3.] Let $\kappa=|C|$, then $G$ is $\kappa$ \emph{colorable} if there is a good coloring $f:V\rightarrow C$. ($C$ need not be well orderable.)
\item[4.] The least $\kappa$ for which $G$ is $\kappa$ colorable (if it exists) is the \emph{chromatic number} of $G$.
\item[5.] A good coloring $f:G\rightarrow C$ is \emph{irreducible} if $f^{-1}(c_{1})\cup
f^{-1}(c_{2})$ is dependent whenever $c_{1},\,c_{2}\in C$ and $c_{1}\neq c_{2}$.

\noindent (So a graph admits an irreducible good coloring just in case its vertex set can be partitioned into independent sets such that the union of any two is dependent.)
\item[6.] For any two sets $A$ and $B$ define the graph $G(A,\,B)$ with vertex set $A\times B$ by $(a_{1},\,b_{1})\neq(a_{2},\,b_{2})$ are adjacent if $a_{1}=a_{2}$ or $b_{1}=b_{2}$.
\end{enumerate}
\end{definition*}

\section*{NOTE 119}

Definitions from Shannon \cite{bib:1992} for form [86($\alpha$) A].

\begin{definition*}
Assume $(P,\,\leq)$ is a quasi-order (i.e., reflexive and transitive).
\begin{enumerate}
\item[1.] $x,\,y\in P$ are \emph{incompatible} if there is no $z\in P$ such that $z\leq x$ and $z\leq y$.
\item[2.] $I\subseteq P$ is an \emph{antichain} if any two different elements of $I$ are incompatible.
\item[3.] A partially ordered set $(P,\,\leq)$ is a \emph{tree} iff $\forall x\in P,\,\{y\in P:y\leq x\}$ is well ordered by $\leq$ and $(P,\,\leq)$ has a unique minimal element.
\end{enumerate}
Assume $(P,\,\leq)$ is a tree\index{Subject}{tree}.
\begin{enumerate}
\item[4.] If\index{Subject}{chromatic number}\index{Subject}{incompatible elements}\index{Subject}{$\kappa$ colorable graph}\index{Author}{Shannon} $x\in P$ then $ht(x)$ (the \emph{height} of $x$) is the order type of $\{y\in P:y<x\}$.
\item[5.] $\mathrm{lev}_{\alpha}(P)=\{x\in P:ht(x)=\alpha\}$.
\item[6.] The \emph{height}\index{Subject}{height of a tree} of $P$ is the least $\alpha$ such that $\mathrm{lev}_{\alpha}(P)=\emptyset$.
\item[7.] A \emph{branch}\index{Subject}{branch in a tree} in $(P,\,\leq)$ is a maximal chain.
\item[8.] If $(P,\,\geq)$ is an upside down tree then a \emph{strong antichain in}\index{Subject}{antichain}\index{Subject}{strong antichain} $(P,\,\geq)$ is an antichain that has at most one element from each level of $(P,\,\leq)$.
\end{enumerate}
\end{definition*}

\section*{NOTE 120}

In this note we list some relationships between forms which are provable in $\mathrm{ZF}^{0}$ and are not readily available from table 1. In each case, we give the numbers of the forms, the relationship, and the reference.
\begin{enumerate}
\item[1.] $8+163\leftrightarrow 1$, Brunner\index{Author}{Brunner} \cite{bib:1982a}.
\item[2.] $67+164\leftrightarrow 1$, Brunner \cite{bib:1982a}.
\item[3.] $49+51\leftrightarrow 1$, Morris\index{Author}{Morris} \cite{bib:1969}, note 121.
\item[4.] $64+133\leftrightarrow 1$, clear.
\item[5.] $88+133\leftrightarrow 1$, Brunner \cite{bib:1982a}.
\item[6.] $15+91\leftrightarrow 1$, Brunner \cite{bib:1982a}.
\item[7.] $8+191+193\leftrightarrow 1$, Blass\index{Author}{Blass} \cite{bib:1979}, note 60.
\item[8.] $191+192\leftrightarrow 1$, Blass \cite{bib:1979}, note 60.
\item[9.] $30+126\rightarrow 8$, K. Keremedis\index{Author}{Keremedis}.
\item[10.] $106+154\leftrightarrow 43$, Brunner \cite{bib:1983c}\cite{bib:1984b}.
\item[11.] $14+297\leftrightarrow 50$, Banaschewski\index{Author}{Banaschewski} \cite{bib:1988}, Bell\index{Author}{Bell}
\cite{bib:1988}.
\item[12.] $14+299\leftrightarrow 50$, Bell \cite{bib:1988}.
\item[13.] $14+300\leftrightarrow 50$, Bell \cite{bib:1988}.
\item[14.] $14+301\leftrightarrow 50$, Bell \cite{bib:1988}.
\item[15.] $14+302\leftrightarrow 50$, Bell \cite{bib:1988}.
\item[16.] $30+118\rightarrow 15$, Krom\index{Author}{Krom} \cite{bib:1986}.
\item[17.] $127+133\rightarrow 132$, Brunner \cite{bib:1984b}.
\item[18.] $106+127\rightarrow 64$, Brunner \cite{bib:1984b}.
\item[19.] $17+133\rightarrow 134$, Brunner \cite{bib:1984b}.
\item[20.] $9+70\rightarrow 63$, clear.
\item[21.] $39+41\rightarrow 135$, van Douwen\index{Author}{Douwen} \cite{bib:1980}.
\item[22.] $128+130\rightarrow 18$, Brunner \cite{bib:1984f}.
\item[23.] $10+163\rightarrow 231$, note 123.
\item[24.] $\neg 206\rightarrow 253$, Howard\index{Author}{Howard} \cite{bib:1975}.
\item[25.] DUSS $+131\leftrightarrow 8$, Keremedis \cite{bib:1996d} and note 132.
\item[26.] $52+286\leftrightarrow 1$, H. Rubin/J. Rubin\index{Author}{H. Rubin/ J. Rubin} \cite{bib:1985}, p179.
\item[27.] $30+296\rightarrow 295$, Pincus\index{Author}{Pincus} \cite{bib:1997}.
\item[28.] $30+144+294\leftrightarrow 1$, clear. (See note 25 for definitions.) \item[29.] $62+99\leftrightarrow 14$, note 33.
\item[30.] $163+10\rightarrow 165$, Brunner \cite{bib:1982a}.
\item[31.] $62+67\leftrightarrow 1$, clear.
\item[32.] $122+328\leftrightarrow 40$, clear.
\item[33.] $327+328\rightarrow 322$, clear.
\item[34.] $122+330\leftrightarrow 165$, clear.
\item[35.] $327+330\rightarrow 324$, clear.
\item[36.] $62+329\leftrightarrow 60$, clear.
\item[37.] $333+334\leftrightarrow 1$, Keremedis\index{Author}{Keremedis} \cite{bib:1996a}.
\item[38.] $8+(\forall n\geq 2)335(n)\leftrightarrow 1$,\index{Author}{Brunner} Keremedis \cite{bib:1996c}.
\item[39.] $61+67\leftrightarrow 218$, Bleicher\index{Author}{Bleicher} \cite{bib:1965}.
\item[40.] $336(n)\rightarrow 342(n)$ for $n=2,3$ and 4, note 111.
\item[41.] $67+122\leftrightarrow 1$. (See [67 B].)
\item[42.] $67+346\leftrightarrow 1$, note 127.
\item[43.] $20+97\rightarrow 359$, Higasikawa\index{Author}{Higasikawa} \cite{bib:1995}.
\item[44.] $20+3\rightarrow 359$, Higasikawa \cite{bib:1995}.
\item[45.] $14+253\leftrightarrow 1$, Howard\index{Author}{Howard} \cite{bib:1975}.
\item[46.] $40+152\leftrightarrow 1$, clear.
\item[47.] $151+152\leftrightarrow 1$, clear.
\item[48.] $214+329\rightarrow 67$, clear.
\item[49.] $32+131\leftrightarrow 8$, note 132.
\item[50.] DUCCC $+131\leftrightarrow 8$, Keremedis \cite{bib:1996d} and note 132.
\item[51.] DUSC $+131\leftrightarrow 8$, Keremedis \cite{bib:1996d} and note 132.
\item[52.] SCDUL $+131\leftrightarrow 8$, Keremedis \cite{bib:1996d} Jand note 132.
\item[53.] $131+350\leftrightarrow 126$, Keremedis \cite{bib:1996d} and note 132.
\item[54.] FSCB $+131\leftrightarrow 126$, Keremedis \cite{bib:1996d} and note 132.
\item[55.] $10+106\leftrightarrow 43$, Blass\index{Author}{Blass} \cite{bib:1979}.
\item[56.] $70+344\leftrightarrow 123$, Morillon\index{Author}{Morillon} \cite{bib:1988}.
\end{enumerate}


\section*{NOTE 121}

This note contains an argument that the order extension principle (form 49) together with ``Every linear ordering has a cofinal well ordered subset.'' (form 51) implies AC (form 1). This result was announced in Morris\index{Author}{Morris} \cite{bib:1969}.

Assume that $Y$ is a set of non-empty sets. Let $P$ be the set of choice functions $f$ such that the domain of $f$ is a subset of $Y$. We define the following partial ordering on $P$:
\begin{equation*}
f\preceq_{0} g\Leftrightarrow \mathrm{dom}(f)\subsetneq \mathrm{dom}(g)\vee f=g.
\end{equation*}
Extend $\preceq_{0}$ to a linear ordering $\preceq$ on $P$. Let $\{f_{\alpha}:\alpha\in\beta\}$ be a cofinal subset of the linear ordering $(P,\,\preceq)$ where $f_{\alpha}\prec f_{\gamma}\leftrightarrow\alpha\in\gamma$. ($\in$ is the usual well ordering on the ordinals.) The set $W= \bigcup_{\alpha\in\beta} \mathrm{dom}(f_{\alpha})$ has a choice function $F$ defined by $F(x)=f_{\alpha}(x)$ where $\alpha$ is the least ordinal for which $x\in \mathrm{dom}(f_{\alpha})$. We complete the proof by showing that $W=Y$. If not, then there is some $x\in Y-W$. Let $t$ be any element of $x$, then $G=F\cup \{(x,\,t)\}$ is in $P$ and has the property that $\forall\alpha\in\beta,\,f_{\alpha}\prec_{0}G$ and hence, since $\preceq$ extends $\preceq_{0},\,f_{\alpha}\prec G$. This contradicts the cofinality of $\{f_{\alpha}:\alpha\in\beta\}$.

\section*{NOTE 122}

In $\mathcal{N}43$, both $C(\aleph_{0},2)$ (form 80) and the Multiple Choice Axiom $(MC(\infty,\,\infty)$, form 67) are false.

\noindent $C(\aleph_{0},2)$ is false in $\mathcal{N}43$: For each $n\in \omega$, let $w_{n}=(a_{0}^{n},\,a_{1}^{n},\,\ldots,\, a_{n}^{n})$ be a fixed well ordering of $P_{n}$. (The function choosing these well orderings will not be in the model.) Let $C_{0}^{n}=$\{$\phi(w_{n}):\phi$ is an even permutation of $P_{n}$\} and $C_{1}^{n}=\{\phi(w_{n}):\phi$ is an odd permutation of $P_{n}$\}. Then for every $n\in \omega$, $\{C_{0}^{n},\,C_{1}^{n}\}$ is in $\mathcal{N}43$ and has empty support. In fact if $\psi|P_{n}$ is even then $\psi$ fixes $\{C_{0}^{n},\,C_{1}^{n}\}$ pointwise and if $\psi|P_{n}$ is odd then $\psi$ interchanges $C_{0}^{n}$ and $C_{1}^{n}$.

If $E$ is any support, then there is a $\psi\in G$ that fixes $E$ pointwise and an $n\in\omega$ such that $\psi|P_{n}$ is odd. Since such a $\psi$ fixes $\{C_{0}^{n},\,C_{1}^{n}\}$ but moves both its elements, $E$ is not a support of a choice function on $\{\{C_{0}^{n},\,C_{1}^{n}\}:n\in \omega\}$.

\noindent $MC(\infty,\,\infty)$ is false in $\mathcal{N}43$: Let $F$ be the set of all choice functions $f:\omega\rightarrow A$ such that $\forall n\in\omega,\,f(n)\in P_{n}$. Suppose that $E$ is a support of a multiple choice function $H$ on $\mathcal{P}(F)-\{\emptyset\}$ in $\mathcal{N}43$. Let $k= \sup\{|E\cap P_{n}|:n\in \omega\}$ and define $T\subseteq F$ by $T=\{f\in F:(\forall i>k)(f(i)\in P_{n}-E)\}$. Since $E$ is a support of $T$ and of $H,\, E$ is a support of $H(T)$ which must be a finite, non-empty subset of $T$. Fix $f$ in $H(T)$ then for every $g\in T$ for which $g(i)=f(i)$ for $i\leq k$, there is a $\phi\in G$ which fixes $E$ pointwise such that $\phi(f)=g$. Therefore $g\in H(T)$. Since there are infinitely many such $g$'s, $H(T)$ is infinite. This is a contradiction.

\section*{NOTE 123}

We shall prove that form 10 $(C(\aleph_{0},\,<\aleph_{0}))\ +$ form 163 (Every non-well orderable set has an infinite Dedekind finite subset.) implies form 231 $(UT(WO,\,WO,\,WO))$ and also that form 133 (Every set is either well orderable or has an infinite amorphous set.) implies form 231. (It is shown in Brunner\index{Author}{Brunner} \cite{bib:1982a} that 10 + 163 implies 165 $(C(WO,\,WO)).)$

Suppose $X$ is a well ordered set of pairwise disjoint well orderable sets. Assume $Y=\bigcup X$ cannot be well ordered. Form 163 implies that $Y$ has an infinite Dedekind finite set $Z$. Let $T=\{w\cap\,z:w\in X\wedge w\cap z\neq\emptyset\}$ For each $w\in Z,\,w\cap Z$ is a well ordered Dedekind finite set so it must be finite. If $T$ is finite, then $Z$ is a finite union of finite sets and is therefore finite. This contradicts the definition of $Z$. If $T$ is infinite, then $T$ is well ordered because $X$ is well ordered. Therefore $T$ must have a denumerable subset, $U$. $U$ is a denumerable set of finite sets, so by 10, $U$ has a choice set. This choice set is a denumerable subset of $Z$. This contradicts the definition of $Z$. Thus, the union of $X$ can be well ordered.

Since form 133 (every set is either well orderable or has an infinite amorphous subset) implies form 10 (Brunner \cite{bib:1984b}) and 133 implies 163 (clear), it follows that 133 implies 231.

\section*{NOTE 124}

A proof that form 314 (Every permutation on a non-empty set can be expressed as a product of two reflections.) implies form 119 ($C(\aleph_{0}$,uniformly orderable with order type of the integers)).

Assume that $B=\{A_{i}:i\in \omega\}$ is a denumerable set of denumerable sets and $f$ is a function defined on $\omega$ such that $\forall i\in \omega,\,f(i)$ is an ordering of $A_{i}$ with order type that of the integers. We want to construct a choice function $g$ for $B$ so we may assume that the elements of $B$ are pairwise disjoint.

Let $\eta$ be the permutation of $\bigcup B$ defined by $\eta(a)=$ the immediate successor of $a$ under the ordering $f(A_{i})$ where $A_{i}$ is the unique element of $B$ such that $a\in A_{i}$. By form 314, $\eta$ is the product of two reflections: $\eta =\phi\circ\psi$.

Fix an $i\in \omega$ and for each permutation $\beta$ of $\bigcup B$, let $\beta^{\prime}=\beta\ |\ A_{i}$. Then $\eta^{\prime}=\phi^{\prime}\circ\psi^{\prime}$ and $\phi^{\prime}$ and $\psi^{\prime}$ are reflections. We complete the proof by showing that exactly one of $\phi^{\prime}$ or $\psi^{\prime}$ has exactly one fixed point. (This fixed point will be $g(A_{i})$.)

Since $\phi^{\prime}$ and $\psi^{\prime}$ are both reflections, $\phi^{\prime}$ and $\psi^{\prime}$ can both be written as a product of disjoint cycles, $\phi'=\prod_{i\in I}(s_{i},\,t_{i})$ and $\psi^{\prime}=\prod_{j\in J}(c_{j},\,d_{j})$. Assume that $A_{i}=\{\ldots,\,a_{-1},\,a_{0},\,a_{1},\,a_{2},\,\ldots\}$ where the elements are listed in the order determined by the ordering $f(A_{i})$. Then $\eta^{\prime}$ is the infinite cycle $\eta^{\prime}= (\ldots,\,a_{-1},\, a_{0},\,a_{1},\,a_{2},\,\ldots)$. Let $r$ be the smallest positive integer such that for some integer $k$ the transposition $(a_{k},\,a_{k+r})$ is one of the $(s_{i},\,t_{i})$ or one of the $(c_{j},\,d_{j})$

We first show that $r=1$. Assume that $r>1$ and that $(a_{k},\,a_{k+r})$ is one of the $(c_{j},\,b_{j})$. (The argument is similar if $(a_{k},\,a_{k+r})$ is one of the $(s_{i},\,t_{i})$.) This means that $\psi^{\prime}(a_{k})=a_{k+r}$. Since $\eta^{\prime}(a_{k})=a_{k+1}$ (and $r\neq 1)$, $\phi^{\prime}(a_{k+r})$ must be $a_{k+1}$, hence the transposition $(a_{k+1},\,a_{k+r})$ must be one of the $(s_{i},\,t_{i})$. This contradicts the definition of $r$.

Assume that $(a_{k},\,a_{k+1})$ is one of the $(c_{j},\,d_{j})$. (The argument is similar if $(a_{k},\, a_{k+1})$ is one of the $(s_{i},\,t_{i}).)$ This means that $\psi^{\prime}(a_{k})=a_{k+1}$. Since $\eta^{\prime}(a_{k})= a_{k+1}$ we conclude that $a_{k+1}$ must be a fixed point of $\phi^{\prime}$. Further, since $\psi^{\prime}(a_{k+1})=a_{k}$ and $\eta^{\prime}(a_{k+1})=a_{k+2}$ we conclude that $\phi^{\prime}(a_{k})=a_{k+2}$, that is, $(a_{k},\,a_{k+2})$ is one of the $(s_{i},\,t_{i})$. By this last conclusion, $\psi^{\prime}(a_{k+2})=a_{k}$. It follows that $\psi^{\prime}(a_{k-1})=a_{k+2}$, that is $(a_{k-1},\, a_{k+2})$ is one of the $(c_{j},\,d_{j})$. Continuing in this way, we see that for every natural number $j,\,\psi^{\prime}$ interchanges $a_{k-j}$ and $a_{k+j+1}$ and $\phi^{\prime}$ interchanges $a_{k-j}$ and $a_{k+j+2}$. Therefore, $\psi^{\prime}$ has no fixed points and $\phi^{\prime}$ has exactly one fixed point.

\section*{NOTE 125}

In this note we survey forms of the Hahn-Banach Theorem\index{Subject}{Hahn-Banach theorem} discussed in Buskes\index{Author}{Buskes}
\cite{bib:1993}.

\begin{definition}\label{part04:def1}
Let $E$ be a vector space over $\mathbb{R}$, and let $p$ be a mapping from $E$ to $\mathbb{R}$.
\begin{enumerate}
\item[(a)] $p$ is called \emph{sublinear}\index{Subject}{sublinear function} if $p(x+y)\leq p(x)+p(y)$, for all $x,\,y\in E$, and $p(rx)=rp(x)$, for all $x\in E$ and all $r\in \mathbb{R}^{+}$.
\item[(b)] $p$ is called a \emph{seminorm} if $p$ is sublinear and $p(rx)=|r|p(x)$, for all $x\in E$ and all $r\in \mathbb{R}$.
\item[(c)] $p$ is called a \emph{norm} if $p$ is a seminorm and $p(x)=0$ iff $x=0$, for all $x\in E$.
\item[(d)] $p$ is called \emph{superlinear} if $p(x)+p(y)\leq p(x+y)$ for
all $x,\,y\in E$, and $p(rx)=rp(x)$, for all $x\in E$ and all $r\in \mathbb{R}^{+}$.
\end{enumerate}
\end{definition}

\begin{unsec}\label{part04:se1}
(Form 52) Assume $E$ is a vector space over $\mathbb{R}$ and $p:E\rightarrow \mathbb{R}$ is a sublinear function. Suppose $E_{0}$ is a subspace of $E$ and $f_{0}$ is a linear function from $E_{0}$ to $\mathbb{R}$ such that $f_{0}(x)\leq p(x)$, for all $x\in E_{0}$. Then there is a linear function $f:E\rightarrow \mathbb{R}$ that extends $f_{0}$ and $f(x)\leq p(x)$, for all $x\in E$.
\end{unsec}

\begin{unsec}\label{part04:se2}
Assume $E$ is a vector space over $\mathbb{R}$ and $p$ is a norm on E. Suppose $E_{0}$ is a subspace of $E$ and $f_{0}$ is a continuous linear function from $E_{0}$ to $\mathbb{R}$. Then there is a continuous linear function $f:E\rightarrow \mathbb{R}$ that extends $f_{0}$ and $\Vert f_{0}\Vert=\Vert f\Vert$. (For any subspace $E_{0}\subseteq E$ and any linear function $f_{0}:E_{0}\rightarrow \mathbb{R}$,
\begin{equation*}
\Vert f_{0}\Vert=sup\{|f_{0}(x)|:x\in E_{0}\wedge p(x)\leq 1\}.\text{)}
\end{equation*}
\end{unsec}

\begin{unsec}\label{part04:se3}
Assume $E$ is a vector space over $\mathbb{R},\,p$ a sublinear function on $E$, and $q$ a superlinear function on E. If $q(x)\leq p(x)$ for all $x\in E$, then there exists a linear function $f: E\rightarrow \mathbb{R}$ such that $q(x)\leq f(x)\leq p(x)$ for all $x\in E$.
\end{unsec}

\begin{unsec}\label{part04:se4}
(Complex form of $\mathbf{H\text{-}B\ 1}$) Assume $E$ is a vector space over $\mathbb{C}$ and $p:E\rightarrow \mathbb{R}$ a seminorm (a sublinear function for which $p(\alpha x)=|\alpha|p(x)$ for all $x\in E$ and all $\alpha\in \mathbb{C}$). Let $E_{0}$ be a subspace of $E$ and $f_{0}:E_{0}\rightarrow \mathbb{C}$ be linear such that $|f_{0}(x)|\leq p(x)$ for all $x\in E$. Then there is a linear function $f:E\rightarrow \mathbb{C}$ such that $|f(x)|\leq p(x)$ for all $x\in E$.
\end{unsec}

\begin{unsec}\label{part04:se5}
(Complex form of $\mathbf{H\text{-}B\ 2}$)\index{Subject}{seminorm}\index{Subject}{superlinear function} (Change ``$\mathbb{R}$'' in $\mathbf{H}\text{-}\mathbf{B\ 2}$ to ``$\mathbb{C}$''. $\mathbf{H\text{-}B\ 5}$ is equivalent to $\mathbf{H\text{-}B\ 2}$. Let $f_{0}=f_{1}+if_{2}$, where $f_{1}$ and $f_{2}$ are real continuous function from $E_{0}$ to $\mathbb{R}$. Let $f_{1}^{\ast}$ be an extension of $f_{1}$ to $E$ so that $\Vert f_{1}^{\ast}\Vert\leq\Vert f_{0}\Vert$. Then the desired extension $f$ of $f_{0}$ can be defined as $f(x)=f_{1}^{\ast}(x)-if_{1}^{\ast}(x)$ for all $x\in E$.)
\end{unsec}

\begin{definition}\label{part04:def2}
A metric space $E$ is \emph{spherically complete} if the set of its closed
balls has the binary intersection property (if the intersection of each pair of balls is non-empty, the intersection of the whole set is non-empty).
\end{definition}

\begin{unsec}\label{part04:se6}
(Ingleton' s Theorem) Assume $K$ is a spherically complete, non-Archimedean valued field. Let $E$ be a normed vector space over $K$ and let $E_{0}$ be a subspace of $E$.

\noindent Let $f_{0}:E_{0}\rightarrow K$ be a continuous linear function. Then there exists a continuous linear function $f:E\rightarrow K$ that extends $f_{0}$ and for which $\Vert f\Vert=\Vert f_{0}\Vert$.
\end{unsec}

\begin{definition}\label{part04:def3}
Let $(E,\,\Vert\ \Vert)$ be a normed vector space.
\begin{enumerate}
\item[(a)] The norm\index{Subject}{norm} is said to be \emph{Gateaux differentiable}\index{Subject}{Gateaux differentiable} at $x\in E$ if for all $y\in E$,
\begin{equation*}
\lim\limits_{t\rightarrow 0}\frac{\Vert x+ty\Vert-\Vert x\Vert}{t}
\end{equation*}
exists.
\item[(b)] $(E,\,\Vert\ \Vert)$ is said to be \emph{uniformly convex}\index{Subject}{uniformly convex} if for every $\epsilon>0$ there is a $\delta>0$ such that for all $x,\,y\in E$ with $\Vert x\Vert=\Vert y\Vert=1,\,\Vert x-y\Vert\geq\epsilon$ implies $\frac{1}{2}\Vert x+y\Vert\leq 1-\delta$.
\end{enumerate}
\end{definition}

\begin{unsec}\label{part04:se7}
(Ishihara's Theorem) Let $E_{0}$ be a subspace of a uniformly convex Banach space $E$ with a Gateaux differentiable norm. Let $f_{0}:E_{0}\rightarrow \mathbb{R}$ be a non-zero continuous linear function. Then there is a unique linear continuous extension $f$ from $E$ to $\mathbb{R}$ such that $\Vert f\Vert=\Vert f_{0}\Vert$.
\end{unsec}

\begin{unsec}\label{part04:se8}
Suppose $E$ is a normed vector space in which the unit ball is weakly compact and the norm in $E$ is Gateaux differentiable. Let $E_{0}$ be a subspace of $E$. Then every continuous linear function $f_{0}:E_{0}\rightarrow \mathbb{R}$ can be uniquely extended to a continuous linear function $f: E\rightarrow \mathbb{R}$ such that $\Vert f\Vert=\Vert f_{0}\Vert$.
\end{unsec}

\begin{unsec}\label{part04:se9}
(Mazur-Orlicz Theorem) Assume $E$ is a vector space over $\mathbb{R}$ and $p:E\rightarrow \mathbb{R}$ is a sublinear function. Let $T$ be a set and let $\phi$ and $\alpha$ be functions such that $\phi:T\rightarrow E$ and $\alpha:T\rightarrow \mathbb{R}$. Then the following are equivalent.
\begin{enumerate}
\item[\emph{(1)}] There exists a linear $f:E\rightarrow \mathbb{R}$ such that
\begin{enumerate}
\item[\emph{(i)}] $f\leq p$ and
\item[\emph{(ii)}] $\alpha(t)\leq f(\phi(t))$ for all $t\in T$.
\end{enumerate}
\item[\emph{(2)}] For all $k\in \mathbb{N}$, for all $r_{1},\,r_{2},\cdots,\,r_{k}\in \mathbb{R}^{+}$, and for all $t_{1},\,t_{2},\,\cdots,\,t_{k}\in T$, we have
\begin{equation*}
\sum\limits_{n=1}^{k}r_{n}\alpha(t_{n})\leq p\bigg(\sum\limits_{n=1}^{k}r_{n}\phi(t_{n})\bigg).
\end{equation*}
$(\mathbf{H\text{-}B\ 9}$ \emph{implies} $\mathbf{H\text{-}B\ 1}$ \emph{and} $\mathbf{H\text{-}B\ 3}.)$
\end{enumerate}
\end{unsec}

The last three forms are examples of Hahn-Banach theorems in which the range of the function that is to be extended may not be commutative.

\begin{unsec}\label{part04:se10}
(Kantorovich) Suppose $E$ is an ordered vector space and $E_{0}$ a cofinal subspace. Let $F$ be an order complete vector space. Then each positive linear map from $E_{0}$ to $F$ can be extended to a positive linear map from $E$ to $F$.
\end{unsec}

\begin{unsec}\label{part04:se11} (Lipecki-Plachky-Thomsen, Luxemburg-Schep) Suppose $E$ is a vector lattice and $E_{0}$ is a cofinal vector sublattice. Let $F$ be an order complete vector lattice. Then each vector lattice homomorphism from $E_{0}$ to $F$ can be extended to a vector lattice homomorphism from $E$ to $F$.
\end{unsec}

\begin{definition}\label{part04:def4}\
\begin{enumerate}
\item[(a)] A \emph{Banach algebra}\index{Subject}{Banach algebra} is a complex algebra together with a complete norm satisfying the condition $\Vert xy\Vert\leq\Vert x\Vert\Vert y\Vert$, for all $x$ and $y$. If the algebra contains an identity element $e,\,\Vert e\Vert=1$.
\item[(b)] A $C^{\ast}$-\emph{algebra}\index{Subject}{$C^{\ast}$-algebra} is a Banach algebra with an involution $x\mapsto x^{\ast}$ on it satisfying
\begin{enumerate}
\item[(i)] $x^{\ast\ast}=x$ for all vectors $x$.
\item[(ii)] $(ax+bx)^{\ast}=\overline{a}x^{\ast}+\overline{b}y^{\ast}$ for all vectors $x$ and $y$ and all complex numbers $a$ and $b$.
\item[(iii)] $(xy)^{\ast}=y^{\ast}x^{\ast}$ for all vectors $x$ and $y$.
\item[(iv)] $\Vert x^{\ast}x\Vert=\Vert x\Vert^{2}$ for all vectors $x$.
\end{enumerate}
\end{enumerate}
\end{definition}

\begin{unsec}\label{part04:se12}
(Arveson) Suppose $E$ is a $C^{\ast}$-algebra and $E_{0}$ is a cofinal $\ast$-invariant subspace of E. Let $H$ be a Hilbert space and $B(H)$ the space of all bounded operators on H. Then each completely positive linear map from $E_{0}$ to $B(H)$ extends to a completely positive linear map from $E$ to $B(H)$.
\end{unsec}

For additional results see Buskes\index{Author}{Buskes} \cite{bib:1993} and his references.

\section*{NOTE 126}

We give definitions for form [43 B] from Goldblatt\index{Author}{Goldblatt} \cite{bib:1985}.

Let $\Phi$ be a class of formulas in a language that includes the language of the sentential calculus and suppose the deducibility relation, $\vdash$, satisfies the following:
\begin{enumerate}
\item[D1.] $\emptyset \vdash\phi$, for every tautology $\phi$.
\item[D2.] $\Gamma \cup\{\phi\}\vdash\psi$ iff $\Gamma\vdash\phi\rightarrow\psi$, where $\Gamma\subseteq\Phi$.
\item[D3.] If $\Gamma\vdash\phi$ and $\Gamma\subseteq\Delta$, then $\Delta\vdash\phi$, where $\Gamma,\, \Delta\subseteq\Phi$.
\end{enumerate}

\begin{definition}\label{part04:def1a}
$\Gamma\subseteq\Phi$ is \emph{finitely consistent} if every finite subset of $\Gamma$ is consistent.
\end{definition}

\begin{definition}\label{part04:def2a}
An ordered pair $(\Pi,\,\psi)$ is called an \emph{inference} if $\Pi\subseteq\Phi$ and $\psi\in\Phi$.
\end{definition}

\begin{definition}\label{part04:def3a}
If $\Gamma\subseteq\Phi,\,\Gamma$ \emph{respects} the inference $(\Pi,\,\psi)$ if
\begin{equation*}
(\Gamma\vdash\phi,\,\forall\phi\in\Pi)\ \text{implies}\ \Gamma\vdash\psi.
\end{equation*}
\end{definition}

\begin{definition}\label{part04:def4a}
If $\Gamma\subseteq\Phi,\,\Gamma$ \emph{decides} the inference $(\Pi,\,\psi)$ if
\begin{equation*}
\text{either}\ \Gamma\vdash\psi\ \text{or}\ \exists\phi\in\Pi\ \text{such that}\ \Gamma\vdash\neg\phi.
\end{equation*}
(We say $\Gamma$ respects (decides) a set of inferences $I$ if $\Gamma$ respects (decides) each inference in $I$.)
\begin{enumerate}
\item[[43 B{]}] Countable Henkin Principle: Let $\Phi$ be the class of formulas in a language that includes the language of the sentential calculus. Let $I$ be a countable set of inferences, and $\Gamma$ a consistent subset of $\Phi$ such that $\Gamma \cup\{\psi\}$ respects $I$ for each $\psi\in\Phi$. Then for any $\phi\in\Phi$, if $\phi$ is consistent with $\Gamma$ then $\Gamma$ has a finitely consistent extension that contains $\phi$ and decides $I$.
\end{enumerate}

Goldblatt\index{Author}{Goldblatt} \cite{bib:1985} proves that [43 B] is equivalent to the principle of dependent choices (form 43) and Goldblatt \cite{bib:1984} shows how [43
B] may be used to prove completeness and omitting-types theorems for countable first-order languages.
\end{definition}

\section*{NOTE 127}

Forms [1 BZ] (Vector space multiple choice) and 346 (Vector space Kinna-Wagner principle) were suggested by K. Keremedis. It is clear that [1 BZ] implies 346. In this note we prove that 346 implies the Kinna-Wagner principle $KW(\infty, < \aleph_{0})$ (form [62 E]). Since the axiom of choice is implied by the conjunction of forms 62 and 67, we obtain a proof that $346+67$ implies the axiom of choice. (Form 62 is $C(\infty,\,<\aleph_{0})$ and 67 is the axiom of multiple choice.) Similarly, since [1 BZ] implies form 67, we obtain the result: [1 BZ] implies the axiom of choice.

Let $X=\{y_{i}:i\in K\}$ be a family of finite sets. For each $y_{i}$ let $U_{i}$ be the real vector space $\mathbb{R}^{y_{i}}$ with pointwise addition and scalar multiplication. (If $y_{i}= \{a_{1},\,\ldots,\, a_{n}\}$ we could think of $U_{i}$ as being the set of all formal sums $k_{1}a_{1}+\cdots+k_{n}a_{n}$ where the $k_{i}$'s are real.) Let $S_{i}$ be the subspace \{ $g\in U_{i}:g$ is constant\}. (Or in terms of formal sums, $S_{i}$ is all formal sums $ka_{1}+\cdots+ka_{n}$.) Let $V_{i}$ be the quotient space $V_{i}=U_{i}/S_{i}$. That is, $V_{i}$ consists of all equivalence classes $[g]$ of elements of $U_{i}$ under the relation $g\sim f\Leftrightarrow g-f\in S_{i}$. By form 346 there is a family $\{F_{i}:i\in K\}$ such that for each $i\in K,\,F_{i}$ is an independent subset of $V_{i}$. Since $U_{i}$ is finite dimensional $F_{i}$ must be finite. Say $F_{i}=\{b_{1},\,\ldots,\,b_{r}\}$ then since $F_{i}$ is independent, the element $w_{i}=b_{1}+\cdots+b_{r}$ is not zero. The vector $w_{i}=[g]$ for some $g\in U_{i}$. Assume $f\in[g]$. Then if $a,\,a^{\prime}\in y_{i}$ and $g(a)\leq g(a^{\prime})$ it follows from the fact that $f-g$ is constant that $f(a)\leq f(a^{\prime})$. Therefore the set $K_{i}=\{a\in y_{i}:g(a)$ is minimum among the numbers $g(a^{\prime})$ for $a^{\prime}\in y_{i}$\} is independent of the choice of $g\in w_{i}$. It is also true that $K_{i}\neq y_{i}$. This follows from the fact that $w_{i}\neq 0$ which implies that $g$ is not constant. The family $\{K_{i}:i\in K\}$ is therefore a Kinna-Wagner function for X.

\section*{NOTE 128}

In this note we prove that in $\mathcal{N}51$, form 64 (There are no amorphous sets.) and form 90 (Every linearly orderable set is well orderable.) are true and form 221 (Every infinite set has a non-principal measure.) is false. The Hahn-Banach Theorem (form 52) implies 221, so it is also false. (See Brunner/Howard/Rubin \cite{bib:1996}.)

We first note that for each finite ordered partition $P= (B_{1},\,\ldots,\,B_{n})$ of $A,\,G_{P}\in \mathcal{F}$, where $G_{P}=\{\psi\in G:\psi(P)=P\}$ and $\mathcal{F}$ is the filter of subgroups which determine $\mathcal{N}51$. Further, if $H\in \mathcal{F}$ then there is a finite ordered partition $P$ of $A$ such that $G_{P}\subseteq H$. Therefore for each $x\in \mathcal{N}51$, there is a finite ordered partition $P$ of $A$ such that $(\forall\psi\in G_{P})(\psi(x)=x)$. We will refer to such a partition $P$ as a support of $x$.

In order to prove that form 221 is false we shall give a proof that in $\mathcal{N}51$ there is no non-principal measure on $A$, the set of atoms.

Suppose, $m$ is a non-principal measure on $A$ supported by the partition $P= (B_{1},\,\ldots,\,B_{n})$ of $A$. Since $m$ is finitely additive on disjoint sets, $m(A)=1$, and $m(\{a\})=0$, for all $a\in A$. It follows that there must be an $i\leq n$ such that $B_{i}$ is infinite and $m(B_{i})\neq 0$. There exist pairwise disjoint infinite subsets, $S_{1},\,S_{2}$, and $S_{3}$, of $B_{i}$ such that $B_{i}=S_{1}\cup S_{2}\cup S_{3}$. Suppose $\sigma_{1},\,\sigma_{2}\in G_{P}$ are such that
\setcounter{equation}{0}
\begin{equation}\label{part04:eqn1a}
\sigma_{1}(S_{1})=S_{2}\cup S_{3}\ \mathrm{and}\ \sigma(S_{2}\cup S_{3})=S_{1}
\end{equation}
\begin{equation}
\label{part04:eqn2} \sigma_{2}(S_{1})=S_{2},\,\sigma_{2}(S_{2})=S_{3},\ \text{and}\
\sigma_{2}(S_{3})=S_{1}
\end{equation}
Then,\index{Author}{Brunner/Howard/Rubin} we would have
\begin{enumerate}
\item[(a)] $m(S_{1})=m(S_{2}\cup S_{3})=m(S_{2})+m(S_{3})$, by (\ref{part04:eqn1a}).
\item[(b)] $m(S_{1})=m(S_{2})=m(S_{3})$, by (\ref{part04:eqn2}).
\end{enumerate}

Consequently, $m(S_{1})=2m(S_{1})$, which implies that $m(S_{1})=0$. Thus, by (b), $m(B_{i})=0$, which is a contradiction.

In order to prove that forms 90 and 64 are true, we need the following lemma which is a slight strengthening of form 314 (from Degen\index{Author}{Degen} \cite{bib:1988}). The proof requires the axiom of choice but we intend to use the lemma in the ground model where AC holds.

\begin{lemma*}
If $\psi$ is a permutation of a non-empty set $A$ then there are two reflections $\rho_{1}$ and $\rho_{2}$ (a reflection on a set $A$ is a permutation $\rho$ for which $\rho^{2}=1_{A}$.) satisfying
\begin{enumerate}
\item[\emph{(i)}] $\psi=\rho_{2}\circ\rho_{1}$ and
\item[\emph{(ii)}] For all $B\subseteq A$, if $\psi(B)=B$ then $\rho_{1}(B)=B$ and $\rho_{2}(B)=B$.
\end{enumerate}
\end{lemma*}

\begin{proof}
Assume $\psi$ is a permutation of $A$. Write $\psi=\prod_{i\in K}c_{i}$ as a product of disjoint cycles. Each $c_{i}$ can be written as a product of reflections $c_{i}=\rho_{2,i}\circ\rho_{1,i}$ as follows. If $c_{i}= (a_{-n},\,a_{-(n-1)},\,\ldots,\,a_{-1},\,a_{0},\,a_{1},\,\ldots,\,a_{n})$ has odd length then
\begin{gather*}
\rho_{1,i}=(a_{-1},\,a_{1})(a_{-2},\,a_{2})\cdots(a_{-n},\,a_{n})\ \mathrm{and}\\
\rho_{2,i}=(a_{0},\,a_{1})(a_{-1},\,a_{2})\cdots(a_{-(n-1)},\,a_{n}).
\end{gather*}
If $c_{i}= (a_{-(n-1)},\,a_{-(n-2)},\,\ldots,\,a_{-1},\,a_{0},\,a_{1},\,\ldots,\,a_{n})$ has even length then
\begin{gather*}
\rho_{1,i}=(a_{-1},\,a_{1})(a_{-2},\,a_{2})\cdots(a_{-(n-1)},\,a_{(n-1)})\ \mathrm{and}\\
\rho_{2,i}=(a_{0},\,a_{1})(a_{-1},\,a_{2})\cdots(a_{-(n-1)},\,a_{n}).
\end{gather*}
If $c_{i}= (\ldots,\,a_{-2},\,a_{-1},\,a_{0},\,a_{1},\,a_{2},\,\ldots)$ is infinite then
\begin{equation*}
\rho_{1,i}=\prod_{n=1}^{\infty}(a_{-n},\,a_{n})\ \mathrm{and}\ \rho_{2,i}=\prod_{n=0}^{\infty}(a_{-n},\, a_{n+1}).
\end{equation*}
Then if $\rho_{1}=\prod_{i\in K}\rho_{1,i}$ and $\rho_{2}=\prod_{i\in K}\rho_{2,i}$ we have $\rho_{1}^{2}=1_{A},\,\rho_{2}^{2}=1_{A}$ and $\psi=\rho_{2}\circ\rho_{1}$. For the argument that (ii) holds we first note that if $\eta$ is a permutation whose expression as a product of disjoint cycles is $\eta =\prod_{j\in J}\sigma_{j}$ and $B\subseteq\mathrm{dom}(\eta)$ then
\begin{equation*}\tag{$\ast$}
\eta(B)=B\Leftrightarrow(\forall j\in J)(\{x:\sigma_{j}(x)\neq x\}\subseteq B\ \mathrm{or}\ \{x:\sigma_{j}(x)\neq x\}\cap\,B=\emptyset).
\end{equation*}
Now assume that $B\subseteq A$ and that $\psi(B)=B$. By $(\ast)$ for each $i\in K$, either $\{a\in A:c_{i}(a)\neq a\}\subseteq B$ or $\{a\in A:c_{i}(a)\neq a\}\cap B=\emptyset$. When $\rho_{1}$ is written as a product of disjoint cycles each cycle is a transposition $(a,\,b)$ which is a cycle of $\rho_{1,i}$ for some $i\in K$. Hence by the definition of $\rho_{1,i},\,\{a,\,b\}\subseteq\{a\in A:c_{i}(a)\neq a\}$. Therefore $\{a,\,b\}\subseteq B$ (if $\{a\in A:c_{i}(a)\neq a\}\subseteq B)$ or $\{a,\,b\}\cap B=\emptyset$ (if $\{a\in A: c_{i}(a)\neq a\}\cap B=\emptyset)$. It follows from $(\ast)$ that $\rho_{1}(B)=B$. Similarly, $\rho_{2}(B)=B$. 
\end{proof}

\begin{theorem*}
In $\mathcal{N}51$ every linearly orderable set is well orderable.
\end{theorem*}

\begin{proof}
Let $X$ be linearly orderable in $\mathcal{N}51$ and assume that the ordered partition $P= (B_{1},\,\ldots,\, B_{k})$ is a support of a linear ordering $<$ on $X$. We shall show that if $\psi\in G_{P}$, then $\psi$ fixes $X$ pointwise and hence $P$ is a support of a well ordering of $X$. Assume $\psi\in G_{P}$. By the lemma, $\psi=\rho_{2}\circ\rho_{1}$ where $\rho_{1}$ and $\rho_{2}$ are reflections and (by part (ii)) $\rho_{1}$ and $\rho_{2}$ are in $G_{P}$. The permutations $\rho_{1}$ and $\rho_{2}$ are therefore order automorphisms of $(X,\,<)$. Since no order automorphism of a linearly ordered set can have a finite cycle of length greater than 1, we conclude that for all $x\in X,\,\rho_{1}(x)=x$ and $\rho_{2}(x)=x$ hence $\forall x\in X,\,\psi(x)=x$. 
\end{proof}

\begin{theorem*}
In $\mathcal{N}51$ there are no infinite amorphous sets.
\end{theorem*}

\begin{proof}
Assume that $X\in \mathcal{N}51$ is infinite and that the ordered partition $P= (B_{1},\,\ldots,\,B_{k})$ of $A$ is a support of $X$. We will show that there is an infinite set $R\in \mathcal{N}51$ such that $R\subset X$ and $X-R$ is infinite.

If every element of $X$ has support $P$ then $X$ is well orderable in $\mathcal{N}51$ and the proof is easy. We therefore assume that $y\in X$ and there is a $\phi\in G_{P}$ such that $\phi(y)\neq y$.

\textsc{Step 1.}
We show that we may assume that $y$ has a support
\begin{equation*}
Q=(C_{1},\,C_{2},\,C_{3},\,\ldots,\,C_{n_{0}},\,C_{n_{0}+1},\,\ldots,\,C_{n})
\end{equation*}
where
\begin{enumerate}
\item[a.] $n_{0}\geq 3$ and $C_{3}$ is infinite.
\item[b.] For $1\leq j\leq n_{0},\,C_{j}\subseteq B_{1}$ and for $n_{0}<j\leq n,\, C_{j}\cap B_{1}=\emptyset$.
\item[c.] $\phi(C_{1})=C_{2},\,\phi(C_{2})=C_{1}$, and for every $a\in A-(C_{1}\cup
C_{2})$, $\phi(a)=a$.
\end{enumerate}

First, we may assume (by refining $P$ if necessary) that each $B_{i}$ is either infinite or a singleton. Secondly, by writing $\phi=\prod_{i=1}^{k}\phi_{i}$ where $\phi_{i}(a)=a$ for $a\in A-B_{i}$ and replacing $\phi$ by one of the $\phi_{i}$, we may assume that there is an $i\leq k$ such that $\phi(a)=a$ for $a\in A-B_{i}$. (If $\prod_{i=1}^{k}\phi_{i}(y)\neq y$ then for some $i,\,\phi_{i}(y)\neq y$.) Further it is no loss of generality to assume that $i=1$.

Let $Q= (C_{1},\,\ldots,\,C_{n})$ be a support of $y$. By replacing each $C_{j}$ with the two sets $C_{j}\cap B_{1}$ and $C_{j}-B_{1}$ (if necessary) and reordering the $C_{j}$'s we may assume that there is an $n_{0}\leq n$ such that $C_{j}\subseteq B_{1}$ for $j\leq n_{0}$ and $C_{j}\cap B_{1}=\emptyset$ for $j>n_{0}$. As in the proof of the previous theorem (using the lemma) we may assume that $\phi$ is a reflection. By replacing $Q$ by a partition whose cells are $\{C_{i}\cap \phi(C_{j}):1\leq i,\, j\leq n\}-\{\emptyset\}$ we may assume (since $\phi$ is a reflection) that $\forall i,\,1\leq i\leq n,\,\exists j,\,1\leq j\leq n$ such that $\phi(C_{i})=C_{j}$ and $\phi(C_{j})=C_{i}$. (By our simplifying assumptions so far, $\phi(C_{j})=C_{j}$ for $j>n_{0}$.) By decomposing $\phi$ into a product $\phi=\prod_{r=1}^{m}\gamma_{r}$ where for each $r,\,1\leq r\leq m$ there are $i$ and $j$ such that $\gamma_{r}(a)=a$ for $a\in A-(C_{i}\cup C_{j})$, $\gamma_{r}(C_{i})=C_{j}$ and $\gamma)r(C_{j})=C_{i}$ and replacing $\phi$ by one of the $\gamma_{r}$, we may assume that $\phi(a)=a$ for $a\in A-(C_{i}\cup C_{j})$, $\phi(C_{i})=C_{j}$ and $\phi(C_{j})=C_{i}$. It is no loss of generality to assume that $i=1$ and $j=2$.

Finally, we may assume that for some $i,\,3\leq i\leq n_{0},\,C_{i}$ is infinite. Assuming that this is not the case and using the facts that $C_{1}\cup C_{2}\subseteq B_{1}$ and that $B_{1}$ is infinite we conclude that $C_{1}$ and $C_{2}$ are infinite. Let $C^{\prime}$ be any infinite subset of $C_{1}$ whose complement in $C_{1}$ is infinite. Note that $C_{2}$ is the disjoint union $C_{2}= \phi(C^{\prime})\cup\phi(C_{1}-C^{\prime})$. Therefore, $Q^{\prime}=(C^{\prime},\,\phi(C^{\prime}),\,C_{1}-C^{\prime},\,\phi(C_{1}-C^{\prime}),\,C_{3},\,\ldots,\,C_{n})$ is a partition of $Q$ which supports $y$. Define $\phi^{\prime}$ and $\phi^{\prime\prime}$ (both in $G_{P}$) as follows:
\begin{equation*}
\phi^{\prime}(a)=\left\{\begin{array}{ll}
\phi(a) & a\in C^{\prime}\cup\phi(C^{\prime})\\
a & \mbox{otherwise}
\end{array}\right.
\end{equation*}
and
\begin{equation*}
\phi^{\prime\prime}(a)=\left\{\begin{array}{ll}
\phi(a) & a\in(C_{1}-C^{\prime})\cup\phi(C_{1}-C^{\prime})\\
a & \mbox{otherwise}
\end{array}\right.\ .
\end{equation*}
Clearly $\phi=\phi^{\prime}\circ\phi^{\prime\prime}$ and therefore one of $\phi^{\prime}$ or $\phi^{\prime\prime}$ moves $y$. Assume without loss of generality that $\phi^{\prime}(y)\neq y$ then replace $Q$ with $Q^{\prime}$ and $\phi$ by $\phi^{\prime}$.

This completes step 1.

\textsc{Step 2.} We use $y$ to obtain an infinite subset $R$ of $X$ whose complement in $X$ is infinite. Our primary tool will be the following lemma.

\begin{lemma*}
Assume that $\delta,\,\lambda\in G_{P}$, that $\lambda(a)=a$ for all $a\in A-(\delta(C_{1})\cup\delta(C_{2}))$, that $\lambda(\delta(C_{1}))=\delta(C_{2})$, and that $\lambda(\delta(C_{2}))=\delta(C_{1})$. Then $\lambda(\delta(y))\neq\delta(y)$.
\end{lemma*}

\begin{proof}
By the hypotheses, $\delta^{-1}\lambda\delta(C_{1})=C_{2}$ and $\delta^{-1}\lambda\delta(C_{2})=C_{1}$ and for all $a\in A-(C_{1}\cup C_{2})$, $\delta^{-1}\lambda\delta(a)=a$. This means that $\delta^{-1}\lambda\delta(Q)=\phi(Q)$ and since $Q$ is a support of $y,\,\delta^{-1}\lambda\delta(y)=\phi(y)\neq y$. It follows that $\lambda(\delta(y))\neq\delta(y)$. 
\end{proof}

Now by using permutations in $G_{P}$ which move $C_{1}$ and $C_{2}$ we can obtain infinitely many ``copies'' of $y$ all of which are in $X$.

Partition $C_{3}$ into sets $\{S_{j},\,T_{j},\,U_{j},\,V_{j}:j\in \omega\}$ so that $|S_{j}|=|T_{j}|=|U_{j}|=|V_{j}|= |C_{1}|(=|C_{2}|)$. Let $C_{3}^{\prime}$ and $C_{3}^{\prime\prime}$ denote the sets $\bigcup_{j\in\omega}(S_{j}\cup T_{j})$ and $\bigcup_{j\in\omega}(U_{j}\cup V_{j})$ respectively. Let $\tau \in G_{P}$ satisfy $\tau(a)=a$ for all $a\in A-(C_{1}\cup C_{2}\cup S_{0}\cup
T_{0}),\tau(C_{1})=S_{0},\,\tau(S_{0})=C_{1},\,\tau(C_{2})=T_{0}$ and $\tau(T_{0})=C_{2}$. Let $\nu\in G_{P}$ satisfy $\nu(a) =a$ for all $a \in A -(C_{1}\cup C_{2}\cup U_{0}\cup V_{0})$, $\nu(C_{1})=U_{0},\,\tau(U_{0})=C_{1},\,\tau(C_{2})=V_{0}$ and $\tau(V_{0})=C_{2}$. For $j\in \omega,\,j\geq 1$, let $\sigma_{j}\in G_{P}$ satisfy $\sigma_{j}(a)=a$ for $a\in A-(S_{0}\cup T_{0}\cup S_{j}\cup
T_{j})$, $\sigma_{j}(S_{0})=S_{j},\,\sigma_{j}(S_{j})=S_{0},\,\sigma_{j}(T_{0})=T_{j}$ and $\sigma_{j}(T_{j})=T_{0}$. Similarly, for $j\geq 1$, let $\gamma_{j}\in G_{P}$ satisfy $\gamma_{j}(a)=a$ for all $a\in A-(U_{0}\cup V_{0}\cup U_{j}\cup
V_{j})$, $\gamma_{j}(U_{0})=U_{j},\,\gamma_{j}(U_{j})=U_{0},\, \gamma_{j}(V_{0})=V_{j}$, and $\gamma_{j}(V_{j})=V_{0}$. The partition $P^{\prime}= (C_{1},\,C_{2},\,C_{3}^{\prime},\, C_{3}^{\prime\prime},\,\ldots,\,C_{n_{0}},\,B_{2},\,\ldots,\,B_{k})$ is a refinement of $P$ and therefore is a support of $X$. Hence the two sets $R_{1}=\{\psi(\tau(y)):\psi\in G_{P^{\prime}}\}$ and $R_{2}=\{\psi(\nu(y)):\psi\in G_{P^{\prime}}\}$ are subsets of $X$. Further, $P^{\prime}$ is a support of $R_{1}$ and $R_{2}$ so $R_{1}$ and $R_{2}$ are in $\mathcal{N}51$. We shall also prove the following:
\begin{enumerate}
\item[(i)] $(\forall j\in \omega,\,j\geq 1)(\sigma_{j}(\tau(y))\in R_{1}\wedge\gamma_{j}(\nu(y))\in R_{2})$.
\item[(ii)] $(\forall j\in \omega,j\geq 1)(\sigma_{j}(\tau(y))\not\in R_{2}\wedge\gamma_{j}(\nu(y))\not\in R_{1})$.
\item[(iii)] $(\forall j,\,r\in \omega,j,\,r\geq 1)(j\neq r\rightarrow[\sigma_{j}(\tau(y))\neq\sigma_{r}(\tau(y))\wedge(\gamma_{j}(\nu(y))\neq\gamma_{r}(\nu(y))])$.
\end{enumerate}

It will follow that $R=R_{1}$ is an infinite subset of $X$ with infinite complement in $X$ and the proof will be complete. To prove (i) we observe that since $\sigma_{j}$ and $\gamma_{j}$ agree with the identity permutation outside of $C_{3}^{\prime}$ and $C_{3}^{\prime\prime}$ respectively for $j\geq 1$, they are both in $G_{P^{\prime}}$.

For (ii) we first note that since $Q$ is a support of $y$, a support of $\nu(y)$ is given by
\begin{align*}
\nu(Q)&=(\nu(C_{1}),\,\nu\,(C_{2}),\,\nu(C_{3}),\,\ldots,\,\nu(C_{n}))\\
&=(U_{0},\,V_{0},\,[C_{3}-(U_{0}\cup V_{0})]\cup(C_{1}\cup C_{2}),\,C_{4},\,\ldots,\,C_{n}).
\end{align*}
Therefore if $\psi$ fixes $P^{\prime}$ then $\psi(\nu(y))$ has support $\psi(\nu(Q))=(\psi(U_{0}), \psi(V_{0})$, $[\psi(C_{3})-\psi(U_{0}\cup V_{0})]\cup\psi(C_{1}\cup C_{2})$, $\psi(C_{4}),\,\ldots,\, \psi(C_{n}))=(\psi(U_{0})$, $\psi(V_{0})$, $[C_{3}-\psi(U_{0}\cup
V_{0})]\cup(C_{1}\cup C_{2})$, $\psi(C_{4})$, $\ldots,\,\psi(C_{n}))$. Since $C_{3}^{\prime}\subseteq C_{3}-\psi(U_{0}\cup V_{0})$ we may conclude that for any $\lambda \in G_{P}$, if $\lambda(a) =a$ for all $a \in A -C_{3}^{\prime}$ then $\lambda(\psi(\nu(Q)) =\psi(\nu(Q))$ and therefore $\lambda(\psi(\nu(y)))=\psi(\nu(y))$. Hence such an $\lambda$ fixes $R_{2}$ pointwise. In particular, if we choose $\lambda$ so that $\lambda(a)=a$ for all $a\in A-(S_{j}\cup T_{j})$, $\lambda(S_{j})=T_{j}$ and $\lambda(T_{j})=S_{j}$ then $\lambda$ fixes $R_{2}$ pointwise and by the lemma (with $\delta=\sigma_{j}\circ\tau)\lambda(\sigma_{j}(\tau(y)))\neq\sigma_{j}(\tau(y))$. It follows that $\sigma_{j}(\tau(y))\not\in R_{2}$. A similar argument shows that $\gamma_{j}(\nu(y))\not\in R_{1}$.

The argument for (iii) is similar. We shall show
\begin{equation*}
(\forall j,\,r\in \omega,\,j,\,r\geq 1)(j\neq r\rightarrow\sigma_{j}(\tau(y))\neq\sigma_{r}(\tau(y)))
\end{equation*}
and leave the similar argument that $\gamma_{j}(\nu(y))\neq\gamma_{r}(\nu(y))$ to the reader. First note that $\sigma_{r}(\tau(y))$ has support
\begin{equation*}
\sigma_{r}(\tau(Q))=(S_{r},\,T_{r},\,[C_{3}-(S_{r}\cup T_{r})]\cup (C_{1}\cup
C_{2}),\,C_{4},\,\ldots,\,C_{n}).
\end{equation*}
Call this support $Q^{\prime}$. Since $S_{j}\cup T_{j}\subseteq C_{3}-(S_{r}\cup
T_{r})$ any $\lambda$ satisfying $\lambda(a)=a$ for all $a\in A-(S_{j}\cup
T_{j})$ fixes $Q^{\prime}$ and therefore fixes $\sigma_{r}(\tau(y))$. In particular if we choose such a $\lambda$ for which $\lambda(S_{j})=T_{j}$ and $\lambda(T_{j})=S_{j}$ then $\lambda$ fixes $\sigma_{r}(\tau(y))$ and by the lemma (with $\delta=\sigma_{j}\circ \tau)\lambda(\sigma_{j}(\tau(y)))\neq\sigma_{j}(\tau(y))$. We can therefore conclude that $\sigma_{r}(\tau(y))\neq\sigma_{j}(\tau(y))$. This completes the proof of the theorem. 
\end{proof}

\section*{NOTE 129}

In this note we give definitions concerning free groups for forms [1 BB], [1 CA], 68 and 348.

\begin{definition*}
Assume that $(G,\,\circ)$ is a group freely generated by $X$, that $A\subseteq G$, and that $a,\,b$ and $c$ are in $G$.
\begin{enumerate}
\item[1.] The subgroup of $G$ generated by $A$ is denoted $\langle A\rangle$.
\item[2.] The \emph{length of a with respect to} $X$ (denoted by $L_{X}(a)$ or $L(a)$) is the unique natural number $n$ such that there are elements $a_{1},\,\ldots,\,a_{n}$ of $X\cup X^{-1}$ such that $a_{j}\neq a_{j+1}^{-1}$ and $a=a_{1}\circ\cdots \circ a_{n}$.
\item[3.] The subset $A$ is \emph{level} with respect to $X$ if $\langle A\rangle$ is freely generated by $A$ and for each $b\in\langle A\rangle,\,b\in\langle\{a:a\in A\wedge L(a)\leq L(b)\}$.
\item[4.] A subset $A$ of $G$ has the \emph{Nielsen property} with respect to $X$ if
\begin{enumerate}
\item[a.] $A\cap A^{-1}=\emptyset$.
\item[b.] If $a,\,b\in A\cup A^{-1}$ and $L(a\circ b)<L(a)$ then $b=a^{-1}$, and
\item[c.] If $a,\,b$ and $c\in A\cup A^{-1}$ and $L(abc)\leq L(a)-L(b)+L(c)$ then either $b=a^{-1}$ or $c=b^{-1}$.
\end{enumerate}
\end{enumerate}

In Federer/Jonsson \cite{bib:1950} it is shown that a set with the Nielsen property is level.
\end{definition*}

\section*{NOTE 130}

In this note we prove that form [30 G] (Every $\subseteq$-chain in $\mathcal{P}(X)$ is contained in a maximal $\subseteq$-chain) is equivalent to form 30 (The ordering principle). Form [30 G] was suggested by Adam Kolany.

Assume [30 G].\index{Subject}{level set}\index{Subject}{Nielsen property}\index{Author}{Federer/Jonsson} Let $X$ be a set and let $Y$ be a maximal $\subseteq$-chain in $\mathcal{P}(X)$. Define $\leq$ on $X$ by $a\leq b$ iff $(\exists A\in Y)(a\in A\wedge b\not\in A)$. $\leq$ is linear because if not there are $a,\,b\in X$ such that $a\not\leq b$ and $b\not\leq a$. If this happens then $Y\cup\{A_{0}\}$ (where $A_{0}=(\bigcup\{A\in Y:a\not\in A\})\cup\{b\})$ is a chain in $\mathcal{P}(X)$ which properly extends $Y$.

Assume form 30. Let $X$ be a set and $Y$ a $\subseteq$-chain in $\mathcal{P}(X)$. Let $\leq$ be a linear ordering of $X$. Define another linear ordering $R$ on $X$ by $a\, R\,b$ iff $[(\exists A\in Y)(a\in A\wedge b\not\in A)]$ or $[(\forall A\in Y)(a\in A\leftrightarrow b\in A)\wedge a\leq b]$. Then $Y^{\prime}=\{A\subseteq X:(\forall a\in A)(\forall x\,R\,a)(x\in A)\}$ is a maximal $\subseteq$-chain containing $Y$.

\section*{NOTE 131}

In this note we give a proof of form [ $0$ AL] (For every set $A$, there is an $\aleph$ which cannot be covered by fewer than $|A|$ sets each of cardinality $<|A|.)$ in $\mathrm{ZF}^{0}$. Form [0 AL] was suggested by K. Keremedis.

Let $A$ be a set and suppose that $K$ is an aleph which is covered by fewer than $|A|$ sets of cardinality $<|A|$. Let $\leq$ be the usual ordering on $K$ and let $\bigcup_{b\in B}C_{b}=K$ where $|B|<|A|$ and for each $b\in B,\,|C_{b}|<|A|$. Since each $C_{b}$ is well orderable, $|C_{b}|<H(A)$. (The
function $H$ is Hartogs' aleph function, $H(A)$ is the least aleph that is not $\leq|A|.)$ Therefore the ordering $(C_{b},\,\leq)$ has order type that of an ordinal $\alpha_{b}<H(A)$. Let $f_{b}$ be the unique order automorphism from $\alpha_{b}$ onto $C_{b}$. Let $\gamma$ be a fixed element of $K$ and define $F:A\times H(A)\rightarrow K$ by $F(b,\,\beta)=f_{b}(\beta)$ if $b\in B$ and $\beta\in\alpha_{b}$ and $F(b,\,\beta)=\gamma$ otherwise. $F$ is onto and $|K|=|\mathcal{E}|$ where $\mathcal{E}$ is the set of equivalence classes $\subseteq A\times H(A)$ under the relation $q\sim r\Leftrightarrow F(q)=F(r)$. Therefore $|K|=\aleph_{\sigma}$ where $A\times H(A)$ has a partition of cardinality $\aleph_{\sigma}$. Therefore, if we let $\aleph_{\rho}$ be the smallest well ordered cardinal greater than every cardinal in the set $\{\aleph_{\sigma}:A\times H(A)$ has a partition of cardinality $\aleph_{\sigma}$\}, then $\aleph_{\rho}$ cannot be covered by fewer than $|A|$ sets each of cardinality $|A|$.

\section*{NOTE 132}

The following definitions and results for forms [8 M], [8 N], [43 Q], [126
B] through [126 F] are from Keremedis \cite{bib:1996d}.

\begin{definition*}
Let $(X,\,T)$ be a topological space. $B\subseteq T-\{\emptyset\}$ is said to be a $\pi$-\emph{base for} $X$ if every non-empty element of $T$ includes an element of $B$.
\end{definition*}

\begin{definition*}
Let $(P,\,\leq)$ be a partially ordered set (poset).
\begin{enumerate}
\item[(a)] A subset $D\subseteq P$ is set to be \emph{dense} if every element of $P$ has a lower bound in $D$. $D$ is said to be \emph{open} if it is open in the order topology.
\item[(b)] An element $a\in P$ is said to be an \emph{atom} if every pair of lower bounds of $a$ are compatible (have a common lower bound).
\item[(c)] If $P$ has no atoms it is called \emph{atomless}. A pairwise incompatible subset of $P$ is called an \emph{antichain}.
\item[(d)] Let $D=\{D_{i}:i\in \omega\}$ be a denumerable family of dense open sets in $P$. $C=\{C_{i}:i\in \omega\},\,\bigcup C=\omega$ is called a \emph{matrix} for $D$ if the following three conditions are satisfied:
\end{enumerate}
\begin{enumerate}
\item[(i)] $\forall i\in \omega\, C_{i}$ is a countable antichain in $P$.
\item[(ii)] $\forall i\in \omega,\,\forall j>i,\,\forall c\in C_{i}$ there exist infinitely many $t\in C_{j}$ such that $t\leq c$.
\item[(iii)] $\forall i\in \omega,\,\forall j\in \omega$, all but finitely many element of $C_{j}$ belong to $D_{i}$.
\end{enumerate}
\end{definition*}

\begin{definition*}
Let $A=\{A_{i}:i\in I\}$ be a family of non-empty sets.
\begin{enumerate}
\item[(a)] $A$ has the \emph{strong finite intersection property, sfip}\index{Subject}{strong finite intersection property
(sfip)} if for each finite $Q\subseteq A,\,|\bigcap Q|\geq \omega $.
\item[(b)] An infinite subset $S$ of $A$ is a \emph{pseudo-intersection}\index{Subject}{pseudo-intersection} if $|S-A_{i}|<\omega$ for all $i\in I$.
\end{enumerate}

The following statements are considered:

\noindent $MC_{\omega}(\aleph_{0},\,\infty)$: For every denumerable family $A$ of pairwise disjoint non-empty sets, there exists a set $C$ such that for all $a\in A,\,0<|C\cap a|\leq \omega$.

\noindent \emph{pseudo}$(w)([126\ \mathrm{F}])$: Every denumerable family of sets having the sfip also has a pseudo-intersection.

\noindent $DP$: If $(X,\,T)$ is a topological space having a countable $\pi$-base\index{Subject}{$\pi$-base}, then for every family $D=\{D_{i}:i\in \omega\}$ of dense open sets of $X$, there is a countable dense set $S\subseteq X$ such that for all $i\in \omega$ and for all but finitely many $s\in S,\,s\in D_{i}$. ($DP$ is form [8 N]).

\noindent$DP^{\ast}\ ([126\ \mathrm{E}])$: If $(X,\,T)$ is a topological space, then for every family $D=\{D_{i}: i\in \omega\}$ of dense open sets of $X$, there is a countable dense set $S\subseteq X$ such that for all $i\in \omega$ and for all but finitely many $s\in S,\,s\in D_{i}$.

\noindent $WCMC\ ([126\ \mathrm{D}])$: For every denumerable family $A$ of disjoint non-empty sets there is a set $C\subset\bigcup A$ such that for every $a\in A,\, 0\leq|C\cap a|<\omega$.

\noindent $FSCB$: Every separable first countable space $(X,\,T)$ has a countable $\pi$-base.

\noindent $Matrix(\omega)$: For every atomless poset\index{Subject}{atomless poset} $(P,\,\leq)$ and every family $D=\{D_{i}:i\in \omega\}$ of dense open subsets of $P$, there exists a matrix $C=\{C_{i}:i\in \omega\},\, |\bigcup C|=\omega$ of $P$ for $D$. ($Matrix(\omega)$ is form [43 Q].)

Results:
\begin{enumerate}
\item[1.] $MC(\aleph_{0},\,\infty)$ (form 126) $\leftrightarrow MC_{\omega}(\aleph_{0},\,\infty)+MC(\aleph_{0},\, \aleph_{0})$ (form 350). (The conjunction $MC_{\omega}(\aleph_{0},\,\infty)+MC(\aleph_{0},\,\aleph_{0})$ is form [126 B].)
\item[2.] $C(\aleph_{0},\,\infty)$ (form 8) $\leftrightarrow MC_{\omega}(\aleph_{0},\,\infty)+C(\aleph_{0},\, \aleph_{0})$ (form 32). (The conjunction $MC_{\omega}(\aleph_{0},\,\infty)+C(\aleph_{0},\,\aleph_{0})$ is form [8 M].)
\item[3.] $C(\aleph_{0},\,\infty)$ (form 8) $\leftrightarrow DP$ (form [8
N]).
\item[4.] $MC(\aleph_{0},\,\infty)$ (form 126) $\leftrightarrow WCMC$ ([126
D]).

\noindent (Clearly, $MC(\aleph_{0},\,\infty)\rightarrow WCMC$. Let $X=\{A_{n}:n\in \omega\}$ be a denumerable set of pairwise disjoint non-empty sets. For each $n\in \omega$, let $B_{n}=A_{0}\times A_{1}\times\cdots\times A_{n}$ and let $B=\{B_{n}:n\in \omega\}$. Let $C=\{B_{n_{i}}:n_{i}\in \omega\}$ be a denumerable subset of $B$ which has a multiple choice function $f$. The function $f$ can be used to construct a choice function on $A$. This result was noted by K. Keremedis\index{Author}{Keremedis}.)
\item[5.] $WCMC$ ([126 D])$\leftrightarrow DP^{\ast}$ ([126 E]) $\leftrightarrow pseudo(\omega)$ ([126 F]).
\item[6.] $MC(\aleph_{0},\,\infty)$ (form 126) $\leftrightarrow FSCB\, +\, MC_{\omega}(\aleph_{0},\,\infty)$. (The conjunction $FSCB+ MC_{\omega}(\aleph_{0},\,\infty)$ is form [126 C].)
\item[7.] $DC(\omega)$ (the Principle of Dependent Choices (43)) $\leftrightarrow Matrix(\omega)$ (form [43 Q]).
\end{enumerate}
\end{definition*}

\section*{NOTE 133}

In this note we prove that forms [1 F] and [1 CC] are equivalent to form 1.

Form [1 F] is ``For all ordinals $\alpha,\,DC(\aleph_{\alpha})$: Given a relation $R$ such that for every subset $Y$ of a set $X$ with $|Y|<\aleph_{\alpha}$, there is and $x\in X$ with $Y\,R\,x$, then there is a function $f:\aleph_{\alpha}\rightarrow X$ such that $\forall\beta<\aleph_{\alpha},\, \{f(\gamma):\gamma<\beta\}\
R\ f(\beta).)$''. It is clear that form 1 implies [1 F]. To prove the converse, well order $X$ by using $Y\,R\,x\Leftrightarrow x\not\in Y$ and use $DC(\aleph_{\alpha})$ for sufficiently large $\aleph_{\alpha}$.

The question of the equivalence of [1 CC] and 1 was first raised by K. Keremedis. Form [1 CC] is the conjunction of $PC(\aleph_{0},\,\mathrm{odd},\,\infty)$: ``Every denumerable family of odd sized (finite) sets has an infinite subfamily with a choice function.'' and $MC_{\omega^{+}}$: ``For every natural number $n\geq 1$ and for every set of non-empty sets, there is a function $f$ such that for each $u\in x,\,f(u)$ is a non-empty finite subset of $u$ such that $|f(u)|$ and $n$ are relatively prime.'' We first prove that $PC(\aleph_{0},\,\mathrm{odd},\,\infty)$ implies ``every denumerable family of odd sized sets has a choice function''. Let $X=\{x_{n}:n\in \omega\}$ be a denumerable family of odd sized sets. For each $n\in \omega$, let $y_{n}=\{f:f$ is a choice function for $\{x_{i}:i\leq n\}$\}. $|y_{n}|=|x_{1}|\cdots|x_{n}|$ and is therefore odd. Applying $PC(\aleph_{0},\,\mathrm{odd},\,\infty)$ to $Y=\{y_{n}:n\in \omega\}$ gives an infinite subset $Y^{\prime}$ of $Y$ with a choice function $F$. Then defining $G$ by $G(x_{n})=F(y_{k})(x_{n})$ where $k$ is the least natural number $\geq n$ for which $y_{k}\in\mathrm{dom}(F)$ gives a choice function for $X$.

To prove that $PC(\aleph_{0}, \mathrm{odd},\infty)$ and $MC_{\omega^{+}}$ imply $AC$. It suffices to show that $\forall n\in \omega,\,C(\infty,\,n)(form
61)+MC(form 67)+MC_{\omega^{+}}+$ ``Every denumerable family of odd sized sets has a choice function'' implies $AC$. This is because $MC_{\omega^{+}}$ is equivalent to $\forall n\in \omega,\, C(\infty,\,n)+MC$ (by lemma 3  of Howard/A. Rubin/J. Rubin\index{Author}{Howard/A. Rubin/J. Rubin} \cite{bib:1997}) and because $PC(\aleph_{0},\,\mathrm{odd},\, \infty)$ is equivalent to ``Every denumerable of odd sized sets has a choice
function''

Let $X$ be a set of pairwise disjoint, non-empty sets. By $MC$ there is an $f:X\rightarrow \mathcal{P}(\bigcup X)$ such that $\forall x\in X,\,\emptyset \neq f(x)\subseteq x$ and $|f(x)|$ is finite. For each $n\in \omega$, let $Y_{n}=\{x\in X:|f(x)|=n\}$. By $C(\infty,\,n)$, $Y_{n}$ has a choice function (even if $Y_{n}$ happens to be empty, but it might be clearer to exclude $n$ for which $Y_{n}$ is empty). Therefore the set $C_{n}=\{g:g$ is a choice function for $Y_{n}$\} is non-empty. By $MC_{\omega^{+}}$ there is a function $F$ defined on $\omega$ such that $F(n)$ is an odd sized subset of $C_{n}$ for each $n\in \omega$. Since every denumerable of odd sized sets has a choice function, there is a function $G$ defined on $\omega$ such that for every $n\in \omega,\,G(n)\in F(n)$. We can now define a choice function $Ch$ for $X$ by $Ch(x)=G(n)(x)$ where $|f(x)|=n$.

\section*{NOTE 134}

In this note we prove the result from Luzin/Sierpi\'{n}ski \cite{bib:1917} that form 100 (If $n$ and $m$ are cardinal numbers, then $m\leq^{\ast}n\rightarrow\neg(n<m).)$ implies form 369 (If $\mathbb{R}$ is partitioned into two sets, at least one has cardinality $2^{\aleph_{0}}$.). We give a proof here because Luzin/Sierpi\'{n}ski \cite{bib:1917} is not readily available.

Since $2^{\aleph_{0}}=|\mathcal{P}(\omega)|=|2^{\omega}|$ it suffices to show that for any partition $\{A,\,B\}$ of $\mathcal{P}(\omega)$, either $|A|=|2^{\omega}|$ or $|B|=|2^{\omega}|$. Since $|A|\leq|2^{\omega}|$ and $|B|\leq|2^{\omega}|$, by 100 it suffices to show that there is either a function from $A$ onto $2^{\omega}$ or a function from $B$ onto $2^{\omega}$. Define $f:\mathcal{P}(\omega)\rightarrow 2^{\omega}$ by the following: Assume $X=\{a_{0},\,a_{1},\,a_{2},\,\ldots\}\in \mathcal{P}(\omega)$ where $a_{0}<a_{1}<a_{2}<\ldots$. Then
\begin{equation*}
(f(X))(n)=\left\{\begin{array}{ll}
0 &\mathrm{if}\ a_{n+1}-a_{n}\ \text{is even or}\ |X|\leq n+1\\
1 & \mathrm{if}\ a_{n+1}-a_{n}\ \text{is odd}
\end{array}\right. .
\end{equation*}
Assume that $f|B$ is not onto $2^{\omega}$. Then for some $\gamma$ in $2^{\omega}$, the non-empty set $C= \{X\in \mathcal{P}(\omega):f(X)=\gamma\}\subseteq A$. Then we can define a function $g$ from $C$ onto $2^{\omega}$ as follows: Assume $X=\{a_{0},\,a_{1},\,a_{2},\,\ldots\} \in C$ where $a_{0}<a_{1}<a_{2}<\ldots$ and assume $n\in \omega$. If $\gamma(n)=0$, then $a_{n+1}-a_{n}$ is even or $|X|\leq n+1$, and we define
\begin{equation*}
(g(X))(n)=\left\{\begin{array}{ll}
0 &\mathrm{if}\ a_{n+1}-a_{n}\equiv 0\mod4\ \text{or}\ |X|\leq n+1\\
1& \text{if}\ a_{n+1}-a_{n}\equiv 2\mod4
\end{array}\right.
\end{equation*}
If $\gamma(n)=1$,\index{Author}{Luzin/Sierpi\'{n}ski} then $a_{n+1}-a_{n}$ is odd and we define
\begin{equation*}
(g(X))(n)=\left\{\begin{array}{ll}
0& \mathrm{if}\ a_{n+1}-a_{n}\equiv 1\mod4\\
1& \mathrm{if}\ a_{n+1}-a_{n}\equiv 3\mod4.
\end{array}\right.
\end{equation*}
The proof that $g$ is onto $2^{\omega}$ could be done by induction, that is, given any $\beta\in 2^{\omega}$ we could define by induction a sequence $\langle a_{n}\rangle_{n\in\omega}$ so that $X=\{a_{1},\,a_{2},\, \ldots\}\in C$ and $g(X)=\beta$.

\section*{NOTE 135}

In this note we prove that form 18 $(PUT(\aleph_{0},2,\,\aleph_{0}))$ and [18 B] (The union of a denumerable family of pairs has a denumerable subset.) are equivalent. It suffices to show that [18 A] $(PC(\aleph_{0},2,\,\infty))$ implies [18 B].

Let $X$ be a denumerable family of pairs. By [18 A] there is a denumerable subset $Y\subseteq X$ and a choice function $f$ for $Y$. If the range of $f$ is infinite then it is a denumerable subset of $\bigcup X$. If the range of $f$ is finite, then for some $t\in\bigcup X$ there are infinitely many elements $y\in Y$ such that $f(y)=t$. Each such $y$ has the form $y=\{a_{y},\,t\}$. In this case the set $\{a_{y}:f(y)=t\}$ is a denumerable subset of $\bigcup X$.

\section*{NOTE 136}

In this note we prove that the Boolean prime ideal theorem (form 14) is true in $\mathcal{N}19$. We shall use Blass' theorem (Blass\index{Author}{Blass} \cite{bib:1986}, Theorem~\ref{part04:thm2c}) which characterizes Fraenkel-Mostowski models in which BPI is true. According to this theorem BPI holds in a Fraenkel-Mostowski model if and only if the model is determined by a filter of subgroups $\Gamma$ with the Ramsey property described below.

\begin{definition*}
\begin{enumerate}
\item[1.] A subgroup $H$ of a group $G$ is a \emph{Ramsey subgroup} if for every finite $F\subseteq G/H$ (the set of left cosets of $H$ in $G$) there is a finite $Y\subseteq G/H$ such that whenever $Y$ is partitioned into two pieces there is a $g\in G$ such that $gF$ is included in one of the pieces.
\item[2.] A filter $\Gamma$ of subgroups of $G$ has the \emph{Ramsey property} if it has a basis consisting of Ramsey subgroups of $G$.
\end{enumerate}

Let $E$ be a support. By Blass' theorem it suffices to show that $\mathrm{fix}_{G}(E)$ is a Ramsey subgroup of $G$. We shall argue that this follows from Rado's corollary which is given in Halpern \cite{bib:1964}, page 61. Rado's corollary is
\end{definition*}

\textsc{Rado's Corollary.}
\emph{Assume that} $n,p\in \omega$ \emph{and that} $\langle m_{i}\rangle_{i\in n}\in\omega^{n}$. \emph{Then there is a} $q\in \omega$ \emph{such that for all sequences} $\langle X_{i}\rangle_{i\in n}$ \emph{of disjoint sets such that} $|X_{i}|\geq q$ \emph{for every} $i\in n$, \emph{if} $\Delta$ \emph{is a partition of} $\{Z\subseteq\bigcup_{i\in n}X_{i}:(\forall i\in n)(|Z\cap\,X_{i}|=m_{i})\}$ \emph{into two pieces there then is a sequence} $\langle Y_{i}\rangle_{i\in n}$ \emph{such that for all} $i\in n,\,Y_{i}\subset X_{i}$ \emph{and} $|Y_{i}|=p$ \emph{and} $\{Z\subseteq\bigcup_{i\in n}Y_{i}:(\forall i\in n)(|Z\cap\,Y_{i}|=m_{i})\}$\
\emph{is a subset of one of the pieces}.

Let $F=\{\phi_{1}\circ \mathrm{fix}_{G}(E),\ldots,\,\phi_{k}\circ \mathrm{fix}_{G}(E)\}$ be a finite subset of $G/\mathrm{fix}_{G}(E)$. We will find a $Y$ satisfying the condition from definition 1 above. We first note that it is easy to verify that
\setcounter{equation}{0}
\begin{equation}
\label{part04:eqn1b} (\forall\phi,\,\psi\in G)(\phi\circ \mathrm{fix}_{G}(E)=\psi\circ \mathrm{fix}_{G}(E)\Leftrightarrow\phi(E)=\psi(E))
\end{equation}
Let $E^{\prime}=\phi_{1}(E)\cup\cdots \cup\phi_{k}(E)$.\index{Author}{Halpern} $E$ and $E^{\prime}$ are finite and therefore there are natural numbers $j_{0},\,\ldots,\,j_{n-1}$ such that $E\cup E^{\prime}\subseteq B_{j_{0}}\cup\cdots \cup B_{j_{n-1}}$. For $i\in n$, define $E_{i}=B_{j_{i}}\cap E,\,E_{i}^{\prime}=B_{j_{i}}\cap E^{\prime},\,m_{i}=|E_{i}|$ and let $p= \max\{|E_{i}^{\prime}|:i\in n\}$. The integers $n,\,p$ and the sequence $\langle m_{i}\rangle_{i\in n}\in \omega^{n}$ satisfy the hypothesis of Rado's corollary. Therefore there is a $q\in \omega$ satisfying the conclusion. For $i\in n$, choose a set $X_{i}\subseteq B_{j_{i}}$ so that $|X_{i}|=q$. Let $Y= \{\phi\circ \mathrm{fix}_{G}(E):\phi(E)\subseteq\bigcup_{i\in n}X_{i}\}$. We now show that this $Y$ satisfies the condition of definition~\ref{part04:def1}.

Let $\Delta=\{\Delta_{1},\,\Delta_{2}\}$ be a partition of $Y$ into two pieces. Let $\Gamma_{1}=\{\phi(E):\phi\circ \mathrm{fix}_{G}(E)\in\Delta_{1}\}$ and $\Gamma_{2}=\{\phi(E):\phi\circ \mathrm{fix}_{G}(E)\in\Delta_{2}\}$. By (\ref{part04:eqn1b}), $\Gamma=\{\Gamma_{1},\,\Gamma_{2}\}$ is a partition of $\{Z\subseteq\bigcup_{i\in n}X_{i}:(\forall i\in n)\,|Z\cap\,X_{i}|=m_{i}\}$. By Rado's corollary, there is a sequence $\langle Y_{i}\rangle_{i\in n}$ such that $|Y_{i}|=p$ for $i\in n$ and $W= \{Z\subseteq\bigcup_{i\in n}Y_{i}:(\forall i\in n)(|Z\cap Y_{i}| =m_{i}\}$ is either a subset of $\Gamma_{1}$ or a subset of $\Gamma_{2}$. We assume without loss of generality that $W\subseteq\Gamma_{1}$. Let $\eta \in G$ be a permutation such that $\eta(E_{i}^{\prime})\subseteq Y_{i}$ for $i\in n$. (This is possible since $|Y_{i}|=p\geq|E_{i}^{\prime}|$.) Then (since $\phi_{t}(E_{i})\subseteq E_{i}^{\prime}$ for $1\leq t\leq k) \eta\phi_{t}(E_{i})\subseteq Y_{i}$ for $i\in n$ and $1\leq t\leq k$. Therefore, for $1\leq t\leq k,\, \eta\phi_{t}(E)\subseteq\bigcup_{i\in n}Y_{i}$. Using this and the fact that $|\eta\phi_{t}(E)\cap Y_{i}|=|\eta\phi_{t}(E_{i})|=|E_{i}|=m_{i}$ for $1\leq t\leq k$ and $i\in n$ we conclude that $\eta\phi_{t}(E)\in\Gamma_{1}$. Hence $\eta\phi_{t}\circ\mathrm{fix}_{G}(E)\in\Delta_{1}$.

\section*{NOTE 137}

We shall give a proof that [1 BU] (In every vector space over $\mathbb{Q}$, every generating set includes a basis.) implies 1, AC. Bleicher\index{Author}{Bleicher} \cite{bib:1964} has shown that

\noindent [1 BU]implies $C(\infty,\,<\aleph_{0})$ (62), so it is sufficient to prove that [1 BU] implies MC (67). (The proof is due to K. Keremedis.)

Let $X=\{x_{i}:i\in I\}$ be a set of pairwise disjoint infinite sets. For each $i\in I$, let $V_{i}$ be the set of all functions from $x_{i}$ to $\mathbb{Q}$ such that for all $f\in V_{i},\,\exists r\in \mathbb{Q}$ such that $\{u\in x_{i}:f(u)\neq r\}$ is finite. Then $V_{i}$ is a vector space. Let $V$ be the weak direct product of the $V_{i}$'s. For each $i\in I$, define
\begin{gather*}
G_{i}=\{F\in V\,:\,(\forall j)[(\forall u\in x_{j})(j\neq i\rightarrow F_{j}(u)=0)\wedge\\
(j=i\rightarrow(\exists u\in x_{j})(F_{j}(u)=1\wedge(\forall v\in x_{j})(v\neq u\rightarrow F_{j}(v)=0)\vee\\
(F_{j}(u)=0\wedge(\forall v\in x_{j})(v\neq u\rightarrow F_{j}(v)=1)))]\}
\end{gather*}
Let $G=\bigcup\{G_{i}:i\in I\}$. Then $G$ is a generating set for $V$. Let $B\subseteq G$ be a basis for $V$. For each $i\in I$, let $H^{i}$ be the function in $V$ such that $\forall u\in x_{j}$,
\begin{equation*}
H^{i}(u)=\left\{\begin{array}{l}
0,\,j\neq i\\
1,\,j=i
\end{array}\right.
\end{equation*}
For each $i\in I,\,H^{i}$ can be uniquely represented as a linear combination of a finite number of basis elements. It follows from the definition of $G_{i}$ that each basis element uniquely specifies an element in $x_{i}$. Thus, each $H^{i}$ uniquely specifies a finite subset of $x_{i}$.

\section*{NOTE 138}

In this note we include a proof due to K. Keremedis that form 375 (Tietze-Urysohn Extension Theorem) implies form 119 ($C(\aleph_{0}$,uniformly orderable with order type of the integers).

Fix $\mathcal{A}=\{A_{i}:i\in \omega\}$ as in form 19, where the $A_{i}$ are pairwise disjoint. Let $\bar{A}_{i}=\{a_{i}\}\cup A_{i}\cup \{b_{i}\}$, where $a_{i},\,b_{i}$ are chosen so that $a_{i},\, b_{i}\not\in A_{i}$ and the $\bar{A}_{i}$ are pairwise disjoint. Declare $a_{i}$ to be less than every member of $A_{i}$ and $b_{i}$ to be greater than every member of $A_{i}$. Put $X=\cup\{\bar{A}_{i}:i\in \omega \}$. Order $X$ by requiring: elements of $\bar{A}_{i}$ are ordered in $X$ the way they are ordered in $\bar{A}_{i}$, otherwise if $x\in\bar{A}_{i}$ and $y\in\bar{A}_{j}$ then $x$ is less than $y$ iff $i<j$. Put $\bar{X}=X\cup \{a\},\,a\not\in X$ and declare $a$ to be greater than every member of $X$. Clearly $\bar{X}$ with the order topology is a compact $\mathrm{T}_{2}$ space and consequently a $\mathrm{T}_{4}$ space. Furthermore, $G=A(=\{a_{i}: i\in \omega \})\cup B(=\{b_{i}:i\in \omega \})\cup
\{a\}$ is a closed subset of $\bar{X}$ and the function $f:G\rightarrow[0,\,1]$ given by $f(a)=0$ and for $i\in \omega\  f(a_{i})=1/(i+2)$, $f(b_{i})=1/(i+1)$ is clearly continuous. Hence, $f$ has a continuous extension $\bar{f}: \bar{X}\rightarrow[0,1]$. Put $f_{i}=\bar{f}|\bar{A}_{i}$. As $\bar{f}$ is continuous it follows that $f_{i}$ is also continuous. Since $\bar{A}_{i}$ is compact it follows that $f_{i}$ attains both its minimum and maximum values. We may assume that $f_{i}$ attains its minimum at $a_{i}$, its maximum at $b_{i}$ and, at every other point of $A_{i}$, it takes on values other than $f_{i}(a_{i})$ and $f_{i}(b_{i})$. Let $C_{i}=f_{i}^{-1}((y_{i},\,1/(1+i)))$ where $y_{i}= \frac{2i+3}{2(i+1)(i+2)}$ is the midpoint of the interval $[1/(2+i),\,1/(1+i)]$. Clearly $C_{i}\neq\emptyset$ and by the completeness of $\bar{A}_{i}$ and the continuity of $f_{i}$ it follows that $c_{i}= \inf(C_{i})\in A_{i}$. Hence, $c=\{c_{i}:i\in \omega \}$ is a choice set for $\mathcal{A}$, finishing the proof of the theorem.

\section*{NOTE 139}

We give a proof that form 67 (MC, Multiple Choice Axiom) implies form 375 (Tietze-Urysohn Extension Theorem). To the best of our knowledge, this result is due to David Pincus.

Suppose $X$ is a normal topological space. Then, if $A$ and $B$ are disjoint closed sets, there exist disjoint open sets $C$ and $D$ such that $A\subseteq C$ and $B\subseteq D$. MC gives us a rule for choosing a finite number of pairs $(C,\,D)$ with the above property. Since the intersection of a finite number of open sets is open, MC implies that there is a function $F$ such that for each pair of disjoint closed sets $(A,\,B)$, $F(A,\,B)=(C,\,D)$, where $C$ and $D$ are disjoint open sets such that $A\subseteq C$ and $B\subseteq D$. Thus, using the standard proof of Urysohn's lemma (see, for example, Kelley \cite{bib:1955}, p 115), there is a rule for choosing a function $f:X\rightarrow[0,\,1]$ which is zero on $A$ and one on $B$. Given this rule for choosing Urysohn's function, the standard proof of the Tietze-Urysohn Extension Theorem can be given in $\mathrm{ZF}^{0}$ (See, for example, Kelley \cite{bib:1955}, p 242.).

\section*{NOTE 140}

In this note we give a proof that form 214 (For every family $A$ of infinite sets, there is a function $f$ such that $\forall y\in A,\,f(y)\subseteq y$ and $|f(y)|=\aleph_{0}.)$ implies form 152 (Every non-well orderable set is the union of a pairwise disjoint, well orderable family of denumerable sets.)

Let $X$ be a set which cannot be well ordered and let $Y$ be the set of all infinite subsets of $X$. Form 214 implies that there is a function $f$ such that for each $y\in Y,\,f(y)\subseteq y$ and $|f(y)|=\aleph_{0}$. We define a function $g$ on the ordinals so that $g(0)=f(X)$ and $g( \alpha)=f(X\,\backslash\, \bigcup_{\beta<\alpha}g(\beta))$. Since $X$ is a set, there is an ordinal $\gamma$ such that $X\, \backslash\, \bigcup_{\beta<\gamma}g(\beta)$ is finite. We can use this function $g$ to cover $X$ with a well ordered family of pairwise disjoint denumerable sets.

\section*{NOTE 141}

In this\index{Author}{Kelley} note we give the definitions for forms [1 CE] through [1 CH], [67
K] through [67 T], [126 G] through [126 K], 173, [173 A], 232, [232 A] through [232 G], 383, and [383 A] through [383 C].

\begin{definition*} Let $(X,\,T)$ be a topological space.
\begin{enumerate}
\item[1.] $X$ is \emph{normal}\index{Subject}{normal space} if for every pair of disjoint closed sets $A$ and $B$, there exist disjoint open sets $C$ and $D$ such that $A\subseteq C$ and $B\subseteq D$.
\item[2.] $X$ satisfies \emph{Urysohn's lemma}\index{Subject}{Urysohn's lemma} iff it satisfies: If $A,\,B$ are closed and disjoint in $X$ then there exists a continuous function $f:X\rightarrow[0,\,1]$ separating $A$ and $B$. i.e., $A\subseteq f^{-1}(0)$ and $B\subseteq f^{-1}(1)$. (In the remainder of this note we shall refer to a space that satisfies Urysohn's lemma as a $U$ \emph{space}\index{Subject}{U space}.)
\item[3.] $X$ satisfies the \emph{Tietze-Urysohn extension theorem}\index{Subject}{Tietze-Urysohn extension theorem} iff it satisfies: If $A$ is closed in $X$ and $f:A\rightarrow[0,\,1]$ is a continuous function, then there exists a continuous extension $\bar{f}$ of $f$ to all of $X$ taking values in $[0,1]$. (We shall refer to a space satisfying the Tietze-Urysohn extension theorem as a $T$ \emph{space}\index{Subject}{T space}.)
\item[4.] The \emph{disjoint union}\index{Subject}{disjoint union} of the disjoint family $\{(X_{i},\,T_{i}):i\in k\}$ of topological spaces is the set $X=\bigcup\{X_{i}:i\in k\}$ such that $O$ is open in $X$ iff $O\cap X_{i}$ is open in $X_{i}$ for all $i\in k$.
\item[5.] A family $\mathcal{K}$ of subsets of a topological space $(X,\,T)$ is \emph{locally finite} (respectively \emph{locally countable}) iff each point of $X$ has a neighborhood meeting only a finite number (respectively countable number) of elements of $\mathcal{K}$.
\item[6.] $X$ is \emph{paracompact} (respectively \emph{para-Lindel\"{o}f}) iff $X$ is $\mathrm{T}_{2}$ and every open cover $\mathcal{U}$ of $X$ has a locally finite (respectively locally countable $\rangle$ open refinement $\mathcal{V}$, (that is, $\mathcal{V}$ is a locally finite open cover of $X$ and every member of $\mathcal{V}$ is included in a member of $\mathcal{U}$).
\item[7.] A family $\mathcal{K}$ of subsets of $X$ is \emph{point finite} (respectively \emph{point countable}) iff each element of $X$ belongs to only finitely many (respectively countably many) members of $\mathcal{K}$.
\item[8.] $X$ is \emph{metacompact} (respectively \emph{meta-Lindel\"{o}f}) iff each open cover $\mathcal{U}$ of $X$ has an open point finite (respectively point countable) refinement.
\item[9.] A set $C\subseteq \mathcal{P}(X)$ is $\sigma$-\emph{locally finite} (respectively $\sigma$-\emph{point finite}, $\sigma$-\emph{disjoint}, $\sigma$-\emph{discrete}) if $C= \bigcup_{n\in\omega}C_{n}$ where each $C_{n}$ is locally finite (respectively, point finite, pairwise disjoint, discrete).
\item[10.] Assume that $\gamma$ is a well ordered cardinal number, that is, $\gamma$ is an initial ordinal. A set $C\subseteq \mathcal{P}(X)$ is $\gamma$-\emph{locally finite} (respectively $\gamma$-\emph{point finite}, $\gamma$-\emph{disjoint}, $\gamma$-\emph{discrete}) if $C= \bigcup_{\alpha\in\gamma}C_{\alpha}$ where each $C_{\alpha}$ is locally finite (respectively, point finite, pairwise disjoint, discrete).
\item[11.] An open cover $\mathcal{U} =\{U_{i}:i\in k\}$ of $X$ is \emph{shrinkable} iff there exists an open cover $\mathcal{V}=\{V_{i}:i\in k\}$ of non-empty sets such that $\bar{V}_{i}\subseteq U_{i}$ for all $i\in k$. $\mathcal{V}$ is also called a \emph{shrinking} of $\mathcal{U}$.
\item[12.] $X$ is a \emph{PFCS space} iff every point finite open cover of $X$ is shrinkable.
\item[13.] $D\subseteq X$ is \emph{discrete} if every $x\in X$ has a neighborhood $U$ such that $|U\,\cap\, D|=1$. A set $C\subseteq \mathcal{P}(X)$ is \emph{discrete} if $\forall x\in X$, there is a neighborhood $U$ of $x$ such that $U\cap A\neq\emptyset$ for at most one element $A\in C$. (Thus, $D\subseteq X$ is discrete\index{Subject}{$\sigma$-discrete} iff $\{\{d\}:d\in D\}$ is discrete.)
\item[14.] $X$ is\index{Subject}{locally finite}\index{Subject}{meta-Lindel\"{o}f}\index{Subject}{metacompact}\index{Subject}{para-Lindel\"{o}f}\index{Subject}{paracompact}\index{Subject}{point countable}\index{Subject}{point finite}\index{Subject}{$\sigma$-point finite}\index{Subject}{shrinkable} cwH (\emph{collectionwise Hausdorff})\index{Subject}{collectionwise Hausdorff} iff if $X$ is $\mathrm{T}_{2}$ and for every closed and discrete subset $D$ of $X$ there exists a family $U=\{U_{d}:d\in D\}$ of disjoint\index{Subject}{$\sigma$-disjoint} open sets such that $d\in U_{d}$ for all $d\in D$.
\item[15.] $X$ is cwH(B) (\emph{collectionwise Hausdorff with respect to the base} $B$) iff if $X$ is $\mathrm{T}_{2}$ and for every closed and discrete subset $D$ of $X$ there exists a family $U=\{U_{d}:d\in D\}$ of disjoint open sets such that $U\subseteq B$ and $d\in U_{d}$ for all $d\in D$.
\item[16.] $X$ is cwN (\emph{collectionwise normal})\index{Subject}{collectionwise normal} if $X$ is $T_{2}$ and for every discrete $C$ of pairwise disjoint closed subsets of $X$, there exists a family $U=\{U_{A}:A\in C\}$ of disjoint open sets such that $A\subseteq U_{A}$ for all $A\in C$. (cwN($B$) is defined similarly to cwH$(B)$ in 15.)
\end{enumerate}
\end{definition*}

\section*{NOTE 142}

In this note we prove that form 234 (There is a non-Ramsey set.) implies form 282 $(w\not\rightarrow (\omega)^{\omega}$).

Assume that there is a non-Ramsey set $A$. The proof will be by contradiction. We therefore assume that $\omega \rightarrow(\omega)^{\omega}$. Then (using the definitions in note 97) for every $F:\omega^{\omega}\rightarrow\{0,1\}$ there is an infinite subset $X$ of $\omega$ such that either
\begin{enumerate}
\item[(i)] $(\forall g\in X^{\omega})(F(g)=0)$ or
\item[(ii)] $(\forall g\in X^{\omega})(F(g)=1)$.
\end{enumerate}

Using $A$ we define a function $F:\omega^{\omega}\rightarrow\{0,1\}$ by
\begin{equation*}
F(g)=\left\{\begin{array}{l}
0\ \mbox{if}\ g\ \mbox{is strictly increasing and range}(g)\in A\\
1\ \mbox{otherwise}
\end{array}\right.
\end{equation*}
Let $X$ be the infinite subset of $\omega$ guaranteed by $\omega\rightarrow(\omega)^{\omega}$. For each infinite subset $Y$ of $\omega$, let $g_{Y}\in \omega^{\omega}$ be the unique strictly increasing function whose range is $Y$. Then $F(g_{Y})=0$ if and only if $Y\in A$. Either (i) or (ii) must be true of $X$ and $F$. If (i) holds then for every infinite $Y\subseteq X,\,g_{Y}\in X^{\omega}$ so $F(g_{Y})=0$ and hence $Y\in A$. On the other hand, if (ii) holds then for every infinite $Y\subseteq X,\,g_{Y}\in X^{\omega}$ so $F(g_{Y})=1$ and hence $Y\not\in A$. This contradicts the assumption that $A$ is non-Ramsey.

\section*{NOTE 143}

Form [0 AM] is a form of the Hahn-Banach Theorem for separable vector spaces. We shall prove by induction that it is provable in $\mathrm{ZF}^{0}$.

Assume $V$ is a separable vector space, $D=\{d_{0},\,d_{1},\,\cdots\}$ is a countable dense subset of $V$, and $p:V\rightarrow \mathbb{R}$ satisfies $p(x+y)\leq p(x)+p(y)$, and $(\forall t\geq 0)(\forall x\in V)(p(tx)=tp(x))$. Also assume that $p$ satisfies the following condition:
\begin{equation*}
\tag{$\ast$} (\forall\epsilon>0)(\forall x\in V-S)(\exists y\in D)\ \mbox{such\ that}\ (p(x-y)<\epsilon\ \mbox{and}\ p(y-x)<\epsilon.
\end{equation*}
Let $S$ be $a$ subspace of $V$ and let $f$ be a linear function from $S$ to $\mathbb{R}$ such that for all $x\in S,\,f(x)\leq p(x)$. We shall define by induction a countable\index{Subject}{locally countable} set of linear functions $G=\{g_{n}:n\in \omega\}$ such that $g_{n}\subseteq g_{n+1}$ and such that the linear extension of $f,\,f^{\ast}=\bigcup\, G$, can be extended to all of $V$ by using the condition $(\ast)$.

We define $g_{0}=f$. Suppose $g_{n}$ has been defined and $S_{n}$ is the domain of $G_{n}$. In what follows, we shall let $d=d_{n+1}$ and $g=g_{n+1}$. If $d$ is in the domain of $g_{n}$, then define $g=g_{n}$. Suppose that $d$ is not in the domain of $g_{n}$. We would like to define $g$ so that for any $t\in \mathbb{R}$ and any $x\in S_{n},\ g(td+x)=tg(d)+g(x)\leq p(td+x)$. Then $tg(d)\leq p(td+x)-g(x)$. If $t=1,\, g(d)\leq p(d+x)-g(x)$ and if $t=-1,\,g(d)\geq g(x)-p(x-d)$. Let $s= \sup_{x\in S}(g(x)-p(x-d))$ and let $i= \inf_{x\in S}(p(d+x)-g(x))$. Then $s\leq g(d)\leq i$. Define $g(d)$ to be the mid-point of the interval $[s,\,i]$. For any $x\in S_{n}$, We define $g(x)=g_{n}(x)$ and, for each $t\in \mathbb{R}$, we define $g(td+x)=tg(d)+g(x)$. Let $f^{\ast}=\bigcup\{g_{n}:n\in \omega\}$. It is clear that $f^{\ast}$ is a linear function, $f^{\ast}(x)\leq p(x)$ for all $x$ in the domain of $f^{\ast}$, and $D$ is a subset of the domain of $f^{\ast}$.

Let $\epsilon>0$ and $x\in V$. Then, by $(\ast)$, there is a $y\in D$ such that $p(x-y)<\epsilon$ and $p(y-x)<\epsilon$. If $f^{\ast}(x-y)\geq 0$, then $0\leq f^{\ast}(x-y)\leq p(x-y)<\epsilon$. If $f^{\ast}(x-y)<0$, then $0<-f^{\ast}(x-y)=f^{\ast}(y-x)\leq p(y-x)<\epsilon$. Consequently, $|f^{\ast}(x-y)|<\epsilon$ and $f^{\ast}$ can be extended to all of $V$.

\section*{NOTE 144}

A proof that dependent choice (form 43) holds in FM models of $\mathrm{ZF}^{0}$ in which the set of supports $S$ is closed under countable unions. This includes models where $S$ consists of all countable subsets of the set $A$ of atoms.

Assume that $\mathcal{N}$ is such a model and that $R$ is a relation in $\mathcal{N}$ such that for all $x\in\mathrm{dom}(R)$ there is a $y$ such that $x\ R\ y$. Then, in the ground model (where AC holds) there is a sequence $(x_{n})_{n\in\omega}$ such that $(\forall n\in\omega)(x_{n}\  R\ x_{n+1})$. If for each $n\in \omega,\,x_{n}$ has support $E_{n}$, then $\bigcup_{n\in \omega}E_{n}$ is a support for $(x_{n})_{n\in\omega}$. Hence, this sequence is in $\mathcal{N}$.

\section*{NOTE 145}

A proof that form 287 (the Hahn-Banach theorem for separable normed linear spaces) implies form 222 (there is a non-principal measure on $\omega$). This result and its proof are due to M. Morillon.

Let $E$ be the vector space of bounded sequences of reals. Let $M$ be a norm on $E$ such that the normed vector space $(E,\,M)$ is separable: for example, if $x= (x_{k})_{k\in \mathbb{N}}\in E$ let $M(x)= \sum\nolimits_{k\in\mathbb{N}}\frac{|x_{k}|}{(k+1)^{2}}$. Here, the normed vector space $(E,\,M)$ is separable because the set of sequences in $\mathbb{Q}^{(\mathbb{N})}$ which end in an infinite sequence of zeros is dense in E.

For every $x=(x_{k})_{k\in \mathbb{N}}\in E$, let $p(x)=\varlimsup\{x_{k};k\in \mathbb{N}\}$. Hence, we obtain a sublinear functional $p:E\rightarrow \mathbb{R}$. Using form 287, let $g:E\rightarrow \mathbb{R}$ be a linear functional such that $g\leq p$. We have $p(1)=1$ and $p(-1)=-1$. Hence, $g(1)=1$. For every finite subset $A\subseteq \mathbb{N}$ we have $p(1_{A})=p(-1_{A})=0$, hence, $g(1_{A})=0$. Here $1_{A}$ denotes the characteristic function of the set $A$.

For every $A\subseteq \mathbb{N}$, lf $m(A)=g(1_{A})$, then $m:\mathcal{P}(\omega)\rightarrow[0,\,1]$ is a non-principal measure on $\omega$.

\section*{NOTE 146}

Form [43 W] is due to Claus-Peter Wirth. In this note we prove that it is equivalent to form 43. Form [43 W] is: Let $M$ be a non-empty set partially ordered by $<$ without $<$-minimal elements. Then there is a subset $C$ of $M$ which is linearly ordered by $<$ such that $C$ has no $<$-minimal element. Form 43 is the principle of dependent choice, DC. To prove [43 W] from DC let $X=M$ and define $R$ by $x\,R\,y$ if and only if $y<x$. Then the sequence $(x_{n})_{n\in \omega}$ which is guaranteed by DC is a linearly ordered subset of $M$ with no $<$-minimal elements.

To prove [43 W] implies DC, let $R$ be a relation on a set $X$ such that $\forall x\in X,\,\exists y\in X$ such that $x\,R\,y$. Let $M$ be the set of all finite sequences $(x_{1},\,x_{2},\,\ldots,\,x_{n})$ of elements of $X$ such that $x_{1}\ R\ x_{2}\ R\ \ldots\ R\ x_{n}$ and define a relation $<$ on $M$ by $(x_{1},\, \ldots,\,x_{n})<(y_{1},\,\ldots,\,y_{m})$ if and only if $(x_{1},\,\ldots,\,x_{n})$ is $a$ proper extension of $(y_{1},\,\ldots,\,y_{m})$. $M$ has no $<$-minimal elements so by (ii) there is a $<$ chain $C$ in $M$ with no minimal element. From $C$ an infinite sequence $(x_{n})_{n\in \omega}$ can be constructed with the property that $x_{n}\ R\,x_{n+1}$ for all $n\in \omega$.

\section*{NOTE 147}

Definitions for measures on Boolean algebras\index{Subject}{measure on a Boolean algebra}.

\begin{definition*} Assume that $\mathcal{B} =(B,\,\wedge,\,\vee,\,\mathbf{1,\,0})$ is a Boolean algebra.
\begin{enumerate}
\item[1.] A \emph{measure} or \emph{additive real valued measure} on $\mathcal{B}$ is a function $m:\mathcal{B} \rightarrow \mathbb{R}$ such that $m(\mathbf{1})=1,\,(\forall x\in B)(m(x)\geq 0)$, and $(\forall x,\,y\in B)(x\wedge y=\mathbf{0}\rightarrow m(x\vee y)= m(x)+m(y))$. (It follows that $m(\mathbf{0})=0$.)
\item[2.] A \emph{two valued measure}\index{Subject}{two valued measure} on $\mathcal{B}$ is a measure on $\mathcal{B}$ whose range is $\{0,1\}$.
\item[3.] If $\mathcal{B} = (\mathcal{P}(X),\,\cap,\,\cup,\,X,\,\emptyset)$ then a measure $m$ on $\mathcal{B}$ is \emph{non-principal} if $m(\{x\})=0$ for every $x\in X$.
\end{enumerate}
\end{definition*}

\section*{NOTE 148}

We give a proof that forms 154 and [154 A] are equivalent. First, it is clear that 154 implies [154 A]. To prove the converse, assume [154 A] holds. Let $\{X_{i}:i\in \omega\}$ be a countable family of pairwise disjoint compact $\mathrm{T}_{2}$ spaces and let $Y$ be the one point compactification of $\bigcup_{i\in\omega}X_{i}$. Form [154 A] implies that $Y^{\omega}$ is compact. Since $\prod_{i\in \omega}X_{i}$ is a closed subspace of $Y^{\omega}$, it is also compact.

%%%%%%part05

\chapter{References for Relationships Between Forms}\label{pt05:part05}

\section*{Introduction}

To explain this part of the book, it is helpful to consider some examples. Five typical lines are below.

\begin{table}[h]
\begin{tabular}{llll}
43 & 63 & (3) & Felgner\index{Author}{Felgner} [\textbf{1971a}] (M2)\\
85 & 78 & (3) & Brunner\index{Author}{Brunner} [\textbf{1980}] (N3T)\\
67 & 80 & (5) & Fraenkel\index{Author}{Fraenkel} [\textbf{1922}], Jech \cite{bib:1973b} (N2)\\
87 & 86 & (1) & Jech\index{Author}{Jech} \cite{bib:1973b} p 120 thm 8.1, Levy\index{Author}{Levy} [\textbf{1963a}], note 9\\
\end{tabular}
\end{table}

The first two numbers are form numbers. The third number which occurs in parentheses is a code number for the status of the implication. For example, in the first line above, the number 3 means that the implication ``form 43 implies form 63'' is not provable in Zermelo-Fraenkel (ZF) set theory. The other code numbers that we use are
\begin{enumerate}
\item[0] for ``The status of the implication is unknown.''
\item[1] for ``The implication is provable.''
\item[2] for ``The implication is provable and this fact follows from implications whose code numbers are 1.''
\item[3] for ``The implication is not provable in ZF.''
\item[4] for ``The implication is not provable in ZF and this fact follows from implications whose code numbers are 1 or 3.''
\item[5] for ``The implication is not provable in $\mathrm{ZF}^{0}$.''
\item[6] for ``The implication is not provable in $\mathrm{ZF}^{0}$ and this fact follows from implications whose code numbers are odd.''
\item[7] for ``The status of the implication is unknown and we have found the question asked somewhere in the literature.''
\end{enumerate}
(See the introduction for the definition of $\mathrm{ZF}^{0}$.)

The implication codes are odd numbers for the most part. In a typical line following the code number, is a reference to a bibliography item and possibly a note and, in the case of a non-implication, a model number. (We normally only list one model even if there are several models with this property.) In the first example above, Felgner [\textbf{1971a}] refers to the Bibliography and M2 refers to the model $\mathcal{M}2$ of ZF from our models section. ($\mathcal{M}2$ is a model in which form 43 is true and form 63 is false.)

In the second example above we see that form 85 implies form 78 is not provable in ZF set theory; the reference is Brunner [\textbf{1980}]; and the model is the permutation model $\mathcal{N}3$. $\mathcal{N}3$ is a model of $\mathrm{ZF}^{0}$. However, many of the results obtained for permutation models can be transferred to ZF models because of the work of Jech/Sochor\index{Author}{Jech/Sochor} [\textbf{1966a}], [\textbf{1966b}], and Pincus\index{Author}{Pincus} [\textbf{1971}], [\textbf{1972a}], [\textbf{1974c}], [\textbf{1977a}]. If a result is transferable, the numeral 3 will occur in the line to mean that the implication is not provable in ZF set theory. Following the reference, we give the permutation model with a ``$\mathrm{T}$'' following its number to indicate the result is transferable. (See note 18 for a list of forms that are transferable and note 103 for a summary of the transfer results of Jech/Sochor\index{Author}{Jech}\index{Author}{Jech/Sochor} [\textbf{1966a}] and [\textbf{1966b}] and Pincus\index{Author}{Pincus} [\textbf{1971}], [\textbf{1972a}], and [\textbf{1977a}].)

The letters ``IC'' following a model number indicate that an inaccessible cardinal was used to construct the model. In some cases, the reference is just a model. In these cases it is clear from the model itself why the implication does not hold. (See, for example, the line at position (317,3) and model $\mathcal{N}$38.)

It follows from example 3 above that form 67 implies form 80 is not provable in $\mathrm{ZF}^{0}$, the references are Fraenkel\index{Author}{Fraenkel} [\textbf{1922}] and Jech \cite{bib:1973b}, and the permutation model is $\mathcal{N}2$. The last example above tells us that form 87 implies form 86 in $\mathrm{ZF}^{0}$. The references are Jech \cite{bib:1973b} (p 120 theorem 8.1), Levy\index{Author}{Levy} [\textbf{1964}], and additional information can be found in note 9 in the notes section of the book. When note 120 is the reference, we state which part of the note applies. For example, at position (152,67) the reference is note 120(41). This means that we are using item 41 in note 120 ($67+122\leftrightarrow 1$, where form 1 is AC) to obtain the independence result. Form 122 is true in $\mathcal{N}$15, so by note 120 part 41, form 67 must be false. (AC is false in all our models.) In addition, form 152 is true in this model. Consequently, there is a 5 at position (152,67) below.

If the reference is note 18, we mean that either the first form listed is true in every permutation (FM) model, the second form is false in each such model, or the second form is false in every Cohen model. We have several forms that depend on parameters. (See note 18 for a list of these forms.) Sometimes the relationships between such forms depend on the value of the parameters. For example, we have the following line below:
\begin{equation*}
120\qquad 33\qquad (0)\qquad \mathrm{note}15\ (\mathrm{see}\ \mathrm{M}47(n, M))
\end{equation*}
Both forms 120 and 33 depend upon parameters. The notation means that for certain values of the parameters, form 120 is true in the model $M47(n, M)$ and form 33 is false. The reader should look at the model to determine when this holds.

Each of our weak forms of the Axiom of Choice is implied by AC (form 1), but for some of our forms, it is unknown whether they imply AC. Each one of these forms, $n$, is listed below in the form:
\begin{equation*}
n\qquad 1\qquad (7)\qquad \mathrm{reference}
\end{equation*}
Similarly, for some forms, it is unknown whether they are provable in $\mathrm{ZF}^{0}$. Each one of these forms, $m$, is listed below as:
\begin{equation*}
0\qquad m\qquad (7)\qquad \mathrm{reference}
\end{equation*}
In these two cases, for those forms for which we have not found a reference in the literature we use ``This project'' as the reference.

\section*{References for relationships between forms}

\begin{table}[h]
\begin{tabular}{llll}
0 & 1 & (4) & Cohen\index{Author}{Cohen} [\textbf{1963}],[\textbf{1965}],[\textbf{1966}], Vop\u{e}nka\index{Author}{Vop\u{e}nka} [\textbf{1967}](M1)\\
2 & 1 & (7) & G. Moore\index{Author}{G. Moore} [\textbf{1982}], p 322 , Mathias\index{Author}{Mathias} [\textbf{1979}] prob 1354\\
3 & 1 & (3) & Sageev\index{Author}{Sageev} [\textbf{1975}], Halpern/Howard\index{Author}{Halpern}\index{Author}{Halpern/Howard}\index{Author}{Howard} [\textbf{1974}] (M6, N9)\\
7 & 1 & (7) & G. Moore [\textbf{1982}] p 149, Banaschewski/Moore\index{Author}{Banaschewski}\index{Author}{Banaschewski/Moore} [\textbf{1990}]\\
14 & 1 & (4) & Halpern [\textbf{1964}], Halpern/Levy\index{Author}{Halpern/Levy}\index{Author}{Levy} [\textbf{1971}], Pincus\index{Author}{Pincus} [\textbf{1977b}](M1)\\
15 & 1 & (4) & Halpern/Levy [\textbf{1964}],[\textbf{1971}] (M1)\\
20 & 1 & (7) & G. Moore [\textbf{1982}]\\
21 & 1 & (7) & G. Moore [\textbf{1982}]\\
22 & 1 & (7) & This project\\
30 & 1 & (4) & Jech\index{Author}{Jech} \cite{bib:1973b} p 71, Halpern/Levy [\textbf{1964}] (M1)\\
40 & 1 & (3) & Jech \cite{bib:1973b} p 127, Pincus [\textbf{1969}] (M2)\\
43 & 1 & (4) & Jech \cite{bib:1973b} p 83 prob 5.26, Feferman\index{Author}{Feferman} [\textbf{1964}], Mostowski\index{Author}{Mostowski}\\
&&&[\textbf{1948}], Pincus [\textbf{1977b}] (M2)\\
50 & 1 & (7) & Banaschewski \cite{bib:1988a}, Bell\index{Author}{Bell} \cite{bib:1983a}, Mathias [\textbf{1979}] p 124\\
54 & 1 & (7) & This project\\
55 & 1 & (7) & Mathias [\textbf{1979}] prob 1340\\
56 & 1 & (7) & This project\\
59 & 1 & (7) & This project\\
60 & 1 & (3) & L\"{a}uchli\index{Author}{L\"{a}uchli} [\textbf{1964}], Jech \cite{bib:1973b} p 82  prob 5.22 (M3)\\
61 & 1 & (4) & Levy [\textbf{1962}] (N6T)\\
65 & 1 & (6) & Pincus [\textbf{1972b}],\cite{bib:1972c1} (N2)\\
66 & 1 & (7) & 66 implies AC in ZF\\
67 & 1 & (6) & 67 implies AC in ZF - Jech \cite{bib:1973b} p 133,\,Felgner/Jech\index{Author}{Felgner/Jech} [\textbf{1973}],\\
&&&Levy [\textbf{1962}] (N2)\\
68 & 1 & (7) & Howard [\textbf{1985}]\\
81 & 1 & (7) & This project\\
89 & 1 & (6) & 89 implies AC in ZF - Jech \cite{bib:1973b} p 133  (N2)\\
90 & 1 & (6) & 90 implies AC in ZF - Jech \cite{bib:1973b} p 133  (N2)\\
91 & 1 & (6) & 91 implies AC in ZF - Jech \cite{bib:1973b} p 133  note 75 (N2)\\
95 & 1 & (6) & 95 implies AC in ZF (see (95,67) entry) (N6)\\
100 & 1 & (7) & Banaschewski/Moore [\textbf{1990}]\\
101 & 1 & (7) & G. Moore [\textbf{1982}] p 10,\,H.\,Rubin/J. Rubin\index{Author}{H. Rubin/ J. Rubin} \cite{bib:1985q} p xiii, Levy\\
&&&[\textbf{1963a}], Pelc\index{Author}{Pelc} [\textbf{1978}], Banaschewski/Moore [\textbf{1990}]\\
103 & 1 & (7) & This project\\
105 & 1 & (7) & This project\\
109 & 1 & (7) & 109 implies AC in ZF - Blass\index{Author}{Blass} [\textbf{1984a}]\\
110 & 1 & (7) & This project\\
115 & 1 & (6) & Brunner\index{Author}{Brunner} [\textbf{1985a}] (N1)\\
116 & 1 & (6) & Brunner [\textbf{1985a}] (N1)\\
117 & 1 & (7) & This project\\
118 & 1 & (3) & Haddad/Morillon\index{Author}{Haddad/Morillon} [\textbf{1990}] (M1)\\
129 & 1 & (7) & This project\\
133 & 1 & (6) & 133 implies AC in ZF - Brunner [\textbf{1982a}] (N1)\\
136 & 1 & (7) & This project\\
138 & 1 & (7) & Truss\index{Author}{Truss} [\textbf{1984}]\\
143 & 1 & (7) & This project\\
148 & 1 & (7) & This project\\
\end{tabular}
\end{table}

\clearpage

\begin{table}[h]
\begin{tabular}{llll}
149 & 1 & (7) & 149 implies AC in ZF\\
156 & 1\index{Author}{David}\index{Author}{Howard}\index{Author}{Mathias}\index{Author}{Monro}\index{Author}{Pelc} & (7) & This project\\
157 & 1 & (7) & This project\\
158 & 1 & (7) & This project\\
164 & 1 & (6) & 164 implies AC in ZF - Brunner\index{Author}{Brunner} [\textbf{1982a}] (N3)\\
168 & 1 & (7) & Banaschewski/Moore\index{Author}{Banaschewski/Moore} [\textbf{1990}], Pelc [\textbf{1978}]\\
172 & 1 & (7) & This project\\
175 & 1 & (7) & This project\\
176 & 1 & (7) & This project\\
177 & 1 & (7) & This project\\
179 & 1 & (3) & David [\textbf{1980}] (M35$(\epsilon)$)\\
180 & 1 & (7) & This project\\
184 & 1 & (7) & This project\\
187 & 1 & (5) & A. Rubin\index{Author}{A. Rubin} [\textbf{1978}] (N12$(\aleph_{1})$)\\
188 & 1 & (7) & This project\\
192 & 1 & (7) & Mathias [\textbf{1977b}] (due to Blass,\index{Author}{Blass} see [192 A])\\
193 & 1 & (7) & This project\\
195 & 1 & (7) & This project\\
202 & 1 & (5) & 202 implies AC in ZF - Truss [\textbf{1978}] (N12$(\aleph_{1})$\\
205 & 1 & (7) & This project\\
210 & 1 & (7) & This project\\
215 & 1 & (7) & This project\\
220 & 1 & (7) & This project\\
224 & 1 & (7) & This project\\
230 & 1 & (7) & This project\\
234 & 1 & (7) & This project\\
236 & 1 & (7) & This project\\
237 & 1 & (7) & This project\\
238 & 1 & (7) & This project\\
239 & 1 & (7) & This project\\
240 & 1 & (7) & This project\\
244 & 1 & (7) & This project\\
245 & 1 & (7) & This project\\
246 & 1 & (7) & This project\\
247 & 1 & (7) & This project\\
248 & 1 & (7) & This project\\
253 & 1 & (3) & Howard [\textbf{1975}], Blass [\textbf{1977b}] (M15)\\
255 & 1 & (7) & This project\\
256 & 1 & (7) & This project\\
257 & 1 & (7) & This project\\
258 & 1 & (7) & This project\\
259 & 1 & (7) & This project\\
261 & 1 & (7) & This project\\
262 & 1 & (7) & This project\\
264 & 1 & (7) & 264 implies AC in ZF\\
265 & 1 & (7) & This project\\
266 & 1 & (7) & This project\\
267 & 1 & (7) & This project\\
274 & 1 & (7) & This project\\
\end{tabular}
\end{table}

\clearpage

\begin{table}[h]
\begin{tabular}{llll}
275 & 1 & (7) & This project\\
277 & 1\index{Author}{Jech} & (3) & Monro\index{Author}{Monro} [\textbf{1973a}] (M1)\\
278 & 1 & (7) & This project\\
281 & 1 & (7) & This project\\
284 & 1 & (7) & This project\\
286 & 1 & (7) & This project\\
297 & 1 & (7) & This project\\
299 & 1 & (7) & This project\\
300 & 1 & (7) & This project\\
301 & 1 & (7) & This project\\
302 & 1 & (7) & This project\\
303 & 1 & (7) & This project\\
310 & 1 & (7) & This project\\
318 & 1 & (7) & This project\\
319 & 1 & (7) & This project\\
320 & 1 & (7) & This project\\
335 & 1 & (7) & 335 implies AC in ZF\\
346 & 1 & (7) & This project\\
348 & 1 & (7) & Howard\index{Author}{Howard} [\textbf{1985}]\\
359 & 1 & (7) & This project\\
0 & 2 & (3) & Mathias\index{Author}{Mathias} [\textbf{1979}] thm 1353, Jech [\textbf{1966a}] (M1)\\
2 & 2 & (1) & Tarski\index{Author}{Tarski} [\textbf{1954a}]\\
3 & 2 & (7) & G. Moore\index{Author}{G. Moore} [\textbf{1982}] p 322 \\
2 & 3 & (1)& Truss\index{Author}{Truss} [\textbf{1973c}]\\
3 & 3 & (1) & Tarski [\textbf{1954a}]\\
4 & 3 & (4) & Brunner/Howard\index{Author}{Brunner/Howard} [\textbf{1992}], note 27 (M29,N15)\\
9 & 3 & (3) & G. Moore [\textbf{1982}] p 322 (N12$(\aleph_{1}$)T)\\
87 & 3 & (3) & Monro [\textbf{1974}], note 104 (N12$(\aleph_{\alpha})$T)\\
152 & 3 & (4) & Brunner/Howard [\textbf{1992}], note 27 (M29,N15)\\
187 & 3 & (5) & A. Rubin\index{Author}{A. Rubin} [\textbf{1978}], Hickman\index{Author}{Hickman} [\textbf{1979b}] (N12$(\aleph_{1})$)\\
200 & 3 & (3) & Pincus\index{Author}{Pincus} [\textbf{1974a}], [\textbf{1978}] (M29)\\
202 & 3 & (5) & Truss [\textbf{1978}] (N12$(\aleph_{1})$)\\
214 & 3 & (3) & Pincus [\textbf{1974a}] (M29)\\
290 & 3 & (3) & Pincus [\textbf{1978}], note 64 (M29)\\
291 & 3 & (3) & Pincus [\textbf{1978}], note 64 (M29)\\
317 & 3 & (5) & (N38)\\
345 & 3 & (5) & (N40)\\
3 & 4 & (1) & Halpern/Howard\index{Author}{Halpern/Howard} [\textbf{1974}], Brunner/Howard [\textbf{1992}], note 27\\
4 & 4 & (1) & Whitehead\index{Author}{Whitehead} [\textbf{1902}]\\
129 & 4 & (1) & clear\\
152 & 4 & (1) & Brunner/Howard [\textbf{1992}], note 27\\
3 & 5 & (3) & Sageev\index{Author}{Sageev} [\textbf{1975}] (M6)\\
5 & 5 & (1) & Sageev [\textbf{1975}]\\
6 & 5 & (1) & Sierpi\'{n}ski\index{Author}{Sierpi\'{n}ski} [\textbf{1918}]\\
30 & 5 & (3) & Sageev [\textbf{1975}] (M6)\\
32 & 5 & (1) & clear\\
60 & 5 & (1) & clear\\
94 & 5 & (1) & clear\\
126 & 5 & (3) & Sageev [\textbf{1975}] (M6)\\
\end{tabular}
\end{table}

\clearpage

\begin{table}[h]
\begin{tabular}{llll}
295 & 5 & (3) & (M6)\\
6 & 6 & (1) & G. Moore\index{Author}{G. Moore} [\textbf{1982}]\\
16 & 6 & (1) & Sierpi\'{n}ski\index{Author}{Sierpi\'{n}ski} [\textbf{1918}]\\
30 & 6 & (3) & G. Moore [\textbf{1982}] p 324, Sageev\index{Author}{Sageev} [\textbf{1975}] (M6)\\
31 & 6 & (1) & clear\\
7 & 7 & (1) & Schoenflies\index{Author}{Schoenflies} [\textbf{1908}]\\
9 & 7 & (7) & G. Moore [\textbf{1982}] p 322 \\
101 & 7 & (7) & Banaschewski/Moore\index{Author}{Banaschewski/Moore} [\textbf{1990}]\\
185 & 7 & (3) & (N6T)\\
8 & 8 & (1) & Sierpi\'{n}ski [\textbf{1918}]\\
9 & 8 & (4) & Levy\index{Author}{Levy} [\textbf{1964}], Pincus\index{Author}{Pincus} [\textbf{1969}], Jech\index{Author}{Jech} \cite{bib:1973b} p 124 (N16T)\\
27 & 8 & (6) & G. Moore [\textbf{1982}] (N39)\\
29 & 8 & (7) & G. Moore [\textbf{1982}]\\
39 & 8 & (1) & clear\\
41 & 8 & (3) & Levy [\textbf{1964}], Pincus [\textbf{1969}], Jech \cite{bib:1973b} p 124 thm 8.6 (M6)\\
43 & 8 & (1) & clear\\
71 & 8 & (3) & Jech \cite{bib:1973b} p 124 , Levy [\textbf{1964}], Pincus [\textbf{1969}] (N16$(\aleph_{\gamma})$T)\\
86 & 8 & (1) &clear\\
113 & 8 & (1) & Brunner\index{Author}{Brunner} [\textbf{1985a}], Kelley\index{Author}{Kelley} [\textbf{1950}]\\
181 & 8 & (1) & clear\\
0 &  9 &(4) &Jech \cite{bib:1973b} p 81  prob 5.18 (M1)\\
3 & 9 & (1) & Halpern/Howard\index{Author}{Halpern}\index{Author}{Halpern/Howard} [\textbf{1974}]\\
4 & 9 & (1) & clear\\
7 & 9 & (1) & Jech \cite{bib:1973b} p 161  prob 11.6, Tarski\index{Author}{Tarski} [\textbf{1965}]\\
8 & 9 & (1) & Dedekind\index{Author}{Dedekind} [\textbf{1888}]\\
9 & 9 & (1) & Dedekind [\textbf{1888}]\\
10 & 9 & (3) & Halpern [\textbf{1964}], Halpern/Levy\index{Author}{Halpern/Levy} [\textbf{1971}], Cohen\index{Author}{Cohen} [\textbf{1963}],[\textbf{1966}]\\
&&&(M1)\\
15 & 9 & (3) & Jech \cite{bib:1973b} probs 5.18, 5.21 (M1)\\
17 & 9 & (3) & Blass\index{Author}{Blass} [\textbf{1977a}], Pincus [\textbf{1972a}] (N1T)\\
18 & 9 & (4) & Brunner [\textbf{1984f}], Blass [\textbf{1977a}] (M1)\\
31 & 9 & (3) & note 88 (N1T)\\
41 & 9 & (1) & clear\\
57 & 9 & (3) & Mathias\index{Author}{Mathias} [\textbf{1979}] prob 1347, Jech \cite{bib:1973b} p 162  prob 11.7,\\
&&&Tarski [\textbf{1965}], Sageev \cite{bib:1981b}, note 44, (M32 IC)\\
62 & 9 & (4) & Sageev \cite{bib:1981b}, note 44 (M1)\\
67 & 9 & (6) & Brunner [\textbf{1983c}] (N2)\\
71 & 9 & (1) & clear\\
100 & 9 & (1) & Jech \cite{bib:1973b} p 162  prob 11.8, Tarski [\textbf{1965}]\\
118 & 9 & (3) & Haddad/Morillon\index{Author}{Haddad/Morillon} [\textbf{1990}] (M1)\\
128 & 9 & (3) & Brunner [\textbf{1984f}] (M1)\\
129 & 9 & (1) & clear\\
161 & 9 & (1) & H\"{a}ussler\index{Author}{H\"{a}ussler} [\textbf{1983}]\\
174 & 9 & (1) & von Rimscha\index{Author}{Rimscha} [\textbf{1982}], note 49\\
185 & 9 & (3) & (N6T)\\
214 & 9 & (1) & clear\\
217 & 9 & (3) & Pincus [\textbf{1972a}], note 105 (N1T)\\
0 &  10 & (3) & Pincus \cite{bib:1972c1}, note 35 (M4)\\
7 & 10 & (2) & Tarski [\textbf{1965}]\\
\end{tabular}
\end{table}

\clearpage

\begin{table}[h]
\begin{tabular}{llll}
9 & 10 & (1) & G. Moore\index{Author}{G. Moore} [\textbf{1982}] p 322 \\
30 & 10\index{Author}{Mathias} & (1) & clear\\
32 & 10 & (1) & clear\\
52 & 10 & (3) & (N6T)\\
60 & 10 & (1) & clear\\
61 & 10 & (3) & Jech\index{Author}{Jech} \cite{bib:1973b} thm 7.11, note 35 (N6T)\\
62 & 10 & (1) & clear\\
64 & 10 & (4) & Brunner\index{Author}{Brunner} [\textbf{1984b}](N6T)\\
119 & 10 & (3) & (N6T)\\
122 & 10 & (1) & clear\\
132 & 10 & (1) & Brunner [\textbf{1984b}]\\
133 & 10 & (1) & Brunner [\textbf{1984b}]\\
153 & 10 & (1) & Fossy/Morillon\index{Author}{Fossy/Morillon} [\textbf{1995}]\\
171 & 10 & (7) & Shannon\index{Author}{Shannon} \cite{bib:1990f}\\
185 & 10 & (3) & (N6T)\\
218 & 10 & (5) & Bleicher\index{Author}{Bleicher} [\textbf{1965}] (N6)\\
219 & 10 & (3) & Pincus\index{Author}{Pincus} [\textbf{1972a}] p 737 (N6T)\\
308 & 10 & (3) & (N22$(K,p)$T)\\
322 & 10 & (5) & Brunner [\textbf{1990}] (N50$(E)$)\\
357 & 10 & (3) & (N50$(E)$T)\\
358 & 10 & (3) & (N50$(E)$T)\\
361 & 10 & (3) & note 18 (N6T)\\
363 & 10 & (3) & note 18 (N6T)\\
11 & 11 & (1) & Lindel\"{o}f\index{Author}{Lindel\"{o}f} [\textbf{1905}]\\
12 & 11 & (7) & G. Moore [\textbf{1982}] p 322 \\
11 & 12 & (1) & clear\\
12 & 12 & (1) & Young\index{Author}{Young} [\textbf{1903}]\\
13 & 12 & (7) & G. Moore [\textbf{1982}] p 322 \\
9 & 13 & (1) & clear\\
12 & 13 & (1) & Jech [\textbf{1968b}], G. Moore [\textbf{1982}]\\
13 & 13 & (1) & G. Moore [\textbf{1982}]\\
14 & 13 & (3) & Jech \cite{bib:1973b} p 97 p 81 prob 5.18, Halpern/Levy\index{Author}{Halpern/Levy} [\textbf{1971}], (M1)\\
15 & 13 & (3) & Jech \cite{bib:1973b} p 81 prob 5.18, p 82 prob 5.21, (M1)\\
16 & 13 & (1) & G. Moore [\textbf{1982}]\\
18 & 13 & (4) & Brunner [\textbf{1984f}] (M1)\\
31 & 13 & (3) & Jech \cite{bib:1973b} (M1)\\
60 & 13 & (3) & Jech \cite{bib:1973b} p 82 prob 5.22, p 81 prob 5.18 (M1)\\
73 & 13 & (1) & Jech \cite{bib:1973b} p 141 \\
82 & 13 & (3) & Jech \cite{bib:1973b} p 81  probs 5.18, 5.20 (M1)\\
92 & 13 & (1) & Jech \cite{bib:1973b} p 141 \\
94 & 13 & (1) & Jech \cite{bib:1973b} p 148 prob 10.1\\
118 & 13 & (3) & Haddad/Morillon\index{Author}{Haddad/Morillon} [\textbf{1990}] (M1)\\
128 & 13 & (3) & Brunner\index{Author}{Brunner} [\textbf{1984f}] (M1)\\
163 & 13 & (3) & Brunner [\textbf{1982a}] (M1)\\
165 & 13 & (3) & (M1)\\
175 & 13 & (1) & von Rimscha\index{Author}{Rimscha} [\textbf{1982}], note 49\\
185 & 13 & (1) & clear\\
191 & 13 & (3) & Blass\index{Author}{Blass} [\textbf{1979}], note 59 (M1)\\
211 & 13 & (1) & clear\\
\end{tabular}
\end{table}

\clearpage

\begin{table}[h]
\begin{tabular}{llll}
277 & 13 & (3) & (M1)\\
15 & 14 & (3) & Jech\index{Author}{Jech} \cite{bib:1973b} p 113 prob 7.8 (M3)\\
30 & 14 & (4) & Mathias\index{Author}{Mathias} [\textbf{1967}] (M3)\\
49 & 14 & (5) & Felgner/Truss\index{Author}{Felgner}\index{Author}{Felgner/Truss} [\textbf{1996}] (N52)\\
50 & 14 & (1) & Mathias [\textbf{1979}] p 124\\
52 & 14 & (3) & Pincus\index{Author}{Pincus} [\textbf{1972b}] (M4)\\
63 & 14 & (3) & Jech \cite{bib:1973b} p 132  prob 8.5 (N2$(\aleph_{1})$T)\\
65 & 14 & (6) & note 99 (N2)\\
99 & 14 & (6) & note 78 (N2)\\
101 & 14 & (7) & Banaschewski/Moore\index{Author}{Banaschewski/Moore} [\textbf{1990}]\\
141 & 14 & (5) & Howard\index{Author}{Howard} [\textbf{1984a}] $(\mathrm{N}2^{\ast}(3))$\\
153 & 14 & (7) & Fossy/Morillon\index{Author}{Fossy/Morillon}\index{Author}{Morillon} [\textbf{1995}]\\
191 & 14 & (4) & note 59 (M7)\\
317 & 14 & (1) & Howard \cite{bib:1973h}\\
345 & 14 & (1) & Morillon \cite{bib:1988h}\\
3 & 15 & (5) & Mostowski\index{Author}{Mostowski} [\textbf{1958}], note 18 (N9)\\
4 & 15 & (4) & Mostowski [\textbf{1958}], note 18 (M47$(n, M)$,N9)\\
9 & 15 & (4) & Mostowski [\textbf{1958}], G. Moore\index{Author}{G. Moore} [\textbf{1982}], note 18 (M47$(n, M)$)\\
14 & 15 & (3) & Pincus [\textbf{1969}],[\textbf{1974b}], Jech \cite{bib:1973b} p 113  prob 7.6, note 18\\
&&&(N3T)\\
15 & 15 & (1) & Kinna/Wagner\index{Author}{Kinna/Wagner} [\textbf{1955}]\\
30 & 15 & (3) & Pincus [\textbf{1972a}], Felgner [\textbf{1971b}], Jech \cite{bib:1973b} p 95  prob 6.11,\\
&&&note 73 note 18 (M44, N3T)\\
43 & 15 & (4) & note 73 note 18 (M47$(n, M)$)\\
134 & 15 & (5) & note 18 (N1)\\
214 & 15 & (4) & note 73, note 18 (M47$(n, M)$)\\
295 & 15 & (3) & note 18 (N48T)\\
345 & 15 & (5) & note 18 (N40)\\
8 & 16 & (1) & clear\\
16 & 16 & (1) & Sierpi\'{n}ski\index{Author}{Sierpi\'{n}ski} [\textbf{1918}]\\
1 & 17 & (1) & Kleinberg\index{Author}{Kleinberg} [\textbf{1969}]\\
9 & 17 & (1) & Kleinberg [\textbf{1969}]\\
10 & 17 & (4) & note 51 (M1)\\
14 & 17 & (3) & Blass\index{Author}{Blass} [\textbf{1977a}] (M1)\\
15 & 17 & (3) & Blass [\textbf{1977a}] (M1)\\
17 & 17 & (1) & Jech \cite{bib:1973b} p 164  prob 11.20, Ramsey\index{Author}{Ramsey} [\textbf{1929}]\\
18 & 17 & (3) & Blass [\textbf{1977a}] (M1)\\
31 & 17 & (3) &Blass [\textbf{1977a}], Jech \cite{bib:1973b} (M1)\\
60 & 17 & (3) & Blass [\textbf{1977a}], Jech \cite{bib:1973b} (M1)\\
82 & 17 & (3) & (M1)\\
118 & 17 & (3) & Blass [\textbf{1977a}], Haddad/Morillon\index{Author}{Haddad/Morillon} [\textbf{1990}] (M1)\\
128 & 17 & (3) & Blass [\textbf{1977a}], Brunner [\textbf{1984f}] (M1)\\
277 & 17 & (3) & (M1)\\
325 & 17 & (1) & clear\\
6 & 18 & (3) & note 18 (N2T)\\
11 & 18 & (3) & note 18 (N2T)\\
17 & 18 & (1) & Jech \cite{bib:1973b} p 164  prob 11.20, Blass [\textbf{1977a}]\\
37 & 18 & (3) & note 18 (N2T)\\
52 & 18 & (3) & (N2T)\\
\end{tabular}
\end{table}

\clearpage

\begin{table}[h]
\begin{tabular}{llll}
65 & 18 & (5) & (N2)\\
67 & 18\index{Author}{Howard} & (5) & (N2)\\
73 & 18 & (3) & note 18 (N2T)\\
35 & 18 & (3) & note 18 (N2T)\\
80 & 18 & (1) & clear\\
82 & 18 & (3) & (N2T)\\
83 & 18 & (3) & (N2T)\\
92 & 18 & (3) & note 18 (N2T)\\
99 & 18 & (5) & (N2)\\
119 & 18 & (3) & (N2T)\\
130 & 18 & (3) & Brunner\index{Author}{Brunner} [\textbf{1984f}] (N2T)\\
163 & 18 & (3) & (N2T)\\
166 & 18 & (1) & clear\\
167 & 18 & (1) & clear\\
185 & 18 & (3) & (N2T)\\
191 & 18 & (3) & Blass\index{Author}{Blass} [\textbf{1979}], note 59 (M7)\\
231 & 18 & (1) & clear\\
273 & 18 & (3) & note 18 (N2T)\\
305 & 18 & (3) & note 18 (N2T)\\
313 & 18 & (3) & note 18 (N2T)\\
334 & 18 & (5) & (N2)\\
350 & 18 & (3) & (N2T)\\
361 & 18 & (3) & note 18 (N2T)\\
363 & 18 & (3) & note 18 (N2T)\\
0 &  19 &(7) &This project\\
19 & 19 & (1) & Lebesgue\index{Author}{Lebesgue} [\textbf{1905}]\\
34 & 19 & (1) & Lebesgue [\textbf{1905}]\\
21 & 20 & (7) & G. Moore\index{Author}{G. Moore} [\textbf{1982}]\\
97 & 20 & (3) & note 120(43) (M1)\\
359 & 20 & (1) & clear\\
20 & 21 & (1) & clear\\
21 & 21 & (1) & G. Moore [\textbf{1982}]\\
23 & 21 & (6) & G. Moore [\textbf{1982}] (N39)\\
22 & 22 & (1) & Borel\index{Author}{Borel} [\textbf{1898}], G. Moore [\textbf{1982}]\\
21 & 23 & (1) & G. Moore [\textbf{1982}]\\
23 & 23 & (1) & G. Moore [\textbf{1982}]\\
25 & 23 & (6) & G. Moore [\textbf{1982}] (N6)\\
27 & 23 & (4) & G. Moore [\textbf{1982}] (N15T)\\
29 & 23 & (4) & G. Moore [\textbf{1982}] (N15T)\\
8 & 24 & (1) & clear\\
24 & 24 & (1) & G. Moore [\textbf{1982}]\\
0 &  25 & (4) & Apter\index{Author}{Apter} [\textbf{1985a}] (M8 IC, M36)\\
23 & 25 & (1) & Hausdorff\index{Author}{Hausdorff} [\textbf{1907}]\\
25 & 25 & (1) & Feferman/Levy\index{Author}{Feferman/Levy} [\textbf{1963}], Gitik\index{Author}{Gitik} [\textbf{1980}]\\
51 & 25 & (1) & Morris\index{Author}{Morris} [\textbf{1969}]\\
22 & 26 & (1) & clear\\
24 & 26 & (1) & Borel [\textbf{1898}], G. Moore [\textbf{1982}] p 66\\
26 & 26 & (1) & Borel [\textbf{1898}]\\
8 & 27 & (1) & clear\\
\end{tabular}
\end{table}

\clearpage

\begin{table}[h]
\begin{tabular}{llll}
23 & 27 & (1) & clear\\
27 & 27\index{Author}{Jech}\index{Author}{Keremedis} & (1) & G. Moore\index{Author}{G. Moore} [\textbf{1982}] p 36 \\
29 & 27 & (1) & Howard\index{Author}{Howard} [\textbf{1994}]\\
31 & 27 & (3) & (N17T)\\
8 & 28 & (1) & clear\\
16 & 28 & (1) & clear\\
28 & 28 & (1) & Sierpi\'{n}ski\index{Author}{Sierpi\'{n}ski} [\textbf{1918}]\\
31 & 28 & (3) & Brunner/Howard\index{Author}{Brunner}\index{Author}{Brunner/Howard} [\textbf{1992}], note 27 (N17T)\\
37 & 28 & (4) & G. Moore [\textbf{1982}] p 324 (N6T)\\
8 & 29 & (1) & G. Moore [\textbf{1982}] p 324 \\
16 & 29 & (5) & (N39)\\
20 & 29 & (1) & clear\\
21 & 29 & (1) & clear\\
23 & 29 & (5) & Howard [\textbf{1994}] (N39)\\
24 & 29 & (5) & (N39)\\
29 & 29 & (1) & G. Moore [\textbf{1982}] p 324 \\
31 & 29 & (3) & (N6T)\\
15 & 30 & (1) & Jech \cite{bib:1973b} p 53  prob 4.12\\
23 & 30 & (5) & (N7)\\
30 & 30 & (1) & Schr\"{o}der\index{Author}{Schr\"{o}der} [\textbf{1898}]\\
49 & 30 & (1) & clear\\
60 & 30 & (3) & L\"{a}uchli\index{Author}{L\"{a}uchli} [\textbf{1964}], Pincus\index{Author}{Pincus} [\textbf{1968}], Jech \cite{bib:1973b} p 104 (N7T)\\
62 & 30 & (4) & L\"{a}uchli [\textbf{1964}], Pincus [\textbf{1968}], Jech \cite{bib:1973b} p 104, note 73\\
&&&(N7T)\\
101 & 30 & (7) & Banaschewski/Moore\index{Author}{Banaschewski/Moore} [\textbf{1990}]\\
123 & 30 & (3) & Pincus [\textbf{1976}] (N7T)\\
214 & 30 & (4) & note 73 (M47$(n, M)$)\\
295 & 30 & (1) & Gonzalez\index{Author}{Gonzalez} [\textbf{1995a}]\\
6 & 31 & (3) & Truss\index{Author}{Truss} [\textbf{1974b}] (M12$(\aleph)$)\\
27 & 31 & (1) & clear\\
28 & 31 & (1) & Sierpi\'{n}ski [\textbf{1918}], Brunner/Howard [\textbf{1992}], note 27\\
29 & 31 & (1) & G. Moore [\textbf{1982}] p 31 \\
31 & 31 & (1) & Cantor\index{Author}{Cantor} [\textbf{1878}]\\
32 & 31 & (3) & Howard [\textbf{1992}] (N18T)\\
60 & 31 & (3) & Felgner\index{Author}{Felgner} [\textbf{1974}] (M20)\\
165 & 31 & (3) & (M20)\\
9 &32 & (4) & Jech \cite{bib:1973b} p 132  prob 8.4, Pincus [\textbf{1969}], [\textbf{1974a}], [\textbf{1978}]\\
&&&(M29)\\
30 & 32 & (4) & Jech \cite{bib:1973b} p 115  prob 7.12 (M6)\\
31 & 32 & (1) & Sierpi\'{n}ski [\textbf{1918}]\\
32 & 32 & (1) & Sierpi\'{n}ski [\textbf{1918}]\\
60 & 32 & (1) & clear\\
85 & 32 & (1) & clear\\
150 & 32 & (1) & clear (see Brunner [\textbf{1983b}] p91)\\
165 & 32 & (1) & clear\\
200 & 32 & (3) & Pincus [\textbf{1974a}] [\textbf{1978}] (M29)\\
214 & 32 & (3) & (M29)\\
290 & 32 & (3) & Pincus [\textbf{1978}], note 64 (M29)\\
291 & 32 & (3) & Pincus [\textbf{1978}], note 64 (M29)\\
\end{tabular}
\end{table}

\clearpage

\begin{table}[h]
\begin{tabular}{llll}
352 & 32 & (1) & Keremedis [\textbf{1996d}]\\
354 & 32\index{Author}{Banaschewski/Moore}\index{Author}{Higasikawa}\index{Author}{Morillon}\index{Author}{Pelc} & (1) & Keremedis [\textbf{1996d}]\\
33 & 33 & (1) & Truss\index{Author}{Truss} [\textbf{1973a}]\\
45 & 33 & (1) & clear\\
52 & 33 & (3) & (N2$(n)$T)\\
120 & 33 & (0) & note 15 (see M47$(n, M)$)\\
121 & 33 & (1) & clear\\
122 & 33 & (3) & Truss [\textbf{1973a}] (N24$(n, LO)$T)\\
313 & 33 & (3) & note 18 (N24$(n, LO)$T)\\
6 & 34 & (3) & (M12$(\aleph)$)\\
25 & 34 & (1) & clear\\
31 & 34 & (1) & clear\\
38 & 34 & (4) & G. Moore\index{Author}{G. Moore} [\textbf{1982}] p 324  (M36)\\
170 & 34 & (3) & Figura\index{Author}{Figura} \cite{bib:1981b}, note 3 (M36)\\
172 & 34 & (1) & Jech\index{Author}{Jech} [\textbf{1982}]\\
231 & 34 & (3) & (M12$(\aleph)$\\
245 & 34 & (1) & Litman\index{Author}{Litman} [\textbf{1976}]\\
6 & 35 & (1) & clear\\
31 & 35 & (1) & clear\\
35 & 35 & (1) & Feferman/Levy\index{Author}{Feferman/Levy} [\textbf{1963}]\\
6 & 36 & (1) & clear\\
36 & 36 & (1) & G. Moore [\textbf{1982}] p 36 \\
38 & 36 & (7) & G. Moore [\textbf{1982}] p 324 \\
28 & 37 & (1) & Sierpi\'{n}ski\index{Author}{Sierpi\'{n}ski} [\textbf{1916}]\\
37 & 37 & (1) & Sierpi\'{n}ski [\textbf{1916}]\\
38 & 37 & (7) & G. Moore [\textbf{1982}] p 324 \\
0 &  38 & (3) & Jech \cite{bib:1973b} p 142  thm 10.6 (M9)\\
6 & 38 &(1) & clear\\
34 & 38 & (1) & Jech \cite{bib:1973b} p 148 prob 10.2\\
36 & 38 & (1) & G. Moore [\textbf{1982}] p 324, note 110\\
37 & 38 & (1) & Sierpi\'{n}ski [\textbf{1918}]\\
79 & 38 & (1) & clear\\
170 & 38 & (1) & clear\\
39 & 39 & (1) & Sierpi\'{n}ski [\textbf{1918}], G. Moore [\textbf{1982}] p 202 \\
40 & 39 & (1) & clear\\
43 & 39 & (3) & Jech [\textbf{1966b}], Mostowski\index{Author}{Mostowski} [\textbf{1948}], Pincus\index{Author}{Pincus} [\textbf{1977a}] (M1$(<\omega_{1}>)$,\\
&&&(M5$(\aleph)$ IC)\\
44 & 39 & (1) & Jech \cite{bib:1973b} p 120  thm 8.1\\
97 & 39 & (3) & (M1$(<(\omega_{1}>)$)\\
135 & 39 & (5) & Brunner\index{Author}{Brunner} [\textbf{1984b}] (N1)\\
40 & 40 & (1) & G. Moore [\textbf{1982}] p 325 \\
44 & 40 & (3) & Jech \cite{bib:1973b} p123 thm 8.3, Mostowski [\textbf{1948}], Jech [\textbf{1966b}]\\
&&&(N2$(\aleph_{2})$T)\\
51 & 40 & (6) & Harper/Rubin\index{Author}{Harper/Rubin} [\textbf{1976}] (N1)\\
87 & 40 & (3) & Jech \cite{bib:1973b} p123 thm 8.3, Mostowski [\textbf{1948}], Jech [\textbf{1966b}]\\
&&&(N2$(\aleph_{\alpha})$T)\\
101 & 40 & (1) & Pelc\index{Author}{Pelc} [\textbf{1978}]\\
202 & 40 & (1) & clear\\
260 & 40 & (1) & Harper/Rubin [\textbf{1976}]\\
\end{tabular}
\end{table}

\clearpage

\begin{table}[h]
\begin{tabular}{llll}
286 & 40 & (1) & Morillon [\textbf{1991b}]\\
317 & 40\index{Author}{Gauntt}\index{Author}{Kolany/Wojtylak}\index{Author}{Zuckerman} & (5) & (N19)\\
345 & 40 & (5) & N40\\
347 & 40 & (1) & Higasikawa [\textbf{1995}]\\
7 &41 & (7) & Banaschewski/Moore [\textbf{1990}]\\
40 & 41 & (3) & Jech\index{Author}{Jech} \cite{bib:1973b} p 127 , Mostowski\index{Author}{Mostowski} [\textbf{1948}], Pincus\index{Author}{Pincus} [\textbf{1969}]\\
&&&(N12$(\aleph_{1})$T)\\
41 & 41 & (1) & G. Moore\index{Author}{G. Moore} [\textbf{1982}] p 326 \\
44 & 41 & (1) & Jech \cite{bib:1973b} p 120  thm 8.1\\
135 & 41 & (5) & Brunner\index{Author}{Brunner} [\textbf{1984b}] (N1)\\
187 & 41 & (5) & (N12$(\aleph_{1})$)\\
202 & 41 & (5) & (N12$(\aleph_{1})$)\\
0 &  42 & (7) & This project\\
42 & 42 & (1) & G. Moore [\textbf{1982}] p 251 \\
43 & 42 & (1) & clear (see G. Moore [\textbf{1982}] p 325)\\
8 &43 & (4) & Jech \cite{bib:1973b} p 130  thm 8.12, Jensen\index{Author}{Jensen} [\textbf{1967}], Brunner [\textbf{1984b}]\\
&&&(N21$(\aleph_{2})$T)\\
14 & 43 & (4) & Brunner [\textbf{1984b}] (M1)\\
39 & 43 & (3) & Jensen [\textbf{1966}], [\textbf{1967}], Jech \cite{bib:1973b} p 130, Felgner\index{Author}{Felgner} [\textbf{1971a}]\\
&&&pp 151ff, note 13 (N21$(\aleph_{2})$T)\\
40 & 43 & (1) & Jech [\textbf{1966a}] p 121  thm 8.2, Jensen [\textbf{1967}]\\
43 & 43 & (1) & Tarski\index{Author}{Tarski} [\textbf{1948}] p 96, Levy\index{Author}{Levy} [\textbf{1964}] p 136\\
44 & 43 & (1) & Levy [\textbf{1964}], Jech \cite{bib:1973b} p 120 thm 8.1\\
67 & 43 & (6) & Pincus [\textbf{1972b}] (N2)\\
71 & 43 & (4) & Jech \cite{bib:1973b} p 130, Jensen [\textbf{1967}] ((N21$(\aleph_{\alpha+1})$T)\\
86 & 43 & (3) & Jech \cite{bib:1973b} p 130, Jensen [\textbf{1967}], (N21$(\aleph_{\alpha+1})$T)\\
87 & 43 & (1) & clear\\
101 & 43 & (1) & Pelc\index{Author}{Pelc} [\textbf{1978}]\\
106 & 43 & (6) & Blass\index{Author}{Blass} [\textbf{1979}], Goldblatt\index{Author}{Goldblatt} \cite{bib:1985m} (N6)\\
113 & 43 & (5) & Brunner [\textbf{1984b}], Howard/Rubin\index{Author}{Howard/Rubin} [\textbf{1996}] (N38)\\
174 & 43 & (1) & H\"{o}ft/Howard\index{Author}{H\"{o}ft/Howard} \cite{bib:1981b}\\
192 & 43 & (1) & Blass [\textbf{1979}]\\
317 & 43 & (5) &(N38)\\
345 & 43 & (1) & Morillon\index{Author}{Morillon} \cite{bib:1988h}\\
40 & 44 & (3) & Jech \cite{bib:1973b} p 127  thm 8.9, Pincus [\textbf{1969}], Levy [\textbf{1964}],\\
&&&note 14 (N12$(\aleph_{1})$T)\\
44 & 44 & (1) & Levy [\textbf{1964}]\\
0 & 45 & (3) & Cohen\index{Author}{Cohen} [\textbf{1963}],[\textbf{1965}],[\textbf{1966}], Jech \cite{bib:1973b} p 71 (M7$(n)$)\\
23 & 45 & (5) & (N12$(\aleph_{1})$)\\
33 & 45 & (4) & note 15 (N24$(n)$T)\\
40 & 45 & (3) & Jech \cite{bib:1973b} p 131  prob 8.1 (N12$(\aleph_{1})$T)\\
43 & 45 & (4) & Pincus [\textbf{1977a}], note 73 (M47$(n, M)$)\\
44 & 45 & (3) & (N2$(\aleph_{\alpha})$T)\\
45 & 45 & (1) & Mostowski [\textbf{1945}], Truss\index{Author}{Truss} [\textbf{1973a}]\\
46 & 45 & (0) & notes 15,73 (see M47$(n, M)$)\\
60 & 45 & (1) & clear\\
61 & 45 & (1) & clear\\
63 & 45 & (3) & (N12$(\aleph_{1})$T)\\
65 & 45 & (5) & (N2)\\
\end{tabular}
\end{table}

\clearpage

\begin{table}[h]
\begin{tabular}{llll}
67 & 45 & (5) & Zuckerman\index{Author}{Zuckerman} [\textbf{1971}], note 15 (N2$(n)$)\\
87 & 45 & (3) & Jech\index{Author}{Jech} \cite{bib:1973b} p 131  prob 8.1 (N2$(\aleph_{\alpha})$T)\\
99 & 45 & (5) & (N2)\\
121 & 45 & (3) & Gauntt [\textbf{1968}], Truss\index{Author}{Truss} [\textbf{1973a}] (N24$(n)$T)\\
133 & 45 & (5) & (N24$(n)$)\\
135 & 45 & (5) & (N2$(\aleph_{\alpha})$\\
163 & 45 & (3) & (N2T)\\
187 & 45 & (5) & (N12$(\aleph_{1})$)\\
202 & 45 & (5) & (N12$(\aleph_{1})$)\\
214 & 45 & (4) & note 73 (M47$(n, M)$)\\
271 & 45 & (0) & Kolany/Wojtylak [\textbf{1991}], (see note 92)\\
288 & 45 & (4) & note 15 (N24$(n)$T)\\
313 & 45 & (3) & note 18 (N2T)\\
334 & 45 & (5) & (N2)\\
6 & 46 &(3) & note 18 (N2$(n)$T)\\
11 & 46 & (3) & note 18 (N2$(n)$T)\\
35 & 46 & (3) & note 18 (N2$(n)$T)\\
37 & 46 & (3) & note 18 (N2$(n)$T)\\
43 & 46 & (3) & (M47$(n, M)$)\\
45 & 46 & (0) & (see note 15)\\
46 & 46 & (1) & Mostowski\index{Author}{Mostowski} [\textbf{1945}], Truss [\textbf{1973a}]\\
52 & 46 & (3) & (N2$(n)$T)\\
60 & 46 & (1) & clear\\
61 & 46 & (1) & clear\\
63 & 46 & (3) & (M47$(n, m)$)\\
67 & 46 & (5) & Zuckerman [\textbf{1971}], note 15 (N2$(n)$)\\
73 & 46 & (3) & note 18 (N2$(n)$T)\\
92 & 46 & (3) & (note 18N2$(n)$T)\\
130 & 46 & (3) & note 18 (N2$(n)$T)\\
191 & 46 & (5) & note 18 ($A$ is a set) (N2$(n)$)\\
214 & 46 & (3) & (M47$(n, M)$)\\
273 & 46 & (3) & note 18 (N2$(n)$T)\\
305 & 46 & (3) & note 18 (N2$(n)$T)\\
313 & 46 & (3) & note 18 (N2$(n)$T)\\
0 & 47 & (3) & Cohen\index{Author}{Cohen} [\textbf{1963}],[\textbf{1965}],[\textbf{1966}], Jech \cite{bib:1973b} p 71 (M7$(n)$)\\
11 & 47 & (3) & note 18 (N2$(n)$T)\\
33 & 47 & (1) & clear\\
43 & 47 & (3) & Pincus\index{Author}{Pincus} [\textbf{1977a}], note 73 (M47$(n, M)$)\\
45 & 47 & (1) & clear\\
46 & 47 & (0) & note 15, note 73 (see M47$(n, M)$)\\
47 & 47 & (1) & Mostowski [\textbf{1945}], Truss [\textbf{1973a}]\\
48 & 47 & (0) & note 15 (see M47$(n, M)$)\\
52 & 47 & (3) & (N2$(n)$T)\\
60 & 47 & (2) & clear\\
63 & 47 & (3) & (M47$(n, M)$)\\
67 & 47 & (5) & note 73 (N2$(n)$)\\
88 & 47 & (3) & ($\mathrm{N}2^{\ast}$(3)T)\\
91 & 47 & (5) & note 18 (N2$(n)$)\\
122 & 47 & (1) & clear\\
\end{tabular}
\end{table}

\clearpage

\begin{table}[h]
\begin{tabular}{llll}
130 & 47 & (3) & note 18 (N2$(n)$T)\\
191 & 47 & (5) & note 18 ($A$ is a set) (N2$(n)$)\\
214 & 47 & (3) & note 73 (M47$(n, M)$)\\
250 & 47 & (1) & clear\\
273 & 47 & (3) & note 18 (N2$(n)$T)\\
305 & 47 & (3) & note 18 (N2$(n)$T)\\
0 &  48 & (7) & This project\\
45 & 48 & (0) & (see note 15)\\
46 & 48 & (1) & clear\\
47 & 48 & (0) & (see note 15)\\
48 & 48 & (1) & Mostowski [\textbf{1945}], Truss [\textbf{1973a}]\\
120 & 48 & (1) & clear\\
122 & 48 & (1) & clear\\
250 & 48 & (1) & clear\\
14 & 49 & (1) & Mathias\index{Author}{Mathias} [\textbf{1979}] p 124 \\
15 & 49 & (3) & Jech\index{Author}{Jech} \cite{bib:1973b} p 113  prob 7.8 (M3)\\
30 & 49 & (3) & Mathias [\textbf{1967}], Jech \cite{bib:1973b} prob 7.8, Felgner\index{Author}{Felgner} [\textbf{1971a}],\\
&&&note 73 (M3)\\
43 & 49 & (4) & note 73 (M47$(n, M)$)\\
49 & 49 & (1) & Tarski\index{Author}{Tarski} [\textbf{1924}] p 78 \\
51 & 49 & (3) & Morris\index{Author}{Morris} [\textbf{1969}] (M14)\\
214 & 49 & (4) & note 73 (M47(n,M))\\
14 & 50 & (3) & Bell\index{Author}{Bell} \cite{bib:1983a} (M1)\\
118 & 50 & (3) & (M1)\\
277 & 50 & (3) & (M1)\\
303 & 50 & (1) & Bell \cite{bib:1988b}\\
51 & 51 & (1) & Sierpi\'{n}ski\index{Author}{Sierpi\'{n}ski} [\textbf{1918}] p 117 \\
87 & 51 & (3) & (M40$(\kappa)$)\\
90 & 51 & (1) & Harper/Rubin\index{Author}{Harper/Rubin} [\textbf{1976}]\\
259 & 51 & (1) & Harper/Rubin [\textbf{1976}]\\
316 & 51 & (5) & (N37)\\
317 & 51 & (5) & (N19)\\
345 & 51 & (3) & M40$(\kappa)$\\
0 & 52 & (3) & Solovay\index{Author}{Solovay} [\textbf{1970}], Pincus\index{Author}{Pincus} \cite{bib:1972c1} (M5$(\aleph)$ IC)\\
14 & 52 & (1) & Lo\'{s}/Ryll-Nardzewski\index{Author}{Lo\'{s}/Ryll-Nardzewski} [\textbf{1951}], Luxemburg\index{Author}{Luxemburg} \cite{bib:1969d},[\textbf{1970}]\\
52 & 52 & (1) & Lo\'{s}/Ryll-Nardzewski [\textbf{1951}]\\
67 & 52 & (1) & Pincus [\textbf{1972b}]\\
142 & 52 & (3) & Pincus \cite{bib:1972c1} (M2$(\kappa))$\\
153 & 52 & (7) & Fossy/Morillon\index{Author}{Fossy/Morillon} [\textbf{1995}]\\
222 & 52 & (3) & Pincus \cite{bib:1972c1} (M2$(\kappa)$)\\
372 & 52 & (1) & clear\\
0 &  53 & (4) & Mathias [\textbf{1979}] prob 1336 (N1T)\\
3 & 53 & (1) & note 68\\
11 & 53 & (3) & note 18 (N1T)\\
16 & 53 & (5) & (N1)\\
17 & 53 & (3) & (N1T)\\
23 & 53 & (5) & (N1)\\
24 & 53 & (5) & (N1)\\
31 & 53 & (3) & (N1T)\\
\end{tabular}
\end{table}

\clearpage

\begin{table}[h]
\begin{tabular}{llll}
37 & 53 & (3) & note 18 (N1T)\\
47 & 53\index{Author}{Jech}\index{Author}{Jech/Sochor}\index{Author}{Tsukada} & (3) & (N1T)\\
48 & 53 & (3) & (N1T)\\
53 & 53 & (1) & Mathias\index{Author}{Mathias} [\textbf{1979}] prob 1336\\
63 & 53 & (5) & (N1)\\
73 & 53 & (3) & note 18 (N1T)\\
89 & 53 & (5) & (N1)\\
92 & 53 & (3) & note 18 (N1T)\\
112 & 53 & (5) & (N1)\\
114 & 53 & (5) & (N1)\\
115 & 53 & (5) & (N1)\\
116 & 53 & (5) & (N1)\\
122 & 53 & (3) & (N1T)\\
127 & 53 & (5) & (N1)\\
130 & 53 & (3) & note 18 (N1T)\\
133 & 53 & (5) & (N1)\\
134 & 53 & (5) & (N1)\\
135 & 53 & (5) & (N1)\\
147 & 53 & (5) & (N1)\\
163 & 53 & (3) & (N1T)\\
165 & 53 & (3) & (N1T)\\
185 & 53 & (3) & (N1T)\\
191 & 53 & (5) & note 18 (N1)\\
217 & 53 & (3) & (N1T)\\
233 & 53 & (5) & (N1)\\
263 & 53 & (5) & (N1)\\
273 & 53 & (3) & note 18 (N1T)\\
304 & 53 & (5) & (N1)\\
305 & 53 & (3) & note 18 (N1T)\\
309 & 53 & (3) & note 18 (N1T)\\
313 & 53 & (3) & note 18 (N1T)\\
322 & 53 & (5) & (N1)\\
325 & 53 & (3) & (N1T)\\
361 & 53 & (3) & note 18 (N1T)\\
363 & 53 & (3) & note 18 (N1T)\\
380 & 53 & (5) & (N1)\\
0 &  54 &(7)& Mathias [\textbf{1979}] prob 1337\\
54 & 54 & (1) & Mathias [\textbf{1979}] prob 1337, thm 1338, prob 1339\\
0 &  55 & (7) & This project\\
55 & 55 & (1) & Mathias [\textbf{1979}] p 125\\
0 & 56 & (3) & Derrick/Drake\index{Author}{Derrick/Drake} [\textbf{1967b}], Mathias [\textbf{1979}] p 125  (M10)\\
56 & 56 & (1) & Mathias [\textbf{1979}] p 125\\
9 & 57 &(1) &clear\\
57 & 57 & (1) & Mathias [\textbf{1979}] p 125 \\
163 & 57 & (3) & (N24T)\\
165 & 57 & (3) & (N24T)\\
185 & 57 & (3) & (N24T)\\
0 & 58 & (7) & Mathias [\textbf{1979}] prob 1349\\
58 & 58 & (1) & Mathias [\textbf{1979}] p 126 \\
\end{tabular}
\end{table}

\clearpage

\begin{table}[h]
\begin{tabular}{llll}
208 & 58 & (1) & clear\\
0 &59\index{Author}{Monro}\index{Author}{Morillon} &(3)\index{Author}{Felgner} &Mathias [\textbf{1979}] p 1350, Jech\index{Author}{Jech} \cite{bib:1973b} p 151, Jech/Sochor\\
&&&[\textbf{1966a}],[\textbf{1966b}]\,(N11T)\\
6 & 59 & (3) & note 18 (N11T)\\
11 & 59 & (3) & note 18 (N11T)\\
14 & 59 & (3) & Tsukada\index{Author}{Tsukada} [\textbf{1977}] (N19T)\\
16 & 59 & (5) & (N19$(\precsim)$)\\
23 & 59 & (5) & (N19)\\
24 & 59 & (5) & (N19$(\precsim)$)\\
31 & 59 & (3) & (N19T)\\
35 & 59 & (3) & note 18 (N11T)\\
37 & 59 & (3) & note 18 (N11T)\\
59 & 59 & (1) & Mathias [\textbf{1979}] p 120 \\
60 & 59 & (4) & Tsukada [\textbf{1977}] (M11)\\
73 & 59 & (3) & note 18 (N11T)\\
91 & 59 & (5) & note 18 (N11)\\
92 & 59 & (3) & note 18 (N11T)\\
130 & 59 & (3) & note 18 (N11T)\\
191 & 59 & (5) & note 18 (N11)\\
231 & 59 & (4) & Tsukada [\textbf{1977}] (M11)\\
273 & 59 & (3) & note 18 (N11T)\\
305 & 59 & (3) & note 18 (N11T)\\
317 & 59 & (5) & (N19)\\
361 & 59 & (3) & note 18 (N11T)\\
363 & 59 & (3) & note 18 (N11T)\\
60 & 60 & (1) & G. Moore\index{Author}{G. Moore} [\textbf{1982}] p 125 \\
62 & 61 & (1) & clear\\
218 & 61 & (1) & clear\\
284 & 61 & (1) & note 36\\
17 & 62 & (4) & Brunner\index{Author}{Brunner} [\textbf{1982a}] (N1T)\\
30 & 62 & (1) & clear\\
60 & 62 & (1) & clear\\
61 & 62 & (3) & Levy\index{Author}{Levy} [\textbf{1962}], Pincus\index{Author}{Pincus} [\textbf{1969}], \cite{bib:1972c1}, Jech \cite{bib:1973b} p 107\\
&&&(N6T)\\
67 & 62 & (6) & Levy [\textbf{1962}] (N6)\\
68 & 62 & (1) & Howard\index{Author}{Howard} [\textbf{1985}]\\
85 & 62 & (1) & clear\\
89 & 62 & (6) & Brunner [\textbf{1982a}] (N1)\\
107 & 62 & (1) & clear\\
123 & 62 & (1) & Pincus [\textbf{1976}]\\
127 & 62 & (4) & Brunner [\textbf{1984b}] (N2$(n,, M)\mathrm{T}^{\ast}$)\\
133 & 62 & (6) & Brunner [\textbf{1982a}] (N1)\\
270 & 62 & (1) & Kolany/Wojtylak\index{Author}{Kolany/Wojtylak} [\textbf{1991}]\\
323 & 62 & (1) & note 70\\
343 & 62 & (1) & clear\\
346 & 62 & (1) & note 127\\
0 &  63 & (3) &Feferman\index{Author}{Feferman} [\textbf{1965}] (M2)\\
14 & 63 & (1) & clear\\
40 & 63 & (4) & Feferman [\textbf{1965}], Truss\index{Author}{Truss} [\textbf{1978}] (M2)\\
\end{tabular}
\end{table}

\clearpage

\begin{table}[h]
\begin{tabular}{llll}
43 & 63 & (3) & Felgner [\textbf{1971a}] (M2)\\
63 & 63\index{Author}{Bell/Jellett}\index{Author}{Blass}\index{Author}{Haddad/Morillon}\index{Author}{Halpern/Levy}\index{Author}{Hamel}\index{Author}{Jech/Sochor}\index{Author}{H. Rubin/ J. Rubin} & (1) & Jech\index{Author}{Jech} \cite{bib:1973b} p 172 prob 8.5\\
88 & 63 & (7) & Morillon \cite{bib:1988h}\\
123 & 63 & (1) & clear\\
133 & 63 & (1) & Brunner\index{Author}{Brunner} [\textbf{1982a}]\\
331 & 63 & (7) & Morillon \cite{bib:1988h}\\
332 & 63 & (7) & Morillon \cite{bib:1988h}\\
343 & 63 & (7) & Morillon \cite{bib:1988h}\\
344 & 63 & (7) & Morillon \cite{bib:1988h}\\
0 &  64 &(3) &Levy\index{Author}{Levy} [\textbf{1958}], Jech \cite{bib:1973b} p 95  prob 6.5, Monro [\textbf{1972}]\\
&&&(N1T,M37)\\
9 & 64 &(1) &clear, Levy [\textbf{1958}]\\
16 & 64 & (5) & (N1)\\
17 & 64 & (3) & (N1T)\\
23 & 64 & (5) & (N1)\\
24 & 64 & (5) & (N1)\\
31 & 64 & (3) & (N1T)\\
45 & 64 & (1) & Truss\index{Author}{Truss} [\textbf{1974a}]\\
57 & 64 & (1) & Truss [\textbf{1974a}]\\
62 & 64 & (1) & Brunner [\textbf{1984b}]\\
63 & 64 & (6) & (N1)\\
64 & 64 & (1) & Levy [\textbf{1958}]\\
83 & 64 & (1) & Jech \cite{bib:1973b} p 52  prob 4.10\\
88 & 64 & (1) & Truss [\textbf{1974a}]\\
89 & 64 & (5) & (N1)\\
112 & 64 & (5) & (N1)\\
114 & 64 & (5) & (N1)\\
115 & 64 & (5) & (N1)\\
116 & 64 & (5) & (N1)\\
121 & 64 & (3) & (N24T)\\
127 & 64 & (5) & Brunner [\textbf{1984b}] (N1)\\
133 & 64 & (5) & Brunner [\textbf{1982a}],[\textbf{1984b}] (N1)\\
134 & 64 & (5) & (N1)\\
135 & 64 & (5) & (N1)\\
147 & 64 & (5) & (N1)\\
163 & 64 & (3) & (N1T)\\
165 & 64 & (3) & (N1T)\\
185 & 64 & (3) & (N24T)\\
217 & 64 & (3) & (N1T)\\
233 & 64 & (5) & (N1)\\
263 & 64 & (5) & (N1)\\
296 & 64 & (1) & clear\\
304 & 64 & (5) & (N1)\\
322 & 64 & (5) & (N1)\\
325 & 64 & (3) & (N1T)\\
336 & 64 & (1) & De la Cruz/Di Prisco\index{Author}{De la Cruz/Di Prisco} [\textbf{1996}]\\
361 & 64 & (3) & note 18 (N1T)\\
363 & 64 & (3) & note 18 (N1T)\\
371 & 64 & (3) & (M37)\\
\end{tabular}
\end{table}

\clearpage

\begin{table}[h]
\begin{tabular}{llll}
380 & 64 & (5) & (N1)\\
14 & 65 & (3) & H.\,Rubin/J. Rubin \cite{bib:1985q}, Halpern/Levy\index{Author}{Levy} [\textbf{1971}] (M1)\\
15 & 65 & (3) & H.\,Rubin/J. Rubin \cite{bib:1985q}, Halpern/Levy [\textbf{1971}] (M1)\\
31 & 65 & (3) & Jech\index{Author}{Jech} \cite{bib:1973b},\,H.\,Rubin/J. Rubin \cite{bib:1985q} p 179  (M1)\\
60 & 65 & (3) & Jech \cite{bib:1973b}, H. Rubin/J. Rubin \cite{bib:1985q} p 179  (M1)\\
82 & 65 & (3) & (M1)\\
118 & 65 & (3) & Haddad/Morillon [\textbf{1990}], H. Rubin/J. Rubin \cite{bib:1985q} p 179\\
&&&(M1)\\
128 & 65 & (3) & Brunner\index{Author}{Brunner} [\textbf{1984f}], H. Rubin/J. Rubin \cite{bib:1985q} p 179  (M1)\\
163 & 65 & (3) & Brunner [\textbf{1982a}], H. Rubin/J. Rubin \cite{bib:1985q} p 179 (M1)\\
165 & 65 & (3) & (M1)\\
191 & 65 & (3) & Blass\index{Author}{Blass} [\textbf{1979}], H. Rubin/J. Rubin \cite{bib:1985q} p 179  (M1)\\
277 & 65 & (3) & (M1)\\
286 & 65 & (1) & Bell/Jellett [\textbf{1971}]\\
0 &  66 & (3) &Jech/Sochor\index{Author}{Jech/Sochor} [\textbf{1966a}],[\textbf{1966b}], Blass [\textbf{1984a}], Jech \cite{bib:1973b}\\
&&&p 145  (66 implies AC in ZF)\\
66 & 66 & (1) & Hamel [\textbf{1905}]\\
109 & 66 & (1) & clear\\
44 & 67 & (3) & (N12$(\aleph_{2})$T)\\
66 & 67 & (1) & Blass [\textbf{1984a}]\\
67 & 67 & (1) & Levy [\textbf{1962}]\\
89 & 67 & (5) & Jech \cite{bib:1973b} p 134,\,Felgner/Jech\index{Author}{Felgner/Jech} [\textbf{1973}], (N1)\\
95 & 67 & (1) & Jech \cite{bib:1973b} p 148 prob 10.4, Bleicher [\textbf{1964}]\\
110 & 67 & (7) & Blass [\textbf{1984a}]\\
114 & 67 & (5) & Brunner [\textbf{1985a}] (N1)\\
115 & 67 & (5) & (N1)\\
116 & 67 & (5) & (N1)\\
133 & 67 & (5) & (N1)\\
134 & 67 & (5) & (N1)\\
135 & 67 & (5) & (N1)\\
147 & 67 & (5) & Brunner [\textbf{1983d}], note 26 (N1)\\
149 & 67 & (1) & Brunner [\textbf{1983d}], note 26\\
232 & 67 & (5) & Howard/Keremedis/Rubin/Stanley\index{Author}{Howard/Keremedis/Rubin/Stanley} [\textbf{1997}] (N1)\\
187 & 67 & (5) & (N12$(\aleph_{1})$)\\
202 & 67 & (5) & (N12$(\aleph_{1})$)\\
218 & 67 & (1) & clear\\
263 & 67 & (5) & (N1)\\
325 & 67 & (3) & (N1T)\\
333 & 67 & (1) & clear\\
334 & 67 & (1) & clear\\
0 & 68 & (3) & L\"{a}uchli\index{Author}{L\"{a}uchli} [\textbf{1962}], Jech/Sochor [\textbf{1966a}], Jech \cite{bib:1973b} p 145 (M4)\\
68 & 68 & (1) & Jech \cite{bib:1973b} p 12\\
0 & 69 & (6) & Jech \cite{bib:1973b} p 147 , Jech/Sochor [\textbf{1966a}], [\textbf{1966b}] (N1)\\
11 & 69 & (5) & note 18 (N1)\\
14 & 69 & (1) & Banaschewski\index{Author}{Banaschewski} [\textbf{1992}]\\
16 & 69 & (5) & (N1)\\
17 & 69 & (5) & (N1)\\
23 & 69 & (5) & (N1)\\
24 & 69 & (5) & (N1)\\
\end{tabular}
\end{table}

\clearpage

\begin{table}[h]
\begin{tabular}{llll}
31 & 69 & (5) & (N1)\\
37 & 69 & (5) & note 18 (N1)\\
63 & 69 & (5) & (N1)\\
69 & 69 & (1) & Jech\index{Author}{Jech} \cite{bib:1973b} p 13 \\
73 & 69 & (5) & note 18 (N1)\\
89 & 69 & (5) & (N1)\\
92 & 69 & (5) & note 18 (N1)\\
112 & 69 & (5) & (N1)\\
114 & 69 & (5) & (N1)\\
115 & 69 & (5) & (N1)\\
116 & 69 & (5) & (N1)\\
127 & 69 & (5) & (N1)\\
130 & 69 & (5) & note 18 (N1)\\
133 & 69 & (5) & (N1)\\
134 & 69 & (5) & (N1)\\
135 & 69 & (5) & (N1)\\
147 & 69 & (5) & (N1)\\
191 & 69 & (5) & note 18 (N1)\\
217 & 69 & (5) & (N1)\\
233 & 69 & (5) & (due to Pincus) Mathias\index{Author}{Mathias} [\textbf{1977b}], note 41 (N1)\\
263 & 69 & (5) & (N1)\\
273 & 69 & (5) & note 18 (N1)\\
304 & 69 & (5) & (N1)\\
305 & 69 & (5) & note 18 (N1)\\
309 & 69 & (5) & note 18 (N1)\\
313 & 69 & (5) & note 18 (N1)\\
325 & 69 & (5) & (N1)\\
361 & 69 & (5) & note 18 (N1)\\
363 & 69 & (5) & note 18 (N1)\\
380 & 69 & (5) & (N1)\\
0 & 70 & (3) & Feferman\index{Author}{Feferman} [\textbf{1965}], Jech \cite{bib:1973b} p 82  prob 5.24 (M2)\\
14 & 70 & (2) & clear\\
40 & 70 & (3) & Truss\index{Author}{Truss} [\textbf{1978}], note 65 (M2)\\
63 & 70 & (1) & clear\\
70 & 70 & (1) & Jech \cite{bib:1973b} prob 5.24\\
79 & 70 & (1) & clear\\
99 & 70 & (1) & Schechter\index{Author}{Schechter} [\textbf{1996a}]\\
204 & 70 & (3) & Truss [\textbf{1978}], note 65 (M2)\\
225 & 70 & (1) & clear\\
344 & 70 & (7) & Morillon\index{Author}{Morillon} \cite{bib:1988h}\\
0 &  71 & (3) & Jech \cite{bib:1973b} p 123  thm 8.3, p 124  thm 8.6, Jech [\textbf{1966c}],\\
&&&Pincus [\textbf{1969}] (M4)\\
86 & 71 & (0) & note 9\\
87 & 71 & (1) & Jech \cite{bib:1973b} p 120  thm 8.1, Levy\index{Author}{Levy} [\textbf{1964}], note 9\\
0 &  72 & (7) & This project\\
14 & 72 & (1) & Tarski\index{Author}{Tarski} [\textbf{1954b}]\\
72 & 72 & (1) & Tarski [\textbf{1954b}]\\
8 &73 &(1) &Jech \cite{bib:1973b} p 21 \\
73 & 73 & (1) & Jech \cite{bib:1973b} p 21 \\
\end{tabular}
\end{table}

\clearpage

\begin{table}[h]
\begin{tabular}{llll}
0 &  74 & (7) & This project\\
73 & 74\index{Author}{Blass}\index{Author}{Cohen}\index{Author}{Haddad/Morillon}\index{Author}{Hilbert}\index{Author}{Urysohn} & (1) & notes 4,5\\
74 & 74 & (1) & Jech\index{Author}{Jech} \cite{bib:1973b} p 21 \\
8&75 &(1) &Jech \cite{bib:1973b} p 21 \\
75 & 75 & (1) & Jech \cite{bib:1973b} p 21 \\
6&76 &(3) &note 18 (N3T)\\
11 & 76 & (3) & note 18 (N3T)\\
16 & 76 & (5) & (N3)\\
23 & 76 & (5) & (N3)\\
24 & 76 & (5) & (N3)\\
31 & 76 & (3) & note 18 (N3T)\\
37 & 76 & (3) & note 18 (N3T)\\
60 & 76 & (3) & note 18 (N3T)\\
67 & 76 & (1) & clear\\
73 & 76 & (3) & note 18 (N3T)\\
89 & 76 & (5) & (N1)\\
91 & 76 & (5) & (N3)\\
92 & 76 & (3) & note 18 (N3T)\\
112 & 76 & (5) & (N1)\\
114 & 76 & (5) & (N1)\\
115 & 76 & (5) & (N1)\\
116 & 76 & (5) & (N1)\\
128 & 76 & (3) & note 18 (N3T)\\
130 & 76 & (3) & note 18 (N3T)\\
133 & 76 & (5) & (N1)\\
134 & 76 & (5) & (N1)\\
135 & 76 & (5) & (N1)\\
147 & 76 & (5) & (N1)\\
164 & 76 & (5) & (N3)\\
165 & 76 & (3) & (N3T)\\
185 & 76 & (3) & note 18 (N1T)\\
214 & 76 & (1) & clear\\
217 & 76 & (3) & note 18 (N1T)\\
232 & 76 & (5) & (N1)\\
263 & 76 & (5) & (N1)\\
304 & 76 & (5) & (N1)\\
305 & 76 & (3) & note 18 (N3T)\\
309 & 76 & (3) & note 18 (N3T)\\
313 & 76 & (3) & note 18 (N3T)\\
317 & 76 & (5) & (N3)\\
361 & 76 & (3) & note 18 (N3T)\\
363 & 76 & (3) & note 18 (N3T)\\
381 & 76 & (5) & (N3)\\
9 & 77 & (3) & (N10T)\\
43 & 77 & (1) & Jech \cite{bib:1973b} p 23 \\
51 & 77 & (1) & H\"{o}ft/Howard\index{Author}{H\"{o}ft/Howard} [\textbf{1994}]\\
77 & 77 & (1) & Jech \cite{bib:1973b} p 23 \\
316 & 77 & (1) & H\"{o}ft/Howard \cite{bib:1981b}\\
0 & 78 & (3) & L\"{a}uchli\index{Author}{L\"{a}uchli} [\textbf{1962}] pp 15-18, Jech/Sochor\index{Author}{Jech/Sochor} [\textbf{1966a}] p 354 (N8T)\\
\end{tabular}
\end{table}

\clearpage

\begin{table}[h]
\begin{tabular}{llll}
43 & 78 & (1) & Jech\index{Author}{Jech} \cite{bib:1973b} p 29 prob 2.26, Kelly 1 p 115\\
78 & 78\index{Author}{H\"{o}ft/Howard} & (1) & Urysohn [\textbf{1925}] pp 290-292\\
85 & 78 & (3) & Brunner\index{Author}{Brunner} [\textbf{1980}] (N3T)\\
106 & 78 & (1) & Blass [\textbf{1979}]\\
163 & 78 & (3) & (N3T)\\
165 & 78 & (3) & (N3T)\\
375 & 78 & (1) & clear\\
0 &  79 &(3) &Jech \cite{bib:1973b} p 67 , Cohen [\textbf{1963}], [\textbf{1965}],[\textbf{1966}] (M1)\\
14 & 79 & (3) & Jech \cite{bib:1973b} pp 67, 97ff (M1)\\
15 & 79 & (3) & Jech \cite{bib:1973b} p 67,\,p 82  prob 5.21 (M1)\\
30 & 79 & (4) & Jech \cite{bib:1973b} pp 71,67 (M1)\\
60 & 79 & (3) & Jech \cite{bib:1973b} p 67,\,p 82  prob 5.22 (M1)\\
79 & 79 & (1) & Hilbert [\textbf{1900}] p 263 \\
82 & 79 & (3) & Jech \cite{bib:1973b} p 67, 81 prob 5.20 (M1)\\
91 & 79 & (1) & clear\\
118 & 79 & (3) & Haddad/Morillon [\textbf{1990}], Jech \cite{bib:1973b} p 67  (M1)\\
130 & 79 & (1) & clear\\
0 &  80 & (3) &Jech \cite{bib:1973b} p 71 , Cohen [\textbf{1963}], [\textbf{1965}],[\textbf{1966}] (M7)\\
10 & 80 & (1) & clear\\
61 & 80 & (1) & clear\\
65 & 80 & (6) & (N2)\\
67 & 80 & (5) & Fraenkel\index{Author}{Fraenkel} [\textbf{1922}], Jech \cite{bib:1973b}, (N2)\\
80 & 80 & (1) & Mostowski\index{Author}{Mostowski} [\textbf{1948}]\\
88 & 80 & (1) & clear\\
99 & 80 & (5) & Howard\index{Author}{Howard} [\textbf{1984b}] (N2)\\
111 & 80 & (1) & clear\\
122 & 80 & (1) & clear\\
163 & 80 & (3) & Brunner [\textbf{1982a}], Fraenkel [\textbf{1922}] (N2T)\\
171 & 80 & (7) & Shannon\index{Author}{Shannon} [\textbf{1990}]\\
250 & 80 & (1) & clear\\
0 & 81 & (3) & Monro\index{Author}{Monro} [\textbf{1972}], Jech \cite{bib:1973b} p 81  prob 5.17 (M16$(n)$)\\
0 &  82 &(3) &Pincus [\textbf{1969}], Jech \cite{bib:1973b} p 95  prob 6.6 (N1T)\\
9 & 82 & (1) & clear\\
14 & 82 & (3) & Pincus [\textbf{1972a}],[\textbf{1974b}](N3T)\\
15 & 82 & (1) & Monro [\textbf{1972}]\\
30 & 82 & (4) & Howard/Yorke\index{Author}{Howard/Yorke} [\textbf{1989}] (N3T)\\
82 & 82 & (1) & Levy\index{Author}{Levy} [\textbf{1958}]\\
126 & 82 & (1) & note 76\\
0 & 83 & (4) & Jech \cite{bib:1973b} p 95  prob 6.8 (N4T)\\
9 & 83 & (1) &Jech \cite{bib:1973b} p 52  prob 9\\
30 & 83 & (1) & note 81\\
64 & 83 & (3) & Pincus [\textbf{1972a}], note 105 (N4T)\\
82 & 83 & (1) & Howard/Yorke [\textbf{1989}]\\
83 & 83 & (1) & Tarski\index{Author}{Tarski} [\textbf{1924}]\\
215 & 83 & (1) & clear\\
0 &  84 &(4) &Jech \cite{bib:1973b} p 95  prob 6.9, Howard/Yorke [\textbf{1989}] (N3T)\\
9 &84 &(1) &Howard/Yorke [\textbf{1989}]\\
11 & 84 & (3) & note 18 (N3T)\\
14 & 84 & (3) & Levy [\textbf{1958}], Pincus [\textbf{1974b}] (N3T)\\
\end{tabular}
\end{table}

\clearpage

\begin{table}[h]
\begin{tabular}{llll}
16 & 84 & (5) & (N3)\\
23 & 84\index{Author}{Felgner/Jech}\index{Author}{Howard/Rubin}\index{Author}{Kurepa}\index{Author}{H. Rubin/ J. Rubin} & (5) & (N3)\\
24 & 84 & (5) & (N3)\\
31 & 84 & (3) & Howard/Yorke [\textbf{1989}], note 88 (N3T)\\
37 & 84 & (3) & note 18 (N3T)\\
60 & 84 & (3) & (N3T)\\
73 & 84 & (3) & note 18 (N3T)\\
82 & 84 & (1) & Howard/Yorke [\textbf{1989}]\\
84 & 84 & (1) & Howard/Yorke [\textbf{1989}]\\
92 & 84 & (3) & note 18 (N3T)\\
128 & 84 & (3) & (N3T)\\
130 & 84 & (3) & note 18 (N3T)\\
163 & 84 & (3) & (N3T)\\
164 & 84 & (5) & Brunner\index{Author}{Brunner} [\textbf{1982a}], Howard/Yorke [\textbf{1989}] (N3)\\
165 & 84 & (5) & (N3)\\
185 & 84 & (1) & H\"{o}ft/Howard\index{Author}{H\"{o}ft/Howard} [\textbf{1994}]\\
191 & 84 & (5) & note 18 (N3)\\
273 & 84 & (3) & note 18 (N3T)\\
305 & 84 & (3) & note 18 (N3T)\\
317 & 84 & (5) & (N3)\\
361 & 84 & (3) & note 18 (N3T)\\
363 & 84 & (3) & note 18 (N3T)\\
381 & 84 & (5) & Howard/Keremedis/Rubin/Rubin\index{Author}{Howard/Keremedis/Rubin/Rubin} [\textbf{1997b}] (N3)\\
60 & 85 & (1) & clear\\
62 & 85 & (4) & note 74 (M6)\\
85 & 85 & (1) & Jech\index{Author}{Jech} \cite{bib:1973b} p 115  prob 7.13\\
97 & 85 & (3) & (M1$(<\omega_{1}>)$)\\
213 & 85 & (1) & clear\\
40 & 86 & (1) & clear\\
86 & 86 & (1) & Levy\index{Author}{Levy} [\textbf{1964}]\\
87 & 86 & (1) & Jech \cite{bib:1973b} p 120  thm 8.1, Levy [\textbf{1964}], note 9\\
174 & 86 & (1) & H\"{o}ft/Howard \cite{bib:1981b}\\
87 & 87 & (1) & Levy [\textbf{1964}]\\
174 & 87 & (7) & H\"{o}ft/Howard \cite{bib:1981b}\\
3 & 88 & (5) &(N9)\\
40 & 88 & (3) & Jech \cite{bib:1973b} p 131  prob 8.1, Truss\index{Author}{Truss} [\textbf{1978}], note 65 (M2)\\
44 & 88 & (3) & Jech \cite{bib:1973b} thm 8.3 (N2$(\aleph_{2})$T)\\
61 & 88 & (1) & clear\\
63 & 88 & (3) & (N9T)\\
67 & 88 & (6) & Zuckerman\index{Author}{Zuckerman} [\textbf{1971}], note 15 (N2)\\
87 & 88 & (3) & Jech \cite{bib:1973b} thm 8.3 (N2$(\aleph_{\alpha})$T)\\
187 & 88 & (5) & (N12$(\aleph_{1})$)\\
201 & 88 & (1) & Schrijver\index{Author}{Schrijver} \cite{bib:1978j}\\
202 & 88 & (5) & (N12$(\aleph_{1})$)\\
204 & 88 & (3) & Truss [\textbf{1978}], note 65 (M2)\\
231 & 88 & (4) & Felgner\index{Author}{Felgner} [\textbf{1974}] (M2)\\
326 & 88 & (1) & Wojtylak\index{Author}{Wojtylak} [\textbf{1995}]\\
333 & 88 & (1) & clear\\
67 & 89 & (1) & Jech [\textbf{1966a}] p 133\\
\end{tabular}
\end{table}

\clearpage

\begin{table}[h]
\begin{tabular}{llll}
89 & 89 & (1) & Jech\index{Author}{Jech} \cite{bib:1973b} p 133, Kurepa [\textbf{1953}]\\
90 & 89\index{Author}{Jech/Sochor}\index{Author}{Mirsky} & (5) & Jech \cite{bib:1973b} p 134  (N4)\\
89 & 90 & (1) & Jech \cite{bib:1973b} p 133 \\
90 & 90 & (1) & Jech \cite{bib:1973b} p 133, Felgner/Jech [\textbf{1973}]\\
91 & 90 & (5) & Jech \cite{bib:1973b} p 134, Felgner/Jech [\textbf{1973}] note 18 (N3)\\
112 & 90 & (1) & H.\,Rubin/J. Rubin \cite{bib:1985q} p 79 \\
114 & 90 & (1) & Brunner [\textbf{1985a}]\\
133 & 90 & (1) & Brunner [\textbf{1982a}]\\
292 & 90 & (1) & Howard/Rubin [\textbf{1977}]\\
51 & 91 & (3) & (91 implies AC in ZF) (M14)\\
87 & 91 & (3) & note 18 (N40$(\kappa)$)\\
90 & 91 & (1) & Jech \cite{bib:1973b} p 133\\
91 & 91 & (1) & Jech \cite{bib:1973b} p 133, H. Rubin/J. Rubin \cite{bib:1985q}\\
97 & 91 & (3) & (91 implies AC in ZF) (M1$(<\omega_{1}>)$)\\
133 & 91 & (1) & Brunner [\textbf{1984b}]\\
147 & 91 & (1) & Brunner [\textbf{1983d}], note 26\\
148 & 91 & (7) & Brunner [\textbf{1983d}], note 26\\
164 & 91 & (1) & Brunner [\textbf{1982a}]\\
179 & 91 & (3) & (91 implies AC in ZF) (M35$(\epsilon)$)\\
200 & 91 & (3) & (91 implies AC in ZF) (M29)\\
202 & 91 & (1) & note 75 (see (202,1) entry)\\
214 & 91 & (3) & (91 implies AC in ZF) (M1($\langle w_{1}\rangle$))\\
277 & 91 & (3) & (91 implies AC in ZF) (M1)\\
290 & 91 & (3) & (91 implies AC in ZF) (M29)\\
291 & 91 & (3) & (91 implies AC in ZF) (M29)\\
345 & 91 & (3) & note 18 (N40$(\kappa)$)\\
75 & 92 & (1) & clear\\
43 & 93 & (3) & Jech \cite{bib:1973b} p 144, Solovay\index{Author}{Solovay} [\textbf{1965}] (M$(\aleph)$5 IC)\\
52 & 93 & (1) & Foreman/Wehrung\index{Author}{Foreman/Wehrung} [\textbf{1991}]\\
70 & 93 & (1) & Jech \cite{bib:1973b} p 7  prob 10\\
88 & 93 & (1) & Jech \cite{bib:1973b} p 7  prob 10\\
100 & 93 & (1) & Lindenbaum/Tarski\index{Author}{Lindenbaum/Tarski} [\textbf{1926}], G. Moore\index{Author}{G. Moore} [\textbf{1982}]\\
101 & 93 & (1) & Sierpi\'{n}ski\index{Author}{Sierpi\'{n}ski} [\textbf{1918}],[\textbf{1947}], G. Moore [\textbf{1982}]\\
142 & 93 & (3) & Shelah\index{Author}{Shelah} [\textbf{1985}] (M38 IC)\\
169 & 93 & (3) & Shelah [\textbf{1985}] (M38 IC)\\
170 & 93 & (1) & Shelah (cited in G. Moore [\textbf{1982}] p 118) Raisonnier\index{Author}{Raisonnier} [\textbf{1982}]\\
184 & 93 & (1) & Kaniewski/Rogers\index{Author}{Kaniewski/Rogers} [\textbf{1980}]\\
195 & 93 & (1) & Harazisville\index{Author}{Kharazishvili} (see Kharazishvili [\textbf{1979}])\\
222 & 93 & (7) & Pincus \cite{bib:1972c1}, Foreman/Wehrung [\textbf{1991}]\\
251 & 93 & (1) & Ash\index{Author}{Ash} [\textbf{1975}]\\
252 & 93 & (1) & Ash [\textbf{1975}]\\
306 & 93 & (1) & Sierpi\'{n}ski [\textbf{1918}]\\
307 & 93 & (3) & (M$(\aleph)$5 IC)\\
309 & 93 & (1) & clear\\
366 & 93 & (1) & G. Moore [\textbf{1982}] table 5\\
368 & 93 & (1) & Sierpi\'{n}ski [\textbf{1917}]\\
8 & 94 & (1) & clear\\
79 & 94 & (1) & clear\\
94 & 94 & (1) & Jech \cite{bib:1973b} p 148  prob 10.1\\
\end{tabular}
\end{table}

\clearpage

\begin{table}[h]
\begin{tabular}{llll}
203 & 94 & (1) & note 67\\
212 & 94 & (1) & clear\\
218 & 95 & (1) & (see [218 A])\\
6 & 96 & (3) & note 18 (N13T)\\
11 & 96 & (3) & note 18 (N13T)\\
35 & 96 & (3) & note 18 (N13T)\\
37 & 96 & (3) & note 18 (N13T)\\
73 & 96 & (3) & note 18 (N13T)\\
91 & 96 & (5) & note 18 (N13)\\
92 & 96 & (3) & note 18 (N13T)\\
107 & 96 & (1) & Mirsky [\textbf{1971}]\\
130 & 96 & (3) & note 18 (N13T)\\
191 & 96 & (5) & note 18 (N13)\\
273 & 96 & (3) & note 18 (N13T)\\
305 & 96 & (3) & note 18 (N13T)\\
309 & 96 & (3) & note 18 (N13T)\\
313 & 96 & (3) & note 18 (N13T)\\
0 & 97 &(6) &Jech\index{Author}{Jech} \cite{bib:1973b} p 154, Pincus [\textbf{1969}], Jech/Sochor [\textbf{1966a}],\\
&&&[\textbf{1966b}] (N3)\\
11 & 97 & (5) & note 18 (N3)\\
14 & 97 & (5) & Jech \cite{bib:1973b} p 154, Pincus [\textbf{1969}], [\textbf{1974b}] (N3)\\
16 & 97 & (5) & (N3)\\
23 & 97 & (5) & (N3)\\
24 & 97 & (5) & (N3)\\
31 & 97 & (5) & (N3)\\
37 & 97 & (5) & note 18 (N3)\\
60 & 97 &(5) & (N3)\\
73 & 97 & (5) & note 18 (N3)\\
91 & 97 & (5) & note 18 (N3)\\
92 & 97 & (5) & note 18 (N3)\\
128 & 97 & (5) & (N3)\\
130 & 97 & (5) & note 18 (N3)\\
164 & 97 & (5) & (N3)\\
97 & 97 & (1) & Jech \cite{bib:1973b} p 154 \\
165 & 97 & (5) & (N3)\\
191 & 97 & (5) & note 18 (N3)\\
273 & 97 & (5) & note 18 (N3)\\
305 & 97 & (5) & note 18 (N3)\\
317 & 97 & (5) & (N3)\\
361 & 97 & (5) & note 18 (N3)\\
363 & 97 & (5) & note 18 (N3)\\
381 & 97 & (5) & Howard/Keremedis/Rubin/Rubin\index{Author}{Howard/Keremedis/Rubin/Rubin} [\textbf{1997b}] (N3)\\
0 & 98 & (3) & Jech \cite{bib:1973b} p 161  prob 11.5 (N2T)\\
6&98 & (3) & note 18 (N2T)\\
9 & 98 & (1) & clear\\
11 & 98 & (3) & note 18 (N2T)\\
35 & 98 & (3) & note 18 (N2T)\\
37 & 98 & (3) & note 18 (N2T)\\
52 & 98 & (3) & (N2T)\\
\end{tabular}
\end{table}

\clearpage

\begin{table}[h]
\begin{tabular}{llll}
65 & 98 & (5) & (N2)\\
67 & 98\index{Author}{Jech/Sochor}\index{Author}{Rado} & (5) & (N2)\\
73 & 98 & (3) & note 18 (N2T)\\
82 & 98 & (3) & (N2T)\\
83 & 98 & (3) & (N2T)\\
92 & 98 & (3) & note 18 (N2T)\\
98 & 98 & (1) & Jech\index{Author}{Jech} \cite{bib:1973b} p 161  prob 11.5\\
99 & 98 & (5) & (N2)\\
119 & 98 & (3) & (N2T)\\
130 & 98 & (3) & note 18 (N2T)\\
163 & 98 & (3) & (N2T)\\
185 & 98 & (3) & (N2T)\\
191 & 98 & (5) & note 18 (N2)\\
273 & 98 & (3) & note 18 (N2T)\\
305 & 98 & (3) & note 18 (N2T)\\
313 & 98 & (3) & note 18 (N2T)\\
334 & 98 & (5) & (N2)\\
350 & 98 & (3) & (N2T)\\
361 & 98 & (3) & note 18 (N2T)\\
363 & 98 & (3) & note 18 (N2T)\\
11 & 99 & (3) & note 18 (N7T)\\
14 & 99 & (1) & Rav\index{Author}{Rav} [\textbf{1977}]\\
15 & 99 & (3) & (M3)\\
23 & 99 & (5) & (N7)\\
49 & 99 & (5) & (N52)\\
51 & 99 & (3) & (M14)\\
60 & 99 & (3) & L\"{a}uchli\index{Author}{L\"{a}uchli} [\textbf{1964}], Rav [\textbf{1977}] (N7T\\
63 & 99 & (3) & (N7T)\\
91 & 99 & (5) & note 18 (N7)\\
99 & 99 & (1) & Rado [\textbf{1949}]\\
123 &99 &(3) &Pincus [\textbf{1976}], Rav [\textbf{1977}] (N7T)\\
130 & 99 & (3) & note 18 (N7T)\\
165 & 99 & (3) & (M3)\\
191 & 99 & (5) & note 18 (N7)\\
214 & 99 & (3) & note 120(31) (M43)\\
273 & 99 & (3) & note 18 (N7T)\\
305 & 99 & (3) & note 18 (N7T)\\
313 & 99 & (3) & note 18 (N5T)\\
7 & 100 & (7) & Banaschewski/Moore\index{Author}{Banaschewski/Moore} [\textbf{1990}]\\
100 & 100 & (1) & Lindenbaum/Tarski\index{Author}{Lindenbaum/Tarski} [\textbf{1926}]\\
101 & 100 & (1) & clear\\
168 & 100 & (1) & clear\\
187 & 100 & (5) & (N12$(\aleph_{1})$)\\
202 & 100 & (5) & (N12$(\aleph_{1})$)\\
7 & 101 & (7) & Banaschewski/Moore [\textbf{1990}]\\
101 & 101 & (1) & Sierpi\'{n}ski\index{Author}{Sierpi\'{n}ski} [\textbf{1947}]\\
168 & 101 & (7) & Levy\index{Author}{Levy} [\textbf{1963a}], Pelc\index{Author}{Pelc} [\textbf{1978}], Banaschewski/Moore [\textbf{1990}]\\
359 & 101 & (1) & clear\\
0 &  102 &(7) & This project\\
\end{tabular}
\end{table}

\clearpage

\begin{table}[h]
\begin{tabular}{llll}
62 & 102 & (1) & Jech\index{Author}{Jech} \cite{bib:1973b} p 162  prob 11.12\\
102 & 102\index{Author}{R. L. Moore} & (1) & Jech \cite{bib:1973b} p 162  prob 11.12\\
0 &  103 &(3) & Jech \cite{bib:1973b} p 164  prob 11.21, Jech/Sochor [\textbf{1966a}],[\textbf{1966b}]\\ &&&(N14T)\\
6 & 103 & (3) & note 18 (N14T)\\
11 & 103 & (3) & note 18 (N14T)\\
35 & 103 & (3) & note 18 (N14T)\\
37 & 103 & (3) & note 18 (N14T)\\
43 & 103 & (3) & (N14T)\\
63 & 103 & (5) & (N14)\\
73 & 103 & (3) & note 18 (N14T)\\
91 & 103 & (5) & note 18 (N14)\\
92 & 103 & (3) & note 18 (N14T)\\
103 & 103 & (1) & Jech \cite{bib:1973b} p 164  prob 11.21\\
130 & 103 & (3) & note 18 (N14T)\\
191 & 103 & (5) & note 18 (N14)\\
273 & 103 & (3) & note 18 (N14T)\\
305 & 103 & (3) & note 18 (N14T)\\
309 & 103 & (3) & note 18 (N14T)\\
313 & 103 & (3) & note 18 (N14T)\\
361 & 103 & (3) & note 18 (N14T)\\
363 & 103 & (3) & note 18 (N14T)\\
0 &  104 &(3) &Jech \cite{bib:1973b} p 165  prob 11.26, Gitik\index{Author}{Gitik} [\textbf{1980}],[\textbf{1985}] (M17 IC)\\
34 &104 &(1) &clear\\
104 & 104 & (1) & Jech \cite{bib:1973b} p 165  prob 11.26\\
159 & 104 & (1) & clear\\
15 & 105 & (3) & Forti/Honsell\index{Author}{Forti/Honsell} [\textbf{1985}] (M11)\\
59 & 105 & (1) & clear\\
60 & 105 & (3) & Forti/Honsell [\textbf{1985}] (M11)\\
105 & 105 & (1) & Forti/Honsell [\textbf{1985}], Jech [\textbf{1966c}]\\
165 & 105 & (3) & (M11)\\
8 & 106 & (3) & Brunner\index{Author}{Brunner} [\textbf{1983c}],[\textbf{1984b}](N38T)\\
11 & 106 & (3) & note 18 (N1T)\\
14 & 106 & (3) & Brunner [\textbf{1983c}], Pincus [\textbf{1974b}](N3T)\\
15 & 106 & (3) & (M1)\\
17 & 106 & (3) & (N1T)\\
23 & 106 & (5) & (N3)\\
31 & 106 & (3) & (N3T)\\
37 & 106 & (3) & note 18 (N1T)\\
39 & 106 & (5) & (N21$(\aleph_{\alpha+1})$)\\
41 & 106 & (5) & (N16)\\
43 & 106 & (1) & Blass\index{Author}{Blass} [\textbf{1979}], Goldblatt\index{Author}{Goldblatt} \cite{bib:1985m}\\
60 & 106 & (3) & (N3T)\\
67 & 106 & (1) & Blass [\textbf{1979}]\\
71 & 106 & (5) & (N16$(\aleph_{\gamma})$)\\
73 & 106 & (3) & note 18 (N1T)\\
78 & 106 & (7) & Fossy/Morillon\index{Author}{Fossy/Morillon} [\textbf{1995}]\\
86 & 106 & (5) & (N21$(\aleph_{\alpha+1})$)\\
89 & 106 & (5) & (N1)\\
\end{tabular}
\end{table}

\clearpage

\begin{table}[h]
\begin{tabular}{llll}
91 & 106 & (5) & note 18 (N3)\\
92 & 106\index{Author}{Birkhoff}\index{Author}{Blass}\index{Author}{Douwen}\index{Author}{Gitik}\index{Author}{Howard/Rubin}\index{Author}{Jech}\index{Author}{Krom} & (3) & note 18 (N1T)\\
106 & 106 & (1) & R.L. Moore [\textbf{1924}], Goldblatt \cite{bib:1985m}\\
112 & 106 & (5) & Brunner\index{Author}{Brunner} [\textbf{1984b}],[\textbf{1985a}] (N1)\\
113 & 106 & (5) & (N38)\\
114 & 106 & (5) & (N1)\\
115 & 106 & (5) & (N1)\\
116 & 106 & (5) & (N1)\\
118 & 106 & (3) & (M1)\\
128 & 106 & (3) & (N3T)\\
130 & 106 & (3) & note 18 (N1T)\\
133 & 106 & (5) & (N1)\\
134 & 106 & (5) & (N1)\\
135 & 106 & (5) & (N1)\\
147 & 106 & (5) & (N1)\\
163 & 106 & (3) & (M1)\\
164 & 106 & (5) & (N3)\\
188 & 106 & (1) & Blass [\textbf{1979}]\\
191 & 106 & (3) & (M1)\\
263 & 106 & (5) & (N1)\\
273 & 106 & (3) & note 18 (N1T)\\
277 & 106 & (3) & (M1)\\
305 & 106 & (3) & note 18 (N1T)\\
317 & 106 & (5) & (N3)\\
325 & 106 & (3) & (N1T)\\
363 & 106 & (3) & note 18 (N1T)\\
1 & 107 & (1) & M. Hall\index{Author}{Hall} [\textbf{1948}]\\
14 & 107 & (1) & (see [14 H])\\
0 & 108 & (7) & Gitik [\textbf{1985}]\\
38 & 108 & (1) & clear\\
108 & 108 & (1) & Gitik [\textbf{1985}]\\
207 & 108 & (1) & clear\\
0 &  110 &(7) &This project\\
110 & 110 & (1) & Blass [\textbf{1984a}]\\
250 & 111 & (1) & clear\\
67 & 112 & (1) & clear\\
112 & 112 & (1) & Howard/Rubin [\textbf{1977}]\\
8 & 113 & (7) & Brunner [\textbf{1985a}]\\
43 & 113 & (1) & note 7\\
113 & 113 & (1) & Brunner [\textbf{1985a}]\\
67 & 114 & (1) & Brunner [\textbf{1985a}]\\
90 & 114 & (5) & Brunner [\textbf{1985a}], (N4)\\
114 & 114 & (1) & Brunner [\textbf{1985a}]\\
0 &  115 &(7) &This project\\
115 & 115 & (1) & Brunner [\textbf{1985a}]\\
0 &  116 &(7) &This project\\
116 & 116 & (1) & Brunner [\textbf{1985a}]\\
0 &  117 & (4) & Apter\index{Author}{Apter} [\textbf{1985b}], Jech [\textbf{1968a}] (M($\aleph$)5 IC)\\
117 & 117 & (1) & Keisler/Tarski\index{Author}{Keisler/Tarski} [\textbf{1964}]
\end{tabular}
\end{table}

\clearpage

\begin{table}[h]
\begin{tabular}{llll}
11 & 118 & (3) & note 18 (N3T)\\
14 & 118\index{Author}{Brunner}\index{Author}{Brunner/Howard}\index{Author}{Halpern}\index{Author}{Morillon}\index{Author}{Rav} & (3) & (N3T)\\
23 & 118 & (5) & (N3)\\
30 & 118 & (4) & Krom [\textbf{1986}] (N3T)\\
31&118&(3)&(N3T)\\
43 & 118 & (3) & (N40T)\\
60 & 118 & (3) & (N3T)\\
90 & 118 & (1) & van Douwen\index{Author}{Douwen} [\textbf{1985}]\\
91 & 118 & (5) & (N3)\\
118 & 118 & (1) & Birkhoff [\textbf{1967}] p 241 \\
130 & 118 & (3) & note 18 (N3T)\\
163 &118 &(3) &(N3T)\\
164 & 118 & (5) & (N3)\\
165 & 118 & (5) & (N3)\\
191 & 118 & (5) & note 18\\
&&&(N3, N19, N38, N40)\\
200 & 118 & (3) & note 120(16) (M29)\\
214 & 118 & (3) & note 120(16) (M29)\\
273 & 118 & (3) & note 18 (N3T)\\
290 & 118 & (3) & note 120(16) (M29)\\
291 & 118 & (3) & note 120(16) (M29)\\
295 & 118 & (3) & (N48T)\\
305 & 118 & (3) & note 18 (N3T)\\
317 & 118 & (5) & (N3)\\
345 & 118 & (5) & N40\\
381 & 118 & (5) & Howard/Keremedis/Rubin/Rubin\index{Author}{Howard/Keremedis/Rubin/Rubin} [\textbf{1997b}] (N3)\\
 0 &  119 & (3) &van Douwen [\textbf{1985}] (N2(LO)T)\\
6 & 119 & (3) & note 18 (N2(LO)T)\\
11 & 119 & (3) & note 18 (N2(LO)T)\\
32 & 119 & (1) & clear\\
35 & 119 & (3) & note 18 (N2(LO)T)\\
37 & 119 & (3) & note 18 (N2(LO)T)\\
73 & 119 & (3) & note 18 (N2(LO)T)\\
91 & 119 & (5) & note 18 (N2(LO))\\
92 & 119 & (3) & note 18 (N2(LO)T)\\
118 & 119 & (1) & van Douwen [\textbf{1985}]\\
119 & 119 & (1) & van Douwen [\textbf{1985}]\\
130 & 119 & (3) & note 18 (N2(LO)T)\\
191 & 119 & (5) & note 18 (N2(LO))\\
273 & 119 & (3) & note 18 (N2(LO)T)\\
305 & 119 & (3) & note 18 (N2(LO)T)\\
309 & 119 & (3) & note 18 (N2(LO)T)\\
313 & 119 & (3) & note 18 (N2(LO)T)\\
361 & 119 & (3) & note 18 (N2(LO)T)\\
363 & 119 & (3) & note 18 (N2(LO)T)\\
375 & 119 & (1) & note 138\\
0 &  120 &(7)& This project\\
46 & 120 & (1) & clear\\
120 & 120 & (1) & Truss\index{Author}{Truss} [\textbf{1973a}]\\
\end{tabular}
\end{table}

\clearpage

\begin{table}[h]
\begin{tabular}{llll}
121 & 120 & (1) & clear\\
62 & 121\index{Author}{Howard/Keremedis/Rubin/Rubin} & (1) & clear\\
121 & 121 & (1) & Truss [\textbf{1973a}]\\
202 & 121 & (1) & clear\\
40 & 122 & (1) & clear\\
44 & 122 & (3) & (N2$(\aleph_{\alpha})$T)\\
87 & 122 & (3) & (N2$(\aleph_{\alpha})$T)\\
121 & 122 & (1) & clear\\
122 & 122 & (1) & Truss [\textbf{1973a}]\\
135 & 122 & (5) & (N2 $(\aleph_{\alpha})$)\\
151 & 122 & (1) & Brunner/Howard\index{Author}{Brunner} [\textbf{1992}], note 27\\
165 & 122 & (1) & clear\\
14 & 123 & (1) & Rav [\textbf{1977}]\\
123 & 123 & (1) & Halpern [\textbf{1972}]\\
344 & 123 & (7) & Morillon \cite{bib:1988h}\\
0 &  124 &(3) &Brunner [\textbf{1984a}](N24T)\\
6 & 124 & (3) & note 18 (N24T)\\
10 & 124 & (3) & (N24T)\\
11 & 124 & (3) & note 18 (N24T)\\
16 & 124 & (5) & (N24)\\
17 & 124 & (1) & Brunner [\textbf{1984a}]\\
23 & 124 & (5) & (N24)\\
24 & 124 & (5) & (N24)\\
31 & 124 & (3) & (N24T)\\
32 & 124 & (3) & (N24T)\\
35 & 124 & (3) & note 18 (N24T)\\
37 & 124 & (3) & note 18 (N24T)\\
73 & 124 & (3) & note 18 (N24T)\\
92 & 124 & (3) & note 18 (N24T)\\
121 & 124 & (3) & (N24T)\\
124 & 124 & (1) & Brunner [\textbf{1984a}]\\
130 & 124 & (3) & note 18 (N24T)\\
133 & 124 & (5) & (N24)\\
163 & 124 & (3) & (N24T)\\
165 & 124 & (3) & (N24T)\\
185 & 124 & (3) & (N24T)\\
191 & 124 & (5) & note 18 (N24)\\
273 & 124 & (3) & note 18 (N24T)\\
305 & 124 & (3) & note 18 (N24T)\\
309 & 124 & (3) & note 18 (N24T)\\
313 & 124 & (3) & note 18 (N24T)\\
361 & 124 & (3) & note 18 (N24T)\\
363 & 124 & (3) & note 18 (N24T)\\
0 &  125 & (3) & Brunner [\textbf{1984c}](N3T)\\
11 & 125 & (3) & note 18 (N3T)\\
14 & 125 & (3) & Brunner [\textbf{1984c}], Pincus\index{Author}{Pincus} [\textbf{1974b}] (N3T)\\
16 & 125 & (5) & (N3)\\
23 & 125 & (5) & (N3)\\
24 & 125 & (5) & (N3)\\
\end{tabular}
\end{table}

\clearpage

\begin{table}[h]
\begin{tabular}{llll}
31 & 125 & (3) & (N3T)\\
37 & 125 & (3) & note 18 (N3T)\\
60 & 125 & (3) & (N3T)\\
73 & 125 & (3) & note 18 (N3T)\\
91 & 125 & (5) & note 18 (N3)\\
92 & 125 & (3) & note 18 (N3T)\\
125 & 125 & (1) & Brunner\index{Author}{Brunner} [\textbf{1984c}]\\
128 & 125 & (3) & (N3T)\\
130 & 125 & (3) & note 18 (N3T)\\
144 & 125 & (1) & Brunner [\textbf{1984c}]\\
163 & 125 & (3) & (N3T)\\
164 & 125 & (5) & (N3)\\
165 & 125 & (5) & (N3)\\
191 & 125 & (5) & note 18 (N3)\\
273 & 125 & (3) & note 18 (N3T)\\
305 & 125 & (3) & note 18 (N3T)\\
317 & 125 & (5) & (N3)\\
361 & 125 & (3) & note 18 (N3T)\\
363 & 125 & (3) & note 18 (N3T)\\
381 & 125 & (5) & Howard/Keremedis/Rubin/Rubin\index{Author}{Howard/Keremedis/Rubin/Rubin} [\textbf{1997b}] (N3)\\
0 &  126 &(4) &Brunner\index{Author}{Brunner} [\textbf{1984d}] (N1T)\\
8 & 126 & (1) & clear\\
11 & 126 & (3) & note 18 (N1T)\\
14 & 126 & (3) & (N3T)\\
17 & 126 & (3) & (N1T)\\
23 & 126 & (5) & (N1)\\
31 & 126 & (3) & (N1T)\\
37 & 126 & (3) & note 18 (N1T)\\
47 & 126 & (3) & (N1T)\\
48 & 126 & (3) & (N1T)\\
60 & 126 & (3) & (N3T)\\
67 & 126 & (1) & clear\\
73 & 126 & (3) & note 18 (N1T)\\
89 & 126 & (5) & (N1)\\
92 & 126 & (3) & note 18 (N1T)\\
106 & 126 & (1) & Blass\index{Author}{Blass} [\textbf{1979}]\\
112 & 126 & (5) & (N1)\\
114 & 126 & (5) & (N1)\\
115 & 126 & (5) & (N1)\\
116 & 126 & (5) & (N1)\\
122 & 126 & (3) & (N1T)\\
128 & 126 & (3) & (N3T)\\
130 & 126 & (3) & note 18 (N1T)\\
133 & 126 & (5) & (N1)\\
134 & 126 & (5) & (N1)\\
135 & 126 & (5) & (N1)\\
147 & 126 & (5) & (N1)\\
163 & 126 & (3) & (N1T)\\
165 & 126 & (3) & (N1T)\\
\end{tabular}
\end{table}

\clearpage

\begin{table}[h]
\begin{tabular}{llll}
185 & 126 & (3) & (N1T)\\
217 & 126 & (3) & (N1T)\\
263 & 126 & (5) & (N1)\\
273 & 126 & (3) & note 18 (N1T)\\
305 & 126 & (3) & note 18 (N1T)\\
317 & 126 & (5) & (N3)\\
325 & 126 & (3) & (N1T)\\
328 & 126 & (1) & clear\\
361 & 126 & (3) & note 18 (N1T)\\
363 & 126 & (3) & note 18 (N1T)\\
381 & 126 & (5) & Howard/Keremedis/Rubin/Rubin [\textbf{1997b}] (N3)\\
0 &  127 &(3) &Brunner [\textbf{1984b}](N24T)\\
6 & 127 & (3) & note 18 (N24T)\\
10 & 127 & (3) & (N24T)\\
11 & 127 & (3) & note 18 (N24T)\\
16 & 127 & (5) & (N24)\\
17 & 127 & (7) & Brunner [\textbf{1984b}]\\
23 & 127 & (5) & (N24)\\
24 & 127 & (5) & (N24)\\
31 & 127 & (3) & (N24T)\\
32 & 127 & (3) & (N24T)\\
35 & 127 & (3) & note 18 (N24T)\\
37 & 127 & (3) & note 18 (N24T)\\
64 & 127 & (1) & Brunner [\textbf{1984b}]\\
73 & 127 & (3) & note 18 (N24T)\\
92 & 127 & (3) & note 18 (N24T)\\
121 & 127 & (3) & (N24T))\\
127 & 127 & (1) & Brunner [\textbf{1984b}]\\
130 & 127 & (3) & note 18 (N24T)\\
133 & 127 & (5) & (N24)\\
163 & 127 & (3) & (N24T)\\
165 & 127 & (3) & (N24T)\\
185 & 127 & (3) & (N24T)\\
191 & 127 & (5) & note 18 (N24)\\
273 & 127 & (3) & note 18 (N24T)\\
305 & 127 & (3) & note 18 (N24T)\\
309 & 127 & (3) & note 18 (N24T)\\
313 & 127 & (3) & note 18 (N24T)\\
361 & 127 & (3) & note 18 (N24T)\\
363 & 127 & (3) & note 18 (N24T)\\
9 & 128 & (1) & Brunner [\textbf{1984f}]\\
11 & 128 & (3) & note 18 (N1T)\\
16 & 128 & (5) & (N1)\\
17 & 128 & (3) & Blass\index{Author}{Blass} [\textbf{1977a}], Brunner [\textbf{1984f}], Pincus\index{Author}{Pincus} [\textbf{1972a}] (N1T)\\
23 & 128 & (5) & (N1)\\
24 & 128 & (5) & (N1)\\
31 & 128 & (3) & (N1T)\\
37 & 128 & (3) & note 18 (N1T)\\
47 & 128 & (3) & (N1T)\\
\end{tabular}
\end{table}

\clearpage

\begin{table}[h]
\begin{tabular}{llll}
48 & 128 & (3) & (N1T)\\
52 & 128 & (3) & (N2T)\\
63 & 128 & (5) & (N1)\\
65 & 128 & (5) & (N2)\\
67 & 128 & (5) & (N2)\\
73 & 128 & (3) & note 18 (N1T)\\
82 & 128 & (3) & (N2T)\\
83 & 128 & (3) & (N2T)\\
89 & 128 & (5) & (N1)\\
92 & 128 & (3) & note 18 (N1T)\\
99 & 128 & (5) & (N2)\\
112 & 128 & (5) & (N1)\\
114 & 128 & (5) & (N1)\\
115 & 128 & (5) & (N1)\\
116 & 128 & (5) & (N1)\\
122 & 128 & (3) & (N1T)\\
128 & 128 & (1) & Brunner\index{Author}{Brunner} [\textbf{1984f}]\\
130 & 128 & (3) & Brunner [\textbf{1984f}] (N1T)\\
133 & 128 & (5) & (N1)\\
134 & 128 & (5) & (N1)\\
135 & 128 & (5) & (N1)\\
147 & 128 & (5) & (N1)\\
163 & 128 & (3) & (N1T)\\
165 & 128 & (3) & (N1T)\\
185 & 128 & (3) & (N1T)\\
191 & 128 & (5) & note 18 (N1)\\
217 & 128 & (3) & (N1T)\\
233 & 128 & (5) & (N1)\\
263 & 128 & (5) & (N1)\\
273 & 128 & (3) & note 18 (N1T)\\
304 & 128 & (5) & (N1)\\
305 & 128 & (3) & note 18 (N1T)\\
313 & 128 & (3) & note 18 (N1T)\\
322 & 128 & (5) & (N1)\\
325 & 128 & (3) & (N1T)\\
334 & 128 & (5) & (N2)\\
361 & 128 & (3) & note 18 (N1T)\\
363 & 128 & (3) & note 18 (N1T)\\
380 & 128 & (5) & (N1)\\
0 &  129 & (4) & (N28T)\\
129 & 129 & (1) & Brunner [\textbf{1984f}]\\
316 & 129 & (5) & (N37)\\
130 & 130 & (1) & Brunner [\textbf{1984f}]\\
11 & 131 & (3) & (N1T)\\
14 & 131 & (3) & (N3T)\\
15 & 131 & (3) & (M1)\\
16 & 131 & (5) & (N1)\\
23 & 131 & (5) & (N1)\\
24 & 131 & (5) & (N1)\\
\end{tabular}
\end{table}

\clearpage

\begin{table}[h]
\begin{tabular}{llll}
31 & 131 & (3) & (N1T)\\
37 & 131\index{Author}{Shelah}\index{Author}{Solovay}\index{Author}{Stern} & (3) & (N1T)\\
60 & 131 & (3) & (N3T)\\
73 & 131 & (3) & (N1T)\\
76 & 131 & (1) & clear\\
92 & 131 & (3) & (N1T)\\
118 & 131 & (3) & (M1)\\
126 & 131 & (1) & clear\\
128 & 131 & (3) & (N3T)\\
130 & 131 & (3) & (N1T)\\
133 & 131 & (5) & (N1)\\
163 & 131 & (3) & (M1)\\
164 & 131 & (5) & (N3)\\
165 & 131 & (5) & (N3)\\
185 & 131 & (3) & (N1T)\\
191 & 131 & (3) & (M1)\\
277 & 131 & (3) & (M1)\\
305 & 131 & (3) & (N1T)\\
317 & 131 & (5) & (N3)\\
361 & 131 & (3) & (N1T)\\
363 & 131 & (3) & (N1T)\\
381 & 131 & (5) & (N3)\\
10 & 132 & (3) & Brunner\index{Author}{Brunner} [\textbf{1983b}](N24T)\\
17 & 132 & (1) & Brunner [\textbf{1984b}]\\
23 & 132 & (5) & (N24)\\
31 & 132 & (3) & (N24T)\\
32 & 132 & (3) & (N24T)\\
62 & 132 & (1) & Brunner [\textbf{1983b}]\\
121 & 132 & (3) & (N24T)\\
132 & 132 & (1) & Blass\index{Author}{Blass} [\textbf{1977a}], Kleinberg\index{Author}{Kleinberg} [\textbf{1969}]\\
150 & 132 & (1) & clear\\
165 & 132 & (3) & (N24T)\\
378 & 132 & (1) & De la Cruz/Di Prisco\index{Author}{De la Cruz/Di Prisco} [\textbf{1996}]\\
380 & 132 & (1) & clear\\
60 & 133 & (4) & Brunner [\textbf{1982a}], Jech\index{Author}{Jech} \cite{bib:1973b} p 82  prob 5.22, p81 prob 5.18\\
&&&(M1)\\
133 & 133 & (1) & Brunner [\textbf{1982a}]\\
164 & 133 & (6) & Brunner [\textbf{1982a}] (N3)\\
187 & 133 & (5) & (N12$(\aleph_{1})$)\\
202 & 133 & (5) & (N12$(\aleph_{1})$)\\
0 & 134 & (7) & This project\\
0 & 135 & (7) & This project\\
0& 136 & (3) & Truss\index{Author}{Truss} [\textbf{1984}], Jech/Sochor\index{Author}{Jech/Sochor} [\textbf{1966a}],\,[\textbf{1966b}](N20T)\\
6& 136 & (3) & note 18 (N20T)\\
11 & 136 & (3) & note 18 (N20T)\\
35 & 136 & (3) & note 18 (N20T)\\
37 & 136 & (3) & note 18 (N20T)\\
73 & 136 & (3) & note 18 (N20T)\\
91 & 136 & (5) & note 18 (N20)\\
\end{tabular}
\end{table}

\clearpage

\begin{table}[h]
\begin{tabular}{llll}
92 & 136 & (3) & note 18 (N20T)\\
130 & 136 & (3) & note 18 (N20T)\\
138 & 136 & (1) & Truss [\textbf{1984}]\\
191 & 136 & (5) & note 18 (N20)\\
273 & 136 & (3) & note 18 (N20T)\\
305 & 136 & (3) & note 18 (N20T)\\
309 & 136 & (3) & note 18 (N20T)\\
313 & 136 & (3) & note 18 (N20T)\\
361 & 136 & (3) & note 18 (N20T)\\
363 & 136 & (3) & note 18 (N20T)\\
0 &  137 &(7) &This project\\
139 & 137 & (1) & Truss [\textbf{1984}]\\
0 &  139 &(7) &This project\\
14 & 139 & (1) & (see [14 K])\\
0 &  140 &(7) &This project\\
88 & 140 & (1) & clear\\
0 &  141 &(7) &This project\\
14 & 141 & (1) & Cowen\index{Author}{Cowen} [\textbf{1977b}]\\
43 & 142 & (3) & Shelah [\textbf{1984}], Solovay [\textbf{1970}] (M18, M$(\aleph)$5 IC)\\
52 & 142 & (1) & Solovay (see Pincus\index{Author}{Pincus} \cite{bib:1972c1})\\
70 & 142 & (1) & Jech \cite{bib:1973b} p 7  prob 11\\
88 & 142 & (1) & Jech \cite{bib:1973b} p 7  prob 11\\
93 & 142 & (3) & Stern [\textbf{1985}] (M18)\\
222 & 142 & (1) & Pincus \cite{bib:1972c1}\\
310 & 142 & (1) & Wagon\index{Author}{Wagon} \cite{bib:1985u}\\
311 & 142 & (1) & Wagon \cite{bib:1985u}\\
0 &  143 &(7) &This project\\
265 & 143 & (1) & Harper/Rubin\index{Author}{Harper/Rubin} [\textbf{1976}]\\
0 &  144 &(4) &David\index{Author}{David} [\textbf{1980}] (M1)\\
8 & 144 & (3) & note 120(28) (N40T)\\
15 & 144 & (3) & (M1)\\
43 & 144 & (3) & note 120(28) (N40T)\\
51 & 144 & (3) & (M14)\\
67 & 144 & (1) & Levy\index{Author}{Levy} [\textbf{1962}]\\
82 & 144 & (3) & (M1)\\
87 & 144 & (3) & (M40$(\kappa)$)\\
118 & 144 & (3) & (M1)\\
128 & 144 & (3) & (M1)\\
163 & 144 & (3) & (M1)\\
165 & 144 & (3) & (M1)\\
179 & 144 & (1) & clear\\
191 & 144 & (3) & (M1)\\
277 & 144 & (3) & (M1)\\
345 & 144 & (3) & (M40$(\kappa)$)\\
0 &  145 & (7) &This project\\
91 & 145 & (1) & Brunner\index{Author}{Brunner} [\textbf{1984c}]\\
0 &  146 &(3) &Brunner [\textbf{1983d}], note 26 (N1T)\\
11 & 146 & (3) & note 18 (N1T)\\
16 & 146 & (5) & (N1)\\
\end{tabular}
\end{table}

\clearpage

\begin{table}[h]
\begin{tabular}{llll}
17 & 146 & (3) & (N1T)\\
23 & 146\index{Author}{Brunner/Howard}\index{Author}{Howard/Keremedis/Rubin/Rubin}\index{Author}{Jech}\index{Author}{Jech/Sochor}\index{Author}{Pincus} & (5) & (N1)\\
24 & 146 & (5) & (N1)\\
31 & 146 & (3) & (N1T)\\
37 & 146 & (3) & note 18 (N1T)\\
47 & 146 & (3) & (N1T)\\
48 & 146 & (3) & (N1T)\\
62 & 146 & (1) & Brunner\index{Author}{Brunner} [\textbf{1983d}], note 26\\
63 & 146 & (5) & (N1)\\
73 & 146 & (3) & note 18 (N1T)\\
89 & 146 & (5) & (N1)\\
92 & 146 & (3) & note 18 (N1T)\\
112 & 146 & (5) & (N1)\\
114 & 146 & (5) & (N1)\\
115 & 146 & (5) & (N1)\\
116 & 146 & (5) & (N1)\\
122 & 146 & (3) & (N1T)\\
127 & 146 & (5) & (N1)\\
130 & 146 & (3) & note 18 (N1T)\\
133 & 146 & (5) & (N1)\\
134 & 146 & (5) & (N1)\\
135 & 146 & (5) & (N1)\\
147 & 146 & (5) & (N1)\\
163 & 146 & (3) & (N1T)\\
165 & 146 & (3) & (N1T)\\
185 & 146 & (3) & (N1T)\\
191 & 146 & (5) & note 18 (N1T)\\
217 & 146 & (3) & (N1T)\\
233 & 146 & (5) & (N1)\\
263 & 146 & (5) & (N1)\\
273 & 146 & (3) & note 18 (N1T)\\
304 & 146 & (5) & (N1)\\
305 & 146 & (3) & note 18 (N1T)\\
309 & 146 & (3) & note 18 (N1T)\\
313 & 146 & (3) & note 18 (N1T)\\
322 & 146 & (5) & (N1)\\
325 & 146 & (3) & (N1T)\\
361 & 146 & (3) & note 18 (N1T)\\
363 & 146 & (3) & note 18 (N1T)\\
380 & 146 & (5) & (N1)\\
11 & 147 & (3) & note 18 (N3T)\\
23 & 147 & (5) & (N3)\\
67 & 147 & (1) & Brunner [\textbf{1983d}], note 26\\
91 & 147 & (5) & Brunner [\textbf{1983d}], notes 18,26 (N3)\\
130 & 147 & (3) & note 18 (N3T)\\
164 & 147 & (5) & (N3)\\
273 & 147 & (3) & note 18 (N3T)\\
305 & 147 & (3) & note 18 (N3T)\\
317 & 147 & (5) & (N3)\\
\end{tabular}
\end{table}

\clearpage

\begin{table}[h]
\begin{tabular}{llll}
381 & 147 & (5) & Howard/Keremedis/Rubin/Rubin [\textbf{1997b}] (N3)\\
0 &  148 &(7) &Brunner\index{Author}{Brunner} [\textbf{1983d}], note 26\\
8 & 150 & (1) & clear (8 implies 9)\\
85 & 150 & (1) & clear\\
8 & 151 & (3) & Brunner/Howard\index{Author}{Brunner/Howard} [\textbf{1992}], note 27, Jech/Sochor [\textbf{1966a}],\\
&&&[\textbf{1966b}], Pincus [\textbf{1969}], Jech \cite{bib:1973b} pp 128-130 (N15T)\\
11 & 151 & (3) & note 18 (N15T)\\
23 & 151 & (1) & clear\\
62 & 151 & (3) & (M43)\\
63 & 151 & (3) & (N15T)\\
91 & 151 & (5) & note 18 (N15)\\
97 & 151 & (3) & (M1 $(\langle\omega_{1}\rangle)$)\\
122 & 151 & (3) & (N15T)\\
130 & 151 & (3) & note 18 (N15T)\\
200 & 151 & (3) & (M29)\\
231 & 151 & (1) & clear\\
273 & 151 & (3) & note 18 (N15T)\\
290 & 151 & (3) & (M29)\\
291 & 151 & (3) & (M29)\\
295 & 151 & (3) & (M1 $(\langle\omega_{1}\rangle)$)\\
305 & 151 & (3) & note 18 (N15T)\\
309 & 151 & (3) & note 18 (N15T)\\
313 & 151 & (3) & note 18 (N15T)\\
363 & 151 & (3) & note 18 (N15T)\\
4 & 152 & (7) & Brunner/Howard [\textbf{1992}], note 27\\
8 & 152 & (3) & Brunner/Howard [\textbf{1992}], note 27 (N36$(\beta)$T)\\
11 & 152 & (3) & note 18 (N36 $(\beta)$T)\\
39 & 152 & (3) & Brunner/Howard [\textbf{1992}], note 27 (N36$(\beta)$T)\\
40 & 152 & (3) & (M2)\\
43 & 152 & (3) & (N36$(\beta)$T)\\
44 & 152 & (5) & (N12$(\aleph_{2})$)\\
63 & 152 & (3) & (N36$(\beta)$T)\\
86 & 152 & (3) & Brunner/Howard [\textbf{1992}], note 27 (N36$(\beta)$T)\\
87 & 152 & (3) & (M40$(\kappa)$)\\
91 & 152 & (5) & note 18 (N36$(\beta)$)\\
130 & 152 & (3) & note 18 (N36$(\beta)$T)\\
187 & 152 & (5) & (N12$(\aleph_{1})$)\\
202 & 152 & (5) & (N12$(\aleph_{1})$)\\
204 & 152 & (3) & (M2)\\
214 & 152 & (1) & note 140\\
273 & 152 & (3) & note 18 (N36$(\beta)$T)\\
305 & 152 & (3) & note 18 (N36$(\beta)$T)\\
345 & 152 & (3) & (M40$(\kappa)$)\\
14 & 153 & (1) & Fossy/Morillon\index{Author}{Fossy/Morillon} [\textbf{1995}]\\
6 & 154 & (3) & note 18 (N2T)\\
11 & 154 & (3) & note 18 (N2T)\\
14 & 154 & (1) & Brunner [\textbf{1983c}]\\
35 & 154 & (3) & note 18 (N2T)\\
37 & 154 & (3) & note 18 (N2T)\\
\end{tabular}
\end{table}

\clearpage

\begin{table}[h]
\begin{tabular}{llll}
43 & 154 & (1) & Brunner\index{Author}{Brunner} [\textbf{1983c}]\\
52 & 154\index{Author}{Howard/Keremedis/Rubin/Rubin} & (3) & (N2T)\\
61 & 154 & (3) & (N6T)\\
65 & 154 & (5) & (N2)\\
67 & 154 & (5) & (N2)\\
73 & 154 & (3) & note 18 (N2T)\\
82 & 154 & (3) & (N2T)\\
83 & 154 & (3) & (N2T)\\
92 & 154 & (3) & note 18 (N2T)\\
99 & 154 & (5) & (N2)\\
113 & 154 & (1) & clear\\
119 & 154 & (3) & (N2T)\\
130 & 154 & (3) & note 18 (N2T)\\
163 & 154 & (3) & (N2T)\\
185 & 154 & (3) & (N2T)\\
191 & 154 & (5) & note 18 (N2)\\
218 & 154 & (5) & (N6)\\
273 & 154 & (3) & note 18 (N2T)\\
305 & 154 & (3) & note 18 (N2T)\\
313 & 154 & (3) & note 18 (N2T)\\
322 & 154 & (5) & note 120(10) (N50($E$))\\
334 & 154 & (5) & (N2)\\
350 & 154 & (3) & (N2T)\\
357 & 154 & (3) & (N50$(E)$T)\\
358 & 154 & (3) & (N50$(E)$T)\\
361 & 154 & (3) & note 18 (N2T)\\
363 & 154 & (3) & note 18 (N2T)\\
11 & 155 & (3) & note 18 (N3T)\\
14 & 155 & (3) & Brunner [\textbf{1983a}], Pincus\index{Author}{Pincus} [\textbf{1974b}](N3T)\\
16 & 155 & (5) & (N3)\\
23 & 155 & (5) & (N3)\\
24 & 155 & (5) & (N3)\\
31 & 155 & (3) & Brunner [\textbf{1983a}], note 88 (N3T)\\
37 & 155 & (3) & note 18 (N3T)\\
60 & 155 & (3) & Brunner [\textbf{1983a}], L\"{a}uchli\index{Author}{L\"{a}uchli} [\textbf{1964}] (N3T)\\ 73 &155 &(3) &note 18 (N3T)\\
78 & 155 & (1) & Brunner [\textbf{1983a}]\\
91 & 155 & (5) & Brunner [\textbf{1983a}], note 18 (N3)\\
92 & 155 & (3) & note 18 (N3T)\\
128 & 155 & (3) & (N3T)\\
130 & 155 & (3) & note 18 (N3T)\\
163 & 155 & (3) & (N3T)\\
164 & 155 & (5) & Brunner [\textbf{1982a}],[\textbf{1983a}] (N3)\\
165 & 155 & (5) & (N3)\\
191 & 155 & (5) & note 18 (N3)\\
273 & 155 & (3) & note 18 (N3T)\\
305 & 155 & (3) & note 18 (N3T)\\
317 & 155 & (5) & (N3)\\
361 & 155 & (3) & note 18 (N3T)\\
\end{tabular}
\end{table}

\clearpage

\begin{table}[h]
\begin{tabular}{llll}
363 & 155 & (3) & note 18 (N3T)\\
381 & 155 & (5) & Howard/Keremedis/Rubin/Rubin\index{Author}{Howard/Keremedis/Rubin/Rubin} [\textbf{1997b}] (N3)\\
11 & 156 & (3) & note 18 (N3T)\\
14 & 156 & (3) & Brunner\index{Author}{Brunner} [\textbf{1983a}], Pincus\index{Author}{Pincus} [\textbf{1974b}](N3T)\\
16 & 156 & (5) & (N3)\\
23 & 156 & (5) & (N3)\\
24 & 156 & (5) & (N3)\\
31 & 156 & (3) & Brunner [\textbf{1983a}], note 88 (N3T)\\
37 & 156 & (3) & note 18 (N3T)\\
60 & 156 & (3) & Brunner [\textbf{1983a}], L\"{a}uchli\index{Author}{L\"{a}uchli} [\textbf{1964}] (N3T)\\
73 & 156 & (3) & note 18 (N3T)\\
91 & 156 & (5) & Brunner [\textbf{1983a}], note 18 (N3)\\
92 & 156 & (3) & note 18 (N3T)\\
128 & 156 & (3) & (N3T)\\
130 & 156 & (3) & note 18 (N3T)\\
163 & 156 & (3) & (N3T)\\
164 & 156 & (5) & Brunner [\textbf{1982a}],[\textbf{1983a}] (N3)\\
165 & 156 & (5) & (N3)\\
191 & 156 & (5) & note 18 (N3)\\
273 & 156 & (3) & note 18 (N3T)\\
305 & 156 & (3) & note 18 (N3T)\\
317 & 156 & (5) & (N3)\\
361 & 156 & (3) & note 18 (N3T)\\
363 & 156 & (3) & note 18 (N3T)\\
381 & 156 & (5) & Howard/Keremedis/Rubin/Rubin [\textbf{1997b}] (N3)\\
11 & 157 & (3) & note 18 (N3T)\\
14 & 157 & (3) & Brunner [\textbf{1983a}], Pincus [\textbf{1974b}] (N3T)\\
16 & 157 & (5) & (N3)\\
23 & 157 & (5) & (N3)\\
24 & 157 & (5) & (N3)\\
31 & 157 & (3) & Brunner [\textbf{1983a}], note 88 (N3T)\\
37 & 157 & (3) & note 18 (N3T)\\
60 & 157 & (3) & Brunner [\textbf{1983a}], L\"{a}uchli [\textbf{1964}] (N3T)\\
73 & 157 & (3) & note 18 (N3T)\\
91 & 157 & (5) & Brunner [\textbf{1983a}], note 18 (N3)\\
92 & 157 & (3) & note 18 (N3T)\\
128 & 157 & (3) & (N3T)\\
130 & 157 & (3) & note 18 (N3T)\\
163 & 157 & (3) & (N3T)\\
164 & 157 & (5) & Brunner [\textbf{1982a}],[\textbf{1983a}] (N3)\\
165 & 157 & (5) & (N3)\\
191 & 157 & (5) & note 18 (N3)\\
273 & 157 & (3) & note 18 (N3T)\\
305 & 157 & (3) & note 18 (N3T)\\
317 & 157 & (5) & (N3)\\
361 & 157 & (3) & note 18 (N3T)\\
363 & 157 & (3) & note 18 (N3T)\\
381 & 157 & (5) & Howard/Keremedis/Rubin/Rubin [\textbf{1997b}] (N3)\\
0 &  158 &(4) &Brunner [\textbf{1983b}](N25T)\\
\end{tabular}
\end{table}

\clearpage

\begin{table}[h]
\begin{tabular}{llll}
6 & 158 & (3) & note 18 (N25T)\\
11 & 158\index{Author}{Brunner}\index{Author}{H\"{o}ft/Howard}\index{Author}{Plotkin} & (3) & note 18 (N25T)\\
35 & 158 & (3) & note 18 (N25T)\\
37 & 158 & (3) & note 18 (N25T)\\
73 & 158 & (3) & note 18 (N25T)\\
8 & 163 & (3) & (N14T)\\
11 & 163 & (5) & (N14)\\
12 & 163 & (3) & (N14T)\\
43 & 163 & (5) & (N14)\\
79 & 163 & (3) & (N14T)\\
91 & 163 & (5) & (N14)\\
130 & 163 & (5) & (N14)\\
165 & 163 & (3) & (M2)\\
191 & 163 & (5) & (N14)\\
305 & 163 & (5) & (N14)\\
309 & 163 & (5) & (N14)\\
313 & 163 & (5) & (N14)\\
363 & 163 & (3) & (N14T)\\
11 & 164 & (5) & (N2)\\
12 & 164 & (3) & (N2T)\\
65 & 164 & (5) & (N2)\\
67 & 164 & (5) & (N2$(n)$)\\
79 & 164 & (3) & (N2T)\\
130 & 164 & (5) & (N2)\\
218 & 164 & (5) & (N6)\\
305 & 164 & (5) & (N2)\\
334 & 164 & (5) & (N2)\\
63 & 167 & (5) & ((N49))\\
165 & 167 & (5) & (N29)\\
82 & 171 & (3) & (N6T)\\
92 & 171 & (3) & note 18 (N6T)\\
119 & 171 & (3) & (N6T)\\
130 & 171 & (3) & note 18 (N6T)\\
185 & 171 & (3) & (N6T)\\
191 & 171 & (5) & note 18 (N6)\\
218 & 171 & (5) & (N6)\\
273 & 171 & (3) & note 18 (N6T)\\
305 & 171 & (3) & note 18 (N6T)\\
313 & 171 & (3) & note 18 (N6T)\\
350 & 171 & (3) & (N6T)\\
361 & 171 & (3) & note 18 (N6T)\\
363 & 171 & (3) & note 18 (N6T)\\
0 & 173 & (3) & Howard/Keremedis/Rubin/Stanley\index{Author}{Howard/Keremedis/Rubin/Stanley} [\textbf{1997}] (N54T)\\
6 & 173 & (3) &note 18 (N54T)\\
11 & 173 & (3) & note 18 (N54T)\\
35 & 173 & (3) & note 18 (N54T)\\
37 & 173 & (3) & note 18 (N54T)\\
73 & 173 & (3) & note 18 (N54T)\\
76 & 173 & (1) & Howard/Keremedis/Rubin/Stanley [\textbf{1997}]\\
\end{tabular}
\end{table}

\clearpage

\begin{table}[h]
\begin{tabular}{llll}
91 & 173 & (5) & note 18 (N54)\\
92 & 173\index{Author}{David}\index{Author}{Freyd}\index{Author}{Gitik}\index{Author}{Hickman}\index{Author}{Hodges} & (3) & note 18 (N54T)\\
130 & 173 & (3) & note 18 (N54T)\\
191 & 173 & (3) & note 18 (N54T)\\
273 & 173 & (3) & note 18 (N54T)\\
305 & 173 & (3) & note 18 (N54T)\\
309 & 173 & (3) & note 18 (N54T)\\
313 & 173 & (3) & note 18 (N54T)\\
361 & 173 & (3) & note 18 (N54T)\\
363 & 173 & (3) & note 18 (N54T)\\
383 & 173 & (1) & Howard/Keremedis/Rubin/Stanley [\textbf{1997}]\\
87 & 174 & (1) & H\"{o}ft/Howard \cite{bib:1981b}\\
0 &  176 &(3) &Plotkin \cite{bib:1981b} (M22)\\
11 & 177 & (3) & note 18 (N1T)\\
16 & 177 & (5) & (N1)\\
17 & 177 & (3) & Brunner [\textbf{1981a}], note 52 (N1T)\\
23 & 177 & (5) & (N1)\\
24 & 177 & (5) & (N1)\\
31 & 177 & (3) & (N1T)\\
37 & 177 & (3) & note 18 (N1T)\\
47 & 177 & (3) & (N1T)\\
48 & 177 & (3) & (N1T)\\
63 & 177 & (5) & (N1)\\
73 & 177 & (3) & note 18 (N1T)\\
89 & 177 & (5) & (N1)\\
92 & 177 & (3) & note 18 (N1T)\\
112 & 177 & (5) & (N1)\\
114 & 177 & (5) & (N1)\\
115 & 177 & (5) & (N1)\\
116 & 177 & (5) & (N1)\\
122 & 177 & (3) & (N1T)\\
127 & 177 & (5) & (N1)\\
130 & 177 & (3) & note 18 (N1T)\\
133 & 177 & (5) & (N1)\\
134 & 177 & (5) & (N1)\\
135 & 177 & (5) & (N1)\\
147 & 177 & (5) & (N1)\\
163 & 177 & (3) & (N1T)\\
165 & 177 & (3) & (N1T)\\
185 & 177 & (3) & (N1T)\\
191 & 177 & (5) & note 18 (N1)\\
217 & 177 & (3) & (N1T)\\
233 & 177 & (5) & (N1)\\
263 & 177 & (5) & (N1)\\
273 & 177 & (3) & note 18 (N1T)\\
304 & 177 & (5) & (N1)\\
305 & 177 & (3) & note 18 (N1T)\\
309 & 177 & (3) & note 18 (N1T)\\
313 & 177 & (3) & note 18 (N1T)\\
\end{tabular}
\end{table}

\clearpage

\begin{table}[h]
\begin{tabular}{llll}
322 & 177 & (5) & (N1)\\
325 & 177 & (3) & (N1T)\\
361 & 177 & (3) & note 18 (N1T)\\
363 & 177 & (3) & note 18 (N1T)\\
380 & 177 & (5) & (N1)\\
0 &  178 &(7) & This project\\
62 & 178 & (1) & clear\\
144 & 179 & (3) & David [\textbf{1980}] (M35$(\epsilon)$)\\
0 &  180 &(3)& Hodges [\textbf{1980}] (M23)\\
40 & 181 & (3) & Freyd [\textbf{1980}] (M25)\\
202 & 181 & (1) & clear\\
345 & 181 & (5) & (N40)\\
0 &  182 & (3) &Gitik [\textbf{1980}],[\textbf{1985}] (M17 IC)\\
104 & 182 & (1) & clear\\
191 & 182 & (1) & Blass\index{Author}{Blass} [\textbf{1979}], note 59\\
0 &  183 & (3) & Hickman\index{Author}{Hickman} [\textbf{1980b}], note 56 (N27T)\\
6 & 183 & (3) & note 18 (N27T)\\
11 & 183 & (3) & note 18 (N27T)\\
35 & 183 & (3) & note 18 (N27T)\\
37 & 183 & (3) & note 18 (N27T)\\
43 & 183 & (3) & (N27T)\\
63 & 183 & (5) & (N27)\\
71 & 183 & (1) & clear\\
73 & 183 & (3) & note 18 (N27T)\\
91 & 183 & (5) & note 18 (N27)\\
92 & 183 & (3) & note 18 (N27T)\\
130 & 183 & (3) & note 18 (N27T)\\
191 & 183 & (5) & note 18 (N27)\\
273 & 183 & (3) & note 18 (N27T)\\
305 & 183 & (3) & note 18 (N27T)\\
309 & 183 & (3) & note 18 (N27T)\\
313 & 183 & (3) & note 18 (N27T)\\
361 & 183 & (3) & note 18 (N27T)\\
363 & 183 & (3) & note 18 (N27T)\\
185 & 185 & (1) & Hickman [\textbf{1979a}]\\
9 & 185 & (1) & clear\\
77 & 185 & (1) & H\"{o}ft/Howard\index{Author}{H\"{o}ft/Howard} [\textbf{1994}]\\
126 & 185 & (1) & note 76\\
0 &  186 &(7) &This project\\
3 & 186 & (1) & J. Rubin\index{Author}{J. Rubin} [\textbf{1967}]\\
187 & 186 & (1) & J. Rubin [\textbf{1967}]\\
0 &  187 &(7) &This project\\
3 & 187 & (7) & Hickman [\textbf{1979b}]\\
186 & 187 & (7) & Hickman [\textbf{1979b}], note 58\\
193 & 188 & (1) & Blass [\textbf{1979}]\\
191 & 189 & (1) & Blass [\textbf{1979}]\\
0 &  190 &(3) &Blass [\textbf{1979}] (M24)\\
6 & 190 & (5) & note 18 (N28)\\
11 & 190 & (5) & note 18 (N28)\\
\end{tabular}
\end{table}

\clearpage

\begin{table}[h]
\begin{tabular}{llll}
35 & 190 & (5) & note 18 (N28)\\
37 & 190\index{Author}{Brunner}\index{Author}{Dawson/Howard} & (5) & note 18 (N28)\\
73 & 190 & (5) & note 18 (N28)\\
91 & 190 & (5) & note 18 (N28)\\
92 & 190 & (5) & note 18 (N28)\\
130 & 190 & (5) & note 18 (N28)\\
189 & 190 & (1) & Blass [\textbf{1979}]\\
273 & 190 & (5) & note 18 (N28)\\
305 & 190 & (5) & note 18 (N28)\\
309 & 190 & (5) & note 18 (N28)\\
313 & 190 & (5) & note 18 (N28)\\
361 & 190 & (5) & note 18 (N28)\\
363 & 190 & (5) & note 18 (N28)\\
40 & 192 & (5) & note 18 (N12$(\aleph_{2})$)\\
44 & 192 & (5) & note 18 (N12$(\aleph_{2})$)\\
87 & 192 & (5) & note 18 (N12$(\aleph_{\alpha})$)\\
187 & 192 & (5) & note 18 (N12$(\aleph_{1})$)\\
202 & 192 & (5) & note 18 (N12$(\aleph_{1})$)\\
316 & 192 & (5) & note 18 (N37)\\
345 & 192 & (5) & note 18 (N40)\\
40 & 193 & (5) & (N12$(\aleph_{1})$)\\
44 & 193 & (5) & (N12$(\aleph_{2})$)\\
71 & 193 & (5) & (see N12$(\aleph_{\alpha})$)\\
87 & 193 & (5) & (see N2$(\aleph_{\alpha})$)\\
187 & 193 & (5) & (N12$(\aleph_{1})$)\\
202 & 193 & (5) & (N12$(\aleph_{1})$)\\
345 & 193 & (5) & N40\\
0 &  194 & (3) & Kanovei\index{Author}{Kanovei} [\textbf{1976}] (M41)\\
16 & 194 & (1) & clear\\
94 & 194 & (1) & clear\\
0 &  196 &(7) & This project\\
86 & 196 & (1) & Bull\index{Author}{Bull} [\textbf{1978}]\\
0 &  197 &(7) &This project\\
79 & 197 & (1) & Davies\index{Author}{Davies} [\textbf{1978}]\\
0 &  198 &(4) &Hickman\index{Author}{Hickman} [\textbf{1978b}](N2T)\\
6 & 198 & (3) & note 18 (N2T)\\
9 & 198 & (1) & clear\\
11 & 198 & (3) & note 18 (N2T)\\
30 & 198 & (1) & clear\\
35 & 198 & (3) & note 18 (N2T)\\
37 & 198 & (3) & note 18 (N2T)\\
52 & 198 & (3) & (N2T)\\
65 & 198 & (5) & Hickman [\textbf{1978b}], Pincus\index{Author}{Pincus} [\textbf{1972b}] (N2)\\
67 & 198 & (5) & Felgner/Jech\index{Author}{Felgner/Jech} [\textbf{1973}], Hickman [\textbf{1978b}] (N2)\\
73 & 198 & (3) & note 18 (N2T)\\
82 & 198 & (3) & (N2T)\\
83 & 198 & (3) & (N2T)\\
92 & 198 & (3) & note 18 (N2T)\\
99 & 198 & (5) & (N2)\\
\end{tabular}
\end{table}

\clearpage

\begin{table}[h]
\begin{tabular}{llll}
119 & 198 & (3) & (N2T)\\
130 & 198\index{Author}{Blass}\index{Author}{Figura}\index{Author}{Harper/Rubin}\index{Author}{Howard/Keremedis/Rubin/Rubin}\index{Author}{Pelc}\index{Author}{Schrijver}\index{Author}{Truss} & (3) & note 18 (N2T)\\
163 & 198 & (3) & Brunner [\textbf{1982a}], Hickman [\textbf{1978b}] (N2T)\\
185 & 198 & (3) & (N2T)\\
191 & 198 & (5) & note 18 (N2)\\
215 & 198 & (1) & clear\\
273 & 198 & (3) & note 18 (N2T)\\
305 & 198 & (3) & note 18 (N2T)\\
313 & 198 & (3) & note 18 (N2T)\\
334 & 198 & (5) & (N2)\\
350 & 198 & (3) & (N2T)\\
361 & 198 & (3) & note 18 (N2T)\\
363 & 198 & (3) & note 18 (N2T)\\
0 &  199 & (3) & Kanovei [\textbf{1978}] (M26)\\
13 & 199 & (1) & clear\\
0 &  200 & (3) &Dawson/Howard\index{Author}{Howard} [\textbf{1976}], note 64 (N29T)\\
11 & 200 & (3) & note 18 (N1T)\\
14 & 200 & (3) & Dawson/Howard [\textbf{1976}], Pincus\index{Author}{Pineus/Solovay} [\textbf{1974b}], note 64 (N3T)\\
16&200&(5)&(N1)\\
17 & 200 & (3) & Dawson/Howard [\textbf{1976}], note 64 (N1T)\\
23 & 200 & (5) & (N3)\\
24 & 200 & (5) & (N1)\\
31 & 200 & (3) & (N1T)\\
37 & 200 & (3) & note 18 (N1T)\\
60 & 200 & (3) & (N3T)\\
73 & 200 & (3) & note 18 (N1T)\\
89 & 200 & (5) & (N1)\\
91 & 200 & (5) & note 18 (N3)\\
92 & 200 & (3) & note 18 (N1T)\\
112 & 200 & (5) & (N1)\\
114 & 200 & (5) & (N1)\\
115 & 200 & (5) & (N1)\\
116 & 200 & (5) & (N1)\\
128 & 200 & (3) & (N3T)\\
130 & 200 & (3) & note 18 (N1T)\\
133 & 200 & (5) & (N1)\\
134 & 200 & (5) & (N1)\\
135 & 200 & (5) & (N1)\\
147 & 200 & (5) & (N1)\\
163 & 200 & (3) & (N3T)\\
164 & 200 & (5) & (N3)\\
165 & 200 & (5) & (N3)\\
185 & 200 & (3) & (N1T)\\
191 & 200 & (5) & note 18 (N1, N3, N29)\\
217 & 200 & (3) & (N1T)\\
263 & 200 & (5) & (N1)\\
273 & 200 & (3) & note 18 (N1T)\\
304 & 200 & (5) & (N1)\\
305 & 200 & (3) & note 18 (N1T)\\
\end{tabular}
\end{table}

\clearpage

\begin{table}[h]
\begin{tabular}{llll}
317 & 200 & (5) & (N3)\\
322 & 200 & (5) & (N1)\\
325 & 200 & (3) & (N1T)\\
361 & 200 & (3) & note 18 (N1T)\\
363 & 200 & (3) & note 18 (N1T)\\
381 & 200 & (5) & Howard/Keremedis/Rubin/Rubin [\textbf{1997b}] (N3)\\
49 & 201 & (1) & Schrijver \cite{bib:1978j}\\
181 & 202 & (3) & Truss [\textbf{1978}] (M2$(<\omega_{2}>)$)\\
264 & 202 & (1) & Harper/Rubin [\textbf{1976}]\\
40 & 203 & (3) & Truss [\textbf{1978}], note 65 (M2)\\
79 & 203 & (1) & clear\\
181 & 203 & (3) & Truss [\textbf{1978}] (M2$(<\omega_{2}>)$)\\
204 & 203 & (3) & Truss [\textbf{1978}], note 65 (M2)\\
0 &  204 & (7) & This project\\
3 & 204 & (1) & clear\\
0 &  205 & (7) & This project\\
101 & 205 & (1) & Pelc [\textbf{1978}]\\
43 & 206 & (4) & Pincus/Solovay\index{Author}{Pincus} [\textbf{1977}], note 77 (M27)\\
52 & 206 & (3) & Pincus/Solovay [\textbf{1977}], note 77 (M27)\\
70 & 206 & (1) & clear\\
144 & 206 & (3) & Blass [\textbf{1977b}] (M15)\\
253 & 206 & (3) & Blass [\textbf{1977b}], Howard [\textbf{1975}] (M15)\\
0 & 207 & (4) & Figura\index{Author}{Figura} [\textbf{1977}] (M28)\\
23 & 207 & (1) & clear\\
0 & 208 & (3) & Figura [\textbf{1977}] (M$(\aleph)$5 IC)\\
40 & 208 & (1) & note 69\\
51 & 208 & (1) & Morris\index{Author}{Morris} [\textbf{1969}]\\
0 & 209 & (3) & Morris [\textbf{1970}] (M28)\\
26 & 209 & (1) & note 72\\
31 & 209 & (1) & note 72\\
207 & 209 & (1) & note 72\\
0 & 210 & (3) & L\"{a}uchli\index{Author}{L\"{a}uchli} [\textbf{1962}] Jech/Sochor\index{Author}{Jech/Sochor} [\textbf{1966a}] (N30T)\\
6 & 210 & (3) & note 18 (N30T)\\
11 & 210 & (3) & note 18 (N30T)\\
35 & 210 & (3) & note 18 (N30T)\\
37 & 210 & (3) & note 18 (N30T)\\
68 & 210 & (1) & clear\\
73 & 210 & (3) & note 18 (N30T)\\
91 & 210 & (5) & note 18 (N30)\\
92 & 210 & (3) & note 18 (N30T)\\
130 & 210 & (3) & note 18 (N30T)\\
191 & 210 &(5) & note 18 (N30)\\
273 & 210 & (3) & note 18 (N30T)\\
305 & 210 & (3) & note 18 (N30T)\\
309 & 210 & (3) & note 18 (N30T)\\
313 & 210 & (3) & note 18 (N30T)\\
361 & 210 & (3) & note 18 (N30T)\\
363 & 210 & (3) & note 18 (N30T)\\
43 & 211 & (1) & clear\\
\end{tabular}
\end{table}

\clearpage

\begin{table}[h]
\begin{tabular}{llll}
212 & 211 & (1) & clear\\
79 & 212\index{Author}{Solovay}\index{Author}{Stern}\index{Author}{Truss}\index{Author}{Wagon} & (1) & clear\\
15 & 213 & (3) & Pincus\index{Author}{Pincus}\index{Author}{Pineus/Solovay} [\textbf{1977a}] (M34)\\
60 & 213 & (1) & clear\\
85 & 213 & (3) & Pincus [\textbf{1977a}], note 74 (M34)\\
44 & 214 & (3) & (N12$(\aleph_{2})$T)\\
83 & 215 & (3) & Pincus [\textbf{1972a}], note 105 (N2T)\\
0 &  216 &(7) &This project\\
10 & 216 & (1) & (see [10 F])\\
217 & 216 & (1) & clear\\
6 & 217 & (3) & note 18 (N34T)\\
9 & 217 & (1) & clear\\
11 & 217 & (3) & note 18 (N34T)\\
35 & 217 & (3) & note 18 (N34T)\\
37 & 217 & (3) & note 18 (N34T)\\
73 & 217 & (3) & note 18 (N34T)\\
91 & 217 & (5) & note 18 (N34)\\
92 & 217 & (3) & note 18 (N34T)\\
130 & 217 & (3) & note 18 (N34T)\\
191 & 217 & (5) & note18 ($A$ is a set) (N34)\\
216 & 217 & (3) & Pincus [\textbf{1972a}], note 105 (N34T)\\
273 & 217 & (3) & note 18 (N34T)\\
305 & 217 & (3) & note 18 (N34T)\\
309 & 217 & (3) & note 18 (N34T)\\
313 & 217 & (3) & note 18 (N34T)\\
361 & 217 & (3) & note 18 (N34T)\\
363 & 217 & (3) & note 18 (N34T)\\
109 & 218 & (1) & H.\,Rubin/J. Rubin\index{Author}{H. Rubin/ J. Rubin} \cite{bib:1985q} p 120 (see [218 A])\\
0 &  219 &(7) & This project\\
0 &  220 &(3) & Hickman\index{Author}{Hickman} [\textbf{1977b}] (N42$(p)$T, N45$(p)$T)\\
6 & 220 & (3) & note 18 (N42$(p)$T, N45$(p)$T)\\
11 & 220 & (3) & note 18 (N42$(p)$T, N45$(p)$T)\\
35 & 220 & (3) & note 18 (N42$(p)$T, N45$(p)$T)\\
37 & 220 & (3) & note 18 (N42$(p)$T, N45$(p)$T)\\
73 & 220 & (3) & note 18 (N42$(p)$T, N45$(p)$T)\\
91 & 220 & (5) & note 18 (N42$(p)$, N45$(p)$)\\
92 & 220 & (3) & note 18 (N42$(p)$T, N45$(p)$T)\\
130 & 220 & (3) & note 18 (N42$(p)$T, N45$(p)$T)\\
191 & 220 & (5) & note 18 (N42$(p)$, N45$(p)$)\\
273 & 220 & (3) & note 18 (N42$(p)$T, N45$(p)$T)\\
305 & 220 & (3) & note 18 (N42$(p)$T, N45$(p)$T)\\
309 & 220 & (3) & note 18 (N42$(p)$T, N45$(p)$T)\\
313 & 220 & (3) & note 18 (N42$(p)$T, N45$(p)$T)\\
361 & 220 & (3) & note 18 (N42$(p)$T, N45$(p)$T)\\
363 & 220 & (3) & note 18 (N42$(p)$T, N45$(p)$T)\\
11 & 221 & (3) & note 18 (N51T)\\
52 & 221 & (1) & clear (see [52 A])\\
63 & 221 & (1) & clear\\
79 & 221 & (3) & (N51T)\\
\end{tabular}
\end{table}

\clearpage

\begin{table}[h]
\begin{tabular}{llll}
90 & 221 & (5) & (N51)\\
130 & 221 & (3) & note 18 (N51T)\\
191 & 221 & (5) & note 18 (N51)\\
305 & 221 & (3) & note 18 (N51T)\\
309 & 221 & (3) & note 18 (N51T)\\
313 & 221 & (3) & note 18 (N51T)\\
40 & 222 & (3) & Pincus \cite{bib:1972c1}, Truss [\textbf{1978}] (M2)\\
70 & 222 & (1) & Pincus \cite{bib:1972c1}\\
142 & 222 & (3) & Pincus \cite{bib:1972c1} (M2)\\
165 & 222 & (3) & (M2)\\
204 & 222 & (3) & (M2)\\
221 & 222 & (1) & clear\\
313 & 222 & (1) & Wagon \cite{bib:1985u}\\
43 & 223 & (3) & Pincus/Solovay [\textbf{1977}] (M30)\\
206 & 223 & (1) & clear\\
222 & 223 & (1) & clear\\
287 & 222 & (1) & (due to Morillon), note 145\\
253 & 223 & (3) & (M7)\\
312 & 223 & (3) & (M30)\\
0 &  224 & (3) & Stern [\textbf{1979}] (M$(\aleph)$5 IC)\\
43 & 224 & (3) & Solovay [\textbf{1970}], Stern [\textbf{1979}] (M$(\aleph)$5 IC)\\
307 & 224 & (3) & (M$(\aleph)$5 IC)\\
99 & 225 & (1) & Rav\index{Author}{Rav} [\textbf{1977}]\\
0 &  226 & (7) &This project\\
14 & 226 & (1) & Rav [\textbf{1977}]\\
0 &  227 & (7) &This project\\
14 & 227 & (1) & Rav [\textbf{1977}]\\
0 &  228 &(7) &This project\\
14 & 228 & (1) & Rav [\textbf{1977}]\\
0 &  229 &(7) &This project\\
14 & 229 & (1) & Rav [\textbf{1977}]\\
0 &  230 &(3)& Szczepaniak\index{Author}{Szczepaniak} [\textbf{1977}] (M31)\\
8 & 231 & (3) & Felgner\index{Author}{Felgner} [\textbf{1974}] (M21)\\
40 & 231 & (1) & Felgner [\textbf{1974}]\\
60 & 231 & (1) & Tsukada\index{Author}{Tsukada} [\textbf{1977}]\\
133 & 231 & (1) & note 123\\
67 & 232 & (1) & Howard/Keremedis/Rubin/Stanley\index{Author}{Howard/Keremedis/Rubin/Stanley} [\textbf{1997}]\\
14 & 233 & (1) & Banaschewski\index{Author}{Banaschewski} [\textbf{1992}]\\
0 &  234 &(3)& Prikry\index{Author}{Prikry} [\textbf{1976}], Mathias\index{Author}{Mathias} [\textbf{1968}] (M$(\aleph)$5 IC)\\
43 & 234 & (3) & (M$(\aleph)$5 IC)\\
307 & 234 & (3) & (M$(\aleph)$5 IC)\\
0 &  235 & (3) & Jech\index{Author}{Jech} \cite{bib:1973b} p 149  prob 10.5, L\"{a}uchli\index{Author}{L\"{a}uchli} [\textbf{1962}], Jech/Sochor\index{Author}{Jech/Sochor}\\
&&&[\textbf{1966a}](N13T)\\
6 & 235 & (3) & note 18 (N13T)\\
11 & 235 & (3) & note 18 (N13T)\\
35 & 235 & (3) & note 18 (N13T)\\
37 & 235 & (3) & note 18 (N13T)\\
73 & 235 & (3) & note 18 (N13T)\\
91 & 235 & (5) & note 18 (N13)\\
\end{tabular}
\end{table}

\clearpage

\begin{table}[h]
\begin{tabular}{llll}
92 & 235 & (3) & note 18 (N13T)\\
96 & 235 & (1) & clear\\
130 & 235 & (3) & note 18 (N13T)\\
191 & 235 & (5) & note 18 (N13)\\
273 & 235 & (3) & note 18 (N13T)\\
305 & 235 & (3) & note 18 (N13T)\\
309 & 235 & (3) & note 18 (N13T)\\
313 & 235 & (3) & note 18 (N13T)\\
361 & 235 & (3) & note 18 (N13T)\\
363 & 235 & (3) & note 18 (N13T)\\
0 &  236 & (3) & Gross\index{Author}{Gross} [\textbf{1975}], [\textbf{1976}] (N44T)\\
6 & 236 & (3) & note 18 (N44T)\\
11 & 236 & (3) & note 18 (N44T)\\
35 & 236 & (3) & note 18 (N44T)\\
37 & 236 & (3) & note 18 (N44T)\\
73 & 236 & (3) & note 18 (N44T)\\
91 & 236 & (5) & note 18 (N44)\\
92 & 236 & (3) & note 18 (N44T)\\
109 & 236 & (1) & clear\\
130 & 236 & (3) & note 18 (N44T)\\
191 & 236 & (5) & note 18 (N44)\\
273 & 236 & (3) & note 18 (N44T)\\
305 & 236 & (3) & note 18 (N44T)\\
309 & 236 & (3) & note 18 (N44T)\\
313 & 236 & (3) & note 18 (N44T)\\
361 & 236 & (3) & note 18 (N44T)\\
363 & 236 & (3) & note 18 (N44T)\\
0 &  237 &(3)& Hickman\index{Author}{Hickman} [\textbf{1976}] (N32T)\\
6 & 237 & (3) & note 18 (N32T)\\
11 & 237 & (3) & note 18 (N32T)\\
35 & 237 & (3) & note 18 (N32T)\\
37 & 237 & (3) & note 18 (N32T)\\
73 & 237 & (3) & note 18 (N32T)\\
91 & 237 & (5) & note 18 (N32)\\
92 & 237 & (3) & note 18 (N32T)\\
130 & 237 & (3) & note 18 (N32T)\\
191 & 237 & (5) & note 18 (N32)\\
273 & 237 & (3) & note 18 (N32T)\\
305 & 237 & (3) & note 18 (N32T)\\
309 & 237 & (3) & note 18 (N32T)\\
313 & 237 & (3) & note 18 (N32T)\\
361 & 237 & (3) & note 18 (N32T)\\
363 & 237 & (3) & note 18 (N32T)\\
0 &  238 & (3) & Hickman [\textbf{1976}] (N32T)\\
6 & 238 & (3) & note 18 (N32T)\\
11 & 238 & (3) & note 18 (N32T)\\
35 & 238 & (3) & note 18 (N32T)\\
37 & 238 & (3) & note 18 (N32T)\\
73 & 238 & (3) & note 18 (N32T)\\
\end{tabular}
\end{table}

\clearpage

\begin{table}[h]
\begin{tabular}{llll}
91 & 238 & (5) & note 18 (N32)\\
92 & 238\index{Author}{L\"{a}uchli} & (3) & note 18 (N32T)\\
130 & 238 & (3) & note 18 (N32T)\\
191 & 238 & (5) & note 18 (N32)\\
273 & 238 & (3) & note 18 (N32T)\\
305 & 238 & (3) & note 18 (N32T)\\
309 & 238 & (3) & note 18 (N32T)\\
313 & 238 & (3) & note 18 (N32T)\\
361 & 238 & (3) & note 18 (N32T)\\
363 & 238 & (3) & note 18 (N32T)\\
0 &  239 & (4) & Hickman\index{Author}{Hickman} [\textbf{1976}] (N1T)\\
11 & 239 & (3) & note 18 (N1T)\\
16 & 239 & (5) & (N1)\\
17 & 239 & (3) & (N1T)\\
23 & 239 & (5) & (N1)\\
24 & 239 & (5) & (N1)\\
31 & 239 & (3) & (N1T)\\
37 & 239 & (3) & note 18 (N1T)\\
47 & 239 & (3) & (N1T)\\
48 & 239 & (3) & (N1T)\\
63 & 239 & (5) & (N1)\\
73 & 239 & (3) & note 18 (N1T)\\
89 & 239 & (5) & (N1)\\
92 & 239 & (3) & note 18 (N1T)\\
112 & 239 & (5) & (N1)\\
114 & 239 & (5) & (N1)\\
115 & 239 & (5) & (N1)\\
116 & 239 & (5) & (N1)\\
122 & 239 & (3) & (N1T)\\
127 & 239 & (5) & (N1)\\
130 & 239 & (3) & note 18 (N1T)\\
133 & 239 & (5) & (N1)\\
134 & 239 & (5) & (N1)\\
135 & 239 & (5) & (N1)\\
147 & 239 & (5) & (N1)\\
163 & 239 & (3) & (N1T)\\
165 & 239 & (3) & (N1T)\\
185 & 239 & (3) & (N1T)\\
191 & 239 & (5) & note 18 (N1)\\
217 & 239 & (3) & (N1T)\\
233 & 239 & (5) & (N1)\\
263 & 239 & (5) & (N1)\\
273 & 239 & (3) & note 18 (N1T)\\
304 & 239 & (5) & (N1)\\
305 & 239 & (3) & note 18 (N1T)\\
309 & 239 & (3) & note 18 (N1T)\\
313 & 239 & (3) & note 18 (N1T)\\
322 & 239 & (5) & (N1)\\
325 & 239 & (3) & (N1T)\\
\end{tabular}
\end{table}

\clearpage

\begin{table}[h]
\begin{tabular}{llll}
361 & 239 & (3) & note 18 (N1T)\\
363 & 239\index{Author}{Fraleigh}\index{Author}{Hodges}\index{Author}{Litman}\index{Author}{Plotkin} & (3) & note 18 (N1T)\\
380 & 239 & (5) & (N1)\\
0 & 240 & (3) & Hickman [\textbf{1976}] (N32T)\\
6 & 240 & (3) & note 18 (N32T)\\
11 & 240 & (3) & note 18 (N32T)\\
35 & 240 & (3) & note 18 (N32T)\\
37 & 240 & (3) & note 18 (N32T)\\
73 & 240 & (3) & note 18 (N32T)\\
91 & 240 & (5) & note 18 (N32)\\
92 & 240 & (3) & note 18 (N32T)\\
130 & 240 & (3) & note 18 (N32T)\\
191 & 240 & (5) & note 18 (N32)\\
273 & 240 & (3) & note 18 (N32T)\\
305 & 240 & (3) & note 18 (N32T)\\
309 & 240 & (3) & note 18 (N32T)\\
313 & 240 & (3) & note 18 (N32T)\\
361 & 240 & (3) & note 18 (N32T)\\
363 & 240 & (3) & note 18 (N32T)\\
0 &  241&(3)&L\"{a}uchli [\textbf{1962}] (N31T)\\
6 & 241 & (3) & note 18 (N31T)\\
11 & 241 & (3) & note 18 (N31T)\\
35 & 241 & (3) & note 18 (N31T)\\
37 & 241 & (3) & note 18 (N31T)\\
73 & 241 & (3) & note 18 (N31T)\\
91 & 241 & (5) & note 18 (N31)\\
92 & 241 & (3) & note 18 (N31T)\\
130 & 241 & (3) & note 18 (N31T)\\
191 & 241 & (5) & note 18 (N31)\\
242 & 241 & (1) & clear\\
273 & 241 & (3) & note 18 (N31T)\\
305 & 241 & (3) & note 18 (N31T)\\
309 & 241 & (3) & note 18 (N31T)\\
313 & 241 & (3) & note 18 (N31T)\\
361 & 241 & (3) & note 18 (N31T)\\
363 & 241 & (3) & note 18 (N31T)\\
233 & 242 & (1) & clear\\
0 & 243 & (3) & Hodges [\textbf{1976a}] (N31T)\\
6 & 243 & (3) & note 18 (N31T)\\
11 & 243 & (3) & note 18 (N31T)\\
35 & 243 & (3) & note 18 (N31T)\\
37 & 243 & (3) & note 18 (N31T)\\
43 & 243 & (1) & Fraleigh [\textbf{1989}] pp 333-336\\
73 & 243 & (3) & note 18 (N31T)\\
91 & 243 & (5) & note 18 (N31)\\
92 & 243 & (3) & note 18 (N31T)\\
130 & 243 & (3) & note 18 (N31T)\\
191 & 243 & (5) & note 18 (N31)\\
273 & 243 & (3) & note 18 (N31T)\\
\end{tabular}
\end{table}

\clearpage

\begin{table}[h]
\begin{tabular}{llll}
305 & 243 & (3) & note 18 (N31T)\\
309 & 243\index{Author}{Ash}\index{Author}{Blass}\index{Author}{Brunner}\index{Author}{Haddad/Morillon}\index{Author}{Halpern/Levy}\index{Author}{Harper/Rubin}\index{Author}{Howard/Keremedis/Rubin/Rubin}\index{Author}{Jech} & (3) & note 18 (N31T)\\
313 & 243 & (3) & note 18 (N31T)\\
361 & 243 & (3) & note 18 (N31T)\\
363 & 243 & (3) & note 18 (N31T)\\
0 &  244&(3)&Hodges [\textbf{1976a}],\,(N31T)\\
6 & 244 & (3) & note 18 (N31T)\\
11 & 244 & (3) & note 18 (N31T)\\
35 & 244 & (3) & note 18 (N31T)\\
37 & 244 & (3) & note 18 (N31T)\\
73 & 244 & (3) & note 18 (N31T)\\
91 & 244 & (5) & note 18 (N31)\\
92 & 244 & (3) & note 18 (N31T)\\
130 & 244 & (3) & note 18 (N31T)\\
191 & 244 & (5) & note 18 (N31)\\
273 & 244 & (3) & note 18 (N31T)\\
278 & 244 & (1) & clear\\
305 & 244 & (3) & note 18 (N31T)\\
309 & 244 & (3) & note 18 (N31T)\\
313 & 244 & (3) & note 18 (N31T)\\
361 & 244 & (3) & note 18 (N31T)\\
363 & 244 & (3) & note 18 (N31T)\\
0 &  246 &(7)& This project\\
245 & 246 & (1) & Litman [\textbf{1976}]\\
0 &  247 & (3) & Plotkin\index{Author}{Plotkin} [\textbf{1976}], note 86 (M33)\\
0 &  248 & (3) & Plotkin \cite{bib:1976p}, note 86 (M33) \\
6 & 249 & (3) & note 18 (N35T)\\
10 & 249 & (1) & (see [10 F])\\
11 & 249 & (3) & note 18 (N35T)\\
35 & 249 & (3) & note 18 (N35T)\\
37 & 249 & (3) & note 18 (N35T)\\
73 & 249 & (3) & note 18 (N35T)\\
91 & 249 & (5) & note 18 (N35)\\
92 & 249 & (3) & note 18 (N35T)\\
130 & 249 & (3) & note 18 (N35T)\\
191 & 249 & (5) & note 18 (N35)\\
250 & 249 & (3) & Truss\index{Author}{Truss} [\textbf{1975}] (N35T)\\
273 & 249 & (3) & note 18 (N35T)\\
305 & 249 & (3) & note 18 (N35T)\\
309 & 249 & (3) & note 18 (N35T)\\
313 & 249 & (3) & note 18 (N35T)\\
361 & 249 & (3) & note 18 (N35T)\\
363 & 249 & (3) & note 18 (N35T)\\
40 & 250 & (2) & clear\\
61 & 250 & (1) & clear\\
88 & 250 & (3) & $(\mathrm{N}2^{\ast}(3)\mathrm{T})$\\
122 & 250 & (1) & clear\\
141 & 250 & (5) & Howard\index{Author}{Howard} [\textbf{1984a}] $(\mathrm{N}2^{\ast}(3))$\\
327 & 250 & (1) & clear\\
\end{tabular}
\end{table}

\clearpage

\begin{table}[h]
\begin{tabular}{llll}
333 & 250 & (5) & (N2$^{\ast}$(3))\\
79 & 251\index{Author}{Howard/H. Rubin / J . Rubin} & (1) & Ash [\textbf{1975}]\\
79 & 252 & (1) & Ash [\textbf{1975}]\\
14 & 253 & (3) & Howard [\textbf{1975}] (M1)\\
15 & 253 & (3) & Halpern/Levy [\textbf{1964}],[\textbf{1971}], Howard [\textbf{1975}] (M1)\\
23 & 253 & (5) & note 120(45) (N3)\\
31 & 253 & (3) & Howard [\textbf{1975}], Jech \cite{bib:1973b} (M1)\\
60 & 253 & (3) & Howard [\textbf{1975}], Jech \cite{bib:1973b} (M1)\\
79 & 253 & (3) & (N3T)\\
82 & 253 & (3) & (M1)\\
87 & 253 & (3) & note 120(45) (M40$(\kappa)$\\
118 & 253 & (3) & Haddad/Morillon [\textbf{1990}], Howard [\textbf{1975}] (M1)\\
128 & 253 & (3) & Brunner [\textbf{1984f}], Howard [\textbf{1975}] (M1)\\
163 & 253 & (3) & Brunner [\textbf{1982a}], Howard [\textbf{1975}] (M1)\\
164 & 253 & (5) & note 120(45) (N3)\\
165 & 253 & (3) & (M1)\\
191 & 253 & (3) & Blass [\textbf{1979}], Howard [\textbf{1975}] (M1)\\
277 & 253 & (3) & (M1)\\
317 & 253 & (5) & note 120(45) (N3)\\
345 & 253 & (3) & note 120(45) (M40$(\kappa)$\\
381 & 253 & (5) & Howard/Keremedis/Rubin/Rubin [\textbf{1997b}] (N3)\\
0 &  254 & (7) & This project\\
89 & 254 & (1) & Harper/Rubin\index{Author}{Harper/Rubin} [\textbf{1976}]\\
256 & 255 & (1) & Harper/Rubin [\textbf{1976}]\\
258 & 255 & (1) & Harper/Rubin [\textbf{1976}]\\
262 & 255 & (1) & Harper/Rubin [\textbf{1976}]\\
261 & 256 & (1) & Harper/Rubin [\textbf{1976}]\\
262 & 257 & (1) & Harper/Rubin [\textbf{1976}]\\
256 & 259 & (1) & Harper/Rubin [\textbf{1976}]\\
40 & 260 & (3) & Howard/H.Rubin/J.Rubin [\textbf{1973}] (N33T)\\
202 & 260 & (1) & Harper/Rubin [\textbf{1976}]\\
255 & 260 & (1) & Harper/Rubin [\textbf{1976}]\\
257 & 260 & (1) & Harper/Rubin [\textbf{1976}]\\
259 & 260 & (1) & Harper/Rubin [\textbf{1976}]\\
261 & 262 & (1) & Harper/Rubin [\textbf{1976}]\\
0 &  263 & (7) & This project\\
143 & 263 & (1) & Harper/Rubin [\textbf{1976}]\\
264 & 263 & (1) & Harper/Rubin [\textbf{1976}]\\
266 & 263 & (1) & Harper/Rubin [\textbf{1976}]\\
263 & 264 & (5) & Harper/Rubin [\textbf{1976}] (N1)\\
0 &  265 & (7) & This project\\
0 &  266 & (7) & This project\\
265 & 266 & (1) & Harper/Rubin [\textbf{1976}]\\
0 &  267 &(6) & Stavi\index{Author}{Stavi} \cite{bib:1975k} (N1)\\
11 & 267 & (5) & note 18 (N1)\\
16 & 267 & (5) & (N1)\\
17 & 267 & (5) & (N1)\\
23 & 267 & (5) & (N1)\\
24 & 267 & (5) & (N1)\\
\end{tabular}
\end{table}

\clearpage

\begin{table}[h]
\begin{tabular}{llll}
31 & 267 & (5) & (N1)\\
37 & 267\index{Author}{Hodges}\index{Author}{Howard/Spi\u{s}iak}\index{Author}{Monro} & (5) & note 18 (N1)\\
63 & 267 & (5) & (N1)\\
73 & 267 & (5) & note 18 (N1)\\
89 & 267 & (5) & (N1)\\
92 & 267 & (5) & note 18 (N1)\\
112 & 267 & (5) & (N1)\\
114 & 267 & (5) & (N1)\\
115 & 267 & (5) & (N1)\\
116 & 267 & (5) & (N1)\\
127 & 267 & (5) & (N1)\\
130 & 267 & (5) & note 18 (N1)\\
133 & 267 & (5) & (N1)\\
134 & 267 & (5) & (N1)\\
135 & 267 & (5) & (N1)\\
147 & 267 & (5) & (N1)\\
191 & 267 & (5) & note 18 (N1)\\
217 & 267 & (5) & (N1)\\
233 & 267 & (5) & (N1)\\
263 & 267 & (5) & (N1)\\
273 & 267 & (5) & note 18 (N1)\\
304 & 267 & (5) & (N1)\\
305 & 267 & (5) & note 18 (N1)\\
309 & 267 & (5) & note 18 (N1)\\
313 & 267 & (5) & note 18 (N1)\\
322 & 267 & (5) & (N1)\\
325 & 267 & (5) & (N1)\\
361 & 267 & (5) & note 18 (N1)\\
363 & 267 & (5) & note 18 (N1)\\
380 & 267 & (5) & (N1)\\
88 & 268 & (1) & Lampe\index{Author}{Lampe} \cite{bib:1974m}\\
0 &  269 & (5) & note 91 (N1)\\
268 & 269 & (1) & Lampe \cite{bib:1974m}\\
14 & 270 & (1) & Kolany/Wojtylak\index{Author}{Kolany/Wojtylak} [\textbf{1991}]\\
62 & 270 & (7) & Kolany/Wojtylak [\textbf{1991}]\\
85 & 270 & (1) & Kolany/Wojtylak [\textbf{1991}]\\
0 &  271 & (7) &This project\\
270 & 271 & (1) & Kolany/Wojtylak [\textbf{1991}]\\
79 & 272 & (1) & Truss\index{Author}{Truss} [\textbf{1974b}]\\
6 & 273 & (3) & Truss [\textbf{1974b}](M12$(\aleph)$)\\
231 & 273 & (3) & Truss [\textbf{1974b}](M12$(\aleph)$)\\
364 & 273 & (1) & clear\\
0 &  274 & (3) & Truss [\textbf{1974b}] (M2)\\
40 & 274 & (3) & (M2)\\
142 & 274 & (3) & Pincus\index{Author}{Pincus} \cite{bib:1972c1}, Truss [\textbf{1974b}] (M2)\\
165 & 274 & (3) & (M2)\\
204 & 274 & (3) & (M2)\\
0 &  275 & (3) & Truss [\textbf{1974b}] (M10)\\
0 &  276 & (7) & This project\\
\end{tabular}
\end{table}

\clearpage

\begin{table}[h]
\begin{tabular}{llll}
88 & 276 & (1) & Howard/Spi\u{s}iak [\textbf{1994}]\\
0 & 277\index{Author}{Makowski/Wi\'{s}niewski}\index{Author}{Truss}\index{Author}{Wi\'{s}niewski} & (3) & Monro [\textbf{1973a}] (M$(\aleph)$5 IC)\\
43 & 277 & (3) & (M$(\aleph)$5 IC)\\
307 & 277 & (3) & (M$(\aleph)$5 IC)\\
0 &  278 & (3) & Hodges [\textbf{1974}] (N1T)\\
16 & 278 & (5) & (N1)\\
17 & 278 & (3) & (N1T)\\
23 & 278 & (5) & (N1)\\
24 & 278 & (5) & (N1)\\
31 & 278 & (3) & (N1T)\\
47 & 278 & (3) & (N1T)\\
48 & 278 & (3) & (N1T)\\
63 & 278 & (5) & (N1)\\
89 & 278 & (5) & (N1)\\
112 & 278 & (5) & (N1)\\
114 & 278 & (5) & (N1)\\
115 & 278 & (5) & (N1)\\
116 & 278 & (5) & (N1)\\
122 & 278 & (3) & (N1T)\\
127 & 278 & (5) & (N1)\\
133 & 278 & (5) & (N1)\\
134 & 278 & (5) & (N1)\\
135 & 278 & (5) & (N1)\\
147 & 278 & (5) & (N1)\\
163 & 278 & (3) & (N1T)\\
165 & 278 & (3) & (N1T)\\
185 & 278 & (3) & (N1T)\\
217 & 278 & (3) & (N1T)\\
233 & 278 & (5) & (N1)\\
263 & 278 & (5) & (N1)\\
304 & 278 & (5) & (N1)\\
322 & 278 & (5) & (N1)\\
325 & 278 & (3) & (N1T)\\
361 & 278 & (3) & note 18 (N1T)\\
363 & 278 & (3) & note 18 (N1T)\\
380 & 278 & (5) & (N1)\\
0 &  279 & (7) & This project\\
43 & 279 & (1) & Wright\index{Author}{Wright} [\textbf{1973}]\\
0 & 280 & (3) & Wright [\textbf{1973}], Solovay\index{Author}{Solovay} [\textbf{1970}] (M$(\aleph)$5 IC)\\
43 & 280 & (3) & (M$(\aleph)$5 IC)\\
142 & 280 & (1) & clear\\
307 & 280 & (3) & (M$(\aleph)$5 IC)\\
0 &  281 & (3) & Wright [\textbf{1973}], note 96 (M$(\aleph)$5 IC)\\
43 & 281 & (3) & (M$(\aleph)$5 IC)\\
307 & 281 & (3) & (M$(\aleph)$5 IC)\\
0 &  282 & (7) & This project\\
8 & 282 & (1) & Kleinberg/Seiferas\index{Author}{Kleinberg/Seiferas} [\textbf{1973}], note 97\\
234 & 282 & (1) & note 142\\
0 &  283 & (7) &This project\\
\end{tabular}
\end{table}

\clearpage

\begin{table}[h]
\begin{tabular}{llll}
62 & 283 & (1) & Truss [\textbf{1973d}]\\
6 & 285\index{Author}{Keremedis}\index{Author}{Pincus}\index{Author}{Smullyan} & (3) & note 18 (N22($K,p$)T)\\
11 & 285 & (3) & note 18 (N22($K,p$)T)\\
35 & 285 & (3) & note 18 (N22($K,p$)T)\\
37 & 285 & (3) & note 18 (N22($K,p$)T)\\
45 & 285 & (3) & Makowski/Wi\'{s}niewski [\textbf{1969}] (N22($\{n\}$)T)\\ 46 & 285 &(3) & Makowski/Wi\'{s}niewski [\textbf{1969}] (N22($K,p$)T) \\
62 & 285 & (1) & Wi\'{s}niewski [\textbf{1973}]\\
64 & 285 & (3) & (N22($\{2\}$)T)\\
73 & 285 & (3) & note 18 (N22($K,p$)T)\\
80 & 285 & (3) & (N22($\{2\}$)T)\\
88 & 285 & (3) & Makowski/Wi\'{s}niewski [\textbf{1969}] (N22($\{2\}$)T)\\
91 & 285 & (5) & note 18 (N22($K,p$))\\
92 & 285 & (3) & note 18 (N22($K,p$)T)\\
130 & 285 & (3) & note 18 (N22($K,p$)T)\\
191 & 285 & (5) & note 18 (N22($K,p$))\\
273 & 285 & (3) & note 18 (N22($K,p$)T)\\
305 & 285 & (3) & note 18 (N22($K,p$))\\
308 & 285 & (3) & (N22($K,p$)T)\\
309 & 285 & (3) & note 18 (N22($K,p$)T)\\
313 & 285 & (3) & note 18 (N22($K,p$)T)\\
361 & 285 & (3) & note 18 (N22($K,p$) T)\\
363 & 285 & (3) & note 18 (N22($K,p$)T)\\
52 & 286 & (4) & note 99 (M1)\\
253 & 286 & (3) & note 120(26) (M27)\\
345 & 286 & (3) & note 120(26) (M40($\kappa)$)\\
52 & 287 & (1) & clear\\
6 & 288 & (3) & note 18 (N22($K,p$)T)\\
10 & 288 & (1) & clear\\
11 & 288 & (3) & note 18 (N22($K,p$)T)\\
33 & 288 & (2) & note 15\\
35 & 288 & (3) & note 18 (N22($K,p$)T)\\
37 & 288 & (3) & note 18 (N22($K,p$)T)\\
45 & 288 & (2) & note 15\\
47 & 288 & (1) & clear\\
52 & 288 & (3) & (N2$^{\ast}$(3)T)\\
73 & 288 & (3) & note 18 (N22($K,p$)T)\\
82 & 288 & (3) & (N2$^{\ast}$(3)T)\\
88 & 288 & (3) & (N2$^{\ast}$(3)T)\\
91 & 288 & (5) & note18 (N122($K,p$)\\
92 & 288 & (3) & note 18 (N22($K,p$)T)\\
119 & 288 & (3) & (N2$^{\ast}$(3)T)\\
130 & 288 & (3) & note 18 (N22($K,p$)T)\\
141 & 288 & (5) & (N2$^{\ast}$(3))\\
185 & 288 & (3) & (N2$^{\ast}$(3)T)\\
191 & 288 & (5) & note 18 (N22($K,p$))\\
273 & 288 & (3) & note 18 (N22($K,p$)T)\\
305 & 288 & (3) & note 18 (N22($K,p$)T)\\
308 & 288 & (5) & Howard/Yorke\index{Author}{Howard/Yorke} [\textbf{1987}] (N22($K,p$))\\
\end{tabular}
\end{table}

\clearpage

\begin{table}[h]
\begin{tabular}{llll}
309 & 288 & (3) & note 18 (N22($K,p$)T)\\
313 & 288 & (3) & note 18 (N22($K,p$)T)\\
333 & 288 & (5) & (N2$^{\ast}$(3))\\
335 & 288 & (1) & Keremedis \cite{bib:1996c1}\\
361 & 288 & (3) & note 18 (N2$^{\ast}$(3)T)\\
363 & 288 & (3) & note 18 (N2$^{\ast}$(3)T)\\
374 & 288 & (1) & clear\\
0 &  289 & (4) & Smullyan [\textbf{1994}], note 102 (M1)\\
14 & 289 & (3) & (M1)\\
15 & 289 & (3) & (M1)\\
31 & 289 & (3) & (M1)\\
60 & 289 & (3) & (M1)\\
79 & 289 & (1) & clear\\
82 & 289 & (3) & (M1)\\
118 & 289 & (3) & (M1)\\
128 & 289 & (3) & (M1)\\
163 & 289 & (3) & (M1)\\
165 & 289 & (3) & (M1)\\
191 & 289 & (3) & (M1)\\
277 & 289 & (3) & (M1)\\
11 & 290 & (3) & note 18 (N3T)\\
14 & 290 & (3) & Pincus [\textbf{1974b}] (N3T)\\
16 & 290 & (5) & (N3)\\
23 & 290 & (5) & (N3)\\
24 & 290 & (5) & (N3)\\
31 & 290 & (3) & (N3T)\\
37 & 290 & (3) & note 18 (N3T)\\
60 & 290 & (3) & (N3T)\\
73 & 290 & (3) & note 18 (N3T)\\
83 & 290 & (4) & (N3T)\\
91 & 290 & (5) & (N3)\\
92 & 290 & (3) & note 18 (N3T)\\
128 & 290 & (3) & (N3T)\\
130 & 290 & (3) & note 18 (N3T)\\
163 & 290 & (3) & (N3T)\\
164 & 290 & (5) & (N3)\\
165 & 290 & (5) & (N3)\\
191 & 290 & (5) & note 18 (N3)\\
273 & 290 & (3) & note 18 (N3T)\\
305 & 290 & (3) & note 18 (N3T)\\
317 & 290 & (5) & (N3)\\
361 & 290 & (3) & note 18 (N3T)\\
363 & 290 & (3) & note 18 (N3T)\\
381 & 290 & (5) & Howard/Keremedis/Rubin/Rubin\index{Author}{Howard/Keremedis/Rubin/Rubin} [\textbf{1997b}] (N3)\\
0 &  291 & (3) & (N29T)\\
6 & 291 & (3) & note 18 (N29T)\\
11 & 291 & (3) & note 18 (N29T)\\
14 & 291 & (3) & (N29T)\\
16 & 291 & (5) & (N29)\\
\end{tabular}
\end{table}

\clearpage

\begin{table}[h]
\begin{tabular}{llll}
23 & 291 & (5) & note 2(1) (N29)\\
24 & 291\index{Author}{Morris} & (5) & (N29)\\
30 & 291 & (3) & (N29T)\\
31 & 291 & (3) & (N29T)\\
35 & 291 & (3) & note 18 (N29T)\\
37 & 291 & (3) & note 18 (N29T)\\
60 & 291 & (3) & (N29T)\\
73 & 291 & (3) & note 18 (N29T)\\
91 & 291 & (5) & note 18 (N29)\\
92 & 291 & (3) & note 18 (N29T)\\
130 & 291 & (3) & note 18 (N29T)\\
165 & 291 & (5) & (N29)\\
191 & 291 & (5) & note 18 (N29)\\
200 & 291 & (5) & (N29)\\
273 & 291 & (3) & note 18 (N29T)\\
305 & 291 & (3) & note 18 (N29T)\\
309 & 291 & (3) & note 18 (N29T)\\
313 & 291 & (3) & note 18 (N29T)\\
317 & 291 & (5) & note 2(4) (N29)\\
361 & 291 & (3) & note 18 (N29T)\\
363 & 291 & (3) & note 18 (N29T)\\
17 & 292 & (3) & (N1T)\\
31 & 292 & (3) & (N1T)\\
67 & 292 & (1) & clear\\
89 & 292 & (5) & (N1)\\
112 & 292 & (5) & Brunner\index{Author}{Brunner} [\textbf{1984b}], Howard/Rubin\index{Author}{Howard/Rubin} [\textbf{1977}] (N1)\\
114 & 292 & (5) & (N1)\\
115 & 292 & (5) & (N1)\\
116 & 292 & (5) & (N1)\\
133 & 292 & (5) & (N1)\\
134 & 292 & (5) & (N1)\\
135 & 292 & (5) & (N1)\\
147 & 292 & (5) & (N1)\\
263 & 292 & (5) & (N1)\\
325 & 292 & (3) & (N1T)\\
0 &  293 & (3) & (N37T)\\
6 & 293 & (3) & note 18 (N37T)\\
11 & 293 & (3) & note 18 (N37T)\\
30 & 293 & (1) & clear\\
35 & 293 & (3) & note 18 (N37T)\\
37 & 293 & (3) & note 18 (N37T)\\
73 & 293 & (3) & note 18 (N37T)\\
91 & 293 & (5) & note 18 (N37)\\
92 & 293 & (3) & note 18 (N37T)\\
130 & 293 & (3) & note 18 (N37T)\\
185 & 293 & (3) & (N37T)\\
191 & 293 & (5) & note 18 (N37)\\
273 & 293 & (3) & note 18 (N37T)\\
305 & 293 & (3) & note 18 (N37T)\\
\end{tabular}
\end{table}

\clearpage

\begin{table}[h]
\begin{tabular}{llll}
309 & 293 & (3) & note 18 (N37T)\\
313 & 293\index{Author}{Brunner}\index{Author}{Howard/Keremedis/Rubin/Rubin} & (3) & note 18 (N37T)\\
316 & 293 & (5) & (N37)\\
361 & 293 & (3) & note 18 (N28T)\\
363 & 293 & (3) & note 18 (N28T)\\
6 & 294 & (3) & note 18 (N41T)\\
11 & 294 & (3) & note 18 (N41T)\\
35 & 294 & (3) & note 18 (N41T)\\
37 & 294 & (3) & note 18 (N41T)\\
51 & 294 & (1) & Morris [\textbf{1969}]\\
63 & 294 & (3) & (N41T)\\
73 & 294 & (3) & note 18 (N41T)\\
91 & 294 & (5) & note 18 (N41)\\
92 & 294 & (3) & note 18 (N41T)\\
130 & 294 & (3) & note 18 (N41T)\\
122 & 294 & (3) & (N41T)\\
191 & 294 & (5) & note 18 (N41)\\
231 & 294 & (1) & clear\\
273 & 294 & (3) & note 18 (N41T)\\
305 & 294 & (3) & note 18 (N41T)\\
309 & 294 & (3) & note 18 (N41T)\\
313 & 294 & (3) & note 18 (N41T)\\
361 & 294 & (3) & note 18 (N41T)\\
363 & 294 & (3) & note 18 (N41T)\\
14 & 295 & (3) & (N3T)\\
15 & 295 & (1) & Pincus\index{Author}{Pincus} [\textbf{1997}]\\
317 & 295 & (5) & (N3)\\
381 & 295 & (5) & Howard/Keremedis/Rubin/Rubin\index{Author}{Howard/Keremedis/Rubin/Rubin} [\textbf{1997b}] (N3)\\
9 & 296 & (1) & clear\\
14 & 296 & (3) & (N48T)\\
15 & 296 & (1) & Pincus [\textbf{1997}]\\
60 & 296 & (3) & (N3T)\\
128 & 296 & (3) & (N3T)\\
295 & 296 & (3) & (N48T)\\
317 & 296 & (5) & (N3)\\
381 & 296 & (5) & Howard/Keremedis/Rubin/Rubin [\textbf{1997b}] (N3)\\
0 &  298 & (7) & This project\\
14 & 298 & (1) & Gleason\index{Author}{Gleason} [\textbf{1958}], Bell\index{Author}{Bell} \cite{bib:1988b}\\
14 & 299 & (3) & Bell \cite{bib:1988b}, Monro\index{Author}{Monro} [\textbf{1972}] (M1)\\
15 & 299 & (3) & (M1)\\
31 & 299 & (3) & (M1)\\
60 & 299 & (3) & (M1)\\
118 & 299 & (3) & (M1)\\
128 & 299 & (3) & (M1)\\
163 & 299 & (3) & (M1)\\
165 & 299 & (3) & (M1)\\
191 & 299 & (3) & (M1)\\
277 & 299 & (3) & (M1)\\
297 & 299 & (1) & Bell \cite{bib:1988b}\\
\end{tabular}
\end{table}

\clearpage

\begin{table}[h]
\begin{tabular}{llll}
14 & 300 & (3) & Bell \cite{bib:1988b}, Monro [\textbf{1972}] (M1)\\
15 & 300\index{Author}{Degen}\index{Author}{Pawlikowski}\index{Author}{Wagon} & (3) & (M1)\\
31 & 300 & (3) & (M1)\\
60 & 300 & (3) & (M1)\\
118 & 300 & (3) & (M1)\\
128 & 300 & (3) & (M1)\\
163 & 300 & (3) & (M1)\\
165 & 300 & (3) & (M1)\\
191 & 300 & (3) & (M1)\\
277 & 300 & (3) & (M1)\\
301 & 300 & (1) & Bell \cite{bib:1988b}\\
302 & 301 & (1) & Bell \cite{bib:1988b}\\
0 &  304 & (4) & Brunner [\textbf{1985c}](N3T)\\

6 & 304 & (3) & note 18 (N3T)\\
9 & 304 & (1) & clear\\
11 & 304 & (3) & note 18 (N3T)\\
14 & 304 & (3) & (N3T)\\
16 & 304 & (5) & (N3)\\
23 & 304 & (5) & (N3)\\
24 & 304 & (5) & (N3)\\
31 & 304 & (3) & (N3T)\\
35 & 304 & (3) & note 18 (N3T)\\
37 & 304 & (3) & note 18 (N3T)\\
60 & 304 & (3) & (N3T)\\
73 & 304 & (3) & note 18 (N3T)\\
92 & 304 & (3) & note 18 (N3T)\\
128 & 304 & (3) & (N3T)\\
130 & 304 & (3) & note 18 (N3T)\\
163 & 304 & (3) & (N3T)\\
164 & 304 & (5) & (N3)\\
165 & 304 & (5) & (N3)\\
191 & 304 & (5) & note 18 (N3)\\
273 & 304 & (3) & note 18 (N3T)\\
305 & 304 & (3) & note 18 (N3T)\\
317 & 304 & (5) & (N3)\\
361 & 304 & (3) & note 18 (N3T)\\
363 & 304 & (3) & note 18 (N3T)\\
381 & 304 & (5) & Howard/Keremedis/Rubin/Rubin\index{Author}{Howard/Keremedis/Rubin/Rubin} [\textbf{1997b}] (N3)\\
203 & 306 & (1) & clear\\
305 & 306 & (1) & Kanovei\index{Author}{Kanovei} [\textbf{1991}]\\
43 & 307 & (3) & Kanovei [\textbf{1991}] (M39$(\kappa, \lambda)$)\\
305 & 307 & (1) & clear\\
6 & 308 & (3) & note 18 (N6T)\\
11 & 308 & (3) & note 18 (N6T)\\
35 & 308 & (3) & note 18 (N6T)\\
37 & 308 & (3) & note 18 (N6T)\\
52 & 308 & (3) & (N6T)\\
61 & 308 & (3) & (N6T)\\
62 & 308 & (1) & Howard/Yorke\index{Author}{Howard/Yorke} [\textbf{1987}]\\
\end{tabular}
\end{table}

\clearpage

\begin{table}[h]
\begin{tabular}{llll}
67 & 308 & (5) & (N6)\\
73 & 308\index{Author}{Bull}\index{Author}{H\"{o}ft/Howard}\index{Author}{Howard} & (3) & note 18 (N6T)\\
82 & 308 & (3) & (N6T)\\
92 & 308 & (3) & note 18 (N6T)\\
119 & 308 & (3) & (N6T)\\
130 & 308 & (3) & note 18 (N6T)\\
185 & 308 & (3) & (N6T)\\
191 & 308 & (5) & note 18 (N6)\\
218 & 308 & (5) & (N6)\\
273 & 308 & (3) & note 18 (N6T)\\
305 & 308 & (3) & note 18 (N6T)\\
313 & 308 & (3) & note 18 (N6T)\\
350 & 308 & (3) & (N6T)\\
361 & 308 & (3) & note 18 (N6T)\\
363 & 308 & (3) & note 18 (N6T)\\
52 & 309 & (1) & Pawlikowski [\textbf{1991}]\\
14 & 311 & (1) & Wagon \cite{bib:1985u}\\
43 & 312 & (3) & (M5$(\aleph)$ IC)\\
307 & 312 & (3) & (M5$(\aleph)$ IC)\\
311 & 313 & (1) & clear\\
0 &  314 & (4) & Degen\index{Author}{Degen} [\textbf{1988}] (N6T)\\
5 & 314 & (3) & (N6T)\\
6 & 314 & (3) & note 18 (N6T)\\
11 & 314 & (3) & note 18 (N6T)\\
13 & 314 & (3) & (N6T)\\
35 & 314 & (3) & note 18 (N6T)\\
37 & 314 & (3) & note 18 (N6T)\\
38 & 314 & (3) & (N6T)\\
52 & 314 & (3) & (N6T)\\
61 & 314 & (3) & (N6T)\\
70 & 314 & (3) & (N6T)\\
73 & 314 & (3) & note 18 (N6T)\\
79 & 314 & (3) & (N6T)\\
82 & 314 & (3) & (N6T)\\
85 & 314 & (1) & Degen [\textbf{1988}], note 124\\
92 & 314 & (3) & note 18 (N6T)\\
119 & 314 & (3) & (N6T)\\
130 & 314 & (3) & note 18 (N6T)\\
185 & 314 & (3) & (N6T)\\
191 & 314 & (5) & note 18 (N6)\\
218 & 314 & (5) & (N6)\\
273 & 314 & (3) & note 18 (N6T)\\
305 & 314 & (3) & note 18 (N6T)\\
313 & 314 & (3) & note 18 (N6T)\\
350 & 314 & (3) & (N6T)\\
361 & 314 & (3) & note 18 (N6T)\\
363 & 314 & (3) & note 18 (N6T)\\
0 &  315 & (3) & Gitik\index{Author}{Gitik} [\textbf{1980}] (M17 IC)\\
25 & 315 & (1) & clear\\
\end{tabular}
\end{table}

\clearpage

\begin{table}[h]
\begin{tabular}{llll}
51 & 316 & (1) & H\"{o}ft/Howard [\textbf{1994}]\\
77 & 316 & (5) & (N47)\\
50 & 317 & (1) & Howard \cite{bib:1973h}\\
43 & 318 & (3) & (M$(\aleph)$5 IC)\\
307 & 318 & (3) & (M$(\aleph)$5 IC)\\
320 & 318 & (1) & clear\\
117 & 319 & (1) & clear\\
319 & 320 & (1) & note 20\\
39 & 321 & (3) & Bull [\textbf{1978}] (M42 IC)\\
86 & 321 & (3) & Bull [\textbf{1978}] (M42 IC)\\
320 & 321 & (1) &clear\\
6 & 322 & (3) & note 18 (N3T)\\
11 & 322 & (3) & note 18 (N3T)\\
14 & 322 & (3) & (N3T)\\
15 & 322 & (1) & clear\\
23 & 322 & (5) & (N3)\\
31 & 322 & (3) & (N3T)\\
37 & 322 & (3) & note 18 (N3T)\\
40 & 322 & (1) & clear\\
60 & 322 & (3) & (N29T)\\
73 & 322 & (3) & note 18 (N3T)\\
91 & 322 & (5) & note 18 (N3)\\
92 & 322 & (3) & note 18 (N3T)\\
130 & 322 & (3) & note 18 (N3T)\\
133 & 322 & (5) & (N26)\\
273 & 322 & (3) & note 18 (N3T)\\
305 & 322 & (3) & note 18 (N3T)\\
6 & 323 & (3) & note 18 (N1T)\\
9 & 323 & (3) & (N12$(\aleph_{1})$T)\\
11 & 323 & (3) & note 18 (N1T)\\
15 & 323 & (1) & clear\\
23 & 323 & (5) & (N1)\\
31 & 323 & (3) & (N1T)\\
37 & 323 & (3) & note 18 (N1T)\\
40 & 323 & (3) & (N12$(\aleph_{1})$T)\\
60 & 323 & (1) & clear\\
89 & 323 & (5) & (N1)\\
112 & 323 & (5) & (N1)\\
114 & 323 & (5) & (N1)\\
115 & 323 & (5) & (N1)\\
116 & 323 & (5) & (N1)\\
121 & 323 & (3) & (N12$(\aleph_{1})$T)\\
130 & 323 & (3) & note 18 (N1T)\\
133 & 323 & (5) & (N1)\\
134 & 323 & (5) & (N1)\\
135 & 323 & (5) & (N1)\\
147 & 323 & (5) & (N1)\\
187 & 323 & (5) & (N12$(\aleph_{1})$)\\
202 & 323 & (5) & (N12$(\aleph_{1})$)\\
\end{tabular}
\end{table}

\clearpage

\begin{table}[h]
\begin{tabular}{llll}
233 & 323 & (5) & (N1)\\
263 & 323 & (5) & (N1)\\
273 & 323 & (3) & note 18 (N1T)\\
305 & 323 & (3) & note 18 (N1T)\\
309 & 323 & (3) & note 18 (N1T)\\
313 & 323 & (3) & note 18 (N1T)\\
6 & 324 & (3) & note 18 (N2T)\\
11 & 324 & (3) & note 18 (N2T)\\
30 & 324 & (3) & (M29)\\
35 & 324 & (3) & note 18 (N2T)\\
37 & 324 & (3) & note 18 (N2T)\\
52 & 324 & (3) & (N2T)\\
65 & 324 & (5) & (N2)\\
67 & 324 & (5) & (N2)\\
73 & 324 & (3) & note 18 (N2T)\\
83 & 324 & (3) & (N2T)\\
92 & 324 & (3) & note 18 (N2T)\\
99 & 324 & (5) & (N2)\\
119 & 324 & (3) & (N2T)\\
130 & 324 & (3) & note 18 (N2T)\\
163 & 324 & (3) & (N2T)\\
165 & 324 & (1) & clear\\
185 & 324 & (3) & (N2T)\\
200 & 324 & (3) & (M29)\\
273 & 324 & (3) & note 18 (N2T)\\
290 & 324 & (3) & (M29)\\
295 & 324 & (3) & (M29)\\
305 & 324 & (3) & note 18 (N2T)\\
309 & 324 & (3) & note 18 (N2T)\\
313 & 324 & (3) & note 18 (N2T)\\
322 & 324 & (1) & clear\\
323 & 324 & (1) & clear\\
9 & 325 & (1) & note 46\\
49 & 326 & (1) & Wojtylak\index{Author}{Wojtylak} [\textbf{1995}]\\
44 & 327 & (3) & (N2$(\aleph_{\alpha})$T)\\
61 & 327 & (3) & (N6T)\\
122 & 327 & (1) & clear\\
135 & 327 & (5) & (N2$(\aleph_{\alpha})$)\\
218 & 327 & (5) & (N6)\\
324 & 327 & (1) & clear\\
8 & 328 & (3) & (N38T)\\
23 & 328 & (5) & (N26)\\
31 & 328 & (3) & (N26T)\\
40 & 328 & (1) & clear\\
43 & 328 & (3) & (N40T)\\
67 & 328 & (1) & clear\\
79 & 328 & (3) & (N26T)\\
97 & 328 & (3) & note 120(32) (M1)\\
163 & 328 & (3) & (N26T)\\
\end{tabular}
\end{table}

\clearpage

\begin{table}[h]
\begin{tabular}{llll}
165 & 328 & (3) & (N26T)\\
185 & 328 & (3) & (N26T)\\
321 & 328 & (5) & (N1)\\
345 & 328 & (5) & N40\\
23 & 329 & (5) & (N33)\\
40 & 329 & (3) & (N33T)\\
60 & 329 & (1) & clear\\
67 & 329 & (1) & clear\\
70 & 329 & (3) & (N33T)\\
97 & 329 & (3) & (M1$(<\omega_{1}>)$)\\
191 & 329 & (5) & note 18 (N33)\\
309 & 329 & (3) & note 18 (N33T)\\
3 & 330 & (3) & (M6)\\
6 & 330 & (3) & note 18 (N41T)\\
8 & 330 & (3) & note 120(34) (N15T)\\
11 & 330 & (3) & note 18 (N41T)\\
15 & 330 & (3) & (M34)\\
35 & 330 & (3) & note 18 (N41T)\\
37 & 330 & (3) & note 18 (N41T)\\
43 & 330 & (3) & (N15T)\\
63 & 330 & (3) & (N41T)\\
73 & 330 & (3) & note 18 (N41T)\\
91 & 330 & (5) & note 18 (N41)\\
92 & 330 & (3) & note 18 (N41T)\\
126 & 330 & (3) & (M6)\\
130 & 330 & (3) & note 18 (N41T)\\
165 & 330 & (1) & clear\\
191 & 330 & (5) & note 18 (N41)\\
200 & 330 & (3) & (M29)\\
214 & 330 & (3) & (M29)\\
273 & 330 & (3) & note 18 (N41T)\\
291 & 330 & (3) & (M29)\\
305 & 330 & (3) & note 18 (N41T)\\
309 & 330 & (3) & note 18 (N41T)\\
313 & 330 & (3) & note 18 (N41T)\\
328 & 330 & (1) & clear\\
329 & 330 & (1) & clear\\
123 & 331 & (1) & Morillon\index{Author}{Morillon} \cite{bib:1988h}\\
332 & 331 & (7) & Morillon \cite{bib:1988h}\\
331 & 332 & (1) & Morillon \cite{bib:1988h}\\
343 & 332 & (7) & Morillon \cite{bib:1988h}\\
95 & 333 & (1) & clear\\
218 & 333 & (1) & clear\\
335 & 333 & (1) & Keremedis\index{Author}{Keremedis} [\textbf{1996a}]\\
218 & 334 & (5) & (N6)\\
9 & 336 & (1) & clear\\
61 & 336 & (1) & clear\\
378 & 336 & (1) & clear\\
0 &  337 & (4) & Good/Tree\index{Author}{Good/Tree} [\textbf{1995}] (M1)\\
\end{tabular}
\end{table}

\clearpage

\begin{table}[h]
\begin{tabular}{llll}
8 & 337 & (1) & clear\\
14 & 337\index{Author}{Brunner}\index{Author}{Morillon} & (3) & (M1)\\
15 & 337 & (3) & (M1)\\
31 & 337 & (3) & (M1)\\
60 & 337 & (3) & (M1)\\
118 & 337 & (3) & (M1)\\
128 & 337 & (3) & (M1)\\
163 & 337 & (3) & (M1)\\
165 & 337 & (3) & (M1)\\
191 & 337 & (3) & (M1)\\
277 & 337 & (3) & (M1)\\
0 & 338 & (4) & Good/Tree\index{Author}{Good/Tree} [\textbf{1995}] (M1)\\
8 & 338 & (1) & clear\\
14 & 338 & (3) & (M1)\\
15 & 338 & (3) & (M1)\\
31 & 338 & (3) & (M1)\\
60 & 338 & (3) & (M1)\\
118 & 338 & (3) & (M1)\\
128 & 338 & (3) & (M1)\\
163 & 338 & (3) & (M1)\\
165 & 338 & (3) & (M1)\\
191 & 338 & (3) & (M1)\\
277 & 338 & (3) & (M1)\\
0 &  339 & (4) & Good/Tree [\textbf{1995}] (M1)\\
14 & 339 & (3) & (M1)\\
15 & 339 & (3) & (M1)\\
31 & 339 & (3) & (M1)\\
60 & 339 & (3) & (M1)\\
118 & 339 & (3) & (M1)\\
128 & 339 & (3) & (M1)\\
163 & 339 & (3) & (M1)\\
165 & 339 & (3) & (M1)\\
191 & 339 & (3) & (M1)\\
277 & 339 & (3) & (M1)\\
360 & 339 & (1) & clear\\
0 & 340 & (4) & Good/Tree [\textbf{1995}] (M7)\\
8 & 340 & (1) & clear\\
191 & 340 & (3) & (M7)\\
0 & 341 & (4) & Good/Tree [\textbf{1995}] (M7)\\
8 & 341 & (1) & clear\\
191 & 341 & (3) & (M7)\\
6 & 342 &(3) & note 18 (N50$(E)$T)\\
8 & 342 & (1) & clear\\
9 & 342 & (1) & clear\\
11 & 342 & (3) & note 18 (N50$(E)$T)\\
35 & 342 & (3) & note 18 (N50$(E)$T)\\
37 & 342 & (3) & note 18 (N50$(E)$T)\\
52 & 342 &(3) &(N50$(E)$T)\\
61 & 342 & (1) & clear\\
\end{tabular}
\end{table}

\clearpage

\begin{table}[h]
\begin{tabular}{llll}
67 & 342 & (5) & Brunner [\textbf{1990}] (N50$(E)$)\\
73 & 342\index{Author}{Higasikawa}\index{Author}{Morillon} & (3) & note 18 (N50$(E)$T)\\
79 & 342 & (3) & (N50$(E)$T)\\
82 & 342 & (3) & (N50$(E)$T)\\
92 & 342 & (3) & note 18 (N50$(E)$T)\\
119 & 342 & (3) & (N50$(E)$T)\\
130 & 342 & (3) & note 18 (N50$(E)$T)\\
132 & 342 & (1) & clear\\
185 & 342 & (3) & (N50$(E)$T)\\
191 & 342 & (5) & note 18 (N50$(E)$)\\
250 & 342 & (3) & (N50$(E)$T)\\
305 & 342 & (3) & note 18 (N50$(E)$T)\\
313 & 342 & (3) & note 18 (N50$(E)$T)\\
322 & 342 & (5) & Brunner [\textbf{1990}] (N50$(E)$)\\
350 & 342 &(3) &(N50$(E)$T)\\
357 & 342 & (3) & (N50$(E)$T)\\
358 & 342 & (3) & (N50$(E)$T)\\
361 & 342 & (3) & note 18 (see N50$(E)$T)\\
363 & 342 & (3) & note 18 (see N50$(E)$T)\\
332 & 343 & (1) & Morillon\index{Author}{Morillon} \cite{bib:1988h}\\
3 & 344 & (5) & (N9)\\
11 & 344 & (5) & (N1)\\
12 & 344 & (3) & (N1T)\\
23 & 344 & (5) & (N1)\\
39 & 344 & (3) & (N2$(\aleph_{\alpha})$T)\\
43 & 344 & (3) & (M47(n,M))\\
44 & 344 & (5) & (N2$(\aleph_{\alpha})$)\\
46 & 344 & (3) & (M47(n,M))\\
52 & 344 & (3) & (N2(n)T)\\
63 & 344 & (3) & (M47(n,M))\\
65 & 344 & (5) & (N2)\\
67 & 344 & (5) & (N2(n))\\
71 & 344 & (3) & (N2$(\aleph_{\alpha})$T)\\
79 & 344 & (3) & (N1T)\\
86 & 344 & (3) & (N2$(\aleph_{\alpha})$T)\\
87 & 344 & (5) & (N2$(\aleph_{\alpha})$)\\
99 & 344 & (5) & (N2)\\
115 & 344 & (5) & (N1)\\
116 & 344 & (5) & (N1)\\
122 & 344 & (3) & (N1T)\\
123 & 344 & (1) & Morillon \cite{bib:1988h}\\
130 & 344 & (5) & (N1)\\
133 & 344 & (5) & (N1)\\
134 & 344 & (5) & (N1)\\
135 & 344 & (5) & (N1)\\
141 & 344 & (5) & (N2$^{\ast}$(3))\\
163 & 344 & (5) & (N1)\\
187 & 344 & (5) & (N12$(\aleph_{1})$)\\
191 & 344 & (5) & (N1)\\
\end{tabular}
\end{table}

\clearpage

\begin{table}[h]
\begin{tabular}{llll}
202 & 344 & (5) & (N12$(\aleph_{1})$)\\
214 & 344\index{Author}{Keremedis} & (3) & (M47(n,M))\\
218 & 344 & (5) & (N6)\\
233 & 344 & (5) & (N1)\\
263 & 344 & (5) & (N1)\\
305 & 344 & (5) & (N1)\\
308 & 344 & (5) & (N22(K,p))\\
313 & 344 & (5) & (N1)\\
331 & 344 & (7) & Morillon \cite{bib:1988h}\\
334 & 344 & (5) & (N2)\\
40 & 347 & (3) & Higasikawa [\textbf{1995}] (N12$(\aleph_{1})$T)\\
101 & 347 & (1) &clear\\
0 &  348 & (7) & Howard\index{Author}{Howard} [\textbf{1985}]\\
11 & 349 & (3) & note 18 (N33T)\\
23 & 349 & (5) & (N33)\\
40 & 349 & (3) &(N33T)\\
63 & 349 & (5) & (N33)\\
70 & 349 & (3) & (N33T)\\
85 & 349 & (1) & clear\\
91 & 349 & (5) & note 18 (N33)\\
130 & 349 & (3) & note 18 (N33T)\\
165 & 349 & (5) & (N33)\\
191 & 349 & (5) & note 18 (N33)\\
305 & 349 & (3) & note 18 (N33T)\\
309 & 349 & (3) & note 18 (N33T)\\
313 & 349 & (3) & note 18 (N33T)\\
329 & 349 & (1) & clear\\
6 & 350 & (3) & note 18 (N41T)\\
9 & 350 & (3) & (M29)\\
11 & 350 & (3) & note 18 (N41T)\\
30 & 350 & (3) & (M29)\\
32 & 350 & (1) & clear\\
35 & 350 & (3) & note 18 (N41T)\\
37 & 350 & (3) & note 18 (N41T)\\
63 & 350 & (5) & (N41)\\
70 & 350 & (3) & (N41T)\\
73 & 350 & (3) & note 18 (N41T)\\
91 & 350 & (5) & note 18 (N41)\\
92 & 350 & (3) & note 18 (N41T)\\
126 & 350 & (1) & clear\\
130 & 350 & (3) & note 18 (N41T)\\
191 & 350 & (5) & note 18 (N41)\\
200 & 350 & (3) & (M29)\\
214 & 350 & (3) & (M29)\\
290 & 350 & (3) & (M29)\\
291 & 350 & (3) & (M29)\\
295 & 350 & (3) & (M29)\\
305 & 350 & (3) & note 18 (N41T)\\
309 & 350 & (3) & note 18 (N41T)\\
\end{tabular}
\end{table}

\clearpage

\begin{table}[h]
\begin{tabular}{llll}
313 & 350 & (3) & note 18 (N41T)\\
330 & 350 & (1) & clear\\
349 & 350 & (1) & clear\\
351 & 350 & (1) & Keremedis [\textbf{1996d}]\\
353 & 350 & (1) & Keremedis [\textbf{1996d}]\\
361 & 350 & (3) & note 18 (N41T)\\
363 & 350 & (3) & note 18 (N41T)\\
8 & 351 & (1) & Keremedis [\textbf{1996d}]\\
8 & 352 & (1) & Keremedis [\textbf{1996d}]\\
8 & 353 & (1) & Keremedis [\textbf{1996d}]\\
8 & 354 & (1) & Keremedis [\textbf{1996d}]\\
8 & 355 & (1) & clear\\
14 & 355 & (3) & (N3T)\\
16 & 355 & (5) & (N3)\\
23 & 355 & (5) & (N3)\\
24 & 355 & (5) & (N3)\\
31 & 355 & (3) & (N3T)\\
60 & 355 & (3) & (N3T)\\
79 & 355 & (3) & (N3T)\\
91 & 355 & (5) & (N3)\\
128 & 355 & (3) & (N3T)\\
163 & 355 & (3) & (N3T)\\
164 & 355 & (5) & (N3)\\
165 & 355 & (5) & (N3)\\
317 & 355 & (5) & (N3)\\
322 & 355 & (1) & clear\\
381 & 355 & (5) & Howard/Keremedis/Rubin/Rubin\index{Author}{Howard/Keremedis/Rubin/Rubin} [\textbf{1997b}] (N3)\\
11 & 356 & (3) & note 18 (N33T)\\
23 & 356 & (5) & (N33)\\
40 & 356 & (3) & (N33T)\\
63 & 356 & (5) & (N33)\\
70 & 356 & (3) & (N33T)\\
85 & 356 & (1) & clear\\
91 & 356 & (5) & note 18 (N33)\\
130 & 356 & (3) & note 18 (N33T)\\
165 & 356 & (5) & (N33)\\
191 & 356 & (5) & note 18 (N33)\\
305 & 356 & (3) & note 18 (N33T)\\
309 & 356 & (3) & note 18 (N33T)\\
313 & 356 & (3) & note 18 (N33T)\\
323 & 356 & (1) & clear\\
6 & 357 & (3) & note 18 (N41T)\\
9 & 357 & (3) & (M29)\\
11 & 357 & (3) & note 18 (N41T)\\
30 & 357 & (3) & (M29)\\
32 & 357 & (1) & clear\\
35 & 357 & (3) & note 18 (N41T)\\
37 & 357 & (3) & note 18 (N41T)\\
63 & 357 & (5) & (N41)\\
\end{tabular}
\end{table}

\clearpage

\begin{table}[h]
\begin{tabular}{llll}
70 & 357 & (3) & (N41T)\\
73 & 357\index{Author}{Hamel}\index{Author}{Lindenbaum/Tarski}\index{Author}{G. Moore}\index{Author}{Morillon} & (3) & note 18 (N41T)\\
91 & 357 & (5) & note 18 (N41)\\
92 & 357 & (3) & note 18 (N41T)\\
130 & 357 & (3) & note 18 (N41T)\\
191 & 357 & (5) & note 18 (N41)\\
200 & 357 & (3) & (M29)\\
214 & 357 & (3) & (M29)\\
290 & 357 & (3) & (M29)\\
291 & 357 & (3) & (M29)\\
295 & 357 & (3) & (M29)\\
305 & 357 & (3) & note 18 (N41T)\\
309 & 357 & (3) & note 18 (N41T)\\
313 & 357 & (3) & note 18 (N41T)\\
324 & 357 & (1) & clear\\
355 & 357 & (1) & clear\\
356 & 357 & (1) & clear\\
361 & 357 & (3) & note 18 (N41T)\\
363 & 357 & (3) & note 18 (N41T)\\
6 & 358 & (3) & note 18 (N2T)\\
10 & 358 & (1) & clear\\
11 & 358 & (3) & note 18 (N2T)\\
35 & 358 & (3) & note 18 (N2T)\\
37 & 358 & (3) & note 18 (N2T)\\
52 & 358 & (3) & (N2$^{\ast}$(3)T)\\
61 & 358 & (3) & (N6T)\\
65 & 358 & (5) & (N2)\\
70 & 358 & (3) & (N2T)\\
73 & 358 & (3) & note 18 (N2T)\\
79 & 358 & (3) & (N2$^{\ast}$(3)T)\\
82 & 358 & (3) & (N2$^{\ast}$(3)T)\\
88 & 358 & (3) & (N2$^{\ast}$(3)T)\\
92 & 358 & (3) & note 18 (N2T)\\
95 & 358 & (5) &(N6)\\
99 & 358 & (5) & (N2)\\
119 & 358 & (3) & (N2$^{\ast}$(3)T)\\
130 & 358 & (3) & note 18 (N2T)\\
141 & 358 & (5) & (N2$^{\ast}$(3))\\
163 & 358 & (3) & (N2T)\\
185 & 358 & (3) & (N2$^{\ast}$(3)T)\\
191 & 358 & (3) & (M7)\\
218 & 358 & (5) & (N6)\\
305 & 358 & (3) & note 18 (N2T)\\
313 & 358 & (3) & note 18 (N2T)\\
327 & 358 & (1) & clear\\
333 & 358 & (5) & (N2$^{\ast}$(3))\\
334 & 358 & (5) & (N2)\\
350 & 358 & (3) & (N2T)\\
355 & 358 & (1) & clear\\
\end{tabular}
\end{table}

\clearpage

\begin{table}[h]
\begin{tabular}{llll}
361 & 358 & (3) & note 18 (N2T)\\
363 & 358\index{Author}{Howard/Keremedis/Rubin/Rubin} & (3) & note 18 (N2T)\\
8 & 360 & (1) & clear\\
8 & 361 & (1) & G. Moore [\textbf{1982}] p 325 \\
0 &  362 & (7) & This project\\
361 & 362 & (1) & G. Moore [\textbf{1982}] p 325 \\
43 & 363 & (1) & G. Moore [\textbf{1982}] p 325 \\
363 & 364 & (1) & G. Moore [\textbf{1982}] p 325 \\
41 & 365 & (1) & G. Moore [\textbf{1982}] table 5 (clear)\\
367 & 366 & (1) & Hamel [\textbf{1905}]\\
66 & 367 & (1) & clear\\
79 & 367 & (1) & Hamel [\textbf{1905}]\\
79 & 368 & (1) & clear\\
79 & 369 & (1) & clear\\
100 & 369 & (1) & Lindenbaum/Tarski [\textbf{1926}]\\
14 & 370 & (1) & clear\\
14 & 371 & (1) & Morillon [\textbf{1993}]\\
42 & 371 & (3) & (M27)\\
53 & 371 & (3) & (M27)\\
79 & 371 & (1) & Morillon [\textbf{1993}]\\
253 & 371 & (3) & (M27)\\
6 & 373 & (3) & note 18 (N2$(n)$T)\\
11 & 373 & (3) & note 18 (N2$(n)$T)\\
35 & 373 & (3) & note 18 (N2$(n)$T)\\
37 & 373 & (3) & note 18 (N2$(n)$T)\\
47 & 373 & (1) & clear\\
52 & 373 & (3) & (N2$(n)$T)\\
67 & 373 & (5) & (N2$(n)$)\\
73 & 373 & (3) & note 18 (N2$(n)$T)\\
79 & 373 & (3) & (N2$(n)$T)\\
82 & 373 & (3) & (N2$(n)$T)\\
92 & 373 & (3) & note 18 (N2$(n)$T)\\
119 & 373 & (3) & (N2$(n)$T)\\
130 & 373 & (3) & note 18 (N2$(n)$T)\\
185 & 373 & (3) & (N2$(n)$T)\\
191 & 373 & (5) & note 18 (N2$(n)$)\\
273 & 373 & (3) & note 18 (N2$(n)$T)\\
288 & 373 & (1) & clear\\
305 & 373 & (3) & note 18 (N2$(n)$T)\\
313 & 373 & (3) & note 18 (N2$(n)$T)\\
350 & 373 & (3) & (N2$(n)$T\\
361 & 373 & (3) & note 18 (N2$(n)$T)\\
363 & 373 & (3) & note 18 (N2$(n)$T)\\
10 & 374 & (1) & clear\\
43 & 375 & (1) & (see Kelley\index{Author}{Kelley} [\textbf{1955}] p 242 and entry (43,78))\\
67 & 375 & (1) & note 139\\
9 & 376 & (1) & clear\\
15 & 376 & (1) & clear\\
376 & 377 & (1) & De la Cruz/Di Prisco\index{Author}{De la Cruz/Di Prisco} [\textbf{1996}]\\
\end{tabular}
\end{table}

\clearpage

\begin{table}[h]
\begin{tabular}{llll}
30 & 377 & (1) & clear\\
62 & 378 & (1) & clear\\
377 & 378 & (1) & clear\\
15 & 379 & (1) & clear\\
8 & 380 & (1) & clear\\
9 & 380 & (3) & (N49T)\\
60 & 380 & (1) & clear\\
63 & 380 & (5) & ((N49))\\
67 & 381 & (1) & Howard/Keremedis/Rubin/Rubin\index{Author}{Howard/Keremedis/Rubin/Rubin} [\textbf{1997b}]\\
6 & 382 & (3) & note 18 (N53T)\\
9 & 382 & (3) & Howard/Keremedis/Rubin/Rubin [\textbf{1997b}] (N53T)\\
11 & 382 & (3) & note 18 (N53T)\\
35 & 382 & (3) & note 18 (N53T)\\
37 & 382 & (3) & note 18 (N53T)\\
63 & 382 & (5) & (N53)\\
73 & 382 & (3) & note 18 (N53T)\\
91 & 382 & (5) & note 18 (N53)\\
92 & 382 & (3) & note 18 (N53T)\\
130 & 382 & (3) & note 18 (N53T)\\
191 & 382 & (3) & note 18 (N53T)\\
273 & 382 & (3) & note 18 (N53T)\\
305 & 382 & (3) & note 18 (N53T)\\
309 & 382 & (3) & note 18 (N53T)\\
313 & 382 & (3) & note 18 (N53T)\\
361 & 382 & (3) & note 18 (N53T)\\
363 & 382 & (3) & note 18 (N53T)\\
381 & 382 & (1) & clear\\
9 & 383 & (3) & Howard/Keremedis/Rubin/Rubin [\textbf{1997b}] (N53T)\\
63 & 383 & (5) & (N53)\\
232 & 383 & (1) & Howard/Keremedis/Rubin/Stanley\index{Author}{Howard/Keremedis/Rubin/Stanley} [\textbf{1997}]\\
\end{tabular}
\end{table}

\clearpage

%%%%%%%part06

\chapter{Bibliography}
\label{pt06:part06}

\section*{Abbreviations of Journal Names}

\begin{enumerate}
\item[] Acta Fac. Rerum Natur. Univ. Comenian. Math. - Acta Facultatis Rerum Naturalium Universitatis Comenainae. Mathematica.

\item[] Acta Math. - Acta Mathematica

\item[] Acta Math. Acad. Sci. Hungar. - Acta mathematica Academiae Scientiarum Hungaricae.

\item[] Advances in Math. - Advances in Mathematics.

\item[] Algebra Universalis. - Algebra Universalis

\item[] Amer. Math. Monthly. - American Mathematical Monthly.

\item[] Amer. Math. Soc. Transl. - American Mathematical Society Translations.

\item[] Ann. Math. Logic. - Annals of Mathematical Logic.

\item[] Ann. Math. Ser. 2 - Annals of Mathematics. Series 2.

\item[] Ann. Pure Appl. Logic. - Annals of Pure and Applied Logic.

\item[] Ann. Soc. Polon. Math. - Soci\'{e}t\'{e} Polonaise de Mathematique. Annales. (Continued as Ann. Polon. Math.)

\item[] Ann. Polon. Math. - Annales Polonici Mathematici.

\item[] Anz. \"{O}sterreich. Akad. Wiss. Math.-Nat. Kl. \"{O}sterreichische Akademie der Wissenschaften. Mathematisch-Naturwissenschaftliche Klasse. Anzeiger.

\item[] Arch. Math. (Basel). - Archiv der Mathematik. Archives of Mathematics. Archives Math\'{e}matiques. 

\item[] Arch. Math. (Brno). - Spisy P\v{r}\v{i}rodov\v{e}deck\'{e} Fakulty University J. E. Purkyn\v{e} v Brn\v{e}. Archivum Mathematicum.

\item[] Arch. Math. Logik Grundlagenforschung - Archiv f\"{u}r Mathematische Logik und Grundlagenforschung.

\item[] Bol. Soc. Paran. Mat. (2). - Boletim da Sociedade Paranaemse de Matem\'{a}tica. $2^{a}$ S\'{e}rie. Soc. Paranaese de Mat.

\item[] Bull. Acad. Polon. Sci. S\'{e}r. Sci. Math. Astronom. Phys. - Bulletin de L'Acad\'{e}mie Polonaise des Sciences. S\'{e}rie des Sciences Math\'{e}matiques, Astronomiques et Physiques.

\item[] Bull. Amer. Math. Soc. - Bulletin of the American Mathematical Society.

\item[] Bull. Austral. Math. Soc. - Bulletin of the Australian Mathematical Society.

\item[] Bull. Greek Math. Soc. - Bulletin of the Greek Mathematical Society.

\item[] Bull. Int. Acad. Sci. Cracovie CI. Math. Nat. - Bulletin International de l'Acad\'{e}mie des Sciences de Cracovie. Classe de Math\'{e}matiques-Sciences Naturelles.

\item[] Bull. London Math. Soc. - Bulletin of the London Mathematical Society.

\item[] Bull. Malaysian Math. Soc. - Bulletin of the Malaysian Mathematical Society.

\item[] Bull. Sci. Math. - Bulletin des Sciences Math\'{e}matiques.

\item[] Bull. Sec. Logic Polish Acad. Sci. - Bulletin of the Section on Logic of the Polish Academy of Sciences.

\item[] C. R. Acad. Sci. Paris S\'{e}r. A--B. - Comptes Rendus Hebdomadaires des S\'{e}ances de l'Acad\'{e}mie des Sciences. S\'{e}ries A et B.

\item[] C. R. Acad. Sci. Paris S\'{e}r. I. - Comptes Rendus des S\'{e}ances de l'Academie des Sciences. Serie I: Mathematique.

\item[] C. R. Soc. Sci. Lett. Varsovie. - Soci\'{e}t\'{e} des Sciences et des lettres de Varsovie. Comptes Rendus des Seances. Classe III: Sciences Mathematiques et Physiques.

\item[] Canad. J. Math. - Canadian Journal of Mathematics.

\item[] Canad. Math. Bull. - Canadian Mathematical Bulletin.

\item[] \v{C}asopis P\v{e}st. Mat. - \v{C}eskoslovensk\'{a} Akademie V\v{e}d. \v{C}asopis Pro P\v{e}stov\'{a}n\'{i} Matematiky.

\item[] Colloq. Math. - Colloquium Mathematicum.

\item[] Comment. Math. Ann. Soc. Math. Polon. S\'{e}r. 1 - Annales Societatis Mathematicae Polonae. Series I. Commentationes Mahtematicae.

\item[] Comment. Math. Univ. Carolinae. - Commentationes Mathematicae Universitatis Carolinae.

\item[] Comment. Math. Helv. - Commentarii Mathematici Helvetici.

\item[] Comment. Math. Prace Mat. - Roczniki Polskiego Towarzystwa matematycznego. Seria I. Commentationes Mathematicae. Prace Matematyczne.

\item[] Compositio Math. - Compositio Mathematica.

\item[] Czechoslovak Math. J.- \v{C}ehoslovackaja Akademija Nauk. \v{C}ehoslovacki\v{i}
Matemati\v{c}eski\v{i} \v{Z}urnal. Czechoslovak Mathematical Journal.

\item[] Deutch. Akad. Naturforscher, nova Acta Leopoldina. - Deutsche Akademie der Naturforsch\d{e}r, Nova Acta Leopoldina.

\item[] Dissertationes Math. (Rozprawy Mat.) - Polska Akademia Nauk. Instytut Matematyczny. Dissertationes Mathematicae. Rozprawy Matematycany.

\item[] Dokl. Akad. Nauk SSSR. - Dokaldy Akademii Nauk SSSR.

\item[] Fund. Math. - Fundamenta Mathematicae. Polska Akademia Nauk.

\item[] Illinois J. Math. - Illinois Journal of Mathematics.

\item[] Indag. Math. - Indagationes Mathematicae.

\item[] Irish Math. Soc. Bull. - Irish Mathematical Society Bulletin.

\item[] Israel J. Math. - Israel Journal of Mathematics.

\item[] Izv. Akad. Nauk SSSR Ser. Mat. - Izvestija Akademii Nauk SSSR. Serija Mathmati\v{c}eskaja. (Translated as: Math. USSR-Izv.)

\item[] J. Austral. Math. Soc. Ser. A. - Journal of the Australian Mathematical Society. Series A.

\item[] J. Austral. Math. Soc. Ser. B. - Journal of the Australian Mathematical Society. Series B.

\item[] J. Combinatorial Theory Ser. A. - Journal of Combinatorial Theory. Series A.

\item[] J. Combinatorial Theory Ser. B. - Journal of Combinatorial Theory. Series B.

\item[] J. London Math. Soc. Ser. 2. - The Journal of the London Mathematical Society. Series 2.

\item[] J. Math. Pures Appl. - Journal de Math\'{e}matiques Pures et Appliqu\'{e}es.

\item[] J. Pure Appl. Algebra. - Journal of Pure and Applied Algebra.

\item[] J. Reine Angew. Math. - Journal f\"{u}r die Reine und Angewandte Mathematik.

\item[] J. Symbolic Logic. - The Journal of Symbolic Logic.

\item[] J. Univ. Kuwait Sci. - The Journal of the University of Kuwait.

\item[] Jber. Deutsch. Math.-Verein. - Jahresbericht der Deutschen Mathematiker-Vereinigung.

\item[] K\"{o}niglich S\"{a}chsischen Gesellschaft der Wissenschaften zu Leipzig, Math.-Phys. Klasse, Sitzungsberichte. - K\"{o}niglich S\"{a}chsischen Gesellschaft der Wissenschaften zu Leipzig, Math.-Phys. Klasse, Sitzungsberichte

\item[] Kyungpook Math. J. - Kyungpook Mathematical Journal. Kyungpook National University.

\item[] Logique et Analyse. - Logique et Analyse

\item[] Manuscripta Math. - Manuscripta Mathematica.

\item[] Mat. Zametki - Akademija Nauk Sojuza SSR. Matemati\v{c}eskie Zametki. (Translated as: Math. Notes)

\item[] Matematiche (Catania). - Le Matematiche.

\item[] Math. Ann. - Mathematische Annalen.

\item[] Math. Japon. - Mathematica Japonica.

\item[] Math. Logic Quart. - Mathematical Logic Quarterly. (Formerly Zeitschrift f\"{u}r Mathematische Logik und Grundlagen der Mathematik.)

\item[] Math. Nachr. - Mathematische Nachrichten.

\item[] Math. Proc. Cambridge Philos. Soc. - Mathematical Proceedings of the Cambridge Philosophical Society.

\item[] Math. Sem. Notes Kobe Univ. - Kobe University. Mathematics Seminar Notes.

\item[] Math. Z. - Mathematische Zeitschrift.

\item[] Monatsh. Math. - Monatshefte f\"{u}r Mathematik.

\item[] Notices Amer. Math. Soc. - Notices of the American Mathematical Society.

\item[] Notre Dame J. Formal Logic. - Notre Dame Journal of Formal Logic.

\item[] \"{O}sterreich. Akad. Wiss. Math.-Natur. Kl. S.-B II. - \"{O}sterreichische Akademie der Wissenschaften. Mathematisch-Naturwissenschaftliche Klasse. Sitzungsberichte. Abteilung II. Mathematik, Astronomie, Physik, Meteorologie und Technik.

\item[] Pacific J. Math. - Pacific Journal of Mathematics.

\item[] Period. Math. Hungar. - Periodica Mathematica Hungarica. Journal of the J\'{a}nos Bolyai Mathematical Society.

\item[] Proc. Amer. Math. Soc. - Proceedings of the American Mathematical Society.

\item[] Proc. Nat. Acad. Sci. U.S.A. - Proceedings of the National Academy of Sciences of the United States of America.

\item[] Proc. London Math. Soc. Ser. 2 - Proceedings of the London Mathematical Society. Series 2.

\item[] Proc. London Math. Soc. Ser. 3. - Proceedings of the London Mathematical Society. Series 3.

\item[] Publ. Inst. Math. (Beograd) - Institut Math\'{e}matique. Publications. Nouvelle S\'{e}rie.

\item[] Publ. Sci. Univ. Alger Ser. A Math. - Universit\'{e} d'Alger. Publications Scientifiques. Series A. Mathematiques.

\item[] Quart. J. Math. Oxford Ser. - The Quarterly Journal of Mathematics. Oxford. Second Series.

\item[] Real Anal. Exchange. - Real Analysis Exchange.

\item[] Rend. Circ. Mat. Palermo. - Rendiconti del Circolo Matematico di Palermo.

\item[] Rend. Sem. Mat. Univ. Padova. - Rendiconti del Seminario Matematico dell'Universit\`{a} di Padova.

\item[] Rep. Math. Logic. - Reports on Mathematical Logic.

\item[] Sci. Rep. Tokyo Kyoiku Daigaku Sect. A - Science Reports of the Tokyo Kyoiku Daigaku (Tokyo University of Education). Section A.

\item[] Simon Stevin. - Simon Stevin.

\item[] Sitzungsber. Preussichen Akad. Wiss., Math. Kl. - Die Preussichen Akademie der Wissenschaften. Sitzungsberichte. Physikalisch-Mathematische Klasse.

\item[] Soob\v{s}\v{c}. Akak. Nauk. Gruzin SSR - Soob\v{s}\v{c}enija Akademii Nauk Gruzinsko\v{i} SSR.

\item[] Studia Math. - Polska Akademia Nauk. Institut Matematyczny. Studia Mathematica.

\item[] Topology and its Applications. - Topology and its Applications.

\item[] Trans. Amer. Math. Soc. - Transactions of the American Mathematical Society.

\item[] Z. Math. Logik Grundlagen Math. - Zeitschrift f\"{u}r Mathematische Logik und Grundlagen der Mathematik.

\item[] Z. Philos. Philos. Kritik. - Zeitschrift f\"{u}r Philosophic und philosophische Kritik.
\end{enumerate}

%%%%%%%ref

\section*{References}

\begin{thebibliography}{[\textbf{1968}]}
\item[]  \textbf{Abian, A.}

\bibitem[\textbf{1968}]{bib:1968} \emph{A fixed point theorem}, Nieuw Arch. Wisk. (3) \textbf{16}, 184--185.

\bibitem[\textbf{1972}]{bib:1968a} \emph{On solvability of infinite systems of polynomial equations over a finite ring}, Simon Stevin \textbf{46}/73), 33--37.

\bibitem[\textbf{1980}]{bib:1968b} \emph{A fundamental fixed point theorem revisited}, Bull. Greek Math. Soc. \textbf{21}, 12 and 4-11.

\bibitem[\textbf{1985}]{bib:1985} \emph{A fixed point theorem equivalent to the axiom of choice}, Arch. Math. Logik Grundlagenforschung \textbf{25}, 173-174.

\item[]  \textbf{Addison, J. W./De Bouvere, Karel/Pitt, W.B.}

\bibitem[\textbf{1965}]{bib:1965} \emph{A bibliography of the theory of models}, in The Theory of Models, J. Addison, L. Henkin and A. Tarski ed., North-Holland, Amsterdam, 442-492.

\item[]  \textbf{Alexander, J. W.}

\bibitem[\textbf{1939}]{bib:1939} \emph{Ordered sets, complexes, and the problem of bicompactification} Proc. Nat. Acad. Sci. U.S.A. \textbf{22}, 296-298.

\item[]  \textbf{Andreka, H./Nemeti, I.}

\bibitem[\textbf{1980}]{bib:1980} \emph{Does SPK $\supseteq$ PSK imply the axiom of choice?}, Comment. Math. Univ. Carolinae \textbf{21}, 699-706.

\bibitem[\textbf{1981}]{bib:1981} \emph{HSP K is an equational class without the axiom of choice}, Algebra Universalis \textbf{13}, 164-166.

\item[]  \textbf{Apter, A.W.}

\bibitem[\textbf{1985a}]{bib:1985a} \emph{An AD-like model}, J. Symbolic Logic \textbf{50}, 531-543.

\bibitem[\textbf{1985b}]{bib:1985b} \emph{Successors of singular cardinals and measurability}, Advances in Math. \textbf{55}, 228-241.

\bibitem[\textbf{1988}]{bib:1988} \emph{On a problem inspired by determinancy}, Israel J. Math. \textbf{61}, 256-270.

\bibitem[\textbf{1990}]{bib:1990} \emph{Successors of singular cardinals and measurability revisited}, J. Symbolic Logic \textbf{55}, 492-501.

\bibitem[\textbf{1992}]{bib:1992} \emph{On the class of measurable cardinals without the axiom of choice}, Israel J. Math. \textbf{79}, 367-379.

\item[]  \textbf{Apter, A.W./Henle, J.M.}

\bibitem[\textbf{1991}]{bib:1991} \emph{Relative consistency results via strong compactness}, Fund. Math. \textbf{139}, 133-149.

\item[]  \textbf{Armbrust, M.}

\bibitem[\textbf{1972}]{bib:1972a} \emph{An algebraic equivalent of the multiple choice axiom}, Fund. Math. \textbf{74}, 145--146.

\item[]  \textbf{Ash, L. J.}

\bibitem[\textbf{1975}]{bib:1975} \emph{A consequence of the axiom of choice}, J. Austral. Math. Soc. Ser. A \textbf{19}, 306-308.

\item[]  \textbf{Bacsich, P. D.}

\bibitem[\textbf{1972}]{bib:1972b} \emph{Extension of Boolean homomorphisms with bounding semi-morphisms}, J. Reine Angew. Math. \textbf{253}, 24--27.

\item[]  \textbf{Baker, I. N.}

\bibitem[\textbf{1964}]{bib:1964} Solution to problem 5077 Amer. Math. Monthly \textbf{71}, 219-220.

\item[]  \textbf{Baldwin, J.T./Blass, A.}

\bibitem[\textbf{1974}]{bib:1974} \emph{An axiomatic approach to rank in model theory}, Ann. Math. Logic \textbf{7}, 295-324.

\item[]  \textbf{Banach, S. /Tarski, A.}

\bibitem[\textbf{1924}]{bib:1924} \emph{Sur la decomposition des ensembles de points en parties respectivement congruents}, Fund. Math. \textbf{6}, 244--257.

\item[]  \textbf{Banaschewski, B.}

\bibitem[\textbf{1981}]{bib:1981b} \emph{Coherent frames, in Continuous Lattices}, Lecture Notes in Mathematics, \textbf{871}, Springer-Verlag, Berlin, 1-11.

\bibitem[\textbf{1983}]{bib:1983} \emph{The power of the ultra-filter theorem}, J. London Math. Soc. Ser. 2 \textbf{27}, 193-202. Mr84f\#03043, ZBL 523\#03037.

\bibitem[\textbf{1985}]{bib:1985c} \emph{Prime elements from prime ideals}, Order \textbf{2}, 211-213.

\bibitem[\textbf{1988}]{bib:1988a} \emph{Compact regular frames and the Sikorski theorem}, Kyungpook Math. J. \textbf{28}, 1--14.

\bibitem[\textbf{1992}]{bib:1992a} \emph{Algebraic closures without choice}, Z. Math. Logik Grundlagen Math. \textbf{38}, 383--385.

\bibitem[\textbf{1997}]{bib:1997} \emph{Choice principles and compactness conditions}, preprint, May.

\item[]  \textbf{Banaschewski, B./Harting, R.}

\bibitem[\textbf{1985}]{bib:1985d} \emph{Lattice aspects of radical ideals and choice principles}, Proc. London Math. Soc. Ser. 3 \textbf{50}, 385-404.

\item[]  \textbf{Banaschewski, B./Moore, G.H.}

\bibitem[\textbf{1990}]{bib:1990a} \emph{The dual Cantor-Bernstein theorem and the partition principle}, Notre Dame J. Formal Logic \textbf{31}, 375-381.

\item[]  \textbf{Bell, J.L.}

\bibitem[\textbf{1977}]{bib:1977} \emph{Boolean valued models and independence proofs in set theory}, Oxford Logic Guides, Oxford Univ. Press, Oxford.

\bibitem[\textbf{1983}]{bib:1983a} \emph{On the strength of the Sikorski extension theorem for Boolean algebras}, J. Symbolic Logic \textbf{48}, 841-845.

\bibitem[\textbf{1988}]{bib:1988b} \emph{Some propositions equivalent to the Sikorski extension theorem for Boolean algebras}, Fund. Math. \textbf{130}, 51--55.

\item[]  \textbf{Bell, J. L./Fremlin, D. H.}

\bibitem[\textbf{1972}]{bib:1972c} \emph{The maximal ideal theorem for lattices of sets}, Bull. London Math. Soc. \textbf{4}, 1--2.

\item[]  \textbf{Bell, J. L./Jellett, F.}

\bibitem[\textbf{1971}]{bib:1971} \emph{On the relationship between the Boolean prime ideal theorem and two principles in functional analysis}, Bull. Acad. Polon. Sci. S\'{e}r. Sci. Math. Astronom. Phys. \textbf{19}, 191--194.

\item[]  \textbf{B\'{e}n\'{e}jam, J. P.}

\bibitem[\textbf{1970}]{bib:1970} \emph{Un remarque sur des th\'{e}or\`{e}mes de Rado et Fra\"{i}ss\'{e}}, C. R. Acad. Sci. Paris S\'{e}r. A--B \textbf{270}, A1656--A1658.

\item[]  \textbf{van Benthem, J. F. A. K.}

\bibitem[\textbf{1975}]{bib:1975a} \emph{A set theoretical equivalent for the prime ideal theorem for Boolean algebras}, Fund. Math. \textbf{89}, 151-153.

\item[]  \textbf{Bernstein, Felix,}

\bibitem[\textbf{1908}]{bib:1908} \emph{Zur Theorie der trigonometrischen Reihe}, K\"{o}niglich S\"{a}chsischen Gesellschaft der Wissenschaften zu Leipzig, Math.-Phys. Klasse, Sitzungsberichte, \textbf{60}, 325-338.

\item[]  \textbf{Birkhoff, G.}

\bibitem[\textbf{1944}]{bib:1944} \emph{Subdirect unions in universal algebra}, Bull. Amer. Math. Soc. \textbf{50}, 764-768.

\bibitem[\textbf{1967}]{bib:1967} \emph{Lattice Theory}, third edition, AMS Collog. Publications \textbf{25}, AMS, Providence, R.I..

\item[]  \textbf{Blair, C.E.}

\bibitem[\textbf{1977}]{bib:1977a} \emph{The Baire Category Theorem implies the principle of dependent choices}, Bull. Acad. Polon. Sci. S\'{e}r. Math. Astronom. Phys. \textbf{25}, 933-934.

\item[]  \textbf{Blass, A.}

\bibitem[\textbf{1975}]{bib:1975b} \emph{A forcing proof of the Kechris-Moschovakis constructibility theorem}, Proc. Amer. Math. Soc. \textbf{47}, 195-197.

\bibitem[\textbf{1977a}]{bib:1977aa} \emph{Ramsey's theorem in the hierarchy of choice principles}, J. Symbolic Logic \textbf{42}, 387-390.

\bibitem[\textbf{1977b}]{bib:1977b} \emph{A model without ultrafilters}, Bull. Acad. Polon. Sci. S\'{e}r. Math. Astronom. Phys. \textbf{25}, 329-331.

\bibitem[\textbf{1979}]{bib:1979} \emph{Injectivity, projectivity and the axiom of choice}, Trans. Amer. Math. Soc. \textbf{255}, 31-59.

\bibitem[\textbf{1981}]{bib:1981c} \emph{The model of set theory generated by countably many generic reals}, J. Symbolic Logic \textbf{46}, 732-752.

\bibitem[\textbf{1983a}]{bib:1983b} \emph{Cohomology detects failures of the axiom of choice}, Trans. Amer. Math. Soc. \textbf{279}, 257-269.

\bibitem[\textbf{1983b}]{bib:1983c} \emph{Words, free algebras and coequalizers}, Fund. Math. \textbf{117}, 117-160.

\bibitem[\textbf{1984a}]{bib:1984a} \emph{Existence of a basis implies the axiom of choice}, Axiomatic Set Theory, J. Baumgartner, D.A. Martin and S. Shelah ed., Contemporary Mathematics \textbf{31}, AMS, Providence, 31-33. MR 86a:04001, ZB1 557\#18001.

\bibitem[\textbf{1984b}]{bib:1984b} \emph{The interaction between category theory and set theory}, Contemporary Mathematics \textbf{30}, Amer. Math. Soc., Providence, 5-29.

\bibitem[\textbf{1986}]{bib:1986} \emph{Prime ideals yield almost maximal ideals}, Fund. Math. \textbf{127}, 57-66.

\item[]  \textbf{Bleicher, M.}

\bibitem[\textbf{1964}]{bib:1964a} \emph{Some theorems on vector spaces and the axiom of choice}, Fund. Math. \textbf{54}, 95-107.

\bibitem[\textbf{1965}]{bib:1965c} \emph{Multiple choice axioms and the axiom of choice for finite sets}, Fund. Math. \textbf{57}, 247-252.

\item[]  \textbf{Borel, E.}

\bibitem[\textbf{1898}]{bib:1898} \emph{Le\c{c}ons sur la th\'{e}orie des fonctions}, Gauthier-Villars, Paris.

\item[]  \textbf{Brunner, N.}

\bibitem[\textbf{1980}]{bib:1980c} \emph{Beitr\"{a}ge zur modelltheoretischen Topologie}, Anz. \"{O}sterr. Akad. Wiss., Math.-Nat. Kl. \textbf{117}, 168-172.

\bibitem[\textbf{1981a}]{bib:1981a} \emph{Boxprodukte und Auswahlaxiom}, Anz. \"{O}sterreich. Akad. Wiss. Math.-Nat. Kl. \textbf{118}, 167-169.

\bibitem[\textbf{1981b}]{bib:1981b1} \emph{Topologie und offene Mengen}, \"{O}sterreich. Akad. Wiss. Math.-Natur. Kl. S.-B II \textbf{190}, 317-325.

\bibitem[\textbf{1982a}]{bib:1982a} \emph{Dedekind-Endlichkeit und Wohlordenbarkeit}, Monatsh. Math. \textbf{94}, 9-31.

\bibitem[\textbf{1982b}]{bib:1982b} $\sigma$-\emph{kompakte R\"{a}ume}, Manuscripta Math. \textbf{38}, 375-379.

\bibitem[\textbf{1982c}]{bib:1982c} \emph{Sequential continuity}, Kyungpook Math. J. \textbf{22}, 233 pp.

\bibitem[\textbf{1982d}]{bib:1982d} \emph{Lindel\"{o}f Raume und Auswahlaxiom (English Summary)}, Ariz. \"{O}sterreich. Akad. Wiss. Math.-Nat. Kl. \textbf{119}, 161-165.

\bibitem[\textbf{1983a}]{bib:1983d} \emph{Geordnete L\"{a}uchli Kontinuen}, Fund. Math. \textbf{117}, 67-73.

\bibitem[\textbf{1983b}]{bib:1983e} \emph{Sequential compactness and the axiom of choice}, Notre Dame J. Formal Logic \textbf{24}, 89-92.

\bibitem[\textbf{1983c}]{bib:1983f} \emph{Kategories\"{a}tze und multiples Auswahlaxiom}, Z. Math. Logik Grundlagen Math. \textbf{29}, 435-443.

\bibitem[\textbf{1983d}]{bib:1983g} \emph{The axiom of choice in topology}, Notre Dame J. Formal Logic \textbf{24}, 305-317.

\bibitem[\textbf{1984a}]{bib:1984a1} \emph{Hilbertr\"{a}ume mit amorphen Basen (English summary)}, Compositio Math. \textbf{52}, 381-387. MR 85i: 03113, ZBl 554\#47001.

\bibitem[\textbf{1984b}]{bib:1984b1} \emph{Amorphe Potenzen kompakter R\"{a}ume}, Arch. Math. Logik Grundlagenforsch \textbf{24}, 119-135. MR 86k:03043, ZBl 553\#47001.

\bibitem[\textbf{1984c}]{bib:1984c} \emph{P--Ra\"{u}me und Auswahlaxiom}, Rend. Circ. Mat. Palermo, Ser. 2 \textbf{33}, 34-36. MR 85f:03045, ZBl 538\#54004.

\bibitem[\textbf{1984d}]{bib:1984d} \emph{Positive functional and the axiom of choice}, Rend. Sem. Mat. Univ. Padova \textbf{72}, 9-12. MR 86f:03080, ZBl 554\#46007.

\bibitem[\textbf{1984e}]{bib:1984e} \emph{Lattice points in metric spaces}, Rep. Math. Logic \textbf{18}, 45-46. ZBl 557\#03031.

\bibitem[\textbf{1984f}]{bib:1984f} \emph{Realisierung und Auswahlaxiom}, Arch. Math. (Brno) \textbf{20}, 39-42. MR 86m:04005, ZBl 551 \#54004.

\bibitem[\textbf{1985a}]{bib:1985e} \emph{Products of compact spaces in the least permutation model}, Z. Math. Logik Grundlagen Math. \textbf{31}, 441-448.

\bibitem[\textbf{1985b}]{bib:1985f} \emph{Well-ordering theorems in topology}, Foundations of Logic and Linguistics. Problems and Their solutions, G. Dorn and P. Weingartner ed. Plenum Publishing, New York, 281-289.

\bibitem[\textbf{1985c}]{bib:1985g} \emph{Spaces of urelements}, Rend. Sem. Mat. Univ. Padova, \textbf{74}, 7-13.

\bibitem[\textbf{1986}]{bib:1986d} \emph{Ultraproducts and the axiom of choice}, Arch. Math. (Brno) \textbf{22} 175-180

\bibitem[\textbf{1987a}]{bib:1987a} \emph{Topologische Maximalprinzipien}, Z. Math. Logik Grundlagen Math. \textbf{33}, 135-139.

\bibitem[\textbf{1987b}]{bib:1987b} \emph{Spaces of urelements. II}, Rend. Sem. Mat. Univ. Padova, \textbf{77}, 305-315.

\bibitem[\textbf{1989}]{bib:1989} \emph{Set mappings on Dedekind sets}, Notre Dame J. Formal Logic \textbf{30}, 268-270.

\bibitem[\textbf{1990}]{bib:1990b} \emph{The Fraenkel-Mostowski method, revisited}, Notre Dame J. Formal Logic \textbf{31}, 64-75.

\item[]  \textbf{Brunner, N./Howard, P.}

\bibitem[\textbf{1992}]{bib:1992b} \emph{Russell's alternative to the axiom of choice}, Z. Math. Logik Grundlagen Math. \textbf{38} 529-534.

\item[]  \textbf{Brunner, N./Howard, P. /Rubin, J.}

\bibitem[\textbf{1996}]{bib:1996} \emph{Choice principles in W\c{e}glorz' models} Accepted Fund. Math. July.

\item[]  \textbf{Brunner, N./Rubin, A./Rubin, J.}

\bibitem[\textbf{1985}]{bib:1985h} \emph{Maximal principles for sets and classes}, \"{O}sterreich. Akad. Wiss. Math.-Natur. Kl. S.-B II \textbf{194}, 203-214.

\item[]  \textbf{Brunner, N./Rubin, J.}

\bibitem[\textbf{1986}]{bib:1986e} \emph{Permutation models and topological groups}, Rend. Sem. Mat. Univ. Padova. \textbf{76}, 149-161.

\item[]  \textbf{Bull Jr., E. L.}

\bibitem[\textbf{1978}]{bib:1978} \emph{Successive large cardinals}, Ann. Math. Logic \textbf{15}, 161-191.

\item[]  \textbf{Buskes, G.}

\bibitem[\textbf{1993}]{bib:1993} \emph{The Hahn-Banach Theorem surveyed}, Dissertationes Math. (Rozprawy Mat.) \textbf{327}, 1-49.

\item[]  \textbf{Cantor, G.}

\bibitem[\textbf{1878}]{bib:1878} \emph{Ein Beitrag zur Mannigfaltigkeitslehre}, J. Reine Angew. Math. \textbf{84}, 242-258.

\bibitem[\textbf{1883}]{bib:1883} \emph{\"{U}ber unendliche, lineare Punktmannichfaltigkeiten V}, Math. Ann. \textbf{21}, 545-591.

\bibitem[\textbf{1887}]{bib:1887} \emph{Mitteilungen zur Lehre vom Transfiniten}, Z. Philos. Philos. Kritik. \textbf{91}, 81-125, 252-270.

\bibitem[\textbf{1895}]{bib:1895} \emph{Beitr\"{a}ge zur Begrundung der transfinitem Megenlehre I} Math. Ann. \textbf{46} 207-246; translated by P. Jourdain in 1915.

\item[]  \textbf{Caristi, J.}

\bibitem[\textbf{1976}]{bib:1976} \emph{Fixed point theorems for mappings satisfying inwardness conditions}, Trans. Amer. Math. Soc. \textbf{215}, 241-251.

\item[]  \textbf{Chandler, R. E.}

\bibitem[\textbf{1972}]{bib:1972d} \emph{An alternative construction of} $\beta x$ \emph{and} $\mu x$, Proc. Amer. Math. Soc. \textbf{32}, 315--318.

\item[]  \textbf{Ciesielski, K./Guzicki, W.}

\bibitem[\textbf{1989}]{bib:1989c} \emph{Generic families and models of set theory with the axiom of choice}, Proc. Amer. Math. Soc. \textbf{106}, 199-206.

\item[]  \textbf{Cohen, P. J.}

\bibitem[\textbf{1963}]{bib:1963} \emph{The independence of the continuum hypothesis} I \& II, Proc. Nat. Acad. Sci. U.S.A. \textbf{50}, 1143-1148 \& \textbf{51} (1964), 105-110.

\bibitem[\textbf{1965}]{bib:1965d} \emph{Independence results in set theory}, Theory of Models; J. Addison, L. Henkin and A. Tarski ed., Studies in Logic and the Foundations of Mathematics, North-Holland, Amsterdam, 39-54.

\bibitem[\textbf{1966}]{bib:1966} \emph{Set Theory and the Continuum Hypothesis}, Benjamin, New York.

\bibitem[\textbf{1971}]{bib:1971c} \emph{Automorphisms of set theory}, Tarski Symposium; Berkeley; L. Henkin, J. Addison, C.C. Chang, W. Craig, D.S. Scott, and R. Vaught \emph{ed}, Proceedings of Symposia in Pure Mathematics \textbf{25}, AMS, 1974, 325-330.

\item[]  \textbf{Comfort, W. W.}

\bibitem[\textbf{1968}]{bib:1968c} \emph{A theorem of Stone-\v{C}ech type, and a theorem of Tychonoff type, without the axiom of choice, and their real compact analogues}, Fund. Math. \textbf{63}, 97-110.

\item[]  \textbf{Conway, J. H.}

\bibitem[\textbf{1971}]{bib:1971d} \emph{Effective implications between the finite choice axioms}, Cambridge Summer School in Mathematical Logic (\emph{Lecture Notes in Mathematics}, \textbf{337}, Springer-Verlag, Berlin, 1973), 439--458.

\item[]  \textbf{Cowen, R.H.}

\bibitem[\textbf{1973}]{bib:1973} \emph{Some combinatorial theorems equivalent to the prime ideal theorem}, Proc. Amer. Math. Soc. \textbf{41}, 268-273.

\bibitem[\textbf{1977a}]{bib:1977a1} \emph{Binary consistent choice on triples}, Notre Dame J. Formal Logic \textbf{18}, 310-312.

\bibitem[\textbf{1977b}]{bib:1977b1} \emph{Generalizing K\"{o}nigs infinity lemma}, Notrd Dame J. Formal Logic, \textbf{18}, 243-247.

\bibitem[\textbf{1982}]{bib:1982} \emph{Partition principles for properties of finite character}, Rep. Math. Logic \textbf{14}, 23-28.

\bibitem[\textbf{1983}]{bib:1983h} \emph{Compactness via prime semilattices}, Notre Dame J. Formal Logic \textbf{24}, 199-204.

\bibitem[\textbf{1990}]{bib:1990c} \emph{Two hypergraph theorems equivalent to BPI}, Notre Dame J. Formal Logic \textbf{31}, 232--240.

\item[]  \textbf{Cowen, R. H./Emerson, W.}

\bibitem[\textbf{1995}]{bib:1995} \emph{A compactness theorem for linear equations}, preprint, April 17.

\item[]  \textbf{Dawson, J. W.}

\bibitem[\textbf{1973}]{bib:1973c} \emph{Ordinal definability in the rank hierarchy}, Ann. Math. Logic \textbf{6}, 1-39.

\bibitem[\textbf{1984}]{bib:1984} \emph{Discussion of the Foundations of Mathematics}, translated and edited by J. Dawson, History and Philosophy of Logic \textbf{5}, 111-129.

\item[]  \textbf{Dawson, J. W./Howard, P.}

\bibitem[\textbf{1976}]{bib:1976c} \emph{Factorials of infinite cardinals}, Fund. Math. \textbf{93}, 185-195.

\item[]  \textbf{David, R.}

\bibitem[\textbf{1975}]{bib:1975c} \emph{Un r\'{e}sultat de consistence relative concernant l'\'{e}nonc\'{e}} $V = WO^{\delta}$, C. R. Acad. Sci. Paris S\'{e}r. A-B \textbf{280}, A981-A983.

\bibitem[\textbf{1980}]{bib:1980d} \emph{A result of relative consistency about the predicate} $WO(\delta$, \emph{x}), J. Symbolic Logic \textbf{45}, 483-492.

\item[]  \textbf{Davies, R. O.}

\bibitem[\textbf{1978}]{bib:1978c} \emph{The plane is the union of three rectilinearly accessible sets}, Real Anal. Exchange \textbf{4}, 97-101.

\item[]  \textbf{Day, M. M.}

\bibitem[\textbf{1973}]{bib:1973d} \emph{Normed Linear Spaces}, Ergebnisse der mathematik und ihrer Grenzgebiete \textbf{21}, Springer Verlag, Berlin.

\item[]  \textbf{Dedekind, R.}

\bibitem[\textbf{1888}]{bib:1888} \emph{Was sind und was sollen die Zahlen?}, Vieweg, Braunschweig.

\item[]  \textbf{Degen, J.W.}

\bibitem[\textbf{1988}]{bib:1988c} \emph{There can be a permutation which is not the product of two reflections}, Z. Math. Logik Grundlagen Math. \textbf{34}, 65-66.

\item[]  \textbf{De la Cruz, O./Di Prisco, C. A.}

\bibitem[\textbf{1996}]{bib:1996c} \emph{Weak choice principles}, Accepted Proc. Amer. Math. Soc. August.

\item[]  \textbf{Derrick, J./Drake, F.}

\bibitem[\textbf{1967a}]{bib:1967a} \emph{Independence of the axiom of choice from variants of the generalized continuum hypothesis}, Proc. Leicester Colloq. Math. Logic, 75-84.

\bibitem[\textbf{1967b}]{bib:1967b} \emph{A model of Z-F set theory with} $\aleph(2^{\omega})=\aleph_{\omega}$, J. Symbolic Logic \textbf{32} 559.

\item[]  \textbf{Diel, J.H.}

\bibitem[\textbf{1974}]{bib:1974c} \emph{Two definitions of finiteness}, Notices Amer. Math. Soc. \textbf{21}, 554-555.

\item[]  \textbf{Diener, Karl-Heinz}

\bibitem[\textbf{1989}]{bib:1989d} \emph{Two new algebraic equivalents to the axiom of choice}, Algebra Universalis \textbf{26}, 189--195.

\bibitem[\textbf{1994}]{bib:1994} \emph{A remark on ascending chain conditions, the countable axiom of choice, and the principle of dependent choices}, Math. Logic Quart. \textbf{40} 415-421.

\item[]  \textbf{van Douwen, E.K.}

\bibitem[\textbf{1980}]{bib:1980e} \emph{Covering and separation properties of box products}, in G. M. Reed, ed., Surveys in General Topology, Academic Press, New York, 55-129.

\bibitem[\textbf{1985}]{bib:1985i} \emph{Horrors of topology without AC}: \emph{A non-normal orderable space}, Proc. Amer. Math. Soc. \textbf{95}, 101-105.

\item[]  \textbf{Drake, F. R.}

\bibitem[\textbf{1974}]{bib:1974d} \emph{Set Theory, An Introduction to Large Cardinals}, North Holland, Amsterdam.

\item[]  \textbf{Dunford, N./Schwartz, J.T.}

\bibitem[\textbf{1957}]{bib:1957} \emph{Linear Operators I}, Interscience Publishers, New York.

\item[]  \textbf{Easton, W. B.}

\bibitem[\textbf{1970}]{bib:1970c} \emph{Powers of regular cardinals}, Ann. Math. Logic \textbf{1}, 139-178.

\item[]  \textbf{Ellentuck, E.M.}

\bibitem[\textbf{1962}]{bib:1962} \emph{The theory of Dedekind finite cardinals}, Doc. dissertation, U. of California, Berkeley.

\bibitem[\textbf{1965}]{bib:1965e} The universal properties of Dedekind finite cardinals, Ann. Math. Ser. 2 \textbf{82}, 225-248.

\bibitem[\textbf{1970}]{bib:1970d} \emph{Extension methods in cardinal arithmetic}, Trans. Amer. Math. Soc. \textbf{149}, 307--325.

\bibitem[\textbf{1973}]{bib:1973e} \emph{The structure of Dedekind cardinals}, Trans. Amer. Math. Soc. \textbf{180}, 109--125.

\bibitem[\textbf{1974}]{bib:1974f} \emph{A model of arithmetic}, Bull. Acad. Polon. Sci. S\'{e}r. Sci. Math. Astronom. Phys. \textbf{22}, 353--355.

\bibitem[\textbf{1976}]{bib:1976d} \emph{Combinatorial functions and indecomposable cardinals}, Proc. Amer. Math. Soc. \textbf{57}, 311--317.

\item[]  \textbf{Erd\"{o}s, P., /Hajnal, A., /M\'{a}t\'{e}, A. /Rado, R.}

\bibitem[\textbf{1984}]{bib:1984k} \emph{Combinatorial Set Theory: Partitions Relations for Cardinals}, Studies in Logic \textbf{106}, North Holland, Amsterdam.

\item[]  \textbf{Erd\"{o}s, P./Rado, R.}

\bibitem[\textbf{1952}]{bib:1952} \emph{Combinatorial theorems on classification of subsets of a given set}, Proc. London Math. Soc. \textbf{2}, 417--439.

\item[]  \textbf{Esenin-Vol'pin, A.S.}

\bibitem[\textbf{1963}]{bib:1963c} \emph{Unprovability of Souslin's hypothesis without the aid of the axiom of choice in the Bernays-Mostowski axiom system}, (translated from Russian), Amer. Math. Soc. Transl. Ser. 2 \textbf{23}, 83-87.

\item[]  \textbf{Federer, H./Jonsson, B.}

\bibitem[\textbf{1950}]{bib:1950} \emph{Some properties of free groups}, Trans. Amer. Math. Soc. \textbf{68}, 1--27.

\item[]  \textbf{Feferman, S.}

\bibitem[\textbf{1964}]{bib:1964b} \emph{Independence of the axiom of choice from the axiom of dependent choices}, J. Symbolic Logic \textbf{29}, 226.

\bibitem[\textbf{1965}]{bib:1965f} \emph{Applications of forcing and generic sets}, Fund. Math. \textbf{56}, 325-345.

\item[]  \textbf{Feferman, S./Levy, A.}

\bibitem[\textbf{1963}]{bib:1963d} \emph{Independence results in set theory by Cohen's method II}, Notices Amer. Math. Soc. \textbf{10}, 593.

\item[]  \textbf{Felgner, U.}

\bibitem[\textbf{1969}]{bib:1969} \emph{Die Existenz wohlgenordneter konfinaler Teilmengen in Ketten und das Auswahlaxiom}, Math. Z. \textbf{111}, 221-232.

\bibitem[\textbf{1971a}]{bib:1971a} \emph{Models of ZF Set Theory}, Springer-Verlag, Berlin.

\bibitem[\textbf{1971b}]{bib:1971b} \emph{\"{U}ber das Ordnungstheorem}, Z. Math. Logik Grundlagen Math. \textbf{17}, 257-272.

\bibitem[\textbf{1971c}]{bib:1971c1} \emph{Comparison of the axioms of local and universal choice}, Fund. Math. \textbf{81}, 43-62.

\bibitem[\textbf{1972}]{bib:1972e} \emph{Independence of the Boolean prime ideal theorem from the order extension principle}, Habilitationsschrift, Heidelberg.

\bibitem[\textbf{1974}]{bib:1974g} \emph{Abz\"{a}hlbarkeit und Wohlordenbarkeit}, Comment. Math. Helv. \textbf{49}, 114-124.

\bibitem[\textbf{1975}]{bib:1975d} \emph{On} $\aleph_{0}$-\emph{categorical extra-special} $p$-\emph{groups}, Logique et Analyse \textbf{71-71}, 407-428.

\bibitem[\textbf{1976}]{bib:1976e} \emph{Choice functions on sets and classes}, Sets and Classes (Studies in Logic and the Foundations of Mathematics, North Holland, Amsterdam, 217-255.

\item[]  \textbf{Felgner, U./Flannagan, T. B.}

\bibitem[\textbf{1978}]{bib:1978d} \emph{Well ordered subclasses of proper classes}, in Higher Set Theory (Lecture notes in Mathematics \textbf{699}, Springer-Verlag, Berlin, 1-14.

\item[]  \textbf{Felgner, U./Jech, T.}

\bibitem[\textbf{1973}]{bib:1973f} \emph{Variants of the axiom of choice in set theory with atoms}, Fund. Math. \textbf{79}, 79-85.

\item[]  \textbf{Felgner, U./Truss, J. K.}

\bibitem[\textbf{1996}]{bib:1996d} \emph{The independence of the prime ideal theorem from the order-extension principle}, preprint December.

\item[]  \textbf{Figura, Adam}

\bibitem[\textbf{1977}]{bib:1977e} \emph{Collapsing functions and dependent choices}, Bull. Acad. Polon. Sci. S\'{e}r. Sci. Math. Astronom. Phys. \textbf{25}, 99-103.

\bibitem[\textbf{1981}]{bib:1981f} \emph{On some properties of cardinal numbers}, in Open days in Model Theory and Set Theory (Proceedings of a conference held in Sept. at Jadwisin near Warsaw, Poland) Ed. W. Guzicki, W. Marek, A. Pelc and C. Rauszer.

\item[]  \textbf{Flannagan, T. B.}

\bibitem[\textbf{1974}]{bib:1974h} \emph{Axioms of choice in Morse-Kelley class theory}, Logic Conference in Kiel (Lecture Notes in Math. \textbf{499}, Springer-Verlag, Berlin, 1975), 190-247.

\item[]  \textbf{Fleischer, I.}

\bibitem[\textbf{1973}]{bib:1973g} \emph{Das Zornsche Lemma f\"{u}r Klassen}, Z. Math. Logik Grundlagen Math. \textbf{19}, 205-206.

\item[]  \textbf{Foreman, M./Wehrung, F.}

\bibitem[\textbf{1991}]{bib:1991c} \emph{The Hahn-Banach theorem implies the existence of a non-Lebesgue measurable set}, Fund. Math. \textbf{138}, 13-19.

\item[]  \textbf{Forti, M./Furio, H.}

\bibitem[\textbf{1985}]{bib:1985j} \emph{A model where cardinal ordering is universal}, Z. Math. Logik Grundlagen Math. \textbf{31}, 533-536.

\item[]  \textbf{Forti, M./Honsell, F.}

\bibitem[\textbf{1985}]{bib:1985k} \emph{The consistency of the axiom of universality for the ordering of cardinals}, J. Symbolic Logic \textbf{50}, 502-509.

\item[]  \textbf{Fossy, J./Morillon, M.}

\bibitem[\textbf{1995}]{bib:1995c} \emph{The Baire category property and some notions of compactness}, preprint.

\item[]  \textbf{Fourman, M. P.}

\bibitem[\textbf{1980}]{bib:1980f} \emph{Sheaf models for set theory}, J. Pure Appl. Algebra \textbf{19}, 91-101.

\item[]  \textbf{Fra\"{i}ss\'{e}, R. J.}

\bibitem[\textbf{1958}]{bib:1958} \emph{Un mod\`{e}le d\'{e}finissant une th\'{e}orie aberrant des ensembles ou sont nies les axiomes du choix et d'extensionalite}, Publ. Sci. Univ. Alger Ser. A Math., \textbf{5}, 17-98.

\item[]  \textbf{Fraenkel, A.}

\bibitem[\textbf{1922}]{bib:1922} \emph{Der Bergriff ``definit'' und der Unabhangigkeit des Auswahlaxioms}, Sitzungsber. Preussichen Akad. Wiss., Math. Kl., 253-257.

\item[]  \textbf{Fraleigh, J. B.}

\bibitem[\textbf{1989}]{bib:1989p} A First Course in Abstract Algebra fourth edition, Addison Welley, New York.

\item[]  \textbf{Freyd, P.}

\bibitem[\textbf{1980}]{bib:1980g} \emph{The axiom of choice}, J. Pure Appl. Algebra \textbf{19}, 103-125.

\item[]  \textbf{Fuchs, L.}

\bibitem[\textbf{1970}]{bib:1970e} \emph{Infinite Abelian groups I}, Academic Press, New York.

\item[]  \textbf{Gaifman, H.}

\bibitem[\textbf{1970}]{bib:1970f} \emph{Infinite Boolean polynomials I}, Academic Press, New York.

\item[]  \textbf{Galvin, F./ Komj\'{a}th, P.}

\bibitem[\textbf{1991}]{bib:1991d} \emph{Graph colorings and the axiom of choice}, Period. Math. Hungar. \textbf{22}, 592--607.

\item[]  \textbf{Gonzalez, C.}

\bibitem[\textbf{1995a}]{bib:1995a} \emph{Dense orderings, partitions, and weak forms of choice}, Fund. Math. \textbf{147}, 11-25.

\bibitem[\textbf{1995b}]{bib:1995b} \emph{Partitioning an infinite set}, preprint.

\item[]  \textbf{Gardiner, G.}

\bibitem[\textbf{1974}]{bib:1974i} \emph{The equivalence of the Boolean prime ideal theorem and a theorem of functional analysis}, Fund. Math. \textbf{84}, 81-86.

\item[]  \textbf{Gauntt, R.J.}

\bibitem[\textbf{1968}]{bib:1968d} \emph{Some restricted versions of the axiom of choice}, Notices Amer. Math. Soc. \textbf{68T-176}, 351.

\bibitem[\textbf{1970}]{bib:1970g} \emph{Axiom of choice for finite sets - a solution of a problem of Mostowski}, Notices Amer. Math. Soc. \textbf{17}, 454.

\item[]  \textbf{Gelfand, I./Kolmogoroff, A.}

\bibitem[\textbf{1939}]{bib:1939c} \emph{On rings of continuous functions on topological spaces}, Dokl. Akad. Nauk SSSR \textbf{22}, 25-38.

\item[]  \textbf{Gitik, M.}

\bibitem[\textbf{1980}]{bib:1980h} \emph{All uncountable cardinals can be singular}, Israel J. Math. \textbf{35}, 61-88.

\bibitem[\textbf{1985}]{bib:1985l} \emph{Regular cardinals in models of ZF}, Trans. Amer. Math. Soc. \textbf{290}, 41-68.

\item[]  \textbf{Gleason, A. M.}

\bibitem[\textbf{1958}]{bib:1958c} \emph{Projective topological spaces}, Illinois J. Math. \textbf{2}, 482--489.

\item[]  \textbf{Good, C./Tree, I. J.}

\bibitem[\textbf{1995}]{bib:1995d} \emph{Continuing horrors of topology without choice}, Topology and its Applications \textbf{63}, 79-90.

\item[]  \textbf{Good, C./Tree, I. J./Watson, W. S.}

\bibitem[\textbf{1996}]{bib:1996e} \emph{On Stone's theorem and the axiom of choice}, preprint.

\item[]  \textbf{Goodner, D. A.}

\bibitem[\textbf{1950}]{bib:1950c} \emph{Projections in normed linear spaces}, Trans. Amer. Math. Soc. \textbf{69}, 89-108.

\item[]  \textbf{Goldblatt, R.}

\bibitem[\textbf{1984}]{bib:1984l} \emph{An abstract setting for Henkin proofs}, Topoi, \textbf{3}, 37-41.

\bibitem[\textbf{1985}]{bib:1985m} \emph{On the role of the Baire category theorem and dependent choice in the foundations of logic}, J. Symbolic Logic \textbf{50}, 412-422.

\item[]  \textbf{Gottschalk, W. H.}

\bibitem[\textbf{1951}]{bib:1951} \emph{Choice functions and Tychonoff's theorem}, Proc. Amer. Math. Soc. \textbf{2}, 172. MR 12.683, Zbl 43.54.

\item[]  \textbf{Gr\"{a}tzer, G.}

\bibitem[\textbf{1968}]{bib:1968e} \emph{Universal Algebra}, Van Nostrand, Princeton.

\bibitem[\textbf{1979}]{bib:1979c} \emph{Universal Algebra}, Springer Verlag, New York.

\bibitem[\textbf{1986}]{bib:1986f} \emph{Birkhoff's representation theorem is equivalent to the axiom of choice}, Algebra Universalis \textbf{23}, 58-60.

\item[]  \textbf{Gross, W. F.}

\bibitem[\textbf{1975}]{bib:1975e} \emph{Models with dimension}, Ph. D. Thesis, Monash Univerity.

\bibitem[\textbf{1976}]{bib:1976f} \emph{Models with dimension}, Bull. Austral. Math. Soc. \textbf{14}, 153-154.

\item[]  \textbf{Guzicki, W.}

\bibitem[\textbf{1976}]{bib:1976g} \emph{On weaker forms of choice in second order arithmetic}, Fund. Math. \textbf{93}, 131-144.

\item[]  \textbf{Haddad, L./Morillon, M.}

\bibitem[\textbf{1990}]{bib:1990d} \emph{L'axiome de normalit\'{e} pour les espaces totalement ordonn\'{e}s}, J. Symbolic Logic \textbf{55}, 277-283.

\item[]  \textbf{Hales, A. W.}

\bibitem[\textbf{1964}]{bib:1964c} \emph{On the non-existence of free complete Boolean algebras}, Fund. Math. \textbf{54}, 45-66.

\item[]  \textbf{Hall, M.}

\bibitem[\textbf{1948}]{bib:1948} \emph{Distinct representatives of subsets}, Bull. Amer. Math. Soc. \textbf{54}, 922-926.

\item[]  \textbf{Hall, P.}

\bibitem[\textbf{1935}]{bib:1935} \emph{On representatives of subsets}, J. London Math. Soc. \textbf{10}, 26-30.

\item[]  \textbf{Halpern, J. D.}

\bibitem[\textbf{1962}]{bib:1962c} \emph{Some structure in the hierarchy of maximal principles in set theory}, Notices Amer. Math. Soc. \textbf{9}, 315.

\bibitem[\textbf{1964}]{bib:1964e} \emph{The independence of the axiom of choice from the Boolean prime ideal theorem}, Fund. Math. \textbf{55}, 57-66.

\bibitem[\textbf{1966}]{bib:1966c} \emph{Bases for vector spaces and the axiom of choice}, Proc. Amer. Math. Soc. \textbf{17}, 670-673.

\bibitem[\textbf{1972}]{bib:1972f} \emph{On a question of Tarski and a maximal theorem of Kurepa}, Pacific J. Math. \textbf{41}, 111-121.

\bibitem[\textbf{1975}]{bib:1975f} \emph{Non-standard combinatorics}, Proc. London Math. Soc. Ser. 3 \textbf{30}, 40-54.

\item[]  \textbf{Halpern, J.D./Howard, P.}

\bibitem[\textbf{1970}]{bib:1970g1} \emph{Cardinals m such that $2m=m$}, Proc. Amer. Math. Soc. \textbf{26}, 487-490.

\bibitem[\textbf{1974}]{bib:1974j} \emph{Cardinal addition and the axiom of choice}, Bull. Amer. Math. Soc. \textbf{80}, 584-586.

\item[]  \textbf{Halpern, J.D./L\"{a}uchli, H.}

\bibitem[\textbf{1966}]{bib:1966d} \emph{A partition theorem}, Trans. Amer. Math. Soc. \textbf{124}, 360-367.

\item[]  \textbf{Halpern, J.D./Levy, A.}

\bibitem[\textbf{1964}]{bib:1964f} \emph{The ordering theorem does not imply the axiom of choice}, Notices Amer. Math. Soc. \textbf{11}, 56.

\bibitem[\textbf{1971}]{bib:1971e} \emph{The Boolean prime ideal theorem does not imply the axiom of choice}, Axiomatic Set Theory, Proceedings of Symposia in Pure Mathematics, 13 part 1, D. Scott ed., AMS, Providence, 83-134.

\item[]  \textbf{Hamel, G.}

\bibitem[\textbf{1905}]{bib:1905} \emph{Eine Basis aller Zahlen und die unstetigen Losungen der Functionalgleichung}: $f(x+y)=f(x)+f(y)$, Math. Ann. \textbf{60}, 459-462.

\item[]  \textbf{Harazisville, A. B. - see Kharazishvili, A. B.}

\item[]  \textbf{Harper, J./Rubin, J.}

\bibitem[\textbf{1976}]{bib:1976h} \emph{Variations of Zorn's lemma, principles of cofinality, and Hausdorff's maximal principle, Part I and II}, Notre Dame J. Formal Logic \textbf{27}, 565-588 and Notre Dame J. Formal Logic \textbf{28} (1977), 151-163.

\item[]  \textbf{Hausdorff, F.}

\bibitem[\textbf{1907}]{bib:1907} \emph{\"{U}ber dichte Ordnungstypen}, Jber. Deutsch. Math.-Verein. \textbf{16}, 541-546.

\item[]  \textbf{H\"{a}ussler, A.}

\bibitem[\textbf{1983}]{bib:1983i} \emph{Defining cardinal addition by $\leq$--formulas}, Fund. Math. \textbf{115}, 195-205. MR 85e03111, Zbl 517\#03017.

\item[]  \textbf{van Heijenoort, J.}

\bibitem[\textbf{1879}]{bib:1879} \emph{From Frege to G\"{o}del, A Source Book in Mathematical Logic}-1931, Harvard University Press, Cambridge, 1967.

\item[]  \textbf{Henkin, L.}

\bibitem[\textbf{1954a}]{bib:1954a} \emph{Metamathematical theorems equivalent to the prime ideal theorem for Boolean algebras}, Bull. Amer. Math. Soc. \textbf{60}, 378-388.

\bibitem[\textbf{1954b}]{bib:1954b} \emph{Boolean representations through propositional calculus}, Fund. Math. \textbf{41}, 89-96.

\item[]  \textbf{Hickin, K. K./Plotkin, J. M.}

\bibitem[\textbf{1976}]{bib:1976i} \emph{A patching lemma}, Notre Dame J. Formal Logic \textbf{17}, 158-160.

\item[]  \textbf{Hickman, J.L.}

\bibitem[\textbf{1972}]{bib:1972g} \emph{On a cardinal equation in set theory}, Bull. Austral. Math. Soc. \textbf{6}, 447--457.

\bibitem[\textbf{1975}]{bib:1975g} \emph{General well-ordered sets}, J. Austral. Math. Soc. \textbf{19} ser. A, 7-20.

\bibitem[\textbf{1976}]{bib:1976j} \emph{Groups in models of set theory that fail the axiom of choice}, Bull. Austral. Math. Soc. \textbf{14}, 199-232.

\bibitem[\textbf{1977a}]{bib:1977a2} \emph{Quasi-minimal posets and lattices}, J. Reine Angew. Math. \textbf{296}, 10-13.

\bibitem[\textbf{1977b}]{bib:1977b2} \emph{A remark on elementary Abelian groups}, Bull. Austral. Math. Soc. \textbf{16}, 213-217.

\bibitem[\textbf{1978a}]{bib:1978a} \emph{Dedekind finite fields}, Bull. Austral. Math. Soc. \textbf{19}, 117-124.

\bibitem[\textbf{1978b}]{bib:1978b} \emph{The ideal of orderable subsets of a set}, Notre Dame J. Formal Logic \textbf{19}, 593-598.

\bibitem[\textbf{1979a}]{bib:1979a} \emph{A class of ``near-finite'' order types}, Z. Math. Logik Grundlagen Math. \textbf{25}, 79-92.

\bibitem[\textbf{1979b}]{bib:1979b} \emph{Boundedness properties of cardinals}, Z. Math. Logik Grundlagen Math. \textbf{25}, 485-486.

\bibitem[\textbf{1979c}]{bib:1979c1} \emph{On a remark by Sageev}, Z. Math. Logik Grundlagen Math. \textbf{25}, 373-374.

\bibitem[\textbf{1980a}]{bib:1980a} \emph{Some observations concerning medial fields}, Z. Math. Logik Grundlagen Math. \textbf{26}, 471-475.

\bibitem[\textbf{1980b}]{bib:1980b} $\Lambda$ -\emph{minimal lattices}, Z. Math. Logik Grundlagen Math. \textbf{26}, 181-191.

\bibitem[\textbf{1982}]{bib:1982f} \emph{Automorphisms of medial fields}, Z. Math. Logik Grundlagen Math. \textbf{28}, 263-267.

\item[]  \textbf{Higasikawa, M.}

\bibitem[\textbf{1995}]{bib:1995e} \emph{Partition principles and infinite sums of cardinal numbers}, Notre Dame J. of Formal Logic \textbf{36}, 425-434.

\item[]  \textbf{Hilbert, D.}

\bibitem[\textbf{1900}]{bib:1900} \emph{Mathematische Probleme. Vortrag, gehalten auf dem internationalem Mathematiker-Kongress zu Paris}, Nachrichten, Akademie der Wissenschafter G\"{o}ttingen, 1900, 253-397, translated in Bulletin of the A.M.S. (2) \textbf{8} (1902), 437-479.

\item[]  \textbf{Hodges, W.}

\bibitem[\textbf{1974}]{bib:1974k} \emph{Six impossible rings}, J. of Algebra \textbf{31}, 218-244.

\bibitem[\textbf{1976a}]{bib:1976a} \emph{L\"{a}uchli's algebraic closure of} $\mathbb{Q}$, Math. Proc. Cambridge Philos. Soc. \textbf{79}, 289-297.

\bibitem[\textbf{1976b}]{bib:1976b} \emph{On the effectivity of some field constructions}, Proc. London Math. Soc. \textbf{32}, 133-162.

\bibitem[\textbf{1979}]{bib:1979d} \emph{Krull implies Zorn}, J. London Math. Soc. Ser. 2 \textbf{19}, 285-287.

\bibitem[\textbf{1980}]{bib:1980j} \emph{Constructing pure infective hulls}, J. Symbolic Logic \textbf{45}, 544-548.

\item[]  \textbf{H\"{o}ft, H./Howard, P.}

\bibitem[\textbf{1981}]{bib:1981g} \emph{Representing multi-algebras by algebras, the axiom of choice and the axiom of dependent choice}, Algebra Universalis \textbf{13}, 69-77.

\bibitem[\textbf{1994}]{bib:1994c} \emph{Well ordered subsets of linearly ordered sets}, Notre Dame J. Formal Logic \textbf{35}, 413--425.

\item[]  \textbf{Howard, P.}

\bibitem[\textbf{1973}]{bib:1973h} \emph{Limitations on the Fraenkel-Mostowski method of independence proofs}, J. Symbolic Logic \textbf{38}, 416--422.

\bibitem[\textbf{1975}]{bib:1975h} \emph{Los theorem and the Boolean prime ideal theorem imply the axiom of choice}, Proc. Amer. Math. Soc. \textbf{49}, 426-428.

\bibitem[\textbf{1984a}]{bib:1984a2} \emph{Binary consistent choice on pairs and a generalization of K\"{o}nig's infinity lemma}, Fund. Math. \textbf{121}, 17-31.

\bibitem[\textbf{1984b}]{bib:1984b2} \emph{Rado's selection lemma does not imply the Boolean prime ideal theorem}, Z. Math. Logik Grundlagen Math. \textbf{30}, 129-132.

\bibitem[\textbf{1985}]{bib:1985n} \emph{Subgroups of a free group and the axiom of choice}, J. Symbolic Logic \textbf{50}, 458-467.

\bibitem[\textbf{1987}]{bib:1987} \emph{The existence of level sets in a free group implies the axiom of choice}, Z. Math. Logik Grundlagen Math. \textbf{33}, 315-316.

\bibitem[\textbf{1992}]{bib:1992c} \emph{The axiom of choice for countable sets of countable sets does not imply the countable union theorem}, Notre Dame J. Formal Logic \textbf{13}, 236-243.

\bibitem[\textbf{1993}]{bib:1993c} \emph{Variations of Rado's lemma}, Math. Logic Quart. \textbf{39}, 353-356.

\bibitem[\textbf{1994}]{bib:1994d} \emph{Unions of well-ordered sets}, J. Austral. Math. Soc. Ser. A \textbf{56}, 117-124.

\item[]  \textbf{Howard, P./Keremedis, K.,/Rubin, H./Rubin, J.}

\bibitem[\textbf{1997a}]{bib:1997a} \emph{Versions of normality and some weak forms of the axiom of choice}, accepted June Math. Logic Quart.

\bibitem[\textbf{1997b}]{bib:1997b} \emph{Disjoint unions of topological spaces and choice}, Accepted September Math. Logic Quart.

\item[]  \textbf{Howard, P./Keremedis, K.,/Rubin, J./Stanley, A.}

\bibitem[\textbf{1997}]{bib:1997d} \emph{Paracompactness of metric spaces and the axiom of choice}, preprint September.

\item[]  \textbf{Howard, P./Rubin, A./Rubin, J.}

\bibitem[\textbf{1978}]{bib:1978e} \emph{Independence results for class forms of the axiom of choice}, J. Symbolic Logic \textbf{43}, 673-684.

\bibitem[\textbf{1997}]{bib:1997e} \emph{Kinna-Wagner selection principles, axioms of choice and multiple choice}, Monashafte f\"{u}r Math. \textbf{123}, 309-319.

\item[]  \textbf{Howard, P./Rubin, H./Rubin, J.}

\bibitem[\textbf{1973}]{bib:1973i} \emph{The relationship between two weak forms of the axiom of choice}, Fund. Math. \textbf{80}, 75-79.

\item[]  \textbf{Howard, P./Rubin, J.}

\bibitem[\textbf{1977}]{bib:1977e1} \emph{The axiom of choice and linearly ordered sets}, Fund. Math. \textbf{98}, 111-122.

\bibitem[\textbf{1995a}]{bib:1995a1} \emph{Elementary Abelian} $p$-\emph{groups revisited}, Bull. Austral. Math. Soc. \textbf{52}, 373-376.

\bibitem[\textbf{1995b}]{bib:1995b1} \emph{The axiom of choice for well ordered families and for families of well-orderable sets}, J. Symbolic Logic, \textbf{60}, 1115-1117.

\bibitem[\textbf{1996}]{bib:1996f} \emph{The Boolean prime ideal theorem plus countable choice do not imply dependent choice}, Math. Logic Quart., \textbf{42}, 410-420.

\item[]  \textbf{Howard, P./Solski, J.}

\bibitem[\textbf{1993}]{bib:1993d} \emph{The strength of the $\Delta$-system lemma}, Notre Dame J. Formal Logic \textbf{34}, 100--106.

\item[]  \textbf{Howard, P./Spi\u{s}iak, L.}

\bibitem[\textbf{1994}]{bib:1994e} \emph{Definitions of finite and the power set operation}, preprint.

\item[]  \textbf{Howard, P./Yorke, M.}

\bibitem[\textbf{1987}]{bib:1987d} \emph{Maximal $p$-subgroups and the axiom of choice}, Notre Dame J. Formal Logic \textbf{28}, 276--283.

\bibitem[\textbf{1989}]{bib:1989q} \emph{Definitions of finite}, Fund. Math. \textbf{133}, 169-177.

\item[]  \textbf{Jech, T.}

\bibitem[\textbf{1966a}]{bib:1966a} \emph{On cardinals and their successors}, Bull. Acad. Polon. Sci. S\'{e}r. Sci. Math. Astronom. Phys. \textbf{14}, 533-537.

\bibitem[\textbf{1966b}]{bib:1966b} \emph{Interdependence of weakened forms of the axiom of choice}, Comment. Math. Univ. Carolinae \textbf{7}, 359-371.

\bibitem[\textbf{1966c}]{bib:1966c1} \emph{On the ordering of cardinalities}, Bull. Acad. Polon. Sci. S\'{e}r. Sci. Math. Astronom. Phys. \textbf{14}, 293-296.

\bibitem[\textbf{1967}]{bib:1967c} \emph{On models of set theory without AC}, Axiomatic Set Theory: Lo\'{s} Angeles, D. Scott and T. Jech ed., Proc. Symp. Pure Math \textbf{13}, AMS, Providence, 1971, 135-141.

\bibitem[\textbf{1968a}]{bib:1968c1} $\omega_{1}$ \emph{can be measurable}, Israel J. Math. \textbf{6}, 363-367.

\bibitem[\textbf{1968b}]{bib:1968d1} \emph{Bemerkung zum Auswahlaxiom}, \v{C}asopis P\v{e}st. Mat. \textbf{93}, 30-31.

\bibitem[\textbf{1973a}]{bib:1973a} \emph{Some combinatorial problems concerning uncountable cardinals}, Ann. Math. Logic \textbf{5}, 165-198.

\bibitem[\textbf{1973b}]{bib:1973b} \emph{The Axiom of Choice}, North-Holland, Amsterdam.

\bibitem[\textbf{1977}]{bib:1977f} \emph{About the axiom of choice}, Handbook of Mathematical Logic (Studies in Logic and the Foundations of Mathematics, \textbf{90}, North-Holland, Amsterdam, 345-370.

\bibitem[\textbf{1982}]{bib:1982g} \emph{On hereditarily countable sets}, J. Symbolic Logic \textbf{47}, 43-47.

\item[]  \textbf{Jech, T./Sochor, A.}

\bibitem[\textbf{1966a}]{bib:1966a1} \emph{Applications of the $\Theta$-model}, Bull. Acad. Polon. Sci. S\'{e}r. Sci. Math. Astronom. Phys. \textbf{14}, 351-355.

\bibitem[\textbf{1966b}]{bib:1966b1} \emph{On $\Theta$-model of the set theory}, Bull. Acad. Polon. Sci. S\'{e}r. Sci. Math. Astronom. Phys. \textbf{14}, 297-303.

\item[]  \textbf{Jensen, R.B.}

\bibitem[\textbf{1966}]{bib:1966e} \emph{Independence of the axiom of dependent choices from the countable axiom of choice}, J. Symbolic Logic \textbf{31}, 294.

\bibitem[\textbf{1967}]{bib:1967d} \emph{Consistency results for ZF}, Notices Amer. Math. Soc. \textbf{14}, 137.

\bibitem[\textbf{1968}]{bib:1968f} \emph{Definable sets of minimal degree}, Mathematical Logic and Foundations of Set Theory, Jerusalem (Studies in Logic and the Foundations of Mathematics, North-Holland, Amsterdam, 1970), 122--128.

\item[]  \textbf{Johnstone, P.T.}

\bibitem[\textbf{1977}]{bib:1977g} \emph{Topos Theory}, L.M.S. Monographs, Academic Press.

\bibitem[\textbf{1984}]{bib:1984m} \emph{Almost maximal ideals}, Fund. Math. \textbf{123}, 197-209.

\item[]  \textbf{Kanamori, A./Magidor, M.}

\bibitem[\textbf{1977}]{bib:1977h} \emph{The evolution of large cardinal axioms in set theory}, Higher set theory (Oberwolfach, Lecture Notes in Math., vol. 669, Springer-Verlag, Berlin, 1978, 99-275.

\item[]  \textbf{Kaniewski, J./Rogers, C.A.}

\bibitem[\textbf{1980}]{bib:1980k} \emph{Double uniformization}, J. London Math. Soc. \textbf{22}, 521-533.

\item[]  \textbf{Kanovei, V.G.}

\bibitem[\textbf{1976}]{bib:1976k} \emph{Definability with the help of degrees of constructibility}, Studies in set theory and non-classical logics (Russian), Izdat. Nauka, Moskow, 5-95.

\bibitem[\textbf{1978}]{bib:1978f} \emph{The non-emptyness of classes in axiomatic set theory}, Izv. Akad. Nauk SSSR Ser. Mat. \textbf{42}, 550-579. (Translated in Math. USSR-Izv. \textbf{42} (1978), 550--579.)

\bibitem[\textbf{1979}]{bib:1979e} \emph{Descriptive forms of the countable axiom of choice}, Studies in nonclassical logics and set theory, (Russian) Moscow, Nauka, 3-136.

\bibitem[\textbf{1991}]{bib:1991e} \emph{Cardinality of the set of Vitali equivalence classes}, Mat. Zametki \textbf{49}, 63-68. (Translated in Math. Notes \textbf{49} (1991), 370--374.)

\item[]  \textbf{Kasahara, S.}

\bibitem[\textbf{1976}]{bib:1976l} \emph{Remarks on some fixed point theorems}, Math. Sem. Notes Kobe Univ. \textbf{4}, 43-50.

\item[]  \textbf{Keisler, H. J.}

\bibitem[\textbf{1970}]{bib:1970i} \emph{Logic with the quantifier ``there exist uncountably many}'', Ann. Math. Logic \textbf{1} 1-93.

\bibitem[\textbf{1971}]{bib:1971f} \emph{Model Theory for Infinitary Logic}, North Holland.

\item[]  \textbf{Keisler, H.J./Tarski, A.}

\bibitem[\textbf{1964}]{bib:1964h} \emph{From accessible to inaccessible cardinals}, Fund. Math. \textbf{53}, 225-308.

\item[]  \textbf{Kelley, J.L.}

\bibitem[\textbf{1950}]{bib:1950d} \emph{Tychonoff's theorem implies AC}, Fund. Math. \textbf{37} 75-76.

\bibitem[\textbf{1955}]{bib:1955} \emph{General Topology}, Van Nostrand, New York.

\item[]  \textbf{Kenyon, H.}

\bibitem[\textbf{1963}]{bib:1963e} Problem 5077 Amer. Math. Monthly \textbf{70}, 216.

\item[]  \textbf{Keremedis, K.}

\bibitem[\textbf{1995}]{bib:1995f} \emph{Continuous real valued functions in} $T_{4}$ \emph{spaces}, Accepted Math. Japon.

\bibitem[\textbf{1996a}]{bib:1996a} \emph{Bases for vector spaces over the two element field and the axiom of choice}, Proc. Amer. Math. Soc.\textbf{124} 2527-2531.

\bibitem[\textbf{1996b}]{bib:1996b} \emph{Filters, antichains and towers in topological spaces and the axiom of choice}, preprint.

\bibitem[\textbf{1996c}]{bib:1996c1} \emph{Some equivalents of AC in algebra}, Accepted Algebra Universalis.

\bibitem[\textbf{1996d}]{bib:1996d1} \emph{Disasters in topology without the axiom of choice}, preprint.

\item[]  \textbf{Kharazishvili, A.B.}

\bibitem[\textbf{1979}]{bib:1979f} \emph{The axiom of choice in the general theory of systems}, Soob\v{s}\v{c}. Akak. Nauk. Gruzin SSR \textbf{94}, 533-536.

\item[]  \textbf{Kim, S.M.}

\bibitem[\textbf{1988}]{bib:1988d} \emph{Variants of the well-ordering theorem and the axiom of multiple choice}, Bol. Soc. Paran. Mat. (2) \textbf{9}, 177-181.

\item[]  \textbf{Kinna, W./Wagner, K.}

\bibitem[\textbf{1955}]{bib:1955c} \emph{\"{U}ber eine Abschw\"{a}chung des Auswahlpostulates}, Fund. Math. \textbf{42}, 75-82.

\item[]  \textbf{Kleinberg, E.}

\bibitem[\textbf{1969}]{bib:1969c} \emph{The independence of Ramsey's theorem}, J. Symbolic Logic \textbf{34}, 205-206.

\item[]  \textbf{Kleinberg, E. M./Seiferas, J. I.}

\bibitem[\textbf{1973}]{bib:1973j} \emph{Infinite exponent partition relations and well-ordered choice}, J. Symbolic Logic \textbf{38}, 299--308.

\item[]  \textbf{Knight, C. J.}

\bibitem[\textbf{1964}]{bib:1964g} \emph{Boxtopologies}, Quart. J. Math. Oxford Ser. \textbf{15}, 41-54.

\item[]  \textbf{Kolany, A.}

\bibitem[\textbf{1992}]{bib:1992d} \emph{Equivalents of the compactness theorem for locally finite sets of sentences}, Bull. Sec. Logic Polish Acad. Sci. \textbf{21}, 12--18.

\item[]  \textbf{Kolany, A./Wojtylak, P.}

\bibitem[\textbf{1991}]{bib:1991f} \emph{Restricted versions of the compactness theorem}, Rep. Math. Logic \textbf{25}, 91-103.

\item[]  \textbf{Krom, M.}

\bibitem[\textbf{1981}]{bib:1981h} \emph{Equivalents of a weak axiom of choice}, Notre Dame J. Formal Logic \textbf{22}, 283-285.

\bibitem[\textbf{1986}]{bib:1986j} \emph{A linearly ordered topological space that is not normal}, Notre Dame J. Formal Logic \textbf{27}, 12-13.

\item[]  \textbf{Krull, W.}

\bibitem[\textbf{1929}]{bib:1929} \emph{Die Idealtheorie in Ringen ohne Endlichkeitsbedingungen}, Math. Ann. \textbf{101}, 729-744.

\item[]  \textbf{Kruse, A. H.}

\bibitem[\textbf{1974}]{bib:1974l} \emph{Some results on partitions and Cartesian products in the absence of the axiom of choice}, Z. Math. Logik Grundlagen Math. \textbf{20}, 149-172.

\item[]  \textbf{Kunen, K.}

\bibitem[\textbf{1980}]{bib:1980m} \emph{Set Theory. An Introduction to Independence Proofs}, Studies in Logic and the Foundations of Mathematics \textbf{102}, North Holland, Amsterdam.

\item[]  \textbf{Kuratowski, K.}

\bibitem[\textbf{1966}]{bib:1966f} \emph{Topology I}, Academic Press, New York.

\item[]  \textbf{Kurepa, D.}

\bibitem[\textbf{1953}]{bib:1953} \emph{\"{U}ber das Auswahlaxiom}, Math. Ann. \textbf{126}, 381-384.

\bibitem[\textbf{1971}]{bib:1971g} \emph{Factorials and the general continuum hypothesis}, General Topology and its Relationship to Modern Analysis and Algebra, III (Proc. Third Prague Topological Sympos., Academia, Prague, 1972, 281--282.

\item[]  \textbf{Lampe, W. A.}

\bibitem[\textbf{1974}]{bib:1974m} \emph{Subalgebra lattices of unary algebras and an axiom of choice}, Colloq. Math. \textbf{30}, 41-55.

\item[]  \textbf{Lang, S.}

\bibitem[\textbf{1993}]{bib:1993e} \emph{Algebra}, third edition, Addison-Wesley, Reading Mass..

\item[]  \textbf{L\"{a}uchli, H.}

\bibitem[\textbf{1962}]{bib:1962d} \emph{Auswahlaxiom in der Algebra}, Comment. Math. Helv. \textbf{37}, 1-18.

\bibitem[\textbf{1964}]{bib:1964i} \emph{The independence of the ordering principle from a restricted axiom of choice}, Fund. Math. \textbf{54}, 31-43.

\bibitem[\textbf{1971}]{bib:1971h} \emph{Coloring infinite graphs and the Boolean prime ideal theorem}, Israel J. Math. \textbf{9}, 422-429.

\item[]  \textbf{Lebesgue, H.}

\bibitem[\textbf{1905}]{bib:1905c} \emph{Sur les fonctions representables analytiquement}, J. Math. Pures Appl. \textbf{60}, 139-216.

\item[]  \textbf{Lembcke, J.}

\bibitem[\textbf{1979}]{bib:1979g} \emph{Two extension theorems effectively equivalent to the axiom of choice}, Bull. London Math. Soc. \textbf{11}, 285-288.

\item[]  \textbf{Levy, A.}

\bibitem[\textbf{1958}]{bib:1958d} \emph{The independence of various definitions of finiteness}, Fund. Math. \textbf{46}, 1-13.

\bibitem[\textbf{1960}]{bib:1960} \emph{On models of set theory with urelements}, Bull. Acad. Polon. Sci. S\'{e}r. Sci. Math. Astronom. Phys. \textbf{8}, 463-465.

\bibitem[\textbf{1962}]{bib:1962e} \emph{Axioms of multiple choice}, Fund. Math. \textbf{50}, 475-483.

\bibitem[\textbf{1963a}]{bib:1963a} \emph{The Fraenkel-Mostowski method for independence proofs in set theory}, Theory of Models:, Berkeley, Studies in Logic and the Foundations of Mathematics, North Holland, Amsterdam, 1965, 221-228.

\bibitem[\textbf{1963b}]{bib:1963b} \emph{Remarks on a paper by J. Mycielski}, Acta Math. Acad. Sci. Hungar. \textbf{14}, 125-130.

\bibitem[\textbf{1964}]{bib:1964k} \emph{The interdependence of certain consequences of the axiom of choice}, Fund. Math. \textbf{54}, 137-157.

\bibitem[\textbf{1971}]{bib:1971i} \emph{Parameters in the comprehension axiom schemas of set theory}, Tarski Symposium: Berkeley, L. Henkin, J. Addison, C.C. Chang, W. Craig, D.S. Scott and R. Vaught ed., Proc Symp Pure Math \textbf{25}, AMS, Providence, 1974, 309-324.

\item[]  \textbf{Lewis, A.A.}

\bibitem[\textbf{1983}]{bib:1983j} \emph{A note on the relationship between the theorems of Tychonoff, Alexander and Rado}, Bull. Malaysian Math. Soc. \textbf{6}, 69-72.

\bibitem[\textbf{1986a}]{bib:1986a1} \emph{A $\Pi_{1}^{0}$-class characterization of $\omega$-selective choice functions}, preprint.

\bibitem[\textbf{1986b}]{bib:1986b1} \emph{On the effectiveness of $\omega$-selective families}, preprint.

\item[]  \textbf{Lindel\"{o}f, E.}

\bibitem[\textbf{1905}]{bib:1905d} \emph{Remarque sur un th\'{e}or\`{e}me fondamental de la th\'{e}orie des ensembles}, Acta Math. \textbf{29}, 183-190.

\item[]  \textbf{Lindenbaum, A./Tarski, A.}

\bibitem[\textbf{1926}]{bib:1926} \emph{Communication sur les recherches de la th\'{e}orie des ensembles}, C. R. Soc. Sci. Lett. Varsovie \textbf{19}, 299-330.

\item[]  \textbf{Litman, A.}

\bibitem[\textbf{1976}]{bib:1976m} \emph{The monadic theory of} $\omega_{1}$, Israel J. Math. \textbf{23}, 251-266.

\item[]  \textbf{Lolli, G.}

\bibitem[\textbf{1977}]{bib:1977i} \emph{On Ramsey's theorem and the axiom of choice}, Notre Dame J. Formal Logic \textbf{18}, 599-601.

\item[]  \textbf{{\L}o\'{s}, J./Ryll-Nardzewski, C.}

\bibitem[\textbf{1951}]{bib:1951c} \emph{On the application of Tychonoff's theorem in mathematical proofs}, Fund. Math. \textbf{38}, 233-237.

\bibitem[\textbf{1954}]{bib:1954} \emph{Effectiveness of the representation theory for Boolean algebras}, Fund. Math. \textbf{41}, 49-56.

\item[]  \textbf{L\"{o}wig, L.}

\bibitem[\textbf{1934}]{bib:1934} \emph{\"{U}ber die Dimension linearer R\"{a}ume}, Studia Math. \textbf{5} 18-23.

\item[]  \textbf{Luxemburg, W.A.J.}

\bibitem[\textbf{1969}]{bib:1969d} \emph{Reduced powers of the real number system and equivalents of the Hahn-Banach theorem}, Applications of Model Theory to Algebra, Analysis and Probability, W.A. Luxemburg ed., Holt, Rinehart and Winston, New York, 123-137.

\bibitem[\textbf{1970}]{bib:1970j} \emph{Two applications of the method of construction by ultrapowers to analysis}, Amer. Math. Soc. Proc. Symp. Pure Math. \textbf{18} part 1, 83-134.

\item[]  \textbf{Luzin, N./Sierpi\'{n}ski, W.}

\bibitem[\textbf{1917}]{bib:1917} \emph{Sur une propri\'{e}t\'{e} du continu}, C. R. Acad. Sci. Paris S\'{e}r. A--B \textbf{165}, 498-500.

\item[]  \textbf{MacLane, S.}

\bibitem[\textbf{1971}]{bib:1971j} \emph{Categories for the Working Mathematician}, Graduate Texts in Math. no. 5, Springer-Verlag

\item[]  \textbf{Makowski, A./Wi\'{s}niewski, K.}

\bibitem[\textbf{1969}]{bib:1969e} \emph{Generalization of Abian's fixed point theorem}, Comment. Math. Ann. Soc. Math. Polon. S\'{e}r. 1 \textbf{13}, 63--65.

\bibitem[\textbf{1972}]{bib:1972h} \emph{On the axiom of choice for families of finite sets}, Fund. Math.\textbf{73},187--192.

\item[]  \textbf{Malitz, J.}

\bibitem[\textbf{1992}]{bib:1992e} \emph{The axiom of choice, the Lowenheim-Skolem theorem and Borel models}, Fund. Math. \textbf{140}, 197.

\item[]  \textbf{Manka, Roman}

\bibitem[\textbf{1988a}]{bib:1988e} \emph{Some forms of the axiom of choice}, Jahrbuch der Kurt-G\"{o}del-Gesellschaft, 24-34.

\bibitem[\textbf{1988b}]{bib:1988f} \emph{Turinici's fixed point theorem and the axiom of choice}, Rep. Math. Logic, no. 12, 15--19.

\item[]  \textbf{Mathias, A.R.D.}

\bibitem[\textbf{1967}]{bib:1967e} \emph{The order extension principle}, Axiomatic Set Theory, Los Angeles, T. Jech ed., Proc. Symposia in Pure Math \textbf{13}, AMS, Providence, 1974, 179-183.

\bibitem[\textbf{1968}]{bib:1968g} \emph{On a generalization of Ramsey's theorem}, Thesis, Bonn.

\bibitem[\textbf{1977a}]{bib:1977a3} \emph{Happy families}, Ann. Math. Logic \textbf{12}, 59-111.

\bibitem[\textbf{1977b}]{bib:1977b3} \emph{The real line and the universe}, Logic Colloquium 76 (Studies in Logic and the Foundations of Math. \textbf{87}, North-Holland, Amsterdam, 531-546.

\bibitem[\textbf{1979}]{bib:1979h} \emph{A survey of recent results in set theory}, Period. Math. Hungar. \textbf{10}, 109-175.

\item[]  \textbf{McCarten, S. D.}

\bibitem[\textbf{1988}]{bib:1988g} \emph{Topological equivalents of the axiom of choice}, Irish Math. Soc. Bull. \textbf{21} 45-48.

\item[]  \textbf{Morillon, M.}

\bibitem[\textbf{1985}]{bib:1985p} \emph{Axiomes de choix}, S\'{e}minaire d'Analyse de Clermont-Ferrand--1986.

\bibitem[\textbf{1986}]{bib:1986g} \emph{Une forme g\'{e}om\'{e}trique de l'axiome de l'id\'{e}al premier bool\'{e}en}, S\'{e}minaire d'Analyse de Clermont-Ferrand II-1987.

\bibitem[\textbf{1987}]{bib:1987e} \emph{Ensembles totalement ordonn\'{e}s et axiomes de choix}, S\'{e}minaire d'analyse, Universit\'{e} Blaise Pascal (Clermont II)-8, Exp. 5, 15 pp.

\bibitem[\textbf{1988}]{bib:1988h} \emph{Topologie, Analyse Nonstandard et Axiome du Choix}, Thesis, Universit\'{e} Blaise Pascal.

\bibitem[\textbf{1989}]{bib:1989r} \emph{Alg\`{e}bres de Boole mesur\'{e}es et axiome du choix}, S\'{e}minaire d'Analyse, Universit\'{e} Blaise-Pascal, Clermont-Ferrand-1990.

\bibitem[\textbf{1990}]{bib:1990e} \emph{Alg\`{e}bres de Gelfand mesur\'{e}es}, S\'{e}minaire d'Analyse, Universit\'{e} Blaise-Pascal, Clermont-Ferrand-1991.

\bibitem[\textbf{1991a}]{bib:1991a} \emph{Extreme choices on complete lexicographic orders}, Z. Math. Logik Grundlagen Math. \textbf{37}, 353--355.

\bibitem[\textbf{1991b}]{bib:1991b} \emph{Duaux continus et axiome du choix}, S\'{e}minaire d'Analyse, Universit\'{e} Blaise-Pascal, Clermont-Ferrand-1992.

\bibitem[\textbf{1993}]{bib:1993f} \emph{Les compacts extr\^{e}mement discontinus sont finis}, S\'{e}minaire d'Analyse, Universit\'{e} Blaise-Pascal, Clermont-Ferrand-1994.

\item[]  \textbf{Mirsky, L.}

\bibitem[\textbf{1971}]{bib:1971k} \emph{Transversal Theory}, Academic Press, New York.

\item[]  \textbf{Monro, G.P.}

\bibitem[\textbf{1964}]{bib:1964m} \emph{On the axiom of determinateness}, Fund. Math. \textbf{53}, 205-224.

\bibitem[\textbf{1972}]{bib:1972i} \emph{Models of ZF with the same sets of ordinals}, Notices Amer. Math. Soc. \textbf{19}, 534.

\bibitem[\textbf{1973a}]{bib:1973a1} \emph{Decomposable cardinals}, Fund. Math. \textbf{80}, 101-104.

\bibitem[\textbf{1973b}]{bib:1973b1} \emph{Models of ZF with the same sets of sets of ordinals}, Fund. Math. \textbf{80}, 105-110.

\bibitem[\textbf{1973c}]{bib:1973c1} \emph{Small sets with large power sets}, Bull. Austral. Math. Soc. \textbf{8}, 413--421.

\bibitem[\textbf{1974}]{bib:1974n} \emph{The cardinal equation} $2m=m$, Colloq. Math. \textbf{29}, 1-5.

\bibitem[\textbf{1975}]{bib:1975i} \emph{Independence results concerning Dedekind finite sets}, J. Austral. Math. Soc. \textbf{19}, 35-46.

\bibitem[\textbf{1983}]{bib:1983k} \emph{On generic extensions without the axiom of choice}, J. Symbolic Logic \textbf{48}, 39-52.

\item[]  \textbf{Monteiro, A.}

\bibitem[\textbf{1965}]{bib:1965g} \emph{G\'{e}n\'{e}ralization d'un th\'{e}or\`{e}me de R. Sikorski sur les alg\`{e}bres de Boole}, Bull. Sci. Math. \textbf{89}, 65--74.

\item[]  \textbf{Moore, G.H.}

\bibitem[\textbf{1982}]{bib:1982h} \emph{Zermelo's Axiom of Choice}, Springer-Verlan, New York.

\item[]  \textbf{Moore, R.L.}

\bibitem[\textbf{1924}]{bib:1924c} \emph{An extension of the theorem that no countable point set is perfect}, Proc. Nat. Acad. Sci. U.S.A. \textbf{10}, 168-170.

\item[]  \textbf{Morris, D.B.}

\bibitem[\textbf{1969}]{bib:1969f} \emph{Choice and cofinal well-ordered subsets}, Notices Amer. Math. Soc. \textbf{16}, 1088.

\bibitem[\textbf{1970}]{bib:1970k} \emph{A model of ZF which cannot be extended to a model of ZFC without adding ordinals}, Notices Amer. Math. Soc. \textbf{17}, 577

\item[]  \textbf{Mostowski, A.}

\bibitem[\textbf{1938}]{bib:1938} \emph{On the independence of the definitions of finiteness in a system of logic}, Ann. Soc. Polon. Math. \textbf{11}, 1-54.

\bibitem[\textbf{1939}]{bib:1939d} \emph{\"{U}ber die Unabh\"{a}ngigkeit des Wohlordnungssatzes vom Ordnungsprinzip}, Fund. Math. \textbf{32}, 201-252.

\bibitem[\textbf{1945}]{bib:1945} \emph{Axiom of choice for finite sets}, Fund. Math. \textbf{33}, 137-168.

\bibitem[\textbf{1948}]{bib:1948c} \emph{On the principle of dependent choices}, Fund. Math. \textbf{35}, 127-130.

\bibitem[\textbf{1958}]{bib:1958e} \emph{On a problem of W. Kinna and K. Wagner}, Colloq. Math. \textbf{6}, 207-208.

\item[]  \textbf{Mycielski, J.}

\bibitem[\textbf{1961}]{bib:1961} \emph{Some remarks and problems on the coloring of infinite graphs and the theorem of Kuratowski}, Acta. Math. Acad. Sci. Hungar. \textbf{12}, 125-129.

\bibitem[\textbf{1964}]{bib:1964n} \emph{Two remarks on Tychonoff's product theorem}, Bull. Acad. Polon. Sci. S\'{e}r. Sci. Math. Astronom. Phys. \textbf{8}, 439-441.

\item[]  \textbf{Myhill, J. R./Scott, D.}

\bibitem[\textbf{1971}]{bib:1971m} \emph{Ordinal definability}, Axiomatic Set Theory (Proc. Symp. Pure Math.\textbf{13}, A.M.S., Providence.), 271-278.

\item[]  \textbf{Parovi\v{c}enko, I.I.}

\bibitem[\textbf{1969}]{bib:1969g} \emph{Topological equivalents of the Tichonov theorem}, Dokl. Akad. Nauk SSSR \textbf{184}, 33-34.

\item[]  \textbf{Paseka, J.}

\bibitem[\textbf{1989}]{bib:1989s} \emph{The Tychonoff product theorem for finite spaces and the Boolean ultrafilter theorem are equivalent}, J. Univ. Kuwait Sci. \textbf{16}, 223--227.

\item[]  \textbf{Pawlikowski, J.}

\bibitem[\textbf{1991}]{bib:1991g} \emph{The Hahn-Banach theorem implies the Banach-Tarski paradox}, Fund. Math. \textbf{138}, 21--22.

\item[]  \textbf{Pelc, A.}

\bibitem[\textbf{1978}]{bib:1978g} \emph{On some weak forms of the axiom of choice in set theory}, Bull. Acad. Polon. Sci. S\'{e}r. Sci. Math. Astronom. Phys. \textbf{26}, 585-589.

\item[]  \textbf{Pincus, D.}

\bibitem[\textbf{1968}]{bib:1968h} \emph{Comparison of independence results in Mostowski's system (G) and ZF set theory}, Notices Amer. Math. Soc. \textbf{15}, 234.

\bibitem[\textbf{1969}]{bib:1969h} \emph{Individuals in Zermelo-Fraenkel set theory}, Doctoral Dissertation, Harvard University.

\bibitem[\textbf{1971}]{bib:1971n} \emph{Support structures for the axiom of choice}, J. Symbolic Logic \textbf{36}, 28-38.

\bibitem[\textbf{1972a}]{bib:1972} \emph{Zermelo-Fraenkel consistency results by Fraenkel-Mostowski methods}, J. Symbolic Logic \textbf{37}, 721-743.

\bibitem[\textbf{1972b}]{bib:1972b1} \emph{Independence of the prime ideal theorem from the Hahn Banach theorem}, Bull. Amer. Math. Soc. \textbf{78}, 766-770.

\bibitem[\textbf{1972c}]{bib:1972c1} \emph{The strength of the Hahn-Banach theorem}, Proc. of the Victoria Symp. on Nonstandard Analysis, Lecture Notes in Math. \textbf{369}, Springer, Heidelberg, 1973, 203-248.

\bibitem[\textbf{1974a}]{bib:1974a} \emph{Cardinal representatives}, Israel J. Math. \textbf{18}, 321-344.

\bibitem[\textbf{1974b}]{bib:1974b} \emph{On the independence of the Kinna Wagner principle}, Z. Math. Logik Grundlagen Math. \textbf{20}, 503-516.

\bibitem[\textbf{1974c}]{bib:1974c1} \emph{Postscript to}: \emph{On the independence of the Kinna Wagner principle}, preprint.

\bibitem[\textbf{1976}]{bib:1976n} \emph{Two model theoretic ideas in independence proofs}, Fund. Math. \textbf{92}, 113-130.

\bibitem[\textbf{1977a}]{bib:1977a4} \emph{Adding dependent choice}, Ann. Math. Logic \textbf{11}, 105-145.

\bibitem[\textbf{1977b}]{bib:1977b4} \emph{Adding dependent choice to the prime ideal theorem}, Logic Colloquium 76, R.O. Gandy and J.M.E. Hyland eds., Studies in Logic and the Foundations of Mathematics \textbf{87}, North Holland, Amsterdam, 547-565.

\bibitem[\textbf{1978}]{bib:1978h} \emph{A note on the cardinal factorial}, Fund. Math. \textbf{98}, 21-24.

\bibitem[\textbf{1997}]{bib:1997f} \emph{The dense linear ordering principle}, J. Symbolic Logic \textbf{62}, 438--456.

\item[]  \textbf{Pincus, D./Solovay, R.M.}

\bibitem[\textbf{1977}]{bib:1977j} \emph{Definability of measures and ultrafilters}, J. Symbolic Logic \textbf{42}, 179-190.

\item[]  \textbf{Plotkin, J. M.}

\bibitem[\textbf{1969}]{bib:1969i} \emph{Generic embeddings}, J. Symbolic Logic, \textbf{34}, 388-394.

\bibitem[\textbf{1976}]{bib:1976p} \emph{ZF and Boolean Algebras}, Israel J. Math. \textbf{23}, 298-308.

\bibitem[\textbf{1981}]{bib:1981i} \emph{ZF and locally finite groups}, Z. Math. Logik Grundlagen Math. \textbf{27}, 375-379.

\item[]  \textbf{Pogorzelski, W.A./Prucnal, T.}

\bibitem[\textbf{1974}]{bib:1974m1} \emph{Equivalence of the structural completeness theorem for propositional calculus and the Boolean representation theorem}, Rep. Math. Logic \textbf{3}, 37-40.

\item[]  \textbf{Prikry, K.}

\bibitem[\textbf{1976}]{bib:1976q} \emph{Determinateness and partitions}, Proc. Amer. Math. Soc. \textbf{54}, 303-306.

\item[]  \textbf{Rado, R.}

\bibitem[\textbf{1949}]{bib:1949} \emph{Axiomatic treatment of rank in infinite sets}, Canad. J. Math. \textbf{1}, 337-343.

\bibitem[\textbf{1971}]{bib:1971p} \emph{A selection lemma}, J. Combinatorial Theory Ser. A \textbf{10}, 176--177.

\item[]  \textbf{Raisonnier, Jean}

\bibitem[\textbf{1982}]{bib:1982i} \emph{Ensembles non mesurables et filtres rapides}, C. R. Acad. Sci. Paris S\'{e}r. I \textbf{294}, 285-287.

\bibitem[\textbf{1984}]{bib:1984n} \emph{A mathematical proof of S. Shelah's theorem on the measure problem and related results}, Israel J. Math. \textbf{48} no. 1, 48-56.

\item[]  \textbf{Ramsey, Frank}

\bibitem[\textbf{1929}]{bib:1929c} \emph{On a problem of formal logic}, Proc. London Math. Soc. \textbf{30}, 264-286.

\item[]  \textbf{Rav, Y.}

\bibitem[\textbf{1977}]{bib:1977k} \emph{Variants of Rado's selection lemma and their applications}, Math. Nachr. \textbf{79}, 145-165.

\bibitem[\textbf{1988}]{bib:1988i} \emph{Subdirect decomposition of rings and the axiom of choice}, Arch. Math. (Basel) \textbf{51}, 125-127.

\bibitem[\textbf{1989}]{bib:1989t} \emph{Lattice theoretical equivalencies of the ultrafilter principle}, Z. Math. Logik Grundlagen Math. \textbf{35}, 131-136.

\item[]  \textbf{von Rimscha, M.}

\bibitem[\textbf{1982}]{bib:1982j} \emph{Transitivit\"{a}tsbedingungen (Transitivity conditions)}, Z. Math. Logik Grundlagen Math., \textbf{28}, 67-74.

\item[]  \textbf{Royden, H. L.}

\bibitem[\textbf{1963}]{bib:1963f} \emph{Real Analysis}, The Macmillan Co., New York.

\item[]  \textbf{Rubin, A.}

\bibitem[\textbf{1978}]{bib:1978i} \emph{Independence of the axiom of choice from the greatest lower bound property of cardinals}, Notices Amer. Math. Soc. \textbf{25}, A-14.

\item[]  \textbf{Rubin, A./Rubin, J.}

\bibitem[\textbf{1974}]{bib:1974n1} \emph{The cardinality of the set of Dedekind finite cardinals in Fraenkel-Mostowski models}, Z. Math. Logik Grundlagen Math. \textbf{20}, 517-528.

\item[]  \textbf{Rubin, H.}

\bibitem[\textbf{1960}]{bib:1960a} \emph{Two propositions equivalent to the axiom of choice only under the axioms of extensionality and regularity}, Notices Amer. Math. Soc. \textbf{7}, 381.

\item[]  \textbf{Rubin, H./Rubin, J.}

\bibitem[\textbf{1963}]{bib:1963g} \emph{Equivalents of the Axiom of Choice}, North-Holland, Amsterdam.

\bibitem[\textbf{1985}]{bib:1985q} \emph{Equivalents of the Axiom of Choice, II}, North-Holland, Amsterdam.

\item[]  \textbf{Rubin, H./Scott, D.}

\bibitem[\textbf{1954}]{bib:1954c} \emph{Some topological theorems equivalent to the Boolean prime ideal theorem}, Bull. Amer. Math. Soc. \textbf{60}, 389.

\item[]  \textbf{Rubin, J.}

\bibitem[\textbf{1967}]{bib:1967f} \emph{Set theory for the mathematician}, Holden-Day, San Francisco.

\bibitem[\textbf{1976}]{bib:1976r} \emph{Properties of finite character and set character}, preprint.

\item[]  \textbf{Rudin, W.}

\bibitem[\textbf{1991}]{bib:1991h} \emph{Functional Analysis} 2nd edition, MeGraw-hill, New York.

\item[]  \textbf{Russell, B.}

\bibitem[\textbf{1906}]{bib:1906} \emph{On some difficulties in the theory of transfinite numbers and order types}, Proc. London Math. Soc. Ser. 2 \textbf{4},29-53.

\item[]  \textbf{Russell, B./Whitehead, A. N.}

\bibitem[\textbf{1925}]{bib:1925} \emph{Principia Mathematica}, 2nd ed, 3 volumes, Cambridge University Press, Cambridge, 1927.

\item[]  \textbf{Sacks, G. E.}

\bibitem[\textbf{1969}]{bib:1969j} \emph{Measure theoretic uniformity in recursion theory and axiomatic set theory} Trans. Amer. Math. Soc. \textbf{142}, 381-420.

\item[]  \textbf{Sageev, G.}

\bibitem[\textbf{1975}]{bib:1975j} \emph{An independence result concerning the axiom of choice}, Ann. Math. Logic \textbf{8}, 1-184.

\bibitem[\textbf{1981}]{bib:1981j} \emph{A model of ZF + there exists an inaccessible, in which the Dedekind cardinals constitute a natural, non-standard model of arithmetic}, Ann. Math. Logic \textbf{21}, 221-281.

\item[]  \textbf{Schechter, E.}

\bibitem[\textbf{1992}]{bib:1992f} \emph{Two topological equivalents of the axiom of choice}, Z. Math. Logic Grundlagen Math. \textbf{38}, 555-557.

\bibitem[\textbf{1996a}]{bib:1996a1} \emph{Handbook of Analysis and its Applications}, Academic Press, New York.

\bibitem[\textbf{1996b}]{bib:1996b1} \emph{Equivalents of the Hahn-Banach theorem and generalizations to vector} \emph{valued functions}, preprint, March.

\item[]  \textbf{Schnare, P. S.}

\bibitem[\textbf{1968}]{bib:1968j} \emph{The maximal $T_{0}$ (respectively $T_{1}$) lemma is equivalent to the axiom of choice}, Amer. Math. Monthly, \textbf{65} 761.

\item[]  \textbf{Schoenflies, A.}

\bibitem[\textbf{1908}]{bib:1908c} \emph{Die Entwickelung der Lehre von den Punktmannichfaltigkeiten. Zweiten Teil}, Jber. Deutsch. Math.-Verein. \textbf{2}, 1-331.

\item[]  \textbf{Schrijver, A.}

\bibitem[\textbf{1978}]{bib:1978j} \emph{The dependence of some logical axioms on disjoint transversals and linked systems}, Colloq. Math. \textbf{39}, 191-199.

\item[]  \textbf{Schr\"{o}der, E.}

\bibitem[\textbf{1898}]{bib:1898c} \emph{\"{U}ber zwei Definitionen der Endlichkeit und G. Cantor'sche S\"{a}tze}, Deutch. Akad. Naturforscher, nova Acta Leopoldina \textbf{71}, 303-362.

\item[]  \textbf{Scott, D.}

\bibitem[\textbf{1954}]{bib:1954d} \emph{Prime ideal theorems for rings lattices and Boolean algebras}, Bull. Amer. Math. Soc. \textbf{60}, 390.

\item[]  \textbf{Shannon, G.P.}

\bibitem[\textbf{1988}]{bib:1988j} \emph{Equivalent versions of a weak form of the axiom of choice}, Notre Dame J. Formal Logic \textbf{29}, 569-573.

\bibitem[\textbf{1990}]{bib:1990f} \emph{Provable forms of Martin's axiom}, Notre Dame J. Formal Logic \textbf{31}, 382-388.

\bibitem[\textbf{1992}]{bib:1992g} \emph{A note on some weak forms of the axiom of choice}, Notre Dame J. Formal Logic \textbf{33}, 144--147.

\item[]  \textbf{Shelah, S.}

\bibitem[\textbf{1984}]{bib:1984p} \emph{Can you take Solovay's inaccessible away?}, Israel J. Math. \textbf{48}, 1-47.

\bibitem[\textbf{1985}]{bib:1985r} \emph{On measure and category}, Israel J. Math. \textbf{52} 110-114.

\item[]  \textbf{Shoenfield, J.R.}

\bibitem[\textbf{1967}]{bib:1967g} \emph{Unramified forcing}, Axiomatic Set Theory; Los Angeles, vol. 1, D. Scott ed. (Proc. Symp. Pure Math. 13), AMS, Providence, 1971, 357-381.

\item[]  \textbf{Sierpi\'{n}ski, W.}

\bibitem[\textbf{1916}]{bib:1916} \emph{Sur le r\^{o}le de l'axiome de M. Zermelo dans l'analyse moderne}, C. R. Acad. Sci. Paris S\'{e}r. A-B \textbf{163}, 688-691.

\bibitem[\textbf{1917}]{bib:1917c} \emph{Sur quelques probl\`{e}mes qui impliquent des fonctions non-mesurables}, C. R. Acad. Sci. Paris S\'{e}r. A--B \textbf{164}, 882-884.

\bibitem[\textbf{1918}]{bib:1918} \emph{L'axiome de M. Zermelo et son r\^{o}le dans la th\'{e}orie des ensembles et l'analyse}, Bull. Int. Acad. Sci. Cracovie Cl. Math. Nat., 97-152.

\bibitem[\textbf{1928}]{bib:1928} Zarys teorii mnogosci, part I, 3rd ed., Kasa im. Mianowskiego, Wassay.

\bibitem[\textbf{1947}]{bib:1947} \emph{Sur une proposition qui entraine l'existence des ensembles non mesurables}, Fund. Math. \textbf{34}, 157-162.

\bibitem[\textbf{1954}]{bib:1954e} \emph{Sur un proposition \'{e}quivalente \`{a} l'existence d'un ensemble de nombres r\'{e}els de puissance} $\aleph_{1}$, Bull. Acad. Polon. Sci. S\'{e}r. Sci. Math., Astronom. et Phys. \textbf{2}, 53-54.

\item[]  \textbf{Sikorski, R.}

\bibitem[\textbf{1964}]{bib:1964p} \emph{Boolean Algrbras}, 2nd ed., Springer-Verlag, Berlin.

\item[]  \textbf{Smith, Perry}

\bibitem[\textbf{1982}]{bib:1982k} \emph{Three propositions equivalent to the axiom of choice}, Publ. Inst. Math. (Beograd) \textbf{32}, 165-166.

\item[]  \textbf{Smullyan, Raymond M.}

\bibitem[\textbf{1994}]{bib:1994f} \emph{Double superinduction and Zorn's lemma}, Zorn Memorial Lectures, Indiana University Logic Group, preprint no. IULG94-29 1-9.

\item[]  \textbf{Solovay, R.M.}

\bibitem[\textbf{1965}]{bib:1965h} \emph{The measure problem}, Notices Amer. Math. Soc. \textbf{12}, 217.

\bibitem[\textbf{1970}]{bib:1970l} \emph{A model of set theory in which every set of reals is Lebesgue measurable}, Ann. Math. Ser. 2, \textbf{92}, 1-56.

\bibitem[\textbf{1971}]{bib:1971q} \emph{Real-valued measurable cardinals}, Axiomatic Set Theory, Proc. of Symposia in Pure Math., Vol. XIII, Part I, 397-428.

\item[]  \textbf{Solovay, R.M./Reinhardt, W./Kanamori, A.}

\bibitem[\textbf{1978}]{bib:1978k} \emph{Strong axioms of infinity and elementary embeddings}, Ann. Math. Logic \textbf{13} 73-116.

\item[]  \textbf{Solovay, R.M./Tennenbaum, S.}

\bibitem[\textbf{1971}]{bib:1971r} \emph{Iterated Cohen extensions and Souslin's problem}, Ann. Math. Ser. 2 \textbf{94}, 201-245.

\item[]  \textbf{Specker, E.}

\bibitem[\textbf{1957}]{bib:1957c} \emph{Zur Axiomatik der Mengenlehre}, Z. Math. Logik Grundlagen Math. \textbf{3}, 173-210.

\item[]  \textbf{Spector, M.}

\bibitem[\textbf{1980}]{bib:1980n} \emph{A measurable cardinal with non-well founded ultra-power}, J. Symbolic Logic \textbf{45}, 623-628.

\bibitem[\textbf{1991}]{bib:1991i} \emph{Extended ultrapowers and the Vop\v{e}nka-Hrb\'{a}\v{c}ek theorem without choice}, J. Symbolic Logic \textbf{56}, 592--607.

\item[]  \textbf{Spi\u{s}iak, L.}

\bibitem[\textbf{1993}]{bib:1993g} \emph{Dependencies between definitions of finiteness II}, Czechoslovak Math. J. \textbf{43} (118), 391-407.

\item[]  \textbf{Spi\u{s}iak, L./Vojt\'{a}\u{s}, P.}

\bibitem[\textbf{1988}]{bib:1988k} \emph{Dependencies between definitions of finiteness}, Czechoslovak Math. J. \textbf{38}, 389-397.

\item[]  \textbf{Stavi, J.}

\bibitem[\textbf{1975}]{bib:1975k} \emph{A model of ZF with an infinite free complete Boolean algebra}, Israel J. Math. \textbf{20}, 149-164.

\item[]  \textbf{Stern, J.}

\bibitem[\textbf{1979}]{bib:1979i} \emph{Suites transfinies d'ensemble boreliens}, C. R. Acad. Sci. Paris Ser. A-B \textbf{288}, A527-A529.

\bibitem[\textbf{1985}]{bib:1985s} \emph{Generic extensions which do not add random reals}, Lecture Notes in Math., 1130, Springer, Berlin-New York, 395-407.

\item[]  \textbf{Stone, A.H.}

\bibitem[\textbf{1948}]{bib:1948d} \emph{Paracompactness and product space}, Bull. Amer. Math. Soc. \textbf{54}, 977-982.

\item[]  \textbf{Stone, M.}

\bibitem[\textbf{1936}]{bib:1936} \emph{The theory of representations for Boolean algebras}, Trans. Amer. Math. Soc. \textbf{40}, 37-111.

\item[]  \textbf{Strauch, O.}

\bibitem[\textbf{1980}]{bib:1980p} \emph{A theorem equivalent to the axiom of choice}, Acta Fac. Rerum Natur. Univ. Comenian. Math. \textbf{36}, 121-123.

\item[]  \textbf{Szczepaniak, Z.}

\bibitem[\textbf{1977}]{bib:1977l} \emph{The consistency of the Theory} $ZF + L^{1}\neq HOD$, Set Theory and Hierarchy Theory, V (Lecture Notes in Mathematics \textbf{619}), Springer-Verlag, Berlin, 285-290.

\item[]  \textbf{Tarski, A.}

\bibitem[\textbf{1924}]{bib:1924d} \emph{Sur les ensembles finis}, Fund. Math. \textbf{6}, 45-95.

\bibitem[\textbf{1930}]{bib:1930} \emph{Une contribution \'{a} la th\'{e}orie de la mesure}, Fund. Math. \textbf{15}, 42--50.

\bibitem[\textbf{1938a}]{bib:1938a} \emph{Review of Mostowski 5} in J. Symbolic Logic \textbf{3}, 115.

\bibitem[\textbf{1938b}]{bib:1938b} \emph{Ein \"{U}berdeckungsatz f\"{u}r endliche Mengen nebst enigen Bemerkungen \"{u}ber die Definitionen der Endlichkeit}, Fund. Math. \textbf{30}, 156-163.

\bibitem[\textbf{1939}]{bib:1939e} \emph{On well-ordered subsets of any set}, Fund. Math. \textbf{32}, 176--183.

\bibitem[\textbf{1948}]{bib:1948e} \emph{Algebraic and axiomatic aspects of two theorems on sums of cardinals}, Fund. Math. \textbf{35}, 79-104.

\bibitem[\textbf{1949a}]{bib:1949a} \emph{Cardinal Algebras}, Oxford Univ. Press.

\bibitem[\textbf{1949b}]{bib:1949b} \emph{Cancellation laws in the arithmetic of cardinals}, Fund. Math. \textbf{36}, 77-92.

\bibitem[\textbf{1954a}]{bib:1954a1} \emph{Theorems on the existence of successors of cardinals and the axiom of choice}, Indag. Math. \textbf{16}, 26-32.

\bibitem[\textbf{1954b}]{bib:1954b1} \emph{Prime ideal theorems for Boolean algebras and the axiom of choice}, Bull. Amer. Math. Soc. \textbf{60}, 390-391.

\bibitem[\textbf{1965}]{bib:1965i} \emph{On the existence of large sets of Dedekind cardinals}, Notices Amer. Math. Soc. \textbf{12}, 719.

\item[]  \textbf{Taskovi\'{c}, M.}

\bibitem[\textbf{1988}]{bib:1988l} \emph{Characterizations of inductive posets with applications}, Proc. Amer. Math. Soc. \textbf{104}, 650--660.

\bibitem[\textbf{1992a}]{bib:1992h} \emph{The axiom of choice, fixed point theorems and inductive ordered sets}, Proc. Amer. Math. Soc. \textbf{116}, 897--904.

\bibitem[\textbf{1992b}]{bib:1992i} \emph{New maximal principles}, Math. Japon. \textbf{37}, 549--554.

\item[]  \textbf{Truss, J. K.}

\bibitem[\textbf{1973a}]{bib:1973a2} \emph{Finite axioms of choice}, Ann. Math. Logic \textbf{6}, 147-176.

\bibitem[\textbf{1973b}]{bib:1973b2} \emph{On certain arbitrarily long sequences of cardinals}, Z. Math. Logik Grundlagen Math. \textbf{19}, 209-210.

\bibitem[\textbf{1973c}]{bib:1973c2} \emph{On successors in cardinal arithmetic}, Fund. Math. \textbf{78}, 7-21.

\bibitem[\textbf{1973d}]{bib:1973d2} \emph{The well-ordered and well-orderable subsets of a set}, Z. Math. Logik Grundlagen Math. \textbf{19}, 211--214.

\bibitem[\textbf{1973e}]{bib:1973e2} \emph{Convex sets of cardinals}, Proc. London Math. Soc. Ser. 3 \textbf{27}, 577--599.

\bibitem[\textbf{1974a}]{bib:1974a1} \emph{Classes of Dedekind finite cardinals}, Fund. Math. \textbf{84}, 187-208.

\bibitem[\textbf{1974b}]{bib:1974b1} \emph{Models of set theory containing many perfect sets}, Ann. Math. Logic \textbf{7}, 197-219.

\bibitem[\textbf{1975}]{bib:1975l} \emph{Some cases of K\"{o}nig's lemma}, Set Theory and Hierarchy Theory, Proc. Second Conf. Bierutowice, (Lecture Notes in Math., \textbf{537}) Springer-Verlag, Berlin, 1975, 273-384.

\bibitem[\textbf{1978}]{bib:1978l} \emph{The axiom of choice for linearly ordered families}, Fund. Math. \textbf{99}, 133-139.

\bibitem[\textbf{1984}]{bib:1984q} \emph{Cancellation laws for surjective cardinals}, Ann. Pure Appl. Logic \textbf{27}, 165-207.

\bibitem[\textbf{1994}]{bib:1994g} \emph{Permutations and the axiom of choice}, Automorphisms of First Order Structures, R. Kaye and D. Macpherson ed., 131-152, 351-368.

\bibitem[\textbf{1995}]{bib:1995g} \emph{The structure of amorphous sets}, Ann. Pure Appl. Logic, \textbf{73}, 191-233.

\item[]  \textbf{Tsukada, N.}

\bibitem[\textbf{1977}]{bib:1977m} \emph{Cardinals and the Boolean prime ideal theorem}, Sci. Rep. Tokyo Kyoiku Daigaku Sect. A \textbf{13}, 276-283.

\item[]  \textbf{Tychonoff, A.}

\bibitem[\textbf{1930}]{bib:1930c} \emph{\"{U}ber die topologische Erweiterung von R\"{a}umen}, Math. Ann. \textbf{102}, 544-561.

\item[]  \textbf{Urysohn, P.}

\bibitem[\textbf{1924}]{bib:1924e} \emph{\"{U}ber die Metrisation der kompakten topologischen R\"{a}ume}, Math. Ann. 92, 275-293.

\bibitem[\textbf{1925}]{bib:1925c} \emph{\"{U}ber die M\"{a}chtigkeit der zusammenh\"{a}ngenden Mengen}, Math. Ann. \textbf{94}, 262-295.

\item[]  \textbf{Vickers, S.}

\bibitem[\textbf{1989}]{bib:1989u} \emph{Topology via Logic}, Cambridge University Press, Cambridge.

\item[]  \textbf{Vop\u{e}nka, P.}

\bibitem[\textbf{1967}]{bib:1967h} \emph{General Theory of $\nabla$-models}, Comment. Math. Univ. Carolinae \textbf{8}, 145-170.

\item[]  \textbf{Wagon, S.}

\bibitem[\textbf{1985}]{bib:1985u} \emph{The Banach-Tarski Paradox}, Encyclopedia of Mathematics and its Applications, \textbf{24}), Cambridge University Press.

\item[]  \textbf{Ward, L.E.J.}

\bibitem[\textbf{1962}]{bib:1962f} \emph{A weak Tychonoff theorem and the axiom of choice}, Proc. Amer. Math. Soc. \textbf{13}, 757-758.

\item[]  \textbf{W\c{e}glorz, B.}

\bibitem[\textbf{1969}]{bib:1969k} \emph{A model of set theory $\mathcal{S}$ over a given Boolean algebra}, Bull. Acad. Polon. Sci. S\'{e}r. Sci. Math. Astronom. Phys., \textbf{17}, 201-202.

\item[]  \textbf{Welsh, P. J.}

\bibitem[\textbf{1973}]{bib:1973k} \emph{On the number of overlapping subsets of a set}, Notre Dame J. Formal Logic \textbf{14}, 269--270.

\item[]  \textbf{Whitehead, A.N.}

\bibitem[\textbf{1902}]{bib:1902} \emph{On cardinal numbers}, Amer. J. Math. \textbf{24}, 367-394.

\item[]  \textbf{Wi\'{s}niewski, K.}

\bibitem[\textbf{1972}]{bib:1972j} \emph{On the axiom of choice for families of finite sets}, Fund. Math. \textbf{73}, 187--192.

\bibitem[\textbf{1973}]{bib:1973m} \emph{On functions without fixed points}, Comment. Math. Prace Mat. \textbf{17}, 227--228.

\item[]  \textbf{Wojtylak, P.}

\bibitem[\textbf{1995}]{bib:1995h} \emph{2-SAT is not equivalent to} $BPI$, preprint June.

\item[]  \textbf{Wolk, E.S.}

\bibitem[\textbf{1967}]{bib:1967i} \emph{On the theorems of Tychonoff, Alexander and Rado}, Proc. Amer. Math. Soc. \textbf{18}, 113-115.

\bibitem[\textbf{1983}]{bib:1983m} \emph{On the principle of dependent choice and some forms of Zorn's Lemma}, Canad. Math. Bull. \textbf{26}, 365-367.

\item[]  \textbf{Wright, J.D.M.}

\bibitem[\textbf{1973}]{bib:1973n} \emph{All operators on a Hilbert space are bounded}, Bull. Amer. Math. Soc. \textbf{79}, 1247--1250.

\item[]  \textbf{Young, W.H.}

\bibitem[\textbf{1903}]{bib:1903} \emph{Overlapping intervals}, Proc. London Math. Soc. \textbf{35}, 384-388.

\item[]  \textbf{Zarach, Andrej}

\bibitem[\textbf{1977}]{bib:1977n} \emph{Extensions of $ZF$ models to models with the scheme of choice}, Fund. Math. \textbf{96}, 87-89.

\item[]  \textbf{Zermelo, E.}

\bibitem[\textbf{1904}]{bib:1904} \emph{Beweis das jede Menge wohlgeordnet kann}, Math. Ann. \textbf{59} 514-516.

\bibitem[\textbf{1908a}]{bib:1908a} \emph{Neuer Beweis f\"{u}r die wohlordnung}, Math. Ann. \textbf{65} 107-128.

\bibitem[\textbf{1908b}]{bib:1908b} \emph{Untersuchungen \"{u}ber die Grundlagen der Mengenlehre} $I$, Math. Ann. \textbf{65}, 261-281. (Translated in van Heijenoort 1, p 199-215.)

\item[]  \textbf{Zuckerman, M.M.}

\bibitem[\textbf{1968}]{bib:1968k} \emph{On a theorem of Sierpi\'{n}ski}, Matematiche (Catania) \textbf{23}, 469-500.

\bibitem[\textbf{1969a}]{bib:1969a} \emph{On choosing subsets of $n$-element sets}, Fund. Math. \textbf{64}, 163-179.

\bibitem[\textbf{1969b}]{bib:1969b} \emph{A unifying condition for implications among the axioms of choice for finite sets}, Pacific J. Math. \textbf{28}, 233-242.

\bibitem[\textbf{1971}]{bib:1971s} \emph{Choices from finite sets and choices of finite subsets}, Proc. Amer. Math. Soc. \textbf{27}, 133--138.

\bibitem[\textbf{1981}]{bib:1981k} \emph{Choosing $l$ element subsets of $n$-element sets}, Pacific J. Math. \textbf{96}, 247-250.
\end{thebibliography}

%\Printindex{Subject}{SUBJECT INDEX}
%\Printindex{Author}{AUTHOR INDEX}

\end{document}
